\PassOptionsToPackage{dvipsnames}{xcolor}
\documentclass[UTF8,12pt]{ctexart}

% 显式给 Fandol 中文字体登记 small-caps 形状。
% 模板的 \ChapterSectionLabel 与 titlepage 都用到 \textsc / \scshape，
% 默认查不到 CJK 字体的 sc 形状会出现 "Font shape ... sc undefined" 警告；
% 这里把 sc 显式落到 Regular/Hei，警告消失，输出无可见变化（中文本无大小写之分）。
\setCJKmainfont{FandolSong-Regular}[
  Extension  = .otf,
  ItalicFont = FandolKai-Regular,
  BoldItalicFont = FandolKai-Regular,
  Script = Default,
  BoldFont   = FandolHei-Regular,
  FontFace   = {m}{sc}{Font=FandolSong-Regular},
  FontFace   = {bx}{sc}{Font=FandolHei-Regular},
]

\usepackage[paper=a4paper,left=2.5cm,right=2.5cm,top=3cm,bottom=3cm]{geometry}
\usepackage{amsmath,amsthm,amssymb,mathrsfs,bbm,needspace}
\usepackage{graphicx,tikz-cd,titlesec}
\usepackage[dvipsnames]{xcolor}
% longtable 与 fvextra 在 hyperref 之前加载，使长表脚注与代码环境应用正确的链接补丁。
\usepackage{longtable,booktabs,enumitem,array,xurl,float,multirow,placeins,fvextra,multicol,makeidx}
\usepackage[
    colorlinks=true,
    linkcolor=MidnightBlue,
    citecolor=BrickRed,
    urlcolor=MidnightBlue,
    linktoc=all
]{hyperref}

\pagestyle{plain}

% 设置环境
\usepackage{aliascnt}
\theoremstyle{plain}
\newtheorem{theorem}{定理}[subsection]
\newaliascnt{definition}{theorem}
\newtheorem{definition}[definition]{定义}
\aliascntresetthe{definition}
\newaliascnt{lemma}{theorem}
\newtheorem{lemma}[lemma]{引理}
\aliascntresetthe{lemma}
\newaliascnt{corollary}{theorem}
\newtheorem{corollary}[corollary]{推论}
\aliascntresetthe{corollary}
\newaliascnt{remark}{theorem}
\newtheorem{remark}[remark]{注}
\aliascntresetthe{remark}
\newaliascnt{example}{theorem}
\newtheorem{example}[example]{例}
\aliascntresetthe{example}
\newaliascnt{proposition}{theorem}
\newtheorem{proposition}[proposition]{命题}
\aliascntresetthe{proposition}
\newaliascnt{problem}{theorem}
\newtheorem{problem}[problem]{问题}
\aliascntresetthe{problem}
\newtheorem{exercise}{习题}[section]
\numberwithin{equation}{subsection}
\numberwithin{figure}{subsection}
\numberwithin{table}{subsection}

% autoref names
\providecommand{\sectionautorefname}{章}
\renewcommand{\sectionautorefname}{章}

\providecommand{\subsectionautorefname}{节}
\renewcommand{\subsectionautorefname}{节}

\providecommand{\definitionautorefname}{定义}
\renewcommand{\definitionautorefname}{定义}

\providecommand{\theoremautorefname}{定理}
\renewcommand{\theoremautorefname}{定理}

\providecommand{\lemmaautorefname}{引理}
\renewcommand{\lemmaautorefname}{引理}

\providecommand{\propositionautorefname}{命题}
\renewcommand{\propositionautorefname}{命题}

\providecommand{\corollaryautorefname}{推论}
\renewcommand{\corollaryautorefname}{推论}

\providecommand{\exampleautorefname}{例}
\renewcommand{\exampleautorefname}{例}

\providecommand{\remarkautorefname}{注}
\renewcommand{\remarkautorefname}{注}

\providecommand{\exerciseautorefname}{习题}
\renewcommand{\exerciseautorefname}{习题}

\providecommand{\problemautorefname}{问题}
\renewcommand{\problemautorefname}{问题}

\providecommand{\figureautorefname}{图}
\renewcommand{\figureautorefname}{图}

\providecommand{\tableautorefname}{表}
\renewcommand{\tableautorefname}{表}

% 设置 proof 样式
\makeatletter
\renewcommand{\proofname}{\textbf{证明}}
\renewenvironment{proof}[1][\proofname]{\par
  \pushQED{\qed}%
  \normalfont
  \topsep6\p@\@plus6\p@\relax
  \trivlist
  \item[\hskip\labelsep
        \normalfont\bfseries
    {#1\@addpunct{.}}]\ignorespaces
}{%
  \popQED\endtrivlist\@endpefalse
}
\makeatother
\renewcommand{\qedsymbol}{\ensuremath{\square}}

% 设置 section 样式
\newlength{\chapterNumSpace}
\setlength{\chapterNumSpace}{1.4cm}

% 附录判断：正文显示 Chapter，\appendix 后显示 Appendix
\newif\ifInAppendix
\InAppendixfalse

% 目录中的附录判断
\newif\ifInAppendixToc
\InAppendixTocfalse

% 目录可调参数
\newlength{\tocSectionLabelWidth}
\setlength{\tocSectionLabelWidth}{4.2em}

\newlength{\tocAppendixGap}
\setlength{\tocAppendixGap}{0.2em}

\newlength{\tocSubIndent}
\setlength{\tocSubIndent}{1.2em}

\newlength{\tocSubNumWidth}
\setlength{\tocSubNumWidth}{3.0em}

% 目录 section 编号文字
\newcommand{\TocSectionNumber}[1]{%
  \ifInAppendixToc
    附录\hspace{\tocAppendixGap}#1%
  \else
    第~#1~章%
  \fi
}

\makeatletter

% section 目录格式：
% 正文：第 1 章 群的定义与性质
% 附录：附录A 集合论补遗
\renewcommand*\l@section[2]{%
  \ifnum\c@tocdepth>\z@
    \addvspace{0.55em}%
    \begingroup
      \parindent\z@
      \rightskip\@pnumwidth
      \parfillskip-\@pnumwidth
      \leavevmode
      \bfseries
      \setlength{\@tempdima}{\tocSectionLabelWidth}%
      \advance\leftskip\@tempdima
      \hskip-\leftskip
      \def\numberline##1{%
        \makebox[\@tempdima][l]{\TocSectionNumber{##1}}%
      }%
      #1\nobreak\hfill\nobreak\makebox[\@pnumwidth][r]{#2}\par
    \endgroup
  \fi
}

% subsection 目录格式：
% 1.1 循环群 / A.1 集合基础
\renewcommand*\l@subsection{%
  \@dottedtocline{2}{\tocSubIndent}{\tocSubNumWidth}%
}

\makeatother

\let\OriginalAppendix\appendix
\renewcommand{\appendix}{%
  \OriginalAppendix
  \InAppendixtrue
  \addtocontents{toc}{\protect\InAppendixToctrue}%
}

\newcommand{\ChapterSectionLabel}{%
  \thispagestyle{plain}%
  \vspace*{-0.5cm}%
  \makebox[\textwidth][r]{%
    {\fontsize{18pt}{18pt}\selectfont\scshape
      \ifInAppendix 附录\else 第\fi
    }\enspace
    {\fontsize{62pt}{62pt}\selectfont\bfseries\thesection}\enspace{\fontsize{18pt}{18pt}\selectfont\ifInAppendix\else 章\fi}%
  }%
  \par\vspace{-0.25cm}%
  \noindent\rule{\dimexpr\textwidth-\chapterNumSpace\relax}{0.45pt}%
}

\titleformat{\section}[display]
  {\normalfont}
  {\ChapterSectionLabel}
  {0cm}
  {\raggedleft\rightskip=\chapterNumSpace\fontsize{30pt}{35pt}\selectfont\bfseries}
  [\vspace{0.3cm}]

\newcommand{\sectionbreak}{\clearpage}

% 不编号 section，例如前言、目录
\titleformat{name=\section,numberless}[block]
  {\normalfont\Large\bfseries\centering}
  {}
  {0pt}
  {}

% section 整体间距
\titlespacing{\section}
  {0pt}
  {0pt}
  {0pt}

% subsection 样式
\ctexset{
  subsection = {
    format = \large\bfseries
  }
}



\makeindex
\DefineVerbatimEnvironment{sourcecode}{Verbatim}{fontsize=\small,breaklines=true,breakanywhere=true}
\setlist[enumerate,1]{label=(\alph*),leftmargin=2.6em}
\setlength{\emergencystretch}{2em}
\setcounter{tocdepth}{3}
\setcounter{secnumdepth}{3}
\renewcommand{\subsubsectionautorefname}{节}
\edef\SourcePDF{\detokenize{../../books/George Casella, Roger Berger - Statistical Inference (Chapman & Hall_CRC Texts in Statistical Science) (2024, Chapman and Hall_CRC) - libgen.li.pdf}}
\newcommand{\sourcepageref}[2]{\href[page=#1]{\SourcePDF}{#2}}
\newcommand{\sourcepageanchor}[1]{\phantomsection\label{src:#1}}
\newcommand{\semanticindexanchor}[1]{\phantomsection\label{idxocc:#1}}
\newcommand{\chapterdoi}[1]{\par\noindent{\footnotesize DOI：\href{https://doi.org/10.1201/9781003456285-#1}{10.1201/9781003456285-#1}}\par\smallskip}
\hypersetup{pdftitle={统计推断：第二版中文译稿},pdfauthor={George Casella; Roger L. Berger},bookmarksnumbered=true}
\begin{document}
\hypersetup{pageanchor=false}
\begin{titlepage}
\noindent\rule{\textwidth}{1.2pt}\\[0.1cm]\rule{\textwidth}{.8pt}\\[1.4cm]
{\fontsize{40pt}{52pt}\selectfont\bfseries 统计推断}\\[.8cm]
{\fontsize{22pt}{30pt}\selectfont\textit{Statistical Inference}}\\[1cm]
\rule{\textwidth}{.6pt}\\[1cm]
{\Large 第二版}\\[1cm]
{\large George Casella\\[.15cm]Roger L. Berger}\\[1.2cm]
{\small 2001年第二版的2024年CRC重印本.\\
中文译稿：全12章及前后材料.\\
原书版权归原权利人. 仅作学习参考.\\
原编号与原疑误保留；疑误及适用条件以译注说明.}
\vfill\begin{flushright}\small 2026年11月4日\end{flushright}
\end{titlepage}
\pagenumbering{roman}\hypersetup{pageanchor=true}

% PDF2–7；装饰性标志不重复描摹，文字信息保留
\section*{统计推断}\phantomsection
这部经典教材从概率论的基本原理建立理论统计学. 作者从概率基础出发，运用具有统计意义的技术、定义和概念发展统计推断理论；这些内容是此前概念的自然延伸与结果. 本书覆盖标准统计推断课程的全部主题，包括分布、随机变量、数据约简、点估计、假设检验与区间估计.

\textbf{特色}
\begin{itemize}
\item 经典的研究生层次统计推断教材；
\item 从概率论的基本原理发展统计理论的基本要素；
\item 文笔清晰，具备一定微积分基础即可阅读；
\item 覆盖标准推断课程的全部关键主题；
\item 全书包含数百个例子，帮助理解；
\item 每章均有大量难度递进的习题.
\end{itemize}
《统计推断》第二版主要面向统计学研究生，也适合具有扎实数学基础的统计学高年级本科生. 本书强调统计理论更实用的一面，更关心理解基本统计概念和推导合理的统计程序，而较少侧重形式上的最优性考察.

本书是Cengage Learning, Inc.于2001年首次出版的第二版的重印本.
\clearpage
\section*{Chapman \& Hall/CRC统计科学教材丛书}\phantomsection
\noindent Joseph K. Blitzstein，美国哈佛大学\\
Julian J. Faraway，英国巴斯大学\\
Martin Tanner，美国西北大学\\
Jim Zidek，加拿大不列颠哥伦比亚大学

\textbf{近期出版书目}
\begin{description}[style=nextline,font=\normalfont\bfseries,itemsep=5pt]
\item[Python贝叶斯建模与计算] Osvaldo A. Martin、Raviv Kumar和Junpeng Lao
\item[贝叶斯规则！应用贝叶斯建模导论] Alicia Johnson、Miles Ott和Mine Dogucu
\item[用R学习随机过程：导论] Olga Korosteleva
\item[使用R设计与分析实验和观察性研究] Nathan Taback
\item[数据科学中的时间序列：分析与预测] Wayne A. Woodward、Bivin Philip Sadler和Stephen Robertson
\item[统计理论：简明导论，第二版] Felix Abramovich和Ya'acov Ritov
\item[纵向数据的应用线性回归：着重缺失观测] Frans E.S. Tan和Shahab Jolani
\item[数理统计基础] Steffen Lauritzen
\item[医学研究中的生存数据建模，第四版] David Collett
\item[应用分类数据与计数数据分析，第二版] Wan Tang、Hua He和Xin M. Tu
\item[Python地理数据科学] Sergio Rey、Dani Arribas-Bel和Levi John Wolf
\item[多状态生存数据模型：率、风险与伪值] Per Kragh Andersen和Henrik Ravn
\item[用R处理环境流行病学中的时空方法，第二版] Gavin Shaddick、James V. Zidek和Alex Schmidt
\item[统计推断的大样本理论教程] W. Jackson Hall和David Oakes
\item[统计推断，第二版] George Casella和Roger L. Berger
\end{description}
丛书详情见：\href{https://www.routledge.com/Chapman-HallCRCTexts-in-Statistical-Science/book-series/CHTEXSTASCI}{\nolinkurl{https://www.routledge.com/Chapman--HallCRC-Texts-in-Statistical-Science/book-series/CHTEXSTASCI}}.
\clearpage
\begin{center}
\vspace*{3cm}
{\Huge\bfseries 统计推断}\\[1cm]
{\Large 第二版}\\[3cm]
George Casella\\[0.3cm]Roger L. Berger
\vfill
CRC Press，Taylor \& Francis Group\\
博卡拉顿\quad 伦敦\quad 纽约\\
CRC Press为Taylor \& Francis Group（Informa旗下企业）的出版品牌\\
Chapman \& Hall图书
\end{center}
\clearpage
\begingroup\small
本第二版最初于2001年由Cengage Learning, Inc.出版.

第二版于2024年由CRC Press出版：\\
2385 NW Executive Center Drive, Suite 320, Boca Raton FL 33431；\\
以及CRC Press，4 Park Square, Milton Park, Abingdon, Oxon, OX14 4RN.

CRC Press是Taylor \& Francis Group, LLC旗下的出版品牌.

\textcopyright\ 2024 George Casella和Roger L. Berger.

我们已尽合理努力保证所出版数据和信息的可靠性，但作者与出版者不能对所有材料的有效性或使用后果承担责任. 作者与出版者曾努力查找本书转载材料的版权所有人；对于未取得以此形式出版的许可的情况，谨向版权所有人致歉. 如有版权材料未获注明，请来函告知，以便在今后重印时更正.

除美国版权法允许的情形外，未经出版者书面许可，不得以任何现有或今后发明的电子、机械或其他方式，翻印、复制、传播或使用本书的任何部分，包括影印、缩微摄影、录音，以及任何信息存储或检索系统.

如需影印或以电子方式使用本书材料的许可，请访问\url{https://www.copyright.com}，或联系Copyright Clearance Center, Inc.（CCC），222 Rosewood Drive, Danvers, MA 01923，电话978-750-8400. 对于CCC不提供许可的作品，请联系\href{mailto:mpkbookspermissions@tandf.co.uk}{mpkbookspermissions@tandf.co.uk}.

商标说明：产品或企业名称可能是商标或注册商标；此处仅用于识别与说明，无意侵权.

\begin{quote}\textbf{美国国会图书馆出版编目资料}\\
姓名：Casella, George，作者；Berger, Roger W.，作者.\footnote{译注：原编目写Roger W.，本书扉页及版权正文写Roger L.，此处保留原编目姓名.}\\
题名：Statistical inference / George Casella and Roger Berger.\\
说明：第二版；Boca Raton: CRC Press, 2024；丛书：CRC texts in statistical science series；含索引.\\
摘要：“这部经典教材从概率论的基本原理建立理论统计学. 作者从概率基础出发，运用具有统计意义且为此前概念的自然延伸与结果的技术、定义和概念，发展统计推断理论. 它覆盖标准推断课程的全部主题，包括分布、随机变量、数据约简、点估计、假设检验、区间估计和回归.”——出版者提供.\\
标识符：LCCN 2023053302（印刷本），LCCN 2023053303（电子书）；ISBN 9781032593036（精装），9781032597942（平装），9781003456285（电子书）.\\
主题：LCSH: Experimental design（实验设计）.\\
分类：LCC QA279 .C375 2024（印刷本），LCC QA279（电子书）；DDC 519.5/7--dc23/eng/20231120.\\
印刷本记录：\url{https://lccn.loc.gov/2023053302}\\
电子书记录：\url{https://lccn.loc.gov/2023053303}
\end{quote}
ISBN：978-1-032-59303-6（精装）；978-1-032-59794-2（平装）；978-1-003-45628-5（电子书）.\\
DOI：\href{https://doi.org/10.1201/9781003456285}{10.1201/9781003456285}.

原书由KnowledgeWorks Global Ltd.使用Nimbus Roman字体排版.
\endgroup
\clearpage
\vspace*{5cm}\begin{center}\Large 献给Anne和Vicki\end{center}
\vfill\begin{center}\small Taylor \& Francis\\Taylor \& Francis Group\\\url{http://taylorandfrancis.com}\end{center}\clearpage

% 用户授权疑误脚注：2026-11-04



\section*{原书目录（原书页码）}\phantomsection\label{front:source-contents}
本节及后面的原书表、图、例目录保留原书页码；点击条目可跳到本译稿的对应位置. 译稿自己的页码见后面的“中文译稿目录”.

\begingroup\small
\begin{longtable}{@{}p{.86\textwidth}r@{}}
\hyperref[front:preface-second]{第二版前言} & \hyperref[front:preface-second]{xi} \\
\hyperref[front:preface-first]{第一版前言} & \hyperref[front:preface-first]{xiii} \\
\hyperref[front:authors]{作者} & \hyperref[front:authors]{xvii} \\
\hyperref[front:tables]{表目录} & \hyperref[front:tables]{xix} \\
\hyperref[front:figures]{图目录} & \hyperref[front:figures]{xxi} \\
\hyperref[front:examples]{例目录} & \hyperref[front:examples]{xxiii} \\
\hyperref[sec:1]{1 概率论} & \hyperref[sec:1]{1} \\
\hyperref[sec:1.1]{1.1 集合论} & \hyperref[sec:1.1]{1} \\
\hyperref[sec:1.2]{1.2 概率论} & \hyperref[sec:1.2]{4} \\
\hyperref[sec:1.2.1]{1.2.1 公理基础} & \hyperref[sec:1.2.1]{5} \\
\hyperref[sec:1.2.2]{1.2.2 概率运算} & \hyperref[sec:1.2.2]{8} \\
\hyperref[sec:1.2.3]{1.2.3 计数} & \hyperref[sec:1.2.3]{11} \\
\hyperref[sec:1.2.4]{1.2.4 枚举结果} & \hyperref[sec:1.2.4]{13} \\
\hyperref[sec:1.3]{1.3 条件概率与独立性} & \hyperref[sec:1.3]{16} \\
\hyperref[sec:1.4]{1.4 随机变量} & \hyperref[sec:1.4]{22} \\
\hyperref[sec:1.5]{1.5 分布函数} & \hyperref[sec:1.5]{24} \\
\hyperref[sec:1.6]{1.6 密度函数与质量函数} & \hyperref[sec:1.6]{28} \\
\hyperref[sec:1.7]{1.7 习题} & \hyperref[sec:1.7]{31} \\
\hyperref[sec:1.8]{1.8 杂录} & \hyperref[sec:1.8]{37} \\
\hyperref[sec:2]{2 变换与期望} & \hyperref[sec:2]{38} \\
\hyperref[sec:2.1]{2.1 随机变量的函数的分布} & \hyperref[sec:2.1]{38} \\
\hyperref[sec:2.2]{2.2 期望值} & \hyperref[sec:2.2]{45} \\
\hyperref[sec:2.3]{2.3 矩与矩母函数} & \hyperref[sec:2.3]{48} \\
\hyperref[sec:2.4]{2.4 积分号下求导} & \hyperref[sec:2.4]{56} \\
\hyperref[sec:2.5]{2.5 习题} & \hyperref[sec:2.5]{62} \\
\hyperref[sec:2.6]{2.6 杂录} & \hyperref[sec:2.6]{68} \\
\hyperref[sec:3]{3 常见分布族} & \hyperref[sec:3]{70} \\
\hyperref[sec:3.1]{3.1 引言} & \hyperref[sec:3.1]{70} \\
\hyperref[sec:3.2]{3.2 离散分布} & \hyperref[sec:3.2]{70} \\
\hyperref[sec:3.3]{3.3 连续分布} & \hyperref[sec:3.3]{81} \\
\hyperref[sec:3.4]{3.4 指数族} & \hyperref[sec:3.4]{90} \\
\hyperref[sec:3.5]{3.5 位置族与尺度族} & \hyperref[sec:3.5]{95} \\
\hyperref[sec:3.6]{3.6 不等式与恒等式} & \hyperref[sec:3.6]{100} \\
\hyperref[sec:3.6.1]{3.6.1 概率不等式} & \hyperref[sec:3.6.1]{100} \\
\hyperref[sec:3.6.2]{3.6.2 恒等式} & \hyperref[sec:3.6.2]{102} \\
\hyperref[sec:3.7]{3.7 习题} & \hyperref[sec:3.7]{105} \\
\hyperref[sec:3.8]{3.8 杂录} & \hyperref[sec:3.8]{111} \\
\hyperref[sec:4]{4 多个随机变量} & \hyperref[sec:4]{114} \\
\hyperref[sec:4.1]{4.1 联合分布与边缘分布} & \hyperref[sec:4.1]{114} \\
\hyperref[sec:4.2]{4.2 条件分布与独立性} & \hyperref[sec:4.2]{121} \\
\hyperref[sec:4.3]{4.3 二元变换} & \hyperref[sec:4.3]{128} \\
\hyperref[sec:4.4]{4.4 分层模型与混合分布} & \hyperref[sec:4.4]{133} \\
\hyperref[sec:4.5]{4.5 协方差与相关} & \hyperref[sec:4.5]{138} \\
\hyperref[sec:4.6]{4.6 多元分布} & \hyperref[sec:4.6]{145} \\
\hyperref[sec:4.7]{4.7 不等式} & \hyperref[sec:4.7]{152} \\
\hyperref[sec:4.7.1]{4.7.1 数值不等式} & \hyperref[sec:4.7.1]{152} \\
\hyperref[sec:4.7.2]{4.7.2 函数不等式} & \hyperref[sec:4.7.2]{155} \\
\hyperref[sec:4.8]{4.8 习题} & \hyperref[sec:4.8]{157} \\
\hyperref[sec:4.9]{4.9 杂录} & \hyperref[sec:4.9]{166} \\
\hyperref[sec:5]{5 随机样本的性质} & \hyperref[sec:5]{169} \\
\hyperref[sec:5.1]{5.1 随机样本的基本概念} & \hyperref[sec:5.1]{169} \\
\hyperref[sec:5.2]{5.2 随机样本中的随机变量之和} & \hyperref[sec:5.2]{172} \\
\hyperref[sec:5.3]{5.3 从正态分布抽样} & \hyperref[sec:5.3]{177} \\
\hyperref[sec:5.3.1]{5.3.1 样本均值与方差的性质} & \hyperref[sec:5.3.1]{178} \\
\hyperref[sec:5.3.2]{5.3.2 导出分布：Student的$t$与Snedecor的$F$} & \hyperref[sec:5.3.2]{181} \\
\hyperref[sec:5.4]{5.4 次序统计量} & \hyperref[sec:5.4]{184} \\
\hyperref[sec:5.5]{5.5 收敛概念} & \hyperref[sec:5.5]{188} \\
\hyperref[sec:5.5.1]{5.5.1 依概率收敛} & \hyperref[sec:5.5.1]{189} \\
\hyperref[sec:5.5.2]{5.5.2 几乎处处收敛} & \hyperref[sec:5.5.2]{190} \\
\hyperref[sec:5.5.3]{5.5.3 依分布收敛} & \hyperref[sec:5.5.3]{191} \\
\hyperref[sec:5.5.4]{5.5.4 Delta方法} & \hyperref[sec:5.5.4]{195} \\
\hyperref[sec:5.6]{5.6 随机样本的生成} & \hyperref[sec:5.6]{200} \\
\hyperref[sec:5.6.1]{5.6.1 直接方法} & \hyperref[sec:5.6.1]{201} \\
\hyperref[sec:5.6.2]{5.6.2 间接方法} & \hyperref[sec:5.6.2]{204} \\
\hyperref[sec:5.6.3]{5.6.3 接受／拒绝算法} & \hyperref[sec:5.6.3]{206} \\
\hyperref[sec:5.7]{5.7 习题} & \hyperref[sec:5.7]{208} \\
\hyperref[sec:5.8]{5.8 杂录} & \hyperref[sec:5.8]{217} \\
\hyperref[sec:6]{6 数据约简原则} & \hyperref[sec:6]{221} \\
\hyperref[sec:6.1]{6.1 引言} & \hyperref[sec:6.1]{221} \\
\hyperref[sec:6.2]{6.2 充分性原则} & \hyperref[sec:6.2]{221} \\
\hyperref[sec:6.2.1]{6.2.1 充分统计量} & \hyperref[sec:6.2.1]{222} \\
\hyperref[sec:6.2.2]{6.2.2 最小充分统计量} & \hyperref[sec:6.2.2]{228} \\
\hyperref[sec:6.2.3]{6.2.3 辅助统计量} & \hyperref[sec:6.2.3]{230} \\
\hyperref[sec:6.2.4]{6.2.4 充分、辅助与完备统计量} & \hyperref[sec:6.2.4]{231} \\
\hyperref[sec:6.3]{6.3 似然原则} & \hyperref[sec:6.3]{236} \\
\hyperref[sec:6.3.1]{6.3.1 似然函数} & \hyperref[sec:6.3.1]{236} \\
\hyperref[sec:6.3.2]{6.3.2 形式化的似然原则} & \hyperref[sec:6.3.2]{238} \\
\hyperref[sec:6.4]{6.4 等变性原则} & \hyperref[sec:6.4]{241} \\
\hyperref[sec:6.5]{6.5 习题} & \hyperref[sec:6.5]{244} \\
\hyperref[sec:6.6]{6.6 杂录} & \hyperref[sec:6.6]{250} \\
\hyperref[sec:7]{7 点估计} & \hyperref[sec:7]{252} \\
\hyperref[sec:7.1]{7.1 引言} & \hyperref[sec:7.1]{252} \\
\hyperref[sec:7.2]{7.2 寻找估计量的方法} & \hyperref[sec:7.2]{253} \\
\hyperref[sec:7.2.1]{7.2.1 矩估计法} & \hyperref[sec:7.2.1]{253} \\
\hyperref[sec:7.2.2]{7.2.2 最大似然估计量} & \hyperref[sec:7.2.2]{255} \\
\hyperref[sec:7.2.3]{7.2.3 贝叶斯估计量} & \hyperref[sec:7.2.3]{262} \\
\hyperref[sec:7.2.4]{7.2.4 EM算法} & \hyperref[sec:7.2.4]{264} \\
\hyperref[sec:7.3]{7.3 评价估计量的方法} & \hyperref[sec:7.3]{267} \\
\hyperref[sec:7.3.1]{7.3.1 均方误差} & \hyperref[sec:7.3.1]{267} \\
\hyperref[sec:7.3.2]{7.3.2 最优无偏估计量} & \hyperref[sec:7.3.2]{270} \\
\hyperref[sec:7.3.3]{7.3.3 充分性与无偏性} & \hyperref[sec:7.3.3]{277} \\
\hyperref[sec:7.3.4]{7.3.4 损失函数意义下的最优性} & \hyperref[sec:7.3.4]{282} \\
\hyperref[sec:7.4]{7.4 习题} & \hyperref[sec:7.4]{287} \\
\hyperref[sec:7.5]{7.5 杂录} & \hyperref[sec:7.5]{298} \\
\hyperref[sec:8]{8 假设检验} & \hyperref[sec:8]{302} \\
\hyperref[sec:8.1]{8.1 引言} & \hyperref[sec:8.1]{302} \\
\hyperref[sec:8.2]{8.2 寻找检验的方法} & \hyperref[sec:8.2]{303} \\
\hyperref[sec:8.2.1]{8.2.1 似然比检验} & \hyperref[sec:8.2.1]{303} \\
\hyperref[sec:8.2.2]{8.2.2 贝叶斯检验} & \hyperref[sec:8.2.2]{307} \\
\hyperref[sec:8.2.3]{8.2.3 并交检验与交并检验} & \hyperref[sec:8.2.3]{308} \\
\hyperref[sec:8.3]{8.3 评价检验的方法} & \hyperref[sec:8.3]{310} \\
\hyperref[sec:8.3.1]{8.3.1 错误概率与功效函数} & \hyperref[sec:8.3.1]{310} \\
\hyperref[sec:8.3.2]{8.3.2 最强检验} & \hyperref[sec:8.3.2]{314} \\
\hyperref[sec:8.3.3]{8.3.3 并交检验与交并检验的大小} & \hyperref[sec:8.3.3]{319} \\
\hyperref[sec:8.3.4]{8.3.4 $p$值} & \hyperref[sec:8.3.4]{322} \\
\hyperref[sec:8.3.5]{8.3.5 损失函数意义下的最优性} & \hyperref[sec:8.3.5]{324} \\
\hyperref[sec:8.4]{8.4 习题} & \hyperref[sec:8.4]{326} \\
\hyperref[sec:8.5]{8.5 杂录} & \hyperref[sec:8.5]{335} \\
\hyperref[sec:9]{9 区间估计} & \hyperref[sec:9]{337} \\
\hyperref[sec:9.1]{9.1 引言} & \hyperref[sec:9.1]{337} \\
\hyperref[sec:9.2]{9.2 寻找区间估计量的方法} & \hyperref[sec:9.2]{339} \\
\hyperref[sec:9.2.1]{9.2.1 反演检验统计量} & \hyperref[sec:9.2.1]{339} \\
\hyperref[sec:9.2.2]{9.2.2 枢轴量} & \hyperref[sec:9.2.2]{345} \\
\hyperref[sec:9.2.3]{9.2.3 对累积分布函数作枢轴变换} & \hyperref[sec:9.2.3]{348} \\
\hyperref[sec:9.2.4]{9.2.4 贝叶斯区间} & \hyperref[sec:9.2.4]{352} \\
\hyperref[sec:9.3]{9.3 评价区间估计量的方法} & \hyperref[sec:9.3]{356} \\
\hyperref[sec:9.3.1]{9.3.1 大小与覆盖概率} & \hyperref[sec:9.3.1]{356} \\
\hyperref[sec:9.3.2]{9.3.2 与检验相关的最优性} & \hyperref[sec:9.3.2]{359} \\
\hyperref[sec:9.3.3]{9.3.3 贝叶斯最优性} & \hyperref[sec:9.3.3]{362} \\
\hyperref[sec:9.3.4]{9.3.4 损失函数意义下的最优性} & \hyperref[sec:9.3.4]{363} \\
\hyperref[sec:9.4]{9.4 习题} & \hyperref[sec:9.4]{365} \\
\hyperref[sec:9.5]{9.5 杂录} & \hyperref[sec:9.5]{375} \\
\hyperref[sec:10]{10 渐近评价} & \hyperref[sec:10]{377} \\
\hyperref[sec:10.1]{10.1 点估计} & \hyperref[sec:10.1]{377} \\
\hyperref[sec:10.1.1]{10.1.1 相合性} & \hyperref[sec:10.1.1]{377} \\
\hyperref[sec:10.1.2]{10.1.2 效率} & \hyperref[sec:10.1.2]{379} \\
\hyperref[sec:10.1.3]{10.1.3 计算与比较} & \hyperref[sec:10.1.3]{382} \\
\hyperref[sec:10.1.4]{10.1.4 自助法标准误} & \hyperref[sec:10.1.4]{385} \\
\hyperref[sec:10.2]{10.2 稳健性} & \hyperref[sec:10.2]{388} \\
\hyperref[sec:10.2.1]{10.2.1 均值与中位数} & \hyperref[sec:10.2.1]{389} \\
\hyperref[sec:10.2.2]{10.2.2 M估计量} & \hyperref[sec:10.2.2]{391} \\
\hyperref[sec:10.3]{10.3 假设检验} & \hyperref[sec:10.3]{394} \\
\hyperref[sec:10.3.1]{10.3.1 似然比检验的渐近分布} & \hyperref[sec:10.3.1]{394} \\
\hyperref[sec:10.3.2]{10.3.2 其他大样本检验} & \hyperref[sec:10.3.2]{397} \\
\hyperref[sec:10.4]{10.4 区间估计} & \hyperref[sec:10.4]{401} \\
\hyperref[sec:10.4.1]{10.4.1 近似最大似然区间} & \hyperref[sec:10.4.1]{401} \\
\hyperref[sec:10.4.2]{10.4.2 其他大样本区间} & \hyperref[sec:10.4.2]{403} \\
\hyperref[sec:10.5]{10.5 习题} & \hyperref[sec:10.5]{408} \\
\hyperref[sec:10.6]{10.6 杂录} & \hyperref[sec:10.6]{417} \\
\hyperref[sec:11]{11 方差分析与回归} & \hyperref[sec:11]{421} \\
\hyperref[sec:11.1]{11.1 引言} & \hyperref[sec:11.1]{421} \\
\hyperref[sec:11.2]{11.2 单因素方差分析} & \hyperref[sec:11.2]{421} \\
\hyperref[sec:11.2.1]{11.2.1 模型与分布假设} & \hyperref[sec:11.2.1]{423} \\
\hyperref[sec:11.2.2]{11.2.2 经典的方差分析假设} & \hyperref[sec:11.2.2]{424} \\
\hyperref[sec:11.2.3]{11.2.3 关于均值线性组合的推断} & \hyperref[sec:11.2.3]{426} \\
\hyperref[sec:11.2.4]{11.2.4 方差分析的$F$检验} & \hyperref[sec:11.2.4]{428} \\
\hyperref[sec:11.2.5]{11.2.5 对比的同时估计} & \hyperref[sec:11.2.5]{431} \\
\hyperref[sec:11.2.6]{11.2.6 平方和的分解} & \hyperref[sec:11.2.6]{433} \\
\hyperref[sec:11.3]{11.3 简单线性回归} & \hyperref[sec:11.3]{435} \\
\hyperref[sec:11.3.1]{11.3.1 最小二乘：一种数学解法} & \hyperref[sec:11.3.1]{438} \\
\hyperref[sec:11.3.2]{11.3.2 最优线性无偏估计量：一种统计解法} & \hyperref[sec:11.3.2]{440} \\
\hyperref[sec:11.3.3]{11.3.3 模型与分布假设} & \hyperref[sec:11.3.3]{443} \\
\hyperref[sec:11.3.4]{11.3.4 正态误差下的估计与检验} & \hyperref[sec:11.3.4]{445} \\
\hyperref[sec:11.3.5]{11.3.5 在给定$x=x_0$处的估计与预测} & \hyperref[sec:11.3.5]{450} \\
\hyperref[sec:11.3.6]{11.3.6 同时估计与置信带} & \hyperref[sec:11.3.6]{452} \\
\hyperref[sec:11.4]{11.4 习题} & \hyperref[sec:11.4]{455} \\
\hyperref[sec:11.5]{11.5 杂录} & \hyperref[sec:11.5]{463} \\
\hyperref[sec:12]{12 回归模型} & \hyperref[sec:12]{467} \\
\hyperref[sec:12.1]{12.1 引言} & \hyperref[sec:12.1]{467} \\
\hyperref[sec:12.2]{12.2 变量含误差的回归} & \hyperref[sec:12.2]{467} \\
\hyperref[sec:12.2.1]{12.2.1 函数关系与结构关系} & \hyperref[sec:12.2.1]{469} \\
\hyperref[sec:12.2.2]{12.2.2 一种最小二乘解法} & \hyperref[sec:12.2.2]{470} \\
\hyperref[sec:12.2.3]{12.2.3 最大似然估计} & \hyperref[sec:12.2.3]{472} \\
\hyperref[sec:12.2.4]{12.2.4 置信集} & \hyperref[sec:12.2.4]{476} \\
\hyperref[sec:12.3]{12.3 逻辑回归} & \hyperref[sec:12.3]{478} \\
\hyperref[sec:12.3.1]{12.3.1 模型} & \hyperref[sec:12.3.1]{478} \\
\hyperref[sec:12.3.2]{12.3.2 估计} & \hyperref[sec:12.3.2]{480} \\
\hyperref[sec:12.4]{12.4 稳健回归} & \hyperref[sec:12.4]{483} \\
\hyperref[sec:12.5]{12.5 习题} & \hyperref[sec:12.5]{487} \\
\hyperref[sec:12.6]{12.6 杂录} & \hyperref[sec:12.6]{492} \\
\hyperref[app:computer-algebra]{计算机代数} & \hyperref[app:computer-algebra]{497} \\
\hyperref[app:distributions]{常见分布表} & \hyperref[app:distributions]{503} \\
\hyperref[app:references]{参考文献} & \hyperref[app:references]{509} \\
\hyperref[app:index]{索引} & \hyperref[app:index]{525} \\
\end{longtable}\endgroup



% 原书印刷xi–xii / PDF12–13
\section*{第二版前言}\phantomsection\label{front:preface-second}
\addcontentsline{toc}{section}{第二版前言}
本书的大部分引语出自阿瑟·柯南·道尔爵士，不过，若要描述这本书的历程，最贴切的也许是感恩至死乐队的那句感叹：“这一路真是漫长而奇异.”

第二版的计划大约始于六年前. 很长一段时间里，我们都在纠结应该增添什么、删去什么. 所幸，随着时间推移，统计学的发展方向日渐清晰，这些问题的答案也逐渐明朗. 我们看到，学科的趋势正在从特殊情形的优美证明，转向更复杂、更实际问题的算法解法. 这并不削弱数学与严谨性的重要性；事实上，我们发现它们变得更加重要了. 改变的是它们的应用方式.

对于熟悉第一版的读者，主要变化可简要概括如下. 渐近方法的讨论大幅扩充，独立成章. 我们更重视计算与模拟（见\hyperref[sec:5.5]{第5.5节}\footnote{译注：原文指向第5.5节；本版目录中“随机样本的生成”实际列于第5.6节.原引用保留.}及\hyperref[app:computer-algebra]{“计算机代数”附录}）；更具应用性的技术得到扩充或新增，例如自助法（bootstrapping）、EM算法、$p$值、逻辑回归与稳健回归；“杂录”和习题也新增了许多内容. 我们降低了等变性与决策论等较专门理论主题的比重，并调整了\hyperref[sec:3]{第3章}至\hyperref[sec:11]{第11章}的一些材料，使叙述更清楚.

有两点需要特别说明. 首先，关于计算机代数程序，虽然我们认为它们正成为越来越有价值的工具，但并不想强迫不持这种看法的教师使用它们. 因此，这部分以“不打扰”的方式处理：只放在附录中，并在书中一些可能有帮助的地方给出提示. 其次，我们采用了更便于查找的编号体系. 现在，定理、引理、例、定义等合用一个编号序列；例如，\hyperref[def:7.2.4]{定义7.2.4}之后是\hyperref[ex:7.2.5]{例7.2.5}，\hyperref[thm:10.1.3]{定理10.1.3}之后是\hyperref[ex:10.1.4]{例10.1.4}.

前四章只作了小幅改动. 我们重排了部分材料，特别是将不等式与恒等式分开，增添了一些例子和习题，并作了一般性的更新. \hyperref[sec:5]{第5章}也调整了顺序，把收敛一节移到较后的位置，并新增了随机变量生成一节. 第一版第7至9章中有关不变性的内容大幅缩减，合并到了\hyperref[sec:6]{第6章}；该章其余部分只作了少量编辑，主要是新增习题. \hyperref[sec:7]{第7章}扩充并更新，新增了EM算法一节. \hyperref[sec:8]{第8章}也作了少量编辑与更新，并新增$p$值一节. 在\hyperref[sec:9]{第9章}中，我们更加重视枢轴方法，因为我们意识到，“保证得到一个区间”实际上不过是“对累积分布函数作枢轴变换”. 此外，第一版第10章的决策论内容有所缩减，点估计、假设检验和区间估计在损失函数意义下的最优性分别以小节形式加入相应章节.

\hyperref[sec:10]{第10章}是全新的，力图阐明大样本推断的基础，包括Delta方法、相合性与渐近正态性、自助法、稳健估计量、得分检验等. \hyperref[sec:11]{第11章}讨论经典的单因素方差分析与线性回归；第一版中这些内容分属两章. 遗憾的是，受篇幅所限，我们删去了随机区组设计. \hyperref[sec:12]{第12章}讨论变量含误差的回归，并新增了稳健回归与逻辑回归的材料.

使用第一版授课多年之后，我们大致知道一年课程可以讲授多少内容. 例如，第二版中的下列内容应当可以在一年内讲完：
\begin{center}\begin{tabular}{ll@{\qquad}ll}
\hyperref[sec:1]{第1章} & 第1至7节 & \hyperref[sec:6]{第6章} & 第1至3节\\
\hyperref[sec:2]{第2章} & 第1至3节 & \hyperref[sec:7]{第7章} & 第1至3节\\
\hyperref[sec:3]{第3章} & 第1至6节 & \hyperref[sec:8]{第8章} & 第1至3节\\
\hyperref[sec:4]{第4章} & 第1至7节 & \hyperref[sec:9]{第9章} & 第1至3节\\
\hyperref[sec:5]{第5章} & 第1至6节 & \hyperref[sec:10]{第10章} & 第1、3、4节
\end{tabular}\end{center}
若学生开始上课时已有一定概率论基础，就可以多讲一些后面章节的材料.

最后，要感谢所有以某种方式促成第二版问世，并帮助我们改正第一版错误的人，几乎是不可能的. 感谢所有花时间寄来便条或电子邮件的学生、朋友和同事. 有些人的关键建议使我们的叙述方式发生了重大改变. 有时这些建议只是简短的便条或评论，有时则是较长的评阅意见. 有些事已过去很久，他们可能已经忘了，但我们没有忘记. 因此，感谢Arthur Cohen、David Cox爵士、Steve Samuels、Rob Strawderman和Tom Wehrly. 我们也十分感谢Jay Beder，他多年来寄来无数评论与建议，对第一版的熟悉程度可能比我们还高；还要感谢Michael Perlman及其班上的学生，就在我们写下这些话时，他们仍不断寄来意见与勘误.

这本书经历了几位编辑. 我们感谢Alex Kugashev，他在20世纪90年代中期最先建议编写第二版；感谢一直鼓励我们的编辑Carolyn Crockett. 除了我们之外，对这本书贡献最大的人也许就是第一位编辑John Kimmel，他鼓励、出版并推广了第一版. 谢谢你，John.
\begin{flushright}George Casella\\Roger L. Berger\end{flushright}

% 用户授权疑误脚注：2026-11-04



% 原书印刷xiii–xvi / PDF14–17，旧版章节号忠实保留，不能套用第二版目录
\section*{第一版前言}\phantomsection\label{front:preface-first}
\addcontentsline{toc}{section}{第一版前言}
当别人得知你正在写一本教科书时，往往会问两个问题中的一个，或者两个都问. 第一个是：“你为什么要写一本书？”第二个是：“你的书与已有的书有什么不同？”第一个问题相当容易回答：因为你对现有教材并不完全满意. 第二个问题则难得多. 答案无法用几句话说清；为了不使听众厌烦——他们也许只是出于礼貌才问——你会试着给出一个简短而机智的回答. 通常并不奏效.

本书的目的是从概率论的基本原理出发，建立理论统计学（theoretical statistics，有别于数理统计学）. 逻辑的发展、证明、思想与主题等，都通过统计论证逐步展开. 因而，我们从概率基础出发，用具有统计意义的技术、定义与概念发展统计推断理论；这些内容又是此前概念的自然延伸与推论. 开始这项工作时，我们并不知道这种做法会有多成功. 当然，最终的评判留给读者.

本书面向统计学专业，或需要专门学习统计学的相关领域的一年级研究生. 先修要求是一年的微积分. 熟悉一些矩阵运算会有帮助，但不是必需的. 本书可用于两个学期或三个学季的统计学入门课程.

前四章介绍概率论基础，并引入后文必需的许多基本内容. 第5章和第6章是最先讨论统计学的两章. 第5章起到概率与统计之间的过渡作用，已有概率基础的学生可由此开始统计理论课程. 第6章颇有特色，详细讨论充分性、似然与不变性三个统计原则，并说明它们在数据建模中的重要性. 并非所有教师都会详细讲授这一章，但我们强烈建议在此花一些时间. 其中，似然原则与不变性原则得到了详细讨论. 它们与充分性原则，以及这些原则背后的思考方式，都是全面理解统计学的基础.

第7至9章构成统计推断的核心：估计，包括点估计与区间估计，以及假设检验. 这些章节的一个主要特点，是区分寻找适当统计技术的方法与评价这些技术的方法. 理论工作者和实践工作者都关心寻找与评价，但我们认为把二者分开十分重要，因为它们关注的问题不同，遵循的规则也不同. 各章题为“其他考虑”的部分也可能引起读者兴趣. 我们在其中说明，统计推断的规则如何像日常实践中那样适当放宽，而仍能产生有意义的推断. 这些部分介绍的许多技术用于统计咨询，有助于分析实际问题并作出推断.

最后三章可以看作专题，但我们认为，任何人的统计教育都应包含对这些材料的一定了解. 第10章全面介绍决策论，并纳入了我们能够包含的最新内容. 第11章讨论方差分析，包括单因素与随机区组情形，从较简单的处理对比理论出发，建立完整分析的理论. 根据我们的经验，实验工作者最关心由对比作出的推断；运用前面发展出的原则，大多数检验与区间都可以由对比推导. 最后，第12章讨论回归理论，先介绍简单线性回归，再讨论“变量含误差”的回归. 后一个主题十分重要，既展示其自身的用途和固有困难，也揭示普通回归推断的局限.

关于如何用本书安排一年课程，我们给出以下更具体的建议. 本书可以用于两类不同的课程. 第一类可称为“偏重数学”，适合有扎实数学基础的统计学专业学生，即至少学过一年半微积分、一些矩阵代数，也许还学过实分析. 对这些学生，我们建议完整讲授第1至9章，大约需要22周，其余时间可从第10至12章中选取专题，按课程需要安排. 学完前九章之后，最后三章各自的内容都是自足的，可以按任意顺序讲授.

另一类课程“偏重实用”. 这种课程也可作为数学基础较强学生的第一门课程，但主要面向学过一年微积分、可能并非统计学专业的学生. 它强调统计理论中更实用的用途，更关心理解基本统计概念并针对不同情形推导合理的统计程序，而较少关注形式上的最优性研究. 这样的课程必然要省略部分材料，但以下各节可以在一年内讲授：
\begin{center}\begin{tabular}{cl}
章 & 节\\\hline
1 & 全部\\
2 & 2.1，2.2，2.3\\
3 & 3.1，3.2\\
4 & 4.1，4.2，4.3，4.5\\
5 & 5.1，5.2，5.3.1，5.4\\
6 & 6.1.1，6.2.1\\
7 & 7.1，7.2.1，7.2.2，7.2.3，7.3.1，7.3.3，7.4\\
8 & 8.1，8.2.1，8.2.3，8.2.4，8.3.1，8.3.2，8.4\\
9 & 9.1，9.2.1，9.2.2，9.2.4，9.3.1，9.4\\
11 & 11.1，11.2\\
12 & 12.1，12.2
\end{tabular}\end{center}
若时间允许，可以讨论第4.4、5.5、6.1.2、6.1.3和6.1.4节的材料，而不必过分强调细节. 第11.3和12.3节的内容也可以考虑.

习题取自许多来源，数量相当丰富. 我们认为，掌握这些材料的唯一途径也许就是练习，因此提供了充足的练习机会. 我们力求习题形式多样，其中许多题目说明了正文之外的新观点，或补充正文内容. 有些习题甚至取自研究论文. 当你能把自己学生时代还是新研究的论文改编成习题时，难免会觉得自己老了！习题虽然不像章节那样细分，但其排列大体遵循章节顺序，因为分组往往会给出过多提示. 此外，随着题号增大，习题的挑战性大体上也会增加.

本书是一本范围较广的入门教材，因此对各主题的讨论并不十分深入. 但我们认为，有责任引导读者在感兴趣的主题上再向前一步. 因此，书中收入了大量参考文献，指出深入理解各个主题的途径. Kotz、Johnson和Read主编的《统计科学百科全书》（\textit{Encyclopedia of Statistical Sciences}）为许多主题提供了很好的入门介绍.

写作本书时，我们借鉴了过去的教学经验与当前的研究工作，也得到了许多人的帮助，对此深怀感激. 感谢康奈尔大学、北卡罗来纳州立大学和普渡大学的同事，尤其是Jim Berger、Larry Brown、David Cox爵士、Ziding Feng、Janet Johnson、Leon Gleser、Costas Goutis、Dave Lansky、George McCabe、Chuck McCulloch、Myra Samuels、Steve Schwager和Shayle Searle. 他们付出时间与专业知识，阅读本书部分手稿，提供帮助，并在多次讨论中提出建设性意见. 我们还感谢Shanti Gupta的热情接待，以及对本书至关重要的普渡大学图书馆. 感谢Shayle Searle，以及匿名和具名审稿人Jim Albert、Dan Coster、Tom Wehrly的细致阅读与有益建议. 感谢David Moore和George McCabe允许我们使用他们的表格，感谢Steve Hirdt提供数据. 由于两位作者在写书的大部分时间里相隔至少600英里，最后，我们感谢Bitnet，使这一切成为可能.
\begin{flushright}George Casella\\Roger L. Berger\end{flushright}
\clearpage
\begin{quote}
“我们已经谈到演绎和推断了，”雷斯垂德向我眨眨眼说，“福尔摩斯，我觉得光是处理事实就够难了，可不愿再飞到理论和幻想中去.”
\begin{flushright}雷斯垂德探长对夏洛克·福尔摩斯\\《博斯科姆比溪谷秘案》\end{flushright}
\end{quote}



% 原书印刷xvii / PDF18
\section*{作者}\phantomsection\label{front:authors}\addcontentsline{toc}{section}{作者}
George Casella教授在福特汉姆大学完成本科教育，在普渡大学完成研究生教育. 他先后任教于罗格斯大学、康奈尔大学和佛罗里达大学. 他的统计学贡献包括蒙特卡洛方法、模型选择与基因组分析. 他尤其活跃于贝叶斯和经验贝叶斯方法领域：相关工作涉及Stein现象，通过Rao--Blackwell化技术等方法评估并加快马尔可夫链蒙特卡洛方法的收敛，以及把Lasso重新表述为采用独立Laplace先验的贝叶斯后验众数估计.

Casella博士于1988年当选美国统计学会和数理统计学会会士，1989年当选国际统计学会会士，2009年当选西班牙皇家科学院外籍院士.

Roger L. Berger教授在普渡大学取得统计学博士学位后，先后在佛罗里达州立大学与北卡罗来纳州立大学任教. 在2004年赴亚利桑那州立大学之前，他还曾在美国国家科学基金会工作两年. Berger教授是教材《统计推断》的合著者，该书现已出版第二版，并已译成中文与葡萄牙文. 他的论文发表于\textit{Journal of the American Statistical Association}、\textit{Statistical Science}、\textit{Biometrics}和\textit{Statistical Methods in Medical Research}等刊物. 他的专长包括假设检验、（生物）等效性、广义线性模型、生物统计学与统计教育.

Berger教授当选为美国统计学会和数理统计学会会士.



\section*{表目录（原书页码）}\phantomsection\label{front:tables}
\addcontentsline{toc}{section}{表目录（原书页码）}
\begingroup\small
\begin{longtable}{@{}p{.86\textwidth}r@{}}
\hyperref[tab:1.2.1]{1.2.1 排列方式的数目} & \hyperref[tab:1.2.1]{13} \\
\hyperref[tab:4.1.1]{4.1.1 联合概率质量函数$f(x,y)$的取值} & \hyperref[tab:4.1.1]{115} \\
\hyperref[tab:7.3.1]{7.3.1 二项分布参数$p$的三个估计量} & \hyperref[tab:7.3.1]{286} \\
\hyperref[tab:7.4.2]{7.4.2 白血病病例计数} & \hyperref[tab:7.4.2]{291} \\
\hyperref[tab:8.3.1]{8.3.1 假设检验中的两类错误} & \hyperref[tab:8.3.1]{310} \\
\hyperref[tab:9.2.1]{9.2.1 位置–尺度枢轴量} & \hyperref[tab:9.2.1]{346} \\
\hyperref[tab:9.2.2]{9.2.2 Sterne的接受域与置信集} & \hyperref[tab:9.2.2]{349} \\
\hyperref[tab:9.3.3]{9.3.3 三个90\%正态置信区间} & \hyperref[tab:9.3.3]{357} \\
\hyperref[tab:10.1.1]{10.1.1 自助法与Delta方法的方差} & \hyperref[tab:10.1.1]{387} \\
\hyperref[tab:10.2.2]{10.2.2 中位数相对于均值的渐近相对效率} & \hyperref[tab:10.2.2]{391} \\
\hyperref[tab:10.2.3]{10.2.3 Huber估计量} & \hyperref[tab:10.2.3]{391} \\
\hyperref[tab:10.2.4]{10.2.4 Huber估计量的渐近相对效率，$k=1.5$} & \hyperref[tab:10.2.4]{393} \\
\hyperref[tab:10.3.5]{10.3.5 Poisson似然比检验统计量} & \hyperref[tab:10.3.5]{396} \\
\hyperref[tab:10.3.6]{10.3.6 稳健检验的功效} & \hyperref[tab:10.3.6]{401} \\
\hyperref[tab:10.4.7]{10.4.7 枢轴区间的置信系数} & \hyperref[tab:10.4.7]{404} \\
\hyperref[tab:10.4.8]{10.4.8 基于Huber的M估计量的区间的置信系数} & \hyperref[tab:10.4.8]{407} \\
\hyperref[tab:11.3.1]{11.3.1 图11.3.1中绘制的数据} & \hyperref[tab:11.3.1]{439} \\
\hyperref[tab:11.3.2]{11.3.2 简单线性回归的方差分析表} & \hyperref[tab:11.3.2]{449} \\
\hyperref[tab:12.3.1]{12.3.1 挑战者号数据} & \hyperref[tab:12.3.1]{481} \\
\hyperref[tab:12.4.2]{12.4.2 长鼻袋鼠数据} & \hyperref[tab:12.4.2]{485} \\
\hyperref[tab:12.4.3]{12.4.3 回归M估计量的渐近相对效率} & \hyperref[tab:12.4.3]{487} \\
\hyperref[tab:12.5.4]{12.5.4 按年龄分组的雀鹰存活情况} & \hyperref[tab:12.5.4]{491} \\
\end{longtable}\endgroup



\section*{图目录（原书页码）}\phantomsection\label{front:figures}
\addcontentsline{toc}{section}{图目录（原书页码）}
\begingroup\small
\begin{longtable}{@{}p{.86\textwidth}r@{}}
\hyperref[fig:1.2.1]{1.2.1 例1.2.7的飞镖靶} & \hyperref[fig:1.2.1]{7} \\
\hyperref[fig:1.2.2]{1.2.2 平均值的直方图} & \hyperref[fig:1.2.2]{16} \\
\hyperref[fig:1.5.1]{1.5.1 例1.5.2的累积分布函数} & \hyperref[fig:1.5.1]{25} \\
\hyperref[fig:1.5.2]{1.5.2 几何分布的累积分布函数，$p=.3$} & \hyperref[fig:1.5.2]{27} \\
\hyperref[fig:1.6.1]{1.6.1 逻辑斯蒂曲线下的面积} & \hyperref[fig:1.6.1]{30} \\
\hyperref[fig:2.1.1]{2.1.1 例2.1.2的变换} & \hyperref[fig:2.1.1]{40} \\
\hyperref[fig:2.1.2]{2.1.2 递增与非递增的累积分布函数}\footnote{译注：原图目录写“increasing and nonincreasing cdfs”，这里照译；累积分布函数应为非递减函数.} & \hyperref[fig:2.1.2]{44} \\
\hyperref[fig:2.3.1]{2.3.1 指数密度} & \hyperref[fig:2.3.1]{49} \\
\hyperref[fig:2.3.2]{2.3.2 具有相同矩的两个概率密度函数} & \hyperref[fig:2.3.2]{53} \\
\hyperref[fig:2.3.3]{2.3.3 二项分布的Poisson近似} & \hyperref[fig:2.3.3]{56} \\
\hyperref[fig:3.3.1]{3.3.1 标准正态密度} & \hyperref[fig:3.3.1]{86} \\
\hyperref[fig:3.3.2]{3.3.2 二项分布的正态近似} & \hyperref[fig:3.3.2]{86} \\
\hyperref[fig:3.3.3]{3.3.3 Beta密度} & \hyperref[fig:3.3.3]{88} \\
\hyperref[fig:3.3.4]{3.3.4 对称Beta密度} & \hyperref[fig:3.3.4]{88} \\
\hyperref[fig:3.3.5]{3.3.5 标准正态密度与Cauchy密度} & \hyperref[fig:3.3.5]{89} \\
\hyperref[fig:3.3.6]{3.3.6 对数正态与Gamma概率密度函数} & \hyperref[fig:3.3.6]{91} \\
\hyperref[fig:3.5.1]{3.5.1 位置密度} & \hyperref[fig:3.5.1]{96} \\
\hyperref[fig:3.5.2]{3.5.2 指数位置密度} & \hyperref[fig:3.5.2]{98} \\
\hyperref[fig:3.5.3]{3.5.3 同一尺度族的成员} & \hyperref[fig:3.5.3]{98} \\
\hyperref[fig:3.5.4]{3.5.4 位置–尺度族} & \hyperref[fig:3.5.4]{99} \\
\hyperref[fig:4.1.1]{4.1.1 积分区域（一）} & \hyperref[fig:4.1.1]{120} \\
\hyperref[fig:4.5.1]{4.5.1 积分区域（二）} & \hyperref[fig:4.5.1]{142} \\
\hyperref[fig:4.5.2]{4.5.2 积分区域（三）} & \hyperref[fig:4.5.2]{143} \\
\hyperref[fig:4.7.1]{4.7.1 凸函数} & \hyperref[fig:4.7.1]{155} \\
\hyperref[fig:4.7.2]{4.7.2 Jensen不等式} & \hyperref[fig:4.7.2]{156} \\
\hyperref[fig:5.4.1]{5.4.1 例5.4.7中$f_{R,V}(r,v)>0$的区域} & \hyperref[fig:5.4.1]{189} \\
\hyperref[fig:5.6.1]{5.6.1 指数概率密度函数的直方图} & \hyperref[fig:5.6.1]{202} \\
\hyperref[fig:5.6.2]{5.6.2 Poisson样本方差的直方图} & \hyperref[fig:5.6.2]{204} \\
\hyperref[fig:5.6.3]{5.6.3 接受–拒绝抽样中的Beta分布} & \hyperref[fig:5.6.3]{205} \\
\hyperref[fig:7.3.1]{7.3.1 二项分布的均方误差比较} & \hyperref[fig:7.3.1]{269} \\
\hyperref[fig:7.3.2]{7.3.2 方差估计量的风险函数} & \hyperref[fig:7.3.2]{284} \\
\hyperref[fig:8.2.1]{8.2.1 似然比检验统计量} & \hyperref[fig:8.2.1]{305} \\
\hyperref[fig:8.3.1]{8.3.1 例8.3.2的功效函数} & \hyperref[fig:8.3.1]{311} \\
\hyperref[fig:8.3.2]{8.3.2 例8.3.3的功效函数} & \hyperref[fig:8.3.2]{312} \\
\hyperref[fig:8.3.3]{8.3.3 例8.3.19中三个检验的功效函数} & \hyperref[fig:8.3.3]{319} \\
\hyperref[fig:8.3.4]{8.3.4 例8.3.31中检验的风险函数} & \hyperref[fig:8.3.4]{325} \\
\hyperref[fig:9.2.1]{9.2.1 置信区间与接受域的关系} & \hyperref[fig:9.2.1]{340} \\
\hyperref[fig:9.2.2]{9.2.2 例9.2.3的接受域与置信区间} & \hyperref[fig:9.2.2]{342} \\
\hyperref[fig:9.2.3]{9.2.3 例9.2.16中的可信区间与置信区间} & \hyperref[fig:9.2.3]{353} \\
\hyperref[fig:9.2.4]{9.2.4 例9.2.16中的可信概率与区间} & \hyperref[fig:9.2.4]{354} \\
\hyperref[fig:9.2.5]{9.2.5 例9.2.16中可信区间的覆盖概率} & \hyperref[fig:9.2.5]{355} \\
\hyperref[fig:9.3.1]{9.3.1 例9.2.16中的三个区间估计量} & \hyperref[fig:9.3.1]{363} \\
\hyperref[fig:10.1.1]{10.1.1 Gamma均值估计量的渐近相对效率} & \hyperref[fig:10.1.1]{386} \\
\hyperref[fig:10.3.1]{10.3.1 Poisson似然的似然比检验} & \hyperref[fig:10.3.1]{396} \\
\hyperref[fig:10.4.1]{10.4.1 二项比例的似然比检验区间} & \hyperref[fig:10.4.1]{406} \\
\hyperref[fig:10.4.2]{10.4.2 名义置信水平为$.9$的二项置信程序的覆盖概率} & \hyperref[fig:10.4.2]{407} \\
\hyperref[fig:11.3.1]{11.3.1 残差平方和所度量的竖直距离} & \hyperref[fig:11.3.1]{438} \\
\hyperref[fig:11.3.2]{11.3.2 最优线性无偏估计量的几何描述} & \hyperref[fig:11.3.2]{443} \\
\hyperref[fig:11.3.3]{11.3.3 Scheffé带、$t$区间与Bonferroni区间} & \hyperref[fig:11.3.3]{455} \\
\hyperref[fig:12.2.1]{12.2.1 正交最小二乘所最小化的距离} & \hyperref[fig:12.2.1]{470} \\
\hyperref[fig:12.2.2]{12.2.2 三条回归直线} & \hyperref[fig:12.2.2]{472} \\
\hyperref[fig:12.2.3]{12.2.3 Creasy–Williams $F$统计量} & \hyperref[fig:12.2.3]{478} \\
\hyperref[fig:12.3.1]{12.3.1 挑战者号数据的逻辑斯蒂曲线} & \hyperref[fig:12.3.1]{481} \\
\hyperref[fig:12.4.1]{12.4.1 最小二乘、最小绝对偏差与M估计拟合} & \hyperref[fig:12.4.1]{485} \\
\end{longtable}\endgroup

% 用户授权疑误脚注：2026-11-04



% 原书例目录：印刷xxiii–xxix／PDF物理24–30；原编号与原页保留.
\clearpage\phantomsection
\section*{例目录（原书页码）}\phantomsection\label{front:examples}
\addcontentsline{toc}{section}{例目录（原书页码）}
\begin{longtable}{@{}p{.82\textwidth}r@{}}
例 1.1.3（事件运算） & \hyperref[ex:1.1.3]{2} \\
例 1.2.2（$\sigma$代数—I） & \hyperref[ex:1.2.2]{5} \\
例 1.2.3（$\sigma$代数—II） & \hyperref[ex:1.2.3]{5} \\
例 1.2.5（定义概率—I） & \hyperref[ex:1.2.5]{6} \\
例 1.2.7（定义概率—II） & \hyperref[ex:1.2.7]{7} \\
例 1.2.10（Bonferroni不等式） & \hyperref[ex:1.2.10]{9} \\
例 1.2.12（彩票—I） & \hyperref[ex:1.2.12]{11} \\
例 1.2.13（锦标赛） & \hyperref[ex:1.2.13]{11} \\
例 1.2.15（彩票—II） & \hyperref[ex:1.2.15]{11} \\
例 1.2.18（扑克牌） & \hyperref[ex:1.2.18]{13} \\
例 1.2.19（有放回抽样） & \hyperref[ex:1.2.19]{14} \\
例 1.2.20（计算平均值） & \hyperref[ex:1.2.20]{15} \\
例 1.3.1（四张A） & \hyperref[ex:1.3.1]{16} \\
例 1.3.3（例1.3.1续） & \hyperref[ex:1.3.3]{17} \\
例 1.3.4（三个囚犯） & \hyperref[ex:1.3.4]{18} \\
例 1.3.6（编码） & \hyperref[ex:1.3.6]{19} \\
例 1.3.8（de Méré骑士） & \hyperref[ex:1.3.8]{20} \\
例 1.3.10（掷两枚骰子） & \hyperref[ex:1.3.10]{21} \\
例 1.3.11（信件） & \hyperref[ex:1.3.11]{21} \\
例 1.3.13（抛掷硬币三次—I） & \hyperref[ex:1.3.13]{22} \\
例 1.4.2（随机变量） & \hyperref[ex:1.4.2]{23} \\
例 1.4.3（抛掷硬币三次—II） & \hyperref[ex:1.4.3]{23} \\
例 1.4.4（随机变量的分布） & \hyperref[ex:1.4.4]{24} \\
例 1.5.2（抛掷三枚硬币） & \hyperref[ex:1.5.2]{24} \\
例 1.5.4（抛掷直到正面出现） & \hyperref[ex:1.5.4]{25} \\
例 1.5.5（连续的累积分布函数） & \hyperref[ex:1.5.5]{26} \\
例 1.5.6（有跳跃的累积分布函数） & \hyperref[ex:1.5.6]{27} \\
例 1.5.9（同分布的随机变量） & \hyperref[ex:1.5.9]{28} \\
例 1.6.2（几何分布概率） & \hyperref[ex:1.6.2]{28} \\
例 1.6.4（Logistic分布概率） & \hyperref[ex:1.6.4]{30} \\
\addlinespace[.5em]
例 2.1.1（二项分布的变换） & \hyperref[ex:2.1.1]{39} \\
例 2.1.2（均匀分布的变换） & \hyperref[ex:2.1.2]{39} \\
例 2.1.4（均匀分布与指数分布的关系—I） & \hyperref[ex:2.1.4]{41} \\
例 2.1.6（逆伽马概率密度函数） & \hyperref[ex:2.1.6]{42} \\
例 2.1.7（平方变换） & \hyperref[ex:2.1.7]{42} \\
例 2.1.9（正态分布与卡方分布的关系） & \hyperref[ex:2.1.9]{43} \\
例 2.2.2（指数分布的均值） & \hyperref[ex:2.2.2]{45} \\
例 2.2.3（二项分布的均值） & \hyperref[ex:2.2.3]{45} \\
例 2.2.4（Cauchy分布的均值） & \hyperref[ex:2.2.4]{46} \\
例 2.2.6（使距离最小） & \hyperref[ex:2.2.6]{47} \\
例 2.2.7（均匀分布与指数分布的关系—II） & \hyperref[ex:2.2.7]{48} \\
例 2.3.3（指数分布的方差） & \hyperref[ex:2.3.3]{49} \\
例 2.3.5（二项分布的方差） & \hyperref[ex:2.3.5]{50} \\
例 2.3.8（伽马分布的矩母函数） & \hyperref[ex:2.3.8]{51} \\
例 2.3.9（二项分布的矩母函数） & \hyperref[ex:2.3.9]{52} \\
例 2.3.10（不能唯一确定分布的矩） & \hyperref[ex:2.3.10]{53} \\
例 2.3.13（Poisson近似） & \hyperref[ex:2.3.13]{54} \\
例 2.4.5（交换积分与微分—I） & \hyperref[ex:2.4.5]{58} \\
例 2.4.6（交换积分与微分—II） & \hyperref[ex:2.4.6]{59} \\
例 2.4.7（交换求和与微分） & \hyperref[ex:2.4.7]{60} \\
例 2.4.9（习题2.4.7续）\footnote{译注：原例目录写“Exercise 2.4.7”，但2.4.7是例的编号，原目录用词保留.} & \hyperref[ex:2.4.9]{61} \\
\addlinespace[.5em]
例 3.2.1（验收抽样） & \hyperref[ex:3.2.1]{72} \\
例 3.2.3（掷骰子的概率） & \hyperref[ex:3.2.3]{74} \\
例 3.2.4（等待时间） & \hyperref[ex:3.2.4]{76} \\
例 3.2.5（Poisson近似） & \hyperref[ex:3.2.5]{77} \\
例 3.2.6（逆二项抽样） & \hyperref[ex:3.2.6]{79} \\
例 3.2.7（失效时间） & \hyperref[ex:3.2.7]{80} \\
例 3.3.1（伽马分布与Poisson分布的关系） & \hyperref[ex:3.3.1]{82} \\
例 3.3.2（正态近似） & \hyperref[ex:3.3.2]{86} \\
例 3.4.1（二项指数族） & \hyperref[ex:3.4.1]{92} \\
例 3.4.3（二项分布的均值与方差） & \hyperref[ex:3.4.3]{93} \\
例 3.4.4（正态指数族） & \hyperref[ex:3.4.4]{93} \\
例 3.4.6（例3.4.4续） & \hyperref[ex:3.4.6]{94} \\
例 3.4.8（一个弯曲指数族） & \hyperref[ex:3.4.8]{95} \\
例 3.4.9（正态近似） & \hyperref[ex:3.4.9]{95} \\
例 3.5.3（指数分布的位置族） & \hyperref[ex:3.5.3]{97} \\
例 3.6.2（Chebychev不等式的例示） & \hyperref[ex:3.6.2]{101} \\
例 3.6.3（一个正态概率不等式） & \hyperref[ex:3.6.3]{101} \\
例 3.6.6（正态分布的高阶矩） & \hyperref[ex:3.6.6]{103} \\
例 3.6.9（Poisson分布的高阶矩） & \hyperref[ex:3.6.9]{104} \\
\addlinespace[.5em]
例 4.1.2（掷骰子的样本空间） & \hyperref[ex:4.1.2]{114} \\
例 4.1.4（例4.1.2续） & \hyperref[ex:4.1.4]{116} \\
例 4.1.5（掷骰子的联合概率质量函数） & \hyperref[ex:4.1.5]{116} \\
例 4.1.7（掷骰子的边缘概率质量函数） & \hyperref[ex:4.1.7]{117} \\
例 4.1.8（掷骰子的概率） & \hyperref[ex:4.1.8]{118} \\
例 4.1.9（边缘分布相同、联合概率质量函数不同） & \hyperref[ex:4.1.9]{118} \\
例 4.1.11（计算联合概率—I） & \hyperref[ex:4.1.11]{119} \\
例 4.1.12（计算联合概率—II） & \hyperref[ex:4.1.12]{120} \\
例 4.2.2（计算条件概率） & \hyperref[ex:4.2.2]{122} \\
例 4.2.4（计算条件概率密度函数） & \hyperref[ex:4.2.4]{123} \\
例 4.2.6（检验独立性—I） & \hyperref[ex:4.2.6]{125} \\
例 4.2.8（检验独立性—II） & \hyperref[ex:4.2.8]{126} \\
例 4.2.9（联合概率模型） & \hyperref[ex:4.2.9]{126} \\
例 4.2.11（独立变量的期望） & \hyperref[ex:4.2.11]{127} \\
例 4.2.13（正态变量之和的矩母函数） & \hyperref[ex:4.2.13]{128} \\
例 4.3.1（Poisson变量之和的分布） & \hyperref[ex:4.3.1]{129} \\
例 4.3.3（贝塔变量之积的分布） & \hyperref[ex:4.3.3]{130} \\
例 4.3.4（正态变量的和与差） & \hyperref[ex:4.3.4]{131} \\
例 4.3.6（正态变量之比的分布） & \hyperref[ex:4.3.6]{132} \\
例 4.4.1（二项–Poisson层次模型） & \hyperref[ex:4.4.1]{133} \\
例 4.4.2（例4.4.1续） & \hyperref[ex:4.4.2]{134} \\
例 4.4.5（例4.4.1的推广） & \hyperref[ex:4.4.5]{135} \\
例 4.4.6（贝塔–二项层次模型） & \hyperref[ex:4.4.6]{137} \\
例 4.4.8（例4.4.6续） & \hyperref[ex:4.4.8]{138} \\
例 4.5.4（相关性—I） & \hyperref[ex:4.5.4]{139} \\
例 4.5.8（相关性—II） & \hyperref[ex:4.5.8]{141} \\
例 4.5.9（相关性—III） & \hyperref[ex:4.5.9]{143} \\
例 4.6.1（多元概率密度函数） & \hyperref[ex:4.6.1]{146} \\
例 4.6.3（多元概率质量函数） & \hyperref[ex:4.6.3]{148} \\
例 4.6.8（伽马变量之和的矩母函数） & \hyperref[ex:4.6.8]{150} \\
例 4.6.13（多元变量替换） & \hyperref[ex:4.6.13]{152} \\
例 4.7.4（协方差不等式） & \hyperref[ex:4.7.4]{153} \\
例 4.7.8（一个关于均值的不等式） & \hyperref[ex:4.7.8]{156} \\
\addlinespace[.5em]
例 5.1.2（样本概率密度函数：指数分布） & \hyperref[ex:5.1.2]{170} \\
例 5.1.3（有限总体模型） & \hyperref[ex:5.1.3]{171} \\
例 5.2.8（均值的分布） & \hyperref[ex:5.2.8]{175} \\
例 5.2.10（Cauchy随机变量之和） & \hyperref[ex:5.2.10]{176} \\
例 5.2.12（Bernoulli随机变量之和） & \hyperref[ex:5.2.12]{177} \\
例 5.3.5（方差比的分布） & \hyperref[ex:5.3.5]{182} \\
例 5.3.7（例5.3.5续） & \hyperref[ex:5.3.7]{183} \\
例 5.4.5（均匀分布次序统计量的概率密度函数） & \hyperref[ex:5.4.5]{187} \\
例 5.4.7（中程数与极差的分布） & \hyperref[ex:5.4.7]{188} \\
例 5.5.3（$S^2$的相合性） & \hyperref[ex:5.5.3]{190} \\
例 5.5.5（$S$的相合性） & \hyperref[ex:5.5.5]{190} \\
例 5.5.7（几乎处处收敛） & \hyperref[ex:5.5.7]{191} \\
例 5.5.8（依概率收敛但不几乎处处收敛） & \hyperref[ex:5.5.8]{191} \\
例 5.5.11（均匀随机变量的最大值） & \hyperref[ex:5.5.11]{192} \\
例 5.5.16（负二项分布的正态近似） & \hyperref[ex:5.5.16]{194} \\
例 5.5.18（使用估计方差的正态近似） & \hyperref[ex:5.5.18]{195} \\
例 5.5.19（估计胜算） & \hyperref[ex:5.5.19]{195} \\
例 5.5.22（例5.5.19续） & \hyperref[ex:5.5.22]{197} \\
例 5.5.23（近似均值与方差） & \hyperref[ex:5.5.23]{197} \\
例 5.5.25（例5.5.23续） & \hyperref[ex:5.5.25]{198} \\
例 5.5.27（比率估计量的矩） & \hyperref[ex:5.5.27]{199} \\
例 5.6.1（指数分布的寿命） & \hyperref[ex:5.6.1]{200} \\
例 5.6.2（例5.6.1续） & \hyperref[ex:5.6.2]{200} \\
例 5.6.3（概率积分变换） & \hyperref[ex:5.6.3]{201} \\
例 5.6.4（Box–Muller，1958） & \hyperref[ex:5.6.4]{202} \\
例 5.6.5（二项随机变量的生成） & \hyperref[ex:5.6.5]{203} \\
例 5.6.6（Poisson方差的分布） & \hyperref[ex:5.6.6]{203} \\
例 5.6.7（贝塔随机变量的生成—I） & \hyperref[ex:5.6.7]{204} \\
例 5.6.9（贝塔随机变量的生成—II） & \hyperref[ex:5.6.9]{207} \\
\addlinespace[.5em]
例 6.2.3（二项分布的充分统计量） & \hyperref[ex:6.2.3]{223} \\
例 6.2.4（正态分布的充分统计量） & \hyperref[ex:6.2.4]{224} \\
例 6.2.5（充分的次序统计量） & \hyperref[ex:6.2.5]{224} \\
例 6.2.7（例6.2.4续） & \hyperref[ex:6.2.7]{225} \\
例 6.2.8（均匀分布的充分统计量） & \hyperref[ex:6.2.8]{226} \\
例 6.2.9（两个参数均未知时正态分布的充分统计量） & \hyperref[ex:6.2.9]{227} \\
例 6.2.12（两个正态充分统计量） & \hyperref[ex:6.2.12]{228} \\
例 6.2.14（正态分布的极小充分统计量） & \hyperref[ex:6.2.14]{229} \\
例 6.2.15（均匀分布的极小充分统计量） & \hyperref[ex:6.2.15]{229} \\
例 6.2.17（均匀分布的辅助统计量） & \hyperref[ex:6.2.17]{230} \\
例 6.2.18（位置族的辅助统计量） & \hyperref[ex:6.2.18]{231} \\
例 6.2.19（尺度族的辅助统计量） & \hyperref[ex:6.2.19]{231} \\
例 6.2.20（辅助精度） & \hyperref[ex:6.2.20]{232} \\
例 6.2.22（二项分布的完备充分统计量） & \hyperref[ex:6.2.22]{232} \\
例 6.2.23（均匀分布的完备充分统计量） & \hyperref[ex:6.2.23]{233} \\
例 6.2.26（使用Basu定理—I） & \hyperref[ex:6.2.26]{235} \\
例 6.2.27（使用Basu定理—II） & \hyperref[ex:6.2.27]{235} \\
例 6.3.2（负二项似然） & \hyperref[ex:6.3.2]{236} \\
例 6.3.3（正态信仰分布） & \hyperref[ex:6.3.3]{237} \\
例 6.3.4（证据函数） & \hyperref[ex:6.3.4]{238} \\
例 6.3.5（二项／负二项试验） & \hyperref[ex:6.3.5]{238} \\
例 6.3.7（例6.3.5续） & \hyperref[ex:6.3.7]{240} \\
例 6.4.1（二项等变性） & \hyperref[ex:6.4.1]{242} \\
例 6.4.3（例6.4.1续） & \hyperref[ex:6.4.3]{242} \\
例 6.4.5（例6.4.1的结论） & \hyperref[ex:6.4.5]{243} \\
例 6.4.6（正态位置不变性） & \hyperref[ex:6.4.6]{243} \\
\addlinespace[.5em]
例 7.2.1（正态分布的矩方法） & \hyperref[ex:7.2.1]{253} \\
例 7.2.2（二项分布的矩方法） & \hyperref[ex:7.2.2]{254} \\
例 7.2.3（Satterthwaite近似） & \hyperref[ex:7.2.3]{254} \\
例 7.2.5（正态似然） & \hyperref[ex:7.2.5]{256} \\
例 7.2.6（例7.2.5续） & \hyperref[ex:7.2.6]{257} \\
例 7.2.7（Bernoulli最大似然估计量） & \hyperref[ex:7.2.7]{257} \\
例 7.2.8（取值范围受限的最大似然估计量） & \hyperref[ex:7.2.8]{257} \\
例 7.2.9（试验次数未知的二项最大似然估计量） & \hyperref[ex:7.2.9]{258} \\
例 7.2.11（$\mu$与$\sigma$未知时的正态最大似然估计量） & \hyperref[ex:7.2.11]{260} \\
例 7.2.12（例7.2.11续） & \hyperref[ex:7.2.12]{260} \\
例 7.2.13（例7.2.2续） & \hyperref[ex:7.2.13]{261} \\
例 7.2.14（二项Bayes估计） & \hyperref[ex:7.2.14]{262} \\
例 7.2.16（正态Bayes估计量） & \hyperref[ex:7.2.16]{263} \\
例 7.2.17（多个Poisson发生率） & \hyperref[ex:7.2.17]{264} \\
例 7.2.18（例7.2.17续） & \hyperref[ex:7.2.18]{265} \\
例 7.2.19（例7.2.17的结论） & \hyperref[ex:7.2.19]{266} \\
例 7.3.3（正态均方误差） & \hyperref[ex:7.3.3]{267} \\
例 7.3.4（例7.3.3续） & \hyperref[ex:7.3.4]{268} \\
例 7.3.5（二项Bayes估计量的均方误差） & \hyperref[ex:7.3.5]{269} \\
例 7.3.6（等变估计量的均方误差） & \hyperref[ex:7.3.6]{270} \\
例 7.3.8（Poisson无偏估计） & \hyperref[ex:7.3.8]{271} \\
例 7.3.12（例7.3.8的结论） & \hyperref[ex:7.3.12]{274} \\
例 7.3.13（尺度均匀分布的无偏估计量） & \hyperref[ex:7.3.13]{275} \\
例 7.3.14（正态方差界） & \hyperref[ex:7.3.14]{276} \\
例 7.3.16（例7.3.14续） & \hyperref[ex:7.3.16]{276} \\
例 7.3.18（以非充分统计量为条件） & \hyperref[ex:7.3.18]{278} \\
例 7.3.21（零的无偏估计量） & \hyperref[ex:7.3.21]{279} \\
例 7.3.22（例7.3.13续） & \hyperref[ex:7.3.22]{280} \\
例 7.3.24（二项最佳无偏估计） & \hyperref[ex:7.3.24]{281} \\
例 7.3.25（二项风险函数） & \hyperref[ex:7.3.25]{283} \\
例 7.3.26（正态方差的风险） & \hyperref[ex:7.3.26]{283} \\
例 7.3.27（使用Stein损失的方差估计） & \hyperref[ex:7.3.27]{284} \\
例 7.3.28（两个Bayes规则） & \hyperref[ex:7.3.28]{285} \\
例 7.3.29（正态Bayes估计） & \hyperref[ex:7.3.29]{286} \\
例 7.3.30（二项Bayes估计） & \hyperref[ex:7.3.30]{286} \\
\addlinespace[.5em]
例 8.2.2（正态似然比检验） & \hyperref[ex:8.2.2]{304} \\
例 8.2.3（指数分布的似然比检验） & \hyperref[ex:8.2.3]{304} \\
例 8.2.5（似然比检验与充分性） & \hyperref[ex:8.2.5]{306} \\
例 8.2.6（方差未知的正态似然比检验） & \hyperref[ex:8.2.6]{306} \\
例 8.2.7（正态Bayes检验） & \hyperref[ex:8.2.7]{307} \\
例 8.2.8（正态并交检验） & \hyperref[ex:8.2.8]{308} \\
例 8.2.9（验收抽样） & \hyperref[ex:8.2.9]{309} \\
例 8.3.2（二项功效函数） & \hyperref[ex:8.3.2]{310} \\
例 8.3.3（正态功效函数） & \hyperref[ex:8.3.3]{311} \\
例 8.3.4（例8.3.3续） & \hyperref[ex:8.3.4]{312} \\
例 8.3.7（似然比检验的大小） & \hyperref[ex:8.3.7]{313} \\
例 8.3.8（并交检验的大小） & \hyperref[ex:8.3.8]{313} \\
例 8.3.10（例8.3.3的结论） & \hyperref[ex:8.3.10]{314} \\
例 8.3.14（二项一致最强检验） & \hyperref[ex:8.3.14]{316} \\
例 8.3.15（正态一致最强检验） & \hyperref[ex:8.3.15]{317} \\
例 8.3.18（例8.3.15续） & \hyperref[ex:8.3.18]{318} \\
例 8.3.19（一致最强检验不存在） & \hyperref[ex:8.3.19]{318} \\
例 8.3.20（无偏检验） & \hyperref[ex:8.3.20]{319} \\
例 8.3.22（一个等价关系） & \hyperref[ex:8.3.22]{320} \\
例 8.3.25（交并检验） & \hyperref[ex:8.3.25]{321} \\
例 8.3.28（双侧正态$p$值） & \hyperref[ex:8.3.28]{322} \\
例 8.3.29（单侧正态$p$值） & \hyperref[ex:8.3.29]{323} \\
例 8.3.30（Fisher精确检验） & \hyperref[ex:8.3.30]{323} \\
例 8.3.31（一致最强检验的风险） & \hyperref[ex:8.3.31]{325} \\
\addlinespace[.5em]
例 9.1.2（区间估计量） & \hyperref[ex:9.1.2]{337} \\
例 9.1.3（例9.1.2续） & \hyperref[ex:9.1.3]{338} \\
例 9.1.6（尺度均匀分布的区间估计量） & \hyperref[ex:9.1.6]{338} \\
例 9.2.1（反演正态检验） & \hyperref[ex:9.2.1]{339} \\
例 9.2.3（反演似然比检验） & \hyperref[ex:9.2.3]{342} \\
例 9.2.4（正态单侧置信界） & \hyperref[ex:9.2.4]{343} \\
例 9.2.5（二项单侧置信界） & \hyperref[ex:9.2.5]{344} \\
例 9.2.7（位置–尺度枢轴量） & \hyperref[ex:9.2.7]{345} \\
例 9.2.8（伽马枢轴量） & \hyperref[ex:9.2.8]{346} \\
例 9.2.9（例9.2.8续） & \hyperref[ex:9.2.9]{347} \\
例 9.2.10（正态枢轴区间） & \hyperref[ex:9.2.10]{347} \\
例 9.2.11（最短长度二项集合） & \hyperref[ex:9.2.11]{348} \\
例 9.2.13（位置指数分布的区间） & \hyperref[ex:9.2.13]{350} \\
例 9.2.15（Poisson区间估计量） & \hyperref[ex:9.2.15]{351} \\
例 9.2.16（Poisson可信集） & \hyperref[ex:9.2.16]{352} \\
例 9.2.17（Poisson可信概率与覆盖概率） & \hyperref[ex:9.2.17]{354} \\
例 9.2.18（正态可信集的覆盖概率） & \hyperref[ex:9.2.18]{354} \\
例 9.3.1（使长度最优） & \hyperref[ex:9.3.1]{356} \\
例 9.3.3（使期望长度最优） & \hyperref[ex:9.3.3]{358} \\
例 9.3.4（最短枢轴区间） & \hyperref[ex:9.3.4]{359} \\
例 9.3.6（一致最准确置信界） & \hyperref[ex:9.3.6]{360} \\
例 9.3.8（例9.3.6续） & \hyperref[ex:9.3.8]{361} \\
例 9.3.11（Poisson最高后验密度区域） & \hyperref[ex:9.3.11]{362} \\
例 9.3.12（正态最高后验密度区域） & \hyperref[ex:9.3.12]{363} \\
例 9.3.13（正态区间估计量） & \hyperref[ex:9.3.13]{364} \\
\addlinespace[.5em]
例 10.1.2（$\bar X$的相合性） & \hyperref[ex:10.1.2]{378} \\
例 10.1.4（例10.1.2续） & \hyperref[ex:10.1.4]{379} \\
例 10.1.8（极限方差） & \hyperref[ex:10.1.8]{379} \\
例 10.1.10（大样本混合方差） & \hyperref[ex:10.1.10]{380} \\
例 10.1.13（渐近正态性与相合性） & \hyperref[ex:10.1.13]{381} \\
例 10.1.14（近似二项方差） & \hyperref[ex:10.1.14]{382} \\
例 10.1.15（例10.1.14续） & \hyperref[ex:10.1.15]{383} \\
例 10.1.17（Poisson估计量的渐近相对效率） & \hyperref[ex:10.1.17]{384} \\
例 10.1.18（估计伽马分布的均值） & \hyperref[ex:10.1.18]{385} \\
例 10.1.19（用自助法估计方差） & \hyperref[ex:10.1.19]{386} \\
例 10.1.20（用自助法估计二项方差） & \hyperref[ex:10.1.20]{387} \\
例 10.1.21（例10.1.20的结论） & \hyperref[ex:10.1.21]{387} \\
例 10.1.22（参数自助法） & \hyperref[ex:10.1.22]{388} \\
例 10.2.1（样本均值的稳健性） & \hyperref[ex:10.2.1]{389} \\
例 10.2.3（中位数的渐近正态性） & \hyperref[ex:10.2.3]{390} \\
例 10.2.4（中位数相对于均值的渐近相对效率） & \hyperref[ex:10.2.4]{391} \\
例 10.2.5（Huber估计量） & \hyperref[ex:10.2.5]{391} \\
例 10.2.6（Huber估计量的极限分布） & \hyperref[ex:10.2.6]{392} \\
例 10.2.7（Huber估计量的渐近相对效率） & \hyperref[ex:10.2.7]{393} \\
例 10.3.2（Poisson似然比检验） & \hyperref[ex:10.3.2]{395} \\
例 10.3.4（多项似然比检验） & \hyperref[ex:10.3.4]{396} \\
例 10.3.5（大样本二项检验） & \hyperref[ex:10.3.5]{398} \\
例 10.3.6（二项得分检验） & \hyperref[ex:10.3.6]{400} \\
例 10.3.7（基于Huber估计量的检验） & \hyperref[ex:10.3.7]{400} \\
例 10.4.1（例10.1.14续） & \hyperref[ex:10.4.1]{402} \\
例 10.4.2（二项得分区间） & \hyperref[ex:10.4.2]{403} \\
例 10.4.3（二项似然比检验区间） & \hyperref[ex:10.4.3]{403} \\
例 10.4.4（近似区间） & \hyperref[ex:10.4.4]{404} \\
例 10.4.5（近似Poisson区间） & \hyperref[ex:10.4.5]{404} \\
例 10.4.6（二项得分区间的进一步讨论） & \hyperref[ex:10.4.6]{405} \\
例 10.4.7（二项区间的比较） & \hyperref[ex:10.4.7]{406} \\
例 10.4.8（基于Huber估计量的区间） & \hyperref[ex:10.4.8]{406} \\
例 10.4.9（负二项区间） & \hyperref[ex:10.4.9]{408} \\
例 10.6.2（均值与中位数的影响函数） & \hyperref[ex:10.6.2]{419} \\
\addlinespace[.5em]
例 11.2.1（单因素方差分析） & \hyperref[ex:11.2.1]{422} \\
例 11.2.3（方差分析假设） & \hyperref[ex:11.2.3]{424} \\
例 11.2.6（方差分析中的对比） & \hyperref[ex:11.2.6]{427} \\
例 11.2.9（成对差） & \hyperref[ex:11.2.9]{431} \\
例 11.2.12（例11.2.1续） & \hyperref[ex:11.2.12]{434} \\
例 11.3.1（预测葡萄产量） & \hyperref[ex:11.3.1]{436} \\
例 11.3.4（例11.3.1续） & \hyperref[ex:11.3.4]{449} \\
\addlinespace[.5em]
例 12.2.1（估计大气压） & \hyperref[ex:12.2.1]{468} \\
例 12.3.1（挑战者号数据） & \hyperref[ex:12.3.1]{481} \\
例 12.3.2（挑战者号数据续） & \hyperref[ex:12.3.2]{482} \\
例 12.4.1（最小二乘估计的稳健性） & \hyperref[ex:12.4.1]{484} \\
例 12.4.2（灾难性观测） & \hyperref[ex:12.4.2]{484} \\
例 12.4.3（最小绝对偏差估计量的渐近正态性） & \hyperref[ex:12.4.3]{484} \\
例 12.4.4（回归M估计量） & \hyperref[ex:12.4.4]{487} \\
例 12.4.5（回归渐近相对效率的模拟） & \hyperref[ex:12.4.5]{487} \\
例 12.6.1（无序抽样） & \hyperref[ex:12.6.1]{497} \\
例 12.6.2（一元变换） & \hyperref[ex:12.6.2]{498} \\
例 12.6.3（二元变换） & \hyperref[ex:12.6.3]{498} \\
例 12.6.4（正态概率） & \hyperref[ex:12.6.4]{499} \\
例 12.6.5（和的密度） & \hyperref[ex:12.6.5]{500} \\
例 12.6.6（均匀随机变量之和的四阶矩） & \hyperref[ex:12.6.6]{501} \\
例 12.6.7（伽马均值的渐近相对效率） & \hyperref[ex:12.6.7]{501} \\
例 12.6.8（卡方矩母函数的极限） & \hyperref[ex:12.6.8]{502} \\
\end{longtable}

% 用户授权疑误脚注：2026-11-04


\clearpage\renewcommand{\contentsname}{中文译稿目录}\tableofcontents\clearpage
\pagenumbering{arabic}
\section{概率论}\label{sec:1}
\chapterdoi{1}

% 原书印刷页1；PDF物理页32.章标题由主文件给出.

\sourcepageanchor{32}

\begin{quote}
例如，你永远无法预言某一个人会做什么，却可以准确地说出一群人平均而言会怎样行动.个体各不相同，但百分比保持不变.统计学家就是这样说的.
\begin{flushright}夏洛克·福尔摩斯\\《四签名》\end{flushright}
\end{quote}
概率论（probability theory）是一切统计学赖以建立的基础，它提供了对总体、试验，以及几乎任何可以视为随机现象的事物进行建模的方法.借助这些模型，统计学家能够仅考察整体的一部分，就对总体作出推断.

概率论有着悠久而丰富的历史，至少可以追溯到十七世纪.当时，应朋友德·梅雷骑士之请，帕斯卡和费马用数学形式描述了赌博中的胜算.

本章并不打算全面介绍概率论；在如此有限的篇幅内作这样的尝试未免鲁莽.我们的目标是概述概率论中一些对学习统计学至关重要的基本思想.

正如统计学建立在概率论的基础之上，概率论又建立在集合论的基础之上，因此我们从集合论开始.



% 原书印刷页1–4；PDF物理页32–35.逐页核对原图.
\subsection{集合论}\label{sec:1.1}
统计学家的主要目标之一，是通过开展试验对一个对象总体作出结论.第一步是确定所有可能的结果，用统计学的术语说，就是确定样本空间.
\begin{definition}\label{def:1.1.1}\index{SI0948@样本空间}
一个特定试验所有可能结果组成的集合$S$，称为该试验的\emph{样本空间}（sample space）.
\end{definition}
如果试验是抛掷一枚硬币，那么样本空间包含正面和反面这两个结果；因此，
$$S=\{\mathrm H,\mathrm T\}.$$
另一方面，如果试验是观察某所大学随机选出学生所报告的SAT成绩，那么样本空间就是200与800之间、为10的倍数的正整数集合，即$S=\{200,210,220,\ldots,780,790,800\}$.最后，考虑观察对某种刺激的反应时间的试验.此时样本空间由所有正数组成，即$S=(0,\infty)$.

按照元素个数，可以把样本空间分成两类：可数的（countable）和不可数的（uncountable）.如果样本空间的元素能够与整数的某个子集建立一一对应，那么样本空间就是可数的.当然，如果它只包含有限个元素，它便是可数的.因此，抛硬币和SAT成绩的样本空间都是可数的（实际上都是有限的），而反应时间的样本空间是不可数的，因为正实数无法与整数建立一一对应.不过，如果把反应时间测量到最接近的整数秒，那么样本空间（以秒为单位）就是$S=\{0,1,2,3,\ldots\}$，此时它便是可数的.


\sourcepageanchor{33}

可数与不可数样本空间的区别，重要之处只在于它决定了赋予概率的方式.尽管数学处理方式不同，但这在大多数情况下并不造成问题.从哲学层面看，人们或许会认为只能存在可数样本空间，因为测量不可能具有无限精度.（例如，由所有十位数字组成的样本空间就是可数的.）实际情况固然如此，但一般说来，与不可数样本空间相联系的概率和统计方法，比可数情形的方法更简便，而且能很好地逼近真实的（可数）情形.

一旦定义了样本空间，就可以考虑试验可能结果的各种集合.
\begin{definition}\label{def:1.1.2}\index{SI0418@事件}
\emph{事件}（event）是试验可能结果的任意集合，即$S$的任意子集（包括$S$本身）.
\end{definition}
设$A$是一个事件，即$S$的一个子集.如果试验结果属于集合$A$，就说事件$A$发生.讨论概率时，我们一般说事件的概率，而不是集合的概率.不过，这两个术语可以互换使用.

首先需要正式定义下面两种关系，它们使我们能够比较集合的包含关系并判定相等：
$$\begin{aligned}
A\subset B&\ \Longleftrightarrow\ x\in A\Rightarrow x\in B;&&\text{（包含）}\\
A=B&\ \Longleftrightarrow\ A\subset B\text{ 且 }B\subset A.&&\text{（相等）}
\end{aligned}$$
给定任意两个事件（或集合）$A$和$B$，有以下基本集合运算：
\begin{description}
\item[并（union）：]\index{SI0969@并}\semanticindexanchor{969:33} $A$与$B$的并记为$A\cup B$，是属于$A$或$B$，或同时属于二者的元素组成的集合：
$$A\cup B=\{x:x\in A\text{ 或 }x\in B\}.$$
\item[交（intersection）：]\index{SI0966@交}\semanticindexanchor{966:33} $A$与$B$的交记为$A\cap B$，是同时属于$A$和$B$的元素组成的集合：
$$A\cap B=\{x:x\in A\text{ 且 }x\in B\}.$$
\item[补（complementation）：]\index{SI0962@求补}\semanticindexanchor{962:33} $A$的补集记为$A^c$，是所有不属于$A$的元素组成的集合：
$$A^c=\{x:x\notin A\}.$$
\end{description}
\begin{example}[事件运算]\label{ex:1.1.3}
考虑从一副标准扑克牌中随机选一张并记录其花色的试验：梅花（C）、方块（D）、红桃（H）或黑桃（S）.样本空间为
$$S=\{\mathrm C,\mathrm D,\mathrm H,\mathrm S\},$$
其中一些可能事件为
$$A=\{\mathrm C,\mathrm D\}\quad\text{和}\quad B=\{\mathrm D,\mathrm H,\mathrm S\}.$$
由这些事件可构成
$$A\cup B=\{\mathrm C,\mathrm D,\mathrm H,\mathrm S\},\quad A\cap B=\{\mathrm D\},\quad A^c=\{\mathrm H,\mathrm S\}.$$
此外，注意$A\cup B=S$（事件$S$），且$(A\cup B)^c=\emptyset$，其中$\emptyset$表示\emph{空集}（empty set，即不含任何元素的集合）\index{SI0397@空集}\semanticindexanchor{397:33}.
\end{example}

\sourcepageanchor{34}

基本集合运算可以组合起来，这与加法和乘法的组合有些相似.只要小心，我们便可以像处理数一样处理集合.现在陈述集合运算的一些有用性质.
\begin{theorem}\label{thm:1.1.4}\index{SI0960@结合律}\index{SI0961@交换律}\index{SI0963@德摩根律}\index{SI0964@分配律}
对定义在样本空间$S$上的任意三个事件$A,B,C$，有
\begin{enumerate}
\item[(a)] 交换律：$A\cup B=B\cup A$，$A\cap B=B\cap A$；
\item[(b)] 结合律：$A\cup(B\cup C)=(A\cup B)\cup C$，$A\cap(B\cap C)=(A\cap B)\cap C$；
\item[(c)] 分配律：$A\cap(B\cup C)=(A\cap B)\cup(A\cap C)$，$A\cup(B\cap C)=(A\cup B)\cap(A\cup C)$；
\item[(d)] 德摩根律（DeMorgan's Laws）：$(A\cup B)^c=A^c\cap B^c$，$(A\cap B)^c=A^c\cup B^c$.
\end{enumerate}
\end{theorem}
\begin{proof}
本定理大部分结论的证明留作习题\ref{exercise:1.3}.习题\ref{exercise:1.9}和\ref{exercise:1.10}还给出了推广.为了说明证明方法，这里证明分配律
$$A\cap(B\cup C)=(A\cap B)\cup(A\cap C).$$
（读者可能熟悉用维恩图（Venn diagram）\index{SI1070@韦恩图}\semanticindexanchor{1070:34}“证明”集合论定理的做法.我们提醒读者，维恩图有时有助于直观理解，但并不构成严格证明.）要证明两个集合相等，必须说明它们彼此包含.正式地说，
$$\begin{aligned}
A\cap(B\cup C)&=\{x\in S:x\in A\text{ 且 }x\in(B\cup C)\};\\
(A\cap B)\cup(A\cap C)&=\{x\in S:x\in(A\cap B)\text{ 或 }x\in(A\cap C)\}.
\end{aligned}$$
先证明$A\cap(B\cup C)\subset(A\cap B)\cup(A\cap C)$.设$x\in(A\cap(B\cup C))$.由交的定义，必有$x\in(B\cup C)$，即$x\in B$或$x\in C$.由于$x$也必须属于$A$，所以$x\in(A\cap B)$或$x\in(A\cap C)$；因此，
$$x\in\bigl((A\cap B)\cup(A\cap C)\bigr),$$
包含关系得证.

反过来，设$x\in((A\cap B)\cup(A\cap C))$.于是$x\in(A\cap B)$或$x\in(A\cap C)$.若$x\in(A\cap B)$，则$x$同时属于$A$和$B$.因为$x\in B$，所以$x\in(B\cup C)$，从而$x\in(A\cap(B\cup C))$.另一方面，若$x\in(A\cap C)$，论证类似，仍可得出$x\in(A\cap(B\cup C))$.这样就证明了$(A\cap B)\cup(A\cap C)\subset A\cap(B\cup C)$，即另一个方向的包含关系，因此分配律得证.
\end{proof}
并和交的运算也可以推广到无穷集合族.\index{SI0967@可数}\semanticindexanchor{967:34}\index{SI0970@可数}\semanticindexanchor{970:34}若$A_1,A_2,A_3,\ldots$是一族都定义在样本空间$S$上的集合，则
$$\begin{aligned}
\bigcup_{i=1}^{\infty}A_i&=\{x\in S:\text{对某个 }i,\ x\in A_i\},\\
\bigcap_{i=1}^{\infty}A_i&=\{x\in S:\text{对所有 }i,\ x\in A_i\}.
\end{aligned}$$

\sourcepageanchor{35}

例如，令$S=(0,1]$并定义$A_i=[(1/i),1]$.那么
$$\begin{aligned}
\bigcup_{i=1}^{\infty}A_i&=\bigcup_{i=1}^{\infty}[(1/i),1]=\{x\in(0,1]:\text{对某个 }i,\ x\in[(1/i),1]\}\\
&=\{x\in(0,1]\}=(0,1];\\
\bigcap_{i=1}^{\infty}A_i&=\bigcap_{i=1}^{\infty}[(1/i),1]=\{x\in(0,1]:\text{对所有 }i,\ x\in[(1/i),1]\}\\
&=\{x\in(0,1]:x\in[1,1]\}=\{1\}.\qquad\text{（点1）}
\end{aligned}$$
也可以对不可数的集合族定义并和交.\index{SI0968@不可数}\semanticindexanchor{968:35}\index{SI0971@不可数}\semanticindexanchor{971:35}如果$\Gamma$是一个指标集（其中的元素用作指标），则
$$\begin{aligned}
\bigcup_{a\in\Gamma}A_a&=\{x\in S:\text{对某个 }a,\ x\in A_a\},\\
\bigcap_{a\in\Gamma}A_a&=\{x\in S:\text{对所有 }a,\ x\in A_a\}.
\end{aligned}$$
例如，取$\Gamma=\{\text{所有正实数}\}$且$A_a=(0,a]$，那么$\bigcup_{a\in\Gamma}A_a=(0,\infty)$就是一个不可数并.虽然不可数并和交在统计学中并不占主要地位，但有时它们提供了求解问题的有效手段（见\hyperref[sec:8.2.3]{第8.2.3节}）.

最后讨论样本空间的划分.
\begin{definition}\label{def:1.1.5}\index{SI0249@不交集}\index{SI0745@两两不交集}
若两个事件$A$和$B$满足$A\cap B=\emptyset$，就称它们\emph{不相交}（disjoint）或\emph{互斥}（mutually exclusive）.若对所有$i\ne j$都有$A_i\cap A_j=\emptyset$，就称事件$A_1,A_2,\ldots$\emph{两两不相交}或\emph{两两互斥}.
\end{definition}
不相交的集合没有公共点.若给两个不相交集合画维恩图，则它们不重叠.集合族
$$A_i=[i,i+1),\qquad i=0,1,2,\ldots$$
中的集合两两不相交.进一步注意到$\bigcup_{i=0}^{\infty}A_i=[0,\infty)$.
\begin{definition}\label{def:1.1.6}\index{SI0753@划分}
若$A_1,A_2,\ldots$两两不相交且$\bigcup_{i=1}^{\infty}A_i=S$，则称集合族$A_1,A_2,\ldots$构成$S$的一个\emph{划分}（partition）.
\end{definition}
集合$A_i=[i,i+1)$构成$[0,\infty)$的一个划分.一般说来，划分十分有用，它使我们能够把样本空间分成互不重叠的小块.



% 原书印刷页4–16（至1.3节前）；PDF物理页35–47.逐页核对原图.
\subsection{概率论}\label{sec:1.2}
进行一次试验时，其实现是样本空间中的一个结果.若把试验重复多次，每次可能出现不同的结果，也可能重复出现某些结果.一个结果的这种“出现频率”可以看作一种概率.概率较大的结果出现得更频繁.如果能够用概率来描述试验的结果，我们就迈出了用统计方法分析试验的一步.

本节介绍概率论的一些基础知识.我们不以频率来定义概率，而是采用数学上更简单的公理化方法.正如将要看到的，公理化方法不关心概率的解释，只关心概率是否由满足公理的函数定义.概率的解释是另一回事.事件的“出现频率”就是概率的一种特定解释.另一种可能的解释是主观解释：不把概率看作频率，而把它看作对某事件发生机会的信念.


\sourcepageanchor{36}

\subsubsection{公理化基础}\label{sec:1.2.1}
对样本空间$S$中的每个事件$A$，我们希望给它赋予一个介于0和1之间的数，称为$A$的概率，记作$P(A)$.似乎很自然，可以把$P$的定义域（函数$P(\cdot)$的自变量所在的集合）定义为$S$的所有子集；也就是说，对每个$A\subset S$，定义$P(A)$为$A$发生的概率.不幸的是，事情没有这么简单.有一些技术困难需要克服.我们不详谈这些技术细节；虽然它们很重要，但概率论学家通常比统计学家更关心它们.不过，要扎实地理解统计学，至少应对以下内容略知一二.
\begin{definition}\label{def:1.2.1}\index{SI0097@博雷尔域（见$\sigma$代数）}\index{SI0972@$\sigma$代数}
$S$的一个子集族若满足以下三个性质，就称为一个\emph{$\sigma$代数}（sigma algebra，或Borel field），记作$\mathscr B$：
\begin{enumerate}
\item[(a)]\phantomsection\label{itm:def:1.2.1:a} $\emptyset\in\mathscr B$（空集是$\mathscr B$的一个元素）.
\item[(b)]\phantomsection\label{itm:def:1.2.1:b} 若$A\in\mathscr B$，则$A^c\in\mathscr B$（$\mathscr B$对取补运算封闭）.
\item[(c)] 若$A_1,A_2,\ldots\in\mathscr B$，则$\bigcup_{i=1}^{\infty}A_i\in\mathscr B$（$\mathscr B$对可数并封闭）.
\end{enumerate}
\end{definition}
空集$\emptyset$是任何集合的子集，因此$\emptyset\subset S$.性质\hyperref[itm:def:1.2.1:a]{(a)}说明，这个子集总属于一个$\sigma$代数.因为$S=\emptyset^c$，性质\hyperref[itm:def:1.2.1:a]{(a)}与\hyperref[itm:def:1.2.1:b]{(b)}还说明$S$总属于$\mathscr B$.此外，由德摩根律可知\index{SI0247@德摩根律（见集合运算）}\semanticindexanchor{247:36}，$\mathscr B$对可数交也封闭.若$A_1,A_2,\ldots\in\mathscr B$，则由性质\hyperref[itm:def:1.2.1:b]{(b)}有$A_1^c,A_2^c,\ldots\in\mathscr B$，因此$\bigcup_{i=1}^{\infty}A_i^c\in\mathscr B$.不过，利用德摩根律（如习题\ref{exercise:1.9}所示），有
\begin{equation}\label{eq:1.2.1}
\left(\bigcup_{i=1}^{\infty}A_i^c\right)^c=\bigcap_{i=1}^{\infty}A_i.
\end{equation}
因此，再用性质\hyperref[itm:def:1.2.1:b]{(b)}，得$\bigcap_{i=1}^{\infty}A_i\in\mathscr B$.

与样本空间$S$相联系的$\sigma$代数可以有许多个.例如，由两个集合组成的集合族$\{\emptyset,S\}$就是一个$\sigma$代数，通常称为平凡$\sigma$代数.\index{SI0976@平凡}\semanticindexanchor{976:36}我们只关心包含给定样本空间$S$中所有开集的最小$\sigma$代数.
\begin{example}[$\sigma$代数——I]\label{ex:1.2.2}\index{SI0973@可数}
若$S$有限或可数，这些技术困难其实不会出现，因为对给定样本空间$S$，我们定义
$$\mathscr B=\{S\text{ 的所有子集，包括 }S\text{ 本身}\}.$$
若$S$有$n$个元素，则$\mathscr B$中有$2^n$个集合（见习题\ref{exercise:1.14}）.例如，若$S=\{1,2,3\}$，则$\mathscr B$是以下$2^3=8$个集合组成的集合族：
$$\begin{array}{ccc}
\{1\}&\{1,2\}&\{1,2,3\}\\[3pt]
\{2\}&\{1,3\}&\emptyset\\[3pt]
\{3\}&\{2,3\}&
\end{array}$$
\end{example}
一般说来，若$S$不可数，描述$\mathscr B$并不容易.不过，我们选取的$\mathscr B$会包含所有感兴趣的集合.
\begin{example}[$\sigma$代数——II]\label{ex:1.2.3}\index{SI0977@不可数}
令$S=(-\infty,\infty)$为实数轴.选取$\mathscr B$使其包含对所有实数$a,b$所形成的下列各类集合：
$$[a,b],\quad(a,b],\quad(a,b),\quad\text{和}\quad[a,b).$$
此外，由$\mathscr B$的性质可知，它还包含对上述各类集合取并和交（可能为可数无穷次）所能得到的一切集合.
\end{example}
现在可以定义概率函数了.

\sourcepageanchor{37}

\begin{definition}\label{def:1.2.4}\index{SI0596@柯尔莫哥洛夫公理}\index{SI0843@概率函数}
给定样本空间$S$及与之相联系的$\sigma$代数$\mathscr B$，\emph{概率函数}（probability function）是以$\mathscr B$为定义域、满足下列条件的函数$P$：
\begin{enumerate}[label=\arabic*.]
\item 对所有$A\in\mathscr B$，$P(A)\ge0$.
\item $P(S)=1$.
\item 若$A_1,A_2,\ldots\in\mathscr B$两两不相交，则$P(\bigcup_{i=1}^{\infty}A_i)=\sum_{i=1}^{\infty}P(A_i)$.
\end{enumerate}
\end{definition}
定义\ref{def:1.2.4}中的三个性质通常称为概率公理（Axioms of Probability），或柯尔莫哥洛夫公理（Kolmogorov Axioms），以概率论奠基人之一A. Kolmogorov命名.满足概率公理的任何函数$P$都称为概率函数.公理化定义不打算告诉我们应该选择哪个具体的函数$P$；它只要求$P$满足这些公理.对任意样本空间，都可以定义许多不同的概率函数.哪些函数能够反映某一特定试验中可能观察到的情形，还有待讨论.
\begin{example}[定义概率——I]\label{ex:1.2.5}
考虑抛一枚公平硬币的简单试验，此时$S=\{\mathrm H,\mathrm T\}$.所谓“公平”，是指硬币平衡，正面朝上和反面朝上的可能性相同.因此，合理的概率函数应赋予正面和反面相等的概率，即
\begin{equation}\label{eq:1.2.2}
P(\{\mathrm H\})=P(\{\mathrm T\}).
\end{equation}
注意，式\eqref{eq:1.2.2}并非概率公理的推论，而是公理之外的要求.我们使用了概率的对称性解释（或者仅凭直觉），要求正面与反面等可能.因为$S=\{\mathrm H\}\cup\{\mathrm T\}$，由公理2得$P(\{\mathrm H\}\cup\{\mathrm T\})=1$.又因为$\{\mathrm H\}$和$\{\mathrm T\}$不相交，$P(\{\mathrm H\}\cup\{\mathrm T\})=P(\{\mathrm H\})+P(\{\mathrm T\})$，于是
\begin{equation}\label{eq:1.2.3}
P(\{\mathrm H\})+P(\{\mathrm T\})=1.
\end{equation}
联立式\eqref{eq:1.2.2}和\eqref{eq:1.2.3}，得$P(\{\mathrm H\})=P(\{\mathrm T\})=\frac12$.

由于式\eqref{eq:1.2.2}依据的是我们对具体试验的了解，而不是公理，因此，只要$P(\{\mathrm H\})$和$P(\{\mathrm T\})$取非负值并满足式\eqref{eq:1.2.3}，就定义了一个合法的概率函数.例如，可以取$P(\{\mathrm H\})=\frac19$和$P(\{\mathrm T\})=\frac89$.
\end{example}
我们需要一些定义概率函数的一般方法，并且知道它们总能满足柯尔莫哥洛夫公理.我们不希望每引入一个新的概率函数，都像例\ref{ex:1.2.5}那样逐一检查公理.下面给出一种定义合法概率函数的常用方法.
\begin{theorem}\label{thm:1.2.6}\index{SI0844@有限集}
设$S=\{s_1,\ldots,s_n\}$是有限集，$\mathscr B$是$S$的子集构成的任意$\sigma$代数，$p_1,\ldots,p_n$是和为1的非负数.对任意$A\in\mathscr B$，定义
$$P(A)=\sum_{\{i:s_i\in A\}}p_i.$$
（空集上的求和定义为0.）则$P$是$\mathscr B$上的概率函数.若$S=\{s_1,s_2,\ldots\}$为可数集，结论仍然成立.
\end{theorem}
\begin{proof}
我们证明$S$有限的情形.对任意$A\in\mathscr B$，$P(A)=\sum_{\{i:s_i\in A\}}p_i\ge0$，因为每个$p_i\ge0$.所以公理1成立.
$$P(S)=\sum_{\{i:s_i\in S\}}p_i=\sum_{i=1}^n p_i=1.$$
所以公理2成立.令$A_1,\ldots,A_k$表示两两不相交的事件.（$\mathscr B$只含有限个集合，因此只需考虑有限个互不相交集合的并.）于是，
$$P\left(\bigcup_{i=1}^k A_i\right)=\sum_{\{j:s_j\in\bigcup_{i=1}^k A_i\}}p_j
=\sum_{i=1}^k\sum_{\{j:s_j\in A_i\}}p_j=\sum_{i=1}^k P(A_i).$$
第一个和第三个等号来自$P(A)$的定义.$A_i$互不相交保证了第二个等号，因为等号两边相同的$p_j$各出现恰好一次.因此公理3成立，柯尔莫哥洛夫公理全部满足.
\end{proof}
试验的物理实际可能决定概率如何分配，下面的例子说明了这一点.

\sourcepageanchor{38}

\begin{figure}[htbp]\centering
\includegraphics[width=.52\textwidth]{assets/chapter1/fig1-2-1.pdf}
\caption{例\ref{ex:1.2.7}中的飞镖靶}\label{fig:1.2.1}
\end{figure}
\begin{example}[定义概率——II]\label{ex:1.2.7}
飞镖游戏中，选手向靶上投掷飞镖，所得分数等于飞镖落入区域所标的数.对初学者而言，假设飞镖命中某个区域的概率与该区域面积成正比，似乎是合理的.因此，面积较大的区域被命中的概率较高.

如图\ref{fig:1.2.1}所示，飞镖靶的半径为$r$，相邻圆环之间的距离为$r/5$.若假设总能命中靶面（习题\ref{exercise:1.7}给出一个变式），则
$$P(\text{得到 }i\text{ 分})=\frac{\text{区域 }i\text{ 的面积}}{\text{飞镖靶面积}}.$$
例如，
$$P(\text{得到1分})=\frac{\pi r^2-\pi(4r/5)^2}{\pi r^2}=1-\left(\frac45\right)^2.$$
一般公式很容易导出，结果为
$$P(\text{得到 }i\text{ 分})=\frac{(6-i)^2-(5-i)^2}{5^2},\qquad i=1,\ldots,5,$$
与$\pi$和$r$无关.各个互不相交区域的面积之和等于靶面积.因此赋予这五个结果的概率之和为1，由定理\ref{thm:1.2.6}可知，这确实是一个概率函数（习题\ref{exercise:1.8}）.
\end{example}
在结束概率的公理化讨论之前，还有一点需要考虑.定义\ref{def:1.2.4}中的公理3，通常称为可数可加性公理（Axiom of Countable Additivity）\index{SI0074@可数可加性公理}\semanticindexanchor{74:38}，并未得到所有统计学家的认可.的确，人们可以认为公理应是简单且不言自明的陈述.与其他简单且不言自明的公理相比，公理3可能使我们怀疑：假定它为真是否合理？

以\hyperref[bib:source95]{De Finetti（1972）}为代表的一个统计学派拒绝可数可加性公理，选择以有限可加性公理取代它.

\noindent\emph{有限可加性公理（Axiom of Finite Additivity）：}\index{SI0075@有限可加性公理}\semanticindexanchor{75:39}若$A\in\mathscr B$与$B\in\mathscr B$不相交，则
$$P(A\cup B)=P(A)+P(B).$$
虽然这个公理未必完全不言自明，但它当然比可数可加性公理简单（并且可由后者推出，见习题\ref{exercise:1.12}）.

只假定有限可加性也许更可信，但会在统计理论中导致意想不到的复杂情况；在本书的层次上，这些复杂情况不一定有助于理解.因此，以下均假定可数可加性公理成立.

\subsubsection{概率的运算}\label{sec:1.2.2}
从概率公理出发，可以建立概率函数的许多性质，它们对计算较复杂的概率很有帮助.本节详细讨论其中一些运算，另一些留作习题.

先从概率函数应用于单个事件时的一些（相当明显的）性质开始.

\sourcepageanchor{39}

\begin{theorem}\label{thm:1.2.8}\index{SI0846@性质}
若$P$是概率函数，$A$是$\mathscr B$中的任意集合，则
\begin{enumerate}
\item[(a)]\phantomsection\label{itm:thm:1.2.8:a} $P(\emptyset)=0$，其中$\emptyset$是空集；
\item[(b)]\phantomsection\label{itm:thm:1.2.8:b} $P(A)\le1$；
\item[(c)]\phantomsection\label{itm:thm:1.2.8:c} $P(A^c)=1-P(A)$.
\end{enumerate}
\end{theorem}
\begin{proof}
先证明\hyperref[itm:thm:1.2.8:c]{(c)}最容易.集合$A$与$A^c$构成样本空间的一个划分，即$S=A\cup A^c$.因此，由第二条公理，
\begin{equation}\label{eq:1.2.4}P(A\cup A^c)=P(S)=1.\end{equation}
又因为$A$与$A^c$不相交，由第三条公理，
\begin{equation}\label{eq:1.2.5}P(A\cup A^c)=P(A)+P(A^c).\end{equation}
合并式\eqref{eq:1.2.4}和\eqref{eq:1.2.5}即得\hyperref[itm:thm:1.2.8:c]{(c)}.

因为$P(A^c)\ge0$，\hyperref[itm:thm:1.2.8:c]{(c)}立即推出\hyperref[itm:thm:1.2.8:b]{(b)}.为证明\hyperref[itm:thm:1.2.8:a]{(a)}，对$S=S\cup\emptyset$作类似论证.（回想$S$与$\emptyset$总在$\mathscr B$中.）由于$S$与$\emptyset$不相交，有
$$1=P(S)=P(S\cup\emptyset)=P(S)+P(\emptyset),$$
因此$P(\emptyset)=0$.
\end{proof}
定理\ref{thm:1.2.8}中的性质非常基本，几乎也带有公理的意味，尽管我们已仅用最初的三条柯尔莫哥洛夫公理严格证明了它们.下面的定理与定理\ref{thm:1.2.8}精神相似，但其结论并不那么显然.
\begin{theorem}\label{thm:1.2.9}
若$P$是概率函数，$A$和$B$是$\mathscr B$中的任意集合，则
\begin{enumerate}
\item[(a)]\phantomsection\label{itm:thm:1.2.9:a} $P(B\cap A^c)=P(B)-P(A\cap B)$；
\item[(b)]\phantomsection\label{itm:thm:1.2.9:b} $P(A\cup B)=P(A)+P(B)-P(A\cap B)$；
\item[(c)]\phantomsection\label{itm:thm:1.2.9:c} 若$A\subset B$，则$P(A)\le P(B)$.
\end{enumerate}
\end{theorem}
\begin{proof}
为证明\hyperref[itm:thm:1.2.9:a]{(a)}，注意对任意集合$A,B$，有
$$B=\{B\cap A\}\cup\{B\cap A^c\},$$
因此
\begin{equation}\label{eq:1.2.6}
P(B)=P(\{B\cap A\}\cup\{B\cap A^c\})=P(B\cap A)+P(B\cap A^c),
\end{equation}
其中最后一个等号来自$B\cap A$与$B\cap A^c$不相交.移项即得\hyperref[itm:thm:1.2.9:a]{(a)}.

为证明\hyperref[itm:thm:1.2.9:b]{(b)}，使用恒等式
\begin{equation}\label{eq:1.2.7}A\cup B=A\cup\{B\cap A^c\}.\end{equation}
维恩图可以说明式\eqref{eq:1.2.7}为何成立，严格证明也不困难（见习题\ref{exercise:1.2}）.利用式\eqref{eq:1.2.7}及$A$与$B\cap A^c$不相交这一事实（因为$A$与$A^c$不相交），由\hyperref[itm:thm:1.2.9:a]{(a)}得到
\begin{equation}\label{eq:1.2.8}
P(A\cup B)=P(A)+P(B\cap A^c)=P(A)+P(B)-P(A\cap B).
\end{equation}
若$A\subset B$，则$A\cap B=A$.所以利用\hyperref[itm:thm:1.2.9:a]{(a)}，有
$$0\le P(B\cap A^c)=P(B)-P(A),$$
从而\hyperref[itm:thm:1.2.9:c]{(c)}成立.
\end{proof}
定理\ref{thm:1.2.9}的\hyperref[itm:thm:1.2.9:b]{(b)}给出了交事件概率的一个有用不等式.由于$P(A\cup B)\le1$，对式\eqref{eq:1.2.8}作简单移项，得
\begin{equation}\label{eq:1.2.9}\index{SI0560@Bonferroni}P(A\cap B)\ge P(A)+P(B)-1.\end{equation}
它是邦费罗尼不等式（Bonferroni's Inequality）的一个特例［\hyperref[bib:source221]{Miller（1981）}是很好的参考文献］.邦费罗尼不等式使我们可以用单个事件的概率来界定同时发生事件（即交事件）的概率.

\sourcepageanchor{40}

\begin{example}[邦费罗尼不等式]\label{ex:1.2.10}
当交事件的概率很难甚至无法计算，但我们仍希望了解该概率的大小时，邦费罗尼不等式尤其有用.设$A,B$是两个事件，各自的概率都是$.95$，则二者同时发生的概率有下界
$$P(A\cap B)\ge P(A)+P(B)-1=.95+.95-1=.90.$$
注意，除非各个事件的概率足够大，否则邦费罗尼界只是一个没有用处（但正确！）的负数.
\end{example}
本小节最后给出一个定理，其中的结论有助于处理集合族.
\begin{theorem}\label{thm:1.2.11}\index{SI0563@Boole}
若$P$是概率函数，则
\begin{enumerate}
\item[(a)]\phantomsection\label{itm:thm:1.2.11:a} 对任意划分$C_1,C_2,\ldots$，有$P(A)=\sum_{i=1}^{\infty}P(A\cap C_i)$；
\item[(b)]\phantomsection\label{itm:thm:1.2.11:b} 对任意集合$A_1,A_2,\ldots$，有$P(\bigcup_{i=1}^{\infty}A_i)\le\sum_{i=1}^{\infty}P(A_i)$.（布尔不等式，Boole's Inequality）
\end{enumerate}
\end{theorem}
\begin{proof}
由于$C_1,C_2,\ldots$构成划分，对所有$i\ne j$有$C_i\cap C_j=\emptyset$，且$S=\bigcup_{i=1}^{\infty}C_i$.于是
$$A=A\cap S=A\cap\left(\bigcup_{i=1}^{\infty}C_i\right)=\bigcup_{i=1}^{\infty}(A\cap C_i),$$

\sourcepageanchor{41}

最后一个等号来自分配律（定理\ref{thm:1.1.4}）.因此
$$P(A)=P\left(\bigcup_{i=1}^{\infty}(A\cap C_i)\right).$$
由于$C_i$互不相交，集合$A\cap C_i$也互不相交，由概率函数的性质，
$$P\left(\bigcup_{i=1}^{\infty}(A\cap C_i)\right)=\sum_{i=1}^{\infty}P(A\cap C_i),$$
故\hyperref[itm:thm:1.2.11:a]{(a)}成立.

为证明\hyperref[itm:thm:1.2.11:b]{(b)}，先构造一族互不相交的集合$A_1^*,A_2^*,\ldots$，使$\bigcup_{i=1}^{\infty}A_i^*=\bigcup_{i=1}^{\infty}A_i$.定义
$$A_1^*=A_1,\qquad A_i^*=A_i\setminus\left(\bigcup_{j=1}^{i-1}A_j\right),\qquad i=2,3,\ldots,$$
其中$A\setminus B$表示$A$中不与$B$相交的部分.用熟悉的符号表示就是$A\setminus B=A\cap B^c$.不难看出$\bigcup_{i=1}^{\infty}A_i^*=\bigcup_{i=1}^{\infty}A_i$，因此
$$P\left(\bigcup_{i=1}^{\infty}A_i\right)=P\left(\bigcup_{i=1}^{\infty}A_i^*\right)=\sum_{i=1}^{\infty}P(A_i^*),$$
最后一个等号成立是因为$A_i^*$互不相交.为说明这一点，写成
$$\begin{aligned}
A_i^*\cap A_k^*
&=\left\{A_i\setminus\left(\bigcup_{j=1}^{i-1}A_j\right)\right\}\cap\left\{A_k\setminus\left(\bigcup_{j=1}^{k-1}A_j\right)\right\}\quad\text{（$A_i^*$的定义）}\\
&=\left\{A_i\cap\left(\bigcup_{j=1}^{i-1}A_j\right)^c\right\}\cap\left\{A_k\cap\left(\bigcup_{j=1}^{k-1}A_j\right)^c\right\}\quad\text{（$\setminus$的定义）}\\
&=\left\{A_i\cap\bigcap_{j=1}^{i-1}A_j^c\right\}\cap\left\{A_k\cap\bigcap_{j=1}^{k-1}A_j^c\right\}.\quad\text{（德摩根律）}
\end{aligned}$$
若$i>k$，上式第一个交中包含集合$A_k^c$，它与$A_k$的交为空集.若$k>i$，论证相同.此外，由构造可知$A_i^*\subset A_i$，所以$P(A_i^*)\le P(A_i)$，于是
$$\sum_{i=1}^{\infty}P(A_i^*)\le\sum_{i=1}^{\infty}P(A_i),$$
故\hyperref[itm:thm:1.2.11:b]{(b)}成立.
\end{proof}
布尔不等式和邦费罗尼不等式之间有相似之处.\index{SI0559@Bonferroni与Boole的相似性}\semanticindexanchor{559:41}事实上，它们本质上是同一回事.我们本来可以用布尔不等式推出式\eqref{eq:1.2.9}.把布尔不等式用于$A^c$，得
$$P\left(\bigcup_{i=1}^n A_i^c\right)\le\sum_{i=1}^n P(A_i^c),$$
再利用$\bigcup A_i^c=(\bigcap A_i)^c$及$P(A_i^c)=1-P(A_i)$，得到
$$1-P\left(\bigcap_{i=1}^n A_i\right)\le n-\sum_{i=1}^n P(A_i).$$

\sourcepageanchor{42}

移项后变为
\begin{equation}\label{eq:1.2.10}\index{SI0560@Bonferroni}
P\left(\bigcap_{i=1}^n A_i\right)\ge\sum_{i=1}^n P(A_i)-(n-1),
\end{equation}
这就是式\eqref{eq:1.2.9}中邦费罗尼不等式的更一般形式.

\subsubsection{计数}\label{sec:1.2.3}\index{SI0219@计数方法}
计数这一基本过程，在统计学家手中可能变得相当精巧.计数方法最常用于在有限样本空间上构造概率分配，当然它们也能用来回答其他问题.
\begin{example}[彩票——I]\label{ex:1.2.12}
纽约州彩票曾多年采用如下规则：每个人可从数字$1,2,\ldots,44$中任选六个作为彩票号码.中奖号码则通过从44个数中随机选出六个来决定.要计算中奖概率，首先必须数清从44个数中能够选出多少组不同的六个数.
\end{example}
\begin{example}[锦标赛]\label{ex:1.2.13}
在单败淘汰赛（single-elimination tournament）中，例如美国网球公开赛，选手只有获胜才能晋级（这不同于双败淘汰赛或循环赛）.如果有16名参赛者，我们可能关心某位选手夺冠的路径有多少条，这里的路径指一系列对手组成的序列.
\end{example}
计数问题一般听起来很复杂，而且经常必须在许多限制条件下计数.解决方法是把问题拆成一系列容易计数的简单任务，再运用已知的任务组合规则.下面的定理是这个过程的第一步，有时称为计数基本定理（Fundamental Theorem of Counting）.
\begin{theorem}\label{thm:1.2.14}\index{SI0449@计数基本定理}
若一项工作由$k$个分开的任务组成，第$i$个任务可以用$n_i$种方式完成，$i=1,\ldots,k$，则整项工作可以用$n_1\times n_2\times\cdots\times n_k$种方式完成.
\end{theorem}
\begin{proof}
只需证明$k=2$的情形（见习题\ref{exercise:1.15}）.证明只是仔细计数.第一个任务可以用$n_1$种方式完成，而对其中每种方式，第二个任务都有$n_2$种选择.因此，完成整项工作的方式数为
$$\underbrace{(1\times n_2)+(1\times n_2)+\cdots+(1\times n_2)}_{n_1\text{ 项}}=n_1\times n_2,$$
这就证明了$k=2$时的定理.
\end{proof}
\begin{example}[彩票——II]\label{ex:1.2.15}
虽然计数基本定理提供了合理的起点，但应用中通常还有更多方面需要考虑.例如，纽约州彩票的第一个号码有44种选法，第二个有43种选法，因此前两个号码共有$44\times43=1,892$种选法.不过，如果允许同一个号码选两次，那么前两个号码就有$44\times44=1,936$种选法.
\end{example}
例\ref{ex:1.2.15}区分了\emph{有放回}（with replacement）与\emph{无放回}（without replacement）计数.任何计数问题还有第二个关键因素：任务的次序是否重要.仍以彩票为例，假设中奖号码按$12,37,35,9,13,22$的顺序选出，那么选了$9,12,13,22,35,37$的人是否中奖？换言之，完成任务的顺序究竟是否重要？把这些因素都考虑进去，可以构成一个$2\times2$的可能情形表：
\begin{center}
\emph{可能的计数方法}\\[3pt]
\begin{tabular}{c|c|c|}
 & 无放回 & 有放回\\\cline{2-3}
有序 & \hspace{2cm} & \hspace{2cm}\\\cline{2-3}
无序 & & \\\cline{2-3}
\end{tabular}
\end{center}

\sourcepageanchor{43}

开始计数之前，下面的定义提供了非常有用的记号.
\begin{definition}\label{def:1.2.16}\index{SI0442@阶乘}
对正整数$n$，$n!$（读作$n$的阶乘，$n$ factorial）是所有小于或等于$n$的正整数之积，即
$$n!=n\times(n-1)\times(n-2)\times\cdots\times3\times2\times1.$$
此外，定义$0!=1$.
\end{definition}
现在分别考虑四种情形下所有可能彩票的计数.
\begin{enumerate}[label=\arabic*.]
\item \emph{有序、无放回.}由计数基本定理，第一个号码有44种选法，第二个有43种选法，依此类推.因此，可能的彩票共有
$$44\times43\times42\times41\times40\times39=\frac{44!}{38!}=5,082,517,440$$
张.
\item \emph{有序、有放回.}由于每个号码现在都可用44种方式选取（选中的号码被放回），可能的彩票共有
$$44\times44\times44\times44\times44\times44=44^6=7,256,313,856$$
张.
\item \emph{无序、无放回.}我们已知道必须考虑次序时的彩票数，因此要做的是除去重复的排列.再次运用基本定理，六个号码可按$6\times5\times4\times3\times2\times1$种方式排列，所以无序彩票总数为
$$\frac{44\times43\times42\times41\times40\times39}{6\times5\times4\times3\times2\times1}=\frac{44!}{6!38!}=7,059,052.$$
\end{enumerate}
这种计数方式在统计学的许多领域中具有核心地位，因此有了专门的记号.
\begin{definition}\label{def:1.2.17}\index{SI0091@二项式系数}
对非负整数$n,r$，$n\ge r$，定义符号$\binom nr$（读作“$n$选$r$”，$n$ choose $r$）为
$$\binom nr=\frac{n!}{r!(n-r)!}.$$
\end{definition}
在彩票例子中，无序、无放回的可能彩票数是$\binom{44}{6}$.这些数也称为\emph{二项式系数}（binomial coefficients），其原因将在\hyperref[sec:3]{第3章}中说明.
\begin{enumerate}[label=\arabic*.]
\setcounter{enumi}{3}
\item \emph{无序、有放回.}这是最难计数的一种情形.你也许首先猜想答案是$44^6/(6\times5\times4\times3\times2\times1)$，但这不正确（它太小了）.


\sourcepageanchor{44}

对此计数时，最简单的想法是在44个数上放置六个标记.实际上，可以把44个数看作44个格子，把六个标记M放入其中，例如下图所示.
\begin{center}\includegraphics[width=.88\textwidth]{assets/chapter1/bins.pdf}\end{center}
可能的彩票数等于把六个标记放入44个格子的方法数.不过还可以进一步化简：只需记录标记与格子隔墙的排列.再注意，最外侧的两面墙不起作用.因此必须数清43面墙（44个格子产生45面墙，但忽略两端的墙）和六个标记的所有排列.一共有$43+6=49$个对象，可按$49!$种方式排列.然而，为除去重复排列，必须同时除以$6!$和$43!$，所以排列总数为
$$\frac{49!}{6!43!}=13,983,816.$$
\end{enumerate}
虽然上述推导都是通过一个例子完成的，但不难看出它们在一般情形下也成立.为求完整，表\ref{tab:1.2.1}汇总了这些情况.
\begin{table}[htbp]\centering
\caption{从$n$个对象中取$r$个的各种可能排列数}\label{tab:1.2.1}\index{SI0218@计数公式}
\renewcommand{\arraystretch}{1.8}
\begin{tabular}{ccc}\hline
 & 无放回 & 有放回\\
有序 & $\displaystyle\frac{n!}{(n-r)!}$ & $n^r$\\
无序 & $\displaystyle\binom nr$ & $\displaystyle\binom{n+r-1}{r}$\\\hline
\end{tabular}
\end{table}

\subsubsection{枚举结果}\label{sec:1.2.4}
当样本空间$S$是有限集，并且$S$中的所有结果等可能时，上一小节的计数技巧很有用.此时，只需数清事件中的结果数，就可以计算事件概率.为说明这一点，设$S=\{s_1,\ldots,s_N\}$是有限样本空间.所有结果等可能意味着对每个结果$s_i$都有$P(\{s_i\})=1/N$.由概率函数定义中的公理3，对任意事件$A$，有
$$P(A)=\sum_{s_i\in A}P(\{s_i\})=\sum_{s_i\in A}\frac1N=\frac{A\text{ 中的元素个数}}{S\text{ 中的元素个数}}.$$
对大型样本空间，可以使用计数技巧分别计算该式的分子和分母.
\begin{example}[扑克]\label{ex:1.2.18}
考虑从一副52张的标准扑克牌中选取一手五张牌.显然，这是无放回抽样.但为了明确可能结果（可能的牌型），必须决定把发牌看作逐张发出（有序）还是一次发出（无序）.如果要计算的事件依赖次序，例如前两张牌中出现一张A的概率，就必须使用有序结果.若事件不依赖次序，则可以使用无序结果.本例使用无序结果，因此样本空间由从52张牌中选出的所有五张牌组合构成.可能的牌组有$\binom{52}{5}=2,598,960$种.如果洗牌充分且随机发牌，则赋予每种可能牌组概率$1/2,598,960$是合理的.

现在通过计数事件中的结果来计算一些概率.拿到四张A的概率是多少？含四张A的牌组有多少种？四张牌已指定为A后，第五张牌有48种选择.因此，
$$P(\text{四张A})=\frac{48}{2,598,960},$$
概率小于五万分之一.使用定理\ref{thm:1.2.14}，只需稍微复杂一点的计数，就能计算拿到四张同点数牌的概率.四张同点数牌的点数有13种选择；确定这四张牌后，第五张牌有48种选择.因此，含四张同点数牌的牌组总数为$(13)(48)$，并且
$$P(\text{四张同点数牌})=\frac{(13)(48)}{2,598,960}=\frac{624}{2,598,960}.$$
要计算恰好一对的概率（不是两对，也不是三张同点数等），需要组合使用几种计数技巧.恰好一对的牌组数为
\begin{equation}\label{eq:1.2.11}
13\binom42\binom{12}{3}4^3=1,098,240.
\end{equation}
式\eqref{eq:1.2.11}来自定理\ref{thm:1.2.14}，因为
$$\begin{aligned}
13&=\text{指定一对的点数的方法数},\\
\binom42&=\text{从该点数中指定两张牌的方法数},\\
\binom{12}{3}&=\text{指定其余三个点数的方法数},\\
4^3&=\text{从这些点数中指定其余三张牌的方法数}.
\end{aligned}$$
因此，
$$P(\text{恰好一对})=\frac{1,098,240}{2,598,960}.$$
\end{example}
在例\ref{ex:1.2.18}这样的无放回抽样中，如果要计算不依赖次序的事件概率，则有序和无序样本空间都可以使用.无序样本空间中的每个结果对应有序样本空间中的$r!$个结果.因此，在有序样本空间中计数时，分子、分母各有一个$r!$因子，相消后得到的概率与在无序样本空间中计数相同.

有放回抽样则不同.无序样本空间中的每个结果对应有序样本空间中的一些结果，但所对应的个数各不相同.

\sourcepageanchor{45}

\begin{example}[有放回抽样]\label{ex:1.2.19}
考虑从$n=3$个对象中有放回地抽取$r=2$个对象.有序和无序样本空间中的结果如下：
\begin{center}\small
\begin{tabular}{ccccccc}\hline
无序 & $\{1,1\}$ & $\{2,2\}$ & $\{3,3\}$ & $\{1,2\}$ & $\{1,3\}$ & $\{2,3\}$\\
有序 & $(1,1)$ & $(2,2)$ & $(3,3)$ & $(1,2),(2,1)$ & $(1,3),(3,1)$ & $(2,3),(3,2)$\\
概率 & $1/9$ & $1/9$ & $1/9$ & $2/9$ & $2/9$ & $2/9$\\\hline
\end{tabular}
\end{center}
这些概率是把有序样本空间中的九个结果视为等可能而得到的.这对应于“有放回抽样”的通常解释：三个对象各以$1/3$的概率被选中；记录该对象后将其放回；混匀后再次从三个对象中选择一个，每个仍以$1/3$的概率被选中.可以看出，在这种抽样方式下，无序样本空间中的六个结果并非等可能.无序样本空间的结果数公式可用于枚举结果，但为了正确计算概率，必须对有序结果计数.
\end{example}
有些作者认为，在“把$r$个不可区分的球随机分配到$n$个可区分的罐子中”时，应给无序结果赋予相等概率.也就是说，随机选一个罐子并放入一球，重复$r$次.不记录放球的顺序，因此最终的结果$\{1,3\}$表示1号罐和3号罐各有一球.

但这种解释存在问题.假设两个人观察同一过程，观察者1记录放球顺序，观察者2不记录.观察者1会给事件$\{1,3\}$赋予概率$2/9$.观察者2看到的是完全相同的过程，也应给这个事件赋予概率$2/9$.然而，若把六个无序结果写在相同的纸片上，随机抽一片来决定球的放置方式，那么每个无序结果的概率就是$1/6$.于是观察者2会给事件$\{1,3\}$赋予概率$1/6$.

混淆的原因在于，“有放回”通常按上述依次抽样的方式理解，因此给事件$\{1,3\}$赋予概率$2/9$.这是正确的处理方式，因为概率应由抽样机制决定，而不是由球可区分还是不可区分来决定.

\sourcepageanchor{46}

\begin{example}[计算平均数]\label{ex:1.2.20}
为说明可区分／不可区分的处理方法，假设要计算从四个数
$$2,4,9,12$$
中有放回抽取四个数所得的所有可能平均数.例如，可能抽得$\{2,4,4,9\}$，平均数为4.75；也可能抽得$\{4,4,9,9\}$，平均数为6.5.如果只关心所抽数值的平均数，次序便不重要，因此不同样本总数可按无序、有放回抽样的方式计数.

不同样本总数为$\binom{n+n-1}{n}$.不过，要计算样本平均数的概率分布，必须数清某个平均数可以出现的不同方式.只有样本包含一个2、两个4和一个9时，才会出现4.75这个值.具有这种构成的可能样本如下表所示.
\begin{center}
\begin{tabular}{cl}
无序 & 有序\\\hline
$\{2,4,4,9\}$ & $(2,4,4,9),(2,4,9,4),(2,9,4,4),(4,2,4,9),$\\
 & $(4,2,9,4),(4,4,2,9),(4,4,9,2),(4,9,2,4),$\\
 & $(4,9,4,2),(9,2,4,4),(9,4,2,4),(9,4,4,2)$
\end{tabular}
\end{center}
有序样本总数为$n^n=4^4=256$，因此抽得无序样本$\{2,4,4,9\}$的概率为$12/256$.与此相比，如果认为无序样本等可能，就会给$\{2,4,4,9\}$及每个其他无序样本赋予概率$1/\binom{n+n-1}{n}=1/\binom74=1/35$.

为计算会产生$\{2,4,4,9\}$的有序样本数，可作如下论证.需要枚举四个数$\{2,4,4,9\}$的所有可能顺序，因此本质上属于第\ref{sec:1.2.3}节的计数方法1.样本可按$4\times3\times2\times1=24$种方式排序.但这里有重复计数，因为两个4的不同排列无法区分.例如，24种方式会把$\{9,4,2,4\}$数两次（若两个4不同，这样计数就没有问题）.为纠正这一点，除以$2!$（两个4有$2!$种排列），得到$24/2=12$个有序样本.一般说来，若有$k$个位置，$m$个不同的数分别重复$k_1,k_2,\ldots,k_m$次，则有序样本数为
$$\frac{k!}{k_1!k_2!\cdots k_m!}.$$
这种计数与\emph{多项分布}（multinomial distribution）有关，我们将在\hyperref[sec:4.6]{第4.6节}中讨论.图\ref{fig:1.2.2}是样本平均数概率分布的直方图，反映了对样本所作的多项式计数.\footnote{译注：原图注明写总体$\{2,4,4,9\}$，正文及直方图对应的却是$\{2,4,9,12\}$；下方保留原图注.}

\sourcepageanchor{47}

\begin{figure}[htbp]\centering
\includegraphics[width=.88\textwidth]{assets/chapter1/fig1-2-2.pdf}
\caption{从四个数$\{2,4,4,9\}$中有放回抽样所得平均数的直方图}\label{fig:1.2.2}
\end{figure}

图\ref{fig:1.2.2}还反映了一个需要进一步细化之处：两个不同的无序样本可能具有相同的平均数.例如，无序样本$\{4,4,12,12\}$和$\{2,9,9,12\}$的平均数都是8.第一个样本的概率为$3/128$，第二个样本的概率为$3/64$，因此取值8的概率为$9/128=.07$.\footnote{译注：原文使用等号；精确值为$0.0703125$，$.07$只是近似值.}关于如何构造这样的直方图，见\hyperref[ex:12.6.1]{例12.6.1}.本例中的计算是一种非常重要的统计技术——\emph{自助法}（bootstrap，\hyperref[bib:source114]{Efron与Tibshirani，1993}）的初级形式.我们将在\hyperref[sec:10.1.4]{第10.1.4节}再次讨论自助法.
\end{example}

% 用户授权疑误脚注：2026-11-04



% 原书印刷16–22页，PDF物理47–53页；从§1.3起至§1.4前.
\subsection{条件概率与独立性}\label{sec:1.3}
到目前为止，我们讨论的概率都是无条件概率.我们先定义一个样本空间，再相对于该样本空间计算所有概率.不过，在许多情形中，可以依据新的信息更新样本空间.此时，我们希望能够更新概率计算，或者计算条件概率（conditional probability）.

\begin{example}[四张A]\label{ex:1.3.1}
从充分洗匀的一副扑克牌顶端发出四张牌.这四张恰好是四张A的概率是多少？可以用上一节的方法计算.由四张牌组成的不同组合共有
$$\binom{52}{4}=270,725.$$
其中只有一个组合由四张A组成，而且每个组合等可能，因此发到四张A的概率为$1/270,725$.

也可以用如下的“更新”论证来计算这一概率.第一张牌是A的概率为$4/52$.已知第一张牌是A，第二张牌是A的概率为$3/51$（还剩三张A和51张牌）.继续这一论证，所求概率为
$$\frac4{52}\times\frac3{51}\times\frac2{50}\times\frac1{49}=\frac1{270,725}.$$
\end{example}

在第二种解法中，每抽出一张牌，我们就更新一次样本空间；我们计算的是条件概率.

\sourcepageanchor{48}

\begin{definition}\label{def:1.3.2}\index{SI0122@条件概率}
若$A$和$B$是$S$中的事件，且$P(B)>0$，则给定$B$时$A$的条件概率记为$P(A\mid B)$，定义为
\begin{equation}\label{eq:1.3.1}
P(A\mid B)=\frac{P(A\cap B)}{P(B)}.
\end{equation}
\end{definition}

注意，在条件概率计算中，$B$成了样本空间：$P(B\mid B)=1$.直观地说，原来的样本空间$S$已更新为$B$，此后所有事件的发生都按它们与$B$的关系重新衡量.尤其要注意不相交集合的条件概率.设$A$与$B$不相交，则$P(A\cap B)=0$，从而$P(A\mid B)=P(B\mid A)=0$.

\begin{example}[续例\ref{ex:1.3.1}]\label{ex:1.3.3}
虽然得到四张A的概率很小，但我们来看看，已知已经得到若干张A时，条件概率会怎样变化.仍从充分洗匀的一副牌中发出四张牌，现在计算
$$P(\text{四张牌中有四张A}\mid\text{$i$张牌中有$i$张A}),\qquad i=1,2,3.$$
事件$\{\text{四张牌中有四张A}\}$是事件$\{\text{$i$张牌中有$i$张A}\}$的子集.因此，由条件概率的定义\eqref{eq:1.3.1}，有
\begin{equation*}
\begin{aligned}
&P(\text{四张牌中有四张A}\mid\text{$i$张牌中有$i$张A})\\
&\quad=\frac{P(\{\text{四张牌中有四张A}\}\cap\{\text{$i$张牌中有$i$张A}\})}{P(\text{$i$张牌中有$i$张A})}\\
&\quad=\frac{P(\text{四张牌中有四张A})}{P(\text{$i$张牌中有$i$张A})}.
\end{aligned}
\end{equation*}
分子已经算出，分母也可以用类似的论证计算.由$i$张牌组成的不同组合共有$\binom{52}{i}$个，且
$$P(\text{$i$张牌中有$i$张A})=\frac{\binom4i}{\binom{52}i}.$$
因此，条件概率为
\begin{equation*}
\begin{aligned}
&P(\text{四张牌中有四张A}\mid\text{$i$张牌中有$i$张A})\\
&\quad=\frac{\binom{52}i}{\binom{52}4\binom4i}
=\frac{(4-i)!48!}{(52-i)!}=\frac1{\binom{52-i}{4-i}}.
\end{aligned}
\end{equation*}
当$i=1,2,3$时，条件概率分别为$.00005$、$.00082$和$.02041$.
\end{example}

对任何满足$P(B)>0$的$B$，直接验证即可知道概率函数$P(\cdot\mid B)$满足Kolmogorov公理\index{SI0597@条件概率}\semanticindexanchor{597:49}（见习题\ref{exercise:1.35}）.你也许觉得$P(B)>0$这个要求是多余的：谁会想在概率为零的事件上取条件呢？有趣的是，有时这恰恰是一种特别有用的思考方式.不过，我们将把这些问题留到\hyperref[sec:4]{第4章}.

条件概率有时很容易令人误判，需要仔细思考.请看下面这个经常讲述的故事.

\sourcepageanchor{49}

\begin{example}[三名囚犯]\label{ex:1.3.4}\index{SI1071@看守问题}
三名囚犯A、B、C正在等候执行死刑.州长决定赦免其中一人，并随机选出要赦免的囚犯.他把选择告诉典狱长，但要求对姓名保密几天.

第二天，A试图让典狱长说出谁已获赦免，典狱长拒绝了.A于是问B和C中谁将被处死.典狱长想了一会儿，随后告诉A，B将被处死.
\begin{description}
\item[典狱长的推理：]每名囚犯获赦免的机会都是$\frac13$.显然，B和C中至少有一人必须被处死，因此我并没有给A提供任何关于A是否会获赦免的信息.
\item[A的推理：]已知B将被处死，获赦免的就只可能是A或C.我获赦免的机会已上升到$\frac12$.
\end{description}
典狱长的推理应当是正确的，不过我们来看看原因.令$A$、$B$、$C$分别表示A、B、C获赦免的事件.我们知道$P(A)=P(B)=P(C)=\frac13$.令$\mathscr W$表示典狱长说B将死去这一事件.利用\eqref{eq:1.3.1}，A可将自己获赦免的概率更新为
$$P(A\mid\mathscr W)=\frac{P(A\cap\mathscr W)}{P(\mathscr W)}.$$
整个情形可以归纳为下表：
\begin{center}
\begin{tabular}{ccc}\hline
获赦免的囚犯&典狱长告诉A&\\\hline
A&B将死去&\multirow{2}{*}{两种说法等概率}\\
A&C将死去&\\
B&C将死去&\\
C&B将死去&\\\hline
\end{tabular}
\end{center}
根据该表，可得
\begin{equation*}\begin{aligned}
P(\mathscr W)&=P(\text{典狱长说B将死去})\\
&=P(\text{典狱长说B将死去且A获赦免})\\
&\quad+P(\text{典狱长说B将死去且C获赦免})\\
&\quad+P(\text{典狱长说B将死去且B获赦免})\\
&=\frac16+\frac13+0=\frac12.
\end{aligned}\end{equation*}
因此，按典狱长的推理，
\begin{equation}\label{eq:1.3.2}\begin{aligned}
P(A\mid\mathscr W)&=\frac{P(A\cap\mathscr W)}{P(\mathscr W)}\\
&=\frac{P(\text{典狱长说B将死去且A获赦免})}{P(\text{典狱长说B将死去})}
=\frac{1/6}{1/2}=\frac13.
\end{aligned}\end{equation}
然而，A错误地把事件$\mathscr W$当成事件$B^c$，并计算
$$P(A\mid B^c)=\frac{P(A\cap B^c)}{P(B^c)}=\frac{1/3}{2/3}=\frac12.$$
可见，条件概率很容易令人误判，需要仔细解释.这个问题的其他变体见习题\ref{exercise:1.37}.
\end{example}

将\eqref{eq:1.3.1}改写，便得到计算交事件概率的一个有用形式：
\begin{equation}\label{eq:1.3.3}\index{SI0123@乘法法则}P(A\cap B)=P(A\mid B)P(B),\end{equation}
它实质上就是例\ref{ex:1.3.1}所用的公式.利用\eqref{eq:1.3.3}的对称性，也可写成
\begin{equation}\label{eq:1.3.4}P(A\cap B)=P(B\mid A)P(A).\end{equation}
面对看起来很困难的计算时，可以按\eqref{eq:1.3.3}或\eqref{eq:1.3.4}分解计算，选其中较容易的一种.进一步令这两个式子的右端相等，整理可得
\begin{equation}\label{eq:1.3.5}P(A\mid B)=P(B\mid A)\frac{P(A)}{P(B)},\end{equation}
这给出了“反转”条件概率的公式.式\eqref{eq:1.3.5}通常以其发现者Thomas Bayes爵士命名，称为贝叶斯公式（Bayes' Rule）［不过，参见\hyperref[bib:source306]{Stigler（1983）}］.

贝叶斯公式还有比\eqref{eq:1.3.5}更一般的形式，可以用于样本空间的划分.因此，我们以下面的形式作为贝叶斯公式的定义.

\sourcepageanchor{50}

\begin{theorem}[贝叶斯公式]\label{thm:1.3.5}\index{SI0086@贝叶斯公式}\index{SI0753@划分}
设$A_1,A_2,\ldots$是样本空间的一个划分，$B$为任意集合.则对每个$i=1,2,\ldots$，\footnote{译注：按本章的条件概率比值定义，须有$P(B)>0$，且使用$P(B\mid A_j)$的各项须有$P(A_j)>0$；零概率分块应剔除或另约定其贡献为0.原定理没有显写这些条件.}
$$P(A_i\mid B)=\frac{P(B\mid A_i)P(A_i)}{\sum_{j=1}^\infty P(B\mid A_j)P(A_j)}.$$
\end{theorem}

\begin{example}[编码]\label{ex:1.3.6}
发送编码信息时，传输过程中有时会出错.具体来说，莫尔斯电码使用“点”和“划”，已知它们出现的比例是$3:4$.这意味着，对任意一个符号，
$$P(\text{发送点})=\frac37,\qquad P(\text{发送划})=\frac47.$$
假设传输线路受到干扰，一个点以$\frac18$的概率被误收为划，反之亦然.如果接收到一个点，能否确定发出的就是点？利用贝叶斯公式，可写成
$$P(\text{发送点}\mid\text{接收点})=P(\text{接收点}\mid\text{发送点})\frac{P(\text{发送点})}{P(\text{接收点})}.$$
根据给定信息，$P(\text{发送点})=\frac37$，且$P(\text{接收点}\mid\text{发送点})=\frac78$.此外，
\begin{equation*}\begin{aligned}
P(\text{接收点})&=P(\text{接收点}\cap\text{发送点})+P(\text{接收点}\cap\text{发送划})\\
&=P(\text{接收点}\mid\text{发送点})P(\text{发送点})\\
&\quad+P(\text{接收点}\mid\text{发送划})P(\text{发送划})\\
&=\frac78\times\frac37+\frac18\times\frac47=\frac{25}{56}.
\end{aligned}\end{equation*}
合并上述结果，可知正确接收到点的概率为
$$P(\text{发送点}\mid\text{接收点})=\frac{(7/8)\times(3/7)}{25/56}=\frac{21}{25}.$$
\end{example}

有些时候，某个事件$B$的发生对另一事件$A$的概率没有影响.用符号表示，就是
\begin{equation}\label{eq:1.3.6}P(A\mid B)=P(A).\end{equation}
若此式成立，则由贝叶斯公式\eqref{eq:1.3.5}及\eqref{eq:1.3.6}，有
\begin{equation}\label{eq:1.3.7}P(B\mid A)=P(A\mid B)\frac{P(B)}{P(A)}=P(A)\frac{P(B)}{P(A)}=P(B),\end{equation}
所以$A$的发生也不影响$B$.又因$P(B\mid A)P(A)=P(A\cap B)$，故
$$P(A\cap B)=P(A)P(B),$$
我们以此定义统计独立性（statistical independence）.

\sourcepageanchor{51}

\begin{definition}\label{def:1.3.7}\index{SI0550@事件}
若两个事件$A$、$B$满足
\begin{equation}\label{eq:1.3.8}P(A\cap B)=P(A)P(B),\end{equation}
则称它们统计独立.
\end{definition}
注意，也可以等价地用\eqref{eq:1.3.6}或\eqref{eq:1.3.7}来定义独立性（只要$P(A)>0$或\allowbreak\mbox{$P(B)>0$）.}\allowbreak 式\eqref{eq:1.3.8}的优点在于它对称地处理两个事件，也更容易推广到两个以上的事件.

许多赌博游戏都提供了独立事件的模型.旋转轮盘或投掷一对骰子，都形成独立事件的序列.
\begin{example}[德梅雷骑士]\label{ex:1.3.8}\index{SI0686@德梅雷骰子问题}
本章开头介绍的赌徒德梅雷骑士，特别感兴趣的是投掷一枚骰子四次至少出现一次6的事件.有
\begin{equation*}\begin{aligned}
P(\text{四次中至少一次为6})&=1-P(\text{四次都没有6})\\
&=1-\prod_{i=1}^4 P(\text{第$i$次没有6}),
\end{aligned}\end{equation*}
最后一个等号来自各次投掷的独立性.每次没有投出6的概率为$\frac56$，所以
$$P(\text{四次中至少一次为6})=1-\left(\frac56\right)^4=.518.$$
\end{example}

$A$与$B$独立还蕴含其补集之间的独立性.事实上，有如下定理.
\begin{theorem}\label{thm:1.3.9}\index{SI0552@性质}
若$A$和$B$是独立事件，则以下各对也独立：
\begin{enumerate}[label=\alph*.]
\item\label{item:1.3.9a}$A$和$B^c$；
\item\label{item:1.3.9b}$A^c$和$B$；
\item\label{item:1.3.9c}$A^c$和$B^c$.
\end{enumerate}
\end{theorem}
\begin{proof}
只证明\hyperref[item:1.3.9a]{(a)}，其余留作习题\ref{exercise:1.40}.要证明\hyperref[item:1.3.9a]{(a)}，必须说明$P(A\cap B^c)=P(A)P(B^c)$.由定理\ref{thm:1.2.9}\hyperref[itm:thm:1.2.9:a]{(a)}，
\begin{equation*}\begin{aligned}
P(A\cap B^c)&=P(A)-P(A\cap B)\\
&=P(A)-P(A)P(B)&&\text{（$A$与$B$独立）}\\
&=P(A)(1-P(B))\\
&=P(A)P(B^c).
\end{aligned}\end{equation*}
\end{proof}

两个以上事件的独立性可以用类似\eqref{eq:1.3.8}的方式定义，不过必须谨慎.例如，或许会认为只要$P(A\cap B\cap C)=P(A)P(B)P(C)$，就可以说$A$、$B$、$C$独立.然而，这并不是正确的条件.

\sourcepageanchor{52}

\begin{example}[投掷两枚骰子]\label{ex:1.3.10}
考虑投掷两枚骰子的试验.其样本空间为
$$S=\{(1,1),(1,2),\ldots,(1,6),(2,1),\ldots,(2,6),\ldots,(6,1),\ldots,(6,6)\},$$
即$S$由数$1$到$6$组成的36个有序对构成.定义事件
\begin{equation*}\begin{aligned}
A&=\{\text{出现相同点数}\}=\{(1,1),(2,2),(3,3),(4,4),(5,5),(6,6)\},\\
B&=\{\text{点数之和在7和10之间}\},\\
C&=\{\text{点数之和为2、7或8}\}.
\end{aligned}\end{equation*}
在36个可能结果中计数，即可算出各概率.有
$$P(A)=\frac16,\qquad P(B)=\frac12,\qquad P(C)=\frac13.$$
此外，
\begin{equation*}\begin{aligned}
P(A\cap B\cap C)&=P(\text{点数之和为8，且两枚都是4})\\
&=\frac1{36}=\frac16\times\frac12\times\frac13=P(A)P(B)P(C).
\end{aligned}\end{equation*}
但是
$$P(B\cap C)=P(\text{点数之和等于7或8})=\frac{11}{36}\ne P(B)P(C).$$
同样，可以证明$P(A\cap B)\ne P(A)P(B)$.因此，要求$P(A\cap B\cap C)=P(A)P(B)P(C)$，还不足以保证两两独立.
\end{example}

鉴于前例，第二次尝试可能是：若所有事件对都独立，就定义$A$、$B$、$C$独立.可惜，这个条件仍然不行.
\begin{example}[字母]\label{ex:1.3.11}
令样本空间$S$由字母a、b、c的$3!$种排列，再加上每个字母重复三次所得的三个三元组组成.即
$$S=\left\{\begin{array}{ccc}\mathrm{aaa}&\mathrm{bbb}&\mathrm{ccc}\\\mathrm{abc}&\mathrm{bca}&\mathrm{cba}\\\mathrm{acb}&\mathrm{bac}&\mathrm{cab}\end{array}\right\}.$$
并令$S$中每个元素的概率都是$\frac19$.定义
$$A_i=\{\text{三元组的第$i$个位置是a}\}.$$
简单计数即得
$$P(A_i)=\frac13,\qquad i=1,2,3,$$
以及
$$P(A_1\cap A_2)=P(A_1\cap A_3)=P(A_2\cap A_3)=\frac19,$$
所以这些$A_i$两两独立.但是
$$P(A_1\cap A_2\cap A_3)=\frac19\ne P(A_1)P(A_2)P(A_3),$$
故这些$A_i$并不满足所需的概率条件.
\end{example}

前两个例子说明，一组事件的同时独立或相互独立需要一个很强的定义.下面的定义是可行的.

\sourcepageanchor{53}

\begin{definition}\label{def:1.3.12}\index{SI0551@相互}
事件族$A_1,\ldots,A_n$称为相互独立（mutually independent），如果对任意子族$A_{i_1},\ldots,A_{i_k}$都有
$$P\left(\bigcap_{j=1}^k A_{i_j}\right)=\prod_{j=1}^k P(A_{i_j}).$$
\end{definition}

\begin{example}[抛硬币三次——I]\label{ex:1.3.13}
考虑抛一枚硬币三次的试验.一个样本点必须说明每次投掷的结果.例如，HHT可以表示先观察到两次正面，再观察到一次反面.本试验的样本空间有八个点，即
$$\{\mathrm{HHH,HHT,HTH,THH,TTH,THT,HTT,TTT}\}.$$
令$H_i$（$i=1,2,3$）表示第$i$次为正面的事件.例如，
\begin{equation}\label{eq:1.3.9}H_1=\{\mathrm{HHH,HHT,HTH,HTT}\}.\end{equation}
若给每个样本点赋予概率$\frac18$，则利用\eqref{eq:1.3.9}这样的枚举，可以看出$P(H_1)=P(H_2)=P(H_3)=\frac12$.这意味着硬币是公平的，每次落地时正反面的概率相同.

在这个概率模型中，事件$H_1,H_2,H_3$也相互独立.为验证这一点，注意
$$P(H_1\cap H_2\cap H_3)=P(\{\mathrm{HHH}\})=\frac18=\frac12\cdot\frac12\cdot\frac12=P(H_1)P(H_2)P(H_3).$$
为验证定义\ref{def:1.3.12}的条件，还必须检查每一对.例如，
$$P(H_1\cap H_2)=P(\{\mathrm{HHH,HHT}\})=\frac28=\frac12\cdot\frac12=P(H_1)P(H_2).$$
另外两对也满足这一等式.因此，$H_1,H_2,H_3$相互独立，也就是说，任一次出现正面对其他各次投掷都没有影响.

可以验证，给每个样本点赋予概率$\frac18$，是唯一使$P(H_1)=P(H_2)=P(H_3)=\frac12$且$H_1,H_2,H_3$相互独立的概率模型.
\end{example}

% 用户授权疑误脚注：2026-11-04



% 原书印刷22–24页，PDF物理53–55页.
\subsection{随机变量}\label{sec:1.4}
在许多试验中，处理一个概括性变量比处理原来的概率结构容易.例如，在民意调查中，我们可能决定询问50个人是否赞同某个议题.若用“1”记录赞同，用“0”记录不赞同，则这个试验的样本空间有$2^{50}$个元素，每个元素都是由1和0组成的长度为50的有序串.我们应当能把它缩减到合理的规模！也许我们关心的唯一数量，就是50人中赞同的人数（等价地，不赞同的人数）.如果定义变量$X=\text{50次记录中1的个数}$，就抓住了问题的实质.注意，$X$的样本空间为整数集合$\{0,1,2,\ldots,50\}$，它比原样本空间容易处理得多.

定义量$X$，就定义了从原样本空间到一个新样本空间的映射（函数）；这个新空间通常是实数的某个集合.一般地，有如下定义.

\sourcepageanchor{54}

\begin{definition}\label{def:1.4.1}\index{SI0864@随机变量}
随机变量（random variable）是从样本空间$S$到实数集的函数.
\end{definition}

\begin{example}[随机变量]\label{ex:1.4.2}
在一些试验中，随机变量被隐含地使用.下面列举若干例子.
\begin{center}
随机变量的例子\par\smallskip
\begin{tabular}{p{.40\linewidth}p{.47\linewidth}}\hline
试验&随机变量\\\hline
投掷两枚骰子&$X=\text{点数之和}$\\
抛一枚硬币25次&$X=\text{25次中正面的次数}$\\
对玉米植株施用不同量的肥料&$X=\text{每英亩产量}$\\\hline
\end{tabular}
\end{center}
\end{example}

在定义随机变量时，我们也定义了一个新的样本空间（随机变量的值域）.现在必须正式检查，定义在原样本空间上的概率函数，能否用于这个随机变量.

设样本空间
$$S=\{s_1,\ldots,s_n\}$$
上有概率函数$P$，并定义随机变量$X$，其值域为$\mathscr X=\{x_1,\ldots,x_m\}$.可按下述方式在$\mathscr X$上定义概率函数$P_X$：观察到$X=x_i$，当且仅当随机试验的结果是满足$X(s_j)=x_i$的某个$s_j\in S$.因此，
\begin{equation}\label{eq:1.4.1}\index{SI0598@诱导概率}\index{SI0845@诱导}
P_X(X=x_i)=P(\{s_j\in S:X(s_j)=x_i\}).
\end{equation}
注意，式\eqref{eq:1.4.1}左边的函数$P_X$是$\mathscr X$上的诱导概率函数，它用原来的函数$P$定义.式\eqref{eq:1.4.1}正式给随机变量$X$定义了概率函数$P_X$.当然，还需验证$P_X$满足Kolmogorov公理，不过这并不困难（见习题\ref{exercise:1.45}）.由于\eqref{eq:1.4.1}中的等价关系，今后将简写为$P(X=x_i)$，而不写$P_X(X=x_i)$.

\emph{记号说明：}随机变量始终用大写字母表示，而变量的实现值（或其值域中的值）用相应的小写字母表示.因此，随机变量$X$可以取值$x$.

\begin{example}[抛硬币三次——II]\label{ex:1.4.3}
再次考虑例\ref{ex:1.3.13}中抛公平硬币三次的试验.定义随机变量$X$为三次投掷中正面的次数.每个样本点对应的$X$值完整列举如下：\footnote{译注：原表第六、七列重复写成HTH、THH，却令$X=1$；按正面次数，这两列应为THT、TTH.原表保留，下文事件集合使用的是正确的三个单正面结果.}
\begin{center}\begin{tabular}{c|cccccccc}\hline
$s$&HHH&HHT&HTH&THH&HTT&HTH&THH&TTT\\\hline
$X(s)$&3&2&2&2&1&1&1&0\\\hline
\end{tabular}\end{center}
随机变量$X$的值域为$\mathscr X=\{0,1,2,3\}$.假设$S$中的八个点各有概率$\frac18$，则对上表简单计数可得，$\mathscr X$上的诱导概率函数为
\begin{center}\begin{tabular}{c|cccc}
$x$&0&1&2&3\\\hline
$P_X(X=x)$&$\frac18$&$\frac38$&$\frac38$&$\frac18$\\\hline
\end{tabular}\end{center}
例如，$P_X(X=1)=P(\{\mathrm{HTT,THT,TTH}\})=\frac38$.
\end{example}
% 原表第六、七列确为HTH、THH；不在译文中静默替换.见核验说明.


\sourcepageanchor{55}

\begin{example}[随机变量的分布]\label{ex:1.4.4}
即使不能像例\ref{ex:1.4.3}那样完整列举，也可能确定$P_X$.如本节开头所述，令$S$为50个0和1组成的$2^{50}$个串，$X=\text{1的个数}$，$\mathscr X=\{0,1,2,\ldots,50\}$.假设这$2^{50}$个串等可能.在原样本空间中，数出含27个1的所有串，就能求得$X=27$的概率.由于每个串等可能，
$$P_X(X=27)=\frac{\text{含27个1的串的个数}}{\text{串的总数}}=\frac{\binom{50}{27}}{2^{50}}.$$
一般地，对任意$i\in\mathscr X$，
$$P_X(X=i)=\frac{\binom{50}{i}}{2^{50}}.$$
\end{example}

前面的例子中，$S$与$\mathscr X$都有限，$P_X$的定义很直接.$\mathscr X$可数时也是如此.若$\mathscr X$不可数，则用类似\eqref{eq:1.4.1}的方式定义诱导概率函数$P_X$.\index{SI0598@诱导概率}\semanticindexanchor{598:55}对任意集合$A\subset\mathscr X$，
\begin{equation}\label{eq:1.4.2}
P_X(X\in A)=P(\{s\in S:X(s)\in A\}).
\end{equation}
这确实定义了一个合法的概率函数，可以验证它满足Kolmogorov公理.（精确地说，我们只对$\mathscr X$的子集所构成的某个$\sigma$代数，利用\eqref{eq:1.4.2}定义概率.不过，这里不讨论这些技术细节.）

% 用户授权疑误脚注：2026-11-04



% 原书印刷24–28页，PDF物理55–59页.
\subsection{分布函数}\label{sec:1.5}
对每个随机变量$X$，都关联一个称为$X$的累积分布函数的函数.
\begin{definition}\label{def:1.5.1}\index{SI0237@累积分布函数}
随机变量$X$的累积分布函数（cumulative distribution function，cdf）记为$F_X(x)$，定义为
$$F_X(x)=P_X(X\le x),\qquad\text{对所有$x$}.$$
\end{definition}

\begin{example}[抛三枚硬币]\label{ex:1.5.2}
考虑抛三枚公平硬币的试验，令$X=\text{观察到正面的次数}$.$X$的cdf为
\begin{equation}\label{eq:1.5.1}
F_X(x)=\begin{cases}
0,&-\infty<x<0,\\
\frac18,&0\le x<1,\\
\frac12,&1\le x<2,\\
\frac78,&2\le x<3,\\
1,&3\le x<\infty.
\end{cases}
\end{equation}
阶梯函数$F_X(x)$的图像见图\ref{fig:1.5.1}.从该图有几点值得注意.$F_X$对所有$x$都有定义，而不只对$\mathscr X=\{0,1,2,3\}$中的值有定义.例如，
$$F_X(2.5)=P(X\le2.5)=P(X=0,1\text{或}2)=\frac78.$$
注意，$F_X$在$x_i\in\mathscr X$处发生跳跃，且在$x_i$处的跳跃大小等于$P(X=x_i)$.又因为$X$不可能为负，所以$x<0$时$F_X(x)=0$；因为$X$必然不超过这样的值，所以$x\ge3$时$F_X(x)=1$.
\end{example}


\sourcepageanchor{56}

\begin{figure}[htbp]\centering
\includegraphics[width=.88\textwidth]{assets/chapter1/fig-1-5-1.pdf}
\caption{例\ref{ex:1.5.2}的cdf}\label{fig:1.5.1}
\end{figure}

从图\ref{fig:1.5.1}可以清楚地看到，$F_X$可能不连续，在某些$x$值处有跳跃.不过，由$F_X$的定义可知，在跳跃点处，$F_X$取跳跃上端的值（注意\eqref{eq:1.5.1}中不同的不等号）.这称为右连续（right-continuity）：从右侧趋近某一点时函数连续.右连续性是cdf定义的一个结果.相比之下，若定义$F_X(x)=P_X(X<x)$（注意严格不等号），则$F_X$将是左连续的.任一点$x$处的跳跃大小都等于$P(X=x)$.

每个cdf都满足一些性质；从概率的角度理解$F_X(x)$的定义时，其中一些性质是显然的.
\begin{theorem}\label{thm:1.5.3}\index{SI0242@性质}
函数$F(x)$是cdf，当且仅当下列三个条件成立：
\begin{enumerate}[label=\alph*.]
\item\label{item:1.5.3a}$\lim_{x\to-\infty}F(x)=0$且$\lim_{x\to\infty}F(x)=1$.
\item\label{item:1.5.3b}$F(x)$是$x$的非减函数.
\item\label{item:1.5.3c}$F(x)$右连续，即对每个数$x_0$，$\lim_{x\downarrow x_0}F(x)=F(x_0)$.
\end{enumerate}
\end{theorem}
\begin{proof}[证明提纲]
为证明必要性，将$F$写成概率函数的形式即可验证这三个性质（习题\ref{exercise:1.48}）.充分性的证明，即由函数$F$满足定理的三个条件，推出它是某个随机变量的cdf，要困难得多.必须证明存在样本空间$S$、$S$上的概率函数$P$以及定义在$S$上的随机变量$X$，使得$F$是$X$的cdf.
\end{proof}

\begin{example}[抛掷直到出现正面]\label{ex:1.5.4}\index{SI0319@累积分布函数}
设试验是不断抛一枚硬币，直到出现正面.令$p=\text{任一次投掷出现正面的概率}$，定义随机变量$X=\text{得到正面所需的投掷次数}$.则对任意$x=1,2,\ldots$，
\begin{equation}\label{eq:1.5.2}P(X=x)=(1-p)^{x-1}p,\end{equation}
因为该事件要求先出现$x-1$次反面，再出现一次正面，而所有试验都独立.由\eqref{eq:1.5.2}，对任意正整数$x$，
\begin{equation}\label{eq:1.5.3}P(X\le x)=\sum_{i=1}^xP(X=i)=\sum_{i=1}^x(1-p)^{i-1}p.\end{equation}
几何级数的部分和为
\begin{equation}\label{eq:1.5.4}\index{SI0460@部分和}\sum_{k=1}^nt^{k-1}=\frac{1-t^n}{1-t},\qquad t\ne1,\end{equation}
这一事实可以用归纳法证明（见习题\ref{exercise:1.50}）.将\eqref{eq:1.5.4}用于上述概率，得随机变量$X$的cdf为
\begin{equation*}\begin{aligned}
F_X(x)&=P(X\le x)\\
&=\frac{1-(1-p)^x}{1-(1-p)}p\\
&=1-(1-p)^x,\qquad x=1,2,\ldots.
\end{aligned}\end{equation*}
与例\ref{ex:1.5.2}类似，cdf $F_X(x)$在相邻非负整数之间保持平坦.

容易证明，若$0<p<1$，则$F_X(x)$满足定理\ref{thm:1.5.3}的条件.首先，
$$\lim_{x\to-\infty}F_X(x)=0,$$
因为对所有$x<0$都有$F_X(x)=0$；并且
$$\lim_{x\to\infty}F_X(x)=\lim_{x\to\infty}[1-(1-p)^x]=1,$$
此处取极限时，$x$只遍历整数值.为验证性质\hyperref[item:1.5.3b]{(b)}，只需注意\eqref{eq:1.5.3}中的和随着$x$增加，包含更多的正项.最后，为验证\hyperref[item:1.5.3c]{(c)}，注意对任意$x$，只要$\varepsilon>0$足够小，就有$F_X(x+\varepsilon)=F_X(x)$.因此，
$$\lim_{\varepsilon\downarrow0}F_X(x+\varepsilon)=F_X(x),$$
故$F_X(x)$右连续.$F_X(x)$是几何分布（geometric distribution，因几何级数而得名）的cdf，其图像见图\ref{fig:1.5.2}.
\end{example}


\sourcepageanchor{57}

\begin{example}[连续的cdf]\label{ex:1.5.5}\index{SI0329@累积分布函数}
连续cdf的一个例子是函数
\begin{equation}\label{eq:1.5.5}F_X(x)=\frac1{1+e^{-x}},\end{equation}
它满足定理\ref{thm:1.5.3}的条件.例如，
\begin{equation*}\begin{aligned}
\lim_{x\to-\infty}F_X(x)&=0,&&\text{因为}\quad\lim_{x\to-\infty}e^{-x}=\infty;\\
\lim_{x\to\infty}F_X(x)&=1,&&\text{因为}\quad\lim_{x\to\infty}e^{-x}=0.
\end{aligned}\end{equation*}
对$F_X(x)$求导，得
$$\frac d{dx}F_X(x)=\frac{e^{-x}}{(1+e^{-x})^2}>0,$$
说明$F_X(x)$递增.$F_X$不仅右连续，而且连续.这是Logistic分布（logistic distribution）的一个特例.
\end{example}


\sourcepageanchor{58}

\begin{figure}[htbp]\centering
\includegraphics[width=.88\textwidth]{assets/chapter1/fig-1-5-2.pdf}
\caption{几何分布的cdf，$p=.3$}\label{fig:1.5.2}
\end{figure}

\begin{example}[有跳跃的cdf]\label{ex:1.5.6}\index{SI0240@混合}
若$F_X$不是$x$的连续函数，它可以由连续片段与跳跃混合而成.例如，对某个满足$1>\varepsilon>0$的$\varepsilon$，将\eqref{eq:1.5.5}中的$F_X(x)$修改为
\begin{equation}\label{eq:1.5.6}
F_Y(y)=\begin{cases}
\dfrac{1-\varepsilon}{1+e^{-y}},&y<0,\\[5pt]
\varepsilon+\dfrac{1-\varepsilon}{1+e^{-y}},&y\ge0,
\end{cases}
\end{equation}
则$F_Y(y)$是某个随机变量$Y$的cdf（见习题\ref{exercise:1.47}）.函数$F_Y$在$y=0$处有高度为$\varepsilon$的跳跃，其余地方都连续.若观察某个仪表的读数，而该读数理论上可处于$-\infty$与$\infty$之间的任何位置，则这个模型可能合适.不过，这台仪表有时会卡在0处.于是可以用$F_Y$对观测建模，其中$\varepsilon$为仪表卡住的概率.
\end{example}

cdf是连续还是有跳跃，对应于相关随机变量是否连续.事实上，根据这种对应关系来定义连续随机变量很方便.
\begin{definition}\label{def:1.5.7}\index{SI0867@连续}\index{SI0868@离散}
若$F_X(x)$是$x$的连续函数，则随机变量$X$称为连续型（continuous）.若$F_X(x)$是$x$的阶梯函数，则随机变量$X$称为离散型（discrete）.
\end{definition}

本节最后给出一个定理，正式说明$F_X$完全确定随机变量$X$的概率分布.这个结论成立的前提是，仅对$\mathscr B^1$中的事件$A$定义$P(X\in A)$；其中$\mathscr B^1$是包含所有形如$(a,b)$、$[a,b)$、$(a,b]$和$[a,b]$的实数区间的最小$\sigma$代数.若对更大的事件类定义概率，则两个随机变量可能具有相同的分布函数，却并非对每个事件都有相同的概率（见\hyperref[bib:source78]{Chung（1974），第27页}）.与大多数统计应用一样，本书只关心区间、区间的可数并或交等事件，因此不考虑这类病态情况.我们首先需要两个随机变量同分布的概念.

\sourcepageanchor{59}

\begin{definition}\label{def:1.5.8}\index{SI0869@同分布}
若对每个集合$A\in\mathscr B^1$都有$P(X\in A)=P(Y\in A)$，则随机变量$X$与$Y$称为同分布（identically distributed）.
\end{definition}
注意，同分布的两个随机变量未必相等.也就是说，定义\ref{def:1.5.8}并不是说$X=Y$.

\begin{example}[同分布的随机变量]\label{ex:1.5.9}
考虑例\ref{ex:1.4.3}中抛一枚公平硬币三次的试验.定义随机变量$X$与$Y$为
\begin{equation*}\begin{aligned}
X&=\text{观察到正面的次数},\\
Y&=\text{观察到反面的次数}.
\end{aligned}\end{equation*}
$X$的分布已在例\ref{ex:1.4.3}中给出，容易验证$Y$的分布与之完全相同.也就是说，对每个$k=0,1,2,3$，都有$P(X=k)=P(Y=k)$.因此$X$与$Y$同分布.然而，没有任何样本点满足$X(s)=Y(s)$.
\end{example}

\begin{theorem}\label{thm:1.5.10}
下列两个陈述等价：
\begin{enumerate}[label=\alph*.]
\item\label{item:1.5.10a}随机变量$X$与$Y$同分布.
\item\label{item:1.5.10b}对每个$x$都有$F_X(x)=F_Y(x)$.
\end{enumerate}
\end{theorem}
\begin{proof}
为证明等价性，必须说明两个陈述相互蕴含.先证明\hyperref[item:1.5.10a]{(a)}$\implies$\hyperref[item:1.5.10b]{(b)}.

因为$X$与$Y$同分布，对任意集合$A\in\mathscr B^1$，有$P(X\in A)=P(Y\in A)$.特别地，对每个$x$，集合$(-\infty,x]$属于$\mathscr B^1$，且
$$F_X(x)=P(X\in(-\infty,x])=P(Y\in(-\infty,x])=F_Y(x).$$
反向的蕴含关系，即\hyperref[item:1.5.10b]{(b)}$\implies$\hyperref[item:1.5.10a]{(a)}，要难证得多.上面的论证说明，若$X$与$Y$的概率在所有集合上相同，则它们在区间上也相同.现在必须证明相反方向，即若$X$与$Y$的概率在所有区间上相同，则它们在所有集合上都相同.证明这一点需要大量使用\index{SI0974@生成的}\semanticindexanchor{974:59}$\sigma$代数，这里不展开细节.只需指出，只要证明两个概率函数在所有区间上一致就够了（\hyperref[bib:source78]{Chung，1974，第2.2节}）.
\end{proof}



% 原书印刷28–31页，PDF物理59–62页；至§1.7前.
\subsection{密度函数与质量函数}\label{sec:1.6}
与随机变量$X$及其cdf $F_X$相关联的，还有另一个函数，称为概率密度函数（probability density function，pdf）或概率质量函数（probability mass function，pmf）.pdf与pmf分别对应连续和离散情形.它们都与随机变量的“点概率”有关.
\begin{definition}\label{def:1.6.1}\index{SI0850@概率质量函数}
离散随机变量$X$的概率质量函数（pmf）为
$$f_X(x)=P(X=x),\qquad\text{对所有$x$}.$$
\end{definition}

\begin{example}[几何分布的概率]\label{ex:1.6.2}\index{SI0321@概率质量函数}
对于例\ref{ex:1.5.4}中的几何分布，其pmf为
$$f_X(x)=P(X=x)=\begin{cases}(1-p)^{x-1}p,&x=1,2,\ldots,\\0,&\text{其他情形}.
\end{cases}$$
回忆一下，$P(X=x)$，等价地$f_X(x)$，就是cdf在$x$处的跳跃大小.我们可以用pmf计算概率.既然现在可以度量单个点的概率，只需对所考虑事件中的所有点求和即可.因此，对满足$a\le b$的正整数$a,b$，有
$$P(a\le X\le b)=\sum_{k=a}^b f_X(k)=\sum_{k=a}^b(1-p)^{k-1}p.$$
特别地，
\begin{equation}\label{eq:1.6.1}P(X\le b)=\sum_{k=1}^b f_X(k)=F_X(b).\end{equation}
\end{example}

我们采用一个广泛接受的约定：用大写字母表示cdf，用相应的小写字母表示pmf或pdf.

在连续情形下定义pdf时，需要稍加谨慎.若直接尝试计算连续随机变量的$P(X=x)$，会得到如下结果.由于对任意$\varepsilon>0$，有$\{X=x\}\subset\{x-\varepsilon<X\le x\}$，由定理\ref{thm:1.2.9}\hyperref[itm:thm:1.2.9:c]{(c)}得
$$P(X=x)\le P(x-\varepsilon<X\le x)=F_X(x)-F_X(x-\varepsilon),$$
对任意$\varepsilon>0$均成立.因此，由$F_X$的连续性，
$$0\le P(X=x)\le\lim_{\varepsilon\downarrow0}[F_X(x)-F_X(x-\varepsilon)]=0.$$
不过，只要理解pdf的用途，它的定义便会清楚.

由例\ref{ex:1.6.2}可见，pmf给出“点概率”.在离散情形中，对pmf的值求和就得到cdf（如\eqref{eq:1.6.1}）.连续情形中相应的做法，是以积分代替求和，得到
$$P(X\le x)=F_X(x)=\int_{-\infty}^x f_X(t)\,dt.$$
利用微积分基本定理，若$f_X(x)$连续，则进一步有
\begin{equation}\label{eq:1.6.2}\index{SI0448@微积分基本定理}\frac d{dx}F_X(x)=f_X(x).\end{equation}
注意，这与离散情形的类比几乎完全一致：将“点概率”$f_X(x)$“加起来”，便得到区间概率.

\sourcepageanchor{60}

\begin{definition}\label{def:1.6.3}\index{SI0828@概率密度函数}
连续随机变量$X$的概率密度函数（pdf）$f_X(x)$是满足下式的函数：
\begin{equation}\label{eq:1.6.3}F_X(x)=\int_{-\infty}^x f_X(t)\,dt,\qquad\text{对所有$x$}.\end{equation}
\end{definition}

\emph{记号说明：}“$X$的分布由$F_X(x)$给出”可简记为“$X\sim F_X(x)$”，其中符号“$\sim$”读作“服从”.类似地，可写$X\sim f_X(x)$；若$X$与$Y$同分布，也可写$X\sim Y$.

在连续情形中，描述区间概率时可以略微随意一些.因为连续随机变量满足$P(X=x)=0$，所以
$$P(a<X<b)=P(a<X\le b)=P(a\le X<b)=P(a\le X\le b).$$
应当清楚，pdf（或pmf）包含的信息与cdf相同.因此，解题时可以使用其中任何一个，并应尽量选择较简单的那个.


\sourcepageanchor{61}

\begin{figure}[htbp]\centering
\includegraphics[width=.88\textwidth]{assets/chapter1/fig-1-6-1.pdf}
\caption{Logistic曲线下的面积}\label{fig:1.6.1}
\end{figure}
\FloatBarrier
\begin{example}[Logistic分布的概率]\label{ex:1.6.4}\index{SI0331@概率密度函数}
对于例\ref{ex:1.5.5}中的Logistic分布，有
$$F_X(x)=\frac1{1+e^{-x}},$$
因此
$$f_X(x)=\frac d{dx}F_X(x)=\frac{e^{-x}}{(1+e^{-x})^2}.$$
曲线$f_X(x)$下的面积给出区间概率（见图\ref{fig:1.6.1}）：
\begin{equation*}\begin{aligned}
P(a<X<b)&=F_X(b)-F_X(a)\\
&=\int_{-\infty}^b f_X(x)\,dx-\int_{-\infty}^a f_X(x)\,dx\\
&=\int_a^b f_X(x)\,dx.
\end{aligned}\end{equation*}
\end{example}

对于pdf（或pmf），实际上只有两个要求，它们都是定义的直接结果.
\begin{theorem}\label{thm:1.6.5}\index{SI0838@性质}\index{SI0856@性质}
函数$f_X(x)$是某个随机变量$X$的pdf（或pmf），当且仅当
\begin{enumerate}[label=\alph*.]
\item 对所有$x$都有$f_X(x)\ge0$.
\item $\sum_x f_X(x)=1$（pmf），或$\int_{-\infty}^{\infty}f_X(x)\,dx=1$（pdf）.
\end{enumerate}
\end{theorem}
\begin{proof}
若$f_X(x)$是pdf（或pmf），两个性质由定义立即得到.特别地，对pdf，利用\eqref{eq:1.6.3}和定理\ref{thm:1.5.3}，有
$$1=\lim_{x\to\infty}F_X(x)=\int_{-\infty}^{\infty}f_X(t)\,dt.$$
反向蕴含同样容易证明.\index{SI0831@构造}\semanticindexanchor{831:61}\index{SI0852@构造}\semanticindexanchor{852:61}一旦有了$f_X(x)$，就可以定义$F_X(x)$，并应用定理\ref{thm:1.5.3}.
\end{proof}


\sourcepageanchor{62}

从纯数学的观点看，任何具有有限正积分（或和）的非负函数，都能转化为pdf或pmf.例如，若$h(x)$为任意非负函数，在集合$A$上为正，其他地方为0，且对某个常数$K>0$有
$$\int_{\{x\in A\}}h(x)\,dx=K<\infty,$$
则$f_X(x)=h(x)/K$是某个取值于$A$的随机变量$X$的pdf.

事实上，关系\eqref{eq:1.6.3}并非总成立，因为$F_X(x)$可能连续却不可微.确实存在这样的连续随机变量：无论选择什么$f_X(x)$，该积分关系都不成立.这些情形相当病态，我们将忽略它们.因此，本书假设每个连续随机变量都满足\eqref{eq:1.6.3}.在更高等的教材中（例如\hyperref[bib:source33]{Billingsley，1995，第31节}），若随机变量满足\eqref{eq:1.6.3}，则称其为绝对连续（absolutely continuous）.\index{SI0865@绝对连续}\semanticindexanchor{865:62}



% 原书印刷31–37页／PDF物理62–68页；习题1.1–1.55.
\subsection{习题}\label{sec:1.7}
\setcounter{exercise}{0}

\begin{exercise}\label{exercise:1.1}
对下列每个试验，描述其样本空间.
\begin{enumerate}
\item[(a)] 抛掷一枚硬币四次.
\item[(b)] 数一株植物上受到虫害的叶片数.
\item[(c)] 测量某一品牌灯泡的寿命（以小时为单位）.
\item[(d)] 记录出生十天的大鼠的体重.
\item[(e)] 观察一批运来的电子元件中次品所占的比例.
\end{enumerate}
\end{exercise}

\begin{exercise}\label{exercise:1.2}\index{SI0965@恒等式}
验证下列恒等式.
\begin{enumerate}
\item[(a)] $A\setminus B=A\setminus(A\cap B)=A\cap B^c$.
\item[(b)] $B=(B\cap A)\cup(B\cap A^c)$.
\item[(c)] $B\setminus A=B\cap A^c$.
\item[(d)] $A\cup B=A\cup(B\cap A^c)$.
\end{enumerate}
\end{exercise}

\begin{exercise}\label{exercise:1.3}
完成定理\ref{thm:1.1.4}的证明.对样本空间$S$上定义的任意事件$A$、$B$和$C$，证明：
\begin{enumerate}
\item[(a)] $A\cup B=B\cup A$且$A\cap B=B\cap A$.（交换律）
\item[(b)] $A\cup(B\cup C)=(A\cup B)\cup C$且$A\cap(B\cap C)=(A\cap B)\cap C$.（结合律）
\item[(c)] $(A\cup B)^c=A^c\cap B^c$且$(A\cap B)^c=A^c\cup B^c$.（De Morgan定律）
\end{enumerate}
\end{exercise}

\begin{exercise}\label{exercise:1.4}
对事件$A$和$B$，用$P(A)$、$P(B)$和$P(A\cap B)$表示下列事件的概率.
\begin{enumerate}
\item[(a)] $A$发生，或$B$发生，或二者都发生.
\item[(b)] $A$与$B$中恰有一个发生.
\item[(c)] $A$与$B$中至少有一个发生.
\item[(d)] $A$与$B$中至多有一个发生.
\end{enumerate}
\end{exercise}

\begin{exercise}\label{exercise:1.5}
所有人类双胞胎中，约三分之一是同卵双胞胎，三分之二是异卵双胞胎.同卵双胞胎必定性别相同，男性与女性的可能性相等.异卵双胞胎中，约四分之一都是女性，四分之一都是男性，一半是一男一女.最后，在美国全部分娩中，约每90次有1次诞下双胞胎.定义下列事件：
$$
\begin{aligned}
A&=\{\text{美国的一次分娩诞下女性双胞胎}\},\\
B&=\{\text{美国的一次分娩诞下同卵双胞胎}\},\\
C&=\{\text{美国的一次分娩诞下双胞胎}\}.
\end{aligned}
$$
\begin{enumerate}
\item[(a)] 用文字描述事件$A\cap B\cap C$.
\item[(b)] 求$P(A\cap B\cap C)$.
\end{enumerate}
\end{exercise}


\sourcepageanchor{63}

\begin{exercise}\label{exercise:1.6}
将两枚硬币一起独立地抛掷，其中一枚满足$P(\text{正面})=u$，另一枚满足$P(\text{正面})=w$.定义
$$
\begin{aligned}
p_0&=P(\text{出现0个正面}),\\
p_1&=P(\text{出现1个正面}),\\
p_2&=P(\text{出现2个正面}).
\end{aligned}
$$
能否选择$u$和$w$，使$p_0=p_1=p_2$？证明你的答案.
\end{exercise}

\begin{exercise}\label{exercise:1.7}
参见例\ref{ex:1.2.7}中的飞镖游戏.假定我们不再假设击中镖盘的概率为1，而是假设该概率与镖盘面积成正比.假定镖盘安装在一面墙上，击中这面墙的概率为1，墙的面积为$A$.
\begin{enumerate}
\item[(a)] 利用击中某一区域的概率与其面积成正比这一事实，构造概率函数$P(\text{得}i\text{分})$，$i=0,\ldots,5$.（若未击中镖盘，则不得分.）
\item[(b)] 证明条件概率分布$P(\text{得}i\text{分}\mid\text{击中镖盘})$恰好就是例\ref{ex:1.2.7}中的概率分布.
\end{enumerate}
\end{exercise}

\begin{exercise}\label{exercise:1.8}
再次考虑例\ref{ex:1.2.7}所说明的飞镖游戏.
\begin{enumerate}
\item[(a)] 推导得$i$分的概率的一般公式.
\item[(b)] 证明$P(\text{得}i\text{分})$是$i$的减函数，即分值越高，得到该分值的概率越小.
\item[(c)] 证明$P(\text{得}i\text{分})$是符合柯尔莫哥洛夫公理的概率函数.
\end{enumerate}
\end{exercise}

\needspace{8\baselineskip}
\begin{exercise}\label{exercise:1.9}\index{SI0247@德摩根律（见集合运算）}
证明De Morgan定律的一般形式.设$\{A_\alpha:\alpha\in\Gamma\}$是一个（可能不可数的）集合族.证明：
\begin{enumerate}
\item[(a)] $(\bigcup_\alpha A_\alpha)^c=\bigcap_\alpha A_\alpha^c$.
\item[(b)] $(\bigcap_\alpha A_\alpha)^c=\bigcup_\alpha A_\alpha^c$.
\end{enumerate}
\end{exercise}

\begin{exercise}\label{exercise:1.10}
陈述并证明适用于有限集合族$A_1,\ldots,A_n$的De Morgan定律.
\end{exercise}

\begin{exercise}\label{exercise:1.11}\index{SI0975@性质}
设$S$为样本空间.
\begin{enumerate}
\item[(a)] 证明集合族$\mathscr{B}=\{\emptyset,S\}$是一个$\sigma$代数.
\item[(b)] 设$\mathscr{B}=\{S\text{的所有子集，包括}S\text{本身}\}$.证明$\mathscr{B}$是一个$\sigma$代数.
\item[(c)] 证明两个$\sigma$代数的交仍是一个$\sigma$代数.
\end{enumerate}
\end{exercise}

\begin{exercise}\label{exercise:1.12}\index{SI0073@连续性公理}\index{SI0074@可数可加性公理}\index{SI0075@有限可加性公理}
第\ref{sec:1.2.1}节曾提到，遵从De Finetti学派的统计学家不接受可列可加性公理，而坚持有限可加性公理.
\begin{enumerate}
\item[(a)] 证明可列可加性公理蕴含有限可加性.
\item[(b)] 有限可加性公理本身并不蕴含可列可加性，但假定我们补充如下条件.设$A_1\supset A_2\supset\cdots\supset A_n\supset\cdots$是一个无限的嵌套集合序列，其极限为空集，记为$A_n\downarrow\emptyset$.考虑下述公理：
$$
\textbf{连续性公理：}\qquad
\text{若 }A_n\downarrow\emptyset,\text{ 则 }P(A_n)\to0.
$$
证明连续性公理与有限可加性公理一起蕴含可列可加性.
\end{enumerate}
\end{exercise}

\begin{exercise}\label{exercise:1.13}
若$P(A)=\frac13$且$P(B^c)=\frac14$，$A$与$B$能否不相交？请解释.
\end{exercise}

\begin{exercise}\label{exercise:1.14}
假定样本空间$S$含有$n$个元素.证明由$S$中元素所能组成的子集数为$2^n$.
\end{exercise}

\begin{exercise}\label{exercise:1.15}
完成定理\ref{thm:1.2.14}的证明.以$k=2$时已经建立的结果作为归纳论证的基础.
\end{exercise}

\begin{exercise}\label{exercise:1.16}
若每个人有一个姓，并且有如下数量的名字，可以形成多少种不同的姓名首字母组合？
\begin{enumerate}
\item[(a)]\phantomsection\label{itm:exercise:1.16:a} 恰好两个名字.
\item[(b)]\phantomsection\label{itm:exercise:1.16:b} 一个或两个名字.
\item[(c)]\phantomsection\label{itm:exercise:1.16:c} 一个、两个或三个名字.
\end{enumerate}
（答案：\hyperref[itm:exercise:1.16:a]{(a)} $26^3$；\hyperref[itm:exercise:1.16:b]{(b)} $26^3+26^2$；\hyperref[itm:exercise:1.16:c]{(c)} $26^4+26^3+26^2$.）
\end{exercise}


\sourcepageanchor{64}

\begin{exercise}\label{exercise:1.17}
在多米诺骨牌游戏中，每张骨牌上标有两个数.骨牌是对称的，因此这两个数没有次序之分（例如$(2,6)=(6,2)$）.使用数$1,2,\ldots,n$，能够构成多少张不同的骨牌？

（答案：$n(n+1)/2$.）
\end{exercise}

\begin{exercise}\label{exercise:1.18}
将$n$个球随机放入$n$个格子中，求恰好有一个格子为空的概率.

（答案：$\binom n2 n!/n^n$.）
\end{exercise}

\begin{exercise}\label{exercise:1.19}
若多元函数具有连续偏导数，则求导顺序不影响结果.例如，二元函数$f(x,y)$的下列三阶偏导数相等：
$$
\frac{\partial^3}{\partial x^2\partial y}f(x,y)
=\frac{\partial^3}{\partial y\partial x^2}f(x,y).
$$
\begin{enumerate}
\item[(a)] 三元函数有多少个四阶偏导数？
\item[(b)] 证明$n$元函数有$\binom{n+r-1}{r}$个$r$阶偏导数.
\end{enumerate}
\end{exercise}

\begin{exercise}\label{exercise:1.20}
我的电话每周响12次，这些电话随机分布在七天之中.我每天至少接到一个电话的概率是多少？

（答案：$.2285$.）
\end{exercise}

\begin{exercise}\label{exercise:1.21}
衣柜里有$n$双鞋.若随机选出$2r$只鞋（$2r<n$），样本中没有任何一双相配的鞋的概率是多少？

（答案：$\binom n{2r}2^{2r}/\binom{2n}{2r}$.）
\end{exercise}

\begin{exercise}\label{exercise:1.22}\leavevmode
\begin{enumerate}
\item[(a)]\phantomsection\label{itm:exercise:1.22:a} 一次征兵抽签包含一年中的366天（包括2月29日）.不放回地抽取日期，最先抽出的180天恰好均匀分布在12个月中的概率是多少？
\item[(b)]\phantomsection\label{itm:exercise:1.22:b} 最先抽出的30天中没有任何一天属于9月的概率是多少？
\end{enumerate}
（答案：\hyperref[itm:exercise:1.22:a]{(a)} $.167\times10^{-8}$；\hyperref[itm:exercise:1.22:b]{(b)} $\binom{336}{30}/\binom{366}{30}$.）
\end{exercise}

\begin{exercise}\label{exercise:1.23}
两个人各将一枚均匀硬币抛掷$n$次.求他们得到相同正面次数的概率.

（答案：$(\frac14)^n\binom{2n}{n}$.）
\end{exercise}

\begin{exercise}\label{exercise:1.24}
两名玩家A和B轮流且独立地抛掷一枚硬币，先得到正面者获胜.假定玩家A先抛.
\begin{enumerate}
\item[(a)]\phantomsection\label{itm:exercise:1.24:a} 若硬币均匀，A获胜的概率是多少？
\item[(b)]\phantomsection\label{itm:exercise:1.24:b} 假定$P(\text{正面})=p$，不一定等于$\frac12$.A获胜的概率是多少？
\item[(c)] 证明对所有$0<p<1$，都有$P(\mathrm{A}\text{获胜})>\frac12$.（提示：尝试用事件$E_1,E_2,\ldots$表示$P(\mathrm{A}\text{获胜})$，其中$E_i=\{\text{第}i\text{次抛掷首次出现正面}\}$.）
\end{enumerate}
（答案：\hyperref[itm:exercise:1.24:a]{(a)} $2/3$；\hyperref[itm:exercise:1.24:b]{(b)} $\dfrac{p}{1-(1-p)^2}$.）
\end{exercise}

\begin{exercise}\label{exercise:1.25}
Smith夫妇有两个孩子.其中至少有一个是男孩.两个孩子都是男孩的概率是多少？（关于此问题的完整讨论，见\hyperref[bib:source132]{Gardner（1961）}.）
\end{exercise}

\begin{exercise}\label{exercise:1.26}
反复掷一枚均匀骰子，直至出现6点.必须投掷超过五次的概率是多少？
\end{exercise}

\begin{exercise}\label{exercise:1.27}\index{SI0531@二项式系数}
对$n\ge2$，验证下列恒等式.
\begin{enumerate}
\item[(a)] $\sum_{k=0}^n(-1)^k\binom nk=0$.
\item[(b)] $\sum_{k=1}^n k\binom nk=n2^{n-1}$.
\item[(c)] $\sum_{k=1}^n(-1)^{k+1}k\binom nk=0$.
\end{enumerate}
\end{exercise}

\begin{exercise}\label{exercise:1.28}\index{SI0571@对数}\index{SI0985@Stirling公式}
近似计算大数的阶乘的一种方法是使用Stirling公式（Stirling's Formula）：
$$
n!\approx\sqrt{2\pi}\,n^{n+(1/2)}e^{-n},
$$
其完整推导比较困难.作为替代，证明下述较容易的事实：
$$
\lim_{n\to\infty}\frac{n!}{n^{n+(1/2)}e^{-n}}=\text{某个常数}.
$$
（提示：\hyperref[bib:source118]{Feller（1968）}首先利用对数函数的单调性得到
$$
\int_{k-1}^{k}\log x\,dx<\log k<\int_k^{k+1}\log x\,dx,
\qquad k=1,\ldots,n,
$$
从而
$$
\int_0^n\log x\,dx<\log n!<\int_1^{n+1}\log x\,dx.
$$
现在将$\log n!$与这两个积分的平均值比较.另一种推导见\hyperref[exercise:5.35]{习题5.35}.）
\end{exercise}


\sourcepageanchor{65}

\begin{exercise}\label{exercise:1.29}\leavevmode
\begin{enumerate}
\item[(a)] 在例\ref{ex:1.2.20}的情形下，列出构成无序样本$\{4,4,12,12\}$与$\{2,9,9,12\}$的有序样本.
\item[(b)] 假定有六个数$\{1,2,7,8,14,20\}$.有放回抽样得到无序样本$\{2,7,7,8,14,14\}$的概率是多少？
\item[(c)] 验证：一个容量为$k$的无序样本中，$m$个不同的数分别重复$k_1,k_2,\ldots,k_m$次，则该样本包含$\dfrac{k!}{k_1!k_2!\cdots k_m!}$种有序排列，其中$k_1+k_2+\cdots+k_m=k$.
\item[(d)] 利用上一小题的结果，建立恒等式\footnote{译注：以下等式原样保留.若$k_1,\ldots,k_m$遍历非负整数且和为$k$，多项式定理给出左端为$m^k$；右端组合数是这类重数配置的个数，并非多项系数之和.}
$$
\sum_{k_1,k_2,\ldots,k_m:\,k_1+k_2+\cdots+k_m=k}
\frac{k!}{k_1!k_2!\cdots k_m!}=\binom{k+m-1}{k}.
$$
\end{enumerate}
\end{exercise}

\begin{exercise}\label{exercise:1.30}
从六个数$\{1,2,7,8,14,20\}$中作有放回抽样，画出所有可能样本均值的分布直方图.\footnote{译注：原题未给样本容量，因此不能唯一确定该分布；需要先补充样本容量.}
\end{exercise}

\needspace{8\baselineskip}
\begin{exercise}\label{exercise:1.31}
在例\ref{ex:1.2.20}的情形下，原数集$\{2,4,9,12\}$的平均值为$\frac{29}{4}$，其概率最高.
\begin{enumerate}
\item[(a)] 证明：一般地，从集合$\{x_1,x_2,\ldots,x_n\}$中作有放回抽样，平均值为$(x_1+x_2+\cdots+x_n)/n$的结果最有可能出现，其概率为$\dfrac{n!}{n^n}$.\footnote{译注：该值是每个原数恰出现一次的无序样本的概率，不能一般地直接当作样本均值的概率.例如从$\{0,1,2\}$有放回抽取3次，均值1还可由$\{1,1,1\}$产生，概率为$7/27$而非$6/27$.}
\item[(b)] 利用Stirling公式（习题\ref{exercise:1.28}），证明$n!/n^n\approx\sqrt{2n\pi}/e^n$（\hyperref[bib:source154]{Hall 1992，附录I}）.
\item[(c)] 证明某个特定的$x_i$没有出现在抽样结果中的概率为$(1-\frac1n)^n$，并且当$n\to\infty$时，该概率趋于$e^{-1}$.
\end{enumerate}
\end{exercise}

\begin{exercise}\label{exercise:1.32}
雇主打算从$N$名候选人中聘用一名新员工，这些候选人的未来潜力可以按1到$N$评分.雇主按如下规则行事：
\begin{enumerate}
\item[i.] 按随机顺序依次面试每位候选人，并决定是否聘用该候选人.
\item[ii.]\phantomsection\label{itm:exercise:1.32:ii} 在拒绝了$m-1$名候选人后（$m>1$），仅当第$m$名候选人优于前面的$m-1$名时，雇主才能聘用他.
\end{enumerate}
假定在第$i$次尝试时聘用了某位候选人.聘用到最优秀候选人的概率是多少？\footnote{译注：规则\hyperref[itm:exercise:1.32:ii]{(ii)}只给出了允许聘用的必要条件，未完全规定何时停止及聘用.若需要唯一的条件概率，还须明确决策机制；译文未替原题补写.}
\end{exercise}

\begin{exercise}\label{exercise:1.33}
假定男性中有$5\%$、女性中有$.25\%$患有色盲.随机选择一人，发现此人患有色盲.此人为男性的概率是多少？（假定男性与女性人数相等.）
\end{exercise}

\begin{exercise}\label{exercise:1.34}
某种啮齿动物产下了两窝幼崽，一窝有两只棕毛、一只灰毛（第1窝），另一窝有三只棕毛、两只灰毛（第2窝）.我们随机选一窝，再从所选的一窝中随机选一只幼崽.
\begin{enumerate}
\item[(a)] 所选动物为棕毛的概率是多少？
\item[(b)] 已知选中的是棕毛幼崽，抽样来自第1窝的概率是多少？
\end{enumerate}
\end{exercise}

\begin{exercise}\label{exercise:1.35}
证明：若$P(\cdot)$是一个合法的概率函数，且集合$B$满足$P(B)>0$，则$P(\cdot\mid B)$也满足柯尔莫哥洛夫公理.
\end{exercise}

\begin{exercise}\label{exercise:1.36}
若命中目标的概率为$\frac15$，独立射击十次，则至少命中两次的概率是多少？已知至少命中一次，至少命中两次的条件概率是多少？
\end{exercise}


\sourcepageanchor{66}

\begin{exercise}\label{exercise:1.37}\index{SI0711@蒙提霍尔问题}\index{SI1071@看守问题}
这里考察例\ref{ex:1.3.4}的若干变式.
\begin{enumerate}
\item[(a)] 在例\ref{ex:1.3.4}中狱长的计算里，假定若A将获赦免，则狱长告诉A“B将被处死”或“C将被处死”的概率相等.然而，情况未必如此.狱长可以分别给这两个事件赋予概率$\gamma$和$1-\gamma$，如下表所示：
\begin{center}
\begin{tabular}{lll}
\hline
获赦免的囚犯 & 狱长告诉A & \\
\hline
A & B将被处死 & 概率为$\gamma$\\
A & C将被处死 & 概率为$1-\gamma$\\
B & C将被处死 & \\
C & B将被处死 & \\
\hline
\end{tabular}
\end{center}
将$P(A\mid\mathscr{W})$作为$\gamma$的函数求出.$\gamma$取何值时，$P(A\mid\mathscr{W})$小于、等于或大于$\frac13$？
\item[(b)] 再假定像例中一样，$\gamma=\frac12$.狱长告诉A，B将被处死之后，A想了一会儿，意识到自己原来的计算是错的.但随后A灵机一动，问狱长能否让自己与C交换命运.狱长认为并未传递任何信息，便同意了.证明A此时的推理正确，他的存活概率已经跃升为$\frac23$！
\end{enumerate}
\hyperref[bib:source285]{Selvin（1975）}讨论了一个类似而略为复杂的问题，即“Monty Hall问题”.当这个形式的问题出现在一份星期日杂志上（\hyperref[bib:source322]{vos Savant 1990}），并附有正确答案但解释存在疑问时，它引起了相当大的关注.随后的争论甚至登上了星期日《纽约时报》的头版\mbox{（\hyperref[bib:source314]{Tierney 1991}）.}\hyperref[bib:source223]{Morgan等（1991）}作了完整且颇有趣味的讨论（另见\hyperref[bib:source323]{vos Savant 1991}\mbox{的回应）.}\hyperref[bib:source77]{Chun（1999）}给出的分析非常详尽，几乎穷尽了这一问题.
\end{exercise}

\begin{exercise}\label{exercise:1.38}\index{SI0124@性质}
证明以下各条陈述.（假定所有作为条件的事件都具有正概率.）
\begin{enumerate}
\item[(a)] 若$P(B)=1$，则对任意$A$，都有$P(A\mid B)=P(A)$.
\item[(b)] 若$A\subset B$，则$P(B\mid A)=1$且$P(A\mid B)=P(A)/P(B)$.
\item[(c)] 若$A$与$B$互斥，则
$$
P(A\mid A\cup B)=\frac{P(A)}{P(A)+P(B)}.
$$
\item[(d)] $P(A\cap B\cap C)=P(A\mid B\cap C)P(B\mid C)P(C)$.
\end{enumerate}
\end{exercise}

\begin{exercise}\label{exercise:1.39}\index{SI0250@独立性}\index{SI0549@不交}
两个事件$A$和$B$不能既互斥又独立.证明：若$P(A)>0$且$P(B)>0$，则
\begin{enumerate}
\item[(a)] 若$A$与$B$互斥，它们就不可能独立.
\item[(b)] 若$A$与$B$独立，它们就不可能互斥.
\end{enumerate}
\end{exercise}

\begin{exercise}\label{exercise:1.40}
证明定理\ref{thm:1.3.9}的\hyperref[item:1.3.9b]{(b)}和\hyperref[item:1.3.9c]{(c)}两部分，完成该定理的证明.
\end{exercise}

\begin{exercise}\label{exercise:1.41}
如例\ref{ex:1.3.6}，考虑按$3:4$比例发送的电报信号“点”和“划”.传输中的差错使“点”以$\frac14$的概率变成“划”，使“划”以$\frac13$的概率变成“点”.
\begin{enumerate}
\item[(a)] 若收到的是“划”，发送的也是“划”的概率是多少？
\item[(b)] 假定各信号彼此独立.若收到的消息是“点—点”，则四种可能发送的消息的概率分布是什么？
\end{enumerate}
\end{exercise}

\begin{exercise}\label{exercise:1.42}\index{SI0540@容斥}
杂录\ref{sec:1.8.1}中的容斥恒等式（inclusion-exclusion identity），因其用容斥法证明而得名（\hyperref[bib:source118]{Feller 1968，第IV.1节}）.这里给出细节.概率$P(\bigcup_{i=1}^n A_i)$等于至少属于一个$A_i$的所有样本点的概率之和.容斥法就是对这些点计数的一套方法.
\begin{enumerate}
\item[(a)]\phantomsection\label{itm:exercise:1.42:a} 若$E_k$表示恰好属于事件$A_1,A_2,\ldots,A_n$中$k$个事件的所有样本点构成的集合，证明$P(\bigcup_{i=1}^n A_i)=\sum_{i=1}^n P(E_i)$.
\item[(b)] 假定$x\in E_1$.证明$P(x\in E_1)=\sum_{i=1}^n P(x\in A_i)$.
\item[(c)]\phantomsection\label{itm:exercise:1.42:c} 假定$x\in E_k$.不失一般性，可以假设$x\in A_1,A_2,\ldots,A_k$且$x\notin A_{k+1},A_{k+2},\ldots,A_n$.证明$P(x\in E_k)$在和$P_1$中出现$k$次，在和$P_2$中出现$\binom k2$次，在和$P_3$中出现$\binom k3$次，依此类推.
\item[(d)] 证明
$$
k-\binom k2+\binom k3-+\cdots\pm\binom kk=1.
$$
（见习题\ref{exercise:1.27}.）\footnote{译注：原式在$\binom k3$后排成“$-+\cdots$”，保留该排印；意图应为继续交替加减.}
\item[(e)] 证明\hyperref[itm:exercise:1.42:a]{(a)}–\hyperref[itm:exercise:1.42:c]{(c)}蕴含$\sum_{i=1}^n P(E_i)=P_1-P_2=\cdots\pm P_n$，从而建立容斥恒等式.\footnote{译注：原题在$P_1-P_2$后写等号；标准容斥和应为$P_1-P_2+P_3-\cdots+(-1)^{n+1}P_n$.}
\end{enumerate}
\end{exercise}


\sourcepageanchor{67}

\begin{exercise}\label{exercise:1.43}
对于杂录\ref{sec:1.8.1}中的容斥恒等式：
\begin{enumerate}
\item[(a)] 从容斥恒等式推导Boole不等式和Bonferroni不等式.
\item[(b)] 证明当$i\ge j$时，$P_i$满足$P_i\ge P_j$；并证明杂录\ref{sec:1.8.1}中的界限序列随着项数增加而得到改进.\footnote{译注：原题的不等号方向与杂录相反，且这些交集概率和一般不具单调性.例如4个事件都等于$A$时，$P_1=4P(A)$、$P_2=6P(A)$；更高阶又会下降.另见杂录处译注.}
\item[(c)] 通常，界限中的项数越多，界限就越有用.然而，\hyperref[bib:source279]{Schwager（1984）}提醒说，在某些情形下改进并不大，尤其是各$A_i$高度相关时.考察极端情形：对每个$i$，都有$A_i=A$，此时这一界限序列会怎样？（见\hyperref[bib:source332]{Worsley（1982）}，以及\hyperref[bib:source333]{Worsley（1985）}与\hyperref[bib:source280]{Schwager（1985）}之间的往来讨论.）
\end{enumerate}
\end{exercise}

\begin{exercise}\label{exercise:1.44}
标准化考试为概率论提供了一个有趣的应用.先假定一次考试包含20道选择题，每道题有四个可能答案.如果学生每题都猜，那么参加考试可以建模为20个独立事件组成的序列.在学生完全猜测的条件下，求他至少答对十题的概率.
\end{exercise}

\begin{exercise}\label{exercise:1.45}
证明式\eqref{eq:1.4.1}中定义的诱导概率函数是一个合法的概率函数，即它满足柯尔莫哥洛夫公理.
\end{exercise}

\begin{exercise}\label{exercise:1.46}
将七个球随机分配到七个格子中.令$X_i=$恰好含有$i$个球的格子数.$X_3$的概率分布是什么？（即对每个可能的$x$，求$P(X_3=x)$.）
\end{exercise}

\begin{exercise}\label{exercise:1.47}
证明下列函数都是累积分布函数（cdf）.
\begin{enumerate}
\item[(a)] $\frac12+\frac1\pi\tan^{-1}(x)$，$x\in(-\infty,\infty)$.
\item[(b)] $(1+e^{-x})^{-1}$，$x\in(-\infty,\infty)$.
\item[(c)] $e^{-e^{-x}}$，$x\in(-\infty,\infty)$.
\item[(d)] $1-e^{-x}$，$x\in(0,\infty)$.\footnote{译注：要作为整个实数轴上的cdf，还须在$x\le0$时取0；原题只写出了正半轴上的形式.}
\item[(e)] 式\eqref{eq:1.5.6}所定义的函数.
\end{enumerate}
\end{exercise}

\begin{exercise}\label{exercise:1.48}
证明定理\ref{thm:1.5.3}的必要性部分.
\end{exercise}

\begin{exercise}\label{exercise:1.49}\index{SI0988@较大}
若对所有$t$都有$F_X(t)\le F_Y(t)$，并且对某个$t$有$F_X(t)<F_Y(t)$，则称累积分布函数$F_X$在随机意义下大于（stochastically greater than）累积分布函数$F_Y$.证明：若$X\sim F_X$且$Y\sim F_Y$，则
$$
P(X>t)\ge P(Y>t)\quad\text{对每个 }t\text{ 都成立},
$$
并且
$$
P(X>t)>P(Y>t)\quad\text{对某个 }t\text{ 成立},
$$
也就是说，$X$往往比$Y$大.
\end{exercise}

\begin{exercise}\label{exercise:1.50}
验证公式\eqref{eq:1.5.4}，即几何级数部分和的公式.
\end{exercise}

\begin{exercise}\label{exercise:1.51}
一家电器商店收到一批30台微波炉，其中五台有缺陷（经理并不知情）.经理从中不放回地随机抽取四台，检查是否有缺陷.令$X=$发现的有缺陷微波炉台数.计算$X$的概率质量函数（pmf）和累积分布函数，并画出累积分布函数.
\end{exercise}

\begin{exercise}\label{exercise:1.52}\index{SI0841@截断}
设$X$是连续随机变量，概率密度函数为$f(x)$，累积分布函数为$F(x)$.对一个固定的数$x_0$，定义函数
$$
g(x)=\begin{cases}
f(x)/[1-F(x_0)],&x\ge x_0,\\
0,&x<x_0.
\end{cases}
$$
证明$g(x)$是一个概率密度函数.（假定$F(x_0)<1$.）
\end{exercise}


\sourcepageanchor{68}

\begin{exercise}\label{exercise:1.53}
某条河流每年都会泛滥.假定最低水位定为1，最高水位$Y$的分布函数为
$$
F_Y(y)=P(Y\le y)=1-\frac1{y^2},\qquad 1\le y<\infty.
$$
\begin{enumerate}
\item[(a)] 验证$F_Y(y)$是累积分布函数.\footnote{译注：原题只给$y\ge1$上的表达式；在全实数轴上定义cdf时，需在$y<1$处取0.}
\item[(b)] 求$Y$的概率密度函数$f_Y(y)$.
\item[(c)] 若把最低水位重新定为零，并使用原测量单位的$\frac1{10}$作为新单位，则最高水位变成$Z=10(Y-1)$.求$F_Z(z)$.
\end{enumerate}
\end{exercise}

\begin{exercise}\label{exercise:1.54}
对下列各函数，确定使$f(x)$成为概率密度函数的$c$值.
\begin{enumerate}
\item[(a)] $f(x)=c\sin x$，$0<x<\pi/2$.
\item[(b)] $f(x)=ce^{-|x|}$，$-\infty<x<\infty$.
\end{enumerate}
\end{exercise}

\begin{exercise}\label{exercise:1.55}
某电子设备的寿命记为$T$.若它在时刻$t=3$之前失效，其价值为$V=5$；否则，其价值为$V=2T$.若$T$的概率密度函数为
$$
f_T(t)=\frac1{1.5}e^{-t/(1.5)},\qquad t>0,
$$
求$V$的累积分布函数.
\end{exercise}

% 用户授权疑误脚注：2026-11-04



% 原书印刷37页／PDF物理68页；本节不补写原文未给出的证明.
\subsection{杂录}\label{sec:1.8}
\subsubsection{Bonferroni界及其推广}\label{sec:1.8.1}\index{SI0560@Bonferroni}\index{SI0563@Boole}
式\eqref{eq:1.2.10}中的Bonferroni界，或Boole不等式（定理\ref{thm:1.2.11}），给出了交事件或并事件概率的简单界限.利用下面的展开，可以使这些界限越来越精确.

对集合$A_1,A_2,\ldots,A_n$，按如下方式构造一组新的嵌套交集.令
$$
\begin{aligned}
P_1&=\sum_{i=1}^n P(A_i),\\
P_2&=\sum_{1\le i<j\le n}^{n}P(A_i\cap A_j),\\
P_3&=\sum_{1\le i<j<k\le n}^{n}P(A_i\cap A_j\cap A_k),\\
&\mathrel{\phantom{=}}\vdots\\
P_n&=P(A_1\cap A_2\cap\cdots\cap A_n).
\end{aligned}
$$
那么，容斥恒等式表明\index{SI0540@容斥}\semanticindexanchor{540:68}
$$
P(A_1\cup A_2\cup\cdots\cup A_n)=P_1-P_2+P_3-P_4+\cdots\pm P_n.
$$
此外，$P_i$是有序的：当$i\le j$时，$P_i\ge P_j$；并且有下列一系列上界和下界：\footnote{译注：这里关于$P_i$单调的断言一般不成立.例如4个相同事件给出$P_1=4P(A)<6P(A)=P_2$（若$P(A)>0$）.以下Bonferroni上下界本身成立，但不同截断阶数的界不必单调变紧.}
$$
\begin{aligned}
P_1&\ge P\left(\bigcup_{i=1}^{n}A_i\right)\ge P_1-P_2,\\
P_1-P_2+P_3&\ge P\left(\bigcup_{i=1}^{n}A_i\right)\ge P_1-P_2+P_3-P_4,\\
&\mathrel{\phantom{\ge}}\vdots
\end{aligned}
$$
详情见习题\ref{exercise:1.42}和\ref{exercise:1.43}.

这些界限随项数增加而越来越紧，从而改进了最初的Bonferroni界.\footnote{译注：此一般性断言需要附加条件.例如6个事件都等于$A$时，一阶上界是$6P(A)$，三阶上界却是$(6-15+20)P(A)=11P(A)$，反而更松.}这些界限的应用包括游程概率的近似（\hyperref[bib:source187]{Karlin和Ost 1988}）以及多重比较程序（\hyperref[bib:source230]{Naiman和Wynn 1992}）.

% 用户授权疑误脚注：2026-11-04



\section{变换与期望}\label{sec:2}
\chapterdoi{2}

\sourcepageanchor{69}
\begin{flushright}\itshape
不过，我们现在需要的可不只是理论和说教.\\
——Sherlock Holmes，《血字的研究》
\end{flushright}
在许多问题中，我们用具有给定累积分布函数 $F_X$ 的随机变量 $X$ 对某种现象建立模型，但真正关心的却是 $X$ 的某个函数. 本章介绍取得这种函数的信息的若干方法；这些信息既可以是完整的分布，也可以是关于其平均行为的较粗略描述.



\subsection{随机变量的函数的分布}\label{sec:2.1}
设 $X$ 是具有累积分布函数 $F_X(x)$ 的随机变量. 那么 $X$ 的任意函数，比如 $g(X)$，也是随机变量. 我们常常关心 $g(X)$ 本身，并写 $Y=g(X)$ 来表示这个新随机变量. 由于 $Y$ 是 $X$ 的函数，可以用 $X$ 的概率行为描述 $Y$ 的概率行为. 即对任意集合 $A$，
$$P(Y\in A)=P(g(X)\in A)$$
这表明 $Y$ 的分布取决于函数 $F_X$ 和 $g$. 根据 $g$ 的选择，有时可以得到这个概率的易于处理的表达式.

形式上，若写 $y=g(x)$，则函数 $g(x)$ 定义了从 $X$ 的原样本空间 $\mathscr X$ 到新样本空间 $\mathscr Y$（随机变量 $Y$ 的样本空间）的映射，即\index{SI0664@映射}\semanticindexanchor{664:69}
$$g(x):\mathscr X\longrightarrow\mathscr Y.$$
我们给 $g$ 配上一个逆映射（inverse mapping）\index{SI0665@逆}\semanticindexanchor{665:69}，记为 $g^{-1}$；它把 $\mathscr Y$ 的子集映到 $\mathscr X$ 的子集，定义为
\begin{equation}\label{eq:2.1.1}
g^{-1}(A)=\{x\in\mathscr X:g(x)\in A\}.
\end{equation}
请注意，$g^{-1}$ 把集合映到集合；$g^{-1}(A)$ 就是 $\mathscr X$ 中被 $g(x)$ 映入集合 $A$ 的点组成的集合. $A$ 可以是单点集 $\{y\}$，此时
$$g^{-1}(\{y\})=\{x\in\mathscr X:g(x)=y\}.$$
我们通常简写为 $g^{-1}(y)$，但若不止一个 $x$ 满足 $g(x)=y$，这个量仍是一个集合. 如果只有一个 $x$ 满足 $g(x)=y$，则 $g^{-1}(y)$ 为单点集 $\{x\}$，我们写作 $g^{-1}(y)=x$. 现在若定义 $Y=g(X)$，则对任意集合 $A\subset\mathscr Y$，可写成
\begin{equation}\label{eq:2.1.2}
\begin{aligned}
P(Y\in A)&=P(g(X)\in A)\\
&=P(\{x\in\mathscr X:g(x)\in A\})\\
&=P(X\in g^{-1}(A)).
\end{aligned}
\end{equation}
\sourcepageanchor{70}
这就定义了 $Y$ 上的概率分布，容易验证它满足 Kolmogorov 公理.\index{SI0598@诱导概率}\semanticindexanchor{598:70}

如果 $X$ 是离散随机变量，$\mathscr X$ 为可数集，那么 $\mathscr Y=\{y:y=g(x),\ x\in\mathscr X\}$ 也是可数集，所以 $Y$ 是离散随机变量. 对 $y\in\mathscr Y$，$Y$ 的概率质量函数为
$$
f_Y(y)=P(Y=y)=\sum_{x\in g^{-1}(y)}P(X=x)
=\sum_{x\in g^{-1}(y)}f_X(x),
$$
在其他地方为 $0$. 因而，我们只需找出所有映到 $y$ 的 $x$，将相应概率相加即可.

\begin{example}[二项变换]\label{ex:2.1.1}\index{SI0261@二项}\index{SI0750@参数}
如果离散随机变量 $X$ 的概率质量函数具有下列形式，就称它服从二项分布（binomial distribution）：
\begin{equation}\label{eq:2.1.3}
f_X(x)=P(X=x)=\binom nx p^x(1-p)^{n-x},\qquad x=0,1,\ldots,n,
\end{equation}
其中 $n$ 是正整数，$0\le p\le1$. 数 $n$ 和 $p$ 称为参数（parameters）；改变它们，就得到不同的二项分布.

考虑 $Y=g(X)$，其中 $g(x)=n-x$，即 $Y=n-X$. 这里 $\mathscr X=\{0,1,\ldots,n\}$，$\mathscr Y=\{y:y=g(x),\ x\in\mathscr X\}=\{0,1,\ldots,n\}$. 对任意 $y\in\mathscr Y$，$n-x=g(x)=y$ 当且仅当 $x=n-y$，所以 $g^{-1}(y)$ 是单点 $x=n-y$. 因而
$$
\begin{aligned}
f_Y(y)&=\sum_{x\in g^{-1}(y)}f_X(x)=f_X(n-y)\\
&=\binom n{n-y}p^{n-y}(1-p)^{n-(n-y)}\\
&=\binom ny(1-p)^y p^{n-y}.
\end{aligned}
$$
最后一个等式使用了\autoref{def:1.2.17}蕴含的二项系数恒等式 $\binom ny=\binom n{n-y}$. 因此，$Y$ 也具有二项分布，但参数为 $n$ 和 $1-p$.
\end{example}

如果 $X$ 和 $Y$ 都是连续随机变量，有时可以写出求变换后随机变量的累积分布函数和概率密度函数的公式. 本节余下部分讨论其中几种情况. 设 $X$ 的概率密度函数为 $f_X$，$Y=g(X)$，则
\begin{equation}\label{eq:2.1.4}
\begin{aligned}
F_Y(y)&=P(Y\le y)=P(g(X)\le y)\\
&=P(\{x\in\mathscr X:g(x)\le y\})\\
&=\int_{\{x\in\mathscr X:g(x)\le y\}}f_X(x)\,dx.
\end{aligned}
\end{equation}
应用这个公式的困难在于确定积分区域并算出积分.

\begin{example}[均匀变换]\label{ex:2.1.2}
设 $X$ 在 $(0,2\pi)$ 上服从均匀分布（uniform distribution），即
$$f_X(x)=\begin{cases}1/(2\pi),&0<x<2\pi,\\0,&\text{其他}.\end{cases}$$
考虑 $Y=\sin^2X$. 如\autoref{fig:2.1.1}所示，
\begin{equation}\label{eq:2.1.5}
P(Y\le y)=P(X\le x_1)+P(x_2\le X\le x_3)+P(X\ge x_4).
\end{equation}
\sourcepageanchor{71}

\begin{figure}[htbp]\centering
\includegraphics[width=.8\textwidth]{assets/chapter2/fig2-1-1.pdf}
\caption{\autoref{ex:2.1.2}中变换 $y=\sin^2(x)$ 的图形.}\label{fig:2.1.1}
\end{figure}
由图形的对称性和均匀分布可知，$P(X\le x_1)=P(X\ge x_4)$，而 $P(x_2\le X\le x_3)=2P(x_2\le X\le\pi)$. 因此，\eqref{eq:2.1.5}化为
\begin{equation}\label{eq:2.1.6}
P(Y\le y)=2P(X\le x_1)+2P(x_2\le X\le\pi),
\end{equation}
其中 $x_1$ 和 $x_2$ 是方程 $\sin^2x=y$ 在 $0<x<\pi$ 内的两个解. 虽然这个例子看起来很简单，所得的累积分布函数却并不简单.
\end{example}

为了避免混淆，必须留意 $X$ 和 $Y$ 的取值集合. 我们记
\begin{equation}\label{eq:2.1.7}
\mathscr X=\{x:f_X(x)>0\},\qquad
\mathscr Y=\{y:y=g(x),\ \text{其中某个 }x\in\mathscr X\}.
\end{equation}
集合 $\mathscr X$ 称为 $f_X$ 的支撑集（support set）\index{SI1015@支持集}\semanticindexanchor{1015:71}，或非正式地称为支撑（support）. 也就是说，$f_X$ 在支撑集上为正，在其他地方为零. 概率质量函数以及任意非负函数也有支撑集.

如果函数 $g$ 满足 $u>v$ 蕴含 $g(u)>g(v)$，或者满足 $u<v$ 蕴含 $g(u)>g(v)$，就称 $g$ 为单调函数（monotone function）\index{SI1040@单调}\semanticindexanchor{1040:71}；前一种情况称为递增，后一种情况称为递减. 单调函数 $g:\mathscr X\to\mathscr Y$ 是一一且满的. 每个 $x$ 对应一个 $y$，而每个 $y$ 至多对应一个 $x$，所以它是一一的；\index{SI0729@一一变换}\semanticindexanchor{729:71}由于 $\mathscr Y$ 是 $g$ 的像，它也是满的.\index{SI0730@满射变换}\semanticindexanchor{730:71}
\sourcepageanchor{72}
因此，单调函数 $g$ 把 $x$ 与 $y$ 唯一配对，$g^{-1}$ 是单值函数，并且 $g^{-1}(y)=x$ 当且仅当 $y=g(x)$.

如果 $g$ 递增，那么
\begin{equation}\label{eq:2.1.8}
\begin{aligned}
\{x\in\mathscr X:g(x)\le y\}
&=\{x\in\mathscr X:g^{-1}(g(x))\le g^{-1}(y)\}\\
&=\{x\in\mathscr X:x\le g^{-1}(y)\}.
\end{aligned}
\end{equation}
同样，如果 $g$ 递减，那么
\begin{equation}\label{eq:2.1.9}
\begin{aligned}
\{x\in\mathscr X:g(x)\le y\}
&=\{x\in\mathscr X:g^{-1}(g(x))\ge g^{-1}(y)\}\\
&=\{x\in\mathscr X:x\ge g^{-1}(y)\}.
\end{aligned}
\end{equation}
画图就能看清，为什么递减变换会使不等号反向.

如果 $g$ 递增，结合\eqref{eq:2.1.4}和\eqref{eq:2.1.8}，得到
$$F_Y(y)=\int_{\{x\in\mathscr X:x\le g^{-1}(y)\}}f_X(x)\,dx
=\int_{-\infty}^{g^{-1}(y)}f_X(x)\,dx=F_X(g^{-1}(y)).$$
如果 $g$ 递减，则
$$F_Y(y)=\int_{g^{-1}(y)}^\infty f_X(x)\,dx=1-F_X(g^{-1}(y)),$$
其中最后一步用到了 $X$ 的连续性. 这些结果可以总结如下.

\begin{theorem}\label{thm:2.1.3}\index{SI0241@单调变换}
设 $X$ 的累积分布函数为 $F_X(x)$，$Y=g(X)$，并按\eqref{eq:2.1.7}定义 $\mathscr X$ 与 $\mathscr Y$.
\begin{enumerate}
\item[(a)] 如果 $g$ 在 $\mathscr X$ 上递增，则 $F_Y(y)=F_X(g^{-1}(y))$，$y\in\mathscr Y$.
\item[(b)] 如果 $g$ 在 $\mathscr X$ 上递减且 $X$ 是连续随机变量，则 $F_Y(y)=1-F_X(g^{-1}(y))$，$y\in\mathscr Y$.
\end{enumerate}
\end{theorem}

\begin{example}[均匀分布与指数分布的关系——I]\label{ex:2.1.4}\index{SI0296@与均匀的关系}\index{SI0390@与指数的关系}
设 $X$ 在 $(0,1)$ 上服从均匀分布，概率密度函数为 $f_X(x)=1$，$0<x<1$，在其他地方为 $0$，所以 $F_X(x)=x$，$0<x<1$. 考虑随机变量 $Y=-\log X$. 因为
$$\frac d{dx}g(x)=\frac d{dx}(-\log x)=-\frac1x<0,\qquad 0<x<1,$$
$g$ 在 $\mathscr X$ 上递减. 当 $X$ 在 $0$ 与 $1$ 之间变化时，$-\log x$ 在 $0$ 与 $\infty$ 之间变化，即 $\mathscr Y=(0,\infty)$. 对 $y>0$，由 $y=-\log x$ 得 $x=e^{-y}$，所以逆函数为 $g^{-1}(y)=e^{-y}$. 因此，对 $y>0$，
$$F_Y(y)=1-F_X(g^{-1}(y))=1-F_X(e^{-y})=1-e^{-y},$$
因为 $F_X(x)=x$. 对 $y\le0$，$F_Y(y)=0$. 注意，只要求 $g$ 在 $X$ 的支撑集 $(0,1)$ 上单调即可.
\end{example}

如果 $Y$ 有连续的概率密度函数，对它的累积分布函数求导便得到该密度.
\begin{theorem}\label{thm:2.1.5}\index{SI0835@单调变换}
设 $X$ 的概率密度函数为 $f_X(x)$，$Y=g(X)$，其中 $g$ 是单调函数，$\mathscr X$ 和 $\mathscr Y$ 按\eqref{eq:2.1.7}定义. 假定 $f_X$ 在 $\mathscr X$ 上连续，$g^{-1}$ 在 $\mathscr Y$ 上有连续导数，则
\begin{equation}\label{eq:2.1.10}
f_Y(y)=\begin{cases}
f_X(g^{-1}(y))\left|\dfrac d{dy}g^{-1}(y)\right|,&y\in\mathscr Y,\\
0,&\text{其他}.
\end{cases}
\end{equation}
\end{theorem}
\sourcepageanchor{73}
\begin{proof}
由\autoref{thm:2.1.3}和链式法则，
$$f_Y(y)=\frac d{dy}F_Y(y)=\begin{cases}
f_X(g^{-1}(y))\dfrac d{dy}g^{-1}(y),&g\text{ 递增},\\
-f_X(g^{-1}(y))\dfrac d{dy}g^{-1}(y),&g\text{ 递减},
\end{cases}$$
这可以简写为\eqref{eq:2.1.10}.
\end{proof}

\begin{example}[逆伽马概率密度函数]\label{ex:2.1.6}\index{SI0313@与逆伽马的关系}\index{SI0326@与伽马的关系}
设 $f_X(x)$ 是伽马概率密度函数（gamma pdf）
$$f(x)=\frac1{(n-1)!\beta^n}x^{n-1}e^{-x/\beta},\qquad0<x<\infty,$$
其中 $\beta$ 为正常数，$n$ 为正整数. 假设我们要计算 $g(X)=1/X$ 的概率密度函数. 注意，这里支撑集 $\mathscr X$ 和 $\mathscr Y$ 都是区间 $(0,\infty)$. 若令 $y=g(x)$，则 $g^{-1}(y)=1/y$ 且 $\frac d{dy}g^{-1}(y)=-1/y^2$. 应用上述定理，对 $y\in(0,\infty)$，
$$\begin{aligned}
f_Y(y)&=f_X(g^{-1}(y))\left|\frac d{dy}g^{-1}(y)\right|\\
&=\frac1{(n-1)!\beta^n}\left(\frac1y\right)^{n-1}e^{-1/(\beta y)}\frac1{y^2}\\
&=\frac1{(n-1)!\beta^n}\left(\frac1y\right)^{n+1}e^{-1/(\beta y)},
\end{aligned}$$
这是称为逆伽马概率密度函数（inverted gamma pdf）的密度的一种特殊情形.
\end{example}

在许多应用中，函数 $g$ 可能既不递增也不递减，因而不能应用上述结果. 不过，$g$ 常常在某些区间上单调，这就使我们能够得到 $Y=g(X)$ 的分布表达式. （如果 $g$ 在任何区间上都不单调，我们就有大麻烦了.）

\begin{example}[平方变换]\label{ex:2.1.7}
设 $X$ 是连续随机变量. 对 $y>0$，$Y=X^2$ 的累积分布函数为
$$F_Y(y)=P(Y\le y)=P(X^2\le y)=P(-\sqrt y\le X\le\sqrt y).$$
由于 $X$ 连续，可以去掉左端点的等号，得到
$$\begin{aligned}
F_Y(y)&=P(-\sqrt y<X\le\sqrt y)\\
&=P(X\le\sqrt y)-P(X\le-\sqrt y)=F_X(\sqrt y)-F_X(-\sqrt y).
\end{aligned}$$
对累积分布函数求导便得到 $Y$ 的概率密度函数：
$$\begin{aligned}
f_Y(y)&=\frac d{dy}F_Y(y)\\
&=\frac d{dy}[F_X(\sqrt y)-F_X(-\sqrt y)]\\
&=\frac1{2\sqrt y}f_X(\sqrt y)+\frac1{2\sqrt y}f_X(-\sqrt y),
\end{aligned}$$
其中对 $F_X(\sqrt y)$ 和 $F_X(-\sqrt y)$ 使用了链式法则. 因此，概率密度函数为
\begin{equation}\label{eq:2.1.11}
f_Y(y)=\frac1{2\sqrt y}\bigl(f_X(\sqrt y)+f_X(-\sqrt y)\bigr).
\end{equation}
注意，\eqref{eq:2.1.11}把 $Y$ 的密度表示成两部分之和，分别对应 $g(x)=x^2$ 单调的两个区间. 一般情形也是如此.
\end{example}

\sourcepageanchor{74}
\begin{theorem}\label{thm:2.1.8}\index{SI0837@分段单调变换}
设 $X$ 的概率密度函数为 $f_X(x)$，$Y=g(X)$，并按\eqref{eq:2.1.7}定义样本空间 $\mathscr X$. 假定存在 $\mathscr X$ 的划分 $A_0,A_1,\ldots,A_k$，使得 $P(X\in A_0)=0$，且 $f_X(x)$ 在每个 $A_i$ 上连续. 再假定存在分别定义在 $A_1,\ldots,A_k$ 上的函数 $g_1(x),\ldots,g_k(x)$，满足
\begin{enumerate}
\item[(i)] 对 $x\in A_i$，$g(x)=g_i(x)$；
\item[(ii)] $g_i(x)$ 在 $A_i$ 上单调；
\item[(iii)] 对所有 $i=1,\ldots,k$，集合 $\mathscr Y=\{y:y=g_i(x),\ \text{其中某个 }x\in A_i\}$ 都相同；
\item[(iv)] 对每个 $i=1,\ldots,k$，$g_i^{-1}(y)$ 在 $\mathscr Y$ 上有连续导数.
\end{enumerate}
则
$$f_Y(y)=\begin{cases}
\displaystyle\sum_{i=1}^k f_X(g_i^{-1}(y))\left|\dfrac d{dy}g_i^{-1}(y)\right|,&y\in\mathscr Y,\\
0,&\text{其他}.
\end{cases}$$
\end{theorem}

\autoref{thm:2.1.8}的关键在于，可以把 $\mathscr X$ 分成集合 $A_1,\ldots,A_k$，使 $g(x)$ 在每个 $A_i$ 上单调. 因为 $P(X\in A_0)=0$，我们可以忽略这个“例外集”$A_0$. 这是一种技术处理，例如可用于处理区间端点. 要注意，每个 $g_i(x)$ 都是一一地把 $A_i$ 映到整个 $\mathscr Y$ 的变换. 此外，$g_i^{-1}(y)$ 是从 $\mathscr Y$ 到整个 $A_i$ 的一一函数；对 $y\in\mathscr Y$，$g_i^{-1}(y)$ 给出唯一的 $x=g_i^{-1}(y)\in A_i$，使得 $g_i(x)=y$. （推广见\autoref{exercise:2.7}.）

\begin{example}[正态分布与卡方分布的关系]\label{ex:2.1.9}\index{SI0284@与正态的关系}\index{SI0366@与卡方的关系}
设 $X$ 具有标准正态分布（standard normal distribution），
$$f_X(x)=\frac1{\sqrt{2\pi}}e^{-x^2/2},\qquad-\infty<x<\infty.$$
考虑 $Y=X^2$. 函数 $g(x)=x^2$ 在 $(-\infty,0)$ 和 $(0,\infty)$ 上单调，集合 $\mathscr Y=(0,\infty)$. 应用\autoref{thm:2.1.8}，取
$$\begin{aligned}
A_0&=\{0\};\\
A_1&=(-\infty,0),&g_1(x)&=x^2,&g_1^{-1}(y)&=-\sqrt y;\\
A_2&=(0,\infty),&g_2(x)&=x^2,&g_2^{-1}(y)&=\sqrt y.
\end{aligned}$$
$Y$ 的概率密度函数为
$$\begin{aligned}
f_Y(y)&=\frac1{\sqrt{2\pi}}e^{-(-\sqrt y)^2/2}\left|-\frac1{2\sqrt y}\right|
+\frac1{\sqrt{2\pi}}e^{-(\sqrt y)^2/2}\left|\frac1{2\sqrt y}\right|\\
&=\frac1{\sqrt{2\pi}}\frac1{\sqrt y}e^{-y/2},\qquad0<y<\infty.
\end{aligned}$$
这是我们将经常遇到的密度，即自由度为 $1$ 的卡方随机变量（chi squared random variable）的概率密度函数.
\end{example}

本节最后介绍一种特殊而又非常有用的变换.
\begin{theorem}[概率积分变换]\label{thm:2.1.10}\index{SI0847@概率积分变换}
设 $X$ 具有连续的累积分布函数 $F_X(x)$，并定义随机变量 $Y=F_X(X)$. 则 $Y$ 在 $(0,1)$ 上服从均匀分布，即 $P(Y\le y)=y$，$0<y<1$.
\end{theorem}
在证明这个定理之前，我们暂且岔开，详细考察 $F_X^{-1}$，也就是累积分布函数 $F_X$ 的逆. 如果 $F_X$ 严格递增，则可用下式良好地定义 $F_X^{-1}$：
\begin{equation}\label{eq:2.1.12}
F_X^{-1}(y)=x\quad\Longleftrightarrow\quad F_X(x)=y.
\end{equation}
\sourcepageanchor{75}
\begin{figure}[htbp]\centering
\includegraphics[width=\textwidth]{assets/chapter2/fig2-1-2.pdf}
\caption{(a) $F(x)$ 严格递增；(b) $F(x)$ 非递减.}\label{fig:2.1.2}
\end{figure}
然而，如果 $F_X$ 在某个区间上为常数，\eqref{eq:2.1.12}就不能良好地定义 $F_X^{-1}$，如\autoref{fig:2.1.2}所示. 任意满足 $x_1\le x\le x_2$ 的 $x$ 都满足 $F_X(x)=y$.

为避免这一问题，对 $0<y<1$，定义
\begin{equation}\label{eq:2.1.13}
F_X^{-1}(y)=\inf\{x:F_X(x)\ge y\}.
\end{equation}
当 $F_X$ 不为常数时，这一定义与\eqref{eq:2.1.12}一致；即使 $F_X$ 并非严格递增，它也给出了单值的 $F_X^{-1}$. 使用这一定义，\autoref{fig:2.1.2}(b) 中有 $F_X^{-1}(y)=x_1$. 在 $y$ 的取值范围的端点，也可以定义 $F_X^{-1}(y)$：如果对所有 $x$ 都有 $F_X(x)<1$，则 $F_X^{-1}(1)=\infty$；对任意 $F_X$，均有 $F_X^{-1}(0)=-\infty$.

\begin{proof}[\autoref{thm:2.1.10}的证明]
对于 $Y=F_X(X)$，当 $0<y<1$ 时，
$$\begin{aligned}
P(Y\le y)&=P(F_X(X)\le y)\\
&=P(F_X^{-1}[F_X(X)]\le F_X^{-1}(y))&&\text{（$F_X^{-1}$ 递增）}\\
&=P(X\le F_X^{-1}(y))&&\text{（见下文）}\\
&=F_X(F_X^{-1}(y))&&\text{（$F_X$ 的定义）}\\
&=y.&&\text{（$F_X$ 的连续性）}
\end{aligned}$$
在端点处，当 $y\ge1$ 时 $P(Y\le y)=1$，当 $y\le0$ 时 $P(Y\le y)=0$，这表明 $Y$ 服从均匀分布.

等式
$$P\bigl(F_X^{-1}(F_X(X))\le F_X^{-1}(y)\bigr)=P(X\le F_X^{-1}(y))$$
背后的理由略有些微妙，值得进一步说明. 如果 $F_X$ 严格递增，确实有 $F_X^{-1}(F_X(x))=x$. （参见\autoref{fig:2.1.2}(a).）然而，如果 $F_X$ 有平坦部分，就可能有 $F_X^{-1}(F_X(x))\ne x$. 假设 $F_X$ 如\autoref{fig:2.1.2}(b) 所示，且 $x\in[x_1,x_2]$，则对该区间内的任意 $x$，都有 $F_X^{-1}(F_X(x))=x_1$.
\sourcepageanchor{76}
即使如此，概率的等式仍然成立，因为对任意 $x\in[x_1,x_2]$，都有 $P(X\le x)=P(X\le x_1)$. 累积分布函数的平坦部分表示一个概率为 $0$ 的区域，即 $P(x_1<X\le x)=F_X(x)-F_X(x_1)=0$.
\end{proof}

\autoref{thm:2.1.10}的一个应用是从给定分布产生随机样本.\index{SI0861@随机数生成}\semanticindexanchor{861:76} 如果需要从累积分布函数为 $F_X$ 的总体产生一个观测值 $X$，只需在 $0$ 与 $1$ 之间产生一个均匀随机数 $U$，然后由方程 $F_X(x)=u$ 解出 $x$. （对许多分布，存在耗费更少计算时间的其他观测值生成方法；不过，本方法具有普遍适用性，因而仍然有用.）



\subsection{期望值}\label{sec:2.2}
随机变量的期望值（expected value），或期望（expectation），就是它的平均值；这里的“平均”是按照概率分布加权的平均. 我们把平均值理解为居中的值，因此分布的期望值可以视为其中心的一种度量. 按概率分布对随机变量的各个值加权，我们希望得到一个数，以概括该随机变量的一次观测的典型值或预期值.

\begin{definition}\label{def:2.2.1}\index{SI0422@期望值}
随机变量 $g(X)$ 的期望值或均值（mean），记为 $\mathrm E g(X)$，定义为
$$\mathrm E g(X)=\begin{cases}
\displaystyle\int_{-\infty}^\infty g(x)f_X(x)\,dx,&X\text{ 连续},\\
\displaystyle\sum_{x\in\mathscr X}g(x)f_X(x)=\sum_{x\in\mathscr X}g(x)P(X=x),&X\text{ 离散},
\end{cases}$$
前提是该积分或级数存在. 如果 $\mathrm E|g(X)|=\infty$，我们就说 $\mathrm E g(X)$ 不存在. （\hyperref[bib:source259]{Ross (1988)} 将此称为“无意识统计学家定律”. 我们并不觉得这有什么好笑的.）
\end{definition}

\begin{example}[指数分布的均值]\label{ex:2.2.2}\index{SI0292@均值}
设 $X$ 具有指数分布 $\operatorname{exponential}(\lambda)$，也就是说，它的概率密度函数为
$$f_X(x)=\frac1\lambda e^{-x/\lambda},\qquad0\le x<\infty,\quad\lambda>0.$$
那么 $\mathrm EX$ 为
$$\begin{aligned}
\mathrm EX&=\int_0^\infty\frac1\lambda xe^{-x/\lambda}\,dx\\
&=\left.-xe^{-x/\lambda}\right|_0^\infty+\int_0^\infty e^{-x/\lambda}\,dx&&\text{（分部积分）}\\
&=\int_0^\infty e^{-x/\lambda}\,dx=\lambda.
\end{aligned}$$
\end{example}

\begin{example}[二项分布的均值]\label{ex:2.2.3}\index{SI0264@均值}
如果 $X$ 具有二项分布，其概率质量函数为
$$P(X=x)=\binom nx p^x(1-p)^{n-x},\qquad x=0,1,\ldots,n,$$
其中 $n$ 为正整数，$0\le p\le1$；对每一对固定的 $n$ 和 $p$，概率质量函数之和为 $1$. 二项随机变量的期望为
$$\mathrm EX=\sum_{x=0}^n x\binom nx p^x(1-p)^{n-x}
=\sum_{x=1}^n x\binom nx p^x(1-p)^{n-x}$$
\sourcepageanchor{77}
（$x=0$ 的项为 $0$）. 利用恒等式 $x\binom nx=n\binom{n-1}{x-1}$，得到
$$\begin{aligned}
\mathrm EX&=\sum_{x=1}^n n\binom{n-1}{x-1}p^x(1-p)^{n-x}\\
&=\sum_{y=0}^{n-1}n\binom{n-1}y p^{y+1}(1-p)^{n-(y+1)}&&\text{（令 $y=x-1$）}\\
&=np\sum_{y=0}^{n-1}\binom{n-1}y p^y(1-p)^{n-1-y}\\
&=np,
\end{aligned}$$
因为最后的求和是 $\operatorname{binomial}(n-1,p)$ 概率质量函数对所有可能值求和，必定等于 $1$.
\end{example}

\begin{example}[柯西分布的均值]\label{ex:2.2.4}\index{SI0278@均值}
期望不存在的随机变量的一个经典例子是柯西随机变量（Cauchy random variable），其概率密度函数为
$$f_X(x)=\frac1\pi\frac1{1+x^2},\qquad-\infty<x<\infty.$$
容易验证 $\int_{-\infty}^\infty f_X(x)\,dx=1$，但 $\mathrm E|X|=\infty$. 写成
$$\mathrm E|X|=\int_{-\infty}^\infty\frac{|x|}\pi\frac1{1+x^2}\,dx
=\frac2\pi\int_0^\infty\frac{x}{1+x^2}\,dx.$$
对任意正数 $M$，
$$\int_0^M\frac{x}{1+x^2}\,dx
=\left.\frac{\log(1+x^2)}2\right|_0^M=\frac{\log(1+M^2)}2.$$
因此，
$$\mathrm E|X|=\lim_{M\to\infty}\frac2\pi\int_0^M\frac{x}{1+x^2}\,dx
=\frac1\pi\lim_{M\to\infty}\log(1+M^2)=\infty,$$
所以 $\mathrm EX$ 不存在.
\end{example}

求期望是线性运算，这意味着只要注意到对任意常数 $a$ 和 $b$ 有
\begin{equation}\label{eq:2.2.1}
\mathrm E(aX+b)=a\mathrm EX+b,
\end{equation}
就能容易地算出 $X$ 的线性函数的期望. 例如，若 $X$ 服从 $\operatorname{binomial}(n,p)$，则 $\mathrm EX=np$，所以
$$\mathrm E(X-np)=\mathrm EX-np=np-np=0.$$
实际上，期望算子有许多能够减轻计算工作的性质. 其中多数来自积分或求和的性质，概括于下面的定理.

\begin{theorem}\label{thm:2.2.5}\index{SI0432@性质}
设 $X$ 是随机变量，$a,b,c$ 是常数. 对任意期望存在的函数 $g_1(x)$ 和 $g_2(x)$，有
\begin{enumerate}
\item[(a)]\phantomsection\label{itm:thm:2.2.5:a} $\mathrm E(ag_1(X)+bg_2(X)+c)=a\mathrm Eg_1(X)+b\mathrm Eg_2(X)+c$.
\item[(b)] 如果对所有 $x$ 都有 $g_1(x)\ge0$，则 $\mathrm Eg_1(X)\ge0$.
\item[(c)] 如果对所有 $x$ 都有 $g_1(x)\ge g_2(x)$，则 $\mathrm Eg_1(X)\ge\mathrm Eg_2(X)$.
\item[(d)] 如果对所有 $x$ 都有 $a\le g_1(x)\le b$，则 $a\le\mathrm Eg_1(X)\le b$.
\end{enumerate}
\end{theorem}
\sourcepageanchor{78}
\begin{proof}
我们只详细证明连续情形，离散情形类似. 由定义和积分的可加性，
$$\begin{aligned}
&\mathrm E(ag_1(X)+bg_2(X)+c)\\
&\quad=\int_{-\infty}^\infty(ag_1(x)+bg_2(x)+c)f_X(x)\,dx\\
&\quad=\int_{-\infty}^\infty ag_1(x)f_X(x)\,dx
+\int_{-\infty}^\infty bg_2(x)f_X(x)\,dx+\int_{-\infty}^\infty cf_X(x)\,dx.
\end{aligned}$$
因为 $a,b,c$ 为常数，可以分别提出积分号，于是
$$\begin{aligned}
&\mathrm E(ag_1(X)+bg_2(X)+c)\\
&\quad=a\int_{-\infty}^\infty g_1(x)f_X(x)\,dx
+b\int_{-\infty}^\infty g_2(x)f_X(x)\,dx+c\int_{-\infty}^\infty f_X(x)\,dx\\
&\quad=a\mathrm Eg_1(X)+b\mathrm Eg_2(x)+c,
\end{aligned}$$
这就证明了 \hyperref[itm:thm:2.2.5:a]{(a)}.\footnote{译注：原文最后一行写作 $b\mathrm Eg_2(x)$，按定理陈述和前面的积分，应为 $b\mathrm Eg_2(X)$；此处保留原文大小写.} 其余三个性质可用类似的方法证明.
\end{proof}

\begin{example}[距离最小化]\label{ex:2.2.6}\index{SI0430@最小化性质}
随机变量的期望还有另一个性质，它与把 $\mathrm EX$ 解释为对 $X$ 的值的一个良好猜测有关.

假设用 $(X-b)^2$ 衡量随机变量 $X$ 与常数 $b$ 之间的距离. $b$ 越接近 $X$，这个量就越小. 现在可以确定使 $\mathrm E(X-b)^2$ 最小的 $b$，从而得到 $X$ 的一个良好预测量. （注意，寻找使 $(X-b)^2$ 最小的 $b$ 没有用处，因为答案依赖于 $X$，无法用来预测 $X$.）

我们可以用微积分求 $\mathrm E(X-b)^2$ 的最小值，但有一个更简单的方法. （基于微积分的证明见\autoref{exercise:2.19}.）相信 $\mathrm EX$ 应该有某种特殊性，写出
$$\begin{aligned}
\mathrm E(X-b)^2&=\mathrm E(X-\mathrm EX+\mathrm EX-b)^2&&\text{（加上 $\pm\mathrm EX$，值不变）}\\
&=\mathrm E((X-\mathrm EX)+(\mathrm EX-b))^2&&\text{（合并项）}\\
&=\mathrm E(X-\mathrm EX)^2+(\mathrm EX-b)^2\\
&\qquad+2\mathrm E((X-\mathrm EX)(\mathrm EX-b)),
\end{aligned}$$
其中展开了平方. 现在注意，
$$\mathrm E((X-\mathrm EX)(\mathrm EX-b))=(\mathrm EX-b)\mathrm E(X-\mathrm EX)=0,$$
因为 $(\mathrm EX-b)$ 是常数，可以提出期望号，而 $\mathrm E(X-\mathrm EX)=\mathrm EX-\mathrm EX=0$. 这意味着
\begin{equation}\label{eq:2.2.2}
\mathrm E(X-b)^2=\mathrm E(X-\mathrm EX)^2+(\mathrm EX-b)^2.
\end{equation}
我们不能改变\eqref{eq:2.2.2}右端第一项；第二项总是大于等于 $0$，选择 $b=\mathrm EX$ 就能使其等于 $0$. 因此，
\begin{equation}\label{eq:2.2.3}
\min_b\mathrm E(X-b)^2=\mathrm E(X-\mathrm EX)^2.
\end{equation}
关于中位数的类似结果见\autoref{exercise:2.18}.
\end{example}

\sourcepageanchor{79}
计算 $X$ 的非线性函数的期望时，可以采用两种方法之一. 根据 $\mathrm Eg(X)$ 的定义，可以直接计算
\begin{equation}\label{eq:2.2.4}
\mathrm Eg(X)=\int_{-\infty}^\infty g(x)f_X(x)\,dx.
\end{equation}
也可以先求出 $Y=g(X)$ 的密度 $f_Y(y)$，于是
\begin{equation}\label{eq:2.2.5}\index{SI0434@两种计算方法}
\mathrm Eg(X)=\mathrm EY=\int_{-\infty}^\infty yf_Y(y)\,dy.
\end{equation}

\begin{example}[均匀分布与指数分布的关系——II]\label{ex:2.2.7}
设 $X$ 服从 $\operatorname{uniform}(0,1)$ 分布，即 $X$ 的概率密度函数为
$$f_X(x)=\begin{cases}1,&0\le x\le1,\\0,&\text{其他},\end{cases}$$
并定义新随机变量 $g(X)=-\log X$. 那么
$$\mathrm Eg(X)=\mathrm E(-\log X)=\int_0^1-\log x\,dx
=\left.x-x\log x\right|_0^1=1.$$
但我们在\autoref{ex:2.1.4}中也已看到，$Y=-\log X$ 的累积分布函数为 $1-e^{-y}$，所以其密度为 $f_Y(y)=\frac d{dy}(1-e^{-y})=e^{-y}$，$0<y<\infty$，这是 $\lambda=1$ 时指数密度的特殊情形. 因而由\autoref{ex:2.2.2}，$\mathrm EY=1$.
\end{example}



\subsection{矩与矩母函数}\label{sec:2.3}
分布的各种矩是一类重要的期望.
\begin{definition}\label{def:2.3.1}\index{SI0697@矩}
对每个整数 $n$，$X$（或 $F_X(x)$）的 $n$ 阶矩（$n$th moment）$\mu'_n$ 为
$$\mu'_n=\mathrm EX^n.$$
$X$ 的 $n$ 阶中心矩（$n$th central moment）$\mu_n$ 为
$$\mu_n=\mathrm E(X-\mu)^n,$$
其中 $\mu=\mu'_1=\mathrm EX$.
\end{definition}
除随机变量的均值 $\mathrm EX$ 以外，最重要的矩或许就是二阶中心矩，通常称为方差.
\begin{definition}\label{def:2.3.2}\index{SI0980@标准差}\index{SI1062@方差}
随机变量 $X$ 的方差（variance）就是它的二阶中心矩，$\operatorname{Var}X=\mathrm E(X-\mathrm EX)^2$. $\operatorname{Var}X$ 的正平方根称为 $X$ 的标准差（standard deviation）.
\end{definition}
方差衡量分布在其均值周围的分散程度. 前面在\autoref{ex:2.2.6}中已经看到，选择 $b=\mathrm EX$ 可使 $\mathrm E(X-b)^2$ 最小. 现在我们考虑这个最小值本身的大小. 方差的解释是，其值越大，$X$ 的变异越大. 极端地，如果 $\operatorname{Var}X=\mathrm E(X-\mathrm EX)^2=0$，那么 $X$ 以概率 $1$ 等于 $\mathrm EX$，$X$ 完全没有变异. 标准差有相同的定性解释：值较小表示 $X$ 很可能接近 $\mathrm EX$，值较大表示 $X$ 的变异很大. 标准差更容易解释，因为它的度量单位与原变量 $X$ 相同，而方差的度量单位是原单位的平方.

\sourcepageanchor{80}
\begin{figure}[htbp]\centering
\includegraphics[width=.74\textwidth]{assets/chapter2/fig2-3-1.pdf}
\caption{$\lambda=1,\frac13,\frac15$ 时的指数密度.}\label{fig:2.3.1}
\end{figure}
\begin{example}[指数分布的方差]\label{ex:2.3.3}\index{SI0297@方差}
设 $X$ 具有\autoref{ex:2.2.2}中定义的 $\operatorname{exponential}(\lambda)$ 分布. 在那里我们算出了 $\mathrm EX=\lambda$，现在可以计算方差：
$$\begin{aligned}
\operatorname{Var}X&=\mathrm E(X-\lambda)^2=\int_0^\infty(x-\lambda)^2\frac1\lambda e^{-x/\lambda}\,dx\\
&=\int_0^\infty(x^2-2x\lambda+\lambda^2)\frac1\lambda e^{-x/\lambda}\,dx.
\end{aligned}$$
为了完成积分，可以对各项分别积分，对含有 $x$ 和 $x^2$ 的项使用分部积分. 这样就得到 $\operatorname{Var}X=\lambda^2$.
\end{example}
可见，指数分布的方差与参数 $\lambda$ 直接相关. \autoref{fig:2.3.1}给出了不同 $\lambda$ 值对应的几个指数分布. 注意，$\lambda$ 越小，分布就越集中在均值周围. 指数分布的方差作为 $\lambda$ 的函数所表现出的性质，是下列定理所概括的方差性质的一个特殊情形.
\begin{theorem}\label{thm:2.3.4}\index{SI1066@性质}
如果随机变量 $X$ 的方差有限，那么对任意常数 $a$ 和 $b$，
$$\operatorname{Var}(aX+b)=a^2\operatorname{Var}X.$$
\end{theorem}
\begin{proof}
由定义，
$$\begin{aligned}
\operatorname{Var}(aX+b)&=\mathrm E((aX+b)-\mathrm E(aX+b))^2\\
&=\mathrm E(aX-a\mathrm EX)^2&&\text{（$\mathrm E(aX+b)=a\mathrm EX+b$）}\\
&=a^2\mathrm E(X-\mathrm EX)^2\\
&=a^2\operatorname{Var}X.
\end{aligned}$$
\end{proof}

\sourcepageanchor{81}
有时使用方差的另一个公式更容易：
\begin{equation}\label{eq:2.3.1}
\operatorname{Var}X=\mathrm EX^2-(\mathrm EX)^2.
\end{equation}
只需注意下式便容易证明它：
$$\begin{aligned}
\operatorname{Var}X&=\mathrm E(X-\mathrm EX)^2=\mathrm E[X^2-2X\mathrm EX+(\mathrm EX)^2]\\
&=\mathrm EX^2-2(\mathrm EX)^2+(\mathrm EX)^2\\
&=\mathrm EX^2-(\mathrm EX)^2,
\end{aligned}$$
其中利用了 $\mathrm E(X\mathrm EX)=(\mathrm EX)(\mathrm EX)=(\mathrm EX)^2$，因为 $\mathrm EX$ 是常数. 现在用一个离散分布说明矩的计算.

\begin{example}[二项分布的方差]\label{ex:2.3.5}\index{SI0272@方差}
设 $X\sim\operatorname{binomial}(n,p)$，即
$$P(X=x)=\binom nx p^x(1-p)^{n-x},\qquad x=0,1,\ldots,n.$$
我们已经知道 $\mathrm EX=np$. 为计算 $\operatorname{Var}X$，先计算 $\mathrm EX^2$. 有
\begin{equation}\label{eq:2.3.2}
\mathrm EX^2=\sum_{x=0}^n x^2\binom nx p^x(1-p)^{n-x}.
\end{equation}
为求出这个级数之和，必须先对二项系数作变换，方法与求 $\mathrm EX$ 时类似（\autoref{ex:2.2.3}）. 写成
\begin{equation}\label{eq:2.3.3}
x^2\binom nx=x\frac{n!}{(x-1)!(n-x)!}=xn\binom{n-1}{x-1}.
\end{equation}
\eqref{eq:2.3.2}中 $x=0$ 对应的项为零，利用\eqref{eq:2.3.3}，得到
$$\begin{aligned}
\mathrm EX^2&=n\sum_{x=1}^n x\binom{n-1}{x-1}p^x(1-p)^{n-x}\\
&=n\sum_{y=0}^{n-1}(y+1)\binom{n-1}y p^{y+1}(1-p)^{n-1-y}&&\text{（令 $y=x-1$）}\\
&=np\sum_{y=0}^{n-1}y\binom{n-1}y p^y(1-p)^{n-1-y}\\
&\qquad+np\sum_{y=0}^{n-1}\binom{n-1}y p^y(1-p)^{n-1-y}.
\end{aligned}$$
现在容易看出，第一个和等于 $(n-1)p$，因为它是 $\operatorname{binomial}(n-1,p)$ 分布的均值；第二个和等于 $1$. 因此，
\begin{equation}\label{eq:2.3.4}
\mathrm EX^2=n(n-1)p^2+np.
\end{equation}
由\eqref{eq:2.3.1}，得到
$$\operatorname{Var}X=n(n-1)p^2+np-(np)^2=-np^2+np=np(1-p).$$
\end{example}

高阶矩的计算方式类似，但数学运算通常会变得相当复杂. 在应用中，有时会关心三阶或四阶矩；但从统计学角度看，通常没有多少理由考察更高阶的矩.

现在介绍与概率分布相关的一个新函数——矩母函数（moment generating function，mgf）. 顾名思义，矩母函数可以用来生成矩. 实际上，很多时候直接计算矩比使用矩母函数更容易. 然而，矩母函数的主要用途并不是生成矩，而是帮助刻画分布. 恰当地使用这一性质，可以得到极为有力的结果.

\sourcepageanchor{82}
\begin{definition}\label{def:2.3.6}\index{SI0689@矩母函数}\index{SI0720@邻域}
设 $X$ 是具有累积分布函数 $F_X$ 的随机变量. $X$（或 $F_X$）的矩母函数，记为 $M_X(t)$，定义为
$$M_X(t)=\mathrm Ee^{tX},$$
前提是这个期望对 $0$ 的某个邻域内的 $t$ 存在. 也就是说，存在 $h>0$，使得对所有 $-h<t<h$，$\mathrm Ee^{tX}$ 都存在. 如果这个期望在 $0$ 的任何邻域内都不是处处存在的，就说矩母函数不存在.
\end{definition}
更明确地说，$X$ 的矩母函数可写为
$$M_X(t)=\int_{-\infty}^\infty e^{tx}f_X(x)\,dx,\qquad X\text{ 连续},$$
或者
$$M_X(t)=\sum_x e^{tx}P(X=x),\qquad X\text{ 离散}.$$
很容易看出矩母函数如何生成矩. 我们将这一结果概括为下列定理.
\begin{theorem}\label{thm:2.3.7}
如果 $X$ 的矩母函数为 $M_X(t)$，那么
$$\mathrm EX^n=M_X^{(n)}(0),$$
其中定义
$$M_X^{(n)}(0)=\left.\frac{d^n}{dt^n}M_X(t)\right|_{t=0}.$$
也就是说，$n$ 阶矩等于 $M_X(t)$ 的 $n$ 阶导数在 $t=0$ 处的值.
\end{theorem}
\begin{proof}
假定可以在积分号下求导（见\autoref{sec:2.4}），则
$$\begin{aligned}
\frac d{dt}M_X(t)&=\frac d{dt}\int_{-\infty}^\infty e^{tx}f_X(x)\,dx\\
&=\int_{-\infty}^\infty\left(\frac d{dt}e^{tx}\right)f_X(x)\,dx\\
&=\int_{-\infty}^\infty(xe^{tx})f_X(x)\,dx\\
&=\mathrm EXe^{tX}.
\end{aligned}$$
于是，
$$\left.\frac d{dt}M_X(t)\right|_{t=0}=\left.\mathrm EXe^{tX}\right|_{t=0}=\mathrm EX.$$
用类似方式可得
$$\left.\frac{d^n}{dt^n}M_X(t)\right|_{t=0}=\left.\mathrm EX^ne^{tX}\right|_{t=0}=\mathrm EX^n.$$
\end{proof}

\begin{example}[伽马分布的矩母函数]\label{ex:2.3.8}\index{SI0308@矩母函数}
在\autoref{ex:2.1.6}中，我们遇到了伽马密度的一种特殊情形；一般的伽马密度为
$$f(x)=\frac1{\Gamma(\alpha)\beta^\alpha}x^{\alpha-1}e^{-x/\beta},\qquad0<x<\infty,\quad\alpha>0,\quad\beta>0,$$
\sourcepageanchor{83}
其中 $\Gamma(\alpha)$ 表示伽马函数（gamma function），它的一些性质将在\autoref{sec:3.3}中给出. 矩母函数为
\begin{equation}\label{eq:2.3.5}
\begin{aligned}
M_X(t)&=\frac1{\Gamma(\alpha)\beta^\alpha}\int_0^\infty e^{tx}x^{\alpha-1}e^{-x/\beta}\,dx\\
&=\frac1{\Gamma(\alpha)\beta^\alpha}\int_0^\infty x^{\alpha-1}e^{-(\frac1\beta-t)x}\,dx\\
&=\frac1{\Gamma(\alpha)\beta^\alpha}\int_0^\infty x^{\alpha-1}e^{-x/(\frac\beta{1-\beta t})}\,dx.
\end{aligned}
\end{equation}
现在可以认出，\eqref{eq:2.3.5}中的被积函数是另一个伽马密度的核（kernel）\index{SI0595@核}\semanticindexanchor{595:83}. （函数的核是函数的主要部分，即忽略常数后余下的部分.）利用对任意正常数 $a,b$，
$$f(x)=\frac1{\Gamma(a)b^a}x^{a-1}e^{-x/b}$$
是概率密度函数这一事实，有
$$\int_0^\infty\frac1{\Gamma(a)b^a}x^{a-1}e^{-x/b}\,dx=1,$$
从而
\begin{equation}\label{eq:2.3.6}
\int_0^\infty x^{a-1}e^{-x/b}\,dx=\Gamma(a)b^a.
\end{equation}
将\eqref{eq:2.3.6}应用于\eqref{eq:2.3.5}，得到
$$M_X(t)=\frac1{\Gamma(\alpha)\beta^\alpha}\Gamma(\alpha)
\left(\frac\beta{1-\beta t}\right)^\alpha
=\left(\frac1{1-\beta t}\right)^\alpha,\qquad t<\frac1\beta.$$
如果 $t\ge1/\beta$，\eqref{eq:2.3.5}被积函数中的 $(1/\beta)-t$ 非正，而\eqref{eq:2.3.6}中的积分为无穷大. 因此，伽马分布的矩母函数只在 $t<1/\beta$ 时存在. （我们将在\autoref{sec:3.3}中更详细地讨论伽马函数.）

伽马分布的均值为
$$\mathrm EX=\left.\frac d{dt}M_X(t)\right|_{t=0}
=\left.\frac{\alpha\beta}{(1-\beta t)^{\alpha+1}}\right|_{t=0}=\alpha\beta.$$
其他矩可用类似方式计算.
\end{example}

\begin{example}[二项分布的矩母函数]\label{ex:2.3.9}\index{SI0265@矩母函数}
作为计算矩母函数的第二个例子，我们考虑一个离散分布——二项分布. $\operatorname{binomial}(n,p)$ 的概率质量函数由\eqref{eq:2.1.3}给出，因此
$$M_X(t)=\sum_{x=0}^n e^{tx}\binom nx p^x(1-p)^{n-x}
=\sum_{x=0}^n\binom nx(pe^t)^x(1-p)^{n-x}.$$
二项式公式（见\autoref{thm:3.2.2}）给出
\begin{equation}\label{eq:2.3.7}
\sum_{x=0}^n\binom nx u^xv^{n-x}=(u+v)^n.
\end{equation}
因此，令 $u=pe^t$、$v=1-p$，得到
$$M_X(t)=[pe^t+(1-p)]^n.$$
\end{example}

\sourcepageanchor{84}
\begin{figure}[htbp]\centering
\includegraphics[width=.76\textwidth]{assets/chapter2/fig2-3-2.pdf}
\caption{具有相同各阶矩的两个概率密度函数：$f_1(x)=\frac1{\sqrt{2\pi}x}e^{-(\log x)^2/2}$，$f_2(x)=f_1(x)[1+\sin(2\pi\log x)]$.}\label{fig:2.3.2}
\end{figure}
如前所述，矩母函数的主要用途并不是生成矩，而在于它在许多情况下能够刻画分布. 不过，用矩来刻画分布会遇到一些技术困难，现在就来研究这些困难.

如果矩母函数存在，它就刻画了无穷多个矩. 一个自然的问题是：刻画这无穷多个矩，是否就唯一确定了分布函数？遗憾的是，答案是否定的. 仅仅刻画所有矩，不足以唯一确定分布，因为两个不同的随机变量可能具有相同的各阶矩.

\begin{example}[矩的不唯一性]\label{ex:2.3.10}\index{SI0334@矩}\index{SI0698@非唯一性}
考虑两个概率密度函数
$$\begin{aligned}
f_1(x)&=\frac1{\sqrt{2\pi}x}e^{-(\log x)^2/2},&0\le x<\infty,\\
f_2(x)&=f_1(x)[1+\sin(2\pi\log x)],&0\le x<\infty.
\end{aligned}$$
（密度 $f_1$ 是对数正态密度（lognormal pdf）的一种特殊情形.）
可以证明，若 $X_1\sim f_1(x)$，则
$$\mathrm EX_1^r=e^{r^2/2},\qquad r=0,1,\ldots,$$
所以 $X_1$ 的所有阶矩都存在. 现设 $X_2\sim f_2(x)$，有
$$\begin{aligned}
\mathrm EX_2^r&=\int_0^\infty x^rf_1(x)[1+\sin(2\pi\log x)]\,dx\\
&=\mathrm EX_1^r+\int_0^\infty x^rf_1(x)\sin(2\pi\log x)\,dx.
\end{aligned}$$
但是，变换 $y=\log x-r$ 表明，最后的积分是一个奇函数在 $(-\infty,\infty)$ 上的积分，因而对 $r=0,1,\ldots$ 都等于 $0$. 所以，尽管 $X_1$ 与 $X_2$ 的密度不同，它们的所有阶矩却都相同. 两个密度画在\autoref{fig:2.3.2}中.

细节见\autoref{exercise:2.35}；有关矩母函数和分布的更多内容，另见\autoref{exercise:2.34}、\autoref{exercise:2.36}及\autoref{exercise:2.37}.
\end{example}

\sourcepageanchor{85}
如果随机变量具有有界支撑，就不会出现矩的唯一性问题. 此时，无穷矩序列确实唯一确定分布（例如，见 \hyperref[bib:source33]{Billingsley 1995}，第~30~节）. 而且，如果矩母函数在零的某个邻域内存在，无论支撑集怎样，分布都被唯一确定. 因而，所有阶矩都存在与矩母函数存在并不等价. 下列定理说明如何刻画分布.
\needspace{6\baselineskip}
\begin{theorem}\label{thm:2.3.11}\index{SI0691@同分布}\index{SI0870@矩}
设 $F_X(x)$ 和 $F_Y(y)$ 是两个各阶矩都存在的累积分布函数.
\begin{enumerate}
\item[(a)] 若 $X$ 和 $Y$ 具有有界支撑，则对所有 $u$ 有 $F_X(u)=F_Y(u)$，当且仅当对所有整数 $r=0,1,2,\ldots$ 有 $\mathrm EX^r=\mathrm EY^r$.
\item[(b)] 若矩母函数存在，并且对 $0$ 的某个邻域内的所有 $t$ 有 $M_X(t)=M_Y(t)$，则对所有 $u$ 有 $F_X(u)=F_Y(u)$.
\end{enumerate}
\end{theorem}
下一个定理讨论一个收敛的矩母函数序列，我们不再单独处理有界支撑的情形. 注意，如果极限矩母函数在 $0$ 的某个邻域内存在，则唯一性假设自动满足（\hyperref[sec:2.6.1]{杂录~2.6.1}）.
\begin{theorem}[矩母函数的收敛]\label{thm:2.3.12}\index{SI0690@收敛}
设 $\{X_i,\ i=1,2,\ldots\}$ 是随机变量序列，每个随机变量的矩母函数为 $M_{X_i}(t)$. 再设
$$\lim_{i\to\infty}M_{X_i}(t)=M_X(t),\qquad\text{对 $0$ 的某个邻域内的所有 $t$},$$
而 $M_X(t)$ 是一个矩母函数. 则存在唯一的累积分布函数 $F_X$，其矩由 $M_X(t)$ 确定；并且对 $F_X(x)$ 连续的所有 $x$，有
$$\lim_{i\to\infty}F_{X_i}(x)=F_X(x).$$
也就是说，矩母函数在 $|t|<h$ 时收敛到一个矩母函数，蕴含累积分布函数的收敛.
\end{theorem}
\autoref{thm:2.3.11}和\autoref{thm:2.3.12}的证明依赖于拉普拉斯变换（Laplace transform）理论. （经典参考文献是 \hyperref[bib:source329]{Widder (1946)}，\hyperref[bib:source119]{Feller (1971)} 也对拉普拉斯变换作了全面介绍.）$M_X(t)$ 的定义式，即
\begin{equation}\label{eq:2.3.8}\index{SI0601@拉普拉斯变换}
M_X(t)=\int_{-\infty}^\infty e^{tx}f_X(x)\,dx,
\end{equation}
定义了一个拉普拉斯变换（$M_X(t)$ 是 $f_X(x)$ 的拉普拉斯变换）. 拉普拉斯变换的一个关键性质是唯一性. 若对满足 $|t|<h$ 的所有 $t$，\eqref{eq:2.3.8}都成立，其中 $h$ 为某个正数，那么给定 $M_X(t)$ 后，只有一个函数 $f_X(x)$ 满足\eqref{eq:2.3.8}. 有了这个事实，前面两个定理就相当合理了. 尽管它们的严格证明没有超出本书的范围，但证明较为技术化，并不能带来真正的理解，因此这里略去.

矩序列可能不具有唯一性，是件麻烦事. 如果证明某个矩序列收敛，我们还不能正式断言随机变量收敛. 为了得出这一结论，还必须核实矩序列的唯一性，而这通常是一项可怕的工作（\hyperref[sec:2.6.1]{杂录~2.6.1}）. 但是，如果矩母函数序列在 $0$ 的某个邻域内收敛，那么随机变量就收敛. 因此，矩母函数的收敛可以视为随机变量序列收敛的一个充分但非必要条件.

\begin{example}[泊松近似]\label{ex:2.3.13}\index{SI0060@以泊松近似二项}
初等统计课程通常介绍的一种近似是：二项概率（见\autoref{ex:2.3.5}）可以用通常更容易计算的泊松概率（Poisson probabilities）来近似. 二项分布由两个量 $n$ 和 $p$ 刻画. 通常的说法是，当“$n$ 大而 $np$ 小”时，泊松近似有效，有时还会给出经验规则.
\sourcepageanchor{86}
$\operatorname{Poisson}(\lambda)$ 的概率质量函数为
$$P(X=x)=\frac{e^{-\lambda}\lambda^x}{x!},\qquad x=0,1,2,\ldots,$$
其中 $\lambda$ 为正常数. 这个近似是说，如果 $X\sim\operatorname{binomial}(n,p)$，$Y\sim\operatorname{Poisson}(\lambda)$，且 $\lambda=np$，则对大的 $n$ 和小的 $np$，
\begin{equation}\label{eq:2.3.9}
P(X=x)\approx P(Y=x).
\end{equation}
现在证明矩母函数收敛，从而为这一近似提供依据. 回顾
\begin{equation}\label{eq:2.3.10}
M_X(t)=[pe^t+(1-p)]^n.
\end{equation}
对于 $\operatorname{Poisson}(\lambda)$ 分布，可以算出（见\autoref{exercise:2.33}）
$$M_Y(t)=e^{\lambda(e^t-1)},$$
而如果定义 $p=\lambda/n$，当 $n\to\infty$ 时就有 $M_X(t)\to M_Y(t)$. \eqref{eq:2.3.9}中近似的有效性于是由\autoref{thm:2.3.12}得到.

我们先稍作离题，介绍一个在统计学中有广泛应用的重要极限结果. 这个引理的证明可在许多标准微积分教材中找到.
\begin{lemma}\label{lem:2.3.14}\index{SI0372@矩母函数}
设数列 $a_1,a_2,\ldots$ 收敛到 $a$，即 $\lim_{n\to\infty}a_n=a$. 则
$$\lim_{n\to\infty}\left(1+\frac{a_n}n\right)^n=e^a.$$
\end{lemma}
回到例子，有
$$\begin{aligned}
M_X(t)&=[pe^t+(1-p)]^n\\
&=\left[1+\frac1n(e^t-1)(np)\right]^n
=\left[1+\frac1n(e^t-1)\lambda\right]^n,
\end{aligned}$$
因为 $\lambda=np$. 现在令 $a_n=a=(e^t-1)\lambda$，应用\autoref{lem:2.3.14}，得到
$$\lim_{n\to\infty}M_X(t)=e^{\lambda(e^t-1)}=M_Y(t),$$
即泊松分布的矩母函数.

即使 $p$ 和 $n$ 的大小适中，泊松近似也可能相当好. \autoref{fig:2.3.3}给出了一个二项质量函数及其 $\lambda=np$ 的泊松近似. 这个近似看起来令人满意.
\end{example}

本节最后给出关于矩母函数的一个有用结果.
\begin{theorem}\label{thm:2.3.15}\index{SI0693@性质}
对任意常数 $a,b$，随机变量 $aX+b$ 的矩母函数为
$$M_{aX+b}(t)=e^{bt}M_X(at).$$
\end{theorem}
\begin{proof}
由定义，
$$\begin{aligned}
M_{aX+b}(t)&=\mathrm E\bigl(e^{(aX+b)t}\bigr)\\
&=\mathrm E\bigl(e^{(aX)t}e^{bt}\bigr)&&\text{（指数函数的性质）}\\
&=e^{bt}\mathrm E\bigl(e^{(at)X}\bigr)&&\text{（$e^{bt}$ 是常数）}\\
&=e^{bt}M_X(at),&&\text{（矩母函数的定义）}
\end{aligned}$$
定理得证.
\end{proof}
\sourcepageanchor{87}
\begin{figure}[htbp]\centering
\includegraphics[width=.85\textwidth]{assets/chapter2/fig2-3-3.pdf}
\caption{二项分布（实线）的泊松近似（虚线），$n=15$，$p=.3$.}\label{fig:2.3.3}
\end{figure}



\subsection{积分号下求导}\label{sec:2.4}
上一节遇到了一种希望交换积分与求导次序的情况. 在理论统计学中，这种情况经常出现. 本节的目的是刻画这种运算成立的条件. 我们也将讨论交换求导与求和次序的问题.

许多条件可由微积分中的标准定理得到，详细证明可在大多数微积分教材中找到，因此这里不再给出详细证明.

我们首先要确立计算
\begin{equation}\label{eq:2.4.1}
\frac d{d\theta}\int_{a(\theta)}^{b(\theta)}f(x,\theta)\,dx
\end{equation}
的方法，其中对所有 $\theta$ 都有 $-\infty<a(\theta),b(\theta)<\infty$. 对\eqref{eq:2.4.1}求导的法则称为 Leibnitz 法则（Leibnitz's Rule），它是微积分基本定理和链式法则的一个应用.
\begin{theorem}[Leibnitz 法则]\label{thm:2.4.1}\index{SI0618@Leibnitz法则}
如果 $f(x,\theta)$、$a(\theta)$ 和 $b(\theta)$ 关于 $\theta$ 可微，则
$$\begin{aligned}
\frac d{d\theta}\int_{a(\theta)}^{b(\theta)}f(x,\theta)\,dx
&=f(b(\theta),\theta)\frac d{d\theta}b(\theta)
-f(a(\theta),\theta)\frac d{d\theta}a(\theta)\\
&\qquad+\int_{a(\theta)}^{b(\theta)}\frac\partial{\partial\theta}f(x,\theta)\,dx.
\end{aligned}$$
\end{theorem}
注意，如果 $a(\theta)$ 和 $b(\theta)$ 为常数，就得到 Leibnitz 法则的一个特殊情形：
$$\frac d{d\theta}\int_a^b f(x,\theta)\,dx
=\int_a^b\frac\partial{\partial\theta}f(x,\theta)\,dx.$$
因此，一般而言，如果是在有限区间上对可微函数积分，对积分求导不会有问题. 然而，如果积分区间无穷，就可能出现问题.\footnote{译注：原文此处以及\autoref{thm:2.4.1}仅列“关于 $\theta$ 可微”，一般不足以保证交换求导与积分. 常见充分条件还要求 $f$ 及其关于 $\theta$ 的偏导在相关紧区域上连续，或满足后文的可积控制条件；正文保留原述.}

注意，上式交换求导与积分时，一边是偏导数，另一边是普通导数. 形式上必须如此，因为左边只是 $\theta$ 的函数，而右边的被积函数同时依赖 $\theta$ 和 $x$.

\sourcepageanchor{88}
交换求导与积分次序是否成立，实际上就是能否交换极限与积分的问题，因为导数是一种特殊的极限. 回顾，若 $f(x,\theta)$ 可微，则
$$\frac\partial{\partial\theta}f(x,\theta)
=\lim_{\delta\to0}\frac{f(x,\theta+\delta)-f(x,\theta)}\delta,$$
所以
$$\int_{-\infty}^\infty\frac\partial{\partial\theta}f(x,\theta)\,dx
=\int_{-\infty}^\infty\lim_{\delta\to0}\left[\frac{f(x,\theta+\delta)-f(x,\theta)}\delta\right]dx,$$
而
$$\frac d{d\theta}\int_{-\infty}^\infty f(x,\theta)\,dx
=\lim_{\delta\to0}\int_{-\infty}^\infty\left[\frac{f(x,\theta+\delta)-f(x,\theta)}\delta\right]dx.$$
因此，如果能证明可以交换极限与积分次序，就能证明可以在积分号下求导. 遗憾的是，要完全一般地处理这个问题，必须使用本书不涉及的测度论. 不过，我们可以给出一些重要结果的陈述和结论. 以下定理都是 Lebesgue 控制收敛定理（Dominated Convergence Theorem）\index{SI0616@勒贝格控制收敛定理}\semanticindexanchor{616:88}的推论（例如，见 \hyperref[bib:source267]{Rudin (1976)}）.

\begin{theorem}\label{thm:2.4.2}
设对每个 $x$，函数 $h(x,y)$ 都在 $y_0$ 处连续，并且存在函数 $g(x)$，满足
\begin{enumerate}
\item[(i)] 对所有 $x,y$，$|h(x,y)|\le g(x)$；
\item[(ii)] $\int_{-\infty}^\infty g(x)\,dx<\infty$.
\end{enumerate}
则
$$\lim_{y\to y_0}\int_{-\infty}^\infty h(x,y)\,dx
=\int_{-\infty}^\infty\lim_{y\to y_0}h(x,y)\,dx.$$
\end{theorem}
这个定理的关键条件是存在积分有限的控制函数 $g(x)$，它保证积分的行为不会过于恶劣. 令 $h(x,y)$ 对应差商 $(f(x,\theta+\delta)-f(x,\theta))/\delta$，便可以把这个定理用于我们正在考虑的情形.

\begin{theorem}\label{thm:2.4.3}
设 $f(x,\theta)$ 在 $\theta=\theta_0$ 处可微，即对每个 $x$，
$$\lim_{\delta\to0}\frac{f(x,\theta_0+\delta)-f(x,\theta_0)}\delta
=\left.\frac\partial{\partial\theta}f(x,\theta)\right|_{\theta=\theta_0}$$
存在，并且存在函数 $g(x,\theta_0)$ 和常数 $\delta_0>0$，使得
\begin{enumerate}
\item[(i)]\phantomsection\label{itm:thm:2.4.3:i} 对所有 $x$ 和 $|\delta|\le\delta_0$，
$$\left|\frac{f(x,\theta_0+\delta)-f(x,\theta_0)}\delta\right|\le g(x,\theta_0);$$
\item[(ii)]\phantomsection\label{itm:thm:2.4.3:ii} $\int_{-\infty}^\infty g(x,\theta_0)\,dx<\infty$.
\end{enumerate}
则
\begin{equation}\label{eq:2.4.2}
\left.\frac d{d\theta}\int_{-\infty}^\infty f(x,\theta)\,dx\right|_{\theta=\theta_0}
=\int_{-\infty}^\infty\left[\left.\frac\partial{\partial\theta}f(x,\theta)\right|_{\theta=\theta_0}\right]dx.
\end{equation}
\end{theorem}
条件 \hyperref[itm:thm:2.4.3:i]{(i)} 类似于所谓的 Lipschitz 条件（Lipschitz condition），这是一种对函数施加光滑性要求的条件. 这里条件 \hyperref[itm:thm:2.4.3:i]{(i)} 实际上约束了一阶导数的变化；其他光滑性约束可能用常数而不是函数 $g$ 来约束这种变化，或者约束 $f$ 的二阶导数的变化.

\autoref{thm:2.4.3}的结论写起来略显繁琐，但必须认识到，尽管看起来把 $\theta$ 当作变量，定理的陈述实际上针对 $\theta$ 的一个取值.
\sourcepageanchor{89}
也就是说，对每个使 $f(x,\theta)$ 在该处可微且满足条件 \hyperref[itm:thm:2.4.3:i]{(i)}、\hyperref[itm:thm:2.4.3:ii]{(ii)} 的 $\theta_0$，都可以交换积分与求导次序. 通常不特别强调 $\theta$ 与 $\theta_0$ 的区别，而将\eqref{eq:2.4.2}写为
\begin{equation}\label{eq:2.4.3}
\frac d{d\theta}\int_{-\infty}^\infty f(x,\theta)\,dx
=\int_{-\infty}^\infty\frac\partial{\partial\theta}f(x,\theta)\,dx.
\end{equation}
通常，$f(x,\theta)$ 在所有 $\theta$ 处都可微，而不仅在一个值 $\theta_0$ 处可微. 这时，可以用一个通常更容易验证的条件代替\autoref{thm:2.4.3}的条件 \hyperref[itm:thm:2.4.3:i]{(i)}. 应用中值定理，对于固定的 $x,\theta_0$ 和 $|\delta|\le\delta_0$，有
$$\frac{f(x,\theta_0+\delta)-f(x,\theta_0)}\delta
=\left.\frac\partial{\partial\theta}f(x,\theta)\right|_{\theta=\theta_0+\delta^*(x)},$$
其中某个数 $\delta^*(x)$ 满足 $|\delta^*(x)|\le\delta_0$. 因此，只要找到满足条件 \hyperref[itm:thm:2.4.3:ii]{(ii)} 和下式的 $g(x,\theta)$，条件 \hyperref[itm:thm:2.4.3:i]{(i)} 就成立：
\begin{equation}\label{eq:2.4.4}
\left|\left.\frac\partial{\partial\theta}f(x,\theta)\right|_{\theta=\theta'}\right|
\le g(x,\theta),\qquad\text{对所有满足 $|\theta'-\theta|\le\delta_0$ 的 $\theta'$}.
\end{equation}
注意，与\autoref{thm:2.4.3}一样，\eqref{eq:2.4.4}中的 $\delta_0$ 暗含地依赖于 $\theta$. 这是允许的，因为定理是逐个应用于 $\theta$ 的每个取值的. 由\eqref{eq:2.4.4}得到下列推论.
\begin{corollary}\label{cor:2.4.4}
设 $f(x,\theta)$ 关于 $\theta$ 可微. 若存在函数 $g(x,\theta)$ 满足\eqref{eq:2.4.4}，且 $\int_{-\infty}^\infty g(x,\theta)\,dx<\infty$，则\eqref{eq:2.4.3}成立.
\end{corollary}
注意，\autoref{thm:2.4.3}的条件 \hyperref[itm:thm:2.4.3:i]{(i)} 和\eqref{eq:2.4.4}都对被控制的函数施加了一致性要求；交换求导与积分之前，一般需要某种一致性.

\begin{example}[交换积分与求导——I]\label{ex:2.4.5}
设 $X$ 具有指数密度 $f(x)=(1/\lambda)e^{-x/\lambda}$，$0<x<\infty$，即服从 $\operatorname{exponential}(\lambda)$ 分布. 假设我们要对正整数 $n$ 计算
\begin{equation}\label{eq:2.4.5}
\frac d{d\lambda}\mathrm EX^n=\frac d{d\lambda}\int_0^\infty x^n\left(\frac1\lambda\right)e^{-x/\lambda}\,dx.
\end{equation}
如果可以把求导移到积分号内，就有
\begin{equation}\label{eq:2.4.6}
\begin{aligned}
\frac d{d\lambda}\mathrm EX^n&=\int_0^\infty\frac\partial{\partial\lambda}x^n\left(\frac1\lambda\right)e^{-x/\lambda}\,dx\\
&=\int_0^\infty\frac{x^n}{\lambda^2}\left(\frac x\lambda-1\right)e^{-x/\lambda}\,dx\\
&=\frac1{\lambda^2}\mathrm EX^{n+1}-\frac1\lambda\mathrm EX^n.
\end{aligned}
\end{equation}
为证明交换积分与求导的合理性，我们要控制 $x^n(1/\lambda)e^{-x/\lambda}$ 的导数. 现在有\footnote{译注：原文左端被求导的指数印为 $e^{-x\lambda}$，而本例其余各处及右端计算均使用 $e^{-x/\lambda}$；左端疑漏除号，正文保留原式.}
$$\begin{aligned}
\left|\frac\partial{\partial\lambda}\left(\frac{x^ne^{-x\lambda}}\lambda\right)\right|
&=\frac{x^ne^{-x/\lambda}}{\lambda^2}\left|\frac x\lambda-1\right|\\
&\le\frac{x^ne^{-x/\lambda}}{\lambda^2}\left(\frac x\lambda+1\right),\qquad\text{（因为 $x/\lambda>0$）}.
\end{aligned}$$

对某个满足 $0<\delta_0<\lambda$ 的常数 $\delta_0$，取
$$g(x,\lambda)=\frac{x^ne^{-x/(\lambda+\delta_0)}}{(\lambda-\delta_0)^2}
\left(\frac{x}{\lambda-\delta_0}+1\right).$$
\sourcepageanchor{90}
于是有
$$\left|\left.\frac\partial{\partial\lambda}\left(\frac{x^ne^{-x/\lambda}}\lambda\right)\right|_{\lambda=\lambda'}\right|
\le g(x,\lambda),\qquad\text{对所有满足 $|\lambda'-\lambda|\le\delta_0$ 的 $\lambda'$}.$$
由于指数分布的各阶矩都存在，只要 $\lambda-\delta_0>0$，就有 $\int_{-\infty}^\infty g(x,\lambda)\,dx<\infty$，所以可以交换积分与求导.
\end{example}
上面用指数分布说明的性质，对一大类密度都成立，\autoref{sec:3.4}将讨论这一类密度.

注意，\eqref{eq:2.4.6}给出了指数分布的矩的递推关系，
\begin{equation}\label{eq:2.4.7}\index{SI0294@递推关系}
\mathrm EX^{n+1}=\lambda\mathrm EX^n+\lambda^2\frac d{d\lambda}\mathrm EX^n,
\end{equation}
使 $(n+1)$ 阶矩的计算相对容易. 其他分布也有这种关系.\index{SI0364@递推关系}\semanticindexanchor{364:90} 特别地，如果 $X$ 具有均值为 $\mu$、方差为 $1$ 的正态分布，概率密度函数为 $f(x)=(1/\sqrt{2\pi})e^{-(x-\mu)^2/2}$，\footnote{译注：原文递推式中导数前为负号. 由 $\partial f(x;\mu)/\partial\mu=(x-\mu)f(x;\mu)$ 可得 $\frac d{d\mu}\mathrm EX^n=\mathrm EX^{n+1}-\mu\mathrm EX^n$，因此该负号应为正号；取 $n=1$ 时原式也会给出 $\mathrm EX^2=\mu^2-1$，与方差为 $1$ 矛盾.}则
$$\mathrm EX^{n+1}=\mu\mathrm EX^n-\frac d{d\mu}\mathrm EX^n.$$


我们再举一个交换求导与积分的例子，这次涉及矩母函数.
\begin{example}[交换积分与求导——II]\label{ex:2.4.6}
仍设 $X$ 具有均值为 $\mu$、方差为 $1$ 的正态分布，考虑 $X$ 的矩母函数，
$$M_X(t)=\mathrm Ee^{tX}=\frac1{\sqrt{2\pi}}\int_{-\infty}^\infty e^{tx}e^{-(x-\mu)^2/2}\,dx.$$
\autoref{sec:2.3}指出，可以通过对 $M_X(t)$ 求导来计算矩，且积分号下求导是成立的：
\begin{equation}\label{eq:2.4.8}
\frac d{dt}M_X(t)=\frac d{dt}\mathrm Ee^{tX}
=\mathrm E\frac\partial{\partial t}e^{tX}=\mathrm E(Xe^{tX}).
\end{equation}
可以应用本节的结果证明\eqref{eq:2.4.8}中运算的合理性. 注意，这里无论应用\autoref{thm:2.4.3}还是\autoref{cor:2.4.4}，都是把 $t$ 对应于\autoref{thm:2.4.3}中的变量 $\theta$，而参数 $\mu$ 作为常数处理.

由\autoref{cor:2.4.4}，必须找到一个积分有限的函数 $g(x,t)$，满足
\begin{equation}\label{eq:2.4.9}
\left|\left.\frac\partial{\partial t}e^{tx}e^{-(x-\mu)^2/2}\right|_{t=t'}\right|
\le g(x,t),\qquad\text{对所有满足 $|t'-t|\le\delta_0$ 的 $t'$}.
\end{equation}
直接计算，得到
$$\left|\frac\partial{\partial t}e^{tx}e^{-(x-\mu)^2/2}\right|
=\left|xe^{tx}e^{-(x-\mu)^2/2}\right|
\le|x|e^{tx}e^{-(x-\mu)^2/2}.$$
最简单的办法是对 $x\ge0$ 和 $x<0$ 分别定义 $g(x,t)$. 取
$$g(x,t)=\begin{cases}
|x|e^{(t-\delta_0)x}e^{-(x-\mu)^2/2},&x<0,\\
|x|e^{(t+\delta_0)x}e^{-(x-\mu)^2/2},&x\ge0.
\end{cases}$$
显然，这个函数满足\eqref{eq:2.4.9}；只需检验它的积分有限.
\sourcepageanchor{91}
对 $x\ge0$，有
$$g(x,t)=xe^{-(x^2-2x(\mu+t+\delta_0)+\mu^2)/2}.$$
现在对指数配方，即写成
$$\begin{aligned}
&x^2-2x(\mu+t+\delta_0)+\mu^2\\
&\quad=x^2-2x(\mu+t+\delta_0)+(\mu+t+\delta_0)^2-(\mu+t+\delta_0)^2+\mu^2\\
&\quad=(x-(\mu+t+\delta_0))^2+\mu^2-(\mu+t+\delta_0)^2,
\end{aligned}$$
因而，对 $x\ge0$，
$$g(x,t)=xe^{-[x-(\mu+t+\delta_0)]^2/2}e^{-[\mu^2-(\mu+t+\delta_0)^2]/2}.$$
由于这个表达式中最后一个指数因子不依赖 $x$，$\int_0^\infty g(x,t)\,dx$ 实质上是在计算均值为 $\mu+t+\delta_0$ 的正态分布的均值，只不过积分仅限于 $[0,\infty)$. 然而，正态分布的均值有限（将在\autoref{sec:3}中证明），因而这个积分有限. 对 $x<0$ 作类似推导，得到
$$g(x,t)=|x|e^{-[x-(\mu+t-\delta_0)]^2/2}e^{-[\mu^2-(\mu+t-\delta_0)^2]/2},$$
所以 $\int_{-\infty}^0g(x,t)\,dx<\infty$. 因此，我们找到了满足\eqref{eq:2.4.9}的可积函数，\eqref{eq:2.4.8}中的运算成立.
\end{example}

现在转向何时可以交换求导与求和的问题，这种运算对离散分布具有重要作用. 当然，我们只关心无穷和，因为有限和总可以逐项求导.
\begin{example}[交换求和与求导]\label{ex:2.4.7}\index{SI0459@求导}
设 $X$ 是具有几何分布（geometric distribution）的离散随机变量，
$$P(X=x)=\theta(1-\theta)^x,\qquad x=0,1,\ldots,\quad0<\theta<1.$$
有 $\sum_{x=0}^\infty\theta(1-\theta)^x=1$；如果下列运算成立，则
$$\begin{aligned}
\frac d{d\theta}\sum_{x=0}^\infty\theta(1-\theta)^x
&=\sum_{x=0}^\infty\frac d{d\theta}\theta(1-\theta)^x\\
&=\sum_{x=0}^\infty[(1-\theta)^x-\theta x(1-\theta)^{x-1}]\\
&=\frac1\theta\sum_{x=0}^\infty\theta(1-\theta)^x
-\frac1{1-\theta}\sum_{x=0}^\infty x\theta(1-\theta)^x.
\end{aligned}$$
因为对所有 $0<\theta<1$，$\sum_{x=0}^\infty\theta(1-\theta)^x=1$，其导数为零，所以
\begin{equation}\label{eq:2.4.10}
\frac1\theta\sum_{x=0}^\infty\theta(1-\theta)^x
-\frac1{1-\theta}\sum_{x=0}^\infty x\theta(1-\theta)^x=0.
\end{equation}
\eqref{eq:2.4.10}中的第一个和等于 $1$，第二个和为 $\mathrm EX$，因而该式变为
$$\frac1\theta-\frac1{1-\theta}\mathrm EX=0,$$
或者
$$\mathrm EX=\frac{1-\theta}\theta.$$
实质上，我们通过求导求出了级数 $\sum_{x=0}^\infty x\theta(1-\theta)^x$ 的和.
\end{example}

\sourcepageanchor{92}
证明可以在求和号内求导，比积分的情况更直接. 下列定理给出详细条件.
\begin{theorem}\label{thm:2.4.8}\index{SI0585@求导与求和}
设级数 $\sum_{x=0}^\infty h(\theta,x)$ 对实数区间 $(a,b)$ 中的所有 $\theta$ 都收敛，并且
\begin{enumerate}
\item[(i)] 对每个 $x$，$\frac\partial{\partial\theta}h(\theta,x)$ 关于 $\theta$ 连续；
\item[(ii)] $\sum_{x=0}^\infty\frac\partial{\partial\theta}h(\theta,x)$ 在 $(a,b)$ 的每个有界闭子区间上一致收敛.
\end{enumerate}
则
\begin{equation}\label{eq:2.4.11}
\frac d{d\theta}\sum_{x=0}^\infty h(\theta,x)
=\sum_{x=0}^\infty\frac\partial{\partial\theta}h(\theta,x).
\end{equation}
\end{theorem}
要证明可以在求和号内求导，关键是验证一致收敛条件. 回顾，如果一个级数的部分和序列一致收敛，则该级数一致收敛. 下例将使用这一事实.

\begin{example}[续\autoref{ex:2.4.7}]\label{ex:2.4.9}
\footnote{译注：原文标题写作“Continuation of Exercise 2.4.7”，但本书习题采用两段编号，且下文承接的是\autoref{ex:2.4.7}；这里将引用目标链接到该例.}
为应用\autoref{thm:2.4.8}，令
$$h(\theta,x)=\theta(1-\theta)^x,$$
以及
$$\frac\partial{\partial\theta}h(\theta,x)=(1-\theta)^x-\theta x(1-\theta)^{x-1},$$
并验证 $\sum_{x=0}^\infty\frac\partial{\partial\theta}h(\theta,x)$ 一致收敛. 定义
$$S_n(\theta)=\sum_{x=0}^n[(1-\theta)^x-\theta x(1-\theta)^{x-1}].$$
如果给定 $\varepsilon>0$，可以找到 $N$，使得
$$n>N\ \Longrightarrow\ |S_n(\theta)-S_\infty(\theta)|<\varepsilon
\qquad\text{对所有 $\theta\in[c,d]$},$$
那么收敛在 $[c,d]\subset(0,1)$ 上就是一致的. 回顾几何级数的部分和\eqref{eq:1.5.3}，若 $y\ne1$，则
$$\sum_{k=0}^n y^k=\frac{1-y^{n+1}}{1-y}.$$
应用它，得到
$$\begin{aligned}
\sum_{x=0}^n(1-\theta)^x&=\frac{1-(1-\theta)^{n+1}}\theta,\\
\sum_{x=0}^n\theta x(1-\theta)^{x-1}
&=\theta\sum_{x=0}^n-\frac\partial{\partial\theta}(1-\theta)^x\\
&=-\theta\frac d{d\theta}\sum_{x=0}^n(1-\theta)^x\\
&=-\theta\frac d{d\theta}\left[\frac{1-(1-\theta)^{n+1}}\theta\right].
\end{aligned}$$
\sourcepageanchor{93}
这里把导数移到有限求和号外是合理的. 计算导数，得到
$$\sum_{x=0}^n\theta x(1-\theta)^{x-1}
=\frac{(1-(1-\theta)^{n+1})-(n+1)\theta(1-\theta)^n}\theta,$$
从而
$$\begin{aligned}
S_n(\theta)&=\frac{1-(1-\theta)^{n+1}}\theta
-\frac{(1-(1-\theta)^{n+1})-(n+1)\theta(1-\theta)^n}\theta\\
&=(n+1)(1-\theta)^n.
\end{aligned}$$
显然，对于 $0<\theta<1$，$S_\infty=\lim_{n\to\infty}S_n(\theta)=0$. 由于 $S_n(\theta)$ 连续，收敛在任意有界闭区间上都是一致的.\footnote{译注：连续函数序列逐点收敛，并不能仅凭连续性推出一致收敛. 本例的结论成立，但还需用显式上界：对 $[c,d]\subset(0,1)$，$\sup_{\theta\in[c,d]}|S_n(\theta)|\le(n+1)(1-c)^n\to0$. 原文理由在此不够充分.} 因此，导数级数一致收敛，可以交换求导与求和.
\end{example}

本节最后给出一个类似于\autoref{thm:2.4.8}的定理，但它处理的是交换求和与积分次序的情况.
\begin{theorem}\label{thm:2.4.10}\index{SI0586@积分与求和}
设级数 $\sum_{x=0}^\infty h(\theta,x)$ 在 $[a,b]$ 上一致收敛，并且对每个 $x$，$h(\theta,x)$ 是 $\theta$ 的连续函数. 则
$$\int_a^b\sum_{x=0}^\infty h(\theta,x)\,d\theta
=\sum_{x=0}^\infty\int_a^b h(\theta,x)\,d\theta.$$
\end{theorem}



\subsection{习题}\label{sec:2.5}
\setcounter{exercise}{0}
\begin{exercise}\label{exercise:2.1}
在下列各情形中求 $Y$ 的概率密度函数，并证明其积分为 $1$.
\begin{enumerate}
\item[(a)] $Y=X^3$，$f_X(x)=42x^5(1-x)$，$0<x<1$.
\item[(b)] $Y=4X+3$，$f_X(x)=7e^{-7x}$，$0<x<\infty$.
\item[(c)] $Y=X^2$，$f_X(x)=30x^2(1-x)^2$，$0<x<1$.
\end{enumerate}
（见计算机代数附录中的\autoref{ex:12.6.2}.）
\end{exercise}
\begin{exercise}\label{exercise:2.2}
在下列各情形中求 $Y$ 的概率密度函数.
\begin{enumerate}
\item[(a)] $Y=X^2$，$f_X(x)=1$，$0<x<1$.
\item[(b)] $Y=-\log X$，$X$ 的密度为
$$f_X(x)=\frac{(n+m+1)!}{n!m!}x^n(1-x)^m,\qquad0<x<1,\quad m,n\text{ 为正整数}.$$
\item[(c)] $Y=e^X$，$X$ 的密度为
$$f_X(x)=\frac1{\sigma^2}xe^{-(x/\sigma)^2/2},\qquad0<x<\infty,\quad\sigma^2\text{ 为正常数}.$$
\end{enumerate}
\end{exercise}
\begin{exercise}\label{exercise:2.3}
设 $X$ 具有几何概率质量函数 $f_X(x)=\frac13(\frac23)^x$，$x=0,1,2,\ldots$. 求 $Y=X/(X+1)$ 的概率分布. 注意，这里 $X$ 和 $Y$ 都是离散随机变量. 要确定 $Y$ 的概率分布，应给出它的概率质量函数.
\end{exercise}
\begin{exercise}\label{exercise:2.4}
设 $\lambda$ 为固定正常数，定义函数 $f(x)$：当 $x\ge0$ 时，$f(x)=\frac12\lambda e^{-\lambda x}$；当 $x<0$ 时，$f(x)=\frac12\lambda e^{\lambda x}$.
\begin{enumerate}
\item[(a)] 验证 $f(x)$ 是概率密度函数.
\item[(b)] 若随机变量 $X$ 的密度为 $f(x)$，求所有 $t$ 对应的 $P(X<t)$. 算出所有积分.
\item[(c)] 求所有 $t$ 对应的 $P(|X|<t)$. 算出所有积分.
\end{enumerate}
\end{exercise}
\sourcepageanchor{94}
\begin{exercise}\label{exercise:2.5}
用\autoref{thm:2.1.8}求\autoref{ex:2.1.2}中 $Y$ 的概率密度函数. 证明，对\eqref{eq:2.1.6}给出的累积分布函数求导会得到相同答案.
\end{exercise}
\begin{exercise}\label{exercise:2.6}
在下列各情形中求 $Y$ 的概率密度函数，并证明其积分为 $1$.
\begin{enumerate}
\item[(a)] $f_X(x)=\frac12e^{-|x|}$，$-\infty<x<\infty$；$Y=|X|^3$.
\item[(b)] $f_X(x)=\frac38(x+1)^2$，$-1<x<1$；$Y=1-X^2$.
\item[(c)] $f_X(x)=\frac38(x+1)^2$，$-1<x<1$；当 $X\le0$ 时 $Y=1-X^2$，当 $X>0$ 时 $Y=1-X$.
\end{enumerate}
\end{exercise}
\begin{exercise}\label{exercise:2.7}
设 $X$ 的概率密度函数为 $f_X(x)=\frac29(x+1)$，$-1\le x\le2$.
\begin{enumerate}
\item[(a)]\phantomsection\label{itm:exercise:2.7:a} 求 $Y=X^2$ 的概率密度函数. 注意，\autoref{thm:2.1.8}不能直接用于本题.
\item[(b)] 证明，当集合 $A_0,A_1,\ldots,A_k$ 包含 $\mathscr X$ 时，\autoref{thm:2.1.8}仍然成立；并取 $A_0=\varnothing$、$A_1=(-2,0)$、$A_2=(0,2)$，用这一推广解答 \hyperref[itm:exercise:2.7:a]{(a)}.
\end{enumerate}
\end{exercise}
\begin{exercise}\label{exercise:2.8}
对下列各情形，证明给定函数是累积分布函数，并求 $F_X^{-1}(y)$.
\begin{enumerate}
\item[(a)] $F_X(x)=\begin{cases}0,&x<0,\\1-e^{-x},&x\ge0.\end{cases}$
\item[(b)] $F_X(x)=\begin{cases}e^x/2,&x<0,\\1/2,&0\le x<1,\\1-(e^{1-x}/2),&1\le x.\end{cases}$
\item[(c)]\phantomsection\label{itm:exercise:2.8:c} $F_X(x)=\begin{cases}e^x/4,&x<0,\\1-(e^{-x}/4),&x\ge0.\end{cases}$
\end{enumerate}
注意，\hyperref[itm:exercise:2.8:c]{(c)} 中 $F_X(x)$ 不连续，但\eqref{eq:2.1.13}仍然是 $F_X^{-1}(y)$ 的合适定义.
\end{exercise}
\begin{exercise}\label{exercise:2.9}
若随机变量 $X$ 的概率密度函数为
$$f(x)=\begin{cases}\dfrac{x-1}2,&1<x<3,\\0,&\text{其他},\end{cases}$$
求一个单调函数 $u(x)$，使随机变量 $Y=u(X)$ 具有 $\operatorname{uniform}(0,1)$ 分布.
\end{exercise}
\begin{exercise}\label{exercise:2.10}\index{SI0849@离散}\index{SI0988@较大}
\autoref{thm:2.1.10}证明了概率积分变换，把均匀累积分布函数与任意连续累积分布函数联系起来. 本题研究离散随机变量与均匀随机变量之间的关系. 设 $X$ 是累积分布函数为 $F_X(x)$ 的离散随机变量，定义 $Y=F_X(X)$.
\begin{enumerate}
\item[(a)] 证明，$Y$ 在随机意义下大于 $\operatorname{uniform}(0,1)$；即若 $U\sim\operatorname{uniform}(0,1)$，则
$$\begin{aligned}
P(Y>y)&\ge P(U>y)=1-y,&&\text{对所有 $0<y<1$},\\
P(Y>y)&>P(U>y)=1-y,&&\text{对某个 $0<y<1$}.
\end{aligned}$$
（回顾，\autoref{exercise:1.49}定义了“在随机意义下大于”.）
\item[(b)] 等价地，证明 $Y$ 的累积分布函数满足：对所有 $0<y<1$，$F_Y(y)\le y$，并且对某个 $0<y<1$，$F_Y(y)<y$. （提示：设 $x_0$ 为 $F_X$ 的一个跳跃点，定义 $y_0=F_X(x_0)$. 证明 $P(Y\le y_0)=y_0$. 然后考虑 $y=y_0+\varepsilon$，建立所需不等式. 画累积分布函数的图会有帮助.）
\end{enumerate}
\end{exercise}
\begin{exercise}\label{exercise:2.11}
设 $X$ 具有标准正态密度 $f_X(x)=(1/\sqrt{2\pi})e^{-x^2/2}$.
\begin{enumerate}
\item[(a)] 直接求 $\mathrm EX^2$；然后利用\autoref{ex:2.1.7}中 $Y=X^2$ 的密度，计算 $\mathrm EY$，再次求出结果.
\item[(b)] 求 $Y=|X|$ 的概率密度函数，并求它的均值和方差.
\end{enumerate}
\end{exercise}
\begin{exercise}\label{exercise:2.12}\index{SI0277@生成}
可以按下列方式构造随机直角三角形. 设 $X$ 是在 $(0,\pi/2)$ 上均匀分布的随机角. 对每个 $X$，如下图构造三角形，其中 $Y$ 为该随机三角形的高. 对固定常数 $d$，求 $Y$ 的分布和 $\mathrm EY$.
\sourcepageanchor{95}
\begin{center}\includegraphics[width=.7\textwidth]{assets/chapter2/exercise2-12.pdf}\end{center}
\end{exercise}
\begin{exercise}\label{exercise:2.13}
考虑一列相互独立的抛硬币试验，每次正面朝上的概率为 $p$. 定义随机变量 $X$ 为从第一次试验开始的游程（正面或反面连续出现）的长度. （例如，观察到 TTTH 或 HHHT 时，都有 $X=3$.）求 $X$ 的分布及 $\mathrm EX$.
\end{exercise}
\needspace{5\baselineskip}
\begin{exercise}\label{exercise:2.14}\index{SI0536@连续}\index{SI0537@离散}\leavevmode
\begin{enumerate}
\item[(a)]\phantomsection\label{itm:exercise:2.14:a} 设 $X$ 是连续的非负随机变量［当 $x<0$ 时，$f(x)=0$］. 证明
$$\mathrm EX=\int_0^\infty[1-F_X(x)]\,dx,$$
其中 $F_X(x)$ 为 $X$ 的累积分布函数.
\item[(b)] 设 $X$ 是取值范围为非负整数的离散随机变量. 证明
$$\mathrm EX=\sum_{k=0}^\infty(1-F_X(k)),$$
其中 $F_X(k)=P(X\le k)$. 将此与 \hyperref[itm:exercise:2.14:a]{(a)} 比较.
\end{enumerate}
\end{exercise}
\begin{exercise}\label{exercise:2.15}\index{SI0431@极大值与极小值的}
\hyperref[bib:source30]{Betteley (1977)} 给出了一个有趣的期望加法法则. 设 $X,Y$ 为任意两个随机变量，定义
$$X\wedge Y=\min(X,Y),\qquad X\vee Y=\max(X,Y).$$
仿照概率法则 $P(A\cup B)=P(A)+P(B)-P(A\cap B)$，证明
$$\mathrm E(X\vee Y)=\mathrm EX+\mathrm EY-\mathrm E(X\wedge Y).$$
（提示：证明 $X+Y=(X\vee Y)+(X\wedge Y)$.）
\end{exercise}
\begin{exercise}\label{exercise:2.16}
用\autoref{exercise:2.14}的结果求某类电话通话的平均时长. 假设一次通话的时长 $T$ 可用下列概率关系描述：$P(T>t)=ae^{-\lambda t}+(1-a)e^{-\mu t}$，其中 $a,\lambda,\mu$ 为常数，$0<a<1$，$\lambda>0$，$\mu>0$.
\end{exercise}
\begin{exercise}\label{exercise:2.17}\index{SI0684@中位数}
分布的中位数（median）是满足 $P(X\le m)\ge\frac12$ 且 $P(X\ge m)\ge\frac12$ 的数 $m$. （若 $X$ 连续，则 $m$ 满足 $\int_{-\infty}^m f(x)\,dx=\int_m^\infty f(x)\,dx=\frac12$.）求下列分布的中位数.
\begin{enumerate}
\item[(a)] $f(x)=3x^2$，$0<x<1$.
\item[(b)] $f(x)=1/[\pi(1+x^2)]$，$-\infty<x<\infty$.
\end{enumerate}
\end{exercise}
\begin{exercise}\label{exercise:2.18}\index{SI0685@最小化性质}
证明，如果 $X$ 是连续随机变量，则
$$\min_a\mathrm E|X-a|=\mathrm E|X-m|,$$
其中 $m$ 为 $X$ 的中位数（见\autoref{exercise:2.17}）.
\end{exercise}
\sourcepageanchor{96}
\begin{exercise}\label{exercise:2.19}
通过对积分求导，证明
$$\frac d{da}\mathrm E(X-a)^2=0\ \Longleftrightarrow\ \mathrm EX=a.$$
用微积分验证 $a=\mathrm EX$ 确实给出最小值. 列出关于 $F_X$ 和 $f_X$ 所需的假设.
\end{exercise}
\begin{exercise}\label{exercise:2.20}
一对夫妇决定不断生育，直到生出一个女儿为止. 这对夫妇的子女数的期望是多少？（提示：见\autoref{ex:1.5.4}.）
\end{exercise}
\begin{exercise}\label{exercise:2.21}
证明期望的“双向”法则\eqref{eq:2.2.5}，即 $\mathrm Eg(X)=\mathrm EY$，其中 $Y=g(X)$. 假定 $g(x)$ 为单调函数.
\end{exercise}
\begin{exercise}\label{exercise:2.22}
设 $X$ 的概率密度函数为
$$f(x)=\frac4{\beta^3\sqrt\pi}x^2e^{-x^2/\beta^2},\qquad0<x<\infty,\quad\beta>0.$$
\begin{enumerate}
\item[(a)] 验证 $f(x)$ 是概率密度函数.
\item[(b)] 求 $\mathrm EX$ 和 $\operatorname{Var}X$.
\end{enumerate}
\end{exercise}
\begin{exercise}\label{exercise:2.23}
设 $X$ 的概率密度函数为
$$f(x)=\frac12(1+x),\qquad-1<x<1.$$
\begin{enumerate}
\item[(a)] 求 $Y=X^2$ 的概率密度函数.
\item[(b)] 求 $\mathrm EY$ 和 $\operatorname{Var}Y$.
\end{enumerate}
\end{exercise}
\begin{exercise}\label{exercise:2.24}
对下列每个概率分布，计算 $\mathrm EX$ 和 $\operatorname{Var}X$.
\begin{enumerate}
\item[(a)] $f_X(x)=ax^{a-1}$，$0<x<1$，$a>0$.
\item[(b)] $f_X(x)=1/n$，$x=1,2,\ldots,n$，$n>0$ 为整数.
\item[(c)] $f_X(x)=\frac32(x-1)^2$，$0<x<2$.
\end{enumerate}
\end{exercise}
\begin{exercise}\label{exercise:2.25}\index{SI0832@偶}
设随机变量 $X$ 的概率密度函数 $f_X(x)$ 为偶函数. （若对每个 $x$，都有 $f_X(x)=f_X(-x)$，则 $f_X(x)$ 称为偶函数.）证明
\begin{enumerate}
\item[(a)] $X$ 与 $-X$ 同分布.
\item[(b)] $M_X(t)$ 关于零对称.
\end{enumerate}
\end{exercise}
\begin{exercise}\label{exercise:2.26}\index{SI0840@对称}
设 $f(x)$ 为概率密度函数，$a$ 为一个数，使得对所有 $\varepsilon>0$ 都有 $f(a+\varepsilon)=f(a-\varepsilon)$. 称这样的密度关于点 $a$ 对称（symmetric）.
\begin{enumerate}
\item[(a)] 举出三个对称密度的例子.
\item[(b)] 证明，如果 $X\sim f(x)$ 且该密度对称，则 $X$ 的中位数（见\autoref{exercise:2.17}）是 $a$.
\item[(c)] 证明，如果 $X\sim f(x)$ 且该密度对称，$\mathrm EX$ 存在，则 $\mathrm EX=a$.
\item[(d)]\phantomsection\label{itm:exercise:2.26:d} 证明，$f(x)=e^{-x}$，$x\ge0$，不是对称密度.
\item[(e)] 证明，对 \hyperref[itm:exercise:2.26:d]{(d)} 中的密度，中位数小于均值.
\end{enumerate}
\end{exercise}
\begin{exercise}\label{exercise:2.27}\index{SI0842@单峰}
设 $f(x)$ 为概率密度函数，$a$ 为一个数，使得当 $a\ge x\ge y$ 时，$f(a)\ge f(x)\ge f(y)$；当 $a\le x\le y$ 时，也有 $f(a)\ge f(x)\ge f(y)$. 称这样的密度为单峰（unimodal）密度，$a$ 为一个众数（mode）.
\begin{enumerate}
\item[(a)] 举出一个众数唯一的单峰密度.
\item[(b)] 举出一个众数不唯一的单峰密度.
\item[(c)] 证明，如果 $f(x)$ 既对称（见\autoref{exercise:2.26}）又单峰，则其对称中心是一个众数.
\item[(d)] 考虑密度 $f(x)=e^{-x}$，$x\ge0$. 证明该密度单峰. 它的众数是什么？
\end{enumerate}
\end{exercise}
\begin{exercise}\label{exercise:2.28}\index{SI0599@峰度}\index{SI0978@偏度}
令 $\mu_n$ 表示随机变量 $X$ 的 $n$ 阶中心矩. 除均值与方差外，另外两个值得关注的量是
$$\alpha_3=\frac{\mu_3}{(\mu_2)^{3/2}},\qquad\alpha_4=\frac{\mu_4}{\mu_2^2}.$$
$\alpha_3$ 称为偏度（skewness），$\alpha_4$ 称为峰度（kurtosis）. 偏度衡量密度的不对称性（见\autoref{exercise:2.26}）. 峰度虽然更难解释，但它衡量密度的尖峭或平坦程度.
\begin{enumerate}
\item[(a)] 证明，如果一个密度关于点 $a$ 对称，则 $\alpha_3=0$.
\item[(b)] 计算 $f(x)=e^{-x}$，$x\ge0$ 的 $\alpha_3$；这是一个右偏（skewed to the right）的密度.
\sourcepageanchor{97}
\item[(c)] 对下列每个密度计算 $\alpha_4$，并评论其尖峭程度.
$$\begin{aligned}
f(x)&=\frac1{\sqrt{2\pi}}e^{-x^2/2},&-\infty<x<\infty,\\
f(x)&=\frac12,&-1<x<1,\\
f(x)&=\frac12e^{-|x|},&-\infty<x<\infty.
\end{aligned}$$
\end{enumerate}
\hyperref[bib:source268]{Ruppert (1987)} 用影响函数（influence function，\autoref{sec:10.6.4}）进一步研究峰度的含义，\hyperref[bib:source152]{Groeneveld (1991)} 用它研究偏度；关于 $\alpha_4$ 的解释，还可参见 \hyperref[bib:source6]{Balanda and MacGillivray (1988)}.
\end{exercise}
\begin{exercise}\label{exercise:2.29}
计算离散分布的矩时，使用阶乘矩（factorial moments）往往更容易（见\hyperref[sec:2.6.2]{杂录~2.6.2}）.
\begin{enumerate}
\item[(a)]\phantomsection\label{itm:exercise:2.29:a} 计算二项分布和泊松分布的阶乘矩 $\mathrm E[X(X-1)]$.
\item[(b)] 用 \hyperref[itm:exercise:2.29:a]{(a)} 的结果计算二项分布和泊松分布的方差.
\item[(c)] 一个特别难处理的离散分布是贝塔二项分布（beta-binomial），其概率质量函数为
$$P(y=y)=a(y+a)\frac{\binom ny\binom{a+b-1}a}{\binom{n+a+b-1}{y+a}},$$
其中 $n,a,b$ 为整数，$y=0,1,2,\ldots,n$. 用阶乘矩计算贝塔二项分布的方差. （这一计算的另一种方法见\autoref{exercise:4.34}.）\footnote{译注：原文左端写作 $P(y=y)$，随机变量通常应大写为 $Y$. 原式前因子 $a(y+a)$ 也有疑误：例如 $n=a=b=1$ 时，两项质量为 $1/2$ 与 $2$，和不为 $1$. 将该因子改为 $a/(y+a)$，才得到通常的贝塔二项质量函数；此外通常要求 $a,b>0$. 此处保留原式与原假设.}
\end{enumerate}
\end{exercise}
\begin{exercise}\label{exercise:2.30}
求下列分布的矩母函数.
\begin{enumerate}
\item[(a)] $f(x)=1/c$，$0<x<c$.
\item[(b)] $f(x)=2x/c^2$，$0<x<c$.
\item[(c)] $f(x)=\frac1{2\beta}e^{-|x-\alpha|/\beta}$，$-\infty<x<\infty$，$-\infty<\alpha<\infty$，$\beta>0$.
\item[(d)] $P(X=x)=\binom{r+x-1}x p^r(1-p)^x$，$x=0,1,\ldots$，$0<p<1$，$r>0$ 为整数.
\end{enumerate}
\end{exercise}
\begin{exercise}\label{exercise:2.31}
是否存在一个分布，其矩母函数为 $M_X(t)=t/(1-t)$，$|t|<1$？如果存在，求出它；如果不存在，证明之.
\end{exercise}
\begin{exercise}\label{exercise:2.32}
设 $M_X(t)$ 为 $X$ 的矩母函数，定义 $S(t)=\log(M_X(t))$. 证明
$$\left.\frac d{dt}S(t)\right|_{t=0}=\mathrm EX,\qquad
\left.\frac{d^2}{dt^2}S(t)\right|_{t=0}=\operatorname{Var}X.$$
\end{exercise}
\begin{exercise}\label{exercise:2.33}
对下列各情形，验证给出的矩母函数表达式，并利用它计算 $\mathrm EX$ 和 $\operatorname{Var}X$.
\begin{enumerate}
\item[(a)] $P(X=x)=e^{-\lambda}\lambda^x/x!$，$M_X(t)=e^{\lambda(e^t-1)}$，$x=0,1,\ldots$；$\lambda>0$.
\item[(b)] $P(X=x)=p(1-p)^x$，$M_X(t)=p/[1-(1-p)e^t]$，$x=0,1,\ldots$；$0<p<1$.
\item[(c)] $f_X(x)=e^{-(x-\mu)^2/(2\sigma^2)}/(\sqrt{2\pi}\sigma)$，$M_X(t)=e^{\mu t+\sigma^2t^2/2}$，$-\infty<x<\infty$；$-\infty<\mu<\infty$，$\sigma>0$.
\end{enumerate}
\end{exercise}
\begin{exercise}\label{exercise:2.34}
有限个矩不能唯一确定一个分布，\hyperref[bib:source258]{Romano and Siegel (1986)} 中的这个例子说明了这一点. 设 $X$ 具有正态分布，即密度为
$$f_X(x)=\frac1{\sqrt{2\pi}}e^{-x^2/2},\qquad-\infty<x<\infty.$$
定义离散随机变量 $Y$：
$$P(Y=\sqrt3)=P(Y=-\sqrt3)=\frac16,\qquad P(Y=0)=\frac23.$$
证明
$$\mathrm EX^r=\mathrm EY^r,\qquad r=1,2,3,4,5.$$
（\hyperref[bib:source258]{Romano and Siegel (1986)} 指出，对任意有限的 $n$，都存在一个离散的、从而非正态的随机变量，其前 $n$ 阶矩与 $X$ 的对应矩相同.）
\end{exercise}
\sourcepageanchor{98}
\begin{exercise}\label{exercise:2.35}\index{SI0698@非唯一性}
补全\autoref{ex:2.3.10}中的细节.
\begin{enumerate}
\item[(a)] 证明，若 $X_1\sim f_1(x)$，则
$$\mathrm EX_1^r=e^{r^2/2},\qquad r=0,1,\ldots.$$
所以 $f_1(x)$ 的所有阶矩都存在，并且全都有限.
\item[(b)] 证明，对所有正整数 $r$，
$$\int_0^\infty x^rf_1(x)\sin(2\pi\log x)\,dx=0,$$
因而对所有 $r$，$\mathrm EX_1^r=\mathrm EX_2^r$. （\hyperref[bib:source258]{Romano and Siegel (1986)} 讨论了这个例子的极端版本，即一整类不同的密度具有相同的各阶矩. 此外，\hyperref[bib:source16]{Berg (1988)} 表明，正态分布的较简单变换，例如 $X^3$，也会出现这种矩的现象.）
\end{enumerate}
\end{exercise}
\begin{exercise}\label{exercise:2.36}\index{SI0333@矩母函数}
\autoref{ex:2.3.10}所依据的对数正态分布具有一个有趣性质. 若密度为
$$f(x)=\frac1{\sqrt{2\pi}x}e^{-(\log x)^2/2},\qquad0\le x<\infty,$$
则\autoref{exercise:2.35}表明所有阶矩都存在且有限. 然而，这个分布没有矩母函数，即
$$M_X(t)=\int_0^\infty\frac{e^{tx}}{\sqrt{2\pi}x}e^{-(\log x)^2/2}\,dx$$
不存在. 证明之.
\end{exercise}
\begin{exercise}\label{exercise:2.37}
参照\hyperref[sec:2.6.3]{杂录~2.6.3}所述的情形：
\begin{enumerate}
\item[(a)] 画出密度 $f_1$ 和 $f_2$，说明它们的差异.
\item[(b)] 画出累积量母函数 $K_1$ 和 $K_2$，说明它们的相似性.
\item[(c)] 计算密度 $f_1$ 和 $f_2$ 的矩母函数. 它们相似还是不同？
\item[(d)] 密度 $f_1$ 和 $f_2$ 与\autoref{ex:2.3.10}中的密度有什么关系？
\end{enumerate}
\end{exercise}
\begin{exercise}\label{exercise:2.38}\index{SI0283@与负二项的关系}\index{SI0354@与卡方的关系}
设 $X$ 具有负二项分布（negative binomial distribution），概率质量函数为
$$f(x)=\binom{r+x-1}x p^r(1-p)^x,\qquad x=0,1,2,\ldots,$$
其中 $0<p<1$，$r>0$ 为整数.
\begin{enumerate}
\item[(a)] 计算 $X$ 的矩母函数.
\item[(b)] 定义新随机变量 $Y=2pX$. 证明，当 $p\downarrow0$ 时，$Y$ 的矩母函数收敛到自由度为 $2r$ 的卡方随机变量的矩母函数，即证明
$$\lim_{p\to0}M_Y(t)=\left(\frac1{1-2t}\right)^r,\qquad |t|<\frac12.$$
\end{enumerate}
\end{exercise}
\begin{exercise}\label{exercise:2.39}
对下列各情形计算所示导数，并说明所有运算的合理性.
\begin{enumerate}
\item[(a)] $\displaystyle\frac d{dx}\int_0^x e^{-\lambda t}\,dt$.
\item[(b)] $\displaystyle\frac d{d\lambda}\int_0^\infty e^{-\lambda t}\,dt$.
\item[(c)] $\displaystyle\frac d{dt}\int_t^1\frac1{x^2}\,dx$.
\item[(d)] $\displaystyle\frac d{dt}\int_1^\infty\frac1{(x-t)^2}\,dx$.
\end{enumerate}
\end{exercise}
\begin{exercise}\label{exercise:2.40}\index{SI0257@与二项的关系}\index{SI0268@与贝塔的关系}
证明
$$\sum_{k=0}^x\binom nk p^k(1-p)^{n-k}
=(n-x)\binom nx\int_0^{1-p}t^{n-x-1}(1-t)^x\,dt.$$
（提示：分部积分，或两边对 $p$ 求导.）
\end{exercise}



\sourcepageanchor{99}
\subsection{杂录}\label{sec:2.6}
\subsubsection{矩序列的唯一性}\label{sec:2.6.1}\index{SI0699@唯一性}
分布不一定由它的矩确定. 但如果 $\sum_{r=1}^\infty\mu^r r^k/k!$ 具有正的收敛半径，其中 $X\sim F_X$，$\mathrm EX^r=\mu'_r$，那么矩序列具有唯一性，从而分布被唯一确定（\hyperref[bib:source33]{Billingsley 1995}，第~30~节）.\footnote{译注：原文级数印为 $\sum_{r=1}^\infty\mu^r r^k/k!$，求和指标、幂指数与后面的矩记号不一致，也未给出作为幂级数变量的量. 从上下文看，可能意指矩母函数的幂级数 $\sum_{r=0}^\infty\mu'_r t^r/r!$；此处保留原式，具体排印修正待核实.} 这个级数的收敛还意味着矩母函数在某个区间内存在，因此矩母函数确定分布.

矩序列具有唯一性的一个充分条件是 Carleman 条件\index{SI0101@Carleman条件}\semanticindexanchor{101:99}（Carleman's Condition；\hyperref[bib:source78]{Chung 1974）}. 若 $X\sim F_X$，并记 $\mathrm EX^r=\mu'_r$，则当
$$\sum_{r=1}^\infty\frac1{(\mu'_{2r})^{1/(2r)}}=+\infty$$
时，矩序列具有唯一性. 这个条件一般不容易验证.

\hyperref[bib:source119]{Feller (1971)} 对拉普拉斯变换作了非常完整的论述，矩母函数是其特殊情形. 特别地，Feller 证明了（与 Billingsley 类似），只要
$$M_X(t)=\sum_{r=0}^\infty\frac{(-1)^r\mu'_r t^r}{r!}$$
在区间 $-t_0\le t<t_0$（$t_0>0$）上收敛，分布 $F_X$ 就被唯一确定.\footnote{译注：本书定义 $M_X(t)=\mathrm Ee^{tX}$，其泰勒系数应为 $\mu'_r/r!$；原式含 $(-1)^r$，对应的是 $\mathrm Ee^{-tX}=M_X(-t)$ 的展开. 这不影响所述唯一性思路，但记号与前文定义不一致.} 因此，当矩母函数存在时，矩序列唯一确定分布 $F_X$.

应当清楚，用矩母函数确定分布是件困难的工作. 更好的办法是使用下面介绍的特征函数. 虽然特征函数简化了分布的刻画，但需要理解复分析. 有得必有失.

\subsubsection{其他母函数}\label{sec:2.6.2}
除了矩母函数，还有其他多种母函数可用. 在多数情况下，其中最有用的是特征函数. 除了少数情形，其他母函数的用途较少，但有时它们可以简化计算.

\paragraph{累积量母函数.}\index{SI0454@累积量}\semanticindexanchor{454:99}
对随机变量 $X$，累积量母函数（cumulant generating function）是函数 $\log[M_X(t)]$. 它可用来生成 $X$ 的累积量（cumulants），累积量被定义为（颇为迂回地）累积量母函数的泰勒级数中的系数（见\autoref{exercise:2.32}）.\footnote{译注：通常第 $r$ 阶累积量定义为 $\log M_X(t)$ 在 $0$ 处的 $r$ 阶导数；它是以 $t^r/r!$ 为基底时的系数，而不是普通 $t^r$ 项的系数. 原文“泰勒级数中的系数”省略了这一阶乘约定.}

\paragraph{阶乘矩母函数.}\index{SI0455@阶乘矩}\semanticindexanchor{455:99}
若期望存在，$X$ 的阶乘矩母函数（factorial moment generating function）定义为 $\mathrm Et^X$. 这个名称来自于它满足
$$\left.\frac{d^r}{dt^r}\mathrm Et^X\right|_{t=1}
=\mathrm E\{X(X-1)\cdots(X-r+1)\},$$
其中右端是一个阶乘矩. 若 $X$ 为离散随机变量，则可以写成
$$\mathrm Et^X=\sum_x t^xP(X=x),$$
阶乘矩母函数又称为概率母函数（probability generating function）\index{SI0456@概率}\semanticindexanchor{456:99}，因为幂级数的系数给出了概率. 也就是说，为求 $X=k$ 的概率，计算\footnote{译注：原文最后一式在 $t=1$ 处求值，应为在 $t=0$ 处求值才能抽取 $t^k$ 的系数. 这里的普通概率母函数通常还要求 $X$ 取非负整数值；任意离散取值不一定形成关于 $t$ 的普通幂级数.}
$$\left.\frac1{k!}\frac{d^k}{dt^k}\mathrm Et^X\right|_{t=1}=P(X=k).$$


\sourcepageanchor{100}
\paragraph{特征函数.}\index{SI0103@特征函数}\semanticindexanchor{103:100}
这些函数中最有用的也许是特征函数（characteristic function）. $X$ 的特征函数定义为
$$\phi_X(t)=\mathrm Ee^{itX},$$
其中 $i$ 是复数 $\sqrt{-1}$，所以上面的期望需要复数积分.

特征函数的作用远超过矩母函数. 当 $F_X$ 的矩存在时，$\phi_X$ 可以像矩母函数一样生成这些矩. 特征函数总是存在，并且完全确定分布. 也就是说，每个累积分布函数都有唯一的特征函数. 因而，例如我们可以陈述一个类似于\autoref{thm:2.3.11}的定理，却不必附加限制条件.

\begin{theorem}[特征函数的收敛]\label{thm:2.6.1}\index{SI0104@收敛}
设 $X_k$，$k=1,2,\ldots$，是随机变量序列，每个随机变量的特征函数为 $\phi_{X_k}(t)$. 再设
$$\lim_{k\to\infty}\phi_{X_k}(t)=\phi_X(t),\qquad\text{对 $0$ 的某个邻域内的所有 $t$},$$
且 $\phi_X(t)$ 是一个特征函数.\footnote{译注：原文把条件写成仅在 $0$ 的某个邻域内收敛；一般特征函数不一定解析，仅有局部相同不足以确定分布，因此该条件一般不够. 通常的连续性定理要求对所有实数 $t$ 收敛，并要求极限函数在 $0$ 处连续（此处若已假设极限为特征函数，连续性自动满足）. 原文结论中的大写 $X$ 也应为自变量 $x$，这里一并保留原述.} 则对 $F_X(x)$ 连续的所有 $X$，
$$\lim_{k\to\infty}F_{X_k}(x)=F_X(x).$$
\end{theorem}


\hyperref[bib:source118]{Feller (1968)} 对母函数有完整论述. 几乎所有高等概率论教材都介绍特征函数，例如 \hyperref[bib:source33]{Billingsley (1995)} 或 \hyperref[bib:source251]{Resnick (1999)}.

\subsubsection{矩母函数能否刻画分布？}\label{sec:2.6.3}\index{SI0689@矩母函数}
在以上述问题为题的文章中，\hyperref[bib:source215]{McCullagh (1994)} 考察了一对与\autoref{ex:2.3.10}类似、但具有矩母函数的密度，
$$f_1=\mathrm n(0,1),\qquad f_2=f_1(x)\left[1+\frac12\sin(2\pi x)\right],$$
其累积量母函数为
$$K_1(t)=t^2/2,\qquad K_2(t)=K_1(t)+\log\left[1+\frac12e^{-2\pi^2}\sin(2\pi t)\right].$$
他指出，虽然两个密度明显不同，它们的累积量母函数实际上几乎相同，整个取值范围内的最大差小于 $1.34\times10^{-9}$（小于一个像素的大小）. 因此，对标题中问题的回答是：“从数学目的看，是；但从数值目的看，是响亮的否.” 相比之下，\hyperref[bib:source325]{Waller (1995)} 说明，尽管矩母函数不能在数值上区分这两个分布，特征函数却做得很好. （\hyperref[bib:source326]{Waller et al. (1995)} 和 \hyperref[bib:source211]{Luce\~no (1997)} 进一步研究了用特征函数数值求取累积分布函数的作用.）细节见\autoref{exercise:2.37}.





\section{常见分布族}\label{sec:3}
\chapterdoi{3}

\sourcepageanchor{101}
\begin{flushright}\itshape
华生，你怎么看这一切反常之处？\\
它们合在一起的影响当然很大，可是每一件本身又都完全可能发生.\\
——Sherlock Holmes 与 Watson 医生，《格兰其庄园》
\end{flushright}



\subsection{引言}\label{sec:3.1}
统计分布用来为总体建立模型；因此，我们通常处理的是一个分布族（family of distributions），而不是单个分布. 这个族由一个或多个参数标定，使我们可以在保持函数形式不变的同时，改变分布的某些特征. 例如，我们可能认为正态分布适合为某个总体建模，但不能准确指定其均值. 此时处理的是一个参数族：均值为 $\mu$ 的正态分布，其中 $\mu$ 是尚未指定的参数，$-\infty<\mu<\infty$.

本章罗列许多较常见的统计分布，其中一些已经遇到过. 对每个分布，我们将给出均值、方差以及许多有助于理解的其他有用或描述性的度量. 我们还会介绍这些分布的一些典型应用，以及它们之间某些有趣而有用的联系. 部分事实汇总于书末的表格中. 本章对统计分布的介绍远非全面；这一任务已由 \hyperref[bib:source177]{Johnson and Kotz (1969–1972)} 的多卷本著作 \textit{Distributions in Statistics}，以及更新版 Johnson, Kotz, and Balakrishnan (\hyperref[bib:source178]{1994}, \hyperref[bib:source179]{1995}) 和 \hyperref[bib:source180]{Johnson, Kotz, and Kemp (1992)} 完成.



\subsection{离散分布}\label{sec:3.2}
如果随机变量 $X$ 的取值范围，也就是样本空间，是可数的，就称 $X$ 具有离散分布（discrete distribution）. 在大多数情形中，随机变量的结果取整数值.

\paragraph{离散均匀分布.}
如果随机变量 $X$ 满足
\begin{equation}\label{eq:3.2.1}\index{SI0287@离散均匀}
P(X=x\mid N)=\frac1N,\qquad x=1,2,\ldots,N,
\end{equation}
其中 $N$ 为给定整数，就称 $X$ 具有离散均匀分布（discrete uniform distribution）$\operatorname{uniform}(1,N)$. 这个分布把相等的概率质量赋给每个结果 $1,2,\ldots,N$.

关于记号的说明：我们几乎总是处理参数分布，而这时分布依赖参数的值. 为强调这一点并明确参数，我们把参数写在概率质量函数中，并在前面加上“$\mid$”（给定）. 这一约定也用于累积分布函数、概率密度函数、期望以及其他需要明确参数的地方. 在不会产生混淆时，可以略去参数，以免记号过于繁杂.

\sourcepageanchor{102}
为计算 $X$ 的均值和方差，回顾下列恒等式（可用归纳法证明）：
$$\sum_{i=1}^k i=\frac{k(k+1)}2,\qquad
\sum_{i=1}^k i^2=\frac{k(k+1)(2k+1)}6.$$
于是
$$\mathrm EX=\sum_{x=1}^N xP(X=x\mid N)=\sum_{x=1}^N x\frac1N=\frac{N+1}2,$$
并且
$$\mathrm EX^2=\sum_{x=1}^N x^2\frac1N=\frac{(N+1)(2N+1)}6,$$
所以
$$\begin{aligned}
\operatorname{Var}X&=\mathrm EX^2-(\mathrm EX)^2\\
&=\frac{(N+1)(2N+1)}6-\left(\frac{N+1}2\right)^2\\
&=\frac{(N+1)(N-1)}{12}.
\end{aligned}$$
这个分布可以推广到样本空间为任意连续整数范围 $N_0,N_0+1,\ldots,N_1$ 的情形，其概率质量函数为 $P(X=x\mid N_0,N_1)=1/(N_1-N_0+1)$.

\paragraph{超几何分布.}
超几何分布（hypergeometric distribution）在有限总体抽样中有许多应用，借助经典的瓮模型例子最容易理解它.

设一个大瓮内有 $N$ 个球，除其中 $M$ 个为红色、$N-M$ 个为绿色外，各方面都完全相同. 我们蒙上眼睛，伸手随机取出 $K$ 个球（一次取出全部 $K$ 个球，属于无放回抽样）. 恰有 $x$ 个红球的概率是多少？

如\autoref{sec:1.2.3}所讨论的，从 $N$ 个球中抽取大小为 $K$ 的样本，共有 $\binom NK$ 种. 要求其中 $x$ 个为红球，有 $\binom Mx$ 种选法，余下用 $K-x$ 个绿球补足样本，有 $\binom{N-M}{K-x}$ 种选法. 因此，若用 $X$ 表示大小为 $K$ 的样本中的红球数，$X$ 就具有下列超几何分布：
\begin{equation}\label{eq:3.2.2}\index{SI0323@超几何}
P(X=x\mid N,M,K)=\frac{\binom Mx\binom{N-M}{K-x}}{\binom NK},\qquad x=0,1,\ldots,K.
\end{equation}
注意，\eqref{eq:3.2.2}中还暗含了关于 $X$ 取值范围的附加假设. 形如 $\binom nr$ 的二项系数只在 $n\ge r$ 时作了定义，因此 $X$ 的取值范围还受到下列一对不等式的限制：
$$M\ge x,\qquad N-M\ge K-x,$$
合并起来就是
$$M-(N-K)\le x\le M.$$
许多情况下，$K$ 相比于 $M,N$ 很小，所以 $0\le x\le K$ 会包含在上述范围内，因而是适用的. 超几何概率函数的公式通常相当难处理. 实际上，甚至验证
$$\sum_{x=0}^K P(X=x)=\sum_{x=0}^K\frac{\binom Mx\binom{N-M}{K-x}}{\binom NK}=1$$
也并非易事.
\sourcepageanchor{103}
超几何分布说明，从统计学角度看，处理有限总体（有限的 $N$）是一项困难的任务.

超几何分布的均值为
$$\mathrm EX=\sum_{x=0}^K x\frac{\binom Mx\binom{N-M}{K-x}}{\binom NK}
=\sum_{x=1}^K x\frac{\binom Mx\binom{N-M}{K-x}}{\binom NK},\qquad\text{（$x=0$ 的项为零）}.$$
为计算这个表达式，使用已经在\autoref{sec:2.3}中遇到的恒等式\index{SI0531@二项式系数}\semanticindexanchor{531:103}
$$x\binom Mx=M\binom{M-1}{x-1},\qquad
\binom NK=\frac NK\binom{N-1}{K-1},$$
得到
$$\mathrm EX=\sum_{x=1}^K\frac{M\binom{M-1}{x-1}\binom{N-M}{K-x}}{\frac NK\binom{N-1}{K-1}}
=\frac{KM}N\sum_{x=1}^K\frac{\binom{M-1}{x-1}\binom{N-M}{K-x}}{\binom{N-1}{K-1}}.$$
现在可以认出，上面第二个和是另一个超几何分布的概率之和，其参数为 $N-1,M-1,K-1$. 令 $y=x-1$，就能清楚地看出这一点：
$$\begin{aligned}
\sum_{x=1}^K\frac{\binom{M-1}{x-1}\binom{N-M}{K-x}}{\binom{N-1}{K-1}}
&=\sum_{y=0}^{K-1}\frac{\binom{M-1}y\binom{(N-1)-(M-1)}{K-1-y}}{\binom{N-1}{K-1}}\\
&=\sum_{y=0}^{K-1}P(Y=y\mid N-1,M-1,K-1)=1,
\end{aligned}$$
其中 $Y$ 是参数为 $N-1,M-1,K-1$ 的超几何随机变量. 因此，对于超几何分布，
$$\mathrm EX=\frac{KM}N.$$
类似但更长的计算可以得到
$$\operatorname{Var}X=\frac{KM}N\left(\frac{(N-M)(N-K)}{N(N-1)}\right).$$
注意这里计算 $\mathrm EX$ 时所用的变换：把求和变成参数不同的另一个超几何分布，认出这一事实后，就能够求出级数之和.

\begin{example}[验收抽样]\label{ex:3.2.1}\index{SI0003@超几何}
本例说明超几何分布在验收抽样（acceptance sampling）中的应用. 设某零售商按批购买货物，每件产品或为合格品，或为次品. 令
$$\begin{aligned}
N&=\text{一批中的产品数},\\
M&=\text{一批中的次品数}.
\end{aligned}$$
便可计算大小为 $K$ 的样本含有 $x$ 件次品的概率. 具体地，假设运来一批 $25$ 个机械零件，只有通过公差检验的零件才被视为合格.
\sourcepageanchor{104}
我们抽取 $10$ 个零件，发现没有次品（全部在公差范围内）. 如果这批 $25$ 个零件中有 $6$ 个次品，这一事件的概率是多少？应用参数为 $N=25,M=6,K=10$ 的超几何分布，得到
$$P(X=0)=\frac{\binom60\binom{19}{10}}{\binom{25}{10}}=.028,$$
这表明，如果该批货物中有 $6$ 个（或更多！）次品，我们观察到的事件就相当不可能发生.
\end{example}

\paragraph{二项分布.}\index{SI0261@二项}\semanticindexanchor{261:104}
二项分布是一种十分有用的离散分布，其基础是 Bernoulli 试验（Bernoulli trial）的思想. Bernoulli 试验以概率论奠基人之一 James Bernoulli 命名，是恰有两个可能结果的试验. 如果随机变量 $X$ 满足\index{SI0088@伯努利试验（见分布：伯努利）}\semanticindexanchor{88:104}
\begin{equation}\label{eq:3.2.3}\index{SI0252@伯努利}
X=\begin{cases}1,&\text{概率为 }p,\\0,&\text{概率为 }1-p,\end{cases}
\qquad0\le p\le1,
\end{equation}
就称它具有 $\operatorname{Bernoulli}(p)$ 分布. 通常把 $X=1$ 称为“成功”，$p$ 称为成功概率；把 $X=0$ 称为“失败”. 容易看出，$\operatorname{Bernoulli}(p)$ 随机变量的均值与方差为
$$\begin{aligned}
\mathrm EX&=1p+0(1-p)=p,\\
\operatorname{Var}X&=(1-p)^2p+(0-p)^2(1-p)=p(1-p).
\end{aligned}$$
许多试验都可以建模为一列 Bernoulli 试验，最简单的是反复抛硬币：$p$ 为正面朝上的概率，正面朝上时 $X=1$. 其他例子包括赌博游戏（例如轮盘中，出现红色时令 $X=1$，所以 $p$ 为红色出现的概率）、选举民意调查（候选人 A 获得一票时 $X=1$）以及疾病发病情况（$p$ 为随机选取的一人受到感染的概率）.

若进行 $n$ 次相同的 Bernoulli 试验，定义事件
$$A_i=\{\text{第 $i$ 次试验中 $X=1$}\},\qquad i=1,2,\ldots,n.$$
若假定 $A_1,\ldots,A_n$ 为一组相互独立的事件（抛硬币便是如此），就容易推导出 $n$ 次试验中成功总次数的分布. 定义随机变量
$$Y=\text{$n$ 次试验中的成功总次数}.$$
事件 $\{Y=y\}$ 发生，当且仅当 $A_1,\ldots,A_n$ 中恰有 $y$ 个发生，而其余 $n-y$ 个不发生. $n$ 次 Bernoulli 试验的某个特定结果（发生与不发生的一种特定排列）可能是 $A_1\cap A_2\cap A_3^c\cap\cdots\cap A_{n-1}\cap A_n^c$. 它发生的概率为
$$\begin{aligned}
P(A_1\cap A_2\cap A_3^c\cap\cdots\cap A_{n-1}\cap A_n^c)
&=pp(1-p)\cdots p(1-p)\\
&=p^y(1-p)^{n-y},
\end{aligned}$$
这里用到了 $A_i$ 的独立性. 注意，这个计算并不依赖于发生的是哪 $y$ 个 $A_i$，只依赖于有某一组 $y$ 个事件发生. 而且，无论发生的是哪一组 $y$ 个 $A_i$，事件 $\{Y=y\}$ 都发生. 综上，
\sourcepageanchor{105}
$n$ 次试验中恰有 $y$ 次成功的某个特定序列，发生概率为 $p^y(1-p)^{n-y}$. 因为这样的序列有 $\binom ny$ 个（即 $y$ 个 $1$ 与 $n-y$ 个 $0$ 的排列数），所以
$$P(Y=y\mid n,p)=\binom ny p^y(1-p)^{n-y},\qquad y=0,1,2,\ldots,n,$$
而 $Y$ 称为 $\operatorname{binomial}(n,p)$ 随机变量.

也可以用另一种等价方式定义 $Y$：在 $n$ 次相同且相互独立、每次成功概率为 $p$ 的 Bernoulli 试验中，定义随机变量 $X_1,\ldots,X_n$ 为
$$X_i=\begin{cases}1,&\text{概率为 }p,\\0,&\text{概率为 }1-p.\end{cases}$$
随机变量
$$Y=\sum_{i=1}^nX_i$$
就具有 $\operatorname{binomial}(n,p)$ 分布.

$\sum_{y=0}^nP(Y=y)=1$ 这一事实来自下列一般定理.
\begin{theorem}[二项式定理]\label{thm:3.2.2}\index{SI0092@二项式定理}
对任意实数 $x,y$ 和整数 $n\ge0$，
\begin{equation}\label{eq:3.2.4}
(x+y)^n=\sum_{i=0}^n\binom ni x^iy^{n-i}.
\end{equation}
\end{theorem}
\begin{proof}
写成
$$(x+y)^n=(x+y)(x+y)\cdots(x+y),$$
并考虑右端如何计算. 从每个因子 $(x+y)$ 中选取 $x$ 或 $y$，再把 $n$ 次选取的结果相乘. 对每个 $i=0,1,\ldots,n$，其中 $x$ 恰好出现 $i$ 次的项数为 $\binom ni$. 因而该项的形式为 $\binom ni x^iy^{n-i}$，结果得证.
\end{proof}
在\eqref{eq:3.2.4}中取 $x=p$、$y=1-p$，得到
$$1=(p+(1-p))^n=\sum_{i=0}^n\binom ni p^i(1-p)^{n-i},$$
可见和中的每一项都是一个二项概率. 另一个特殊情形是在\autoref{thm:3.2.2}中取 $x=y=1$，得到恒等式
$$2^n=\sum_{i=0}^n\binom ni.$$
二项分布的均值与方差已经在\autoref{ex:2.2.3}和\autoref{ex:2.3.5}中推导过，这里不再重复推导. 为了完整起见，列出结果：若 $X\sim\operatorname{binomial}(n,p)$，则
$$\mathrm EX=np,\qquad\operatorname{Var}X=np(1-p).$$
二项分布的矩母函数已经在\autoref{ex:2.3.9}中算出，即
$$M_X(t)=[pe^t+(1-p)]^n.$$

\begin{example}[掷骰子的概率]\label{ex:3.2.3}\index{SI0686@德梅雷骰子问题}
假设我们要计算掷一枚公平骰子四次至少出现一次 $6$ 点的概率. 这个试验可以建模为一列
\sourcepageanchor{106}
四次 Bernoulli 试验，成功概率为 $p=\frac16=P(\text{骰子出现 $6$ 点})$. 定义随机变量
$$X=\text{四次投掷中出现 $6$ 点的总次数}.$$
则 $X\sim\operatorname{binomial}(4,\frac16)$，并且
$$\begin{aligned}
P(\text{至少一次 $6$ 点})&=P(X>0)=1-P(X=0)\\
&=1-\binom40\left(\frac16\right)^0\left(\frac56\right)^4\\
&=1-\left(\frac56\right)^4\\
&=.518.
\end{aligned}$$
现在考虑另一种游戏：掷一对骰子 $24$ 次，求至少一次同时出现两个 $6$ 点的概率. 同样可以用成功概率为 $p$ 的二项分布建模，其中
$$p=P(\text{掷出两个 $6$ 点})=\frac1{36}.$$
所以，若 $Y$ 为 $24$ 次投掷中出现双 $6$ 点的次数，则 $Y\sim\operatorname{binomial}(24,\frac1{36})$，且
$$\begin{aligned}
P(\text{至少一次双 $6$ 点})&=P(Y>0)\\
&=1-P(Y=0)\\
&=1-\binom{24}0\left(\frac1{36}\right)^0\left(\frac{35}{36}\right)^{24}\\
&=1-\left(\frac{35}{36}\right)^{24}\\
&=.491.
\end{aligned}$$
这项计算最初由 Pascal 在十八世纪应赌徒 de M\'er\'e 的请求完成，后者原以为两个事件的概率相同.\footnote{译注：原文写作“十八世纪”，但 Pascal 活跃于十七世纪；此处年代疑误，正文保留原文.}（他在第二种赌局中开始输钱以后，才开始相信自己错了.）
\end{example}

\paragraph{泊松分布.}
泊松分布是一种应用广泛的离散分布，可以为许多不同类型的试验建立模型. 例如，如果要对等待某个事件发生的现象建模（如等公共汽车，或在银行等待顾客到达），那么在给定时间区间内发生的次数，有时可用泊松分布建模. 泊松分布的一个基本假设是：对较短的时间区间，一次到达的概率与等待时间的长度成正比. 这使它成为前述情形的合理模型. 例如，等待越久，越可能有顾客进入银行，这一假设是有道理的. 更正式的处理见\autoref{sec:3.8}.

另一个应用领域是空间分布. 例如，可用泊松分布对一个区域中的炸弹落点分布，或湖中鱼的分布建模.

泊松分布只有一个参数 $\lambda$，有时称为强度参数（intensity parameter）. 如果取非负整数值的随机变量 $X$ 满足
\begin{equation}\label{eq:3.2.5}\index{SI0371@泊松}
P(X=x\mid\lambda)=\frac{e^{-\lambda}\lambda^x}{x!},\qquad x=0,1,\ldots,
\end{equation}
就称它具有 $\operatorname{Poisson}(\lambda)$ 分布.
\sourcepageanchor{107}

为说明$\sum_{x=0}^\infty P(X=x\mid\lambda)=1$，回顾$e^y$的Taylor级数展开
$$e^y=\sum_{i=0}^\infty\frac{y^i}{i!}.$$
因此，
$$\sum_{x=0}^\infty P(X=x\mid\lambda)=e^{-\lambda}\sum_{x=0}^\infty\frac{\lambda^x}{x!}=e^{-\lambda}e^\lambda=1.$$
容易看出，$X$的均值为
\begin{align*}
EX&=\sum_{x=0}^\infty x\frac{e^{-\lambda}\lambda^x}{x!}
=\sum_{x=1}^\infty x\frac{e^{-\lambda}\lambda^x}{x!}\\
&=\lambda e^{-\lambda}\sum_{x=1}^\infty\frac{\lambda^{x-1}}{(x-1)!}
=\lambda e^{-\lambda}\sum_{y=0}^\infty\frac{\lambda^y}{y!}\qquad(\text{令 }y=x-1)\\
&=\lambda.
\end{align*}
类似计算可得
$$\operatorname{Var}X=\lambda,$$
所以参数$\lambda$同时是Poisson分布的均值和方差.

同样利用$e^y$的Taylor级数，也可直接算出矩母函数：
$$M_X(t)=e^{\lambda(e^t-1)}.$$
（见\autoref{exercise:2.33}和\autoref{ex:2.3.13}.）

\begin{example}[等待时间]\label{ex:3.2.4}
作为等待事件发生的一个应用例子，考虑一位平均每三分钟处理五通电话的接线员. 下一分钟没有电话的概率是多少？至少有两通电话的概率是多少？

若令$X=$一分钟内的电话数，则$X$具有Poisson分布，且$EX=\lambda=5/3$. 因而
\begin{align*}
P(\text{下一分钟没有电话})&=P(X=0)
=\frac{e^{-5/3}(5/3)^0}{0!}\\
&=e^{-5/3}=.189;\\
P(\text{下一分钟至少两通电话})&=P(X\ge2)
=1-P(X=0)-P(X=1)\\
&=1-.189-\frac{e^{-5/3}(5/3)^1}{1!}=.496.
\end{align*}
\end{example}
\sourcepageanchor{108}

利用以下递推关系，可以迅速计算Poisson概率：
\begin{equation}\label{eq:3.2.6}\index{SI0375@递推关系}
P(X=x)=\frac{\lambda}{x}P(X=x-1),\qquad x=1,2,\ldots.
\end{equation}
把Poisson分布的pmf写出来，就容易证明这个关系. 其他离散分布也有类似关系. 例如，若$Y\sim\operatorname{binomial}(n,p)$，则
\begin{equation}\label{eq:3.2.7}\index{SI0267@递推关系}
P(Y=y)=\frac{n-y+1}{y}\frac{p}{1-p}P(Y=y-1).
\end{equation}
递推关系\eqref{eq:3.2.6}和\eqref{eq:3.2.7}可用来建立二项分布的Poisson近似. 我们已在\autoref{sec:2.3}见过这种近似，当时使用矩母函数加以说明. 令$\lambda=np$，当$p$很小时，可写成
$$\frac{n-y+1}{y}\frac{p}{1-p}=\frac{np-p(y-1)}{y-py}\approx\frac{\lambda}{y},$$
因为$p$很小时，$p(y-1)$与$py$两项可以忽略. 因此，在这一近似精度下，\eqref{eq:3.2.7}成为
\begin{equation}\label{eq:3.2.8}\index{SI0060@以泊松近似二项}
P(Y=y)=\frac{\lambda}{y}P(Y=y-1),
\end{equation}
即Poisson递推关系. 因而只须再证明$P(X=0)\approx P(Y=0)$就完成了这个近似，因为其他概率都由\eqref{eq:3.2.8}得到. 现在令$np=\lambda$，则
$$P(Y=0)=(1-p)^n=\left(1-\frac{np}{n}\right)^n=\left(1-\frac{\lambda}{n}\right)^n.$$
回顾\autoref{sec:2.3}，固定$\lambda$时，$\lim_{n\to\infty}(1-(\lambda/n))^n=e^{-\lambda}$，故$n$很大时有近似
$$P(Y=0)=\left(1-\frac{\lambda}{n}\right)^n\approx e^{-\lambda}=P(X=0),$$
从而完成了二项分布的Poisson近似.

这一近似在$n$很大而$p$很小时有效，此时它也恰好最有用，使我们免于计算$n$很大时的二项系数及幂.

\begin{example}[Poisson近似]\label{ex:3.2.5}
一位排字工平均每排500个词出现一个错误. 一般每页有300个词. 五页中不超过两个错误的概率是多少？

若假定排一个词是成功概率$p=1/500$的Bernoulli试验（注意我们把出错称为“成功”），且各次试验独立，则$X=$五页（1500个词）中的错误数服从$\operatorname{binomial}(1500,1/500)$分布. 因此，
\begin{align*}
P(\text{不超过两个错误})&=P(X\le2)\\
&=\sum_{x=0}^2\binom{1500}{x}\left(\frac1{500}\right)^x\left(\frac{499}{500}\right)^{1500-x}\\
&=.4230,
\end{align*}
这个计算相当繁琐. 若使用$\lambda=1500(1/500)=3$的Poisson近似，则
$$P(X\le2)\approx e^{-3}\left(1+3+\frac{3^2}{2}\right)=.4232.$$
\end{example}
\sourcepageanchor{109}

\paragraph{负二项分布.}
二项分布计算固定次数Bernoulli试验中的成功次数. 假设我们转而计算获得固定次数成功所需的Bernoulli试验次数. 后一种表述就导出了负二项分布（negative binomial distribution）.

在一列独立的Bernoulli$(p)$试验中，令随机变量$X$表示第$r$次成功发生时的试验序号，其中$r$为固定整数. 则
\begin{equation}\label{eq:3.2.9}
P(X=x\mid r,p)=\binom{x-1}{r-1}p^r(1-p)^{x-r},\qquad x=r,r+1,\ldots,
\end{equation}
称$X$具有负二项$(r,p)$分布.

由二项分布可迅速推得\eqref{eq:3.2.9}. 事件$\{X=x\}$发生，当且仅当前$x-1$次试验中恰有$r-1$次成功，并且第$x$次试验成功. 在$x-1$次试验中有$r-1$次成功的概率为二项概率$\binom{x-1}{r-1}p^{r-1}(1-p)^{x-r}$，而第$x$次试验成功的概率为$p$. 把这两个概率相乘就得到\eqref{eq:3.2.9}.

负二项分布有时也用随机变量$Y=$第$r$次成功之前的失败次数来定义. 由于$Y=X-r$，这种表述在统计上与前面用$X=$第$r$次成功发生时的试验序号给出的表述等价. 利用$Y$与$X$的关系，负二项分布的另一种形式为
\begin{equation}\label{eq:3.2.10}\index{SI0351@负二项}
P(Y=y)=\binom{r+y-1}{y}p^r(1-p)^y,\qquad y=0,1,\ldots.
\end{equation}
除非另有说明，今后提到负二项$(r,p)$分布时，我们都使用这个pmf.

负二项分布因以下关系而得名：
$$\binom{r+y-1}{y}=(-1)^y\binom{-r}{y}=(-1)^y\frac{(-r)(-r-1)(-r-2)\cdots(-r-y+1)}{(y)(y-1)(y-2)\cdots(2)(1)},$$
这实际上就是负整数二项系数的定义式［完整论述见\hyperref[bib:source118]{Feller（1968）}］. 代入\eqref{eq:3.2.10}得
$$P(Y=y)=(-1)^y\binom{-r}{y}p^r(1-p)^y,$$
它与二项分布具有惊人的相似性.

$\sum_{y=0}^\infty P(Y=y)=1$这一事实不容易直接验证，但可由二项式定理中包含负指数的推广得到. 这里不再深入讨论. \hyperref[bib:source118]{Feller（1968）}对二项系数有极好的论述.

用类似于二项分布的技巧，可以计算$Y$的均值与方差：
\begin{align*}
EY&=\sum_{y=0}^\infty y\binom{r+y-1}{y}p^r(1-p)^y\\
&=\sum_{y=1}^\infty\frac{(r+y-1)!}{(y-1)!(r-1)!}p^r(1-p)^y\\
&=\sum_{y=1}^\infty r\binom{r+y-1}{y-1}p^r(1-p)^y.
\end{align*}
\sourcepageanchor{110}

现在令$z=y-1$，上述和变为
\begin{align*}
EY&=\sum_{z=0}^\infty r\binom{r+z}{z}p^r(1-p)^{z+1}\\
&=\frac{r(1-p)}p\sum_{z=0}^\infty\binom{(r+1)+z-1}{z}p^{r+1}(1-p)^z\quad(\text{求和项为负二项pmf})\\
&=\frac{r(1-p)}p.
\end{align*}
由于此和遍历负二项$(r+1,p)$分布的全部取值，所以它等于1. 类似计算给出
$$\operatorname{Var}Y=\frac{r(1-p)}{p^2}.$$
用均值对负二项分布重新参数化，是一种有趣且有时很有用的做法. 若定义参数$\mu=r(1-p)/p$，则$EY=\mu$，且稍作代数运算便可得到\index{SI0352@二次方差}\semanticindexanchor{352:110}\index{SI1067@均值的二次函数}\semanticindexanchor{1067:110}
$$\operatorname{Var}Y=\mu+\frac1r\mu^2.$$
方差是均值的二次函数. 这种关系在数据分析和理论研究中都有用（\hyperref[bib:source224]{Morris 1982）}.

负二项分布族以Poisson分布为其极限情形. 若$r\to\infty$且$p\to1$，使得$r(1-p)\to\lambda$，其中$0<\lambda<\infty$，则\index{SI0356@与泊松的关系}\semanticindexanchor{356:110}\index{SI0378@与负二项的关系}\semanticindexanchor{378:110}
\begin{align*}
EY&=\frac{r(1-p)}p\to\lambda,\\
\operatorname{Var}Y&=\frac{r(1-p)}{p^2}\to\lambda,
\end{align*}
这与Poisson分布的均值和方差一致. 要说明负二项$(r,p)\to\operatorname{Poisson}(\lambda)$，可以证明所有概率都收敛. 矩母函数收敛这一事实使我们预期会有这样的结果（见\autoref{exercise:3.15}）.

\begin{example}[逆二项抽样]\label{ex:3.2.6}\index{SI0592@逆二项抽样}
一种称为逆二项抽样（inverse binomial sampling）的技术，在生物总体的抽样中很有用. 若具有某一特征的个体比例为$p$，且我们持续抽样直到见到$r$个这样的个体，则抽取的个体总数就是一个负二项随机变量.

例如，假设我们关心果蝇总体中残翅个体所占的比例，决定持续抽样，直到找到100只这样的果蝇为止. 必须检查至少$N$只果蝇的概率为（利用\eqref{eq:3.2.9}）
\begin{align*}
P(X\ge N)&=\sum_{x=N}^\infty\binom{x-1}{99}p^{100}(1-p)^{x-100}\\
&=1-\sum_{x=100}^{N-1}\binom{x-1}{99}p^{100}(1-p)^{x-100}.
\end{align*}
给定$p$和$N$，就可计算这个表达式，确定我们很可能需要观察多少只果蝇. （虽然计算繁琐，但使用递推关系可以加快计算.）
\end{example}
\autoref{ex:3.2.6}表明，与Poisson分布一样，负二项分布也可用来对等待某事件发生的现象建模. 在负二项情形，我们等待的是指定次数的成功.
\sourcepageanchor{111}

\paragraph{几何分布.}\index{SI0318@几何}\semanticindexanchor{318:111}
几何分布（geometric distribution）是最简单的等待时间分布，也是负二项分布的特例. 在\eqref{eq:3.2.9}中令$r=1$，得到
$$P(X=x\mid p)=p(1-p)^{x-1},\qquad x=1,2,\ldots,$$
这就定义了成功概率为$p$的几何随机变量$X$的pmf. $X$可以解释为首次成功发生时的试验序号，所以我们是在“等待一次成功”. $\sum_{x=1}^\infty P(X=x)=1$由几何级数的性质得到. 对满足$|a|<1$的任何数$a$，有
$$\sum_{x=1}^\infty a^{x-1}=\frac1{1-a},$$
我们已在\autoref{ex:1.5.4}中遇到这个公式.

令$X=Y+1$，利用负二项分布的公式就可算出$X$的均值和方差：
$$EX=EY+1=\frac1p\qquad\text{及}\qquad\operatorname{Var}X=\frac{1-p}{p^2}.$$
几何分布具有一个有趣的性质，称为“无记忆性”（memoryless property）. 对整数$s>t$，有
\begin{equation}\label{eq:3.2.11}\index{SI0320@无记忆性}
P(X>s\mid X>t)=P(X>s-t);
\end{equation}
也就是说，几何分布“忘记”已经发生的事情. 已经观察到$t$次失败后，再发生$s-t$次失败的概率，与从序列开始就观察到$s-t$次失败的概率相同. 换言之，出现连续一段失败的概率只取决于这段的长度，而不取决于它的位置.

为证明\eqref{eq:3.2.11}，先注意到，对任何整数$n$，
\begin{equation}\label{eq:3.2.12}
\begin{aligned}
P(X>n)&=P(n\text{次试验中没有成功})\\
&=(1-p)^n,
\end{aligned}
\end{equation}
因此
\begin{align*}
P(X>s\mid X>t)&=\frac{P(X>s\text{ 且 }X>t)}{P(X>t)}\\
&=\frac{P(X>s)}{P(X>t)}=(1-p)^{s-t}=P(X>s-t).
\end{align*}
\begin{example}[失效时间]\label{ex:3.2.7}
几何分布有时用来对元件的“寿命”或“失效前的时间”建模. 例如，若一个灯泡在任一天失效的概率为$.001$，那么它至少持续30天的概率为
$$P(X>30)=\sum_{x=31}^\infty .001(1-.001)^{x-1}=(.999)^{30}=.970.$$
\end{example}
几何分布的无记忆性描述了一种很特殊的“不老化”性质. 它表明，几何分布不适用于对预期失效概率随时间增加的寿命建模. 还有其他分布可用于对不同类型的老化建模，例如见\hyperref[bib:source9]{Barlow和Proschan（1975）}.



\sourcepageanchor{112}
\subsection{连续分布}\label{sec:3.3}
本节讨论一些更常见的、有熟悉名称的连续分布族. 这里提到的分布绝不涵盖统计学中使用的所有分布. 实际上，如\autoref{sec:1.6}所见，任何非负、可积的函数都可变换成一个pdf.

\paragraph{均匀分布.}
把质量均匀分布在区间$[a,b]$上，就定义了连续均匀分布（uniform distribution）. 它的pdf为
\begin{equation}\label{eq:3.3.1}\index{SI0387@均匀}
f(x\mid a,b)=\begin{cases}\dfrac1{b-a},&x\in[a,b],\\0,&\text{其他情形}.\end{cases}
\end{equation}
容易验证$\int_a^b f(x)\,dx=1$. 还有
\begin{align*}
EX&=\int_a^b\frac{x}{b-a}\,dx=\frac{b+a}2,\\
\operatorname{Var}X&=\int_a^b\frac{(x-\frac{b+a}2)^2}{b-a}\,dx=\frac{(b-a)^2}{12}.
\end{align*}

\paragraph{Gamma分布.}
Gamma分布族是$[0,\infty)$上一个灵活的分布族，可用\autoref{sec:1.6}讨论的构造方法导出. 若$\alpha$为正常数，则积分
$$\int_0^\infty t^{\alpha-1}e^{-t}\,dt$$
有限. 若$\alpha$为正整数，该积分可用闭式表示，否则不能.\footnote{译注：原文“否则不能”过于绝对；例如本页随后即指出$\Gamma(1/2)=\sqrt\pi$，再由递推关系可得所有正半整数处的闭式值. 此处保留原文表述.} 无论哪种情形，它的值都定义了Gamma函数（gamma function），
\begin{equation}\label{eq:3.3.2}\index{SI0450@伽马函数（另见分布：伽马）}
\Gamma(\alpha)=\int_0^\infty t^{\alpha-1}e^{-t}\,dt.
\end{equation}
Gamma函数满足许多有用的关系，特别是
\begin{equation}\label{eq:3.3.3}
\Gamma(\alpha+1)=\alpha\Gamma(\alpha),\qquad\alpha>0,
\end{equation}
这可用分部积分验证. 结合\eqref{eq:3.3.3}与容易验证的事实$\Gamma(1)=1$，对任意整数$n>0$有
\begin{equation}\label{eq:3.3.4}
\Gamma(n)=(n-1)!.
\end{equation}
（另一个有用的特例是$\Gamma(1/2)=\sqrt\pi$，将在\eqref{eq:3.3.15}中看到.）

\eqref{eq:3.3.3}和\eqref{eq:3.3.4}给出的递推关系，使Gamma函数值的计算更为容易. 借助递推关系，只须知道$\Gamma(c)$在$0<c\le1$时的值，就能计算Gamma函数的任意值.

由于\eqref{eq:3.3.2}中的被积函数为正，立即得到
\begin{equation}\label{eq:3.3.5}
f(t)=\frac{t^{\alpha-1}e^{-t}}{\Gamma(\alpha)},\qquad0<t<\infty
\end{equation}
是一个pdf. 不过，完整的Gamma族具有两个参数. 取正常数$\beta$，在\eqref{eq:3.3.5}中作变量变换，求出随机变量$X=\beta T$的pdf，就得到Gamma$(\alpha,\beta)$族：
\begin{equation}\label{eq:3.3.6}\index{SI0306@伽马}
f(x\mid\alpha,\beta)=\frac1{\Gamma(\alpha)\beta^\alpha}x^{\alpha-1}e^{-x/\beta},\qquad0<x<\infty,\quad\alpha>0,\quad\beta>0.
\end{equation}
\sourcepageanchor{113}

参数$\alpha$称为形状参数（shape parameter），因为它主要影响分布峰的形状；参数$\beta$称为尺度参数（scale parameter），因为它主要影响分布的离散程度.

Gamma$(\alpha,\beta)$分布的均值为
\begin{equation}\label{eq:3.3.7}
EX=\frac1{\Gamma(\alpha)\beta^\alpha}\int_0^\infty xx^{\alpha-1}e^{-x/\beta}\,dx.
\end{equation}
为计算\eqref{eq:3.3.7}，注意被积函数是Gamma$(\alpha+1,\beta)$的pdf的核. 由\eqref{eq:3.3.6}知道，对任意$\alpha,\beta>0$，
\begin{equation}\label{eq:3.3.8}
\int_0^\infty x^{\alpha-1}e^{-x/\beta}\,dx=\Gamma(\alpha)\beta^\alpha,
\end{equation}
所以
\begin{align*}
EX&=\frac1{\Gamma(\alpha)\beta^\alpha}\int_0^\infty x^\alpha e^{-x/\beta}\,dx\\
&=\frac1{\Gamma(\alpha)\beta^\alpha}\Gamma(\alpha+1)\beta^{\alpha+1}\\
&=\frac{\alpha\Gamma(\alpha)\beta}{\Gamma(\alpha)}\qquad(\text{由\eqref{eq:3.3.3}})\\
&=\alpha\beta.
\end{align*}
注意，计算$EX$时，我们再次使用了把积分中的函数识别为另一个pdf的核这一技巧. （此前已用此技巧在\autoref{ex:2.3.8}中计算Gamma矩母函数，并在离散情形下于\autoref{ex:2.2.3}和\autoref{ex:2.3.5}中进行二项分布的计算.）

Gamma$(\alpha,\beta)$分布方差的计算方法类似于均值. 特别地，计算$EX^2$时处理的是Gamma$(\alpha+2,\beta)$分布的核. 结果为
$$\operatorname{Var}X=\alpha\beta^2.$$
在\autoref{ex:2.3.8}中，我们算出了Gamma$(\alpha,\beta)$分布的矩母函数：
$$M_X(t)=\left(\frac1{1-\beta t}\right)^\alpha,\qquad t<\frac1\beta.$$

\begin{example}[Gamma与Poisson的关系]\label{ex:3.3.1}\index{SI0315@与泊松的关系}\index{SI0377@与伽马的关系}
Gamma分布与Poisson分布之间有一个有趣的关系. 若$X$是Gamma$(\alpha,\beta)$随机变量，其中$\alpha$为整数，则对任意$x$，
\begin{equation}\label{eq:3.3.9}
P(X\le x)=P(Y\ge\alpha),
\end{equation}
其中$Y\sim\operatorname{Poisson}(x/\beta)$. 连续使用分部积分可证明\eqref{eq:3.3.9}，如下所示. 因为$\alpha$为整数，写$\Gamma(\alpha)=(\alpha-1)!$，得到
\begin{align*}
P(X\le x)&=\frac1{(\alpha-1)!\beta^\alpha}\int_0^x t^{\alpha-1}e^{-t/\beta}\,dt\\
&=\frac1{(\alpha-1)!\beta^\alpha}\left[\left.-t^{\alpha-1}\beta e^{-t/\beta}\right|_0^x+\int_0^x(\alpha-1)t^{\alpha-2}\beta e^{-t/\beta}\,dt\right],
\end{align*}
其中使用分部积分代换$u=t^{\alpha-1}$，$dv=e^{-t/\beta}\,dt$. 继续计算，有
\begin{align*}
P(X\le x)&=\frac{-1}{(\alpha-1)!\beta^{\alpha-1}}x^{\alpha-1}e^{-x/\beta}+\frac1{(\alpha-2)!\beta^{\alpha-1}}\int_0^x t^{\alpha-2}e^{-t/\beta}\,dt\\
&=\frac1{(\alpha-2)!\beta^{\alpha-1}}\int_0^x t^{\alpha-2}e^{-t/\beta}\,dt-P(Y=\alpha-1),
\end{align*}
其中$Y\sim\operatorname{Poisson}(x/\beta)$. 如此继续，就可建立\eqref{eq:3.3.9}. （见\autoref{exercise:3.19}.）
\end{example}
\sourcepageanchor{114}

Gamma分布有若干重要特例. 若令$\alpha=p/2$，其中$p$为整数，并令$\beta=2$，则Gamma的pdf成为
\begin{equation}\label{eq:3.3.10}\index{SI0245@自由度}\index{SI0280@卡方}
f(x\mid p)=\frac1{\Gamma(p/2)2^{p/2}}x^{(p/2)-1}e^{-x/2},\qquad0<x<\infty,
\end{equation}
即具有$p$个自由度的卡方pdf（chi squared pdf with $p$ degrees of freedom）. 卡方分布的均值、方差和矩母函数，都可以利用前面导出的Gamma公式计算.

卡方分布在统计推断中起着重要作用，特别是在从正态分布抽样时. \autoref{sec:5}将详细讨论这一主题.

令$\alpha=1$，就得到Gamma分布的另一个重要特例. 此时
\begin{equation}\label{eq:3.3.11}\index{SI0289@指数}
f(x\mid\beta)=\frac1\beta e^{-x/\beta},\qquad0<x<\infty,
\end{equation}
即尺度参数为$\beta$的指数pdf（exponential pdf）. 它的均值与方差已在\autoref{ex:2.2.2}和\autoref{ex:2.3.3}中计算过.

指数分布可以用来对寿命建模，类似于离散情形中几何分布的用途. 实际上，指数分布也具有几何分布的“无记忆”性质. 若$X\sim\operatorname{exponential}(\beta)$，即pdf由\eqref{eq:3.3.11}给出，则对$s>t\ge0$有\index{SI0293@无记忆性}\semanticindexanchor{293:114}
$$P(X>s\mid X>t)=P(X>s-t),$$
因为
\begin{align*}
P(X>s\mid X>t)&=\frac{P(X>s,X>t)}{P(X>t)}\\
&=\frac{P(X>s)}{P(X>t)}\qquad(\text{因为 }s>t)\\
&=\frac{\int_s^\infty\frac1\beta e^{-x/\beta}\,dx}{\int_t^\infty\frac1\beta e^{-x/\beta}\,dx}
=\frac{e^{-s/\beta}}{e^{-t/\beta}}\\
&=e^{-(s-t)/\beta}=P(X>s-t).
\end{align*}
另一个同时与指数族和Gamma族有关的分布是Weibull分布. 若$X\sim\operatorname{exponential}(\beta)$，则$Y=X^{1/\gamma}$具有Weibull$(\gamma,\beta)$分布：
\begin{equation}\label{eq:3.3.12}\index{SI0391@威布尔}
f_Y(y\mid\gamma,\beta)=\frac\gamma\beta y^{\gamma-1}e^{-y^\gamma/\beta},\quad0<y<\infty,\quad\gamma>0,\quad\beta>0.
\end{equation}
显然，我们也可以从Weibull分布出发，再把指数分布作为特例（$\gamma=1$）推导出来. 这只是取舍偏好的问题. Weibull分布在失效时间数据分析中起着极为重要的作用［这一主题的全面论述见\hyperref[bib:source186]{Kalbfleisch和Prentice（1980）}］. 特别地，Weibull分布非常适合对风险函数（hazard function）建模［见\autoref{exercise:3.25}和\autoref{exercise:3.26}］.\index{SI0469@风险率函数}\semanticindexanchor{469:114}

\paragraph{正态分布.}\index{SI0317@高斯}\semanticindexanchor{317:114}\index{SI0358@正态}\semanticindexanchor{358:114}
正态分布（normal distribution，有时称为Gaussian分布）在统计学的大量内容中起着核心作用. 主要有三个原因. 第一，正态分布及与它相关的分布非常便于解析处理（虽然乍看之下可能并非如此）. 第二，正态分布具有熟悉的钟形，其对称性使它成为许多总体模型中颇有吸引力的选择. 尽管还有许多其他分布也呈钟形，
\sourcepageanchor{115}
其中大多数却没有正态分布那样便于解析处理. 第三，存在中心极限定理（Central Limit Theorem，详见\autoref{sec:5}）；它表明，在温和条件下，可以在大样本中用正态分布近似种类繁多的分布.

正态分布有两个参数，通常记作$\mu$和$\sigma^2$，它们分别是均值与方差. 均值为$\mu$、方差为$\sigma^2$的正态分布（通常记作$\operatorname n(\mu,\sigma^2)$）的pdf为
\begin{equation}\label{eq:3.3.13}
f(x\mid\mu,\sigma^2)=\frac1{\sqrt{2\pi}\sigma}e^{-(x-\mu)^2/(2\sigma^2)},\qquad-\infty<x<\infty.
\end{equation}
若$X\sim\operatorname n(\mu,\sigma^2)$，则随机变量$Z=(X-\mu)/\sigma$具有$\operatorname n(0,1)$分布，也称为标准正态分布（standard normal distribution）. 写出以下各式便容易证明：\index{SI0368@标准}\semanticindexanchor{368:115}
\begin{align*}
P(Z\le z)&=P\left(\frac{X-\mu}\sigma\le z\right)=P(X\le z\sigma+\mu)\\
&=\frac1{\sqrt{2\pi}\sigma}\int_{-\infty}^{z\sigma+\mu}e^{-(x-\mu)^2/(2\sigma^2)}\,dx\\
&=\frac1{\sqrt{2\pi}}\int_{-\infty}^z e^{-t^2/2}\,dt\qquad\left(\text{令 }t=\frac{x-\mu}\sigma\right),
\end{align*}
这说明$P(Z\le z)$就是标准正态cdf.

因此，所有正态概率都可以利用标准正态分布计算. 此外，计算期望值时可先在$\operatorname n(0,1)$情形下完成具体运算，再把结果变换到$\operatorname n(\mu,\sigma^2)$情形，从而简化计算. 例如，若$Z\sim\operatorname n(0,1)$，则
$$EZ=\frac1{\sqrt{2\pi}}\int_{-\infty}^\infty ze^{-z^2/2}\,dz=\left.-\frac1{\sqrt{2\pi}}e^{-z^2/2}\right|_{-\infty}^\infty=0,$$
所以若$X\sim\operatorname n(\mu,\sigma^2)$，由\autoref{thm:2.2.5}可得
$$EX=E(\mu+\sigma Z)=\mu+\sigma EZ=\mu.$$
类似地，$\operatorname{Var}Z=1$，再用\autoref{thm:2.3.4}，得到$\operatorname{Var}X=\sigma^2$.

我们尚未证明\eqref{eq:3.3.13}在整个实数轴上的积分等于1. 应用标准化变换，只须证明
$$\frac1{\sqrt{2\pi}}\int_{-\infty}^\infty e^{-z^2/2}\,dz=1.$$
注意，上面的被积函数关于0对称，故在$(-\infty,0)$上的积分等于在$(0,\infty)$上的积分. 因此，问题归结为证明
\begin{equation}\label{eq:3.3.14}
\int_0^\infty e^{-z^2/2}\,dz=\frac{\sqrt{2\pi}}2=\sqrt{\frac\pi2}.
\end{equation}
函数$e^{-z^2/2}$没有可用初等函数显式表示（即以闭式表示）的原函数，所以我们不能直接积分. 事实上，这是一类积分的例子：要么你知道怎样做，要么可能花很长时间仍毫无进展. 因为\eqref{eq:3.3.14}的两边都为正，只要证明它们的平方相等，就能得到该等式. 对\eqref{eq:3.3.14}中的积分取平方，得到
\begin{align*}
\left(\int_0^\infty e^{-z^2/2}\,dz\right)^2
&=\left(\int_0^\infty e^{-t^2/2}\,dt\right)\left(\int_0^\infty e^{-u^2/2}\,du\right)\\
&=\int_0^\infty\int_0^\infty e^{-(t^2+u^2)/2}\,dt\,du.
\end{align*}
\sourcepageanchor{116}

积分变量只是哑变量，所以可以改变它们的名称. 现在转为极坐标. 定义
$$t=r\cos\theta\qquad\text{及}\qquad u=r\sin\theta.$$
则$t^2+u^2=r^2$，且$dt\,du=r\,d\theta\,dr$，积分范围变为$0<r<\infty$，$0<\theta<\pi/2$（因为$t$和$u$均限制为正，所以$\theta$的上限为$\pi/2$）. 于是
\begin{align*}
\int_0^\infty\int_0^\infty e^{-(t^2+u^2)/2}\,dt\,du
&=\int_0^\infty\int_0^{\pi/2}re^{-r^2/2}\,d\theta\,dr\\
&=\frac\pi2\int_0^\infty re^{-r^2/2}\,dr
=\frac\pi2\left[\left.-e^{-r^2/2}\right|_0^\infty\right]=\frac\pi2,
\end{align*}
从而证明了\eqref{eq:3.3.14}.

这个积分与Gamma函数密切相关；事实上，在\eqref{eq:3.3.14}中作代换$w=\frac12z^2$，就会看到该积分实质上就是$\Gamma(1/2)$. 如果仔细处理常数，就会看出\eqref{eq:3.3.14}蕴含
\begin{equation}\label{eq:3.3.15}
\Gamma\left(\frac12\right)=\int_0^\infty w^{-1/2}e^{-w}\,dw=\sqrt\pi.
\end{equation}
正态分布在以下意义上有些特殊：它的两个参数$\mu$（均值）和$\sigma^2$（方差）提供了有关分布确切形状与位置的全部信息. 分布由$\mu$和$\sigma^2$确定这一性质，并非正态pdf所独有，而是一族称为位置—尺度族（location--scale families）的pdf所共有的性质；将在\autoref{sec:3.5}讨论.

直接作微积分运算可知，正态pdf \eqref{eq:3.3.13}在$x=\mu$处取最大值，在$\mu\pm\sigma$处有拐点（曲线由凹变凸的位置）. 此外，离均值1、2或3个标准差以内的概率为
\begin{align*}
P(|X-\mu|\le\sigma)&=P(|Z|\le1)=.6826,\\
P(|X-\mu|\le2\sigma)&=P(|Z|\le2)=.9544,\\
P(|X-\mu|\le3\sigma)&=P(|Z|\le3)=.9974,
\end{align*}
其中$X\sim\operatorname n(\mu,\sigma^2)$，$Z\sim\operatorname n(0,1)$，这些数值可由许多计算机软件包或数表得到. 通常报告的两位小数值分别为$.68$、$.95$和$.99$. 尽管这些并非四舍五入值，却是常用的数值. \autoref{fig:3.3.1}展示了正态pdf及这些关键特征.

在正态分布的众多应用中，一个重要应用是近似其他分布（中心极限定理为此提供了部分依据）. 例如，若$X\sim\operatorname{binomial}(n,p)$，则$EX=np$且$\operatorname{Var}X=np(1-p)$；在适当条件下，$X$的分布可用均值$\mu=np$、方差$\sigma^2=np(1-p)$的正态随机变量的分布近似. “适当条件”是$n$应很大，且$p$不应过于极端（接近0或1）. 我们希望$n$足够大，使$X$有足够多的（离散）取值，从而用连续分布近似是合理的；而$p$应处于“中间”，使二项分布像正态分布一样近乎对称. 与大多数近似一样，这里没有绝对准则，每次应用都应检查近似是否足以满足预定用途. 一个保守的经验规则是：若$\min(np,n(1-p))\ge5$，则近似较好.
\sourcepageanchor{117}

\begin{figure}[htbp]
\centering\includegraphics[width=.72\textwidth]{assets/chapter3/fig3-3-1.pdf}
\caption{标准正态密度}\label{fig:3.3.1}
\end{figure}
\begin{example}[正态近似]\label{ex:3.3.2}\index{SI0059@以正态近似二项}\index{SI0198@连续性校正}
设$X\sim\operatorname{binomial}(25,.6)$. 可用均值$\mu=25(.6)=15$、标准差$\sigma=((25)(.6)(.4))^{1/2}=2.45$的正态随机变量$Y$来近似$X$. 因此，
$$P(X\le13)\approx P(Y\le13)=P\left(Z\le\frac{13-15}{2.45}\right)=P(Z\le-.82)=.206,$$
而精确的二项计算给出
$$P(X\le13)=\sum_{x=0}^{13}\binom{25}{x}(.6)^x(.4)^{25-x}=.267,$$
这表明正态近似还不错，但并不十分出色. 不过，通过“连续性校正”（continuity correction）可大大改善近似. 为了解其原理，请看\autoref{fig:3.3.2}，图中展示了$\operatorname{binomial}(25,.6)$的pmf与$\operatorname n(15,(2.45)^2)$的pdf. 我们以宽为1、高等于概率的条形画出二项pmf. 因而条形面积就给出二项概率. 注意在上面的近似中，近似正态的面积小于二项面积（正态面积是13处直线左边的全部面积，而二项面积包含13处直到13.5的整根条形）. 连续性校正通过给截断点加$1/2$，把这块面积补回去. 所以我们不直接近似$P(X\le13)$，而是近似与它等价的表达式（由于离散性）$P(X\le13.5)$，得到
$$P(X\le13)=P(X\le13.5)\approx P(Y\le13.5)=P(Z\le-.61)=.271,$$
\begin{figure}[htbp]
\centering\includegraphics[width=.8\textwidth]{assets/chapter3/fig3-3-2.pdf}
\caption{用正态$(15,(2.45)^2)$近似二项$(25,.6)$}\label{fig:3.3.2}
\end{figure}
\sourcepageanchor{118}
这是好得多的近似. 一般而言，带连续性校正的正态近似远优于没有连续性校正的近似.

下端点也要作校正. 若$X\sim\operatorname{binomial}(n,p)$，$Y\sim\operatorname n(np,np(1-p))$，则我们采用近似
\begin{align*}
P(X\le x)&\approx P(Y\le x+1/2),\\
P(X\ge x)&\approx P(Y\ge x-1/2).
\end{align*}
\end{example}

\paragraph{Beta分布.}
Beta分布族是$(0,1)$上由两个参数标识的连续分布族. Beta$(\alpha,\beta)$的pdf为
\begin{equation}\label{eq:3.3.16}\index{SI0254@贝塔}
f(x\mid\alpha,\beta)=\frac1{B(\alpha,\beta)}x^{\alpha-1}(1-x)^{\beta-1},\quad0<x<1,\quad\alpha>0,\quad\beta>0,
\end{equation}
其中$B(\alpha,\beta)$表示Beta函数（beta function），
$$B(\alpha,\beta)=\int_0^1x^{\alpha-1}(1-x)^{\beta-1}\,dx.$$
Beta函数通过以下恒等式与Gamma函数联系起来：
\begin{equation}\label{eq:3.3.17}\index{SI0089@贝塔函数（另见分布：贝塔）}
B(\alpha,\beta)=\frac{\Gamma(\alpha)\Gamma(\beta)}{\Gamma(\alpha+\beta)}.
\end{equation}
处理Beta函数时，\eqref{eq:3.3.17}非常有用，使我们可以利用Gamma函数的性质. 实际上，我们不会直接处理Beta函数，而会在所有计算中使用\eqref{eq:3.3.17}.

Beta分布是少数几个把概率1赋予一个有限区间的常见“有名”分布之一，这里取的区间是$(0,1)$. 因此，Beta分布常用来对自然位于0与1之间的比例建模. \autoref{sec:4}将给出这方面的例子.

由于pdf的特殊形式，Beta分布的矩很容易计算. 对$n>-\alpha$，有
\begin{align*}
EX^n&=\frac1{B(\alpha,\beta)}\int_0^1x^nx^{\alpha-1}(1-x)^{\beta-1}\,dx\\
&=\frac1{B(\alpha,\beta)}\int_0^1x^{(\alpha+n)-1}(1-x)^{\beta-1}\,dx.
\end{align*}
现在可识别出被积函数是Beta$(\alpha+n,\beta)$的pdf的核，因此
\begin{equation}\label{eq:3.3.18}
EX^n=\frac{B(\alpha+n,\beta)}{B(\alpha,\beta)}=\frac{\Gamma(\alpha+n)\Gamma(\alpha+\beta)}{\Gamma(\alpha+\beta+n)\Gamma(\alpha)}.
\end{equation}
在\eqref{eq:3.3.18}中取$n=1$和$n=2$，并用\eqref{eq:3.3.3}，可算出Beta$(\alpha,\beta)$分布的均值与方差：
$$EX=\frac\alpha{\alpha+\beta}\qquad\text{及}\qquad\operatorname{Var}X=\frac{\alpha\beta}{(\alpha+\beta)^2(\alpha+\beta+1)}.$$
随着参数$\alpha$和$\beta$变化，Beta分布呈现许多形状，如\autoref{fig:3.3.3}所示. pdf可以严格递增（$\alpha>1,\beta=1$）、严格递减（$\alpha=1,\beta>1$）、呈U形（$\alpha<1,\beta<1$）或单峰\mbox{（$\alpha>1,\beta>1$）.} 当$\alpha=\beta$时，pdf关于$1/2$对称，
\sourcepageanchor{119}

\begin{figure}[htbp]
\centering\includegraphics[width=.72\textwidth]{assets/chapter3/fig3-3-3.pdf}
\caption{Beta密度}\label{fig:3.3.3}
\end{figure}
\begin{figure}[htbp]
\centering\includegraphics[width=.72\textwidth]{assets/chapter3/fig3-3-4.pdf}
\caption{对称Beta密度}\label{fig:3.3.4}
\end{figure}
其均值（必然）为$1/2$，方差为$(4(2\alpha+1))^{-1}$. 随$\alpha$增加，pdf更集中，但保持对称，如\autoref{fig:3.3.4}所示. 最后，若$\alpha=\beta=1$，Beta分布就化为均匀$(0,1)$分布，说明均匀分布可视为Beta族的成员. Beta分布还通过一个变换与$F$分布相关，后者在统计分析中起着极为重要的作用（见\autoref{sec:5.3}）.

\paragraph{Cauchy分布.}
Cauchy分布是在$(-\infty,\infty)$上对称的钟形分布，其pdf为
\begin{equation}\label{eq:3.3.19}\index{SI0276@柯西}
f(x\mid\theta)=\frac1\pi\frac1{1+(x-\theta)^2},\qquad-\infty<x<\infty,\quad-\infty<\theta<\infty.
\end{equation}
（更一般的Cauchy pdf见\autoref{exercise:3.39}.）从外观上看，Cauchy分布似乎与正态分布没有很大差别. 然而，二者确实有非常大的区别.
\sourcepageanchor{120}

\begin{figure}[htbp]
\centering\includegraphics[width=.72\textwidth]{assets/chapter3/fig3-3-5.pdf}
\caption{标准正态密度与Cauchy密度}\label{fig:3.3.5}
\end{figure}
如\autoref{sec:2}中所见，Cauchy分布的均值不存在，即
\begin{equation}\label{eq:3.3.20}
E|X|=\int_{-\infty}^\infty\frac1\pi\frac{|x|}{1+(x-\theta)^2}\,dx=\infty.
\end{equation}
\needspace{5\baselineskip}容易看出，\eqref{eq:3.3.19}对所有$\theta$都定义了一个恰当的pdf. 回顾$\frac d{dt}\arctan(t)=(1+t^2)^{-1}$，因此
$$\int_{-\infty}^\infty\frac1\pi\frac1{1+(x-\theta)^2}\,dx=\left.\frac1\pi\arctan(x-\theta)\right|_{-\infty}^\infty=1,$$
因为$\arctan(\pm\infty)=\pm\pi/2$.

由于$E|X|=\infty$，Cauchy分布的各阶矩都不存在；换言之，所有绝对矩都等于$\infty$.\footnote{译注：此处“各阶矩”“所有绝对矩”应按上下文理解为正整数阶矩；若允许非整数阶，则$0<q<1$时$E|X|^q$有限.} 特别地，矩母函数不存在.

\eqref{eq:3.3.19}中的参数$\theta$确实衡量分布的中心：它是中位数. 若$X$具有参数为$\theta$的Cauchy分布，则由\autoref{exercise:3.37}可得$P(X\ge\theta)=1/2$，这说明$\theta$是分布的中位数. \autoref{fig:3.3.5}把Cauchy$(0)$分布与$\operatorname n(0,1)$分布画在一起；可以看到二者形状相似，但Cauchy的尾部厚得多.

Cauchy分布在统计理论中扮演特殊角色. 它代表一种可用来检验猜想的极端情形. 但不要误以为Cauchy分布只是一个病态例子，因为它会在最意想不到的地方出现. 例如，实验人员经常计算观测值之比，即随机变量之比. （衡量生长时，通常把体重与身高合成一个量——身高对应体重，即体重/身高.）一个令人惊讶的事实是，两个标准正态变量之比具有Cauchy分布（见\autoref{ex:4.3.6}）. 取比值可能导致表现不良的分布.

\paragraph{对数正态分布.}
若随机变量$X$的对数服从正态分布（即$\log X\sim\operatorname n(\mu,\sigma^2)$），则$X$具有对数正态分布（lognormal distribution）. 用\autoref{thm:2.1.5}对正态pdf作直接变换，就得到$X$的pdf：
\begin{equation}\label{eq:3.3.21}\index{SI0332@对数正态}
f(x\mid\mu,\sigma^2)=\frac1{\sqrt{2\pi}\sigma}\frac1x e^{-(\log x-\mu)^2/(2\sigma^2)},\quad0<x<\infty,\quad-\infty<\mu<\infty,\quad\sigma>0,
\end{equation}
\sourcepageanchor{121}

即对数正态pdf. $X$的矩可以直接用\eqref{eq:3.3.21}计算，也可利用与正态分布的关系，写成
\begin{align*}
EX&=Ee^{\log X}=Ee^Y\qquad(Y=\log X\sim\operatorname n(\mu,\sigma^2))\\
&=e^{\mu+(\sigma^2/2)}.
\end{align*}
最后一个等式通过识别正态分布的矩母函数而得到（令$t=1$；见\autoref{exercise:2.33}）. 用类似技巧可计算$EX^2$，得到
$$\operatorname{Var}X=e^{2(\mu+\sigma^2)}-e^{2\mu+\sigma^2}.$$
对数正态分布的外观与Gamma分布相似，如\autoref{fig:3.3.6}所示. 当感兴趣的变量右偏时，这个分布在建模应用中很常见. 例如，收入必然右偏，而用对数正态分布建模，就可以对$\log(\text{收入})$使用正态理论统计方法，这十分方便.

\paragraph{双指数分布.}
把指数分布关于其均值反射，就形成双指数分布（double exponential distribution）.\footnote{译注：原文说“关于其均值反射”；就下面给出的密度而言，准确的构造是以$\mu$为中心，将指数形状的右半支及其镜像各赋予一半质量. 普通指数分布直接关于其自身均值反射，并不会得到此式.} 它的pdf为
\begin{equation}\label{eq:3.3.22}
f(x\mid\mu,\sigma)=\frac1{2\sigma}e^{-|x-\mu|/\sigma},\qquad-\infty<x<\infty,\quad-\infty<\mu<\infty,\quad\sigma>0.
\end{equation}
双指数给出一个具有“肥”尾（比正态尾部肥得多）的对称分布，但它的所有矩仍然存在. 直接计算可得
$$EX=\mu\qquad\text{及}\qquad\operatorname{Var}X=2\sigma^2.$$
双指数分布并非钟形. 事实上，它在$x=\mu$处有一个尖峰（更正式地说，是一个不可微点）. 对这一分布作解析处理时，牢记这一点很重要. 绝对值符号在积分时也可能带来麻烦，最好以$x=\mu$为界把积分分成两个区域：
\begin{equation}\label{eq:3.3.23}\index{SI0288@双指数}
\begin{aligned}
EX&=\int_{-\infty}^\infty\frac{x}{2\sigma}e^{-|x-\mu|/\sigma}\,dx\\
&=\int_{-\infty}^\mu\frac{x}{2\sigma}e^{(x-\mu)/\sigma}\,dx+\int_\mu^\infty\frac{x}{2\sigma}e^{-(x-\mu)/\sigma}\,dx.
\end{aligned}
\end{equation}
注意，在这两个积分区域内可以去掉绝对值符号. （一般而言，这个策略对处理含绝对值的积分很有用：把积分区域分开，使绝对值符号可以去掉.）对每个积分使用分部积分，就可完成\eqref{eq:3.3.23}的计算.

还有许多连续分布可用于不同的统计应用，其中很多会在本书后续内容中出现. \hyperref[bib:source177]{Johnson和Kotz（1969、1970a、1970b）}的综合性著作，是查阅多数常用统计分布的宝贵参考资料.
\sourcepageanchor{122}

\begin{figure}[htbp]
\centering\includegraphics[height=.78\textheight]{assets/chapter3/fig3-3-6.pdf}
\caption{(a)若干对数正态密度；(b)若干Gamma密度}\label{fig:3.3.6}
\end{figure}



\subsection{指数族}\label{sec:3.4}
如果一个pdf或pmf族可表示为
\begin{equation}\label{eq:3.4.1}\index{SI0435@指数族}
f(x\mid\theta)=h(x)c(\theta)\exp\left(\sum_{i=1}^k w_i(\theta)t_i(x)\right),
\end{equation}
就称它为指数族（exponential family）.
\sourcepageanchor{123}

这里$h(x)\ge0$和$t_1(x),\ldots,t_k(x)$是观测值$x$的实值函数（它们不能依赖于$\theta$），而$c(\theta)\ge0$和$w_1(\theta),\ldots,w_k(\theta)$是可能为向量值的参数$\theta$的实值函数（它们不能依赖\mbox{于$x$）.} 上一节介绍的许多常见分布族都是指数族，包括连续的正态族、Gamma族和Beta族，以及离散的二项族、Poisson族和负二项族.

要验证一个pdf或pmf族是指数族，必须识别出函数$h(x)$、$c(\theta)$、$w_i(\theta)$和$t_i(x)$，并说明该族具有\eqref{eq:3.4.1}的形式. 下例说明这一点.

\begin{example}[二项指数族]\label{ex:3.4.1}\index{SI0262@指数族}
设$n$为正整数，考虑$0<p<1$的二项$(n,p)$族. 对$x=0,\ldots,n$及$0<p<1$，这个族的pmf为
\begin{equation}\label{eq:3.4.2}
\begin{aligned}
f(x\mid p)&=\binom nx p^x(1-p)^{n-x}\\
&=\binom nx(1-p)^n\left(\frac p{1-p}\right)^x\\
&=\binom nx(1-p)^n\exp\left(\log\left(\frac p{1-p}\right)x\right).
\end{aligned}
\end{equation}
定义
$$h(x)=\begin{cases}\binom nx,&x=0,\ldots,n,\\0,&\text{其他情形},\end{cases}\qquad c(p)=(1-p)^n,\quad0<p<1,$$
$$w_1(p)=\log\left(\frac p{1-p}\right),\quad0<p<1,\qquad\text{及}\qquad t_1(x)=x.$$
于是
\begin{equation}\label{eq:3.4.3}
f(x\mid p)=h(x)c(p)\exp[w_1(p)t_1(x)],
\end{equation}
这就是$k=1$时的\eqref{eq:3.4.1}. 特别注意，只有$x=0,\ldots,n$时$h(x)>0$，且$c(p)$只在$0<p<1$时有定义. 这一点很重要，因为对所有$x$值，\eqref{eq:3.4.3}都必须与\eqref{eq:3.4.2}一致，而且只有$0<p<1$时它才是指数族（因此参数的函数只在这里定义）. 此外，二项模型有时包含参数值$p=0$和1，但这里没有把它们包括进来，因为当$p=0$和1时，使$f(x\mid p)>0$的$x$值集合与其他$p$值时不同.
\end{example}
\eqref{eq:3.4.1}的特殊形式使指数族具有许多良好的数学性质. 对统计模型更重要的是，这种形式还带来许多良好的统计性质. 下面说明计算指数族的矩的一条捷径.

\begin{theorem}\label{thm:3.4.2}\index{SI0439@矩}\index{SI0538@指数族}
若$X$是pdf或pmf具有\eqref{eq:3.4.1}形式的随机变量，则\footnote{译注：上述公式需要相应参数导数存在，并且允许把求导与积分或求和交换. 原定理未在陈述中写出这些正则性条件；其证明所用的运算须在这些条件成立处理解.}
\begin{equation}\label{eq:3.4.4}
E\left(\sum_{i=1}^k\frac{\partial w_i(\theta)}{\partial\theta_j}t_i(X)\right)=-\frac\partial{\partial\theta_j}\log c(\theta)
\end{equation}
以及
\begin{equation}\label{eq:3.4.5}
\operatorname{Var}\left(\sum_{i=1}^k\frac{\partial w_i(\theta)}{\partial\theta_j}t_i(X)\right)=-\frac{\partial^2}{\partial\theta_j^2}\log c(\theta)-E\left(\sum_{i=1}^k\frac{\partial^2w_i(\theta)}{\partial\theta_j^2}t_i(X)\right).
\end{equation}
\end{theorem}

这些等式虽然看起来令人生畏，但用于具体情形时却可能相当好用. 它们的优点在于，可把积分或求和替换为通常更直接的求导.
\sourcepageanchor{124}

\begin{example}[二项均值和方差]\label{ex:3.4.3}\index{SI0266@矩}
由\autoref{ex:3.4.1}有
\begin{align*}
\frac d{dp}w_1(p)&=\frac d{dp}\log\frac p{1-p}=\frac1{p(1-p)},\\
\frac d{dp}\log c(p)&=\frac d{dp}n\log(1-p)=\frac{-n}{1-p}.
\end{align*}
因而由\autoref{thm:3.4.2}得到
$$E\left(\frac1{p(1-p)}X\right)=\frac n{1-p},$$
稍加整理就得到$E(X)=np$. 方差恒等式也以类似方式起作用.
\end{example}
\autoref{thm:3.4.2}的证明是一段微积分练习，留在\autoref{exercise:3.31}中. 一个特殊情形还见\autoref{exercise:3.32}.

现在来看另一个例子，以及指数族的其他一些特征.

\begin{example}[正态指数族]\label{ex:3.4.4}\index{SI0360@指数族}
设$f(x\mid\mu,\sigma^2)$为$\operatorname n(\mu,\sigma^2)$的pdf族，其中$\theta=(\mu,\sigma)$，$-\infty<\mu<\infty$，$\sigma>0$. 则
\begin{equation}\label{eq:3.4.6}
\begin{aligned}
f(x\mid\mu,\sigma^2)&=\frac1{\sqrt{2\pi}\sigma}\exp\left(-\frac{(x-\mu)^2}{2\sigma^2}\right)\\
&=\frac1{\sqrt{2\pi}\sigma}\exp\left(-\frac{\mu^2}{2\sigma^2}\right)\exp\left(-\frac{x^2}{2\sigma^2}+\frac{\mu x}{\sigma^2}\right).
\end{aligned}
\end{equation}
定义
\begin{align*}
h(x)&=1\quad\text{对所有 }x,\\
c(\theta)=c(\mu,\sigma)&=\frac1{\sqrt{2\pi}\sigma}\exp\left(\frac{-\mu^2}{2\sigma^2}\right),\quad-\infty<\mu<\infty,\ \sigma>0,\\
w_1(\mu,\sigma)&=\frac1{\sigma^2},\quad\sigma>0,\qquad w_2(\mu,\sigma)=\frac\mu{\sigma^2},\quad\sigma>0,\\
t_1(x)&=-x^2/2,\qquad\text{及}\qquad t_2(x)=x.
\end{align*}
于是
$$f(x\mid\mu,\sigma^2)=h(x)c(\mu,\sigma)\exp[w_1(\mu,\sigma)t_1(x)+w_2(\mu,\sigma)t_2(x)],$$
这就是$k=2$时的\eqref{eq:3.4.1}. 再次注意，参数的函数只在参数的取值范围内定义.
\end{example}
一般而言，在指数族中，使$f(x\mid\theta)>0$的$x$值集合不能依赖于$\theta$. pdf或pmf的整个定义都必须包含在\eqref{eq:3.4.1}的形式中. 最容易的实现方法是使用示性函数，把$x$的取值范围纳入$f(x\mid\theta)$的表达式.

\begin{definition}\label{def:3.4.5}\index{SI0556@示性函数}
集合$A$的示性函数（indicator function）通常记作$I_A(x)$，定义为
$$I_A(x)=\begin{cases}1,&x\in A,\\0,&x\notin A.\end{cases}$$
另一种记号是$I(x\in A)$.
\end{definition}
\sourcepageanchor{125}

因此，\autoref{ex:3.4.4}中的正态pdf可写为
$$f(x\mid\mu,\sigma^2)=h(x)c(\mu,\sigma)\exp[w_1(\mu,\sigma)t_1(x)+w_2(\mu,\sigma)t_2(x)]I_{(-\infty,\infty)}(x).$$
因为示性函数只是$x$的函数，可以把它并入$h(x)$，从而表明这个pdf具有\eqref{eq:3.4.1}的形式.

由\eqref{eq:3.4.1}，因子$\exp(\cdot)$恒为正，所以对任意$\theta\in\Theta$，即对使$c(\theta)>0$的任意$\theta$，有$\{x:f(x\mid\theta)>0\}=\{x:h(x)>0\}$，而这个集合不依赖于$\theta$. 所以，例如由$f(x\mid\theta)=\theta^{-1}\exp(1-(x/\theta))$，$0<\theta<x<\infty$给出的pdf族不是指数族，尽管可写成$\theta^{-1}\exp(1-(x/\theta))=h(x)c(\theta)\exp(w(\theta)t(x))$，其中$h(x)=e^1$，$c(\theta)=\theta^{-1}$，$w(\theta)=\theta^{-1}$且$t(x)=-x$. 用示性函数把pdf写出来，就能看得很清楚：
$$f(x\mid\theta)=\theta^{-1}\exp\left(1-\left(\frac x\theta\right)\right)I_{[\theta,\infty)}(x).$$
这个示性函数无法并入\eqref{eq:3.4.1}中的任何函数，因为它既不是仅关于$x$的函数，也不是仅关于$\theta$的函数，而且不能表示成指数形式. 因此，这不是指数族.

指数族有时参数化为
\begin{equation}\label{eq:3.4.7}
f(x\mid\eta)=h(x)c^*(\eta)\exp\left(\sum_{i=1}^k\eta_i t_i(x)\right).
\end{equation}
这里$h(x)$和$t_i(x)$与原参数化\eqref{eq:3.4.1}中的函数相同. 集合
$$\mathcal H=\left\{\eta=(\eta_1,\ldots,\eta_k):\int_{-\infty}^\infty h(x)\exp\left(\sum_{i=1}^k\eta_i t_i(x)\right)\,dx<\infty\right\}$$
称为该族的自然参数空间（natural parameter space）. 如果$X$为离散型，就把积分换成对使$h(x)>0$的$x$值求和. 对$\eta\in\mathcal H$，必须有\index{SI0440@自然参数}\semanticindexanchor{440:125}
$$c^*(\eta)=\left[\int_{-\infty}^\infty h(x)\exp\left(\sum_{i=1}^k\eta_i t_i(x)\right)\,dx\right]^{-1},$$
以保证pdf的积分为1. 因为\eqref{eq:3.4.1}中原来的$f(x\mid\theta)$是pdf或pmf，集合$\{\eta=(w_1(\theta),\ldots,w_k(\theta)):\theta\in\Theta\}$必为自然参数空间的子集. 但也可能还有其他$\eta\in\mathcal H$. 自然参数化及自然参数空间有许多有用的数学性质. 例如，$\mathcal H$是凸集.

\begin{example}[续\autoref{ex:3.4.4}]\label{ex:3.4.6}
为确定正态分布族的自然参数空间，在\eqref{eq:3.4.6}中以$\eta_i$替换$w_i(\mu,\sigma)$，得到
\begin{equation}\label{eq:3.4.8}
f(x\mid\eta_1,\eta_2)=\frac{\sqrt{\eta_1}}{\sqrt{2\pi}}\exp\left(-\frac{\eta_2^2}{2\eta_1}\right)\exp\left(-\frac{\eta_1x^2}2+\eta_2x\right).
\end{equation}
积分有限，当且仅当$x^2$的系数为负. 这意味着$\eta_1$必须为正. 若$\eta_1>0$，则无论$\eta_2$取何值，积分都有限. 因此，自然参数空间为$\{(\eta_1,\eta_2):\eta_1>0,-\infty<\eta_2<\infty\}$. 对照\eqref{eq:3.4.8}和\eqref{eq:3.4.6}，可得$\eta_2=\mu/\sigma^2$，$\eta_1=1/\sigma^2$. 自然参数虽然提供了方便的数学表达，但有时缺乏像均值和方差那样简单的解释.
\end{example}
在表示式\eqref{eq:3.4.1}中，向量$\theta$的维数常常等于$k$，即指数中求和项的个数. 但不一定如此，向量$\theta$的维数也可能等于$d<k$. 这样的指数族称为弯曲指数族（curved exponential family）.

\begin{definition}\label{def:3.4.7}\index{SI0436@曲线}
弯曲指数族是具有\eqref{eq:3.4.1}形式的密度族，其中向量$\theta$的维数等于$d<k$. 若$d=k$，则该族是满指数族（full exponential family）. （另见\hyperref[sec:3.8.3]{杂录3.8.3}.）
\end{definition}
\sourcepageanchor{126}

\begin{example}[一个弯曲指数族]\label{ex:3.4.8}
\autoref{ex:3.4.4}中的正态族是满指数族. 但若假定$\sigma^2=\mu^2$，该族就成为弯曲的. （这样的模型可用于方差分析；见\autoref{exercise:11.1}和\autoref{exercise:11.2}.）此时
\begin{equation}\label{eq:3.4.9}
\begin{aligned}
f(x\mid\mu)&=\frac1{\sqrt{2\pi\mu^2}}\exp\left(-\frac{(x-\mu)^2}{2\mu^2}\right)\\
&=\frac1{\sqrt{2\pi\mu^2}}\exp\left(-\frac12\right)\exp\left(-\frac{x^2}{2\mu^2}+\frac x\mu\right).
\end{aligned}
\end{equation}
对正态族，满指数族的参数空间为$(\mu,\sigma^2)=\mathbb R\times(0,\infty)$，而弯曲族的参数空间$(\mu,\sigma^2)=(\mu,\mu^2)$是一条抛物线.
\end{example}
弯曲指数族有许多用途. 下例给出一个简单应用.

\begin{example}[正态近似]\label{ex:3.4.9}\index{SI0373@正态近似}
在\autoref{sec:5}中将看到，若$X_1,\ldots,X_n$为来自Poisson$(\lambda)$总体的样本，则$\bar X=\sum_iX_i/n$的分布近似为
$$\bar X\sim\operatorname n(\lambda,\lambda/n),$$
这是一个弯曲指数族.

中心极限定理（\autoref{thm:5.5.14}）为$\operatorname n(\lambda,\lambda/n)$近似提供了依据. 实际上可以意识到，大多数这样的中心极限定理近似都会导出弯曲正态族. 我们已经见过正态—二项近似（\autoref{ex:3.3.2}）：若$X_1,X_2,\ldots,X_n$独立同分布，服从Bernoulli$(p)$，则近似地
$$\bar X\sim\operatorname n(p,p(1-p)/n).$$
另一个例子见\autoref{ex:5.5.16}.
\end{example}
虽然参数空间维数较低这一事实会对族的性质产生一定影响，但我们将看到，弯曲族仍然具有满族的许多性质. 特别地，\autoref{thm:3.4.2}适用于弯曲指数族. 此外，满指数族与弯曲指数族还具有其他统计性质，本书余下部分将陆续讨论. 例如，假设我们拥有来自一个总体的大量数据值，其pdf或pmf具有\eqref{eq:3.4.1}的形式. 那么，只须从数据中算出$k$个数（$k$为\eqref{eq:3.4.1}中求和项的个数），就可概括数据中关于$\theta$的全部信息. 这种“数据压缩”性质将在\autoref{sec:6}讨论充分统计量时更详细地处理（\autoref{thm:6.2.10}）.

关于指数族的进一步入门介绍，见\hyperref[bib:source198]{Lehmann（1986}，第2.7节），或\hyperref[bib:source200]{Lehmann和Casella（1998}，第1.5节及注1.10.6）. \hyperref[bib:source50]{Brown（1986）}的经典专著提供了更深入且全面的介绍.



\subsection{位置族与尺度族}\label{sec:3.5}\index{SI0647@位置尺度族}
在\autoref{sec:3.3}和\autoref{sec:3.4}中，我们讨论了若干常见连续分布族. 本节讨论构造分布族的三种方法. 所得到的族具有易于理解的实际解释，因而适合建模，同时也有方便的数学性质.

这三类族称为位置族、尺度族和位置—尺度族. 每一类族的构造，都是先指定一个pdf，例如$f(x)$，称为该族的标准pdf（standard pdf）. 然后按指定方式变换这个标准pdf，生成族中的所有其他pdf. 我们从关于pdf的一个简单定理开始.\index{SI0651@标准概率密度函数}\semanticindexanchor{651:126}
\sourcepageanchor{127}

\begin{theorem}\label{thm:3.5.1}
设$f(x)$为任意pdf，$\mu$与$\sigma>0$为任意给定常数. 则函数
$$g(x\mid\mu,\sigma)=\frac1\sigma f\left(\frac{x-\mu}\sigma\right)$$
也是pdf.
\end{theorem}
\begin{proof}
为验证这个变换产生了合法的pdf，需要检查：对代入公式的每个$\mu$与$\sigma$值，$(1/\sigma)f((x-\mu)/\sigma)$作为$x$的函数都是pdf. 即须检查它非负且积分为1. 因为$f(x)$是pdf，所以对所有$x$都有$f(x)\ge0$. 于是对所有$x,\mu,\sigma$都有$(1/\sigma)f((x-\mu)/\sigma)\ge0$. 再注意
\begin{align*}
\int_{-\infty}^\infty\frac1\sigma f\left(\frac{x-\mu}\sigma\right)\,dx
&=\int_{-\infty}^\infty f(y)\,dy\qquad\left(\text{令 }y=\frac{x-\mu}\sigma\right)\\
&=1\qquad(\text{因为 }f(y)\text{是pdf}),
\end{align*}
即得所需结论.
\end{proof}
现在转向第一种构造，即位置族的构造.

\begin{definition}\label{def:3.5.2}\index{SI0648@位置参数}
设$f(x)$为任意pdf. 则由参数$\mu$，$-\infty<\mu<\infty$标识的pdf族$f(x-\mu)$，称为以$f(x)$为标准pdf的位置族（location family）；$\mu$称为该族的位置参数（location parameter）.
\end{definition}
为了解引入位置参数$\mu$的效果，考虑\autoref{fig:3.5.1}. 在$x=\mu$处，$f(x-\mu)=f(0)$；在$x=\mu+1$处，$f(x-\mu)=f(1)$；一般地，在$x=\mu+a$处，$f(x-\mu)=f(a)$. 当然，$\mu=0$时$f(x-\mu)$就是$f(x)$. 因此，位置参数$\mu$只把pdf $f(x)$平移，图形形状保持不变，但$f(x)$中原来在$x=0$上方的图上点，在$f(x-\mu)$中位于$x=\mu$上方. 从\autoref{fig:3.5.1}可清楚看出，$f(x)$图形下$x=-1$与$x=2$之间的面积，等于$f(x-\mu)$图形下$x=\mu-1$与$x=\mu+2$之间的面积. 因此，若$X$是pdf为$f(x-\mu)$的随机变量，就可写成
$$P(-1\le X\le2\mid0)=P(\mu-1\le X\le\mu+2\mid\mu),$$
其中等式左边随机变量$X$的pdf为$f(x-0)=f(x)$，右边的pdf为$f(x-\mu)$.

\autoref{sec:3.2}中介绍的若干族，本身就是位置族，或包含位置族作为子族.\footnote{译注：原文此处指向\hyperref[sec:3.2]{第3.2节}，但随后列出的正态、Cauchy和双指数分布均在\autoref{sec:3.3}介绍；原交叉引用疑应为\hyperref[sec:3.3]{第3.3节}.} 例如，若$\sigma>0$是一个指定的已知数，定义
$$f(x)=\frac1{\sqrt{2\pi}\sigma}e^{-x^2/(2\sigma^2)},\qquad-\infty<x<\infty,$$
\begin{figure}[htbp]
\centering\includegraphics[width=.72\textwidth]{assets/chapter3/fig3-5-1.pdf}
\caption{同一位置族的两个成员：均值为0和2}\label{fig:3.5.1}
\end{figure}
\sourcepageanchor{128}

则以$f(x)$为标准pdf的位置族，就是均值$\mu$未知而方差$\sigma^2$已知的正态分布集合. 为看出这一点，只须验证在上式中以$x-\mu$替换$x$，得到的pdf具有\eqref{eq:3.3.13}定义的形式. 类似地，Cauchy族，以及$\sigma$为指定值、$\mu$为参数的双指数族，都是位置族的例子. 不过，\autoref{def:3.5.2}的关键在于，可以从任意pdf $f(x)$出发，引入位置参数来生成一个pdf族.

若$X$是pdf为$f(x-\mu)$的随机变量，则可表示为$X=Z+\mu$，其中$Z$是pdf为$f(z)$的随机变量. 这个表示是稍后将证明的\autoref{thm:3.5.6}（取$\sigma=1$）的结果. 考察这个表示，可看出何时位置族适合用来对观测变量$X$建模. 下面描述两种情形.

首先，假设一个实验旨在测量某个物理常数$\mu$，例如一种溶液的温度. 但观测中存在测量误差，所以实际观测值$X$为$Z+\mu$，其中$Z$是测量误差. 对这次观测，若$Z>0$，则$X$大于$\mu$；若$Z<0$，则$X$小于$\mu$. 从过去用这个测量装置测量其他溶液的经验中，可能已很清楚随机测量误差的分布. 若该分布的pdf为$f(z)$，则观测值$X$的pdf就是$f(x-\mu)$.

再举一个例子：假设由以往实验已知驾驶员在协调性测试中的反应时间分布. 用随机变量$Z$表示随机选取一位驾驶员的反应时间. 设描述已知分布的$Z$的pdf为$f(z)$. 现在考虑对这个总体“施加处理”. 例如，假设每个人都喝三杯啤酒，会发生什么？可以假定每个人的反应时间都会改变某个未知量$\mu$. （这个非常简单的模型假定每个人的反应时间都改变相同的量$\mu$，可能不是最好的模型. 例如，已知酒精的作用与体重有关，所以较重的人受到啤酒的影响可能较小.）作为思想开放的科学工作者，我们甚至可能允许$\mu<0$，也就是反应时间缩短. 那么，如果我们观察“处理”后一位随机选取的驾驶员的反应时间，这个时间就是$X=Z+\mu$，而$X$的可能分布族由$f(x-\mu)$给出.

若使$f(x)>0$的$x$集合不是整个实数轴，那么使$f(x-\mu)>0$的$x$集合就会依赖于$\mu$. \autoref{ex:3.5.3}说明这一点.

\begin{example}[指数位置族]\label{ex:3.5.3}\index{SI1034@阈值参数}
设$f(x)=e^{-x}$，$x\ge0$；$f(x)=0$，$x<0$. 为构造位置族，以$x-\mu$替换$x$，得到
\begin{align*}
f(x\mid\mu)&=\begin{cases}e^{-(x-\mu)},&x-\mu\ge0,\\0,&x-\mu<0,\end{cases}\\
&=\begin{cases}e^{-(x-\mu)},&x\ge\mu,\\0,&x<\mu.\end{cases}
\end{align*}
不同$\mu$值对应的$f(x\mid\mu)$图形见\autoref{fig:3.5.2}. 与\autoref{fig:3.5.1}一样，图形被平移了. 此时图形的正值部分从$\mu$而非零开始. 若$X$衡量时间，则可能限制$\mu$非负，使每个$\mu$值下$X$都以概率1为正. 在这种模型中，$\mu$表示$X$取值范围的一个界，因而有时称为阈值参数（threshold parameter）.
\end{example}
本节还要讨论的另外两类族是尺度族和位置—尺度族.

\begin{definition}\label{def:3.5.4}\index{SI0650@尺度参数}
设$f(x)$为任意pdf. 对任意$\sigma>0$，由参数$\sigma$标识的pdf族$(1/\sigma)f(x/\sigma)$，称为以$f(x)$为标准pdf的尺度族（scale family），$\sigma$称为该族的尺度参数.
\end{definition}
引入尺度参数$\sigma$的效果是拉伸（$\sigma>1$）或压缩（$\sigma<1$）$f(x)$的图形，同时保持图形的基本形状不变. \autoref{fig:3.5.3}说明了这一点.
\sourcepageanchor{129}

\begin{figure}[htbp]
\centering\includegraphics[width=.72\textwidth]{assets/chapter3/fig3-5-2.pdf}
\caption{指数位置密度}\label{fig:3.5.2}
\end{figure}
\begin{figure}[htbp]
\centering\includegraphics[width=.72\textwidth]{assets/chapter3/fig3-5-3.pdf}
\caption{同一尺度族的成员}\label{fig:3.5.3}
\end{figure}
使用尺度参数时，$f(x)$通常或者关于0对称，或者仅在$x>0$时为正. 此时拉伸或者关于0对称，或者仅沿正方向进行. 但在定义中，任意pdf都可作为标准pdf.

\autoref{sec:3.3}介绍的若干族，本身是尺度族，或含有尺度族作为子族. 这些包括：$\alpha$固定而$\beta$为尺度参数的Gamma族；$\mu=0$而$\sigma$为尺度参数的正态族；指数族；以及$\mu=0$而$\sigma$为尺度参数的双指数族. 在每种情形中，令尺度参数等于1，就得到标准pdf. 然后可证明该族的所有其他成员都具有\autoref{def:3.5.4}中的形式.

\begin{definition}\label{def:3.5.5}\index{SI0648@位置参数}\index{SI0651@标准概率密度函数}
设$f(x)$为任意pdf. 对任意$\mu$，$-\infty<\mu<\infty$，及任意$\sigma>0$，由参数$(\mu,\sigma)$标识的pdf族$(1/\sigma)f((x-\mu)/\sigma)$，称为以$f(x)$为标准pdf的位置—尺度族；$\mu$称为位置参数，$\sigma$称为尺度参数.
\end{definition}
同时引入位置参数和尺度参数的效果是：先像尺度参数那样拉伸（$\sigma>1$）或压缩（$\sigma<1$）图形，再平移图形，使原先在0上方的点现在位于$\mu$上方. \autoref{fig:3.5.4}展示了$f(x)$的这种变换. 正态族和双指数族都是位置—尺度族的例子. \autoref{exercise:3.39}把Cauchy族作为位置—尺度族来讨论.

下面的定理把用于定义位置—尺度族的pdf $f(x)$变换，与pdf为$f(z)$的随机变量$Z$的变换联系起来.
\sourcepageanchor{130}

\begin{figure}[htbp]
\centering\includegraphics[width=.72\textwidth]{assets/chapter3/fig3-5-4.pdf}
\caption{同一位置—尺度族的成员}\label{fig:3.5.4}
\end{figure}
如前面讨论位置族时所述，用$Z$表示是一种有用的数学工具，也能帮助理解何时位置—尺度族适合建模. 在\autoref{thm:3.5.6}中令$\sigma=1$，得到（仅含）位置族的结果；令$\mu=0$，得到（仅含）尺度族的结果.

\begin{theorem}\label{thm:3.5.6}
设$f(\cdot)$为任意pdf，$\mu$为任意实数，$\sigma$为任意正实数. 则$X$是pdf为$(1/\sigma)f((x-\mu)/\sigma)$的随机变量，当且仅当存在pdf为$f(z)$的随机变量$Z$，使$X=\sigma Z+\mu$.
\end{theorem}
\begin{proof}
为证明“如果”部分，定义$g(z)=\sigma z+\mu$. 则$X=g(Z)$，$g$为单调函数，$g^{-1}(x)=(x-\mu)/\sigma$，且$|(d/dx)g^{-1}(x)|=1/\sigma$. 因此，由\autoref{thm:2.1.5}，$X$的pdf为
$$f_X(x)=f_Z(g^{-1}(x))\left|\frac d{dx}g^{-1}(x)\right|=f\left(\frac{x-\mu}\sigma\right)\frac1\sigma.$$
为证明“仅当”部分，定义$g(x)=(x-\mu)/\sigma$，令$Z=g(X)$. 再次应用\autoref{thm:2.1.5}，有$g^{-1}(z)=\sigma z+\mu$，$|(d/dz)g^{-1}(z)|=\sigma$，且$Z$的pdf为
$$f_Z(z)=f_X(g^{-1}(z))\left|\frac d{dz}g^{-1}(z)\right|=\frac1\sigma f\left(\frac{(\sigma z+\mu)-\mu}\sigma\right)\sigma=f(z).$$
此外，
$$\sigma Z+\mu=\sigma g(X)+\mu=\sigma\left(\frac{X-\mu}\sigma\right)+\mu=X.$$
\end{proof}
从\autoref{thm:3.5.6}可以提取出一个重要事实：随机变量$Z=(X-\mu)/\sigma$具有pdf
\footnote{译注：原式最右端写作$f(x)$，但由左边的自变量及本段定义，应为$f(z)$. 这里保留原式并注明.}
$$f_Z(z)=\frac11f\left(\frac{z-0}1\right)=f(x).$$

也就是说，$Z$的分布是位置—尺度族中对应于$\mu=0,\sigma=1$的成员. 对正态族这一特殊情形，我们已在\autoref{sec:3.3}证明过这一点.

通常可先对pdf为$f(z)$的“标准”随机变量$Z$进行计算，再容易地导出pdf为$(1/\sigma)f((x-\mu)/\sigma)$的随机变量$X$的相应结果. 下面的定理给出一个例子，它推广了\autoref{sec:3.3}中对正态族所作的计算.
\sourcepageanchor{131}

\begin{theorem}\label{thm:3.5.7}
设$Z$是pdf为$f(z)$的随机变量，假定$EZ$与$\operatorname{Var}Z$存在. 若$X$是pdf为$(1/\sigma)f((x-\mu)/\sigma)$的随机变量，则
$$EX=\sigma EZ+\mu\qquad\text{及}\qquad\operatorname{Var}X=\sigma^2\operatorname{Var}Z.$$
特别地，若$EZ=0$且$\operatorname{Var}Z=1$，则$EX=\mu$且$\operatorname{Var}X=\sigma^2$.
\end{theorem}
\begin{proof}
由\autoref{thm:3.5.6}，存在pdf为$f(z)$的随机变量$Z^*$，使$X=\sigma Z^*+\mu$. 所以$EX=\sigma EZ^*+\mu=\sigma EZ+\mu$，且$\operatorname{Var}X=\sigma^2\operatorname{Var}Z^*=\sigma^2\operatorname{Var}Z$.
\end{proof}
对任意均值和方差有限的位置—尺度族，可选择标准pdf $f(z)$，使$EZ=0$且$\operatorname{Var}Z=1$. （证明可以这样选择，留作\autoref{exercise:3.40}.）于是$\mu$和$\sigma^2$就分别可方便地解释为$X$的均值和方差. \autoref{sec:3.3}给出的通常正态族定义就是这种情形. 然而，\autoref{sec:3.3}给出的通常双指数族定义并未作这种选择；那里$\operatorname{Var}Z=2$.

位置—尺度族任一成员的概率都可以用标准变量$Z$计算，因为
$$P(X\le x)=P\left(\frac{X-\mu}\sigma\le\frac{x-\mu}\sigma\right)=P\left(Z\le\frac{x-\mu}\sigma\right).$$
因此，如果标准变量$Z$的$P(Z\le z)$有数表可查或容易计算，就可得到$X$的概率. 用标准正态表计算正态概率就是这方面的例子.



\subsection{不等式与恒等式}\label{sec:3.6}
统计理论中确实充满了不等式与恒等式，多到有整本著作专门讨论这个主题. \hyperref[bib:source214]{Marshall和Olkin（1979）}的重要著作包含许多使用主序（majorization）概念的不等式. 较早的\hyperref[bib:source158]{Hardy、Littlewood和Polya（1952）}的著作则是经典不等式的汇编. 本节及\autoref{sec:4.7}将混合讨论一些新旧结果，让读者了解这些结果的类型. 本节专门讨论源于概率问题的恒等式与不等式，而\autoref{sec:4.7}中的结果更多依赖数与函数的基本性质.

\subsubsection{概率不等式}\label{sec:3.6.1}
最著名且可能最有用的概率不等式是Chebyshev不等式. 它的用处来自广泛的适用性. 与许多重要结果一样，它的证明几乎是平凡的.

\begin{theorem}[Chebyshev不等式]\label{thm:3.6.1}\index{SI0565@切比雪夫}
设$X$为随机变量，$g(x)$为非负函数. 则对任意$r>0$，
$$P(g(X)\ge r)\le\frac{Eg(X)}r.$$
\end{theorem}
\begin{proof}
\begin{align*}
Eg(X)&=\int_{-\infty}^\infty g(x)f_X(x)\,dx\\
&\ge\int_{\{x:g(x)\ge r\}}g(x)f_X(x)\,dx\qquad(g\text{非负})\\
&\ge r\int_{\{x:g(x)\ge r\}}f_X(x)\,dx\\
&=rP(g(X)\ge r).\qquad(\text{定义})
\end{align*}
重新整理即得所需不等式.
\end{proof}
\sourcepageanchor{132}

\begin{example}[说明Chebyshev不等式]\label{ex:3.6.2}
Chebyshev不等式最广泛的用途涉及均值和方差. 令$g(x)=(x-\mu)^2/\sigma^2$，其中$\mu=EX$，$\sigma^2=\operatorname{Var}X$. 为方便起见，写$r=t^2$. 则
$$P\left(\frac{(X-\mu)^2}{\sigma^2}\ge t^2\right)\le\frac1{t^2}E\frac{(X-\mu)^2}{\sigma^2}=\frac1{t^2}.$$
作一些显然的代数运算，得到不等式
$$P(|X-\mu|\ge t\sigma)\le\frac1{t^2}$$
及其对应形式
$$P(|X-\mu|<t\sigma)\ge1-\frac1{t^2},$$
这给出了用$\sigma$表示的偏差$|X-\mu|$的通用界. 例如，取$t=2$，得到
$$P(|X-\mu|\ge2\sigma)\le\frac1{2^2}=.25,$$
所以，一个随机变量位于其均值的$2\sigma$范围以内的机会至少有75\%（无论$X$的分布如何）.
\end{example}
Chebyshev不等式虽然适用范围广，但必然是保守的. （例如见\autoref{exercise:3.46}和\hyperref[sec:3.8.2]{杂录3.8.2}.）特别地，对某些特定分布，常常可以得到更紧的界.

\begin{example}[一个正态概率不等式]\label{ex:3.6.3}\index{SI0576@正态}
若$Z$服从标准正态分布，则
\begin{equation}\label{eq:3.6.1}
P(|Z|\ge t)\le\sqrt{\frac2\pi}\frac{e^{-t^2/2}}t,\qquad\text{对所有 }t>0.
\end{equation}
将它与Chebyshev不等式比较. 当$t=2$时，Chebyshev给出$P(|Z|\ge t)\le.25$，而$\sqrt{2/\pi}e^{-2}/2=.054$，改进非常大.

为证明\eqref{eq:3.6.1}，写成
\begin{align*}
P(Z\ge t)&=\frac1{\sqrt{2\pi}}\int_t^\infty e^{-x^2/2}\,dx\\
&\le\frac1{\sqrt{2\pi}}\int_t^\infty\frac xt e^{-x^2/2}\,dx\qquad(\text{当 }x>t\text{ 时，}x/t>1)\\
&=\frac1{\sqrt{2\pi}}\frac{e^{-t^2/2}}t,
\end{align*}
再用$P(|Z|\ge t)=2P(Z\ge t)$即可. 可用类似方法建立$P(|Z|\ge t)$的下界\mbox{（见\autoref{exercise:3.47}）.}
\end{example}
还有许多其他概率不等式，几乎所有这些不等式的思想都与Chebyshev不等式相似. 例如，将看到（\autoref{exercise:3.45}）
$$P(X\ge a)\le e^{-at}M_X(t),$$
但这个不等式当然要求矩母函数存在.\footnote{译注：这一上尾界取$t>0$；该条件在\autoref{exercise:3.45}中明确给出.} 还有一些比Chebyshev不等式更紧、但需要更多假设的不等式（详见\hyperref[sec:3.8.2]{杂录3.8.2}）.
\sourcepageanchor{133}

\subsubsection{恒等式}\label{sec:3.6.2}
本节介绍若干恒等式，它们不仅有助于建立定理，也可简化数值计算. 有一整类恒等式可视为“递推关系”，其中一些我们已见过. 回顾，若$X$服从Poisson$(\lambda)$，则
\begin{equation}\label{eq:3.6.2}\index{SI0375@递推关系}
P(X=x+1)=\frac\lambda{x+1}P(X=x),
\end{equation}
于是可以从$P(X=0)=e^{-\lambda}$出发，递推计算Poisson概率. 几乎所有离散分布都有类似\eqref{eq:3.6.2}的关系（见\autoref{exercise:3.48}）. 有时，对连续分布也存在形式略有不同的关系.

\begin{theorem}\label{thm:3.6.4}\index{SI0310@递推关系}
设$X_{\alpha,\beta}$为Gamma$(\alpha,\beta)$随机变量，pdf为$f(x\mid\alpha,\beta)$，其中$\alpha>1$. 则对任意常数$a$和$b$，
\begin{equation}\label{eq:3.6.3}
P(a<X_{\alpha,\beta}<b)=\beta\bigl(f(a\mid\alpha,\beta)-f(b\mid\alpha,\beta)\bigr)+P(a<X_{\alpha-1,\beta}<b).
\end{equation}
\end{theorem}
\begin{proof}
由定义，
\begin{align*}
P(a<X_{\alpha,\beta}<b)&=\frac1{\Gamma(\alpha)\beta^\alpha}\int_a^b x^{\alpha-1}e^{-x/\beta}\,dx\\
&=\frac1{\Gamma(\alpha)\beta^\alpha}\left[\left.-x^{\alpha-1}\beta e^{-x/\beta}\right|_a^b+\int_a^b(\alpha-1)x^{\alpha-2}\beta e^{-x/\beta}\,dx\right],
\end{align*}
这里作了分部积分，其中$u=x^{\alpha-1}$，$dv=e^{-x/\beta}\,dx$. 继续计算，得到
$$P(a<X_{\alpha,\beta}<b)=\beta\bigl(f(a\mid\alpha,\beta)-f(b\mid\alpha,\beta)\bigr)+\frac{\alpha-1}{\Gamma(\alpha)\beta^{\alpha-1}}\int_a^b x^{\alpha-2}e^{-x/\beta}\,dx.$$
利用$\Gamma(\alpha)=(\alpha-1)\Gamma(\alpha-1)$，可见最后一项为$P(a<X_{\alpha-1,\beta}<b)$.
\end{proof}
若$\alpha$为整数，反复使用\eqref{eq:3.6.3}，最终会得到可以解析计算的积分（当$\alpha=1$时，即指数分布）. 因而可以方便地计算这些Gamma概率.

有一整类恒等式依赖于分部积分. 其中第一个归功于Charles Stein，他在研究多元正态均值的估计时使用了这一恒等式（\hyperref[bib:source303]{Stein 1973}、\hyperref[bib:source304]{1981}）.

\begin{lemma}[Stein引理]\label{lem:3.6.5}\index{SI0541@分部积分}\index{SI0984@Stein引理}
设$X\sim\operatorname n(\theta,\sigma^2)$，$g$为满足$E|g'(X)|<\infty$的可微函数. 则
$$E[g(X)(X-\theta)]=\sigma^2Eg'(X).$$
\end{lemma}
\begin{proof}
左边为
$$E[g(X)(X-\theta)]=\frac1{\sqrt{2\pi}\sigma}\int_{-\infty}^\infty g(x)(x-\theta)e^{-(x-\theta)^2/(2\sigma^2)}\,dx.$$
令$u=g(x)$、$dv=(x-\theta)e^{-(x-\theta)^2/(2\sigma^2)}\,dx$，分部积分得到
$$E[g(X)(X-\theta)]=\frac1{\sqrt{2\pi}\sigma}\left[\left.-\sigma^2g(x)e^{-(x-\theta)^2/(2\sigma^2)}\right|_{-\infty}^\infty+\sigma^2\int_{-\infty}^\infty g'(x)e^{-(x-\theta)^2/(2\sigma^2)}\,dx\right].$$
对$g'$的条件足以保证第一项为0，右边剩下的就是$\sigma^2Eg'(X)$.
\end{proof}
\sourcepageanchor{134}

\begin{example}[高阶正态矩]\label{ex:3.6.6}
Stein引理使高阶矩的计算十分容易. 例如，若$X\sim\operatorname n(\theta,\sigma^2)$，则
\begin{align*}
EX^3&=EX^2(X-\theta+\theta)\\
&=EX^2(X-\theta)+\theta EX^2\\
&=2\sigma^2EX+\theta EX^2\qquad(g(x)=x^2,\ g'(x)=2x)\\
&=2\sigma^2\theta+\theta(\sigma^2+\theta^2)\\
&=3\theta\sigma^2+\theta^3.
\end{align*}
\end{example}
许多分布都有类似的分部积分恒等式（见\autoref{exercise:3.49}及\hyperref[bib:source172]{Hudson（1978）}）. 也可以利用特定分布的性质得到有用的恒等式，如下一个定理所示.

\begin{theorem}\label{thm:3.6.7}\index{SI0532@卡方}
设$\chi_p^2$表示具有$p$个自由度的卡方随机变量. 对任意函数$h(x)$，只要期望存在，就有
\begin{equation}\label{eq:3.6.4}
Eh(\chi_p^2)=pE\left(\frac{h(\chi_{p+2}^2)}{\chi_{p+2}^2}\right).
\end{equation}
\end{theorem}
\begin{proof}
“只要期望存在”这个说法，是避免详细列出$h$所需条件的一种省事做法. 一般而言，合理的函数都会满足\eqref{eq:3.6.4}. 我们有
\begin{align*}
Eh(\chi_p^2)&=\frac1{\Gamma(p/2)2^{p/2}}\int_0^\infty h(x)x^{(p/2)-1}e^{-x/2}\,dx\\
&=\frac1{\Gamma(p/2)2^{p/2}}\int_0^\infty\left(\frac{h(x)}x\right)x^{((p+2)/2)-1}e^{-x/2}\,dx,
\end{align*}
这里把被积函数乘以$x/x$. 现在写成
$$\Gamma\left(\frac p2\right)2^{p/2}=\frac{\Gamma((p+2)/2)2^{(p+2)/2}}p,$$
所以
\begin{align*}
Eh(\chi_p^2)&=\frac p{\Gamma((p+2)/2)2^{(p+2)/2}}\int_0^\infty\left(\frac{h(x)}x\right)x^{((p+2)/2)-1}e^{-x/2}\,dx\\
&=pE\left(\frac{h(\chi_{p+2}^2)}{\chi_{p+2}^2}\right).
\end{align*}
\end{proof}
利用\eqref{eq:3.6.4}，一些矩的计算非常容易. 例如，$\chi_p^2$的均值为
$$E\chi_p^2=pE\left(\frac{\chi_{p+2}^2}{\chi_{p+2}^2}\right)=pE(1)=p,$$
二阶矩为
\begin{align*}
E(\chi_p^2)^2&=pE\left(\frac{(\chi_{p+2}^2)^2}{\chi_{p+2}^2}\right)=pE(\chi_{p+2}^2)\\
&=p(p+2).
\end{align*}
所以$\operatorname{Var}\chi_p^2=p(p+2)-p^2=2p$.

最后，用前述恒等式的几个离散对应结果结束本节. \autoref{thm:3.6.8}中两个恒等式的一般版本归功于\hyperref[bib:source174]{Hwang（1982）}.
\sourcepageanchor{135}

\begin{theorem}[Hwang]\label{thm:3.6.8}\index{SI0473@Hwang引理}
设函数$g(x)$满足$-\infty<Eg(X)<\infty$且$-\infty<g(-1)<\infty$. 则
\begin{enumerate}[label=\alph*.]
\item\phantomsection\label{itm:thm:3.6.8:a} 若$X\sim\operatorname{Poisson}(\lambda)$，
\begin{equation}\label{eq:3.6.5}\index{SI0544@泊松}
E(\lambda g(X))=E(Xg(X-1)).
\end{equation}
\item\phantomsection\label{itm:thm:3.6.8:b} 若$X\sim\text{负二项}(r,p)$，\footnote{译注：当$r=1$且$X=0$时，\eqref{eq:3.6.6}右边分式为$0/0$，而该事件有正概率. 为使恒等式包括几何分布这一情形，应将该点的整个求和项约定为0；或者把此式先限制于$r>1$. 原文未注明这一边界约定.}
\begin{equation}\label{eq:3.6.6}\index{SI0543@负二项}
E((1-p)g(X))=E\left(\frac X{r+X-1}g(X-1)\right).
\end{equation}
\end{enumerate}
\end{theorem}

\begin{proof}
这里只证明\hyperref[itm:thm:3.6.8:a]{(a)}，把\hyperref[itm:thm:3.6.8:b]{(b)}留给\autoref{exercise:3.50}. 有
\begin{align*}
E(\lambda g(X))&=\sum_{x=0}^\infty\lambda g(x)\frac{e^{-\lambda}\lambda^x}{x!}\\
&=\sum_{x=0}^\infty g(x)\frac{e^{-\lambda}\lambda^{x+1}}{x!}\frac{x+1}{x+1}\\
&=\sum_{x=0}^\infty(x+1)g(x)\frac{e^{-\lambda}\lambda^{x+1}}{(x+1)!}.
\end{align*}
现在变换求和指标，令$y=x+1$. 当$x$从0变到$\infty$时，$y$从1变到$\infty$. 因此，
\begin{align*}
E(\lambda g(X))&=\sum_{y=1}^\infty yg(y-1)\frac{e^{-\lambda}\lambda^y}{y!}\\
&=\sum_{y=0}^\infty yg(y-1)\frac{e^{-\lambda}\lambda^y}{y!}\qquad(\text{增加的项为0})\\
&=E(Xg(X-1)),
\end{align*}
因为最后这个和是一个Poisson$(\lambda)$期望.
\end{proof}
\hyperref[bib:source174]{Hwang（1982）}以类似Stein的方式使用这个恒等式，证明了有关多元估计量的结果. 这个恒等式还有其他用途，特别是在矩的计算方面.

\begin{example}[高阶Poisson矩]\label{ex:3.6.9}
对$X\sim\operatorname{Poisson}(\lambda)$，取$g(x)=x^2$，利用\eqref{eq:3.6.5}：
\begin{align*}
E(\lambda X^2)&=E(X(X-1)^2)\\
&=E(X^3-2X^2+X).
\end{align*}
因此，Poisson$(\lambda)$的三阶矩为
\begin{align*}
EX^3&=\lambda EX^2+2EX^2-EX\\
&=\lambda(\lambda+\lambda^2)+2(\lambda+\lambda^2)-\lambda\\
&=\lambda^3+3\lambda^2+\lambda.
\end{align*}
对于负二项分布，在\eqref{eq:3.6.6}中取$g(x)=r+x$，就可计算均值：
$$E((1-p)(r+X))=E\left(\frac X{r+X-1}(r+X-1)\right)=EX,$$
所以整理得到
$$(EX)((1-p)-1)=-r(1-p),$$
即
$$EX=\frac{r(1-p)}p.$$
其他矩可类似计算.
\end{example}



\sourcepageanchor{136}
\subsection{习题}\label{sec:3.7}
\setcounter{exercise}{0}
\begin{exercise}\label{exercise:3.1}
设随机变量$X$具有一般离散均匀$(N_0,N_1)$分布，它给$N_0,N_0+1,\ldots,N_1$各个值赋予相同的概率. 这里$N_0\le N_1$且二者均为整数. 求$EX$与$\operatorname{Var}X$的表达式.
\end{exercise}
\begin{exercise}\label{exercise:3.2}
一个制造商收到供应商交来的一批100个零件. 若其中不合格零件超过五个，这批货就是不可接受的. 制造商将从中随机选择$K$个零件检查，如果样本中没有发现不合格零件，就接受这批货.
\begin{enumerate}
\item[(a)] $K$至少要多大，才能保证制造商接受不可接受批次的概率小于$.10$？
\item[(b)] 假设制造商决定，只要样本中至多有一个不合格零件，就接受这批货. $K$至少要多大，才能保证制造商接受不可接受批次的概率小于$.10$？
\end{enumerate}
\end{exercise}
\begin{exercise}\label{exercise:3.3}
某些街角的交通流量，有时可建模为一列Bernoulli试验：假定任意给定一秒内有车通过的概率为常数$p$，且不同秒内的车辆通过情况互不影响. 把每秒视为不可分的时间单位（试验），就可应用Bernoulli模型. 假设只有未来三秒内没有车通过时，行人才能横穿街道. 求该行人在开始横穿前恰好需要等待四秒的概率.
\end{exercise}
\begin{exercise}\label{exercise:3.4}
一个有$n$把钥匙的人想开门，于是随机试钥匙. 恰有一把钥匙能打开门. 在下列情况下求试验次数的均值：
\begin{enumerate}
\item[(a)] 试过而未成功的钥匙仍可在以后再次选取.
\item[(b)] 试过而未成功的钥匙被排除.
\end{enumerate}
\end{exercise}
\begin{exercise}\label{exercise:3.5}
已知一种标准药物在使用它的病例中有80\%有效. 在100名患者中试验一种新药，发现85例有效. 新药是否更好？（提示：假定新药和旧药同样有效，计算观察到85次或更多成功的概率.）
\end{exercise}
\begin{exercise}\label{exercise:3.6}
预计某个品种的玫瑰会吸引大量昆虫. 一种市售杀虫剂宣称有效率为99\%. 假设2000只昆虫侵入一个已喷洒该杀虫剂的玫瑰园，令$X=$存活昆虫数.
\begin{enumerate}
\item[(a)]\phantomsection\label{itm:exercise:3.6:a} 哪一种概率分布可能为这个实验提供合理模型？
\item[(b)]\phantomsection\label{itm:exercise:3.6:b} 用\hyperref[itm:exercise:3.6:a]{(a)}中的模型，写出存活昆虫少于100只的概率表达式，但不要计算.
\item[(c)] 计算\hyperref[itm:exercise:3.6:b]{(b)}中概率的一个近似值.
\end{enumerate}
\end{exercise}
\begin{exercise}\label{exercise:3.7}
设某种曲奇中巧克力碎片的数量具有Poisson分布. 我们希望随机选取的一块曲奇含有至少两片巧克力的概率大于$.99$. 求能保证这一概率的分布均值的最小值.\footnote{译注：原文要求概率严格“大于”$.99$，在均值为连续参数时满足条件的集合是一个开射线，因而严格来说只有下确界而无最小值；若改为“不小于”$.99$，则存在最小值.}
\end{exercise}
\begin{exercise}\label{exercise:3.8}
两家电影院争夺1000名顾客的生意. 假设每个顾客独立且“无偏好”地选择其中一家. 令$N$表示每家影院的座位数.
\begin{enumerate}
\item[(a)] 使用二项模型，求$N$应满足的一个表达式，以保证因满座而拒绝顾客的概率小于1\%.
\item[(b)] 用正态近似求$N$的数值.
\end{enumerate}
\end{exercise}
\begin{exercise}\label{exercise:3.9}\index{SI0109@巧合}
新闻经常把一些巧合报道成惊人的“百万分之一”事件，但仔细考察后，它们实际上并不罕见，甚至是可以预期发生的.

几年前，纽约州一所小学报告，其新入学的幼儿园儿童中有五对双胞胎. 这件事当然传遍全州，报道还引用校长的话，称这是“统计上不可能的事”.

果真如此吗？还是说，这只是\hyperref[bib:source100]{Diaconis和Mosteller（1989）}所称“真正大数定律”的一个例子？让我们作一些计算.
\sourcepageanchor{137}
\begin{enumerate}
\item[(a)]\phantomsection\label{itm:exercise:3.9:a} 双胞胎出生的概率约为$1/90$，可假定一所小学大约有60名儿童进入幼儿园（三个班，每班20人）. 解释为什么这个“统计上不可能”的事件可视为二项$(60,1/90)$中出现五次或更多成功的概率. 它是否罕见到足以成为新闻？
\item[(b)]\phantomsection\label{itm:exercise:3.9:b} 即使\hyperref[itm:exercise:3.9:a]{(a)}中的概率足够小，值得报道，也应考虑它可能发生在县内任何一所学校，以及州内任何一个县，而报道方式仍然完全一样. （“真正大数定律”开始发挥作用了.）纽约州有62个县，假定每个县有五所小学是合理的. 这个事件仍可称为“统计上不可能”，还是已经变成可以预期发生的事情？
\item[(c)] 如果\hyperref[itm:exercise:3.9:b]{(b)}中的概率仍然显得很小，再进一步考虑：这个事件可能发生在50个州中的任何一个、过去十年中的任何一年，仍会受到同样的新闻报道.
\end{enumerate}
除了\hyperref[bib:source100]{Diaconis和Mosteller（1989）}，还可参见\hyperref[bib:source157]{Hanley（1992）}对巧合的进一步讨论.
\end{exercise}
\begin{exercise}\label{exercise:3.10}
\hyperref[bib:source290]{Shuster（1991）}描述了他为一宗涉及贩卖可卡因的案件所作的一些概率计算. 佛罗里达州一个警察部门查获496包疑似可卡因，从中随机抽取四包检验，确认确实是可卡因. 随后警察又随机挑出两包，假扮毒贩把它们卖给被告. 最后这两包在接受检验、确认确为可卡因之前遗失了.
\begin{enumerate}
\item[(a)]\phantomsection\label{itm:exercise:3.10:a} 若原来的496包由$N$包可卡因和$M=496-N$包非可卡因组成，证明先选出四包可卡因、再选出两包非可卡因的概率，也就是被告没有购买可卡因而无罪的概率，为\footnote{译注：这个式子是“四包检验为可卡因且随后两包为非可卡因”的联合抽样概率. 在已经观察到前四包均为可卡因后，固定$N,M$时后两包均为非可卡因的条件概率是$\binom M2/\binom{N+M-4}2$. 因此，把原联合概率直接称为被告的“无罪概率”需谨慎；此处保留原题的表述.}
$$\frac{\binom N4\binom M2}{\binom{N+M}4\binom{N+M-4}2}.$$

\item[(b)] 对$M,N$最大化\hyperref[itm:exercise:3.10:a]{(a)}中的概率，就最大化了被告的“无罪概率”. 证明这个概率为$.022$，在$M=165$、$N=331$时达到.
\end{enumerate}
\end{exercise}
\begin{exercise}\label{exercise:3.11}
超几何分布可用二项分布或Poisson分布近似. （当然也可用其他分布近似，但本题只关注这两种.）设$X$具有超几何分布
$$P(X=x\mid N,M,K)=\frac{\binom Mx\binom{N-M}{K-x}}{\binom NK},\qquad x=0,1,\ldots,K.$$
\begin{enumerate}
\item[(a)] 证明，当$N\to\infty$、$M\to\infty$且$M/N\to p$时，\index{SI0061@以二项近似超几何}\semanticindexanchor{61:137}
$$P(X=x\mid N,M,K)\to\binom Kx p^x(1-p)^{K-x},\qquad x=0,1,\ldots,K.$$
（Stirling公式，见\autoref{exercise:1.23}，可能有用.）
\item[(b)]\phantomsection\label{itm:exercise:3.11:b}\index{SI0062@以泊松近似超几何} 利用二项分布可用Poisson分布近似这一事实，证明若$N\to\infty$、$M\to\infty$、$K\to\infty$、$M/N\to0$且$KM/N\to\lambda$，则
$$P(X=x\mid N,M,K)\to\frac{e^{-\lambda}\lambda^x}{x!},\qquad x=0,1,\ldots.$$
\item[(c)] 不用二项分布的Poisson近似，直接验证\hyperref[itm:exercise:3.11:b]{(b)}中的近似. （\autoref{lem:2.3.14}有帮助.）
\end{enumerate}
\end{exercise}
\begin{exercise}\label{exercise:3.12}\index{SI0270@与负二项的关系}\index{SI0353@与二项的关系}
设$X$具有二项$(n,p)$分布，$Y$具有负二项$(r,p)$分布. 证明$F_X(r-1)=1-F_Y(n-r)$.
\end{exercise}
\sourcepageanchor{138}

\begin{exercise}\label{exercise:3.13}\index{SI0385@截断}
截断离散分布（truncated discrete distribution）是指某个特定类别无法被观察到，并从样本空间中去除后得到的分布. 特别地，若$X$的取值范围为$0,1,2,\ldots$，且无法观察0这一类别（通常如此），则零截断随机变量$X_T$的pmf为
$$P(X_T=x)=\frac{P(X=x)}{P(X>0)},\qquad x=1,2,\ldots.$$
分别从下列分布出发，求零截断随机变量的pmf、均值和方差：
\begin{enumerate}
\item[(a)] $X\sim\operatorname{Poisson}(\lambda)$.
\item[(b)] $X\sim\text{负二项}(r,p)$，如\eqref{eq:3.2.10}.
\end{enumerate}
\end{exercise}
\begin{exercise}\label{exercise:3.14}\index{SI0327@对数级数}
从零截断负二项分布出发（参见\autoref{exercise:3.13}），令$r\to0$，就得到一个有趣的分布——对数级数分布（logarithmic series distribution）. 若随机变量$X$满足
$$P(X=x)=\frac{-(1-p)^x}{x\log p},\qquad x=1,2,\ldots,\quad0<p<1,$$
则称其具有参数为$p$的对数级数分布.
\begin{enumerate}
\item[(a)] 验证这定义了一个合法的概率函数.
\item[(b)] 求$X$的均值与方差. （对数级数分布已被证明在物种丰度建模中有用. 对这个分布更详细的讨论见\hyperref[bib:source309]{Stuart和Ord（1987）}.）
\end{enumerate}
\end{exercise}
\begin{exercise}\label{exercise:3.15}\index{SI0378@与负二项的关系}
\autoref{sec:3.2}曾指出，当$r\to\infty$、$p\to1$且$r(1-p)\to\lambda$时，Poisson$(\lambda)$分布是负二项$(r,p)$分布的极限. 证明在这些条件下，负二项分布的矩母函数收敛到Poisson分布的矩母函数.
\end{exercise}
\begin{exercise}\label{exercise:3.16}
验证正文中给出的关于Gamma函数的以下两个恒等式.
\begin{enumerate}
\item[(a)] $\Gamma(\alpha+1)=\alpha\Gamma(\alpha)$.
\item[(b)] $\Gamma(1/2)=\sqrt\pi$.
\end{enumerate}
\end{exercise}
\begin{exercise}\label{exercise:3.17}\index{SI0309@矩}
对Gamma分布建立一个类似\eqref{eq:3.3.18}的公式. 若$X\sim\operatorname{gamma}(\alpha,\beta)$，则对任意正常数$\nu$，
$$EX^\nu=\frac{\beta^\nu\Gamma(\nu+\alpha)}{\Gamma(\alpha)}.$$
\end{exercise}
\begin{exercise}\label{exercise:3.18}\index{SI0314@与负二项的关系}\index{SI0355@与伽马的关系}
负二项随机变量与Gamma随机变量之间有一个有趣的关系，有时可提供有用的近似. 设$Y$为参数$r,p$的负二项随机变量，其中$p$为成功概率. 证明当$p\to0$时，随机变量$pY$的矩母函数收敛到参数为$r$和1的Gamma分布的矩母函数.
\end{exercise}
\begin{exercise}\label{exercise:3.19}\index{SI0315@与泊松的关系}\index{SI0377@与伽马的关系}
证明
$$\int_x^\infty\frac1{\Gamma(\alpha)}z^{\alpha-1}e^{-z}\,dz=\sum_{y=0}^{\alpha-1}\frac{x^ye^{-x}}{y!},\qquad\alpha=1,2,3,\ldots.$$
（提示：用分部积分.）把这个公式表述为Poisson随机变量与Gamma随机变量之间的概率关系.
\end{exercise}
\begin{exercise}\label{exercise:3.20}\index{SI0305@折叠正态}
设随机变量$X$具有pdf
$$f(x)=\frac2{\sqrt{2\pi}}e^{-x^2/2},\qquad0<x<\infty.$$
\begin{enumerate}
\item[(a)] 求$X$的均值和方差. （这个分布有时称为折叠正态分布，folded normal.）
\item[(b)] 若$X$具有折叠正态分布，求变换$g(X)=Y$及$\alpha,\beta$的值，使$Y\sim\operatorname{gamma}(\alpha,\beta)$.
\end{enumerate}
\end{exercise}
\begin{exercise}\label{exercise:3.21}
写出定义下列pdf的矩母函数所需的积分：
$$f(x)=\frac1\pi\frac1{1+x^2}.$$
这个积分有限吗？（你预期它会有限吗？）
\end{exercise}
\sourcepageanchor{139}

\begin{exercise}\label{exercise:3.22}
对下列每个分布，验证正文中给出的$EX$与$\operatorname{Var}X$公式.
\begin{enumerate}
\item[(a)] 若$X$具有Poisson$(\lambda)$分布，验证$\operatorname{Var}X$. （提示：计算$EX(X-1)=EX^2-EX$.）
\item[(b)] 若$X$具有负二项$(r,p)$分布，验证$\operatorname{Var}X$.
\item[(c)] 若$X$具有Gamma$(\alpha,\beta)$分布，验证$\operatorname{Var}X$.
\item[(d)] 若$X$具有Beta$(\alpha,\beta)$分布，验证$EX$与$\operatorname{Var}X$.
\item[(e)] 若$X$具有双指数$(\mu,\sigma)$分布，验证$EX$与$\operatorname{Var}X$.
\end{enumerate}
\end{exercise}
\begin{exercise}\label{exercise:3.23}\index{SI0370@帕累托}
参数为$\alpha,\beta$的Pareto分布具有pdf
$$f(x)=\frac{\beta\alpha^\beta}{x^{\beta+1}},\qquad\alpha<x<\infty,\quad\alpha>0,\quad\beta>0.$$
\begin{enumerate}
\item[(a)] 验证$f(x)$是pdf.
\item[(b)] 推导这个分布的均值和方差.
\item[(c)] 证明若$\beta\le2$，则方差不存在.
\end{enumerate}
\end{exercise}
\begin{exercise}\label{exercise:3.24}
许多“有名”分布是前面讨论过的常见分布的特殊情形. 对下列每个有名分布，推导pdf的形式，验证它是pdf，并计算均值和方差.
\begin{enumerate}
\item[(a)] 若$X\sim\operatorname{exponential}(\beta)$，则$Y=X^{1/\gamma}$具有Weibull$(\gamma,\beta)$分布，其中$\gamma>0$为常数.\index{SI0391@威布尔}\semanticindexanchor{391:139}
\item[(b)] 若$X\sim\operatorname{exponential}(\beta)$，则$Y=(2X/\beta)^{1/2}$具有Rayleigh分布.\index{SI0380@瑞利}\semanticindexanchor{380:139}
\item[(c)] 若$X\sim\operatorname{gamma}(a,b)$，则$Y=1/X$具有逆Gamma分布（inverse gamma distribution）$IG(a,b)$. （这个分布在方差的Bayes估计中有用；见\autoref{exercise:7.23}.）
\item[(d)] 若$X\sim\operatorname{gamma}(3/2,\beta)$，则$Y=(X/\beta)^{1/2}$具有Maxwell分布.\index{SI0335@麦克斯韦}\semanticindexanchor{335:139}
\item[(e)] 若$X\sim\operatorname{exponential}(1)$，则$Y=\alpha-\gamma\log X$具有Gumbel$(\alpha,\gamma)$分布，其中$-\infty<\alpha<\infty$且$\gamma>0$. （Gumbel分布也称为极值分布，extreme value distribution.）\index{SI0300@极值}\semanticindexanchor{300:139}\index{SI0322@耿贝尔}\semanticindexanchor{322:139}
\end{enumerate}
\end{exercise}
\begin{exercise}\label{exercise:3.25}\index{SI0469@风险率函数}
假设随机变量$T$为某个对象的寿命（可能是电子元件的寿命，或接受某种治疗的受试者的生存时间）. 与随机变量$T$相关的风险函数$h_T(t)$定义为
$$h_T(t)=\lim_{\delta\to0}\frac{P(t\le T<t+\delta\mid T\ge t)}\delta.$$
因此，可把$h_T(t)$解释为：在对象存活到时间$t$的条件下，该对象存活至略晚于$t$的时刻的概率变化率.\footnote{译注：定义式的分子是给定存活到$t$后在$[t,t+\delta)$内失效的概率，故风险函数更准确地说是条件瞬时失效率；原文将其描述为“存活至略晚于$t$”的概率变化率，表述容易混淆. 极限应理解为$\delta\downarrow0$.} 证明若$T$为连续随机变量，则
$$h_T(t)=\frac{f_T(t)}{1-F_T(t)}=-\frac d{dt}\log(1-F_T(t)).$$
\end{exercise}
\begin{exercise}\label{exercise:3.26}
验证下列pdf具有所给出的风险函数（见\autoref{exercise:3.25}）.
\begin{enumerate}
\item[(a)] 若$T\sim\operatorname{exponential}(\beta)$，则$h_T(t)=1/\beta$.\index{SI0291@风险率函数}\semanticindexanchor{291:139}
\item[(b)] 若$T\sim\operatorname{Weibull}(\gamma,\beta)$，则$h_T(t)=(\gamma/\beta)t^{\gamma-1}$.\index{SI0392@风险率函数}\semanticindexanchor{392:139}
\item[(c)] 若$T\sim\operatorname{logistic}(\mu,\beta)$，即\index{SI0330@风险率函数}\semanticindexanchor{330:139}
$$F_T(t)=\frac1{1+e^{-(t-\mu)/\beta}},$$
则$h_T(t)=(1/\beta)F_T(t)$.
\end{enumerate}
\end{exercise}
\begin{exercise}\label{exercise:3.27}
对下列每个族，说明族中是否所有pdf都是单峰的（见\autoref{exercise:2.27}）.
\begin{enumerate}
\item[(a)] 均匀$(a,b)$.
\item[(b)] Gamma$(\alpha,\beta)$.
\item[(c)] $\operatorname n(\mu,\sigma^2)$.
\item[(d)] Beta$(\alpha,\beta)$.
\end{enumerate}
\end{exercise}
\sourcepageanchor{140}

\begin{exercise}\label{exercise:3.28}\index{SI0435@指数族}
证明下列各族都是指数族.
\begin{enumerate}
\item[(a)] 参数$\mu$或$\sigma$之一已知的正态族.
\item[(b)] 参数$\alpha$或$\beta$之一已知，或二者均未知的Gamma族.
\item[(c)] 参数$\alpha$或$\beta$之一已知，或二者均未知的Beta族.
\item[(d)] Poisson族.
\item[(e)] $r$已知、$0<p<1$的负二项族.
\end{enumerate}
\end{exercise}
\begin{exercise}\label{exercise:3.29}
对\autoref{exercise:3.28}中的每个族，描述其自然参数空间.
\end{exercise}
\begin{exercise}\label{exercise:3.30}
使用\autoref{thm:3.4.2}的恒等式：
\begin{enumerate}
\item[(a)] 计算二项随机变量的方差.
\item[(b)] 计算Beta$(a,b)$随机变量的均值和方差.
\end{enumerate}
\end{exercise}
\begin{exercise}\label{exercise:3.31}
本题将证明\autoref{thm:3.4.2}.
\begin{enumerate}
\item[(a)] 从等式
$$\int f(x\mid\theta)=h(x)c(\theta)\exp\left(\sum_{i=1}^k w_i(\theta)t_i(x)\right)\,dx=1,$$
出发，对两边求导，再整理各项以建立\eqref{eq:3.4.4}. （$\frac d{dx}\log g(x)=g'(x)/g(x)$这一事实有帮助.）\footnote{译注：原式把“$f(x\mid\theta)=$”直接写在积分号与被积表达式之间，形式不完整. 应理解为$\int f(x\mid\theta)\,dx=\int h(x)c(\theta)\exp(\sum_iw_i(\theta)t_i(x))\,dx=1$. 原式照录.}
\item[(b)] 对上面的等式再求一次导数，然后整理以建立\eqref{eq:3.4.5}. （$\frac{d^2}{dx^2}\log g(x)=g''(x)/g(x)-(g'(x)/g(x))^2$这一事实有帮助.）
\end{enumerate}
\end{exercise}
\begin{exercise}\label{exercise:3.32}\leavevmode
\begin{enumerate}
\item[(a)] 若一个指数族可写成\eqref{eq:3.4.7}的形式，证明\autoref{thm:3.4.2}中的恒等式简化为
\begin{align*}
E(t_j(X))&=-\frac\partial{\partial\eta_j}\log c^*(\eta),\\
\operatorname{Var}(t_j(X))&=-\frac{\partial^2}{\partial\eta_j^2}\log c^*(\eta).
\end{align*}
\item[(b)] 用这个恒等式计算Gamma$(a,b)$随机变量的均值与方差.
\end{enumerate}
\end{exercise}
\begin{exercise}\label{exercise:3.33}
对下列每个族：
\begin{enumerate}
\item[(i)] 验证它是指数族.
\item[(ii)] 描述参数向量$\theta$所在的曲线.
\item[(iii)] 画出弯曲参数空间的示意图.
\end{enumerate}
\begin{enumerate}
\item[(a)] $\operatorname n(\theta,\theta)$.
\item[(b)] $\operatorname n(\theta,a\theta^2)$，$a$已知.
\item[(c)] Gamma$(\alpha,1/\alpha)$.
\item[(d)] $f(x\mid\theta)=C\exp(-(x-\theta)^4)$，$C$为归一化常数.
\end{enumerate}
\end{exercise}
\needspace{6\baselineskip}
\begin{exercise}\label{exercise:3.34}\index{SI0437@正态近似}
在\autoref{ex:3.4.9}中，我们看到正态近似可产生弯曲指数族. 对下列每个正态近似：
\begin{enumerate}
\item[(i)] 描述参数向量$\theta$所在的曲线.
\item[(ii)] 画出弯曲参数空间的示意图.
\end{enumerate}
\begin{enumerate}
\item[(a)] Poisson近似：$\bar X\sim\operatorname n(\lambda,\lambda/n)$.
\item[(b)] 二项近似：$\bar X\sim\operatorname n(p,p(1-p)/n)$.
\item[(c)] 负二项近似：$\bar X\sim\operatorname n(r(1-p)/p,r(1-p)/np^2)$.
\end{enumerate}
\end{exercise}
\begin{exercise}\label{exercise:3.35}\leavevmode
\begin{enumerate}
\item[(a)] 近似Poisson的正态族也可参数化为$\operatorname n(e^\theta,e^\theta)$，其中$-\infty<\theta<\infty$. 画出参数空间的图形，并与\hyperref[exercise:3.34]{习题3.34(a)}中的近似比较.\footnote{译注：原式方差为$e^\theta$，而所引用的\hyperref[exercise:3.34]{习题3.34(a)}是样本均值的近似，其方差为$\lambda/n$. 若这里也表示同一固定样本量$n$的样本均值族，在令$\lambda=e^\theta$后方差应为$e^\theta/n$；若指单个Poisson变量的正态近似，则原式可成立. 原文意图待核实.}
\sourcepageanchor{141}
\item[(b)] 假设$X\sim\operatorname{gamma}(\alpha,\beta)$，并假定$EX=\mu$. 画出参数空间的示意图.
\item[(c)] 假设$X_i\sim\operatorname{gamma}(\alpha_i,\beta_i)$，$i=1,2,\ldots,n$，并假定$EX_i=\mu$. 描述参数空间$(\alpha_1,\ldots,\alpha_n,\beta_1,\ldots,\beta_n)$.
\end{enumerate}
\end{exercise}
\begin{exercise}\label{exercise:3.36}
考虑pdf $f(x)=\frac{63}4(x^6-x^8)$，$-1<x<1$. 对下列各情形，在同一坐标轴上画出$(1/\sigma)f((x-\mu)/\sigma)$.
\begin{enumerate}
\item[(a)] $\mu=0,\sigma=1$.
\item[(b)] $\mu=3,\sigma=1$.
\item[(c)] $\mu=3,\sigma=2$.
\end{enumerate}
\end{exercise}
\begin{exercise}\label{exercise:3.37}
证明，若$f(x)$是关于0对称的pdf，则$\mu$是位置—尺度pdf $(1/\sigma)f((x-\mu)/\sigma)$，$-\infty<x<\infty$的中位数.
\end{exercise}
\begin{exercise}\label{exercise:3.38}
设$Z$为pdf为$f(z)$的随机变量. 定义$z_\alpha$为满足下列关系的一个数：
$$\alpha=P(Z>z_\alpha)=\int_{z_\alpha}^\infty f(z)\,dz.$$
证明，若$X$是pdf为$(1/\sigma)f((x-\mu)/\sigma)$的随机变量，且$x_\alpha=\sigma z_\alpha+\mu$，则$P(X>x_\alpha)=\alpha$. （因此，若有$z_\alpha$值的数表，就能容易地算出位置—尺度族任意成员的$x_\alpha$值.）
\end{exercise}
\begin{exercise}\label{exercise:3.39}\index{SI0276@柯西}
考虑\autoref{sec:3.3}定义的Cauchy族. 这个族可推广为位置—尺度族，其pdf形式为
$$f(x\mid\mu,\sigma)=\frac1{\sigma\pi\left(1+\left(\frac{x-\mu}\sigma\right)^2\right)},\qquad-\infty<x<\infty.$$
Cauchy分布的均值和方差不存在，所以参数$\mu$和$\sigma^2$不是均值和方差. 但它们仍有重要意义. 证明，若$X$是参数为$\mu,\sigma$的Cauchy随机变量，则
\begin{enumerate}
\item[(a)] $\mu$是$X$分布的中位数，即$P(X\ge\mu)=P(X\le\mu)=1/2$.
\item[(b)] $\mu+\sigma$和$\mu-\sigma$是$X$分布的四分位数，即$P(X\ge\mu+\sigma)=P(X\le\mu-\sigma)=1/4$. （提示：先对$\mu=0,\sigma=1$证明，再用\autoref{exercise:3.38}.）
\end{enumerate}
\end{exercise}
\begin{exercise}\label{exercise:3.40}
设$f(x)$为任意均值为$\mu$、方差为$\sigma^2$的pdf. 说明如何以$f(x)$为基础构造位置—尺度族，使该族的标准pdf（记为$f^*(x)$）具有均值0、方差1.
\end{exercise}
\begin{exercise}\label{exercise:3.41}\index{SI0989@递增}
若$\theta_1>\theta_2$蕴含$F(x\mid\theta_1)$在随机意义下大于$F(x\mid\theta_2)$，则称cdf族$\{F(x\mid\theta),\theta\in\Theta\}$关于$\theta$随机递增（stochastically increasing）. （随机大于的定义见\autoref{exercise:1.49}.）
\begin{enumerate}
\item[(a)] 证明，固定$\sigma^2$时，$\operatorname n(\mu,\sigma^2)$族关于$\mu$随机递增.
\item[(b)] 证明，固定形状参数$\alpha$时，\eqref{eq:3.3.6}的Gamma$(\alpha,\beta)$族关于尺度参数$\beta$随机递增.
\end{enumerate}
\end{exercise}
\begin{exercise}\label{exercise:3.42}
随机递增族的定义见\autoref{exercise:3.41}.
\begin{enumerate}
\item[(a)] 证明，位置族关于其位置参数随机递增.
\item[(b)] 证明，若样本空间为$[0,\infty)$，则尺度族关于其尺度参数随机递增.
\end{enumerate}
\end{exercise}
\begin{exercise}\label{exercise:3.43}\index{SI0987@递减}
若$\theta_1>\theta_2$蕴含$F(x\mid\theta_2)$在随机意义下大于$F(x\mid\theta_1)$，则称cdf族$\{F(x\mid\theta),\theta\in\Theta\}$关于$\theta$随机递减（stochastically decreasing）. （见\autoref{exercise:3.41}和\autoref{exercise:3.42}.）
\begin{enumerate}
\item[(a)] 证明，若$X\sim F_X(x\mid\theta)$，其中$X$的样本空间为$(0,\infty)$，且$F_X(x\mid\theta)$关于$\theta$随机递增，则$F_Y(y\mid\theta)$关于$\theta$随机递减，其中$Y=1/X$.
\item[(b)] 证明，若$X\sim F_X(x\mid\theta)$，其中$F_X(x\mid\theta)$关于$\theta$随机递增且$\theta>0$，则$F_X(x\mid1/\theta)$关于$\theta$随机递减.
\end{enumerate}
\end{exercise}
\begin{exercise}\label{exercise:3.44}
对任意$EX^2$与$E|X|$存在的随机变量$X$，证明$P(|X|\ge b)$既不超过$EX^2/b^2$，也不超过$E|X|/b$，其中$b$为正常数. 若$f(x)=e^{-x}$，$x>0$，证明当$b=3$时其中一个界更好，而当$b=\sqrt2$时另一个界更好. （注意\hyperref[sec:3.8.2]{杂录3.8.2}中的Markov不等式.）
\end{exercise}
\sourcepageanchor{142}

\begin{exercise}\label{exercise:3.45}\index{SI0574@矩母函数}
设随机变量$X$具有矩母函数$M_X(t)$，$-h<t<h$.
\begin{enumerate}
\item[(a)]\phantomsection\label{itm:exercise:3.45:a} 证明$P(X\ge a)\le e^{-at}M_X(t)$，$0<t<h$. （采用类似Chebyshev不等式证明的方法即可.）
\item[(b)] 类似地，证明$P(X\le a)\le e^{-at}M_X(t)$，$-h<t<0$.
\item[(c)] \hyperref[itm:exercise:3.45:a]{(a)}的一个特例是：对矩母函数有定义的所有$t\ge0$，$P(X\ge0)\le Ee^{tX}$. 对函数$h(t,x)$，什么样的一般条件可以保证：对$Eh(t,X)$存在的所有$t\ge0$，都有$P(X\ge0)\le Eh(t,X)$？［在\hyperref[itm:exercise:3.45:a]{(a)}中，$h(t,x)=e^{tx}$.］
\end{enumerate}
\end{exercise}
\begin{exercise}\label{exercise:3.46}
对$X\sim\operatorname{uniform}(0,1)$和$X\sim\operatorname{exponential}(\lambda)$，计算$P(|X-\mu_X|\ge k\sigma_X)$，并将答案与Chebyshev不等式给出的界比较.
\end{exercise}
\begin{exercise}\label{exercise:3.47}\index{SI0576@正态}
若$Z$为标准正态随机变量，证明\autoref{ex:3.6.3}中不等式的以下对应式：
$$P(|Z|\ge t)\ge\sqrt{\frac2\pi}\frac t{1+t^2}e^{-t^2/2}.$$
\end{exercise}
\begin{exercise}\label{exercise:3.48}
对二项分布、负二项分布和超几何分布，推导类似\eqref{eq:3.6.2}的递推关系.
\end{exercise}
\begin{exercise}\label{exercise:3.49}
在假定函数$g$满足适当条件的前提下，证明下列Stein引理的对应结果.
\begin{enumerate}
\item[(a)] 若$X\sim\operatorname{gamma}(\alpha,\beta)$，则\index{SI0539@伽马}\semanticindexanchor{539:142}
$$E(g(X)(X-\alpha\beta))=\beta E(Xg'(X)).$$
\item[(b)] 若$X\sim\operatorname{beta}(\alpha,\beta)$，则\index{SI0530@贝塔}\semanticindexanchor{530:142}
$$E\left[g(X)\left(\beta-(\alpha-1)\frac{1-X}X\right)\right]=E((1-X)g'(X)).$$
\end{enumerate}
\end{exercise}
\begin{exercise}\label{exercise:3.50}
证明\hyperref[thm:3.6.8]{定理3.6.8(b)}中给出的负二项分布恒等式.
\end{exercise}



\subsection{杂录}\label{sec:3.8}
\subsubsection{Poisson公设}\label{sec:3.8.1}
Poisson分布可从一组基本假设推导出来，这些假设有时称为Poisson公设（Poisson postulates）. 它们与所考察过程的物理性质有关. 一般而言，这些假设并不容易验证，但它们为实验人员判断Poisson是否能提供合理模型给出了一组指导原则. 关于Poisson公设更完整的论述，见\hyperref[bib:source118]{Feller（1968）}或\hyperref[bib:source12]{Barr和Zehna（1983）}的经典著作.

\begin{theorem}\label{thm:3.8.1}\index{SI0374@基本假设}\index{SI0818@泊松基本假设}
对每个$t\ge0$，设$N_t$为具有下列性质的整数值随机变量. （把$N_t$理解为从时刻0到时刻$t$这一时间段内的到达次数.）
\begin{enumerate}
\item[(i)]\phantomsection\label{itm:thm:3.8.1:i} $N_0=0$. \hfill（开始时没有到达）
\item[(ii)] $s<t\Rightarrow N_s$与$N_t-N_s$独立. \hfill（不相交时间段内的到达独立）
\item[(iii)] $N_s$与$N_{t+s}-N_t$同分布. \hfill（到达次数只依赖时间段长度）
\item[(iv)] $\displaystyle\lim_{t\to0}\frac{P(N_t=1)}t=\lambda$. \hfill（时间段很短时，到达概率与其长度成比例）
\item[(v)]\phantomsection\label{itm:thm:3.8.1:v} $\displaystyle\lim_{t\to0}\frac{P(N_t>1)}t=0$. \hfill（没有同时到达）
\end{enumerate}
若\hyperref[itm:thm:3.8.1:i]{(i)}至\hyperref[itm:thm:3.8.1:v]{(v)}成立，则对任意整数$n$，
$$P(N_t=n)=e^{-\lambda t}\frac{(\lambda t)^n}{n!},$$
即$N_t\sim\operatorname{Poisson}(\lambda t)$.
\end{theorem}
\sourcepageanchor{143}

这些公设也可解释为描述对象的空间行为（例如昆虫的运动），从而给出Poisson在空间分布方面的应用.

\subsubsection{Chebyshev及其推广}\label{sec:3.8.2}
\hyperref[bib:source140]{Ghosh和Meeden（1977）}讨论了Chebyshev不等式非常保守，而且几乎从不取到等号这一事实. 若用$\bar X_n$表示随机变量$X_1,X_2,\ldots,X_n$的均值，则Chebyshev不等式给出\footnote{译注：这里以及后面的可达性结论，需要样本的相应假设，通常是$X_1,\ldots,X_n$独立同分布，具有共同均值$\mu$和方差$\sigma^2$. 原段未明说；若只给任意随机变量，则$\operatorname{Var}(\bar X_n)=\sigma^2/n$一般不成立.}
$$P(|\bar X_n-\mu|\ge k\sigma)\le\frac1{nk^2}.$$

他们证明了下述定理.

\begin{theorem}\label{thm:3.8.2}\index{SI0565@切比雪夫}
若$0<\sigma<\infty$，则
\begin{enumerate}
\item[(a)] 若$n=1$，则当$k\ge1$时不等式可取到等号，当$0<k<1$时不可取到等号.
\item[(b)] 若$n=2$，则不等式可取到等号，当且仅当$k=1$.
\item[(c)] 若$n\ge3$，则不等式不可取到等号.
\end{enumerate}
\end{theorem}
对于不等式取到等号的情形，他们给出了例子. 其大部分技术论证基于以下称为Markov不等式的结果.

\begin{lemma}[Markov不等式]\label{lem:3.8.3}\index{SI0572@马尔可夫}
若$P(Y\ge0)=1$且$P(Y=0)<1$，则对任意$r>0$，
$$P(Y\ge r)\le\frac{EY}r,$$
其中等号成立，当且仅当$P(Y=r)=p=1-P(Y=0)$，$0<p\le1$.
\end{lemma}
随后可把Markov不等式用于
$$Y=\frac{(\bar X_n-\mu)^2}{\sigma^2},$$
得到上述结果.

Chebyshev不等式如此宽松的一个原因，是它没有对底层分布施加限制. 加上单峰性（unimodality）的限制后，就可以得到更紧的界，以及Gauss不等式和Vysochanski\u{i}--Petunin不等式. （这些不等式的细节及基于初等微积分的证明见\mbox{\hyperref[bib:source248]{Pukelsheim（1994）}.）}

\begin{theorem}[Gauss不等式]\label{thm:3.8.4}\index{SI0567@高斯}
设$X\sim f$，其中$f$为众数在$\nu$处的单峰分布，定义$\tau^2=E(X-\nu)^2$. 则
$$P(|X-\nu|>\varepsilon)\le\begin{cases}\dfrac{4\tau^2}{9\varepsilon^2},&\text{对所有 }\varepsilon\ge\sqrt{4/3}\,\tau,\\[4pt]1-\dfrac\varepsilon{\sqrt3\tau},&\text{对所有 }\varepsilon\le\sqrt{4/3}\,\tau.\end{cases}$$
\end{theorem}
虽然这个界比Chebyshev更紧，但对众数的依赖限制了它的用途. Vysochanski\u{i}--Petunin的推广消除了这一限制.

\begin{theorem}[Vysochanski\u{i}--Petunin不等式]\label{thm:3.8.5}\index{SI0579@Vysochanskiĭ–Petunin}
设$X\sim f$，其中$f$为单峰分布，并对任意点$\alpha$定义$\xi^2=E(X-\alpha)^2$. 则\footnote{译注：原式第二分支的分母印为$9\varepsilon^2$. 在$\varepsilon=\sqrt{8/3}\,\xi$处，它给出$-1/6$，不可能是一般概率上界，并且与第一分支的$1/6$不一致. 该分支疑应为$4\xi^2/(3\varepsilon^2)-1/3$；原式保留，需与所引\hyperref[bib:source248]{Pukelsheim（1994）}原论文进一步核实.}
$$P(|X-\alpha|>\varepsilon)\le\begin{cases}\dfrac{4\xi^2}{9\varepsilon^2},&\text{对所有 }\varepsilon\ge\sqrt{8/3}\,\xi,\\[4pt]\dfrac{4\xi^2}{9\varepsilon^2}-\dfrac13,&\text{对所有 }\varepsilon\le\sqrt{8/3}\,\xi.\end{cases}$$
\end{theorem}

\sourcepageanchor{144}

Pukelsheim指出，取$\alpha=\mu=E(X)$及$\varepsilon=3\sigma$，其中$\sigma^2=\operatorname{Var}(X)$，得到
$$P(|X-\mu|>3\sigma)\le\frac4{81}<.05,$$
这就是所谓的三西格玛规则（three sigma rule）：$X$与总体均值相差超过三个标准差的概率小于5\%.

\subsubsection{关于指数族的进一步讨论}\label{sec:3.8.3}\index{SI0438@对数正态}
对数正态分布属于指数族吗？\eqref{eq:3.3.21}给出的密度可以写成\eqref{eq:3.4.1}规定的形式，因此我们把对数正态纳入了指数族.

按照\hyperref[bib:source50]{Brown（1986}，第1.1节），定义一个指数分布族时，先从非负函数$\nu(x)$出发，定义集合$\mathcal N$为
$$\mathcal N=\left\{\theta:\int_{\mathcal X}e^{\theta x}\nu(x)\,dx<\infty\right\}.$$
令$\lambda(\theta)=\int_{\mathcal X}e^{\theta x}\nu(x)\,dx$，由
$$f(x\mid\theta)=\frac{e^{\theta x}\nu(x)}{\lambda(\theta)},\qquad x\in\mathcal X,\quad\theta\in\mathcal N$$
定义的概率密度集合，就是一个指数族. $f(x\mid\theta)$的矩母函数为
$$M_X(t)=\int_{\mathcal X}e^{tx}f(x\mid\theta)\,dx=\frac{\lambda(t+\theta)}{\lambda(\theta)},$$
因而根据构造它存在. 若参数空间$\Theta$等于集合$\mathcal N$，则称该指数族为满族. 当$\Theta$为$\mathcal N$的较低维子集时，就得到弯曲指数族.

回到对数正态分布，我们知道它没有矩母函数，所以它不能满足Brown关于指数族的定义. 不过，对数正态满足\autoref{thm:3.4.2}的期望恒等式，也具有\autoref{sec:6.2.1}（\autoref{thm:6.2.10}）中详述的充分性性质. 对本书的目的而言，这些就是所需的主要性质，也是把一个分布识别为指数族成员的主要理由. 更高级的性质可能需要矩母函数存在，但我们不在这里研究这些性质.\footnote{译注：本段比较时需区分观测值$X$与指数表示中的充分统计量. 对数正态密度确实可按\eqref{eq:3.4.1}写成以$\log X$、$(\log X)^2$为统计量的指数族；$X$本身没有正参数邻域上的矩母函数，并不否定这一点. 这里显示的$e^{\theta x}$形式使用$x$本身作为自然统计量，较为特殊. 此外，$\lambda(t+\theta)/\lambda(\theta)$在$0$的邻域内有限，需要$\theta$处于相应自然参数空间的内点；仅有$\theta\in\mathcal N$不足以保证双侧邻域存在. 原文关于Brown定义的具体归属仍需核对所引专著.}





\section{多个随机变量}\label{sec:4}
\chapterdoi{4}

\sourcepageanchor{145}

\begin{quote}
我承认自己曾像鼹鼠一样盲目，但迟些学得智慧，总比永远学不到好.
\hfill 夏洛克·福尔摩斯，《歪唇男人》
\end{quote}



\subsection{联合分布与边缘分布}\label{sec:4.1}
前几章研究的模型只涉及一个随机变量，称为单变量模型（univariate model）. 本章讨论包含多个随机变量的模型，称为多变量模型（multivariate model）.

在实验中，我们很少只观测一个值. 例如，研究某个总体的健康状况时，如果只测量一个人的体重，这个实验的规模实在很小. 若测量若干人的体重，就有若干个随机变量，每个人对应一个. 或者，也可以为每个人测量若干特征，例如体重、体温、身高和血压；每个特征都对应一个随机变量. 因此，需要能够处理多个随机变量的模型. 本章前几节主要讨论两个随机变量的模型，即二元模型（bivariate model）.

回顾\autoref{def:1.4.1}，随机变量是从样本空间 $S$ 到实数集的函数. 由若干随机变量组成的向量也可以类似地定义.
\begin{definition}\label{def:4.1.1}\index{SI0871@随机向量}
一个 $n$ 维随机向量（random vector）是从样本空间 $S$ 到 $\mathbb R^n$（$n$ 维欧几里得空间）的函数.
\end{definition}
例如，若为样本空间的每个点指定平面 $\mathbb R^2$ 中的一个有序对 $(x,y)$，就定义了二维（或二元）随机向量 $(X,Y)$. 下例说明这一过程.
\begin{example}[骰子的样本空间]\label{ex:4.1.2}
掷两颗公平骰子的样本空间包含 36 个等可能的点，已在\autoref{ex:1.3.10} 中介绍. 点 $(3,3)$ 表示两颗骰子均出现 3；点 $(4,1)$ 表示第一颗出现 4，第二颗出现 1，等等. 现在为每个样本点指定两个随机变量 $X$ 和 $Y$ 的值，其中
$$
X=\text{两颗骰子的点数之和},\qquad Y=|\text{两颗骰子的点数之差}|.
$$
\sourcepageanchor{146}

\begin{table}[htbp]
\centering\small
\caption{联合 pmf $f(x,y)$ 的取值}\label{tab:4.1.1}
\setlength{\tabcolsep}{5pt}
\begin{tabular}{c|ccccccccccc}
 &\multicolumn{11}{c}{$x$}\\
$y$&2&3&4&5&6&7&8&9&10&11&12\\\hline
0&$\frac1{36}$&&$\frac1{36}$&&$\frac1{36}$&&$\frac1{36}$&&$\frac1{36}$&&$\frac1{36}$\\[3pt]
1&&$\frac1{18}$&&$\frac1{18}$&&$\frac1{18}$&&$\frac1{18}$&&$\frac1{18}$&\\[3pt]
2&&&$\frac1{18}$&&$\frac1{18}$&&$\frac1{18}$&&$\frac1{18}$&&\\[3pt]
3&&&&$\frac1{18}$&&$\frac1{18}$&&$\frac1{18}$&&&\\[3pt]
4&&&&&$\frac1{18}$&&$\frac1{18}$&&&&\\[3pt]
5&&&&&&$\frac1{18}$&&&&&
\end{tabular}
\end{table}

例如，在样本点 $(3,3)$ 处，$X=6$，$Y=0$；在 $(4,1)$ 处，$X=5$，$Y=3$；在 $(1,4)$ 处也是 $X=5$，$Y=3$. 对全部 36 个样本点进行这样的计算，就定义了随机向量 $(X,Y)$.

有关 $X$ 和 $Y$ 的事件，其概率由样本空间中相应事件的概率来定义. 例如，$P(X=5\text{ 且 }Y=3)$ 是多少？只有样本点 $(4,1)$ 和 $(1,4)$ 对应 $X=5$ 且 $Y=3$. 换言之，$X=5$ 且 $Y=3$ 当且仅当事件 $\{(4,1),(1,4)\}$ 发生. 因为这 36 个样本点等可能，
$$
P(\{(4,1),(1,4)\})=\frac2{36}=\frac1{18}.
$$
所以 $P(X=5\text{ 且 }Y=3)=1/18$. 此后记为 $P(X=5,Y=3)$，逗号表示“且”. 同理，$P(X=6,Y=0)=1/36$，因为只有样本点 $(3,3)$ 对应这组值. 更复杂的事件也可照此处理. 例如，$P(X=7,Y\le4)=4/36=1/9$，因为对应的样本点是 $(4,3)$、$(3,4)$、$(5,2)$ 和 $(2,5)$.
\end{example}

这个随机向量称为离散的（discrete），因为其可能取值的个数是可数的（此处为有限个）. 一旦知道函数 $f(x,y)=P(X=x,Y=y)$，就能计算所有这类事件的概率.\index{SI0873@离散}\semanticindexanchor{873:146}
\begin{definition}\label{def:4.1.3}\index{SI0854@二元}
设 $(X,Y)$ 是离散随机向量. 函数 $f:\mathbb R^2\to\mathbb R$，
$$f(x,y)=P(X=x,Y=y),$$
称为联合概率质量函数（joint probability mass function），简称联合 pmf. 如果需要强调它属于哪个随机向量，就记为 $f_{X,Y}$.
\end{definition}
联合 pmf 完全确定随机向量的概率分布，正如单变量的 pmf 完全确定一个随机变量的分布一样. \autoref{ex:4.1.2} 的随机向量有 21 个可能取值，\autoref{tab:4.1.1} 列出了联合 pmf 的值. $f(5,3)=1/18$ 和 $f(6,0)=1/36$ 已经算过，其余各值可类似计算. 注意，函数 $f$ 定义在整个 $\mathbb R^2$ 上，而不仅仅是这 21 个有序对上. 对所有其他 $(x,y)$，$f(x,y)=P(X=x,Y=y)=0$.
\sourcepageanchor{147}

借助联合 pmf，可以计算任意事件 $A\subset\mathbb R^2$ 的概率：
$$P((X,Y)\in A)=\sum_{(x,y)\in A}f(x,y).$$
因为 $f$ 至多在可数个点上非零，即使 $A$ 不可数，此和也至多包含可数个非零项. 例如，令 $A=\{(x,y):x=7\text{ 且 }y\le4\}$. 这是 $\mathbb R^2$ 中一条半无限直线；\autoref{tab:4.1.1} 中只有 $(7,1)$ 和 $(7,3)$ 属于 $A$ 且具有非零概率. 因此，
$$
P(X=7,Y\le4)=P((X,Y)\in A)=f(7,1)+f(7,3)=\frac1{18}+\frac1{18}=\frac19.
$$
这与\autoref{ex:4.1.2} 通过样本点得到的结果相同. 通常，用联合 pmf 计算更简单.

随机向量的函数的期望也与单变量情形相同. 若 $g(x,y)$ 是在随机向量所有可能取值处定义的实值函数，则 $g(X,Y)$ 是一个随机变量，其期望为\index{SI0423@二元}\semanticindexanchor{423:147}
$$Eg(X,Y)=\sum_{(x,y)\in\mathbb R^2}g(x,y)f(x,y).$$
\begin{example}[续\autoref{ex:4.1.2}]\label{ex:4.1.4}
对于联合 pmf 如\autoref{tab:4.1.1} 所示的随机向量，乘积 $XY$ 的平均值是多少？取 $g(x,y)=xy$，对 21 个可能取值求和，得到
$$
EXY=(2)(0)\frac1{36}+(4)(0)\frac1{36}+\cdots+(8)(4)\frac1{18}+(7)(5)\frac1{18}=13\frac{11}{18}.
$$
\end{example}
将 $X$ 替换为随机向量 $(X,Y)$ 后，\autoref{thm:2.2.5} 的期望性质仍成立. 例如，对于函数 $g_1,g_2$ 和常数 $a,b,c$，
$$E(ag_1(X,Y)+bg_2(X,Y)+c)=aEg_1(X,Y)+bEg_2(X,Y)+c.$$
由求和的性质，可以像单变量情形一样证明这些性质（\autoref{exercise:4.2}）.

由于联合 pmf 的值是概率，必有 $f(x,y)\ge0$，且
$$\sum_{(x,y)\in\mathbb R^2}f(x,y)=P((X,Y)\in\mathbb R^2)=1.$$
反之，任何从 $\mathbb R^2$ 到 $\mathbb R$ 的非负函数，只要至多在可数个点上非零且这些值之和等于 1，就是某个随机向量的联合 pmf. 因此，我们可以直接指定联合 pmf 来构造模型，而不必描述底层样本空间.
\begin{example}[骰子的联合 pmf]\label{ex:4.1.5}
定义
$$
\begin{gathered}
f(0,0)=f(0,1)=\frac16,\\
f(1,0)=f(1,1)=\frac13,\\
f(x,y)=0,\qquad\text{其余所有 }(x,y).
\end{gathered}
$$
\sourcepageanchor{148}

这个函数非负，而且所有值的和为 1，所以它是一个联合 pmf. 无须知道样本空间，便可计算 $P(X=Y)=f(0,0)+f(1,1)=1/2$. 许多不同的样本空间及其上的函数都能产生这个 pmf. 下面是一种构造：取掷两颗骰子的 36 点样本空间. 若第一颗骰子的点数不超过 2，令 $X=0$，否则 $X=1$；若第二颗骰子的点数为奇数，令 $Y=0$，否则 $Y=1$. \autoref{exercise:4.3} 要求验证这一构造确实给出上述联合 pmf.
\end{example}

即使用随机向量建模，有时我们仍只关心其中一个随机变量的概率或期望，例如 $P(X=2)$. $X$ 本身是随机变量，有\autoref{sec:1} 中定义的 pmf $f_X(x)=P(X=x)$. 下标用于区别它与联合 pmf $f_{X,Y}$. 为了强调它是随机向量联合模型中 $X$ 的 pmf，我们称之为 $X$ 的边缘 pmf（marginal pmf）. 下述定理表明，由联合 pmf 很容易求得边缘 pmf.
\begin{theorem}\label{thm:4.1.6}\index{SI0855@边缘}
设离散随机向量 $(X,Y)$ 的联合 pmf 为 $f_{X,Y}$. 则 $X$ 和 $Y$ 的边缘 pmf $f_X(x)=P(X=x)$ 和 $f_Y(y)=P(Y=y)$ 为
$$f_X(x)=\sum_{y\in\mathbb R}f_{X,Y}(x,y),\qquad f_Y(y)=\sum_{x\in\mathbb R}f_{X,Y}(x,y).$$
\end{theorem}
\begin{proof}
只证明 $f_X$ 的公式，$f_Y$ 的证明类似. 对任意 $x\in\mathbb R$，令 $A_x=\{(x,y):-\infty<y<\infty\}$，即第一坐标为 $x$ 的直线. 则
$$
\begin{aligned}
f_X(x)&=P(X=x)\\
&=P(X=x,-\infty<Y<\infty)&&[P(-\infty<Y<\infty)=1]\\
&=P((X,Y)\in A_x)&&[A_x\text{ 的定义}]\\
&=\sum_{(x,y)\in A_x}f_{X,Y}(x,y)\\
&=\sum_{y\in\mathbb R}f_{X,Y}(x,y).
\end{aligned}
$$
\end{proof}
\begin{example}[骰子的边缘 pmf]\label{ex:4.1.7}
将\autoref{thm:4.1.6} 用于\autoref{tab:4.1.1} 的联合 pmf. 为求 $Y$ 的边缘 pmf，对每个 $y$，将所有可能的 $x$ 对应的概率相加. 例如，
$$
\begin{aligned}
f_Y(0)&=f_{X,Y}(2,0)+f_{X,Y}(4,0)+f_{X,Y}(6,0)\\
&\quad+f_{X,Y}(8,0)+f_{X,Y}(10,0)+f_{X,Y}(12,0)=\frac16.
\end{aligned}
$$
同理，$f_Y(1)=5/18$，$f_Y(2)=2/9$，$f_Y(3)=1/6$，$f_Y(4)=1/9$，$f_Y(5)=1/18$. 这六个概率之和为 1，因为 $Y$ 只有这六个可能取值.
\end{example}
\sourcepageanchor{149}

$X$ 或 $Y$ 的边缘 pmf 与\autoref{sec:1} 中的 pmf 完全相同. 它可以用来计算只涉及该随机变量的概率与期望；但同时涉及两个随机变量时，就需要联合 pmf.
\begin{example}[骰子的概率]\label{ex:4.1.8}
使用上一例中 $Y$ 的边缘 pmf，得到
$$P(Y<3)=f_Y(0)+f_Y(1)+f_Y(2)=\frac16+\frac5{18}+\frac29=\frac23.$$
此外，
$$EY^3=0^3f_Y(0)+1^3f_Y(1)+\cdots+5^3f_Y(5)=20\frac{11}{18}.$$
\end{example}
边缘 pmf 不能完全确定联合 pmf. 许多不同的联合 pmf 可以有相同的边缘 pmf，因此，仅凭边缘 pmf 无法恢复联合 pmf.
\begin{example}[相同的边缘 pmf，不同的联合 pmf]\label{ex:4.1.9}
定义
$$
\begin{gathered}
f(0,0)=\frac1{12},\quad f(1,0)=\frac5{12},\quad f(0,1)=f(1,1)=\frac3{12},\\
f(x,y)=0,\qquad\text{其余所有 }(x,y).
\end{gathered}
$$
于是 $f_Y(0)=f(0,0)+f(1,0)=1/2$，$f_Y(1)=f(0,1)+f(1,1)=1/2$. 同样可算得 $f_X(0)=1/3$，$f_X(1)=2/3$. 可以核对，\autoref{ex:4.1.5} 的联合 pmf 虽然与这里不同，边缘 pmf 却完全相同. 联合 pmf 包含边缘 pmf 中没有的额外信息.
\end{example}

到目前为止，我们讨论的都是离散随机向量. 也可以考虑各分量为连续随机变量的随机向量；通常，这类向量可像单个连续随机变量一样，用密度函数描述.
\begin{definition}\label{def:4.1.10}\index{SI0833@联合}\index{SI0872@连续}
函数 $f:\mathbb R^2\to\mathbb R$ 称为连续二元随机向量 $(X,Y)$ 的联合概率密度函数（joint probability density function），简称联合 pdf，如果对每个 $A\subset\mathbb R^2$，有
$$P((X,Y)\in A)=\iint_A f(x,y)\,dx\,dy.$$
\end{definition}
联合 pdf 与单变量 pdf 的用法相同，只是需要在平面集合上做二重积分. 记号 $\iint_A$ 表示选择积分限，使积分遍及整个集合 $A$. 计算函数的期望时，将求和换为积分、pmf 换为 pdf 即可. 若 $g$ 是实值函数，则\index{SI0423@二元}\semanticindexanchor{423:149}
\begin{equation}\label{eq:4.1.1}
Eg(X,Y)=\int_{-\infty}^{\infty}\int_{-\infty}^{\infty}g(x,y)f(x,y)\,dx\,dy.
\end{equation}
联合 pdf 定义在整个 $\mathbb R^2$ 上. 如果 $P((X,Y)\in A)=0$，它可能在一个很大的集合 $A$ 上都为 0，但在这些点上仍然有定义.
\sourcepageanchor{150}

$X$ 和 $Y$ 的边缘概率密度函数也像离散情形那样定义，只是用积分代替求和. 边缘 pdf 可用来计算只涉及 $X$ 或 $Y$ 的概率或期望. 具体说来，边缘 pdf 为
\begin{equation}\label{eq:4.1.2}\index{SI0834@边缘}
\begin{aligned}
f_X(x)&=\int_{-\infty}^{\infty}f(x,y)\,dy,\qquad-\infty<x<\infty,\\
f_Y(y)&=\int_{-\infty}^{\infty}f(x,y)\,dx,\qquad-\infty<y<\infty.
\end{aligned}
\end{equation}
任何函数 $f(x,y)$，只要对所有 $(x,y)\in\mathbb R^2$ 都有 $f(x,y)\ge0$，并且
$$1=\int_{-\infty}^{\infty}\int_{-\infty}^{\infty}f(x,y)\,dx\,dy,$$
就是某个连续二元随机向量 $(X,Y)$ 的联合 pdf. 下面两个例子说明所有这些有关联合 pdf 的概念.
\begin{example}[联合概率的计算——I]\label{ex:4.1.11}
定义联合 pdf
$$f(x,y)=\begin{cases}6xy^2,&0<x<1\text{ 且 }0<y<1,\\0,&\text{其他情形}.\end{cases}$$
（此后，定义中未明确提及的 $(x,y)$ 均默认满足 $f(x,y)=0$.）首先检验 $f(x,y)$ 确实是联合 pdf. 在定义的范围内，对所有 $(x,y)$ 均有 $f(x,y)\ge0$，这是显然的. 要求 $f(x,y)$ 在整个平面上的积分，注意它在单位正方形以外为 0，因此全平面上的积分等于在这个正方形上的积分. 从而
$$
\begin{aligned}
\int_{-\infty}^{\infty}\int_{-\infty}^{\infty}f(x,y)\,dx\,dy
&=\int_0^1\int_0^1 6xy^2\,dx\,dy=\int_0^1\left.3x^2y^2\right|_0^1\,dy\\
&=\int_0^1 3y^2\,dy=\left.y^3\right|_0^1=1.
\end{aligned}
$$
现在考虑计算 $P(X+Y\ge1)$ 这样的概率. 令 $A=\{(x,y):x+y\ge1\}$，便可将它写为 $P((X,Y)\in A)$. 根据\autoref{def:4.1.10}，只需在集合 $A$ 上对联合 pdf 积分. 然而，联合 pdf 在单位正方形以外为 0，因此在 $A$ 上积分，等于只在 $A$ 落入单位正方形的部分上积分. 集合 $A$ 是平面东北方向的一个半平面；它在单位正方形内的部分是由直线 $x=1$、$y=1$ 和 $x+y=1$ 围成的三角形区域. 可以写成\footnote{译注：下列等式的 $A$ 沿用原文记号，实际指原先半平面 $A$ 与单位正方形的交集.}
$$
\begin{aligned}
A&=\{(x,y):x+y\ge1,\ 0<x<1,\ 0<y<1\}\\
&=\{(x,y):x\ge1-y,\ 0<x<1,\ 0<y<1\}\\
&=\{(x,y):1-y\le x<1,\ 0<y<1\}.
\end{aligned}
$$
由此得到计算该概率所需的积分限：
$$P(X+Y\ge1)=\iint_A f(x,y)\,dx\,dy=\int_0^1\int_{1-y}^1 6xy^2\,dx\,dy=\frac9{10}.$$
利用\eqref{eq:4.1.2}，还可以求 $X$ 或 $Y$ 的边缘 pdf. 例如，为计算 $f_X(x)$，注意当 $x\ge1$ 或 $x\le0$ 时，对所有 $y$ 都有 $f(x,y)=0$. 因此，这时
$$f_X(x)=\int_{-\infty}^{\infty}f(x,y)\,dy=0.$$
\par\pagebreak[4]
\sourcepageanchor{151}

当 $0<x<1$ 时，$f(x,y)$ 仅在 $0<y<1$ 时非零. 因此，对 $0<x<1$，
$$f_X(x)=\int_{-\infty}^{\infty}f(x,y)\,dy=\int_0^1 6xy^2\,dy=\left.2xy^3\right|_0^1=2x.$$
现在可以用 $X$ 的这个边缘 pdf 计算只涉及 $X$ 的概率. 例如，
$$P\left(\frac12<X<\frac34\right)=\int_{1/2}^{3/4}2x\,dx=\frac5{16}.$$
\end{example}
\begin{figure}[htbp]
\centering\includegraphics[width=.54\linewidth]{assets/chapter4/fig4-1-1.pdf}
\caption{\autoref{ex:4.1.12} 中的区域}\label{fig:4.1.1}
\end{figure}

\begin{example}[联合概率的计算——II]\label{ex:4.1.12}
再看一个联合 pdf 的例子：令 $f(x,y)=e^{-y}$，$0<x<y<\infty$. 虽然 $e^{-y}$ 不依赖 $x$，但 $f(x,y)$ 显然是 $x$ 的函数，因为使 $f(x,y)$ 非零的集合依赖 $x$. 用示性函数写成
$$f(x,y)=e^{-y}I_{\{(u,v):0<u<v<\infty\}}(x,y),$$
这一点就更清楚了. 为计算 $P(X+Y\ge1)$，可以在集合 $A=\{(x,y):x+y\ge1\}$ 与 $f(x,y)$ 非零区域的交集上对联合 pdf 积分. 画出这些集合，可见该区域是一个无界区域（\autoref{fig:4.1.1} 中较浅的阴影），三条边分别位于直线 $x=y$、$x+y=1$ 和 $x=0$ 上. 为写出合适的积分限，至少需要把它拆成两部分.

在集合 $B=\{(x,y):x+y<1\}$ 与 $f(x,y)$ 非零区域的交集上积分会更容易. 这个交集是由直线 $x=y$、$x+y=1$ 和 $x=0$ 围成的三角形区域（\autoref{fig:4.1.1} 中较深的阴影）. 因此，
$$
\begin{aligned}
P(X+Y\ge1)&=1-P(X+Y<1)=1-\int_0^{1/2}\int_x^{1-x}e^{-y}\,dy\,dx\\
&=1-\int_0^{1/2}(e^{-x}-e^{-(1-x)})\,dx=2e^{-1/2}-e^{-1}.
\end{aligned}
$$
这说明，在这类问题中，为确定恰当的积分限，把相关集合画出来几乎总是有帮助的.
\end{example}
\sourcepageanchor{152}

$(X,Y)$ 的联合概率分布也可以用联合累积分布函数（joint cumulative distribution function，联合 cdf）来完全描述，而不使用联合 pmf 或联合 pdf. 联合 cdf 是定义在所有 $(x,y)\in\mathbb R^2$ 上的函数\index{SI0239@联合}\semanticindexanchor{239:152}
$$F(x,y)=P(X\le x,Y\le y).$$
对离散随机向量而言，联合 cdf 通常不太方便. 但对于连续二元随机向量，像单变量情形一样，有重要关系
$$F(x,y)=\int_{-\infty}^x\int_{-\infty}^y f(s,t)\,dt\,ds.$$
由二元微积分基本定理，在 $f(x,y)$ 的连续点处，
\begin{equation}\label{eq:4.1.3}\index{SI0448@微积分基本定理}
\frac{\partial^2F(x,y)}{\partial x\partial y}=f(x,y).
\end{equation}
当能够求出 $F(x,y)$ 的表达式时，这个关系十分有用：计算混合偏导数即可得到联合 pdf.



\subsection{条件分布与独立性}\label{sec:4.2}
观测两个随机变量 $(X,Y)$ 时，两者的取值往往有关系. 例如，从一个人群中抽样，令 $X$ 表示某人的身高，$Y$ 表示同一个人的体重. 相比于得知 $X=41$ 英寸，若得知 $X=73$ 英寸，我们自然会认为 $Y>200$ 磅更有可能. 即使知道 $X$ 的值不能确切告诉我们 $Y$ 的值，它也提供了关于 $Y$ 的一些信息. 利用 $(X,Y)$ 的联合分布，可以计算已知 $X$ 值时关于 $Y$ 的条件概率. 然而，有时关于 $X$ 的知识并不提供任何关于 $Y$ 的信息. 本节讨论这些有关条件概率的问题.

若 $(X,Y)$ 是离散随机向量，则 $P(Y=y\mid X=x)$ 这样的条件概率完全按照\autoref{def:1.3.2} 来解释. 对可数个（也许有限个）$x$，有 $P(X=x)>0$. 对这些 $x$，根据定义，$P(Y=y\mid X=x)$ 就是 $P(X=x,Y=y)/P(X=x)$. 事件 $\{Y=y\}$ 对应公式中的 $A$，事件 $\{X=x\}$ 对应 $B$. 固定 $x$ 后，可以对所有可能的 $y$ 计算 $P(Y=y\mid X=x)$，从而在已观测到 $X=x$ 的条件下，评价各个 $y$ 值的概率. 用 $X$ 和 $Y$ 的联合与边缘 pmf 表示这些概率，即 $P(X=x,Y=y)=f(x,y)$，$P(X=x)=f_X(x)$，便可简化计算. 由此得到下面的定义.
\begin{definition}\label{def:4.2.1}\index{SI0851@条件}
设 $(X,Y)$ 是离散二元随机向量，联合 pmf 为 $f(x,y)$，边缘 pmf 为 $f_X(x)$ 和 $f_Y(y)$. 对任意满足 $P(X=x)=f_X(x)>0$ 的 $x$，给定 $X=x$ 时 $Y$ 的条件 pmf（conditional pmf）是关于 $y$ 的函数，记为 $f(y\mid x)$，定义为
$$f(y\mid x)=P(Y=y\mid X=x)=\frac{f(x,y)}{f_X(x)}.$$
对任意满足 $P(Y=y)=f_Y(y)>0$ 的 $y$，给定 $Y=y$ 时 $X$ 的条件 pmf 是关于 $x$ 的函数，记为 $f(x\mid y)$，定义为
$$f(x\mid y)=P(X=x\mid Y=y)=\frac{f(x,y)}{f_Y(y)}.$$
\end{definition}
\sourcepageanchor{153}

既然把 $f(y\mid x)$ 称为 pmf，就应验证这个关于 $y$ 的函数确实定义了某个随机变量的 pmf. 首先，因为 $f(x,y)\ge0$ 且 $f_X(x)>0$，所以对每个 $y$，$f(y\mid x)\ge0$. 其次，
$$\sum_y f(y\mid x)=\frac{\sum_y f(x,y)}{f_X(x)}=\frac{f_X(x)}{f_X(x)}=1.$$
因此，$f(y\mid x)$ 确实是 pmf，可以按通常方式计算已知 $X=x$ 发生时关于 $Y$ 的概率.
\begin{example}[条件概率的计算]\label{ex:4.2.2}
定义 $(X,Y)$ 的联合 pmf 为
$$
\begin{gathered}
f(0,10)=f(0,20)=\frac2{18},\qquad f(1,10)=f(1,30)=\frac3{18},\\
f(1,20)=\frac4{18},\qquad f(2,30)=\frac4{18}.
\end{gathered}
$$
可以用\autoref{def:4.2.1}，分别对 $X$ 的每个可能取值 $x=0,1,2$，计算给定 $X$ 时 $Y$ 的条件 pmf. 首先，$X$ 的边缘 pmf 为
$$
\begin{aligned}
f_X(0)&=f(0,10)+f(0,20)=\frac4{18},\\
f_X(1)&=f(1,10)+f(1,20)+f(1,30)=\frac{10}{18},\\
f_X(2)&=f(2,30)=\frac4{18}.
\end{aligned}
$$
当 $x=0$ 时，$f(0,y)$ 仅在 $y=10$ 和 $y=20$ 时为正. 因此，$f(y\mid0)$ 也仅在这两个值处为正，而且
$$
f(10\mid0)=\frac{f(0,10)}{f_X(0)}=\frac{2/18}{4/18}=\frac12,\qquad
f(20\mid0)=\frac{f(0,20)}{f_X(0)}=\frac12.
$$
也就是说，已知 $X=0$ 时，$Y$ 的条件概率分布是在 $y=10$ 和 $y=20$ 两点各赋予概率 $1/2$ 的离散分布.

当 $x=1$ 时，$f(y\mid1)$ 在 $y=10,20,30$ 处为正，而且
$$
\begin{aligned}
f(10\mid1)&=f(30\mid1)=\frac{3/18}{10/18}=\frac3{10},\\
f(20\mid1)&=\frac{4/18}{10/18}=\frac4{10};
\end{aligned}
$$
当 $x=2$ 时，
$$f(30\mid2)=\frac{4/18}{4/18}=1.$$
最后这个结果反映了一个从联合 pmf 中也能看出的事实：如果知道 $X=2$，就知道 $Y$ 必定等于 30.

利用这些条件 pmf，还能计算其他条件概率. 例如，
$$
\begin{aligned}
P(Y>10\mid X=1)&=f(20\mid1)+f(30\mid1)=\frac7{10},\\
\text{或}\qquad P(Y>10\mid X=0)&=f(20\mid0)=\frac12.
\end{aligned}
$$
\end{example}
\sourcepageanchor{154}

若 $X$ 和 $Y$ 是连续随机变量，则对每个 $x$ 都有 $P(X=x)=0$. 因为分母 $P(X=73)$ 为 0，不能用\autoref{def:1.3.2} 计算 $P(Y>200\mid X=73)$ 这样的条件概率. 可是在实际中，确实观测到了 $X$ 的一个值. 如果在测量精度范围内看到 $X=73$，这仍可能提供有关 $Y$ 的信息，正如本节开头的身高与体重例子所说明的那样. 当 $X$ 和 $Y$ 都连续时，给定 $X=x$ 时 $Y$ 的条件概率分布的恰当定义与离散情形类似，只需用 pdf 替代 pmf（见\hyperref[misc:4.9.3]{杂录 4.9.3}）.
\begin{definition}\label{def:4.2.3}\index{SI0830@条件}
设 $(X,Y)$ 为连续二元随机向量，联合 pdf 为 $f(x,y)$，边缘 pdf 为 $f_X(x)$ 和 $f_Y(y)$. 对任意满足 $f_X(x)>0$ 的 $x$，给定 $X=x$ 时 $Y$ 的条件 pdf（conditional pdf）是关于 $y$ 的函数，记为 $f(y\mid x)$，定义为
$$f(y\mid x)=\frac{f(x,y)}{f_X(x)}.$$
对任意满足 $f_Y(y)>0$ 的 $y$，给定 $Y=y$ 时 $X$ 的条件 pdf 是关于 $x$ 的函数，记为 $f(x\mid y)$，定义为
$$f(x\mid y)=\frac{f(x,y)}{f_Y(y)}.$$
\end{definition}
验证 $f(x\mid y)$ 和 $f(y\mid x)$ 确实是 pdf，只需重复前面验证\autoref{def:4.2.1} 给出了真正 pmf 的步骤，并用积分代替求和.

条件 pdf 或 pmf 除了可用来计算概率，也可用来计算期望. 只要记住 $f(y\mid x)$ 作为 $y$ 的函数是 pdf 或 pmf，就可以像之前使用无条件 pdf 或 pmf 那样使用它. 若 $g(Y)$ 是 $Y$ 的函数，则给定 $X=x$ 时 $g(Y)$ 的条件期望（conditional expected value）记为 $E(g(Y)\mid x)$；在离散和连续情形下分别为\index{SI0425@条件}\semanticindexanchor{425:154}
$$E(g(Y)\mid x)=\sum_y g(y)f(y\mid x),\qquad E(g(Y)\mid x)=\int_{-\infty}^{\infty}g(y)f(y\mid x)\,dy.$$
条件期望具有\autoref{thm:2.2.5} 所列通常期望的全部性质. 此外，$E(Y\mid X)$ 给出了基于 $X$ 的知识对 $Y$ 的最佳猜测，这推广了\autoref{ex:2.2.6} 的结论（见\autoref{exercise:4.13}）.
\begin{example}[条件 pdf 的计算]\label{ex:4.2.4}
与\autoref{ex:4.1.12} 一样，设连续随机向量 $(X,Y)$ 的联合 pdf 为 $f(x,y)=e^{-y}$，$0<x<y<\infty$. 假定要计算给定 $X=x$ 时 $Y$ 的条件 pdf. $X$ 的边缘 pdf 可如下求得：如果 $x\le0$，则对所有 $y$ 均有 $f(x,y)=0$，所以 $f_X(x)=0$；如果 $x>0$，则 $f(x,y)>0$ 当且仅当 $y>x$. 因此，
$$f_X(x)=\int_{-\infty}^{\infty}f(x,y)\,dy=\int_x^\infty e^{-y}\,dy=e^{-x}.$$
于是，$X$ 的边缘分布是指数分布. 根据\autoref{def:4.2.3}，对任何 $x>0$ 都可以计算给定 $X=x$ 时 $Y$ 的条件分布，因为恰好对这些值有 $f_X(x)>0$. 对任意这样的 $x$，
$$f(y\mid x)=\frac{f(x,y)}{f_X(x)}=\frac{e^{-y}}{e^{-x}}=e^{-(y-x)},\qquad y>x,$$
而
$$f(y\mid x)=\frac{f(x,y)}{f_X(x)}=\frac0{e^{-x}}=0,\qquad y\le x.$$
\sourcepageanchor{155}

因此，给定 $X=x$ 时，$Y$ 服从指数分布，其中 $x$ 是位置参数，$\beta=1$ 是尺度参数. 每个不同的 $x$ 值都对应 $Y$ 的一个不同条件分布. 进而有
$$E(Y\mid X=x)=\int_x^\infty ye^{-(y-x)}\,dy=1+x.$$
由 $f(y\mid x)$ 描述的概率分布的方差，称为给定 $X=x$ 时 $Y$ 的条件方差（conditional variance）. 将其记为 $\operatorname{Var}(Y\mid x)$，根据通常的方差定义，有\index{SI1064@条件}\semanticindexanchor{1064:155}
$$\operatorname{Var}(Y\mid x)=E(Y^2\mid x)-(E(Y\mid x))^2.$$
将该定义用于本例，得到
$$\operatorname{Var}(Y\mid x)=\int_x^\infty y^2e^{-(y-x)}\,dy-\left(\int_x^\infty ye^{-(y-x)}\,dy\right)^2=1.$$
这里，给定 $X=x$ 时 $Y$ 的条件方差对所有 $x$ 都相同. 但在其他情形，它可能随 $x$ 而异. 可以将此条件方差与 $Y$ 的无条件方差比较. $Y$ 的边缘分布是 $\operatorname{gamma}(2,1)$，其方差为 $\operatorname{Var}Y=2$. 得知 $X=x$ 后，$Y$ 的变异程度显著减小.
\end{example}

\autoref{ex:4.2.4} 的模型可用于如下实际情形. 假设有两只灯泡，它们的发光寿命分别为随机变量 $X$ 和 $Z$. 寿命 $X$ 与 $Z$ 相互独立，且 pdf 均为 $e^{-x}$，$x>0$. 先点亮第一只灯泡；它一烧坏，便立即点亮第二只. 现在考虑观测第一只烧坏的时刻 $X$，以及第二只烧坏的时刻 $Y=X+Z$. 已知第一只在 $X=x$ 时烧坏且第二只开始发光，则 $Y=Z+x$. 这类似于\autoref{ex:3.5.3}，其中 $x$ 起位置参数的作用，而 $Y$ 的 pdf（这里是给定 $X=x$ 时的条件 pdf）为 $f(y\mid x)=f_Z(y-x)=e^{-(y-x)}$，$y>x$.

给定 $X=x$ 时 $Y$ 的条件分布，对不同的 $x$ 可能是不同的概率分布. 因而，实际上我们有一个 $Y$ 的概率分布族，每个 $x$ 对应一个分布. 要描述整个分布族时，就使用“$Y\mid X$ 的分布”这一说法. 例如，若 $X$ 是取正整数值的随机变量，给定 $X=x$ 时 $Y$ 的条件分布为 $\operatorname{binomial}(x,p)$，则可以说 $Y\mid X$ 的分布为 $\operatorname{binomial}(X,p)$，或写成 $Y\mid X\sim\operatorname{binomial}(X,p)$. 每当使用符号 $Y\mid X$，或以随机变量作为概率分布的参数时，我们描述的都是一个条件概率分布族. 有时通过指定条件分布 $f(y\mid x)$ 和边缘分布 $f_X(x)$ 来定义联合 pdf 或 pmf. 此时由定义有 $f(x,y)=f(y\mid x)f_X(x)$. \autoref{sec:4.4} 将进一步讨论这类模型.

还要注意，$E(g(Y)\mid x)$ 是 $x$ 的函数. 即对每个 $x$，$E(g(Y)\mid x)$ 是通过适当的积分或求和得到的一个实数. 所以 $E(g(Y)\mid X)$ 是随机变量，其值取决于 $X$ 的值. 当 $X=x$ 时，随机变量 $E(g(Y)\mid X)$ 的值就是 $E(g(Y)\mid x)$. 因此，在\autoref{ex:4.2.4} 中可以写成 $E(Y\mid X)=1+X$.

在前面所有例子中，给定 $X=x$ 时 $Y$ 的条件分布都随 $x$ 改变. 某些情形下，知道 $X=x$ 并不比原有知识提供更多有关 $Y$ 的信息. $X$ 与 $Y$ 的这一重要关系称为独立性（independence）. 和\autoref{sec:1} 中独立事件的情形一样，先以对称方式定义独立性，再推导刚刚提到的条件性质，会更方便. 下面就这样做.
\begin{definition}\label{def:4.2.5}\index{SI0554@随机变量}
设 $(X,Y)$ 为二元随机向量，其联合 pdf 或 pmf 为 $f(x,y)$，边缘 pdf 或 pmf 为 $f_X(x)$ 和 $f_Y(y)$. 若对每个 $x\in\mathbb R$ 和 $y\in\mathbb R$，均有
\begin{equation}\label{eq:4.2.1}
f(x,y)=f_X(x)f_Y(y),
\end{equation}
则称 $X$ 和 $Y$ 是相互独立的随机变量（independent random variables）.
\end{definition}
\sourcepageanchor{156}

如果 $X$ 和 $Y$ 相互独立，则给定 $X=x$ 时 $Y$ 的条件 pdf 为
$$
\begin{aligned}
f(y\mid x)&=\frac{f(x,y)}{f_X(x)}&&\text{（定义）}\\
&=\frac{f_X(x)f_Y(y)}{f_X(x)}&&\text{（由\eqref{eq:4.2.1}）}\\
&=f_Y(y),
\end{aligned}
$$
与 $x$ 的值无关. 因此，对任意 $A\subset\mathbb R$ 和 $x\in\mathbb R$，$P(Y\in A\mid x)=\int_A f(y\mid x)\,dy=\int_A f_Y(y)\,dy=P(Y\in A)$. 知道 $X=x$ 并未提供关于 $Y$ 的额外信息.\footnote{译注：依本节条件密度的定义，上述比值及由它得到的结论限于 $f_X(x)>0$ 的 $x$；原文“任意 $x\in\mathbb R$”应作此限制.}

\autoref{def:4.2.5} 有两种用法. 一种是从联合 pdf 或 pmf 出发，检查 $X$ 与 $Y$ 是否独立. 为此，必须验证\eqref{eq:4.2.1} 对所有 $x,y$ 都成立. 另一种是希望定义一个 $X$ 和 $Y$ 独立的模型. 根据 $X$ 和 $Y$ 的实际含义，可能认为知道 $X=x$ 不应提供任何有关 $Y$ 的信息. 此时，可以指定 $X$ 和 $Y$ 的边缘分布，再按\eqref{eq:4.2.1} 将联合分布定义为它们的乘积.
\begin{example}[独立性的检验——I]\label{ex:4.2.6}
考虑离散二元随机向量 $(X,Y)$，其联合 pmf 为
$$
\begin{gathered}
f(10,1)=f(20,1)=f(20,2)=\frac1{10},\\
f(10,2)=f(10,3)=\frac15,\qquad f(20,3)=\frac3{10}.
\end{gathered}
$$
容易算出边缘 pmf：
$$f_X(10)=f_X(20)=\frac12,\qquad f_Y(1)=\frac15,\quad f_Y(2)=\frac3{10},\quad f_Y(3)=\frac12.$$
$X$ 和 $Y$ 不独立，因为\eqref{eq:4.2.1} 并非对所有 $x,y$ 都成立. 例如，
$$f(10,3)=\frac15\ne\frac12\frac12=f_X(10)f_Y(3).$$
若 $X$ 和 $Y$ 独立，则\eqref{eq:4.2.1} 必须对每一组 $x,y$ 都成立. 注意，$f(10,1)=1/10=(1/2)(1/5)=f_X(10)f_Y(1)$. 仅对某些 $x,y$ 满足\eqref{eq:4.2.1}，并不能保证独立；必须检查所有取值.
\end{example}

直接用\eqref{eq:4.2.1} 验证独立性，需要先知道 $f_X(x)$ 和 $f_Y(y)$. 下面的引理可使验证稍为容易些.
\begin{lemma}\label{lem:4.2.7}
设 $(X,Y)$ 为联合 pdf 或 pmf 为 $f(x,y)$ 的二元随机向量. 则 $X$ 和 $Y$ 独立，当且仅当存在函数 $g(x)$ 和 $h(y)$，使对每个 $x\in\mathbb R$ 和 $y\in\mathbb R$，都有
$$f(x,y)=g(x)h(y).$$
\end{lemma}
\begin{proof}
必要性由取 $g(x)=f_X(x)$、$h(y)=f_Y(y)$ 并利用\eqref{eq:4.2.1} 即得. 对连续随机变量，证明充分性时，假设 $f(x,y)=g(x)h(y)$. 定义
$$\int_{-\infty}^{\infty}g(x)\,dx=c,\qquad\int_{-\infty}^{\infty}h(y)\,dy=d.$$
常数 $c$ 和 $d$ 满足
\begin{equation}\label{eq:4.2.2}
\begin{aligned}
cd&=\left(\int_{-\infty}^{\infty}g(x)\,dx\right)\left(\int_{-\infty}^{\infty}h(y)\,dy\right)\\
&=\int_{-\infty}^{\infty}\int_{-\infty}^{\infty}g(x)h(y)\,dx\,dy\\
&=\int_{-\infty}^{\infty}\int_{-\infty}^{\infty}f(x,y)\,dx\,dy\\
&=1.\qquad\text{（$f(x,y)$ 是联合 pdf）}
\end{aligned}
\end{equation}
\sourcepageanchor{157}

此外，边缘 pdf 为
\begin{equation}\label{eq:4.2.3}
\begin{aligned}
f_X(x)&=\int_{-\infty}^{\infty}g(x)h(y)\,dy=g(x)d,\\
f_Y(y)&=\int_{-\infty}^{\infty}g(x)h(y)\,dx=h(y)c.
\end{aligned}
\end{equation}
因此，由\eqref{eq:4.2.2} 和\eqref{eq:4.2.3}，
$$f(x,y)=g(x)h(y)=g(x)h(y)cd=f_X(x)f_Y(y),$$
可见 $X$ 和 $Y$ 独立. 将积分换为求和，即得离散随机向量情形的证明.
\end{proof}
\begin{example}[独立性的检验——II]\label{ex:4.2.8}
$f(x,y)=\frac1{384}x^2y^4e^{-y-(x/2)}$，$x>0$ 且 $y>0$，是一个联合 pdf. 如果定义
$$g(x)=\begin{cases}x^2e^{-x/2},&x>0,\\0,&x\le0,\end{cases}\qquad
h(y)=\begin{cases}y^4e^{-y}/384,&y>0,\\0,&y\le0,\end{cases}$$
则对所有 $x\in\mathbb R$ 和所有 $y\in\mathbb R$，都有 $f(x,y)=g(x)h(y)$. 由\autoref{lem:4.2.7}，可断定 $X$ 与 $Y$ 独立，无须计算边缘 pdf.
\end{example}

若 $X$ 与 $Y$ 独立，则由\eqref{eq:4.2.1} 可见，$f(x,y)>0$ 的集合为 $\{(x,y):x\in A\text{ 且 }y\in B\}$，其中 $A=\{x:f_X(x)>0\}$，$B=\{y:f_Y(y)>0\}$. 这种集合称为笛卡尔积（cross-product），通常记为 $A\times B$. 判断一个点是否属于笛卡尔积，可以分别考察其 $x$ 值和 $y$ 值. 若 $f(x,y)$ 为联合 pdf 或 pmf，而使 $f(x,y)>0$ 的集合不是笛卡尔积，则具有该联合 pdf 或 pmf 的随机变量 $X$ 和 $Y$ 不独立. 在\autoref{ex:4.2.4} 中，集合 $0<x<y<\infty$ 不是笛卡尔积. 要判断成员资格，不仅要检查 $0<x<\infty$ 和 $0<y<\infty$，还必须检查 $x<y$. 因此，\autoref{ex:4.2.4} 中的随机变量不独立. \autoref{ex:4.2.2} 也给出了一个例子：联合 pmf 为正的集合不是笛卡尔积.\index{SI0236@直积集}\semanticindexanchor{236:157}
\begin{example}[联合概率模型]\label{ex:4.2.9}
作为利用独立性定义联合概率模型的例子，考虑如下情形. 从堪萨斯城一所小学中随机抽取一名学生，记录 $X=$ 该学生仍在世的父母人数. 假设 $X$ 的边缘分布为
$$f_X(0)=.01,\qquad f_X(1)=.09,\qquad f_X(2)=.90.$$
从太阳城随机抽取一位退休者，记录 $Y=$ 该退休者仍在世的父母人数. 假设 $Y$ 的边缘分布为
$$f_Y(0)=.70,\qquad f_Y(1)=.25,\qquad f_Y(2)=.05.$$
假定这两个随机变量独立似乎是合理的. 知道该学生的父母人数，并不能告诉我们该退休者的父母人数. 唯一能反映这种独立性的 $X,Y$ 联合分布，就是由\eqref{eq:4.2.1} 定义的分布. 例如，
$$f(0,0)=f_X(0)f_Y(0)=.0070,\qquad f(0,1)=f_X(0)f_Y(1)=.0025.$$
此联合分布可用于计算诸如
$$
\begin{aligned}
P(X=Y)&=f(0,0)+f(1,1)+f(2,2)\\
&=(.01)(.70)+(.09)(.25)+(.90)(.05)=.0745
\end{aligned}
$$
这样的量.
\end{example}
\sourcepageanchor{158}

如果 $X$ 与 $Y$ 独立，有些概率和期望很容易计算，如下述定理所示.
\begin{theorem}\label{thm:4.2.10}\index{SI0427@独立性}\index{SI0553@期望值}
设 $X$ 和 $Y$ 为相互独立的随机变量.
\begin{enumerate}
\item[(a)]\phantomsection\label{itm:thm:4.2.10:a} 对任意 $A\subset\mathbb R$ 和 $B\subset\mathbb R$，$P(X\in A,Y\in B)=P(X\in A)P(Y\in B)$；即事件 $\{X\in A\}$ 和 $\{Y\in B\}$ 相互独立.
\item[(b)]\phantomsection\label{itm:thm:4.2.10:b} 设 $g(x)$ 仅为 $x$ 的函数，$h(y)$ 仅为 $y$ 的函数，则
$$E(g(X)h(Y))=(Eg(X))(Eh(Y)).$$
\end{enumerate}
\end{theorem}
\begin{proof}
对于连续随机变量，\hyperref[itm:thm:4.2.10:b]{(b)} 由下式得到：
$$
\begin{aligned}
E(g(X)h(Y))&=\int_{-\infty}^{\infty}\int_{-\infty}^{\infty}g(x)h(y)f(x,y)\,dx\,dy\\
&=\int_{-\infty}^{\infty}\int_{-\infty}^{\infty}g(x)h(y)f_X(x)f_Y(y)\,dx\,dy\quad\text{（由\eqref{eq:4.2.1}）}\\
&=\int_{-\infty}^{\infty}h(y)f_Y(y)\int_{-\infty}^{\infty}g(x)f_X(x)\,dx\,dy\\
&=\left(\int_{-\infty}^{\infty}g(x)f_X(x)\,dx\right)\left(\int_{-\infty}^{\infty}h(y)f_Y(y)\,dy\right)\\
&=(Eg(X))(Eh(Y)).
\end{aligned}
$$
将积分换为求和，即得离散随机变量的结论. \hyperref[itm:thm:4.2.10:a]{(a)} 可以用类似步骤证明，也可作如下论证. 令 $g(x)$ 为集合 $A$ 的示性函数，$h(y)$ 为集合 $B$ 的示性函数. 注意，$g(x)h(y)$ 是集合 $C=\{(x,y):x\in A,y\in B\}\subset\mathbb R^2$ 的示性函数. 还要注意，对 $g(x)$ 这样的示性函数，有 $Eg(X)=P(X\in A)$. 因此，利用刚刚证明的期望等式，
$$
\begin{aligned}
P(X\in A,Y\in B)&=P((X,Y)\in C)=E(g(X)h(Y))\\
&=(Eg(X))(Eh(Y))=P(X\in A)P(Y\in B).
\end{aligned}
$$
\end{proof}
\begin{example}[独立变量的期望]\label{ex:4.2.11}
设 $X$ 和 $Y$ 为独立的 $\operatorname{exponential}(1)$ 随机变量. 由\autoref{thm:4.2.10}，
$$P(X\ge4,Y<3)=P(X\ge4)P(Y<3)=e^{-4}(1-e^{-3}).$$
令 $g(x)=x^2$，$h(y)=y$，得到
$$E(X^2Y)=(EX^2)(EY)=\bigl(\operatorname{Var}X+(EX)^2\bigr)EY=(1+1^2)1=2.$$
\end{example}
下面关于独立随机变量之和的结论，是\autoref{thm:4.2.10} 的简单推论.
\begin{theorem}\label{thm:4.2.12}\index{SI0555@和}\index{SI0692@独立随机变量}
设 $X$ 和 $Y$ 是相互独立的随机变量，矩母函数分别为 $M_X(t)$ 和 $M_Y(t)$. 则随机变量 $Z=X+Y$ 的矩母函数为
$$M_Z(t)=M_X(t)M_Y(t).$$
\end{theorem}
\begin{proof}
根据 mgf 的定义及\autoref{thm:4.2.10}，
$$M_Z(t)=Ee^{tZ}=Ee^{t(X+Y)}=E(e^{tX}e^{tY})=(Ee^{tX})(Ee^{tY})=M_X(t)M_Y(t).$$
\end{proof}
\sourcepageanchor{159}

\begin{example}[正态变量之和的 mgf]\label{ex:4.2.13}\index{SI0369@和}
有时利用\autoref{thm:4.2.12}，可以从 $X$ 和 $Y$ 的分布轻易推出 $Z$ 的分布. 例如，设 $X\sim n(\mu,\sigma^2)$ 与 $Y\sim n(\gamma,\tau^2)$ 是独立的正态随机变量. 由\autoref{exercise:2.33}，它们的 mgf 为
$$M_X(t)=\exp(\mu t+\sigma^2t^2/2),\qquad M_Y(t)=\exp(\gamma t+\tau^2t^2/2).$$
所以根据\autoref{thm:4.2.12}，$Z=X+Y$ 的 mgf 为
$$M_Z(t)=M_X(t)M_Y(t)=\exp\bigl((\mu+\gamma)t+(\sigma^2+\tau^2)t^2/2\bigr).$$
这是均值为 $\mu+\gamma$、方差为 $\sigma^2+\tau^2$ 的正态随机变量的 mgf. 这个结论很重要，值得列为定理.
\end{example}
\begin{theorem}\label{thm:4.2.14}
设 $X\sim n(\mu,\sigma^2)$ 与 $Y\sim n(\gamma,\tau^2)$ 是独立的正态随机变量. 则随机变量 $Z=X+Y$ 服从 $n(\mu+\gamma,\sigma^2+\tau^2)$ 分布.
\end{theorem}

若 $f(x,y)$ 是连续随机向量 $(X,Y)$ 的联合 pdf，\eqref{eq:4.2.1} 可能在某个满足 $\iint_A dx\,dy=0$ 的集合 $A$ 上不成立. 即便如此，仍称 $X$ 和 $Y$ 独立. 这是因为，仅在这样的集合 $A$ 上不同的两个 pdf，定义了同一个 $(X,Y)$ 概率分布. 为说明这一点，设两个 pdf $f(x,y)$ 和 $f^*(x,y)$ 除了在满足 $\iint_A dx\,dy=0$ 的集合 $A$ 上以外，处处相等. 设 $(X,Y)$ 的 pdf 为 $f(x,y)$，$(X^*,Y^*)$ 的 pdf 为 $f^*(x,y)$，$B$ 为 $\mathbb R^2$ 的任意子集. 则
$$
\begin{aligned}
P((X,Y)\in B)&=\iint_B f(x,y)\,dx\,dy=\iint_{B\cap A^c}f(x,y)\,dx\,dy\\
&=\iint_{B\cap A^c}f^*(x,y)\,dx\,dy=\iint_B f^*(x,y)\,dx\,dy\\
&=P((X^*,Y^*)\in B).
\end{aligned}
$$
因此，$(X,Y)$ 和 $(X^*,Y^*)$ 有相同的概率分布. 例如，$f(x,y)=e^{-x-y}$，$x>0$ 且 $y>0$，是两个独立指数随机变量的 pdf，并满足\eqref{eq:4.2.1}. 但令 $f^*(x,y)$ 除在 $x=y$ 时取 0 外都等于 $f(x,y)$，所得 $f^*$ 仍是两个独立指数随机变量的 pdf，尽管\eqref{eq:4.2.1} 在集合 $A=\{(x,x):x>0\}$ 上不成立.



\subsection{二元变换}\label{sec:4.3}
\autoref{sec:2.1} 讨论了求随机变量的函数的分布的方法. 本节将这些思想推广到二元随机向量.

设 $(X,Y)$ 为概率分布已知的二元随机向量. 考虑由 $U=g_1(X,Y)$ 和 $V=g_2(X,Y)$ 定义的新二元随机向量 $(U,V)$，其中 $g_1(x,y)$ 和 $g_2(x,y)$ 为给定函数. 若 $B$ 是 $\mathbb R^2$ 的任意子集，则 $(U,V)\in B$ 当且仅当 $(X,Y)\in A$，其中 $A=\{(x,y):(g_1(x,y),g_2(x,y))\in B\}$. 因此，$P((U,V)\in B)=P((X,Y)\in A)$，$(U,V)$ 的概率分布完全由 $(X,Y)$ 的概率分布确定.

若 $(X,Y)$ 是离散二元随机向量，其联合 pmf 仅在一个可数集合上为正. 将这个集合记为 $\mathscr A$. 定义 $\mathscr B=\{(u,v):\text{存在 }(x,y)\in\mathscr A,\ u=g_1(x,y)\text{ 且 }v=g_2(x,y)\}$. 则 $\mathscr B$ 是离散随机向量 $(U,V)$ 的可能取值组成的可数集合. 对任意 $(u,v)\in\mathscr B$，定义 $A_{uv}=\{(x,y)\in\mathscr A:g_1(x,y)=u\text{ 且 }g_2(x,y)=v\}$，则 $(U,V)$ 的联合 pmf 可由 $(X,Y)$ 的联合 pmf 计算：
\begin{equation}\label{eq:4.3.1}
f_{U,V}(u,v)=P(U=u,V=v)=P((X,Y)\in A_{uv})=\sum_{(x,y)\in A_{uv}}f_{X,Y}(x,y).
\end{equation}
\sourcepageanchor{160}

\begin{example}[Poisson 变量之和的分布]\label{ex:4.3.1}
设 $X$ 和 $Y$ 是相互独立的 Poisson 随机变量，参数分别为 $\theta$ 和 $\lambda$. 因此，$(X,Y)$ 的联合 pmf 为
$$f_{X,Y}(x,y)=\frac{\theta^xe^{-\theta}}{x!}\frac{\lambda^ye^{-\lambda}}{y!},\qquad x=0,1,2,\ldots,\ y=0,1,2,\ldots.$$
集合 $\mathscr A=\{(x,y):x=0,1,2,\ldots\text{ 且 }y=0,1,2,\ldots\}$. 现在定义 $U=X+Y$ 和 $V=Y$，即 $g_1(x,y)=x+y$、$g_2(x,y)=y$. 下面描述所有可能 $(u,v)$ 值组成的集合 $\mathscr B$. $v$ 的可能值为非负整数，因为 $v=y$，两者的可能取值集合相同. 给定 $v$ 后，由于 $x$ 是非负整数，$u=x+y=x+v$ 必须是大于或等于 $v$ 的整数. 因此，所有可能 $(u,v)$ 值的集合为 $\mathscr B=\{(u,v):v=0,1,2,\ldots\text{ 且 }u=v,v+1,v+2,\ldots\}$. 对任意 $(u,v)\in\mathscr B$，唯一满足 $x+y=u$、$y=v$ 的取值是 $x=u-v$、$y=v$. 所以本例中 $A_{uv}$ 总是只包含单点 $(u-v,v)$. 根据\eqref{eq:4.3.1}，得到
$$f_{U,V}(u,v)=f_{X,Y}(u-v,v)=\frac{\theta^{u-v}e^{-\theta}}{(u-v)!}\frac{\lambda^ve^{-\lambda}}{v!},\quad
\begin{matrix}v=0,1,2,\ldots,\\u=v,v+1,v+2,\ldots.\end{matrix}$$
本例中求 $U$ 的边缘 pmf 很有意思. 对固定的非负整数 $u$，$f_{U,V}(u,v)>0$ 仅在 $v=0,1,\ldots,u$ 时成立，这给出了求边缘 pmf 时对 $v$ 的求和范围. 得到
$$
\begin{aligned}
f_U(u)&=\sum_{v=0}^u\frac{\theta^{u-v}e^{-\theta}}{(u-v)!}\frac{\lambda^ve^{-\lambda}}{v!}\\
&=e^{-(\theta+\lambda)}\sum_{v=0}^u\frac{\theta^{u-v}}{(u-v)!}\frac{\lambda^v}{v!},\qquad u=0,1,2,\ldots.
\end{aligned}
$$
将每项乘除以 $u!$，再利用二项式定理，可以化简为
$$f_U(u)=\frac{e^{-(\theta+\lambda)}}{u!}\sum_{v=0}^u\binom uv\lambda^v\theta^{u-v}
=\frac{e^{-(\theta+\lambda)}}{u!}(\theta+\lambda)^u,\quad u=0,1,2,\ldots.$$
这是参数为 $\theta+\lambda$ 的 Poisson 随机变量的 pmf. 此结论很重要，值得列为定理.
\end{example}
\begin{theorem}\label{thm:4.3.2}\index{SI0379@和}
若 $X\sim\operatorname{Poisson}(\theta)$，$Y\sim\operatorname{Poisson}(\lambda)$，且 $X$ 与 $Y$ 独立，则 $X+Y\sim\operatorname{Poisson}(\theta+\lambda)$.
\end{theorem}

若 $(X,Y)$ 是联合 pdf 为 $f_{X,Y}(x,y)$ 的连续随机向量，则可用类似\eqref{eq:2.1.8} 的方式，通过 $f_{X,Y}(x,y)$ 表示 $(U,V)$ 的联合 pdf. 与前面一样，令 $\mathscr A=\{(x,y):f_{X,Y}(x,y)>0\}$，$\mathscr B=\{(u,v):\text{存在 }(x,y)\in\mathscr A,\ u=g_1(x,y)\text{ 且 }v=g_2(x,y)\}$. 联合 pdf $f_{U,V}(u,v)$ 在 $\mathscr B$ 上为正. 在此结论最简单的版本中，假设 $u=g_1(x,y)$、$v=g_2(x,y)$ 定义了从 $\mathscr A$ 到 $\mathscr B$ 的一一变换. 由 $\mathscr B$ 的定义，该变换是满射. 我们假定，对每个 $(u,v)\in\mathscr B$，只有一个 $(x,y)\in\mathscr A$ 满足 $(u,v)=(g_1(x,y),g_2(x,y))$. 对这样的一一满射，可以解方程 $u=g_1(x,y)$、$v=g_2(x,y)$，用 $u,v$ 表示 $x,y$. 将逆变换记为 $x=h_1(u,v)$、$y=h_2(u,v)$. 在单变量情形中由导数发挥的作用，现在由称为变换的 Jacobian（雅可比行列式）的量来发挥. 这是一个关于 $(u,v)$ 的函数，记为 $J$，是偏导数矩阵的行列式，定义为\index{SI0248@行列式}\semanticindexanchor{248:160}\index{SI0593@雅可比行列式}\semanticindexanchor{593:160}\index{SI0729@一一变换}\semanticindexanchor{729:160}\index{SI0730@满射变换}\semanticindexanchor{730:160}
$$J=\begin{vmatrix}\dfrac{\partial x}{\partial u}&\dfrac{\partial x}{\partial v}\\[4pt]\dfrac{\partial y}{\partial u}&\dfrac{\partial y}{\partial v}\end{vmatrix}
=\frac{\partial x}{\partial u}\frac{\partial y}{\partial v}-\frac{\partial y}{\partial u}\frac{\partial x}{\partial v},$$
\sourcepageanchor{161}

其中
$$
\frac{\partial x}{\partial u}=\frac{\partial h_1(u,v)}{\partial u},\quad
\frac{\partial x}{\partial v}=\frac{\partial h_1(u,v)}{\partial v},\quad
\frac{\partial y}{\partial u}=\frac{\partial h_2(u,v)}{\partial u},\quad
\frac{\partial y}{\partial v}=\frac{\partial h_2(u,v)}{\partial v}.
$$
假设 $J$ 在 $\mathscr B$ 上不恒为零. 则 $(U,V)$ 的联合 pdf 在 $\mathscr B$ 外为 0，在 $\mathscr B$ 上为
\begin{equation}\label{eq:4.3.2}\index{SI0829@二元变换}
f_{U,V}(u,v)=f_{X,Y}(h_1(u,v),h_2(u,v))|J|,
\end{equation}
其中 $|J|$ 是 $J$ 的绝对值. 使用\eqref{eq:4.3.2} 时，确定 $\mathscr B$ 并验证变换为一一变换，有时与代入公式本身同样困难. 阅读下面例子的说明时，请留意这些部分.
\begin{example}[Beta 变量乘积的分布]\label{ex:4.3.3}\index{SI0256@乘积}
设 $X\sim\operatorname{beta}(\alpha,\beta)$ 和 $Y\sim\operatorname{beta}(\alpha+\beta,\gamma)$ 相互独立. $(X,Y)$ 的联合 pdf 为
$$
\begin{aligned}
f_{X,Y}(x,y)&=\frac{\Gamma(\alpha+\beta)}{\Gamma(\alpha)\Gamma(\beta)}x^{\alpha-1}(1-x)^{\beta-1}
\frac{\Gamma(\alpha+\beta+\gamma)}{\Gamma(\alpha+\beta)\Gamma(\gamma)}y^{\alpha+\beta-1}(1-y)^{\gamma-1},\\
&\hspace{7cm}0<x<1,\quad0<y<1.
\end{aligned}
$$
考虑变换 $U=XY$、$V=X$. 因为 $V=X$，所以 $V$ 的可能取值为 $0<v<1$. 固定 $V=v$ 后，由于 $X=V=v$ 且 $Y$ 介于 0 与 1 之间，$U$ 必须介于 0 与 $v$ 之间. 因而该变换将 $\mathscr A$ 映满到 $\mathscr B=\{(u,v):0<u<v<1\}$. 对每个 $(u,v)\in\mathscr B$，由 $u=xy$、$v=x$ 可唯一解出 $x=h_1(u,v)=v$、$y=h_2(u,v)=u/v$. 注意，若将此变换视为定义在整个 $\mathbb R^2$ 上的变换，它并不是一一的，因为每个点 $(0,y)$ 都映到 $(0,0)$. 但作为仅定义在 $\mathscr A$ 上的函数，它是到 $\mathscr B$ 的一一满射. Jacobian 为
$$J=\begin{vmatrix}\dfrac{\partial x}{\partial u}&\dfrac{\partial x}{\partial v}\\[4pt]\dfrac{\partial y}{\partial u}&\dfrac{\partial y}{\partial v}\end{vmatrix}
=\begin{vmatrix}0&1\\[2pt]1/v&-u/v^2\end{vmatrix}=-\frac1v.$$
于是，根据\eqref{eq:4.3.2} 得到联合 pdf
\begin{equation}\label{eq:4.3.3}
\begin{aligned}
f_{U,V}(u,v)&=\frac{\Gamma(\alpha+\beta+\gamma)}{\Gamma(\alpha)\Gamma(\beta)\Gamma(\gamma)}
v^{\alpha-1}(1-v)^{\beta-1}\left(\frac uv\right)^{\alpha+\beta-1}\left(1-\frac uv\right)^{\gamma-1}\frac1v,\\
&\hspace{7cm}0<u<v<1.
\end{aligned}
\end{equation}
$V=X$ 的边缘分布当然是 $\operatorname{beta}(\alpha,\beta)$，但 $U$ 的分布也是 Beta 分布：
$$
\begin{aligned}
f_U(u)&=\int_u^1 f_{U,V}(u,v)\,dv\\
&=\frac{\Gamma(\alpha+\beta+\gamma)}{\Gamma(\alpha)\Gamma(\beta)\Gamma(\gamma)}u^{\alpha-1}
\int_u^1\left(\frac uv-u\right)^{\beta-1}\left(1-\frac uv\right)^{\gamma-1}\frac u{v^2}\,dv.
\end{aligned}
$$
\sourcepageanchor{162}

这里使用了\eqref{eq:4.3.3}，并重新排列了一些因子. 现在作一元变量替换 $y=(u/v-u)/(1-u)$，于是 $dy=-u/[v^2(1-u)]\,dv$，得到
$$
\begin{aligned}
f_U(u)&=\frac{\Gamma(\alpha+\beta+\gamma)}{\Gamma(\alpha)\Gamma(\beta)\Gamma(\gamma)}u^{\alpha-1}(1-u)^{\beta+\gamma-1}
\int_0^1 y^{\beta-1}(1-y)^{\gamma-1}\,dy\\
&=\frac{\Gamma(\alpha+\beta+\gamma)}{\Gamma(\alpha)\Gamma(\beta)\Gamma(\gamma)}u^{\alpha-1}(1-u)^{\beta+\gamma-1}
\frac{\Gamma(\beta)\Gamma(\gamma)}{\Gamma(\beta+\gamma)}\\
&=\frac{\Gamma(\alpha+\beta+\gamma)}{\Gamma(\alpha)\Gamma(\beta+\gamma)}u^{\alpha-1}(1-u)^{\beta+\gamma-1},\qquad0<u<1.
\end{aligned}
$$
第二个等号利用了被积函数是 Beta pdf 的核，并使用了\hyperref[eq:3.3.17]{式 (3.2.17)}.\footnote{译注：原文式号 (3.2.17) 不存在；按所用的 Beta--Gamma 恒等式，应指\eqref{eq:3.3.17}.} 因此，$U$ 的边缘分布是 $\operatorname{beta}(\alpha,\beta+\gamma)$.
\end{example}
\begin{example}[正态变量的和与差]\label{ex:4.3.4}\index{SI0369@和}
设 $X$ 和 $Y$ 为独立的标准正态随机变量. 考虑变换 $U=X+Y$、$V=X-Y$. 沿用上面的记号，$U=g_1(X,Y)$，其中 $g_1(x,y)=x+y$；$V=g_2(X,Y)$，其中 $g_2(x,y)=x-y$. $X,Y$ 的联合 pdf 为 $f_{X,Y}(x,y)=(2\pi)^{-1}\exp(-x^2/2)\exp(-y^2/2)$，$-\infty<x<\infty$，$-\infty<y<\infty$. 因而 $\mathscr A=\mathbb R^2$. 要确定使 $f_{U,V}(u,v)$ 为正的集合 $\mathscr B$，必须求出当 $(x,y)$ 遍及 $\mathscr A=\mathbb R^2$ 时，
\begin{equation}\label{eq:4.3.4}
u=x+y,\qquad v=x-y
\end{equation}
能取到的全部值. 然而，可以将 $u$ 和 $v$ 分别取为任意实数，再由\eqref{eq:4.3.4} 唯一解出
\begin{equation}\label{eq:4.3.5}
x=h_1(u,v)=\frac{u+v}{2},\qquad y=h_2(u,v)=\frac{u-v}{2}.
\end{equation}
这说明两件事. 第一，对任何 $(u,v)\in\mathbb R^2$，都存在由\eqref{eq:4.3.5} 定义的 $(x,y)\in\mathscr A$，使 $u=x+y$、$v=x-y$；因此全部可能 $(u,v)$ 的集合 $\mathscr B$ 就是 $\mathbb R^2$. 第二，由于\eqref{eq:4.3.5} 的解唯一，所考虑的变换是一一的. 只有其中给出的 $(x,y)$ 才能产生 $u=x+y$ 和 $v=x-y$. 根据\eqref{eq:4.3.5} 很容易计算 $x,y$ 的偏导数，得到
$$J=\begin{vmatrix}\dfrac{\partial x}{\partial u}&\dfrac{\partial x}{\partial v}\\[4pt]\dfrac{\partial y}{\partial u}&\dfrac{\partial y}{\partial v}\end{vmatrix}
=\begin{vmatrix}1/2&1/2\\1/2&-1/2\end{vmatrix}=-\frac12.$$
将\eqref{eq:4.3.5} 中 $x,y$ 的表达式代入 $f_{X,Y}(x,y)$，并利用 $|J|=1/2$，由\eqref{eq:4.3.2} 得到
$$f_{U,V}(u,v)=f_{X,Y}(h_1(u,v),h_2(u,v))|J|=\frac1{2\pi}e^{-((u+v)/2)^2/2}e^{-((u-v)/2)^2/2}\frac12,$$
其中 $-\infty<u<\infty$，$-\infty<v<\infty$. 展开指数中的平方，可见含 $uv$ 的项相互抵消. 化简并重新排列后得到
$$f_{U,V}(u,v)=\left(\frac1{\sqrt{2\pi}\sqrt2}e^{-u^2/4}\right)\left(\frac1{\sqrt{2\pi}\sqrt2}e^{-v^2/4}\right).$$
联合 pdf 分解成了只含 $u$ 的函数与只含 $v$ 的函数之积. 由\autoref{lem:4.2.7}，$U$ 与 $V$ 独立. 由\autoref{thm:4.2.14}，$U=X+Y$ 的边缘分布是 $n(0,2)$. 同理，可用\autoref{thm:4.2.12} 求得 $V$ 的边缘分布也是 $n(0,2)$. 独立正态随机变量的和与差仍为独立正态随机变量，这一重要事实不依赖于 $X,Y$ 的均值，只要求 $\operatorname{Var}X=\operatorname{Var}Y$. 这一结论留作\autoref{exercise:4.27}. \autoref{thm:4.2.12} 和\autoref{thm:4.2.14} 给出了 $U,V$ 的边缘分布；但要确定二者独立，则需要这里更细致的分析.
\end{example}
\sourcepageanchor{163}

在\autoref{ex:4.3.4} 中，我们发现 $U$ 与 $V$ 独立. 对由原变量 $X,Y$ 定义的新变量 $U,V$，还有一种简单得多却十分重要的独立情形，见下述定理.
\begin{theorem}\label{thm:4.3.5}\index{SI0554@随机变量}
设 $X$ 和 $Y$ 为独立随机变量，$g(x)$ 仅为 $x$ 的函数，$h(y)$ 仅为 $y$ 的函数. 则随机变量 $U=g(X)$ 和 $V=h(Y)$ 相互独立.
\end{theorem}
\begin{proof}
在假设 $U$ 和 $V$ 为连续随机变量的条件下证明. 对任意 $u\in\mathbb R$ 和 $v\in\mathbb R$，定义
$$A_u=\{x:g(x)\le u\},\qquad B_v=\{y:h(y)\le v\}.$$
则 $(U,V)$ 的联合 cdf 为
$$
\begin{aligned}
F_{U,V}(u,v)&=P(U\le u,V\le v)&&\text{（cdf 的定义）}\\
&=P(X\in A_u,Y\in B_v)&&\text{（$U,V$ 的定义）}\\
&=P(X\in A_u)P(Y\in B_v).&&\text{（\autoref{thm:4.2.10}）}
\end{aligned}
$$
$(U,V)$ 的联合 pdf 为
$$
\begin{aligned}
f_{U,V}(u,v)&=\frac{\partial^2}{\partial u\partial v}F_{U,V}(u,v)&&\text{（\eqref{eq:4.1.3}）}\\
&=\left(\frac d{du}P(X\in A_u)\right)\left(\frac d{dv}P(Y\in B_v)\right).
\end{aligned}
$$
正如记号所表明的，第一个因子只依赖 $u$，第二个只依赖 $v$. 所以由\autoref{lem:4.2.7}，$U$ 和 $V$ 独立.
\end{proof}

有时只关心一个函数，例如 $U=g_1(X,Y)$. 这时仍可利用上述方法求 $U$ 的分布. 如果能够选择另一个方便的函数 $V=g_2(X,Y)$，使从 $(X,Y)$ 到 $(U,V)$ 的变换在 $\mathscr A$ 上为一一变换，就能用\eqref{eq:4.3.2} 求出 $(U,V)$ 的联合 pdf，再由它求出 $U$ 的边缘 pdf. 在前面的例子中，我们也许只对 $U=XY$ 感兴趣. 可以选择 $V=X$，因为这样得到的变换在 $\mathscr A$ 上是一一的，然后像例中那样求出 $U$ 的边缘 pdf. 其他选择，如 $V=Y$，也同样可行（见\autoref{exercise:4.23}）.\footnote{译注：此段的 $U=XY$ 对应\autoref{ex:4.3.3}，并非紧接在前的\autoref{ex:4.3.4}；原文泛称“前面的例子”.}

当然，在许多情形中，所关心的变换并不是一一的. 正如\autoref{thm:2.1.8} 将一元方法推广到多对一函数，这里也可作同样推广. 仍令 $\mathscr A=\{(x,y):f_{X,Y}(x,y)>0\}$. 假设 $A_0,A_1,\ldots,A_k$ 是 $\mathscr A$ 的一个划分，满足以下性质. 集合 $A_0$ 可以为空，并满足 $P((X,Y)\in A_0)=0$. 对每个 $i=1,2,\ldots,k$，变换 $U=g_1(X,Y)$、$V=g_2(X,Y)$ 从 $A_i$ 一一映满到 $\mathscr B$. 那么，对每个 $i$，都可以求出从 $\mathscr B$ 到 $A_i$ 的逆函数. 记第 $i$ 个逆变换为 $x=h_{1i}(u,v)$、$y=h_{2i}(u,v)$. 对 $(u,v)\in\mathscr B$，该逆变换给出唯一的 $(x,y)\in A_i$，使 $(u,v)=(g_1(x,y),g_2(x,y))$. 以 $J_i$ 表示根据第 $i$ 个逆变换计算的 Jacobian. 假设这些 Jacobian 在 $\mathscr B$ 上都不恒为零，则联合 pdf $f_{U,V}$ 可表示为
\begin{equation}\label{eq:4.3.6}
f_{U,V}(u,v)=\sum_{i=1}^k f_{X,Y}(h_{1i}(u,v),h_{2i}(u,v))|J_i|.
\end{equation}
\begin{example}[正态变量之比的分布]\label{ex:4.3.6}\index{SI0279@与正态的关系}\index{SI0363@比值}\index{SI0365@与柯西的关系}
设 $X$ 和 $Y$ 为独立的 $n(0,1)$ 随机变量. 考虑变换 $U=X/Y$ 和 $V=|Y|$. （若 $Y=0$，可将 $(U,V)$ 任意定义，例如取为 $(1,1)$，因为 $P(Y=0)=0$.）这一变换不是一一的，因为 $(x,y)$ 和 $(-x,-y)$ 都映到同一个 $(u,v)$.
\sourcepageanchor{164}

但若只考虑 $y$ 的正值或只考虑负值，则变换是一一的. 用上面的记号，令
$$A_1=\{(x,y):y>0\},\quad A_2=\{(x,y):y<0\},\quad A_0=\{(x,y):y=0\}.$$
$A_0,A_1,A_2$ 构成 $\mathscr A=\mathbb R^2$ 的划分，且 $P((X,Y)\in A_0)=P(Y=0)=0$. 对 $A_1$ 或 $A_2$，若 $(x,y)\in A_i$，则 $v=|y|>0$；固定 $v=|y|$ 后，由于 $x$ 可取任意实数，$u=x/y$ 也可取任意实数. 因此，$\mathscr B=\{(u,v):v>0\}$ 是 $A_1$ 和 $A_2$ 在该变换下共同的像. 此外，从 $\mathscr B$ 到 $A_1$ 及从 $\mathscr B$ 到 $A_2$ 的逆变换分别为 $x=h_{11}(u,v)=uv$、$y=h_{21}(u,v)=v$，以及 $x=h_{12}(u,v)=-uv$、$y=h_{22}(u,v)=-v$. 注意，第一个逆变换给出正的 $y$，第二个给出负的 $y$. 两个逆变换的 Jacobian 为 $J_1=J_2=v$. 利用
$$f_{X,Y}(x,y)=\frac1{2\pi}e^{-x^2/2}e^{-y^2/2},$$
由\eqref{eq:4.3.6} 得到
$$
\begin{aligned}
f_{U,V}(u,v)&=\frac1{2\pi}e^{-(uv)^2/2}e^{-v^2/2}|v|+\frac1{2\pi}e^{-(-uv)^2/2}e^{-(-v)^2/2}|v|\\
&=\frac v\pi e^{-(u^2+1)v^2/2},\qquad-\infty<u<\infty,\quad0<v<\infty.
\end{aligned}
$$
由此可求 $U$ 的边缘 pdf：
$$
\begin{aligned}
f_U(u)&=\int_0^\infty\frac v\pi e^{-(u^2+1)v^2/2}\,dv\\
&=\frac1{2\pi}\int_0^\infty e^{-(u^2+1)z/2}\,dz&&\text{（换元 $z=v^2$）}\\
&=\frac1{2\pi}\frac2{u^2+1}&&\left(\begin{gathered}\text{被积函数是指数 pdf 的核，}\\\beta=2/(u^2+1)\end{gathered}\right)\\
&=\frac1{\pi(u^2+1)},\qquad-\infty<u<\infty.
\end{aligned}
$$
所以，两个独立标准正态随机变量之比为 Cauchy 随机变量. 关于正态与 Cauchy 随机变量的更多关系，见\autoref{exercise:4.28}.
\end{example}



\subsection{分层模型与混合分布}\label{sec:4.4}
到目前为止，一个随机变量都只有一个分布，这个分布可能依赖于参数. 虽然一般说来，随机变量只能有一个分布，但以层次结构来思考问题，常常更容易建立模型.
\begin{example}[二项--Poisson 层次]\label{ex:4.4.1}\index{SI0340@二项–泊松}
下面也许是最经典的分层模型（hierarchical model）. 一只昆虫产下大量卵，每个卵存活的概率为 $p$. 平均会有多少卵存活？

产卵的“大量”是一个随机变量，常假定服从 $\operatorname{Poisson}(\lambda)$ 分布. 此外，若假定各个卵独立存活，就得到了 Bernoulli 试验. 因此，令 $X=$ 存活卵数，$Y=$ 产卵数，则有
$$\begin{aligned}X\mid Y&\sim\operatorname{binomial}(Y,p),\\Y&\sim\operatorname{Poisson}(\lambda),\end{aligned}$$
这就是一个分层模型. 回顾 $X\mid Y\sim\operatorname{binomial}(Y,p)$ 这样的记号，表示给定 $Y=y$ 时 $X$ 的条件分布为 $\operatorname{binomial}(y,p)$.
\end{example}
\sourcepageanchor{165}

层次结构的优点在于，可以用一系列处于不同层次的相对简单的模型，来描述复杂过程. 此外，处理层次结构，并不比处理条件分布和边缘分布更困难.
\begin{example}[续\autoref{ex:4.4.1}]\label{ex:4.4.2}
所关心的随机变量 $X=$ 存活卵数的分布为
$$
\begin{aligned}
P(X=x)&=\sum_{y=0}^{\infty}P(X=x,Y=y)\\
&=\sum_{y=0}^{\infty}P(X=x\mid Y=y)P(Y=y)&&\text{（条件概率的定义）}\\
&=\sum_{y=x}^{\infty}\left[\binom yx p^x(1-p)^{y-x}\right]\left[\frac{e^{-\lambda}\lambda^y}{y!}\right],&&\text{（$y<x$ 时条件概率为 0）}
\end{aligned}
$$
因为 $X\mid Y=y$ 服从 $\operatorname{binomial}(y,p)$，而 $Y$ 服从 $\operatorname{Poisson}(\lambda)$. 对最后的表达式约去可约的因子，并乘以 $\lambda^x/\lambda^x$，得到
$$
\begin{aligned}
P(X=x)&=\frac{(\lambda p)^xe^{-\lambda}}{x!}\sum_{y=x}^{\infty}\frac{((1-p)\lambda)^{y-x}}{(y-x)!}\\
&=\frac{(\lambda p)^xe^{-\lambda}}{x!}\sum_{t=0}^{\infty}\frac{((1-p)\lambda)^t}{t!}&&\text{（$t=y-x$）}\\
&=\frac{(\lambda p)^xe^{-\lambda}}{x!}e^{(1-p)\lambda}&&\text{（和式来自 Poisson 分布的核）}\\
&=\frac{(\lambda p)^x}{x!}e^{-\lambda p}.
\end{aligned}
$$
所以 $X\sim\operatorname{Poisson}(\lambda p)$. 因而，所有关于 $X$ 的边缘推断都是针对 $\operatorname{Poisson}(\lambda p)$ 分布进行的，$Y$ 完全不参与. 在层次结构中引入 $Y$，主要是为了帮助理解模型. 还有一个额外好处：$X$ 的分布参数是两个参数的乘积，而这两个参数各自都比较容易理解.

现在很容易回答最初的问题：
$$EX=\lambda p,$$
所以平均有 $\lambda p$ 个卵存活. 如果只关心这个均值而不需要整个分布，也可以利用条件期望的性质.
\end{example}
有时，下述定理能够大大简化计算. 回顾\autoref{sec:4.2}，$E(X\mid y)$ 是 $y$ 的函数，而 $E(X\mid Y)$ 是取值依赖于 $Y$ 的随机变量.
\begin{theorem}\label{thm:4.4.3}\index{SI0428@迭代}
若 $X$ 和 $Y$ 是任意两个随机变量，则只要有关期望存在，就有
\begin{equation}\label{eq:4.4.1}
EX=E(E(X\mid Y)).
\end{equation}
\end{theorem}
\begin{proof}
记 $X,Y$ 的联合 pdf 为 $f(x,y)$. 由定义，
\begin{equation}\label{eq:4.4.2}
EX=\iint xf(x,y)\,dx\,dy=\int\left[\int xf(x\mid y)\,dx\right]f_Y(y)\,dy,
\end{equation}
其中 $f(x\mid y)$ 和 $f_Y(y)$ 分别是给定 $Y=y$ 时 $X$ 的条件 pdf 与 $Y$ 的边缘 pdf. 注意，\eqref{eq:4.4.2} 中的内层积分就是条件期望 $E(X\mid y)$，所以
$$EX=\int E(X\mid y)f_Y(y)\,dy=E(E(X\mid Y)),$$
即所欲证. 将积分换为求和，就得到离散情形的证明.
\end{proof}
\sourcepageanchor{166}

注意，\eqref{eq:4.4.1} 中存在记号的混用，因为同一个等式里的“$E$”代表不同的期望. 左端的 $E$ 是关于 $X$ 的边缘分布取期望；右端第一个 $E$ 是关于 $Y$ 的边缘分布取期望，第二个则是关于 $X\mid Y$ 的条件分布取期望. 不过，这其实不会引起混淆，因为这些就是符号 $E$ 在各处唯一可能的解释！

现在很容易计算\autoref{ex:4.4.1} 中存活卵数的期望. 根据\autoref{thm:4.4.3}，
$$\begin{aligned}
EX&=E(E(X\mid Y))\\
&=E(pY)&&\text{（因为 $X\mid Y\sim\operatorname{binomial}(Y,p)$）}\\
&=p\lambda.&&\text{（因为 $Y\sim\operatorname{Poisson}(\lambda)$）}
\end{aligned}$$
本节标题中的混合分布（mixture distribution），指由分层结构产生的分布. 虽然这一术语没有统一的定义，我们将采用下面这个似乎比较常用的定义.
\begin{definition}\label{def:4.4.4}\index{SI0336@混合}
如果随机变量 $X$ 的分布依赖于一个本身也有分布的量，就称 $X$ 具有混合分布.
\end{definition}
因此，\autoref{ex:4.4.1} 中的 $\operatorname{Poisson}(\lambda p)$ 分布是混合分布，因为它是将 $\operatorname{binomial}(Y,p)$ 与 $Y\sim\operatorname{Poisson}(\lambda)$ 结合的结果. 一般说来，分层模型产生混合分布.

层次结构不必止于两层，但也容易看出，理论上任何更复杂的层次结构都可以作为两层结构来处理. 不过，用多层结构为某种现象建模可能有好处，因为它更容易理解.
\begin{example}[\autoref{ex:4.4.1} 的推广]\label{ex:4.4.5}\index{SI0341@二项–泊松–指数}
考虑\autoref{ex:4.4.1} 的一个推广：现在有大量母虫，从中随机选择一只，而不是只有一只母虫. 我们仍然想知道平均存活卵数，但不能再明确断言，每只母虫的产卵数都服从同一个 Poisson 分布. 下列三层结构也许更合适. 令 $X=$ 一窝中的存活卵数，则
$$\begin{aligned}
X\mid Y&\sim\operatorname{binomial}(Y,p),\\
Y\mid\Lambda&\sim\operatorname{Poisson}(\Lambda),\\
\Lambda&\sim\operatorname{exponential}(\beta),
\end{aligned}$$
其中最后一层用于描述不同母虫之间的变异. $X$ 的均值很容易计算：
$$\begin{aligned}
EX&=E(E(X\mid Y))\\
&=E(pY)&&\text{（同前）}\\
&=E(E(pY\mid\Lambda))\\
&=E(p\Lambda)\\
&=p\beta,&&\text{（指数分布的期望）}
\end{aligned}$$
计算完成.
\end{example}
本例的模型与之前略有不同：其中两个随机变量是离散的，一个是连续的. 使用这类模型不会造成困难. 仍可像以前那样定义联合密度 $f(x,y,\lambda)$，条件密度 $f(x\mid y)$、$f(x\mid y,\lambda)$ 等，以及边缘密度 $f(x)$、$f(x,y)$ 等. 只需理解，计算概率或期望时，对离散变量求和，对连续变量积分.
\sourcepageanchor{167}

注意，将最后两层合并，就可以把这个三层模型看作两层结构. 若 $Y\mid\Lambda\sim\operatorname{Poisson}(\Lambda)$ 且 $\Lambda\sim\operatorname{exponential}(\beta)$，则
$$\begin{aligned}
P(Y=y)&=P(Y=y,0<\Lambda<\infty)\\
&=\int_0^\infty f(y,\lambda)\,d\lambda\\
&=\int_0^\infty f(y\mid\lambda)f(\lambda)\,d\lambda\\
&=\int_0^\infty\left[\frac{e^{-\lambda}\lambda^y}{y!}\right]\frac1\beta e^{-\lambda/\beta}\,d\lambda\\
&=\frac1{\beta y!}\int_0^\infty\lambda^ye^{-\lambda(1+\beta^{-1})}\,d\lambda&&\text{（Gamma pdf 的核）}\\
&=\frac1{\beta y!}\Gamma(y+1)\left(\frac1{1+\beta^{-1}}\right)^{y+1}\\
&=\frac1{1+\beta}\left(\frac1{1+\beta^{-1}}\right)^y.
\end{aligned}$$
$Y$ 的这个 pmf 就是负二项 pmf 的\hyperref[eq:3.2.10]{式 (3.1.10)} 形式.\footnote{译注：原文引用 (3.1.10)，对应公式实际编号为\eqref{eq:3.2.10}.} 因此，\autoref{ex:4.4.5} 的三层结构等价于两层结构\index{SI0339@二项–负二项}\semanticindexanchor{339:167}
$$\begin{aligned}
X\mid Y&\sim\operatorname{binomial}(Y,p),\\
Y&\sim\text{负二项}\left(p=\frac1{1+\beta},r=1\right).
\end{aligned}$$
不过，从理解模型的角度说，三层模型要容易理解得多！

一个有用的推广是 Poisson--Gamma 混合，它推广了上述模型的一部分. 若有层次结构\index{SI0346@泊松–伽马}\semanticindexanchor{346:167}
$$\begin{aligned}Y\mid\Lambda&\sim\operatorname{Poisson}(\Lambda),\\\Lambda&\sim\operatorname{gamma}(\alpha,\beta),\end{aligned}$$
则 $Y$ 的边缘分布是负二项分布（见\autoref{exercise:4.32}）. 负二项分布的这个模型表明，可以将它看成“变异更大”的 Poisson 分布. \hyperref[bib:source297]{Solomon (1983)} 解释了这些以及其他会产生负二项分布的生物学和数学模型（见\autoref{exercise:4.33}）.

除了有助于理解，分层模型还常能简化计算. 例如，统计学中常见的一个分布是非中心卡方分布（noncentral chi squared distribution）. 自由度为 $p$、非中心参数为 $\lambda$ 时，其 pdf 为
\begin{equation}\label{eq:4.4.3}\index{SI0342@卡方–泊松}\index{SI0357@非中心卡方}
f(x\mid\lambda,p)=\sum_{k=0}^{\infty}\frac{x^{p/2+k-1}e^{-x/2}}{\Gamma(p/2+k)2^{p/2+k}}\frac{\lambda^ke^{-\lambda}}{k!},
\end{equation}
这是一个极为繁杂的表达式. 例如，计算 $EX$ 看起来颇费周折. 但仔细观察 pdf，就能发现这是一个由中心卡方密度（如\hyperref[eq:3.3.10]{式 (3.2.10)} 所给出的密度）和 Poisson 分布组成的混合分布.\footnote{译注：此处原文引用 (3.2.10)，中心卡方密度实际见\eqref{eq:3.3.10}；\eqref{eq:3.2.10} 是负二项 pmf.} 也就是说，如果建立层次结构
$$\begin{aligned}X\mid K&\sim\chi^2_{p+2K},\\K&\sim\operatorname{Poisson}(\lambda),\end{aligned}$$
则 $X$ 的边缘分布就是\eqref{eq:4.4.3}. 从而
$$\begin{aligned}EX&=E(E(X\mid K))\\&=E(p+2K)\\&=p+2\lambda,\end{aligned}$$
计算相当简单. $\operatorname{Var}X$ 也能用这种方法求出.
\sourcepageanchor{168}

本节最后再介绍一个分层模型，并再说明一种条件期望计算.
\begin{example}[Beta--二项层次]\label{ex:4.4.6}\index{SI0337@伯努利–贝塔}
二项分布的一种推广，是允许成功概率按某个分布变化. 描述这种情形的一个标准模型是\footnote{译注：原文第一行写为 $X\mid P\sim\operatorname{binomial}(P)$，并带有未使用的 $i=1,\ldots,n$. 由紧随其后的 $E(X\mid P)=nP$ 及\autoref{ex:4.4.8} 可知，意指 $X\mid P\sim\operatorname{binomial}(n,P)$. 此处保留原式.}
$$\begin{aligned}X\mid P&\sim\operatorname{binomial}(P),\qquad i=1,\ldots,n,\\P&\sim\operatorname{beta}(\alpha,\beta).\end{aligned}$$
通过迭代取期望，计算 $X$ 的均值为
$$EX=E[E(X\mid P)]=E[nP]=n\frac\alpha{\alpha+\beta}.$$
\end{example}
计算 $X$ 的方差只稍微复杂一些. 我们可以使用一个条件方差公式，它与\autoref{thm:4.4.3} 的期望恒等式具有相似的思想.
\begin{theorem}[条件方差恒等式]\label{thm:4.4.7}\index{SI1064@条件}
对任意两个随机变量 $X$ 和 $Y$，只要有关期望存在，就有
\begin{equation}\label{eq:4.4.4}
\operatorname{Var}X=E(\operatorname{Var}(X\mid Y))+\operatorname{Var}(E(X\mid Y)).
\end{equation}
\end{theorem}
\begin{proof}
由定义，
$$\operatorname{Var}X=E([X-EX]^2)=E([X-E(X\mid Y)+E(X\mid Y)-EX]^2),$$
最后一步加上又减去了 $E(X\mid Y)$. 展开最后一个期望中的平方，得到
\begin{equation}\label{eq:4.4.5}
\begin{aligned}
\operatorname{Var}X&=E([X-E(X\mid Y)]^2)+E([E(X\mid Y)-EX]^2)\\
&\quad+2E([X-E(X\mid Y)][E(X\mid Y)-EX]).
\end{aligned}
\end{equation}
然而，最后一项为 0，这可通过迭代取期望轻易看出：
\begin{equation}\label{eq:4.4.6}
\begin{aligned}
&E([X-E(X\mid Y)][E(X\mid Y)-EX])\\
&\qquad=E(E\{[X-E(X\mid Y)][E(X\mid Y)-EX]\mid Y\}).
\end{aligned}
\end{equation}
在条件分布 $X\mid Y$ 中，$X$ 是随机变量. 所以在表达式
$$E\{[X-E(X\mid Y)][E(X\mid Y)-EX]\mid Y\}$$
中，$E(X\mid Y)$ 与 $EX$ 都是常数. 因此，\footnote{译注：下式第三行第一个因子后多出的一个右括号为原文排印错误；正确配对为 $(E(X\mid Y)-EX)(E(X\mid Y)-E(X\mid Y))$.}
$$\begin{aligned}
&E\{[X-E(X\mid Y)][E(X\mid Y)-EX]\mid Y\}\\
&\quad=(E(X\mid Y)-EX)(E\{[X-E(X\mid Y)]\mid Y\})\\
&\quad=(E(X\mid Y)-EX))(E(X\mid Y)-E(X\mid Y))\\
&\quad=(E(X\mid Y)-EX)(0)\\
&\quad=0.
\end{aligned}$$

因此，由\eqref{eq:4.4.6} 得到 $E((X-E(X\mid Y))(E(X\mid Y)-EX))=E(0)=0$. 再看\eqref{eq:4.4.5}，有
$$E([X-E(X\mid Y)]^2)=E\bigl(E\{[X-E(X\mid Y)]^2\mid Y\}\bigr)=E(\operatorname{Var}(X\mid Y)),$$
以及
$$E([E(X\mid Y)-EX]^2)=\operatorname{Var}(E(X\mid Y)),$$
从而证明了\eqref{eq:4.4.4}.
\end{proof}
\sourcepageanchor{169}

\begin{example}[续\autoref{ex:4.4.6}]\label{ex:4.4.8}\index{SI0337@伯努利–贝塔}
为计算 $X$ 的方差，由\eqref{eq:4.4.4}，
$$\operatorname{Var}X=\operatorname{Var}(E(X\mid P))+E(\operatorname{Var}(X\mid P)).$$
现在 $E(X\mid P)=nP$，而且 $P\sim\operatorname{beta}(\alpha,\beta)$，所以
$$\operatorname{Var}(E(X\mid P))=\operatorname{Var}(nP)=n^2\frac{\alpha\beta}{(\alpha+\beta)^2(\alpha+\beta+1)}.$$
又因为 $X\mid P$ 服从 $\operatorname{binomial}(n,P)$，有 $\operatorname{Var}(X\mid P)=nP(1-P)$. 因而\footnote{译注：原文下面的积分前漏写因子 $n$；左端与下一行均包含 $n$，故应在 $\Gamma(\alpha+\beta)/(\Gamma(\alpha)\Gamma(\beta))$ 前乘以 $n$. 此处保留原式.}
$$E[\operatorname{Var}(X\mid P)]=nE[P(1-P)]=\frac{\Gamma(\alpha+\beta)}{\Gamma(\alpha)\Gamma(\beta)}\int_0^1 p(1-p)p^{\alpha-1}(1-p)^{\beta-1}\,dp.$$
注意，被积函数是另一个 Beta pdf 的核，其参数为 $\alpha+1$ 和 $\beta+1$. 所以
$$E(\operatorname{Var}(X\mid P))=n\frac{\Gamma(\alpha+\beta)}{\Gamma(\alpha)\Gamma(\beta)}\left[\frac{\Gamma(\alpha+1)\Gamma(\beta+1)}{\Gamma(\alpha+\beta+2)}\right]
=n\frac{\alpha\beta}{(\alpha+\beta)(\alpha+\beta+1)}.$$
将两部分相加并化简，得到
$$\operatorname{Var}X=n\frac{\alpha\beta(\alpha+\beta+n)}{(\alpha+\beta)^2(\alpha+\beta+1)}.$$
\end{example}



\subsection{协方差与相关}\label{sec:4.5}
前面几节讨论了两个随机变量之间有无关系，即独立或不独立. 但即便存在关系，也可能有强有弱. 本节讨论度量两个随机变量关系强弱的两个数值指标：协方差与相关.

为说明两个随机变量之间关系的强弱是什么意思，考虑两个不同的实验. 在第一个实验中，测量随机变量 $X,Y$：$X$ 是一份水样的重量，$Y$ 是同一水样的体积. 显然，$X$ 与 $Y$ 有很强的关系. 对若干水样测量 $(X,Y)$，并将观测数据对画出来，由于二者的物理关系，数据点应落在一条直线上. 测量误差、水中杂质等因素会使它们不能精确地落在直线上，但只要实验操作足够细致，这些点就会非常接近一条直线. 再考虑另一个实验：$X$ 是某人的体重，$Y$ 是同一个人的身高. 二者显然也存在关系，但远没有前一个实验那么强. 我们不会期望不同人的 $(X,Y)$ 数据对排成一条直线，尽管可能预期图中呈现上升趋势. 协方差与相关就是将这种关系强弱的差别数量化的两个指标.

本节将经常涉及 $X$ 的均值和方差，以及 $Y$ 的均值和方差. 使用记号 $EX=\mu_X$、$EY=\mu_Y$、$\operatorname{Var}X=\sigma_X^2$、$\operatorname{Var}Y=\sigma_Y^2$. 始终假定 $0<\sigma_X^2<\infty$，$0<\sigma_Y^2<\infty$.
\begin{definition}\label{def:4.5.1}\index{SI0223@协方差（另见相关）}
$X$ 和 $Y$ 的协方差（covariance）定义为
$$\operatorname{Cov}(X,Y)=E((X-\mu_X)(Y-\mu_Y)).$$
\end{definition}
\sourcepageanchor{170}

\begin{definition}\label{def:4.5.2}\index{SI0216@系数}\index{SI0217@相关（另见协方差）}
$X$ 和 $Y$ 的相关（correlation）定义为
$$\rho_{XY}=\frac{\operatorname{Cov}(X,Y)}{\sigma_X\sigma_Y}.$$
数值 $\rho_{XY}$ 也称为相关系数（correlation coefficient）.
\end{definition}
如果 $X$ 的大值往往与 $Y$ 的大值一起出现，$X$ 的小值往往与 $Y$ 的小值一起出现，则 $\operatorname{Cov}(X,Y)$ 为正. 若 $X>\mu_X$，则很可能 $Y>\mu_Y$，于是乘积 $(X-\mu_X)(Y-\mu_Y)$ 为正. 若 $X<\mu_X$，则很可能 $Y<\mu_Y$，这个乘积仍为正. 所以 $\operatorname{Cov}(X,Y)=E((X-\mu_X)(Y-\mu_Y))>0$. 如果 $X$ 的大值往往与 $Y$ 的小值一起出现，而 $X$ 的小值与 $Y$ 的大值一起出现，则协方差为负，因为 $X>\mu_X$ 时，$Y$ 往往小于 $\mu_Y$，反之亦然，所以 $(X-\mu_X)(Y-\mu_Y)$ 往往为负. 因而，协方差的符号提供了关于 $X,Y$ 关系的信息.

但是，$\operatorname{Cov}(X,Y)$ 可以取任意数值；仅仅知道一个给定值，例如 $\operatorname{Cov}(X,Y)=3$，本身不能说明关系有多强. 另一方面，相关系数总在 $-1$ 与 1 之间，取 $-1$ 或 1 表示 $X$ 与 $Y$ 之间有完全线性关系. \autoref{thm:4.5.7} 将证明这一点.

在考察协方差与相关的这些性质之前，先在一个具体例子中计算它们. 下述结论将简化计算.
\begin{theorem}\label{thm:4.5.3}
对任意随机变量 $X$ 和 $Y$，
$$\operatorname{Cov}(X,Y)=EXY-\mu_X\mu_Y.$$
\end{theorem}
\begin{proof}
$$\begin{aligned}
\operatorname{Cov}(X,Y)&=E((X-\mu_X)(Y-\mu_Y))\\
&=E(XY-\mu_XY-\mu_YX+\mu_X\mu_Y)&&\text{（展开乘积）}\\
&=EXY-\mu_XEY-\mu_YEX+\mu_X\mu_Y&&\text{（$\mu_X,\mu_Y$ 为常数）}\\
&=EXY-\mu_X\mu_Y-\mu_Y\mu_X+\mu_X\mu_Y\\
&=EXY-\mu_X\mu_Y.
\end{aligned}$$
\end{proof}
\begin{example}[相关——I]\label{ex:4.5.4}
设 $(X,Y)$ 的联合 pdf 为 $f(x,y)=1$，$0<x<1$，$x<y<x+1$. $X$ 的边缘分布是 $\operatorname{uniform}(0,1)$，所以 $\mu_X=1/2$，$\sigma_X^2=1/12$. $Y$ 的边缘 pdf 为 $f_Y(y)=y$（$0<y<1$）及 $f_Y(y)=2-y$（$1\le y<2$），且 $\mu_Y=1$，$\sigma_Y^2=1/6$. 还有
$$\begin{aligned}
EXY&=\int_0^1\int_x^{x+1}xy\,dy\,dx=\int_0^1\left.\frac12xy^2\right|_x^{x+1}\,dx\\
&=\int_0^1\left(x^2+\frac12x\right)\,dx=\frac7{12}.
\end{aligned}$$
由\autoref{thm:4.5.3}，$\operatorname{Cov}(X,Y)=7/12-(1/2)(1)=1/12$. 相关系数为
$$\rho_{XY}=\frac{\operatorname{Cov}(X,Y)}{\sigma_X\sigma_Y}=\frac{1/12}{\sqrt{1/12}\sqrt{1/6}}=\frac1{\sqrt2}.$$
\end{example}
接下来三个定理描述协方差与相关的一些基本性质.
\sourcepageanchor{171}

\begin{theorem}\label{thm:4.5.5}\index{SI0221@独立性}\index{SI0548@协方差}
如果 $X$ 和 $Y$ 是独立随机变量，则 $\operatorname{Cov}(X,Y)=0$，$\rho_{XY}=0$.
\end{theorem}
\begin{proof}
因为 $X$ 与 $Y$ 独立，由\autoref{thm:4.2.10}，$EXY=(EX)(EY)$. 所以
$$\operatorname{Cov}(X,Y)=EXY-(EX)(EY)=(EX)(EY)-(EX)(EY)=0,$$
且
$$\rho_{XY}=\frac{\operatorname{Cov}(X,Y)}{\sigma_X\sigma_Y}=\frac0{\sigma_X\sigma_Y}=0.$$
\end{proof}
因此，在某种意义上，$\operatorname{Cov}(X,Y)=\rho_{XY}=0$ 表示 $X$ 与 $Y$ 没有关系. 但必须注意，\autoref{thm:4.5.5} 并没有说协方差为 0 就意味着独立. 例如，设 $X\sim f(x-\theta)$，关于 0 对称，且 $EX=\theta$；令 $Y$ 为示性函数 $Y=I(|X-\theta|<2)$，则 $X$ 和 $Y$ 显然不独立.\footnote{译注：原文“关于 0 对称”应指基准密度 $f(t)$；$X$ 的密度 $f(x-\theta)$ 是关于 $\theta$ 对称. 此外，断言不独立还需 $0<P(|X-\theta|<2)<1$，否则 $Y$ 为常数.} 然而，
$$\begin{aligned}
E(XY)&=\int_{-\infty}^{\infty}xI(|x-\theta|<2)f(x-\theta)\,dx\\
&=\int_{-2}^2(t+\theta)f(t)\,dt=\theta\int_{-2}^2f(t)\,dt=EXEY,
\end{aligned}$$
其中利用了对称性给出的 $\int_{-2}^2tf(t)\,dt=0$. 可见，找到不相关但不独立的随机变量并不困难.

协方差与相关只度量一种特定的线性关系，\autoref{thm:4.5.7} 将进一步说明. 另见\autoref{ex:4.5.9}：两个随机变量关系很强，但由于关系不是线性的，协方差与相关都为 0.

协方差还在理解随机变量之和的变异方面起重要作用，下面这个对\autoref{thm:2.3.4} 的推广说明了这一点. 更进一步的推广见\autoref{exercise:4.44}.
\begin{theorem}\label{thm:4.5.6}\index{SI1068@和}
若 $X,Y$ 为任意两个随机变量，$a,b$ 为任意两个常数，则
$$\operatorname{Var}(aX+bY)=a^2\operatorname{Var}X+b^2\operatorname{Var}Y+2ab\operatorname{Cov}(X,Y).$$
若 $X$ 和 $Y$ 独立，则
$$\operatorname{Var}(aX+bY)=a^2\operatorname{Var}X+b^2\operatorname{Var}Y.$$
\end{theorem}
\begin{proof}
$aX+bY$ 的均值为 $E(aX+bY)=aEX+bEY=a\mu_X+b\mu_Y$. 因此，
$$\begin{aligned}
\operatorname{Var}(aX+bY)&=E\bigl[((aX+bY)-(a\mu_X+b\mu_Y))^2\bigr]\\
&=E\bigl[(a(X-\mu_X)+b(Y-\mu_Y))^2\bigr]\\
&=E\bigl[a^2(X-\mu_X)^2+b^2(Y-\mu_Y)^2+2ab(X-\mu_X)(Y-\mu_Y)\bigr]\\
&=a^2E(X-\mu_X)^2+b^2E(Y-\mu_Y)^2+2abE((X-\mu_X)(Y-\mu_Y))\\
&=a^2\operatorname{Var}X+b^2\operatorname{Var}Y+2ab\operatorname{Cov}(X,Y).
\end{aligned}$$
若 $X$ 与 $Y$ 独立，由\autoref{thm:4.5.5} 有 $\operatorname{Cov}(X,Y)=0$，第二个等式立即由第一个得到.
\end{proof}
根据\autoref{thm:4.5.6}，如果 $X$ 和 $Y$ 正相关（$\operatorname{Cov}(X,Y)>0$），则 $X+Y$ 的变异大于 $X,Y$ 各自变异之和. 如果二者负相关，则 $X+Y$ 的变异小于这个和. 对负相关随机变量，一个的大值往往伴随另一个的小值，求和时两个极端相互抵消. 所得 $X+Y$ 往往没有那么多极端值，因而方差较小. 取 $a=1$、$b=-1$，便得到两个随机变量之差的方差表达式，同样可以作类似解释.

下述定理在一定程度上说明了协方差与相关所度量的线性关系的性质.
\sourcepageanchor{172}

\begin{theorem}\label{thm:4.5.7}\index{SI0217@相关（另见协方差）}\index{SI0564@Cauchy–Schwarz}
对任意随机变量 $X$ 和 $Y$，
\begin{enumerate}
\item[(a)] $-1\le\rho_{XY}\le1$.
\item[(b)] $|\rho_{XY}|=1$ 当且仅当存在数 $a\ne0$ 与 $b$，使 $P(Y=aX+b)=1$. 如果 $\rho_{XY}=1$，则 $a>0$；如果 $\rho_{XY}=-1$，则 $a<0$.
\end{enumerate}
\end{theorem}
\begin{proof}
考虑函数
$$h(t)=E\bigl[((X-\mu_X)t+(Y-\mu_Y))^2\bigr].$$
展开得到
$$\begin{aligned}
h(t)&=t^2E(X-\mu_X)^2+2tE((X-\mu_X)(Y-\mu_Y))+E(Y-\mu_Y)^2\\
&=t^2\sigma_X^2+2t\operatorname{Cov}(X,Y)+\sigma_Y^2.
\end{aligned}$$
这个关于 $t$ 的二次函数对所有 $t$ 均大于或等于 0，因为它是非负随机变量的期望. 所以它至多有一个实根，判别式必不为正，即
$$(2\operatorname{Cov}(X,Y))^2-4\sigma_X^2\sigma_Y^2\le0.$$
这等价于
$$-\sigma_X\sigma_Y\le\operatorname{Cov}(X,Y)\le\sigma_X\sigma_Y.$$
除以 $\sigma_X\sigma_Y$，得到
$$-1\le\frac{\operatorname{Cov}(X,Y)}{\sigma_X\sigma_Y}=\rho_{XY}\le1.$$
此外，$|\rho_{XY}|=1$ 当且仅当判别式为 0，也就是 $h(t)$ 恰有一个实根. 由于 $((X-\mu_X)t+(Y-\mu_Y))^2\ge0$，期望 $h(t)=E[((X-\mu_X)t+(Y-\mu_Y))^2]=0$ 当且仅当
$$P([(X-\mu_X)t+(Y-\mu_Y)]^2=0)=1.$$
这又等价于
$$P((X-\mu_X)t+(Y-\mu_Y)=0)=1.$$
取 $a=-t$、$b=\mu_Xt+\mu_Y$，其中 $t$ 是 $h(t)$ 的根，即得 $P(Y=aX+b)=1$. 根据二次方程求根公式，这个根为 $t=-\operatorname{Cov}(X,Y)/\sigma_X^2$. 因此，$a=-t$ 与 $\rho_{XY}$ 同号，最后一个断言也得证.
\end{proof}
在\autoref{sec:4.7} 中，将证明一个称为 Cauchy--Schwarz 不等式的定理. 直接应用它就能得到 $\rho_{XY}$ 介于 $-1$ 和 1 之间；我们会看到，使用它可以缩短上面的证明.

如果存在直线 $y=ax+b$，$a\ne0$，使 $(X,Y)$ 的取值有很高概率位于该直线附近，那么 $X,Y$ 的相关系数就接近 1 或 $-1$. 如果不存在这样的直线，相关系数就接近 0. 这是对相关所度量的线性关系的一种直观理解. 下面两个例子将进一步说明.
\begin{example}[相关——II]\label{ex:4.5.8}
本例与\autoref{ex:4.5.4} 类似，但采用不同的推导方式，以说明其他建模和计算技巧. 设 $X$ 服从 $\operatorname{uniform}(0,1)$，$Z$ 服从 $\operatorname{uniform}(0,1/10)$，且 $X$ 与 $Z$ 独立. 令 $Y=X+Z$，考虑随机向量 $(X,Y)$. 可以利用\autoref{sec:4.3} 的方法，从 $(X,Z)$ 的联合分布推得 $(X,Y)$ 的联合分布. $(X,Y)$ 的联合 pdf 为
$$f(x,y)=10,\qquad0<x<1,\quad x<y<x+\frac1{10}.$$
\sourcepageanchor{173}

\begin{figure}[htbp]
\centering\includegraphics[width=.94\linewidth]{assets/chapter4/fig4-5-1.pdf}
\caption{(a) \autoref{ex:4.5.4} 中 $f(x,y)>0$ 的区域；(b) \autoref{ex:4.5.8} 中 $f(x,y)>0$ 的区域}\label{fig:4.5.1}
\end{figure}

无须使用\autoref{sec:4.3} 的正式方法，也可如下说明. 给定 $X=x$，有 $Y=x+Z$. 由于 $X$ 与 $Z$ 独立，给定 $X=x$ 时 $Z$ 的条件分布仍为 $\operatorname{uniform}(0,1/10)$. 因此，$x$ 在 $Y\mid X=x$ 的条件分布中充当位置参数，这个条件分布就是 $\operatorname{uniform}(x,x+1/10)$. 将此条件 pdf 与 $X$ 的边缘 pdf（$\operatorname{uniform}(0,1)$）相乘，便得到上述联合 pdf. 表示式 $Y=X+Z$ 使协方差与相关的计算很容易. $X,Y$ 的期望分别为 $EX=1/2$ 和 $EY=E(X+Z)=EX+EZ=1/2+1/20=11/20$，因此
$$\begin{aligned}
\operatorname{Cov}(X,Y)&=EXY-(EX)(EY)\\
&=E[X(X+Z)]-(EX)(E(X+Z))\\
&=EX^2+EXZ-(EX)^2-(EX)(EZ)\\
&=EX^2-(EX)^2+(EX)(EZ)-(EX)(EZ)&&\text{（$X,Z$ 独立）}\\
&=\sigma_X^2=\frac1{12}.
\end{aligned}$$
根据\autoref{thm:4.5.6}，$Y$ 的方差为 $\sigma_Y^2=\operatorname{Var}(X+Z)=\operatorname{Var}X+\operatorname{Var}Z=1/12+1/1200$. 所以
$$\rho_{XY}=\frac{1/12}{\sqrt{1/12}\sqrt{1/12+1/1200}}=\sqrt{\frac{100}{101}}.$$
这比\autoref{ex:4.5.4} 中的 $\rho_{XY}=1/\sqrt2$ 大得多. \autoref{fig:4.5.1} 展示了\autoref{ex:4.5.4} 和本例中 $f(x,y)$ 为正的集合. 回顾，这个集合称为分布的支撑集（support）. 两种情形中，$(X,Y)$ 都是相应集合中的随机点，且 $X,Y$ 都有线性递增关系. 但\hyperref[fig:4.5.1]{图 4.5.1(b)} 中的关系强得多. 还可以这样理解：本例中，$Y\mid X=x$ 的条件分布为 $\operatorname{uniform}(x,x+1/10)$；\autoref{ex:4.5.4} 中则为 $\operatorname{uniform}(x,x+1)$. 在本模型中，知道 $X=x$ 能提供比\autoref{ex:4.5.4} 多得多的关于 $Y$ 值的信息，因此相关系数更接近 1.
\end{example}
下例说明，即便 $X,Y$ 的关系很强，如果关系不是线性的，相关系数仍可能很小.
\sourcepageanchor{174}

\begin{figure}[htbp]
\centering\includegraphics[width=.65\linewidth]{assets/chapter4/fig4-5-2.pdf}
\caption{\autoref{ex:4.5.9} 中 $f(x,y)>0$ 的区域}\label{fig:4.5.2}
\end{figure}

\begin{example}[相关——III]\label{ex:4.5.9}
设 $X$ 服从 $\operatorname{uniform}(-1,1)$，$Z$ 服从 $\operatorname{uniform}(0,1/10)$，且 $X$ 与 $Z$ 独立. 令 $Y=X^2+Z$，考虑随机向量 $(X,Y)$. 与\autoref{ex:4.5.8} 一样，给定 $X=x$ 时，$Y=x^2+Z$，所以 $Y\mid X=x$ 的条件分布是 $\operatorname{uniform}(x^2,x^2+1/10)$. 此条件 pdf 与 $X$ 的边缘 pdf 之积给出联合 pdf：
$$f(x,y)=5,\qquad-1<x<1,\quad x^2<y<x^2+\frac1{10}.$$
\autoref{fig:4.5.2} 展示了 $f(x,y)>0$ 的集合. 从给定 $X=x$ 时 $Y$ 的条件分布可见，$X,Y$ 的关系很强，但不是线性的. $(X,Y)$ 的可能值聚集在一条抛物线附近，而不是直线附近. 相关系数不度量这种非线性关系；事实上，$\rho_{XY}=0$. 因为 $X\sim\operatorname{uniform}(-1,1)$，有 $EX=EX^3=0$；又因为 $X,Z$ 独立，$EXZ=(EX)(EZ)$. 因此，
$$\begin{aligned}
\operatorname{Cov}(X,Y)&=E[X(X^2+Z)]-(EX)(E(X^2+Z))\\
&=EX^3+EXZ-0E(X^2+Z)\\
&=0+(EX)(EZ)=0(EZ)=0,
\end{aligned}$$
且 $\rho_{XY}=\operatorname{Cov}(X,Y)/(\sigma_X\sigma_Y)=0$.
\end{example}
本节最后引入一个非常重要的二元分布，其中相关系数自然地作为参数出现.
\begin{definition}\label{def:4.5.10}\index{SI0273@二元正态}
设 $-\infty<\mu_X<\infty$、$-\infty<\mu_Y<\infty$、$0<\sigma_X$、$0<\sigma_Y$、$-1<\rho<1$ 为五个实数. 均值为 $\mu_X,\mu_Y$，方差为 $\sigma_X^2,\sigma_Y^2$，相关系数为 $\rho$ 的二元正态 pdf（bivariate normal pdf）定义为
$$\begin{aligned}
f(x,y)&=\left(2\pi\sigma_X\sigma_Y\sqrt{1-\rho^2}\right)^{-1}\\
&\,\times\exp\left\{-\frac1{2(1-\rho^2)}\left[\left(\frac{x-\mu_X}{\sigma_X}\right)^2
-2\rho\left(\frac{x-\mu_X}{\sigma_X}\right)\left(\frac{y-\mu_Y}{\sigma_Y}\right)
+\left(\frac{y-\mu_Y}{\sigma_Y}\right)^2\right]\right\},
\end{aligned}$$
其中 $-\infty<x<\infty$，$-\infty<y<\infty$.
\end{definition}
\sourcepageanchor{175}

虽然二元正态 pdf 的公式看起来很吓人，它却是使用最频繁的二元分布之一. 实际上，公式的推导完全可以不那么吓人，见\autoref{exercise:4.46}.

此分布有许多良好性质，包括：
\begin{enumerate}
\item[(a)]\phantomsection\label{itm:bivnormal:a} $X$ 的边缘分布为 $n(\mu_X,\sigma_X^2)$.
\item[(b)]\phantomsection\label{itm:bivnormal:b} $Y$ 的边缘分布为 $n(\mu_Y,\sigma_Y^2)$.
\item[(c)]\phantomsection\label{itm:bivnormal:c} $X,Y$ 之间的相关系数为 $\rho_{XY}=\rho$.
\item[(d)]\phantomsection\label{itm:bivnormal:d} 对任意常数 $a,b$，$aX+bY$ 的分布为 $n(a\mu_X+b\mu_Y,a^2\sigma_X^2+b^2\sigma_Y^2+2ab\rho\sigma_X\sigma_Y)$.
\end{enumerate}
性质 \hyperref[itm:bivnormal:a]{(a)}、\hyperref[itm:bivnormal:b]{(b)}、\hyperref[itm:bivnormal:d]{(d)} 的验证留作习题（\autoref{exercise:4.45}）. 假设 \hyperref[itm:bivnormal:a]{(a)}、\hyperref[itm:bivnormal:b]{(b)} 成立，下面证明 \hyperref[itm:bivnormal:c]{(c)}. 由定义，
$$\begin{aligned}
\rho_{XY}&=\frac{\operatorname{Cov}(X,Y)}{\sigma_X\sigma_Y}
=\frac{E((X-\mu_X)(Y-\mu_Y))}{\sigma_X\sigma_Y}\\
&=E\left[\left(\frac{X-\mu_X}{\sigma_X}\right)\left(\frac{Y-\mu_Y}{\sigma_Y}\right)\right]\\
&=\int_{-\infty}^{\infty}\int_{-\infty}^{\infty}\left(\frac{x-\mu_X}{\sigma_X}\right)\left(\frac{y-\mu_Y}{\sigma_Y}\right)f(x,y)\,dx\,dy.
\end{aligned}$$
作变量替换
$$s=\left(\frac{x-\mu_X}{\sigma_X}\right)\left(\frac{y-\mu_Y}{\sigma_Y}\right),\qquad t=\frac{x-\mu_X}{\sigma_X}.$$
则 $x=\sigma_Xt+\mu_X$，$y=\sigma_Ys/t+\mu_Y$，变换的 Jacobian 为 $J=\sigma_X\sigma_Y/t$.\footnote{译注：按\autoref{sec:4.3} 的行列式次序，$\partial(x,y)/\partial(s,t)=-\sigma_X\sigma_Y/t$，原文漏了负号；下文取绝对值，故计算不受影响. 逆变换在 $t=0$ 处无定义，但对应 $X=\mu_X$ 为零概率事件，可在换元时略去.} 由此得到
$$\begin{aligned}
\rho_{XY}&=\int_{-\infty}^{\infty}\int_{-\infty}^{\infty}s f\left(\sigma_Xt+\mu_X,\frac{\sigma_Ys}t+\mu_Y\right)
\left|\frac{\sigma_X\sigma_Y}t\right|\,ds\,dt\\
&=\int_{-\infty}^{\infty}\int_{-\infty}^{\infty}s\left(2\pi\sigma_X\sigma_Y\sqrt{1-\rho^2}\right)^{-1}\\
&\qquad\times\exp\left[-\frac1{2(1-\rho^2)}\left(t^2-2\rho s+\left(\frac st\right)^2\right)\right]
\frac{\sigma_X\sigma_Y}{|t|}\,ds\,dt.
\end{aligned}$$
注意，$|t|=\sqrt{t^2}$，并且
$$t^2-2\rho s+\left(\frac st\right)^2=\left(\frac{s-\rho t^2}t\right)^2+(1-\rho^2)t^2,$$
所以可改写为
$$\rho_{XY}=\int_{-\infty}^{\infty}\frac1{\sqrt{2\pi}}\exp\left(-\frac{t^2}2\right)
\left[\int_{-\infty}^{\infty}\frac{s}{\sqrt{2\pi}\sqrt{(1-\rho^2)t^2}}
\exp\left(-\frac{(s-\rho t^2)^2}{2(1-\rho^2)t^2}\right)\,ds\right]\,dt.$$
内层积分是 $ES$，其中 $S$ 为均值 $ES=\rho t^2$、方差 $\operatorname{Var}S=(1-\rho^2)t^2$ 的正态随机变量. 因而内层积分为 $\rho t^2$，从而
$$\rho_{XY}=\int_{-\infty}^{\infty}\frac{\rho t^2}{\sqrt{2\pi}}\exp\left(-\frac{t^2}2\right)\,dt.$$
此积分又是 $\rho ET^2$，其中 $T\sim n(0,1)$. 因此 $ET^2=1$，得到 $\rho_{XY}=\rho$.
\sourcepageanchor{176}

给定 $X=x$ 时 $Y$ 的条件分布，以及给定 $Y=y$ 时 $X$ 的条件分布，也全都是正态分布. 利用上面的联合与边缘 pdf，容易验证，$Y\mid X=x$ 的条件分布为
$$n\bigl(\mu_Y+\rho(\sigma_Y/\sigma_X)(x-\mu_X),\sigma_Y^2(1-\rho^2)\bigr).$$
当 $\rho$ 趋于 1 或 $-1$ 时，条件方差 $\sigma_Y^2(1-\rho^2)$ 趋于 0. 因而，$Y\mid X=x$ 的条件分布越来越集中在点 $\mu_Y+\rho(\sigma_Y/\sigma_X)(x-\mu_X)$ 附近，$(X,Y)$ 的联合概率分布也越来越集中在直线 $y=\mu_Y+\rho(\sigma_Y/\sigma_X)(x-\mu_X)$ 附近. 这再次说明前面的观点：相关系数接近 1 或 $-1$，意味着存在直线 $y=ax+b$，使 $(X,Y)$ 的值以很高概率聚集在它附近.

注意一个重要事实：所有这些正态边缘与条件 pdf，都是从联合二元正态这一前提出发推导出来的. 反方向的推导并不成立，即边缘正态性并不蕴含联合正态性. \autoref{exercise:4.47} 给出了说明.



\subsection{多元分布}\label{sec:4.6}
本章开头讨论过在实验中观测两个以上随机变量的情形. 前几节主要研究二元随机向量 $(X,Y)$；本节讨论多元随机向量 $(X_1,\ldots,X_n)$. 在章首的例子中，测量一个人的体温、身高、体重和血压. 这里 $n=4$，观测的随机向量为 $(X_1,X_2,X_3,X_4)$，其中 $X_1$ 是体温，$X_2$ 是身高，等等. 前几节的概念，包括边缘分布与条件分布，都能从二元推广到多元情形. 本节介绍其中一些推广.

\textit{记号说明：}使用粗体字母表示多个变量. 因此，用 $\mathbf X$ 表示随机变量 $X_1,\ldots,X_n$，用 $\mathbf x$ 表示样本 $x_1,\ldots,x_n$.

随机向量 $\mathbf X=(X_1,\ldots,X_n)$ 的样本空间是 $\mathbb R^n$ 的一个子集. 若 $(X_1,\ldots,X_n)$ 为离散随机向量（样本空间可数），其联合 pmf 定义为 $f(\mathbf x)=f(x_1,\ldots,x_n)=P(X_1=x_1,\ldots,X_n=x_n)$，$(x_1,\ldots,x_n)\in\mathbb R^n$. 对任意 $A\subset\mathbb R^n$，
\begin{equation}\label{eq:4.6.1}\index{SI0853@联合}
P(\mathbf X\in A)=\sum_{\mathbf x\in A}f(\mathbf x).
\end{equation}
若 $(X_1,\ldots,X_n)$ 为连续随机向量，其联合 pdf 是满足
\begin{equation}\label{eq:4.6.2}\index{SI0833@联合}
P(\mathbf X\in A)=\int\cdots\int_A f(\mathbf x)\,d\mathbf x
=\int\cdots\int_A f(x_1,\ldots,x_n)\,dx_1\cdots dx_n
\end{equation}
的函数 $f(x_1,\ldots,x_n)$. 这些是 $n$ 重积分，积分限应使积分遍及 $A$ 中所有点 $\mathbf x$.

设 $g(\mathbf x)=g(x_1,\ldots,x_n)$ 是定义在 $\mathbf X$ 的样本空间上的实值函数，则 $g(\mathbf X)$ 是随机变量，其期望在连续和离散情形下分别为
\begin{equation}\label{eq:4.6.3}\index{SI0429@联合}
\begin{aligned}
Eg(\mathbf X)&=\int_{-\infty}^{\infty}\cdots\int_{-\infty}^{\infty}g(\mathbf x)f(\mathbf x)\,d\mathbf x,\\
Eg(\mathbf X)&=\sum_{\mathbf x\in\mathbb R^n}g(\mathbf x)f(\mathbf x).
\end{aligned}
\end{equation}
这些以及其他定义都类似于二元情形，只是积分或求和现在是在 $\mathbb R^n$ 的适当子集上进行，而不是 $\mathbb R^2$.

$(X_1,\ldots,X_n)$ 中任意一部分坐标的边缘 pdf 或 pmf，都可以通过对其他坐标的所有可能值积分或求和得到.\index{SI0834@边缘}\semanticindexanchor{834:176}\index{SI0855@边缘}\semanticindexanchor{855:176}
\sourcepageanchor{177}

例如，前 $k$ 个坐标 $(X_1,\ldots,X_k)$ 的边缘 pdf 或 pmf 为
\begin{equation}\label{eq:4.6.4}
f(x_1,\ldots,x_k)=\int_{-\infty}^{\infty}\cdots\int_{-\infty}^{\infty}f(x_1,\ldots,x_n)\,dx_{k+1}\cdots dx_n,
\end{equation}
或
\begin{equation}\label{eq:4.6.5}
f(x_1,\ldots,x_k)=\sum_{(x_{k+1},\ldots,x_n)\in\mathbb R^{n-k}}f(x_1,\ldots,x_n),
\end{equation}
其中 $(x_1,\ldots,x_k)\in\mathbb R^k$. 给定其余坐标的值时，某一部分坐标的条件 pdf 或 pmf，可由联合 pdf 或 pmf 除以其余坐标的边缘 pdf 或 pmf 得到. 例如，若 $f(x_1,\ldots,x_k)>0$，给定 $X_1=x_1,\ldots,X_k=x_k$ 时 $(X_{k+1},\ldots,X_n)$ 的条件 pdf 或 pmf，是关于 $(x_{k+1},\ldots,x_n)$ 的函数
\begin{equation}\label{eq:4.6.6}\index{SI0830@条件}\index{SI0851@条件}
f(x_{k+1},\ldots,x_n\mid x_1,\ldots,x_k)=\frac{f(x_1,\ldots,x_n)}{f(x_1,\ldots,x_k)}.
\end{equation}
下例说明这些思想.
\begin{example}[多元 pdf]\label{ex:4.6.1}
令 $n=4$，并取
$$f(x_1,x_2,x_3,x_4)=\begin{cases}\dfrac34(x_1^2+x_2^2+x_3^2+x_4^2),&0<x_i<1,\ i=1,2,3,4,\\0,&\text{其他情形}.\end{cases}$$
此非负函数是随机向量 $(X_1,X_2,X_3,X_4)$ 的联合 pdf，可以验证
$$\begin{aligned}
&\int_{-\infty}^{\infty}\int_{-\infty}^{\infty}\int_{-\infty}^{\infty}\int_{-\infty}^{\infty}f(x_1,x_2,x_3,x_4)\,dx_1\,dx_2\,dx_3\,dx_4\\
&\quad=\int_0^1\int_0^1\int_0^1\int_0^1\frac34(x_1^2+x_2^2+x_3^2+x_4^2)\,dx_1\,dx_2\,dx_3\,dx_4\\
&\quad=1.
\end{aligned}$$
此联合 pdf 可用于计算如下概率：
$$\begin{aligned}
&P\left(X_1<\frac12,X_2<\frac34,X_4>\frac12\right)\\
&\quad=\int_{1/2}^1\int_0^1\int_0^{3/4}\int_0^{1/2}\frac34(x_1^2+x_2^2+x_3^2+x_4^2)\,dx_1\,dx_2\,dx_3\,dx_4.
\end{aligned}$$
注意，这些积分限把积分限制在所求事件内且 $f(x_1,x_2,x_3,x_4)>0$ 的点上. 四项 $\frac34x_1^2$、$\frac34x_2^2$ 等可分别积分，再将结果相加. 例如，
$$\int_{1/2}^1\int_0^1\int_0^{3/4}\int_0^{1/2}\frac34x_1^2\,dx_1\,dx_2\,dx_3\,dx_4=\frac3{256}.$$
其余三个积分分别为 $7/1024$、$3/64$ 和 $21/256$.\footnote{译注：第二项即 $\frac34x_2^2$ 的积分应为 $(3/4)(1/2)(9/64)(1)(1/2)=27/1024$，原文写为 $7/1024$. 因而下面所求总概率应为 $171/1024$，而非原文的 $151/1024$. 此处保留原数值.} 所以
$$P\left(X_1<\frac12,X_2<\frac34,X_4>\frac12\right)=\frac3{256}+\frac7{1024}+\frac3{64}+\frac{21}{256}=\frac{151}{1024}.$$
\sourcepageanchor{178}

利用\eqref{eq:4.6.4}，将变量 $x_3,x_4$ 积掉，即可求得 $(X_1,X_2)$ 的边缘 pdf：
$$\begin{aligned}
f(x_1,x_2)&=\int_{-\infty}^{\infty}\int_{-\infty}^{\infty}f(x_1,x_2,x_3,x_4)\,dx_3\,dx_4\\
&=\int_0^1\int_0^1\frac34(x_1^2+x_2^2+x_3^2+x_4^2)\,dx_3\,dx_4
=\frac34(x_1^2+x_2^2)+\frac12,
\end{aligned}$$
其中 $0<x_1<1$，$0<x_2<1$. 任何只涉及 $X_1,X_2$ 的概率或期望，都可用此边缘 pdf 计算. 例如，
$$\begin{aligned}
EX_1X_2&=\int_{-\infty}^{\infty}\int_{-\infty}^{\infty}x_1x_2 f(x_1,x_2)\,dx_1\,dx_2\\
&=\int_0^1\int_0^1x_1x_2\left(\frac34(x_1^2+x_2^2)+\frac12\right)\,dx_1\,dx_2\\
&=\int_0^1\int_0^1\left(\frac34x_1^3x_2+\frac34x_1x_2^3+\frac12x_1x_2\right)\,dx_1\,dx_2\\
&=\int_0^1\left(\frac3{16}x_2+\frac38x_2^3+\frac14x_2\right)\,dx_2
=\frac3{32}+\frac3{32}+\frac18=\frac5{16}.
\end{aligned}$$
对任意满足 $0<x_1<1$、$0<x_2<1$ 的 $(x_1,x_2)$，有 $f(x_1,x_2)>0$，可以用\eqref{eq:4.6.6} 求给定 $X_1=x_1,X_2=x_2$ 时 $(X_3,X_4)$ 的条件 pdf. 对这样的 $(x_1,x_2)$，当 $0<x_3<1$、$0<x_4<1$ 时，$f(x_1,x_2,x_3,x_4)>0$，相应的条件 pdf 为
$$\begin{aligned}
f(x_3,x_4\mid x_1,x_2)&=\frac{f(x_1,x_2,x_3,x_4)}{f(x_1,x_2)}\\
&=\frac{\frac34(x_1^2+x_2^2+x_3^2+x_4^2)}{\frac34(x_1^2+x_2^2)+\frac12}
=\frac{x_1^2+x_2^2+x_3^2+x_4^2}{x_1^2+x_2^2+\frac23}.
\end{aligned}$$
例如，给定 $X_1=1/3$、$X_2=2/3$ 时 $(X_3,X_4)$ 的条件 pdf 为
$$f\left(x_3,x_4\mid X_1=\frac13,X_2=\frac23\right)
=\frac{(1/3)^2+(2/3)^2+x_3^2+x_4^2}{(1/3)^2+(2/3)^2+2/3}
=\frac5{11}+\frac9{11}x_3^2+\frac9{11}x_4^2.$$
用它可计算
$$\begin{aligned}
&P\left(X_3>\frac34,X_4<\frac12\mathrel{\Big|}X_1=\frac13,X_2=\frac23\right)\\
&\quad=\int_0^{1/2}\int_{3/4}^1\left(\frac5{11}+\frac9{11}x_3^2+\frac9{11}x_4^2\right)\,dx_3\,dx_4\\
&\quad=\int_0^{1/2}\left(\frac5{44}+\frac{111}{704}+\frac9{44}x_4^2\right)\,dx_4\\
&\quad=\frac5{88}+\frac{111}{1408}+\frac3{352}=\frac{203}{1408}.
\end{aligned}$$
\end{example}
在给出离散多元随机向量的条件与边缘分布计算例子之前，先介绍一个重要的离散多元分布族. 它将二项分布族推广到每次试验有 $n$ 个（而不是两个）不同可能结果的情形.
\sourcepageanchor{179}

\begin{definition}\label{def:4.6.2}\index{SI0348@多项}
设 $n,m$ 为正整数，$p_1,\ldots,p_n$ 满足 $0\le p_i\le1$，$i=1,\ldots,n$，且 $\sum_{i=1}^np_i=1$. 如果随机向量 $(X_1,\ldots,X_n)$ 的联合 pmf 在每个 $x_i$ 均为非负整数且 $\sum_{i=1}^nx_i=m$ 的集合上为
$$f(x_1,\ldots,x_n)=\frac{m!}{x_1!\cdots x_n!}p_1^{x_1}\cdots p_n^{x_n}=m!\prod_{i=1}^n\frac{p_i^{x_i}}{x_i!},$$
则称其服从试验次数为 $m$、格概率（cell probabilities）为 $p_1,\ldots,p_n$ 的多项分布（multinomial distribution）.
\end{definition}
多项分布是下列实验的模型. 实验包含 $m$ 次独立试验，每次试验得到 $n$ 个不同可能结果之一. 在每次试验中，第 $i$ 个结果出现的概率均为 $p_i$，$X_i$ 统计这 $m$ 次试验中第 $i$ 个结果出现的次数. 当 $n=2$ 时，这就是二项实验：每次有 $n=2$ 个可能结果，$X_1$ 统计 $m$ 次试验中的“成功”次数，$X_2=m-X_1$ 统计“失败”次数. 在一般的多项实验中，要统计 $n$ 种不同的可能结果.
\begin{example}[多元 pmf]\label{ex:4.6.3}
考虑掷一颗六面骰子十次. 假设骰子不均匀，出现 1 的概率为 $1/21$，出现 2 的概率为 $2/21$，一般地，出现 $i$ 的概率为 $i/21$. 考虑随机向量 $(X_1,\ldots,X_6)$，其中 $X_i$ 统计十次投掷中 $i$ 出现的次数. 则该向量服从多项分布，试验次数 $m=10$，可能结果数 $n=6$，格概率为 $p_1=1/21,p_2=2/21,\ldots,p_6=6/21$. 用\autoref{def:4.6.2} 中的公式，可计算出现四个 6、三个 5、两个 4 和一个 3 的概率：
$$\begin{aligned}
f(0,0,1,2,3,4)&=\frac{10!}{0!0!1!2!3!4!}\left(\frac1{21}\right)^0\left(\frac2{21}\right)^0
\left(\frac3{21}\right)^1\left(\frac4{21}\right)^2\left(\frac5{21}\right)^3\left(\frac6{21}\right)^4\\
&=.0059.
\end{aligned}$$
\end{example}
因子 $m!/(x_1!\cdots x_n!)$ 称为多项系数（multinomial coefficient）. 它表示将 $m$ 个物体分成 $n$ 组、第一组有 $x_1$ 个、第二组有 $x_2$ 个、……、第 $n$ 组有 $x_n$ 个的方法数. 二项式定理（\autoref{thm:3.2.2}）的一个推广是多项式定理.\index{SI0712@多项式系数}\semanticindexanchor{712:179}
\begin{theorem}[多项式定理]\label{thm:4.6.4}\index{SI0713@多项式定理}
设 $m,n$ 为正整数，$\mathscr A$ 为所有每个 $x_i$ 均为非负整数且 $\sum_{i=1}^nx_i=m$ 的向量 $\mathbf x=(x_1,\ldots,x_n)$ 组成的集合. 对任意实数 $p_1,\ldots,p_n$，
$$(p_1+\cdots+p_n)^m=\sum_{\mathbf x\in\mathscr A}\frac{m!}{x_1!\cdots x_n!}p_1^{x_1}\cdots p_n^{x_n}.$$
\end{theorem}
\autoref{thm:4.6.4} 表明多项 pmf 的和为 1. 集合 $\mathscr A$ 就是\autoref{def:4.6.2} 中概率为正的点集. 根据该定理，在这些点上将 pmf 求和，得到 $(p_1+\cdots+p_n)^m=1^m=1$.

现在考虑多项模型的一些边缘与条件分布. 对单个坐标 $X_i$，若把第 $i$ 个结果称为“成功”，其他结果称为“失败”，则 $X_i$ 统计 $m$ 次独立试验中的成功次数，每次成功概率为 $p_i$. 因而，$X_i$ 应服从 $\operatorname{binomial}(m,p_i)$. 为验证这一点，应根据\eqref{eq:4.6.5} 计算 $X_i$ 的边缘分布. 例如，考虑 $X_n$ 的边缘 pmf. 固定 $x_n\in\{0,1,\ldots,n\}$，\footnote{译注：原文取值上界 $n$ 应为试验次数 $m$；$n$ 是类别数.} 要计算 $f(x_n)$，须对所有可能的 $(x_1,\ldots,x_{n-1})$ 求和.\index{SI0269@与多项的关系}\semanticindexanchor{269:179}\index{SI0350@与二项的关系}\semanticindexanchor{350:179}
\sourcepageanchor{180}

也就是说，对所有 $x_i$ 均为非负整数且 $\sum_{i=1}^{n-1}x_i=m-x_n$ 的 $(x_1,\ldots,x_{n-1})$ 求和. 记此集合为 $\mathscr B$，则
$$\begin{aligned}
f(x_n)&=\sum_{(x_1,\ldots,x_{n-1})\in\mathscr B}\frac{m!}{x_1!\cdots x_n!}p_1^{x_1}\cdots p_n^{x_n}\\
&=\sum_{(x_1,\ldots,x_{n-1})\in\mathscr B}\frac{m!}{x_1!\cdots x_n!}p_1^{x_1}\cdots p_n^{x_n}
\frac{(m-x_n)!}{(m-x_n)!}\frac{(1-p_n)^{m-x_n}}{(1-p_n)^{m-x_n}}\\
&=\frac{m!}{x_n!(m-x_n)!}p_n^{x_n}(1-p_n)^{m-x_n}\\
&\quad\times\sum_{(x_1,\ldots,x_{n-1})\in\mathscr B}\frac{(m-x_n)!}{x_1!\cdots x_{n-1}!}
\left(\frac{p_1}{1-p_n}\right)^{x_1}\cdots\left(\frac{p_{n-1}}{1-p_n}\right)^{x_{n-1}}.
\end{aligned}$$
利用 $x_1+\cdots+x_{n-1}=m-x_n$、$p_1+\cdots+p_{n-1}=1-p_n$ 以及\autoref{thm:4.6.4}，可见最后的和式等于 1. 因而，$X_n$ 的边缘分布是 $\operatorname{binomial}(m,p_n)$. 类似论证表明，其余每个坐标的边缘分布也都是二项分布.

已知 $X_n=x_n$ 时，必有 $m-x_n$ 次试验的结果属于前 $n-1$ 类. 向量 $(X_1,\ldots,X_{n-1})$ 统计这 $m-x_n$ 次试验中每类结果的次数. 因而，给定 $X_n=x_n$ 时，$(X_1,\ldots,X_{n-1})$ 似乎应服从多项分布；事实确实如此. 根据\eqref{eq:4.6.6}，条件 pmf 为
$$\begin{aligned}
f(x_1,\ldots,x_{n-1}\mid x_n)&=\frac{f(x_1,\ldots,x_n)}{f(x_n)}\\
&=\frac{\dfrac{m!}{x_1!\cdots x_n!}p_1^{x_1}\cdots p_n^{x_n}}
{\dfrac{m!}{x_n!(m-x_n)!}p_n^{x_n}(1-p_n)^{m-x_n}}\\
&=\frac{(m-x_n)!}{x_1!\cdots x_{n-1}!}\left(\frac{p_1}{1-p_n}\right)^{x_1}\cdots\left(\frac{p_{n-1}}{1-p_n}\right)^{x_{n-1}}.
\end{aligned}$$
这是试验次数为 $m-x_n$、格概率为 $p_1/(1-p_n),\ldots,p_{n-1}/(1-p_n)$ 的多项分布的 pmf.\footnote{译注：此处除以 $1-p_n$ 的推导要求 $p_n<1$，条件分布还要求 $P(X_n=x_n)>0$. 当 $p_n=1$ 时，$X_n=m$、其余坐标恒为 0，应单独按退化分布处理.} 事实上，给定其余坐标的值时，$(X_1,\ldots,X_n)$ 中任意一部分坐标的条件分布都是多项分布.

从条件分布可以看出，向量 $(X_1,\ldots,X_n)$ 的各坐标有关系. 尤其是 $m$ 固定时，必然有某种负相关. 事实上，所有两两协方差都为负，由下式给出（\autoref{exercise:4.39}）：\footnote{译注：原文期望式中应减去均值 $mp_i,mp_j$，而非 $p_i,p_j$；即 $\operatorname{Cov}(X_i,X_j)=E[(X_i-mp_i)(X_j-mp_j)]=-mp_ip_j$，$i\ne j$. 当某个格概率为 0 时协方差为 0，故“都为负”在允许边界概率的定义下应为“都非正”.}\index{SI0349@协方差}\semanticindexanchor{349:180}
$$\operatorname{Cov}(X_i,X_j)=E[(X_i-p_i)(X_j-p_j)]=-mp_ip_j.$$
因此，成功概率较大的变量之间负相关更强. 这很合理，因为变量总和固定为 $m$，一个开始变大时，另一个就不太可能也变大.
\begin{definition}\label{def:4.6.5}\index{SI0554@随机变量}
设 $\mathbf X_1,\ldots,\mathbf X_n$ 是联合 pdf 或 pmf 为 $f(\mathbf x_1,\ldots,\mathbf x_n)$ 的随机向量，以 $f_{\mathbf X_i}(\mathbf x_i)$ 表示 $\mathbf X_i$ 的边缘 pdf 或 pmf. 如果对每个 $(\mathbf x_1,\ldots,\mathbf x_n)$，均有
$$f(\mathbf x_1,\ldots,\mathbf x_n)=f_{\mathbf X_1}(\mathbf x_1)\cdots f_{\mathbf X_n}(\mathbf x_n)=\prod_{i=1}^n f_{\mathbf X_i}(\mathbf x_i),$$
则称 $\mathbf X_1,\ldots,\mathbf X_n$ 为相互独立的随机向量（mutually independent random vectors）. 如果这些向量都是一维的，则称 $X_1,\ldots,X_n$ 为相互独立的随机变量.
\end{definition}
如果 $X_1,\ldots,X_n$ 相互独立，知道部分坐标的值，不会提供其余坐标取值的信息. 利用\autoref{def:4.6.5}，可以证明：给定其余坐标的值时，任意一部分坐标的条件分布，
\sourcepageanchor{181}

都等于这部分坐标的边缘分布. 相互独立蕴含任意一对 $X_i,X_j$ 两两独立（pairwise independent），即它们的二元边缘 pdf 或 pmf $f(x_i,x_j)$ 满足\autoref{def:4.2.5}. 但相互独立比两两独立要求更多. 和\autoref{ex:1.3.11} 一样，可以指定 $(X_1,\ldots,X_n)$ 的一个概率分布，使每对 $(X_i,X_j)$ 都两两独立，但 $X_1,\ldots,X_n$ 不相互独立.

相互独立的随机变量有许多良好性质. 以下定理的证明类似于\autoref{sec:4.2} 和\autoref{sec:4.3} 中相应结论的证明.
\begin{theorem}[\autoref{thm:4.2.10} 的推广]\label{thm:4.6.6}\index{SI0427@独立性}\index{SI0553@期望值}
设 $X_1,\ldots,X_n$ 相互独立，$g_1,\ldots,g_n$ 为实值函数，且 $g_i(x_i)$ 只依赖于 $x_i$，$i=1,\ldots,n$. 则
$$E(g_1(X_1)\cdots g_n(X_n))=(Eg_1(X_1))\cdots(Eg_n(X_n)).$$
\end{theorem}
\begin{theorem}[\autoref{thm:4.2.12} 的推广]\label{thm:4.6.7}\index{SI0695@和}
设 $X_1,\ldots,X_n$ 相互独立，mgf 分别为 $M_{X_1}(t)$、$\ldots$、$M_{X_n}(t)$. 令 $Z=X_1+\cdots+X_n$，则 $Z$ 的 mgf 为
$$M_Z(t)=M_{X_1}(t)\cdots M_{X_n}(t).$$
特别地，如果 $X_1,\ldots,X_n$ 同分布，mgf 均为 $M_X(t)$，则
$$M_Z(t)=(M_X(t))^n.$$
\end{theorem}
\begin{example}[Gamma 变量之和的 mgf]\label{ex:4.6.8}\index{SI0316@和}
假设 $X_1,\ldots,X_n$ 相互独立，$X_i$ 的分布为 $\operatorname{gamma}(\alpha_i,\beta)$. 由\autoref{ex:2.3.8}，$\operatorname{gamma}(\alpha,\beta)$ 分布的 mgf 为 $M(t)=(1-\beta t)^{-\alpha}$. 因而，若 $Z=X_1+\cdots+X_n$，则
$$\begin{aligned}
M_Z(t)&=M_{X_1}(t)\cdots M_{X_n}(t)\\
&=(1-\beta t)^{-\alpha_1}\cdots(1-\beta t)^{-\alpha_n}
=(1-\beta t)^{-(\alpha_1+\cdots+\alpha_n)}.
\end{aligned}$$
这是 $\operatorname{gamma}(\alpha_1+\cdots+\alpha_n,\beta)$ 分布的 mgf. 所以，具有共同尺度参数 $\beta$ 的独立 Gamma 随机变量之和仍服从 Gamma 分布.
\end{example}
如果考虑独立随机变量的线性函数之和，就可得到\autoref{thm:4.6.7} 的一个推广.
\begin{corollary}\label{cor:4.6.9}
设 $X_1,\ldots,X_n$ 相互独立，mgf 为 $M_{X_1}(t),\ldots,M_{X_n}(t)$；$a_1,\ldots,a_n$ 和 $b_1,\ldots,b_n$ 为固定常数. 令 $Z=(a_1X_1+b_1)+\cdots+(a_nX_n+b_n)$，则
$$M_Z(t)=e^{t(\sum b_i)}M_{X_1}(a_1t)\cdots M_{X_n}(a_nt).$$
\end{corollary}
\begin{proof}
由定义，$Z$ 的 mgf 为
$$\begin{aligned}
M_Z(t)&=Ee^{tZ}\\
&=Ee^{t\sum(a_iX_i+b_i)}\\
&=e^{t(\sum b_i)}E(e^{ta_1X_1}\cdots e^{ta_nX_n})&&\text{（指数与期望的性质）}\\
&=e^{t(\sum b_i)}M_{X_1}(a_1t)\cdots M_{X_n}(a_nt),&&\text{（\autoref{thm:4.6.6}）}
\end{aligned}$$
即所欲证.
\end{proof}
毫无疑问，\autoref{cor:4.6.9} 最重要的应用是正态随机变量的情形：独立正态随机变量的线性组合服从正态分布.
\sourcepageanchor{182}

\begin{corollary}\label{cor:4.6.10}\index{SI0362@线性组合}
设 $X_1,\ldots,X_n$ 相互独立，$X_i\sim n(\mu_i,\sigma_i^2)$，$a_1,\ldots,a_n$ 和 $b_1,\ldots,b_n$ 为固定常数. 则
$$Z=\sum_{i=1}^n(a_iX_i+b_i)\sim n\left(\sum_{i=1}^n(a_i\mu_i+b_i),\sum_{i=1}^na_i^2\sigma_i^2\right).$$
\end{corollary}
\begin{proof}
回顾 $n(\mu,\sigma^2)$ 随机变量的 mgf 为 $M(t)=e^{\mu t+\sigma^2t^2/2}$. 代入\autoref{cor:4.6.9}，得到
$$\begin{aligned}
M_Z(t)&=e^{t(\sum b_i)}e^{\mu_1a_1t+\sigma_1^2a_1^2t^2/2}\cdots e^{\mu_na_nt+\sigma_n^2a_n^2t^2/2}\\
&=e^{(\sum(a_i\mu_i+b_i))t+(\sum a_i^2\sigma_i^2)t^2/2},
\end{aligned}$$
即所述正态分布的 mgf.
\end{proof}
\begin{theorem}[\autoref{lem:4.2.7} 的推广]\label{thm:4.6.11}
设 $\mathbf X_1,\ldots,\mathbf X_n$ 为随机向量. 它们相互独立，当且仅当存在函数 $g_i(\mathbf x_i)$，$i=1,\ldots,n$，使联合 pdf 或 pmf 可以写为
$$f(\mathbf x_1,\ldots,\mathbf x_n)=g_1(\mathbf x_1)\cdots g_n(\mathbf x_n).$$
\end{theorem}
\begin{theorem}[\autoref{thm:4.3.5} 的推广]\label{thm:4.6.12}
设 $\mathbf X_1,\ldots,\mathbf X_n$ 为独立随机向量，$g_i(\mathbf x_i)$ 仅依赖于 $\mathbf x_i$，$i=1,\ldots,n$. 则随机变量 $U_i=g_i(\mathbf X_i)$，$i=1,\ldots,n$，相互独立.
\end{theorem}
本节最后介绍求随机向量变换分布的方法的推广. 我们将推广\eqref{eq:4.3.6}，用原随机向量的 pdf 表示新随机向量的 pdf. 要完全理解本节剩余内容，需要一些矩阵代数知识，例如可参阅 \hyperref[bib:source283]{Searle (1982)}. 尤其需要计算矩阵的行列式. 这是本书唯一需要此类知识的地方.

设 $(X_1,\ldots,X_n)$ 的 pdf 为 $f_{\mathbf X}(x_1,\ldots,x_n)$，令 $\mathscr A=\{\mathbf x:f_{\mathbf X}(\mathbf x)>0\}$. 考虑新随机向量 $(U_1,\ldots,U_n)$，由 $U_1=g_1(X_1,\ldots,X_n)$、$U_2=g_2(X_1,\ldots,X_n)$、……、$U_n=g_n(X_1,\ldots,X_n)$ 定义. 假设 $A_0,A_1,\ldots,A_k$ 构成 $\mathscr A$ 的一个划分，并满足下列性质. 集合 $A_0$ 可以为空，且 $P((X_1,\ldots,X_n)\in A_0)=0$. 对每个 $i=1,2,\ldots,k$，变换 $(U_1,\ldots,U_n)=(g_1(\mathbf X),\ldots,g_n(\mathbf X))$ 从 $A_i$ 一一映满到 $\mathscr B$. 对每个 $i$，都可以求出从 $\mathscr B$ 到 $A_i$ 的逆函数，记为 $x_1=h_{1i}(u_1,\ldots,u_n)$、$x_2=h_{2i}(u_1,\ldots,u_n)$、……、$x_n=h_{ni}(u_1,\ldots,u_n)$. 对 $(u_1,\ldots,u_n)\in\mathscr B$，此第 $i$ 个逆变换给出唯一的 $(x_1,\ldots,x_n)\in A_i$，使 $(u_1,\ldots,u_n)=(g_1(x_1,\ldots,x_n),\ldots,g_n(x_1,\ldots,x_n))$. 记根据第 $i$ 个逆变换计算的 Jacobian 为 $J_i$，即\index{SI0248@行列式}\semanticindexanchor{248:182}
$$J_i=\begin{vmatrix}
\dfrac{\partial x_1}{\partial u_1}&\dfrac{\partial x_1}{\partial u_2}&\cdots&\dfrac{\partial x_1}{\partial u_n}\\[4pt]
\dfrac{\partial x_2}{\partial u_1}&\dfrac{\partial x_2}{\partial u_2}&\cdots&\dfrac{\partial x_2}{\partial u_n}\\[4pt]
\vdots&\vdots&\ddots&\vdots\\
\dfrac{\partial x_n}{\partial u_1}&\dfrac{\partial x_n}{\partial u_2}&\cdots&\dfrac{\partial x_n}{\partial u_n}
\end{vmatrix}
=\begin{vmatrix}
\dfrac{\partial h_{1i}(\mathbf u)}{\partial u_1}&\dfrac{\partial h_{1i}(\mathbf u)}{\partial u_2}&\cdots&\dfrac{\partial h_{1i}(\mathbf u)}{\partial u_n}\\[4pt]
\dfrac{\partial h_{2i}(\mathbf u)}{\partial u_1}&\dfrac{\partial h_{2i}(\mathbf u)}{\partial u_2}&\cdots&\dfrac{\partial h_{2i}(\mathbf u)}{\partial u_n}\\[4pt]
\vdots&\vdots&\ddots&\vdots\\
\dfrac{\partial h_{ni}(\mathbf u)}{\partial u_1}&\dfrac{\partial h_{ni}(\mathbf u)}{\partial u_2}&\cdots&\dfrac{\partial h_{ni}(\mathbf u)}{\partial u_n}
\end{vmatrix},$$
这是一个 $n\times n$ 矩阵的行列式. 假设这些 Jacobian 在 $\mathscr B$ 上都不恒为零，那么对于 $\mathbf u\in\mathscr B$，联合 pdf 为
\begin{equation}\label{eq:4.6.7}\index{SI0836@多元变换}
f_{\mathbf U}(u_1,\ldots,u_n)=\sum_{i=1}^k f_{\mathbf X}(h_{1i}(u_1,\ldots,u_n),\ldots,h_{ni}(u_1,\ldots,u_n))|J_i|.
\end{equation}
\sourcepageanchor{183}

\begin{example}[多元变量替换]\label{ex:4.6.13}
设 $(X_1,X_2,X_3,X_4)$ 的联合 pdf 为
$$f_{\mathbf X}(x_1,x_2,x_3,x_4)=24e^{-x_1-x_2-x_3-x_4},\qquad0<x_1<x_2<x_3<x_4<\infty.$$
考虑变换
$$U_1=X_1,\quad U_2=X_2-X_1,\quad U_3=X_3-X_2,\quad U_4=X_4-X_3.$$
此变换将 $\mathscr A$ 映满到 $\mathscr B=\{\mathbf u:0<u_i<\infty,\ i=1,2,3,4\}$. 该变换是一一的，所以 $k=1$，逆变换为
$$X_1=U_1,\quad X_2=U_1+U_2,\quad X_3=U_1+U_2+U_3,\quad X_4=U_1+U_2+U_3+U_4.$$
逆变换的 Jacobian 为
$$J=\begin{vmatrix}1&0&0&0\\1&1&0&0\\1&1&1&0\\1&1&1&1\end{vmatrix}=1.$$
因为矩阵是三角矩阵，其行列式等于对角元素的乘积. 因而，由\eqref{eq:4.6.7} 得到\footnote{译注：原文下面化简后的指数写为 $-2\mu_3$；按上一行的变量替换应为 $-2u_3$，原文未定义 $\mu_3$. 此处保留原式.}
$$\begin{aligned}
f_{\mathbf U}(u_1,\ldots,u_4)&=24e^{-u_1-(u_1+u_2)-(u_1+u_2+u_3)-(u_1+u_2+u_3+u_4)}\\
&=24e^{-4u_1-3u_2-2\mu_3-u_4},\qquad\mathbf u\in\mathscr B.
\end{aligned}$$
由此可以计算 $U_1,U_2,U_3,U_4$ 的边缘 pdf，结果为 $f_U(u_i)=(5-i)e^{-(5-i)u_i}$，$0<u_i$. 即 $U_i\sim\operatorname{exponential}(1/(5-i))$. 从\autoref{thm:4.6.11} 可见，$U_1,U_2,U_3,U_4$ 相互独立.
\end{example}
\autoref{ex:4.6.13} 的模型可以如下产生. 假设 $Y_1,Y_2,Y_3,Y_4$ 相互独立，每个都服从 $\operatorname{exponential}(1)$ 分布. 定义 $X_1=\min(Y_1,Y_2,Y_3,Y_4)$，$X_2=$ 其中第二小的值，$X_3=$ 其中第二大的值，$X_4=\max(Y_1,Y_2,Y_3,Y_4)$. \autoref{sec:5.5} 将这些变量称为次序统计量（order statistics），并会证明它们的联合 pdf 就是\autoref{ex:4.6.13} 所给出的 pdf. 例中定义的 $U_2,U_3,U_4$ 称为次序统计量之间的间距（spacings）. 该例表明，对这些指数随机变量 $(Y_1,\ldots,Y_n)$，次序统计量的间距相互独立，且也服从指数分布.\index{SI0737@间距}\semanticindexanchor{737:183}



\subsection{不等式}\label{sec:4.7}
\autoref{sec:3.6} 介绍了用概率论证导出的不等式. 本节介绍的不等式适用于概率和期望，但论证基于函数与数的性质.
\subsubsection{数值不等式}\label{sec:4.7.1}
本小节的不等式虽然常以期望的形式表述，却主要依赖数的性质. 事实上，它们都基于下面这个简单引理.
\begin{lemma}\label{lem:4.7.1}
设 $a,b$ 为任意正数，$p,q$ 为满足
\begin{equation}\label{eq:4.7.1}
\frac1p+\frac1q=1
\end{equation}
的任意正数（因而必大于 1）. 则
\begin{equation}\label{eq:4.7.2}
\frac1pa^p+\frac1qb^q\ge ab,
\end{equation}
等号成立当且仅当 $a^p=b^q$.
\end{lemma}
\sourcepageanchor{184}

\begin{proof}
固定 $b$，考虑函数
$$g(a)=\frac1pa^p+\frac1qb^q-ab.$$
为求其最小值，求导并令导数为 0：
$$\frac d{da}g(a)=0\quad\Longrightarrow\quad a^{p-1}-b=0\quad\Longrightarrow\quad b=a^{p-1}.$$
检查二阶导数可知，这的确是极小值. 此时函数值为
$$\begin{aligned}
\frac1pa^p+\frac1q(a^{p-1})^q-aa^{p-1}
&=\frac1pa^p+\frac1qa^p-a^p&&\text{（由\eqref{eq:4.7.1}，$(p-1)q=p$）}\\
&=0.&&\text{（再次利用\eqref{eq:4.7.1}）}
\end{aligned}$$
所以最小值为 0，\eqref{eq:4.7.2} 得证. 最小值点唯一（为什么？），因此等号仅当 $a^{p-1}=b$ 时成立；由\eqref{eq:4.7.1}，这又等价于 $a^p=b^q$.
\end{proof}
第一个期望不等式是最常用、最重要的不等式之一，由此引理很容易得到.
\begin{theorem}[H\"older 不等式]\label{thm:4.7.2}\index{SI0568@Hölder}
设 $X,Y$ 为任意两个随机变量，$p,q$ 满足\eqref{eq:4.7.1}. 则
\begin{equation}\label{eq:4.7.3}
|EXY|\le E|XY|\le(E|X|^p)^{1/p}(E|Y|^q)^{1/q}.
\end{equation}
\end{theorem}
\begin{proof}
第一个不等式由 $-|XY|\le XY\le|XY|$ 及\autoref{thm:2.2.5} 得到. 为证明第二个不等式，定义
$$a=\frac{|X|}{(E|X|^p)^{1/p}},\qquad b=\frac{|Y|}{(E|Y|^q)^{1/q}}.$$
应用\autoref{lem:4.7.1}，得到
$$\frac1p\frac{|X|^p}{E|X|^p}+\frac1q\frac{|Y|^q}{E|Y|^q}
\ge\frac{|XY|}{(E|X|^p)^{1/p}(E|Y|^q)^{1/q}}.$$
对两边取期望. 左边的期望为 1，整理即得\eqref{eq:4.7.3}.
\end{proof}
H\"older 不等式最著名的特例也许是 $p=q=2$，称为 Cauchy--Schwarz 不等式.
\begin{theorem}[Cauchy--Schwarz 不等式]\label{thm:4.7.3}\index{SI0564@Cauchy–Schwarz}
对任意两个随机变量 $X,Y$，
\begin{equation}\label{eq:4.7.4}
|EXY|\le E|XY|\le(E|X|^2)^{1/2}(E|Y|^2)^{1/2}.
\end{equation}
\end{theorem}
\begin{example}[协方差不等式]\label{ex:4.7.4}
若 $X,Y$ 的均值分别为 $\mu_X,\mu_Y$，方差分别为 $\sigma_X^2,\sigma_Y^2$，则应用 Cauchy--Schwarz 不等式，
$$E|(X-\mu_X)(Y-\mu_Y)|\le\{E(X-\mu_X)^2\}^{1/2}\{E(Y-\mu_Y)^2\}^{1/2}.$$
两边平方，并使用统计学记号，得到
$$(\operatorname{Cov}(X,Y))^2\le\sigma_X^2\sigma_Y^2.$$
\sourcepageanchor{185}

回顾相关系数 $\rho$ 的定义，我们已经证明 $0\le\rho^2\le1$. 此外，\autoref{lem:4.7.1} 的等号条件仍然适用：这里仅当存在常数 $c$ 使 $X-\mu_X=c(Y-\mu_Y)$ 时等号成立. 也就是说，相关系数为 $\pm1$ 当且仅当 $X$ 和 $Y$ 线性相关. 比较一下，这个证明比尚未掌握 Cauchy--Schwarz 不等式时\autoref{thm:4.5.7} 的证明简单多少.
\end{example}
H\"older 不等式的其他一些特例也常有用. 在\eqref{eq:4.7.3} 中令 $Y\equiv1$，得到
\begin{equation}\label{eq:4.7.5}
E|X|\le\{E(|X|^p)\}^{1/p},\qquad1<p<\infty.
\end{equation}
对 $1<r<p$，将\eqref{eq:4.7.5} 中的 $|X|$ 替换为 $|X|^r$，得到
$$E|X|^r\le\{E(|X|^{pr})\}^{1/p}.$$
再令 $s=pr$（注意 $s>r$），整理各项可得
\begin{equation}\label{eq:4.7.6}\index{SI0570@Liapounov}
\{E|X|^r\}^{1/r}\le\{E(|X|^s)\}^{1/s},\qquad1<r<s<\infty,
\end{equation}
这称为 Liapounov 不等式.\footnote{译注：推导前写的 $r<p$ 是多余限制；若坚持该限制，令 $s=pr$ 只直接覆盖 $s>r^2$. 实际上对任意 $r>0$、$p>1$ 都可作上述替换，因此取 $p=s/r>1$ 即得到所述范围.}

下一个冠名不等式与 H\"older 不等式思想相近，事实上也由它推出.
\begin{theorem}[Minkowski 不等式]\label{thm:4.7.5}\index{SI0575@Minkowski}
设 $X,Y$ 为任意两个随机变量，则对 $1\le p<\infty$，
\begin{equation}\label{eq:4.7.7}
[E|X+Y|^p]^{1/p}\le[E|X|^p]^{1/p}+[E|Y|^p]^{1/p}.
\end{equation}
\end{theorem}
\begin{proof}
写成
\begin{equation}\label{eq:4.7.8}\index{SI0578@三角}
\begin{aligned}
E|X+Y|^p&=E\bigl(|X+Y||X+Y|^{p-1}\bigr)\\
&\le E\bigl(|X||X+Y|^{p-1}\bigr)+E\bigl(|Y||X+Y|^{p-1}\bigr),
\end{aligned}
\end{equation}
这里使用了 $|X+Y|\le|X|+|Y|$（三角不等式，见\autoref{exercise:4.64}）. 对\eqref{eq:4.7.8} 右边的每个期望应用 H\"older 不等式，得到
$$\begin{aligned}
E(|X+Y|^p)&\le[E(|X|^p)]^{1/p}[E|X+Y|^{q(p-1)}]^{1/q}\\
&\quad+[E(|Y|^p)]^{1/p}[E|X+Y|^{q(p-1)}]^{1/q},
\end{aligned}$$
其中 $q$ 满足 $1/p+1/q=1$. 两边除以 $[E(|X+Y|^{q(p-1)})]^{1/q}$，注意 $q(p-1)=p$ 且 $1-1/q=1/p$，即得\eqref{eq:4.7.7}.\footnote{译注：此证明中的有限共轭指数 $q$ 要求 $p>1$；$p=1$ 直接由三角不等式取期望得到. 除法还需所除期望有限且非零；零的情形直接成立，有限性可由 $|X+Y|^p\le2^{p-1}(|X|^p+|Y|^p)$ 保证，右端矩无穷的情形作扩展实数解释.}
\end{proof}
上述定理也适用于没有显式期望记号的数值和式. 例如，对数 $a_i,b_i$，$i=1,\ldots,n$，不等式
\begin{equation}\label{eq:4.7.9}
\sum_{i=1}^n|a_ib_i|\le\left(\sum_{i=1}^na_i^p\right)^{1/p}\left(\sum_{i=1}^nb_i^q\right)^{1/q},\qquad\frac1p+\frac1q=1,
\end{equation}
是 H\"older 不等式的一种形式.\footnote{译注：若 $a_i,b_i$ 是任意实数，右端应为 $\sum|a_i|^p$ 与 $\sum|b_i|^q$；原式缺少绝对值，只在非负等适当限制下可照写. 此处保留原式.} 要证明\eqref{eq:4.7.9}，可以形式化地引入分别取值 $a_1,\ldots,a_n$ 和 $b_1,\ldots,b_n$ 的随机变量，再构造期望. \autoref{ex:4.7.8} 将使用这种做法.

\eqref{eq:4.7.9} 的一个重要特例是 $b_i\equiv1$、$p=q=2$，此时
$$\frac1n\left(\sum_{i=1}^n|a_i|\right)^2\le\sum_{i=1}^na_i^2.$$
\sourcepageanchor{186}

\subsubsection{函数不等式}\label{sec:4.7.2}
本小节的不等式依赖于实值函数的性质，而非任何统计性质；但它们在许多情形下都非常有用. 其中最有用的不等式之一是适用于凸函数的 Jensen 不等式.
\begin{figure}[htbp]
\centering\includegraphics[width=.76\linewidth]{assets/chapter4/fig4-7-1.pdf}
\caption{凸函数及其在 $x_1$、$x_2$ 处的切线}\label{fig:4.7.1}
\end{figure}

\begin{definition}\label{def:4.7.6}\index{SI0121@凹函数}\index{SI0212@凸函数}
若对所有 $x,y$ 及 $0<\lambda<1$，都有 $g(\lambda x+(1-\lambda)y)\le\lambda g(x)+(1-\lambda)g(y)$，则称函数 $g(x)$ 为凸函数（convex function）. 若 $-g(x)$ 为凸函数，则称 $g(x)$ 为凹函数（concave function）.
\end{definition}
直观上，可以把凸函数看成“盛得住水”的函数，即形状像碗（例如 $g(x)=x^2$ 是凸的）；凹函数则会“把水洒出来”（例如 $g(x)=\log x$ 是凹的）. 更正式地说，凸函数位于连接图像上任意两点的直线下方，见\autoref{fig:4.7.1}. 当 $\lambda$ 从 0 变到 1 时，$\lambda g(x_1)+(1-\lambda)g(x_2)$ 描述连接 $g(x_1)$ 与 $g(x_2)$ 的线段. 若 $g(x)$ 为凸函数，此线段位于图像之上. 此外，凸函数位于其所有切线之上，\autoref{fig:4.7.1} 也显示了这一点；这正是 Jensen 不等式的基础.
\begin{theorem}[Jensen 不等式]\label{thm:4.7.7}\index{SI0569@Jensen}
对任意随机变量 $X$，若 $g(x)$ 为凸函数，则
$$Eg(X)\ge g(EX).$$
等号成立当且仅当对每条在 $x=EX$ 处与 $g(x)$ 相切的直线 $a+bx$，都有 $P(g(X)=a+bX)=1$.
\end{theorem}
\begin{proof}
为证明不等式，令 $l(x)$ 是 $g(x)$ 在 $g(EX)$ 处的切线. 回顾 $EX$ 是常数. 将其写为 $l(x)=a+bx$. \autoref{fig:4.7.2} 图示了这一情形.

由 $g$ 的凸性，$g(x)\ge a+bx$. 因为取期望保持不等式，有
$$\begin{aligned}
Eg(X)&\ge E(a+bX)\\
&=a+bEX&&\text{（期望的线性性，\autoref{thm:2.2.5}）}\\
&=l(EX)&&\text{（$l(x)$ 的定义）}\\
&=g(EX)&&\text{（$l$ 在 $EX$ 处相切）}
\end{aligned}$$
即所欲证.
\sourcepageanchor{187}

\begin{figure}[htbp]
\centering\includegraphics[width=.74\linewidth]{assets/chapter4/fig4-7-2.pdf}
\caption{Jensen 不等式的图示}\label{fig:4.7.2}
\end{figure}

如果 $g(x)$ 是线性的，则由期望的性质（\autoref{thm:2.2.5}）可知等号成立. 必要性部分见\autoref{exercise:4.62}.
\end{proof}
Jensen 不等式的一个直接应用是 $EX^2\ge(EX)^2$，因为 $g(x)=x^2$ 为凸函数. 又如，对正数 $x$，$1/x$ 为凸函数，因此 $E(1/X)\ge1/EX$，这也是一个有用的应用.

判断二次可微函数的凸性很容易：若对所有 $x$，$g''(x)\ge0$，则 $g(x)$ 为凸函数；若对所有 $x$，$g''(x)\le0$，则为凹函数. Jensen 不等式同样适用于凹函数：若 $g$ 凹，则 $Eg(X)\le g(EX)$.
\needspace{6\baselineskip}
\begin{example}[均值的不等式]\label{ex:4.7.8}\index{SI0064@算术平均数}\index{SI0457@几何平均数}\index{SI0468@调和平均数}\index{SI0573@均值}
Jensen 不等式可用于证明三种不同均值之间的不等式. 设 $a_1,\ldots,a_n$ 为正数，定义
$$\begin{aligned}
a_A&=\frac1n(a_1+a_2+\cdots+a_n),&&\text{（算术平均，arithmetic mean）}\\
a_G&=[a_1a_2\cdots a_n]^{1/n},&&\text{（几何平均，geometric mean）}\\
a_H&=\frac1{\frac1n\left(\frac1{a_1}+\frac1{a_2}+\cdots+\frac1{a_n}\right)}.&&\text{（调和平均，harmonic mean）}
\end{aligned}$$
它们之间满足
$$a_H\le a_G\le a_A.$$
为应用 Jensen 不等式，令 $X$ 的取值范围为 $a_1,\ldots,a_n$，且 $P(X=a_i)=1/n$，$i=1,\ldots,n$.\footnote{译注：若若干 $a_i$ 相同，严格说应在指标 $i\in\{1,\ldots,n\}$ 上均匀抽取，再令 $X=a_i$；重复取值的概率须合并. 原文逐值写 $P(X=a_i)=1/n$ 默认了取值不同，所给均值不等式对重复取值仍成立.} 因为 $\log x$ 为凹函数，Jensen 不等式给出 $E(\log X)\le\log(EX)$，所以
$$\begin{aligned}
\log a_G&=\frac1n\sum_{i=1}^n\log a_i=E(\log X)\le\log(EX)\\
&=\log\left(\frac1n\sum_{i=1}^na_i\right)=\log a_A,
\end{aligned}$$
故 $a_G\le a_A$. 再次利用 $\log x$ 的凹性，
$$\log\frac1{a_H}=\log\left(\frac1n\sum_{i=1}^n\frac1{a_i}\right)
=\log E\frac1X\ge E\left(\log\frac1X\right)=-E(\log X).$$
因为 $E(\log X)=\log a_G$，故 $\log(1/a_H)\ge\log(1/a_G)$，即 $a_G\ge a_H$.
\end{example}
\sourcepageanchor{188}

下一个不等式仅利用协方差的定义，但有时很有用. 如果随机变量 $X$ 有有限均值 $\mu$，且 $g(x)$ 为非减函数，则
$$E(g(X)(X-\mu))\ge0,$$
因为
$$\begin{aligned}
&E(g(X)(X-\mu))\\
&=E\bigl(g(X)(X-\mu)I_{(-\infty,0)}(X-\mu)\bigr)+E\bigl(g(X)(X-\mu)I_{[0,\infty)}(X-\mu)\bigr)\\
&\ge E\bigl(g(\mu)(X-\mu)I_{(-\infty,0)}(X-\mu)\bigr)\\
&\qquad+E\bigl(g(\mu)(X-\mu)I_{[0,\infty)}(X-\mu)\bigr)\qquad\text{（因为 $g$ 非减）}\\
&=g(\mu)E(X-\mu)\\
&=0.
\end{aligned}$$
推广这一论证，可以建立下面的不等式（见\autoref{exercise:4.65}）.
\begin{theorem}[协方差不等式]\label{thm:4.7.9}\index{SI0566@协方差}
设 $X$ 为任意随机变量，$g,h$ 为使 $Eg(X)$、$Eh(X)$ 和 $E(g(X)h(X))$ 均存在的函数.
\begin{enumerate}
\item[(a)]\phantomsection\label{itm:thm:4.7.9:a} 若 $g(x)$ 非减、$h(x)$ 非增，则
$$E(g(X)h(X))\le(Eg(X))(Eh(X)).$$
\item[(b)]\phantomsection\label{itm:thm:4.7.9:b} 若 $g(x),h(x)$ 同为非减函数或同为非增函数，则
$$E(g(X)h(X))\ge(Eg(X))(Eh(X)).$$
\end{enumerate}
\end{theorem}
不等式背后的直观很简单：\hyperref[itm:thm:4.7.9:a]{(a)} 中 $g$ 与 $h$ 负相关，\hyperref[itm:thm:4.7.9:b]{(b)} 中二者正相关；不等式只是反映了这一事实. 协方差不等式的用途在于，无须使用高阶矩，就可以给期望作界.



\subsection{习题}\label{sec:4.8}
\setcounter{exercise}{0}
\begin{exercise}\label{exercise:4.1}
随机点 $(X,Y)$ 均匀分布在顶点为 $(1,1),(1,-1),(-1,1),(-1,-1)$ 的正方形上，即在此正方形上联合 pdf 为 $f(x,y)=1/4$. 求下列事件的概率.
\begin{enumerate}
\item[(a)] $X^2+Y^2<1$.
\item[(b)] $2X-Y>0$.
\item[(c)] $|X+Y|<2$.
\end{enumerate}
\end{exercise}
\begin{exercise}\label{exercise:4.2}\index{SI0424@性质}
证明二元期望的下列性质（\autoref{thm:2.2.5} 的二元对应结果）. 对随机变量 $X,Y$，函数 $g_1(x,y),g_2(x,y)$ 及常数 $a,b,c$：
\begin{enumerate}
\item[(a)] $E(ag_1(X,Y)+bg_2(X,Y)+c)=aE(g_1(X,Y))+bE(g_2(X,Y))+c$.
\item[(b)] 若 $g_1(x,y)\ge0$，则 $E(g_1(X,Y))\ge0$.
\item[(c)] 若 $g_1(x,y)\ge g_2(x,y)$，则 $E(g_1(X,Y))\ge E(g_2(X,Y))$.
\item[(d)] 若 $a\le g_1(x,y)\le b$，则 $a\le E(g_1(X,Y))\le b$.
\end{enumerate}
\end{exercise}
\begin{exercise}\label{exercise:4.3}
利用\autoref{def:4.1.1}，说明\autoref{ex:4.1.5} 末尾定义的随机向量 $(X,Y)$ 的 pmf 正是该例给出的 pmf.
\end{exercise}
\sourcepageanchor{189}

\begin{exercise}\label{exercise:4.4}
一个 pdf 定义为
$$f(x,y)=\begin{cases}C(x+2y),&0<y<1\text{ 且 }0<x<2,\\0,&\text{其他情形}.\end{cases}$$
\begin{enumerate}
\item[(a)] 求 $C$.
\item[(b)] 求 $X$ 的边缘分布.
\item[(c)] 求 $X,Y$ 的联合 cdf.
\item[(d)] 求随机变量 $Z=9/(X+1)^2$ 的 pdf.
\end{enumerate}
\end{exercise}
\begin{exercise}\label{exercise:4.5}\leavevmode
\begin{enumerate}
\item[(a)] 若 $X,Y$ 的联合 pdf 为 $f(x,y)=x+y$，$0\le x\le1$，$0\le y\le1$，求 $P(X>\sqrt Y)$.
\item[(b)] 若 $X,Y$ 的联合 pdf 为 $f(x,y)=2x$，$0\le x\le1$，$0\le y\le1$，求 $P(X^2<Y<X)$.
\end{enumerate}
\end{exercise}
\begin{exercise}\label{exercise:4.6}
甲乙约定下午 1 点至 2 点之间在某处见面. 假设两人在这一小时内独立、随机地到达会面地点. 求甲等待乙的时间长度的分布. 若乙先于甲到达，定义甲的等待时间为 0.
\end{exercise}
\begin{exercise}\label{exercise:4.7}
一位女士在上午 8 点至 8 点半之间出发上班，路程耗时 40 至 50 分钟. 用随机变量 $X$ 表示出发时刻，$Y$ 表示路程时间. 假设二者独立且均匀分布，求她上午 9 点前到达工作地点的概率.
\end{exercise}
\begin{exercise}\label{exercise:4.8}
参照\hyperref[misc:4.9.1]{杂录 4.9.1}：
\begin{enumerate}
\item[(a)] 证明 $P(X=m\mid M=m)=P(X=2m\mid M=m)=1/2$，并验证 $P(M=x\mid X=x)$ 和 $P(M=x/2\mid X=x)$ 的表达式.
\item[(b)] 验证仅当 $\pi(x/2)<2\pi(x)$ 时应当交换；若 $\pi$ 是 $\operatorname{exponential}(\lambda)$ 密度，证明 $x<2\log2/\lambda$ 时交换最优.\footnote{译注：这里的阈值把 $\lambda$ 当作指数分布的速率参数，而本书通常以 $\operatorname{exponential}(\beta)$ 的参数表示尺度；若仅对原文条件作尺度换写，阈值为 $2\lambda\log2$；不过，连续先验还遗漏了换元因子，修正后的速率参数阈值为 $4\log2/\lambda$. 详见\hyperref[misc:4.9.1]{杂录 4.9.1} 的连续先验译注.}
\item[(c)] 对经典方法，证明 $P(Y=2x\mid X=m)=1$ 和 $P(Y=x/2\mid X=2m)=1$；无论交换还是保留自己的信封，期望所得为 $E(Y)=3m/2$.\footnote{译注：原文条件式混用了随机值记号 $x$ 与固定金额 $m$；应分别理解为在 $x=m$ 与 $x=2m$ 时的另一信封金额. 保留信封的期望应记为 $E(X)=3m/2$，交换的为 $E(Y)=3m/2$.}
\end{enumerate}
\end{exercise}
\begin{exercise}\label{exercise:4.9}\index{SI0238@独立性}
证明：若 $X,Y$ 的联合 cdf 满足 $F_{X,Y}(x,y)=F_X(x)F_Y(y)$，则对任意一对区间 $(a,b)$、$(c,d)$，
$$P(a\le X\le b,c\le Y\le d)=P(a\le X\le b)P(c\le Y\le d).$$
\end{exercise}
\begin{exercise}\label{exercise:4.10}
随机对 $(X,Y)$ 的分布为
$$\begin{array}{c|ccc}
 &\multicolumn{3}{c}{X}\\
Y&1&2&3\\\hline
2&1/12&1/6&1/12\\
3&1/6&0&1/6\\
4&0&1/3&0
\end{array}$$
\begin{enumerate}
\item[(a)] 证明 $X,Y$ 不独立.
\item[(b)] 给出随机变量 $U,V$ 的概率表，使其边缘分布与 $X,Y$ 相同，但 $U,V$ 独立.
\end{enumerate}
\end{exercise}
\begin{exercise}\label{exercise:4.11}
反复投掷一枚公平硬币，令 $U=$ 首次出现正面所需的投掷次数，$V=$ 出现两个正面所需的投掷次数. $U,V$ 是否独立？
\end{exercise}
\begin{exercise}\label{exercise:4.12}
将一根棍子随机折成三段，这三段能拼成三角形的概率是多少？关于此问题的完整讨论，见 \hyperref[bib:source132]{Gardner (1961)}.
\end{exercise}
\sourcepageanchor{190}

\begin{exercise}\label{exercise:4.13}\index{SI0426@最小化}\index{SI0884@条件期望}
设 $X,Y$ 为具有有限均值的随机变量.
\begin{enumerate}
\item[(a)]\phantomsection\label{itm:exercise:4.13:a} 证明
$$\min_{g(x)}E(Y-g(X))^2=E(Y-E(Y\mid X))^2,$$
其中 $g(x)$ 遍历所有函数. $E(Y\mid X)$ 有时称为 $Y$ 关于 $X$ 的回归（regression），即给定 $X$ 时 $Y$ 的“最佳”预测量.
\item[(b)] 说明\eqref{eq:2.2.4} 可作为 \hyperref[itm:exercise:4.13:a]{(a)} 的特例推出.
\end{enumerate}
\end{exercise}
\begin{exercise}\label{exercise:4.14}
设 $X,Y$ 是独立的 $n(0,1)$ 随机变量.
\begin{enumerate}
\item[(a)] 求 $P(X^2+Y^2<1)$.
\item[(b)] 先验证 $X^2\sim\chi_1^2$，再求 $P(X^2<1)$.
\end{enumerate}
\end{exercise}
\begin{exercise}\label{exercise:4.15}\index{SI0271@与泊松的关系}\index{SI0376@与二项的关系}
设 $X\sim\operatorname{Poisson}(\theta)$，$Y\sim\operatorname{Poisson}(\lambda)$，且二者独立. \autoref{thm:4.3.2} 已证明 $X+Y\sim\operatorname{Poisson}(\theta+\lambda)$. 证明 $X\mid X+Y$ 的分布是成功概率为 $\theta/(\theta+\lambda)$ 的二项分布. $Y\mid X+Y$ 的分布是什么？
\end{exercise}
\begin{exercise}\label{exercise:4.16}
设 $X,Y$ 独立且服从相同的几何分布.
\begin{enumerate}
\item[(a)] 证明 $U,V$ 独立，其中 $U=\min(X,Y)$，$V=X-Y$.
\item[(b)] 求 $Z=X/(X+Y)$ 的分布；当 $X+Y=0$ 时，定义 $Z=0$.
\item[(c)] 求 $X$ 与 $X+Y$ 的联合 pdf.\footnote{译注：这里的变量离散，原文 pdf 应为 pmf.}
\end{enumerate}
\end{exercise}
\begin{exercise}\label{exercise:4.17}
设 $X\sim\operatorname{exponential}(1)$，定义 $Y$ 为 $X+1$ 的整数部分，即
$$Y=i+1\quad\text{当且仅当}\quad i\le X<i+1,\qquad i=0,1,2,\ldots.$$
\begin{enumerate}
\item[(a)] 求 $Y$ 的分布. 它是哪一种熟知的分布？
\item[(b)] 求给定 $Y\ge5$ 时 $X-4$ 的条件分布.
\end{enumerate}
\end{exercise}
\begin{exercise}\label{exercise:4.18}
已知 $g(x)\ge0$ 且 $\int_0^\infty g(x)\,dx=1$，证明
$$f(x,y)=\frac{2g(\sqrt{x^2+y^2})}{\pi\sqrt{x^2+y^2}},\qquad x,y>0,$$
是一个 pdf.
\end{exercise}
\begin{exercise}\label{exercise:4.19}\leavevmode
\begin{enumerate}
\item[(a)] 设 $X_1,X_2$ 为独立的 $n(0,1)$ 随机变量. 求 $(X_1-X_2)^2/2$ 的 pdf.
\item[(b)] 若 $X_i$，$i=1,2$，为独立的 $\operatorname{gamma}(\alpha_i,1)$ 随机变量，求 $X_1/(X_1+X_2)$ 和 $X_2/(X_1+X_2)$ 的边缘分布.
\end{enumerate}
\end{exercise}
\begin{exercise}\label{exercise:4.20}
$X_1,X_2$ 为独立的 $n(0,\sigma^2)$ 随机变量.
\begin{enumerate}
\item[(a)] 求 $Y_1,Y_2$ 的联合分布，其中 $Y_1=X_1^2+X_2^2$，$Y_2=X_1/\sqrt{Y_1}$.
\item[(b)] 证明 $Y_1,Y_2$ 独立，并从几何上解释这一结果.
\end{enumerate}
\end{exercise}
\begin{exercise}\label{exercise:4.21}
按以下极坐标方案在平面上随机生成一点. 选择半径 $R$，使 $R^2$ 的分布为自由度 2 的卡方分布. 独立地选择角度 $\theta\sim\operatorname{uniform}(0,2\pi)$. 求 $X=R\cos\theta$ 与 $Y=R\sin\theta$ 的联合分布.
\end{exercise}
\begin{exercise}\label{exercise:4.22}
设 $(X,Y)$ 为联合 pdf 为 $f(x,y)$ 的二元随机向量. 令 $U=aX+b$、$V=cY+d$，其中 $a,b,c,d$ 为固定常数，且 $a>0,c>0$. 证明 $(U,V)$ 的联合 pdf 为
$$f_{U,V}(u,v)=\frac1{ac}f\left(\frac{u-b}a,\frac{v-d}c\right).$$
\end{exercise}
\sourcepageanchor{191}

\begin{exercise}\label{exercise:4.23}
设 $X,Y$ 如\autoref{ex:4.3.3} 所述，分别作 \hyperref[itm:exercise:4.23:a]{(a)}、\hyperref[itm:exercise:4.23:b]{(b)} 中的变换，并将 $V$ 积掉，求 $XY$ 的分布.
\begin{enumerate}
\item[(a)]\phantomsection\label{itm:exercise:4.23:a} $U=XY$，$V=Y$.
\item[(b)]\phantomsection\label{itm:exercise:4.23:b} $U=XY$，$V=X/Y$.
\end{enumerate}
\end{exercise}
\begin{exercise}\label{exercise:4.24}\index{SI0259@与伽马的关系}\index{SI0311@与贝塔的关系}
设 $X,Y$ 独立，$X\sim\operatorname{gamma}(r,1)$，$Y\sim\operatorname{gamma}(s,1)$. 证明 $Z_1=X+Y$ 与 $Z_2=X/(X+Y)$ 独立，并求各自分布. $Z_1$ 为 Gamma 分布，$Z_2$ 为 Beta 分布.
\end{exercise}
\begin{exercise}\label{exercise:4.25}
用\autoref{sec:4.3} 的方法，在\autoref{ex:4.5.8} 和\autoref{ex:4.5.9} 中，由 $(X,Z)$ 的联合分布推导 $(X,Y)$ 的联合分布.
\end{exercise}
\begin{exercise}\label{exercise:4.26}\index{SI0866@删失}
$X,Y$ 独立，$X\sim\operatorname{exponential}(\lambda)$，$Y\sim\operatorname{exponential}(\mu)$. 无法直接观测 $X,Y$，而是观测随机变量 $Z,W$，其中
$$Z=\min\{X,Y\},\qquad W=\begin{cases}1,&Z=X,\\0,&Z=Y.\end{cases}$$
这种情形尤其会出现在医学实验中；变量 $X,Y$ 受到删失（censoring）.
\begin{enumerate}
\item[(a)] 求 $Z,W$ 的联合分布.
\item[(b)] 证明 $Z,W$ 独立. 提示：证明对 $i=0,1$，$P(Z\le z\mid W=i)=P(Z\le z)$.
\end{enumerate}
\end{exercise}
\begin{exercise}\label{exercise:4.27}
设 $X\sim n(\mu,\sigma^2)$，$Y\sim n(\gamma,\sigma^2)$，且二者独立. 定义 $U=X+Y$、$V=X-Y$. 证明 $U,V$ 是独立的正态随机变量，并求各自分布.
\end{exercise}
\begin{exercise}\label{exercise:4.28}\index{SI0279@与正态的关系}\index{SI0365@与柯西的关系}
设 $X,Y$ 为独立的标准正态随机变量.
\begin{enumerate}
\item[(a)] 证明 $X/(X+Y)$ 服从 Cauchy 分布.
\item[(b)]\phantomsection\label{itm:exercise:4.28:b} 求 $X/|Y|$ 的分布.
\item[(c)] \hyperref[itm:exercise:4.28:b]{(b)} 的答案是否出乎意料？你能提出一个一般定理吗？
\end{enumerate}
\end{exercise}
\begin{exercise}\label{exercise:4.29}
\hyperref[bib:source181]{Jones (1999)} 研究了 $X=R\cos\theta$、$Y=R\sin\theta$ 时 $X,Y$ 的函数的分布，其中 $\theta\sim U(0,2\pi)$，$R$ 为正随机变量. 以下是他考虑的许多情形中的两种.\footnote{译注：\hyperref[itm:exercise:4.29:b]{(b)} 一般需要补充 $R$ 与 $\theta$ 独立；只指定均匀角度的边缘分布并不能保证该结论. 例如令 $R=1+\varepsilon\cos2\theta$，$0<\varepsilon<1$，则 $E[R^2\cos^2\theta]=1/2+\varepsilon/2+\varepsilon^2/4$，而 $E[R^2\sin^22\theta]=1/2+\varepsilon^2/8$，二阶矩不同. \hyperref[itm:exercise:4.29:a]{(a)} 中 $R$ 消去，无须此额外条件.}
\begin{enumerate}
\item[(a)]\phantomsection\label{itm:exercise:4.29:a} 证明 $X/Y$ 服从 Cauchy 分布.
\item[(b)]\phantomsection\label{itm:exercise:4.29:b} 证明 $2XY/\sqrt{X^2+Y^2}$ 与 $X$ 同分布. 将此结论特殊化，得到一个关于 $n(0,\sigma^2)$ 随机变量的结论.
\end{enumerate}
\end{exercise}
\begin{exercise}\label{exercise:4.30}
假设给定 $X=x$ 时 $Y$ 的分布为 $n(x,x^2)$，$X$ 的边缘分布为 $\operatorname{uniform}(0,1)$.
\begin{enumerate}
\item[(a)] 求 $EY$、$\operatorname{Var}Y$ 和 $\operatorname{Cov}(X,Y)$.
\item[(b)] 证明 $Y/X$ 与 $X$ 独立.
\end{enumerate}
\end{exercise}
\begin{exercise}\label{exercise:4.31}
假设随机变量 $Y$ 服从试验次数为 $n$、成功概率为 $X$ 的二项分布，其中 $n$ 是给定常数，$X\sim\operatorname{uniform}(0,1)$.
\begin{enumerate}
\item[(a)] 求 $EY$ 与 $\operatorname{Var}Y$.
\item[(b)] 求 $X,Y$ 的联合分布.
\item[(c)] 求 $Y$ 的边缘分布.
\end{enumerate}
\end{exercise}
\begin{exercise}\label{exercise:4.32}\index{SI0346@泊松–伽马}\index{SI0351@负二项}\leavevmode
\begin{enumerate}
\item[(a)]\phantomsection\label{itm:exercise:4.32:a} 对分层模型 $Y\mid\Lambda\sim\operatorname{Poisson}(\Lambda)$，$\Lambda\sim\operatorname{gamma}(\alpha,\beta)$，求 $Y$ 的边缘分布、均值和方差. 证明当 $\alpha$ 为整数时，$Y$ 的边缘分布为负二项分布.
\item[(b)] 证明三层模型
$$Y\mid N\sim\operatorname{binomial}(N,p),\quad N\mid\Lambda\sim\operatorname{Poisson}(\Lambda),\quad\Lambda\sim\operatorname{gamma}(\alpha,\beta)$$
得到 $Y$ 的同一边缘（无条件）分布.\footnote{译注：这里“同一”若指完全相同参数的分布，原文有误. 三层模型给出 $Y\mid\Lambda\sim\operatorname{Poisson}(p\Lambda)$，其边缘分布相当于 \hyperref[itm:exercise:4.32:a]{(a)} 中把尺度 $\beta$ 换为 $p\beta$；两者属于同一负二项分布族，但一般并不同分布.}
\end{enumerate}
\end{exercise}
\sourcepageanchor{192}

\begin{exercise}[负二项分布的另一种推导]\label{exercise:4.33}\index{SI0344@对数级数–泊松}\index{SI0876@随机停止的和}
\hyperref[bib:source297]{Solomon (1983)} 详细描述了如下生物模型. 假设昆虫的数量 $N$ 为随机数，每只昆虫产卵 $X_i$ 个，其中 $X_i$ 独立同分布. 总产卵数为 $H=X_1+\cdots+X_N$. $H$ 的分布是什么？

通常假定 $N\sim\operatorname{Poisson}(\lambda)$. 再假定每个 $X_i$ 服从成功概率为 $p$ 的对数级数分布（见\autoref{exercise:3.14}），得到分层模型
$$\begin{aligned}
H\mid N&=X_1+\cdots+X_N,\qquad P(X_i=t)=\frac{-1}{\log p}\frac{(1-p)^t}t,\\
N&\sim\operatorname{Poisson}(\lambda).
\end{aligned}$$
证明 $H$ 的边缘分布是 $\operatorname{negative\ binomial}(r,p)$，其中 $r=-\lambda/\log p$. 最容易的方法是利用\autoref{thm:4.4.3} 和\autoref{thm:4.6.7} 计算并识别 $H$ 的 mgf. \hyperref[bib:source309]{Stuart and Ord (1987}, Section 5.21) 也提及了这种对数级数分布的推导；他们将 $H$ 称为随机停止和（randomly stopped sum\mbox{）.}
\end{exercise}
\begin{exercise}\label{exercise:4.34}\index{SI0260@贝塔二项}\index{SI0338@二项–贝塔}\leavevmode
\begin{enumerate}
\item[(a)]\phantomsection\label{itm:exercise:4.34:a} 对\autoref{ex:4.4.6} 的层次结构，证明 $X$ 的边缘分布为 Beta--二项分布：
$$P(X=x)=\binom nx\frac{\Gamma(\alpha+\beta)}{\Gamma(\alpha)\Gamma(\beta)}\frac{\Gamma(x+\alpha)\Gamma(n-x+\beta)}{\Gamma(\alpha+\beta+n)}.$$
\item[(b)] \hyperref[itm:exercise:4.34:a]{(a)} 中分层模型的一种变体是
$$X\mid P\sim\operatorname{negative\ binomial}(r,P),\qquad P\sim\operatorname{beta}(\alpha,\beta).$$
求 $X$ 的边缘 pmf、均值与方差. 此分布称为 Beta--Pascal 分布.
\end{enumerate}
\end{exercise}
\begin{exercise}\label{exercise:4.35}\leavevmode
\begin{enumerate}
\item[(a)] 对\autoref{ex:4.4.6} 的层次结构，证明 $X$ 的方差可写为
$$\operatorname{Var}X=nEP(1-EP)+n(n-1)\operatorname{Var}P.$$
第一项反映成功概率为 $EP$ 的二项变异，第二项常称为“超二项”（extra-binomial）变异，表明分层模型的方差大于单纯二项模型的方差.
\item[(b)] 对\autoref{exercise:4.32} 的层次结构，证明 $Y$ 的方差可写为
$$\operatorname{Var}Y=E\Lambda+\operatorname{Var}\Lambda=\mu+\frac1\alpha\mu^2,$$
其中 $\mu=E\Lambda$. 指出层次结构引入的“超 Poisson”（extra-Poisson）变异.
\end{enumerate}
\end{exercise}
\begin{exercise}\label{exercise:4.36}
\autoref{ex:4.4.6} 中 Bernoulli 试验层次结构的一种推广，是允许成功概率随试验而变化，同时保持各次试验独立. 标准模型为
$$\begin{aligned}X_i\mid P_i&\sim\operatorname{Bernoulli}(P_i),\qquad i=1,\ldots,n,\\P_i&\sim\operatorname{beta}(\alpha,\beta).\end{aligned}$$
例如，测量药物对 $n$ 名患者的疗效时，因为患者各不相同，不愿假定成功概率恒定，这个模型可能合适. 它可视为一种经验 Bayes 模型（empirical Bayes model），见\hyperref[sec:7.5.6]{杂录 7.5.6}. 感兴趣的随机变量为总成功次数 $Y=\sum_{i=1}^nX_i$.
\begin{enumerate}
\item[(a)] 证明 $EY=n\alpha/(\alpha+\beta)$.
\item[(b)] 证明 $\operatorname{Var}Y=n\alpha\beta/(\alpha+\beta)^2$，因而 $Y$ 与 $\operatorname{binomial}(n,\alpha/(\alpha+\beta))$ 随机变量具有相同均值和方差. $Y$ 的分布是什么？
\item[(c)] 现在假设模型为
$$\begin{aligned}X_i\mid P_i&\sim\operatorname{binomial}(n_i,P_i),\qquad i=1,\ldots,k,\\P_i&\sim\operatorname{beta}(\alpha,\beta).\end{aligned}$$
证明对 $Y=\sum_{i=1}^kX_i$，有 $EY=\frac\alpha{\alpha+\beta}\sum_{i=1}^kn_i$，$\operatorname{Var}Y=\sum_{i=1}^k\operatorname{Var}X_i$，其中
$$\operatorname{Var}X_i=n_i\frac{\alpha\beta(\alpha+\beta+n_i)}{(\alpha+\beta)^2(\alpha+\beta+1)}.$$
\end{enumerate}
\end{exercise}
\sourcepageanchor{193}

\begin{exercise}\label{exercise:4.37}\index{SI0446@强制二元选择}
D. G. \hyperref[bib:source226]{Morrison (1978)} 描述了\autoref{exercise:4.34} 中层次结构的一个推广，建立了强迫二元选择（forced binary choice）模型. 强迫二元选择是指要求受试者必须在两个选项之间作出选择，例如品尝测试. 受试者可能实际上无法分辨两种选项（你能分清可口可乐与百事可乐吗？），但实验设计要求必须作出选择. 因此，正确辨别与正确猜中混杂在一起. Morrison 为此定义参数
$$\begin{aligned}p&=\text{受试者确实能够分辨的概率},\\c&=\text{受试者判断正确的概率}.\end{aligned}$$
于是
$$c=p+\frac12(1-p)=\frac12(1+p),\qquad\frac12<c<1,$$
其中 $\frac12(1-p)$ 是受试者猜对的概率. 进行实验并观测 $X_1,\ldots,X_n\sim\operatorname{Bernoulli}(c)$，所以
$$P\left(\sum X_i=k\mid c\right)=\binom nk c^k(1-c)^{n-k}.$$
然而，$p$ 很可能因人而异，因此允许 $p$ 按 Beta 分布变化：$P\sim\operatorname{beta}(a,b)$.
\begin{enumerate}
\item[(a)] 证明 $\sum X_i$ 的分布为 Beta--二项分布.\footnote{译注：这里 $C=(1+P)/2$ 是支撑在 $(1/2,1)$ 上的缩放平移 Beta 变量，通常不是\autoref{exercise:4.34} 所用的 $(0,1)$ 上 Beta 混合. 因而题述一般不能直接得到通常的 Beta--二项分布；其 pmf 应由 $\binom nk E[C^k(1-C)^{n-k}]$ 求得. 原文是否将这一推广也称为 Beta--二项分布，需核实 Morrison 原文.}
\item[(b)] 求 $\sum X_i$ 的均值和方差.
\end{enumerate}
\end{exercise}
\begin{exercise}[作为指数混合的 Gamma 分布]\label{exercise:4.38}\index{SI0295@与伽马的关系}\index{SI0312@与指数的关系}\index{SI0343@指数–伽马}
\hyperref[bib:source146]{Gleser (1989)} 证明，在某些情形下 Gamma 分布可以写为指数分布的尺度混合；对同一数据的不同分析启发了这一恒等式. 设 $f(x)$ 为 $\operatorname{gamma}(r,\lambda)$ 的 pdf.
\begin{enumerate}
\item[(a)]\phantomsection\label{itm:exercise:4.38:a} 证明若 $r\le1$，则
$$f(x)=\int_0^\lambda\frac1\nu e^{-x/\nu}p_\lambda(\nu)\,d\nu,$$
其中
$$p_\lambda(\nu)=\frac1{\Gamma(r)\Gamma(1-r)}\frac{\nu^{r-1}}{(\lambda-\nu)^r},\qquad0<\nu<\lambda.$$
提示：作变量替换 $u=x/\nu-x/\lambda$.
\item[(b)]\phantomsection\label{itm:exercise:4.38:b} 通过证明 $\int_0^\lambda p_\lambda(\nu)\,d\nu=1$，说明当 $r\le1$ 时 $p_\lambda(\nu)$ 是 pdf.\footnote{译注：\hyperref[itm:exercise:4.38:a]{(a)}、\hyperref[itm:exercise:4.38:b]{(b)} 所给普通密度公式实际要求 $0<r<1$；$r=1$ 时 $\Gamma(1-r)$ 无定义. 边界 $r=1$ 对应在 $\nu=\lambda$ 处的单位点质量，即一个指数分布，而非此普通密度.}
\item[(c)] 证明限制 $r\le1$ 是 \hyperref[itm:exercise:4.38:a]{(a)} 的表示成立的必要条件，即 $r>1$ 时不存在这种表示. 提示：假设对某个 pdf $q_\lambda(\nu)$，可写成
$$f(x)=\int_0^\infty\frac{e^{-x/\nu}}\nu q_\lambda(\nu)\,d\nu.$$
证明 $\frac\partial{\partial x}\log f(x)>0$，而 $\frac\partial{\partial x}\log\left(\int_0^\infty\frac{e^{-x/\nu}}\nu q_\lambda(\nu)\,d\nu\right)<0$，产生矛盾.\footnote{译注：Gamma 密度的对数导数为 $(r-1)/x-1/\lambda$，只在 $0<x<\lambda(r-1)$ 时为正，而非所有 $x>0$；在该非空区间取得矛盾已足够.}
\end{enumerate}
\end{exercise}
\begin{exercise}\label{exercise:4.39}\index{SI0349@协方差}
设 $(X_1,\ldots,X_n)$ 服从试验次数为 $m$、格概率为 $p_1,\ldots,p_n$ 的多项分布（见\autoref{def:4.6.2}）. 证明对每个 $i,j$，\footnote{译注：这里需 $i\ne j$，条件分布公式还需 $p_j<1$ 及 $P(X_j=x_j)>0$；$i=j$ 时协方差应为方差 $mp_i(1-p_i)$.}
$$\begin{aligned}X_i\mid X_j=x_j&\sim\operatorname{binomial}\left(m-x_j,\frac{p_i}{1-p_j}\right),\\X_j&\sim\operatorname{binomial}(m,p_j),\end{aligned}$$
并且 $\operatorname{Cov}(X_i,X_j)=-mp_ip_j$.
\end{exercise}
\begin{exercise}\label{exercise:4.40}
Beta 分布的一个推广是 Dirichlet 分布. 在二元版本中，$(X,Y)$ 的 pdf 为
$$f(x,y)=Cx^{a-1}y^{b-1}(1-x-y)^{c-1},\quad0<x<1,\quad0<y<1,\quad0<y<1-x<1,$$
其中 $a,b,c$ 为正的常数.
\sourcepageanchor{194}

\begin{enumerate}
\item[(a)] 证明 $C=\Gamma(a+b+c)/(\Gamma(a)\Gamma(b)\Gamma(c))$.
\item[(b)] 证明 $X,Y$ 的边缘分布都是 Beta 分布.
\item[(c)] 求 $Y\mid X=x$ 的条件分布，并证明 $Y/(1-x)$ 服从 $\operatorname{beta}(b,c)$.
\item[(d)] 证明 $E(XY)=ab/((a+b+c+1)(a+b+c))$，并求协方差.
\end{enumerate}
\end{exercise}
\begin{exercise}\label{exercise:4.41}
证明任意随机变量都与常数不相关.\footnote{译注：这里“不相关”应理解为在期望存在时协方差为 0；常数的标准差为 0，按\autoref{def:4.5.2} 定义的相关系数并不存在.}
\end{exercise}
\begin{exercise}\label{exercise:4.42}
设 $X,Y$ 独立，均值为 $\mu_X,\mu_Y$，方差为 $\sigma_X^2,\sigma_Y^2$. 用这些均值和方差表示 $XY$ 与 $Y$ 的相关系数.
\end{exercise}
\begin{exercise}\label{exercise:4.43}
设 $X_1,X_2,X_3$ 不相关，均值均为 $\mu$，方差均为 $\sigma^2$. 用 $\mu,\sigma^2$ 表示 $\operatorname{Cov}(X_1+X_2,X_2+X_3)$ 和 $\operatorname{Cov}(X_1+X_2,X_1-X_2)$.
\end{exercise}
\begin{exercise}\label{exercise:4.44}\index{SI1068@和}
证明\autoref{thm:4.5.6} 的下列推广：对任意随机向量 $(X_1,\ldots,X_n)$，
$$\operatorname{Var}\left(\sum_{i=1}^nX_i\right)=\sum_{i=1}^n\operatorname{Var}X_i+2\sum_{1\le i<j\le n}\operatorname{Cov}(X_i,X_j).$$
\end{exercise}
\begin{exercise}\label{exercise:4.45}\index{SI0273@二元正态}
证明：若 $(X,Y)$ 服从参数为 $\mu_X,\mu_Y,\sigma_X^2,\sigma_Y^2,\rho$ 的二元正态分布，则：
\begin{enumerate}
\item[(a)] $X$ 的边缘分布为 $n(\mu_X,\sigma_X^2)$，$Y$ 的边缘分布为 $n(\mu_Y,\sigma_Y^2)$.
\item[(b)] 给定 $X=x$ 时 $Y$ 的条件分布为
$$n\bigl(\mu_Y+\rho(\sigma_Y/\sigma_X)(x-\mu_X),\sigma_Y^2(1-\rho^2)\bigr).$$
\item[(c)] 对任意常数 $a,b$，$aX+bY$ 的分布为
$$n(a\mu_X+b\mu_Y,a^2\sigma_X^2+b^2\sigma_Y^2+2ab\rho\sigma_X\sigma_Y).$$
\end{enumerate}
\end{exercise}
\begin{exercise}[二元正态分布的一种推导]\label{exercise:4.46}
设 $Z_1,Z_2$ 为独立的 $n(0,1)$ 随机变量，并定义
$$X=a_XZ_1+b_XZ_2+c_X,\qquad Y=a_YZ_1+b_YZ_2+c_Y,$$
其中 $a_X,b_X,c_X,a_Y,b_Y,c_Y$ 为常数.
\begin{enumerate}
\item[(a)] 证明
$$\begin{gathered}EX=c_X,\quad\operatorname{Var}X=a_X^2+b_X^2,\\EY=c_Y,\quad\operatorname{Var}Y=a_Y^2+b_Y^2,\\\operatorname{Cov}(X,Y)=a_Xa_Y+b_Xb_Y.\end{gathered}$$
\item[(b)]\phantomsection\label{itm:exercise:4.46:b} 如果将这些常数定义为
$$\begin{aligned}
a_X&=\sqrt{\frac{1+\rho}2}\sigma_X,&b_X&=\sqrt{\frac{1-\rho}2}\sigma_X,&c_X&=\mu_X,\\
a_Y&=\sqrt{\frac{1+\rho}2}\sigma_Y,&b_Y&=-\sqrt{\frac{1-\rho}2}\sigma_Y,&c_Y&=\mu_Y,
\end{aligned}$$
其中 $\mu_X,\mu_Y,\sigma_X^2,\sigma_Y^2,\rho$ 为常数，$-1\le\rho\le1$，证明
$$EX=\mu_X,\quad\operatorname{Var}X=\sigma_X^2,\qquad EY=\mu_Y,\quad\operatorname{Var}Y=\sigma_Y^2,\qquad\rho_{XY}=\rho.$$
\item[(c)] 证明 $(X,Y)$ 具有参数为 $\mu_X,\mu_Y,\sigma_X^2,\sigma_Y^2,\rho$ 的二元正态 pdf.\footnote{译注：普通二元 pdf 要求 $|\rho|<1$（且两个方差为正）. \hyperref[itm:exercise:4.46:b]{(b)} 允许的 $\rho=\pm1$ 给出集中在直线上的退化联合分布，不具有\autoref{def:4.5.10} 的二维密度.}
\item[(d)] 从二元正态参数 $\mu_X,\mu_Y,\sigma_X^2,\sigma_Y^2,\rho$ 出发，可将常数定义为以下方程的解：
$$\begin{gathered}\mu_X=c_X,\quad\sigma_X^2=a_X^2+b_X^2,\\\mu_Y=c_Y,\quad\sigma_Y^2=a_Y^2+b_Y^2,\\\rho\sigma_X\sigma_Y=a_Xa_Y+b_Xb_Y.\end{gathered}$$
再给出一组解，说明 \hyperref[itm:exercise:4.46:b]{(b)} 中的解并不唯一. 一共有多少组解？
\end{enumerate}
\end{exercise}
\sourcepageanchor{195}

\begin{exercise}[边缘正态性不蕴含二元正态性]\label{exercise:4.47}\index{SI0359@二元}
设 $X,Y$ 为独立的 $n(0,1)$ 随机变量，定义
$$Z=\begin{cases}X,&XY>0,\\-X,&XY<0.\end{cases}$$
\begin{enumerate}
\item[(a)] 证明 $Z$ 服从正态分布.
\item[(b)] 证明 $Z,Y$ 的联合分布不是二元正态分布. 提示：说明 $Z,Y$ 总是同号.
\end{enumerate}
\end{exercise}
\begin{exercise}\label{exercise:4.48}
\hyperref[bib:source135]{Gelman and Meng (1991)} 给出了一个二元分布族的例子，它不是二元正态分布，但条件分布均为正态. 定义 $(X,Y)$ 的联合 pdf 为
$$f(x,y)\propto\exp\left\{-\frac12[Ax^2y^2+x^2+y^2-2Bxy-2Cx-2Dy]\right\},$$
其中 $A,B,C,D$ 为常数.
\begin{enumerate}
\item[(a)] 证明 $X\mid Y=y$ 服从均值为 $(By+C)/(Ay^2+1)$、方差为 $1/(Ay^2+1)$ 的正态分布. 推导 $Y\mid X=x$ 的相应结论.
\item[(b)] 一个十分有趣的配置是 $A=1,B=0,C=D=8$. 证明这个联合分布是双峰的.
\end{enumerate}
\end{exercise}
\begin{exercise}\label{exercise:4.49}
\hyperref[bib:source15]{Behboodian (1990)} 说明了如何构造不相关但不独立的二元随机变量. 设 $f_1,f_2,g_1,g_2$ 为一元密度，均值分别为 $\mu_1,\mu_2,\xi_1,\xi_2$，二元随机变量 $(X,Y)$ 的密度为
$$(X,Y)\sim af_1(x)g_1(y)+(1-a)f_2(x)g_2(y),$$
其中 $0<a<1$ 已知.
\begin{enumerate}
\item[(a)] 证明边缘分布为 $f_X(x)=af_1(x)+(1-a)f_2(x)$，$f_Y(y)=ag_1(y)+(1-a)g_2(y)$.
\item[(b)] 证明 $X,Y$ 独立，当且仅当 $[f_1(x)-f_2(x)][g_1(y)-g_2(y)]=0$.
\item[(c)] 证明 $\operatorname{Cov}(X,Y)=a(1-a)[\mu_1-\mu_2][\xi_1-\xi_2]$，并据此解释如何构造不独立但不相关的随机变量.
\item[(d)] 令 $f_1,f_2,g_1,g_2$ 为二项 pmf，给出参数组合的例子，分别产生独立的 $(X,Y)$、相关的 $(X,Y)$，以及不相关但不独立的 $(X,Y)$.
\end{enumerate}
\end{exercise}
\begin{exercise}\label{exercise:4.50}
若 $(X,Y)$ 具有二元正态 pdf
$$f(x,y)=\frac1{2\pi(1-\rho^2)^{1/2}}\exp\left(\frac{-1}{2(1-\rho^2)}(x^2-2\rho xy+y^2)\right),$$
证明 $\operatorname{Corr}(X,Y)=\rho$，$\operatorname{Corr}(X^2,Y^2)=\rho^2$. 条件期望可以简化计算.
\end{exercise}
\begin{exercise}\label{exercise:4.51}
设 $X,Y,Z$ 为独立的 $\operatorname{uniform}(0,1)$ 随机变量.
\begin{enumerate}
\item[(a)] 求 $P(X/Y\le t)$ 和 $P(XY\le t)$. 画图会有帮助.
\item[(b)] 求 $P(XY/Z\le t)$.
\end{enumerate}
\end{exercise}
\begin{exercise}\label{exercise:4.52}
向 $(x,y)$ 坐标系的原点射击，命中点 $(X,Y)$ 是随机变量. 假定 $X,Y$ 是独立的 $n(0,1)$ 随机变量. 若独立发射两颗子弹，两命中点间距离的分布是什么？
\end{exercise}
\begin{exercise}\label{exercise:4.53}
设 $A,B,C$ 为独立的 $\operatorname{uniform}(0,1)$ 随机变量. $Ax^2+Bx+C$ 有实根的概率是多少？提示：若 $X\sim\operatorname{uniform}(0,1)$，则 $-\log X$ 为指数随机变量；两个独立指数随机变量之和为 Gamma 随机变量.
\end{exercise}
\begin{exercise}\label{exercise:4.54}
求 $\prod_{i=1}^nX_i$ 的 pdf，其中 $X_i$ 为独立的 $\operatorname{uniform}(0,1)$ 随机变量. 提示：尝试计算 cdf，并记住均匀分布与指数分布的关系.
\end{exercise}
\sourcepageanchor{196}

\begin{exercise}\label{exercise:4.55}\index{SI0749@并联系统}
并联系统（parallel system）只要至少有一个部件工作，就能工作. 某并联系统由三个独立部件组成，每个部件的寿命都服从 $\operatorname{exponential}(\lambda)$ 分布. 系统寿命是各部件寿命的最大值. 求系统寿命的分布.
\end{exercise}
\begin{exercise}\label{exercise:4.56}
大量人员 $N=mk$ 需要验血，可采用两种方案.
\begin{enumerate}
\item[(i)] 每人单独检测，此时需要 $N$ 次检测.
\item[(ii)]\phantomsection\label{itm:exercise:4.56:ii} 将 $k$ 人的血样混合检测. 如果结果为阴性，这一次检测就足够覆盖这 $k$ 人；若结果为阳性，则对这 $k$ 人逐个检测，总共需要 $k+1$ 次.
\end{enumerate}
假设每个人检测阳性的概率均为 $p$，不同人的检测结果统计独立.
\begin{enumerate}
\item[(a)] $k$ 人混合样本检测阳性的概率是多少？
\item[(b)]\phantomsection\label{itm:exercise:4.56:b} 令 $X=$ 方案 \hyperref[itm:exercise:4.56:ii]{(ii)} 所需的检测次数，求 $EX$.
\item[(c)] 若已知 $p$ 接近 0，从最小化这 $N$ 人所需检测次数的期望来看，应选择哪种方案？利用 \hyperref[itm:exercise:4.56:b]{(b)} 的表达式说明理由.
\end{enumerate}
\end{exercise}
\begin{exercise}\label{exercise:4.57}
参照\hyperref[misc:4.9.2]{杂录 4.9.2}.
\begin{enumerate}
\item[(a)] 证明 $A_1$ 是算术平均，$A_{-1}$ 是调和平均，$A_0=\lim_{r\to0}A_r$ 是几何平均.
\item[(b)] 如果能证明 $A_r$ 在 $-\infty<r<\infty$ 上关于 $r$ 非减，就能得到算术--几何--调和平均不等式.
\begin{enumerate}
\item[(i)] 验证：若 $\log A_r$ 关于 $r$ 非减，则 $A_r$ 关于 $r$ 非减.
\item[(ii)] 证明
$$\frac d{dr}\log A_r=\frac1{r^2}\left\{\frac{r\sum_i x_i^r\log x_i}{\sum_i x_i^r}-\log\left(\frac1n\sum_i x_i^r\right)\right\}.$$
\item[(iii)]\phantomsection\label{itm:exercise:4.57:b:iii} 定义 $a_i=x_i^r/\sum_i x_i^r$，将花括号中的量写为 $\log n-\sum a_i\log(1/a_i)$，其中 $\sum a_i=1$. 证明此量非负，从而建立 $A_r$ 的单调性，并将算术--几何--调和平均不等式作为特例得到.
\end{enumerate}
\end{enumerate}
量 $\sum_i a_i\log(1/a_i)$ 称为熵（entropy）\index{SI0398@熵}\semanticindexanchor{398:196}，有时视为不确定性的绝对度量，见 \hyperref[bib:source29]{Bernardo and Smith (1994}, Section 2.7). \hyperref[itm:exercise:4.57:b:iii]{(iii)} 的结论表明，当所有概率相同时（随机性），熵达到最大.

提示：为证明不等式，注意 $a_i$ 构成概率分布，可以写为
$$E\log\left(\frac1a\right)=\sum_i a_i\log(1/a_i),$$
再由 Jensen 不等式，$E\log(1/a)\le\log(E(1/a))$.
\end{exercise}
\begin{exercise}\label{exercise:4.58}
对任意具有有限方差的两个随机变量 $X,Y$，证明：
\begin{enumerate}
\item[(a)] $\operatorname{Cov}(X,Y)=\operatorname{Cov}(X,E(Y\mid X))$.
\item[(b)] $X$ 与 $Y-E(Y\mid X)$ 不相关.
\item[(c)] $\operatorname{Var}(Y-E(Y\mid X))=E(\operatorname{Var}(Y\mid X))$.
\end{enumerate}
\end{exercise}
\begin{exercise}\label{exercise:4.59}\index{SI0533@条件协方差}
对任意具有有限方差的三个随机变量 $X,Y,Z$，按\autoref{thm:4.4.7} 的思路证明协方差恒等式
$$\operatorname{Cov}(X,Y)=E(\operatorname{Cov}(X,Y\mid Z))+\operatorname{Cov}(E(X\mid Z),E(Y\mid Z)),$$
其中 $\operatorname{Cov}(X,Y\mid Z)$ 是在 pdf $f(x,y\mid z)$ 下 $X,Y$ 的协方差.
\end{exercise}
\begin{exercise}\label{exercise:4.60}
参照\hyperref[misc:4.9.3]{杂录 4.9.3}，分别对那里给出的条件 $Y=X$ 的三种解释，求给定 $Y=X$ 时 $Y$ 的条件分布.
\end{exercise}
\sourcepageanchor{197}

\begin{exercise}\label{exercise:4.61}
\hyperref[bib:source96]{DeGroot (1986)} 给出如下 Borel 悖论（\hyperref[misc:4.9.3]{杂录 4.9.3}）的例子. 设 $X_1,X_2$ 为独立同分布的 $\operatorname{exponential}(1)$ 随机变量，定义 $Z=(X_2-1)/X_1$. 零概率集合 $\{Z=0\}$ 与 $\{X_2=1\}$ 看似提供了相同信息，却导出不同条件分布.
\begin{enumerate}
\item[(a)]\phantomsection\label{itm:exercise:4.61:a} 求 $X_1\mid Z=0$ 的分布，并与 $X_1\mid X_2=1$ 比较.
\item[(b)] 对小的 $\varepsilon>0$ 及 $x_1>0,x_2>0$，考虑
$$B_1=\left\{(x_1,x_2):-\varepsilon<\frac{x_2-1}{x_1}<\varepsilon\right\},\qquad
B_2=\{(x_1,x_2):1-\varepsilon<x_2<1+\varepsilon\}.$$
画出这两个集合，并说明为何 $B_1$ 提供了关于 $X_1$ 的信息，而 $B_2$ 没有.
\item[(c)] 计算 $P(X_1\le x\mid B_1)$ 与 $P(X_1\le x\mid B_2)$，说明它们在 $\varepsilon\to0$ 时的极限与 \hyperref[itm:exercise:4.61:a]{(a)} 一致.
\end{enumerate}
（由俄亥俄州立大学 L. Mark Berliner 提供.）
\end{exercise}
\begin{exercise}\label{exercise:4.62}
完成 Jensen 不等式（\autoref{thm:4.7.7}）中等号情形的证明. 设 $g$ 为凸函数，$a+bx$ 是它在 $x=EX$ 处的切线，且除 $x=EX$ 外都有 $g(x)>a+bx$. 则除非 $P(X=EX)=1$，否则 $Eg(X)>g(EX)$.
\end{exercise}
\begin{exercise}\label{exercise:4.63}
随机变量 $X$ 由 $Z=\log X$ 定义，其中 $EZ=0$. $EX$ 大于、小于还是等于 1？
\end{exercise}
\begin{exercise}\label{exercise:4.64}\index{SI0578@三角}
本题讨论著名的三角不等式（triangle inequality），它是 Minkowski 不等式的一个特例.
\begin{enumerate}
\item[(a)]\phantomsection\label{itm:exercise:4.64:a} 不使用 Minkowski 不等式，证明对任意数 $a,b$，$|a+b|\le|a|+|b|$.
\item[(b)] 利用 \hyperref[itm:exercise:4.64:a]{(a)} 证明，对任意具有有限期望的随机变量 $X,Y$，$E|X+Y|\le E|X|+E|Y|$.
\end{enumerate}
\end{exercise}
\begin{exercise}\label{exercise:4.65}
推广正文中紧接协方差不等式之前的论证，证明\autoref{thm:4.7.9}.
\end{exercise}



\subsection{杂录}\label{sec:4.9}
\subsubsection{交换悖论}\label{sec:4.9.1}\index{SI0420@交换悖论}\index{SI0748@交换}\label{misc:4.9.1}
“交换悖论”（exchange paradox；\hyperref[bib:source76]{Christensen and Utts 1992）}在统计学家之间引发了长篇讨论. 问题（或悖论）如下：
\begin{quote}
一位大师在一个信封中放入 $m$ 美元，另一个信封中放入 $2m$ 美元. 你和对手随机各取一个信封. 你打开自己的信封，发现其中有 $x$ 美元，大师随后问你是否愿意交换信封. 你想，如果交换，就会得到 $x/2$ 或 $2x$ 美元，各有 $1/2$ 的概率. 于是交换所得的期望为 $(1/2)(x/2)+(1/2)(2x)=5x/4$，大于手中的 $x$ 美元. 所以你提出交换.
\end{quote}
悖论在于，对手也作了同样的计算. 怎么可能交换对你们双方都有利？
\begin{enumerate}
\item[(i)] Christensen and Utts 说：“交换信封总是最优的这一结论，建立在观测信封内容并未获得任何信息的假设之上.” 他们提出如下解释.

设 $M\sim\pi(m)$，$\pi$ 为放进第一个信封的钱数的 pdf，$X$ 为你信封中的钱数. 则 $P(X=m\mid M=m)=P(X=2m\mid M=m)=1/2$，从而\footnote{译注：若 $\pi$ 为离散 pmf，下面的权重形式可成立；若像原文所说是连续 pdf，则 $X=2M$ 的密度须含换元因子 $1/2$. 此时两个后验概率应分别为 $2\pi(x)/(2\pi(x)+\pi(x/2))$ 与 $\pi(x/2)/(2\pi(x)+\pi(x/2))$. 相应交换条件应为 $\pi(x/2)<4\pi(x)$；若 $\pi(m)=\lambda e^{-\lambda m}$，阈值为 $x<4\log2/\lambda$；若以尺度 $\lambda$ 参数化，则为 $x<4\lambda\log2$. 原式及后面的原结论照录.}
$$P(M=x\mid X=x)=\frac{\pi(x)}{\pi(x)+\pi(x/2)},\qquad
P(M=x/2\mid X=x)=\frac{\pi(x/2)}{\pi(x)+\pi(x/2)}.$$
\sourcepageanchor{198}

于是，交换所得的期望为
$$\frac{\pi(x)}{\pi(x)+\pi(x/2)}2x+\frac{\pi(x/2)}{\pi(x)+\pi(x/2)}\frac x2,$$
所以仅当 $\pi(x/2)<2\pi(x)$ 时应交换. 若 $\pi$ 是 $\operatorname{exponential}(\lambda)$ 密度，$x<2\log2/\lambda$ 时交换最优.
\item[(ii)] 更经典的方法不假设第一个信封中的金额服从某个 pdf. Christensen and Utts 也在这里给出解释：若错误地假设对所有 $X,Y$ 的取值都有 $P(Y=y\mid X=x)=1/2$，就会出现悖论，其中 $X$ 是你信封中的金额，$Y$ 是对手信封中的金额. 他们指出，正确的条件分布是 $P(Y=2x\mid X=m)=1$ 和 $P(Y=x/2\mid X=2m)=1$；交换的期望所得为 $E(Y)=3m/2$，与保留信封的期望所得相同.\footnote{译注：这里应将第一式的 $x$ 理解为 $m$、第二式的 $x$ 理解为 $2m$；不混用记号时，两式为 $P(Y=2m\mid X=m)=1$、$P(Y=m\mid X=2m)=1$.}
\end{enumerate}
这个悖论常伴随着支持或反对 Bayes 推断方法的论辩（见\autoref{sec:7}），但这些争论与底层概率计算并不直接相关. 关于评论、批评及其他分析，见 \hyperref[bib:source34]{Binder (1993)}、\hyperref[bib:source252]{Ridgeway (1993)}（其中包含 Marilyn vos Savant 的一个解答）、\hyperref[bib:source260]{Ross (1994)}、\hyperref[bib:source36]{Blachman (1996)} 的读者来信，以及 Christen 和 Utts 随附的回应.

\subsubsection{再论算术--几何--调和平均不等式}\label{sec:4.9.2}\index{SI0558@算术–几何–调和平均数}\index{SI0577@幂平均数}\index{SI0821@幂平均数}\label{misc:4.9.2}
算术--几何--调和平均不等式是关于幂平均（power means）的一般结果的特例. 对 $x_i\ge0$，幂平均定义为
$$A_r=\left[\frac1n\sum_{i=1}^nx_i^r\right]^{1/r}.$$
\hyperref[bib:source289]{Shier (1988)} 证明 $A_r$ 关于 $r$ 非减，即若 $r\le r'$，则 $A_r\le A_{r'}$，也就是
$$\left[\frac1n\sum_{i=1}^nx_i^r\right]^{1/r}\le\left[\frac1n\sum_{i=1}^nx_i^{r'}\right]^{1/r'},\qquad r\le r'.$$
显然 $A_1$ 是算术平均，$A_{-1}$ 是调和平均. 不那么显然、但仍然正确的是，$A_0=\lim_{r\to0}A_r$ 是几何平均. 因而，算术--几何--调和平均不等式是幂平均不等式的特例，见\autoref{exercise:4.57}.\footnote{译注：当允许 $x_i=0$ 时，负幂以及 $r=0$ 的原式不能直接按普通实数运算使用，需以极限作延拓. 若统一假定 $x_i>0$，则除 $r=0$ 按所述极限定义外，上述公式可直接使用.}

\subsubsection{Borel 悖论}\label{sec:4.9.3}\index{SI0098@博雷尔悖论}\index{SI0747@博雷尔}\label{misc:4.9.3}
本章一直在对连续随机变量 $X,Y$ 使用 $E(Y\mid X=x)$、$P(Y\le y\mid X=x)$ 这样的表达式. 到目前为止没有遇到麻烦，但原本可能会遇到.

形式上，条件期望的条件是一个子 $\sigma$ 代数（\autoref{def:1.2.1}）. 条件期望 $E(Y\mid\mathscr G)$ 被定义为这样一个随机变量：它在子 $\sigma$ 代数 $\mathscr G$ 中任意集合上的积分，与 $X$ 的积分相同.\footnote{译注：原文最后的 $X$ 应为 $Y$. 完整定义还要求 $E(Y\mid\mathscr G)$ 对 $\mathscr G$ 可测；在 $Y$ 可积时，其在每个 $G\in\mathscr G$ 上的积分等于 $Y$ 在 $G$ 上的积分.} 这是概率论中相当深入的概念，见 \hyperref[bib:source33]{Billingsley (1995}, Section 34).

由于条件期望仅通过积分来定义，即使条件明确，它也可能不唯一. 但当我们以零概率集合（如 $\{X=x\}$）为条件时，条件本身可能未明确指定，因此更可能出现不同的条件期望. 要看这会怎样影响我们，最容易的是考察条件分布，这相当于计算 $E[I(Y\le y)\mid X=x]$.
\sourcepageanchor{199}

\hyperref[bib:source247]{Proschan and Presnell (1998)} 讲述了一个统计学考试的故事，题目是：“若 $X,Y$ 是独立的标准正态随机变量，给定 $Y=X$ 时 $Y$ 的条件分布是什么？”不同学生将 $Y=X$ 解释为：
\begin{enumerate}
\item[(1)] $Z_1=0$，其中 $Z_1=Y-X$；
\item[(2)] $Z_2=1$，其中 $Z_2=Y/X$；
\item[(3)]\phantomsection\label{itm:borel:3} $Z_3=1$，其中 $Z_3=I(Y=X)$.
\end{enumerate}
每一种都是对条件 $Y=X$ 的正确解释，但每一种都导出不同的条件分布（\autoref{exercise:4.60}\mbox{）.}\footnote{译注：固定条件 $\sigma$ 代数时，条件期望在几乎处处意义下唯一；只在零测集上可改动. 特别是 \hyperref[itm:borel:3]{(3)} 中 $P(Z_3=1)=0$，通常的离散条件概率比值无法定义该条件分布，其在这一零概率取值处的版本可任意指定，不能仅凭事件本身得到唯一的第三个分布.}

这就是 Borel 悖论（Borel paradox）：对同一个零概率条件集合的不同而正确的解释，会产生不同条件期望. 如何避免？一种方法是避免直接以零概率集合为条件，即只计算 $E(Y\mid X\in B)$，其中 $P(X\in B)>0$. 为计算 $E(Y\mid X=x)$ 这样的量，取一列 $B_n\downarrow x$，并定义 $E(Y\mid X=x)=\lim_{n\to\infty}E(Y\mid X\in B_n)$. 这样就避免了悖论，因为不同的 $E(Y\mid X=x)$ 答案来自不同的集合列，就不应感到意外了（\autoref{exercise:4.61}）.\footnote{译注：此极限法还需明确集合列、每项条件概率为正，以及极限存在；任意递减集合列并不自动保证存在或给出同一极限. 因而该方法是在说明具体的逼近机制，而非为一切零概率条件给出唯一、无条件适用的定义.}





\section{随机样本的性质}\label{sec:5}
\chapterdoi{5}

\sourcepageanchor{200}
\begin{flushright}
“我恐怕一解释，就把自己的底细泄露得太多了，”他说.“只有结果而没有原因，要令人印象深刻得多.”\\
夏洛克·福尔摩斯\\《证券经纪人的书记员》
\end{flushright}




% 原书印刷169–172页，PDF物理200–203页.
\subsection{随机样本的基本概念}\label{sec:5.1}
实验收集的数据往往包含对某个感兴趣变量的若干次观测.我们在\autoref{sec:4}开头讨论过这种例子.本章介绍一种常用于描述这种情形的数据收集模型，称为随机抽样（random sampling）.下面的定义从数学上说明这种数据收集方法的含义.
\begin{definition}\label{def:5.1.1}\index{SI0545@独立同分布}\index{SI0554@随机变量}\index{SI0917@样本}
若随机变量$X_1,\ldots,X_n$相互独立，且每个$X_i$的边际pdf或pmf都是同一个函数$f(x)$，则称$X_1,\ldots,X_n$为来自总体$f(x)$的容量为$n$的随机样本（random sample）.也称$X_1,\ldots,X_n$是具有pdf或pmf $f(x)$的独立同分布（independent and identically distributed）随机变量，通常简称为iid随机变量.
\end{definition}
随机抽样模型描述这样一种实验情形：感兴趣变量的概率分布由$f(x)$描述.若仅对此变量观测一次，得到$X$，则可用$f(x)$计算关于$X$的概率.大多数实验会对此变量重复观测$n>1$次，其中$n$是固定的正整数；第一次观测为$X_1$，第二次为$X_2$，依此类推.在随机抽样模型下，每个$X_i$都是对同一个变量的观测，并且每个$X_i$的边际分布都由$f(x)$给出.此外，观测的取得方式保证某次观测的值对其他任何观测都没有影响，也没有关系；即$X_1,\ldots,X_n$相互独立.（关于独立性的一种推广，见\autoref{exercise:5.4}.）

由\autoref{def:4.6.5}，$X_1,\ldots,X_n$的联合pdf或pmf为
\begin{equation}\label{eq:5.1.1}\index{SI0839@样本}\index{SI0857@样本}\index{SI0932@概率密度函数与概率质量函数}
f(x_1,\ldots,x_n)=f(x_1)f(x_2)\cdots f(x_n)=\prod_{i=1}^n f(x_i).
\end{equation}
可用这一联合pdf或pmf计算涉及样本的概率.由于$X_1,\ldots,X_n$同分布，所有边际密度$f(x)$都是同一个函数.特别地，若总体pdf或pmf属于某个参数族，例如\autoref{sec:3}介绍的某个族，且pdf或pmf为$f(x\mid\theta)$，则联合pdf或pmf为
\begin{equation}\label{eq:5.1.2}
f(x_1,\ldots,x_n\mid\theta)=\prod_{i=1}^n f(x_i\mid\theta),
\end{equation}
\sourcepageanchor{201}
其中乘积的每一项都使用同一个参数值$\theta$.在统计问题中，若假定所观测总体属于某个指定的参数族，但真实参数值未知，则来自该总体的随机样本具有上述形式的联合pdf或pmf，其中$\theta$未知.考察$\theta$的不同可能值，就能研究来自不同总体的随机样本会有怎样的表现.
\begin{example}[样本pdf——指数分布]\label{ex:5.1.2}
设$X_1,\ldots,X_n$是来自$\operatorname{exponential}(\beta)$总体的随机样本.具体来说，$X_1,\ldots,X_n$可以是$n$块相同电路板从投入试验、使用直至失效的时间（以年计）.样本的联合pdf为
$$f(x_1,\ldots,x_n\mid\beta)=\prod_{i=1}^n f(x_i\mid\beta)=\prod_{i=1}^n\frac1\beta e^{-x_i/\beta}=\frac1{\beta^n}e^{-(x_1+\cdots+x_n)/\beta}.$$
这个pdf可用于回答关于样本的问题.例如，所有电路板寿命都超过两年的概率是多少？计算得到
\begin{align*}
&P(X_1>2,\ldots,X_n>2)\\
&\quad=\int_2^\infty\cdots\int_2^\infty\prod_{i=1}^n\frac1\beta e^{-x_i/\beta}\,dx_1\cdots dx_n\\
&\quad=e^{-2/\beta}\int_2^\infty\cdots\int_2^\infty\prod_{i=2}^n\frac1\beta e^{-x_i/\beta}\,dx_2\cdots dx_n&&\text{（对$x_1$积分）}\\
&\qquad\vdots&&\text{（依次对其余$x_i$积分）}\\
&\quad=(e^{-2/\beta})^n\\
&\quad=e^{-2n/\beta}.
\end{align*}
若电路板平均寿命$\beta$相对于$n$很大，可见这个概率接近1.

上述计算说明，如何利用\eqref{eq:5.1.1}、或更具体的\eqref{eq:5.1.2}所定义的随机样本pdf计算关于样本的概率.还应认识到，在这种计算中也可直接使用随机样本独立同分布的性质.例如，上述计算可以写成
\begin{align*}
P(X_1>2,\ldots,X_n>2)&=P(X_1>2)\cdots P(X_n>2)&&\text{（独立性）}\\
&=[P(X_1>2)]^n&&\text{（同分布）}\\
&=(e^{-2/\beta})^n&&\text{（指数分布的计算）}\\
&=e^{-2n/\beta}.
\end{align*}
\end{example}
\autoref{def:5.1.1}中的随机抽样模型有时称为从无限总体中抽样.\index{SI0920@无限总体}\semanticindexanchor{920:201}可以把它想象为依次取得$X_1,\ldots,X_n$的值.首先进行实验，观测到$X_1=x_1$；然后重复实验，观测到$X_2=x_2$.随机抽样中的独立性假定意味着，先观测到$X_1=x_1$这一事实不影响$X_2$的概率分布.从无限总体中“取出”$x_1$并不改变总体，因此$X_2=x_2$仍是来自同一总体的一个随机观测.

从有限总体中抽样时，\index{SI0919@有限总体}\semanticindexanchor{919:201}\autoref{def:5.1.1}是否适用取决于数据收集的方式.有限总体是一个有限数集$\{x_1,\ldots,x_N\}$.要从此总体中抽出样本$X_1,\ldots,X_n$.\autoref{sec:1.2.3}描述了四种抽取方法，这里讨论前两种.

设从总体中选取一个值，使$N$个值中的每一个都以相同概率$1/N$被选中.（可以想象从帽子里抽取数字.）把这个值记为$X_1=x_1$.然后重复这个过程，$N$个值中的每一个仍然等可能被选中.
\sourcepageanchor{202}
把第二次选中的值记为$X_2=x_2$.（若选中同一个数，则$x_1=x_2$.）将从这$N$个值中抽取的过程重复$n$次，就得到样本$X_1,\ldots,X_n$.这种抽样称为有放回抽样（sampling with replacement），\index{SI0942@有放回}\semanticindexanchor{942:202}因为每一阶段选中的值都“放回”总体，在下一阶段仍可被选中.这种抽样满足\autoref{def:5.1.1}的条件.每个$X_i$都是离散随机变量，等概率地取$x_1,\ldots,x_N$中的各个值.随机变量$X_1,\ldots,X_n$相互独立，因为不论其他变量取何值，选取任何$X_i$的过程都相同.（这种抽样用于自助法，见\autoref{sec:10.1.4}.）

从有限总体抽取随机样本的第二种方法称为无放回抽样（sampling without replacement）.\index{SI0943@无放回}\semanticindexanchor{943:202}其方法如下：从$\{x_1,\ldots,x_N\}$中选取一个值，各个值被选中的概率均为$1/N$，记为$X_1=x_1$.然后从剩余的$N-1$个值中选取第二个值，每个值被选中的概率均为$1/(N-1)$，记为$X_2=x_2$.继续以这种方式选取其余值，得到样本$X_1,\ldots,X_n$.但一个值一旦被选中，就不能在之后的任何阶段再次被选中.

从有限总体无放回抽出的样本不满足\autoref{def:5.1.1}的全部条件.随机变量$X_1,\ldots,X_n$并不相互独立.为说明这一点，设$x$和$y$是$\{x_1,\ldots,x_N\}$中的不同元素.则$P(X_2=y\mid X_1=y)=0$，因为若第一个阶段已选中$y$，第二个阶段便不能再选中它.然而，$P(X_2=y\mid X_1=x)=1/(N-1)$.因此$X_2$的概率分布依赖于观测到的$X_1$值，故$X_1$与$X_2$不独立.不过，有趣的是，$X_1,\ldots,X_n$仍然同分布，即每个$i=1,\ldots,n$的$X_i$具有相同的边际分布.显然，对每个$x\in\{x_1,\ldots,x_N\}$，$X_1$的边际分布为$P(X_1=x)=1/N$.为计算$X_2$的边际分布，利用\autoref{thm:1.2.11}\hyperref[itm:thm:1.2.11:a]{(a)}及条件概率的定义，写成
$$P(X_2=x)=\sum_{i=1}^N P(X_2=x\mid X_1=x_i)P(X_1=x_i).$$
对某个指标值（记为$k$），有$x=x_k$，且$P(X_2=x\mid X_1=x_k)=0$.对所有其他$j\ne k$，有$P(X_2=x\mid X_1=x_j)=1/(N-1)$.因此，
\begin{equation}\label{eq:5.1.3}
P(X_2=x)=(N-1)\left(\frac1{N-1}\frac1N\right)=\frac1N.
\end{equation}
类似论证可以说明每个$X_i$都有相同的边际分布.

从有限总体无放回抽样有时称为简单随机抽样（simple random sampling）.\index{SI0934@简单随机}\semanticindexanchor{934:202}必须认识到，它与\autoref{def:5.1.1}描述的抽样情形不同.不过，若总体容量$N$远大于样本容量$n$，则$X_1,\ldots,X_n$近似独立，可以假定它们独立来进行某些近似概率计算.说它们“近似独立”，仅仅意味着给定$X_1,\ldots,X_{i-1}$时$X_i$的条件分布与其边际分布差别不大.例如，给定$X_1$时$X_2$的条件分布是
$$P(X_2=x_1\mid X_1=x_1)=0,\qquad P(X_2=x\mid X_1=x_1)=\frac1{N-1}\quad(x\ne x_1).$$
若$N$很大，这与\eqref{eq:5.1.3}给出的$X_2$边际分布相差不大.给定$X_1,\ldots,X_{i-1}$时$X_i$条件分布中的非零概率为$1/(N-i+1)$；若$i\le n$相对于$N$很小，它就接近$1/N$.
\begin{example}[有限总体模型]\label{ex:5.1.3}
作为利用独立性进行近似计算的例子，设有限总体为$\{1,\ldots,1{,}000\}$，故$N=1{,}000$.从中无放回抽取容量$n=10$的样本.
\sourcepageanchor{203}
全部十个样本值都大于200的概率是多少？若$X_1,\ldots,X_{10}$相互独立，就有
\begin{equation}\label{eq:5.1.4}
\begin{aligned}
P(X_1>200,\ldots,X_{10}>200)&=P(X_1>200)\cdots P(X_{10}>200)\\
&=\left(\frac{800}{1{,}000}\right)^{10}=.107374.
\end{aligned}
\end{equation}
为精确计算这个概率，令随机变量$Y$表示样本中大于200的项数.则$Y$服从$\operatorname{hypergeometric}(N=1{,}000,M=800,K=10)$分布.因此
\begin{align*}
P(X_1>200,\ldots,X_{10}>200)&=P(Y=10)\\
&=\frac{\binom{800}{10}\binom{200}{0}}{\binom{1{,}000}{10}}\\
&=.106164.
\end{align*}
因此，\eqref{eq:5.1.4}是真实值的合理近似.
\end{example}
在本书其余部分，我们都采用\autoref{def:5.1.1}作为来自某总体的随机样本的定义.



% 原书印刷172–177页，PDF物理203–208页.
\subsection{随机样本中的随机变量之和}\label{sec:5.2}
抽得样本$X_1,\ldots,X_n$后，通常会计算这些数值的某种概要.任何定义明确的概要，在数学上都可以表示为函数$T(x_1,\ldots,x_n)$，其定义域包含随机向量$(X_1,\ldots,X_n)$的样本空间.函数$T$可以取实数值，也可以取向量值；因此这种概要是一个随机变量（或向量）$Y=T(X_1,\ldots,X_n)$.\autoref{sec:4}详细讨论了把随机变量定义为其他随机变量的函数的方法；该章的方法可用来根据产生样本的总体分布描述$Y$的分布.由于随机样本$X_1,\ldots,X_n$具有简单的概率结构（$X_i$独立同分布），$Y$的分布尤其容易处理.由于这一分布通常从随机样本中各变量的分布推导而来，称其为$Y$的抽样分布（sampling distribution）.这个名称区分了$Y$的概率分布和总体分布，即每个$X_i$的边际分布.本节讨论抽样分布的一些性质，特别是当函数$T(x_1,\ldots,x_n)$由随机变量之和定义时的性质.
\begin{definition}\label{def:5.2.1}\index{SI0951@抽样分布}\index{SI0981@统计量}
设$X_1,\ldots,X_n$为来自某总体的容量为$n$的随机样本，$T(x_1,\ldots,x_n)$为实值或向量值函数，其定义域包含$(X_1,\ldots,X_n)$的样本空间.则随机变量或随机向量$Y=T(X_1,\ldots,X_n)$称为统计量（statistic）.统计量$Y$的概率分布称为$Y$的抽样分布.
\end{definition}
统计量的定义十分宽泛，唯一的限制是统计量不能是参数的函数.统计量给出的样本概要可包含多种信息.例如，它可以给出样本的最小值或最大值、样本平均值，或样本观测值变异程度的度量.下面定义三个常用且能很好概括样本的统计量.
\begin{definition}\label{def:5.2.2}\index{SI0921@均值}
样本均值（sample mean）是随机样本各数值的算术平均，通常记为
$$\bar X=\frac{X_1+\cdots+X_n}{n}=\frac1n\sum_{i=1}^nX_i.$$
\end{definition}
\sourcepageanchor{204}
\begin{definition}\label{def:5.2.3}\index{SI0935@标准差}\index{SI0936@方差}
样本方差（sample variance）是如下定义的统计量：
$$S^2=\frac1{n-1}\sum_{i=1}^n(X_i-\bar X)^2.$$
样本标准差（sample standard deviation）是统计量$S=\sqrt{S^2}$.
\end{definition}
依照通常做法，上述定义省略了这些统计量的函数记号，即写$S$而不是$S(X_1,\ldots,X_n)$.统计量对样本的依赖是不言而喻的.与前面一样，用小写字母表示统计量的观测值；因此$\bar x$、$s^2$和$s$表示$\bar X$、$S^2$和$S$的观测值.

大家当然都熟悉样本均值.样本方差和标准差是衡量样本变异程度的量，它们与总体方差和标准差的关系将在下面说明.先推导样本均值和方差的一些性质.特别地，\autoref{thm:5.2.4}给出的样本方差关系与总体方差的类似关系\eqref{eq:2.3.1}相对应.
\begin{theorem}\label{thm:5.2.4}
设$x_1,\ldots,x_n$为任意数，$\bar x=(x_1+\cdots+x_n)/n$.则
\begin{enumerate}[label=\textbf{\alph*.}]
\item\label{item:5.2.4a} $\min_a\sum_{i=1}^n(x_i-a)^2=\sum_{i=1}^n(x_i-\bar x)^2$.
\item\label{item:5.2.4b} $(n-1)s^2=\sum_{i=1}^n(x_i-\bar x)^2=\sum_{i=1}^n x_i^2-n\bar x^2$.
\end{enumerate}
\end{theorem}
\begin{proof}
为证明\hyperref[item:5.2.4a]{(a)}，加上并减去$\bar x$得到
\begin{align*}
\sum_{i=1}^n(x_i-a)^2&=\sum_{i=1}^n(x_i-\bar x+\bar x-a)^2\\
&=\sum_{i=1}^n(x_i-\bar x)^2+2\sum_{i=1}^n(x_i-\bar x)(\bar x-a)+\sum_{i=1}^n(\bar x-a)^2\\
&=\sum_{i=1}^n(x_i-\bar x)^2+\sum_{i=1}^n(\bar x-a)^2.\qquad\text{（交叉项为零）}
\end{align*}
现在显然右端在$a=\bar x$时取得最小值.（注意与\autoref{ex:2.2.6}及\autoref{exercise:4.13}的相似性.）

在上式取$a=0$，便得到\hyperref[item:5.2.4b]{(b)}.
\end{proof}
\autoref{thm:5.2.4}\hyperref[item:5.2.4b]{(b)}中的表达式在计算和理论上都很有用，因为它把$s^2$表示为容易处理的和.

我们从考察一些统计量的期望开始研究抽样分布.下面的结果十分有用.
\begin{lemma}\label{lem:5.2.5}\index{SI0433@和}\index{SI1068@和}
设$X_1,\ldots,X_n$为来自某总体的随机样本，函数$g(x)$满足$Eg(X_1)$和$\operatorname{Var}g(X_1)$存在.则
\begin{equation}\label{eq:5.2.1}
E\left(\sum_{i=1}^n g(X_i)\right)=n(Eg(X_1))
\end{equation}
并且
\begin{equation}\label{eq:5.2.2}
\operatorname{Var}\left(\sum_{i=1}^n g(X_i)\right)=n(\operatorname{Var}g(X_1)).
\end{equation}
\end{lemma}
\begin{proof}
为证明\eqref{eq:5.2.1}，注意
$$E\left(\sum_{i=1}^n g(X_i)\right)=\sum_{i=1}^n Eg(X_i)=n(Eg(X_1)).$$
\sourcepageanchor{205}
由于$X_i$同分布，所有$i$的$Eg(X_i)$都相同，故第二个等号成立.注意，\eqref{eq:5.2.1}的成立不需要$X_1,\ldots,X_n$独立.事实上，它对任意$n$个同分布随机变量都成立.

为证明\eqref{eq:5.2.2}，注意
\begin{align*}
\operatorname{Var}\left(\sum_{i=1}^n g(X_i)\right)
&=E\left[\sum_{i=1}^n g(X_i)-E\left(\sum_{i=1}^n g(X_i)\right)\right]^2&&\text{（方差的定义）}\\
&=E\left[\sum_{i=1}^n(g(X_i)-Eg(X_i))\right]^2.&&\text{（期望的性质及重排各项）}
\end{align*}
最后一个表达式中共有$n^2$项.首先，有$n$项$(g(X_i)-Eg(X_i))^2$，$i=1,\ldots,n$，每一项都满足
\begin{align*}
E(g(X_i)-Eg(X_i))^2&=\operatorname{Var}g(X_i)&&\text{（方差的定义）}\\
&=\operatorname{Var}g(X_1).&&\text{（同分布）}
\end{align*}
其余$n(n-1)$项都形如$(g(X_i)-Eg(X_i))(g(X_j)-Eg(X_j))$，其中$i\ne j$.对每一项，
\begin{align*}
&E[(g(X_i)-Eg(X_i))(g(X_j)-Eg(X_j))]\\
&\quad=\operatorname{Cov}(g(X_i),g(X_j))&&\text{（协方差的定义）}\\
&\quad=0.&&\text{（独立性，\autoref{thm:4.5.5}）}
\end{align*}
由此得到\eqref{eq:5.2.2}.
\end{proof}
\begin{theorem}\label{thm:5.2.6}\index{SI0923@期望值}\index{SI0938@期望值}
设$X_1,\ldots,X_n$为来自均值为$\mu$、方差为$\sigma^2<\infty$的总体的随机样本.则
\begin{enumerate}[label=\textbf{\alph*.}]
\item\label{item:5.2.6a} $E\bar X=\mu$；
\item\label{item:5.2.6b} $\operatorname{Var}\bar X=\sigma^2/n$；
\item\label{item:5.2.6c} $ES^2=\sigma^2$.
\end{enumerate}
\end{theorem}
\begin{proof}
为证明\hyperref[item:5.2.6a]{(a)}，令$g(X_i)=X_i/n$，则$Eg(X_i)=\mu/n$.由\autoref{lem:5.2.5}，
$$E\bar X=E\left(\frac1n\sum_{i=1}^nX_i\right)=\frac1nE\left(\sum_{i=1}^nX_i\right)=\frac1n nEX_1=\mu.$$
类似地，对\hyperref[item:5.2.6b]{(b)}，有
$$\operatorname{Var}\bar X=\operatorname{Var}\left(\frac1n\sum_{i=1}^n X_i\right)=\frac1{n^2}\operatorname{Var}\left(\sum_{i=1}^nX_i\right)=\frac1{n^2}n\operatorname{Var}X_1=\frac{\sigma^2}{n}.$$
对样本方差，利用\autoref{thm:5.2.4}，有
\begin{align*}
ES^2&=E\left(\frac1{n-1}\left[\sum_{i=1}^nX_i^2-n\bar X^2\right]\right)=\frac1{n-1}(nEX_1^2-nE\bar X^2)\\
&=\frac1{n-1}\left(n(\sigma^2+\mu^2)-n\left(\frac{\sigma^2}{n}+\mu^2\right)\right)=\sigma^2,
\end{align*}
这就证明了\hyperref[item:5.2.6c]{(c)}及整个定理.
\end{proof}
\sourcepageanchor{206}
\autoref{thm:5.2.6}中的\hyperref[item:5.2.6a]{(a)}和\hyperref[item:5.2.6c]{(c)}是统计量与总体参数之间的关系，也是无偏统计量（unbiased statistics）的例子.\index{SI1044@无偏估计量}\semanticindexanchor{1044:206}\autoref{sec:7}将讨论这些概念.$\bar X$是$\mu$的无偏估计量，$S^2$是$\sigma^2$的无偏估计量.$S^2$定义中使用$n-1$可能显得不够直观；现在可以看出，这样定义使$ES^2=\sigma^2$.若按通常平均数的方式定义$S^2$，使离差平方和的分母为$n$而非$n-1$，则$ES^2$会等于$\frac{n-1}{n}\sigma^2$，$S^2$就不是$\sigma^2$的无偏估计量.

下面更详细地讨论$\bar X$的抽样分布.\autoref{sec:4.3}和\autoref{sec:4.6}的方法可用来从总体分布推导这一抽样分布.但由于随机样本特有的概率结构（iid随机变量），所得$\bar X$的抽样分布具有简单的表达形式.

先指出一些简单关系.由于$\bar X=(X_1+\cdots+X_n)/n$，若$f(y)$是$Y=X_1+\cdots+X_n$的pdf，则$f_{\bar X}(x)=nf(nx)$是$\bar X$的pdf（见\autoref{exercise:5.5}）.因此，关于$Y$的pdf的结果容易转化为关于$\bar X$的pdf的结果.矩母函数也有类似关系：
$$M_{\bar X}(t)=Ee^{t\bar X}=Ee^{t(X_1+\cdots+X_n)/n}=Ee^{(t/n)Y}=M_Y(t/n).$$
由于$X_1,\ldots,X_n$同分布，各$i$的$M_{X_i}(t)$都是同一函数.由\autoref{thm:4.6.7}可得如下结论.
\begin{theorem}\label{thm:5.2.7}\index{SI0694@样本均值}\index{SI0925@矩母函数}
设$X_1,\ldots,X_n$是来自矩母函数为$M_X(t)$的总体的随机样本，则样本均值的矩母函数为
$$M_{\bar X}(t)=[M_X(t/n)]^n.$$
\end{theorem}
当然，只有当$M_{\bar X}(t)$的表达式是某个熟悉的矩母函数时，\autoref{thm:5.2.7}才有用.这种情况比较有限，但下面的例子表明，一旦这种方法适用，就能十分巧妙地推导出$\bar X$的抽样分布.
\begin{example}[均值的分布]\label{ex:5.2.8}\index{SI0926@正态}
设$X_1,\ldots,X_n$为来自$n(\mu,\sigma^2)$总体的随机样本.样本均值的矩母函数为
\begin{align*}
M_{\bar X}(t)&=\left[\exp\left(\mu\frac tn+\frac{\sigma^2(t/n)^2}{2}\right)\right]^n\\
&=\exp\left(n\left(\mu\frac tn+\frac{\sigma^2(t/n)^2}{2}\right)\right)=\exp\left(\mu t+\frac{(\sigma^2/n)t^2}{2}\right).
\end{align*}
故$\bar X$服从$n(\mu,\sigma^2/n)$分布.

另一个简单例子是$\operatorname{gamma}(\alpha,\beta)$随机样本（见\autoref{ex:4.6.8}）.同样容易推导样本均值的分布，其矩母函数为
$$M_{\bar X}(t)=\left[\left(\frac1{1-\beta(t/n)}\right)^\alpha\right]^n=\left(\frac1{1-(\beta/n)t}\right)^{n\alpha},$$
这正是$\operatorname{gamma}(n\alpha,\beta/n)$的矩母函数，因此$\bar X$服从该分布.
\end{example}
若\autoref{thm:5.2.7}不适用——或者因为无法辨认所得$\bar X$的矩母函数，或者因为总体矩母函数不存在——则可尝试用\autoref{sec:4.3}和\autoref{sec:4.6}的变换法求$Y=X_1+\cdots+X_n$及$\bar X$的pdf.这时，下面的卷积公式（convolution formula）很有用.
\begin{theorem}\label{thm:5.2.9}\index{SI0213@卷积}
若$X$和$Y$是独立的连续随机变量，pdf分别为$f_X(x)$及$f_Y(y)$，则$Z=X+Y$的pdf为
\begin{equation}\label{eq:5.2.3}
f_Z(z)=\int_{-\infty}^\infty f_X(w)f_Y(z-w)\,dw.
\end{equation}
\end{theorem}
\sourcepageanchor{207}
\begin{proof}
令$W=X$.从$(X,Y)$到$(Z,W)$的变换的Jacobian为1.因此，由\eqref{eq:4.3.2}得到$(Z,W)$的联合pdf为
$$f_{Z,W}(z,w)=f_{X,Y}(w,z-w)=f_X(w)f_Y(z-w).$$
对$w$积分消去它，就得到\eqref{eq:5.2.3}给出的$Z$的边际pdf.
\end{proof}
若$f_X$、$f_Y$之一或二者只在某些值处为正，\eqref{eq:5.2.3}中的积分限可以相应调整.例如，若$f_X$和$f_Y$都仅在正值处为正，则积分限为0和$z$，因为$w$在此范围之外时被积函数为零.对于加法以外的运算，也可推导类似\eqref{eq:5.2.3}的卷积公式，例如差、积及商的公式（见\autoref{exercise:5.6}）.
\begin{example}[Cauchy随机变量之和]\label{ex:5.2.10}\index{SI0922@柯西}
作为矩母函数方法失效的例子，考虑从Cauchy分布抽样.我们最终要推导$\bar Z$的分布，其中$\bar Z$是iid $\operatorname{Cauchy}(0,1)$观测$Z_1,\ldots,Z_n$的均值.不过，先从两个独立Cauchy随机变量之和的分布入手，应用\eqref{eq:5.2.3}.

设$U$和$V$是独立的Cauchy随机变量，$U\sim\operatorname{Cauchy}(0,\sigma)$，$V\sim\operatorname{Cauchy}(0,\tau)$，即
$$f_U(u)=\frac1{\pi\sigma}\frac1{1+(u/\sigma)^2},\quad f_V(v)=\frac1{\pi\tau}\frac1{1+(v/\tau)^2},\qquad -\infty<u,v<\infty.$$
由\eqref{eq:5.2.3}，$Z=U+V$的pdf为
\begin{equation}\label{eq:5.2.4}
f_Z(z)=\int_{-\infty}^\infty\frac1{\pi\sigma}\frac1{1+(w/\sigma)^2}\frac1{\pi\tau}\frac1{1+((z-w)/\tau)^2}\,dw,\quad -\infty<z<\infty.
\end{equation}
这个积分有些复杂，但利用部分分式分解并仔细求原函数即可算出（见\autoref{exercise:5.7}）.结果为
\begin{equation}\label{eq:5.2.5}
f_Z(z)=\frac1{\pi(\sigma+\tau)}\frac1{1+(z/(\sigma+\tau))^2},\qquad -\infty<z<\infty.
\end{equation}
因此，两个独立Cauchy随机变量之和仍服从Cauchy分布，其尺度参数相加.由此，若$Z_1,\ldots,Z_n$为iid $\operatorname{Cauchy}(0,1)$随机变量，则$\sum Z_i$服从$\operatorname{Cauchy}(0,n)$，而且$\bar Z$仍服从$\operatorname{Cauchy}(0,1)$！样本均值与单个观测具有相同的分布.（这一计算的计算机代数版本，见计算机代数附录中的\autoref{ex:12.6.5}.）
\end{example}
若从位置–尺度族或某些类型的指数族抽样，随机变量之和，特别是$\bar X$的抽样分布，很容易推导.本节最后讨论这两种情形.

先讨论\autoref{sec:3.5}中的位置–尺度情形.\index{SI0924@位置尺度}\semanticindexanchor{924:208}设$X_1,\ldots,X_n$为来自位置–尺度族成员$(1/\sigma)\allowbreak f((x-\mu)/\sigma)$的随机样本.则$\bar X$的分布与$\bar Z$的分布有简单的关系，其中$\bar Z$是来自标准pdf $f(z)$的随机样本的样本均值.为说明这种关系，由\autoref{thm:3.5.6}，存在随机变量$Z_1,\ldots,Z_n$，满足$X_i=\sigma Z_i+\mu$，且每个$Z_i$的pdf为$f(z)$.此外，$Z_1,\ldots,Z_n$相互独立，所以构成来自$f(z)$的随机样本.样本均值$\bar X$与$\bar Z$满足
$$\bar X=\frac1n\sum_{i=1}^nX_i=\frac1n\sum_{i=1}^n(\sigma Z_i+\mu)=\frac1n\left(\sigma\sum_{i=1}^nZ_i+n\mu\right)=\sigma\bar Z+\mu.$$
\sourcepageanchor{208}
因此，再次利用\autoref{thm:3.5.6}，若$g(z)$是$\bar Z$的pdf，则$(1/\sigma)g((x-\mu)/\sigma)$是$\bar X$的pdf.先利用$Z_1,\ldots,Z_n$和$f(z)$求出$\bar Z$的pdf $g(z)$，可能更容易.这样无须处理参数$\mu$和$\sigma$，计算可能不那么繁琐；随后立即可知$\bar X$的pdf为$(1/\sigma)g((x-\mu)/\sigma)$.

\autoref{ex:5.2.10}说明，若$Z_1,\ldots,Z_n$为来自$\operatorname{Cauchy}(0,1)$分布的随机样本，则$\bar Z$也服从$\operatorname{Cauchy}(0,1)$分布.现在可以断定，若$X_1,\ldots,X_n$为来自$\operatorname{Cauchy}(\mu,\sigma)$分布的随机样本，则$\bar X$也服从$\operatorname{Cauchy}(\mu,\sigma)$分布.须特别注意，以$\sigma$衡量的$\bar X$分布的离散程度不随样本容量$n$改变.这与\autoref{thm:5.2.6}中更常见的情形（总体方差有限）形成鲜明对照；在后者中，$\operatorname{Var}\bar X=\sigma^2/n$随样本容量增大而减小.

从指数族抽样时，随机样本中某些和的抽样分布容易推导.下一个定理中的$T_1,\ldots,T_k$是重要的样本概要统计量，这将在\autoref{sec:6.2}中看到.
\begin{theorem}\label{thm:5.2.11}\index{SI0441@和}\index{SI0918@指数族}
设$X_1,\ldots,X_n$为来自pdf或pmf $f(x\mid\theta)$的随机样本，其中
$$f(x\mid\theta)=h(x)c(\theta)\exp\left(\sum_{i=1}^k w_i(\theta)t_i(x)\right)$$
是指数族的一个成员.定义统计量$T_1,\ldots,T_k$为
$$T_i(X_1,\ldots,X_n)=\sum_{j=1}^n t_i(X_j),\qquad i=1,\ldots,k.$$
若集合$\{(w_1(\theta),w_2(\theta),\ldots,w_k(\theta)),\theta\in\Theta\}$包含$\mathbb R^k$的一个开子集，则$(T_1,\ldots,T_k)$的分布是形如下式的指数族：
\begin{equation}\label{eq:5.2.6}
f_T(u_1,\ldots,u_k\mid\theta)=H(u_1,\ldots,u_k)[c(\theta)]^n\exp\left(\sum_{i=1}^k w_i(\theta)u_i\right).
\end{equation}
\end{theorem}
开集条件排除了$n(\theta,\theta^2)$这样的密度；\index{SI0299@曲线}\semanticindexanchor{299:208}一般来说，它把弯曲指数族排除在\autoref{thm:5.2.11}之外.注意，在$(T_1,\ldots,T_k)$的pdf或pmf中，函数$c(\theta)$和$w_i(\theta)$与原来族中的相同，当然函数$H(u_1,\ldots,u_k)$不同于$h(x)$.这里不证明该定理，只在一个简单例子中说明这一结果.
\begin{example}[Bernoulli随机变量之和]\label{ex:5.2.12}
设$X_1,\ldots,X_n$为来自$\operatorname{Bernoulli}(p)$分布的随机样本.由\autoref{ex:3.4.1}（取$n=1$），$\operatorname{Bernoulli}(p)$是指数族，且$k=1$、$c(p)=1-p$、$w_1(p)=\log(p/(1-p))$、$t_1(x)=x$.所以在前一定理中，$T_1=T_1(X_1,\ldots,X_n)=X_1+\cdots+X_n$.由\autoref{sec:3.2}中二项分布的定义可知，$T_1$服从$\operatorname{binomial}(n,p)$分布.由\autoref{ex:3.4.1}还可知，$\operatorname{binomial}(n,p)$是具有同一$w_1(p)$及$c(p)=(1-p)^n$的指数族.因此，在本例中验证了\eqref{eq:5.2.6}.
\end{example}



% 原书印刷177–184页，PDF物理208–215页.
\subsection{从正态分布抽样}\label{sec:5.3}
本节讨论从正态总体抽取的样本量的性质；正态总体仍然是应用最广的统计模型之一.从正态总体抽样会使样本统计量具有许多有用的性质，也会产生许多熟知的抽样分布.
\sourcepageanchor{209}
\subsubsection{样本均值与方差的性质}\label{sec:5.3.1}
前面已经看到，一般情形下如何计算$\bar X$与$S^2$的均值和方差.现在加上正态性假设，就可以推导它们的完整分布以及更多结论.下面的定理概括了$\bar X$与$S^2$的性质.
\begin{theorem}\label{thm:5.3.1}\index{SI0926@正态}\index{SI0940@正态}
设$X_1,\ldots,X_n$为来自$n(\mu,\sigma^2)$分布的随机样本，$\bar X=(1/n)\sum_{i=1}^nX_i$，$S^2=[1/(n-1)]\sum_{i=1}^n(X_i-\bar X)^2$.则
\begin{enumerate}[label=\textbf{\alph*.}]
\item\label{item:5.3.1a} $\bar X$与$S^2$是独立随机变量；
\item\label{item:5.3.1b} $\bar X$服从$n(\mu,\sigma^2/n)$分布；
\item\label{item:5.3.1c} $(n-1)S^2/\sigma^2$服从自由度为$n-1$的卡方分布（chi squared distribution）.
\end{enumerate}
\end{theorem}
\begin{proof}
首先，由\autoref{sec:3.5}关于位置–尺度族的结论，不失一般性，可以假定$\mu=0$、$\sigma=1$.（另见\autoref{thm:5.2.11}之前的讨论.）此外，\hyperref[item:5.3.1b]{(b)}已在\autoref{ex:5.2.8}中证明，所以只须证明\hyperref[item:5.3.1a]{(a)}和\hyperref[item:5.3.1c]{(c)}.

为证明\hyperref[item:5.3.1a]{(a)}，将应用\autoref{thm:4.6.12}，说明$\bar X$与$S^2$是独立随机向量的函数.注意$S^2$可写成$n-1$个离差的函数：
\begin{align*}
S^2&=\frac1{n-1}\sum_{i=1}^n(X_i-\bar X)^2\\
&=\frac1{n-1}\left((X_1-\bar X)^2+\sum_{i=2}^n(X_i-\bar X)^2\right)\\
&=\frac1{n-1}\left(\left[\sum_{i=2}^n(X_i-\bar X)\right]^2+\sum_{i=2}^n(X_i-\bar X)^2\right),
\end{align*}
因为$\sum_{i=1}^n(X_i-\bar X)=0$.所以$S^2$可写成仅依赖于$(X_2-\bar X,\ldots,X_n-\bar X)$的函数.下面说明这些随机变量与$\bar X$独立.样本$X_1,\ldots,X_n$的联合pdf为
$$f(x_1,\ldots,x_n)=\frac1{(2\pi)^{n/2}}e^{-(1/2)\sum_{i=1}^n x_i^2},\qquad -\infty<x_i<\infty.$$
作变换
$$\begin{aligned}
y_1&=\bar x,\\y_2&=x_2-\bar x,\\&\ \vdots\\y_n&=x_n-\bar x.
\end{aligned}$$
这是Jacobian等于$1/n$的线性变换.有
\begin{align*}
&f(y_1,\ldots,y_n)\\
&\quad=\frac n{(2\pi)^{n/2}}e^{-(1/2)(y_1-\sum_{i=2}^n y_i)^2}e^{-(1/2)\sum_{i=2}^n(y_i+y_1)^2},\qquad -\infty<y_i<\infty\\
&\quad=\left[\left(\frac n{2\pi}\right)^{1/2}e^{-ny_1^2/2}\right]
\left[\frac{n^{1/2}}{(2\pi)^{(n-1)/2}}e^{-(1/2)[\sum_{i=2}^n y_i^2+(\sum_{i=2}^n y_i)^2]}\right],\quad -\infty<y_i<\infty.
\end{align*}
由于$Y_1,\ldots,Y_n$的联合pdf可以这样分解，由\autoref{thm:4.6.11}可知$Y_1$与$Y_2,\ldots,Y_n$独立，再由\autoref{thm:4.6.12}可知$\bar X$与$S^2$独立.
\end{proof}
为完成定理的证明，还须推导$S^2$的分布.但在此之前，先稍作离题，讨论卡方分布；其性质在推导$S^2$的pdf时起着重要作用.
\sourcepageanchor{210}
回顾\autoref{sec:3.3}，卡方pdf是gamma pdf的特例，表达式为
$$f(x)=\frac1{\Gamma(p/2)2^{p/2}}x^{(p/2)-1}e^{-x/2},\qquad 0<x<\infty,$$
其中$p$称为自由度（degrees of freedom）.现概括卡方分布的一些相关事实.
\begin{lemma}[关于卡方随机变量的事实]\label{lem:5.3.2}\index{SI0282@性质}
以$\chi_p^2$表示自由度为$p$的卡方随机变量.
\begin{enumerate}[label=\textbf{\alph*.}]
\item\label{item:5.3.2a} 若$Z$是$n(0,1)$随机变量，则$Z^2\sim\chi_1^2$，即标准正态随机变量的平方是卡方随机变量.
\item\label{item:5.3.2b} 若$X_1,\ldots,X_n$独立且$X_i\sim\chi_{p_i}^2$，则$X_1+\cdots+X_n\sim\chi_{p_1+\cdots+p_n}^2$，即独立卡方变量之和仍是卡方变量，自由度也相加.
\end{enumerate}
\end{lemma}
\begin{proof}
这些事实前面已经遇到过.\hyperref[item:5.3.2a]{(a)}已在\autoref{ex:2.1.7}中证明.\hyperref[item:5.3.2b]{(b)}是\autoref{ex:4.6.8}中独立gamma随机变量之和的特例.由于$\chi_p^2$随机变量就是$\operatorname{gamma}(p/2,2)$，应用该例即可得到\hyperref[item:5.3.2b]{(b)}.
\end{proof}
\begin{proof}[{\autoref{thm:5.3.1}\hyperref[item:5.3.1c]{(c)}的证明}]
用归纳法确定$S^2$的分布，以$\bar X_k$和$S_k^2$表示基于前$k$个观测的样本均值和方差.（观测的实际次序无关紧要；只是为了方便证明而把它们看作有序.）容易证明（见\autoref{exercise:5.15}）
\begin{equation}\label{eq:5.3.1}
(n-1)S_n^2=(n-2)S_{n-1}^2+\left(\frac{n-1}{n}\right)(X_n-\bar X_{n-1})^2.
\end{equation}
先考虑$n=2$.定义$0\times S_1^2=0$，则由\eqref{eq:5.3.1}有
$$S_2^2=\frac12(X_2-X_1)^2.$$
由于$(X_2-X_1)/\sqrt2$服从$n(0,1)$，\autoref{lem:5.3.2}\hyperref[item:5.3.2a]{(a)}表明$S_2^2\sim\chi_1^2$.继续归纳，假设$n=k$时$(k-1)S_k^2\sim\chi_{k-1}^2$.当$n=k+1$时，由\eqref{eq:5.3.1}，
\begin{equation}\label{eq:5.3.2}
kS_{k+1}^2=(k-1)S_k^2+\left(\frac k{k+1}\right)(X_{k+1}-\bar X_k)^2.
\end{equation}
由归纳假设，$(k-1)S_k^2\sim\chi_{k-1}^2$.若能证明$(k/(k+1))(X_{k+1}-\bar X_k)^2\sim\chi_1^2$且与$S_k^2$独立，则由\autoref{lem:5.3.2}\hyperref[item:5.3.2b]{(b)}可得$kS_{k+1}^2\sim\chi_k^2$，定理即得证.

$(X_{k+1}-\bar X_k)^2$与$S_k^2$的独立性再次由\autoref{thm:4.6.12}得到.向量$(X_{k+1},\bar X_k)$与$S_k^2$独立，故向量的任何函数也与$S_k^2$独立.此外，$X_{k+1}-\bar X_k$是均值为0、方差为
$$\operatorname{Var}(X_{k+1}-\bar X_k)=\frac{k+1}{k}$$
的正态随机变量，所以$(k/(k+1))(X_{k+1}-\bar X_k)^2\sim\chi_1^2$，定理得证.
\end{proof}
证明$\bar X$与$S^2$独立，还可以采用不同于\autoref{thm:5.3.1}证明的方法.不必验证联合pdf能分解，可以使用下面这个把正态样本的独立性与相关性联系起来的引理.
\sourcepageanchor{211}
\begin{lemma}\label{lem:5.3.3}\index{SI0362@线性组合}
设$X_j\sim n(\mu_j,\sigma_j^2)$，$j=1,\ldots,n$，相互独立.对常数$a_{ij},b_{rj}$（$j=1,\ldots,n$；$i=1,\ldots,k$；$r=1,\ldots,m$），其中$k+m\le n$，定义
$$\begin{aligned}
U_i&=\sum_{j=1}^n a_{ij}X_j,\qquad i=1,\ldots,k,\\
V_r&=\sum_{j=1}^n b_{rj}X_j,\qquad r=1,\ldots,m.
\end{aligned}$$
\begin{enumerate}[label=\textbf{\alph*.}]
\item\label{item:5.3.3a} 随机变量$U_i$与$V_r$独立，当且仅当$\operatorname{Cov}(U_i,V_r)=0$.此外，$\operatorname{Cov}(U_i,V_r)=\sum_{j=1}^n a_{ij}b_{rj}\sigma_j^2$.
\item\label{item:5.3.3b} 随机向量$(U_1,\ldots,U_k)$与$(V_1,\ldots,V_m)$独立，当且仅当所有配对$i,r$（$i=1,\ldots,k$；$r=1,\ldots,m$）的$U_i$与$V_r$均独立.
\end{enumerate}
\end{lemma}
\begin{proof}
只须证明$\mu_i=0$、$\sigma_i^2=1$时的结论，一般形式就能很快得到.此外，由独立性推出协方差为零是直接的（\autoref{thm:4.5.5}），协方差表达式也容易验证（\autoref{exercise:5.14}）.还应注意，\autoref{cor:4.6.10}表明$U_i$与$V_r$都服从正态分布.

所以，只需证明：若常数满足上述限制（等价地，协方差为零），则在正态性假设下有独立性.这里只证明$n=2$的情形；一般$n$的证明类似，但需要详尽的$n$元变换.

为证明\hyperref[item:5.3.3a]{(a)}，从$X_1$与$X_2$的联合pdf出发：
$$f_{X_1,X_2}(x_1,x_2)=\frac1{2\pi}e^{-(1/2)(x_1^2+x_2^2)},\qquad -\infty<x_1,x_2<\infty.$$
作变换（$n=2$时可省略双下标）
$$u=a_1x_1+a_2x_2,\qquad v=b_1x_1+b_2x_2,$$
于是
$$x_1=\frac{b_2u-a_2v}{a_1b_2-b_1a_2},\qquad x_2=\frac{a_1v-b_1u}{a_1b_2-b_1a_2},$$
Jacobian为
$$J=\begin{vmatrix}\dfrac{\partial x_1}{\partial u}&\dfrac{\partial x_1}{\partial v}\\[4pt]\dfrac{\partial x_2}{\partial u}&\dfrac{\partial x_2}{\partial v}\end{vmatrix}=\frac1{a_1b_2-b_1a_2}.$$
因此，$U$与$V$的pdf为
\begin{align*}
f_{U,V}(u,v)&=f_{X_1,X_2}\left(\frac{b_2u-a_2v}{a_1b_2-b_1a_2},\frac{a_1v-b_1u}{a_1b_2-b_1a_2}\right)|J|\\
&=\frac1{2\pi}\exp\left\{\frac{-1}{2(a_1b_2-b_1a_2)^2}[(b_2u-a_2v)^2+(a_1v-b_1u)^2]\right\}|J|,
\end{align*}
其中$-\infty<u,v<\infty$.展开指数中的平方，得
$$\begin{aligned}
&(b_2u-a_2v)^2+(a_1v-b_1u)^2\\
&\quad=(b_1^2+b_2^2)u^2+(a_1^2+a_2^2)v^2-2(a_1b_1+a_2b_2)uv.
\end{aligned}$$
常数所满足的假设表明交叉项恒等于零.因此pdf可分解，由\autoref{lem:4.2.7}，$U$与$V$独立，\hyperref[item:5.3.3a]{(a)}得证.

类似论证可用于\hyperref[item:5.3.3b]{(b)}，这里不详述.作适当变换便可得到向量$(U_1,\ldots,U_k)$与$(V_1,\ldots,V_m)$的联合pdf.由\autoref{thm:4.6.11}，联合pdf可分解时，这两个向量独立.根据正态pdf的形式，这当且仅当所有配对$i,r$（$i=1,\ldots,k$；$r=1,\ldots,m$）的$U_i$与$V_r$独立时成立.
\end{proof}
\sourcepageanchor{212}
这个引理说明，对于独立正态随机变量的线性函数，协方差为零与独立性等价.因此，只须检查协方差项，就能验证这些正态变量的独立性；计算简单得多.这并没有什么神奇之处，只是正态pdf形式的结果.此外，\hyperref[item:5.3.3b]{(b)}允许仅通过两两独立性推断正态向量整体的独立性，这一性质对一般随机变量并不成立.

利用\autoref{lem:5.3.3}可以给出正态抽样中$\bar X$与$S^2$独立的另一证明.由于$S^2$可写成$n-1$个离差$(X_2-\bar X,\ldots,X_n-\bar X)$的函数，只须证明这些随机变量与$\bar X$不相关.正态性假设结合\autoref{lem:5.3.3}便能推出独立性.

为说明\autoref{lem:5.3.3}的应用，写成
$$\begin{aligned}
\bar X&=\sum_{i=1}^n\left(\frac1n\right)X_i,\\
X_j-\bar X&=\sum_{i=1}^n\left(\delta_{ij}-\frac1n\right)X_i,
\end{aligned}$$
其中$i=j$时$\delta_{ij}=1$，否则$\delta_{ij}=0$.于是容易证明
$$\operatorname{Cov}(\bar X,X_j-\bar X)=\sum_{i=1}^n\left(\frac1n\right)\left(\delta_{ij}-\frac1n\right)=0,$$
说明$\bar X$与$X_j-\bar X$独立（只要各$X_i$的方差相同）.
\subsubsection{导出分布：Student的\texorpdfstring{$t$}{t}与Snedecor的\texorpdfstring{$F$}{F}}\label{sec:5.3.2}
从某种意义上说，\autoref{sec:5.3.1}推导的分布是在正态性假设下进行统计分析的第一步.特别地，实际情形中方差$\sigma^2$大多未知，因此，要了解$\bar X$（作为$\mu$的估计）的变异程度，就必须估计这一方差.W. S. Gosset（以Student为笔名发表论文）在20世纪初首次研究这个问题.\index{SI0461@Gosset，W. S.}\semanticindexanchor{461:212}Student的奠基性工作产生了Student的$t$分布，简称$t$分布.

若$X_1,\ldots,X_n$为来自$n(\mu,\sigma^2)$的随机样本，已知
\begin{equation}\label{eq:5.3.3}\frac{\bar X-\mu}{\sigma/\sqrt n}\end{equation}
服从$n(0,1)$分布.若知道$\sigma$的值，并测得$\bar X$，则可用\eqref{eq:5.3.3}作为关于$\mu$推断的依据，因为这时$\mu$是唯一未知量.但多数时候$\sigma$未知.Student做了一个自然的尝试，考察
\begin{equation}\label{eq:5.3.4}\frac{\bar X-\mu}{S/\sqrt n},\end{equation}
这个量可在$\sigma$未知时用作推断$\mu$的依据.

只要先注意一些简化步骤，\eqref{eq:5.3.4}的分布就容易推导.将它乘以$\sigma/\sigma$并稍加整理，得到
\begin{equation}\label{eq:5.3.5}
\frac{\bar X-\mu}{S/\sqrt n}=\frac{(\bar X-\mu)/(\sigma/\sqrt n)}{\sqrt{S^2/\sigma^2}}.
\end{equation}
\eqref{eq:5.3.5}的分子是$n(0,1)$随机变量，分母是$\sqrt{\chi_{n-1}^2/(n-1)}$，且与分子独立.因此，求\eqref{eq:5.3.4}的分布可简化为求$U/\sqrt{V/p}$的分布，其中$U\sim n(0,1)$、$V\sim\chi_p^2$且$U$、$V$独立.由此得到Student的$t$分布.
\sourcepageanchor{213}
\begin{definition}\label{def:5.3.4}\index{SI0381@t}
设$X_1,\ldots,X_n$为来自$n(\mu,\sigma^2)$分布的随机样本，则$(\bar X-\mu)/(S/\sqrt n)$服从自由度为$n-1$的Student $t$分布.等价地，若随机变量$T$具有pdf
\begin{equation}\label{eq:5.3.6}
f_T(t)=\frac{\Gamma\left(\frac{p+1}2\right)}{\Gamma\left(\frac p2\right)}\frac1{(p\pi)^{1/2}}\frac1{(1+t^2/p)^{(p+1)/2}},\qquad -\infty<t<\infty,
\end{equation}
则称它服从自由度为$p$的Student $t$分布，记作$T\sim t_p$.
\end{definition}
注意，$p=1$时，\eqref{eq:5.3.6}就是Cauchy分布的pdf，这出现在样本容量为2时.Cauchy分布再次出现在一个普通情形中.

$t$的pdf的推导很直接.从上面定义的$U$和$V$出发，由\eqref{eq:5.3.5}，其联合pdf为
$$f_{U,V}(u,v)=\frac1{(2\pi)^{1/2}}e^{-u^2/2}\frac1{\Gamma(p/2)2^{p/2}}v^{(p/2)-1}e^{-v/2},\quad -\infty<u<\infty,\quad 0<v<\infty.$$
（记住$U$、$V$独立.）现在作变换
$$t=\frac u{\sqrt{v/p}},\qquad w=v.$$
变换的Jacobian为$(w/p)^{1/2}$，$T$的边际pdf为
\begin{align*}
f_T(t)&=\int_0^\infty f_{U,V}\left(t\left(\frac wp\right)^{1/2},w\right)\left(\frac wp\right)^{1/2}\,dw\\
&=\frac1{(2\pi)^{1/2}}\frac1{\Gamma(p/2)2^{p/2}}\int_0^\infty e^{-(1/2)t^2w/p}w^{(p/2)-1}e^{-w/2}\left(\frac wp\right)^{1/2}\,dw\\
&=\frac1{(2\pi)^{1/2}}\frac1{\Gamma(p/2)2^{p/2}p^{1/2}}\int_0^\infty e^{-(1/2)(1+t^2/p)w}w^{((p+1)/2)-1}\,dw.
\end{align*}
识别出被积函数是$\operatorname{gamma}((p+1)/2,2/(1+t^2/p))$的pdf核，因而
$$f_T(t)=\frac1{(2\pi)^{1/2}}\frac1{\Gamma(p/2)2^{p/2}p^{1/2}}\Gamma\left(\frac{p+1}2\right)\left[\frac2{1+t^2/p}\right]^{(p+1)/2},$$
它等于\eqref{eq:5.3.6}.

Student的$t$分布没有矩母函数，因为它并非所有阶矩都存在.事实上，自由度为$p$时，仅有$p-1$个矩.所以$t_1$没有均值，$t_2$没有方差，依此类推.容易验证（见\autoref{exercise:5.18}），若$T_p$服从$t_p$分布，则
\begin{equation}\label{eq:5.3.7}
\begin{aligned}ET_p&=0,&&p>1,\\ \operatorname{Var}T_p&=\frac p{p-2},&&p>2.\end{aligned}
\end{equation}
另一个重要导出分布是Snedecor的$F$分布，其推导与Student的$t$分布很相似，但动机有所不同.$F$分布以纪念Ronald Fisher爵士而命名，自然地作为方差比的分布出现.
\begin{example}[方差比的分布]\label{ex:5.3.5}
设$X_1,\ldots,X_n$为来自$n(\mu_X,\sigma_X^2)$总体的随机样本，$Y_1,\ldots,Y_m$为来自独立的$n(\mu_Y,\sigma_Y^2)$总体的随机样本.若要比较两个总体的变异程度，感兴趣的量之一是比值$\sigma_X^2/\sigma_Y^2$.
\sourcepageanchor{214}
样本方差之比$S_X^2/S_Y^2$包含关于该比值的信息.$F$分布给出
\begin{equation}\label{eq:5.3.8}
\frac{S_X^2/S_Y^2}{\sigma_X^2/\sigma_Y^2}=\frac{S_X^2/\sigma_X^2}{S_Y^2/\sigma_Y^2}
\end{equation}
的分布，从而可以比较这些量.考察\eqref{eq:5.3.8}便能看出如何推导$F$分布.比值$S_X^2/\sigma_X^2$和$S_Y^2/\sigma_Y^2$各自都是按比例缩放的卡方变量，而且相互独立.
\end{example}
\begin{definition}\label{def:5.3.6}\index{SI0301@F}
设$X_1,\ldots,X_n$为来自$n(\mu_X,\sigma_X^2)$总体的随机样本，$Y_1,\ldots,Y_m$为来自独立的$n(\mu_Y,\sigma_Y^2)$总体的随机样本.随机变量$F=(S_X^2/\sigma_X^2)/(S_Y^2/\sigma_Y^2)$服从自由度为$n-1$和$m-1$的Snedecor $F$分布.等价地，若随机变量$F$具有pdf
\begin{equation}\label{eq:5.3.9}
f_F(x)=\frac{\Gamma\left(\frac{p+q}2\right)}{\Gamma\left(\frac p2\right)\Gamma\left(\frac q2\right)}\left(\frac pq\right)^{p/2}\frac{x^{(p/2)-1}}{[1+(p/q)x]^{(p+q)/2}},\qquad 0<x<\infty,
\end{equation}
则它服从自由度为$p$和$q$的$F$分布.
\end{definition}
$F$分布可以在比这里更一般的条件下导出.即使母总体不是正态，方差比也可能服从$F$分布.\hyperref[bib:source188]{Kelker（1970）}证明，只要母总体具有某种对称性（球对称性，spherical symmetry），方差比就服从$F$分布.

从正态分布出发推导$F$的pdf，与推导Student的$t$分布类似.事实上，在一个特例中，$F$就是$t$的变换（见\autoref{thm:5.3.8}）.与$t$的情形类似，可将推导$F$的pdf化为求$(U/p)/(V/q)$的pdf，其中$U$、$V$独立，$U\sim\chi_p^2$、$V\sim\chi_q^2$（见\autoref{exercise:5.17}）.
\begin{example}[续\autoref{ex:5.3.5}]\label{ex:5.3.7}
为说明如何用$F$分布推断真实的总体方差比，考虑如下情形.$(S_X^2/\sigma_X^2)/(S_Y^2/\sigma_Y^2)$服从$F_{n-1,m-1}$分布.（一般用$F_{p,q}$表示自由度为$p$和$q$的$F$随机变量.）计算得
\begin{align*}
EF_{n-1,m-1}&=E\left(\frac{\chi_{n-1}^2/(n-1)}{\chi_{m-1}^2/(m-1)}\right)&&\text{（定义）}\\
&=E\left(\frac{\chi_{n-1}^2}{n-1}\right)E\left(\frac{m-1}{\chi_{m-1}^2}\right)&&\text{（独立性）}\\
&=\left(\frac{n-1}{n-1}\right)\left(\frac{m-1}{m-3}\right)&&\text{（卡方分布的计算）}\\
&=\frac{m-1}{m-3}.
\end{align*}
注意，最后这个表达式仅在$m>3$时为有限正数.有
$$E\left(\frac{S_X^2/\sigma_X^2}{S_Y^2/\sigma_Y^2}\right)=EF_{n-1,m-1}=\frac{m-1}{m-3},$$
去掉期望，当$m$足够大时，有
$$\frac{S_X^2/S_Y^2}{\sigma_X^2/\sigma_Y^2}\approx\frac{m-1}{m-3}\approx1,$$
正如所预期的那样.\footnote{译注：原文由期望接近1直接写出随机比值接近1，这只是启发式表述.仅令$m$增大而固定$n$，分子的随机波动一般不会消失；要论证该比值集中在1附近，还需控制两个样本方差的波动，例如同时增大两组样本容量.}
\end{example}
\sourcepageanchor{215}
$F$分布具有许多有趣性质，并与多种其他分布相关.下面的定理概括其中一些事实，其证明留作习题（见\autoref{exercise:5.17}和\autoref{exercise:5.18}）.
\begin{theorem}\label{thm:5.3.8}\index{SI0258@与F的关系}\index{SI0302@倒数}\index{SI0303@与贝塔的关系}\index{SI0304@与t的关系}\index{SI0383@与F的关系}\leavevmode
\begin{enumerate}[label=\textbf{\alph*.}]
\item 若$X\sim F_{p,q}$，则$1/X\sim F_{q,p}$；即$F$随机变量的倒数仍是$F$随机变量.
\item 若$X\sim t_q$，则$X^2\sim F_{1,q}$.
\item 若$X\sim F_{p,q}$，则$(p/q)X/(1+(p/q)X)\sim\operatorname{beta}(p/2,q/2)$.
\end{enumerate}
\end{theorem}



% 原书印刷184–188页，PDF物理215–219页；图见物理220页.
\subsection{次序统计量}\label{sec:5.4}
随机样本中的最小、最大或中间观测值等，可以提供额外的概要信息.例如，为未来紧急情况制定计划时，过去50年记录到的最高洪水水位或最低冬季气温可能是有用的数据.上个月出售房屋的价格中位数可能有助于估计生活成本.这些都是次序统计量（order statistics）的例子.
\begin{definition}\label{def:5.4.1}\index{SI0731@次序统计量}
随机样本$X_1,\ldots,X_n$的次序统计量，是将样本值按递增次序排列后所得的值，记为$X_{(1)},\ldots,X_{(n)}$.
\end{definition}
次序统计量是满足$X_{(1)}\le\cdots\le X_{(n)}$的随机变量.特别地，
$$\begin{aligned}
X_{(1)}&=\min_{1\le i\le n}X_i,\\
X_{(2)}&=\text{第二小的$X_i$},\\
&\ \vdots\\
X_{(n)}&=\max_{1\le i\le n}X_i.
\end{aligned}$$
既然它们是随机变量，就可讨论它们取各种值的概率.计算这些概率需要次序统计量的pdf或pmf.连续总体随机样本的次序统计量pdf公式是本节后面的主要内容，但先介绍几种易于用次序统计量定义的统计量.

样本极差（sample range）\index{SI0933@极差}\semanticindexanchor{933:215}$R=X_{(n)}-X_{(1)}$是最小观测值与最大观测值之间的距离.它衡量样本的离散程度，应反映总体的离散程度.

样本中位数（sample median）记为$M$，大约一半的观测小于$M$，另一半大于$M$.用次序统计量定义为
\begin{equation}\label{eq:5.4.1}\index{SI0930@中位数}
M=\begin{cases}X_{((n+1)/2)},&n\text{为奇数},\\[2pt]
\bigl(X_{(n/2)}+X_{(n/2+1)}\bigr)/2,&n\text{为偶数}.
\end{cases}
\end{equation}
中位数是一种位置度量，可以作为样本均值的替代量.样本中位数相对于样本均值的一个优点，是它受极端观测值的影响较小（详见\autoref{sec:10.2}）.

虽然均值与中位数有关，但它们通常衡量不同的东西.例如，最近一次棒球薪资谈判的主要争议点，是老板对球员养老金的缴款.老板的观点可以概括为：“棒球运动员的平均年薪为433,659美元；有这么多钱，现有养老金已经足够.”但球员的观点是：“超过一半的球员年收入不到250,000美元，而且多数球员的职业生涯短暂，需要较高养老金提供保障.”（这些数字属于1988赛季，而非发生争议的年份.）双方的数字都正确，但老板讨论的是均值，而球员讨论的是中位数.
\sourcepageanchor{216}
大约十来位年薪超过200万美元的球员，就能把平均薪资抬到433,659美元，而大多数球员收入不到250,000美元，包括所有年薪62,500美元的新秀.讨论工资、价格或任何有少数极端值的变量时，中位数比均值更能代表“典型”值.还有其他可由次序统计量定义、对极端值较不敏感的统计量，例如\autoref{exercise:10.20}中的$\alpha$截尾均值；\hyperref[bib:source318]{Tukey（1977）}等书对此有所讨论.

对0与1之间任意数$p$，第$100p$样本百分位数（sample percentile）\index{SI0754@百分位数}\semanticindexanchor{754:216}是这样一个观测值：大约$np$个观测值比它小，$n(1-p)$个比它大.第50样本百分位数（$p=.5$）就是样本中位数.对于其他$p$值，可按下述方式用次序统计量更精确地定义样本百分位数.
\begin{definition}\label{def:5.4.2}
出现在下标中的记号$\{b\}$，表示将$b$按通常方法舍入到最近的整数.更精确地，若$i$为整数且$i-.5\le b<i+.5$，则$\{b\}=i$.
\end{definition}
当$1/(2n)<p<.5$时，第$100p$样本百分位数是$X_{(\{np\})}$；当$.5<p<1-1/(2n)$时，则是$X_{(n+1-\{n(1-p)\})}$.例如，若$n=12$，要求第65百分位数，注意$12\times(1-.65)=4.2$，且$12+1-4=9$，所以它是$X_{(9)}$.对$p$的范围有限制，是因为样本容量限制了样本百分位数的范围.

将$p<.5$与$p>.5$两种情形分别定义，是为了让样本百分位数具有如下对称性：若第$100p$样本百分位数为第$i$小的观测值，则第$100(1-p)$样本百分位数应为第$i$大的观测值，上述定义正好满足这一点.例如，$n=11$时，第30样本百分位数为$X_{(3)}$，第70样本百分位数为$X_{(9)}$.

除了中位数，还常用两个样本百分位数：\index{SI0860@四分位数}\semanticindexanchor{860:216}下四分位数（lower quartile，第25百分位数）与上四分位数（upper quartile，第75百分位数）.有时用四分位距（interquartile range），即上、下四分位数之间的距离，衡量离散程度.

由于次序统计量是样本的函数，关于次序统计量的概率可由样本的概率计算.若$X_1,\ldots,X_n$为iid离散随机变量，计算次序统计量的概率主要是一个计数问题.\autoref{thm:5.4.3}推导了相应公式.若$X_1,\ldots,X_n$为来自连续总体的随机样本，\autoref{thm:5.4.4}和\autoref{thm:5.4.6}给出一个或多个次序统计量pdf的方便表达式，之后可用来推导次序统计量的函数的分布.
\begin{theorem}\label{thm:5.4.3}\index{SI0732@累积分布函数}\index{SI0736@概率质量函数}
设$X_1,\ldots,X_n$为来自pmf为$f_X(x_i)=p_i$的离散分布的随机样本，其中$x_1<x_2<\cdots$是$X$按递增顺序排列的可能取值.定义
$$\begin{aligned}P_0&=0,\\P_1&=p_1,\\P_2&=p_1+p_2,\\&\ \vdots\\P_i&=p_1+p_2+\cdots+p_i,\\&\ \vdots\end{aligned}$$
以$X_{(1)},\ldots,X_{(n)}$表示样本的次序统计量.则
\begin{equation}\label{eq:5.4.2}
P(X_{(j)}\le x_i)=\sum_{k=j}^n\binom nk P_i^k(1-P_i)^{n-k}
\end{equation}
并且
\begin{equation}\label{eq:5.4.3}
P(X_{(j)}=x_i)=\sum_{k=j}^n\binom nk\left[P_i^k(1-P_i)^{n-k}-P_{i-1}^k(1-P_{i-1})^{n-k}\right].
\end{equation}
\end{theorem}
\sourcepageanchor{217}
\begin{proof}
固定$i$，令随机变量$Y$表示$X_1,\ldots,X_n$中小于或等于$x_i$的个数.对于每个$X_j$，将事件$\{X_j\le x_i\}$称为“成功”，$\{X_j>x_i\}$称为“失败”.则$Y$是$n$次试验的成功次数.由于$X_1,\ldots,X_n$同分布，每次试验的成功概率相同，都是$P_i=P(X_j\le x_i)$.由于$X_j$与其他$X_i$独立，第$j$次试验成功还是失败与其他试验结果独立.故$Y\sim\operatorname{binomial}(n,P_i)$.

事件$\{X_{(j)}\le x_i\}$等价于$\{Y\ge j\}$，即至少$j$个样本值不超过$x_i$.\eqref{eq:5.4.2}表示这一二项概率：
$$P(X_{(j)}\le x_i)=P(Y\ge j).$$
\eqref{eq:5.4.3}只不过表示下面的差：
$$P(X_{(j)}=x_i)=P(X_{(j)}\le x_i)-P(X_{(j)}\le x_{i-1}).$$
$i=1$是一个例外，此时$P(X_{(j)}=x_1)=P(X_{(j)}\le x_1)$.定义$P_0=0$使\eqref{eq:5.4.3}也包含这个例外.
\end{proof}
若$X_1,\ldots,X_n$为来自连续总体的随机样本，情形会稍简单，因为任意两个$X_j$相等的概率为0，无须担心并列值.故$P(X_{(1)}<X_{(2)}<\cdots<X_{(n)})=1$，且$(X_{(1)},\ldots,X_{(n)})$的样本空间为$\{(x_1,\ldots,x_n):x_1<x_2<\cdots<x_n\}$.\autoref{thm:5.4.4}和\autoref{thm:5.4.6}再次利用二项分布论证，分别给出单个次序统计量的pdf和两个次序统计量的联合pdf.
\begin{theorem}\label{thm:5.4.4}\index{SI0735@概率密度函数}
设$X_{(1)},\ldots,X_{(n)}$是随机样本$X_1,\ldots,X_n$的次序统计量，总体连续且cdf为$F_X(x)$、pdf为$f_X(x)$.则$X_{(j)}$的pdf为
\begin{equation}\label{eq:5.4.4}
f_{X_{(j)}}(x)=\frac{n!}{(j-1)!(n-j)!}f_X(x)[F_X(x)]^{j-1}[1-F_X(x)]^{n-j}.
\end{equation}
\end{theorem}
\begin{proof}
先求$X_{(j)}$的cdf，再求导得到pdf.如\autoref{thm:5.4.3}，令$Y$表示$X_1,\ldots,X_n$中不超过$x$的个数.把事件$\{X_j\le x\}$定义为“成功”，可知$Y\sim\operatorname{binomial}(n,F_X(x))$.（注意，\autoref{thm:5.4.3}中可写$P_i=F_X(x_i)$.另外，虽然$X_1,\ldots,X_n$是连续随机变量，计数变量$Y$却是离散的.）因此
$$F_{X_{(j)}}(x)=P(Y\ge j)=\sum_{k=j}^n\binom nk[F_X(x)]^k[1-F_X(x)]^{n-k},$$
而$X_{(j)}$的pdf为
\begin{align*}
f_{X_{(j)}}(x)&=\frac d{dx}F_{X_{(j)}}(x)\\
&=\sum_{k=j}^n\binom nk\Bigl(k[F_X(x)]^{k-1}[1-F_X(x)]^{n-k}f_X(x)\\
&\hspace{38mm}-(n-k)[F_X(x)]^k[1-F_X(x)]^{n-k-1}f_X(x)\Bigr)\quad\text{（链式法则）}\\
&=\binom nj j f_X(x)[F_X(x)]^{j-1}[1-F_X(x)]^{n-j}\\
&\quad+\sum_{k=j+1}^n\binom nk k[F_X(x)]^{k-1}[1-F_X(x)]^{n-k}f_X(x)\\
&\quad-\sum_{k=j}^{n-1}\binom nk(n-k)[F_X(x)]^k[1-F_X(x)]^{n-k-1}f_X(x),
\end{align*}
最后一项求和中$k=n$的项为零.
\sourcepageanchor{218}
改变求和指标，继续得到
\begin{align*}
f_{X_{(j)}}(x)&=\frac{n!}{(j-1)!(n-j)!}f_X(x)[F_X(x)]^{j-1}[1-F_X(x)]^{n-j}\\
&\quad+\sum_{k=j}^{n-1}\binom n{k+1}(k+1)[F_X(x)]^k[1-F_X(x)]^{n-k-1}f_X(x)\\
&\quad-\sum_{k=j}^{n-1}\binom nk(n-k)[F_X(x)]^k[1-F_X(x)]^{n-k-1}f_X(x).
\end{align*}
注意
\begin{equation}\label{eq:5.4.5}
\binom n{k+1}(k+1)=\frac{n!}{k!(n-k-1)!}=\binom nk(n-k),
\end{equation}
可见上面最后两个和相消.因此pdf $f_{X_{(j)}}(x)$就是\eqref{eq:5.4.4}的表达式.
\end{proof}
\begin{example}[均匀分布次序统计量的pdf]\label{ex:5.4.5}\index{SI0389@次序统计量}
设$X_1,\ldots,X_n$为iid $\operatorname{uniform}(0,1)$，故$x\in(0,1)$时$f_X(x)=1$、$F_X(x)=x$.由\eqref{eq:5.4.4}，第$j$个次序统计量的pdf为
\begin{align*}
f_{X_{(j)}}(x)&=\frac{n!}{(j-1)!(n-j)!}x^{j-1}(1-x)^{n-j},\qquad x\in(0,1)\\
&=\frac{\Gamma(n+1)}{\Gamma(j)\Gamma(n-j+1)}x^{j-1}(1-x)^{(n-j+1)-1}.
\end{align*}
所以，$\operatorname{uniform}(0,1)$样本的第$j$个次序统计量服从$\operatorname{beta}(j,n-j+1)$分布.由此推得
$$EX_{(j)}=\frac j{n+1},\qquad\operatorname{Var}X_{(j)}=\frac{j(n-j+1)}{(n+1)^2(n+2)}.$$
\end{example}
两个或多个次序统计量的联合分布，可用于推导本节开头提到的某些统计量的分布.任意两个次序统计量的联合pdf由下述定理给出，证明留在\autoref{exercise:5.26}中.
\begin{theorem}\label{thm:5.4.6}\index{SI0734@联合概率密度函数}
设$X_{(1)},\ldots,X_{(n)}$是随机样本$X_1,\ldots,X_n$的次序统计量，总体连续且cdf为$F_X(x)$、pdf为$f_X(x)$.则$X_{(i)}$与$X_{(j)}$（$1\le i<j\le n$）的联合pdf为
\begin{equation}\label{eq:5.4.6}
\begin{aligned}
f_{X_{(i)},X_{(j)}}(u,v)&=\frac{n!}{(i-1)!(j-1-i)!(n-j)!}f_X(u)f_X(v)[F_X(u)]^{i-1}\\
&\quad\times[F_X(v)-F_X(u)]^{j-1-i}[1-F_X(v)]^{n-j},
\end{aligned}
\end{equation}
其中$-\infty<u<v<\infty$.
\end{theorem}
用类似但更繁复的论证，可以推导三个或更多次序统计量的联合pdf.另一个最有用的pdf或许是全体次序统计量的联合pdf $f_{X_{(1)},\ldots,X_{(n)}}(x_1,\ldots,x_n)$，它为
$$f_{X_{(1)},\ldots,X_{(n)}}(x_1,\ldots,x_n)=\begin{cases}
n!f_X(x_1)\cdots f_X(x_n),&-\infty<x_1<\cdots<x_n<\infty,\\
0,&\text{其他情形}.
\end{cases}$$
公式中自然出现$n!$，因为对于任意一组数值$x_1,\ldots,x_n$，将它们分配给$X_1,\ldots,X_n$共有$n!$种等可能方式，且都产生相同的次序统计量值.结合这个联合pdf与\autoref{sec:4}的方法，可以推导边际分布、条件分布以及次序统计量其他函数的分布（见\autoref{exercise:5.27}和\autoref{exercise:5.28}\mbox{）.}
\sourcepageanchor{219}
现在利用联合pdf \eqref{eq:5.4.6}推导本节开头提到的某些函数的分布.
\begin{example}[中程数与极差的分布]\label{ex:5.4.7}\index{SI0931@中程数}\index{SI0945@均匀}
设$X_1,\ldots,X_n$为iid $\operatorname{uniform}(0,a)$，次序统计量为$X_{(1)},\ldots,X_{(n)}$.前面已定义极差$R=X_{(n)}-X_{(1)}$.中程数（midrange）与样本中位数或均值一样，是位置度量，定义为$V=(X_{(1)}+X_{(n)})/2$.我们由$X_{(1)}$与$X_{(n)}$的联合pdf推导$R$与$V$的联合pdf.由\eqref{eq:5.4.6}有
\begin{align*}
f_{X_{(1)},X_{(n)}}(x_1,x_n)&=\frac{n(n-1)}{a^2}\left(\frac{x_n}a-\frac{x_1}a\right)^{n-2}\\
&=\frac{n(n-1)(x_n-x_1)^{n-2}}{a^n},\qquad0<x_1<x_n<a.
\end{align*}
解出$X_{(1)}$和$X_{(n)}$，得$X_{(1)}=V-R/2$、$X_{(n)}=V+R/2$.此变换的Jacobian为$-1$.从$(X_{(1)},X_{(n)})$到$(R,V)$的变换把$\{(x_1,x_n):0<x_1<x_n<a\}$映到$\{(r,v):0<r<a,r/2<v<a-r/2\}$.这是因为显然$0<r<a$，而固定$r$时，$v$从$r/2$（对应$x_1=0,x_n=r$）到$a-r/2$（对应$x_1=a-r,x_n=a$）变化.因此$(R,V)$的联合pdf为
$$f_{R,V}(r,v)=\frac{n(n-1)r^{n-2}}{a^n},\qquad0<r<a,\quad r/2<v<a-r/2.$$
于是$R$的边际pdf为
\begin{equation}\label{eq:5.4.7}
\begin{aligned}
f_R(r)&=\int_{r/2}^{a-r/2}\frac{n(n-1)r^{n-2}}{a^n}\,dv\\
&=\frac{n(n-1)r^{n-2}(a-r)}{a^n},\qquad0<r<a.
\end{aligned}
\end{equation}
当$a=1$时，$r$服从$\operatorname{beta}(n-1,2)$分布.或者，对任意$a$，容易从\eqref{eq:5.4.7}推出$R/a$服从beta分布.注意常数$a$是尺度参数.

$f_{R,V}(r,v)>0$的区域如\autoref{fig:5.4.1}所示，可见$r$的积分范围取决于$v>a/2$还是$v\le a/2$.因此$V$的边际pdf为
$$f_V(v)=\int_0^{2v}\frac{n(n-1)r^{n-2}}{a^n}\,dr=\frac{n(2v)^{n-1}}{a^n},\qquad0<v\le a/2,$$
以及
$$f_V(v)=\int_0^{2(a-v)}\frac{n(n-1)r^{n-2}}{a^n}\,dr=\frac{n[2(a-v)]^{n-1}}{a^n},\qquad a/2<v\le a.$$
这个pdf关于$a/2$对称，并在$a/2$处达到峰值.
\end{example}
\begin{figure}[htbp]\centering
\includegraphics[width=.62\textwidth]{assets/chapter5/fig-5-4-1.pdf}
\caption{\autoref{ex:5.4.7}中$f_{R,V}(r,v)>0$的区域}\label{fig:5.4.1}
\end{figure}



% 原书印刷188–200页，PDF物理219–231页.
\subsection{收敛概念}\label{sec:5.5}
本节考虑让样本容量趋于无穷这一有些理想化的想法，并研究在此过程中某些样本量的行为.虽然无限样本容量只是理论构造，但它往往能为有限样本情形提供有用的近似，因为表达式通常会在极限中简化.

我们主要关注三种收敛，讨论的详略有所不同.（例如，\hyperref[bib:source33]{Billingsley 1995}或\hyperref[bib:source251]{Resnick 1999}对收敛有完整论述.）特别地，我们想考察$n\to\infty$时$n$个观测的均值$\bar X_n$的行为.
\sourcepageanchor{220}
\subsubsection{依概率收敛}\label{sec:5.5.1}
这种收敛是较弱的收敛之一，因此通常容易验证.
\begin{definition}\label{def:5.5.1}\index{SI0206@依概率}
若对每个$\varepsilon>0$，都有
$$\lim_{n\to\infty}P(|X_n-X|\ge\varepsilon)=0,\quad\text{或等价地}\quad\lim_{n\to\infty}P(|X_n-X|<\varepsilon)=1,$$
则称随机变量序列$X_1,X_2,\ldots$依概率收敛（converges in probability）到随机变量$X$.
\end{definition}
\autoref{def:5.5.1}（以及本节其他定义）中的$X_1,X_2,\ldots$通常不像随机样本那样是独立同分布的随机变量.$X_n$的分布随下标变化；本节的收敛概念描述了下标增大时，$X_n$的分布以不同方式趋向某个极限分布.

统计学家经常关心这样的情形：极限随机变量是常数，而序列中的随机变量是某种样本均值.这类结果中最著名的是以下定理.
\begin{theorem}[弱大数定律]\label{thm:5.5.2}\index{SI0210@弱大数定律}\index{SI0565@切比雪夫}\index{SI0605@弱}
设$X_1,X_2,\ldots$为iid随机变量，$EX_i=\mu$、$\operatorname{Var}X_i=\sigma^2<\infty$.定义$\bar X_n=(1/n)\sum_{i=1}^nX_i$.则对每个$\varepsilon>0$，
$$\lim_{n\to\infty}P(|\bar X_n-\mu|<\varepsilon)=1,$$
即$\bar X_n$依概率收敛到$\mu$.
\end{theorem}
\begin{proof}
证明很简单，只需直接应用Chebychev不等式.对每个$\varepsilon>0$，
$$P(|\bar X_n-\mu|\ge\varepsilon)=P((\bar X_n-\mu)^2\ge\varepsilon^2)\le\frac{E(\bar X_n-\mu)^2}{\varepsilon^2}=\frac{\operatorname{Var}\bar X}{\varepsilon^2}=\frac{\sigma^2}{n\varepsilon^2}.$$
故当$n\to\infty$时，$P(|\bar X_n-\mu|<\varepsilon)=1-P(|\bar X_n-\mu|\ge\varepsilon)\ge1-\sigma^2/(n\varepsilon^2)\to1$.
\end{proof}
\sourcepageanchor{221}
弱大数定律（Weak Law of Large Numbers，WLLN）简洁地说明，在一般条件下，$n\to\infty$时样本均值趋近总体均值.事实上还有更一般的弱大数定律，只需假设均值有限.不过，\autoref{thm:5.5.2}的版本已适用于多数实际情形.

弱大数定律所概括的性质，即“同一个”样本量构成的序列在$n\to\infty$时趋向一个常数，称为相合性（consistency）.\autoref{sec:7}会更仔细地讨论这一性质.
\begin{example}[$S^2$的相合性]\label{ex:5.5.3}\index{SI0196@相合性（见点估计量的相合性）}
设有iid随机变量序列$X_1,X_2,\ldots$，满足$EX_i=\mu$、$\operatorname{Var}X_i=\sigma^2<\infty$.若定义
$$S_n^2=\frac1{n-1}\sum_{i=1}^n(X_i-\bar X_n)^2,$$
能否证明$S_n^2$的弱大数定律？由Chebychev不等式，
$$P(|S_n^2-\sigma^2|\ge\varepsilon)\le\frac{E(S_n^2-\sigma^2)^2}{\varepsilon^2}=\frac{\operatorname{Var}S_n^2}{\varepsilon^2},$$
所以$S_n^2$依概率收敛到$\sigma^2$的一个充分条件是$n\to\infty$时$\operatorname{Var}S_n^2\to0$.
\end{example}
\autoref{def:5.5.1}的一个自然推广涉及随机变量的函数.即若序列$X_1,X_2,\ldots$依概率收敛到随机变量$X$或常数$a$，对于某个表现较好的函数$h$，能否对随机变量序列$h(X_1),h(X_2),\ldots$作出结论？下一个定理表明可以（证明见\autoref{exercise:5.39}）.
\begin{theorem}\label{thm:5.5.4}
设$X_1,X_2,\ldots$依概率收敛到随机变量$X$，且$h$是连续函数，则$h(X_1),h(X_2),\ldots$依概率收敛到$h(X)$.
\end{theorem}
\begin{example}[$S$的相合性]\label{ex:5.5.5}
若$S_n^2$是$\sigma^2$的相合估计量，则由\autoref{thm:5.5.4}，样本标准差$S_n=\sqrt{S_n^2}=h(S_n^2)$是$\sigma$的相合估计量.注意$S_n$实际上是$\sigma$的有偏估计量（\autoref{exercise:5.11}），但偏差渐近消失.
\end{example}
\subsubsection{几乎处处收敛}\label{sec:5.5.2}
比依概率收敛更强的一种收敛是几乎处处收敛（almost sure convergence；有时也称以概率1收敛，这一名称容易引起混淆）.它类似于函数列的逐点收敛，只是不要求在概率为0的集合上收敛（故称“几乎”处处）.
\begin{definition}\label{def:5.5.6}\index{SI0201@几乎处处}\index{SI0211@以概率1}
若对每个$\varepsilon>0$，都有
$$P\left(\lim_{n\to\infty}|X_n-X|<\varepsilon\right)=1,$$
则称随机变量序列$X_1,X_2,\ldots$几乎处处收敛到随机变量$X$.
\end{definition}
注意\autoref{def:5.5.1}与\autoref{def:5.5.6}的表述相似.虽然外观相近，它们却十分不同，后者强得多.为理解几乎处处收敛，须回顾\autoref{def:1.4.1}中随机变量的基本定义：随机变量是定义在样本空间$S$上的实值函数.若$S$的元素记为$s$，则$X_n(s)$与$X(s)$都是$S$上的函数.\autoref{def:5.5.6}说明，若除可能属于某个$N\subset S$、$P(N)=0$的$s$以外，对所有$s\in S$都有$X_n(s)\to X(s)$，则$X_n$几乎处处收敛到$X$.\autoref{ex:5.5.7}说明几乎处处收敛，\autoref{ex:5.5.8}说明它与依概率收敛的差别.
\sourcepageanchor{222}
\begin{example}[几乎处处收敛]\label{ex:5.5.7}
设样本空间$S$为具有均匀概率分布的闭区间$[0,1]$.定义随机变量$X_n(s)=s+s^n$及$X(s)=s$.对每个$s\in[0,1)$，$n\to\infty$时$s^n\to0$，故$X_n(s)\to s=X(s)$.但对每个$n$都有$X_n(1)=2$，所以$X_n(1)$不收敛到$1=X(1)$.由于收敛发生在集合$[0,1)$上，且$P([0,1))=1$，$X_n$几乎处处收敛到$X$.
\end{example}
\begin{example}[依概率收敛但不几乎处处收敛]\label{ex:5.5.8}
本例给出一个依概率收敛但不几乎处处收敛的序列.仍设$S$为具有均匀概率分布的闭区间$[0,1]$.如下定义序列$X_1,X_2,\ldots$：
$$\begin{array}{lll}
X_1(s)=s+I_{[0,1]}(s),&X_2(s)=s+I_{[0,1/2]}(s),&X_4(s)=s+I_{[0,1/3]}(s),\\
&X_3(s)=s+I_{[1/2,1]}(s),&X_5(s)=s+I_{[1/3,2/3]}(s),\\
&&X_6(s)=s+I_{[2/3,1]}(s),
\end{array}$$
依此类推.令$X(s)=s$.容易看出$X_n$依概率收敛到$X$.$n\to\infty$时，$P(|X_n-X|\ge\varepsilon)$等于某个$s$值区间的概率，而该区间长度趋于0.但$X_n$不几乎处处收敛到$X$.事实上，不存在任何$s\in S$使$X_n(s)\to s=X(s)$.对每个$s$，$X_n(s)$在$s$与$s+1$之间无限次交替.例如，$s=3/8$时，$X_1(s)=1\frac38$、$X_2(s)=1\frac38$、$X_3(s)=\frac38$、$X_4(s)=\frac38$、$X_5(s)=1\frac38$、$X_6(s)=\frac38$，等等.此序列在任何一点都不收敛.
\end{example}
可以预料，几乎处处收敛作为更强的条件，蕴含依概率收敛.逆命题当然不成立，\autoref{ex:5.5.8}已说明这一点.然而，如果序列依概率收敛，就可以找到一个几乎处处收敛的子序列.（\hyperref[bib:source251]{Resnick 1999}，第6.3节详细讨论了这两种收敛的联系.）

统计学家仍常关注到常数的收敛.下面不加证明地给出弱大数定律的更强对应结果——强大数定律（Strong Law of Large Numbers，SLLN）.证明提纲见\autoref{sec:5.8.4}.
\begin{theorem}[强大数定律]\label{thm:5.5.9}\index{SI0209@强大数定律}\index{SI0604@强}
设$X_1,X_2,\ldots$为iid随机变量，$EX_i=\mu$、$\operatorname{Var}X_i=\sigma^2<\infty$，定义$\bar X_n=(1/n)\sum_{i=1}^nX_i$.则对每个$\varepsilon>0$，
$$P\left(\lim_{n\to\infty}|\bar X_n-\mu|<\varepsilon\right)=1,$$
即$\bar X_n$几乎处处收敛到$\mu$.
\end{theorem}
在弱大数定律与强大数定律中，都假设方差有限.虽然多数应用中这一假设成立（也很理想），但事实上它强于实际所需条件.去掉这一假设，弱定律与强定律仍然成立.唯一需要的矩条件是$E|X_i|<\infty$（见\hyperref[bib:source251]{Resnick 1999}，第7章；或\hyperref[bib:source33]{Billingsley 1995}，第22节\mbox{）.}
\subsubsection{依分布收敛}\label{sec:5.5.3}
\autoref{sec:2}已介绍依分布收敛（convergence in distribution）的概念.回顾矩母函数的性质，以及其收敛如何蕴含依分布收敛（\autoref{thm:2.3.12}）.
\begin{definition}\label{def:5.5.10}\index{SI0204@依分布}
若在$F_X(x)$连续的每一点$x$都有
$$\lim_{n\to\infty}F_{X_n}(x)=F_X(x),$$
则称随机变量序列$X_1,X_2,\ldots$依分布收敛到随机变量$X$.
\end{definition}
\sourcepageanchor{223}
\begin{example}[均匀随机变量的最大值]\label{ex:5.5.11}\index{SI0388@最大值}
若$X_1,X_2,\ldots$为iid $\operatorname{uniform}(0,1)$，$X_{(n)}=\max_{1\le i\le n}X_i$，考察$X_{(n)}$是否依分布收敛，以及收敛到哪里.

$n\to\infty$时，我们预期$X_{(n)}$接近1.由于$X_{(n)}$必小于1，对任何$\varepsilon>0$，
\begin{align*}
P(|X_{(n)}-1|\ge\varepsilon)&=P(X_{(n)}\ge1+\varepsilon)+P(X_{(n)}\le1-\varepsilon)\\
&=0+P(X_{(n)}\le1-\varepsilon).
\end{align*}
再利用样本iid，有
$$P(X_{(n)}\le1-\varepsilon)=P(X_i\le1-\varepsilon,\ i=1,\ldots,n)=(1-\varepsilon)^n,$$
它趋于0.所以证明了$X_{(n)}$依概率收敛到1.不过，若取$\varepsilon=t/n$，则
$$P(X_{(n)}\le1-t/n)=(1-t/n)^n\to e^{-t},$$
整理得到
$$P(n(1-X_{(n)})\le t)\to1-e^{-t},$$
即随机变量$n(1-X_{(n)})$依分布收敛到$\operatorname{exponential}(1)$随机变量.
\end{example}
注意，虽然说的是随机变量序列依分布收敛，实际收敛的是cdf，而不是随机变量本身.在这一根本意义上，依分布收敛与依概率收敛、几乎处处收敛都很不同.不过，其他两种收敛都蕴含它.
\begin{theorem}\label{thm:5.5.12}\index{SI0205@与依概率收敛的关系}\index{SI0207@与依分布收敛的关系}
若随机变量序列$X_1,X_2,\ldots$依概率收敛到随机变量$X$，则该序列也依分布收敛到$X$.
\end{theorem}
证明见\autoref{exercise:5.40}.另外，由\autoref{sec:5.5.2}可知，几乎处处收敛也蕴含依分布收敛.

在一个特殊情形中，\autoref{thm:5.5.12}有一个有用的逆命题.例示见\autoref{ex:10.1.13}，证明见\autoref{exercise:5.41}.
\begin{theorem}\label{thm:5.5.13}\index{SI0208@收敛到常数}
随机变量序列$X_1,X_2,\ldots$依概率收敛到常数$\mu$，当且仅当它依分布收敛到$\mu$.即
$$P(|X_n-\mu|>\varepsilon)\to0\quad\text{对每个$\varepsilon>0$}$$
等价于
$$P(X_n\le x)\to\begin{cases}0,&x<\mu,\\1,&x>\mu.\end{cases}$$
\end{theorem}
样本均值是大样本行为极其重要的统计量之一.特别地，我们想研究它的极限分布.统计学中最令人惊叹的定理之一——中心极限定理（Central Limit Theorem，CLT）——概括了这一结果.
\begin{theorem}[中心极限定理]\label{thm:5.5.14}\index{SI0102@中心极限定理}
设$X_1,X_2,\ldots$为iid随机变量序列，其矩母函数在0的某个邻域内存在（即对某个$h>0$，$|t|<h$时$M_{X_i}(t)$存在）.设$EX_i=\mu$、$\operatorname{Var}X_i=\sigma^2>0$.（矩母函数存在，故$\mu$与$\sigma^2$都有限.）定义$\bar X_n=(1/n)\sum_{i=1}^nX_i$，以$G_n(x)$表示$\sqrt n(\bar X_n-\mu)/\sigma$的cdf.则对任意$-\infty<x<\infty$，
$$\lim_{n\to\infty}G_n(x)=\int_{-\infty}^x\frac1{\sqrt{2\pi}}e^{-y^2/2}\,dy,$$
即$\sqrt n(\bar X_n-\mu)/\sigma$具有标准正态极限分布.
\end{theorem}
\sourcepageanchor{224}
在证明此定理之前（证明颇有些平淡），先考察它的含义.几乎不作假设（除了独立性和方差有限），最终竟得到正态性！关键在于，正态性来自“小的”（方差有限的）独立扰动之和.方差有限的假设对收敛到正态分布基本上是必要的.虽然可以略加放宽，却不能完全取消.（回顾讨论Cauchy分布的\autoref{ex:5.2.10}；那里没有到正态分布的收敛.）

在赞叹中心极限定理的同时，也应思考其局限.它虽然给出有用的一般近似，但通常无法自动知道这种近似究竟有多好.事实上，近似精度依赖原始分布，必须逐个情形检查.此外，当今计算能力充裕而廉价，使中心极限定理这类近似的重要性有所降低.尽管如此，它仍是一个非凡的结果.
\begin{proof}[\autoref{thm:5.5.14}的证明]
下面证明，对$|t|<h$，$\sqrt n(\bar X_n-\mu)/\sigma$的矩母函数收敛到$n(0,1)$随机变量的矩母函数$e^{t^2/2}$.

定义$Y_i=(X_i-\mu)/\sigma$，以$M_Y(t)$表示各$Y_i$共同的矩母函数，它在$|t|<\sigma h$时存在，可由\autoref{thm:2.3.15}给出.由于
\begin{equation}\label{eq:5.5.1}
\frac{\sqrt n(\bar X_n-\mu)}\sigma=\frac1{\sqrt n}\sum_{i=1}^nY_i,
\end{equation}
由矩母函数的性质（见\autoref{thm:2.3.15}和\autoref{thm:4.6.7}），有
\begin{equation}\label{eq:5.5.2}
\begin{aligned}
M_{\sqrt n(\bar X_n-\mu)/\sigma}(t)&=M_{\sum_{i=1}^nY_i/\sqrt n}(t)\\
&=M_{\sum_{i=1}^nY_i}\left(\frac t{\sqrt n}\right)&&\text{（\autoref{thm:2.3.15}）}\\
&=\left(M_Y\left(\frac t{\sqrt n}\right)\right)^n.&&\text{（\autoref{thm:4.6.7}）}
\end{aligned}
\end{equation}
把$M_Y(t/\sqrt n)$在0附近展开为Taylor级数\index{SI0696@在中心极限定理中的用途}\semanticindexanchor{696:224}\index{SI1031@在中心极限定理中的用途}\semanticindexanchor{1031:224}（幂级数；见\autoref{def:5.5.20}），得到
\begin{equation}\label{eq:5.5.3}
M_Y\left(\frac t{\sqrt n}\right)=\sum_{k=0}^\infty M_Y^{(k)}(0)\frac{(t/\sqrt n)^k}{k!},
\end{equation}
其中$M_Y^{(k)}(0)=(d^k/dt^k)M_Y(t)|_{t=0}$.由于矩母函数在$|t|<h$时存在，当$t<\sqrt n\sigma h$时幂级数展开有效.\footnote{译注：此处原文写$t<\sqrt n\sigma h$，按前文矩母函数存在的双侧邻域，应理解为$|t|<\sqrt n\sigma h$.原式保留；对固定$t$取充分大的$n$，证明仍适用.}

利用$M_Y^{(0)}=1$、$M_Y^{(1)}=0$和$M_Y^{(2)}=1$（按构造，$Y$的均值和方差为0和1），得到
\begin{equation}\label{eq:5.5.4}
M_Y\left(\frac t{\sqrt n}\right)=1+\frac{(t/\sqrt n)^2}{2!}+R_Y\left(\frac t{\sqrt n}\right),
\end{equation}
其中$R_Y$是Taylor展开的余项，
$$R_Y\left(\frac t{\sqrt n}\right)=\sum_{k=3}^\infty M_Y^{(k)}(0)\frac{(t/\sqrt n)^k}{k!}.$$
由Taylor定理（\autoref{thm:5.5.21}），对固定的$t\ne0$，有
$$\lim_{n\to\infty}\frac{R_Y(t/\sqrt n)}{(t/\sqrt n)^2}=0.$$
由于$t$固定，也有
\begin{equation}\label{eq:5.5.5}
\lim_{n\to\infty}\frac{R_Y(t/\sqrt n)}{(1/\sqrt n)^2}=\lim_{n\to\infty}nR_Y\left(\frac t{\sqrt n}\right)=0.
\end{equation}
\sourcepageanchor{225}
当$t=0$时，由于$R_Y(0/\sqrt n)=0$，\eqref{eq:5.5.5}也成立.因此对任意固定$t$，可以写
\begin{equation}\label{eq:5.5.6}
\begin{aligned}
\lim_{n\to\infty}\left(M_Y\left(\frac t{\sqrt n}\right)\right)^n
&=\lim_{n\to\infty}\left[1+\frac{(t/\sqrt n)^2}{2!}+R_Y\left(\frac t{\sqrt n}\right)\right]^n\\
&=\lim_{n\to\infty}\left[1+\frac1n\left(\frac{t^2}2+nR_Y\left(\frac t{\sqrt n}\right)\right)\right]^n\\
&=e^{t^2/2},
\end{aligned}
\end{equation}
这里应用了\autoref{lem:2.3.14}，取$a_n=(t^2/2)+nR_Y(t/\sqrt n)$.（注意，\eqref{eq:5.5.5}意味着$n\to\infty$时$a_n\to t^2/2$.）由于$e^{t^2/2}$是$n(0,1)$分布的矩母函数，定理得证.
\end{proof}
中心极限定理在比\autoref{thm:5.5.14}更一般的条件下仍成立（\autoref{sec:5.8.1}）.特别地，所有关于矩母函数的假设都不需要，可以改用特征函数（\autoref{sec:2.6.2}）.下面的定理不给证明.这是几乎足以满足全部统计用途的中心极限定理版本.注意对母分布的唯一假设是方差有限.
\begin{theorem}[更强形式的中心极限定理]\label{thm:5.5.15}\index{SI0102@中心极限定理}\index{SI0105@在中心极限定理中的用途}
设$X_1,X_2,\ldots$为iid随机变量序列，$EX_i=\mu$，$0<\operatorname{Var}X_i=\sigma^2<\infty$.定义$\bar X_n=(1/n)\sum_{i=1}^nX_i$，以$G_n(x)$表示$\sqrt n(\bar X_n-\mu)/\sigma$的cdf.则对任意$-\infty<x<\infty$，
$$\lim_{n\to\infty}G_n(x)=\int_{-\infty}^x\frac1{\sqrt{2\pi}}e^{-y^2/2}\,dy,$$
即$\sqrt n(\bar X_n-\mu)/\sigma$具有标准正态极限分布.
\end{theorem}
证明与\autoref{thm:5.5.14}几乎相同，只是用特征函数代替矩母函数.分布的特征函数总存在，所以无须在定理假设中提及.不过证明更细致，因为须处理复变量函数.详见\hyperref[bib:source33]{Billingsley（1995}，第27节）.

中心极限定理提供一种通用近似（但请记住关于近似精度的提醒）.实践中，总可先用它进行粗略计算.
\begin{example}[负二项分布的正态近似]\label{ex:5.5.16}\index{SI0063@以正态近似负二项}
设$X_1,\ldots,X_n$为来自$\operatorname{negative\ binomial}(r,p)$分布的随机样本.回顾
$$EX=\frac{r(1-p)}p,\qquad\operatorname{Var}X=\frac{r(1-p)}{p^2},$$
中心极限定理告诉我们
$$\frac{\sqrt n(\bar X-r(1-p)/p)}{\sqrt{r(1-p)/p^2}}$$
近似服从$n(0,1)$.近似概率计算比精确计算容易得多.例如$r=10$、$p=1/2$、$n=30$时，精确计算为
\begin{align*}
P(\bar X\le11)&=P\left(\sum_{i=1}^{30}X_i\le330\right)\\
&=\sum_{x=0}^{330}\binom{300+x-1}x\left(\frac12\right)^{300}\left(\frac12\right)^x
&&\text{（$\sum X_i$服从负二项$(nr,p)$分布）}\\
&=.8916,
\end{align*}
\sourcepageanchor{226}
这个计算非常困难.（即使借助计算机也不容易，阶乘的巨大数值会带来很大困难.不信可以试试！）中心极限定理给出近似
\begin{align*}
P(\bar X\le11)&=P\left(\frac{\sqrt{30}(\bar X-10)}{\sqrt{20}}\le\frac{\sqrt{30}(11-10)}{\sqrt{20}}\right)\\
&\approx P(Z\le1.2247)=.8888.
\end{align*}
进一步的改进见\autoref{exercise:5.37}.
\end{example}
一种可结合中心极限定理使用的近似工具称为Slutsky定理.
\begin{theorem}[Slutsky定理]\label{thm:5.5.17}\index{SI0204@依分布}\index{SI0979@Slutsky定理}
若$X_n$依分布收敛到$X$，而$Y_n$依概率收敛到常数$a$，则
\begin{enumerate}[label=\textbf{\alph*.}]
\item $Y_nX_n$依分布收敛到$aX$.
\item $X_n+Y_n$依分布收敛到$X+a$.
\end{enumerate}
\end{theorem}
Slutsky定理的证明略去，因为它依赖于尚未讨论的依分布收敛的一种刻画.下面的例子说明一种典型应用.
\begin{example}[以估计方差作正态近似]\label{ex:5.5.18}
设
$$\frac{\sqrt n(\bar X_n-\mu)}\sigma\to n(0,1),$$
但$\sigma$未知.\autoref{ex:5.5.3}表明，若$\lim_{n\to\infty}\operatorname{Var}S_n^2=0$，则$S_n^2$依概率收敛到$\sigma^2$.由\autoref{exercise:5.32}，$\sigma/S_n$依概率收敛到1.于是Slutsky定理告诉我们
$$\frac{\sqrt n(\bar X_n-\mu)}{S_n}=\frac\sigma{S_n}\frac{\sqrt n(\bar X_n-\mu)}\sigma\to n(0,1).$$
\end{example}
\subsubsection{Delta方法}\label{sec:5.5.4}\index{SI0246@Delta方法}
上一小节给出了标准化随机变量具有正态极限分布的条件.不过很多时候，我们关心的并非随机变量本身的分布，而是其某个函数的分布.
\begin{example}[优势的估计]\label{ex:5.5.19}
设观测到独立的$\operatorname{Bernoulli}(p)$随机变量$X_1,X_2,\ldots,X_n$.通常关注的参数是成功概率$p$，另一个常见参数是优势（odds）$p/(1-p)$.例如，若数据代表某种医疗方法的治疗结果，且$p=2/3$，则病情好转的优势为$2:1$.此外，若另一种疗法的成功概率为$r$，生物统计学家常估计优势比（odds ratio）$\frac p{1-p}/\frac r{1-r}$，它表示一种疗法相对于另一种疗法的优势.

通常用观测成功比例$\hat p=\sum_iX_i/n$估计成功概率$p$，所以可考虑用$\hat p/(1-\hat p)$估计$p/(1-p)$.但这一估计量有什么性质？怎样估计$\hat p/(1-\hat p)$的方差？又怎样近似它的抽样分布？

直觉不再帮得上忙，而精确计算也几乎没有希望，只能依赖近似.Delta方法将使我们得到合理的近似答案.
\end{example}
\sourcepageanchor{227}
一种方法是利用Taylor级数近似，它能近似随机变量的函数的均值与方差.还将看到，这些相当直接的近似足以得到一个中心极限定理.先简要回顾Taylor级数.
\begin{definition}\label{def:5.5.20}\index{SI1029@泰勒级数}
若函数$g(x)$有$r$阶导数，即$g^{(r)}(x)=\frac{d^r}{dx^r}g(x)$存在，则对任意常数$a$，在$a$处的$r$阶Taylor多项式（Taylor polynomial）为
$$T_r(x)=\sum_{i=0}^r\frac{g^{(i)}(a)}{i!}(x-a)^i.$$
\end{definition}
Taylor的主要定理说明，近似的余项$g(x)-T_r(x)$趋于零总比显式写出的最高阶项更快.这里不证明.
\begin{theorem}[Taylor]\label{thm:5.5.21}\index{SI1030@泰勒定理}
若$g^{(r)}(a)=\frac{d^r}{dx^r}g(x)|_{x=a}$存在，则
$$\lim_{x\to a}\frac{g(x)-T_r(x)}{(x-a)^r}=0.$$
\end{theorem}
一般来说，我们不关心余项的显式形式.由于关注近似，直接忽略余项即可.不过余项有很多显式形式，其中一个有用的形式为
$$g(x)-T_r(x)=\int_a^x\frac{g^{(r+1)}(t)}{r!}(x-t)^r\,dt.$$
在Taylor定理的统计应用中，我们最关心一阶Taylor级数，即只利用一阶导数的近似（上述公式取$r=1$）.另外还会用到多元Taylor级数.上面只详述了一元情形，因此下面有些结果须暂且接受.

设$T_1,\ldots,T_k$是均值分别为$\theta_1,\ldots,\theta_k$的随机变量，定义$\boldsymbol T=(T_1,\ldots,T_k)$、$\boldsymbol\theta=(\theta_1,\ldots,\theta_k)$.设有可微函数$g(\boldsymbol T)$（某参数的估计量），我们想近似估计它的方差.定义
$$g'_i(\boldsymbol\theta)=\frac\partial{\partial t_i}g(\boldsymbol t)\bigg|_{t_1=\theta_1,\ldots,t_k=\theta_k}.$$
$g$在$\boldsymbol\theta$处的一阶Taylor展开为
$$g(\boldsymbol t)=g(\boldsymbol\theta)+\sum_{i=1}^k g'_i(\boldsymbol\theta)(t_i-\theta_i)+\text{余项}.$$
在统计近似中忽略余项，写成
\begin{equation}\label{eq:5.5.7}
 g(\boldsymbol t)\approx g(\boldsymbol\theta)+\sum_{i=1}^k g'_i(\boldsymbol\theta)(t_i-\theta_i).
\end{equation}
对\eqref{eq:5.5.7}两边取期望，得到
\begin{equation}\label{eq:5.5.8}
\begin{aligned}
E_{\boldsymbol\theta}g(\boldsymbol T)&\approx g(\boldsymbol\theta)+\sum_{i=1}^k g'_i(\boldsymbol\theta)E_{\boldsymbol\theta}(T_i-\theta_i)\\
&=g(\boldsymbol\theta).\qquad\text{（$T_i$的均值为$\theta_i$）}
\end{aligned}
\end{equation}
\sourcepageanchor{228}
现在可以近似$g(\boldsymbol T)$的方差：
\begin{equation}\label{eq:5.5.9}\index{SI0764@近似}\index{SI1069@泰勒近似}
\begin{aligned}
\operatorname{Var}_{\boldsymbol\theta}g(\boldsymbol T)
&\approx E_{\boldsymbol\theta}([g(\boldsymbol T)-g(\boldsymbol\theta)]^2)&&\text{（由\eqref{eq:5.5.8}）}\\
&\approx E_{\boldsymbol\theta}\left(\left(\sum_{i=1}^k g'_i(\boldsymbol\theta)(T_i-\theta_i)\right)^2\right)&&\text{（由\eqref{eq:5.5.7}）}\\
&=\sum_{i=1}^k[g'_i(\boldsymbol\theta)]^2\operatorname{Var}_{\boldsymbol\theta}T_i
+2\sum_{i>j}g'_i(\boldsymbol\theta)g'_j(\boldsymbol\theta)\operatorname{Cov}_{\boldsymbol\theta}(T_i,T_j),
\end{aligned}
\end{equation}
最后的等号由展开平方、利用方差和协方差的定义得到（类似\autoref{exercise:4.44}）.近似\eqref{eq:5.5.9}十分有用，因为仅用简单的方差和协方差就给出一般函数的方差公式.下面给出两个例子.
\begin{example}[续\autoref{ex:5.5.19}]\label{ex:5.5.22}
回顾我们关心$\hat p/(1-\hat p)$作为$p/(1-p)$估计量的性质，其中$p$是二项成功概率.沿用上面的记号，取$g(p)=p/(1-p)$，则$g'(p)=1/(1-p)^2$，且
$$\operatorname{Var}\left(\frac{\hat p}{1-\hat p}\right)\approx[g'(p)]^2\operatorname{Var}(\hat p)=\left[\frac1{(1-p)^2}\right]^2\frac{p(1-p)}n=\frac p{n(1-p)^3},$$
由此得到估计量方差的近似.
\end{example}
\begin{example}[近似均值与方差]\label{ex:5.5.23}
设随机变量$X$满足$E_\mu X=\mu\ne0$.若要估计函数$g(\mu)$，一阶近似给出
$$g(X)=g(\mu)+g'(\mu)(X-\mu).$$
若用$g(X)$估计$g(\mu)$，则近似有
$$\begin{aligned}E_\mu g(X)&\approx g(\mu),\\\operatorname{Var}_\mu g(X)&\approx[g'(\mu)]^2\operatorname{Var}_\mu X.\end{aligned}$$
具体地，取$g(\mu)=1/\mu$，用$1/X$估计$1/\mu$，便有
$$\begin{aligned}E_\mu\left(\frac1X\right)&\approx\frac1\mu,\\
\operatorname{Var}_\mu\left(\frac1X\right)&\approx\left(\frac1\mu\right)^4\operatorname{Var}_\mu X.\end{aligned}$$
\end{example}
利用这些均值、方差的Taylor级数近似，可得到下面这个称为Delta方法的中心极限定理推广.
\begin{theorem}[Delta方法]\label{thm:5.5.24}
设随机变量序列$Y_n$满足$\sqrt n(Y_n-\theta)\to n(0,\sigma^2)$（依分布）.给定函数$g$及特定的$\theta$，假设$g'(\theta)$存在且非零.则
\begin{equation}\label{eq:5.5.10}
\sqrt n[g(Y_n)-g(\theta)]\to n(0,\sigma^2[g'(\theta)]^2)\quad\text{（依分布）}.
\end{equation}
\end{theorem}
\begin{proof}
$g(Y_n)$在$Y_n=\theta$附近的Taylor展开为
\begin{equation}\label{eq:5.5.11}
g(Y_n)=g(\theta)+g'(\theta)(Y_n-\theta)+\text{余项},
\end{equation}
其中$Y_n\to\theta$时余项趋于0.由于$Y_n$依概率收敛到$\theta$，余项也依概率收敛到0.将Slutsky定理（\autoref{thm:5.5.17}）用于
$$\sqrt n[g(Y_n)-g(\theta)]=g'(\theta)\sqrt n(Y_n-\theta)$$
即得结论.详见\autoref{exercise:5.43}.\footnote{译注：原文最后一式省去了余项，直接写成等式.严谨论证须说明乘以$\sqrt n$后的余项依概率趋于0；仅有未缩放余项趋于0还不够.可微性给出余项为$(Y_n-\theta)o(1)$，结合$\sqrt n(Y_n-\theta)$依分布收敛，便可完成这一步.}
\end{proof}
\sourcepageanchor{229}
\begin{example}[续\autoref{ex:5.5.23}]\label{ex:5.5.25}
现在设有随机样本的均值$\bar X$.当$\mu\ne0$时，
$$\sqrt n\left(\frac1{\bar X}-\frac1\mu\right)\to n\left(0,\left(\frac1\mu\right)^4\operatorname{Var}_\mu X_1\right)\quad\text{（依分布）}.$$
若不知道$X_1$的方差，要使用这个近似，就须用某个估计，例如$S^2$.此外，$1/\mu$项也有问题，因为$\mu$同样未知.可以把所有未知量都估计出来，得到近似方差\footnote{译注：原文这一式若指$1/\bar X$本身的方差，右端应再除以$n$；所写右端对应于上式经$\sqrt n$缩放后的渐近方差估计.随后标准化公式与后者一致.}
$$\operatorname{Var}\left(\frac1{\bar X}\right)\approx\left(\frac1{\bar X}\right)^4S^2.$$

此外，$\bar X$和$S^2$都是相合估计量，故可再次用Slutsky定理得出：当$\mu\ne0$时，
$$\frac{\sqrt n\left(\frac1{\bar X}-\frac1\mu\right)}{\left(\frac1{\bar X}\right)^2S}\to n(0,1)\quad\text{（依分布）}.$$
注意后一量的写法：除以估计的标准差，使极限分布成为标准正态分布.如果需要估计极限分布中的某些参数，这是唯一合理的方式.还应指出，当存在待估参数时另有一种方法，这里实际上可以避免在方差中使用$\mu$的估计（见\autoref{sec:10.3.2}的得分检验）.
\end{example}
为了完整讨论，还需要处理基本Delta方法的两个推广.第一个是$g'(\mu)=0$的可能性.例如，估计二项方差的方差时就可能出现这种情况（见\autoref{exercise:5.44}）.

若$g'(\theta)=0$，在Taylor展开中多取一项，得
$$g(Y_n)=g(\theta)+g'(\theta)(Y_n-\theta)+\frac{g''(\theta)}2(Y_n-\theta)^2+\text{余项}.$$
稍加整理（令$g'=0$），有
\begin{equation}\label{eq:5.5.12}
g(Y_n)-g(\theta)=\frac{g''(\theta)}2(Y_n-\theta)^2+\text{余项}.
\end{equation}
回顾$n(0,1)$的平方服从$\chi_1^2$（\autoref{ex:2.1.9}），这意味着
$$\frac{n(Y_n-\theta)^2}{\sigma^2}\to\chi_1^2\quad\text{（依分布）}.$$
因此，类似\autoref{thm:5.5.24}的论证即可证明如下定理.
\begin{theorem}[二阶Delta方法]\label{thm:5.5.26}
设随机变量序列$Y_n$满足$\sqrt n(Y_n-\theta)\to n(0,\sigma^2)$（依分布）.对于给定函数$g$及特定$\theta$，假定$g'(\theta)=0$，而$g''(\theta)$存在且非零.则
\begin{equation}\label{eq:5.5.13}
n[g(Y_n)-g(\theta)]\to\sigma^2\frac{g''(\theta)}2\chi_1^2\quad\text{（依分布）}.
\end{equation}
\end{theorem}
\sourcepageanchor{230}
当待估函数包含多个参数、估计量用到多个随机变量时，近似方法十分有用.生长研究中一个常见例子，是把体重与身高之比作为感兴趣变量.（回顾\autoref{sec:3}，两个正态随机变量之比服从Cauchy分布.比值问题对实验者很重要，但理论上十分棘手.）

这引出Delta方法的第二个推广，即多元情形.既然已有多元Taylor定理，这一推广并不出人意料.
\begin{example}[比值估计量的矩]\label{ex:5.5.27}\index{SI0812@比值}
设随机变量$X$与$Y$的非零均值分别为$\mu_X$和$\mu_Y$.待估参数函数为$g(\mu_X,\mu_Y)=\mu_X/\mu_Y$.直接计算得到
$$\frac\partial{\partial\mu_X}g(\mu_X,\mu_Y)=\frac1{\mu_Y},\qquad
\frac\partial{\partial\mu_Y}g(\mu_X,\mu_Y)=-\frac{\mu_X}{\mu_Y^2}.$$
一阶Taylor近似\eqref{eq:5.5.8}与\eqref{eq:5.5.9}给出
$$E\left(\frac XY\right)\approx\frac{\mu_X}{\mu_Y},$$
以及
\begin{align*}
\operatorname{Var}\left(\frac XY\right)&\approx\frac1{\mu_Y^2}\operatorname{Var}X+\frac{\mu_X^2}{\mu_Y^4}\operatorname{Var}Y-2\frac{\mu_X}{\mu_Y^3}\operatorname{Cov}(X,Y)\\
&=\left(\frac{\mu_X}{\mu_Y}\right)^2\left(\frac{\operatorname{Var}X}{\mu_X^2}+\frac{\operatorname{Var}Y}{\mu_Y^2}-2\frac{\operatorname{Cov}(X,Y)}{\mu_X\mu_Y}\right).
\end{align*}
由此得到比值估计量均值、方差的近似，而且只用到$\bar X$与$Y$的均值、方差和协方差.\footnote{译注：原文此处写$\bar X$与$Y$；按本例设定和前面的公式，应为$X$与$Y$.}精确计算几乎没有希望，也无法得到闭式表达式.
\end{example}
下面给出适用于比值估计量等情形的中心极限定理.虽然最终的中心极限定理是一元的，推导仍须处理多个随机变量.设向量值随机变量$\boldsymbol X=(X_1,\ldots,X_p)$的均值为$\boldsymbol\mu=(\mu_1,\ldots,\mu_p)$，协方差为$\operatorname{cov}(X_i,X_j)=\sigma_{ij}$，观测独立随机样本$\boldsymbol X_1,\ldots,\boldsymbol X_n$，计算均值$\bar X_i=\sum_{k=1}^nX_{ik}$，$i=1,\ldots,p$.\footnote{译注：原文均值定义漏写因子$1/n$；样本均值应为$\bar X_i=n^{-1}\sum_{k=1}^nX_{ik}$.}对函数$g(\boldsymbol x)=g(x_1,\ldots,x_p)$，利用\eqref{eq:5.5.7}之后的推导可写\footnote{译注：原文这一式写成等式，且导数在$\boldsymbol x$处取值.作为在$\boldsymbol\mu$处的一阶Taylor近似，应把等号理解为近似，并在$\boldsymbol\mu$处计算导数，或显式保留相应余项.}
$$g(\bar x_1,\ldots,\bar x_p)=g(\mu_1,\ldots,\mu_p)+\sum_{k=1}^p g'_k(\boldsymbol x)(\bar x_k-\mu_k),$$

于是得到如下定理.
\Needspace{10\baselineskip}
\begin{theorem}[多元Delta方法]\label{thm:5.5.28}
设$\boldsymbol X_1,\ldots,\boldsymbol X_n$为随机样本，$E(X_{ij})=\mu_i$、\allowbreak $\operatorname{cov}(X_{ik},X_{jk})=\sigma_{ij}$.给定具有连续一阶偏导数的函数$g$及特定值$\boldsymbol\mu=(\mu_1,\mu_2,\ldots,\mu_p)$，满足
$$\tau^2=\sum\sum\sigma_{ij}\frac{\partial g(\boldsymbol\mu)}{\partial\mu_i}\cdot\frac{\partial g(\boldsymbol\mu)}{\partial\mu_j}>0,$$
则\footnote{译注：原文左端写$\bar X_s$，但这里的向量维数是$p$，应为$\bar X_p$.}
$$\sqrt n[g(\bar X_1,\ldots,\bar X_s)-g(\mu_1,\ldots,\mu_p)]\to n(0,\tau^2)\quad\text{（依分布）}.$$

\end{theorem}
证明需要处理多元随机变量的收敛.这里不涉及这些多元问题的细节，暂且接受\autoref{thm:5.5.28}.感兴趣的读者可参阅\hyperref[bib:source200]{Lehmann和Casella（1998}，第1.8节）.



% 原书印刷200–208页，PDF物理231–239页.
\sourcepageanchor{231}
\subsection{随机样本的生成}\label{sec:5.6}\index{SI0863@生成}
到目前为止，我们关心的是描述随机变量行为的各种方法：变换、分布、矩的计算和极限定理.实际中，用这些随机变量描述真实现象并建立模型，对它们的观测就是我们收集的数据.

所以，通常我们观测来自分布$f(x\mid\theta)$的随机变量$X_1,\ldots,X_n$，主要关心如何用$f(x\mid\theta)$的性质描述随机变量的行为.本节实质上将这一策略反过来，关心从给定分布$f(x\mid\theta)$生成随机样本$X_1,\ldots,X_n$.
\begin{example}[指数寿命]\label{ex:5.6.1}
设某种电子元件的寿命用$\operatorname{exponential}(\lambda)$建模.制造商想确定$c$个元件中至少$t$个能使用$h$小时的概率.分步考虑，有
\begin{equation}\label{eq:5.6.1}
p_1=P(\text{元件至少使用$h$小时})=P(X\ge h\mid\lambda).
\end{equation}
假定元件相互独立，则可将$c$个元件的结果视为Bernoulli试验，故
\begin{equation}\label{eq:5.6.2}
\begin{aligned}p_2&=P(\text{至少$t$个元件使用$h$小时})\\
&=\sum_{k=t}^c\binom ck p_1^k(1-p_1)^{c-k}.
\end{aligned}
\end{equation}
虽然\eqref{eq:5.6.2}的计算直接，但计算负担可能很重，尤其$t$与$c$都很大时.而且指数模型有个优点：$p_1$能写成闭式，即
\begin{equation}\label{eq:5.6.3}
p_1=\int_h^\infty\frac1\lambda e^{-x/\lambda}\,dx=e^{-h/\lambda}.
\end{equation}
但若每个元件改用例如gamma分布建模，$p_1$可能无法用闭式表示，从而使$p_2$的计算更加复杂.
\end{example}
计算\eqref{eq:5.6.2}这类表达式的一种模拟方法，是生成服从所需分布的随机变量，再用大数定律（\autoref{thm:5.5.2}）验证模拟的有效性.即若$Y_i$（$i=1,\ldots,n$）为iid，则该定理的一个推论（只要假设成立）是
\begin{equation}\label{eq:5.6.4}
\frac1n\sum_{i=1}^n h(Y_i)\longrightarrow Eh(Y)
\end{equation}
在$n\to\infty$时依概率成立.（\eqref{eq:5.6.4}也几乎处处成立，这是强大数定律\autoref{thm:5.5.9}的推论.）
\begin{example}[续\autoref{ex:5.6.1}]\label{ex:5.6.2}
可按下列步骤计算概率$p_2$.对$j=1,\ldots,n$：
\begin{enumerate}[label=\textbf{\alph*.}]
\item 生成iid $X_1,\ldots,X_c\sim\operatorname{exponential}(\lambda)$.
\item 若至少$t$个$X_i\ge h$，令$Y_j=1$；否则令$Y_j=0$.
\end{enumerate}
由于$Y_j\sim\operatorname{Bernoulli}(p_2)$、$EY_j=p_2$，因此
$$\frac1n\sum_{j=1}^nY_j\to p_2,\qquad n\to\infty.$$
\end{example}
\sourcepageanchor{232}
\autoref{ex:5.6.1}和\autoref{ex:5.6.2}突出本节的主要问题：首先研究如何生成所需随机变量，然后用某个版本的大数定律验证模拟近似.

总得有个起点，这里首先假定能够生成iid均匀随机变量$U_1,\ldots,U_m$.（计算机科学家对均匀随机数生成问题已作了卓有成效的研究.）有许多生成伪随机数（pseudo-random numbers）的算法，产生的数可通过几乎所有均匀性检验.此外，多数优秀统计软件包都有合理的均匀随机数生成器.（关于伪随机数生成，更多内容见\hyperref[bib:source98]{Devroye 1985}或\hyperref[bib:source253]{Ripley 1987}.）

既然从均匀随机变量出发，这里真正的问题便不是生成所需随机变量，而是把均匀随机变量变换成所需分布.基本上有两类一般方法，我们把它们称为直接方法和间接方法，尽管这两个名称并没有提供多少信息.
\subsubsection{直接方法}\label{sec:5.6.1}
生成随机变量的直接方法（direct method），是指存在闭式函数$g(u)$，当$U\sim\operatorname{uniform}(0,1)$时，变换后的$Y=g(U)$具有所需分布.回顾概率积分变换\autoref{thm:2.1.10}，它已对连续随机变量完成这一任务：将任意分布变为均匀分布.所以逆变换就能解决这里的问题.
\begin{example}[概率积分变换]\label{ex:5.6.3}\index{SI0290@生成}\index{SI0847@概率积分变换}\index{SI1041@概率积分}
若连续随机变量$Y$的cdf为$F_Y$，则由\autoref{thm:2.1.10}，$U\sim\operatorname{uniform}(0,1)$时随机变量$F_Y^{-1}(U)$具有分布$F_Y$.若$Y\sim\operatorname{exponential}(\lambda)$，则
$$F_Y^{-1}(U)=-\lambda\log(1-U)$$
是$\operatorname{exponential}(\lambda)$随机变量（见\autoref{exercise:5.49}）.

因此，生成iid均匀随机变量$U_1,\ldots,U_n$后，$Y_i=-\lambda\log(1-U_i)$（$i=1,\ldots,n$）就是iid指数随机变量.例如取$n=10{,}000$，生成$u_1,u_2,\ldots,u_{10{,}000}$，计算得到
$$\frac1n\sum u_i=.5019,\qquad\frac1{n-1}\sum(u_i-\bar u)^2=.0842.$$
由弱大数定律（\autoref{thm:5.5.2}）导出的\eqref{eq:5.6.4}，$\bar U\to EU=1/2$；由\autoref{ex:5.5.3}，$S^2\to\operatorname{Var}U=1/12=.0833$，所以这些估计十分接近真实参数.变换后的变量$Y_i$服从$\operatorname{exponential}(2)$分布，计算得到
$$\frac1n\sum y_i=2.0004,\qquad\frac1{n-1}\sum(y_i-\bar y)^2=4.0908,$$
与$EY=2$、$\operatorname{Var}Y=4$很吻合.\autoref{fig:5.6.1}说明样本直方图与总体pdf的一致程度.
\end{example}
\sourcepageanchor{233}
\par\begin{figure}[H]\centering
\includegraphics[width=.60\textwidth]{assets/chapter5/fig-5-6-1.pdf}
\caption{从$\lambda=2$的指数pdf生成的10,000个观测的直方图及其pdf}\label{fig:5.6.1}
\end{figure}\par
指数分布与其他分布的关系，使许多随机变量可以快速生成.例如，若$U_i$为iid $\operatorname{uniform}(0,1)$随机变量，则
\begin{equation}\label{eq:5.6.5}\index{SI0255@生成}\index{SI0281@生成}\index{SI0307@生成}\index{SI0361@生成}
\begin{aligned}
Y&=-2\sum_{j=1}^\nu\log(U_j)\sim\chi_{2\nu}^2,\\
Y&=-\beta\sum_{j=1}^a\log(U_j)\sim\operatorname{gamma}(a,\beta),\\
Y&=\frac{\sum_{j=1}^a\log(U_j)}{\sum_{j=1}^{a+b}\log(U_j)}\sim\operatorname{beta}(a,b).
\end{aligned}
\end{equation}
还可以作许多变形（见\autoref{exercise:5.47}—\autoref{exercise:5.49}），但都由指数—均匀变换导出.

遗憾的是，这一变换有局限.例如，不能用它生成奇数自由度的$\chi^2$随机变量，所以不能得到$\chi_1^2$，进而得到极其有用的$\operatorname{normal}(0,1)$随机变量.下一小节会再讨论此问题.

回顾\autoref{ex:5.6.3}以及\eqref{eq:5.6.5}中各变换的基础，都是概率积分变换，一般可写成
\begin{equation}\label{eq:5.6.6}
F_Y^{-1}(u)=y\iff u=\int_{-\infty}^y f_Y(t)\,dt.
\end{equation}
将此公式用于指数分布尤其方便，因为积分方程有简单解（另见\autoref{exercise:5.59}）.但很多情况下\eqref{eq:5.6.6}没有闭式解，于是每生成一个随机变量都要解一次积分方程，实践中可能过于耗时和复杂.例如，用\eqref{eq:5.6.6}生成$\chi_1^2$就有这个问题.

没有闭式解时，应探索其他办法，包括其他类型的生成方法以及间接方法.下面是前一类方法的例子.
\begin{example}[{\hyperref[bib:source46]{Box–Muller，1958}}]\label{ex:5.6.4}\index{SI0099@Box–Muller变换}\index{SI1038@Box–Muller}
生成两个独立的$\operatorname{uniform}(0,1)$随机变量$U_1$和$U_2$，令
$$R=\sqrt{-2\log U_1},\qquad\theta=2\pi U_2.$$
则
$$X=R\cos\theta,\qquad Y=R\sin\theta$$
\sourcepageanchor{234}
是独立的$\operatorname{normal}(0,1)$随机变量.所以，虽然没有生成单个$n(0,1)$变量的快速变换，却有一次生成两个变量的方法（见\autoref{exercise:5.50}）.
\end{example}
遗憾的是，\autoref{ex:5.6.4}这样的解法并不多.它们还利用了特定分布的结构，因此不太适合作为一般策略.对于其他尚未讨论的连续分布，多数时候最好用间接方法生成.在探索这些方法之前，先以离散随机变量的情形结束本小节；此时\eqref{eq:5.6.6}很有用.

若离散随机变量$Y$的取值为$y_1<y_2<\cdots<y_k$，类似\eqref{eq:5.6.6}可写
\begin{equation}\label{eq:5.6.7}
\begin{aligned}
P[F_Y(y_i)\le U<F_Y(y_{i+1})]&=F_Y(y_{i+1})-F_Y(y_i)\\
&=P(Y=y_{i+1}).
\end{aligned}
\end{equation}
实际用\eqref{eq:5.6.7}生成离散随机变量很直接，可概括如下.为生成$Y_i\sim F_Y(y)$：
\begin{enumerate}[label=\textbf{\alph*.}]
\item 生成$U\sim\operatorname{uniform}(0,1)$.
\item 若$F_Y(y_i)\le U<F_Y(y_{i+1})$，令$Y=y_{i+1}$.
\end{enumerate}
定义$y_0=-\infty$及$F_Y(y_0)=0$.
\begin{example}[二项随机变量的生成]\label{ex:5.6.5}\index{SI0263@生成}\index{SI0286@生成}
例如，为生成$Y\sim\operatorname{binomial}(4,5/8)$，生成$U\sim\operatorname{uniform}(0,1)$并令
\begin{equation}\label{eq:5.6.8}
Y=\begin{cases}0,&0\le U<.020,\\1,&.020\le U<.152,\\2,&.152\le U<.481,\\3,&.481\le U<.847,\\4,&.847\le U\le1.\end{cases}
\end{equation}
\end{example}
即使离散随机变量的值域无限，如Poisson或负二项分布，\eqref{eq:5.6.8}中的算法也适用.理论上这可能需要大量求值，但实际上有简单巧妙的加速方法，因而通常不会如此.例如，不必按$1,2,\ldots$的次序逐个检查$y_i$，从靠近均值的$y_i$开始可能快得多（见\hyperref[bib:source253]{Ripley 1987}，第3.3节，以及\autoref{exercise:5.55}）.

后面将看到模拟方法的许多用途.先考察下面对Poisson分布的探索，它是\autoref{sec:10.1.4}将介绍的参数自助法（parametric bootstrap）的一个版本.
\begin{example}[Poisson方差的分布]\label{ex:5.6.6}\index{SI0096@参数}\index{SI0752@参数自助法}
若$X_1,\ldots,X_n$为iid $\operatorname{Poisson}(\lambda)$，由\autoref{thm:5.2.7}或\autoref{thm:5.2.11}，$\sum X_i$服从$\operatorname{Poisson}(n\lambda)$.因此很容易描述样本均值$\bar X$的分布，但描述样本方差$S^2=\frac1{n-1}\sum(X_i-\bar X)^2$的分布并不容易.

不过，$S^2$的分布很容易模拟.例如，取$n=5$，从$\operatorname{Poisson}(\lambda)$分布生成1,000个容量为5的样本.\autoref{fig:5.6.2}给出这些样本的直方图.\footnote{译注：此处原文的1,000个、容量5的设定，与图及后文实际使用的5,000个、容量13的设定不一致.图题和后续计算均按原文保留.}此外，模拟样本还能计算关于$S^2$的概率.若$S_i^2$为第$i$个模拟样本计算出的值，则
$$\frac1M\sum_{i=1}^M I(S_i^2\ge a)\to P_\lambda(S^2\ge a),\qquad M\to\infty.$$
\sourcepageanchor{235}
\par\begin{figure}[H]\centering
\includegraphics[width=.60\textwidth]{assets/chapter5/fig-5-6-2.pdf}
\caption{从$\lambda=18.69$的Poisson分布抽取5,000个容量为13的样本所得样本方差$S^2$的直方图.这5,000个数值的均值与标准差为18.86和7.68}\label{fig:5.6.2}
\end{figure}\par
为说明这种方法的用途，考虑1984年8月下旬在哈德逊河取得的海湾{\upshape\CJKfontspec{SimSun}鳀}鱼幼体计数样本：
\begin{equation}\label{eq:5.6.9}
19,32,29,13,8,12,16,20,14,17,22,18,23.
\end{equation}
若假定幼体在河中随机且均匀分布，则固定大小渔网捕获的数量应服从Poisson分布.这一论证来自Poisson公设的空间版本（见\autoref{sec:2.6}）.为考察假设是否合理，可检查观测数据的均值和方差是否符合Poisson假设.

对\eqref{eq:5.6.9}的数据，计算得$\bar x=18.69$、$s^2=44.90$.Poisson假设下，预期两者相同.当然，由于抽样变异，它们不会完全相同，可用模拟了解应有的差异.\autoref{fig:5.6.2}中，从$\lambda=18.69$的Poisson分布模拟5,000个容量为$n=13$的样本，构造$S^2$的相对频数直方图.观测值$S^2=44.90$落在分布尾部.事实上，有27个$S^2$值大于44.90，故可估计
$$P(S^2>44.90\mid\lambda=18.69)=\frac1{5000}\sum_{i=1}^{5000}I(S_i^2>44.90)=\frac{27}{5000}=.0054,$$
这使我们怀疑Poisson假设；见\autoref{exercise:5.54}.（这类发现催生了一个很糟糕的双语双关：“哈德逊河里有些‘鱼’常——Poisson失效了.”\footnote{译注：原文为“Something is fishy in the Hudson—the Poisson has failed.”英语fishy兼有“鱼腥的”与“可疑的”之意；Poisson与法语“鱼”poisson同形.}）
\end{example}
\subsubsection{间接方法}\label{sec:5.6.2}
找不到容易的直接变换来生成所需随机变量时，一种非常有力的间接方法——接受／拒绝算法（Accept/Reject Algorithm）——常能提供解法.它的思想或许最好用一个简单例子说明.
\begin{example}[beta随机变量的生成——I]\label{ex:5.6.7}\index{SI0001@接受／拒绝算法}\index{SI0006@接受／拒绝}\index{SI0255@生成}
设目标是生成$Y\sim\operatorname{beta}(a,b)$.若$a,b$都是整数，可用直接变换法\eqref{eq:5.6.5}；若不是整数，该方法就不适用.具体取$a=2.7$、$b=6.3$.\autoref{fig:5.6.3}把beta密度$f_Y(y)$放在边长为1和$c\ge\max_yf_Y(y)$的长方形里.考虑以下计算$P(Y\le y)$的方法.若$(U,V)$为独立的$\operatorname{uniform}(0,1)$随机变量，则阴影区域的概率为
\sourcepageanchor{236}
\begin{equation}\label{eq:5.6.10}
\begin{aligned}
P\left(V\le y,U\le\frac1c f_Y(V)\right)&=\int_0^y\int_0^{f_Y(v)/c}\,du\,dv\\
&=\frac1c\int_0^y f_Y(v)\,dv\\
&=\frac1cP(Y\le y).
\end{aligned}
\end{equation}
\par\begin{figure}[H]\centering
\includegraphics[width=.62\textwidth]{assets/chapter5/fig-5-6-3.pdf}
\caption{$a=2.7$、$b=6.3$的beta分布，$c=\max_yf_Y(y)=2.669$.均匀随机变量$V$给出横坐标，利用$U$判断是否处于密度曲线下方}\label{fig:5.6.3}
\end{figure}\par
所以可以从均匀分布的概率算出beta概率，这提示可从均匀随机变量生成beta随机变量.

在\eqref{eq:5.6.10}中令$y=1$，得$1/c=P(U<(1/c)f_Y(V))$，因此
\begin{equation}\label{eq:5.6.11}
\begin{aligned}
P(Y\le y)&=\frac{P(V\le y,U\le(1/c)f_Y(V))}{P(U\le(1/c)f_Y(V))}\\
&=P(V\le y\mid U\le(1/c)f_Y(V)),
\end{aligned}
\end{equation}
这提示如下算法.为生成$Y\sim\operatorname{beta}(a,b)$：
\begin{enumerate}[label=\textbf{\alph*.}]
\item\label{item:5.6.7a} 生成独立的$\operatorname{uniform}(0,1)$变量$(U,V)$.
\item\phantomsection\label{itm:ex:5.6.7:b} 若$U<(1/c)f_Y(V)$，令$Y=V$；否则返回步骤\hyperref[item:5.6.7a]{a}.
\end{enumerate}
只要$c\ge\max_yf_Y(y)$，这一算法就能生成$\operatorname{beta}(a,b)$随机变量；实际上还可推广到支撑集有界的任意密度（\autoref{exercise:5.59}和\autoref{exercise:5.60}）.
\end{example}
显然，$c$的最佳选择是$c=\max_yf_Y(y)$.为说明理由，注意该算法的运行次数不固定：不知道要获得一个$Y$需要多少对$(U,V)$.但若定义随机变量
\begin{equation}\label{eq:5.6.12}
N=\text{生成一个$Y$所需的$(U,V)$对数},
\end{equation}
回顾$1/c=P(U\le(1/c)f_Y(V))$，有
$$\begin{aligned}P(N=1)&=\frac1c,\\P(N=2)&=\left(1-\frac1c\right)\frac1c,\\&\ \vdots\end{aligned}$$
\sourcepageanchor{237}
所以$N$是几何随机变量.因此，为生成一个$Y$，平均需要$E(N)=c$对$(U,V)$；从这一意义上，最小化$c$会使算法最优.

由\autoref{fig:5.6.3}可见，$U>(1/c)f_Y(V)$区域造成浪费，这是因为用均匀随机变量$V$生成beta随机变量$Y$.为改进，可以从更接近beta的分布开始.

算法的检验步骤\hyperref[itm:ex:5.6.7:b]{(b)}可理解为检查随机变量$V$是否“看起来像”来自密度$f_Y$.设$V\sim f_V$，计算
$$M=\sup_y\frac{f_Y(y)}{f_V(y)}<\infty.$$
步骤\hyperref[itm:ex:5.6.7:b]{(b)}的一般化，是比较$U\sim\operatorname{uniform}(0,1)$与$(1/M)f_Y(V)/f_V(V)$.比值越大，$V$越“像”来自随机密度$f_Y$，越有可能满足$U<(1/M)f_Y(V)/f_V(V)$.这就是一般接受／拒绝算法的基础.
\subsubsection{接受／拒绝算法}\label{sec:5.6.3}
\begin{theorem}\label{thm:5.6.8}\index{SI0001@接受／拒绝算法}\index{SI0006@接受／拒绝}
设$Y\sim f_Y(y)$、$V\sim f_V(V)$，其中$f_Y$和$f_V$有共同支撑集，且
$$M=\sup_y f_Y(y)/f_V(y)<\infty.$$
为生成随机变量$Y\sim f_Y$：
\begin{enumerate}[label=\textbf{\alph*.}]
\item\label{item:5.6.8a} 生成独立的$U\sim\operatorname{uniform}(0,1)$和$V\sim f_V$.
\item 若$U<(1/M)f_Y(V)/f_V(V)$，令$Y=V$；否则返回\hyperref[item:5.6.8a]{第一步}.
\end{enumerate}
\end{theorem}
\begin{proof}
生成的随机变量$Y$具有cdf
\begin{align*}
P(Y\le y)&=P(V\le y\mid\text{停止})\\
&=P(V\le y\mid U<(1/M)f_Y(V)/f_V(V))\\
&=\frac{P(V\le y,U<(1/M)f_Y(V)/f_V(V))}{P(U<(1/M)f_Y(V)/f_V(V))}\\
&=\frac{\displaystyle\int_{-\infty}^y\int_0^{(1/M)f_Y(v)/f_V(v)}\,du\,f_V(v)\,dv}{\displaystyle\int_{-\infty}^\infty\int_0^{(1/M)f_Y(v)/f_V(v)}\,du\,f_V(v)\,dv}\\
&=\int_{-\infty}^y f_Y(v)\,dv,
\end{align*}
正是所需cdf.
\end{proof}
另注意
$$\begin{aligned}
M&=\sup_yf_Y(y)/f_V(y)\\
&=[P(U<(1/M)f_Y(V)/f_V(V))]^{-1}\\
&=\frac1{P(\text{停止})},
\end{aligned}$$
所以生成一个$Y$所需的试验次数服从$\operatorname{geometric}(1/M)$，$M$为期望试验次数.
\sourcepageanchor{238}
\begin{example}[beta随机变量的生成——II]\label{ex:5.6.9}\index{SI0255@生成}
为生成$Y\sim\operatorname{beta}(2.7,6.3)$，考虑算法：
\begin{enumerate}[label=\textbf{\alph*.}]
\item\label{item:5.6.9a} 生成$U\sim\operatorname{uniform}(0,1)$、$V\sim\operatorname{beta}(2,6)$.
\item 若$U<(1/M)f_Y(V)/f_V(V)$，令$Y=V$；否则返回\hyperref[item:5.6.9a]{第一步}.
\end{enumerate}
只要$\sup_y f_Y(y)/f_V(y)\le M<\infty$，该接受／拒绝算法就生成所需的$Y$.对给定密度有$M=1.67$，满足要求（见\autoref{exercise:5.63}）.

这一算法有$EN=1.67$，而\autoref{ex:5.6.7}中使用均匀$V$的算法有$EN=2.67$.虽然这似乎说明后一算法更快，\footnote{译注：原文此处为“latter algorithm”，但按前面的试验次数比较，应指当前使用beta候选的前一算法.随后关于每次生成成本的提醒仍然成立.}别忘了生成一个$\operatorname{beta}(2,6)$变量需要八个均匀随机变量.所以比较算法不总是直接的，还须同时考虑计算机速度和编程便利性.
\end{example}
必须强调$M<\infty$的重要性.它可理解为要求$V$的密度（常称候选密度，candidate density）的尾部比$Y$的密度（常称目标密度，target density）更重.这有助于保证$Y$的各种值、包括尾部值，都得到良好代表.例如，若$V\sim\operatorname{Cauchy}$、$Y\sim n(0,1)$，预期$V$样本的范围宽于$Y$样本，因此基于这两个密度的接受／拒绝算法应有良好表现.反过来，把$n(0,1)$变量变成Cauchy变量就困难得多，因为极端值的代表不足.

但有时目标密度具有重尾，难以找到使$M$有限的候选密度.此时接受／拒绝算法不再适用，便引出另一类方法，称为马尔可夫链Monte Carlo方法（Markov Chain Monte Carlo，MCMC）.其特例包括Gibbs抽样器（Gibbs Sampler）和Metropolis算法.下面给出后者.

\noindent\textbf{Metropolis算法.}\index{SI0015@Metropolis}\semanticindexanchor{15:238}\index{SI0688@Metropolis算法}\semanticindexanchor{688:238} 设$Y\sim f_Y(y)$、$V\sim f_V(v)$，其中$f_Y$和$f_V$有共同支撑集.为生成$Y\sim f_Y$：
\begin{enumerate}[label=\arabic*.,start=0]
\item 生成$V\sim f_V$，令$Z_0=V$.
\end{enumerate}
对$i=1,2,\ldots$：
\begin{enumerate}[label=\arabic*.]
\item 生成$U_i\sim\operatorname{uniform}(0,1)$、$V_i\sim f_V$，计算
$$\rho_i=\min\left\{\frac{f_Y(V_i)}{f_V(V_i)}\cdot\frac{f_V(Z_{i-1})}{f_Y(Z_{i-1})},1\right\}.$$
\item 令
$$Z_i=\begin{cases}V_i,&U_i\le\rho_i,\\Z_{i-1},&U_i>\rho_i.\end{cases}$$
\end{enumerate}
则当$i\to\infty$时，$Z_i$依分布收敛到$Y$.

虽然算法不要求$M$有限，但它并不生成密度恰为$f_Y$的随机变量，而是产生一个收敛序列.实际运行一段时间后（$i$变大），产生的$Z$的表现就很接近来自$f_Y$的变量.（关于Metropolis算法的初步介绍，见Chib和Greenberg.）
\sourcepageanchor{239}
虽然MCMC方法至少可追溯到\hyperref[bib:source219]{Metropolis等（1953）}，但它们真正受到广泛重视，是在\hyperref[bib:source137]{Geman和Geman（1984）}的基础上，随着\hyperref[bib:source134]{Gelfand和Smith（1990）}的工作而发生的.更多内容见\autoref{sec:5.8.5}.



% 原书印刷208–217页，PDF物理239–248页.
\subsection{习题}\label{sec:5.7}
\begin{exercise}\label{exercise:5.1}
某总体有1\%的人是色盲.若要样本中包含色盲者的概率至少为.95，样本容量须多大？（假定总体大到可视为无限总体，抽样可视为有放回抽样.）
\end{exercise}
\begin{exercise}\label{exercise:5.2}
设$X_1,X_2,\ldots$联合连续且独立，各自的边际pdf均为$f(x)$，每个$X_i$表示某地点的年降雨量.
\begin{enumerate}[label=(\alph*)]
\item 求直到第一次超过第一年降雨量$X_1$所需年数的分布.
\item 证明直到第一次超过$X_1$所需年数的均值为无穷.
\end{enumerate}
\end{exercise}
\begin{exercise}\label{exercise:5.3}
设$X_1,\ldots,X_n$为iid随机变量，具有连续cdf $F_X$，且$EX_i=\mu$.定义随机变量$Y_1,\ldots,Y_n$为
$$Y_i=\begin{cases}1,&X_i>\mu,\\0,&X_i\le\mu.\end{cases}$$
求$\sum_{i=1}^nY_i$的分布.
\end{exercise}
\begin{exercise}\label{exercise:5.4}\index{SI0421@可交换}
iid随机变量的一种推广是可交换随机变量（exchangeable random variables），这一概念出自\hyperref[bib:source95]{De Finetti（1972）}.\hyperref[bib:source119]{Feller（1971）}也讨论了可交换性.若随机变量$X_1,\ldots,X_n$中任意容量为$k$（$k\le n$）的子集的任意排列都具有相同分布，则称它们可交换.本题给出可交换但非iid的例子.设$X_i\mid P\sim\operatorname{Bernoulli}(P)$，$i=1,\ldots,n$，条件iid，且$P\sim\operatorname{uniform}(0,1)$.
\begin{enumerate}[label=(\alph*)]
\item 证明任意$k$个$X$的边际分布都与下式相同：
$$P(X_1=x_1,\ldots,X_k=x_k)=\int_0^1p^t(1-p)^{k-t}\,dp=\frac{t!(k-t)!}{(k+1)!},$$
其中$t=\sum_{i=1}^kx_i$.因此这些$X$可交换.
\item 证明在边际意义下，
$$P(X_1=x_1,\ldots,X_n=x_n)\ne\prod_{i=1}^nP(X_i=x_i),$$
所以这些$X$的分布可交换，但并非iid.
\end{enumerate}
（De Finetti证明了无限可交换随机变量序列的一个优美刻画定理：任何这样的序列都是iid随机变量的混合.）
\end{exercise}
\begin{exercise}\label{exercise:5.5}\index{SI0927@概率密度函数}
设$X_1,\ldots,X_n$为iid，pdf为$f_X(x)$，样本均值为$\bar X$.证明即使$X$的矩母函数不存在，也有
$$f_{\bar X}(x)=nf_{X_1+\cdots+X_n}(nx).$$
\end{exercise}
\begin{exercise}\label{exercise:5.6}
若$X$的pdf为$f_X(x)$，与$X$独立的$Y$的pdf为$f_Y(y)$，分别为下列随机变量$Z$建立类似\eqref{eq:5.2.3}的公式.
\begin{enumerate}[label=(\alph*)]\item $Z=X-Y$.\item $Z=XY$.\item $Z=X/Y$.\end{enumerate}
\end{exercise}
\sourcepageanchor{240}
\begin{exercise}\label{exercise:5.7}
\autoref{ex:5.2.10}推导两个独立Cauchy随机变量之和的分布时，需要部分分式分解.本题补充该例省略的细节.
\begin{enumerate}[label=(\alph*)]
\item 求满足下式的常数$A,B,C,D$：
$$\begin{aligned}
&\frac1{1+(w/\sigma)^2}\frac1{1+((z-w)/\tau)^2}\\
&\quad=\frac{Aw}{1+(w/\sigma)^2}+\frac B{1+(w/\sigma)^2}-\frac{Cw}{1+((z-w)/\tau)^2}-\frac D{1+((z-w)/\tau)^2},
\end{aligned}$$
其中$A,B,C,D$可以依赖$z$，但不依赖$w$.
\item\phantomsection\label{itm:exercise:5.7:b} 利用
$$\int\frac t{1+t^2}\,dt=\frac12\log(1+t^2)+\text{常数},\qquad\int\frac1{1+t^2}\,dt=\arctan(t)+\text{常数},$$
计算\eqref{eq:5.2.4}，进而验证\eqref{eq:5.2.5}.
\end{enumerate}
（注意\hyperref[itm:exercise:5.7:b]{(b)}中的积分很微妙.Cauchy均值不存在，所以下列积分不存在：
$$\int_{-\infty}^\infty\frac{Aw}{1+(w/\sigma)^2}\,dw,\qquad\int_{-\infty}^\infty\frac{Cw}{1+((z-w)/\tau)^2}\,dw.$$
但二者之差的积分确实存在，而这正是所需的.）
\end{exercise}
\begin{exercise}\label{exercise:5.8}\index{SI0939@恒等式}
设$X_1,\ldots,X_n$为随机样本，$\bar X$与$S^2$按通常方式计算.
\begin{enumerate}[label=(\alph*)]
\item\phantomsection\label{item:5.8.a} 证明
$$S^2=\frac1{2n(n-1)}\sum_{i=1}^n\sum_{j=1}^n(X_i-X_j)^2.$$
现在假定$X_i$具有有限四阶矩，记$\theta_1=EX_i$、$\theta_j=E(X_i-\theta_1)^j$，$j=2,3,4$.
\item\phantomsection\label{item:5.8.b} 证明$\operatorname{Var}S^2=\frac1n(\theta_4-\frac{n-3}{n-1}\theta_2^2)$.
\item 用$\theta_1,\ldots,\theta_4$表示$\operatorname{Cov}(\bar X,S^2)$.什么条件下$\operatorname{Cov}(\bar X,S^2)=0$？
\end{enumerate}
\end{exercise}
\begin{exercise}\label{exercise:5.9}\index{SI0542@拉格朗日}
证明Lagrange恒等式：对任意数$a_1,a_2,\ldots,a_n$和$b_1,b_2,\ldots,b_n$，
$$\left(\sum_{i=1}^na_i^2\right)\left(\sum_{i=1}^nb_i^2\right)-\left(\sum_{i=1}^na_ib_i\right)^2=\sum_{i=1}^{n-1}\sum_{j=i+1}^n(a_ib_j-a_jb_i)^2.$$
利用此恒等式证明，相关系数等于1，当且仅当全部样本点位于一条直线上（\hyperref[bib:source334]{Wright 1992）}.\footnote{译注：原文写相关系数“等于1”.一般直线上的点可能具有负斜率，正确的共线刻画应是相关系数绝对值为1，并要求两组样本方差非零；等于$+1$还须正斜率.}提示：先证明$n=2$，再用归纳法.
\end{exercise}
\begin{exercise}\label{exercise:5.10}
设$X_1,\ldots,X_n$为来自$n(\mu,\sigma^2)$总体的随机样本.
\begin{enumerate}[label=(\alph*)]
\item\phantomsection\label{itm:exercise:5.10:a} 用$\mu$与$\sigma^2$表示\autoref{exercise:5.8}定义的$\theta_1,\ldots,\theta_4$.
\item 结合\autoref{exercise:5.8}及\hyperref[itm:exercise:5.10:a]{(a)}的结果，计算$\operatorname{Var}S^2$.
\item 用完全不同且更容易的方法计算$\operatorname{Var}S^2$：利用$(n-1)S^2/\sigma^2\sim\chi_{n-1}^2$.
\end{enumerate}
\end{exercise}
\begin{exercise}\label{exercise:5.11}
设$\bar X$与$S^2$由来自方差$\sigma^2$有限的总体的随机样本$X_1,\ldots,X_n$计算.已知$ES^2=\sigma^2$.证明$ES\le\sigma$，且若$\sigma^2>0$，则$ES<\sigma$.
\end{exercise}
\begin{exercise}\label{exercise:5.12}
设$X_1,\ldots,X_n$为来自$n(0,1)$总体的随机样本.定义
$$Y_1=\left|\frac1n\sum_{i=1}^nX_i\right|,\qquad Y_2=\frac1n\sum_{i=1}^n|X_i|.$$
计算$EY_1$、$EY_2$，并建立它们之间的不等式.
\end{exercise}
\begin{exercise}\label{exercise:5.13}
设$X_1,\ldots,X_n$为iid $n(\mu,\sigma^2)$.求样本方差$S^2$的函数$g(S^2)$，满足$Eg(S^2)=\sigma$.（提示：尝试$g(S^2)=c\sqrt{S^2}$，其中$c$为常数.）
\end{exercise}
\sourcepageanchor{241}
\needspace{5\baselineskip}
\begin{exercise}\label{exercise:5.14}\leavevmode
\begin{enumerate}[label=(\alph*)]
\item 证明\autoref{lem:5.3.3}的一般结论可由$\mu_i=0$、$\sigma_i^2=1$的特例推出.即证明：若$X_j=\sigma_jZ_j+\mu_j$、$Z_j\sim n(0,1)$，$j=1,\ldots,n$，全部独立，$a_{ij},b_{rj}$为常数，且
$$\operatorname{Cov}\left(\sum_{j=1}^na_{ij}Z_j,\sum_{j=1}^nb_{rj}Z_j\right)=0\implies\sum_{j=1}^na_{ij}Z_j\text{与}\sum_{j=1}^nb_{rj}Z_j\text{独立},$$
则
$$\operatorname{Cov}\left(\sum_{j=1}^na_{ij}X_j,\sum_{j=1}^nb_{rj}X_j\right)=0\implies\sum_{j=1}^na_{ij}X_j\text{与}\sum_{j=1}^nb_{rj}X_j\text{独立}.$$
\item 验证\autoref{lem:5.3.3}中$\operatorname{Cov}(\sum_{j=1}^na_{ij}X_j,\sum_{j=1}^nb_{rj}X_j)$的表达式.
\end{enumerate}
\end{exercise}
\begin{exercise}\label{exercise:5.15}\index{SI0928@递推关系}\index{SI0941@递推关系}
建立以下均值、方差的递推关系.设$\bar X_n$和$S_n^2$分别为$X_1,\ldots,X_n$的均值与方差，又取得一个观测$X_{n+1}$.证明
\begin{enumerate}[label=(\alph*)]
\item $\bar X_{n+1}=(X_{n+1}+n\bar X_n)/(n+1)$.
\item $nS_{n+1}^2=(n-1)S_n^2+\frac n{n+1}(X_{n+1}-\bar X_n)^2$.
\end{enumerate}
\end{exercise}
\begin{exercise}\label{exercise:5.16}
设$X_i$（$i=1,2,3$）独立且各服从$n(i,i^2)$分布.在每种情形中，用$X_i$构造具有指定分布的统计量.
\begin{enumerate}[label=(\alph*)]\item 自由度为3的卡方分布.\item 自由度为2的$t$分布.\item 自由度为1和2的$F$分布.\end{enumerate}
\end{exercise}
\begin{exercise}\label{exercise:5.17}
设随机变量$X$服从$F_{p,q}$分布.
\begin{enumerate}[label=(\alph*)]\item 推导$X$的pdf.\item 推导$X$的均值与方差.\item 证明$1/X$服从$F_{q,p}$分布.\item 证明$(p/q)X/[1+(p/q)X]$服从参数为$p/2$和$q/2$的beta分布.\end{enumerate}
\end{exercise}
\begin{exercise}\label{exercise:5.18}\index{SI0367@与t的关系}\index{SI0384@与正态的关系}
设$X$服从自由度为$p$的Student $t$分布.
\begin{enumerate}[label=(\alph*)]
\item\phantomsection\label{itm:exercise:5.18:a} 推导$X$的均值和方差.
\item\phantomsection\label{itm:exercise:5.18:b} 证明$X^2$服从自由度为1和$p$的$F$分布.
\item 以$f(x\mid p)$表示$X$的pdf.证明对每个$-\infty<x<\infty$，
$$\lim_{p\to\infty}f(x\mid p)\to\frac1{\sqrt{2\pi}}e^{-x^2/2}.$$
这正确地提示，$p\to\infty$时，$X$依分布收敛到$n(0,1)$随机变量.（提示：使用Stirling公式.）
\item 利用\hyperref[itm:exercise:5.18:a]{(a)}、\hyperref[itm:exercise:5.18:b]{(b)}说明，$p\to\infty$时，$X^2$依分布收敛到$\chi_1^2$随机变量.
\item 关于$p\to\infty$时$qF_{q,p}$的分布极限，可以提出什么猜想？
\end{enumerate}
\end{exercise}
\begin{exercise}\label{exercise:5.19}\index{SI0990@卡方}\index{SI0991@F}\leavevmode
\begin{enumerate}[label=(\alph*)]
\item\phantomsection\label{itm:exercise:5.19:a} 证明$\chi^2$分布随自由度随机递增（stochastically increasing）：若$p>q$，则对任意$a$有$P(\chi_p^2>a)\ge P(\chi_q^2>a)$，且对某个$a$严格不等号成立.
\item 用\hyperref[itm:exercise:5.19:a]{(a)}证明，对任意$\nu$，$kF_{k,\nu}$随$k$随机递增.
\item 证明对任意$k,\nu,\alpha$，$kF_{\alpha,k,\nu}>(k-1)F_{\alpha,k-1,\nu}$.（$F_{\alpha,k-1,\nu}$表示水平为$\alpha$的临界点；见\autoref{sec:8.3.1}.另见\autoref{sec:8.5.1}和\autoref{exercise:11.15}.）
\end{enumerate}
\end{exercise}
\begin{exercise}\label{exercise:5.20}\index{SI0345@正态–伽马}\index{SI0347@t–F}\leavevmode
\begin{enumerate}[label=(\alph*)]
\item 以下论证说明$t$分布是正态分布的混合：
$$P(T_\nu\le t)=P\left(\frac Z{\sqrt{\chi_\nu^2/\nu}}\le t\right)=\int_0^\infty P(Z\le t\sqrt x/\sqrt\nu)P(\chi_\nu^2=x)\,dx,$$
\sourcepageanchor{242}
其中$T_\nu$为自由度为$\nu$的$t$随机变量.利用微积分基本定理，并将$P(\chi_\nu^2=\nu x)$解释为pdf，得到
$$f_{T_\nu}(t)=\int_0^\infty\frac1{\sqrt{2\pi}}e^{-t^2x/(2\nu)}\frac{\sqrt x}{\sqrt\nu}\frac1{\Gamma(\nu/2)2^{\nu/2}}x^{(\nu/2)-1}e^{-x/2}\,dx,$$
它是正态分布的尺度混合.用直接积分验证此公式.
\item\phantomsection\label{itm:exercise:5.20:b} $F$分布也有类似公式，即可写成卡方分布的混合.若$F_{1,\nu}$为自由度为1和$\nu$的$F$随机变量，则
$$P(F_{1,\nu}\le\nu t)=\int_0^\infty P(\chi_1^2\le ty)f_\nu(y)\,dy,$$
其中$f_\nu(y)$为$\chi_\nu^2$的pdf.用微积分基本定理得到$F_{1,\nu}$的pdf的积分表达式，并证明积分等于该pdf.
\item 验证\hyperref[itm:exercise:5.20:b]{(b)}的推广
$$P\left(F_{m,\nu}\le\frac\nu m t\right)=\int_0^\infty P(\chi_m^2\le ty)f_\nu(y)\,dy$$
对所有整数$m>1$成立.
\end{enumerate}
\end{exercise}
\begin{exercise}\label{exercise:5.21}
两个连续iid随机变量中较大者超过总体中位数的概率是多少？推广到容量为$n$的样本.
\end{exercise}
\begin{exercise}\label{exercise:5.22}
设$X,Y$为iid $n(0,1)$随机变量，定义$Z=\min(X,Y)$.证明$Z^2\sim\chi_1^2$.
\end{exercise}
\begin{exercise}\label{exercise:5.23}
设$U_i$（$i=1,2,\ldots$）是独立的$\operatorname{uniform}(0,1)$随机变量，$X$具有分布
$$P(X=x)=\frac c{x!},\qquad x=1,2,3,\ldots,$$
其中$c=1/(e-1)$.求$Z=\min\{U_1,\ldots,U_X\}$的分布.（提示：$Z\mid X=x$的分布是容量为$x$的样本的第一个次序统计量的分布.）
\end{exercise}
\begin{exercise}\label{exercise:5.24}
设随机样本$X_1,\ldots,X_n$来自pdf为
$$f_X(x)=\begin{cases}1/\theta,&0<x<\theta,\\0,&\text{其他情形}\end{cases}$$
的总体，次序统计量为$X_{(1)}<\cdots<X_{(n)}$.证明$X_{(1)}/X_{(n)}$与$X_{(n)}$独立.
\end{exercise}
\begin{exercise}\label{exercise:5.25}
推广上一题，设$X_1,\ldots,X_n$为iid，pdf为
$$f_X(x)=\begin{cases}\dfrac a{\theta^a}x^{a-1},&0<x<\theta,\\0,&\text{其他情形}.\end{cases}$$
次序统计量为$X_{(1)}<\cdots<X_{(n)}$.证明$X_{(1)}/X_{(2)},X_{(2)}/X_{(3)},\ldots,X_{(n-1)}/X_{(n)}$与$X_{(n)}$相互独立，并分别求它们的分布.
\end{exercise}
\begin{exercise}\label{exercise:5.26}\index{SI0733@联合累积分布函数}
完成\autoref{thm:5.4.6}的证明.
\begin{enumerate}[label=(\alph*)]
\item 令$U$表示$X_1,\ldots,X_n$中不超过$u$的个数，$V$表示大于$u$且不超过$v$的个数.证明$(U,V,n-U-V)$是试验次数为$n$、单元概率为$(F_X(u),F_X(v)-F_X(u),1-F_X(v))$的多项随机向量.
\sourcepageanchor{243}
\item 证明$X_{(i)}$与$X_{(j)}$的联合cdf可写成
\begin{align*}
F_{X_{(i)},X_{(j)}}(u,v)&=P(U\ge i,U+V\ge j)\\
&=\sum_{k=i}^{j-1}\sum_{m=j-k}^{n-k}P(U=k,V=m)+P(U\ge j)\\
&=\sum_{k=i}^{j-1}\sum_{m=j-k}^{n-k}\frac{n!}{k!m!(n-k-m)!}[F_X(u)]^k[F_X(v)-F_X(u)]^m\\
&\qquad\times[1-F_X(v)]^{n-k-m}+P(U\ge j).
\end{align*}
\item 按照\eqref{eq:4.1.3}计算混合偏导数，求联合pdf.（$P(U\ge j)$只依赖$u$而不依赖$v$，所以其混合偏导数为0.其余项可利用\eqref{eq:5.4.5}之类的关系大量相消.）
\end{enumerate}
\end{exercise}
\begin{exercise}\label{exercise:5.27}
设$X_1,\ldots,X_n$为iid，pdf为$f_X(x)$、cdf为$F_X(x)$，次序统计量为$X_{(1)}<\cdots<X_{(n)}$.
\begin{enumerate}[label=(\alph*)]
\item 用$f_X$和$F_X$表示给定$X_{(j)}$时$X_{(i)}$的条件pdf.
\item 求$V\mid R=r$的pdf，其中$V,R$见\autoref{ex:5.4.7}.
\end{enumerate}
\end{exercise}
\begin{exercise}\label{exercise:5.28}
设$X_1,\ldots,X_n$为iid，pdf为$f_X(x)$、cdf为$F_X(x)$，$X_{(i_1)}<\cdots<X_{(i_l)}$与$X_{(j_1)}<\cdots<X_{(j_m)}$是任意两组互不重叠的次序统计量.用$f_X,F_X$表示：
\begin{enumerate}[label=(\alph*)]
\item $X_{(i_1)},\ldots,X_{(i_l)}$的边际cdf和pdf.
\item 给定$X_{(j_1)},\ldots,X_{(j_m)}$时，$X_{(i_1)},\ldots,X_{(i_l)}$的条件cdf和pdf.
\end{enumerate}
\end{exercise}
\begin{exercise}\label{exercise:5.29}
一家小册子制造商每箱装100本.已知小册子的平均重量为1盎司，标准差为.05盎司.制造商想计算
$$P(\text{100本小册子的重量超过100.4盎司}),$$
用它帮助检测一箱是否放入过多小册子.说明如何计算这个概率的（近似？）值，并列出所需定理或假设.
\end{exercise}
\begin{exercise}\label{exercise:5.30}
若$\bar X_1$和$\bar X_2$是从方差为$\sigma^2$的总体取得的两个独立、容量均为$n$的样本的均值，求使$P(|\bar X_1-\bar X_2|<\sigma/5)\approx.99$的$n$.说明计算依据.
\end{exercise}
\begin{exercise}\label{exercise:5.31}
设$\bar X$是从均值$\mu$、方差$\sigma^2=9$的总体取得的100个观测的均值.求使$\bar X-\mu$以至少.90的概率落在其中的界限.分别使用Chebychev不等式和中心极限定理，并评论两种方法.
\end{exercise}
\begin{exercise}\label{exercise:5.32}
设随机变量序列$X_1,X_2,\ldots$依概率收敛到常数$a$，且对所有$i$都有$P(X_i>0)=1$.
\begin{enumerate}[label=(\alph*)]\item\phantomsection\label{itm:exercise:5.32:a} 验证$Y_i=\sqrt{X_i}$与$Y'_i=a/X_i$定义的序列依概率收敛.\item 利用\hyperref[itm:exercise:5.32:a]{(a)}证明\autoref{ex:5.5.18}所用的结论：$\sigma/S_n$依概率收敛到1.\end{enumerate}
\end{exercise}
\begin{exercise}\label{exercise:5.33}
设随机变量序列$X_n$依分布收敛到随机变量$X$，而随机变量序列$Y_n$满足：对任意有限数$c$，
$$\lim_{n\to\infty}P(Y_n>c)=1.$$
证明对任意有限数$c$，$\lim_{n\to\infty}P(X_n+Y_n>c)=1$.

（\autoref{sec:10.3.2}讨论检验的功效性质时，会使用这类结果.）
\end{exercise}
\sourcepageanchor{244}
\begin{exercise}\label{exercise:5.34}
设$X_1,\ldots,X_n$为来自均值$\mu$、方差$\sigma^2$的总体的随机样本.证明
$$E\frac{\sqrt n(\bar X_n-\mu)}\sigma=0,\qquad\operatorname{Var}\frac{\sqrt n(\bar X_n-\mu)}\sigma=1.$$
因此，中心极限定理中对$\bar X_n$的标准化，得到与极限$n(0,1)$分布具有相同均值、方差的随机变量.
\end{exercise}
\begin{exercise}\label{exercise:5.35}\index{SI0985@Stirling公式}
Stirling公式给出阶乘的近似，已在\autoref{exercise:1.28}推导，也容易由中心极限定理推出.
\begin{enumerate}[label=(\alph*)]
\item\phantomsection\label{itm:exercise:5.35:a} 说明若$X_i\sim\operatorname{exponential}(1)$（$i=1,2,\ldots$）全部独立，则对每个$x$，
$$P\left(\frac{\bar X_n-1}{1/\sqrt n}\le x\right)\to P(Z\le x),$$
其中$Z$为标准正态随机变量.
\item 说明对\hyperref[itm:exercise:5.35:a]{(a)}中近似式两边求导，提示
$$\frac{\sqrt n}{\Gamma(n)}(x\sqrt n+n)^{n-1}e^{-(x\sqrt n+n)}\approx\frac1{\sqrt{2\pi}}e^{-x^2/2},$$
且令$x=0$可得Stirling公式.
\end{enumerate}
\end{exercise}
\begin{exercise}\label{exercise:5.36}
给定$N=n$，$Y$的条件分布为$\chi_{2n}^2$；$N$的无条件分布为$\operatorname{Poisson}(\theta)$.
\begin{enumerate}[label=(\alph*)]\item 计算无条件矩$EY$和$\operatorname{Var}Y$.\item 证明$\theta\to\infty$时，$(Y-EY)/\sqrt{\operatorname{Var}Y}$依分布收敛到$n(0,1)$.\end{enumerate}
\end{exercise}
\begin{exercise}\label{exercise:5.37}
\autoref{ex:5.5.16}给出负二项分布的正态近似.与\autoref{ex:3.3.2}中二项分布的正态近似一样，可用“连续性修正”改进.对\autoref{ex:5.5.16}中的$X_i$，令$V_n=\sum_{i=1}^nX_i$.取$n=10$、$p=.7$、$r=2$，分别用下列三种方法计算$v=0,1,\ldots,10$时的$P(V_n=v)$.
\begin{enumerate}[label=(\alph*)]\item 精确计算.\item \autoref{ex:5.5.16}给出的正态近似.\item 带连续性修正的正态近似.\end{enumerate}
\end{exercise}
\begin{exercise}\label{exercise:5.38}\index{SI0574@矩母函数}
下面对\autoref{exercise:3.45}不等式的推广，可用于证明强大数定律（见\autoref{sec:5.8.4}）.设$X_1,X_2,\ldots,X_n$为iid，矩母函数为$M_X(t)$，$-h<t<h$，令$S_n=\sum_{i=1}^nX_i$、$\bar X_n=S_n/n$.
\begin{enumerate}[label=(\alph*)]
\item 证明$0<t<h$时$P(S_n>a)\le e^{-at}[M_X(t)]^n$；$-h<t<0$时$P(S_n\le a)\le e^{-at}[M_X(t)]^n$.
\item 利用$M_X(0)=1$及$M'_X(0)=E(X)$证明：若$E(X)<0$，则存在$0<c<1$使$P(S_n>a)\le c^n$.为$P(S_n\le a)$建立类似界.
\item\phantomsection\label{itm:exercise:5.38:c} 定义$Y_i=X_i-\mu-\varepsilon$，在上述论证中取$a=0$，证明$P(\bar X_n-\mu>\varepsilon)\le c^n$.
\item\phantomsection\label{itm:exercise:5.38:d} 现在定义$Y_i=-X_i+\mu-\varepsilon$，建立类似\hyperref[itm:exercise:5.38:c]{(c)}的等式，结合两者得到：对某个$0<c<1$，
$$P(|\bar X_n-\mu|>\varepsilon)\le2c^n.$$
\end{enumerate}
\end{exercise}
\begin{exercise}\label{exercise:5.39}
本题及接下来两题研究收敛的若干数学细节.
\begin{enumerate}[label=(\alph*)]
\item 证明\autoref{thm:5.5.4}.（提示：$h$连续，给定$\varepsilon>0$，可找到$\delta$，使$|x_n-x|<\delta$时$|h(x_n)-h(x)|<\varepsilon$.将其转化为概率表述.）
\item 在\autoref{ex:5.5.8}中，找出$Y_i$的一个几乎处处收敛、即逐点收敛的子序列.\footnote{译注：所引例子使用$X_i$，原题这里写成$Y_i$；应指该例的$X_i$序列.}
\end{enumerate}
\end{exercise}
\begin{exercise}\label{exercise:5.40}
在$X_n$与$X$都是连续随机变量的情形，证明\autoref{thm:5.5.12}.
\begin{enumerate}[label=(\alph*)]
\item 给定$t,\varepsilon$，证明$P(X\le t-\varepsilon)\le P(X_n\le t)+P(|X_n-X|\ge\varepsilon)$.这给出$P(X_n\le t)$的下界.
\item 用类似策略求$P(X_n\le t)$的上界.
\item 用夹逼法推出$P(X_n\le t)\to P(X\le t)$.
\end{enumerate}
\end{exercise}
\sourcepageanchor{245}
\begin{exercise}\label{exercise:5.41}
证明\autoref{thm:5.5.13}，即证明\footnote{译注：右端原文写$x\ge\mu$，但依分布收敛到常数只要求在$x\ne\mu$的连续点成立；应将第二行条件改为$x>\mu$，与所引定理一致.}
$$P(|X_n-\mu|>\varepsilon)\to0\text{（对每个$\varepsilon$）}\iff P(X_n\le x)\to\begin{cases}0,&x<\mu,\\1,&x\ge\mu.\end{cases}$$

\begin{enumerate}[label=(\alph*)]
\item 取$\varepsilon=|x-\mu|$，证明$x>\mu$时$P(X_n\le x)\ge P(|X_n-\mu|\le\varepsilon)$；$x<\mu$时$P(X_n\le x)\le P(|X_n-\mu|\ge\varepsilon)$.推出正向蕴含.
\item 利用$\{x:|x-\mu|>\varepsilon\}=\{x:x-\mu<-\varepsilon\}\cup\{x:x-\mu>\varepsilon\}$推出反向蕴含.
\end{enumerate}
（上述结果的详细论述见\hyperref[bib:source33]{Billingsley 1995}，第25节.）
\end{exercise}
\begin{exercise}\label{exercise:5.42}
类似\autoref{ex:5.5.11}，设$X_1,X_2,\ldots$为iid $f$，$X_{(n)}=\max_{1\le i\le n}X_i$.
\begin{enumerate}[label=(\alph*)]\item 若$f$为$\operatorname{beta}(1,\beta)$，求使$n^\nu(1-X_{(n)})$依分布收敛的$\nu$.\item 若$f$为$\operatorname{exponential}(1)$，求使$X_{(n)}-a_n$依分布收敛的序列$a_n$.\end{enumerate}
\end{exercise}
\begin{exercise}\label{exercise:5.43}
补全\autoref{thm:5.5.24}证明的细节.
\begin{enumerate}[label=(\alph*)]\item 证明若$\sqrt n(Y_n-\mu)\to n(0,\sigma^2)$（依分布），则$Y_n\to\mu$（依概率）.\item 给出应用Slutsky定理（\autoref{thm:5.5.17}）的细节.\end{enumerate}
\end{exercise}
\begin{exercise}\label{exercise:5.44}\index{SI0253@极限}
设$X_i$（$i=1,2,\ldots$）为独立的$\operatorname{Bernoulli}(p)$随机变量，$Y_n=n^{-1}\sum_{i=1}^nX_i$.
\begin{enumerate}[label=(\alph*)]
\item 证明$\sqrt n(Y_n-p)\to n[0,p(1-p)]$（依分布）.
\item 证明$p\ne1/2$时，方差估计$Y_n(1-Y_n)$满足$\sqrt n[Y_n(1-Y_n)-p(1-p)]\to n[0,(1-2p)^2p(1-p)]$（依分布）.
\item 证明$p=1/2$时，$n[Y_n(1-Y_n)-1/4]\to-\frac14\chi_1^2$（依分布）.（若觉得奇怪，注意$Y_n(1-Y_n)\le1/4$，所以左端总为负.等价形式为$4n[1/4-Y_n(1-Y_n)]\to\chi_1^2$.）
\end{enumerate}
\end{exercise}
\begin{exercise}\label{exercise:5.45}
在\autoref{ex:5.6.1}的情形中，计算至少75\%的元件寿命达到150小时的概率，分别取：
\begin{enumerate}[label=(\alph*)]\item\phantomsection\label{itm:exercise:5.45:a} $c=300$，$X\sim\operatorname{gamma}(a,b)$，$a=4,b=5$.\item\phantomsection\label{itm:exercise:5.45:b} $c=100$，$X\sim\operatorname{gamma}(a,b)$，$a=20,b=5$.\item\phantomsection\label{itm:exercise:5.45:c} $c=100$，$X\sim\operatorname{gamma}(a,b)$，$a=20.7,b=5$.\end{enumerate}
提示：\hyperref[itm:exercise:5.45:a]{(a)}、\hyperref[itm:exercise:5.45:b]{(b)}的gamma积分可求闭式，尽管\hyperref[itm:exercise:5.45:b]{(b)}可能不值得这样费力.\hyperref[itm:exercise:5.45:c]{(c)}的积分没有闭式，须用数值积分或模拟计算.
\end{exercise}
\begin{exercise}\label{exercise:5.46}
将\autoref{exercise:5.45}的答案与二项分布正态近似（见\autoref{ex:3.3.2}）所得结果比较.
\end{exercise}
\begin{exercise}\label{exercise:5.47}
验证\eqref{eq:5.6.5}中各随机变量的分布.
\end{exercise}
\begin{exercise}\label{exercise:5.48}
用类似\eqref{eq:5.6.5}的策略说明，如何生成$m,n$均为偶数时的$F_{m,n}$随机变量.
\end{exercise}
\begin{exercise}\label{exercise:5.49}
设$U\sim\operatorname{uniform}(0,1)$.
\begin{enumerate}[label=(\alph*)]\item 证明$-\log U$与$-\log(1-U)$都是指数随机变量.\item 证明$X=\log\frac u{1-u}$是$\operatorname{logistic}(0,1)$随机变量.\item 说明如何生成$\operatorname{logistic}(\mu,\beta)$随机变量.\end{enumerate}
\end{exercise}
\begin{exercise}\label{exercise:5.50}\index{SI0861@随机数生成}
生成正态伪随机变量的Box–Muller方法（\autoref{ex:5.6.4}）基于变换
$$X_1=\cos(2\pi U_1)\sqrt{-2\log U_2},\qquad X_2=\sin(2\pi U_1)\sqrt{-2\log U_2},$$
其中$U_1,U_2$为iid $\operatorname{uniform}(0,1)$.证明$X_1,X_2$是独立的$n(0,1)$随机变量.
\end{exercise}
\sourcepageanchor{246}
\begin{exercise}\label{exercise:5.51}
从均匀随机变量生成伪标准正态随机变量的一种较早方法（并非较好的方法），是取$X=\sum_{i=1}^{12}U_i-6$，其中$U_i$为iid $\operatorname{uniform}(0,1)$.
\begin{enumerate}[label=(\alph*)]\item 说明$X$为何近似服从$n(0,1)$.\item 能否找出这种近似明显失效之处？\item 比较前四阶矩，说明近似有多好（或多差）.（四阶矩为$29/10$，计算较长；矩母函数和计算机代数有帮助，见\autoref{ex:12.6.6}.）\end{enumerate}
\end{exercise}
\begin{exercise}\label{exercise:5.52}\index{SI0861@随机数生成}
对以下各个分布，写出生成相应随机变量的算法.
\begin{enumerate}[label=(\alph*)]\item $Y\sim\operatorname{binomial}(8,2/3)$.\item $Y\sim\operatorname{hypergeometric}(N=10,M=8,K=4)$.\item $Y\sim\operatorname{negative\ binomial}(5,1/3)$.\end{enumerate}
\end{exercise}
\begin{exercise}\label{exercise:5.53}
对上一题中的每个分布：
\begin{enumerate}[label=(\alph*)]\item 生成1,000个服从该分布的变量.\item 将生成变量的均值、方差及直方图与理论值比较.\end{enumerate}
\end{exercise}
\begin{exercise}\label{exercise:5.54}
参见\autoref{ex:5.6.6}，另一份海湾{\upshape\CJKfontspec{SimSun}鳀}鱼幼体计数样本给出数据
$$158,143,106,57,97,80,109,109,350,224,109,214,84.$$
\begin{enumerate}[label=(\alph*)]
\item 用\autoref{ex:5.6.6}的方法构造$S^2$的模拟分布，判断Poisson计数假设是否合理.
\item Poisson假设失效（及方差增大）的一种解释，是幼体在河中均匀分布的假设不成立.若幼体倾向于聚集，则负二项$(r,p)$分布（均值$\mu=r(1-p)/p$，方差$\mu+\mu^2/r$）是合理的替代模型.取$\mu=\bar x$，哪些$r$值产生的模拟分布与数据相符？
\end{enumerate}
\end{exercise}
\begin{exercise}\label{exercise:5.55}
设用\eqref{eq:5.6.7}的方法生成随机变量$Y$，其中$y_i=i$，$i=0,1,2,\ldots$.证明比较次数的期望为$E(Y+1)$.（提示：\autoref{exercise:2.14}.）
\end{exercise}
\begin{exercise}\label{exercise:5.56}
设$Y$服从Cauchy分布，$f_Y(y)=1/(1+y^2)$，$-\infty<y<\infty$.
\begin{enumerate}[label=(\alph*)]\item 证明$F_Y(y)=\tan^{-1}(y)$.\item 说明如何从$\operatorname{uniform}(0,1)$随机变量模拟$\operatorname{Cauchy}(a,b)$随机变量.\end{enumerate}
（相关结果见\autoref{exercise:2.12}.）\footnote{译注：原题的pdf漏掉归一化因子$1/\pi$，cdf也漏掉缩放和加性常数.标准Cauchy的pdf应为$1/[\pi(1+y^2)]$，cdf为$1/2+\arctan(y)/\pi$.}
\end{exercise}
\begin{exercise}\label{exercise:5.57}\index{SI0215@二元}
\hyperref[bib:source239]{Park等（1996）}给出如下生成相关二元变量的方法.设$X_1,X_2,X_3$为均值分别为$\lambda_1,\lambda_2,\lambda_3$的独立Poisson随机变量，构造$Y_1=X_1+X_3$及$Y_2=X_2+X_3$.
\begin{enumerate}[label=(\alph*)]
\item 证明$\operatorname{Cov}(Y_1,Y_2)=\lambda_3$.
\item 定义$Z_i=I(Y_i=0)$、$p_i=e^{-(\lambda_i+\lambda_3)}$.证明$Z_i\sim\operatorname{Bernoulli}(p_i)$，且
$$\operatorname{Corr}(Z_1,Z_2)=\frac{p_1p_2(e^{\lambda_3}-1)}{\sqrt{p_1(1-p_1)}\sqrt{p_2(1-p_2)}}.$$
\item 证明$Z_1,Z_2$的相关系数不能在$[-1,1]$中任意取值，而满足
$$\operatorname{Corr}(Z_1,Z_2)\le\min\left\{\sqrt{\frac{p_2(1-p_1)}{p_1(1-p_2)}},\sqrt{\frac{p_1(1-p_2)}{p_2(1-p_1)}}\right\}.$$
\end{enumerate}
\end{exercise}
\sourcepageanchor{247}
\begin{exercise}\label{exercise:5.58}
设$U_1,U_2,\ldots,U_n$为iid $\operatorname{uniform}(0,1)$随机变量，$S_n=\sum_{i=1}^nU_i$.定义$N=\min\{k:S_k>1\}$.
\begin{enumerate}[label=(\alph*)]
\item 证明$P(S_k\le t)=t^k/k!$.\footnote{译注：这一公式需要$0\le t\le1$；原题未写适用范围.后续使用$t=1$，在该范围内.}
\item\phantomsection\label{itm:exercise:5.58:b} 证明$P(N=n)=P(S_{n-1}<1)-P(S_n<1)$，并证明令人惊讶的$E(N)=e$，其中$e$为自然对数的底.
\item 利用\hyperref[itm:exercise:5.58:b]{(b)}的结果，通过模拟计算$e$.
\item 若希望有95\%的把握保证$e$的前四位数字正确，$n$须多大？
\end{enumerate}
（\hyperref[bib:source270]{Russell 1991}描述过这种模拟实验，并将问题归于\hyperref[bib:source151]{Gnedenko 1978}.）
\end{exercise}
\begin{exercise}\label{exercise:5.59}
证明\autoref{ex:5.6.7}的算法生成$\operatorname{beta}(a,b)$随机变量.
\end{exercise}
\begin{exercise}\label{exercise:5.60}
把\autoref{ex:5.6.7}的算法推广到$[0,1]$上的任意连续pdf，即对$[a,b]$上的任意有界pdf $f(x)$，定义$c=\max_{a\le x\le b}f(x)$.设$X,Y$独立，$X\sim\operatorname{uniform}(a,b)$、$Y\sim\operatorname{uniform}(0,c)$.取$d>b$，定义
$$W=\begin{cases}X,&Y<f(X),\\d,&Y\ge f(X).\end{cases}$$
\begin{enumerate}[label=(\alph*)]\item\phantomsection\label{itm:exercise:5.60:a} 证明$a\le w\le b$时，$P(W\le w)=\int_a^w f(t)\,dt/[c(b-a)]$.\item 利用\hyperref[itm:exercise:5.60:a]{(a)}说明如何生成pdf为$f(x)$的随机变量.（提示：用几何论证，画图会有帮助.）\end{enumerate}
\end{exercise}
\begin{exercise}\label{exercise:5.61}\index{SI0255@生成}\index{SI0307@生成}\leavevmode
\begin{enumerate}[label=(\alph*)]
\item\phantomsection\label{itm:exercise:5.61:a} 要生成$Y\sim\operatorname{beta}(a,b)$，其中$a,b$不是整数.证明取$V\sim\operatorname{beta}([a],[b])$会使$M=\sup_y f_Y(y)/f_V(y)$有限.
\item\phantomsection\label{itm:exercise:5.61:b} 要生成$Y\sim\operatorname{gamma}(a,b)$，其中$a,b$不是整数.证明取$V\sim\operatorname{gamma}([a],b)$会使$M=\sup_y f_Y(y)/f_V(y)$有限.
\item 证明在\hyperref[itm:exercise:5.61:a]{(a)}、\hyperref[itm:exercise:5.61:b]{(b)}中，若$V$的参数为$[a]+1$，则$M$为无穷.
\item 在\hyperref[itm:exercise:5.61:a]{(a)}、\hyperref[itm:exercise:5.61:b]{(b)}中分别求使$E(N)$最小的$V$的最优参数（见\eqref{eq:5.6.12}）.
\end{enumerate}
（回顾$[a]$是不超过$a$的最大整数.）\footnote{译注：候选beta、gamma分布的形状参数必须为正；因此\hyperref[itm:exercise:5.61:a]{(a)}需$[a],[b]\ge1$，\hyperref[itm:exercise:5.61:b]{(b)}需$[a]\ge1$.原题未明说这些条件；参数在$(0,1)$时，直接向下取整会得到无效的零形状参数.}
\end{exercise}
\begin{exercise}\label{exercise:5.62}
求使接受／拒绝算法可以用$U\sim\operatorname{uniform}(0,1)$和下列$V$生成$Y\sim n(0,1)$的$M$.
\begin{enumerate}[label=(\alph*)]\item $V\sim\operatorname{Cauchy}$.\item $V$服从双指数分布.\item 比较两种算法.你推荐哪一种？\end{enumerate}
\end{exercise}
\begin{exercise}\label{exercise:5.63}
用接受／拒绝算法生成$Y\sim n(0,1)$时，可生成均匀$U$及$V\sim\operatorname{exponential}(\lambda)$，再给$V$附上随机符号（正负各等概率）.什么$\lambda$使算法最优？
\end{exercise}
\begin{exercise}\label{exercise:5.64}\index{SI0014@重要性抽样}
与接受／拒绝方法类似的一种技术是重要性抽样（importance sampling），它很适合计算分布的特征.设$X\sim f$，但难以从pdf $f$模拟.从$g$生成iid $Y_1,Y_2,\ldots,Y_m$，对任意函数$h$计算$m^{-1}\sum_{i=1}^m[f(Y_i)/g(Y_i)]h(Y_i)$.假定$f,g$支撑集相同，且$\operatorname{Var}h(X)<\infty$.
\begin{enumerate}[label=(\alph*)]
\item\phantomsection\label{itm:exercise:5.64:a} 证明$E\bigl(m^{-1}\sum_{i=1}^m[f(Y_i)/g(Y_i)]h(Y_i)\bigr)=Eh(X)$.
\item 证明$m^{-1}\sum_{i=1}^m[f(Y_i)/g(Y_i)]h(Y_i)\to Eh(X)$（依概率）.
\item 虽然\hyperref[itm:exercise:5.64:a]{(a)}的估计量有正确期望，实际中更常用
$$\sum_{i=1}^m\left(\frac{f(Y_i)/g(Y_i)}{\sum_{j=1}^m f(Y_j)/g(Y_j)}\right)h(Y_i).$$
证明它依概率收敛到$Eh(X)$.此外，证明若$h$是常数，这一估计量优于\hyperref[itm:exercise:5.64:a]{(a)}中的估计量.（\hyperref[bib:source70]{Casella和Robert 1996}进一步研究了其性质.）
\end{enumerate}
\end{exercise}
\sourcepageanchor{248}
\begin{exercise}\label{exercise:5.65}
\autoref{exercise:5.64}的重要性抽样算法的一种变体，实际上可产生来自$f$的近似样本.仍设$X\sim f$，从$g$生成iid $Y_1,Y_2,\ldots,Y_m$，计算$q_i=[f(Y_i)/g(Y_i)]/[\sum_{j=1}^mf(Y_j)/g(Y_j)]$.然后从取值$Y_1,Y_2,\ldots,Y_m$、概率$P(X^*=Y_k)=q_k$的离散分布生成随机变量$X^*$.证明$X_1^*,X_2^*,\ldots,X_r^*$近似为来自$f$的随机样本.

提示：证明$P(X^*\le x)=\sum_{i=1}^m q_iI(Y_i\le x)$；令$m\to\infty$，在分子、分母中使用弱大数定律.

\hyperref[bib:source266]{Rubin（1988）}称此法为抽样／重要性再抽样（Sampling/Importance Resampling，SIR）算法；\hyperref[bib:source292]{Smith和Gelfand（1992）}称其为加权自助法（weighted bootstrap）.
\end{exercise}
\begin{exercise}\label{exercise:5.66}
若$X_1,\ldots,X_n$为iid $n(\mu,\sigma^2)$，已知样本均值$\bar X$服从$n(\mu,\sigma^2/n)$.若想使用更稳健的位置估计量，例如中位数\eqref{eq:5.4.1}，其分布就更难推导.
\begin{enumerate}[label=(\alph*)]
\item 证明$m$是$X_i$的中位数，当且仅当$(m-\mu)/\sigma$是$(X_i-\mu)/\sigma$的中位数.所以只须考虑$n(0,1)$样本中位数的分布.
\item\phantomsection\label{itm:exercise:5.66:b} 对来自$n(0,1)$、容量$n=15$的样本，模拟中位数$m$的分布.
\item 将\hyperref[itm:exercise:5.66:b]{(b)}的分布与中位数的渐近分布$\sqrt n(m-\mu)\sim n[0,1/(4f^2(0))]$比较，其中$f$为pdf.$n=15$是否足以使渐近结果有效？
\end{enumerate}
\end{exercise}
\begin{exercise}\label{exercise:5.67}
很多情况下Metropolis算法是首选，因为：(i)没有明显满足接受／拒绝上确界条件的候选密度；或(ii)难以验证上确界条件；或(iii)懒惰促使我们用计算力代替脑力.

对以下各情形，说明如何实施Metropolis算法，从指定分布生成容量为100的样本.
\begin{enumerate}[label=(\alph*)]
\item $X\sim\sigma^{-1}f[(x-\mu)/\sigma]$，$f$为自由度$\nu$的Student $t$分布，$\nu,\mu,\sigma$已知.
\item $X\sim\operatorname{lognormal}(\mu,\sigma^2)$，$\mu,\sigma^2$已知.
\item $X\sim\operatorname{Weibull}(\alpha,\beta)$，$\alpha,\beta$已知.
\end{enumerate}
\end{exercise}
\begin{exercise}\label{exercise:5.68}
用Metropolis算法代替接受／拒绝算法，就不必验证后者的上确界条件.当然，也放弃了得到精确服从所需分布的随机变量的性质，只能接受近似.
\begin{enumerate}[label=(\alph*)]
\item 说明如何从$n(0,1)$随机变量出发，用Metropolis算法生成近似服从自由度$\nu$的Student $t$分布的随机变量.
\item 说明如何从Cauchy随机变量出发，用接受／拒绝算法生成自由度$\nu$的Student $t$随机变量.
\item 说明如何用变换直接生成自由度$\nu$的Student $t$随机变量.
\item 对$\nu=2,10,25$，生成容量100的样本来比较各方法.你更喜欢哪一种？为什么？
\end{enumerate}
[\hyperref[bib:source218]{Mengersen和Tweedie 1995}\footnote{译注：正文此处引1995年，原书书后同名论文条目[218]列1996年；这里分别保留两处年份，链接至该条目.}表明，若满足上确界条件，即$\sup f/g<\!\le M<\infty$（$f$为目标密度，$g$为候选密度），Metropolis算法收敛快得多.\footnote{译注：原文此处连写“$<\!\le$”；结合前文上确界条件，可理解为$\sup f/g\le M<\infty$.}]
\end{exercise}
\begin{exercise}\label{exercise:5.69}
证明pdf $f_Y(y)$是Metropolis算法的稳定点（stable point）.即若$Z_i\sim f_Y(y)$，则$Z_{i+1}\sim f_Y(y)$.
\end{exercise}



% 原书印刷217–220页，PDF物理248–251页.
\subsection{杂录}\label{sec:5.8}
\subsubsection{中心极限定理补充}\label{sec:5.8.1}\index{SI0102@中心极限定理}\index{SI0600@Lévy定理}\index{SI0643@Lindeberg–Feller条件}
对于iid随机变量序列，收敛到正态分布的充分必要条件已经知道，其中最著名的结果或许属于Lindeberg与Feller.以下特例出自Lévy.设$X_1,X_2,\ldots$为iid序列，$EX_i=\mu<\infty$，令$V_n=\sum_{i=1}^nX_i$.
\sourcepageanchor{249}
序列$V_n$经适当标准化后收敛到$n(0,1)$随机变量，当且仅当
$$\lim_{t\to\infty}\frac{t^2P(|X_1-\mu|>t)}{E((X_1-\mu)^2I_{[-t,t]}(X_1-\mu))}=0.$$
注意这是一个方差条件.它虽不完全要求方差有限，却要求方差“几乎”有限.这是收敛到正态分布的重要一点：正态性来自小扰动之和.

还有大量其他类型的中心极限定理，尤其是旨在放宽独立性假设的结果.虽然不能完全取消这一假设，但可以使其限制更弱（见\hyperref[bib:source33]{Billingsley 1995}，第27节；或\hyperref[bib:source251]{Resnick 1999}，第8章）.
\subsubsection{\texorpdfstring{$S^2$}{S²}的偏差}\label{sec:5.8.2}\index{SI0937@偏差}
本章多数计算假定观测相互独立，一些期望的计算也依赖这一事实.H. A. \hyperref[bib:source92]{David（1985）}指出，若观测相关，$S^2$可能是$\sigma^2$的有偏估计，即$ES^2=\sigma^2$未必成立.但可能偏差的范围很容易计算.若随机变量$X_1,\ldots,X_n$（不一定独立）都具有均值$\mu$、方差$\sigma^2$，则
$$(n-1)ES^2=E\left(\sum_{i=1}^n(X_i-\mu)^2-n(\bar X-\mu)^2\right)=n\sigma^2-n\operatorname{Var}\bar X.$$
依赖的程度与类型不同时，$\operatorname{Var}\bar X$可从0（所有变量都为常数时）变化到$\sigma^2$（所有变量都是$X_1$的副本时）.\footnote{译注：若固定$\sigma^2>0$，所有$X_i$为常数不符合其方差均为$\sigma^2$的假设.方差下界0可由使样本和恒定的负相关构造实现，例如$n=2$时$X_2=2\mu-X_1$；原文的括号说明应作此区分.}把这些值代入上式，得到存在依赖时$ES^2$的范围：
$$0\le ES^2\le\frac n{n-1}\sigma^2.$$
\subsubsection{再论Chebychev不等式}\label{sec:5.8.3}\index{SI0565@切比雪夫}
\autoref{sec:3.6}讨论过Chebychev不等式（另见\autoref{sec:3.8.2}），\autoref{ex:3.6.2}给出一个特别有用的形式.该形式仍要求知道随机变量的均值与方差，但有时我们想用均值、方差的估计值来给出界.

若$X_1,\ldots,X_n$为来自均值$\mu$、方差$\sigma^2$的总体的随机样本，Chebychev不等式为
$$P(|X-\mu|\ge k\sigma)\le\frac1{k^2}.$$
\hyperref[bib:source273]{Saw等（1984）}证明，用$\bar X$代替$\mu$、$S^2$代替$\sigma^2$，可得到
$$P(|X-\bar X|\ge kS)\le\frac1{n+1}g\left(\frac{n(n+1)k^2}{n-1+(n+1)k^2}\right),$$
其中
$$g(t)=\begin{cases}\nu,&\nu\text{为偶数},\\\nu,&\nu\text{为奇数且$t<a$},\\\nu-1,&\nu\text{为奇数且$t>a$},\end{cases}$$
且
$$\nu=\text{小于}\ \frac{n+1}t\ \text{的最大整数},\qquad a=\frac{(n+1)(n+1-\nu)}{1+\nu(n+1-\nu)}.$$
\sourcepageanchor{250}
\subsubsection{强大数定律补充}\label{sec:5.8.4}\index{SI0202@几乎处处}\index{SI0995@强大数定律}
如前所述，强大数定律\autoref{thm:5.5.9}可在较弱条件下证明，只要求随机变量的均值有限（例如见\hyperref[bib:source251]{Resnick 1999}，第7章；或\hyperref[bib:source33]{Billingsley 1995}，第22节）.不过，在矩母函数存在的假设下，\hyperref[bib:source191]{Koopmans（1993）}给出了只使用微积分的证明.

强大数定律所断言的收敛是我们最熟悉的一种，即随机变量序列$\bar X_n$逐点收敛到共同均值$\mu$.由\autoref{ex:5.5.8}可知，它比弱大数定律的依概率收敛更强.

强大数定律的结论是
$$P\left(\lim_{n\to\infty}|\bar X_n-\mu|<\varepsilon\right)=1,$$
即以概率1，序列$\{\bar X_n\}$的极限为$\mu$.换言之，序列发散的集合概率为0.序列要发散，必存在$\delta>0$，使每个$n$都有某个$k>n$满足$|\bar X_k-\mu|>\delta$.满足此条件的全体$\bar X_k$构成发散序列，可用集合
$$A_\delta=\bigcap_{n=1}^\infty\bigcup_{k=n}^\infty\{|\bar X_k-\mu|>\delta\}$$
表示.删去交集项，就得到$P(A_\delta)$的上界，因此序列$\{\bar X_n\}$发散的集合概率由下式控制：
\begin{align*}
P(A_\delta)&\le P\left(\bigcup_{k=n}^\infty\{|\bar X_k-\mu|>\delta\}\right)\\
&\le\sum_{k=n}^\infty P(\{|\bar X_k-\mu|>\delta\})&&\text{（Boole不等式，\autoref{thm:1.2.11}）}\\
&\le2\sum_{k=n}^\infty c^k,\qquad0<c<1,
\end{align*}
最后一个不等式可用\autoref{exercise:5.38}\hyperref[itm:exercise:5.38:d]{(d)}证明.注意这里求的是几何级数之和，由\eqref{eq:1.5.4}得到
$$P(A_\delta)\le2\sum_{k=n}^\infty c^k=2\frac{c^n}{1-c}\to0,\qquad n\to\infty.$$
因此，$\{\bar X_n\}$发散的集合概率为零，强大数定律得证.
\subsubsection{马尔可夫链Monte Carlo}\label{sec:5.8.5}\index{SI0017@马尔可夫链蒙特卡洛}\index{SI0666@马尔可夫链蒙特卡洛}
统称为马尔可夫链Monte Carlo（MCMC）的方法用于生成随机变量，在复杂计算中非常有用，尤其是涉及积分和最大化的计算.Metropolis算法（见\autoref{sec:5.6}）就是MCMC方法的例子.

顾名思义，这些方法基于马尔可夫链（Markov chains），这是我们尚未研究的概率结构（入门介绍见\hyperref[bib:source78]{Chung 1974}或\hyperref[bib:source259]{Ross 1988）}.若随机变量序列$X_1,X_2,\ldots$满足
$$P(X_{k+1}\in A\mid X_1,\ldots,X_k)=P(X_{k+1}\in A\mid X_k),$$
则称其为马尔可夫链；即当前随机变量的分布至多依赖紧邻的前一个随机变量.注意这是独立性的推广.

遍历定理（Ergodic Theorem）是大数定律的推广，\index{SI0203@遍历定理}\semanticindexanchor{203:250}\index{SI0406@遍历定理}\semanticindexanchor{406:250}它说明若马尔可夫链$X_1,X_2,\ldots$满足某些正则条件（统计问题中常常满足），则只要期望存在，
$$\frac1n\sum_{i=1}^n h(X_i)\to Eh(X),\qquad n\to\infty.$$
\sourcepageanchor{251}
所以，\autoref{sec:5.6}的计算可推广到马尔可夫链和MCMC方法.

要充分理解MCMC方法，确实需要了解更多马尔可夫链知识，这里不展开.MCMC的理论及应用已有大量文献.\hyperref[bib:source312]{Tanner（1993）}\footnote{译注：正文此处引1993年，原书书后同名著作条目[312]列1996年；这里分别保留两处年份，链接至该条目.}提供了统计计算方法的良好入门介绍；\hyperref[bib:source255]{Robert（1994}，第9章）也作了介绍，偏重理论且带有Bayes色彩.\hyperref[bib:source67]{Casella和George（1992）}通过Gibbs抽样器（一种特定MCMC方法）给出了更易读的入门介绍.Gibbs抽样器或许仍是应用最广的MCMC方法，也是这类方法流行的原因之一［源于\hyperref[bib:source134]{Gelfand和Smith（1990）}在\hyperref[bib:source137]{Geman和Geman（1984）}基础上的奠基性工作］.涉及MCMC的参考文献实在太多.其他入门文献包括\hyperref[bib:source136]{Gelman和Rubin（1992）}、\hyperref[bib:source138]{Geyer和Thompson（1992）}、\hyperref[bib:source293]{Smith和Roberts（1993）}的论文；\hyperref[bib:source315]{Tierney（1994）}给出了尤其优美的理论介绍.\hyperref[bib:source256]{Robert和Casella（1999）}则以一整本教材系统讨论这一领域.





\section{数据约简原则}\label{sec:6}
\chapterdoi{6}

\sourcepageanchor{252}
\begin{flushright}
\begin{minipage}{.90\textwidth}
……我们正饱受种种揣测、猜想和假设之苦.困难在于把事实——绝对不可否认的事实——的骨架，从理论家和报道者添上的装饰中剥离出来.
\begin{flushright}夏洛克·福尔摩斯\\《银色马》\end{flushright}
\end{minipage}
\end{flushright}



\subsection{引言}\label{sec:6.1}
实验者用样本$X_1,\ldots,X_n$中的信息推断未知参数$\theta$.若样本容量$n$很大，观测样本$x_1,\ldots,x_n$就是一长串难以解释的数字.实验者可能希望确定样本值的几个关键特征，来概括样本中的信息.通常通过计算统计量，即样本的函数，来完成这一点.例如，样本均值、样本方差、最大观测值和最小观测值，是可以用来概括样本某些关键特征的四个统计量.回顾我们用粗体字母表示多元变量，因而$\boldsymbol X$表示随机变量$X_1,\ldots,X_n$，$\boldsymbol x$表示样本$x_1,\ldots,x_n$.

任何统计量$T(\boldsymbol X)$都确定一种数据约简（data reduction）或数据概括的形式.若实验者只使用统计量的观测值$T(\boldsymbol x)$，那么对满足$T(\boldsymbol x)=T(\boldsymbol y)$的两个样本$\boldsymbol x$和$\boldsymbol y$，他就会等同看待，即使实际样本值在某些方面可能不同.

以某个特定统计量进行的数据约简，可以看作样本空间$\mathcal X$的一个划分.\index{SI0753@划分}\semanticindexanchor{753:252}令$\mathcal T=\{t:t=T(\boldsymbol x)\text{，其中某个}\boldsymbol x\in\mathcal X\}$为$\mathcal X$在$T(\boldsymbol x)$下的像.那么$T(\boldsymbol x)$把样本空间划分为集合$A_t$、$t\in\mathcal T$，其定义为$A_t=\{\boldsymbol x:T(\boldsymbol x)=t\}$.统计量概括数据的含义是：它只报告$T(\boldsymbol x)=t$，或等价地报告$\boldsymbol x\in A_t$，而不报告整个样本$\boldsymbol x$.例如，若$T(\boldsymbol x)=x_1+\cdots+x_n$，则$T(\boldsymbol x)$只报告总和，不报告实际样本值.可能有许多不同的样本点具有相同的总和.这种数据约简的优点和后果是本章讨论的主题.

我们研究三种数据约简原则.我们关心既不丢弃有关未知参数$\theta$的重要信息，又能成功去掉与认识$\theta$无关的信息的约简方法.充分性原则（Sufficiency Principle）提出一种在概括数据的同时不丢弃关于$\theta$的信息的方法.\index{SI0997@充分性原则}\semanticindexanchor{997:252}似然原则（Likelihood Principle）描述一个由观测样本确定的参数函数，这个函数含有样本能提供的全部关于$\theta$的信息.等变性原则（Equivariance Principle）则规定另一种仍保留模型某些重要特征的数据约简方法.



\subsection{充分性原则}\label{sec:6.2}
参数$\theta$的充分统计量（sufficient statistic），是在某种意义上捕获样本中全部关于$\theta$的信息的统计量.除充分统计量的值之外，样本中的其他信息不再含有更多关于$\theta$的信息.这些考虑引出称为充分性原则的数据约简技术.
\sourcepageanchor{253}
\begin{quote}
\textbf{充分性原则：}若$T(\boldsymbol X)$是$\theta$的充分统计量，则关于$\theta$的任何推断都应只通过$T(\boldsymbol X)$的值依赖样本$\boldsymbol X$.即若$\boldsymbol x$与$\boldsymbol y$是满足$T(\boldsymbol x)=T(\boldsymbol y)$的两个样本点，那么无论观测到$\boldsymbol X=\boldsymbol x$还是$\boldsymbol X=\boldsymbol y$，关于$\theta$的推断都应相同.
\end{quote}
本节考察充分统计量和充分性原则的若干方面.
\subsubsection{充分统计量}\label{sec:6.2.1}
充分统计量的形式定义如下.
\begin{definition}\label{def:6.2.1}\index{SI1000@充分统计量}
若给定$T(\boldsymbol X)$的值后，样本$\boldsymbol X$的条件分布不依赖$\theta$，则统计量$T(\boldsymbol X)$是$\theta$的充分统计量.
\end{definition}
若$T(\boldsymbol X)$有连续分布，则对所有$t$都有$P_\theta(T(\boldsymbol X)=t)=0$.此时要充分理解\autoref{def:6.2.1}，需要比\autoref{sec:1}中介绍的条件概率更精细的概念.可参阅\hyperref[bib:source198]{Lehmann（1986）}等更高阶教材中的讨论.我们将在离散情形下进行计算，并指出连续情形中成立的类似结果.

为理解\autoref{def:6.2.1}，设$t$是$T(\boldsymbol X)$的一个可能值，即$P_\theta(T(\boldsymbol X)=t)>0$.我们要考虑条件概率$P_\theta(\boldsymbol X=\boldsymbol x\mid T(\boldsymbol X)=t)$.若样本点$\boldsymbol x$满足$T(\boldsymbol x)\ne t$，显然$P_\theta(\boldsymbol X=\boldsymbol x\mid T(\boldsymbol X)=t)=0$.因此我们关心$P(\boldsymbol X=\boldsymbol x\mid T(\boldsymbol X)=T(\boldsymbol x))$.按定义，若$T(\boldsymbol X)$是充分统计量，这个条件概率对所有$\theta$都相同，所以省略了下标.

充分统计量在以下意义上捕获全部关于$\theta$的信息.考虑实验者1，他观测到$\boldsymbol X=\boldsymbol x$，自然也能计算$T(\boldsymbol X)=T(\boldsymbol x)$.为了推断$\theta$，他可以利用$\boldsymbol X=\boldsymbol x$和$T(\boldsymbol X)=T(\boldsymbol x)$的信息.再考虑实验者2，他不知道$\boldsymbol X$的值，只知道$T(\boldsymbol X)=T(\boldsymbol x)$.实验者2知道$P(\boldsymbol X=\boldsymbol y\mid T(\boldsymbol X)=T(\boldsymbol x))$，它是$A_{T(\boldsymbol x)}=\{\boldsymbol y:T(\boldsymbol y)=T(\boldsymbol x)\}$上的概率分布，因为无须知道$\theta$的真值就能根据模型计算它.因此，实验者2可以用这个分布以及随机数表之类的随机化工具，产生一个观测$\boldsymbol Y$，满足
$$P(\boldsymbol Y=\boldsymbol y\mid T(\boldsymbol X)=T(\boldsymbol x))=P(\boldsymbol X=\boldsymbol y\mid T(\boldsymbol X)=T(\boldsymbol x)).$$
下面会看到，对每个$\theta$，$\boldsymbol X$和$\boldsymbol Y$有相同的无条件概率分布.于是，知道$\boldsymbol X$的实验者1与知道$\boldsymbol Y$的实验者2，具有等价的关于$\theta$的信息.但用随机数表生成$\boldsymbol Y$显然没有增加实验者2对$\theta$的了解.他关于$\theta$的全部知识都包含在$T(\boldsymbol X)=T(\boldsymbol x)$这个事实中.所以，只知道$T(\boldsymbol X)=T(\boldsymbol x)$的实验者2，与知道整个样本$\boldsymbol X=\boldsymbol x$的实验者1，拥有同样多的关于$\theta$的信息.

为完成上述论证，要说明$\boldsymbol X$和$\boldsymbol Y$的无条件分布相同，即对所有$\boldsymbol x$与$\theta$有$P_\theta(\boldsymbol X=\boldsymbol x)=P_\theta(\boldsymbol Y=\boldsymbol x)$.注意事件$\{\boldsymbol X=\boldsymbol x\}$和$\{\boldsymbol Y=\boldsymbol x\}$都是事件$\{T(\boldsymbol X)=T(\boldsymbol x)\}$的子集.还要记住
$$P(\boldsymbol X=\boldsymbol x\mid T(\boldsymbol X)=T(\boldsymbol x))=P(\boldsymbol Y=\boldsymbol x\mid T(\boldsymbol X)=T(\boldsymbol x))$$
且这些条件概率不依赖$\theta$.
\sourcepageanchor{254}
因此，有
\begin{align*}
P_\theta(\boldsymbol X=\boldsymbol x)&=P_\theta(\boldsymbol X=\boldsymbol x\text{且}T(\boldsymbol X)=T(\boldsymbol x))\\
&=P(\boldsymbol X=\boldsymbol x\mid T(\boldsymbol X)=T(\boldsymbol x))P_\theta(T(\boldsymbol X)=T(\boldsymbol x))\quad\text{（条件概率定义）}\\
&=P(\boldsymbol Y=\boldsymbol x\mid T(\boldsymbol X)=T(\boldsymbol x))P_\theta(T(\boldsymbol X)=T(\boldsymbol x))\\
&=P_\theta(\boldsymbol Y=\boldsymbol x\text{且}T(\boldsymbol X)=T(\boldsymbol x))\\
&=P_\theta(\boldsymbol Y=\boldsymbol x).
\end{align*}
用\autoref{def:6.2.1}检验统计量$T(\boldsymbol X)$对$\theta$充分，必须验证：对任意固定的$\boldsymbol x$和$t$，条件概率$P_\theta(\boldsymbol X=\boldsymbol x\mid T(\boldsymbol X)=t)$对所有$\theta$相同.若$T(\boldsymbol x)\ne t$，这个概率对所有$\theta$都是0.故只须验证$P_\theta(\boldsymbol X=\boldsymbol x\mid T(\boldsymbol X)=T(\boldsymbol x))$不依赖$\theta$.由于$\{\boldsymbol X=\boldsymbol x\}$是$\{T(\boldsymbol X)=T(\boldsymbol x)\}$的子集，
\begin{align*}
P_\theta(\boldsymbol X=\boldsymbol x\mid T(\boldsymbol X)=T(\boldsymbol x))
&=\frac{P_\theta(\boldsymbol X=\boldsymbol x\text{且}T(\boldsymbol X)=T(\boldsymbol x))}{P_\theta(T(\boldsymbol X)=T(\boldsymbol x))}\\
&=\frac{P_\theta(\boldsymbol X=\boldsymbol x)}{P_\theta(T(\boldsymbol X)=T(\boldsymbol x))}\\
&=\frac{p(\boldsymbol x\mid\theta)}{q(T(\boldsymbol x)\mid\theta)},
\end{align*}
其中$p(\boldsymbol x\mid\theta)$是样本$\boldsymbol X$的联合pmf，$q(t\mid\theta)$是$T(\boldsymbol X)$的pmf.所以，$T(\boldsymbol X)$对$\theta$充分，当且仅当对每个$\boldsymbol x$，以上pmf之比作为$\theta$的函数是常数.若$\boldsymbol X$和$T(\boldsymbol X)$具有连续分布，则上述条件概率不能按\autoref{sec:1}的意义解释.但仍可用上述准则判断$T(\boldsymbol X)$是否对$\theta$充分.
\begin{theorem}\label{thm:6.2.2}
设$p(\boldsymbol x\mid\theta)$是$\boldsymbol X$的联合pdf或pmf，$q(t\mid\theta)$是$T(\boldsymbol X)$的pdf或pmf.若对样本空间中每个$\boldsymbol x$，比值$p(\boldsymbol x\mid\theta)/q(T(\boldsymbol x)\mid\theta)$作为$\theta$的函数是常数，则$T(\boldsymbol X)$是$\theta$的充分统计量.
\end{theorem}
现在用\autoref{thm:6.2.2}验证某些常用统计量的充分性.
\begin{example}[二项分布的充分统计量]\label{ex:6.2.3}\index{SI1002@伯努利}
设$X_1,\ldots,X_n$为参数$\theta$的iid Bernoulli随机变量，$0<\theta<1$.下面说明$T(\boldsymbol X)=X_1+\cdots+X_n$是$\theta$的充分统计量.注意$T(\boldsymbol X)$计算等于1的$X_i$个数，故$T(\boldsymbol X)$服从$\operatorname{binomial}(n,\theta)$分布.pmf之比为
\begin{align*}
\frac{p(\boldsymbol x\mid\theta)}{q(T(\boldsymbol x)\mid\theta)}
&=\frac{\prod\theta^{x_i}(1-\theta)^{1-x_i}}{\binom nt\theta^t(1-\theta)^{n-t}}\qquad\text{（定义$t=\sum x_i$）}\\
&=\frac{\theta^{\sum x_i}(1-\theta)^{\sum(1-x_i)}}{\binom nt\theta^t(1-\theta)^{n-t}}\qquad\text{（$\prod\theta^{x_i}=\theta^{\sum x_i}$）}\\
&=\frac{\theta^t(1-\theta)^{n-t}}{\binom nt\theta^t(1-\theta)^{n-t}}\\
&=\frac1{\binom nt}=\frac1{\binom n{\sum x_i}}.
\end{align*}
这一比值不依赖$\theta$，故由\autoref{thm:6.2.2}，$T(\boldsymbol X)$是$\theta$的充分统计量.其含义是：这个Bernoulli样本中1的总数，包含数据中全部关于$\theta$的信息.数据的其他特征，例如$X_3$的确切值，不含额外信息.
\end{example}
\sourcepageanchor{255}
\begin{example}[正态分布的充分统计量]\label{ex:6.2.4}\index{SI1011@正态}
设$X_1,\ldots,X_n$为iid $n(\mu,\sigma^2)$，其中$\sigma^2$已知.我们要说明样本均值$T(\boldsymbol X)=\bar X=(X_1+\cdots+X_n)/n$是$\mu$的充分统计量.样本$\boldsymbol X$的联合pdf为
\begin{equation}\label{eq:6.2.1}
\begin{aligned}
f(\boldsymbol x\mid\mu)&=\prod_{i=1}^n(2\pi\sigma^2)^{-1/2}\exp\bigl(-(x_i-\mu)^2/(2\sigma^2)\bigr)\\
&=(2\pi\sigma^2)^{-n/2}\exp\left(-\sum_{i=1}^n(x_i-\mu)^2/(2\sigma^2)\right)\\
&=(2\pi\sigma^2)^{-n/2}\exp\left(-\sum_{i=1}^n(x_i-\bar x+\bar x-\mu)^2/(2\sigma^2)\right)\quad\text{（加、减$\bar x$）}\\
&=(2\pi\sigma^2)^{-n/2}\exp\left(-\left(\sum_{i=1}^n(x_i-\bar x)^2+n(\bar x-\mu)^2\right)/(2\sigma^2)\right).
\end{aligned}
\end{equation}
最后一个等式成立，是因为交叉乘积项$\sum_{i=1}^n(x_i-\bar x)(\bar x-\mu)$可写为$(\bar x-\mu)\sum_{i=1}^n(x_i-\bar x)$，而$\sum_{i=1}^n(x_i-\bar x)=0$.回顾样本均值$\bar X$服从$n(\mu,\sigma^2/n)$分布.因此pdf之比为
\begin{align*}
\frac{f(\boldsymbol x\mid\theta)}{q(T(\boldsymbol x)\mid\theta)}
&=\frac{(2\pi\sigma^2)^{-n/2}\exp\left(-\left(\sum_{i=1}^n(x_i-\bar x)^2+n(\bar x-\mu)^2\right)/(2\sigma^2)\right)}{(2\pi\sigma^2/n)^{-1/2}\exp(-n(\bar x-\mu)^2/(2\sigma^2))}\\
&=n^{-1/2}(2\pi\sigma^2)^{-(n-1)/2}\exp\left(-\sum_{i=1}^n(x_i-\bar x)^2/(2\sigma^2)\right),
\end{align*}
它不依赖$\mu$.由\autoref{thm:6.2.2}，样本均值是$\mu$的充分统计量.
\end{example}
下例考察不能显著约简样本的情形.
\begin{example}[充分的次序统计量]\label{ex:6.2.5}
设$X_1,\ldots,X_n$为来自pdf $f$的iid随机变量，而我们不能对这个pdf规定更多信息（非参数估计，nonparametric estimation，就是这种情形）.则样本密度为
\begin{equation}\label{eq:6.2.2}
f(\boldsymbol x)=\prod_{i=1}^nf(x_i)=\prod_{i=1}^nf(x_{(i)}),
\end{equation}
其中$x_{(i)}\le x_{(2)}\le\cdots\le x_{(n)}$为次序统计量.\footnote{译注：此处原文首项写$x_{(i)}$；按次序统计量的排列，应为$x_{(1)}$.}由\autoref{thm:6.2.2}可说明次序统计量是充分统计量.当然，这种约简幅度不大，但对密度$f$所知如此之少，也不应期望更多.

不过，即使对密度作了更多规定，仍可能无法得到很大的充分性约简.例如，设$f$是Cauchy的pdf $f(x\mid\theta)=1/[\pi(x-\theta)^2]$\footnote{译注：原文此处Cauchy密度分母漏写$1+$；标准位置Cauchy密度应为$1/[\pi(1+(x-\theta)^2)]$.原式保留.}，或logistic的pdf $f(x\mid\theta)=e^{-(x-\theta)}/(1+e^{-(x-\theta)})^2$.此时仍只能得到\eqref{eq:6.2.2}中的约简，不能再进一步.所以在这些族中，约简为次序统计量就是所能达到的极限（更多例子见\autoref{exercise:6.8}和\autoref{exercise:6.9}）.

事实表明，在指数分布族之外，很少能找到维数小于样本容量的充分统计量，因而许多情况下次序统计量就是我们所能做到的最好结果.（进一步细节见\hyperref[bib:source200]{Lehmann和Casella 1998}，第1.6节.）
\end{example}
对特定模型，直接用定义寻找充分统计量可能很繁琐.使用定义时，必须先猜测一个可能充分的统计量$T(\boldsymbol X)$，再求$T(\boldsymbol X)$的pmf或pdf，并检查pdf或pmf之比是否不依赖$\theta$.
\sourcepageanchor{256}
第一步需要相当多的直觉，第二步有时需要繁琐的分析.幸而下面由\hyperref[bib:source155]{Halmos和Savage（1949）}给出的定理，只要观察样本的pdf或pmf就能找到充分统计量.\footnote{不过，按Halmos和Savage的说法，他们的定理“可以改写成一种在精神上更接近先前关于充分性概念研究的形式”.这里的研究指\hyperref[bib:source233]{Neyman（1935）}的工作.（威斯康星大学密尔沃基分校的J. Beder教授指出了这一点.）}
\begin{theorem}[因子分解定理]\label{thm:6.2.6}\index{SI0443@因子分解定理}\index{SI1005@因子分解定理}
以$f(\boldsymbol x\mid\theta)$表示样本$\boldsymbol X$的联合pdf或pmf.统计量$T(\boldsymbol X)$对$\theta$充分，当且仅当存在函数$g(t\mid\theta)$和$h(\boldsymbol x)$，使对所有样本点$\boldsymbol x$和所有参数点$\theta$都有
\begin{equation}\label{eq:6.2.3}
f(\boldsymbol x\mid\theta)=g(T(\boldsymbol x)\mid\theta)h(\boldsymbol x).
\end{equation}
\end{theorem}
\begin{proof}
只给出离散分布的证明.

设$T(\boldsymbol X)$是充分统计量.取$g(t\mid\theta)=P_\theta(T(\boldsymbol X)=t)$及$h(\boldsymbol x)=P(\boldsymbol X=\boldsymbol x\mid T(\boldsymbol X)=T(\boldsymbol x))$.由于$T(\boldsymbol X)$充分，定义$h(\boldsymbol x)$的条件概率不依赖$\theta$.所以$h(\boldsymbol x)$和$g(t\mid\theta)$的这种选择合法，且有
\begin{align*}
f(\boldsymbol x\mid\theta)&=P_\theta(\boldsymbol X=\boldsymbol x)\\
&=P_\theta(\boldsymbol X=\boldsymbol x\text{且}T(\boldsymbol X)=T(\boldsymbol x))\\
&=P_\theta(T(\boldsymbol X)=T(\boldsymbol x))P(\boldsymbol X=\boldsymbol x\mid T(\boldsymbol X)=T(\boldsymbol x))\quad\text{（充分性）}\\
&=g(T(\boldsymbol x)\mid\theta)h(\boldsymbol x).
\end{align*}
这就给出了分解\eqref{eq:6.2.3}.从上面最后两行还能看出
$$P_\theta(T(\boldsymbol X)=T(\boldsymbol x))=g(T(\boldsymbol x)\mid\theta),$$
因此$g(T(\boldsymbol x)\mid\theta)$是$T(\boldsymbol X)$的pmf.

再假设分解\eqref{eq:6.2.3}存在.设$q(t\mid\theta)$是$T(\boldsymbol X)$的pmf.为说明$T(\boldsymbol X)$充分，考察比值$f(\boldsymbol x\mid\theta)/q(T(\boldsymbol x)\mid\theta)$.定义$A_{T(\boldsymbol x)}=\{\boldsymbol y:T(\boldsymbol y)=T(\boldsymbol x)\}$.则
\begin{align*}
\frac{f(\boldsymbol x\mid\theta)}{q(T(\boldsymbol x)\mid\theta)}
&=\frac{g(T(\boldsymbol x)\mid\theta)h(\boldsymbol x)}{q(T(\boldsymbol x)\mid\theta)}&&\text{（因\eqref{eq:6.2.3}成立）}\\
&=\frac{g(T(\boldsymbol x)\mid\theta)h(\boldsymbol x)}{\sum_{A_{T(\boldsymbol x)}}g(T(\boldsymbol y)\mid\theta)h(\boldsymbol y)}&&\text{（$T$的pmf定义）}\\
&=\frac{g(T(\boldsymbol x)\mid\theta)h(\boldsymbol x)}{g(T(\boldsymbol x)\mid\theta)\sum_{A_{T(\boldsymbol x)}}h(\boldsymbol y)}&&\text{（$T$在$A_{T(\boldsymbol x)}$上为常数）}\\
&=\frac{h(\boldsymbol x)}{\sum_{A_{T(\boldsymbol x)}}h(\boldsymbol y)}.
\end{align*}
比值不依赖$\theta$，故由\autoref{thm:6.2.2}，$T(\boldsymbol X)$是$\theta$的充分统计量.
\end{proof}
用因子分解定理（Factorization Theorem）寻找充分统计量时，把样本的联合pdf分解为两部分，其中一部分不依赖$\theta$.不依赖$\theta$的部分构成函数$h(\boldsymbol x)$.另一部分依赖$\theta$，通常只通过某个函数$T(\boldsymbol x)$依赖样本$\boldsymbol x$，这个函数就是$\theta$的充分统计量.下例说明这一点.
\begin{example}[续\autoref{ex:6.2.4}]\label{ex:6.2.7}\index{SI1011@正态}
对前述正态模型，我们已看到pdf可分解为
\begin{equation}\label{eq:6.2.4}
\begin{aligned}
f(\boldsymbol x\mid\mu)={}&(2\pi\sigma^2)^{-n/2}\exp\left(-\sum_{i=1}^n(x_i-\bar x)^2/(2\sigma^2)\right)\\
&\quad\cdot\exp(-n(\bar x-\mu)^2/(2\sigma^2)).
\end{aligned}
\end{equation}
\sourcepageanchor{257}
可以定义
$$h(\boldsymbol x)=(2\pi\sigma^2)^{-n/2}\exp\left(-\sum_{i=1}^n(x_i-\bar x)^2/(2\sigma^2)\right),$$
它不依赖未知参数$\mu$.\eqref{eq:6.2.4}中含$\mu$的因子只通过函数$T(\boldsymbol x)=\bar x$，即样本均值，依赖样本$\boldsymbol x$.故有
$$g(t\mid\mu)=\exp\bigl(-n(t-\mu)^2/(2\sigma^2)\bigr),$$
并注意
$$f(\boldsymbol x\mid\mu)=h(\boldsymbol x)g(T(\boldsymbol x)\mid\mu).$$
所以由因子分解定理，$T(\boldsymbol X)=\bar X$是$\mu$的充分统计量.
\end{example}
因子分解定理要求等式$f(\boldsymbol x\mid\theta)=g(T(\boldsymbol x)\mid\theta)h(\boldsymbol x)$对所有$\boldsymbol x$和$\theta$成立.若使$f(\boldsymbol x\mid\theta)$为正的$\boldsymbol x$集合依赖$\theta$，定义$h$和$g$时必须谨慎，保证在$f$为0的地方乘积也为0.当然，正确地定义$h$和$g$就能使充分统计量显现出来，下例说明这一点.
\begin{example}[均匀分布的充分统计量]\label{ex:6.2.8}\index{SI1003@离散均匀}
设$X_1,\ldots,X_n$为来自$1,\ldots,\theta$上离散均匀分布的iid观测.即未知参数$\theta$为正整数，且$X_i$的pmf为
$$f(x\mid\theta)=\begin{cases}1/\theta,&x=1,2,\ldots,\theta,\\0,&\text{其他情形}.\end{cases}$$
于是$X_1,\ldots,X_n$的联合pmf为
$$f(\boldsymbol x\mid\theta)=\begin{cases}\theta^{-n},&x_i\in\{1,\ldots,\theta\},\quad i=1,\ldots,n,\\0,&\text{其他情形}.\end{cases}$$
限制“$x_i\in\{1,\ldots,\theta\}$，$i=1,\ldots,n$”可重写为“$x_i\in\{1,2,\ldots\}$，$i=1,\ldots,n$（注意这一限制中不含$\theta$），且$\max_i x_i\le\theta$”.若定义$T(\boldsymbol x)=\max_i x_i$，
$$h(\boldsymbol x)=\begin{cases}1,&x_i\in\{1,2,\ldots\},\quad i=1,\ldots,n,\\0,&\text{其他情形},\end{cases}$$
以及
$$g(t\mid\theta)=\begin{cases}\theta^{-n},&t\le\theta,\\0,&\text{其他情形},\end{cases}$$
则容易验证$f(\boldsymbol x\mid\theta)=g(T(\boldsymbol x)\mid\theta)h(\boldsymbol x)$对所有$\boldsymbol x$和$\theta$成立.因此，最大次序统计量$T(\boldsymbol X)=\max_iX_i$是本问题中的充分统计量.

有时使用指示函数能更清楚、简洁地进行这种分析.回顾$I_A(x)$是集合$A$的指示函数，即$x\in A$时为1，否则为0.令$\mathcal N=\{1,2,\ldots\}$为正整数集，$\mathcal N_\theta=\{1,2,\ldots,\theta\}$.那么$X_1,\ldots,X_n$的联合pmf为
$$f(\boldsymbol x\mid\theta)=\prod_{i=1}^n\theta^{-1}I_{\mathcal N_\theta}(x_i)=\theta^{-n}\prod_{i=1}^n I_{\mathcal N_\theta}(x_i).$$
定义$T(\boldsymbol x)=\max_i x_i$，显然有
$$\prod_{i=1}^n I_{\mathcal N_\theta}(x_i)=\left(\prod_{i=1}^n I_{\mathcal N}(x_i)\right)I_{\mathcal N_\theta}(T(\boldsymbol x)).$$
\sourcepageanchor{258}
因而得到分解
$$f(\boldsymbol x\mid\theta)=\theta^{-n}I_{\mathcal N_\theta}(T(\boldsymbol x))\left(\prod_{i=1}^n I_{\mathcal N}(x_i)\right).$$
第一个因子仅通过$T(\boldsymbol x)=\max_i x_i$的值依赖$x_1,\ldots,x_n$，第二个因子不依赖$\theta$.由因子分解定理，$T(\boldsymbol X)=\max_iX_i$是$\theta$的充分统计量.
\end{example}
在前面所有例子中，充分统计量都是样本的实值函数.样本$\boldsymbol x$中全部关于$\theta$的信息都概括在单个数$T(\boldsymbol x)$中.有时一个数无法概括这些信息，需要若干个数.这时充分统计量是一个向量，例如$T(\boldsymbol X)=(T_1(\boldsymbol X),\ldots,T_r(\boldsymbol X))$.当参数也是向量，如$\boldsymbol\theta=(\theta_1,\ldots,\theta_s)$时，常出现这种情形；通常充分统计量向量和参数向量具有相同长度，即$r=s$.不过，长度的不同组合也是可能的，习题和\autoref{ex:6.2.15}、\autoref{ex:6.2.18}及\autoref{ex:6.2.20}会说明这一点.因子分解定理也可用于寻找向量值充分统计量，如\autoref{ex:6.2.9}.
\begin{example}[两个参数均未知时的正态充分统计量]\label{ex:6.2.9}\index{SI1011@正态}
仍设$X_1,\ldots,X_n$为iid $n(\mu,\sigma^2)$，但与\autoref{ex:6.2.4}不同，假设$\mu$和$\sigma^2$都未知，故参数向量为$\boldsymbol\theta=(\mu,\sigma^2)$.此时使用因子分解定理，联合pdf中凡依赖$\mu$或$\sigma^2$的部分都必须纳入函数$g$.由\eqref{eq:6.2.1}显然可见，pdf仅通过$T_1(\boldsymbol x)=\bar x$和$T_2(\boldsymbol x)=s^2=\sum_{i=1}^n(x_i-\bar x)^2/(n-1)$两个值依赖样本$\boldsymbol x$.因此可定义$h(\boldsymbol x)=1$，以及
\begin{align*}
g(\boldsymbol t\mid\boldsymbol\theta)&=g(t_1,t_2\mid\mu,\sigma^2)\\
&=(2\pi\sigma^2)^{-n/2}\exp\left(-\bigl(n(t_1-\mu)^2+(n-1)t_2\bigr)/(2\sigma^2)\right).
\end{align*}
于是可见
\begin{equation}\label{eq:6.2.5}
f(\boldsymbol x\mid\mu,\sigma^2)=g(T_1(\boldsymbol x),T_2(\boldsymbol x)\mid\mu,\sigma^2)h(\boldsymbol x).
\end{equation}
因此由因子分解定理，$T(\boldsymbol X)=(T_1(\boldsymbol X),T_2(\boldsymbol X))=(\bar X,S^2)$是此正态模型中$(\mu,\sigma^2)$的充分统计量.
\end{example}
\autoref{ex:6.2.9}表明，对正态模型，只报告样本均值和方差来概括数据集这一常见做法是有依据的.充分统计量$(\bar X,S^2)$包含样本能提供的全部关于$(\mu,\sigma^2)$的信息.不过，实验者应记住充分统计量的定义依赖模型.对于另一模型，即另一密度族，样本均值和方差可能不再是总体均值和方差的充分统计量.若实验者只计算$\bar X$和$S^2$，而完全忽略其余数据，那就对正态模型假设寄予了很强的信任.

利用因子分解定理，很容易找到指数分布族的充分统计量.以下重要结果的证明留作\autoref{exercise:6.4}.
\begin{theorem}\label{thm:6.2.10}\index{SI1004@指数族}
设$X_1,\ldots,X_n$为来自pdf或pmf $f(x\mid\boldsymbol\theta)$的iid观测，且此分布属于下式给出的指数族：
$$f(x\mid\boldsymbol\theta)=h(x)c(\boldsymbol\theta)\exp\left(\sum_{i=1}^k w_i(\boldsymbol\theta)t_i(x)\right),$$
其中$\boldsymbol\theta=(\theta_1,\theta_2,\ldots,\theta_d)$，$d\le k$.则
$$T(\boldsymbol X)=\left(\sum_{j=1}^nt_1(X_j),\ldots,\sum_{j=1}^nt_k(X_j)\right)$$
是$\boldsymbol\theta$的充分统计量.
\end{theorem}
\sourcepageanchor{259}
\subsubsection{最小充分统计量}\label{sec:6.2.2}
上一小节为每个所考察的模型找到了一个充分统计量.事实上，在任何问题中都有许多充分统计量.

完整样本$\boldsymbol X$总是充分统计量.可把$\boldsymbol X$的pdf或pmf分解为$f(\boldsymbol x\mid\theta)=f(T(\boldsymbol x)\mid\theta)h(\boldsymbol x)$，其中对所有$\boldsymbol x$都有$T(\boldsymbol x)=\boldsymbol x$和$h(\boldsymbol x)=1$.由因子分解定理，$T(\boldsymbol X)=\boldsymbol X$是充分统计量.

还可以推出，充分统计量的任意一一函数仍是充分统计量.设$T(\boldsymbol X)$是充分统计量，对所有$\boldsymbol x$定义$T^*(\boldsymbol x)=r(T(\boldsymbol x))$，其中$r$是一一函数，逆为$r^{-1}$.由因子分解定理，存在$g$和$h$，使
$$f(\boldsymbol x\mid\theta)=g(T(\boldsymbol x)\mid\theta)h(\boldsymbol x)=g(r^{-1}(T^*(\boldsymbol x))\mid\theta)h(\boldsymbol x).$$
定义$g^*(t\mid\theta)=g(r^{-1}(t)\mid\theta)$，可见
$$f(\boldsymbol x\mid\theta)=g^*(T^*(\boldsymbol x)\mid\theta)h(\boldsymbol x).$$
所以由因子分解定理，$T^*(\boldsymbol X)$是充分统计量.

一个问题中有众多充分统计量，因此自然会问其中是否有更好的.回顾充分统计量的目的是在不损失关于参数$\theta$的信息的条件下约简数据；因此，仍保留全部关于$\theta$的信息而约简程度最大的统计量，可以认为更可取.下面正式定义这样的统计量.
\begin{definition}\label{def:6.2.11}\index{SI1006@最小}
若对任意其他充分统计量$T'(\boldsymbol X)$，$T(\boldsymbol x)$都是$T'(\boldsymbol x)$的函数，则称充分统计量$T(\boldsymbol X)$为最小充分统计量（minimal sufficient statistic）.
\end{definition}
说$T(\boldsymbol x)$是$T'(\boldsymbol x)$的函数，只是指$T'(\boldsymbol x)=T'(\boldsymbol y)$必推出$T(\boldsymbol x)=T(\boldsymbol y)$.用本章开头所述的划分集合来说，\index{SI0753@划分}\semanticindexanchor{753:259}若$\{B_{t'}:t'\in\mathcal T'\}$是$T'(\boldsymbol x)$的划分集合，而$\{A_t:t\in\mathcal T\}$是$T(\boldsymbol x)$的划分集合，则\autoref{def:6.2.11}断言每个$B_{t'}$都是某个$A_t$的子集.所以最小充分统计量对应的划分，是充分统计量所能产生的最粗划分；最小充分统计量实现了充分统计量所能达到的最大数据约简.
\begin{example}[两个正态充分统计量]\label{ex:6.2.12}\index{SI1011@正态}
\autoref{ex:6.2.4}的模型是$X_1,\ldots,X_n$为iid $n(\mu,\sigma^2)$且$\sigma^2$已知.用分解\eqref{eq:6.2.4}，我们得出$T(\boldsymbol X)=\bar X$对$\mu$充分.也可以对本问题写出分解\eqref{eq:6.2.5}（现在$\sigma^2$是已知值），同样正确地得出$T'(\boldsymbol X)=(\bar X,S^2)$对$\mu$充分.显然，$T(\boldsymbol X)$比$T'(\boldsymbol X)$的约简更大，因为只知道$T(\boldsymbol X)$就不知道样本方差.定义$r(a,b)=a$，便可把$T(\boldsymbol x)$写成$T'(\boldsymbol x)$的函数：$T(\boldsymbol x)=\bar x=r(\bar x,s^2)=r(T'(\boldsymbol x))$.由于$T(\boldsymbol X)$和$T'(\boldsymbol X)$都是充分统计量，二者包含同样的关于$\mu$的信息.因此，总体方差$\sigma^2$已知时，关于样本方差$S^2$值的额外信息，不会增进我们对$\mu$的了解.当然，若如\autoref{ex:6.2.9}那样$\sigma^2$未知，$T(\boldsymbol X)=\bar X$就不充分，而$T'(\boldsymbol X)$包含的关于参数$(\mu,\sigma^2)$的信息比$T(\boldsymbol X)$更多.
\end{example}
用\autoref{def:6.2.11}寻找最小充分统计量不切实际，正如用\autoref{def:6.2.1}寻找充分统计量一样.需要猜测$T(\boldsymbol X)$最小充分，然后验证定义中的条件.（注意，\autoref{ex:6.2.12}并未说明$\bar X$最小充分.）幸而\hyperref[bib:source201]{Lehmann和Scheffé（1950}，定理6.3）的以下结果，给出了更容易的寻找方法.
\begin{theorem}\label{thm:6.2.13}\index{SI1007@刻画}
设$f(\boldsymbol x\mid\theta)$是样本$\boldsymbol X$的pmf或pdf.若存在函数$T(\boldsymbol x)$，使得对每两个样本点$\boldsymbol x$和$\boldsymbol y$，比值$f(\boldsymbol x\mid\theta)/f(\boldsymbol y\mid\theta)$作为$\theta$的函数为常数，当且仅当$T(\boldsymbol x)=T(\boldsymbol y)$，则$T(\boldsymbol X)$是$\theta$的最小充分统计量.
\end{theorem}
\begin{proof}
为简化证明，假设对所有$\boldsymbol x\in\mathcal X$和$\theta$都有$f(\boldsymbol x\mid\theta)>0$.
\sourcepageanchor{260}
先说明$T(\boldsymbol X)$是充分统计量.令$\mathcal T=\{t:t=T(\boldsymbol x)\text{，其中某个}\boldsymbol x\in\mathcal X\}$为$\mathcal X$在$T(\boldsymbol x)$下的像.定义$T(\boldsymbol x)$诱导的划分集合$A_t=\{\boldsymbol x:T(\boldsymbol x)=t\}$.对每个$A_t$，选定并固定一个元素$\boldsymbol x_t\in A_t$.对任意$\boldsymbol x\in\mathcal X$，$\boldsymbol x_{T(\boldsymbol x)}$是与$\boldsymbol x$处于同一集合$A_t$中的固定元素.由于$\boldsymbol x$和$\boldsymbol x_{T(\boldsymbol x)}$在同一$A_t$中，$T(\boldsymbol x)=T(\boldsymbol x_{T(\boldsymbol x)})$，所以$f(\boldsymbol x\mid\theta)/f(\boldsymbol x_{T(\boldsymbol x)}\mid\theta)$作为$\theta$的函数为常数.因而可在$\mathcal X$上定义$h(\boldsymbol x)=f(\boldsymbol x\mid\theta)/f(\boldsymbol x_{T(\boldsymbol x)}\mid\theta)$，且$h$不依赖$\theta$.在$\mathcal T$上定义$g(t\mid\theta)=f(\boldsymbol x_t\mid\theta)$.于是可见
$$f(\boldsymbol x\mid\theta)=\frac{f(\boldsymbol x_{T(\boldsymbol x)}\mid\theta)f(\boldsymbol x\mid\theta)}{f(\boldsymbol x_{T(\boldsymbol x)}\mid\theta)}=g(T(\boldsymbol x)\mid\theta)h(\boldsymbol x),$$
由因子分解定理，$T(\boldsymbol X)$对$\theta$充分.

再说明$T(\boldsymbol X)$最小.设$T'(\boldsymbol X)$是任意其他充分统计量.由因子分解定理，存在函数$g'$和$h'$，使$f(\boldsymbol x\mid\theta)=g'(T'(\boldsymbol x)\mid\theta)h'(\boldsymbol x)$.设$\boldsymbol x$、$\boldsymbol y$是满足$T'(\boldsymbol x)=T'(\boldsymbol y)$的任意两个样本点.则
$$\frac{f(\boldsymbol x\mid\theta)}{f(\boldsymbol y\mid\theta)}=\frac{g'(T'(\boldsymbol x)\mid\theta)h'(\boldsymbol x)}{g'(T'(\boldsymbol y)\mid\theta)h'(\boldsymbol y)}=\frac{h'(\boldsymbol x)}{h'(\boldsymbol y)}.$$
比值不依赖$\theta$，故由定理假设$T(\boldsymbol x)=T(\boldsymbol y)$.所以$T(\boldsymbol x)$是$T'(\boldsymbol x)$的函数，$T(\boldsymbol x)$最小.
\end{proof}
\begin{example}[正态分布的最小充分统计量]\label{ex:6.2.14}\index{SI1009@正态}
设$X_1,\ldots,X_n$为iid $n(\mu,\sigma^2)$，$\mu$和$\sigma^2$均未知.以$\boldsymbol x$和$\boldsymbol y$表示两个样本点，$(\bar x,s_x^2)$和$(\bar y,s_y^2)$分别是相应样本的均值和方差.由\eqref{eq:6.2.5}，密度之比为
\begin{align*}
\frac{f(\boldsymbol x\mid\mu,\sigma^2)}{f(\boldsymbol y\mid\mu,\sigma^2)}
&=\frac{(2\pi\sigma^2)^{-n/2}\exp\left(-[n(\bar x-\mu)^2+(n-1)s_x^2]/(2\sigma^2)\right)}{(2\pi\sigma^2)^{-n/2}\exp\left(-[n(\bar y-\mu)^2+(n-1)s_y^2]/(2\sigma^2)\right)}\\
&=\exp\left(\bigl[-n(\bar x^2-\bar y^2)+2n\mu(\bar x-\bar y)-(n-1)(s_x^2-s_y^2)\bigr]/(2\sigma^2)\right).
\end{align*}
此比值作为$\mu$和$\sigma^2$的函数为常数，当且仅当$\bar x=\bar y$且$s_x^2=s_y^2$.因此由\autoref{thm:6.2.13}，$(\bar X,S^2)$是$(\mu,\sigma^2)$的最小充分统计量.
\end{example}
若使pdf或pmf为正的$\boldsymbol x$集合依赖参数$\theta$，那么要使\autoref{thm:6.2.13}中的比值作为$\theta$的函数为常数，分子和分母必须恰在同样的$\theta$值处为正.这个限制通常会反映在最小充分统计量中，如下例所示.
\begin{example}[均匀分布的最小充分统计量]\label{ex:6.2.15}\index{SI1010@均匀}
设$X_1,\ldots,X_n$为区间$(\theta,\theta+1)$上的iid均匀观测，$-\infty<\theta<\infty$.则$\boldsymbol X$的联合pdf为
$$f(\boldsymbol x\mid\theta)=\begin{cases}1,&\theta<x_i<\theta+1,\quad i=1,\ldots,n,\\0,&\text{其他情形},\end{cases}$$
可写为
$$f(\boldsymbol x\mid\theta)=\begin{cases}1,&\max_i x_i-1<\theta<\min_i x_i,\\0,&\text{其他情形}.\end{cases}$$
因此，对两个样本点$\boldsymbol x$和$\boldsymbol y$，比值$f(\boldsymbol x\mid\theta)/f(\boldsymbol y\mid\theta)$的分子和分母在同样的$\theta$值处为正，当且仅当$\min_i x_i=\min_i y_i$且$\max_i x_i=\max_i y_i$.最小值和最大值均相等时，比值为常数，且实际上等于1.于是令$X_{(1)}=\min_iX_i$、$X_{(n)}=\max_iX_i$，便有$T(\boldsymbol X)=(X_{(1)},X_{(n)})$是最小充分统计量.此例中最小充分统计量的维数与参数维数不同.
\end{example}
\sourcepageanchor{261}
最小充分统计量并不唯一.它的任何一一函数仍是最小充分统计量.例如，\autoref{ex:6.2.15}中$T'(\boldsymbol X)=(X_{(n)}-X_{(1)},(X_{(n)}+X_{(1)})/2)$也是最小充分统计量；\autoref{ex:6.2.14}中$T'(\boldsymbol X)=(\sum_{i=1}^nX_i,\sum_{i=1}^nX_i^2)$也是最小充分统计量.
\subsubsection{辅助统计量}\label{sec:6.2.3}
前面各小节讨论的充分统计量，在某种意义上包含样本能提供的全部关于$\theta$的信息.本小节引入另一类具有互补作用的统计量.
\begin{definition}\label{def:6.2.16}\index{SI0018@辅助统计量}
分布不依赖参数$\theta$的统计量$S(\boldsymbol X)$称为辅助统计量（ancillary statistic）.
\end{definition}
单独看，辅助统计量不含关于$\theta$的信息.它是一个分布固定且已知、与$\theta$无关的随机变量的观测.看似矛盾的是，辅助统计量与其他统计量联合使用时，有时确实含有关于$\theta$推断的宝贵信息.下一小节将考察这种现象.现在只给出若干辅助统计量的例子.
\begin{example}[均匀分布的辅助统计量]\label{ex:6.2.17}\index{SI0024@均匀}
如\autoref{ex:6.2.15}，设$X_1,\ldots,X_n$为区间$(\theta,\theta+1)$上的iid均匀观测，$-\infty<\theta<\infty$.以$X_{(1)}<\cdots<X_{(n)}$表示样本次序统计量.下面通过说明极差统计量$R=X_{(n)}-X_{(1)}$的pdf不依赖$\theta$，来证明$R$为辅助统计量.回顾每个$X_i$的cdf为
$$F(x\mid\theta)=\begin{cases}0,&x\le\theta,\\x-\theta,&\theta<x<\theta+1,\\1,&\theta+1\le x.\end{cases}$$
因此，由\eqref{eq:5.5.7}给出的$X_{(1)}$与$X_{(n)}$联合pdf为\footnote{译注：原文引用式(5.5.7)，该式属于Delta方法；这里应参见次序统计量联合密度\eqref{eq:5.4.6}及\autoref{ex:5.4.7}.}
$$g(x_{(1)},x_{(n)}\mid\theta)=\begin{cases}n(n-1)(x_{(n)}-x_{(1)})^{n-2},&\theta<x_{(1)}<x_{(n)}<\theta+1,\\0,&\text{其他情形}.\end{cases}$$
作变换$R=X_{(n)}-X_{(1)}$、$M=(X_{(1)}+X_{(n)})/2$，逆变换为$X_{(1)}=(2M-R)/2$、$X_{(n)}=(2M+R)/2$，Jacobian为1，可得$R$和$M$的联合pdf为
$$h(r,m\mid\theta)=\begin{cases}n(n-1)r^{n-2},&0<r<1,\ \theta+r/2<m<\theta+1-r/2,\\0,&\text{其他情形}.\end{cases}$$
（注意$h(r,m\mid\theta)$为正的区域相当复杂.）所以$R$的pdf为
\begin{align*}
h(r\mid\theta)&=\int_{\theta+r/2}^{\theta+1-r/2}n(n-1)r^{n-2}\,dm\\
&=n(n-1)r^{n-2}(1-r),\qquad0<r<1.
\end{align*}
这是$\alpha=n-1$、$\beta=2$的beta密度.更重要的是，对所有$\theta$这个pdf相同.所以$R$的分布不依赖$\theta$，$R$为辅助统计量.
\end{example}
\autoref{ex:6.2.17}中的极差是辅助统计量，是因为所考虑的模型是位置参数模型.$R$的辅助性并不依赖$X_i$服从均匀分布，而是依赖分布参数为位置参数.现在考虑一般的位置参数模型.
\sourcepageanchor{262}
\begin{example}[位置族的辅助统计量]\label{ex:6.2.18}\index{SI0020@位置族}
设$X_1,\ldots,X_n$为来自cdf为$F(x-\theta)$的位置参数族的iid观测，$-\infty<\theta<\infty$.下面说明极差$R=X_{(n)}-X_{(1)}$是辅助统计量.利用\autoref{thm:3.5.6}，取来自$F(x)$（对应$\theta=0$）的iid观测$Z_1,\ldots,Z_n$，满足$X_1=Z_1+\theta,\ldots,X_n=Z_n+\theta$.所以极差统计量$R$的cdf为
\begin{align*}
F_R(r\mid\theta)&=P_\theta(R\le r)\\
&=P_\theta(\max_iX_i-\min_iX_i\le r)\\
&=P_\theta(\max_i(Z_i+\theta)-\min_i(Z_i+\theta)\le r)\\
&=P_\theta(\max_iZ_i-\min_iZ_i+\theta-\theta\le r)\\
&=P_\theta(\max_iZ_i-\min_iZ_i\le r).
\end{align*}
最后一个概率不依赖$\theta$，因为$Z_1,\ldots,Z_n$的分布不依赖$\theta$.故$R$的cdf不依赖$\theta$，从而$R$为辅助统计量.
\end{example}
\begin{example}[尺度族的辅助统计量]\label{ex:6.2.19}\index{SI0023@尺度族}
尺度参数族也有某些类型的辅助统计量.设$X_1,\ldots,X_n$为来自cdf为$F(x/\sigma)$的尺度参数族的iid观测，$\sigma>0$.则任何仅通过$n-1$个值$X_1/X_n,\ldots,X_{n-1}/X_n$依赖样本的统计量，都是辅助统计量.例如，
$$\frac{X_1+\cdots+X_n}{X_n}=\frac{X_1}{X_n}+\cdots+\frac{X_{n-1}}{X_n}+1$$
是辅助统计量.为说明这一点，设$Z_1,\ldots,Z_n$为来自$F(x)$（对应$\sigma=1$）的iid观测，满足$X_i=\sigma Z_i$.$X_1/X_n,\ldots,X_{n-1}/X_n$的联合cdf为
\begin{align*}
F(y_1,\ldots,y_{n-1}\mid\sigma)&=P_\sigma(X_1/X_n\le y_1,\ldots,X_{n-1}/X_n\le y_{n-1})\\
&=P_\sigma(\sigma Z_1/(\sigma Z_n)\le y_1,\ldots,\sigma Z_{n-1}/(\sigma Z_n)\le y_{n-1})\\
&=P_\sigma(Z_1/Z_n\le y_1,\ldots,Z_{n-1}/Z_n\le y_{n-1}).
\end{align*}
最后一个概率不依赖$\sigma$，因为$Z_1,\ldots,Z_n$的分布不依赖$\sigma$.所以$X_1/X_n,\ldots,X_{n-1}/X_n$的分布与$\sigma$无关，这些量的任何函数的分布也如此.

特别地，令$X_1$和$X_2$为iid $n(0,\sigma^2)$观测.由上述结果，$X_1/X_2$对每个$\sigma$都有相同分布.而\autoref{ex:4.3.6}说明$\sigma=1$时，$X_1/X_2$服从$\operatorname{Cauchy}(0,1)$.因此任意$\sigma>0$时，$X_1/X_2$都具有同一Cauchy分布.
\end{example}
本小节给出了各种模型下辅助统计量的例子，其中有些相当一般.下一小节将考虑充分统计量与辅助统计量之间的关系.
\subsubsection{充分、辅助与完备统计量}\label{sec:6.2.4}
最小充分统计量是在仍保留全部关于参数$\theta$的信息的同时，达到最大可能约简程度的统计量.直观地说，最小充分统计量去掉样本中所有无关信息，只保留关于$\theta$的信息.由于辅助统计量的分布不依赖$\theta$，人们或许会猜测最小充分统计量与辅助统计量无关（或用数学语言说，函数上独立）.然而并不一定如此.本小节较详细地考察这种关系.
\sourcepageanchor{263}
我们已经讨论过辅助统计量与最小充分统计量不独立的情形.回顾\autoref{ex:6.2.15}，$X_1,\ldots,X_n$为来自$\operatorname{uniform}(\theta,\theta+1)$的iid观测.在\autoref{sec:6.2.2}末尾，指出$(X_{(n)}-X_{(1)},(X_{(n)}+X_{(1)})/2)$是最小充分统计量；\autoref{ex:6.2.17}则说明$X_{(n)}-X_{(1)}$是辅助统计量.所以此时辅助统计量是最小充分统计量的重要组成部分.二者当然不独立.

为强调辅助统计量有时能为$\theta$的推断提供重要信息，再给出一个例子.
\begin{example}[辅助统计量与精度]\label{ex:6.2.20}\index{SI0025@用途}
设$X_1$和$X_2$为来自满足下式的离散分布的iid观测：
$$P_\theta(X=\theta)=P_\theta(X=\theta+1)=P_\theta(X=\theta+2)=\frac13,$$
其中未知参数$\theta$可为任意整数.令$X_{(1)}\le X_{(2)}$为样本次序统计量.用类似\autoref{ex:6.2.15}的论证可知，$(R,M)$为最小充分统计量，其中$R=X_{(2)}-X_{(1)}$、$M=(X_{(1)}+X_{(2)})/2$.由于这是位置参数族，由\autoref{ex:6.2.17}，$R$为辅助统计量.为看出$R$虽然是辅助统计量却可能给出关于$\theta$的信息，考虑样本点$(r,m)$，其中$m$为整数.先只考虑$m$，要使这个样本点有正概率，$\theta$必须取三个值之一：$\theta=m$、$\theta=m-1$或$\theta=m-2$.仅凭$M=m$这一信息，三个$\theta$值都有可能.但若再获知$R=2$，就必有$X_{(1)}=m-1$、$X_{(2)}=m+1$.有了这一额外信息，$\theta$唯一可能的值就是$\theta=m-1$.所以知道辅助统计量$R$的值，增进了我们对$\theta$的了解.当然，单独知道$R$并不提供关于$\theta$的信息.（辅助统计量能提供关于$\theta$估计精度的信息，这一思想并非新事物.更多观点可参见\hyperref[bib:source85]{Cox（1971）}或\hyperref[bib:source110]{Efron和Hinkley（1978）}.）
\end{example}
不过，在许多重要情形中，最小充分统计量与任意辅助统计量独立这一直觉是正确的.描述何时出现这种情况，要用到下面的定义.
\begin{definition}\label{def:6.2.21}\index{SI0110@完备性}
设$f(t\mid\theta)$是统计量$T(\boldsymbol X)$的pdf或pmf族.若对所有$\theta$有$E_\theta g(T)=0$，就能推出对所有$\theta$有$P_\theta(g(T)=0)=1$，则称这个概率分布族完备（complete）.等价地，称$T(\boldsymbol X)$为完备统计量（complete statistic）.
\end{definition}
注意完备性是概率分布族的性质，并非某一个特定分布的性质.例如，若$X$服从$n(0,1)$，定义$g(x)=x$就有$Eg(X)=EX=0$.但$g(x)=x$满足$P(g(X)=0)=P(X=0)=0$，而非1.不过这只是一个特定分布，不是分布族.若$X$服从$n(\theta,1)$，$-\infty<\theta<\infty$，我们将看到，除对所有$\theta$以概率1等于0的函数之外，没有任何$X$的函数能对所有$\theta$满足$E_\theta g(X)=0$.所以$n(\theta,1)$、$-\infty<\theta<\infty$这一分布族是完备的.
\begin{example}[二项分布的完备充分统计量]\label{ex:6.2.22}\index{SI0114@二项}
设$T$服从$\operatorname{binomial}(n,p)$，$0<p<1$.设$g$为满足$E_pg(T)=0$的函数.则
\begin{align*}
0=E_pg(T)&=\sum_{t=0}^n g(t)\binom nt p^t(1-p)^{n-t}\\
&=(1-p)^n\sum_{t=0}^n g(t)\binom nt\left(\frac p{1-p}\right)^t.
\end{align*}
\sourcepageanchor{264}
对所有$0<p<1$成立.在这个范围内，因子$(1-p)^n$不为0.因此，必有
$$0=\sum_{t=0}^ng(t)\binom nt\left(\frac p{1-p}\right)^t=\sum_{t=0}^ng(t)\binom nt r^t$$
对所有$0<r<\infty$成立.最后一个表达式是$r$的$n$次多项式，其中$r^t$的系数为$g(t)\binom nt$.要使多项式对所有$r$都为0，每个系数都必须为0.由于所有$\binom nt$均非零，故$t=0,1,\ldots,n$时$g(t)=0$.又$T$以概率1取值于$0,1,\ldots,n$，从而对所有$p$都有$P_p(g(T)=0)=1$，这正是所需结论.所以$T$是完备统计量.
\end{example}
\begin{example}[均匀分布的完备充分统计量]\label{ex:6.2.23}\index{SI0120@均匀}
设$X_1,\ldots,X_n$为iid $\operatorname{uniform}(0,\theta)$观测，$0<\theta<\infty$.用类似\autoref{ex:6.2.8}的论证，可说明$T(\boldsymbol X)=\max_iX_i$是充分统计量；由\autoref{thm:5.4.4}，$T(\boldsymbol X)$的pdf为
$$f(t\mid\theta)=\begin{cases}nt^{n-1}\theta^{-n},&0<t<\theta,\\0,&\text{其他情形}.\end{cases}$$
设函数$g(t)$对所有$\theta$满足$E_\theta g(T)=0$.由于$E_\theta g(T)$作为$\theta$的函数为常数，其对$\theta$的导数为0.故有
\begin{align*}
0=\frac d{d\theta}E_\theta g(T)&=\frac d{d\theta}\int_0^\theta g(t)nt^{n-1}\theta^{-n}\,dt\\
&=\theta^{-n}\frac d{d\theta}\int_0^\theta ng(t)t^{n-1}\,dt+\left(\frac d{d\theta}\theta^{-n}\right)\int_0^\theta ng(t)t^{n-1}\,dt\\
&=\theta^{-n}ng(\theta)\theta^{n-1}+0\qquad\text{（乘积求导法则）}\\
&=\theta^{-1}ng(\theta).
\end{align*}
倒数第二行的第一项来自微积分基本定理.第二项为0，是因为其中的积分除一个常数因子外等于$E_\theta g(T)$，而后者为0.由于$\theta^{-1}ng(\theta)=0$且$\theta^{-1}n\ne0$，必有$g(\theta)=0$.这对每个$\theta>0$都成立，故$T$是完备统计量.（稍作严格的提醒：微积分基本定理并非适用于所有函数，而只适用于Riemann可积函数.等式
$$\frac d{d\theta}\int_0^\theta g(t)\,dt=g(\theta)$$
只在Riemann可积函数$g$的连续点成立.因此，严格地说，上述论证没有证明$T$完备，因为完备性条件针对所有函数，而不仅是Riemann可积函数.但从更实用的角度，这一区别并不重要，因为Riemann可积条件非常宽泛，几乎包含我们能想到的任何函数.）
\end{example}
下面利用完备性，给出最小充分统计量与每个辅助统计量独立的一个条件.
\begin{theorem}[Basu定理]\label{thm:6.2.24}\index{SI0076@Basu定理}\index{SI0111@辅助性}\index{SI0112@Basu定理}\index{SI0117@独立性}\index{SI0119@充分性}\index{SI0547@完备性}
若$T(\boldsymbol X)$是完备的最小充分统计量，则$T(\boldsymbol X)$与每个辅助统计量独立.
\end{theorem}
\begin{proof}
只给出离散分布的证明.

设$S(\boldsymbol X)$为任意辅助统计量.由于其辅助性，$P(S(\boldsymbol X)=s)$不依赖$\theta$.条件概率
$$P(S(\boldsymbol X)=s\mid T(\boldsymbol X)=t)=P(\boldsymbol X\in\{\boldsymbol x:S(\boldsymbol x)=s\}\mid T(\boldsymbol X)=t)$$
也不依赖$\theta$，因为$T(\boldsymbol X)$是充分统计量（回顾其定义！）.
\sourcepageanchor{265}
所以，要说明$S(\boldsymbol X)$与$T(\boldsymbol X)$独立，只须说明对所有可能的值$t\in\mathcal T$都有
\begin{equation}\label{eq:6.2.6}
P(S(\boldsymbol X)=s\mid T(\boldsymbol X)=t)=P(S(\boldsymbol X)=s).
\end{equation}
现在，
$$P(S(\boldsymbol X)=s)=\sum_{t\in\mathcal T}P(S(\boldsymbol X)=s\mid T(\boldsymbol X)=t)P_\theta(T(\boldsymbol X)=t).$$
又因为$\sum_{t\in\mathcal T}P_\theta(T(\boldsymbol X)=t)=1$，可写
$$P(S(\boldsymbol X)=s)=\sum_{t\in\mathcal T}P(S(\boldsymbol X)=s)P_\theta(T(\boldsymbol X)=t).$$
因此，若定义统计量
$$g(t)=P(S(\boldsymbol X)=s\mid T(\boldsymbol X)=t)-P(S(\boldsymbol X)=s),$$
以上两式说明
$$E_\theta g(T)=\sum_{t\in\mathcal T}g(t)P_\theta(T(\boldsymbol X)=t)=0,\qquad\text{对所有$\theta$}.$$
由于$T(\boldsymbol X)$是完备统计量，可推出对所有可能的$t\in\mathcal T$都有$g(t)=0$.因此\eqref{eq:6.2.6}得到验证.
\end{proof}
Basu定理的用处在于，无须求两个统计量的联合分布，就能推出其独立性.使用Basu定理，需要说明一个统计量完备，这有时是相当困难的分析问题.幸而下面的定理涵盖了我们关心的大多数问题.这里不证明该定理，但指出其证明依赖Laplace变换的唯一性，\index{SI0601@拉普拉斯变换}\semanticindexanchor{601:265}\autoref{sec:2.3}曾提及这一性质.
\begin{theorem}[指数族中的完备统计量]\label{thm:6.2.25}\index{SI0116@指数族}
设$X_1,\ldots,X_n$为来自指数族的iid观测，其pdf或pmf具有形式
\begin{equation}\label{eq:6.2.7}
f(x\mid\boldsymbol\theta)=h(x)c(\boldsymbol\theta)\exp\left(\sum_{j=1}^k w(\theta_j)t_j(x)\right),
\end{equation}
其中$\boldsymbol\theta=(\theta_1,\theta_2,\ldots,\theta_k)$.则统计量
$$T(\boldsymbol X)=\left(\sum_{i=1}^nt_1(X_i),\sum_{i=1}^nt_2(X_i),\ldots,\sum_{i=1}^nt_k(X_i)\right)$$
只要参数空间$\Theta$包含$\mathbb R^k$中的一个开集，就是完备的.\footnote{译注：本定理的开集条件严格说应加在自然参数$(w(\theta_1),\ldots,w(\theta_k))$的取值集合上；若$w$退化，单凭原参数空间$\Theta$含开集不足以保证结论.原表述保留；通常的满秩指数族形式满足相应自然参数条件.}
\end{theorem}
要求参数空间包含开集，是为避免如下情形.\index{SI0436@曲线}\semanticindexanchor{436:265}$n(\theta,\theta^2)$分布可以写成\eqref{eq:6.2.7}的形式，但参数空间$(\theta,\theta^2)$只由抛物线上的点组成，不含二维开集.于是可找到统计量$T(\boldsymbol X)$的一个变换，它是0的无偏估计量（见\autoref{exercise:6.15}）.（回顾$n(\theta,\theta^2)$这类参数空间为低维曲线的指数族，称为弯曲指数族，curved exponential families；见\autoref{sec:3.4}.）指数族中充分性、完备性与最小性的关系很有意思，简短介绍见\autoref{sec:6.6.3}.

下面举例说明Basu定理、\autoref{thm:6.2.25}以及本章前面许多结果的运用.
\sourcepageanchor{266}
\begin{example}[Basu定理的运用——I]\label{ex:6.2.26}\index{SI0019@指数}\index{SI0115@指数}
设$X_1,\ldots,X_n$为参数$\theta$的iid指数观测.考虑计算
$$g(\boldsymbol X)=\frac{X_n}{X_1+\cdots+X_n}$$
的期望值.先注意指数分布构成尺度参数族，因此由\autoref{ex:6.2.19}，$g(\boldsymbol X)$是辅助统计量.指数分布也构成$t(x)=x$的指数族，所以由\autoref{thm:6.2.25}，
$$T(\boldsymbol X)=\sum_{i=1}^nX_i$$
是完备统计量；由\autoref{thm:6.2.10}，$T(\boldsymbol X)$也是充分统计量.（如下文指出的，无须验证$T(\boldsymbol X)$最小，虽然用\autoref{thm:6.2.13}很容易验证.）故由Basu定理，$T(\boldsymbol X)$与$g(\boldsymbol X)$独立.于是
$$\theta=E_\theta X_n=E_\theta T(\boldsymbol X)g(\boldsymbol X)=(E_\theta T(\boldsymbol X))(E_\theta g(\boldsymbol X))=n\theta E_\theta g(\boldsymbol X).$$
因此，对任意$\theta$有$E_\theta g(\boldsymbol X)=n^{-1}$.
\end{example}
\begin{example}[Basu定理的运用——II]\label{ex:6.2.27}
作为Basu定理的另一个应用，考虑从$n(\mu,\sigma^2)$总体抽样时，样本均值$\bar X$和方差$S^2$的独立性.当然，\autoref{thm:5.3.1}已证明这两个统计量独立，但我们要在这个重要背景下说明Basu定理的用法.先固定$\sigma^2$，让$\mu$变化，$-\infty<\mu<\infty$.由\autoref{ex:6.2.4}，$\bar X$是$\mu$的充分统计量.由\autoref{thm:6.2.25}可推出$n(\mu,\sigma^2/n)$、$-\infty<\mu<\infty$且$\sigma^2/n$已知这一分布族完备.这是$\bar X$的分布，所以$\bar X$是完备统计量.再考虑$S^2$.用类似\autoref{ex:6.2.18}和\autoref{ex:6.2.19}的论证，可说明任意位置参数族中（记住$\sigma^2$固定，$\mu$为位置参数），$S^2$都是辅助统计量.或者，对此正态模型，用\autoref{thm:5.3.1}可知$S^2$的分布依赖固定量$\sigma^2$，却不依赖参数$\mu$.无论采用哪种方式，$S^2$都是辅助统计量，因此由Basu定理，它与完备充分统计量$\bar X$独立.对任意$\mu$和固定的$\sigma^2$，$\bar X$与$S^2$独立.由于$\sigma^2$任意，所以对$\mu$与$\sigma^2$的任意选择，样本均值和方差都独立.注意在$\mu$与$\sigma^2$均未知的模型中，$\bar X$和$S^2$都不是辅助统计量.但通过此论证，我们仍能用Basu定理推出独立性.这类论证有时有用，不过通常证明一个统计量完备，比证明两个统计量独立更难，这一点仍然成立.
\end{example}
应注意，Basu定理的证明未用到充分统计量的“最小性”.事实上，去掉这个词定理仍成立，因为完备统计量的一项基本性质就是最小性.
\begin{theorem}\label{thm:6.2.28}
若最小充分统计量存在，则任何完备统计量也都是最小充分统计量.\footnote{译注：原文这里遗漏“充分”条件.正确结论应针对“完备充分统计量”；例如常数统计量是完备的，却一般不充分，因而不能仅由完备性推出最小充分性.前一句及后面的解释应结合本节充分统计量的语境理解.}
\end{theorem}
所以尽管Basu定理陈述中的“最小”一词是冗余的，保留它是为了提醒：定理中的统计量$T(\boldsymbol X)$是最小充分统计量.（关于完备统计量与最小充分统计量关系的更多内容，见\hyperref[bib:source201]{Lehmann和Scheffé 1950}，以及\hyperref[bib:source276]{Schervish 1995}，第2.1节.）

Basu定理利用完备统计量的概念，给出充分统计量与辅助统计量之间的一种关系.辅助性和完备性还有其他可能的定义.\hyperref[bib:source197]{Lehmann（1981）}讨论了这些定义下充分性与辅助性的一些关系.



\sourcepageanchor{267}
\subsection{似然原则}\label{sec:6.3}
本节研究一个称为似然函数的特定而重要的统计量，它也可以用来概括数据.似然函数有许多用法，本节和后面各章会分别提及其中一些.但本节主要考虑一个论证，它说明如果接受某些其他原则，就必须把似然函数用作数据约简工具.
\subsubsection{似然函数}\label{sec:6.3.1}
\begin{definition}\label{def:6.3.1}\index{SI0619@似然函数}
设$f(\boldsymbol x\mid\theta)$为样本$\boldsymbol X=(X_1,\ldots,X_n)$的联合pdf或pmf.给定观测$\boldsymbol X=\boldsymbol x$，定义为
$$L(\theta\mid\boldsymbol x)=f(\boldsymbol x\mid\theta)$$
的$\theta$的函数，称为似然函数（likelihood function）.
\end{definition}
若$\boldsymbol X$是离散随机向量，则$L(\theta\mid\boldsymbol x)=P_\theta(\boldsymbol X=\boldsymbol x)$.若比较两个参数点的似然函数值，发现
$$P_{\theta_1}(\boldsymbol X=\boldsymbol x)=L(\theta_1\mid\boldsymbol x)>L(\theta_2\mid\boldsymbol x)=P_{\theta_2}(\boldsymbol X=\boldsymbol x),$$
那么我们实际观测的样本在$\theta=\theta_1$时比在$\theta=\theta_2$时更有可能出现.可以将其解释为，$\theta_1$作为$\theta$真值的可能性比$\theta_2$更可信.利用这种信息，人们提出过许多不同方法；不过，考察$\theta$各种可能值下实际观测样本的概率，显然是合理的.这正是似然函数提供的信息.

若$X$是连续实值随机变量，且其pdf关于$x$连续，则当$\varepsilon$很小时，$P_\theta(x-\varepsilon<X<x+\varepsilon)$近似等于$2\varepsilon f(x\mid\theta)=2\varepsilon L(\theta\mid x)$（这来自导数的定义）.因此，
$$\frac{P_{\theta_1}(x-\varepsilon<X<x+\varepsilon)}{P_{\theta_2}(x-\varepsilon<X<x+\varepsilon)}\approx\frac{L(\theta_1\mid x)}{L(\theta_2\mid x)},$$
比较两个参数值处的似然函数，再次给出了观测样本值$x$出现概率的近似比较.

\autoref{def:6.3.1}似乎把似然函数定义成了与pdf或pmf相同的东西.这两个函数唯一的区别在于哪个变量被视为固定，哪个变量在变化.考虑pdf或pmf $f(\boldsymbol x\mid\theta)$时，视$\theta$固定，而$\boldsymbol x$为变量；考虑似然函数$L(\theta\mid\boldsymbol x)$时，视$\boldsymbol x$为观测样本点，而$\theta$遍历所有可能的参数值.
\begin{example}[负二项似然]\label{ex:6.3.2}\index{SI0622@负二项}
设$X$服从$r=3$、成功概率$p$的负二项分布.若观测到$x=2$，似然函数就是$0\le p\le1$上的五次多项式：
$$L(p\mid2)=P_p(X=2)=\binom42p^3(1-p)^2.$$
一般地，若观测到$X=x$，似然函数是次数为$3+x$的多项式，
$$L(p\mid x)=\binom{3+x-1}{x}p^3(1-p)^x.$$
\end{example}
\sourcepageanchor{268}
似然原则规定如何把似然函数用作数据约简工具.
\begin{quote}
\textbf{似然原则：}\index{SI0624@似然原则}\semanticindexanchor{624:268}若两个样本点$\boldsymbol x$和$\boldsymbol y$满足$L(\theta\mid\boldsymbol x)$与$L(\theta\mid\boldsymbol y)$成比例，即存在常数$C(\boldsymbol x,\boldsymbol y)$使
\begin{equation}\label{eq:6.3.1}
L(\theta\mid\boldsymbol x)=C(\boldsymbol x,\boldsymbol y)L(\theta\mid\boldsymbol y)\qquad\text{对所有$\theta$},
\end{equation}
则根据$\boldsymbol x$和$\boldsymbol y$得出的结论应完全一致.
\end{quote}
注意\eqref{eq:6.3.1}中的常数$C(\boldsymbol x,\boldsymbol y)$可随$(\boldsymbol x,\boldsymbol y)$而不同，但不依赖$\theta$.

特别地，当$C(\boldsymbol x,\boldsymbol y)=1$时，似然原则说：若两个样本点产生相同的似然函数，则它们含有相同的关于$\theta$的信息.但似然原则还进一步声称，即使两个样本点的似然函数只是成比例，它们也含有等价的关于$\theta$的信息.理由是：似然函数用来比较各参数值的可信程度；若$L(\theta_2\mid\boldsymbol x)=2L(\theta_1\mid\boldsymbol x)$，则在某种意义上，$\theta_2$比$\theta_1$可信两倍.若\eqref{eq:6.3.1}也成立，则$L(\theta_2\mid\boldsymbol y)=2L(\theta_1\mid\boldsymbol y)$.因此，无论观测到$\boldsymbol x$还是$\boldsymbol y$，都会得出$\theta_2$比$\theta_1$可信两倍的结论.

上一段谨慎地使用“可信”（plausible）一词，而非“有概率”（probable），因为我们往往把$\theta$视为固定值，虽说它未知.此外，尽管$f(\boldsymbol x\mid\theta)$作为$\boldsymbol x$的函数是pdf，却不能保证$L(\theta\mid\boldsymbol x)$作为$\theta$的函数也是pdf.

一种称为信仰推断（fiducial inference）的推断形式，\index{SI0444@信仰推断}\semanticindexanchor{444:268}有时把似然解释为$\theta$的概率.即把$L(\theta\mid\boldsymbol x)$乘以$M(\boldsymbol x)=(\int_{-\infty}^{\infty}L(\theta\mid\boldsymbol x)\,d\theta)^{-1}$（参数空间可数时用求和代替积分），再把$M(\boldsymbol x)L(\theta\mid\boldsymbol x)$解释为$\theta$的pdf（当然须$M(\boldsymbol x)$有限！）.显然，满足\eqref{eq:6.3.1}的$L(\theta\mid\boldsymbol x)$与$L(\theta\mid\boldsymbol y)$给出相同的pdf，因为$C(\boldsymbol x,\boldsymbol y)$只会被吸收入归一化常数.大多数统计学家并不赞同信仰推断理论，但它历史悠久，可追溯至\hyperref[bib:source124]{Fisher（1930）}关于所谓逆概率（inverse probability）的工作，即概率积分变换的一项应用.现在出于历史原因，计算一个信仰分布.
\begin{example}[正态信仰分布]\label{ex:6.3.3}
设$X_1,\ldots,X_n$为iid $n(\mu,\sigma^2)$，$\sigma^2$已知.用\eqref{eq:6.2.4}作为$L(\mu\mid\boldsymbol x)$的表达式，先注意\eqref{eq:6.3.1}成立当且仅当$\bar x=\bar y$，此时
$$C(\boldsymbol x,\boldsymbol y)=\exp\left(-\sum_{i=1}^n(x_i-\bar x)^2/(2\sigma^2)+\sum_{i=1}^n(y_i-\bar y)^2/(2\sigma^2)\right).$$
所以似然原则认为，对于满足$\bar x=\bar y$的任意两个样本点，应得出相同的关于$\mu$的结论.为计算$\mu$的信仰pdf，若定义$M(\boldsymbol x)=n^{n/2}\exp(\sum_{i=1}^n(x_i-\bar x)^2/(2\sigma^2))$，则$M(\boldsymbol x)L(\mu\mid\boldsymbol x)$（作为$\mu$的函数）是$n(\bar x,\sigma^2/n)$的pdf.\footnote{译注：原文此处归一化因子写$n^{n/2}$，与\eqref{eq:6.2.4}不符.按该式积分，指数项之前的因子应为$\sqrt n(2\pi\sigma^2)^{(n-1)/2}$，这样乘积才积分为1.}这就是$\mu$的信仰分布，信仰推断论者可对$\mu$作如下概率计算.

参数$\mu$服从$n(\bar x,\sigma^2/n)$分布，因而$(\mu-\bar x)/(\sigma/\sqrt n)$服从$n(0,1)$.于是
\begin{align*}
.95&=P\left(-1.96<\frac{\mu-\bar x}{\sigma/\sqrt n}<1.96\right)\\
&=P(-1.96\sigma/\sqrt n<\mu-\bar x<1.96\sigma/\sqrt n)\\
&=P(\bar x-1.96\sigma/\sqrt n<\mu<\bar x+1.96\sigma/\sqrt n).
\end{align*}
这些代数运算类似前面的计算，但解释截然不同.这里$\bar x$是已观测数据所给出的固定已知数，而$\mu$是具有正态概率分布的变量.
\end{example}
后面各章讨论具体推断方法时，将介绍似然函数其他更常见的用途.现在则考察一个论证，说明似然原则是另外两个基本原则的必然结果.
\sourcepageanchor{269}
\subsubsection{形式化的似然原则}\label{sec:6.3.2}\index{SI0627@形式化}
对离散分布，似然原则可以从两个直观上更简单的思想推出.对连续分布，在附加一些限制后也成立.本小节只处理离散分布.\hyperref[bib:source22]{Berger和Wolpert（1984）}深入讨论了离散与连续两种情形的似然原则.\hyperref[bib:source35]{Birnbaum（1962）}最早在一篇里程碑式论文中证明这些结果，但这里的叙述更接近Berger和Wolpert.

形式上，把实验$E$定义为三元组$(\boldsymbol X,\theta,\{f(\boldsymbol x\mid\theta)\})$，其中$\boldsymbol X$为随机向量，对参数空间$\Theta$中的某个$\theta$，具有pmf $f(\boldsymbol x\mid\theta)$.实验者知道进行了哪个实验$E$，并观测到特定样本$\boldsymbol X=\boldsymbol x$，便会对$\theta$作出某种推断或结论.以$\operatorname{Ev}(E,\boldsymbol x)$表示这个结论，\index{SI0419@证据}\semanticindexanchor{419:269}含义是由$E$和$\boldsymbol x$产生的关于$\theta$的证据.
\begin{example}[证据函数]\label{ex:6.3.4}
设实验$E$为观测iid $n(\mu,\sigma^2)$变量$X_1,\ldots,X_n$，其中$\sigma^2$已知.因为样本均值$\bar X$是$\mu$的充分统计量，且$E\bar X=\mu$，可以用观测值$\bar X=\bar x$估计$\mu$.为衡量估计精度，通常报告$\bar X$的标准差$\sigma/\sqrt n$.所以可以定义$\operatorname{Ev}(E,\boldsymbol x)=(\bar x,\sigma/\sqrt n)$.这里$\bar x$坐标依赖观测样本$\boldsymbol x$，而$\sigma/\sqrt n$坐标依赖关于$E$的知识.
\end{example}
为把证据函数概念与熟悉的内容联系起来，现在用这些概念重述\autoref{sec:6.2}的充分性原则.
\begin{quote}
\textbf{形式化的充分性原则：}\index{SI0998@形式化}\semanticindexanchor{998:269}考虑实验$E=(\boldsymbol X,\theta,\{f(\boldsymbol x\mid\theta)\})$，设$T(\boldsymbol X)$是$\theta$的充分统计量.若样本点$\boldsymbol x$和$\boldsymbol y$满足$T(\boldsymbol x)=T(\boldsymbol y)$，则$\operatorname{Ev}(E,\boldsymbol x)=\operatorname{Ev}(E,\boldsymbol y)$.
\end{quote}
形式化的充分性原则比\autoref{sec:6.2}的充分性原则略进了一步.那里没有提及实验；这里则约定充分统计量相同时证据相同.由形式化的充分性原则和下面这个十分合理的原则，可以推出似然原则.
\begin{quote}
\textbf{条件性原则（Conditionality Principle）：}\index{SI0125@条件性原则}\semanticindexanchor{125:269}设$E_1=(\boldsymbol X_1,\theta,\{f_1(\boldsymbol x_1\mid\theta)\})$和$E_2=(\boldsymbol X_2,\theta,\{f_2(\boldsymbol x_2\mid\theta)\})$为两个实验，二者只须具有同一未知参数$\theta$.考虑如下混合实验：先观测随机变量$J$，其中$P(J=1)=P(J=2)=1/2$（与$\theta$、$\boldsymbol X_1$或$\boldsymbol X_2$无关），然后进行实验$E_J$.形式上，所进行的实验为$E^*=(\boldsymbol X^*,\theta,\{f^*(\boldsymbol x^*\mid\theta)\})$，其中$\boldsymbol X^*=(j,\boldsymbol X_j)$且$f^*(\boldsymbol x^*\mid\theta)=f^*((j,\boldsymbol x_j)\mid\theta)=\tfrac12f_j(\boldsymbol x_j\mid\theta)$.则
\begin{equation}\label{eq:6.3.2}
\operatorname{Ev}(E^*,(j,\boldsymbol x_j))=\operatorname{Ev}(E_j,\boldsymbol x_j).
\end{equation}
\end{quote}
条件性原则只是说：若随机选取两个实验中的一个，完成所选实验并得到数据$\boldsymbol x$，那么关于$\theta$的信息只依赖实际进行的实验.也就是说，它与从一开始就非随机地决定做该实验、并观测到数据$\boldsymbol x$时获得的信息相同.进行了这个实验而非其他实验这一事实，并没有增加、减少或改变关于$\theta$的知识.
\begin{example}[二项／负二项实验]\label{ex:6.3.5}
设所关心的参数是概率$p$，$0<p<1$，表示某枚硬币投掷后正面朝上的概率.令$E_1$为投掷硬币20次、记录其中正面次数的实验.$E_1$是二项实验，$\{f_1(x_1\mid p)\}$是$\operatorname{binomial}(20,p)$的pmf族.令$E_2$为一直掷到第7次正面出现、记录第7次正面前反面次数的实验.$E_2$是负二项实验.现在假设实验者用随机数表在两个实验间选择，恰好选到$E_2$，得到的数据是第7次正面在第20次试验出现.
\sourcepageanchor{270}
条件性原则认为，实验者此时拥有的关于$\theta$的信息$\operatorname{Ev}(E^*,(2,13))$，与他直接选择做负二项实验、从未考虑过二项实验时会拥有的信息$\operatorname{Ev}(E_2,13)$相同.
\end{example}
现在可由形式化的充分性原则与条件性原则，推出下面的形式化似然原则.
\begin{quote}
\textbf{形式化的似然原则：}\index{SI0627@形式化}\semanticindexanchor{627:270}设有两个实验$E_1=(\boldsymbol X_1,\theta,\{f_1(\boldsymbol x_1\mid\theta)\})$和$E_2=(\boldsymbol X_2,\theta,\{f_2(\boldsymbol x_2\mid\theta)\})$，其未知参数$\theta$相同.设$\boldsymbol x_1^*$、$\boldsymbol x_2^*$分别为$E_1$与$E_2$的样本点，并且
\begin{equation}\label{eq:6.3.3}
L(\theta\mid\boldsymbol x_2^*)=CL(\theta\mid\boldsymbol x_1^*)
\end{equation}
对所有$\theta$成立，其中常数$C$可依赖$\boldsymbol x_1^*$和$\boldsymbol x_2^*$，却不依赖$\theta$.则
$$\operatorname{Ev}(E_1,\boldsymbol x_1^*)=\operatorname{Ev}(E_2,\boldsymbol x_2^*).$$
\end{quote}
形式化似然原则涉及两个实验，\autoref{sec:6.3.1}的似然原则只涉及一个，因此二者不同.不过令$E_2$为$E_1$的精确重复，就能由形式化似然原则推出似然原则.所以形式化似然原则的两个实验设置，在某种程度上是人为安排；重要后果是下面的推论，其证明留作习题（见\autoref{exercise:6.32}）.
\begin{quote}
\textbf{似然原则推论：}若$E=(\boldsymbol X,\theta,\{f(\boldsymbol x\mid\theta)\})$是一个实验，则$\operatorname{Ev}(E,\boldsymbol x)$应只通过$L(\theta\mid\boldsymbol x)$依赖$E$和$\boldsymbol x$.
\end{quote}
现在陈述Birnbaum定理，再考察其有些出人意料的后果.
\begin{theorem}[Birnbaum定理]\label{thm:6.3.6}\index{SI0093@Birnbaum定理}\index{SI0626@Birnbaum定理}
由形式化的充分性原则与条件性原则可推出形式化的似然原则.逆命题也成立.
\end{theorem}
\begin{proof}
只概述证明，细节留作\autoref{exercise:6.33}.设$E_1$、$E_2$、$\boldsymbol x_1^*$、$\boldsymbol x_2^*$如形式化似然原则中所定义，并以$E^*$表示条件性原则中的混合实验.在$E^*$的样本空间上定义统计量
$$T(j,\boldsymbol x_j)=\begin{cases}(1,\boldsymbol x_1^*),&j=1\text{且}\boldsymbol x_1=\boldsymbol x_1^*\text{，或}j=2\text{且}\boldsymbol x_2=\boldsymbol x_2^*,\\(j,\boldsymbol x_j),&\text{其他情形}.\end{cases}$$
利用因子分解定理，可证明$T(J,\boldsymbol X_J)$是$E^*$实验中的充分统计量.于是形式化的充分性原则推出
\begin{equation}\label{eq:6.3.4}
\operatorname{Ev}(E^*,(1,\boldsymbol x_1^*))=\operatorname{Ev}(E^*,(2,\boldsymbol x_2^*)),
\end{equation}
条件性原则推出
\begin{equation}\label{eq:6.3.5}
\begin{aligned}
\operatorname{Ev}(E^*,(1,\boldsymbol x_1^*))&=\operatorname{Ev}(E_1,\boldsymbol x_1^*),\\
\operatorname{Ev}(E^*,(2,\boldsymbol x_2^*))&=\operatorname{Ev}(E_2,\boldsymbol x_2^*),
\end{aligned}
\end{equation}
从而可得$\operatorname{Ev}(E_1,\boldsymbol x_1^*)=\operatorname{Ev}(E_2,\boldsymbol x_2^*)$，即形式化的似然原则.

证明逆命题时，先令一个实验为$E^*$，另一个为$E_j$.可证明$\operatorname{Ev}(E^*,(j,\boldsymbol x_j))=\operatorname{Ev}(E_j,\boldsymbol x_j)$，即条件性原则.然后，若$T(\boldsymbol X)$充分且$T(\boldsymbol x)=T(\boldsymbol y)$，则似然成比例，由形式化似然原则有$\operatorname{Ev}(E,\boldsymbol x)=\operatorname{Ev}(E,\boldsymbol y)$，即形式化的充分性原则.
\end{proof}
\sourcepageanchor{271}
\begin{example}[续\autoref{ex:6.3.5}]\label{ex:6.3.7}\index{SI0126@解释}\index{SI0625@二项与负二项}\index{SI0628@解释}\index{SI0999@解释}
再次考虑二项与负二项实验及两个样本点$x_1=7$（二项实验20次中7次正面）和$x_2=13$（第7次正面在第20次投掷出现）.似然函数分别为
$$L(p\mid x_1=7)=\binom{20}7p^7(1-p)^{13}\qquad\text{（二项实验）},$$
以及
$$L(p\mid x_2=13)=\binom{19}6p^7(1-p)^{13}\qquad\text{（负二项实验）}.$$
这两个似然函数成比例，所以形式化似然原则认为两种情形应得出相同的关于$p$的结论.特别地，该原则断言：第一种情形因完成20次试验而停止抽样，第二种情形因观测到第7次正面而停止抽样，这一事实与我们关于$p$的结论无关.\hyperref[bib:source206]{Lindley和Phillips（1976）}对二项–负二项推断问题作了深入讨论.
\end{example}
不同实验产生等价推断这一点，可以通过考察Birnbaum定理证明中定义的充分统计量$T$及样本点$x_1^*=7$、$x_2^*=13$进一步阐明.对混合实验中除$(1,7)$和$(2,13)$之外的任何样本点，$T$会告诉我们做的是二项还是负二项实验，以及实验结果.但对$(1,7)$和$(2,13)$，有$T(1,7)=T(2,13)=(1,7)$.若只用充分统计量推断且$T=(1,7)$，则只知道观测到20次中有7次正面，却不知道7和20中哪一个是固定量.

许多常用统计程序违反形式化似然原则.用这些程序，会对\autoref{ex:6.3.5}中的两个实验得出不同结论.这似乎奇怪，因为由Birnbaum定理，此时必违反充分性原则或条件性原则之一.下面更仔细地考察这两个原则.

形式化的充分性原则，本质上与\autoref{sec:6.1}所讨论的相同.那里说明全部关于$\theta$的信息都包含在充分统计量中，知道整个样本无法增加任何信息.所以把证据建立在充分统计量上，是非常可信的原则.这个原则一个容易导致违反的弱点，是它强烈依赖模型.如\autoref{ex:6.2.9}之后所述，相信该原则就必须相信模型，而这未必容易做到.

大多数数据分析者分析数据集时都会做某种“模型检查”.多数模型检查必然基于充分统计量以外的统计量.例如，考察模型残差（residuals）是常见做法，残差衡量模型未能解释的数据变异.（\autoref{sec:11}和\autoref{sec:12}将更详细地介绍残差.）这种做法立即违反充分性原则，因为残差并非基于充分统计量.（当然，它也直接违反似然原则.）因此必须认识到，在考虑充分性原则或似然原则之前，我们必须认可模型.

条件性原则非形式地说就是“只有实际进行的实验才重要”.即在\autoref{ex:6.3.5}中，若做的是二项实验，那么未进行的负二项实验就不应以任何方式影响关于$\theta$的结论.这一原则也显得非常可信.

既然违反形式化似然原则意味着违反充分性或条件性原则之一，那么统计实践为何仍会违反它？一些作者探讨过此问题，包括\hyperref[bib:source103]{Durbin（1970）}和\mbox{\hyperref[bib:source185]{Kalbfleisch（1975）}.}\allowbreak{}Kalbfleisch的一种论点是：形式化似然原则的证明并不令人信服，因为它在忽略条件性原则的情况下应用充分性原则.\autoref{thm:6.3.6}证明中所用的充分统计量$T(J,\boldsymbol X_J)$，是在混合实验上定义的.
\sourcepageanchor{272}
若先援引条件性原则，就必须分别为每个实验定义充分统计量.此时不再能推出形式化似然原则.（Birnbaum定理证明的关键，是$T(J,\boldsymbol X_J)$能对来自两个实验的样本点取同一个值.分别定义充分统计量时就不会如此.）

无论如何，许多直观上有吸引力的推断程序确实违反似然原则，所以它并未得到所有统计学家的普遍接受.不过，它在数学上有吸引力，也确实提供一种有用的数据约简技术.



\subsection{等变性原则}\label{sec:6.4}\index{SI0403@等变性（另见不变性）}
前两节以如下方式描述数据约简原则：指定样本的一个函数$T(\boldsymbol x)$，并规定若两个样本点$\boldsymbol x$和$\boldsymbol y$满足$T(\boldsymbol x)=T(\boldsymbol y)$，那么无论观测到哪一个，都应对$\theta$作出相同推断.使用充分性原则时，函数$T(\boldsymbol x)$是充分统计量；使用似然原则时，$T(\boldsymbol x)$的“值”是所有与$L(\theta\mid\boldsymbol x)$成比例的似然函数组成的集合.等变性原则以略有不同的方式描述数据约简技术.每次应用等变性原则时，都指定函数$T(\boldsymbol x)$，但若$T(\boldsymbol x)=T(\boldsymbol y)$，该原则要求观测到$\boldsymbol x$所作的推断与观测到$\boldsymbol y$所作的推断具有某种关系，虽然二者未必相同.这种对推断程序的限制有时能简化分析，正如前面各节的数据约简原则一样.\footnote{与许多其他教材（\hyperref[bib:source276]{Schervish 1995}；\hyperref[bib:source200]{Lehmann和Casella 1998}；\hyperref[bib:source311]{Stuart、Ord和Arnold 1999）}一样，我们区分等变性（equivariance）和不变性（invariance）：前者指数据变换时估计值按规定方式变化，后者指数据变换时估计值保持不变.}

虽然通常合称等变性原则，下面所述的数据约简技术实际上结合了两种不同的等变性考虑.

第一类可称为测量等变性（measurement equivariance）.\index{SI0401@测量}\semanticindexanchor{401:272}它规定推断不应依赖所采用的测量尺度.例如，设两位林业人员要估计一片森林树木的平均直径.第一位使用以英寸表示的树径数据，第二位使用以米表示的同一数据.现在要求两人都给出以英寸表示的估计.（第二位可以先方便地以米估计平均直径，再把估计转换为英寸.）测量等变性要求两人得出相同估计.毫无疑问，几乎所有人都会同意这种等变性是合理的.

第二类等变性实际上是不变性，可称为形式不变性（formal invariance）.\index{SI0589@形式化}\semanticindexanchor{589:272}它规定：若两个推断问题的数学模型具有相同的形式结构，就应使用同一推断程序.必须相同的模型要素是参数空间$\Theta$、样本pdf或pmf集合$\{f(\boldsymbol x\mid\theta):\theta\in\Theta\}$，以及允许作出的推断的集合和错误推断的后果.最后这一要素此前未多作讨论；本节假设可能推断的集合就是$\Theta$，即推断只是从$\Theta$选取一个元素，作为$\theta$真值的估计或猜测.形式不变性只关心涉及的数学对象，不关心实验的物理描述.例如，两个问题的参数空间都可能是$\Theta=\{\theta:\theta>0\}$，但一个问题中$\theta$表示美国一打鸡蛋的平均价格（以美分计），另一个问题中$\theta$表示肯尼亚长颈鹿的平均身高（以米计）.形式不变性仍将这两个参数空间视为相同，因为它们指向同一实数集合.
\begin{quote}
\textbf{等变性原则：}\index{SI0404@等变性原则}\semanticindexanchor{404:272}若$\boldsymbol Y=g(\boldsymbol X)$是测量尺度的改变，且$\boldsymbol Y$的模型与$\boldsymbol X$的模型具有相同形式结构，则推断程序应既满足测量等变性，又满足形式等变性.
\end{quote}
\sourcepageanchor{273}
下面说明这两种等变性概念如何结合，提供有用的数据约简.
\begin{example}[二项等变性]\label{ex:6.4.1}\index{SI0400@二项}
设$X$服从样本容量$n$已知、成功概率$p$未知的二项分布.设$T(x)$为观测到$X=x$时采用的$p$估计.除了用成功次数$X$推断$p$，也可以用失败次数$Y=n-X$.$Y$也服从二项分布，参数为$(n,q=1-p)$.设$T^*(y)$为观测到$Y=y$时采用的$q$估计，则此时$p$的估计为$1-T^*(y)$.若观测到$x$次成功，$p$的估计是$T(x)$；但$x$次成功意味着$n-x$次失败，$1-T^*(n-x)$也是$p$的估计.测量等变性要求两个估计相等，即$T(x)=1-T^*(n-x)$，因为从$X$到$Y$只是测量尺度的改变.此外，基于$X$和$Y$的推断问题形式结构相同，二者都服从$\operatorname{binomial}(n,\theta)$，$0\le\theta\le1$.所以形式不变性要求对所有$z=0,\ldots,n$有$T(z)=T^*(z)$.因此测量与形式不变性共同要求
\begin{equation}\label{eq:6.4.1}
T(x)=1-T^*(n-x)=1-T(n-x).
\end{equation}
若只考虑满足\eqref{eq:6.4.1}的估计量，所愿考虑的估计量集合便大为缩小和简化.指定任意估计量需要指定$T(0),T(1),\ldots,T(n)$；指定满足\eqref{eq:6.4.1}的估计量则只须指定$T(0),T(1),\ldots,T([n/2])$，其中$[n/2]$是不大于$n/2$的最大整数.其余$T(x)$值由已指定的值及\eqref{eq:6.4.1}决定.例如，$T(n)=1-T(0)$、$T(n-1)=1-T(1)$.等变性原则总是实现这种类型的数据约简：对某些样本点应作的推断，决定了对其他样本点应作的推断.

本问题的两个等变估计量是$T_1(x)=x/n$和$T_2(x)=.9(x/n)+.1(.5)$.$T_1(x)$用样本成功比例估计$p$.$T_2(x)$把样本比例向$.5$“收缩”；若有理由认为$p$接近$.5$，这样的估计可能合理.容易验证二者均满足\eqref{eq:6.4.1}，故都等变.$T_3(x)=.8(x/n)+.2(1)$则不等变，因为$T_3(0)=.2\ne0=1-T_3(n-0)$，不满足\eqref{eq:6.4.1}.关于测量与形式不变性的更多内容，见\autoref{exercise:6.39}.
\end{example}
\autoref{ex:6.4.1}以及任何等变性论证的关键之一，是变换的选择.该例使用的数据变换为$Y=n-X$.等变性原则应用中的变换（测量尺度的改变），用样本空间上的一组函数描述，称为变换群（group of transformations）.
\begin{definition}\label{def:6.4.2}\index{SI0466@变换的}
从样本空间$\mathcal X$到自身的满射函数集合$\{g(\boldsymbol x):g\in\mathcal G\}$称为$\mathcal X$上的变换群，若
\begin{enumerate}[label=(\roman*)]
\item\label{item:6.4.2i}（逆）对每个$g\in\mathcal G$，存在$g'\in\mathcal G$，使对所有$\boldsymbol x\in\mathcal X$有$g'(g(\boldsymbol x))=\boldsymbol x$.
\item\label{item:6.4.2ii}（复合）对每个$g\in\mathcal G$和$g'\in\mathcal G$，存在$g''\in\mathcal G$，使对所有$\boldsymbol x\in\mathcal X$有$g'(g(\boldsymbol x))=g''(\boldsymbol x)$.
\item\label{item:6.4.2iii}（单位元）由$e(\boldsymbol x)=\boldsymbol x$定义的恒等映射$e$是$\mathcal G$的元素.
\end{enumerate}
\end{definition}
有时第三个要求被列为群定义的一部分.但\hyperref[item:6.4.2iii]{(iii)}是\hyperref[item:6.4.2i]{(i)}和\hyperref[item:6.4.2ii]{(ii)}的推论，不必单独验证（见\autoref{exercise:6.38}）.
\begin{example}[续\autoref{ex:6.4.1}]\label{ex:6.4.3}
此问题只涉及两个变换，可设$\mathcal G=\{g_1,g_2\}$，其中$g_1(x)=n-x$、$g_2(x)=x$.条件\hyperref[item:6.4.2i]{(i)}和\hyperref[item:6.4.2ii]{(ii)}容易验证.取$g'=g$即验证\hyperref[item:6.4.2i]{(i)}，也就是说每个元素都是自身的逆.例如，
$$g_1(g_1(x))=g_1(n-x)=n-(n-x)=x.$$
\sourcepageanchor{274}
在\hyperref[item:6.4.2ii]{(ii)}中，若$g'=g$，取$g''=g_2$；若$g'\ne g$，则$g''=g_1$满足等式.例如，取$g'\ne g=g_1$，则
$$g_2(g_1(x))=g_2(n-x)=n-x=g_1(x).$$
\end{example}
要使用等变性原则，必须能对变换后的问题应用形式不变性.即改变测量尺度后仍应具有相同形式结构.由于结构不变，我们希望底层模型或分布族保持不变.下一个定义概括这一要求.
\begin{definition}\label{def:6.4.4}\index{SI0588@分布}
设$\mathcal F=\{f(\boldsymbol x\mid\theta):\theta\in\Theta\}$是$\boldsymbol X$的pdf或pmf集合，$\mathcal G$是样本空间$\mathcal X$的变换群.若对每个$\theta\in\Theta$和$g\in\mathcal G$，存在唯一$\theta'\in\Theta$，使得当$\boldsymbol X$具有分布$f(\boldsymbol x\mid\theta)$时，$\boldsymbol Y=g(\boldsymbol X)$具有分布$f(\boldsymbol y\mid\theta')$，则称$\mathcal F$在群$\mathcal G$下不变.
\end{definition}
\begin{example}[\autoref{ex:6.4.1}的结束]\label{ex:6.4.5}
二项问题中须同时检查$g_1$和$g_2$.若$X\sim\operatorname{binomial}(n,p)$，则$g_1(X)=n-X\sim\operatorname{binomial}(n,1-p)$，故$p'=1-p$，这里$p$起\autoref{def:6.4.4}中$\theta$的作用.又$g_2(X)=X\sim\operatorname{binomial}(n,p)$，故此时$p'=p$.因此二项pmf集合在群$\mathcal G=\{g_1,g_2\}$下不变.
\end{example}
\autoref{ex:6.4.1}的变换群只有两个元素.许多情况下变换群是无限的，下例说明这一点（另见\autoref{exercise:6.41}和\autoref{exercise:6.42}）.
\begin{example}[正态位置不变性]\label{ex:6.4.6}\index{SI0405@解释}\index{SI0464@位置}\index{SI0591@正态}
设$X_1,\ldots,X_n$为iid $n(\mu,\sigma^2)$，$\mu$和$\sigma^2$均未知.考虑变换群$\mathcal G=\{g_a(\boldsymbol x):-\infty<a<\infty\}$，其中$g_a(x_1,\ldots,x_n)=(x_1+a,\ldots,x_n+a)$.要验证它是群，须验证\autoref{def:6.4.2}的\hyperref[item:6.4.2i]{(i)}和\hyperref[item:6.4.2ii]{(ii)}.对\hyperref[item:6.4.2i]{(i)}，注意
\begin{align*}
g_{-a}(g_a(x_1,\ldots,x_n))&=g_{-a}(x_1+a,\ldots,x_n+a)\\
&=(x_1+a-a,\ldots,x_n+a-a)\\
&=(x_1,\ldots,x_n).
\end{align*}
所以$g=g_a$时，$g'=g_{-a}$满足\hyperref[item:6.4.2i]{(i)}.对\hyperref[item:6.4.2ii]{(ii)}，注意
\begin{align*}
g_{a_2}(g_{a_1}(x_1,\ldots,x_n))&=g_{a_2}(x_1+a_1,\ldots,x_n+a_1)\\
&=(x_1+a_1+a_2,\ldots,x_n+a_1+a_2)\\
&=g_{a_1+a_2}(x_1,\ldots,x_n).
\end{align*}
所以$g=g_{a_1}$、$g'=g_{a_2}$时，$g''=g_{a_1+a_2}$满足\hyperref[item:6.4.2ii]{(ii)}，从而\autoref{def:6.4.2}得到验证.$\mathcal G$是变换群.

此问题中的$\mathcal F$是由“$X_1,\ldots,X_n$为iid $n(\mu,\sigma^2)$，其中某个$-\infty<\mu<\infty$、$\sigma^2>0$”定义的全部联合密度$f(x_1,\ldots,x_n\mid\mu,\sigma^2)$的集合.对任意$-\infty<a<\infty$，由
$$(Y_1,\ldots,Y_n)=g_a(X_1,\ldots,X_n)=(X_1+a,\ldots,X_n+a)$$
定义的$Y_1,\ldots,Y_n$是iid $n(\mu+a,\sigma^2)$随机变量.故$\boldsymbol Y=g_a(\boldsymbol X)$的联合分布属于$\mathcal F$，从而$\mathcal F$在$\mathcal G$下不变.用\autoref{def:6.4.4}的记号，若$\theta=(\mu,\sigma^2)$，则$\theta'=(\mu+a,\sigma^2)$.
\end{example}
再提醒一次，等变性原则由两类不同的等变性组成.测量等变性直观上合理.许多人想到等变性原则时，认为它只指测量等变性.若真如此，该原则大概会被普遍接受.但另一原则——形式不变性——很不一样.
\sourcepageanchor{275}
它把数学结构相同的任意两个问题等同起来，不考虑它们试图解释的物理现实.它认为即使物理现实很不相同，同一推断程序仍然适用；这一假设有时难以证明合理.

不过，与充分性和似然原则一样，等变性原则是一种数据约简技术，通过规定相关样本点上必须作出的其他推断来限制推断.三种原则都规定不同样本点上推断之间的关系，限制允许推断的集合，并由此简化问题分析.



\subsection{习题}\label{sec:6.5}
\begin{exercise}\label{exercise:6.1}
设$X$是来自$n(0,\sigma^2)$总体的一个观测.$|X|$是否为充分统计量？
\end{exercise}
\begin{exercise}\label{exercise:6.2}
设$X_1,\ldots,X_n$为独立随机变量，其密度为
$$f_{X_i}(x\mid\theta)=\begin{cases}e^{i\theta-x},&x\ge i\theta,\\0,&x<i\theta.\end{cases}$$
证明$T=\min_i(X_i/i)$是$\theta$的充分统计量.
\end{exercise}
\begin{exercise}\label{exercise:6.3}
设$X_1,\ldots,X_n$为来自下列pdf的随机样本：
$$f(x\mid\mu,\sigma)=\frac1\sigma e^{-(x-\mu)/\sigma},\qquad\mu<x<\infty,\quad0<\sigma<\infty.$$
求$(\mu,\sigma)$的二维充分统计量.
\end{exercise}
\begin{exercise}\label{exercise:6.4}
证明\autoref{thm:6.2.10}.
\end{exercise}
\begin{exercise}\label{exercise:6.5}
设$X_1,\ldots,X_n$为独立随机变量，其pdf为
$$f(x_i\mid\theta)=\begin{cases}\dfrac1{2i\theta},&-i(\theta-1)<x_i<i(\theta+1),\\0,&\text{其他情形},\end{cases}$$
其中$\theta>0$.求$\theta$的二维充分统计量.
\end{exercise}
\begin{exercise}\label{exercise:6.6}
设$X_1,\ldots,X_n$为来自$\operatorname{gamma}(\alpha,\beta)$总体的随机样本.求$(\alpha,\beta)$的二维充分统计量.
\end{exercise}
\begin{exercise}\label{exercise:6.7}
设$f(x,y\mid\theta_1,\theta_2,\theta_3,\theta_4)$是$\mathbb R^2$中左下角为$(\theta_1,\theta_2)$、右上角为$(\theta_3,\theta_4)$的矩形上均匀分布的二维pdf.参数满足$\theta_1<\theta_3$和$\theta_2<\theta_4$.设$(X_1,Y_1),\ldots,(X_n,Y_n)$为来自该pdf的随机样本.求$(\theta_1,\theta_2,\theta_3,\theta_4)$的四维充分统计量.
\end{exercise}
\begin{exercise}\label{exercise:6.8}\index{SI0738@充分}
设$X_1,\ldots,X_n$为来自具有位置pdf $f(x-\theta)$的总体的随机样本.说明次序统计量$T(X_1,\ldots,X_n)=(X_{(1)},\ldots,X_{(n)})$是$\theta$的充分统计量，且不能再进一步约简.\footnote{译注：后半句对任意给定位置密度并不成立；例如已知方差的正态位置族可约简到样本均值.若$f$本身未知、考虑整个相应非参数族，则需明确这一额外范围.原题保留，参见\autoref{exercise:6.29}.}
\end{exercise}
\begin{exercise}\label{exercise:6.9}
对以下每个分布，设$X_1,\ldots,X_n$为随机样本.求$\theta$的最小充分统计量.
\begin{enumerate}[label=(\alph*)]
\item $f(x\mid\theta)=\dfrac1{\sqrt{2\pi}}e^{-(x-\theta)^2/2}$，$-\infty<x<\infty$，$-\infty<\theta<\infty$（正态）.
\item $f(x\mid\theta)=e^{-(x-\theta)}$，$\theta<x<\infty$，$-\infty<\theta<\infty$（位置指数）.
\item $f(x\mid\theta)=\dfrac{e^{-(x-\theta)}}{(1+e^{-(x-\theta)})^2}$，$-\infty<x<\infty$，$-\infty<\theta<\infty$（logistic）.
\item $f(x\mid\theta)=\dfrac1{\pi[1+(x-\theta)^2]}$，$-\infty<x<\infty$，$-\infty<\theta<\infty$（Cauchy）.
\item $f(x\mid\theta)=\tfrac12e^{-|x-\theta|}$，$-\infty<x<\infty$，$-\infty<\theta<\infty$（双指数）.
\end{enumerate}
\end{exercise}
\begin{exercise}\label{exercise:6.10}\index{SI1010@均匀}
说明\autoref{ex:6.2.15}中求得的$\operatorname{uniform}(\theta,\theta+1)$的最小充分统计量不完备.
\end{exercise}
\sourcepageanchor{276}
\begin{exercise}\label{exercise:6.11}\index{SI0020@位置族}
参见\autoref{exercise:6.9}中的pdf.对每种情形，以$X_{(1)}<\cdots<X_{(n)}$表示排序后的样本，定义$Y_i=X_{(n)}-X_{(i)}$，$i=1,\ldots,n-1$.
\begin{enumerate}[label=(\alph*)]
\item 对\autoref{exercise:6.9}的每个pdf，验证$(Y_1,\ldots,Y_{n-1})$对$\theta$为辅助统计量.尝试像\autoref{ex:6.2.18}那样，证明一个一次处理所有这些族的一般定理.
\item 对每种情形，确定$(Y_1,\ldots,Y_{n-1})$是否与最小充分统计量独立.
\end{enumerate}
\end{exercise}
\begin{exercise}\label{exercise:6.12}\index{SI0022@样本量}
许多问题中一个自然的辅助统计量是样本容量.例如，设随机变量$N$取值$1,2,\ldots$，已知概率为$p_1,p_2,\ldots$，其中$\sum p_i=1$.观测到$N=n$后，进行$n$次成功概率为$\theta$的Bernoulli试验，得到$X$次成功.
\begin{enumerate}[label=(\alph*)]
\item 证明$(X,N)$最小充分，且$N$对$\theta$为辅助统计量.（注意与\autoref{sec:4.4}所述某些分层模型的相似之处.）
\item 证明估计量$X/N$对$\theta$无偏，方差为$\theta(1-\theta)E(1/N)$.
\end{enumerate}
\end{exercise}
\begin{exercise}\label{exercise:6.13}
设$X_1$和$X_2$为来自pdf $f(x\mid\alpha)=\alpha x^{\alpha-1}e^{-x^\alpha}$、$x>0$、$\alpha>0$的iid观测.说明$(\log X_1)/(\log X_2)$是辅助统计量.
\end{exercise}
\begin{exercise}\label{exercise:6.14}
设$X_1,\ldots,X_n$为来自位置族的随机样本.说明$M-\bar X$为辅助统计量，其中$M$为样本中位数.
\end{exercise}
\begin{exercise}\label{exercise:6.15}
设$X_1,\ldots,X_n$为iid $n(\theta,a\theta^2)$，其中$a$是已知常数且$\theta>0$.
\begin{enumerate}[label=(\alph*)]
\item 说明参数空间不含二维开集.
\item 说明统计量$T=(\bar X,S^2)$是$\theta$的充分统计量，但其分布族不完备.
\end{enumerate}
\end{exercise}
\begin{exercise}\label{exercise:6.16}
遗传建模中的一个著名例子（\hyperref[bib:source312]{Tanner 1996}，或\hyperref[bib:source97]{Dempster、Laird和Rubin 1977）}是遗传连锁多项模型，其中观测多项向量$(x_1,x_2,x_3,x_4)$，单元概率为$(\tfrac12+\tfrac\theta4,\tfrac14(1-\theta),\tfrac14(1-\theta),\tfrac\theta4)$.
\begin{enumerate}[label=(\alph*)]
\item 说明这是弯曲指数族.
\item 求$\theta$的充分统计量.
\item 求$\theta$的最小充分统计量.
\end{enumerate}
\end{exercise}
\begin{exercise}\label{exercise:6.17}
设$X_1,\ldots,X_n$为iid几何分布变量，
$$P_\theta(X=x)=\theta(1-\theta)^{x-1},\qquad x=1,2,\ldots,\quad0<\theta<1.$$
说明$\sum X_i$对$\theta$充分，并求$\sum X_i$的分布族.这个族是否完备？
\end{exercise}
\begin{exercise}\label{exercise:6.18}\index{SI0118@泊松}
设$X_1,\ldots,X_n$为iid $\operatorname{Poisson}(\lambda)$.说明$\sum X_i$的分布族完备.不用\autoref{thm:6.2.25}证明其完备性.
\end{exercise}
\begin{exercise}\label{exercise:6.19}
随机变量$X$按下列某一分布取值0、1、2：
\begin{center}
\begin{tabular}{lcccc}
&$P(X=0)$&$P(X=1)$&$P(X=2)$&\\\hline
分布1&$p$&$3p$&$1-4p$&$0<p<\tfrac14$\\
分布2&$p$&$p^2$&$1-p-p^2$&$0<p<\tfrac12$
\end{tabular}
\end{center}
对每种情形，确定$X$的分布族是否完备.
\end{exercise}
\begin{exercise}\label{exercise:6.20}
对以下各pdf，设$X_1,\ldots,X_n$为iid观测.求一个完备充分统计量，或说明不存在这样的统计量.
\begin{enumerate}[label=(\alph*)]
\item $f(x\mid\theta)=2x/\theta^2$，$0<x<\theta$，$\theta>0$.
\item $f(x\mid\theta)=\theta/(1+x)^{1+\theta}$，$0<x<\infty$，$\theta>0$.
\item $f(x\mid\theta)=(\log\theta)\theta^x/(\theta-1)$，$0<x<1$，$\theta>1$.
\item $f(x\mid\theta)=e^{-(x-\theta)}\exp(-e^{-(x-\theta)})$，$-\infty<x<\infty$，$-\infty<\theta<\infty$.
\item $f(x\mid\theta)=\binom2x\theta^x(1-\theta)^{2-x}$，$x=0,1,2$，$0\le\theta\le1$.
\end{enumerate}
\end{exercise}
\sourcepageanchor{277}
\begin{exercise}\label{exercise:6.21}
设$X$为来自下列pdf的一个观测：
$$f(x\mid\theta)=\left(\frac\theta2\right)^{|x|}(1-\theta)^{1-|x|},\qquad x=-1,0,1,\quad0\le\theta\le1.$$
\begin{enumerate}[label=(\alph*)]
\item $X$是否为完备充分统计量？
\item $|X|$是否为完备充分统计量？
\item $f(x\mid\theta)$是否属于指数类？
\end{enumerate}
\end{exercise}
\begin{exercise}\label{exercise:6.22}
设$X_1,\ldots,X_n$为来自具有下列pdf的总体的随机样本：
$$f(x\mid\theta)=\theta x^{\theta-1},\qquad0<x<1,\quad\theta>0.$$
\begin{enumerate}[label=(\alph*)]
\item $\sum X_i$是否对$\theta$充分？
\item 求$\theta$的完备充分统计量.
\end{enumerate}
\end{exercise}
\begin{exercise}\label{exercise:6.23}
设$X_1,\ldots,X_n$为来自区间$(\theta,2\theta)$上均匀分布的随机样本，$\theta>0$.求$\theta$的最小充分统计量.该统计量是否完备？
\end{exercise}
\begin{exercise}\label{exercise:6.24}
考虑分布族
$$\mathcal P=\{P_\lambda(X=x):P_\lambda(X=x)=\lambda^xe^{-\lambda}/x!;\ x=0,1,2,\ldots;\ \lambda=0\text{或}1\}.$$
这是把$\lambda$限制为0或1的Poisson族.说明$\mathcal P$不完备，从而展示完备性可能依赖参数的取值范围.（见\autoref{exercise:6.15}和\autoref{exercise:6.18}.）
\end{exercise}
\begin{exercise}\label{exercise:6.25}
我们已见到若干关于指数族充分性及相关概念的定理.\autoref{thm:5.2.11}给出了某统计量的分布，\autoref{thm:6.2.10}刻画它的充分性，\autoref{thm:6.2.25}刻画其完备性.但若该族弯曲，则\autoref{thm:6.2.25}的开集条件不满足.此时\autoref{thm:6.2.10}的充分统计量是否仍然最小？对\autoref{thm:6.2.10}中的$T(\boldsymbol x)$应用\autoref{thm:6.2.13}，建立以下结论：
\begin{enumerate}[label=(\alph*)]
\item 在$n(\mu,\mu)$族中，$(\sum X_i,\sum X_i^2)$充分，但非最小充分.
\item 在$n(\mu,\mu)$族中，$\sum X_i^2$最小充分.
\item 在$n(\mu,\mu^2)$族中，$(\sum X_i,\sum X_i^2)$最小充分.
\item 在$n(\mu,\sigma^2)$族中，$(\sum X_i,\sum X_i^2)$最小充分.
\end{enumerate}
\end{exercise}
\begin{exercise}\label{exercise:6.26}\index{SI1006@最小}
利用\autoref{thm:6.6.5}证明，给定样本$X_1,X_2,\ldots,X_n$，以下统计量最小充分.
\begin{center}
\begin{tabular}{cll}
&统计量&分布\\\hline
(a)&$\bar X$&$n(\theta,1)$\\
(b)&$\sum X_i$&$\operatorname{gamma}(\alpha,\beta)$，$\alpha$已知\\
(c)&$\max X_i$&$\operatorname{uniform}(0,\theta)$\\
(d)&$X_{(1)},X_{(2)},\ldots,X_{(n)}$&$\operatorname{Cauchy}(\theta,1)$\\
(e)&$X_{(1)},X_{(2)},\ldots,X_{(n)}$&$\operatorname{logistic}(\mu,\beta)$
\end{tabular}
\end{center}
\end{exercise}
\begin{exercise}\label{exercise:6.27}\index{SI0324@逆高斯}
设$X_1,X_2,\ldots,X_n$为来自逆高斯分布（inverse Gaussian distribution）的随机样本，其pdf为
$$f(x\mid\mu,\lambda)=\left(\frac\lambda{2\pi x^3}\right)^{1/2}e^{-\lambda(x-\mu)^2/(2\mu^2x)},\qquad0<x<\infty.$$
\begin{enumerate}[label=(\alph*)]
\item 说明统计量\footnote{译注：原文$T$的分母写$\sum_i1/X_i-1/\bar X$；与后述独立性和卡方分布相符的常见统计量应为$T=n/(\sum_i1/X_i-n/\bar X)$.原式保留.另外“充分且完备”应理解为二者组成的联合统计量.}
$$\bar X=\frac1n\sum_{i=1}^nX_i\qquad\text{及}\qquad T=\frac n{\sum_{i=1}^n\frac1{X_i}-\frac1{\bar X}}$$
充分且完备.
\item 当$n=2$时，说明$\bar X$服从逆高斯分布，$n\lambda/T$服从$\chi_{n-1}^2$分布，且二者独立.（\hyperref[bib:source281]{Schwarz和Samanta 1991}处理了一般情形.）
\end{enumerate}
逆高斯分布有许多应用，尤其用于寿命建模.参见\hyperref[bib:source74]{Chhikara和Folks（1989）}\footnote{译注：原书正文将第一作者姓氏拼作Chhikara，书后对应条目[74]作Chikkara；两处拼写分别保留，链接至该条目.}或\hyperref[bib:source286]{Seshadri（1993）}的著作.
\end{exercise}
\sourcepageanchor{278}
\begin{exercise}\label{exercise:6.28}
证明\autoref{thm:6.6.5}.提示：先由\autoref{thm:6.2.13}确立$T(\boldsymbol X)$在族
$$\{f_0(\boldsymbol x),f_1(\boldsymbol x),\ldots,f_k(\boldsymbol x)\}$$
中的最小充分性.然后论证，在$\mathcal F$中充分的任一统计量都必须是$T(\boldsymbol x)$的函数.\footnote{译注：最后一句的函数方向在原文中颠倒了.最小充分性要求$T$是任一充分统计量的函数；先把在$\mathcal F$中充分的统计量限制到上述有限子族，再应用该子族中$T$的最小充分性，即得所需方向.}
\end{exercise}
\begin{exercise}\label{exercise:6.29}
最小充分性的概念可以推广到参数分布族之外.证明：若$X_1,X_2,\ldots,X_n$是来自未知密度$f$的随机样本，则次序统计量最小充分.

提示：利用\autoref{thm:6.6.5}，取$\{f_0(\boldsymbol x),f_1(\boldsymbol x),\ldots,f_k(\boldsymbol x)\}$为logistic密度.
\end{exercise}
\begin{exercise}\label{exercise:6.30}\index{SI0115@指数}
设$X_1,\ldots,X_n$为来自pdf $f(x\mid\mu)=e^{-(x-\mu)}$的随机样本，其中$-\infty<\mu<x<\infty$.
\begin{enumerate}[label=(\alph*)]
\item 说明$X_{(1)}=\min_iX_i$是完备充分统计量.
\item 利用Basu定理说明$X_{(1)}$与$S^2$独立.
\end{enumerate}
\end{exercise}
\begin{exercise}\label{exercise:6.31}\index{SI0710@巧算}
\hyperref[bib:source42]{Boos和Hughes-Oliver（1998）}详述了应用Basu定理简化计算的若干实例.下面列出其中几个.
\begin{enumerate}[label=(\alph*)]
\item 设$X_1,X_2,\ldots,X_n$独立同分布于$n(\mu,\sigma^2)$，其中$\sigma^2$已知.
\begin{enumerate}[label=(\roman*)]
\item 说明$\bar X$对于$\mu$完备充分，而$S^2$是辅助统计量.因而由Basu定理，$\bar X$与$S^2$独立.
\item 说明即使$\sigma^2$未知，这一独立性仍成立，因为是否知道$\sigma^2$不影响这些分布.（将此证明与较为复杂的\hyperref[thm:5.3.1]{定理5.3.1(a)}比较.）
\end{enumerate}
\item 蒙特卡罗巧计（Monte Carlo swindle）是一种改进方差估计的技巧.设随机变量$X_1,X_2,\ldots,X_n$独立同分布于$N(\mu,\sigma^2)$，我们希望计算中位数$M$的方差.
\begin{enumerate}[label=(\roman*)]
\item 说明$\operatorname{Var}(M)=\operatorname{Var}(M-\bar X)+\operatorname{Var}(\bar X)$；因此，对于$\operatorname{Var}(M)$，我们只需模拟其中的$\operatorname{Var}(M-\bar X)$部分（因为已知$\operatorname{Var}(\bar X)=\sigma^2/n$）.
\item 通过说明$M$的渐近方差约为$2[\operatorname{Var}(M)]^2/(N-1)$，而$M-\bar X$的渐近方差约为$2[\operatorname{Var}(M-\bar X)]^2/(N-1)$，证明巧计估计更为精确，其中$N$是蒙特卡罗样本数.\footnote{译注：原文这里说“$M$的渐近方差”和“$M-\bar X$的渐近方差”，但随模拟次数$N$减小的量应指用$N$次模拟所得的相应\emph{方差估计量}的方差.所给近似式还采用了被模拟统计量近似正态的条件；对一般分布，样本方差的渐近方差涉及四阶中心矩.}
\end{enumerate}
\item \begin{enumerate}[label=(\roman*)]
\item 若$X/Y$与$Y$是独立随机变量，证明
$$E\left(\frac XY\right)^k=\frac{E(X^k)}{E(Y^k)}.$$
\item 利用该结果及Basu定理说明，若$X_1,X_2,\ldots,X_n$独立同分布于$\operatorname{gamma}(\alpha,\beta)$，其中$\alpha$已知，则
$$E\left(X_{(i)}\,\middle|\,\sum_iX_i\right)=\sum_iX_i\frac{E(X_{(i)})}{E(\sum_iX_i)}.$$
\end{enumerate}
\end{enumerate}
\end{exercise}
\begin{exercise}\label{exercise:6.32}
证明似然原则推论.也就是说，假设形式充分性原则和条件性原则都成立，证明：若$E=(\mathcal X,\theta,\{f(\boldsymbol x\mid\theta)\})$是一个试验，则$\operatorname{Ev}(E,\boldsymbol x)$应当仅通过$L(\theta\mid\boldsymbol x)$依赖于$E$和$\boldsymbol x$.
\end{exercise}
\begin{exercise}\label{exercise:6.33}\index{SI0093@Birnbaum定理}\index{SI0626@Birnbaum定理}
补全\autoref{thm:6.3.6}（Birnbaum定理）证明中的细节.
\begin{enumerate}[label=(\alph*)]
\item 定义$g(t\mid\theta)=g((j,\boldsymbol x_j)\mid\theta)=f^*((j,\boldsymbol x_j)\mid\theta)$及
$$h(j,\boldsymbol x_j)=\begin{cases}C,&\text{若 }(j,\boldsymbol x_j)=(2,\boldsymbol x_2^*),\\1,&\text{其他情形}.\end{cases}$$
通过验证对所有$(j,\boldsymbol x_j)$均有下式，说明$T(j,\boldsymbol x_j)$是$E^*$试验中的充分统计量：\footnote{译注：原文连等式中间一项在$(j,\boldsymbol x_j)=(2,\boldsymbol x_2^*)$处一般不成立；此处应为$g((1,\boldsymbol x_1^*)\mid\theta)C$.由证明中的比例关系，它仍等于最右端的$f^*((2,\boldsymbol x_2^*)\mid\theta)$，因此所需因子分解成立.}
\begin{align*}
g(T(j,\boldsymbol x_j)\mid\theta)h(j,\boldsymbol x_j)
&=g((j,\boldsymbol x_j)\mid\theta)(1)\\
&=f^*((j,\boldsymbol x_j)\mid\theta).
\end{align*}
\item 由于$T$充分，说明形式充分性原则蕴含\eqref{eq:6.3.4}.条件性原则也蕴含\eqref{eq:6.3.5}，由此推出形式似然原则.
\item 为证明逆命题，先取一个试验为$E^*$试验，另一个为$E_j$，推出$\operatorname{Ev}(E^*,(j,\boldsymbol x_j))=\operatorname{Ev}(E_j,\boldsymbol x_j)$，即条件性原则.然后，若$T(\boldsymbol X)$充分且$T(\boldsymbol x)=T(\boldsymbol y)$，说明似然成比例，再利用形式似然原则推出$\operatorname{Ev}(E,\boldsymbol x)=\operatorname{Ev}(E,\boldsymbol y)$，即形式充分性原则.
\end{enumerate}
\end{exercise}
\sourcepageanchor{279}
\begin{exercise}\label{exercise:6.34}\index{SI0629@样本量}
考虑\autoref{exercise:6.12}中的模型.说明形式似然原则意味着：关于$\theta$的任何结论，都不应依赖于样本量$n$是随机选取的这一事实.也就是说，\autoref{exercise:6.12}中样本点$(n,x)$的似然，与固定样本量的$\operatorname{binomial}(n,\theta)$试验中样本点$x$的似然成比例.
\end{exercise}
\begin{exercise}\label{exercise:6.35}
一项有风险的试验性治疗最多用于三名病人.先对一名病人治疗；若成功，则对第二名病人治疗；若又成功，则对第三名病人治疗.将各病人的结果建模为相互独立的$\operatorname{Bernoulli}(p)$随机变量.确定此模型中的四个样本点，并说明，按照形式似然原则，关于$p$的推断不应依赖于样本量由数据决定这一事实.
\end{exercise}
\begin{exercise}\label{exercise:6.36}
如下题所示，使用最小充分统计量的一个优点是无偏估计量将有较小的方差.设$T_1$充分，$T_2$最小充分，$U$是$\theta$的无偏估计量，并定义$U_1=E(U\mid T_1)$和$U_2=E(U\mid T_2)$.
\begin{enumerate}[label=(\alph*)]
\item 说明$U_2=E(U_1\mid T_2)$.
\item 现在利用条件方差公式（\autoref{thm:4.4.7}）证明$\operatorname{Var}U_2\le\operatorname{Var}U_1$.
\end{enumerate}
（关于充分性与无偏性的关系，另见\hyperref[bib:source240]{Pena和Rohatgi 1994}.）
\end{exercise}
\begin{exercise}\label{exercise:6.37}\index{SI0725@问题}
\hyperref[bib:source182]{Joshi和Nabar（1989）}研究了所谓“尼罗河问题”（Problem of the Nile）中参数的线性估计量的性质，其中$(X,Y)$具有联合密度
$$f(x,y\mid\theta)=\exp\{-(\theta x+y/\theta)\},\qquad x>0,\quad y>0.$$
\begin{enumerate}[label=(\alph*)]
\item 对样本量为$n$的独立同分布样本，说明Fisher信息为$I(\theta)=2n/\theta^2$.
\item 对估计量
$$T=\sqrt{\sum Y_i/\sum X_i}\qquad\text{及}\qquad U=\sqrt{\sum X_i\sum Y_i},$$
说明
\begin{enumerate}[label=(\roman*)]
\item 仅$T$所含的信息为$[2n/(2n+1)]I(\theta)$.
\item $(T,U)$所含的信息为$I(\theta)$.
\item $(T,U)$联合充分，但不完备.
\end{enumerate}
\end{enumerate}
\end{exercise}
\begin{exercise}\label{exercise:6.38}
在\autoref{def:6.4.2}中，说明\hyperref[item:6.4.2iii]{(iii)}由\hyperref[item:6.4.2i]{(i)}和\hyperref[item:6.4.2ii]{(ii)}蕴含.
\end{exercise}
\begin{exercise}\label{exercise:6.39}\index{SI0401@测量}
测量等变性要求对两个等价的数据点作出相同的推断：$\boldsymbol x$是在一种尺度下表示的测量值，$\boldsymbol y$是完全相同的测量在另一尺度下的表示.形式不变性最终则给出同一测量尺度下两个不同数据点处的推断之间的关系.

假设一位实验者希望根据一次观测$X$——用摄氏度测量的水的沸点——估计$\theta$，即水的平均沸点.考虑到海拔和水中的杂质，他决定采用估计$T(x)=.5x+.5(100)$.如果测量尺度改成华氏度，实验者会采用$T^*(y)=.5y+.5(212)$，估计以华氏度表示的平均沸点.
\begin{enumerate}[label=(\alph*)]
\item 摄氏度与华氏度之间熟知的关系使我们用变换$\frac59(T^*(y)-32)$把华氏度转为摄氏度.说明这个方法具有测量等变性，即相同数据将得到相同答案.也就是说，$\frac59(T^*(y)-32)=T(x)$.
\item 形式不变性要求对所有$x$均有$T(x)=T^*(x)$.说明上面定义的估计量不满足这一要求.因此，它们在等变性原则的意义下不等变.
\end{enumerate}
\end{exercise}
\begin{exercise}\label{exercise:6.40}\index{SI0021@位置尺度族}
设$X_1,\ldots,X_n$为来自一个位置--尺度族的独立同分布观测.设$T_1(X_1,\ldots,X_n)$和$T_2(X_1,\ldots,X_n)$是两个统计量，均对所有$x_1,\ldots,x_n$和$b$以及任意$a>0$满足
$$T_i(ax_1+b,\ldots,ax_n+b)=aT_i(x_1,\ldots,x_n).$$
\sourcepageanchor{280}
\begin{enumerate}[label=(\alph*)]
\item 说明$T_1/T_2$是辅助统计量.
\item 设$R$是样本极差，$S$是样本标准差.验证$R$和$S$满足上述条件，从而$R/S$是辅助统计量.
\end{enumerate}
\end{exercise}
\begin{exercise}\label{exercise:6.41}\index{SI0402@位置}\index{SI0590@位置}\index{SI0784@位置}\index{SI1054@位置}
假设在\autoref{ex:6.4.6}的模型中，要作出的推断是对均值$\mu$的估计.设当观测到$\boldsymbol X=\boldsymbol x$时所用的估计为$T(\boldsymbol x)$.若观测到$g_a(\boldsymbol X)=\boldsymbol Y=\boldsymbol y$，则令$T^*(\boldsymbol y)$为$\mu+a$（每个$Y_i$的均值）的估计.如果用$T^*(\boldsymbol y)$估计$\mu+a$，就会用$T^*(\boldsymbol y)-a$估计$\mu$.
\begin{enumerate}[label=(\alph*)]
\item 说明测量等变性要求：对所有$\boldsymbol x=(x_1,\ldots,x_n)$和所有$a$，均有$T(\boldsymbol x)=T^*(\boldsymbol y)-a$.
\item 说明形式不变性要求$T(\boldsymbol x)=T^*(\boldsymbol x)$，从而等变性原则要求：对所有$(x_1,\ldots,x_n)$和所有$a$，均有
$$T(x_1,\ldots,x_n)+a=T(x_1+a,\ldots,x_n+a).$$
\item 若$X_1,\ldots,X_n$独立同分布于$f(x-\theta)$，说明只要$E_0X_1=0$，估计量
$$W(X_1,\ldots,X_n)=\bar X$$
就是估计$\theta$的等变估计量，并满足$E_\theta W=\theta$.
\end{enumerate}
\end{exercise}
\begin{exercise}\label{exercise:6.42}\index{SI0465@位置尺度}\index{SI0785@位置尺度}
假设$X_1,\ldots,X_n$是来自位置--尺度pdf $\frac1\sigma f((x-\theta)/\sigma)$的随机样本.我们想估计$\theta$，考虑两个变换群：
$$\mathcal G_1=\{g_{a,c}(\boldsymbol x):-\infty<a<\infty,\ c>0\},$$
其中$g_{a,c}(x_1,\ldots,x_n)=(cx_1+a,\ldots,cx_n+a)$，以及
$$\mathcal G_2=\{g_a(\boldsymbol x):-\infty<a<\infty\},$$
其中$g_a(x_1,\ldots,x_n)=(x_1+a,\ldots,x_n+a)$.
\begin{enumerate}[label=(\alph*)]
\item 说明形如
$$W(x_1,\ldots,x_n)=\bar x+k,$$
其中$k$为非零常数的估计量，关于群$\mathcal G_2$等变，但关于群$\mathcal G_1$不等变.
\item 对每一个群，在什么条件下等变估计量$W$满足$E_\theta W=\theta$，即它对$\theta$的估计是无偏的？
\end{enumerate}
\end{exercise}
\begin{exercise}\label{exercise:6.43}\index{SI0467@尺度}
仍假设$X_1,\ldots,X_n$为来自位置--尺度pdf $\frac1\sigma f((x-\theta)/\sigma)$的随机样本，不过现在我们感兴趣的是估计$\sigma^2$.可以考虑三个变换群：
$$\mathcal G_1=\{g_{a,c}(\boldsymbol x):-\infty<a<\infty,\ c>0\},$$
其中$g_{a,c}(x_1,\ldots,x_n)=(cx_1+a,\ldots,cx_n+a)$；
$$\mathcal G_2=\{g_a(\boldsymbol x):-\infty<a<\infty\},$$
其中$g_a(x_1,\ldots,x_n)=(x_1+a,\ldots,x_n+a)$；以及
$$\mathcal G_3=\{g_c(\boldsymbol x):c>0\},$$
其中$g_c(x_1,\ldots,x_n)=(cx_1,\ldots,cx_n)$.
\begin{enumerate}[label=(\alph*)]
\item 说明形如$kS^2$的$\sigma^2$估计量（$k$为正常数，$S^2$为样本方差）关于$\mathcal G_2$不变，关于另外两个群等变.
\item 说明形如（\hyperref[bib:source48]{Brewster和Zidek，1974）}
$$W(X_1,\ldots,X_n)=\phi\left(\frac{\bar X}S\right)S^2,$$
其中$\phi(x)$是一个函数的、更大一类$\sigma^2$估计量，关于$\mathcal G_3$等变，但关于$\mathcal G_1$和$\mathcal G_2$均不等变，除非$\phi(x)$为常数.

对这种形式估计量的考虑，使\hyperref[bib:source302]{Stein（1964）}以及\hyperref[bib:source48]{Brewster和Zidek（1974）}找到了改进的方差估计量（见\hyperref[bib:source200]{Lehmann和Casella 1998}，第3.3节）.
\end{enumerate}
\end{exercise}



\sourcepageanchor{281}
\subsection{杂录}\label{sec:6.6}
\subsubsection{Basu定理的逆命题}\label{sec:6.6.1}\index{SI0077@逆命题}
一个有趣的统计事实是，Basu定理的逆命题并不成立.也就是说，若一个统计量$T(\boldsymbol X)$独立于每一个辅助统计量，它不一定完备且最小充分.\hyperref[bib:source197]{Lehmann（1981）}对这一问题作了很好的讨论.他指出，部分原因在于辅助性是\emph{整个分布}的性质，而完备性只涉及期望.据此，可以修改辅助性的定义.
\begin{definition}\label{def:6.6.1}
如果$E_\theta V(\boldsymbol X)$与$\theta$无关，则统计量$V(\boldsymbol X)$称为\textbf{一阶辅助统计量}（first-order ancillary statistic）.
\end{definition}
Lehmann证明了如下定理，它在某种意义上是Basu定理的逆命题.
\begin{theorem}\label{thm:6.6.2}\index{SI0111@辅助性}
设$T$为满足$\operatorname{Var}T<\infty$的统计量.$T$完备当且仅当每一个有界的一阶辅助统计量$V$与$T$的每一个有界实值函数均不相关（对所有$\theta$）.\footnote{译注：按本书的完备性定义，单独的完备性并不足以推出这里的结论，原述似缺少充分性等限定. 例如，取单个观测$X\sim\operatorname{uniform}(\theta,\theta+1)$，$0<\theta<1$；$T=\lfloor X\rfloor$服从完备的Bernoulli$(\theta)$族，而$V=X-\lfloor X\rfloor\sim\operatorname{uniform}(0,1)$为有界辅助统计量，但$\operatorname{Cov}_\theta(T,V)=-\theta(1-\theta)/2\ne0$. 正文保留原述；此反例为译者独立核验，精确修订条件仍需对照所引原始论文.}
\end{theorem}
Lehmann还指出，通过修改完备性的概念，而非辅助性的概念，还能得到其他逆命题.
\subsubsection{关于辅助性的困惑}\label{sec:6.6.2}\index{SI0018@辅助统计量}
辅助性有许多不同的定义，各自具有不同性质.若一个定义就已经令人困惑，那么五六个定义简直让人无从着手.

\hyperref[bib:source56]{Buehler（1982）}指出，这个主题始于\hyperref[bib:source123]{Fisher（1925）}，后者“留下了一连串耐人寻味的概念，却没有给出定义，这正是他的一贯风格”.Buehler至少区分出三个不同定义，将其归功于\hyperref[bib:source13]{Basu（1959）}以及\hyperref[bib:source86]{Cox和Hinkley（1974）}，并给出了八条性质和二十五个例子.

这个困难的主题值得理解，因为它在统计推断中起着重要作用.例如，\hyperref[bib:source52]{Brown（1990）}说明辅助性如何影响回归分析.\hyperref[bib:source250]{Reid（1995）}综述了辅助性以及其他条件化形式；\hyperref[bib:source202]{Lehmann和Scholz（1992）}也给出了一篇很好的综述，可作为进入这一主题的起点.
\subsubsection{充分性的进一步讨论}\label{sec:6.6.3}\index{SI0623@充分性}
\paragraph{充分性与似然}
\autoref{thm:6.2.13}与似然原则之间有一个引人注目的相似之处：两者都涉及似然比$L(\theta\mid\boldsymbol x)/L(\theta\mid\boldsymbol y)$，前者用它描述最小充分统计量，后者用它描述似然原则.仔细结合二者可知，$T(\boldsymbol X)$最小充分当且仅当它是似然$L(\theta\mid\boldsymbol x)$的一一函数（把满足\eqref{eq:6.3.1}的两个样本点视为具有相同的似然函数）.参见\autoref{ex:6.3.3}和\autoref{exercise:6.9}.
\paragraph{充分性与必要性}
“既然有充分统计量，为什么没有必要统计量？”实际上是有的，见\hyperref[bib:source104]{Dynkin（1951）}.
\begin{definition}\label{def:6.6.3}\index{SI0719@必要统计量}
如果一个统计量是每一个充分统计量的函数，则称它为\textbf{必要统计量}（necessary statistic）.
\end{definition}
比较必要统计量和最小充分统计量的定义，便立即得到如下定理.
\sourcepageanchor{282}
\begin{theorem}\label{thm:6.6.4}\index{SI1008@必要}
一个统计量最小充分，当且仅当它既必要又充分.
\end{theorem}
\paragraph{最小充分性}
由\autoref{thm:6.2.13}可以得到一个有趣的发展（见\autoref{exercise:6.28}）.这个结果对于处理指数族之外的分布尤为有用.
\begin{theorem}[最小充分统计量]\label{thm:6.6.5}
设$\{f_0(\boldsymbol x),f_1(\boldsymbol x),\ldots,f_k(\boldsymbol x)\}$为具有相同支撑的密度族.
\begin{enumerate}[label=(\alph*)]
\item 统计量
$$T(\boldsymbol X)=\left(\frac{f_1(\boldsymbol X)}{f_0(\boldsymbol X)},\frac{f_2(\boldsymbol X)}{f_0(\boldsymbol X)},\ldots,\frac{f_k(\boldsymbol X)}{f_0(\boldsymbol X)}\right)$$
对于该有限密度族最小充分.
\item 若$\mathcal F$是具有相同支撑的密度族，并且
\begin{enumerate}[label=(\roman*)]
\item $f_i(\boldsymbol x)\in\mathcal F$，$i=0,\ldots,k$；
\item $T(\boldsymbol x)$对于$\mathcal F$充分，
\end{enumerate}
则$T(\boldsymbol x)$对于$\mathcal F$最小充分.
\end{enumerate}
\end{theorem}
虽然可以利用这个定理确立$\bar X$对于$n(\theta,1)$族的最小充分性，但它真正有用的地方在这些简单情形之外.例如，可利用它确立次序统计量对于logistic族和双指数族的最小充分性（见\autoref{exercise:6.26}），甚至可以将其推广到非参数情形（见\autoref{exercise:6.29}）.

关于最小充分性与完备性的进一步讨论，参见\hyperref[bib:source200]{Lehmann和Casella（1998）}第1.6节.





\section{点估计}\label{sec:7}
\chapterdoi{7}

\sourcepageanchor{283}
\begin{flushright}
\begin{minipage}{.9\textwidth}
“什么！你已经解决了？”\\
“唔，这么说就过头了.我只是发现了一个富有启发性的事实，仅此而已.”
\begin{flushright}华生医生与夏洛克·福尔摩斯\\《四签名》\end{flushright}
\end{minipage}
\end{flushright}



\subsection{引言}\label{sec:7.1}
本章分为两部分.第一部分讨论构造估计量的方法，第二部分讨论如何评价这些（以及其他）估计量.一般而言，这两项活动相互交织.评价估计量的方法往往会启发我们提出新的估计量.不过，暂时我们还是把构造估计量与评价估计量区分开来.

点估计（point estimation）的基本理由很简单.如果从pdf或pmf为$f(x\mid\theta)$的总体抽样，知道$\theta$就意味着知道整个总体.因此，自然会希望找到一个方法，构造点$\theta$的良好估计量，即良好的点估计量.参数$\theta$也常常具有有意义的实际解释（例如总体均值），因而获得$\theta$的良好点估计本身就是我们直接关心的事情.还可能感兴趣的是$\theta$的某个函数，比如$\tau(\theta)$.本章介绍的方法也可用于构造$\tau(\theta)$的估计量.

下面对点估计量的定义看起来也许过于含糊.不过，在这个阶段，我们希望谨慎一些，不把任何候选者排除在考虑范围之外.
\begin{definition}\label{def:7.1.1}\index{SI0763@点估计量}
点估计量（point estimator）是样本的任意函数$W(X_1,\ldots,X_n)$.也就是说，任一统计量都是点估计量.
\end{definition}
注意，这个定义完全没有提到估计量与它要估计的参数之间应有什么对应关系.虽然可以认为定义中应包括这样的说明，但这样会限制可用估计量的集合.定义中也没有提及统计量$W(X_1,\ldots,X_n)$的值域.原则上统计量的值域应与参数的值域一致，但我们将看到情况并非总是如此.

有一点必须区分清楚，即估计值（estimate）与估计量（estimator）的差别.\index{SI0762@点估计值}\semanticindexanchor{762:283}估计量是样本的函数，而估计值是在实际抽取一个样本之后获得的估计量的实现值（即一个数）.在符号上，抽取样本时，估计量是随机变量$X_1,\ldots,X_n$的函数，而估计值是实现值$x_1,\ldots,x_n$的函数.

很多时候，某个参数的点估计量会有显而易见或自然的候选者.例如，样本均值就是总体均值点估计量的自然候选者.不过，一旦离开这种简单情形，直觉不仅可能帮不上忙，甚至可能把我们引入歧途.因此，掌握一些至少能给出合理候选者的技巧很有用.但请注意，这些技巧本身不附带任何保证.它们所产生的点估计量仍须经过评价，才能确立其价值.



\sourcepageanchor{284}
\subsection{构造估计量的方法}\label{sec:7.2}
有时，决定如何估计一个参数很容易，仅靠直觉就常常能得到很好的估计量.例如，用参数的样本对应量来估计它，通常是合理的.特别地，样本均值是总体均值的良好估计.在实践中常见的、更复杂的模型里，我们需要更加系统的参数估计方法.本节详细介绍四种构造估计量的方法.
\subsubsection{矩估计法}\label{sec:7.2.1}\index{SI0806@矩估计法}
矩估计法（method of moments）也许是构造点估计量最古老的方法，至少可追溯到19世纪末的Karl Pearson.它的优点是使用十分简单，而且几乎总能给出某种估计.遗憾的是，在许多情形下，这种方法给出的估计量还有改进余地.不过，当其他方法难以运用时，它是一个良好的起点.

设$X_1,\ldots,X_n$是来自pdf或pmf为$f(x\mid\theta_1,\ldots,\theta_k)$的总体的样本.矩估计量通过令前$k$个样本矩等于对应的$k$个总体矩，并求解所得联立方程组来构造.更精确地，定义
\begin{equation}\label{eq:7.2.1}
\begin{aligned}
m_1&=\frac1n\sum_{i=1}^nX_i^1,&\quad\mu_1'&=EX^1,\\
m_2&=\frac1n\sum_{i=1}^nX_i^2,&\mu_2'&=EX^2,\\
&\ \vdots&&\ \vdots\\
m_k&=\frac1n\sum_{i=1}^nX_i^k,&\mu_k'&=EX^k.
\end{aligned}
\end{equation}
总体矩$\mu_j'$通常是$\theta_1,\ldots,\theta_k$的函数，记作$\mu_j'(\theta_1,\ldots,\theta_k)$.通过求解下面的方程组，把$(\theta_1,\ldots,\theta_k)$表示成$(m_1,\ldots,m_k)$的函数，就得到$(\theta_1,\ldots,\theta_k)$的矩估计量$(\widetilde\theta_1,\ldots,\widetilde\theta_k)$：
\begin{equation}\label{eq:7.2.2}
\begin{aligned}
m_1&=\mu_1'(\theta_1,\ldots,\theta_k),\\
m_2&=\mu_2'(\theta_1,\ldots,\theta_k),\\
&\ \vdots\\
m_k&=\mu_k'(\theta_1,\ldots,\theta_k).
\end{aligned}
\end{equation}
\begin{example}[正态分布的矩估计]\label{ex:7.2.1}\index{SI0809@正态}
设$X_1,\ldots,X_n$独立同分布于$n(\theta,\sigma^2)$.用前面的记号，$\theta_1=\theta$、$\theta_2=\sigma^2$.有$m_1=\bar X$、$m_2=(1/n)\sum X_i^2$、$\mu_1'=\theta$、$\mu_2'=\theta^2+\sigma^2$，因而须求解
$$\bar X=\theta,\qquad\frac1n\sum X_i^2=\theta^2+\sigma^2.$$
解出$\theta$和$\sigma^2$，得到矩估计量
$$\widetilde\theta=\bar X\qquad\text{及}\qquad\widetilde{\sigma^2}=\frac1n\sum X_i^2-\bar X^2=\frac1n\sum(X_i-\bar X)^2.$$
\end{example}
在这个简单例子中，矩估计法的解与直觉一致，也许使我们对二者都更有信心.不过，当没有显而易见的估计量可用时，这种方法更有帮助.
\sourcepageanchor{285}
\begin{example}[二项分布的矩估计]\label{ex:7.2.2}\index{SI0807@二项}
设$X_1,\ldots,X_n$独立同分布于$\operatorname{binomial}(k,p)$，即
$$P(X_i=x\mid k,p)=\binom kx p^x(1-p)^{k-x},\qquad x=0,1,\ldots,k.$$
这里假设$k$和$p$都未知，希望获得两个参数的点估计量.（二项模型的这种稍不寻常的应用，曾用于估计那些已知有大量未报案事件的犯罪发生率.对这样的犯罪，真实报案率$p$和总发生次数$k$都未知.）

令前两个样本矩等于相应总体矩，得到方程组
\begin{align*}
\bar X&=kp,\\
\frac1n\sum X_i^2&=kp(1-p)+k^2p^2,
\end{align*}
现在须由此解出$k$和$p$.稍作代数运算，得到矩估计量
$$\widetilde k=\frac{\bar X^2}{\bar X-(1/n)\sum(X_i-\bar X)^2},$$
以及
$$\widetilde p=\frac{\bar X}{\widetilde k}.$$
诚然，这些并非总体参数的最佳估计量.特别地，它们可能给出$k$和$p$的负估计值，而这些参数当然必须为正.（这里估计量的值域与所估计参数的值域不一致.）不过，为矩估计法说句公道话，负估计值只会在样本均值小于样本方差时出现，而这意味着数据的变异程度很大.在这个例子中，矩估计法至少给出了$k$和$p$点估计量的一组候选者.虽然直觉也许能为$p$的估计量提供一个候选者，但想出$k$的估计量困难得多.
\end{example}
矩估计法在获得统计量分布的近似方面很有用.这种有时称为“矩匹配”（moment matching）的技巧，通过匹配分布的矩来给出近似.理论上，任何统计量分布的矩都可与任一分布的矩匹配；但在实践中，最好采用形状相似的分布.下面的例子展示了这种技巧最著名的应用之一，即\hyperref[bib:source272]{Satterthwaite（1946）}近似.它至今仍在使用（见\autoref{exercise:8.42}）.
\begin{example}[Satterthwaite近似]\label{ex:7.2.3}\index{SI0808@卡方}\index{SI0952@Satterthwaite近似}
若$Y_i$（$i=1,\ldots,k$）为独立的$\chi_{r_i}^2$随机变量，我们已经看到（\autoref{lem:5.3.2}），$\sum Y_i$也服从卡方分布，自由度为$\sum r_i$.遗憾的是，$\sum a_iY_i$（$a_i$为已知常数）的分布一般相当难以求出.不过，假定某个自由度$\nu$的$\chi_\nu^2$能提供良好近似，看来是合理的.

这几乎就是Satterthwaite的问题.他关心的是近似一个$t$统计量的分母，而$\sum a_iY_i$表示该统计量分母的平方.因此，给定$a_1,\ldots,a_k$，他希望找到一个$\nu$，使
$$\sum_{i=1}^ka_iY_i\sim\frac{\chi_\nu^2}\nu\qquad\text{（近似地）}.$$
因为$E(\chi_\nu^2/\nu)=1$，为了匹配一阶矩，需要
$$E\left(\sum_{i=1}^ka_iY_i\right)=\sum_{i=1}^ka_iEY_i=\sum_{i=1}^ka_ir_i=1,$$
\sourcepageanchor{286}
这对$a_i$给出了一个约束，但没有提供如何估计$\nu$的信息.为此必须匹配二阶矩，需要
$$E\left(\sum_{i=1}^ka_iY_i\right)^2=E\left(\frac{\chi_\nu^2}\nu\right)^2=\frac2\nu+1.$$
应用矩估计法，去掉第一个期望算子并解出$\nu$，得到
$$\widehat\nu=\frac2{(\sum_{i=1}^ka_iY_i)^2-1}.$$
因此，直接应用矩估计法确实给出了$\nu$的估计量，但它可能取负值.我们可以想象Satterthwaite会对这种可能性感到震惊，因为这并非他提出的估计量.他作了更多努力，按下面的方式对矩估计法进行了专门调整.写成
\begin{align*}
E\left(\sum a_iY_i\right)^2
&=\operatorname{Var}\left(\sum a_iY_i\right)+\left(E\sum a_iY_i\right)^2\\
&=\left(E\sum a_iY_i\right)^2\left[\frac{\operatorname{Var}(\sum a_iY_i)}{(E\sum a_iY_i)^2}+1\right]\\
&=\left[\frac{\operatorname{Var}(\sum a_iY_i)}{(E\sum a_iY_i)^2}+1\right].\qquad\left(E\sum a_iY_i=1\right)
\end{align*}
现在令二阶矩相等，得到
$$\nu=\frac{2(E\sum a_iY_i)^2}{\operatorname{Var}(\sum a_iY_i)}.$$
最后，利用$Y_1,\ldots,Y_k$是独立卡方随机变量这一事实，写成
\begin{align*}
\operatorname{Var}\left(\sum a_iY_i\right)&=\sum a_i^2\operatorname{Var}Y_i\\
&=2\sum\frac{a_i^2(EY_i)^2}{r_i}.\qquad\left(\operatorname{Var}Y_i=\frac{2(EY_i)^2}{r_i}\right)
\end{align*}
代入这个方差表达式，并去掉期望算子，得到Satterthwaite估计量
$$\widehat\nu=\frac{(\sum a_iY_i)^2}{\sum\frac{a_i^2}{r_i}Y_i^2}.$$
这一近似很好，至今仍被广泛使用.注意，Satterthwaite成功得到了一个总为正的估计量，从而消除了直接矩估计量所存在的明显问题.
\end{example}
\subsubsection{极大似然估计量}\label{sec:7.2.2}
极大似然法（method of maximum likelihood）是迄今最常用的构造估计量的技巧.回顾一下，若$X_1,\ldots,X_n$为来自pdf或pmf为$f(x\mid\theta_1,\ldots,\theta_k)$的总体的独立同分布样本，则似然函数定义为
\begin{equation}\label{eq:7.2.3}
L(\boldsymbol\theta\mid\boldsymbol x)=L(\theta_1,\ldots,\theta_k\mid x_1,\ldots,x_n)=\prod_{i=1}^nf(x_i\mid\theta_1,\ldots,\theta_k).
\end{equation}
\begin{definition}\label{def:7.2.4}\index{SI0667@最大似然估计量}
对于每个样本点$\boldsymbol x$，设$\widehat{\boldsymbol\theta}(\boldsymbol x)$是在固定$\boldsymbol x$时，使$L(\boldsymbol\theta\mid\boldsymbol x)$作为$\boldsymbol\theta$的函数达到最大值的一个参数值.基于样本$\boldsymbol X$的参数$\boldsymbol\theta$的极大似然估计量（maximum likelihood estimator，MLE）为$\widehat{\boldsymbol\theta}(\boldsymbol X)$.
\end{definition}
\sourcepageanchor{287}
注意，由其构造可知，MLE的值域与参数的值域一致.此外，在讨论估计量的实现值时，我们也用缩写MLE表示极大似然估计值.

直观上，MLE是估计量的一个合理选择.MLE是使观测样本最有可能出现的参数点.一般而言，MLE是良好的点估计量，具有后面要讨论的一些最优性质.

求函数最大值的一般问题有两个内在困难，因此极大似然估计也有这些困难.第一个问题是实际找到全局最大值，并验证找到的确实是全局最大值.很多时候，这只是一道简单的微分学练习；不过，即使对常见密度，有时也会出现困难.第二个问题是数值敏感性，即估计值对数据的微小变化有多敏感？（严格说，这是与任何极大化方法相关的数学问题，而非统计问题.不过，既然MLE是通过极大化求得的，这个问题就必须处理.）遗憾的是，有时稍有差别的样本会产生差别极大的MLE，使它的使用受到质疑.我们先考虑寻找MLE的问题.

若似然函数关于$\theta_i$可微，则MLE的可能候选者是满足
\begin{equation}\label{eq:7.2.4}
\frac\partial{\partial\theta_i}L(\boldsymbol\theta\mid\boldsymbol x)=0,\qquad i=1,\ldots,k
\end{equation}
的$(\theta_1,\ldots,\theta_k)$值.

注意，\eqref{eq:7.2.4}的解只是MLE的\emph{可能候选者}，因为一阶导数为零只是极大值的必要条件，而非充分条件.而且，一阶导数的零点只能定位函数定义域内部的极值点.若极值在边界取得，一阶导数未必为零.因此，必须另外检查边界上的极值.

一阶导数为零的点，可能是局部或全局极小值点、局部或全局极大值点，或者拐点.我们的任务是找到全局最大值.
\begin{example}[正态似然]\label{ex:7.2.5}\index{SI0679@正态}
设$X_1,\ldots,X_n$独立同分布于$n(\theta,1)$，用$L(\theta\mid\boldsymbol x)$表示似然函数.则
$$L(\theta\mid\boldsymbol x)=\prod_{i=1}^n\frac1{(2\pi)^{1/2}}e^{-(1/2)(x_i-\theta)^2}=\frac1{(2\pi)^{n/2}}e^{(-1/2)\sum_{i=1}^n(x_i-\theta)^2}.$$
方程$(d/d\theta)L(\theta\mid\boldsymbol x)=0$化为
$$\sum_{i=1}^n(x_i-\theta)=0,$$
其解为$\widehat\theta=\bar x$.因此，$\bar x$是MLE的候选者.为了验证$\bar x$确实是似然函数的全局最大值点，可以作如下论证.首先，$\widehat\theta=\bar x$是$\sum(x_i-\theta)=0$的唯一解，所以$\bar x$是一阶导数的唯一零点.其次，验证
$$\left.\frac{d^2}{d\theta^2}L(\theta\mid\boldsymbol x)\right|_{\theta=\bar x}<0.$$
因此，$\bar x$是内部唯一的极值点，且为极大值点.最后，为验证它是全局最大值点，必须检查边界$\pm\infty$.取极限很容易说明似然在$\pm\infty$为零.故$\widehat\theta=\bar x$是全局最大值点，从而$\bar X$是MLE.（实际上，可以稍微巧妙些，免去检查$\pm\infty$.既然已确定$\bar x$是唯一的内部极值点，且为极大值点，则$\pm\infty$不可能再有极大值.否则，内部必有一个极小值点，与唯一性矛盾.）
\end{example}
\sourcepageanchor{288}
另一种求MLE的方法，是不用微分而直接极大化.这种方法在代数上通常更简单，尤其当导数表达式繁杂时；不过，由于没有固定规则可循，有时更难实施.一种一般技巧是为似然函数找到一个全局上界，再确立达到该上界的点是唯一的.
\begin{example}[续\autoref{ex:7.2.5}]\label{ex:7.2.6}
回顾\autoref{thm:5.2.4}：对任意数$a$，
$$\sum_{i=1}^n(x_i-a)^2\ge\sum_{i=1}^n(x_i-\bar x)^2,$$
等号成立当且仅当$a=\bar x$.这意味着对任意$\theta$，
$$e^{-(1/2)\sum(x_i-\theta)^2}\le e^{-(1/2)\sum(x_i-\bar x)^2},$$
等号成立当且仅当$\theta=\bar x$.因此，$\bar X$是MLE.
\end{example}
在大多数情形中，尤其准备使用微分时，处理$L(\boldsymbol\theta\mid\boldsymbol x)$的自然对数$\log L(\boldsymbol\theta\mid\boldsymbol x)$（称为对数似然，log likelihood），\index{SI0621@对数似然}\semanticindexanchor{621:288}比直接处理$L(\boldsymbol\theta\mid\boldsymbol x)$更容易.这样做可行，是因为对数函数在$(0,\infty)$上严格递增，意味着$L(\boldsymbol\theta\mid\boldsymbol x)$和$\log L(\boldsymbol\theta\mid\boldsymbol x)$的极值点相同（见\autoref{exercise:7.3}）.
\begin{example}[Bernoulli分布的MLE]\label{ex:7.2.7}\index{SI0669@伯努利}
设$X_1,\ldots,X_n$独立同分布于$\operatorname{Bernoulli}(p)$.似然函数为
$$L(p\mid\boldsymbol x)=\prod_{i=1}^np^{x_i}(1-p)^{1-x_i}=p^y(1-p)^{n-y},$$
其中$y=\sum x_i$.虽然这个函数求导并不太难，但对数似然的求导容易得多：
$$\log L(p\mid\boldsymbol x)=y\log p+(n-y)\log(1-p).$$
若$0<y<n$，对$\log L(p\mid\boldsymbol x)$求导并令结果为零，得到解$\widehat p=y/n$.也容易验证，此时$y/n$是全局最大值点.若$y=0$或$y=n$，则
$$\log L(p\mid\boldsymbol x)=\begin{cases}n\log(1-p),&\text{若 }y=0,\\n\log p,&\text{若 }y=n.\end{cases}$$
两种情形下，$\log L(p\mid\boldsymbol x)$都是$p$的单调函数，也同样容易验证各自都有$\widehat p=y/n$.因此，我们说明了$\sum X_i/n$是$p$的MLE.
\end{example}
在这个推导中，我们假设参数空间为$0\le p\le1$.为了使$\widehat p=y/n$在$y=0$和$n$时也是MLE，参数空间必须包括$p=0$和$1$.与\autoref{ex:3.4.1}比较：那里为了满足指数族的要求，取的是$0<p<1$.

寻找极大似然估计量时，还须注意，极大化只在参数取值范围内进行.在某些情形中，这一点十分重要.
\begin{example}[参数范围受限的MLE]\label{ex:7.2.8}
设$X_1,\ldots,X_n$独立同分布于$n(\theta,1)$，并且已知$\theta$必须非负.在不限制$\theta$时，我们已看到其MLE为$\bar X$；不过，若$\bar X$为负，它便落在参数范围之外.

若$\bar x$为负，很容易验证（见\autoref{exercise:7.4}），似然函数$L(\theta\mid\boldsymbol x)$在$\theta\ge0$时关于$\theta$递减，在$\widehat\theta=0$处最大.因此，此时$\theta$的MLE为
$$\widehat\theta=\bar X\ \text{若 }\bar X\ge0,\qquad\text{而}\qquad\widehat\theta=0\ \text{若 }\bar X<0.$$
\end{example}
\sourcepageanchor{289}
若不能解析地极大化$L(\boldsymbol\theta\mid\boldsymbol x)$，也许可以借助计算机对它作数值极大化.实际上，这是MLE最重要的特点之一.只要能够写出模型（似然），就有希望通过数值方法将其极大化，从而找到参数的MLE.不过，采用这种方法时，总有一个问题：找到的是局部还是全局最大值？因此，在使用数值极大化之前，始终应尽可能充分地分析似然函数，确定局部极大值的数目和性质.
\begin{example}[试验次数未知的二项分布MLE]\label{ex:7.2.9}\index{SI0671@二项}
设$X_1,\ldots,X_n$是来自$\operatorname{binomial}(k,p)$总体的随机样本，其中$p$已知，$k$未知.例如，投掷一枚已知公平的硬币，观测到$x_i$次正面，却不知道硬币一共投了多少次.似然函数为
$$L(k\mid\boldsymbol x,p)=\prod_{i=1}^n\binom k{x_i}p^{x_i}(1-p)^{k-x_i}.$$
由于阶乘的存在，而且$k$必须是整数，通过微分来极大化$L(k\mid\boldsymbol x,p)$很困难.因此，尝试另一种方法.

当然，当$k<\max_ix_i$时，$L(k\mid\boldsymbol x,p)=0$.故MLE是一个整数$k\ge\max_ix_i$，满足
$$\frac{L(k\mid\boldsymbol x,p)}{L(k-1\mid\boldsymbol x,p)}\ge1\qquad\text{及}\qquad\frac{L(k+1\mid\boldsymbol x,p)}{L(k\mid\boldsymbol x,p)}<1.$$
我们将说明这样的$k$只有一个.似然之比为
$$\frac{L(k\mid\boldsymbol x,p)}{L(k-1\mid\boldsymbol x,p)}=\frac{(k(1-p))^n}{\prod_{i=1}^n(k-x_i)}.$$
因此，极大值的条件为
\begin{gather*}
(k(1-p))^n\ge\prod_{i=1}^n(k-x_i),\\
((k+1)(1-p))^n<\prod_{i=1}^n(k+1-x_i).
\end{gather*}
除以$k^n$，令$z=1/k$，我们希望在$0\le z\le1/\max_ix_i$内求解
$$(1-p)^n=\prod_{i=1}^n(1-x_iz).$$
显然，右端在这个范围内是$z$的严格递减函数，$z=0$时取值1，$z=1/\max_ix_i$时取值0.因此，方程有唯一解$z$，记为$\widehat z$.$1/\widehat z$未必是整数.不过，满足这些不等式、从而为MLE的整数$\widehat k$，就是不超过$1/\widehat z$的最大整数（见\autoref{exercise:7.5}）.因此，这一分析说明似然函数存在唯一的最大值点，并可通过数值求解一个$n$次多项式方程来找到它.\footnote{译注：原文的“唯一”须留意并列最大值：当$1/\widehat z=m$为整数时，$L(m)=L(m-1)$，二者均为MLE；文中的严格第二不等式选取了其中较大的一个.例如$n=1$、$x_1=1$、$p=1/2$时，$k=1,2$并列.此外，若所有$x_i=0$，这里的$1/\max_ix_i$未定义，须另外处理边界情形.}这种关于$k$的MLE的描述由\hyperref[bib:source116]{Feldman和Fox（1968）}给出.关于$k$的估计，另见\autoref{ex:7.2.13}.
\end{example}
极大似然估计量有一个有用性质，通常称为MLE的不变性（invariance property），不要与\autoref{sec:6}讨论的不变性混淆.设分布由参数$\theta$标定，但我们感兴趣的是$\theta$的某个函数，比如$\tau(\theta)$的估计量.非正式地说，MLE的不变性是指：若$\widehat\theta$是$\theta$的MLE，则$\tau(\widehat\theta)$是$\tau(\theta)$的MLE.例如，若$\theta$是正态分布的均值，则$\sin(\theta)$的MLE为$\sin(\bar X)$.这里介绍\hyperref[bib:source337]{Zehna（1966）}的处理；关于MLE不变性的其他处理方法，参见\hyperref[bib:source238]{Pal和Berry（1992）}.

当然，要把MLE不变性的概念形式化，还须克服一些技术问题，它们大多集中在我们想估计的函数$\tau(\theta)$上.若映射$\theta\mapsto\tau(\theta)$是一一的（即每个$\theta$对应唯一的$\tau(\theta)$，反之亦然），就没有问题.此时容易看出，无论把似然视为$\theta$的函数还是$\tau(\theta)$的函数来极大化，
\sourcepageanchor{290}
都得到相同答案.令$\eta=\tau(\theta)$，则反函数$\tau^{-1}(\eta)=\theta$有良好定义，而$\tau(\theta)$的似然函数写成$\eta$的函数为
$$L^*(\eta\mid\boldsymbol x)=\prod_{i=1}^nf(x_i\mid\tau^{-1}(\eta))=L(\tau^{-1}(\eta)\mid\boldsymbol x),$$
并且
$$\sup_\eta L^*(\eta\mid\boldsymbol x)=\sup_\eta L(\tau^{-1}(\eta)\mid\boldsymbol x)=\sup_\theta L(\theta\mid\boldsymbol x).$$
因此，$L^*(\eta\mid\boldsymbol x)$的最大值在$\eta=\tau(\theta)=\tau(\widehat\theta)$处取得，说明$\tau(\theta)$的MLE就是$\tau(\widehat\theta)$.

很多时候，这个简单形式的MLE不变性并不适用，因为我们感兴趣的许多函数不是一一的.例如，估计正态均值的平方$\theta^2$时，映射$\theta\mapsto\theta^2$并非一一映射.因此，需要一个更一般的定理，实际上还需要对$\tau(\theta)$的似然函数给出更一般的定义.

若$\tau(\theta)$不是一一的，则给定$\eta$，可能有不止一个$\theta$满足$\tau(\theta)=\eta$.此时，对$\eta$的极大化与对$\theta$的极大化之间的对应可能失效.例如，若$\widehat\theta$为$\theta$的MLE，可能还有另一个$\theta$值，比如$\theta_0$，满足$\tau(\widehat\theta)=\tau(\theta_0)$.我们需要避免这样的困难.

为此，为$\tau(\theta)$定义诱导似然函数（induced likelihood function）$L^*$：
\begin{equation}\label{eq:7.2.5}\index{SI0620@诱导}
L^*(\eta\mid\boldsymbol x)=\sup_{\{\theta:\tau(\theta)=\eta\}}L(\theta\mid\boldsymbol x).
\end{equation}
使$L^*(\eta\mid\boldsymbol x)$最大的值$\widehat\eta$，称为$\eta=\tau(\theta)$的MLE.由\eqref{eq:7.2.5}可见，$L^*$与$L$的最大值相同.
\begin{theorem}[MLE的不变性]\label{thm:7.2.10}\index{SI0677@不变性}
若$\widehat\theta$是$\theta$的MLE，则对任意函数$\tau(\theta)$，$\tau(\theta)$的MLE为$\tau(\widehat\theta)$.
\end{theorem}
\begin{proof}
令$\widehat\eta$表示使$L^*(\eta\mid\boldsymbol x)$最大的值.须证$L^*(\widehat\eta\mid\boldsymbol x)=L^*[\tau(\widehat\theta)\mid\boldsymbol x]$.如前所述，$L$与$L^*$的最大值相同，因而
\begin{align*}
L^*(\widehat\eta\mid\boldsymbol x)&=\sup_\eta\sup_{\{\theta:\tau(\theta)=\eta\}}L(\theta\mid\boldsymbol x)&&\text{（$L^*$的定义）}\\
&=\sup_\theta L(\theta\mid\boldsymbol x)\\
&=L(\widehat\theta\mid\boldsymbol x),&&\text{（$\widehat\theta$的定义）}
\end{align*}
第二个等号成立，是因为逐次极大化等于无约束地对$\theta$极大化，而后者在$\widehat\theta$处达到.此外，
\begin{align*}
L(\widehat\theta\mid\boldsymbol x)&=\sup_{\{\theta:\tau(\theta)=\tau(\widehat\theta)\}}L(\theta\mid\boldsymbol x)&&\text{（$\widehat\theta$是MLE）}\\
&=L^*[\tau(\widehat\theta)\mid\boldsymbol x].&&\text{（$L^*$的定义）}
\end{align*}
因此，这一连串等式表明$L^*(\widehat\eta\mid\boldsymbol x)=L^*(\tau(\widehat\theta)\mid\boldsymbol x)$，从而$\tau(\widehat\theta)$是$\tau(\theta)$的MLE.
\end{proof}
利用此定理，现在知道正态均值平方$\theta^2$的MLE为$\bar X^2$.也可将\autoref{thm:7.2.10}应用于更复杂的函数；例如，当$p$为二项分布概率时，$\sqrt{p(1-p)}$的MLE为$\sqrt{\widehat p(1-\widehat p)}$.

在结束极大似然估计量的构造问题之前，还有几点需要提及.
\sourcepageanchor{291}
MLE的不变性在多元情形也成立.\autoref{thm:7.2.10}的证明没有任何地方排除$\theta$为向量的可能.若$(\theta_1,\ldots,\theta_k)$的MLE为$(\widehat\theta_1,\ldots,\widehat\theta_k)$，而$\tau(\theta_1,\ldots,\theta_k)$是参数的任意函数，则$\tau(\theta_1,\ldots,\theta_k)$的MLE为$\tau(\widehat\theta_1,\ldots,\widehat\theta_k)$.

若$\boldsymbol\theta=(\theta_1,\ldots,\theta_k)$是多维的，寻找MLE就是求一个多元函数的最大值.若似然函数可微，令一阶偏导数为零便给出内部极值点的必要条件.不过，在多维情形中，用二阶导数条件检查极大值很繁琐，可以先尝试其他方法.先介绍一种通常更简单的技巧，即逐次极大化.
\begin{example}[$\mu$与$\sigma$未知时的正态MLE]\label{ex:7.2.11}\index{SI0679@正态}
设$X_1,\ldots,X_n$独立同分布于$n(\theta,\sigma^2)$，其中$\theta$与$\sigma^2$均未知.则
$$L(\theta,\sigma^2\mid\boldsymbol x)=\frac1{(2\pi\sigma^2)^{n/2}}e^{-(1/2)\sum_{i=1}^n(x_i-\theta)^2/\sigma^2},$$
以及
$$\log L(\theta,\sigma^2\mid\boldsymbol x)=-\frac n2\log2\pi-\frac n2\log\sigma^2-\frac12\sum_{i=1}^n(x_i-\theta)^2/\sigma^2.$$
关于$\theta$和$\sigma^2$的偏导数为
$$\frac\partial{\partial\theta}\log L(\theta,\sigma^2\mid\boldsymbol x)=\frac1{\sigma^2}\sum_{i=1}^n(x_i-\theta),$$
以及
$$\frac\partial{\partial\sigma^2}\log L(\theta,\sigma^2\mid\boldsymbol x)=-\frac n{2\sigma^2}+\frac1{2\sigma^4}\sum_{i=1}^n(x_i-\theta)^2.$$
令这些偏导数为零并求解，得$\widehat\theta=\bar x$、$\widehat{\sigma^2}=n^{-1}\sum_{i=1}^n(x_i-\bar x)^2$.为验证此解确为全局最大值点，首先回顾：若$\theta\ne\bar x$，则$\sum(x_i-\theta)^2>\sum(x_i-\bar x)^2$.因此，对任意$\sigma^2$，
\begin{equation}\label{eq:7.2.6}
\frac1{(2\pi\sigma^2)^{n/2}}e^{-(1/2)\sum_{i=1}^n(x_i-\bar x)^2/\sigma^2}
\ge\frac1{(2\pi\sigma^2)^{n/2}}e^{-(1/2)\sum_{i=1}^n(x_i-\theta)^2/\sigma^2}.
\end{equation}
所以，验证已经找到极大似然估计量的问题，就化为一维问题：验证
$$ (\sigma^2)^{-n/2}\exp\left(-\frac12\sum(x_i-\bar x)^2/\sigma^2\right) $$
在$\sigma^2=n^{-1}\sum(x_i-\bar x)^2$处达到全局最大值.用一元微积分很容易完成验证；事实上，估计量$(\bar X,n^{-1}\sum(X_i-\bar X)^2)$就是MLE.\footnote{译注：这里及下一例的内部极大值计算要求$\sum(x_i-\bar x)^2>0$.若全部观测相同（特别是$n=1$），在通常参数空间$\sigma^2>0$上似然随$\sigma^2\downarrow0$无界，最大值并不取得；不能把零方差直接当作该参数空间内的MLE.}

注意，\eqref{eq:7.2.6}不等式左端称为$\sigma^2$的剖面似然（profile likelihood）.参见\autoref{sec:7.5.5}.
\end{example}
现在考虑用二元微积分求解同一个问题.
\begin{example}[续\autoref{ex:7.2.11}]\label{ex:7.2.12}\index{SI0673@微积分}
要用二元微积分验证函数$H(\theta_1,\theta_2)$在$(\widehat\theta_1,\widehat\theta_2)$处有局部极大值，须说明下列三个条件成立.
\begin{enumerate}[label=(\alph*)]
\item\phantomsection\label{itm:ex:7.2.12:a} 一阶偏导数为零：
$$\left.\frac\partial{\partial\theta_1}H(\theta_1,\theta_2)\right|_{\theta_1=\widehat\theta_1,\theta_2=\widehat\theta_2}=0\quad\text{且}\quad\left.\frac\partial{\partial\theta_2}H(\theta_1,\theta_2)\right|_{\theta_1=\widehat\theta_1,\theta_2=\widehat\theta_2}=0.$$
\item\phantomsection\label{itm:ex:7.2.12:b} 至少一个二阶偏导数为负：
$$\left.\frac{\partial^2}{\partial\theta_1^2}H(\theta_1,\theta_2)\right|_{\theta_1=\widehat\theta_1,\theta_2=\widehat\theta_2}<0\quad\text{或}\quad\left.\frac{\partial^2}{\partial\theta_2^2}H(\theta_1,\theta_2)\right|_{\theta_1=\widehat\theta_1,\theta_2=\widehat\theta_2}<0.$$
\sourcepageanchor{292}
\item\phantomsection\label{itm:ex:7.2.12:c} 二阶偏导数的Jacobian为正：\footnote{译注：这里原文把二阶偏导数矩阵的行列式称为Jacobian；通常术语是Hessian矩阵的行列式，它也可看作梯度映射的Jacobian行列式.原称谓保留.}
\begin{align*}
&\left.\begin{vmatrix}
\dfrac{\partial^2}{\partial\theta_1^2}H(\theta_1,\theta_2)&\dfrac{\partial^2}{\partial\theta_1\partial\theta_2}H(\theta_1,\theta_2)\\[4pt]
\dfrac{\partial^2}{\partial\theta_1\partial\theta_2}H(\theta_1,\theta_2)&\dfrac{\partial^2}{\partial\theta_2^2}H(\theta_1,\theta_2)
\end{vmatrix}\right|_{\theta_1=\widehat\theta_1,\theta_2=\widehat\theta_2}\\
&\quad=\left.\left[\frac{\partial^2H(\theta_1,\theta_2)}{\partial\theta_1^2}\frac{\partial^2H(\theta_1,\theta_2)}{\partial\theta_2^2}-\left(\frac{\partial^2H(\theta_1,\theta_2)}{\partial\theta_1\partial\theta_2}\right)^2\right]\right|_{\theta_1=\widehat\theta_1,\theta_2=\widehat\theta_2}>0.
\end{align*}
\end{enumerate}
对于正态对数似然，二阶偏导数为
\begin{align*}
\frac{\partial^2}{\partial\theta^2}\log L(\theta,\sigma^2\mid\boldsymbol x)&=-\frac n{\sigma^2},\\
\frac{\partial^2}{\partial(\sigma^2)^2}\log L(\theta,\sigma^2\mid\boldsymbol x)&=\frac n{2\sigma^4}-\frac1{\sigma^6}\sum_{i=1}^n(x_i-\theta)^2,\\
\frac{\partial^2}{\partial\theta\partial\sigma^2}\log L(\theta,\sigma^2\mid\boldsymbol x)&=-\frac1{\sigma^4}\sum_{i=1}^n(x_i-\theta).
\end{align*}
容易看出\hyperref[itm:ex:7.2.12:a]{(a)}和\hyperref[itm:ex:7.2.12:b]{(b)}成立，Jacobian为
\begin{align*}
&\left.\begin{vmatrix}
-\dfrac n{\sigma^2}&-\dfrac1{\sigma^4}\sum_{i=1}^n(x_i-\theta)\\[4pt]
-\dfrac1{\sigma^4}\sum_{i=1}^n(x_i-\theta)&\dfrac n{2\sigma^4}-\dfrac1{\sigma^6}\sum_{i=1}^n(x_i-\theta)^2
\end{vmatrix}\right|_{\theta=\bar x,\sigma^2=\widehat{\sigma^2}}\\
&=\left.\frac1{\sigma^6}\left[-\frac{n^2}2+\frac n{\sigma^2}\sum_{i=1}^n(x_i-\theta)^2-\frac1{\sigma^2}\left(\sum_{i=1}^n(x_i-\theta)\right)^2\right]\right|_{\theta=\bar x,\sigma^2=\widehat{\sigma^2}}\\
&=\frac1{\widehat\sigma^6}\left[-\frac{n^2}2+\frac{n^2}{\widehat\sigma^2}\widehat\sigma^2-\frac1{\widehat\sigma^2}\left(\sum_{i=1}^n(x_i-\bar x)\right)^2\right]\\
&=\frac1{\widehat\sigma^6}\frac{n^2}2>0.
\end{align*}
因此，微积分条件得到满足，确实找到了极大值.（当然，严格说来，我们只验证了$(\bar x,\widehat\sigma^2)$是内部极大值点.还须检查它是唯一的，而且无穷远处没有极大值.）即使这么简单的问题，计算量也已经相当可观，往后只会更糟.（想想有三个参数时需要做些什么.）因此，得到的教训是：虽然总要验证确实找到了极大值，但应寻找二阶导数条件以外的验证方法.
\end{example}
最后，前面提到，MLE通过极大化过程求得，因此容易受到该过程相关问题的影响，包括数值不稳定性.现在更详细地考察这个问题.

回顾似然函数是参数$\boldsymbol\theta$的函数，数据$\boldsymbol x$保持不变.不过，数据测量有误差，因此可以问，数据的微小变化会如何影响MLE？即我们根据$L(\boldsymbol\theta\mid\boldsymbol x)$计算$\widehat{\boldsymbol\theta}$，但也可询问：若对于很小的$\boldsymbol\varepsilon$，改为根据$L(\boldsymbol\theta\mid\boldsymbol x+\boldsymbol\varepsilon)$计算，会得到怎样的MLE？直观上，当$\boldsymbol\varepsilon$很小时，新的MLE，记为$\widehat{\boldsymbol\theta}_1$，应当接近$\widehat{\boldsymbol\theta}$.但事实并非总是如此.
\begin{example}[续\autoref{ex:7.2.2}]\label{ex:7.2.13}\index{SI0671@二项}\index{SI0676@不稳定性}
\hyperref[bib:source237]{Olkin等（1981）}说明，在二项抽样中，$k$和$p$的MLE可能极不稳定.他们用下面的例子说明这一点.观测到一个$\operatorname{binomial}(k,p)$试验的五个实现值，其中$k$和$p$都
\sourcepageanchor{293}
未知.第一组数据为$(16,18,22,25,27)$（这是未知次数的二项试验中观测到的成功次数）.对此数据集，$k$的MLE为$\widehat k=99$.若第二组数据为$(16,18,22,25,28)$，唯一的变化是把27换成28，则$k$的MLE为$\widehat k=190$，显示出极大的变异.
\end{example}
当似然函数在最大值附近十分平坦，或不存在有限的最大值点时，就会出现这种现象.若MLE可显式求出——我们后面的例子往往如此——通常不会有这个问题.不过，很多情形像上例一样，MLE不能显式求解，必须用数值方法寻找.面对这种问题时，多花一点时间考察解的稳定性，往往是明智的.
\subsubsection{Bayes估计量}\label{sec:7.2.3}
Bayes统计方法与我们一直采用的经典方法有根本区别.然而，Bayes方法的某些方面对其他统计方法也很有帮助.在介绍构造Bayes估计量的方法之前，先讨论Bayes统计的思路.

在经典方法中，参数$\theta$被视为未知但固定的量.从由$\theta$标定的总体中抽取随机样本$X_1,\ldots,X_n$，再根据样本观测值获得关于$\theta$取值的知识.在Bayes方法中，$\theta$则被视为一个其变动可由概率分布描述的量，这个分布称为先验分布（prior distribution）.\index{SI0824@先验分布}\semanticindexanchor{824:293}它是依据实验者信念形成的主观分布，在看到数据之前制定（因此称为“先验”分布）.随后，从由$\theta$标定的总体中抽取样本，用样本信息更新先验分布.更新后的先验称为后验分布（posterior distribution）.\index{SI0819@后验分布}\semanticindexanchor{819:293}更新使用了\autoref{sec:1}中介绍的Bayes法则，因此称为Bayes统计.

若以$\pi(\theta)$表示先验分布，以$f(\boldsymbol x\mid\theta)$表示抽样分布，则后验分布，即给定样本$\boldsymbol x$时$\theta$的条件分布，为
\begin{equation}\label{eq:7.2.7}
\pi(\theta\mid\boldsymbol x)=\frac{f(\boldsymbol x\mid\theta)\pi(\theta)}{m(\boldsymbol x)},\qquad\bigl(f(\boldsymbol x\mid\theta)\pi(\theta)=f(\boldsymbol x,\theta)\bigr),
\end{equation}
其中$m(\boldsymbol x)$是$\boldsymbol X$的边缘分布，即
\begin{equation}\label{eq:7.2.8}
m(\boldsymbol x)=\int f(\boldsymbol x\mid\theta)\pi(\theta)\,d\theta.
\end{equation}
注意，后验分布是以观测到样本为条件的条件分布.然后用它对$\theta$作出陈述，而$\theta$仍被视为随机量.例如，可以用后验分布的均值作为$\theta$的点估计值.

\emph{关于记号的说明：}讨论参数$\theta$上的分布时，我们将打破此前“大写字母表示随机变量，小写字母表示函数自变量”的记号约定.因此，会说随机量$\theta$具有分布$\pi(\theta)$.这更符合通常用法，不应造成混淆.
\begin{example}[二项分布的Bayes估计]\label{ex:7.2.14}\index{SI0079@伯努利}
设$X_1,\ldots,X_n$独立同分布于$\operatorname{Bernoulli}(p)$.则$Y=\sum X_i$服从$\operatorname{binomial}(n,p)$.假设$p$的先验分布为$\operatorname{beta}(\alpha,\beta)$.$Y$与$p$的联合分布为
\begin{align*}
f(y,p)&=\left[\binom ny p^y(1-p)^{n-y}\right]\left[\frac{\Gamma(\alpha+\beta)}{\Gamma(\alpha)\Gamma(\beta)}p^{\alpha-1}(1-p)^{\beta-1}\right]\\
&\hspace{2cm}\text{（条件分布$\times$边缘分布，即$f(y\mid p)\times\pi(p)$）}\\
&=\binom ny\frac{\Gamma(\alpha+\beta)}{\Gamma(\alpha)\Gamma(\beta)}p^{y+\alpha-1}(1-p)^{n-y+\beta-1}.
\end{align*}
\sourcepageanchor{294}
$Y$的边缘pdf为
\begin{equation}\label{eq:7.2.9}
f(y)=\int_0^1f(y,p)\,dp=\binom ny\frac{\Gamma(\alpha+\beta)}{\Gamma(\alpha)\Gamma(\beta)}\frac{\Gamma(y+\alpha)\Gamma(n-y+\beta)}{\Gamma(n+\alpha+\beta)},
\end{equation}
这种分布称为beta--二项分布（见\autoref{exercise:4.34}和\autoref{ex:4.4.6}）.后验分布，即给定$y$时$p$的分布，为
$$f(p\mid y)=\frac{f(y,p)}{f(y)}=\frac{\Gamma(n+\alpha+\beta)}{\Gamma(y+\alpha)\Gamma(n-y+\beta)}p^{y+\alpha-1}(1-p)^{n-y+\beta-1},$$
也就是$\operatorname{beta}(y+\alpha,n-y+\beta)$分布.（记住，这里$p$是变量，$y$被视为固定.）$p$的一个自然估计值是后验分布的均值，由此得到$p$的Bayes估计量
$$\widehat p_B=\frac{y+\alpha}{\alpha+\beta+n}.$$
\end{example}
考察一下$p$的Bayes估计值如何形成.先验分布的均值为$\alpha/(\alpha+\beta)$，在没有看到数据时，这会是我们对$p$的最佳估计.如果忽略先验信息，我们可能会用$p=y/n$作为$p$的估计值.$p$的Bayes估计值把这些信息全部结合起来.把$\widehat p_B$写成
$$\widehat p_B=\left(\frac n{\alpha+\beta+n}\right)\left(\frac yn\right)+\left(\frac{\alpha+\beta}{\alpha+\beta+n}\right)\left(\frac\alpha{\alpha+\beta}\right),$$
便能清楚看到这种结合方式.因此，$p_B$是先验均值与样本均值的线性组合，权重由$\alpha$、$\beta$和$n$决定.

估计二项分布参数时，并不一定要从beta族中选取先验分布.不过，选择beta族有一定好处，其中一个重要好处就是得到了估计量的闭式表达式.一般而言，对任意抽样分布，都有一个自然的先验分布族，称为共轭族（conjugate family）.
\begin{definition}\label{def:7.2.15}\index{SI0195@共轭族}
用$\mathcal F$表示由$\theta$标定的pdf或pmf $f(\boldsymbol x\mid\theta)$的类.若对所有$f\in\mathcal F$、$\Pi$中的所有先验，以及所有$\boldsymbol x\in\mathcal X$，后验分布都属于类$\Pi$，则先验分布类$\Pi$称为$\mathcal F$的共轭族.
\end{definition}
beta族与二项族共轭.因此，若从beta先验开始，最终就得到beta后验.先验的更新表现为其参数的更新.数学上这很方便，因为通常会使计算很容易.不过，对某一具体问题，共轭族是否为合理选择，应由实验者判断.

以最后一个例子结束本小节.
\begin{example}[正态分布的Bayes估计量]\label{ex:7.2.16}\index{SI0080@正态}
设$X\sim n(\theta,\sigma^2)$，假设$\theta$的先验分布为$n(\mu,\tau^2)$.（这里假设$\sigma^2$、$\mu$和$\tau^2$都已知.）$\theta$的后验分布也为正态，均值和方差为
\begin{equation}\label{eq:7.2.10}
\begin{aligned}
E(\theta\mid x)&=\frac{\tau^2}{\tau^2+\sigma^2}x+\frac{\sigma^2}{\sigma^2+\tau^2}\mu,\\
\operatorname{Var}(\theta\mid x)&=\frac{\sigma^2\tau^2}{\sigma^2+\tau^2}.
\end{aligned}
\end{equation}
（细节见\autoref{exercise:7.22}.）注意，正态族是它自身的共轭族.仍取后验均值，得到$\theta$的Bayes估计量为$E(\theta\mid X)$.
\sourcepageanchor{295}
Bayes估计量再次成为先验均值与样本均值的线性组合.还要注意，当先验方差$\tau^2$趋于无穷时，Bayes估计量趋于样本均值.可以把这解释为：先验信息越模糊，Bayes估计量就越重视样本信息.另一方面，如果先验信息很好，使得$\sigma^2>\tau^2$，那么先验均值就获得更大的权重.
\end{example}
\subsubsection{EM算法}\label{sec:7.2.4}\index{SI0007@EM}\index{SI0013@缺失数据}\index{SI0395@EM算法（见算法：EM）}
我们最后考察的一种构造估计量的方法，在思路上与前面的方法有本质区别，它专门用于寻找MLE.\footnote{本小节的内容较为专门，也更深入.略去本小节不会影响正文的连贯性.}这里不是详细给出一个解出MLE的过程，而是指定一种保证收敛到MLE的算法.\footnote{译注：一般EM算法保证的是观测数据似然的非递减性；这本身不保证参数序列收敛，更不保证收敛到全局MLE.收敛及全局最优性需要附加条件，并可能依赖初值.本节原文的保证性表述保留；参见后面\autoref{thm:7.2.20}所实际陈述的单调性质.}它称为EM算法，即期望--极大化算法（Expectation-Maximization algorithm）.其思想是：把一次困难的似然极大化，换成一系列较容易的极大化，其极限便是原问题的答案.它特别适用于“缺失数据”问题，因为数据缺失本身有时就会使计算繁琐.不过，我们将看到，把“缺失数据”补上，反而常使计算更加顺畅.（还会看到，“缺失数据”有不同解释，例如见\autoref{exercise:7.30}.）

使用EM算法时，要考虑两个不同的似然问题.我们希望解决的是“不完全数据”问题，实际求解的则是“完全数据”问题.视情况而定，可以从任一个问题开始.
\begin{example}[多个Poisson率]\label{ex:7.2.17}
观测到$X_1,X_2,\ldots,X_n$和$Y_1,Y_2,\ldots,Y_n$，它们全部相互独立，其中$Y_i\sim\operatorname{Poisson}(\beta\tau_i)$，$X_i\sim\operatorname{Poisson}(\tau_i)$.例如，这可用于建模某种疾病的发生次数$Y_i$，其潜在发生率由总体效应$\beta$和一个附加因子$\tau_i$共同决定.$\tau_i$可以衡量地区$i$的人口密度，或该地区人群的健康状况.我们看不到$\tau_i$，但通过$X_i$获得关于它的信息.

因此，联合pmf为
\begin{equation}\label{eq:7.2.11}\index{SI0009@完整数据似然}
\begin{aligned}
&f((x_1,y_1),(x_2,y_2),\ldots,(x_n,y_n)\mid\beta,\tau_1,\tau_2,\ldots,\tau_n)\\
&\qquad=\prod_{i=1}^n\frac{e^{-\beta\tau_i}(\beta\tau_i)^{y_i}}{y_i!}\frac{e^{-\tau_i}\tau_i^{x_i}}{x_i!}.
\end{aligned}
\end{equation}
通过直接微分即可求得似然估计量（见\autoref{exercise:7.27}）：
\begin{equation}\label{eq:7.2.12}
\widehat\beta=\frac{\sum_{i=1}^ny_i}{\sum_{i=1}^nx_i},\qquad\widehat\tau_j=\frac{x_j+y_j}{\widehat\beta+1},\quad j=1,2,\ldots,n.
\end{equation}
基于pmf \eqref{eq:7.2.11}的似然称为完全数据似然（complete-data likelihood），而
$$((x_1,y_1),(x_2,y_2),\ldots,(x_n,y_n))$$
称为完全数据.数据缺失很常见，\index{SI0008@完整数据}\semanticindexanchor{8:295}它会使估计更困难.例如，假设$x_1$的值缺失.我们可以把$y_1$也丢弃，用大小为$n-1$的样本继续分析，但这忽略了$y_1$中的信息.使用这些信息可以改进估计.

从pmf \eqref{eq:7.2.11}出发，$x_1$缺失时样本的pmf自然定义为\footnote{译注：原文离散求和末尾多出$dx_1$，应删去该微分符号；此处按原式保留.原式中条件竖线前还多了一个右括号，译文仅配平函数括号.}
\begin{equation}\label{eq:7.2.13}\index{SI0012@不完整数据似然}
\sum_{x_1=0}^\infty f((x_1,y_1),(x_2,y_2),\ldots,(x_n,y_n)\mid\beta,\tau_1,\tau_2,\ldots,\tau_n)\,dx_1.
\end{equation}
基于\eqref{eq:7.2.13}的似然是不完全数据似然（incomplete-data likelihood）.\index{SI0011@不完整数据}\semanticindexanchor{11:295}这正是需要极大化的似然.
\end{example}
\sourcepageanchor{296}
一般而言，可以从完全数据问题走向不完全数据问题，也可以反过来.若$\boldsymbol Y=(Y_1,Y_2,\ldots,Y_n)$是不完全数据，$\boldsymbol X=(X_1,X_2,\ldots,X_m)$为增广数据（augmented data），使$(\boldsymbol Y,\boldsymbol X)$成为完全数据，则$\boldsymbol Y$的密度$g(\cdot\mid\theta)$与$(\boldsymbol Y,\boldsymbol X)$的密度$f(\cdot\mid\theta)$满足
\begin{equation}\label{eq:7.2.14}
g(\boldsymbol y\mid\theta)=\int f(\boldsymbol y,\boldsymbol x\mid\theta)\,d\boldsymbol x.
\end{equation}
把它们看成似然，则$L(\theta\mid\boldsymbol y)=g(\boldsymbol y\mid\theta)$是不完全数据似然，$L(\theta\mid\boldsymbol y,\boldsymbol x)=f(\boldsymbol y,\boldsymbol x\mid\theta)$是完全数据似然.当$L(\theta\mid\boldsymbol y)$难以处理时，完全数据似然有时会更容易处理.
\begin{example}[续\autoref{ex:7.2.17}]\label{ex:7.2.18}
在\eqref{eq:7.2.11}中对$x_1$求和，得到不完全数据似然：
\begin{equation}\label{eq:7.2.15}
\begin{aligned}
&L(\beta,\tau_1,\tau_2,\ldots,\tau_n\mid y_1,(x_2,y_2),\ldots,(x_n,y_n))\\
&\qquad=\left[\prod_{i=1}^n\frac{e^{-\beta\tau_i}(\beta\tau_i)^{y_i}}{y_i!}\right]\left[\prod_{i=2}^n\frac{e^{-\tau_i}\tau_i^{x_i}}{x_i!}\right],
\end{aligned}
\end{equation}
其中$(y_1,(x_2,y_2),\ldots,(x_n,y_n))$是不完全数据.这就是需要极大化的似然.微分后得到MLE方程
\begin{equation}\label{eq:7.2.16}
\begin{aligned}
\widehat\beta&=\frac{\sum_{i=1}^ny_i}{\sum_{i=1}^n\widehat\tau_i},\\
y_1&=\widehat\tau_1\widehat\beta,\\
x_j+y_j&=\widehat\tau_j(\widehat\beta+1),\qquad j=2,3,\ldots,n,
\end{aligned}
\end{equation}
现在用EM算法求解.
\end{example}
EM算法使我们只需处理$L(\theta\mid\boldsymbol y,\boldsymbol x)$以及给定$\boldsymbol y$与$\theta$时$\boldsymbol X$的条件pdf或pmf，就能极大化$L(\theta\mid\boldsymbol y)$.定义
\begin{equation}\label{eq:7.2.17}
\begin{aligned}
L(\theta\mid\boldsymbol y,\boldsymbol x)&=f(\boldsymbol y,\boldsymbol x\mid\theta),\qquad L(\theta\mid\boldsymbol y)=g(\boldsymbol y\mid\theta),\\
k(\boldsymbol x\mid\theta,\boldsymbol y)&=\frac{f(\boldsymbol y,\boldsymbol x\mid\theta)}{g(\boldsymbol y\mid\theta)}.
\end{aligned}
\end{equation}
重新整理\eqref{eq:7.2.17}的最后一个等式，得到恒等式
\begin{equation}\label{eq:7.2.18}
\log L(\theta\mid\boldsymbol y)=\log L(\theta\mid\boldsymbol y,\boldsymbol x)-\log k(\boldsymbol x\mid\theta,\boldsymbol y).
\end{equation}
由于$\boldsymbol x$是缺失数据，并未观测到，我们将\eqref{eq:7.2.18}右端换成它在$k(\boldsymbol x\mid\theta',\boldsymbol y)$下的期望，得到新的恒等式
\begin{equation}\label{eq:7.2.19}
\begin{aligned}
\log L(\theta\mid\boldsymbol y)
&=E[\log L(\theta\mid\boldsymbol y,\boldsymbol X)\mid\theta',\boldsymbol y]\\
&\quad-E[\log k(\boldsymbol X\mid\theta,\boldsymbol y)\mid\theta',\boldsymbol y].
\end{aligned}
\end{equation}
现在开始算法：从初值$\theta^{(0)}$出发，按下式生成序列$\theta^{(r)}$：
\begin{equation}\label{eq:7.2.20}
\theta^{(r+1)}=\text{使 }E[\log L(\theta\mid\boldsymbol y,\boldsymbol X)\mid\theta^{(r)},\boldsymbol y]\text{ 最大的值}.
\end{equation}
算法的“E步”计算期望对数似然，“M步”寻找其最大值.在探讨算法为何确实收敛到MLE之前，先回到我们的例子.
\sourcepageanchor{297}
\begin{example}[\autoref{ex:7.2.17}的结尾]\label{ex:7.2.19}
用$(\boldsymbol x,\boldsymbol y)=((x_1,y_1),(x_2,y_2),\ldots,(x_n,y_n))$表示完全数据，用$(\boldsymbol x_{(-1)},\boldsymbol y)=(y_1,(x_2,y_2),\ldots,(x_n,y_n))$表示不完全数据.完全数据对数似然的期望为\footnote{译注：原文\eqref{eq:7.2.21}最后一次分组中，$[\tau_ix_i\log\tau_i]$和$[\tau_1x_1\log\tau_1]$应分别为$[-\tau_i+x_i\log\tau_i]$和$[-\tau_1+x_1\log\tau_1]$；第二组括号的求和下标$2=1$应为$i=2$.这些原文误排保留，紧随的\eqref{eq:7.2.22}实际采用了正确的$-\tau_1+x_1\log\tau_1$项.}
\begin{equation}\label{eq:7.2.21}
\begin{aligned}
&E[\log L(\beta,\tau_1,\tau_2,\ldots,\tau_n\mid(\boldsymbol x,\boldsymbol y))\mid\boldsymbol\tau^{(r)},(\boldsymbol x_{(-1)},\boldsymbol y)]\\
&=\sum_{x_1=0}^\infty\log\left(\prod_{i=1}^n\frac{e^{-\beta\tau_i}(\beta\tau_i)^{y_i}}{y_i!}\frac{e^{-\tau_i}\tau_i^{x_i}}{x_i!}\right)\frac{e^{-\tau_1^{(r)}}(\tau_1^{(r)})^{x_1}}{x_1!}\\
&=\sum_{i=1}^n[-\beta\tau_i+y_i(\log\beta+\log\tau_i)-\log y_i!]\\
&\quad+\sum_{i=2}^n[-\tau_i+x_i\log\tau_i-\log x_i!]\\
&\quad+\sum_{x_1=0}^\infty[-\tau_1+x_1\log\tau_1-\log x_1!]\frac{e^{-\tau_1^{(r)}}(\tau_1^{(r)})^{x_1}}{x_1!}\\
&=\left(\begin{aligned}&\sum_{i=1}^n[-\beta\tau_i+y_i(\log\beta+\log\tau_i)]+\sum_{i=2}^n[\tau_ix_i\log\tau_i]\\
&+\sum_{x_1=0}^\infty[\tau_1x_1\log\tau_1]\frac{e^{-\tau_1^{(r)}}(\tau_1^{(r)})^{x_1}}{x_1!}\end{aligned}\right)\\
&\quad-\left(\sum_{i=1}^n\log y_i!+\sum_{2=1}^n\log x_i!+\sum_{x_1=0}^\infty\log x_1!\frac{e^{-\tau_1^{(r)}}(\tau_1^{(r)})^{x_1}}{x_1!}\right),
\end{aligned}
\end{equation}
在最后一个等式中，把含$\beta$和$\tau_i$的项以及不含这些参数的项分别归为一组.由于计算期望对数似然的目的是关于$\beta$和$\tau_i$将其极大化，可以忽略第二组括号中的项.因此，只需极大化第一组括号中的项，其中最后一个和可写成
\begin{equation}\label{eq:7.2.22}
-\tau_1+\log\tau_1\sum_{x_1=0}^\infty x_1\frac{e^{-\tau_1^{(r)}}(\tau_1^{(r)})^{x_1}}{x_1!}=-\tau_1+\tau_1^{(r)}\log\tau_1.
\end{equation}
将此代回\eqref{eq:7.2.21}可见，期望完全数据似然与原来的完全数据似然相同，区别仅在于把$x_1$替换为$\tau_1^{(r)}$.所以，第$r$步的MLE只对\eqref{eq:7.2.12}稍作改动，为
\begin{equation}\label{eq:7.2.23}
\begin{aligned}
\widehat\beta^{(r+1)}&=\frac{\sum_{i=1}^ny_i}{\tau_1^{(r)}+\sum_{i=2}^nx_i},\qquad
\widehat\tau_1^{(r+1)}=\frac{\widehat\tau_1^{(r)}+y_1}{\widehat\beta^{(r+1)}+1},\\
\widehat\tau_j^{(r+1)}&=\frac{x_j+y_j}{\widehat\beta^{(r+1)}+1},\qquad j=2,3,\ldots,n.
\end{aligned}
\end{equation}
这既定义了E步（用$\widehat\tau_1^{(r)}$替换$x_1$），也定义了M步（按\eqref{eq:7.2.23}计算第$r$次迭代的\mbox{MLE）.}EM算法的性质保证序列$(\widehat\beta^{(r)},\widehat\tau_1^{(r)},\widehat\tau_2^{(r)},\ldots,\widehat\tau_n^{(r)})$在$r\to\infty$时收敛到不完全数据MLE.更多内容见\autoref{exercise:7.27}.
\end{example}
我们不给出EM序列$\{\widehat\theta^{(r)}\}$收敛到不完全数据MLE的完整证明，但下面的关键性质提示了这一点.证明留作\autoref{exercise:7.31}.
\begin{theorem}[单调EM序列]\label{thm:7.2.20}\index{SI0010@收敛}
由\eqref{eq:7.2.20}定义的序列$\{\widehat\theta^{(r)}\}$满足
\begin{equation}\label{eq:7.2.24}
L(\widehat\theta^{(r+1)}\mid\boldsymbol y)\ge L(\widehat\theta^{(r)}\mid\boldsymbol y),
\end{equation}
等号成立当且仅当相继迭代给出相同的极大化期望完全数据对数似然值，即\footnote{译注：一般情况下，这里“当且仅当”的充分方向还需缺失数据条件分布相同.由\eqref{eq:7.2.19}，观测对数似然的增量等于下面期望差加上$D_{\mathrm{KL}}(k(\cdot\mid\widehat\theta^{(r)},\boldsymbol y)\Vert k(\cdot\mid\widehat\theta^{(r+1)},\boldsymbol y))$.两项均非负，故观测似然相等要求二者都为零；仅期望差为零一般不够.}
\begin{align*}
&E[\log L(\widehat\theta^{(r+1)}\mid\boldsymbol y,\boldsymbol X)\mid\widehat\theta^{(r)},\boldsymbol y]\\
&\qquad=E[\log L(\widehat\theta^{(r)}\mid\boldsymbol y,\boldsymbol X)\mid\widehat\theta^{(r)},\boldsymbol y].
\end{align*}
\end{theorem}



\sourcepageanchor{298}
\subsection{评价估计量的方法}\label{sec:7.3}
上一节介绍了构造参数点估计量的一些合理技巧.不过，一个困难随之出现：在具体情形中通常不止一种方法可用，因此经常需要在不同估计量之间作选择.当然，不同构造方法可能给出相同答案，使评价稍微容易些；但很多时候，不同方法会产生不同估计量.

评价统计程序的一般问题，属于称为决策论（decision theory）的统计学分支，\autoref{sec:7.3.4}将较详细地介绍它.不过，在获得一些关于性能的线索之前，不应考虑采用任何程序.本节将介绍评价估计量的若干基本准则，并据此考察一些估计量.
\subsubsection{均方误差}\label{sec:7.3.1}
先研究在有限样本下衡量估计量质量的方法，从均方误差开始.
\begin{definition}\label{def:7.3.1}\index{SI0683@均方误差}
参数$\theta$的估计量$W$的均方误差（mean squared error，MSE）是定义为$E_\theta(W-\theta)^2$的$\theta$的函数.
\end{definition}
注意，MSE衡量估计量$W$与参数$\theta$之差的平方的平均值，作为点估计量的性能指标是相当合理的.一般而言，绝对距离$|W-\theta|$的任意递增函数都可用来衡量估计量的好坏（平均绝对误差$E_\theta(|W-\theta|)$是一个合理替代）.但MSE比其他距离指标至少有两个优点：首先，解析处理十分方便；其次，它具有如下解释：
\begin{equation}\label{eq:7.3.1}\index{SI0801@方差}
E_\theta(W-\theta)^2=\operatorname{Var}_\theta W+(E_\theta W-\theta)^2=\operatorname{Var}_\theta W+(\operatorname{Bias}_\theta W)^2,
\end{equation}
其中，估计量的偏差定义如下.
\begin{definition}\label{def:7.3.2}\index{SI0090@偏差}\index{SI1044@无偏估计量}
参数$\theta$的点估计量$W$的偏差（bias），是$W$的期望与$\theta$之差，即$\operatorname{Bias}_\theta W=E_\theta W-\theta$.偏差作为$\theta$的函数恒等于零的估计量，称为无偏估计量（unbiased estimator），它对所有$\theta$满足$E_\theta W=\theta$.
\end{definition}
因此，MSE包括两个分量：一个衡量估计量的变异性（精密度，precision），另一个衡量其偏差（准确度，accuracy）.具有良好MSE性质的估计量，其方差与偏差的组合较小.为了找到MSE性质良好的估计量，需要寻找同时控制方差和偏差的估计量.显然，无偏估计量在控制偏差方面做得很好.

对于无偏估计量，有
$$E_\theta(W-\theta)^2=\operatorname{Var}_\theta W,$$
所以，估计量若无偏，其MSE就等于方差.
\begin{example}[正态分布的MSE]\label{ex:7.3.3}\index{SI1055@正态}
设$X_1,\ldots,X_n$独立同分布于$n(\mu,\sigma^2)$.统计量$\bar X$与$S^2$都是无偏估计量，因为
$$E\bar X=\mu,\qquad ES^2=\sigma^2,\qquad\text{对所有 }\mu\text{ 和 }\sigma^2.$$
\sourcepageanchor{299}
（不假设正态性时也成立，见\autoref{thm:5.2.6}.）这些估计量的MSE为
\begin{align*}
E(\bar X-\mu)^2&=\operatorname{Var}\bar X=\frac{\sigma^2}n,\\
E(S^2-\sigma^2)^2&=\operatorname{Var}S^2=\frac{2\sigma^4}{n-1}.
\end{align*}
即使去掉正态性假设，$\bar X$的MSE仍为$\sigma^2/n$.不过，若放宽正态性假设，上面$S^2$的MSE表达式便不再相同（见\autoref{exercise:5.8}）.
\end{example}
虽然从MSE角度看，很多无偏估计量也很合理，但应注意，控制偏差并不保证控制MSE.特别地，有时方差与偏差之间存在权衡：以较小的偏差增加换取较大的方差减少，从而改善MSE.
\begin{example}[续\autoref{ex:7.3.3}]\label{ex:7.3.4}
$\sigma^2$的另一个估计量是极大似然估计量$\widehat\sigma^2=\frac1n\sum_{i=1}^n(X_i-\bar X)^2=\frac{n-1}nS^2$.直接计算得
$$E\widehat\sigma^2=E\left(\frac{n-1}nS^2\right)=\frac{n-1}n\sigma^2,$$
所以$\widehat\sigma^2$是$\sigma^2$的有偏估计量.其方差也可算出：
$$\operatorname{Var}\widehat\sigma^2=\operatorname{Var}\left(\frac{n-1}nS^2\right)=\left(\frac{n-1}n\right)^2\operatorname{Var}S^2=\frac{2(n-1)\sigma^4}{n^2},$$
因此，MSE为
$$E(\widehat\sigma^2-\sigma^2)^2=\frac{2(n-1)\sigma^4}{n^2}+\left(\frac{n-1}n\sigma^2-\sigma^2\right)^2=\left(\frac{2n-1}{n^2}\right)\sigma^4.$$
于是有
$$E(\widehat\sigma^2-\sigma^2)^2=\left(\frac{2n-1}{n^2}\right)\sigma^4<\left(\frac2{n-1}\right)\sigma^4=E(S^2-\sigma^2)^2,$$
表明$\widehat\sigma^2$的MSE小于$S^2$.因此，通过以偏差换取方差的减少，MSE得到改善.
\end{example}
必须立即指出，上例并不意味着应放弃用$S^2$估计$\sigma^2$.上述论证表明，如果用MSE衡量，平均而言$\widehat\sigma^2$比$S^2$更接近$\sigma^2$.不过，$\widehat\sigma^2$有偏，而且平均而言会低估$\sigma^2$.仅此一点就可能使我们对用$\widehat\sigma^2$估计$\sigma^2$感到不安.此外，也可以认为，MSE虽是位置参数的合理准则，却不是尺度参数的合理准则，因此根本不应作上述比较.（问题之一是，MSE对高估与低估施加同等惩罚，在位置参数情形中这没有问题.但对尺度参数，零是一个自然下界，所以估计问题并不对称.此时用MSE往往对低估较为宽容.）最终结果并不是获得绝对的答案，而是收集更多关于估计量的信息，希望对具体情形选择一个好的估计量.
\sourcepageanchor{300}
一般而言，由于MSE是参数的函数，并不存在一个“最佳”估计量.两个估计量的MSE曲线往往相交，说明每个估计量只在参数空间的一部分中优于另一个.不过，即使这种局部信息，有时也能指导我们在估计量之间作选择.
\begin{example}[二项Bayes估计量的MSE]\label{ex:7.3.5}\index{SI0079@伯努利}\index{SI0800@伯努利}
设$X_1,\ldots,X_n$独立同分布于$\operatorname{Bernoulli}(p)$.$\bar X$作为$p$的估计量，其MSE为
$$E_p(\bar X-p)^2=\operatorname{Var}_p\bar X=\frac{p(1-p)}n.$$
令$Y=\sum X_i$，回顾\autoref{ex:7.2.14}中得到的Bayes估计量$\widehat p_B=(Y+\alpha)/(\alpha+\beta+n)$.它作为$p$的估计量，其MSE为
\begin{align*}
E_p(\widehat p_B-p)^2&=\operatorname{Var}_p\widehat p_B+(\operatorname{Bias}_p\widehat p_B)^2\\
&=\operatorname{Var}_p\left(\frac{Y+\alpha}{\alpha+\beta+n}\right)+\left(E_p\left(\frac{Y+\alpha}{\alpha+\beta+n}\right)-p\right)^2\\
&=\frac{np(1-p)}{(\alpha+\beta+n)^2}+\left(\frac{np+\alpha}{\alpha+\beta+n}-p\right)^2.
\end{align*}
在没有关于$p$的良好先验信息时，可以尝试选择$\alpha$和$\beta$，使$\widehat p_B$的MSE为常数.细节不难算出（见\autoref{exercise:7.33}），取$\alpha=\beta=\sqrt{n/4}$便得到
$$\widehat p_B=\frac{Y+\sqrt{n/4}}{n+\sqrt n}\qquad\text{且}\qquad E(\widehat p_B-p)^2=\frac n{4(n+\sqrt n)^2}.$$
\begin{figure}[H]
\centering\includegraphics[width=.95\linewidth]{assets/chapter7/fig-7-3-1.pdf}
\caption{\autoref{ex:7.3.5}中，样本量$n=4$与$n=400$时$\bar X$和$\widehat p_B$的MSE比较}\label{fig:7.3.1}
\end{figure}
若想根据MSE在$\widehat p_B$与$\bar X$之间选择，\autoref{fig:7.3.1}很有帮助.$n$较小时，$\widehat p_B$是更好的选择（除非坚信$p$接近0或1）；$n$较大时，$\bar X$是更好的选择（除非坚信$p$接近$1/2$）.虽然MSE准则没有说明一个估计量一致优于另一个，但仍提供了有用信息.将这些信息与对实际问题的了解结合起来，就能为具体情形选出较好的估计量.
\end{example}
\sourcepageanchor{301}
在某些情形中，尤其估计位置参数时，MSE有助于在一类等变估计量中寻找最佳者（\autoref{sec:6.4}）.对群$\mathcal G$中固定的$g$，用$\bar g(\theta)$表示把$\theta$映到$\theta'$的函数，即$\bar g(\theta)=\theta'$.若$W(\boldsymbol X)$估计$\theta$，则有
\begin{quote}
测量等变性：$W(\boldsymbol x)$估计$\theta$ $\Rightarrow$ $\bar g(W(\boldsymbol x))$估计$\bar g(\theta)=\theta'$.

形式不变性：\index{SI0788@形式化}\semanticindexanchor{788:301}$W(\boldsymbol x)$估计$\theta$ $\Rightarrow$ $W(g(\boldsymbol x))$估计$\bar g(\theta)=\theta'$.
\end{quote}
将这两个要求结合，得到$W(g(\boldsymbol x))=\bar g(W(\boldsymbol x))$.
\begin{example}[等变估计量的MSE]\label{ex:7.3.6}\index{SI0464@位置}\index{SI0789@均方误差}
设$X_1,\ldots,X_n$独立同分布于$f(x-\theta)$.估计量$W(X_1,\ldots,X_n)$若满足$W(g_a(\boldsymbol x))=\bar g_a(W(\boldsymbol x))$，就必须有
\begin{equation}\label{eq:7.3.2}
W(x_1,\ldots,x_n)+a=W(x_1+a,\ldots,x_n+a),
\end{equation}
它规定了关于变换群$\mathcal G=\{g_a(\boldsymbol x):-\infty<a<\infty\}$的等变估计量，其中$g_a(x_1,\ldots,x_n)=(x_1+a,\ldots,x_n+a)$.对这些估计量，有
\begin{equation}\label{eq:7.3.3}
\begin{aligned}
&E_\theta(W(X_1,\ldots,X_n)-\theta)^2\\
&=E_\theta(W(X_1+a,\ldots,X_n+a)-a-\theta)^2\\
&=E_\theta(W(X_1-\theta,\ldots,X_n-\theta))^2\qquad(a=-\theta)\\
&=\int_{-\infty}^\infty\cdots\int_{-\infty}^\infty(W(x_1-\theta,\ldots,x_n-\theta))^2\prod_{i=1}^nf(x_i-\theta)\,dx_i\\
&=\int_{-\infty}^\infty\cdots\int_{-\infty}^\infty(W(u_1,\ldots,u_n))^2\prod_{i=1}^nf(u_i)\,du_i.\quad(u_i=x_i-\theta)
\end{aligned}
\end{equation}
最后一个表达式不依赖$\theta$，所以这些等变估计量的MSE不是$\theta$的函数.因而可以用MSE排列它们的优劣，并找到MSE最小的等变估计量.实际上，这一估计量是如下数学问题的解：在约束\eqref{eq:7.3.2}下，寻找使\eqref{eq:7.3.3}最小的函数$W$.（见\autoref{exercise:7.35}和\autoref{exercise:7.36}.）
\end{example}
\subsubsection{最佳无偏估计量}\label{sec:7.3.2}
如上一小节所述，基于MSE比较估计量，未必能选出明确的优胜者.事实上，不存在一个“最佳MSE”估计量.很多人觉得这令人烦恼，希望得到一个“推荐”的估计量，而不愿逐一比较候选者的MSE.

不存在一个“最佳MSE”估计量，是因为所有估计量构成的类太大.（例如，在$\theta=17$处，估计量$\widehat\theta=17$的MSE无可超越，但在其他地方它是一个很差的估计量.）使寻找“最佳”估计量的问题变得可处理的一种办法，是限制估计量的类.一种常用限制方式，也是本小节要考虑的方式，就是只考虑无偏估计量.

若$W_1$和$W_2$都是参数$\theta$的无偏估计量，即$E_\theta W_1=E_\theta W_2=\theta$，那么它们的均方误差等于方差，因此应选择方差较小的估计量.若能找到方差一致最小的无偏估计量——最佳无偏估计量——任务便完成了.

继续之前要指出，虽然我们将讨论无偏估计量，但本小节及下一小节的结果实际上更一般.设$W^*$是$\theta$的一个估计量，满足$E_\theta W^*=\tau(\theta)\ne\theta$，而我们希望考察$W^*$的价值.考虑估计量类
$$\mathcal C_\tau=\{W:E_\theta W=\tau(\theta)\}.$$
\sourcepageanchor{302}
对任意$W_1,W_2\in\mathcal C_\tau$，有$\operatorname{Bias}_\theta W_1=\operatorname{Bias}_\theta W_2$，所以
$$E_\theta(W_1-\theta)^2-E_\theta(W_2-\theta)^2=\operatorname{Var}_\theta W_1-\operatorname{Var}_\theta W_2,$$
在$\mathcal C_\tau$内部作MSE比较，只需比较方差.因此，虽然使用无偏估计量的措辞，我们实际比较的是具有相同期望$\tau(\theta)$的估计量.

本小节的目标是研究寻找“最佳”无偏估计量的方法，其定义如下.
\begin{definition}\label{def:7.3.7}\index{SI1045@最优}
若估计量$W^*$对所有$\theta$满足$E_\theta W^*=\tau(\theta)$，而且对于任一其他满足$E_\theta W=\tau(\theta)$的估计量$W$，对所有$\theta$都有$\operatorname{Var}_\theta W^*\le\operatorname{Var}_\theta W$，则$W^*$称为$\tau(\theta)$的最佳无偏估计量（best unbiased estimator）.$W^*$也称为$\tau(\theta)$的一致最小方差无偏估计量（uniform minimum variance unbiased estimator，UMVUE）.
\end{definition}
由于许多原因，寻找最佳无偏估计量（假如存在的话！）并不容易.下例展示其中两个原因.
\begin{example}[Poisson分布的无偏估计]\label{ex:7.3.8}\index{SI1057@泊松}
设$X_1,\ldots,X_n$独立同分布于$\operatorname{Poisson}(\lambda)$，$\bar X$和$S^2$分别为样本均值与样本方差.回顾Poisson分布的均值和方差均等于$\lambda$.因此，应用\autoref{thm:5.2.6}得到
$$E_\lambda\bar X=\lambda,\qquad\text{对所有 }\lambda,$$
以及
$$E_\lambda S^2=\lambda,\qquad\text{对所有 }\lambda,$$
所以$\bar X$与$S^2$都是$\lambda$的无偏估计量.

为了确定$\bar X$与$S^2$中哪个更好，应比较方差.同样由\autoref{thm:5.2.6}，有$\operatorname{Var}_\lambda\bar X=\lambda/n$，但$\operatorname{Var}_\lambda S^2$的计算很长（类似\hyperref[exercise:5.10]{习题5.10(b)}）.这是寻找最佳无偏估计量首先遇到的问题之一：计算不仅冗长繁琐，还可能白费力气（本例即如此），因为我们将看到，对所有$\lambda$都有$\operatorname{Var}_\lambda\bar X\le\operatorname{Var}_\lambda S^2$.

即使能够确立$\bar X$优于$S^2$，还要考虑估计量类
$$W_a(\bar X,S^2)=a\bar X+(1-a)S^2.$$
对任意常数$a$，都有$E_\lambda W_a(\bar X,S^2)=\lambda$，所以现在有无穷多个$\lambda$的无偏估计量.即使$\bar X$优于$S^2$，它也优于每个$W_a(\bar X,S^2)$吗？而且，如何确信没有其他更好的无偏估计量隐藏在别处？
\end{example}
这个例子展示了寻找最佳无偏估计量时可能遇到的若干问题，也提示需要一种更全面的方法.假设要估计分布$f(\boldsymbol x\mid\theta)$的参数$\tau(\theta)$，能够给出任一无偏估计量方差的下界，记为$B(\theta)$.若能找到无偏估计量$W^*$满足$\operatorname{Var}_\theta W^*=B(\theta)$，便找到了最佳无偏估计量.这就是使用Cramér--Rao下界（Cramér--Rao lower bound）的思路.
\begin{theorem}[Cramér--Rao不等式]\label{thm:7.3.9}\index{SI0228@Cramér–Rao下界}
设$X_1,\ldots,X_n$是具有pdf $f(\boldsymbol x\mid\theta)$的样本，$W(\boldsymbol X)=W(X_1,\ldots,X_n)$为满足下列条件的任意估计量：
\begin{equation}\label{eq:7.3.4}
\begin{aligned}
\frac d{d\theta}E_\theta W(\boldsymbol X)&=\int_{\mathcal X}\frac\partial{\partial\theta}[W(\boldsymbol x)f(\boldsymbol x\mid\theta)]\,d\boldsymbol x,\\
\text{且}\qquad\operatorname{Var}_\theta W(\boldsymbol X)&<\infty.
\end{aligned}
\end{equation}
\sourcepageanchor{303}
则
\begin{equation}\label{eq:7.3.5}
\operatorname{Var}_\theta(W(\boldsymbol X))\ge\frac{\bigl(\frac d{d\theta}E_\theta W(\boldsymbol X)\bigr)^2}{E_\theta\left(\left(\frac\partial{\partial\theta}\log f(\boldsymbol X\mid\theta)\right)^2\right)}.
\end{equation}
\end{theorem}
\begin{proof}
此定理的证明简洁优美，是对Cauchy--Schwarz不等式的巧妙应用；用统计语言说，就是对任意两个随机变量$X$和$Y$都有
\begin{equation}\label{eq:7.3.6}
[\operatorname{Cov}(X,Y)]^2\le(\operatorname{Var}X)(\operatorname{Var}Y).
\end{equation}
重新整理\eqref{eq:7.3.6}，可得$X$的方差下界：
$$\operatorname{Var}X\ge\frac{[\operatorname{Cov}(X,Y)]^2}{\operatorname{Var}Y}.$$
这个定理的巧妙之处，在于选$X$为估计量$W(\boldsymbol X)$，选$Y$为$\frac\partial{\partial\theta}\log f(\boldsymbol X\mid\theta)$，然后应用Cauchy--Schwarz不等式.

首先注意
\begin{equation}\label{eq:7.3.7}
\begin{aligned}
\frac d{d\theta}E_\theta W(\boldsymbol X)&=\int_{\mathcal X}W(\boldsymbol x)\left[\frac\partial{\partial\theta}f(\boldsymbol x\mid\theta)\right]d\boldsymbol x\\
&=E_\theta\left[W(\boldsymbol X)\frac{\frac\partial{\partial\theta}f(\boldsymbol X\mid\theta)}{f(\boldsymbol X\mid\theta)}\right]\quad\text{（乘以$f/f$）}\\
&=E_\theta\left[W(\boldsymbol X)\frac\partial{\partial\theta}\log f(\boldsymbol X\mid\theta)\right],\quad\text{（对数的性质）}
\end{aligned}
\end{equation}
这提示了$W(\boldsymbol X)$与$\frac\partial{\partial\theta}\log f(\boldsymbol X\mid\theta)$之间的协方差.要成为协方差，还须减去期望的乘积，因此计算$E_\theta(\frac\partial{\partial\theta}\log f(\boldsymbol X\mid\theta))$.若取$W(\boldsymbol x)=1$应用\eqref{eq:7.3.7}，有\footnote{译注：这一步还要求微分与积分的交换对$W\equiv1$成立，而定理原述仅对所考察的$W$列出\eqref{eq:7.3.4}.通常的正则性条件还保证此项交换、共同支撑以及所用信息量为正且有限；应用定理时须一并检查，不能仅凭某个$W$的交换条件推出得分期望为零.}
\begin{equation}\label{eq:7.3.8}
E_\theta\left(\frac\partial{\partial\theta}\log f(\boldsymbol X\mid\theta)\right)=\frac d{d\theta}E_\theta[1]=0.
\end{equation}
因此，$\operatorname{Cov}_\theta(W(\boldsymbol X),\frac\partial{\partial\theta}\log f(\boldsymbol X\mid\theta))$等于乘积的期望，由\eqref{eq:7.3.7}和\eqref{eq:7.3.8}可得
\begin{equation}\label{eq:7.3.9}
\begin{aligned}
\operatorname{Cov}_\theta\left(W(\boldsymbol X),\frac\partial{\partial\theta}\log f(\boldsymbol X\mid\theta)\right)
&=E_\theta\left(W(\boldsymbol X)\frac\partial{\partial\theta}\log f(\boldsymbol X\mid\theta)\right)\\
&=\frac d{d\theta}E_\theta W(\boldsymbol X).
\end{aligned}
\end{equation}
另外，由于$E_\theta(\frac\partial{\partial\theta}\log f(\boldsymbol X\mid\theta))=0$，有
\begin{equation}\label{eq:7.3.10}
\operatorname{Var}_\theta\left(\frac\partial{\partial\theta}\log f(\boldsymbol X\mid\theta)\right)=E_\theta\left(\left(\frac\partial{\partial\theta}\log f(\boldsymbol X\mid\theta)\right)^2\right).
\end{equation}
结合Cauchy--Schwarz不等式、\eqref{eq:7.3.9}和\eqref{eq:7.3.10}，得到
$$\operatorname{Var}_\theta(W(\boldsymbol X))\ge\frac{\bigl(\frac d{d\theta}E_\theta W(\boldsymbol X)\bigr)^2}{E_\theta\left(\left(\frac\partial{\partial\theta}\log f(\boldsymbol X\mid\theta)\right)^2\right)},$$
定理得证.
\end{proof}
\sourcepageanchor{304}
若再加上样本独立的假设，下界的计算便得到简化.分母中的期望化为一元计算，如下述推论所示.
\begin{corollary}[Cramér--Rao不等式，独立同分布情形]\label{cor:7.3.10}
若\autoref{thm:7.3.9}的假设满足，而且$X_1,\ldots,X_n$独立同分布，具有pdf $f(x\mid\theta)$，则
$$\operatorname{Var}_\theta W(\boldsymbol X)\ge\frac{\bigl(\frac d{d\theta}E_\theta W(\boldsymbol X)\bigr)^2}{nE_\theta\left(\left(\frac\partial{\partial\theta}\log f(X\mid\theta)\right)^2\right)}.$$
\end{corollary}
\begin{proof}
只需证明
$$E_\theta\left(\left(\frac\partial{\partial\theta}\log f(\boldsymbol X\mid\theta)\right)^2\right)=nE_\theta\left(\left(\frac\partial{\partial\theta}\log f(X\mid\theta)\right)^2\right).$$
由于$X_1,\ldots,X_n$相互独立，
\begin{equation}\label{eq:7.3.11}
\begin{aligned}
E_\theta\left(\frac\partial{\partial\theta}\log f(\boldsymbol X\mid\theta)\right)^2
&=E_\theta\left(\left(\frac\partial{\partial\theta}\log\prod_{i=1}^nf(X_i\mid\theta)\right)^2\right)\\
&=E_\theta\left(\left(\sum_{i=1}^n\frac\partial{\partial\theta}\log f(X_i\mid\theta)\right)^2\right)\quad\text{（对数性质）}\\
&=\sum_{i=1}^nE_\theta\left(\left(\frac\partial{\partial\theta}\log f(X_i\mid\theta)\right)^2\right)\quad\text{（展开平方）}\\
&\quad+\sum_{i\ne j}E_\theta\left(\frac\partial{\partial\theta}\log f(X_i\mid\theta)\frac\partial{\partial\theta}\log f(X_j\mid\theta)\right).
\end{aligned}
\end{equation}
对于$i\ne j$，有
\begin{align*}
&E_\theta\left(\frac\partial{\partial\theta}\log f(X_i\mid\theta)\frac\partial{\partial\theta}\log f(X_j\mid\theta)\right)\\
&\quad=E_\theta\left(\frac\partial{\partial\theta}\log f(X_i\mid\theta)\right)E_\theta\left(\frac\partial{\partial\theta}\log f(X_j\mid\theta)\right)\quad\text{（独立性）}\\
&\quad=0.\qquad\text{（由\eqref{eq:7.3.8}）}
\end{align*}
所以\eqref{eq:7.3.11}的第二个和为零，第一个和为
\begin{align*}
&\sum_{i=1}^nE_\theta\left(\left(\frac\partial{\partial\theta}\log f(X_i\mid\theta)\right)^2\right)\\
&\qquad=nE_\theta\left(\left(\frac\partial{\partial\theta}\log f(X\mid\theta)\right)^2\right),\qquad\text{（同分布）}
\end{align*}
从而确立推论.
\end{proof}
继续之前指出，虽然Cramér--Rao下界是针对连续随机变量陈述的，但它也适用于离散随机变量.允许交换积分与微分的关键条件\eqref{eq:7.3.4}，只需作显然的修改.若$f(x\mid\theta)$是pmf，就必须能够交换微分与求和.（当然，这假设虽然$f(x\mid\theta)$是pmf，关于$x$不可微，但关于$\theta$可微.大多数常见pmf都如此.）
\sourcepageanchor{305}
量$E_\theta((\frac\partial{\partial\theta}\log f(\boldsymbol X\mid\theta))^2)$称为样本的信息数（information number），或Fisher信息（Fisher information）.\index{SI0582@信息量}\semanticindexanchor{582:305}这个术语反映了如下事实：信息数给出了$\theta$的最佳无偏估计量方差的下界.信息数越大，即关于$\theta$的信息越多，最佳无偏估计量的方差下界就越小.

事实上，Cramér--Rao不等式也称信息不等式（Information Inequality），而信息不等式存在比这里更一般的形式.更一般形式的一个关键区别，是去掉对候选估计量的全部假设，改为对潜在密度作假设.在这种形式下，信息不等式对于比较估计量性能十分有用.细节见\hyperref[bib:source200]{Lehmann和Casella（1998）}第2.6节.

现在，对于任意可微函数$\tau(\theta)$，我们得到了满足\eqref{eq:7.3.4}及$E_\theta W=\tau(\theta)$的任意估计量$W$的方差下界.此下界只依赖于$\tau(\theta)$和$f(\boldsymbol x\mid\theta)$，是方差的一致下界.任何满足$E_\theta W=\tau(\theta)$且达到下界的候选估计量，都是$\tau(\theta)$的最佳无偏估计量.

在看例子之前，先给出一个有助于应用此定理的计算结果.证明留作\autoref{exercise:7.39}.
\begin{lemma}\label{lem:7.3.11}
若$f(x\mid\theta)$满足
$$\frac d{d\theta}E_\theta\left(\frac\partial{\partial\theta}\log f(X\mid\theta)\right)=\int\frac\partial{\partial\theta}\left[\left(\frac\partial{\partial\theta}\log f(x\mid\theta)\right)f(x\mid\theta)\right]dx$$
（指数族满足此条件），则
$$E_\theta\left(\left(\frac\partial{\partial\theta}\log f(X\mid\theta)\right)^2\right)=-E_\theta\left(\frac{\partial^2}{\partial\theta^2}\log f(X\mid\theta)\right).$$
\end{lemma}
利用刚建立的工具，回到Poisson例子并解决它.
\begin{example}[\autoref{ex:7.3.8}的结尾]\label{ex:7.3.12}\index{SI1049@泊松}
这里$\tau(\lambda)=\lambda$，所以$\tau'(\lambda)=1$.又因这是指数族，利用\autoref{lem:7.3.11}得到
\begin{align*}
E_\lambda\left(\left(\frac\partial{\partial\lambda}\log\prod_{i=1}^nf(X_i\mid\lambda)\right)^2\right)
&=-nE_\lambda\left(\frac{\partial^2}{\partial\lambda^2}\log f(X\mid\lambda)\right)\\
&=-nE_\lambda\left(\frac{\partial^2}{\partial\lambda^2}\log\left(\frac{e^{-\lambda}\lambda^X}{X!}\right)\right)\\
&=-nE_\lambda\left(\frac{\partial^2}{\partial\lambda^2}(-\lambda+X\log\lambda-\log X!)\right)\\
&=-nE_\lambda\left(-\frac X{\lambda^2}\right)\\
&=\frac n\lambda.
\end{align*}
因此，$\lambda$的任意无偏估计量$W$都必须满足
$$\operatorname{Var}_\lambda W\ge\frac\lambda n.$$
由于$\operatorname{Var}_\lambda\bar X=\lambda/n$，$\bar X$就是$\lambda$的最佳无偏估计量.
\end{example}
务必记住，Cramér--Rao定理的一个关键假设是能够在积分号下求导，这当然有一定限制.如已看到的，指数类中的密度满足这些假设，但一般情况下仍须检查，否则会出现下面这样的矛盾.
\sourcepageanchor{306}
\begin{example}[尺度均匀分布的无偏估计量]\label{ex:7.3.13}\index{SI1059@均匀}
设$X_1,\ldots,X_n$独立同分布，pdf为$f(x\mid\theta)=1/\theta$，$0<x<\theta$.因为$\frac\partial{\partial\theta}\log f(x\mid\theta)=-1/\theta$，有
$$E_\theta\left(\left(\frac\partial{\partial\theta}\log f(X\mid\theta)\right)^2\right)=\frac1{\theta^2}.$$
Cramér--Rao定理似乎意味着，若$W$是$\theta$的任意无偏估计量，则
$$\operatorname{Var}_\theta W\ge\frac{\theta^2}n.$$
现在希望寻找方差小的无偏估计量.首先考虑充分统计量$Y=\max(X_1,\ldots,X_n)$，即最大次序统计量.$Y$的pdf为$f_Y(y\mid\theta)=ny^{n-1}/\theta^n$，$0<y<\theta$，所以
$$E_\theta Y=\int_0^\theta\frac{ny^n}{\theta^n}\,dy=\frac n{n+1}\theta,$$
说明$\frac{n+1}nY$是$\theta$的无偏估计量.接着计算
\begin{align*}
\operatorname{Var}_\theta\left(\frac{n+1}nY\right)&=\left(\frac{n+1}n\right)^2\operatorname{Var}_\theta Y\\
&=\left(\frac{n+1}n\right)^2\left[E_\theta Y^2-\left(\frac n{n+1}\theta\right)^2\right]\\
&=\left(\frac{n+1}n\right)^2\left[\frac n{n+2}\theta^2-\left(\frac n{n+1}\theta\right)^2\right]\\
&=\frac1{n(n+2)}\theta^2,
\end{align*}
它一致小于$\theta^2/n$.这表明Cramér--Rao定理不适用于这个pdf.为说明这一点，用Leibnitz法则（\autoref{sec:2.4}）计算
\begin{align*}
\frac d{d\theta}\int_0^\theta h(x)f(x\mid\theta)\,dx
&=\frac d{d\theta}\int_0^\theta h(x)\frac1\theta\,dx\\
&=\frac{h(\theta)}\theta+\int_0^\theta h(x)\frac\partial{\partial\theta}\left(\frac1\theta\right)dx\\
&\ne\int_0^\theta h(x)\frac\partial{\partial\theta}f(x\mid\theta)\,dx,
\end{align*}
除非对所有$\theta$都有$h(\theta)/\theta=0$.因此，Cramér--Rao定理不适用.一般而言，若pdf的取值范围依赖参数，该定理就不适用.
\end{example}
用这种方法寻找最佳无偏估计量有一个不足：即使Cramér--Rao定理适用，也不能保证下界是精确的.也就是说，Cramér--Rao下界可能严格小于任一无偏估计量的方差.事实上，即使$f(x\mid\theta)$是单参数指数族这一通常有利的情形，我们最多也只能说，存在某个参数$\tau(\theta)$，其某个无偏估计量达到Cramér--Rao下界.但在其他典型情形中，对其他参数，这个下界可能无法达到.这令人担忧：如果找不到达到下界的估计量，必须判断是任何估计量都无法达到，还是我们需要继续考察更多估计量.
\sourcepageanchor{307}
\begin{example}[正态方差的下界]\label{ex:7.3.14}\index{SI1060@方差界}
设$X_1,\ldots,X_n$独立同分布于$n(\mu,\sigma^2)$，考虑在$\mu$未知时估计$\sigma^2$.正态pdf满足Cramér--Rao定理及\autoref{lem:7.3.11}的假设，因此有
$$\frac{\partial^2}{\partial(\sigma^2)^2}\log\left(\frac1{(2\pi\sigma^2)^{1/2}}e^{-(1/2)(x-\mu)^2/\sigma^2}\right)=\frac1{2\sigma^4}-\frac{(x-\mu)^2}{\sigma^6},$$
以及
\begin{align*}
-E\left(\frac{\partial^2}{\partial(\sigma^2)^2}\log f(X\mid\mu,\sigma^2)\,\middle|\,\mu,\sigma^2\right)
&=-E\left(\frac1{2\sigma^4}-\frac{(X-\mu)^2}{\sigma^6}\,\middle|\,\mu,\sigma^2\right)\\
&=\frac1{2\sigma^4}.
\end{align*}
因此，$\sigma^2$的任意无偏估计量$W$必须满足
$$\operatorname{Var}(W\mid\mu,\sigma^2)\ge\frac{2\sigma^4}n.$$
在\autoref{ex:7.3.3}中，已经看到
$$\operatorname{Var}(S^2\mid\mu,\sigma^2)=\frac{2\sigma^4}{n-1},$$
所以$S^2$没有达到Cramér--Rao下界.
\end{example}
上例留下了一个不完整的答案：是否有比$S^2$更好的$\sigma^2$无偏估计量，还是Cramér--Rao下界本来就无法达到？

达到Cramér--Rao下界的条件其实很简单.回顾此下界来自Cauchy--Schwarz不等式，所以达到下界的条件，就是Cauchy--Schwarz不等式取等号的条件（见\autoref{sec:4.7}）.还要注意，\autoref{cor:7.3.15}是一个有用工具，因为它隐含地给出了寻找最佳无偏估计量的方法.
\begin{corollary}[下界的达到]\label{cor:7.3.15}\index{SI0229@必要条件}
设$X_1,\ldots,X_n$独立同分布于$f(x\mid\theta)$，其中$f(x\mid\theta)$满足Cramér--Rao定理的条件.用$L(\theta\mid\boldsymbol x)=\prod_{i=1}^nf(x_i\mid\theta)$表示似然函数.若$W(\boldsymbol X)=W(X_1,\ldots,X_n)$为$\tau(\theta)$的任意无偏估计量，则$W(\boldsymbol X)$达到Cramér--Rao下界，当且仅当对某个函数$a(\theta)$，有
\begin{equation}\label{eq:7.3.12}
a(\theta)[W(\boldsymbol x)-\tau(\theta)]=\frac\partial{\partial\theta}\log L(\theta\mid\boldsymbol x).
\end{equation}
\end{corollary}
\begin{proof}
\eqref{eq:7.3.6}中的Cramér--Rao不等式可写为
\begin{align*}
&\left[\operatorname{Cov}_\theta\left(W(\boldsymbol X),\frac\partial{\partial\theta}\log\prod_{i=1}^nf(X_i\mid\theta)\right)\right]^2\\
&\qquad\le\operatorname{Var}_\theta W(\boldsymbol X)\operatorname{Var}_\theta\left(\frac\partial{\partial\theta}\log\prod_{i=1}^nf(X_i\mid\theta)\right).
\end{align*}
回顾$E_\theta W=\tau(\theta)$以及$E_\theta(\frac\partial{\partial\theta}\log\prod_{i=1}^nf(X_i\mid\theta))=0$，再利用\autoref{thm:4.5.7}的结果，等号成立当且仅当$W(\boldsymbol x)-\tau(\theta)$与$\frac\partial{\partial\theta}\log\prod_{i=1}^nf(x_i\mid\theta)$成比例.这正是\eqref{eq:7.3.12}所表达的内容.
\end{proof}
\begin{example}[续\autoref{ex:7.3.14}]\label{ex:7.3.16}
这里有
$$L(\mu,\sigma^2\mid\boldsymbol x)=\frac1{(2\pi\sigma^2)^{n/2}}e^{-(1/2)\sum_{i=1}^n(x_i-\mu)^2/\sigma^2},$$
因此
$$\frac\partial{\partial\sigma^2}\log L(\mu,\sigma^2\mid\boldsymbol x)=\frac n{2\sigma^4}\left(\sum_{i=1}^n\frac{(x_i-\mu)^2}n-\sigma^2\right).$$
\sourcepageanchor{308}
取$a(\sigma^2)=n/(2\sigma^4)$，说明$\sigma^2$的最佳无偏估计量为$\sum_{i=1}^n(x_i-\mu)^2/n$，但只有$\mu$已知时才能计算它.若$\mu$未知，下界就无法达到.
\end{example}
本小节建立的理论仍留下一些未解决的问题.首先，若$f(x\mid\theta)$不满足Cramér--Rao定理的假设，该怎么办？（在\autoref{ex:7.3.13}中，仍不知道$\frac{n+1}nY$是否为最佳无偏估计量.）其次，若允许使用的估计量无法达到下界，如\autoref{ex:7.3.14}，又该怎么办？那里仍不知道$S^2$是否最佳无偏.

回答这些问题的一种办法，是寻找适用面更广、能给出更精确（即更大）下界的方法.这方面已有许多研究，其中最著名的下界也许是\hyperref[bib:source72]{Chapman和Robbins（1951）}的下界.\hyperref[bib:source311]{Stuart、Ord和Arnold（1999）}第17章对此有很好的介绍.不过，这里不走这条路，而将利用充分性的概念，从另一角度继续研究最佳无偏估计量.
\subsubsection{充分性与无偏性}\label{sec:7.3.3}
上一小节寻找无偏估计时没有用到充分性的概念.现在将看到，考虑充分性确实是一件有力工具.

本小节的主要定理把充分统计量与无偏估计联系起来；与Cramér--Rao定理一样，它又是对熟知定理的巧妙应用.回顾\autoref{sec:4}：若$X$与$Y$为任意两个随机变量，只要期望存在，就有
\begin{equation}\label{eq:7.3.13}
\begin{aligned}
EX&=E[E(X\mid Y)],\\
\operatorname{Var}X&=\operatorname{Var}[E(X\mid Y)]+E[\operatorname{Var}(X\mid Y)].
\end{aligned}
\end{equation}
利用这些工具，可以证明如下定理.
\begin{theorem}[Rao--Blackwell]\label{thm:7.3.17}\index{SI0817@充分性}\index{SI0877@Rao–Blackwell定理}\index{SI1050@充分性}
设$W$为$\tau(\theta)$的任意无偏估计量，$T$为$\theta$的充分统计量.定义$\phi(T)=E(W\mid T)$.则对所有$\theta$，有$E_\theta\phi(T)=\tau(\theta)$以及$\operatorname{Var}_\theta\phi(T)\le\operatorname{Var}_\theta W$；即$\phi(T)$是$\tau(\theta)$的一致更优的无偏估计量.
\end{theorem}
\begin{proof}
由\eqref{eq:7.3.13}，有
$$\tau(\theta)=E_\theta W=E_\theta[E(W\mid T)]=E_\theta\phi(T),$$
所以$\phi(T)$对$\tau(\theta)$无偏.此外，
\begin{align*}
\operatorname{Var}_\theta W&=\operatorname{Var}_\theta[E(W\mid T)]+E_\theta[\operatorname{Var}(W\mid T)]\\
&=\operatorname{Var}_\theta\phi(T)+E_\theta[\operatorname{Var}(W\mid T)]\\
&\ge\operatorname{Var}_\theta\phi(T).\qquad(\operatorname{Var}(W\mid T)\ge0)
\end{align*}
因此，$\phi(T)$一致优于$W$，只剩说明$\phi(T)$确实是估计量.即必须说明$\phi(T)=E(W\mid T)$仅是样本的函数，特别地，不依赖$\theta$.但由充分性的定义，以及$W$仅是样本的函数，可知$W\mid T$的分布不依赖$\theta$.所以$\phi(T)$是$\tau(\theta)$的一致更优无偏估计量.
\end{proof}
因此，将任一无偏估计量对充分统计量取条件期望，会得到一致改进；所以寻找最佳无偏估计量时，只需考虑充分统计量的函数.

\eqref{eq:7.3.13}的恒等式没有提及充分性，所以乍看之下，似乎以任何量为条件都会带来改进.事实上确实如此，但问题是所得量很可能依赖$\theta$，因而不是估计量.
\sourcepageanchor{309}
\begin{example}[以不充分统计量为条件]\label{ex:7.3.18}
设$X_1,X_2$独立同分布于$n(\theta,1)$.统计量$\bar X=\frac12(X_1+X_2)$满足
$$E_\theta\bar X=\theta\qquad\text{且}\qquad\operatorname{Var}_\theta\bar X=\frac12.$$
考虑以不充分的$X_1$为条件.令$\phi(X_1)=E_\theta(\bar X\mid X_1)$.由\eqref{eq:7.3.13}，有$E_\theta\phi(X_1)=\theta$和$\operatorname{Var}_\theta\phi(X_1)\le\operatorname{Var}_\theta\bar X$，所以$\phi(X_1)$优于$\bar X$.然而，
\begin{align*}
\phi(X_1)&=E_\theta(\bar X\mid X_1)\\
&=\frac12E_\theta(X_1\mid X_1)+\frac12E_\theta(X_2\mid X_1)\\
&=\frac12X_1+\frac12\theta,
\end{align*}
因为独立性给出$E_\theta(X_2\mid X_1)=E_\theta X_2$.所以$\phi(X_1)$不是估计量.
\end{example}
现在知道，寻找$\tau(\theta)$的最佳无偏估计量时，只需考虑基于充分统计量的估计量.于是问题是：若$E_\theta\phi=\tau(\theta)$，且$\phi$基于充分统计量，即$E(\phi\mid T)=\phi$，如何知道$\phi$是最佳无偏的？当然，若$\phi$达到Cramér--Rao下界，它就是最佳无偏；但若没有达到，我们究竟获得了什么？例如，若$\phi^*$是$\tau(\theta)$的另一个无偏估计量，$E(\phi^*\mid T)$与$\phi$相比如何？下一个定理通过说明最佳无偏估计量的唯一性，部分回答这个问题.
\begin{theorem}\label{thm:7.3.19}\index{SI1052@唯一}
若$W$是$\tau(\theta)$的最佳无偏估计量，则$W$唯一.\footnote{译注：这里的唯一性是在每个参数对应分布下几乎处处相等的意义下，且证明使用有限方差.证明末由协方差推出$a(\theta)=1$时默认方差为正；若方差为零，两估计量均几乎处处等于同一均值，唯一性可直接得到.}
\end{theorem}
\begin{proof}
假设$W'$是另一个最佳无偏估计量，考虑$W^*=\frac12(W+W')$.注意$E_\theta W^*=\tau(\theta)$，且
\begin{equation}\label{eq:7.3.14}
\begin{aligned}
\operatorname{Var}_\theta W^*&=\operatorname{Var}_\theta\left(\frac12W+\frac12W'\right)\\
&=\frac14\operatorname{Var}_\theta W+\frac14\operatorname{Var}_\theta W'+\frac12\operatorname{Cov}_\theta(W,W')\quad\text{（\autoref{exercise:4.44}）}\\
&\le\frac14\operatorname{Var}_\theta W+\frac14\operatorname{Var}_\theta W'+\frac12[(\operatorname{Var}_\theta W)(\operatorname{Var}_\theta W')]^{1/2}\\
&\hspace{7cm}\text{（Cauchy--Schwarz）}\\
&=\operatorname{Var}_\theta W.\qquad(\operatorname{Var}_\theta W=\operatorname{Var}_\theta W')
\end{aligned}
\end{equation}
但若上述不等式严格成立，就与$W$最佳无偏矛盾，因此对所有$\theta$必须取等号.由于该不等式是Cauchy--Schwarz不等式的应用，只有$W'=a(\theta)W+b(\theta)$时才能取等号.利用协方差的性质，有
\begin{align*}
\operatorname{Cov}_\theta(W,W')&=\operatorname{Cov}_\theta[W,a(\theta)W+b(\theta)]\\
&=\operatorname{Cov}_\theta[W,a(\theta)W]\\
&=a(\theta)\operatorname{Var}_\theta W,
\end{align*}
但因\eqref{eq:7.3.14}取等号，有$\operatorname{Cov}_\theta(W,W')=\operatorname{Var}_\theta W$.因此$a(\theta)=1$；又因$E_\theta W'=\tau(\theta)$，必须有$b(\theta)=0$和$W=W'$，说明$W$唯一.
\end{proof}
为了判断一个无偏估计量何时最佳无偏，可以问：怎样改进给定的无偏估计量？假设$W$满足$E_\theta W=\tau(\theta)$，另有估计量$U$对所有$\theta$满足$E_\theta U=0$，即$U$是零的无偏估计量.估计量
$$\phi_a=W+aU,$$
\sourcepageanchor{310}
其中$a$为常数，满足$E_\theta\phi_a=\tau(\theta)$，因而也是$\tau(\theta)$的无偏估计量.$\phi_a$能优于$W$吗？它的方差为
$$\operatorname{Var}_\theta\phi_a=\operatorname{Var}_\theta(W+aU)=\operatorname{Var}_\theta W+2a\operatorname{Cov}_\theta(W,U)+a^2\operatorname{Var}_\theta U.$$
现在，若对某个$\theta=\theta_0$有$\operatorname{Cov}_{\theta_0}(W,U)<0$，则选择
$$a\in\left(0,-\frac{2\operatorname{Cov}_{\theta_0}(W,U)}{\operatorname{Var}_{\theta_0}U}\right)$$
即可使$2a\operatorname{Cov}_{\theta_0}(W,U)+a^2\operatorname{Var}_{\theta_0}U<0$.因此，$\phi_a$在$\theta=\theta_0$处优于$W$，而$W$不可能最佳无偏.类似论证说明，若对任意某个$\theta_0$有$\operatorname{Cov}_{\theta_0}(W,U)>0$，$W$也不可能最佳无偏（见\autoref{exercise:7.53}）.所以，判断$W$是否最佳无偏时，它与零的无偏估计量之间的关系至关重要.事实上，这种关系刻画了最佳无偏性.
\begin{theorem}\label{thm:7.3.20}\index{SI1047@刻画}
若$E_\theta W=\tau(\theta)$，则$W$是$\tau(\theta)$的最佳无偏估计量，当且仅当$W$与零的所有无偏估计量都不相关.
\end{theorem}
\begin{proof}
若$W$最佳无偏，上面的论证说明，对于任意满足$E_\theta U=0$的$U$，$W$都必须对所有$\theta$满足$\operatorname{Cov}_\theta(W,U)=0$.必要性得到确立.

反之，设无偏估计量$W$与零的所有无偏估计量都不相关.令$W'$是任一其他满足$E_\theta W'=E_\theta W=\tau(\theta)$的估计量.下面说明$W$优于$W'$.写成
$$W'=W+(W'-W),$$
并计算
\begin{equation}\label{eq:7.3.15}
\begin{aligned}
\operatorname{Var}_\theta W'&=\operatorname{Var}_\theta W+\operatorname{Var}_\theta(W'-W)+2\operatorname{Cov}_\theta(W,W'-W)\\
&=\operatorname{Var}_\theta W+\operatorname{Var}_\theta(W'-W),
\end{aligned}
\end{equation}
最后一个等号成立，是因为$W'-W$是零的无偏估计量，按假设与$W$不相关.由于$\operatorname{Var}_\theta(W'-W)\ge0$，\eqref{eq:7.3.15}意味着$\operatorname{Var}_\theta W'\ge\operatorname{Var}_\theta W$.$W'$是任意的，故$W$是$\tau(\theta)$的最佳无偏估计量.
\end{proof}
注意，零的无偏估计量不过是随机噪声（random noise），即零的估计量中没有信息.（估计零最合理的方法是用零，而不是用随机噪声，这很自然.）因此，若一个估计量可以通过加上随机噪声得到改进，那么它大概有缺陷.（也可以质疑评价估计量所用的准则，但这里的准则看来无可非议.）\autoref{thm:7.3.20}正是把这种直觉形式化.

虽然得到了最佳无偏估计量的一个有趣刻画，但实际应用有限.验证一个估计量与零的\emph{所有}无偏估计量不相关往往很难，因为通常难以描述零的所有无偏估计量.不过，有时可以用它判定一个估计量不是最佳无偏的.
\begin{example}[零的无偏估计量]\label{ex:7.3.21}\index{SI1059@均匀}\index{SI1061@零}
设$X$是来自$\operatorname{uniform}(\theta,\theta+1)$分布的一个观测.则
$$E_\theta X=\int_\theta^{\theta+1}x\,dx=\theta+\frac12,$$
所以$X-\frac12$是$\theta$的无偏估计量，而且容易验证$\operatorname{Var}_\theta X=1/12$.

对此pdf，零的无偏估计量是以1为周期的周期函数.这是因为，若$h(x)$满足
$$\int_\theta^{\theta+1}h(x)\,dx=0,\qquad\text{对所有 }\theta,$$
则
$$0=\frac d{d\theta}\int_\theta^{\theta+1}h(x)\,dx=h(\theta+1)-h(\theta),\qquad\text{对所有 }\theta.$$
\sourcepageanchor{311}
这样的函数之一是$h(x)=\sin(2\pi x)$.现在
\begin{align*}
\operatorname{Cov}_\theta\left(X-\frac12,\sin(2\pi X)\right)
&=\operatorname{Cov}_\theta(X,\sin(2\pi X))\\
&=\int_\theta^{\theta+1}x\sin(2\pi x)\,dx\\
&=\left.-\frac{x\cos(2\pi x)}{2\pi}\right|_\theta^{\theta+1}+\int_\theta^{\theta+1}\frac{\cos(2\pi x)}{2\pi}\,dx\quad\text{（分部积分）}\\
&=-\frac{\cos(2\pi\theta)}{2\pi},
\end{align*}
其中用了$\cos(2\pi(\theta+1))=\cos(2\pi\theta)$及$\sin(2\pi(\theta+1))=\sin(2\pi\theta)$.

因此，$X-\frac12$与零的某个无偏估计量相关，不可能是$\theta$的最佳无偏估计量.事实上，容易验证估计量$X-\frac12+\sin(2\pi X)/(2\pi)$对$\theta$无偏，方差为$.071<1/12$.\footnote{译注：原文给出的方差$.071$并非对所有$\theta$成立.由上面的协方差和$\operatorname{Var}(\sin(2\pi X))=1/2$，该估计量的方差为$1/12+1/(8\pi^2)-\cos(2\pi\theta)/(2\pi^2)$，随$\theta$变化，并不一致小于$1/12$.不过，在某些参数值处可改进，已足以说明$X-1/2$不是最佳无偏估计量.}
\end{example}
要回答最佳无偏估计量是否存在，需要对零的所有无偏估计量作某种刻画.有了这种刻画，就能判断候选的最佳无偏估计量是否确实最优.

刻画零的无偏估计量并不容易，需要对所处理的pdf或pmf加上条件.注意，本小节至今没有规定pdf的条件（不像Cramér--Rao下界所需的那些条件）.为这种一般性付出的代价，就是难以验证最佳无偏估计量的存在性.

若pdf或pmf族$f(x\mid\theta)$具有这样一种性质：除零本身之外没有零的无偏估计量，那么寻找过程就结束了，因为任一统计量$W$都满足$\operatorname{Cov}_\theta(W,0)=0$.回顾\hyperref[def:6.2.21]{定义6.1.4}中定义的完备性，它保证这种情形发生.\footnote{译注：原文此处引用“定义6.1.4”，本版对应的完备性定义实际为\autoref{def:6.2.21}.保留原显示编号，并将链接指向实际定义.}
\begin{example}[续\autoref{ex:7.3.13}]\label{ex:7.3.22}\index{SI1051@均匀}
对于独立同分布于$\operatorname{uniform}(0,\theta)$的$X_1,\ldots,X_n$，已经看到$\frac{n+1}nY$是$\theta$的无偏估计量，其中$Y=\max\{X_1,\ldots,X_n\}$.Cramér--Rao定理的条件不满足，因此尚未确立这个估计量是否最佳无偏.不过，\autoref{ex:6.2.23}已经说明$Y$是完备充分统计量.这意味着$Y$的pdf族完备，不存在基于$Y$的非零的零的无偏估计量.（由Rao--Blackwell定理形式的充分性，只需考虑基于$Y$的零的无偏估计量.）因此，$\frac{n+1}nY$与零的所有无偏估计量不相关（因为唯一这样的估计量就是零本身），从而$\frac{n+1}nY$是$\theta$的最佳无偏估计量.
\end{example}
值得再次指出，重要的是充分统计量的分布族的完备性.原始分布族是否完备并不重要.这来自Rao--Blackwell定理：可以把注意力限制在充分统计量的函数上，所以所有期望都针对该统计量的分布计算.

下面的定理概括了完备性与最佳无偏性之间的关系.
\begin{theorem}\label{thm:7.3.23}\index{SI0113@最优无偏估计量}\index{SI1048@完备性}
设$T$是参数$\theta$的完备充分统计量，$\phi(T)$是任一仅基于$T$的估计量.则$\phi(T)$是其期望值的唯一最佳无偏估计量.
\end{theorem}
以本小节理论的一个有趣而实用的应用结束讨论.很多情况下，函数$\tau(\theta)$的无偏估计量没有显而易见的候选者，更不用说最佳无偏估计量.不过，在有完备性的条件下，本小节的理论告诉我们，只要能找到任意一个无偏估计量，就能找到最佳无偏估计量.若$T$是参数$\theta$的完备充分统计量，$h(X_1,\ldots,X_n)$是$\tau(\theta)$的任意无偏估计量，则$\phi(T)=E(h(X_1,\ldots,X_n)\mid T)$就是$\tau(\theta)$的最佳无偏估计量（见\autoref{exercise:7.56}）.
\sourcepageanchor{312}
\begin{example}[二项分布的最佳无偏估计]\label{ex:7.3.24}\index{SI1046@二项}
设$X_1,\ldots,X_n$独立同分布于$\operatorname{binomial}(k,\theta)$.问题是估计该分布中恰有一次成功的概率，即估计
$$\tau(\theta)=P_\theta(X=1)=k\theta(1-\theta)^{k-1}.$$
现在$\sum_{i=1}^nX_i\sim\operatorname{binomial}(kn,\theta)$是完备充分统计量，但基于它的无偏估计量并不一目了然.遇到这种情形，应尝试最简单的解法.朴素的估计量
$$h(X_1)=\begin{cases}1,&\text{若 }X_1=1,\\0,&\text{其他情形}\end{cases}$$
满足
\begin{align*}
E_\theta h(X_1)&=\sum_{x_1=0}^kh(x_1)\binom k{x_1}\theta^{x_1}(1-\theta)^{k-x_1}\\
&=k\theta(1-\theta)^{k-1},
\end{align*}
因此是$k\theta(1-\theta)^{k-1}$的无偏估计量.理论告诉我们，估计量
$$\phi\left(\sum_{i=1}^nX_i\right)=E\left(h(X_1)\,\middle|\,\sum_{i=1}^nX_i\right)$$
是$k\theta(1-\theta)^{k-1}$的最佳无偏估计量.（注意，并不需要实际计算$\phi(\sum_{i=1}^nX_i)$的期望；由迭代期望的性质，已经知道它有正确的期望值.）不过，必须能够算出$\phi$.假设观测到$\sum_{i=1}^nX_i=t$，则
\begin{align*}
\phi(t)&=E\left(h(X_1)\,\middle|\,\sum_{i=1}^nX_i=t\right)&&\text{（期望不依赖$\theta$）}\\
&=P\left(X_1=1\,\middle|\,\sum_{i=1}^nX_i=t\right)&&\text{（$h$为0或1）}\\
&=\frac{P_\theta(X_1=1,\sum_{i=1}^nX_i=t)}{P_\theta(\sum_{i=1}^nX_i=t)}&&\text{（条件概率定义）}\\
&=\frac{P_\theta(X_1=1,\sum_{i=2}^nX_i=t-1)}{P_\theta(\sum_{i=1}^nX_i=t)}&&\text{（$X_1=1$是多余条件）}\\
&=\frac{P_\theta(X_1=1)P_\theta(\sum_{i=2}^nX_i=t-1)}{P_\theta(\sum_{i=1}^nX_i=t)}.&&\text{（$X_1$与其余变量独立）}
\end{align*}
现在$X_1\sim\operatorname{binomial}(k,\theta)$，并且
$$\sum_{i=2}^nX_i\sim\operatorname{binomial}(k(n-1),\theta),\qquad\sum_{i=1}^nX_i\sim\operatorname{binomial}(kn,\theta).$$
利用这些事实，得到
\begin{align*}
\phi(t)&=\frac{[k\theta(1-\theta)^{k-1}]\left[\binom{k(n-1)}{t-1}\theta^{t-1}(1-\theta)^{k(n-1)-(t-1)}\right]}{\binom{kn}t\theta^t(1-\theta)^{kn-t}}\\
&=k\frac{\binom{k(n-1)}{t-1}}{\binom{kn}t}.
\end{align*}
注意，所有$\theta$都消去了；由于$\sum_{i=1}^nX_i$充分，本来就必须如此.所以$k\theta(1-\theta)^{k-1}$的最佳无偏估计量为
$$\phi\left(\sum_{i=1}^nX_i\right)=k\frac{\binom{k(n-1)}{\sum X_i-1}}{\binom{kn}{\sum X_i}}.$$
无需进行困难的$E_\theta[\phi(\sum_{i=1}^nX_i)]$计算，就可以断言无偏性.
\end{example}
\sourcepageanchor{313}
\subsubsection{损失函数下的最优性}\label{sec:7.3.4}\index{SI0244@决策论}
前面对点估计量的评价，基于它们的均方误差表现.均方误差是一种称为损失函数的函数的特殊情形.通过损失函数研究估计量的性能与最优性，是决策论的一个分支.

观测到数据$\boldsymbol X=\boldsymbol x$之后，其中$\boldsymbol X\sim f(\boldsymbol x\mid\theta)$、$\theta\in\Theta$，要对$\theta$作出决策.允许决策的集合称为行动空间（action space），记为$\mathcal A$.\index{SI0751@参数空间}\semanticindexanchor{751:313}在点估计问题中，$\mathcal A$往往等于参数空间$\Theta$，但其他问题未必如此（例如假设检验，见\autoref{sec:8.3.5}）.

点估计问题中的损失函数反映了这样的事实：\index{SI0653@损失函数}\semanticindexanchor{653:313}若行动$a$接近$\theta$，决策$a$便合理，产生的损失很小；若$a$远离$\theta$，损失便很大.损失函数是非负函数，通常随着$a$与$\theta$之间距离增大而增大.若$\theta$为实数，两种常用的损失函数为
\begin{align*}
\text{绝对误差损失（absolute error loss）：}\quad L(\theta,a)&=|a-\theta|,\\
\text{平方误差损失（squared error loss）：}\quad L(\theta,a)&=(a-\theta)^2.
\end{align*}
两者均随着$\theta$与$a$之间距离增大而增大，最小值为$L(\theta,\theta)=0$.即行动正确时，损失最小.平方误差损失对较大偏离施加相对更多的惩罚，绝对误差损失则对较小偏离施加相对更多的惩罚.平方误差损失的一种变形，对高估的惩罚大于低估：
$$L(\theta,a)=\begin{cases}(a-\theta)^2,&\text{若 }a<\theta,\\10(a-\theta)^2,&\text{若 }a\ge\theta.\end{cases}$$
另一种相对平方误差损失，在$\theta$接近零时比$|\theta|$较大时更严厉地惩罚估计误差：
$$L(\theta,a)=\frac{(a-\theta)^2}{|\theta|+1}.$$
注意，后面这两种平方误差损失的变形，也可以改为基于绝对误差损失.一般而言，实验者必须考虑不同$\theta$值下各种估计错误的后果，并指定能反映这些后果的损失函数.

在损失函数或决策论分析中，用风险函数（risk function）量化估计量的质量.即对于$\theta$的估计量$\delta(\boldsymbol x)$，风险函数是$\theta$的函数：
\begin{equation}\label{eq:7.3.16}\index{SI0914@风险函数}
R(\theta,\delta)=E_\theta L(\theta,\delta(\boldsymbol X)).
\end{equation}
在给定$\theta$处，风险函数就是使用估计量$\delta(\boldsymbol x)$时会产生的平均损失.

由于$\theta$的真值未知，希望采用对所有$\theta$均有较小$R(\theta,\delta)$的估计量.这样，不论$\theta$的真值是什么，估计量的期望损失都小.要比较两个估计量$\delta_1$和$\delta_2$的优劣，就比较其风险函数$R(\theta,\delta_1)$和$R(\theta,\delta_2)$.若对所有$\theta\in\Theta$都有$R(\theta,\delta_1)<R(\theta,\delta_2)$，则偏好$\delta_1$，因为它在所有$\theta$处都更好.更常见的情形是两个风险函数相交，此时哪个估计量更好就未必那么明确.

估计量$\delta$的风险函数是\eqref{eq:7.3.16}定义的期望损失.在平方误差损失下，风险函数就是熟悉的量——\autoref{sec:7.3.1}中使用的均方误差MSE.那里定义$\operatorname{MSE}(\theta)=E_\theta(\delta(\boldsymbol X)-\theta)^2$；若$L(\theta,a)=(a-\theta)^2$，这正是$E_\theta L(\theta,\delta(\boldsymbol X))=R(\theta,\delta)$.如\eqref{eq:7.3.1}，平方误差损失下有
\begin{equation}\label{eq:7.3.17}
\begin{aligned}
R(\theta,\delta)&=\operatorname{Var}_\theta\delta(\boldsymbol X)+(E_\theta\delta(\boldsymbol X)-\theta)^2\\
&=\operatorname{Var}_\theta\delta(\boldsymbol X)+(\operatorname{Bias}_\theta\delta(\boldsymbol X))^2.
\end{aligned}
\end{equation}
\sourcepageanchor{314}
这个平方误差损失的风险函数清楚表明，好的估计量应当既有小方差，也有小偏差.决策论分析会评价估计量在同时使这两个量变小方面做得如何.

像\autoref{sec:7.3.2}那样，把允许估计量集合$\mathcal D$限制为无偏估计量集合，在决策论分析中并不典型.那样，最小化风险就只是最小化方差.决策论分析会更全面，因为方差和偏差都包含在风险中，需要同时考虑.若估计量有小但可能非零的偏差，同时有小方差，就会被认为是好的.
\begin{example}[二项分布的风险函数]\label{ex:7.3.25}
在\autoref{ex:7.3.5}中，考虑了来自$\operatorname{Bernoulli}(p)$总体的随机样本$X_1,\ldots,X_n$，以及两个估计量
$$\widehat p_B=\frac{\sum_{i=1}^nX_i+\sqrt{n/4}}{n+\sqrt n}\qquad\text{及}\qquad\bar X=\frac1n\sum_{i=1}^nX_i.$$
当$n=4$和$n=400$时，这两个估计量的风险函数绘于\autoref{fig:7.3.1}；比较结果如\autoref{ex:7.3.5}所述.根据风险比较，$n$较小时偏好$\widehat p_B$，$n$较大时偏好$\bar X$.
\end{example}
\begin{example}[正态方差估计的风险]\label{ex:7.3.26}\index{SI0778@方差}
设$X_1,\ldots,X_n$为来自$n(\mu,\sigma^2)$总体的随机样本.考虑在平方误差损失下估计$\sigma^2$.考察形如$\delta_b(\boldsymbol X)=bS^2$的估计量，其中$S^2$为样本方差，$b$为任意非负常数.回顾$ES^2=\sigma^2$，并且对正态样本，$\operatorname{Var}S^2=2\sigma^4/(n-1)$.利用\eqref{eq:7.3.17}，可计算$\delta_b$的风险函数：
\begin{align*}
R((\mu,\sigma^2),\delta_b)&=\operatorname{Var}bS^2+(EbS^2-\sigma^2)^2\\
&=b^2\operatorname{Var}S^2+(bES^2-\sigma^2)^2\\
&=\frac{b^22\sigma^4}{n-1}+(b-1)^2\sigma^4\qquad\text{（利用$\operatorname{Var}S^2$）}\\
&=\left[\frac{2b^2}{n-1}+(b-1)^2\right]\sigma^4.
\end{align*}
$\delta_b$的风险函数不依赖$\mu$，并且是$\sigma^2$的二次函数，形式为$c_b(\sigma^2)^2$，其中$c_b$是正常数.要比较两个风险函数，从而比较两个估计量的价值，只须注意：若$c_b<c_{b'}$，则对所有$(\mu,\sigma^2)$都有
$$R((\mu,\sigma^2),\delta_b)=c_b(\sigma^2)^2<c_{b'}(\sigma^2)^2=R((\mu,\sigma^2),\delta_{b'}).$$
因此$\delta_b$优于$\delta_{b'}$.使
\begin{equation}\label{eq:7.3.18}
c_b=\frac{2b^2}{n-1}+(b-1)^2
\end{equation}
达到全局最小值的$b$，给出此类中的最佳估计量$\delta_b$.用标准微积分方法可知，最小化值为$b=(n-1)/(n+1)$.所以，在每个$(\mu,\sigma^2)$处，估计量
$$\widetilde S^2=\frac{n-1}{n+1}S^2=\frac1{n+1}\sum(X_i-\bar X)^2$$
在所有$bS^2$形式的估计量中风险最小.当$n=5$时，此估计量与同类中另两个估计量的风险函数见\autoref{fig:7.3.2}.另两个估计量是无偏估计量$S^2$和$\sigma^2$的MLE $\widehat\sigma^2=\frac{n-1}nS^2$.显然，$\widetilde S^2$的风险函数处处最小.
\end{example}
\sourcepageanchor{315}
\begin{figure}[H]
\centering\includegraphics[width=.66\linewidth]{assets/chapter7/fig-7-3-2.pdf}
\caption{\autoref{ex:7.3.26}中三个方差估计量的风险函数}\label{fig:7.3.2}
\end{figure}
\begin{example}[Stein损失下的方差估计]\label{ex:7.3.27}\index{SI0658@似然}\index{SI0662@Stein损失}
再次考虑用形如$bS^2$的估计量估计总体方差$\sigma^2$.这次分析可以十分一般，只假设$X_1,\ldots,X_n$是来自某个具有正且有限方差$\sigma^2$的总体的随机样本.现在使用归于Stein的损失函数（\hyperref[bib:source176]{James和Stein 1961}；另见\hyperref[bib:source51]{Brown 1990a）}
$$L(\sigma^2,a)=\frac a{\sigma^2}-1-\log\frac a{\sigma^2}.$$
它比平方误差损失复杂，但有一些合理性质.注意，若$a=\sigma^2$，损失为零.此外，对于固定的$\sigma^2$，当$a\to0$或$a\to\infty$时，均有$L(\sigma^2,a)\to\infty$.即严重低估与严重高估受到同样严厉的惩罚.（在方差估计中，对平方误差损失的一项批评是，低估只受有限惩罚，高估则受无限惩罚.）若样本来自正态总体，这个损失函数也可从$\sigma^2$的似然函数导出，因此将良好的决策论性质与良好的似然性质联系起来（见\autoref{exercise:7.61}）.

对估计量$\delta_b=bS^2$，风险函数为
\begin{align*}
R(\sigma^2,\delta_b)&=E\left(\frac{bS^2}{\sigma^2}-1-\log\frac{bS^2}{\sigma^2}\right)\\
&=bE\frac{S^2}{\sigma^2}-1-E\log\frac{bS^2}{\sigma^2}\\
&=b-\log b-1-E\log\frac{S^2}{\sigma^2}.\qquad\left(E\frac{S^2}{\sigma^2}=1\right)
\end{align*}
量$E\log(S^2/\sigma^2)$可能是$\sigma^2$及其他总体参数的函数，但不是$b$的函数.因此，对所有$\sigma^2$，使$b-\log b$最小的$b$也使$R(\sigma^2,\delta_b)$最小，即$b=1$.所以，在$bS^2$形式的估计量中，对所有$\sigma^2$风险最小的是$\delta_1=S^2$.\footnote{译注：要作这里有限风险的比较，还需$S^2>0$几乎处处及相应对数期望有限；仅总体方差正且有限不足以保证这些条件.例如离散总体可能使$P(S^2=0)>0$，此时所有正$b$的Stein风险均为无穷，不能据该计算得到有意义的严格优劣比较.}
\end{example}
\sourcepageanchor{316}
也可用Bayes方法研究损失函数下的最优性，此时给定先验分布$\pi(\theta)$.在Bayes分析中，用该先验计算平均风险
$$\int_\Theta R(\theta,\delta)\pi(\theta)\,d\theta,$$
它称为Bayes风险（Bayes risk）.\index{SI0083@贝叶斯风险}\semanticindexanchor{83:316}对风险函数取平均，得到一个数，用来评估给定损失函数下估计量的性能.而且，可以尝试寻找使Bayes风险最小的估计量，称为关于先验$\pi$的Bayes规则（Bayes rule with respect to a prior $\pi$），\index{SI0085@先验分布}\semanticindexanchor{85:316}常记为$\delta^\pi$.

寻找给定先验$\pi$的Bayes决策规则看起来可能很困难，但实际过程相当机械化，正如下面的定理所指出的.（下面这种寻找Bayes规则的方法，在比此处更一般的情形中也适用，见\hyperref[bib:source54]{Brown和Purves 1973}.）

对于$\boldsymbol X\sim f(\boldsymbol x\mid\theta)$和$\theta\sim\pi$，决策规则$\delta$的Bayes风险可写为
$$\int_\Theta R(\theta,\delta)\pi(\theta)\,d\theta=\int_\Theta\left(\int_{\mathcal X}L(\theta,\delta(\boldsymbol x))f(\boldsymbol x\mid\theta)\,d\boldsymbol x\right)\pi(\theta)\,d\theta.$$
若写成$f(\boldsymbol x\mid\theta)\pi(\theta)=\pi(\theta\mid\boldsymbol x)m(\boldsymbol x)$，其中$\pi(\theta\mid\boldsymbol x)$是$\theta$的后验分布，$m(\boldsymbol x)$是$\boldsymbol X$的边缘分布，则Bayes风险可写为
\begin{equation}\label{eq:7.3.19}\index{SI0660@后验期望损失}\index{SI0820@后验期望损失（另见损失函数）}
\int_\Theta R(\theta,\delta)\pi(\theta)\,d\theta=\int_{\mathcal X}\left[\int_\Theta L(\theta,\delta(\boldsymbol x))\pi(\theta\mid\boldsymbol x)\,d\theta\right]m(\boldsymbol x)\,d\boldsymbol x.
\end{equation}
方括号内是损失函数关于后验分布的期望，称为后验期望损失（posterior expected \mbox{loss）.}它仅是$\boldsymbol x$的函数，不是$\theta$的函数.因此，对每个$\boldsymbol x$，若选行动$\delta(\boldsymbol x)$使后验期望损失最小，就使Bayes风险最小.

注意，现在有了构造Bayes规则的步骤：对于给定观测$\boldsymbol x$，Bayes规则应最小化后验期望损失.这与前面各小节的任何方法都很不一样.例如，考虑前面寻找最佳无偏估计量的方法.要用\autoref{thm:7.3.23}，首先须找到完备充分统计量$T$，然后找到一个函数$\phi(T)$，它是所求参数的无偏估计量.如果已经知道参数的某个无偏估计量，Rao--Blackwell定理（\autoref{thm:7.3.17}）可能有帮助；但如果想不出任何无偏估计量，该方法并不告诉我们怎样构造一个.

即使不能解析地最小化后验期望损失，也可以用数值方法计算积分并进行最小化.事实上，观测到$\boldsymbol X=\boldsymbol x$后，只需对这个特定的$\boldsymbol x$完成最小化.不过，在有些问题中，可以显式描述Bayes规则.
\needspace{7\baselineskip}
\begin{example}[两种Bayes规则]\label{ex:7.3.28}\index{SI0654@绝对误差}\index{SI0659@点估计}\index{SI0661@平方误差}\index{SI0683@均方误差}\index{SI0776@决策论}
考虑实参数$\theta$的点估计问题.
\begin{enumerate}[label=(\alph*)]
\item 对平方误差损失，后验期望损失为
$$\int_\Theta(\theta-a)^2\pi(\theta\mid\boldsymbol x)\,d\theta=E((\theta-a)^2\mid\boldsymbol X=\boldsymbol x).$$
这里$\theta$为分布是$\pi(\theta\mid\boldsymbol x)$的随机变量.由\autoref{ex:2.2.6}，$\delta^\pi(\boldsymbol x)=E(\theta\mid\boldsymbol x)$使该期望最小.所以Bayes规则是后验分布的均值.
\item 对绝对误差损失，后验期望损失为$E(|\theta-a|\mid\boldsymbol X=\boldsymbol x)$.应用\autoref{exercise:2.18}可知，取$\delta^\pi(\boldsymbol x)$为$\pi(\theta\mid\boldsymbol x)$的中位数，即使它最小.
\end{enumerate}
\end{example}
\sourcepageanchor{317}
\begin{table}[H]
\centering
\caption{二项分布参数$p$的三个估计量}\label{tab:7.3.1}
\begin{tabular}{rrrr}
\hline
\multicolumn{4}{c}{$n=10$，先验$\pi(p)\sim\operatorname{uniform}(0,1)$}\\
$y$&MLE&\shortstack{Bayes\\绝对误差}&\shortstack{Bayes\\平方误差}\\\hline
0&.0000&.0611&.0833\\
1&.1000&.1480&.1667\\
2&.2000&.2358&.2500\\
3&.3000&.3238&.3333\\
4&.4000&.4119&.4167\\
5&.5000&.5000&.5000\\
6&.6000&.5881&.5833\\
7&.7000&.6762&.6667\\
8&.8000&.7642&.7500\\
9&.9000&.8520&.8333\\
10&1.0000&.9389&.9137\footnotemark\\\hline
\end{tabular}
\end{table}
\footnotetext{译注：原表最后一项为$.9137$，按$(y+1)/(n+2)=11/12$计算，应为$.9167$.原数值保留.}
在\autoref{sec:7.2.3}中，所讨论的Bayes估计量为后验均值$\delta^\pi(\boldsymbol x)=E(\theta\mid\boldsymbol x)$.现在知道，它是平方误差损失下的Bayes估计量.若认为另一种损失函数比平方误差损失更合适，Bayes估计量可能是不同的统计量.
\begin{example}[正态分布的Bayes估计]\label{ex:7.3.29}
设$X_1,\ldots,X_n$是来自$n(\theta,\sigma^2)$总体的随机样本，先验$\pi(\theta)$为$n(\mu,\tau^2)$.$\sigma^2$、$\mu$和$\tau^2$都已知.在\autoref{ex:7.2.16}及其在\autoref{exercise:7.22}中的推广里，已得到给定$\bar X=\bar x$时$\theta$的后验分布为正态，且
$$E(\theta\mid\bar x)=\frac{\tau^2}{\tau^2+(\sigma^2/n)}\bar x+\frac{\sigma^2/n}{\tau^2+(\sigma^2/n)}\mu,$$
以及
$$\operatorname{Var}(\theta\mid\bar x)=\frac{\tau^2\sigma^2/n}{\tau^2+(\sigma^2/n)}.$$
在平方误差损失下，Bayes估计量为$\delta^\pi(\boldsymbol x)=E(\theta\mid\bar x)$.由于后验分布为正态，它关于均值对称，$\pi(\theta\mid\boldsymbol x)$的中位数等于$E(\theta\mid\boldsymbol x)$.因此，在绝对误差损失下，Bayes估计量也为$\delta^\pi(\boldsymbol x)=E(\theta\mid\bar x)$.
\end{example}
\begin{example}[二项分布的Bayes估计]\label{ex:7.3.30}\index{SI0777@二项}
设$X_1,\ldots,X_n$独立同分布于$\operatorname{Bernoulli}(p)$，令$Y=\sum X_i$.假设$p$的先验为$\operatorname{beta}(\alpha,\beta)$.在\autoref{ex:7.2.14}中，已发现后验分布只通过$Y=y$的观测值依赖样本，且为$\operatorname{beta}(y+\alpha,n-y+\beta)$.因此，$\delta^\pi(y)=E(p\mid y)=(y+\alpha)/(\alpha+\beta+n)$是平方误差损失下$p$的Bayes估计量.

对于绝对误差损失，需要寻找$\pi(p\mid y)=\operatorname{beta}(y+\alpha,n-y+\beta)$的中位数.一般没有简单表达式.中位数隐式定义为满足
$$\int_0^m\frac{\Gamma(\alpha+\beta+n)}{\Gamma(y+\alpha)\Gamma(n-y+\beta)}p^{y+\alpha-1}(1-p)^{n-y+\beta-1}\,dp=\frac12$$
的数$m$.可以数值计算该积分，近似找出满足等式的$m$.对于$n=10$、$\alpha=\beta=1$，即$\operatorname{uniform}(0,1)$先验，我们已作了这种计算.绝对误差损失的Bayes估计量列于\autoref{tab:7.3.1}；表中还列出前面导出的平方误差损失的Bayes估计量，以及MLE $\widehat p=y/n$.

注意\autoref{tab:7.3.1}：与MLE不同，即使$y=0$或$n$，两个Bayes估计量都不会把$p$估为0或1.Bayes估计量通常不会取参数空间中的
\sourcepageanchor{318}
极端值.不论样本量多大，先验始终对估计量有一定影响，倾向于把它从极端值拉开.从上面$E(p\mid y)$的表达式可见，即使$y=0$且$n$很大，Bayes估计量仍是正数.
\end{example}



\subsection{习题}\label{sec:7.4}
\begin{exercise}\label{exercise:7.1}
对离散随机变量$X$作一次观测，其pmf为$f(x\mid\theta)$，其中$\theta\in\{1,2,3\}$.求$\theta$的MLE.
\begin{center}\begin{tabular}{cccc}
$x$&$f(x\mid1)$&$f(x\mid2)$&$f(x\mid3)$\\\hline
0&$1/3$&$1/4$&0\\
1&$1/3$&$1/4$&0\\
2&0&$1/4$&$1/4$\\
3&$1/6$&$1/4$&$1/2$\\
4&$1/6$&0&$1/4$
\end{tabular}\end{center}
\end{exercise}
\begin{exercise}\label{exercise:7.2}
设$X_1,\ldots,X_n$为来自$\operatorname{gamma}(\alpha,\beta)$总体的随机样本.
\begin{enumerate}[label=(\alph*)]
\item\phantomsection\label{itm:exercise:7.2:a} 假设$\alpha$已知，求$\beta$的MLE.
\item 若$\alpha$与$\beta$均未知，其MLE没有显式公式，但可通过数值方法寻找最大值.利用\hyperref[itm:exercise:7.2:a]{(a)}的结果，可以把问题化为一元函数的极大化.对\hyperref[itm:exercise:7.10:c]{习题7.10(c)}的数据，求$\alpha$和$\beta$的MLE.
\end{enumerate}
\end{exercise}
\begin{exercise}\label{exercise:7.3}\index{SI0621@对数似然}
给定来自pdf为$f(x\mid\theta)$的总体的随机样本$X_1,\ldots,X_n$，说明把似然函数$L(\theta\mid\boldsymbol x)$作为$\theta$的函数极大化，等价于极大化$\log L(\theta\mid\boldsymbol x)$.
\end{exercise}
\begin{exercise}\label{exercise:7.4}
证明\autoref{ex:7.2.8}中的断言.即证明，当$\theta$的范围限制为正半轴时，该例给出的$\widehat\theta$为MLE.
\end{exercise}
\begin{exercise}\label{exercise:7.5}
考虑如\autoref{ex:7.2.9}中那样估计二项分布参数$k$.
\begin{enumerate}[label=(\alph*)]
\item 证明如下断言：满足那些不等式且为MLE的整数$\widehat k$，是不超过$1/\widehat z$的最大整数.
\item 令$p=1/2$、$n=4$，且$X_1=0$、$X_2=20$、$X_3=1$、$X_4=19$.$\widehat k$是多少？
\end{enumerate}
\end{exercise}
\begin{exercise}\label{exercise:7.6}
设$X_1,\ldots,X_n$为来自下列pdf的随机样本：
$$f(x\mid\theta)=\theta x^{-2},\qquad0<\theta\le x<\infty.$$
\begin{enumerate}[label=(\alph*)]
\item $\theta$的充分统计量是什么？
\item 求$\theta$的MLE.
\item 求$\theta$的矩估计量.
\end{enumerate}
\end{exercise}
\begin{exercise}\label{exercise:7.7}
设$X_1,\ldots,X_n$独立同分布于两个pdf之一.若$\theta=0$，则
$$f(x\mid\theta)=\begin{cases}1,&\text{若 }0<x<1,\\0,&\text{其他情形};\end{cases}$$
而若$\theta=1$，则
$$f(x\mid\theta)=\begin{cases}1/(2\sqrt x),&\text{若 }0<x<1,\\0,&\text{其他情形}.\end{cases}$$
求$\theta$的MLE.
\end{exercise}
\begin{exercise}\label{exercise:7.8}
从$n(0,\sigma^2)$总体中取得一个观测$X$.
\begin{enumerate}[label=(\alph*)]
\item 求$\sigma^2$的一个无偏估计量.
\item 求$\sigma$的MLE.
\item 讨论如何求$\sigma$的矩估计量.
\end{enumerate}
\end{exercise}
\sourcepageanchor{319}
\begin{exercise}\label{exercise:7.9}
设$X_1,\ldots,X_n$独立同分布，pdf为
$$f(x\mid\theta)=\frac1\theta,\qquad0\le x\le\theta,\quad\theta>0.$$
分别用矩估计法和极大似然法估计$\theta$.计算两个估计量的均值和方差.应选择哪一个，为什么？
\end{exercise}
\begin{exercise}\label{exercise:7.10}
独立随机变量$X_1,\ldots,X_n$具有共同分布
$$P(X_i\le x\mid\alpha,\beta)=\begin{cases}0,&\text{若 }x<0,\\(x/\beta)^\alpha,&\text{若 }0\le x\le\beta,\\1,&\text{若 }x>\beta,\end{cases}$$
其中参数$\alpha$与$\beta$均为正.
\begin{enumerate}[label=(\alph*)]
\item 求$(\alpha,\beta)$的一个二维充分统计量.
\item 求$\alpha$和$\beta$的MLE.
\item\phantomsection\label{itm:exercise:7.10:c} 在篱雀巢中发现的杜鹃卵的长度（毫米）可用此分布建模.对数据
\begin{gather*}
22.0,\ 23.9,\ 20.9,\ 23.8,\ 25.0,\ 24.0,\ 21.7,\\
23.8,\ 22.8,\ 23.1,\ 23.1,\ 23.5,\ 23.0,\ 23.0,
\end{gather*}
求$\alpha$和$\beta$的极大似然估计值.
\end{enumerate}
\end{exercise}
\begin{exercise}\label{exercise:7.11}
设$X_1,\ldots,X_n$独立同分布，pdf为
$$f(x\mid\theta)=\theta x^{\theta-1},\qquad0\le x\le1,\quad0<\theta<\infty.$$
\begin{enumerate}[label=(\alph*)]
\item 求$\theta$的MLE，并说明当$n\to\infty$时，其方差趋于零.
\item 求$\theta$的矩估计量.
\end{enumerate}
\end{exercise}
\begin{exercise}\label{exercise:7.12}
设$X_1,\ldots,X_n$为来自具有下列pmf的总体的随机样本：
$$P_\theta(X=x)=\theta^x(1-\theta)^{1-x},\qquad x=0\text{ 或 }1,\quad0\le\theta\le\frac12.$$
\begin{enumerate}[label=(\alph*)]
\item 求$\theta$的矩估计量和MLE.
\item 求两个估计量各自的均方误差.
\item 更倾向于哪个估计量？说明选择的理由.
\end{enumerate}
\end{exercise}
\begin{exercise}\label{exercise:7.13}
设$X_1,\ldots,X_n$为来自双指数pdf总体的样本：
$$f(x\mid\theta)=\frac12e^{-|x-\theta|},\qquad-\infty<x<\infty,\quad-\infty<\theta<\infty.$$
求$\theta$的MLE.［提示：分别考虑$n$为偶数和奇数的情形，用次序统计量表示MLE.\hyperref[bib:source235]{Norton（1984）}给出了此问题的完整处理.］
\end{exercise}
\begin{exercise}\label{exercise:7.14}
设$X$和$Y$为独立指数随机变量，且
$$f(x\mid\lambda)=\frac1\lambda e^{-x/\lambda},\ x>0,\qquad f(y\mid\mu)=\frac1\mu e^{-y/\mu},\ y>0.$$
观测到$Z$和$W$，其中
$$Z=\min(X,Y)\qquad\text{及}\qquad W=\begin{cases}1,&\text{若 }Z=X,\\0,&\text{若 }Z=Y.\end{cases}$$
\autoref{exercise:4.26}已求出$Z$和$W$的联合分布.现假设$(Z_i,W_i)$，$i=1,\ldots,n$，为$n$个独立同分布观测.求$\lambda$和$\mu$的MLE.
\end{exercise}
\sourcepageanchor{320}
\begin{exercise}\label{exercise:7.15}\index{SI0324@逆高斯}
设$X_1,X_2,\ldots,X_n$为来自逆高斯pdf的样本：
$$f(x\mid\mu,\lambda)=\left(\frac\lambda{2\pi x^3}\right)^{1/2}\exp\{-\lambda(x-\mu)^2/(2\mu^2x)\},\qquad x>0.$$
\begin{enumerate}[label=(\alph*)]
\item 说明$\mu$与$\lambda$的极大似然估计量为\footnote{译注：原式分母中$1/\bar X$前漏了$n$；对数似然求导给出$\widehat\lambda_n=n/(\sum_i1/X_i-n/\bar X)$.与\autoref{exercise:6.27}为同一处公式疑误，原式保留.}
$$\widehat\mu_n=\bar X\qquad\text{及}\qquad\widehat\lambda_n=\frac n{\sum_i\frac1{X_i}-\frac1{\bar X}}.$$
\item \hyperref[bib:source319]{Tweedie（1957）}说明，$\widehat\mu_n$和$\widehat\lambda_n$独立，$\widehat\mu_n$服从参数为$\mu$和$n\lambda$的逆高斯分布，$n\lambda/\widehat\lambda_n$服从$\chi_{n-1}^2$分布.\hyperref[bib:source281]{Schwarz和Samantha（1991）}用归纳论证证明了这些事实.\footnote{译注：此处原文写 Samantha，书末对应文献作者为 Samanta；保留原拼写，链接指向该条文献.}
\begin{enumerate}[label=(\roman*)]
\item 说明$\widehat\mu_2$服从参数为$\mu$与$2\lambda$的逆高斯分布，$2\lambda/\widehat\lambda_2$服从$\chi_1^2$分布，且二者独立.
\item 假设结果在$n=k$时成立，并得到一个新的独立观测$x$.建立\hyperref[bib:source281]{Schwarz和Samantha（1991）}所用的归纳步骤，把pdf $f(x,\widehat\mu_k,\widehat\lambda_k)$变换为$f(x,\widehat\mu_{k+1},\widehat\lambda_{k+1})$.说明该密度以适当方式分解，从而得到Tweedie的结果.
\end{enumerate}
\end{enumerate}
\end{exercise}
\begin{exercise}\label{exercise:7.16}\index{SI0452@广义均值}\index{SI0786@广义均值}\index{SI0811@幂平均数}\index{SI0821@幂平均数}\index{SI1037@Box–Cox}
\hyperref[bib:source26]{Berger和Casella（1992）}还研究了\autoref{exercise:4.57}中见过的幂平均.回顾其定义为$[(1/n)\sum_{i=1}^nx_i^r]^{1/r}$.注意幂函数$x^r$可以替换为任意连续单调函数$h$，于是可进一步推广，得到广义平均$h^{-1}((1/n)\sum_{i=1}^nh(x_i))$.
\begin{enumerate}[label=(\alph*)]
\item 最小二乘问题$\min_a\sum_i(x_i-a)^2$有时通过变换变量求解，即求解$\min_a\sum_i[h(x_i)-h(a)]^2$.说明后一个问题的解为$a=h^{-1}((1/n)\sum_i h(x_i))$.
\item 说明算术平均是未变换的最小二乘问题的解；几何平均是经$h(x)=\log x$变换后问题的解；调和平均是经$h(x)=1/x$变换后问题的解.
\item 说明若用Box--Cox变换（见\autoref{exercise:11.3}）变换最小二乘问题，则解是$h(x)=x^\lambda$所对应的广义平均.
\item 设$X_1,X_2,\ldots,X_n$为来自$\operatorname{lognormal}(\mu,\sigma^2)$总体的样本.说明$\mu$的MLE为几何平均.\footnote{译注：按本书通常的对数正态参数约定，$\log X\sim n(\mu,\sigma^2)$，则$\mu$的MLE是$(1/n)\sum\log X_i$，几何平均是$e^\mu$的MLE.原文把二者混同，原题保留.}
\item 假设$X_1,X_2,\ldots,X_n$为来自单参数指数族$f(x\mid\theta)=\exp\{\theta h(x)-H(\theta)\}g(x)$的样本，其中$h=H'$，且$h$递增.
\begin{enumerate}[label=(\roman*)]
\item 说明$\theta$的极大似然估计量为$\widehat\theta=h^{-1}((1/n)\sum_i h(x_i))$.
\item 说明满足$h=H'$的两个密度是正态密度，以及pdf为$f(x\mid\theta)=\theta x^{-2}\exp\{-\theta/x\}$（$x>0$）的逆gamma密度；对于正态分布，MLE是算术平均，对于逆gamma分布，MLE是调和平均.
\end{enumerate}
\end{enumerate}
\end{exercise}
\begin{exercise}\label{exercise:7.17}
Borel悖论（\autoref{sec:4.9.3}）也可能出现在推断问题中.设$X_1$和$X_2$为独立同分布的$\operatorname{exponential}(\theta)$随机变量.
\begin{enumerate}[label=(\alph*)]
\item\phantomsection\label{itm:exercise:7.17:a} 若只观测$X_2$，说明$\theta$的MLE为$\widehat\theta=X_2$.
\item\phantomsection\label{itm:exercise:7.17:b} 假设改为只观测$Z=(X_2-1)/X_1$.求$(X_1,Z)$的联合分布，并将$X_1$积分消去，得到似然函数.
\item 假设$X_2=1$.比较\hyperref[itm:exercise:7.17:a]{(a)}和\hyperref[itm:exercise:7.17:b]{(b)}中$\theta$的MLE.
\item Bayes分析也不能免受Borel悖论影响.若$\pi(\theta)$是$\theta$的先验密度，说明在$X_2=1$时，\hyperref[itm:exercise:7.17:a]{(a)}与\hyperref[itm:exercise:7.17:b]{(b)}中的后验分布不同.
\end{enumerate}
（由俄亥俄州立大学L. Mark Berliner提供.）
\end{exercise}
\sourcepageanchor{321}
\begin{exercise}\label{exercise:7.18}\index{SI0672@二元正态}
设$(X_1,Y_1),\ldots,(X_n,Y_n)$为独立同分布的二元正态随机变量对，全部五个参数均未知.
\begin{enumerate}[label=(\alph*)]
\item 说明$\mu_X,\mu_Y,\sigma_X^2,\sigma_Y^2,\rho$的矩估计量为
\begin{gather*}
\widetilde\mu_X=\bar x,\quad\widetilde\mu_Y=\bar y,\quad\widetilde\sigma_X^2=\frac1n\sum(x_i-\bar x)^2,\quad\widetilde\sigma_Y^2=\frac1n\sum(y_i-\bar y)^2,\\
\widetilde\rho=\frac{(1/n)\sum(x_i-\bar x)(y_i-\bar y)}{\widehat\sigma_X\widehat\sigma_Y}.
\end{gather*}
\item 推导未知参数的MLE，并说明它们与矩估计量相同.（一种办法是把联合pdf写成条件密度与边缘密度的乘积，即
\begin{align*}
f(x,y\mid\mu_X,\mu_Y,\sigma_X^2,\sigma_Y^2,\rho)
&=f(y\mid x,\mu_X,\mu_Y,\sigma_X^2,\sigma_Y^2,\rho)\\
&\quad\times f(x\mid\mu_X,\sigma_X^2),
\end{align*}
并论证$\mu_X$和$\sigma_X^2$的MLE为$\bar x$与$(1/n)\sum(x_i-\bar x)^2$.然后调换角色，求$\mu_Y$和$\sigma_Y^2$的MLE.最后，处理“部分极大化”的似然函数$L(\bar x,\bar y,\widehat\sigma_X^2,\widehat\sigma_Y^2,\rho\mid\boldsymbol x,\boldsymbol y)$，求出$\rho$的MLE.可以猜到，这是一道难题.）
\end{enumerate}
\end{exercise}
\begin{exercise}\label{exercise:7.19}\index{SI0911@零截距}
假设随机变量$Y_1,\ldots,Y_n$满足
$$Y_i=\beta x_i+\varepsilon_i,\qquad i=1,\ldots,n,$$
其中$x_1,\ldots,x_n$为固定常数，$\varepsilon_1,\ldots,\varepsilon_n$独立同分布于$n(0,\sigma^2)$，$\sigma^2$未知.
\begin{enumerate}[label=(\alph*)]
\item 求$(\beta,\sigma^2)$的一个二维充分统计量.
\item 求$\beta$的MLE，并说明它是$\beta$的无偏估计量.
\item 求$\beta$的MLE的分布.
\end{enumerate}
\end{exercise}
\begin{exercise}\label{exercise:7.20}
设$Y_1,\ldots,Y_n$如\autoref{exercise:7.19}所定义.
\begin{enumerate}[label=(\alph*)]
\item 说明$\sum Y_i/\sum x_i$是$\beta$的无偏估计量.
\item 计算$\sum Y_i/\sum x_i$的精确方差，并与MLE的方差比较.
\end{enumerate}
\end{exercise}
\begin{exercise}\label{exercise:7.21}
仍设$Y_1,\ldots,Y_n$如\autoref{exercise:7.19}所定义.
\begin{enumerate}[label=(\alph*)]
\item 说明$[\sum(Y_i/x_i)]/n$也是$\beta$的无偏估计量.
\item 计算$[\sum(Y_i/x_i)]/n$的精确方差，并与前两题中估计量的方差比较.
\end{enumerate}
\end{exercise}
\begin{exercise}\label{exercise:7.22}\index{SI0080@正态}
本题将证明\autoref{ex:7.2.16}的断言，并作进一步推广.设$X_1,\ldots,X_n$为来自$n(\theta,\sigma^2)$总体的随机样本，$\theta$的先验分布为$n(\mu,\tau^2)$.这里假设$\sigma^2$、$\mu$和$\tau^2$都已知.
\begin{enumerate}[label=(\alph*)]
\item 求$\bar X$与$\theta$的联合pdf.
\item 说明$\bar X$的边缘分布$m(\bar x\mid\sigma^2,\mu,\tau^2)$为$n(\mu,(\sigma^2/n)+\tau^2)$.
\item 说明$\theta$的后验分布$\pi(\theta\mid\bar x,\sigma^2,\mu,\tau^2)$为正态，均值和方差由\eqref{eq:7.2.10}给出.\footnote{译注：对样本均值$\bar X$应用\eqref{eq:7.2.10}时，应将该式的$x$换为$\bar x$，$\sigma^2$换为$\sigma^2/n$；原题未明说此替换.}
\end{enumerate}
\end{exercise}
\begin{exercise}\label{exercise:7.23}\index{SI0081@正态方差}\index{SI0826@伽马}
若$S^2$是正态总体中大小为$n$的样本的样本方差，已知$(n-1)S^2/\sigma^2$服从$\chi_{n-1}^2$分布.$\sigma^2$的共轭先验是逆gamma pdf，记为$\operatorname{IG}(\alpha,\beta)$：
$$\pi(\sigma^2)=\frac1{\Gamma(\alpha)\beta^\alpha}\frac1{(\sigma^2)^{\alpha+1}}e^{-1/(\beta\sigma^2)},\qquad0<\sigma^2<\infty,$$
其中$\alpha$和$\beta$为正常数.说明$\sigma^2$的后验分布为
$$\operatorname{IG}\left(\alpha+\frac{n-1}2,\left[\frac{(n-1)S^2}2+\frac1\beta\right]^{-1}\right).$$
求此分布的均值，即$\sigma^2$的Bayes估计量.
\end{exercise}
\begin{exercise}\label{exercise:7.24}\index{SI0827@泊松}
设$X_1,\ldots,X_n$独立同分布于$\operatorname{Poisson}(\lambda)$，令$\lambda$服从Poisson分布的共轭族$\operatorname{gamma}(\alpha,\beta)$分布.
\begin{enumerate}[label=(\alph*)]
\item 求$\lambda$的后验分布.
\item 计算后验均值和方差.
\end{enumerate}
\end{exercise}
\begin{exercise}\label{exercise:7.25}\index{SI0782@正态}
考察\autoref{ex:7.2.16}和\autoref{exercise:7.22}中层次（Bayes）模型的一个推广.假设观测到$X_1,\ldots,X_n$，其中
\begin{align*}
X_i\mid\theta_i&\sim n(\theta_i,\sigma^2),\qquad i=1,\ldots,n,\quad\text{相互独立},\\
\theta_i&\sim n(\mu,\tau^2),\qquad i=1,\ldots,n,\quad\text{相互独立}.
\end{align*}
\sourcepageanchor{322}
\begin{enumerate}[label=(\alph*)]
\item 说明$X_i$的边缘分布为$n(\mu,\sigma^2+\tau^2)$，且边缘意义下$X_1,\ldots,X_n$独立同分布.（经验Bayes分析会用$X_i$的边缘分布估计先验参数$\mu$和$\tau^2$，参见\autoref{sec:7.5.6}.）
\item 一般地，说明若
\begin{align*}
X_i\mid\theta_i&\sim f(x\mid\theta_i),\qquad i=1,\ldots,n,\quad\text{相互独立},\\
\theta_i&\sim\pi(\theta\mid\tau),\qquad i=1,\ldots,n,\quad\text{相互独立},
\end{align*}
则边缘意义下$X_1,\ldots,X_n$独立同分布.
\end{enumerate}
\end{exercise}
\begin{exercise}\label{exercise:7.26}
在\autoref{ex:7.2.16}中，看到正态分布是自身的共轭族.不过，有时共轭先验不能准确反映先验知识，需要寻找不同的先验.设$X_1,\ldots,X_n$独立同分布于$n(\theta,\sigma^2)$，$\theta$服从双指数分布，即$\pi(\theta)=e^{-|\theta|/a}/(2a)$，$a$已知.求$\theta$的后验分布的均值.
\end{exercise}
\needspace{8\baselineskip}
\begin{exercise}\label{exercise:7.27}
参见\hyperref[ex:7.2.17]{习题7.2.17}.\footnote{译注：原文“习题7.2.17”应指\autoref{ex:7.2.17}；下一题的同一引用也如此.保留原显示称谓，链接指向该例.}
\begin{enumerate}[label=(\alph*)]
\item 说明完全数据似然\eqref{eq:7.2.11}所给的似然估计量为\eqref{eq:7.2.12}.
\item 说明\eqref{eq:7.2.23}中EM序列的极限满足\eqref{eq:7.2.16}.
\item 原始（不完全数据）似然方程可以直接求解.说明\eqref{eq:7.2.16}的解为
$$\widehat\beta=\frac{\sum_{i=2}^ny_i}{\sum_{i=2}^nx_i},\qquad\widehat\tau_1=\frac{y_1}{\widehat\beta},\qquad\widehat\tau_j=\frac{x_j+y_j}{\widehat\beta+1},\quad j=2,3,\ldots,n,$$
并且它就是\eqref{eq:7.2.23}中EM序列的极限.
\end{enumerate}
\end{exercise}
\begin{exercise}\label{exercise:7.28}
把\hyperref[ex:7.2.17]{习题7.2.17}的模型用于\autoref{tab:7.4.2}的数据，数据改编自\hyperref[bib:source193]{Lange等（1994）}.这些是纽约州若干地区的白血病病例数及相应人口数.
\begin{enumerate}[label=(\alph*)]
\item\phantomsection\label{itm:exercise:7.28:a} 对\autoref{tab:7.4.2}的完整数据，以及假设首个人口数$x_1=3540$缺失的“不完全”数据，分别拟合Poisson模型.
\end{enumerate}
\setcounter{table}{1}
\begin{table}[H]
\centering\small
\caption{白血病病例计数}\label{tab:7.4.2}
\setlength{\tabcolsep}{3.2pt}
\begin{tabular}{lrrrrrrrrr}\hline
人口数&3540&3560&3739&2784&2571&2729&3952&993&1908\\
病例数&3&4&1&1&3&1&2&0&2\\\hline
人口数&948&1172&1047&3138&5485&5554&2943&4969&4828\\
病例数&0&1&3&5&4&6&2&5&4\\\hline
\end{tabular}
\end{table}
\begin{enumerate}[label=(\alph*),start=2]
\item 假设缺失的不是$x$值，而是一项白血病计数（设$y_1=3$缺失）.用EM算法求此时的MLE，并与\hyperref[itm:exercise:7.28:a]{(a)}的答案比较.
\end{enumerate}
\end{exercise}
\begin{exercise}\label{exercise:7.29}
\autoref{ex:7.2.17}模型的一个替代如下：观测$(Y_i,X_i)$，$i=1,2,\ldots,n$，其中$Y_i\sim\operatorname{Poisson}(m\beta\tau_i)$，而$(X_1,X_2,\ldots,X_n)\sim\operatorname{multinomial}(m;\boldsymbol\tau)$，$\boldsymbol\tau=(\tau_1,\tau_2,\ldots,\tau_n)$且$\sum_{i=1}^n\tau_i=1$.例如，这里假设人口计数是多项分配，而非Poisson计数.（把$m=\sum x_i$视为已知.）
\begin{enumerate}[label=(\alph*)]
\item\phantomsection\label{itm:exercise:7.29:a} 说明$\boldsymbol Y=(Y_1,Y_2,\ldots,Y_n)$与$\boldsymbol X=(X_1,X_2,\ldots,X_n)$的联合密度为
$$f(\boldsymbol y,\boldsymbol x\mid\beta,\boldsymbol\tau)=\prod_{i=1}^n\frac{e^{-m\beta\tau_i}(m\beta\tau_i)^{y_i}}{y_i!}\,m!\frac{\tau_i^{x_i}}{x_i!}.$$
\item\phantomsection\label{itm:exercise:7.29:b} 若观测到完全数据，说明MLE为
$$\widehat\beta=\frac{\sum_{i=1}^ny_i}{\sum_{i=1}^nx_i}\qquad\text{及}\qquad\widehat\tau_j=\frac{x_j+y_j}{\sum_{i=1}^nx_i+y_i},\quad j=1,2,\ldots,n.$$
\sourcepageanchor{323}
\item\phantomsection\label{itm:exercise:7.29:c} 假设$x_1$缺失.利用$X_1\sim\operatorname{binomial}(m,t_1)$这一事实，计算期望完全数据对数似然.说明EM序列为\footnote{译注：本题原文有几处需辨明：\hyperref[itm:exercise:7.29:a]{(a)}中的多项系数$m!$只能出现一次，\hyperref[itm:exercise:7.29:b]{(b)}分母应理解为$\sum_i(x_i+y_i)$.在\hyperref[itm:exercise:7.29:c]{(c)}中，若$m$确已知且其余$x_i$全部观测到，则$x_1=m-\sum_{i=2}^nx_i$已经确定，不能用边缘分布的均值$m\tau_1^{(r)}$代替条件均值；所给EM迭代一般不适用于原述观测模型.此外，原文$t_1$应为$\tau_1$，且$j=1$的更新分子仍含缺失的$x_1$.原题保留.}
\begin{align*}
\widehat\beta^{(r+1)}&=\frac{\sum_{i=1}^ny_i}{m\widehat\tau_1^{(r)}+\sum_{i=2}^nx_i},\\
\widehat\tau_j^{(r+1)}&=\frac{x_j+y_j}{m\widehat\tau_1^{(r)}+\sum_{i=2}^nx_i+\sum_{i=1}^ny_i},\qquad j=1,2,\ldots,n.
\end{align*}
\item 用此模型求\autoref{tab:7.4.2}数据的MLE，先假设全部数据可用，再假设$x_1=3540$缺失.
\end{enumerate}
\end{exercise}
\begin{exercise}\label{exercise:7.30}
EM算法适用于多种情形，“缺失数据”的定义也可以扩展，以容纳许多不同模型.假设混合密度为$pf(x)+(1-p)g(x)$，其中$p$未知.若观测$\boldsymbol X=(X_1,X_2,\ldots,X_n)$，则样本密度为
$$\prod_{i=1}^n[pf(x_i)+(1-p)g(x_i)],$$
这可能难以处理.（其实，两种分布的混合还不算太糟；但想想当$k$很大时，混合$\sum_{i=1}^kp_if_i(x)$的似然会是什么样.）

EM解法是用$\boldsymbol Z=(Z_1,Z_2,\ldots,Z_n)$增广已观测（不完全）数据，其中$Z_i$说明$X_i$来自哪个混合成分.即
$$X_i\mid z_i=1\sim f(x_i)\qquad\text{且}\qquad X_i\mid z_i=0\sim g(x_i),$$
并且$P(Z_i=1)=p$.
\begin{enumerate}[label=(\alph*)]
\item 说明$(\boldsymbol X,\boldsymbol Z)$的联合密度为$\prod_{i=1}^n[pf(x_i)^{z_i}][(1-p)g(x_i)^{1-z_i}]$.\footnote{译注：原式把$p$和$1-p$漏在幂之外.正确联合密度应为$\prod_i[pf(x_i)]^{z_i}[(1-p)g(x_i)]^{1-z_i}$；原式保留.}
\item 说明缺失数据分布，即$Z_i\mid x_i,p$的分布，是成功概率为$pf(x_i)/(pf(x_i)+(1-p)g(x_i))$的Bernoulli分布.
\item 计算期望完全数据对数似然，并说明EM序列为
$$\widehat p^{(r+1)}=\frac1n\sum_{i=1}^n\frac{\widehat p^{(r)}f(x_i)}{\widehat p^{(r)}f(x_i)+(1-\widehat p^{(r)})g(x_i)}.$$
\end{enumerate}
\end{exercise}
\needspace{9\baselineskip}
\begin{exercise}\label{exercise:7.31}
证明\autoref{thm:7.2.20}.
\begin{enumerate}[label=(\alph*)]
\item\phantomsection\label{itm:exercise:7.31:a} 说明，利用\eqref{eq:7.2.19}，可以写成
\begin{align*}
\log L(\widehat\theta^{(r)}\mid\boldsymbol y)
&=E[\log L(\widehat\theta^{(r)}\mid\boldsymbol y,\boldsymbol X)\mid\widehat\theta^{(r)},\boldsymbol y]\\
&\quad-E[\log k(\boldsymbol X\mid\widehat\theta^{(r)},\boldsymbol y)\mid\widehat\theta^{(r)},\boldsymbol y],
\end{align*}
而由于$\widehat\theta^{(r+1)}$是极大值点，有\footnote{译注：原文下面不等式左端漏了期望算子$E$，应与右端一样在方括号外加$E$.关于定理中等号的充分条件，另见\autoref{thm:7.2.20}处的译注.}
\begin{align*}
&[\log L(\widehat\theta^{(r+1)}\mid\boldsymbol y,\boldsymbol X)\mid\widehat\theta^{(r)},\boldsymbol y]\\
&\qquad\ge E[\log L(\widehat\theta^{(r)}\mid\boldsymbol y,\boldsymbol X)\mid\widehat\theta^{(r)},\boldsymbol y].
\end{align*}
何时不等式取等号？
\item 用Jensen不等式说明
\begin{align*}
&E[\log k(\boldsymbol X\mid\widehat\theta^{(r+1)},\boldsymbol y)\mid\widehat\theta^{(r)},\boldsymbol y]\\
&\qquad\le E[\log k(\boldsymbol X\mid\widehat\theta^{(r)},\boldsymbol y)\mid\widehat\theta^{(r)},\boldsymbol y],
\end{align*}
将其与\hyperref[itm:exercise:7.31:a]{(a)}结合即可证明定理.
\end{enumerate}
提示：若$f$和$g$是密度，因$\log$是凹函数，Jensen不等式\eqref{eq:4.7.7}意味着
\begin{align*}
\int\log\left(\frac{f(x)}{g(x)}\right)g(x)\,dx
&\le\log\left(\int\frac{f(x)}{g(x)}g(x)\,dx\right)\\
&=\log\int f(x)\,dx=0.
\end{align*}
再由对数的性质，得到$\int\log[f(x)]g(x)\,dx\le\int\log[g(x)]g(x)\,dx$.
\end{exercise}
\begin{exercise}\label{exercise:7.32}
\autoref{exercise:5.65}的算法可以调整，使其仅用一个先验分布样本，就能近似模拟后验分布的样本.设$X_1,X_2,\ldots,X_n\sim f(x\mid\theta)$，其中$\theta$的先验分布为$\pi$.从$\pi$生成$\theta_1,\theta_2,\ldots,\theta_m$，并计算$ q_i=L(\theta_i/\mid\boldsymbol x)\sum_jL(\theta_j/\mid\boldsymbol x)$，其中$L(\theta/\mid\boldsymbol x)=\prod_i f(x_i\mid\theta)$是似然函数.\footnote{译注：原文$q_i$把归一化写成了乘法，并在似然的条件竖线前多出斜线.为使概率和为1，应为$q_i=L(\theta_i\mid\boldsymbol x)/\sum_jL(\theta_j\mid\boldsymbol x)$.原记号保留.}
\begin{enumerate}[label=(\alph*)]
\item 生成$\theta_1^*,\theta_2^*,\ldots,\theta_r^*$，其中$P(\theta^*=\theta_i)=q_i$.说明这是一个近似的后验样本，其意义是$P(\theta^*\le t)$收敛到$\int_{-\infty}^t\pi(\theta\mid\boldsymbol x)\,d\theta$.
\item\phantomsection\label{itm:exercise:7.32:b} 说明估计量$\sum_{j=1}^rh(\theta_j^*)/r$收敛到$E[h(\theta)\mid\boldsymbol x]$，其中期望针对后验分布.
\sourcepageanchor{324}
\item \hyperref[bib:source261]{Ross（1996）}建议用Rao--Blackwell化改进\hyperref[itm:exercise:7.32:b]{(b)}中的估计.说明对于任意$j$，
$$E[h(\theta_j^*)\mid\theta_1,\theta_2,\ldots,\theta_m]=\frac1{\sum_{i=1}^mL(\theta_i\mid\boldsymbol x)}\sum_{i=1}^mh(\theta_i)L(\theta_i\mid\boldsymbol x)$$
与\hyperref[itm:exercise:7.32:b]{(b)}的估计量均值相同，方差较小.
\end{enumerate}
\end{exercise}
\begin{exercise}\label{exercise:7.33}
\autoref{ex:7.3.5}计算了成功概率的Bayes估计量$\widehat p_B$的MSE（估计量在\autoref{ex:7.2.14}中导出）.说明取$\alpha=\beta=\sqrt{n/4}$，可使$\widehat p_B$的MSE为常数.
\end{exercise}
\begin{exercise}\label{exercise:7.34}
设$X_1,\ldots,X_n$为来自$\operatorname{binomial}(n,p)$的随机样本.希望利用\autoref{ex:6.4.1}描述的群，寻找$p$的等变点估计量.
\begin{enumerate}[label=(\alph*)]
\item 求关于此群等变的估计量类.
\item\phantomsection\label{itm:exercise:7.34:b} 在\autoref{ex:7.2.14}的Bayes估计量类中，求关于此群等变的估计量.
\item 从\hyperref[itm:exercise:7.34:b]{(b)}的等变Bayes估计量中，求MSE最小者.
\end{enumerate}
\end{exercise}
\begin{exercise}\label{exercise:7.35}\index{SI0755@Pitman估计量（见点估计量：不变）}\index{SI0790@Pitman}
Pitman位置估计量（Pitman estimator of location；见\hyperref[bib:source200]{Lehmann和Casella 1998}，第3.1节，或\hyperref[bib:source242]{Pitman 1939}的原始论文）为
$$d_P(\boldsymbol X)=\frac{\int_{-\infty}^\infty t\prod_{i=1}^nf(x_i-t)\,dt}{\int_{-\infty}^\infty\prod_{i=1}^nf(x_i-t)\,dt},$$
其中观测到来自$f(x-\theta)$的随机样本$X_1,\ldots,X_n$.Pitman说明，该估计量是在位置等变估计量中均方误差最小者（即使\eqref{eq:7.3.3}最小）.本题的目标较为有限.
\begin{enumerate}[label=(\alph*)]
\item 说明$d_P(\boldsymbol X)$关于\autoref{ex:7.3.6}的位置群不变.\footnote{译注：原文此处用“invariant”（不变），实际应为“equivariant”（等变）：平移样本$a$会使估计值也增加$a$.原词保留.}
\item 说明若$f(x-\theta)$为$n(\theta,1)$，则$d_P(\boldsymbol X)=\bar X$.
\item 说明若$f(x-\theta)$为$\operatorname{uniform}(\theta-\frac12,\theta+\frac12)$，则$d_P(\boldsymbol X)=\frac12(X_{(1)}+X_{(n)})$.
\end{enumerate}
\end{exercise}
\begin{exercise}\label{exercise:7.36}
Pitman尺度估计量（Pitman estimator of scale）为
$$d_P^r(\boldsymbol X)=\frac{\int_0^\infty t^{n+r-1}\prod_{i=1}^nf(tx_i)\,dt}{\int_0^\infty t^{n+2r-1}\prod_{i=1}^nf(tx_i)\,dt},$$
其中观测到来自$\frac1\sigma f(x/\sigma)$的随机样本$X_1,\ldots,X_n$.Pitman说明，它是估计$\sigma^r$的尺度等变估计量中，尺度化均方误差最小者（即使$E(d-\sigma^r)^2/\sigma^{2r}$最小）.
\begin{enumerate}[label=(\alph*)]
\item 说明$d_P^r(\boldsymbol X)$关于尺度群等变，即对任意常数$c>0$满足
$$d_P^r(cx_1,\ldots,cx_n)=c^rd_P^r(x_1,\ldots,x_n).$$
\item 若$X_1,\ldots,X_n$独立同分布于$n(0,\sigma^2)$，求$\sigma^2$的Pitman尺度等变估计量.
\item 若$X_1,\ldots,X_n$独立同分布于$\operatorname{exponential}(\beta)$，求$\beta$的Pitman尺度等变估计量.
\item 若$X_1,\ldots,X_n$独立同分布于$\operatorname{uniform}(0,\theta)$，求$\theta$的Pitman尺度等变估计量.
\end{enumerate}
\end{exercise}
\begin{exercise}\label{exercise:7.37}
设$X_1,\ldots,X_n$为来自具有下列pdf的总体的随机样本：
$$f(x\mid\theta)=\frac1{2\theta},\qquad-\theta<x<\theta,\quad\theta>0.$$
若存在，求$\theta$的最佳无偏估计量.
\end{exercise}
\sourcepageanchor{325}
\begin{exercise}\label{exercise:7.38}
对于下列每个分布，设$X_1,\ldots,X_n$为随机样本.是否存在$\theta$的某个函数$g(\theta)$，使其有一个无偏估计量的方差达到Cramér--Rao下界？若存在，求出它；若不存在，说明原因.
\begin{enumerate}[label=(\alph*)]
\item $f(x\mid\theta)=\theta x^{\theta-1}$，$0<x<1$，$\theta>0$.
\item $f(x\mid\theta)=\frac{\log\theta}{\theta-1}\theta^x$，$0<x<1$，$\theta>1$.
\end{enumerate}
\end{exercise}
\begin{exercise}\label{exercise:7.39}
证明\autoref{lem:7.3.11}.
\end{exercise}
\begin{exercise}\label{exercise:7.40}
设$X_1,\ldots,X_n$独立同分布于$\operatorname{Bernoulli}(p)$.说明$\bar X$的方差达到Cramér--Rao下界，从而$\bar X$是$p$的最佳无偏估计量.
\end{exercise}
\begin{exercise}\label{exercise:7.41}\index{SI0792@线性}\index{SI1053@线性}
设$X_1,\ldots,X_n$为来自均值$\mu$、方差$\sigma^2$的总体的随机样本.
\begin{enumerate}[label=(\alph*)]
\item 说明若$\sum_{i=1}^na_i=1$，则估计量$\sum_{i=1}^na_iX_i$对$\mu$无偏.
\item 在所有这种形式的无偏估计量（称为线性无偏估计量，linear unbiased estimators）中，求方差最小者，并计算其方差.
\end{enumerate}
\end{exercise}
\begin{exercise}\label{exercise:7.42}
设$W_1,\ldots,W_k$为参数$\theta$的无偏估计量，$\operatorname{Var}W_i=\sigma_i^2$，且$i\ne j$时$\operatorname{Cov}(W_i,W_j)=0$.
\begin{enumerate}[label=(\alph*)]
\item 说明在所有形如$\sum a_iW_i$（$a_i$为常数）且满足$E_\theta(\sum a_iW_i)=\theta$的估计量中，$W^*=\sum W_i/\sigma_i^2\,/\sum(1/\sigma_i^2)$的方差最小.
\item 说明$\operatorname{Var}W^*=1/\sum(1/\sigma_i^2)$.
\end{enumerate}
\end{exercise}
\begin{exercise}\label{exercise:7.43}
\autoref{exercise:7.42}确立了最优权重$q_i^*=(1/\sigma_i^2)/(\sum_j1/\sigma_j^2)$.Tukey的一个结果（见\hyperref[bib:source38]{Bloch和Moses 1988）}指出：若$W=\sum_iq_iW_i$基于另一组权重$q_i\ge0$、$\sum_iq_i=1$，则
$$\frac{\operatorname{Var}W}{\operatorname{Var}W^*}\le\frac1{1-\lambda^2},$$
其中$\lambda$满足$(1+\lambda)/(1-\lambda)=b_{\max}/b_{\min}$，而$b_{\max}$与$b_{\min}$分别为$b_i=q_i/q_i^*$的最大和最小值.
\begin{enumerate}[label=(\alph*)]
\item 证明Tukey不等式.
\item 利用该不等式，作为$\sigma_{\max}^2/\sigma_{\min}^2$的函数评估通常均值$\sum_iW_i/k$的表现.
\end{enumerate}
\end{exercise}
\begin{exercise}\label{exercise:7.44}
设$X_1,\ldots,X_n$独立同分布于$n(\theta,1)$.说明$\theta^2$的最佳无偏估计量为$\bar X^2-1/n$.计算其方差（使用\autoref{sec:3.6}的Stein恒等式），并说明它大于Cramér--Rao下界.
\end{exercise}
\begin{exercise}\label{exercise:7.45}\index{SI0599@峰度}
设$X_1,X_2,\ldots,X_n$独立同分布于均值为$\mu$、方差为$\sigma^2$的分布，$S^2$为通常的$\sigma^2$无偏估计量.在\autoref{ex:7.3.4}中，看到在正态假设下，MLE的MSE小于$S^2$.本题进一步考察方差估计.
\begin{enumerate}[label=(\alph*)]
\item\phantomsection\label{itm:exercise:7.45:a} 说明对任意形如$aS^2$（$a$为常数）的估计量，有
$$\operatorname{MSE}(aS^2)=E[aS^2-\sigma^2]^2=a^2\operatorname{Var}(S^2)+(a-1)^2\sigma^4.$$
\item\phantomsection\label{itm:exercise:7.45:b} 说明
$$\operatorname{Var}(S^2)=\frac1n\left(\kappa-\frac{n-3}{n-1}\sigma^4\right),$$
其中$\kappa=E[X-\mu]^4$为峰度（kurtosis）.（你可能已在\hyperref[item:5.8.b]{习题5.8(b)}中完成此题.）
\item 说明正态假设下峰度为$3\sigma^4$，并确立此时$aS^2$形式中MSE最小的估计量是$\frac{n-1}{n+1}S^2$.（\autoref{lem:3.6.5}可能有帮助.）
\item\phantomsection\label{itm:exercise:7.45:d} 若不假设正态性，说明$\operatorname{MSE}(aS^2)$在
$$a=\frac{n-1}{(n+1)+\frac{(\kappa-1)(n-1)}n}$$
处最小；由于依赖参数，这个结果没有实际用处.\footnote{译注：原题混用了四阶中心矩与标准化峰度，并且\hyperref[itm:exercise:7.45:d]{(d)}式的常数有误.若按\hyperref[itm:exercise:7.45:b]{(b)}定义$\kappa=E(X-\mu)^4$，由\hyperref[itm:exercise:7.45:a]{(a)}\hyperref[itm:exercise:7.45:b]{(b)}直接求导得$a=(n-1)/[n+1+(\kappa/\sigma^4-3)(n-1)/n]$；\hyperref[itm:exercise:7.45:e]{(e)}中与3比较的也应是$\kappa/\sigma^4$.原公式与原符号保留.}
\needspace{5\baselineskip}
\item\phantomsection\label{itm:exercise:7.45:e} 说明
\begin{enumerate}[label=(\roman*)]
\item 对于$\kappa>3$的分布，最优$a$满足$a<(n-1)/(n+1)$；
\item 对于$\kappa<3$的分布，最优$a$满足$(n-1)/(n+1)<a<1$.
\end{enumerate}
\end{enumerate}
更多细节参见\hyperref[bib:source284]{Searls和Intarapanich（1990）}.
\end{exercise}
\sourcepageanchor{326}
\begin{exercise}\label{exercise:7.46}
设$X_1,X_2,X_3$为来自$\operatorname{uniform}(\theta,2\theta)$分布的大小为3的随机样本，其中$\theta>0$.
\begin{enumerate}[label=(\alph*)]
\item 求$\theta$的矩估计量.
\item 求MLE $\widehat\theta$，并求常数$k$，使$E_\theta(k\widehat\theta)=\theta$.
\item 这两个估计量中，哪个可利用充分性改进？如何改进？
\item 根据数据$1.29,.86,1.33$求$\theta$的矩估计值和MLE.这些是酿酒葡萄平均果粒大小（厘米）的三次观测.
\end{enumerate}
\end{exercise}
\begin{exercise}\label{exercise:7.47}
假设测量一个圆的半径时，误差服从$n(0,\sigma^2)$分布.若作$n$次独立测量，求圆面积的一个无偏估计量.它是最佳无偏的吗？
\end{exercise}
\begin{exercise}\label{exercise:7.48}
假设$X_i$，$i=1,\ldots,n$，独立同分布于$\operatorname{Bernoulli}(p)$.
\begin{enumerate}[label=(\alph*)]
\item 说明$p$的MLE的方差达到Cramér--Rao下界.
\item 对$n\ge4$，说明乘积$X_1X_2X_3X_4$是$p^4$的无偏估计量，并利用这一事实求$p^4$的最佳无偏估计量.
\end{enumerate}
\end{exercise}
\begin{exercise}\label{exercise:7.49}
设$X_1,\ldots,X_n$独立同分布于$\operatorname{exponential}(\lambda)$.
\begin{enumerate}[label=(\alph*)]
\item\phantomsection\label{itm:exercise:7.49:a} 求仅基于$Y=\min\{X_1,\ldots,X_n\}$的$\lambda$无偏估计量.
\item\phantomsection\label{itm:exercise:7.49:b} 求一个优于\hyperref[itm:exercise:7.49:a]{(a)}中估计量的估计量，并证明它更好.
\item 下面数据是航天飞机持续加压环境中使用的Kevlar/环氧树脂球形容器在高应力下的失效时间（小时）：
\begin{gather*}
50.1,\ 70.1,\ 137.0,\ 166.9,\ 170.5,\ 152.8,\\
80.5,\ 123.5,\ 112.6,\ 148.5,\ 160.0,\ 125.4.
\end{gather*}
失效时间常用指数分布建模.用\hyperref[itm:exercise:7.49:a]{(a)}与\hyperref[itm:exercise:7.49:b]{(b)}中的估计量估计平均失效时间.
\end{enumerate}
\end{exercise}
\begin{exercise}\label{exercise:7.50}\index{SI1056@标准差}
设$X_1,\ldots,X_n$独立同分布于$n(\theta,\theta^2)$，$\theta>0$.在此模型中，$\bar X$与$cS$都是$\theta$的无偏估计量，其中
$$c=\frac{\sqrt{n-1}\,\Gamma((n-1)/2)}{\sqrt2\,\Gamma(n/2)}.$$
\begin{enumerate}[label=(\alph*)]
\item 证明对任意数$a$，估计量$a\bar X+(1-a)(cS)$都是$\theta$的无偏估计量.
\item 求使估计量方差最小的$a$值.
\item 说明$(\bar X,S^2)$是$\theta$的充分统计量，但不是完备充分统计量.
\end{enumerate}
\end{exercise}
\begin{exercise}\label{exercise:7.51}
\hyperref[bib:source149]{Gleser和Healy（1976）}详细讨论了$n(\theta,a\theta^2)$族中的估计问题，其中$a$是已知常数（\autoref{exercise:7.50}是其特殊情形）.这里考察他们结果的一小部分.仍设$X_1,\ldots,X_n$独立同分布于$n(\theta,\theta^2)$，$\theta>0$，令$\bar X$与$cS$如\hyperref[exercise:7.50]{习题7.46}.\footnote{译注：原文此处写“习题7.46”，实际应为\autoref{exercise:7.50}.原显示编号保留，链接指向实际定义.}定义估计量类
$$\mathcal T=\{T:T=a_1\bar X+a_2(cS)\},$$
其中不假设$a_1+a_2=1$.
\begin{enumerate}[label=(\alph*)]
\item 求使$E_\theta(\theta-T)^2$最小的估计量$T\in\mathcal T$，记为$T^*$.
\item 说明$T^*$的MSE小于\hyperref[exercise:7.50]{习题7.50(b)}所得估计量的MSE.
\item 说明$T^{*+}=\max\{0,T^*\}$的MSE小于$T^*$的MSE.
\item $\theta$应归为位置参数还是尺度参数？解释理由.
\end{enumerate}
\end{exercise}
\begin{exercise}\label{exercise:7.52}
设$X_1,\ldots,X_n$独立同分布于$\operatorname{Poisson}(\lambda)$，$\bar X$与$S^2$分别表示样本均值和样本方差.现在用另一种方法完成\autoref{ex:7.3.8}.那里用了Cramér--Rao下界，这里使用完备性.
\begin{enumerate}[label=(\alph*)]
\item 不用Cramér--Rao定理，证明$\bar X$是$\lambda$的最佳无偏估计量.
\item\phantomsection\label{itm:exercise:7.52:b} 证明颇为惊人的恒等式$E(S^2\mid\bar X)=\bar X$，并用它明确说明$\operatorname{Var}S^2>\operatorname{Var}\bar X$.
\item 利用完备性，能否表述一个一般定理，使\hyperref[itm:exercise:7.52:b]{(b)}中的恒等式成为其特殊情形？
\end{enumerate}
\end{exercise}
\sourcepageanchor{327}
\begin{exercise}\label{exercise:7.53}
补全\autoref{thm:7.3.20}证明中省略的一些细节.设$W$是$\tau(\theta)$的无偏估计量，$U$是零的无偏估计量.证明：如果对某个$\theta=\theta_0$有$\operatorname{Cov}_{\theta_0}(W,U)\ne0$，则$W$不可能是$\tau(\theta)$的最佳无偏估计量.
\end{exercise}
\needspace{8\baselineskip}
\begin{exercise}\label{exercise:7.54}\index{SI0859@尼罗河问题}
对于“尼罗河问题”（见\autoref{exercise:6.37}），
\begin{enumerate}[label=(\alph*)]
\item 说明$T$是$\theta$的MLE，$U$是辅助统计量，并且
$$E(T)=\frac{\Gamma(n+1/2)\Gamma(n-1/2)}{[\Gamma(n)]^2}\theta,\qquad E(T^2)=\frac{\Gamma(n+1)\Gamma(n-1)}{[\Gamma(n)]^2}\theta^2.$$
\item 令$Z_1=(n-1)/\sum X_i$，$Z_2=\sum Y_i/n$.说明两者均为无偏估计量，方差分别为$\theta^2/(n-2)$和$\theta^2/n$.
\item 求形如$aZ_1+(1-a)Z_2$的最佳无偏估计量，计算其方差，并将它与经偏差校正的MLE比较.
\end{enumerate}
\end{exercise}
\begin{exercise}\label{exercise:7.55}
对下列每个pdf，设$X_1,\ldots,X_n$是来自该分布的样本.分别求$\theta^r$的最佳无偏估计量.（\hyperref[bib:source153]{Guenther（1978）}对这一问题有完整讨论.）
\begin{enumerate}[label=(\alph*)]
\item $f(x\mid\theta)=1/\theta$，$0<x<\theta$，$r<n$.
\item $f(x\mid\theta)=e^{-(x-\theta)}$，$x>\theta$.
\item $f(x\mid\theta)=e^{-x}/(e^{-\theta}-e^{-b})$，$\theta<x<b$，$b$已知.
\end{enumerate}
\end{exercise}
\begin{exercise}\label{exercise:7.56}
证明\autoref{ex:7.3.24}之前正文中的断言：若$T$是参数$\theta$的完备充分统计量，$h(X_1,\ldots,X_n)$是$\tau(\theta)$的任意无偏估计量，则$\varphi(T)=E(h(X_1,\ldots,X_n)\mid T)$是$\tau(\theta)$的最佳无偏估计量.
\end{exercise}
\begin{exercise}\label{exercise:7.57}
设$X_1,\ldots,X_{n+1}$独立同分布于$\operatorname{Bernoulli}(p)$，定义函数$h(p)$为
$$h(p)=P\left(\sum_{i=1}^nX_i>X_{n+1}\mid p\right),$$
即前$n$个观测值之和超过第$n+1$个观测值的概率.
\begin{enumerate}[label=(\alph*)]
\item 说明
$$T(X_1,\ldots,X_{n+1})=\begin{cases}1&\text{若 }\sum_{i=1}^nX_i>X_{n+1},\\0&\text{其他情形}\end{cases}$$
是$h(p)$的无偏估计量.
\item 求$h(p)$的最佳无偏估计量.
\end{enumerate}
\end{exercise}
\begin{exercise}\label{exercise:7.58}
设$X$是来自下列pdf的一个观测：
$$f(x\mid\theta)=\left(\frac\theta2\right)^{|x|}(1-\theta)^{1-|x|},\qquad x=-1,0,1;\quad0\le\theta\le1.$$
\begin{enumerate}[label=(\alph*)]
\item 求$\theta$的MLE.
\item 定义估计量$T(X)$为
$$T(X)=\begin{cases}2&\text{若 }x=1,\\0&\text{其他情形}.\end{cases}$$
说明$T(X)$是$\theta$的无偏估计量.
\item 求一个优于$T(X)$的估计量，并证明它更好.
\end{enumerate}
\end{exercise}
\begin{exercise}\label{exercise:7.59}
设$X_1,\ldots,X_n$独立同分布于$n(\mu,\sigma^2)$.求$\sigma^p$的最佳无偏估计量，其中$p$是已知正常数，不一定是整数.
\end{exercise}
\begin{exercise}\label{exercise:7.60}
设$X_1,\ldots,X_n$独立同分布于$\operatorname{gamma}(\alpha,\beta)$，$\alpha$已知.求$1/\beta$的最佳无偏估计量.
\end{exercise}
\begin{exercise}\label{exercise:7.61}
说明基于观测$S^2\sim\sigma^2\chi_\nu^2/\nu$估计$\sigma^2$时，似然函数的对数可写为
$$\log L(\sigma^2\mid s^2)=K_1\frac{s^2}{\sigma^2}-K_2\log\frac{s^2}{\sigma^2}+K_3,$$
其中$K_1,K_2,K_3$是不依赖于$\sigma^2$的常数.将上述对数似然与\autoref{ex:7.3.27}讨论的损失函数联系起来.Anderson（1984a）讨论了这种关系.
\end{exercise}
\sourcepageanchor{328}
\begin{exercise}\label{exercise:7.62}
设$X_1,\ldots,X_n$是来自$n(\theta,\sigma^2)$总体的随机样本，$\sigma^2$已知.考虑在平方误差损失下估计$\theta$.令$\pi(\theta)$为$\theta$的$n(\mu,\tau^2)$先验分布，$\delta^\pi$为$\theta$的贝叶斯估计量.验证下列风险函数与贝叶斯风险公式.
\begin{enumerate}[label=(\alph*)]
\item 对任意常数$a,b$，估计量$\delta(\boldsymbol x)=a\bar X+b$的风险函数为
$$R(\theta,\delta)=a^2\frac{\sigma^2}{n}+(b-(1-a)\theta)^2.$$
\item 令$\eta=\sigma^2/(n\tau^2+\sigma^2)$.贝叶斯估计量的风险函数为
$$R(\theta,\delta^\pi)=(1-\eta)^2\frac{\sigma^2}{n}+\eta^2(\theta-\mu)^2.$$
\item 贝叶斯估计量的贝叶斯风险为$B(\pi,\delta^\pi)=\tau^2\eta$.
\end{enumerate}
\end{exercise}
\begin{exercise}\label{exercise:7.63}
设$X\sim n(\mu,1)$.令$\delta^\pi$为平方误差损失下$\mu$的贝叶斯估计量.分别对$\pi(\mu)\sim n(0,1)$与$\pi(\mu)\sim n(0,10)$计算风险函数$R(\mu,\delta^\pi)$并作图.讨论先验如何影响贝叶斯估计量的风险函数.
\end{exercise}
\begin{exercise}\label{exercise:7.64}
设$X_1,\ldots,X_n$是独立随机变量，$X_i$的cdf为$F(x\mid\theta_i)$.证明：若对$i=1,\ldots,n$，$\delta_i^{\pi_i}(X_i)$是在损失$L(\theta_i,a_i)$和先验$\pi_i(\theta_i)$下估计$\theta_i$的贝叶斯法则，则$\delta^\pi(\boldsymbol X)=(\delta_1^{\pi_1}(X_1),\ldots,\delta_n^{\pi_n}(X_n))$是在损失$\sum_{i=1}^nL(\theta_i,a_i)$和先验$\pi(\boldsymbol\theta)=\prod_{i=1}^n\pi_i(\theta_i)$下估计$\boldsymbol\theta=(\theta_1,\ldots,\theta_n)$的贝叶斯法则.
\end{exercise}
\begin{exercise}\label{exercise:7.65}
\hyperref[bib:source338]{Zellner（1986）}研究的一种损失函数是LINEX（LINear--EXponential，线性—指数）损失，它能以光滑方式处理非对称性.LINEX损失为
$$L(\theta,a)=e^{c(a-\theta)}-c(a-\theta)-1,$$
其中$c$是正常数.随着$c$变化，损失函数可以从极不对称变为近乎对称.
\begin{enumerate}[label=(\alph*)]
\item 对$c=.2,.5,1$，将$L(\theta,a)$作为$a-\theta$的函数作图.
\item 若$X\sim F(x\mid\theta)$，证明采用先验$\pi$时，$\theta$的贝叶斯估计量为$\delta^\pi(X)=-(1/c)\log E(e^{-c\theta}\mid X)$.
\item 设$X_1,\ldots,X_n$独立同分布于$n(\theta,\sigma^2)$，其中$\sigma^2$已知，并设$\theta$具有无信息先验$\pi(\theta)=1$.证明LINEX损失下的贝叶斯估计量为$\delta^B(\boldsymbol X)=\bar X-c\sigma^2/(2n)$.
\item 在LINEX损失下，计算$\delta^B(\boldsymbol X)$与$\bar X$的后验期望损失.
\item 在平方误差损失下，计算$\delta^B(\boldsymbol X)$与$\bar X$的后验期望损失.
\end{enumerate}
\end{exercise}
\begin{exercise}\label{exercise:7.66}\index{SI0791@刀切法}
刀切法（jackknife）是减少估计量偏差的一种一般方法（\hyperref[bib:source249]{Quenouille 1956）}.一步刀切估计量定义如下.设$X_1,\ldots,X_n$为随机样本，$T_n=T_n(X_1,\ldots,X_n)$为参数$\theta$的某个估计量.为对$T_n$实施“刀切”，我们计算$n$个统计量$T_n^{(i)}$，$i=1,\ldots,n$；$T_n^{(i)}$的计算方法与$T_n$相同，但使用从样本中去掉$X_i$后剩下的$n-1$个观测.$\theta$的刀切估计量记为$\operatorname{JK}(T_n)$，定义为
$$\operatorname{JK}(T_n)=nT_n-\frac{n-1}{n}\sum_{i=1}^nT_n^{(i)}.$$
（一般说来，$\operatorname{JK}(T_n)$的偏差比$T_n$小.\hyperref[bib:source220]{Miller（1974）}对刀切法的性质作了很好的综述.）

具体地，现设$X_1,\ldots,X_n$独立同分布于$\operatorname{Bernoulli}(\theta)$.目标是估计$\theta^2$.
\begin{enumerate}[label=(\alph*)]
\item 说明$\theta^2$的MLE，$(\sum_{i=1}^nX_i/n)^2$，是$\theta^2$的有偏估计量.
\item 推导基于MLE的一步刀切估计量.
\item 说明一步刀切估计量是$\theta^2$的无偏估计量.（一般说来，刀切仅能减小偏差，但在这一特殊情形中，它完全消除了偏差.）
\item 这个刀切估计量是$\theta^2$的最佳无偏估计量吗？若是，请证明；若不是，请求出最佳无偏估计量.
\end{enumerate}
\end{exercise}
\begin{exercise}\label{exercise:7.67}
证明\autoref{thm:10.1.5}.
\end{exercise}



\sourcepageanchor{329}
\subsection{杂录}\label{sec:7.5}
\subsubsection{矩估计量与MLE}\label{sec:7.5.1}\index{SI0680@与矩估计法的关系}\index{SI0810@与最大似然估计量的关系}
一般说来，矩估计量不是充分统计量的函数，因此总可以通过对充分统计量取条件期望加以改进.不过，对指数族而言，修正的矩方法与最大似然估计之间可能存在对应关系.\hyperref[bib:source93]{Davidson和Solomon（1974）}详细讨论了这种对应关系，并介绍了一些有趣的历史.

设有随机样本$\boldsymbol X=(X_1,\ldots,X_n)$，其pdf属于指数族（见\autoref{thm:5.2.11}）
$$f(x\mid\boldsymbol\theta)=h(x)c(\boldsymbol\theta)\exp\left(\sum_{i=1}^kw_i(\boldsymbol\theta)t_i(x)\right),$$
其中$f(x\mid\boldsymbol\theta)$的取值范围不依赖于$\boldsymbol\theta$.（注意$\boldsymbol\theta$可以是向量.）似然函数具有形式
$$L(\boldsymbol\theta\mid\boldsymbol x)=H(\boldsymbol x)[c(\boldsymbol\theta)]^n\exp\left(\sum_{i=1}^kw_i(\boldsymbol\theta)\sum_{j=1}^nt_i(x_j)\right),$$
而修正的矩方法用下面$k$个方程的解$\widehat w_i(\boldsymbol\theta)$估计$w_i(\boldsymbol\theta)$，$i=1,\ldots,k$：
$$\sum_{j=1}^nt_i(x_j)=E_{\boldsymbol\theta}\left(\sum_{j=1}^nt_i(X_j)\right),\qquad i=1,\ldots,k.$$
Davidson和Solomon推广了\hyperref[bib:source173]{Huzurbazar（1949）}的工作，证明估计量$\widehat w_i(\boldsymbol\theta)$实际上就是$w_i(\boldsymbol\theta)$的MLE.若定义$\eta_i=w_i(\boldsymbol\theta)$，$i=1,\ldots,k$，则对任何一一函数$g$，$g(\eta_i)$的MLE等于$g(\widehat\eta_i)=g(\widehat w_i(\boldsymbol\theta))$.利用下列事实，可以简化上述期望的计算（\hyperref[bib:source198]{Lehmann，1986}，第2.7节）.\footnote{译注：按本段所写的密度形式，规范化因子为$c(\boldsymbol\theta)$而非其倒数，下面两个导数式右侧都应带负号.例如泊松族$c(\lambda)=e^{-\lambda}$、$w(\lambda)=\log\lambda$，则$\partial\log c/\partial w=-\lambda$，不能等于$E(X)=\lambda$.原式保留.}
\begin{align*}
E_{\boldsymbol\theta}(t_i(X_j))
&=\frac{\partial}{\partial w_i(\boldsymbol\theta)}\log(c(\boldsymbol\theta)),\qquad i=1,\ldots,k,\quad j=1,\ldots,n;\\
\operatorname{Cov}_{\boldsymbol\theta}(t_i(X_j),t_{i'}(X_j))
&=\frac{\partial^2}{\partial w_i(\boldsymbol\theta)\partial w_{i'}(\boldsymbol\theta)}\log(c(\boldsymbol\theta)),\\[-2pt]
&\hspace{2cm}i,i'=1,\ldots,k,\quad j=1,\ldots,n.
\end{align*}
\subsubsection{无偏贝叶斯估计}\label{sec:7.5.2}\index{SI0082@与无偏估计量的关系}\index{SI1058@与贝叶斯估计量的关系}
如\autoref{sec:7.2.3}所见，进行贝叶斯计算时，通常取后验分布的均值作为点估计量.具体说来，若$X$的pdf为$f(x\mid\theta)$，满足$E_\theta(X)=\theta$，并有先验分布$\pi(\theta)$，则后验均值这一$\theta$的贝叶斯点估计量为
$$E(\theta\mid x)=\int\theta\pi(\theta\mid x)\,d\theta.$$
可以问：$E(\theta\mid X)$能否是$\theta$的无偏估计量，从而满足方程
$$E_\theta[E(\theta\mid X)]=\int\left[\int\theta\pi(\theta\mid x)\,d\theta\right]f(x\mid\theta)\,dx=\theta.$$
\sourcepageanchor{330}
答案是否定的.也就是说，后验均值从来都不是无偏估计量.如果是，对$X$与$\theta$的联合分布取期望，便可写成\footnote{译注：本段结论须在存在适当先验及相应可积性条件下理解，且排除参数可由数据完全确定的退化情形；不能推广到非正常先验，例如正态位置模型的平坦非正常先验给出无偏后验均值$\bar X$.此外，下方第二组计算直接假设$E(\theta\mid X)=X$，这不等同于上一段的$E_\theta[E(\theta\mid X)]=\theta$.要论证一般情形，可令$T=E(\theta\mid X)$；若$E_\theta T=\theta$且相关二阶矩有限，则$E(T^2)=E(T\theta)=E(\theta^2)$，故$E[(T-\theta)^2]=0$.原文推导保留.}
\begin{align*}
E[(X-\theta)^2]&=E[X^2-2X\theta+\theta^2]&&\text{（展开平方）}\\
&=E\bigl(E((X^2-2X\theta+\theta^2)\mid\theta)\bigr)&&\text{（迭代取期望）}\\
&=E\bigl(E(X^2\mid\theta)-2\theta^2+\theta^2\bigr)&&\text{（$E(X\mid\theta)=E_\theta X=\theta$）}\\
&=E\bigl(E(X^2\mid\theta)-\theta^2\bigr)\\
&=E(X^2)-E(\theta^2),&&\text{（期望的性质）}
\end{align*}
这是按一种方式取条件期望.类似地，对$X$取条件期望可得
\begin{align*}
E[(X-\theta)^2]&=E\bigl(E[(X^2-2X\theta+\theta^2)\mid X]\bigr)\\
&=E\bigl(X^2-2X^2+E(\theta^2\mid X)\bigr)&&\text{（由假设$E(\theta\mid X)=X$）}\\
&=E(\theta^2)-E(X^2).
\end{align*}
比较两次计算可知，只有$E(X^2)=E(\theta^2)$时才不会矛盾；这又意味着$E(X-\theta)^2=0$，从而$X=\theta$.只有$P(X=\theta)=1$时才会如此，这是一个无甚意义的情形，因此我们得到了矛盾.于是，要么$E(X\mid\theta)\ne\theta$，要么$E(\theta\mid X)\ne X$，这说明后验均值不可能是无偏估计量.注意，我们隐含地假设了$E(X^2)<\infty$，但事实上，这个结果在更一般的条件下仍成立.\hyperref[bib:source32]{Bickel和Mallows（1988）}对这一主题作了更深入的讨论.\hyperref[bib:source234]{Noorbaloochi和Meeden（1983）}从更高的理论层次刻画了这种联系.

\subsubsection{Lehmann--Scheffé定理}\label{sec:7.5.3}
Lehmann--Scheffé定理是数理统计的一项重大成果，它将充分性、完备性和唯一性联系起来.正文的论述与Lehmann--Scheffé定理有某种互补关系，因此我们一直没有以经典形式陈述它（这一形式与\autoref{thm:7.3.23}相似）.事实上，Lehmann--Scheffé定理包含在\autoref{thm:7.3.19}和\autoref{thm:7.3.23}之中.
\begin{theorem}[Lehmann--Scheffé]\label{thm:7.5.1}\index{SI0617@Lehmann–Scheffé定理}
基于完备充分统计量的无偏估计量是唯一的.
\end{theorem}
\begin{proof}
设$T$是完备充分统计量，$\varphi(T)$是一个满足$E_\theta\varphi(T)=\tau(\theta)$的估计量.由\autoref{thm:7.3.23}，$\varphi(T)$是$\tau(\theta)$的最佳无偏估计量；由\autoref{thm:7.3.19}，最佳无偏估计量是唯一的.
\end{proof}
这个定理也可以不借助\hyperref[ex:7.3.3]{定理7.3.3}，仅用完备性的推论来证明，并由此给出通向\autoref{thm:7.3.23}的一条略有不同的路径.\footnote{译注：原文写“定理7.3.3”，本章该编号对应的是例，而非定理，疑为引用编号误植.此外，上述唯一性指对同一待估函数、在各参数下几乎处处相等；若两个基于$T$的估计量均无偏，则它们之差的期望恒为零，由完备性立即得到这一结论.}

\subsubsection{关于EM算法的进一步讨论}\label{sec:7.5.4}\index{SI0007@EM}\index{SI0395@EM算法（见算法：EM）}
EM算法源于20世纪50年代的工作（\hyperref[bib:source159]{Hartley 1958）}，但真正受到统计学界广泛重视，是在\hyperref[bib:source97]{Dempster、Laird和Rubin（1977）}的奠基性工作之后；他们详细说明了算法的基本结构，并通过多种应用展示了其用法.

EM算法的一个优点是，人们知道它收敛到不完全数据MLE的条件，尽管这一主题也衍生出一些流传的说法.\hyperref[bib:source97]{Dempster、Laird和Rubin（1977）}最初的收敛性证明有一个缺陷，后来\hyperref[bib:source47]{Boyles（1983）}和\hyperref[bib:source335]{Wu（1983）}给出了正确证明；另见\hyperref[bib:source122]{Finch等（1989）}.

我们的论述止于\autoref{thm:7.2.20}，该定理保证每次迭代都会增加似然值.不过，这可能不足以断定序列$\{\widehat{\boldsymbol\theta}^{(r)}\}$收敛到最大似然估计量.要作出这样的保证，还需要
\sourcepageanchor{331}
进一步的条件.下面由\hyperref[bib:source335]{Wu（1983）}给出的定理保证收敛到一个驻点，该点可能是局部极大值点，也可能是鞍点.
\begin{theorem}\label{thm:7.5.2}\index{SI0010@收敛}
若完全数据对数似然的期望$E[\log L(\boldsymbol\theta\mid\boldsymbol y,\boldsymbol x)\mid\boldsymbol\theta',\boldsymbol y]$关于$\boldsymbol\theta$和$\boldsymbol\theta'$都连续，则EM序列$\{\widehat{\boldsymbol\theta}^{(r)}\}$的所有极限点都是$L(\boldsymbol\theta\mid\boldsymbol y)$的驻点，并且$L(\widehat{\boldsymbol\theta}^{(r)}\mid\boldsymbol y)$单调收敛到某个驻点$\widehat{\boldsymbol\theta}$处的$L(\widehat{\boldsymbol\theta}\mid\boldsymbol y)$.
\end{theorem}
在指数族中，计算变得较为简单，因为对数似然对于缺失数据是线性的.计算有所简化，因为我们可以写成\footnote{译注：原文前两句有重复且第二句缺主语，译文保留其内容.按指数族的一般形式，对数似然实际是关于充分统计量$T_i$线性的，不必关于原始缺失数据线性；下面第二行的$E_{\boldsymbol\theta'}[\log h(\boldsymbol y,\boldsymbol X)]$似漏写条件$\mid\boldsymbol y$.}
\begin{align*}
&E[\log L(\boldsymbol\theta\mid\boldsymbol y,\boldsymbol x)\mid\boldsymbol\theta',\boldsymbol y]\\
&\quad=E_{\boldsymbol\theta'}\left[\left.\log\left(h(\boldsymbol y,\boldsymbol X)e^{\sum\eta_i(\boldsymbol\theta)T_i-B(\boldsymbol\theta)}\right)\right|\boldsymbol y\right]\\
&\quad=E_{\boldsymbol\theta'}[\log h(\boldsymbol y,\boldsymbol X)]+\sum\eta_i(\boldsymbol\theta)E_{\boldsymbol\theta'}[T_i\mid\boldsymbol y]-B(\boldsymbol\theta).
\end{align*}
因此，计算完全数据MLE只涉及较简单的期望$E_{\boldsymbol\theta'}[T_i\mid\boldsymbol y]$.\hyperref[bib:source208]{Little和Rubin（1987）}、\hyperref[bib:source312]{Tanner（1996）}以及\hyperref[bib:source274]{Sheaffer（1997）}\footnote{译注：原文作者写 Sheaffer；书末与此处主题、年份相符的书目为 Schafer（1997）的《Analysis of Incomplete Multivariate Data》. 保留原拼写，并链接至该书目.}提供了很好的EM算法综述；另见\hyperref[bib:source200]{Lehmann和Casella（1998}，第6.4节）.\hyperref[bib:source216]{McLachlan和Krishnan（1997）}则用一整本书论述EM算法.

\subsubsection{其他似然}\label{sec:7.5.5}
本章使用了最大似然法，并且看到它不仅提供了寻找估计量的方法，还带来了一套对推断很有用的大样本理论.

似然有许多修正形式.有的用于处理干扰参数（如剖面似然，profile likelihood）；有的用于需要更稳健的设定时（如拟似然，quasi likelihood）；还有的适用于删失数据（如部分似然，partial likelihood）.

此外还有许多变体，它们都能对这里介绍的普通似然作出某种改进.\hyperref[bib:source166]{Hinkley（1980）}的综述，或\hyperref[bib:source167]{Hinkley、Reid和Snell（1991）}编辑的综述文集，可作为进入这个丰富的似然世界的入口.

\subsubsection{其他贝叶斯分析}\label{sec:7.5.6}
\begin{enumerate}[label=\arabic*.]
\item \textbf{稳健贝叶斯分析（robust Bayes analysis）.}\index{SI0768@稳健}\semanticindexanchor{768:331}贝叶斯法则可能对先验分布的（主观）选择相当敏感，这使许多贝叶斯统计学家感到担忧.\hyperref[bib:source17]{Berger（1984）}的论文提出了稳健贝叶斯分析的思想.这种贝叶斯分析寻求对一系列先验分布都具有良好性质的估计量.也就是说，我们寻找一个估计量$\delta^*$，其表现具有稳健性，即不敏感于某个先验类中究竟哪个先验$\pi$是正确的先验.稳健贝叶斯估计量还可能具有良好的频率学性质，因而成为颇有吸引力的方法.Berger（\hyperref[bib:source19]{1990}，\hyperref[bib:source20]{1994}）与\hyperref[bib:source327]{Wasserman（1994）}\footnote{译注：原文在此写 Wasserman（1994），书末对应稳健贝叶斯综述列为1992年；保留原年份，链接至该条文献.}的综述可作为这一主题的入门读物.
\item \textbf{经验贝叶斯分析（empirical Bayes analysis）.}\index{SI0396@经验贝叶斯估计量}\semanticindexanchor{396:331}\index{SI0766@经验}\semanticindexanchor{766:331}标准贝叶斯分析中，先验分布通常含有需要实验者指定的参数.例如，考虑如下设定：
$$\begin{aligned}X\mid\theta&\sim n(\theta,1),\\\theta\mid\tau^2&\sim n(0,\tau^2).\end{aligned}$$
贝叶斯实验者会事先指定$\tau^2$的值，从而进行贝叶斯分析.然而，由于$X$的边缘分布为$n(0,\tau^2+1)$，它包含关于$\tau$的信息，可以用来估计$\tau$.这种从边缘分布估计先验参数的思想，是经验贝叶斯分析的特征.经验贝叶斯方法有助于构造改进的方法，\hyperref[bib:source225]{Morris（1983）}以及\hyperref[bib:source68]{Casella和Hwang（1987）}给出了这方面的例子.\hyperref[bib:source141]{Gianola和Fernando（1986）}已成功地
\sourcepageanchor{332}
将这类方法应用于解决实际问题.\hyperref[bib:source57]{Carlin和Louis（1996）}全面论述了经验贝叶斯方法，而Casella（\hyperref[bib:source63]{1985}，\hyperref[bib:source65]{1992}）提供了技术性较弱的入门介绍.
\item \textbf{层次贝叶斯分析（hierarchical Bayes analysis）.}\index{SI0470@分层贝叶斯估计量}\semanticindexanchor{470:332}\index{SI0767@分层}\semanticindexanchor{767:332}处理上述设定的另一种方法，是不为$\tau^2$事先指定一个值，而采用层次设定，即再为$\tau^2$指定第二层先验.例如，可以采用\footnote{译注：这里的$\operatorname{uniform}(0,\infty)$只是对$\tau^2$取常数密度的非正常先验记号，不是概率分布.而且对于本例单个观测$X$，积分掉$\theta$后，$\tau^2$的后验核为$(1+\tau^2)^{-1/2}\exp\{-x^2/[2(1+\tau^2)]\}$，其在无穷远处的积分发散，因此这一具体层次设定不能产生正常后验.原例保留.}
$$\begin{aligned}X\mid\theta&\sim n(\theta,1),\\\theta\mid\tau^2&\sim n(0,\tau^2),\\\tau^2&\sim\operatorname{uniform}(0,\infty)\qquad\text{（非正常先验）}.\end{aligned}$$
层次建模，无论贝叶斯还是非贝叶斯，都是十分有效的工具，通常给出对于底层模型相当稳健的答案.\hyperref[bib:source207]{Lindley和Smith（1972）}展示了其用途，此后层次模型得到了广泛应用与发展.\hyperref[bib:source134]{Gelfand和Smith（1990）}的奠基性论文将层次模型与计算算法联系起来，贝叶斯方法的适用范围由此迅速扩大.\hyperref[bib:source200]{Lehmann和Casella（1999，第4.5节）}\footnote{译注：此处原文写1999年，书末该书第二版列为1998年；原年份保留.}介绍了层次贝叶斯理论，\hyperref[bib:source256]{Robert和Casella（1999）}则讨论其应用以及与计算算法的联系.
\end{enumerate}





\section{假设检验}\label{sec:8}
\chapterdoi{8}

\sourcepageanchor{333}
\begin{quote}
把奇怪与神秘混为一谈，是一种错误.
\begin{flushright}夏洛克·福尔摩斯\\《血字的研究》\end{flushright}
\end{quote}



\subsection{引言}\label{sec:8.1}
在\autoref{sec:7}中，我们研究了称为点估计的一种推断方法. 现在转向另一种推断方法——假设检验（hypothesis testing\mbox{）.} 为同时满足构造和评价假设检验的需要，本章与\autoref{sec:7}一样分为两部分. 首先定义统计假设.
\begin{definition}\label{def:8.1.1}\index{SI0474@假设}
假设（hypothesis）是关于总体参数的一个陈述.
\end{definition}
假设的定义相当宽泛，但重要的是，假设对总体作出陈述. 假设检验的目标是根据来自总体的一个样本，判断两个互补假设中哪一个为真.
\begin{definition}\label{def:8.1.2}\index{SI0475@备择}\index{SI0477@零}
假设检验问题中的两个互补假设称为原假设（null hypothesis）和备择假设（alternative hypothesis），分别记为 $H_0$ 和 $H_1$.
\end{definition}
若 $\theta$ 表示总体参数，原假设和备择假设的一般形式为 $H_0:\theta\in\Theta_0$ 和 $H_1:\theta\in\Theta_0^c$，其中 $\Theta_0$ 是参数空间的某个子集，$\Theta_0^c$ 是它的补集. 例如，若 $\theta$ 表示患者服药后血压变化的平均值，实验者可能希望检验 $H_0:\theta=0$ 对 $H_1:\theta\ne0$. 原假设说药物平均而言对血压没有影响，备择假设则说存在某种影响. 这种由 $H_0$ 表示处理无效的常见情况，产生了“零”假设这一名称. 再如，消费者可能关心某供应商产品中的次品比例. 若用 $\theta$ 表示次品比例，消费者可能希望检验 $H_0:\theta\ge\theta_0$ 对 $H_1:\theta<\theta_0$. $\theta_0$ 是可接受的最大次品比例，而 $H_0$ 表示次品比例高到不可接受. 假设涉及产品质量的问题称为验收抽样（acceptance sampling）\index{SI0002@验收抽样}\semanticindexanchor{2:333}问题.

在假设检验问题中，观察样本之后，实验者必须作出以下二者之一的决定：接受 $H_0$ 为真，或者拒绝 $H_0$、判定它为假，并判定 $H_1$ 为真.
\begin{definition}\label{def:8.1.3}\index{SI0482@假设检验}
假设检验程序（hypothesis testing procedure）或假设检验（hypothesis test）是规定下列事项的规则：
\begin{enumerate}[label=\roman*.]
\item 对哪些样本值，作出接受 $H_0$ 为真的决定.
\item 对哪些样本值，拒绝 $H_0$ 并接受 $H_1$ 为真.
\end{enumerate}
使 $H_0$ 被拒绝的那些样本点构成的样本空间子集，称为拒绝域（rejection region）\index{SI0506@拒绝域}\semanticindexanchor{506:333}或临界域（critical region\mbox{）.} 拒绝域的补集称为接受域（acceptance region\mbox{）.}\index{SI0484@接受域}\semanticindexanchor{484:333}
\end{definition}
\sourcepageanchor{334}
在哲学层面上，有人关心“拒绝 $H_0$”与“接受 $H_1$”的区别. 前者并不暗示实验者接受什么状态，只说明由 $H_0$ 所定义的状态遭到拒绝. 同样，也可以区分“接受 $H_0$”与“不拒绝 $H_0$”. 前者暗示实验者愿意断言 $H_0$ 所指定的自然状态，后者则暗示实验者实际上并不相信 $H_0$，但缺乏拒绝它的证据. 本书大体上不讨论这些问题. 我们把假设检验视为必须在两种行动中择一的问题，这两种行动就是断言 $H_0$ 或断言 $H_1$.

通常，假设检验通过一个检验统计量（test statistic）\index{SI1033@检验统计量}\semanticindexanchor{1033:334}$W(X_1,\ldots,X_n)=W(\boldsymbol X)$ 来规定；它是样本的函数. 例如，一个检验可以规定：若样本均值 $\bar X$ 大于 $3$，则拒绝 $H_0$. 这时 $W(\boldsymbol X)=\bar X$ 是检验统计量，拒绝域为 $\{(x_1,\ldots,x_n):\bar x>3\}$. \autoref{sec:8.2}讨论选择检验统计量和拒绝域的方法，\autoref{sec:8.3}介绍评价检验的准则. 与点估计量一样，构造检验的方法本身不提供任何保证；必须对所得检验作出评价，才能确定它们的价值.



\subsection{构造检验的方法}\label{sec:8.2}
我们将详细讨论四种构造检验程序的方法，它们适用于不同情形，利用问题的不同方面. 首先介绍一种非常一般的方法，它几乎总能应用，而且在某些情况下还是最优的.
\subsubsection{似然比检验}\label{sec:8.2.1}
假设检验的似然比方法与\autoref{sec:7.2.2}讨论的最大似然估计量有关，似然比检验的适用范围也与最大似然估计一样广泛. 回顾，若 $X_1,\ldots,X_n$ 是来自密度或质量函数为 $f(x\mid\theta)$ 的总体的随机样本（$\theta$ 可以是向量），则似然函数定义为
$$L(\theta\mid x_1,\ldots,x_n)=L(\theta\mid\boldsymbol x)=f(\boldsymbol x\mid\theta)=\prod_{i=1}^n f(x_i\mid\theta).$$
用 $\Theta$ 表示整个参数空间. 似然比检验定义如下.
\begin{definition}\label{def:8.2.1}\index{SI0630@似然比检验}
检验 $H_0:\theta\in\Theta_0$ 对 $H_1:\theta\in\Theta_0^c$ 的似然比检验统计量（likelihood ratio test statistic）为
$$\lambda(\boldsymbol x)=\frac{\sup_{\Theta_0}L(\theta\mid\boldsymbol x)}{\sup_\Theta L(\theta\mid\boldsymbol x)}.$$
似然比检验（likelihood ratio test，LRT）是拒绝域形如 $\{\boldsymbol x:\lambda(\boldsymbol x)\le c\}$ 的任一检验，其中 $c$ 是满足 $0\le c\le1$ 的任意数.
\end{definition}
当 $f(\boldsymbol x\mid\theta)$ 是离散随机变量的概率质量函数时，或许最容易理解似然比检验背后的道理. 此时，$\lambda(\boldsymbol x)$ 的分子是已观察样本的最大概率，最大值在原假设中的参数上计算（见\autoref{exercise:8.4}\mbox{）.} 分母则是已观察样本在所有可能参数上的最大概率. 若备择假设中存在某些参数点，使已观察样本比在原假设任一参数点下都可能得多，那么这两个最大值之比就很小. 此时，似然比检验准则认为应当拒绝 $H_0$，接受 $H_1$ 为真. 选择数值 $c$ 的方法将在\autoref{sec:8.3}讨论.

\sourcepageanchor{335}
如果把最大化区分为在整个参数空间上进行的无约束最大化（unrestricted maximization），以及在参数空间子集上进行的约束最大化（restricted maximization），似然比检验与最大似然估计量的联系就更清楚了. 假设 $\theta$ 的一个最大似然估计量 $\hat\theta$ 存在，它由对 $L(\theta\mid\boldsymbol x)$ 作无约束最大化得到. 也可以把 $\Theta_0$ 看作参数空间，考虑约束最大化所得的 $\theta$ 的最大似然估计量，记为 $\hat\theta_0$.\index{SI0635@最大似然估计量}\semanticindexanchor{635:335}\index{SI0678@似然比检验}\semanticindexanchor{678:335} 即 $\hat\theta_0=\hat\theta_0(\boldsymbol x)$ 是使 $L(\theta\mid\boldsymbol x)$ 最大的 $\theta\in\Theta_0$ 的值. 于是似然比检验统计量为
$$\lambda(\boldsymbol x)=\frac{L(\hat\theta_0\mid\boldsymbol x)}{L(\hat\theta\mid\boldsymbol x)}.$$
\begin{example}[正态总体的似然比检验]\label{ex:8.2.2}\index{SI0638@正态}
设 $X_1,\ldots,X_n$ 是来自 $\mathrm n(\theta,1)$ 总体的随机样本. 考虑检验 $H_0:\theta=\theta_0$ 对 $H_1:\theta\ne\theta_0$，其中 $\theta_0$ 是实验者在实验前固定的数. 由于 $H_0$ 只规定一个 $\theta$ 值，$\lambda(\boldsymbol x)$ 的分子为 $L(\theta_0\mid\boldsymbol x)$. \autoref{ex:7.2.5}已求出 $\theta$ 的无约束最大似然估计量是样本均值 $\bar X$，所以分母为 $L(\bar x\mid\boldsymbol x)$. 因此，似然比检验统计量为
\begin{equation}\label{eq:8.2.1}
\begin{aligned}
\lambda(\boldsymbol x)&=\frac{(2\pi)^{-n/2}\exp[-\sum_{i=1}^n(x_i-\theta_0)^2/2]}{(2\pi)^{-n/2}\exp[-\sum_{i=1}^n(x_i-\bar x)^2/2]}\\
&=\exp\left[\left(-\sum_{i=1}^n(x_i-\theta_0)^2+\sum_{i=1}^n(x_i-\bar x)^2\right)/2\right].
\end{aligned}
\end{equation}
注意到
$$\sum_{i=1}^n(x_i-\theta_0)^2=\sum_{i=1}^n(x_i-\bar x)^2+n(\bar x-\theta_0)^2,$$
就可化简 $\lambda(\boldsymbol x)$，得到
\begin{equation}\label{eq:8.2.2}
\lambda(\boldsymbol x)=\exp[-n(\bar x-\theta_0)^2/2].
\end{equation}
似然比检验在 $\lambda(\boldsymbol x)$ 较小时拒绝 $H_0$. 由\eqref{eq:8.2.2}，拒绝域 $\{\boldsymbol x:\lambda(\boldsymbol x)\le c\}$ 可写为
$$\{\boldsymbol x:|\bar x-\theta_0|\ge\sqrt{-2(\log c)/n}\}.$$
当 $c$ 在 $0$ 与 $1$ 之间变化时，$\sqrt{-2(\log c)/n}$ 在 $0$ 与 $\infty$ 之间变化. 因而，似然比检验就是当样本均值与假设值 $\theta_0$ 的差超过某个指定量时，拒绝 $H_0:\theta=\theta_0$ 的那些检验.
\end{example}
\autoref{ex:8.2.2}中的分析是典型的：先按\autoref{def:8.2.1}求出 $\lambda(\boldsymbol X)$ 的表达式，如\eqref{eq:8.2.1}所做的那样；再尽可能把拒绝域化为含有更简单统计量的表达式，本例中这个统计量就是 $|\bar X-\theta_0|$.
\begin{example}[指数总体的似然比检验]\label{ex:8.2.3}\index{SI0634@指数}
设 $X_1,\ldots,X_n$ 是来自指数总体的随机样本，其密度为
$$f(x\mid\theta)=\begin{cases}e^{-(x-\theta)},&x\ge\theta,\\0,&x<\theta,\end{cases}$$
\sourcepageanchor{336}
其中 $-\infty<\theta<\infty$. 似然函数为
$$L(\theta\mid\boldsymbol x)=\begin{cases}e^{-\sum x_i+n\theta},&\theta\le x_{(1)},\\0,&\theta>x_{(1)},\end{cases}\qquad(x_{(1)}=\min x_i).$$
考虑检验 $H_0:\theta\le\theta_0$ 对 $H_1:\theta>\theta_0$，其中 $\theta_0$ 由实验者指定. 显然，$L(\theta\mid\boldsymbol x)$ 在 $-\infty<\theta\le x_{(1)}$ 上是 $\theta$ 的增函数. 所以 $\lambda(\boldsymbol x)$ 的分母，即似然函数的无约束最大值，为
$$L(x_{(1)}\mid\boldsymbol x)=e^{-\sum x_i+nx_{(1)}}.$$
若 $x_{(1)}\le\theta_0$，分子也是 $L(x_{(1)}\mid\boldsymbol x)$. 但由于分子只在 $\theta\le\theta_0$ 上最大化，当 $x_{(1)}>\theta_0$ 时，它等于 $L(\theta_0\mid\boldsymbol x)$. 因此，似然比检验统计量为
$$\lambda(\boldsymbol x)=\begin{cases}1,&x_{(1)}\le\theta_0,\\e^{-n(x_{(1)}-\theta_0)},&x_{(1)}>\theta_0.\end{cases}$$

\begin{figure}[htbp]\centering
\includegraphics[width=.85\textwidth]{assets/chapter8/fig8-2-1.pdf}
\caption{$\lambda(\boldsymbol x)$，一个仅依赖于 $x_{(1)}$ 的函数.}\label{fig:8.2.1}
\end{figure}

\autoref{fig:8.2.1}画出了 $\lambda(\boldsymbol x)$. 似然比检验在 $\lambda(\boldsymbol X)\le c$ 时拒绝 $H_0$，其拒绝域是 $\{\boldsymbol x:x_{(1)}\ge\theta_0-(\log c)/n\}$.\footnote{译注：这一拒绝域的改写适用于 $0<c<1$. 若取定义允许的端点 $c=1$，由 $\lambda\le1$ 应拒绝全部样本，而所写集合仅含 $x_{(1)}\ge\theta_0$ 的样本；$c=0$ 时则无有限样本被拒绝，涉及 $\log0$ 的表达式应作极限理解.} 注意，拒绝域仅通过充分统计量 $X_{(1)}$ 依赖样本. \autoref{thm:8.2.4}将说明，这一现象普遍成立.
\end{example}
\autoref{ex:8.2.3}再次说明\autoref{sec:7.2.2}指出的一点：对似然函数求导，并不是求最大似然估计量的唯一方法. 在该例中，$L(\theta\mid\boldsymbol x)$ 在 $\theta=x_{(1)}$ 处不可微.

若 $T(\boldsymbol X)$ 是 $\theta$ 的充分统计量，密度或质量函数为 $g(t\mid\theta)$，则可以考虑用 $T$ 及其似然函数 $L^*(\theta\mid t)=g(t\mid\theta)$ 构造似然比检验，而不直接使用样本 $\boldsymbol X$ 及其似然函数 $L(\theta\mid\boldsymbol x)$. 用 $\lambda^*(t)$ 表示基于 $T$ 的似然比检验统计量. 直观上，$\boldsymbol x$ 中关于 $\theta$ 的全部信息都包含在 $T(\boldsymbol x)$ 中，因此基于 $T$ 的检验应当与基于完整样本的检验一样好. 事实上，这两种检验等价.
\begin{theorem}\label{thm:8.2.4}\index{SI0642@充分统计量}
若 $T(\boldsymbol X)$ 是 $\theta$ 的充分统计量，$\lambda^*(t)$ 和 $\lambda(\boldsymbol x)$ 分别为基于 $T$ 与 $\boldsymbol X$ 的似然比检验统计量，则对样本空间中的每个 $\boldsymbol x$，$\lambda^*(T(\boldsymbol x))=\lambda(\boldsymbol x)$.
\end{theorem}
\sourcepageanchor{337}
\begin{proof}
由因子分解定理（\autoref{thm:6.2.6}），$\boldsymbol X$ 的密度或质量函数可写为 $f(\boldsymbol x\mid\theta)=g(T(\boldsymbol x)\mid\theta)h(\boldsymbol x)$，其中 $g(t\mid\theta)$ 是 $T$ 的密度或质量函数，$h(\boldsymbol x)$ 不依赖于 $\theta$. 因此，
$$\begin{aligned}
\lambda(\boldsymbol x)&=\frac{\sup_{\Theta_0}L(\theta\mid\boldsymbol x)}{\sup_\Theta L(\theta\mid\boldsymbol x)}
=\frac{\sup_{\Theta_0}f(\boldsymbol x\mid\theta)}{\sup_\Theta f(\boldsymbol x\mid\theta)}\\
&=\frac{\sup_{\Theta_0}g(T(\boldsymbol x)\mid\theta)h(\boldsymbol x)}{\sup_\Theta g(T(\boldsymbol x)\mid\theta)h(\boldsymbol x)}&&\text{（$T$ 充分）}\\
&=\frac{\sup_{\Theta_0}g(T(\boldsymbol x)\mid\theta)}{\sup_\Theta g(T(\boldsymbol x)\mid\theta)}&&\text{（$h$ 不依赖于 $\theta$）}\\
&=\frac{\sup_{\Theta_0}L^*(\theta\mid T(\boldsymbol x))}{\sup_\Theta L^*(\theta\mid T(\boldsymbol x))}&&\text{（$g$ 是 $T$ 的密度或质量函数）}\\
&=\lambda^*(T(\boldsymbol x)).
\end{aligned}$$
\end{proof}
\autoref{ex:8.2.2}之后指出，求出 $\lambda(\boldsymbol x)$ 的表达式后，应设法将它化简. 根据\autoref{thm:8.2.4}，这句话的一种解释是：若 $T(\boldsymbol X)$ 是 $\theta$ 的充分统计量，则化简后的 $\lambda(\boldsymbol x)$ 应仅通过 $T(\boldsymbol x)$ 依赖于 $\boldsymbol x$.
\begin{example}[似然比检验与充分性]\label{ex:8.2.5}\index{SI0634@指数}
在\autoref{ex:8.2.2}中，可以看出 $\bar X$ 是 $\theta$ 的充分统计量. 利用 $\bar X$ 的似然函数（$\bar X\sim\mathrm n(\theta,1/n)$），可以更容易得到结论：检验 $H_0:\theta=\theta_0$ 对 $H_1:\theta\ne\theta_0$ 的似然比检验在 $|\bar X-\theta_0|$ 较大时拒绝 $H_0$.

类似地，在\autoref{ex:8.2.3}中，$X_{(1)}=\min X_i$ 是 $\theta$ 的充分统计量. $X_{(1)}$ 的似然函数（即其概率密度函数）为
$$L^*(\theta\mid x_{(1)})=\begin{cases}ne^{-n(x_{(1)}-\theta)},&\theta\le x_{(1)},\\0,&\theta>x_{(1)}.\end{cases}$$
也可利用这一似然函数，推出 $H_0:\theta\le\theta_0$ 对 $H_1:\theta>\theta_0$ 的似然比检验在 $X_{(1)}$ 较大时拒绝 $H_0$.
\end{example}
似然比检验对于存在干扰参数（nuisance parameters）\index{SI0727@滋扰参数}\semanticindexanchor{727:337}的情况也很有用；干扰参数是模型中存在、却不是推断直接关心的参数. 它们的存在不影响似然比检验的构造方法，但正如可能预料到的，它们可能导致不同的检验.
\begin{example}[方差未知时正态总体的似然比检验]\label{ex:8.2.6}\index{SI0638@正态}\index{SI0639@滋扰参数}\index{SI1019@似然比检验}
设 $X_1,\ldots,X_n$ 是来自 $\mathrm n(\mu,\sigma^2)$ 的随机样本，实验者只关心关于 $\mu$ 的推断，例如检验 $H_0:\mu\le\mu_0$ 对 $H_1:\mu>\mu_0$. 此时 $\sigma^2$ 是干扰参数. 似然比检验统计量为\footnote{译注：原文在最大化范围中写 $\sigma^2\ge0$，但所用普通正态密度要求 $\sigma^2>0$. $\sigma^2=0$ 对应退化分布，不能直接代入该密度；正文保留原范围.}
$$\begin{aligned}
\lambda(\boldsymbol x)&=\frac{\max_{\{\mu,\sigma^2:\mu\le\mu_0,\,\sigma^2\ge0\}}L(\mu,\sigma^2\mid\boldsymbol x)}{\max_{\{\mu,\sigma^2:-\infty<\mu<\infty,\,\sigma^2\ge0\}}L(\mu,\sigma^2\mid\boldsymbol x)}\\
&=\frac{\max_{\{\mu,\sigma^2:\mu\le\mu_0,\,\sigma^2\ge0\}}L(\mu,\sigma^2\mid\boldsymbol x)}{L(\hat\mu,\hat\sigma^2\mid\boldsymbol x)},
\end{aligned}$$
\sourcepageanchor{338}
其中 $\hat\mu$ 与 $\hat\sigma^2$ 是 $\mu$ 和 $\sigma^2$ 的最大似然估计量（\autoref{ex:7.2.11}\mbox{）.} 进一步，若 $\hat\mu\le\mu_0$，约束最大值与无约束最大值相同；若 $\hat\mu>\mu_0$，则约束最大值为 $L(\mu_0,\hat\sigma_0^2\mid\boldsymbol x)$，其中 $\hat\sigma_0^2=\sum(x_i-\mu_0)^2/n$. 因此，
$$\lambda(\boldsymbol x)=\begin{cases}1,&\hat\mu\le\mu_0,\\\dfrac{L(\mu_0,\hat\sigma_0^2\mid\boldsymbol x)}{L(\hat\mu,\hat\sigma^2\mid\boldsymbol x)},&\hat\mu>\mu_0.\end{cases}$$
经过一些代数运算，可以证明基于 $\lambda(\boldsymbol x)$ 的检验等价于基于 Student 的 $t$ 统计量的检验. 细节留作\autoref{exercise:8.37}. （\autoref{exercise:8.38}至\autoref{exercise:8.42}也讨论干扰参数问题.）
\end{example}

\subsubsection{贝叶斯检验}\label{sec:8.2.2}
假设检验问题也可以用贝叶斯模型表述. 回顾\autoref{sec:7.2.3}，贝叶斯模型不仅包含抽样分布 $f(\boldsymbol x\mid\theta)$，还包含先验分布 $\pi(\theta)$；先验分布反映实验者在抽样前对参数 $\theta$ 的看法.

贝叶斯范式要求用贝叶斯定理把样本信息与先验信息结合起来，得到后验分布 $\pi(\theta\mid\boldsymbol x)$. 此后关于 $\theta$ 的全部推断都以这个后验分布为依据.

在假设检验问题中，可以用后验分布计算 $H_0$ 与 $H_1$ 为真的概率. 要记住，$\pi(\theta\mid\boldsymbol x)$ 是一个随机变量的概率分布，因此可以计算后验概率
$$P(\theta\in\Theta_0\mid\boldsymbol x)=P(H_0\text{ 为真}\mid\boldsymbol x),\qquad
P(\theta\in\Theta_0^c\mid\boldsymbol x)=P(H_1\text{ 为真}\mid\boldsymbol x).$$
对经典统计学家而言，$P(H_0\text{ 为真}\mid\boldsymbol x)$ 和 $P(H_1\text{ 为真}\mid\boldsymbol x)$ 没有意义. 经典统计学家把 $\theta$ 看成固定的数，因此一个假设非真即假. 若 $\theta\in\Theta_0$，则对所有 $\boldsymbol x$，$P(H_0\text{ 为真}\mid\boldsymbol x)=1$，$P(H_1\text{ 为真}\mid\boldsymbol x)=0$. 若 $\theta\in\Theta_0^c$，这两个值就颠倒过来. 因为这些概率未知（$\theta$ 未知），且不依赖样本 $\boldsymbol x$，经典统计学家不使用它们. 在假设检验的贝叶斯表述中，这些概率依赖样本 $\boldsymbol x$，能提供有关 $H_0$ 和 $H_1$ 真伪的有用信息.

贝叶斯检验者使用后验分布的一种方式是：当 $P(\theta\in\Theta_0\mid\boldsymbol X)\ge P(\theta\in\Theta_0^c\mid\boldsymbol X)$ 时，接受 $H_0$ 为真，否则拒绝 $H_0$. 用前几节的术语来说，作为样本函数的检验统计量是 $P(\theta\in\Theta_0^c\mid\boldsymbol X)$，拒绝域是 $\{\boldsymbol x:P(\theta\in\Theta_0^c\mid\boldsymbol x)>1/2\}$. 另一种做法是：若希望防范错误拒绝 $H_0$，可以决定仅在 $P(\theta\in\Theta_0^c\mid\boldsymbol X)$ 超过某个很大的数（例如 $.99$）时才拒绝 $H_0$.
\begin{example}[正态总体的贝叶斯检验]\label{ex:8.2.7}\index{SI0488@正态}
设 $X_1,\ldots,X_n$ 独立同分布于 $\mathrm n(\theta,\sigma^2)$，$\theta$ 的先验分布为 $\mathrm n(\mu,\tau^2)$，其中 $\sigma^2,\mu,\tau^2$ 已知. 考虑检验 $H_0:\theta\le\theta_0$ 对 $H_1:\theta>\theta_0$. 由\autoref{ex:7.2.16}，后验分布 $\pi(\theta\mid\boldsymbol x)$ 为正态分布，均值为 $(n\tau^2\bar x+\sigma^2\mu)/(n\tau^2+\sigma^2)$，方差为 $\sigma^2\tau^2/(n\tau^2+\sigma^2)$.

若决定当且仅当 $P(\theta\in\Theta_0\mid\boldsymbol X)\ge P(\theta\in\Theta_0^c\mid\boldsymbol X)$ 时接受 $H_0$，则接受 $H_0$ 当且仅当
$$\frac12\le P(\theta\in\Theta_0\mid\boldsymbol X)=P(\theta\le\theta_0\mid\boldsymbol X).$$
由于 $\pi(\theta\mid\boldsymbol x)$ 对称，此式成立当且仅当后验均值不大于 $\theta_0$. 因此，若
$$\bar X\le\theta_0+\frac{\sigma^2(\theta_0-\mu)}{n\tau^2},$$
\sourcepageanchor{339}
则接受 $H_0$ 为真，否则接受 $H_1$ 为真. 特别地，若 $\mu=\theta_0$，即实验前赋予 $H_0$ 和 $H_1$ 的概率均为 $1/2$，那么当 $\bar x\le\theta_0$ 时接受 $H_0$ 为真，否则接受 $H_1$.
\end{example}
\autoref{sec:8.3.5}将讨论在假设检验问题中利用后验分布作推断的其他方法.

\subsubsection{并交检验与交并检验}\label{sec:8.2.3}
在某些情形下，可以从简单原假设的检验构造复杂原假设的检验. 我们讨论两种相关方法.

当原假设便于写成一个交集时，并交法（union–intersection method）可能很有用，例如
\begin{equation}\label{eq:8.2.3}H_0:\theta\in\bigcap_{\gamma\in\Gamma}\Theta_\gamma.\end{equation}
这里 $\Gamma$ 是任意指标集，依问题不同可以有限，也可以无限. 假设每个检验 $H_{0\gamma}:\theta\in\Theta_\gamma$ 对 $H_{1\gamma}:\theta\in\Theta_\gamma^c$ 的问题，都已有可用的检验. 设检验 $H_{0\gamma}$ 的拒绝域为 $\{\boldsymbol x:T_\gamma(\boldsymbol x)\in R_\gamma\}$. 则并交检验（union–intersection test）的拒绝域为
\begin{equation}\label{eq:8.2.4}\index{SI0522@并交}\bigcup_{\gamma\in\Gamma}\{\boldsymbol x:T_\gamma(\boldsymbol x)\in R_\gamma\}.\end{equation}
道理很简单：只要任一 $H_{0\gamma}$ 被拒绝，$H_0$ 也必须被拒绝，因为由\eqref{eq:8.2.3}，$H_0$ 为真仅当每个 $H_{0\gamma}$ 都为真. 只有每个 $H_{0\gamma}$ 都被接受为真，交集 $H_0$ 才会被接受为真.

有时，并交检验的拒绝域可以写得很简单. 特别地，假设每个单独检验的拒绝域都形如 $\{\boldsymbol x:T_\gamma(\boldsymbol x)>c\}$，其中 $c$ 不依赖于 $\gamma$. 则\eqref{eq:8.2.4}可写为
$$\bigcup_{\gamma\in\Gamma}\{\boldsymbol x:T_\gamma(\boldsymbol x)>c\}
=\{\boldsymbol x:\sup_{\gamma\in\Gamma}T_\gamma(\boldsymbol x)>c\}.$$
所以，检验 $H_0$ 的统计量是 $T(\boldsymbol x)=\sup_{\gamma\in\Gamma}T_\gamma(\boldsymbol x)$. \autoref{sec:11}给出了一些 $T(\boldsymbol x)$ 具有简单公式的例子.
\begin{example}[正态总体的并交检验]\label{ex:8.2.8}\index{SI0525@正态}\index{SI1021@双侧}
设 $X_1,\ldots,X_n$ 是来自 $\mathrm n(\mu,\sigma^2)$ 总体的随机样本. 考虑检验 $H_0:\mu=\mu_0$ 对 $H_1:\mu\ne\mu_0$，$\mu_0$ 为指定数. 可把 $H_0$ 写成两个集合的交：
$$H_0:\{\mu:\mu\le\mu_0\}\cap\{\mu:\mu\ge\mu_0\}.$$
检验 $H_{0L}:\mu\le\mu_0$ 对 $H_{1L}:\mu>\mu_0$ 的似然比检验为
$$\text{若 }\frac{\bar X-\mu_0}{S/\sqrt n}\ge t_L,\quad\text{则拒绝 $H_{0L}:\mu\le\mu_0$，支持 $H_{1L}:\mu>\mu_0$}$$
（见\autoref{exercise:8.37}\mbox{）.} 类似地，检验 $H_{0U}:\mu\ge\mu_0$ 对 $H_{1U}:\mu<\mu_0$ 的似然比检验为
$$\text{若 }\frac{\bar X-\mu_0}{S/\sqrt n}\le t_U,\quad\text{则拒绝 $H_{0U}:\mu\ge\mu_0$，支持 $H_{1U}:\mu<\mu_0$}.$$
由这两个似然比检验构成的并交检验，检验 $H_0:\mu=\mu_0$ 对 $H_1:\mu\ne\mu_0$，其规则为
$$\text{若 }\frac{\bar X-\mu_0}{S/\sqrt n}\ge t_L\quad\text{或}\quad\frac{\bar X-\mu_0}{S/\sqrt n}\le t_U,\quad\text{则拒绝 }H_0.$$
\sourcepageanchor{340}
若 $t_L=-t_U\ge0$，则可更简洁地写为
$$\text{若 }\frac{|\bar X-\mu_0|}{S/\sqrt n}\ge t_L,\quad\text{则拒绝 }H_0.$$
事实上，这个并交检验也是该问题的似然比检验（见\autoref{exercise:8.38}），称为双侧 $t$ 检验（two-sided $t$ test\mbox{）.}
\end{example}
原假设便于写成交集时，并交构造法很有用. 另一种方法——交并法（intersection–union method）——则可能适用于原假设便于写成并集的情形. 假设要检验的原假设为
\begin{equation}\label{eq:8.2.5}H_0:\theta\in\bigcup_{\gamma\in\Gamma}\Theta_\gamma.\end{equation}
假设对每个 $\gamma\in\Gamma$，$\{\boldsymbol x:T_\gamma(\boldsymbol x)\in R_\gamma\}$ 是检验 $H_{0\gamma}:\theta\in\Theta_\gamma$ 对 $H_{1\gamma}:\theta\in\Theta_\gamma^c$ 的拒绝域. 则检验 $H_0$ 对 $H_1$ 的交并检验（intersection–union test）的拒绝域为
\begin{equation}\label{eq:8.2.6}\index{SI0494@交并}\bigcap_{\gamma\in\Gamma}\{\boldsymbol x:T_\gamma(\boldsymbol x)\in R_\gamma\}.\end{equation}
由\eqref{eq:8.2.5}，$H_0$ 为假当且仅当所有 $H_{0\gamma}$ 都为假，所以仅当每个单独假设 $H_{0\gamma}$ 都可拒绝时，$H_0$ 才可拒绝. 同样，若单独假设的拒绝域都形如 $\{\boldsymbol x:T_\gamma(\boldsymbol x)\ge c\}$（$c$ 与 $\gamma$ 无关），检验就可大为简化. 此时 $H_0$ 的拒绝域为
$$\bigcap_{\gamma\in\Gamma}\{\boldsymbol x:T_\gamma(\boldsymbol x)\ge c\}
=\{\boldsymbol x:\inf_{\gamma\in\Gamma}T_\gamma(\boldsymbol x)\ge c\}.$$
交并检验在检验统计量 $\inf_{\gamma\in\Gamma}T_\gamma(\boldsymbol X)$ 较大时拒绝 $H_0$.
\begin{example}[验收抽样]\label{ex:8.2.9}\index{SI0002@验收抽样}\index{SI0495@伯努利}\index{SI0497@正态}
本例说明，验收抽样是交并检验一个极有用的应用领域. （关于这个问题更详细的讨论，见 \hyperref[bib:source23]{Berger (1982)}.）

评估家具装饰织物质量时，两个重要参数是平均断裂强度 $\theta_1$ 和通过燃烧性能试验的概率 $\theta_2$. 标准可能规定 $\theta_1$ 应超过 $50$ 磅，$\theta_2$ 应超过 $.95$；只有两项均达标，织物才可接受. 可以用下列假设检验建模：
$$H_0:\{\theta_1\le50\text{ 或 }\theta_2\le.95\}\quad\text{对}\quad
H_1:\{\theta_1>50\text{ 且 }\theta_2>.95\},$$
只有接受 $H_1$，这批材料才予以验收.

假设 $X_1,\ldots,X_n$ 是 $n$ 个样品的断裂强度测量值，并假定它们独立同分布于 $\mathrm n(\theta_1,\sigma^2)$. $H_{01}:\theta_1\le50$ 的似然比检验，在 $(\bar X-50)/(S/\sqrt n)>t$ 时拒绝 $H_{01}$. 又假设还获得了 $m$ 次燃烧性能试验的结果 $Y_1,\ldots,Y_m$；第 $i$ 个样品通过试验时 $Y_i=1$，否则 $Y_i=0$. 若将它们建模为独立同分布的 $\operatorname{Bernoulli}(\theta_2)$ 随机变量，则似然比检验在 $\sum_{i=1}^mY_i>b$ 时拒绝 $H_{02}:\theta_2\le.95$（见\autoref{exercise:8.3}\mbox{）.} 合起来，交并检验的拒绝域为
$$\left\{(\boldsymbol x,\boldsymbol y):\frac{\bar x-50}{s/\sqrt n}>t\text{ 且 }\sum_{i=1}^my_i>b\right\}.$$
因此，这个交并检验判定产品可接受、即 $H_1$ 为真，当且仅当它判定每个单独参数都符合相应标准、即 $H_{1i}$ 为真. 若产品质量由两个以上参数定义，可以用交并法组合各参数的单独检验，得到产品质量的总体检验.
\end{example}



\sourcepageanchor{341}
\subsection{评价检验的方法}\label{sec:8.3}
实验者在决定接受或拒绝原假设 $H_0$ 时，可能犯错误. 通常通过犯错概率评价和比较假设检验. 本节讨论如何控制这些错误概率；有时甚至可以确定哪些检验具有尽可能小的错误概率.

\subsubsection{错误概率与功效函数}\label{sec:8.3.1}
检验 $H_0:\theta\in\Theta_0$ 对 $H_1:\theta\in\Theta_0^c$，可能犯两种错误中的一种. 传统上，这两种错误被赋予不易帮助记忆的名称：第一类错误（Type I Error）和第二类错误（Type II Error\mbox{）.} 若 $\theta\in\Theta_0$，检验却错误地决定拒绝 $H_0$，就犯了第一类错误. 反之，若 $\theta\in\Theta_0^c$，检验却决定接受 $H_0$，就犯了第二类错误. \autoref{tab:8.3.1}描述了这两种不同情形.

\begin{table}[htbp]\centering
\caption{假设检验中的两类错误}\label{tab:8.3.1}\index{SI0513@第一类错误}\index{SI0514@第二类错误}
\begin{tabular}{cc|cc}\toprule
&&\multicolumn{2}{c}{决策}\\
&&接受 $H_0$&拒绝 $H_0$\\\midrule
\multirow{2}{*}{真实状态}&$H_0$&正确决策&第一类错误\\
&$H_1$&第二类错误&正确决策\\\bottomrule
\end{tabular}
\end{table}

设 $R$ 表示某检验的拒绝域. 对 $\theta\in\Theta_0$，如果 $\boldsymbol x\in R$，检验就会犯错，故第一类错误的概率为 $P_\theta(\boldsymbol X\in R)$. 对 $\theta\in\Theta_0^c$，第二类错误的概率为 $P_\theta(\boldsymbol X\in R^c)$. 从 $R$ 转到 $R^c$ 略显混乱，但注意到 $P_\theta(\boldsymbol X\in R^c)=1-P_\theta(\boldsymbol X\in R)$，便知 $\theta$ 的函数 $P_\theta(\boldsymbol X\in R)$ 包含了有关拒绝域为 $R$ 的检验的全部信息. 有
$$P_\theta(\boldsymbol X\in R)=\begin{cases}\text{第一类错误的概率},&\theta\in\Theta_0,\\1-\text{第二类错误的概率},&\theta\in\Theta_0^c.\end{cases}$$
这一考虑引出如下定义.
\begin{definition}\label{def:8.3.1}\index{SI0501@功效}
拒绝域为 $R$ 的假设检验的功效函数（power function）是 $\theta$ 的函数 $\beta(\theta)=P_\theta(\boldsymbol X\in R)$.
\end{definition}
理想的功效函数对所有 $\theta\in\Theta_0$ 为零，对所有 $\theta\in\Theta_0^c$ 为一. 除平凡情形外，这一理想无法实现. 定性地说，好的检验应对大多数 $\theta\in\Theta_0^c$ 具有接近一的功效，对大多数 $\theta\in\Theta_0$ 具有接近零的功效.
\begin{example}[二项功效函数]\label{ex:8.3.2}\index{SI0502@二项}
设 $X\sim\operatorname{binomial}(5,\theta)$. 考虑检验 $H_0:\theta\le1/2$ 对 $H_1:\theta>1/2$. 先考虑仅在观察到全部“成功”时才拒绝 $H_0$ 的检验，其功效函数为
$$\beta_1(\theta)=P_\theta(\boldsymbol X\in R)=P_\theta(X=5)=\theta^5.$$
$\beta_1(\theta)$ 的图像见\autoref{fig:8.3.1}. 考察该功效函数，可能认为虽然第一类错误概率低到可以接受（对所有 $\theta\le1/2$，$\beta_1(\theta)\le(1/2)^5=.0312$），但对大多数 $\theta>1/2$，第二类错误概率太高（$\beta_1(\theta)$ 太小\mbox{）.} 只有当 $\theta>(1/2)^{1/5}=.87$ 时，第二类错误概率才小于 $1/2$. 为降低第二类错误概率，可以考虑在 $X=3,4,5$ 时拒绝 $H_0$ 的检验，其功效函数为
$$\begin{aligned}
\beta_2(\theta)&=P_\theta(X=3,4\text{ 或 }5)\\
&=\binom53\theta^3(1-\theta)^2+\binom54\theta^4(1-\theta)^1+\binom55\theta^5(1-\theta)^0.
\end{aligned}$$
\sourcepageanchor{342}

\begin{figure}[htbp]\centering
\includegraphics[width=.75\textwidth]{assets/chapter8/fig8-3-1.pdf}
\caption{\autoref{ex:8.3.2}的功效函数.}\label{fig:8.3.1}
\end{figure}

$\beta_2(\theta)$ 也画在\autoref{fig:8.3.1}中. 可见第二个检验确实减小了第二类错误概率，因为对 $\theta>1/2$，$\beta_2(\theta)$ 更大. 但第二个检验的第一类错误概率也更大，因为对 $\theta\le1/2$，$\beta_2(\theta)$ 也更大. 若要在两者之间选择，研究者必须决定哪一种错误结构——由 $\beta_1(\theta)$ 还是由 $\beta_2(\theta)$ 描述的结构——更可接受.
\end{example}
\begin{example}[正态功效函数]\label{ex:8.3.3}\index{SI0504@正态}
设 $X_1,\ldots,X_n$ 是来自 $\mathrm n(\theta,\sigma^2)$ 总体的随机样本，$\sigma^2$ 已知. 检验 $H_0:\theta\le\theta_0$ 对 $H_1:\theta>\theta_0$ 的一种似然比检验，是当 $(\bar X-\theta_0)/(\sigma/\sqrt n)>c$ 时拒绝 $H_0$（\autoref{exercise:8.37}\mbox{）.} 常数 $c$ 可以是任意正数. 其功效函数为
$$\begin{aligned}
\beta(\theta)&=P_\theta\left(\frac{\bar X-\theta_0}{\sigma/\sqrt n}>c\right)\\
&=P_\theta\left(\frac{\bar X-\theta}{\sigma/\sqrt n}>c+\frac{\theta_0-\theta}{\sigma/\sqrt n}\right)\\
&=P\left(Z>c+\frac{\theta_0-\theta}{\sigma/\sqrt n}\right),
\end{aligned}$$
其中 $Z$ 为标准正态随机变量，因为 $(\bar X-\theta)/(\sigma/\sqrt n)\sim\mathrm n(0,1)$. 容易看出，当 $\theta$ 从 $-\infty$ 增至 $\infty$ 时，这一正态概率从零增至一. 因此 $\beta(\theta)$ 是 $\theta$ 的增函数，并且
$$\lim_{\theta\to-\infty}\beta(\theta)=0,\qquad\lim_{\theta\to\infty}\beta(\theta)=1,\qquad
\beta(\theta_0)=\alpha\ \text{若 }P(Z>c)=\alpha.$$
\autoref{fig:8.3.2}给出了 $c=1.28$ 时的 $\beta(\theta)$ 图像.
\end{example}
通常，检验的功效函数依赖样本量 $n$. 若 $n$ 可由实验者选择，考虑功效函数就可能帮助确定实验中适当的样本量.
\sourcepageanchor{343}

\begin{figure}[htbp]\centering
\includegraphics[width=.78\textwidth]{assets/chapter8/fig8-3-2.pdf}
\caption{\autoref{ex:8.3.3}的功效函数.}\label{fig:8.3.2}
\end{figure}

\begin{example}[续\autoref{ex:8.3.3}]\label{ex:8.3.4}\index{SI0510@正态}\index{SI0513@第一类错误}\index{SI0514@第二类错误}\index{SI0947@假设检验：正态}
假设实验者希望最大第一类错误概率为 $.1$；此外，希望在 $\theta\ge\theta_0+\sigma$ 时，最大第二类错误概率为 $.2$. 现在说明如何选择 $c$ 和 $n$ 达到这些目标，采用的检验是在 $(\bar X-\theta_0)/(\sigma/\sqrt n)>c$ 时拒绝 $H_0:\theta\le\theta_0$. 如上所述，这种检验的功效函数为
$$\beta(\theta)=P\left(Z>c+\frac{\theta_0-\theta}{\sigma/\sqrt n}\right).$$
由于 $\beta(\theta)$ 随 $\theta$ 递增，只要
$$\beta(\theta_0)=.1\quad\text{且}\quad\beta(\theta_0+\sigma)=.8,$$
就满足要求.

选择 $c=1.28$，无论 $n$ 为何，都有 $\beta(\theta_0)=P(Z>1.28)=.1$. 现在要选取 $n$，使 $\beta(\theta_0+\sigma)=P(Z>1.28-\sqrt n)=.8$. 因 $P(Z>-.84)=.8$，令 $1.28-\sqrt n=-.84$，解得 $n=4.49$. 当然，$n$ 必须是整数，所以取 $c=1.28$、$n=5$，即可按实验者的规定控制错误概率.\footnote{译注：这里的正态分位数和概率均经过舍入，原文用等号表示近似值. 实际上 $P(Z>1.28)\approx.10027$，略大于 $.1$；若要求严格不超过 $.1$，应取 $c=z_{.1}\approx1.28155$，此时取 $n=5$ 仍满足第二类错误的要求.}
\end{example}
样本量固定时，通常不可能把两类错误概率都任意缩小. 寻找好的检验时，通常先把考虑范围限制为能将第一类错误概率控制在指定水平的检验，再在这一类中寻找第二类错误概率尽可能小的检验. 讨论第一类错误概率的控制时，下列两个术语很有用.
\begin{definition}\label{def:8.3.5}\index{SI0511@检验大小}
对 $0\le\alpha\le1$，若功效函数为 $\beta(\theta)$ 的检验满足 $\sup_{\theta\in\Theta_0}\beta(\theta)=\alpha$，则称其为大小为 $\alpha$ 的检验（size $\alpha$ test\mbox{）.}
\end{definition}
\begin{definition}\label{def:8.3.6}\index{SI0498@检验水平}
对 $0\le\alpha\le1$，若功效函数为 $\beta(\theta)$ 的检验满足 $\sup_{\theta\in\Theta_0}\beta(\theta)\le\alpha$，则称其为水平为 $\alpha$ 的检验（level $\alpha$ test\mbox{）.}
\end{definition}
有些作者不像我们这样区分大小（size）与水平（level），有时两词互换使用. 但按本书的定义，水平为 $\alpha$ 的检验集合包含大小为 $\alpha$ 的检验集合. 而且，在复杂模型和复杂检验问题中，这一区别很重要，因为计算上往往无法构造大小恰为 $\alpha$ 的检验. 此时实验者必须满足于水平为 $\alpha$ 的检验，意识到可能需要某些折中.
\sourcepageanchor{344}
后面会看到一些例子，特别是在并交检验和交并检验中.

实验者通常指定所用检验的水平，典型选择为 $\alpha=.01,.05,.10$. 要注意，固定检验水平只控制第一类错误概率，并不控制第二类错误. 采用这种方法时，应恰当设置原假设和备择假设，使控制第一类错误概率成为最重要的事情. 例如，某实验者期待实验支持一个特定假设，但除非数据确实给出令人信服的支持，否则不愿作出断言. 可以把检验设置成：备择假设正是她期待数据支持、希望证明的假设. （在这一语境下，备择假设有时称为研究假设（research hypothesis\mbox{）.}\index{SI0479@研究}\semanticindexanchor{479:344}）通过采用 $\alpha$ 较小的水平 $\alpha$ 检验，她就能防范在研究假设为假时声称数据支持它.

\autoref{sec:8.2}的方法通常给出检验统计量和拒绝域的一般形式，却通常不会确定某一个具体检验. 例如，似然比检验（\autoref{def:8.2.1}）在 $\lambda(\boldsymbol X)\le c$ 时拒绝 $H_0$，但 $c$ 未指定，所以定义的是整类似然比检验，每个 $c$ 值对应一个. 现在将范围限制为大小 $\alpha$ 的检验，就可能从中选出一个.
\begin{example}[似然比检验的大小]\label{ex:8.3.7}\index{SI0634@指数}\index{SI0638@正态}
一般地，选择 $c$ 使 $\sup_{\theta\in\Theta_0}P_\theta(\lambda(\boldsymbol X)\le c)=\alpha$，便构造出大小为 $\alpha$ 的似然比检验. 如何确定 $c$，依具体问题而定. 例如，在\autoref{ex:8.2.2}中，$\Theta_0$ 只含单点 $\theta=\theta_0$，且在此点 $\sqrt n(\bar X-\theta_0)\sim\mathrm n(0,1)$. 因而，检验
$$\text{若 }|\bar X-\theta_0|\ge z_{\alpha/2}/\sqrt n,\quad\text{则拒绝 }H_0,$$
就是大小为 $\alpha$ 的似然比检验；其中 $z_{\alpha/2}$ 满足 $P(Z\ge z_{\alpha/2})=\alpha/2$，$Z\sim\mathrm n(0,1)$. 具体而言，这相当于取 $c=\exp(-z_{\alpha/2}^2/2)$，但这一点并不重要.

对\autoref{ex:8.2.3}的问题，原假设 $H_0:\theta\le\theta_0$ 包含不止一个点，因此寻找大小为 $\alpha$ 的似然比检验更加复杂. 检验在 $X_{(1)}\ge c$ 时拒绝 $H_0$，其中 $c$ 应使检验大小为 $\alpha$. 但若 $c=(-\log\alpha)/n+\theta_0$，则
$$P_{\theta_0}(X_{(1)}\ge c)=e^{-n(c-\theta_0)}=\alpha.$$
由于 $\theta$ 是 $X_{(1)}$ 的位置参数，
$$P_\theta(X_{(1)}\ge c)\le P_{\theta_0}(X_{(1)}\ge c),\qquad\text{对任意 }\theta\le\theta_0.$$
因此，
$$\sup_{\theta\in\Theta_0}\beta(\theta)=\sup_{\theta\le\theta_0}P_\theta(X_{(1)}\ge c)
=P_{\theta_0}(X_{(1)}\ge c)=\alpha,$$
这个 $c$ 就给出了大小为 $\alpha$ 的似然比检验.
\end{example}
\emph{记号说明：}上例用 $z_{\alpha/2}$ 表示标准正态密度曲线下右侧概率为 $\alpha/2$ 的点. 我们将普遍使用这种记法，不仅用于正态分布，也用于其他分布（为清楚起见，会按需给出定义\mbox{）.} 例如，$z_\alpha$ 满足 $P(Z>z_\alpha)=\alpha$，其中 $Z\sim\mathrm n(0,1)$；$t_{n-1,\alpha/2}$ 满足 $P(T_{n-1}>t_{n-1,\alpha/2})=\alpha/2$，其中 $T_{n-1}\sim t_{n-1}$；$\chi^2_{p,1-\alpha}$ 满足 $P(\chi_p^2>\chi^2_{p,1-\alpha})=1-\alpha$，其中 $\chi_p^2$ 是自由度为 $p$ 的卡方随机变量. $z_{\alpha/2},z_\alpha,t_{n-1,\alpha/2},\chi^2_{p,1-\alpha}$ 等点称为临界点（cutoff points\mbox{）.}\index{SI0243@临界点}\semanticindexanchor{243:344}
\begin{example}[并交检验的大小]\label{ex:8.3.8}\index{SI0525@正态}\index{SI1028@并交}
在\autoref{ex:8.2.8}中，寻找大小为 $\alpha$ 的并交检验，就是寻找常数 $t_L,t_U$，使
$$\sup_{\theta\in\Theta_0}P_\theta\left(\frac{\bar X-\mu_0}{\sqrt{S^2/n}}\ge t_L\quad\text{或}\quad\frac{\bar X-\mu_0}{\sqrt{S^2/n}}\le t_U\right)=\alpha.$$
\sourcepageanchor{345}
但对任意 $(\mu,\sigma^2)=\theta\in\Theta_0$，都有 $\mu=\mu_0$，因此 $(\bar X-\mu_0)/\sqrt{S^2/n}$ 服从自由度为 $n-1$ 的 Student $t$ 分布. 所以，只要取 $t_U=t_{n-1,1-\alpha_1}$、$t_L=t_{n-1,\alpha_2}$，其中 $\alpha_1+\alpha_2=\alpha$，就得到对所有 $\theta\in\Theta_0$ 第一类错误概率均恰为 $\alpha$ 的检验. 通常取 $t_L=-t_U=t_{n-1,\alpha/2}$.
\end{example}
除了 $\alpha$ 水平，还可能关心检验的其他特征. 例如，希望在 $\theta\in\Theta_0^c$ 时，比在 $\theta\in\Theta_0$ 时更可能拒绝 $H_0$. \autoref{fig:8.3.1}和\autoref{fig:8.3.2}的全部功效函数都有此性质，得到的检验称为无偏检验.
\begin{definition}\label{def:8.3.9}\index{SI0520@无偏}
若对每个 $\theta'\in\Theta_0^c$ 和 $\theta''\in\Theta_0$，都有 $\beta(\theta')\ge\beta(\theta'')$，则功效函数为 $\beta(\theta)$ 的检验称为无偏的（unbiased\mbox{）.}
\end{definition}
\begin{example}[\autoref{ex:8.3.3}的结论]\label{ex:8.3.10}
检验 $H_0:\theta\le\theta_0$ 对 $H_1:\theta>\theta_0$ 的似然比检验具有功效函数
$$\beta(\theta)=P\left(Z>c+\frac{\theta_0-\theta}{\sigma/\sqrt n}\right),$$
其中 $Z\sim\mathrm n(0,1)$. 由于对固定 $\theta_0$，$\beta(\theta)$ 是 $\theta$ 的增函数，故
$$\beta(\theta)>\beta(\theta_0)=\max_{t\le\theta_0}\beta(t),\qquad\text{对所有 }\theta>\theta_0,$$
所以该检验无偏.
\end{example}
多数问题有许多无偏检验（见\autoref{exercise:8.45}\mbox{）.} 同样，大小为 $\alpha$ 的检验、似然比检验等也往往有许多个. 有时施加的限制已足以把考虑范围缩减为一个检验. \autoref{ex:8.3.7}中的两个问题，都只有一个大小为 $\alpha$ 的似然比检验. 在其他情形下，仍有许多检验可选；我们只讨论了在 $T$ 较大时拒绝 $H_0$ 的那一个. 接下来将讨论从一类检验中作选择的其他准则，它们都与检验的功效函数有关.



\subsubsection{最强检验}\label{sec:8.3.2}
前几节描述了各种假设检验类. 其中一些控制第一类错误概率，例如水平为 $\alpha$ 的检验对所有 $\theta\in\Theta_0$ 都使第一类错误概率不超过 $\alpha$. 这类中的好检验还应具有较小的第二类错误概率，即在 $\theta\in\Theta_0^c$ 时功效函数较大. 若一个检验的第二类错误概率比类中所有其他检验都小，它当然有力地竞争着类中最佳检验的地位. 下列定义将这一想法形式化.
\begin{definition}\label{def:8.3.11}\index{SI0515@一致最强}
设 $\mathscr C$ 是检验 $H_0:\theta\in\Theta_0$ 对 $H_1:\theta\in\Theta_0^c$ 的一个检验类. 类 $\mathscr C$ 中功效函数为 $\beta(\theta)$ 的检验，若对每个 $\theta\in\Theta_0^c$，以及类 $\mathscr C$ 中每个检验的功效函数 $\beta'(\theta)$，都有 $\beta(\theta)\ge\beta'(\theta)$，就称为类 $\mathscr C$ 中的一致最强检验（uniformly most powerful，UMP\mbox{）.}
\end{definition}
本节中，$\mathscr C$ 取全部水平 $\alpha$ 检验构成的类，\autoref{def:8.3.11}中的检验于是称为水平 $\alpha$ 的一致最强检验. 要使这种检验有意义，对 $\mathscr C$ 的限制必须包含对第一类错误概率的某种限制. 完全不控制第一类错误概率而使第二类错误概率最小，并没有多大意义. （例如，以概率 $1$ 拒绝 $H_0$ 的检验永远不会犯第二类错误，见\autoref{exercise:8.16}.）

\autoref{def:8.3.11}要求很强，许多实际问题中不存在一致最强检验. 但在存在时，它完全可以被视为类中最佳检验. 因此，我们希望能识别出存在的一致最强检验. 下面著名的定理明确指出：当原假设和备择假设各自只包含样本的一种概率分布，即二者均为简单假设时\index{SI0480@简单}\semanticindexanchor{480:345}，哪些检验是水平 $\alpha$ 的一致最强检验.
\sourcepageanchor{346}
\begin{theorem}[Neyman–Pearson 引理]\label{thm:8.3.12}\index{SI0722@Neyman–Pearson引理}
考虑检验 $H_0:\theta=\theta_0$ 对 $H_1:\theta=\theta_1$，$\theta_i$ 对应的密度或质量函数为 $f(\boldsymbol x\mid\theta_i)$，$i=0,1$. 使用拒绝域为 $R$ 的检验，满足对某个 $k\ge0$，
\begin{equation}\label{eq:8.3.1}
\begin{aligned}
\boldsymbol x\in R&\quad\text{若 }f(\boldsymbol x\mid\theta_1)>kf(\boldsymbol x\mid\theta_0),\\
\text{且}\qquad\boldsymbol x\in R^c&\quad\text{若 }f(\boldsymbol x\mid\theta_1)<kf(\boldsymbol x\mid\theta_0),
\end{aligned}
\end{equation}
以及
\begin{equation}\label{eq:8.3.2}\alpha=P_{\theta_0}(\boldsymbol X\in R).\end{equation}
则
\begin{enumerate}
\item\phantomsection\label{itm:thm:8.3.12:a}（充分性）任何满足\eqref{eq:8.3.1}和\eqref{eq:8.3.2}的检验，都是水平 $\alpha$ 的一致最强检验.
\item\phantomsection\label{itm:thm:8.3.12:b}（必要性）若存在满足\eqref{eq:8.3.1}、\eqref{eq:8.3.2}且 $k>0$ 的检验，则每个水平 $\alpha$ 的一致最强检验大小都为 $\alpha$（满足\eqref{eq:8.3.2}）；并且每个水平 $\alpha$ 的一致最强检验都满足\eqref{eq:8.3.1}，至多在某个满足 $P_{\theta_0}(\boldsymbol X\in A)=P_{\theta_1}(\boldsymbol X\in A)=0$ 的集合 $A$ 上例外.
\end{enumerate}
\end{theorem}
\begin{proof}
我们证明 $f(\boldsymbol x\mid\theta_0)$ 和 $f(\boldsymbol x\mid\theta_1)$ 是连续随机变量密度的情形. 离散情形的证明只需把积分换成求和（见\autoref{exercise:8.21}\mbox{）.}

首先注意，任何满足\eqref{eq:8.3.2}的检验，大小都为 $\alpha$，从而水平也为 $\alpha$，因为 $\Theta_0$ 只有一个点，故 $\sup_{\theta\in\Theta_0}P_\theta(\boldsymbol X\in R)=P_{\theta_0}(\boldsymbol X\in R)=\alpha$.

为简化记号，定义检验函数（test function）\index{SI1032@检验函数}\semanticindexanchor{1032:346}：它是样本空间上的函数，在 $\boldsymbol x\in R$ 时为 $1$，在 $\boldsymbol x\in R^c$ 时为 $0$，即拒绝域的示性函数. 设 $\phi(\boldsymbol x)$ 是满足\eqref{eq:8.3.1}与\eqref{eq:8.3.2}的检验的检验函数，$\phi'(\boldsymbol x)$ 是任一其他水平 $\alpha$ 检验的检验函数，$\beta(\theta)$ 与 $\beta'(\theta)$ 分别为二者的功效函数. 因为 $0\le\phi'(\boldsymbol x)\le1$，\eqref{eq:8.3.1}蕴含对每个 $\boldsymbol x$，
$$[\phi(\boldsymbol x)-\phi'(\boldsymbol x)][f(\boldsymbol x\mid\theta_1)-kf(\boldsymbol x\mid\theta_0)]\ge0$$
（因为当 $f(\boldsymbol x\mid\theta_1)>kf(\boldsymbol x\mid\theta_0)$ 时 $\phi=1$，当反向严格不等式成立时 $\phi=0$\mbox{）.} 因此，
\begin{equation}\label{eq:8.3.3}
\begin{aligned}
0&\le\int[\phi(\boldsymbol x)-\phi'(\boldsymbol x)][f(\boldsymbol x\mid\theta_1)-kf(\boldsymbol x\mid\theta_0)]\,d\boldsymbol x\\
&=\beta(\theta_1)-\beta'(\theta_1)-k[\beta(\theta_0)-\beta'(\theta_0)].
\end{aligned}
\end{equation}
为证明\hyperref[itm:thm:8.3.12:a]{(a)}，注意 $\phi'$ 的水平为 $\alpha$ 而 $\phi$ 的大小为 $\alpha$，故 $\beta(\theta_0)-\beta'(\theta_0)=\alpha-\beta'(\theta_0)\ge0$. 因而\eqref{eq:8.3.3}与 $k\ge0$ 蕴含
$$0\le\beta(\theta_1)-\beta'(\theta_1)-k[\beta(\theta_0)-\beta'(\theta_0)]\le\beta(\theta_1)-\beta'(\theta_1),$$
表明 $\beta(\theta_1)\ge\beta'(\theta_1)$，即 $\phi$ 的功效不低于 $\phi'$. 因 $\phi'$ 是任意水平 $\alpha$ 的检验，$\theta_1$ 又是 $\Theta_0^c$ 中的唯一点，所以 $\phi$ 是水平 $\alpha$ 的一致最强检验.

为证明\hyperref[itm:thm:8.3.12:b]{(b)}，现令 $\phi'$ 为任一水平 $\alpha$ 的一致最强检验的检验函数. 由\hyperref[itm:thm:8.3.12:a]{(a)}，满足\eqref{eq:8.3.1}和\eqref{eq:8.3.2}的 $\phi$ 也是水平 $\alpha$ 的一致最强检验，故 $\beta(\theta_1)=\beta'(\theta_1)$. 这一事实、\eqref{eq:8.3.3}以及 $k>0$ 蕴含
$$\alpha-\beta'(\theta_0)=\beta(\theta_0)-\beta'(\theta_0)\le0.$$
由于 $\phi'$ 的水平为 $\alpha$，又有 $\beta'(\theta_0)\le\alpha$. 所以 $\beta'(\theta_0)=\alpha$，即 $\phi'$ 的大小为 $\alpha$，且\eqref{eq:8.3.3}此时取等号. 然而，非负被积函数 $[\phi(\boldsymbol x)-\phi'(\boldsymbol x)][f(\boldsymbol x\mid\theta_1)-kf(\boldsymbol x\mid\theta_0)]$ 的积分只有在 $\phi'$ 满足\eqref{eq:8.3.1}时才为零，至多在某个满足 $\int_Af(\boldsymbol x\mid\theta_i)\,d\boldsymbol x=0$ 的集合 $A$ 上例外. 这就证明了\hyperref[itm:thm:8.3.12:b]{(b)}的最后一个断言.
\end{proof}
\sourcepageanchor{347}
下面的推论把 Neyman–Pearson 引理与充分性联系起来.
\begin{corollary}\label{cor:8.3.13}\index{SI0512@充分性}
考虑\autoref{thm:8.3.12}中的假设检验问题. 设 $T(\boldsymbol X)$ 是 $\theta$ 的充分统计量，$g(t\mid\theta_i)$ 是 $T$ 在 $\theta_i$ 下的密度或质量函数，$i=0,1$. 则基于 $T$、拒绝域为 $S$（$T$ 的样本空间的一个子集）的任一检验，只要对某个 $k\ge0$ 满足
\begin{equation}\label{eq:8.3.4}
\begin{aligned}t\in S&\quad\text{若 }g(t\mid\theta_1)>kg(t\mid\theta_0),\\
\text{且}\qquad t\in S^c&\quad\text{若 }g(t\mid\theta_1)<kg(t\mid\theta_0),\end{aligned}
\end{equation}
就是水平 $\alpha$ 的一致最强检验，其中
\begin{equation}\label{eq:8.3.5}\alpha=P_{\theta_0}(T\in S).\end{equation}
\end{corollary}
\begin{proof}
用原样本 $\boldsymbol X$ 表示，基于 $T$ 的检验的拒绝域为 $R=\{\boldsymbol x:T(\boldsymbol x)\in S\}$. 由因子分解定理，可将 $\boldsymbol X$ 的密度或质量函数写为 $f(\boldsymbol x\mid\theta_i)=g(T(\boldsymbol x)\mid\theta_i)h(\boldsymbol x)$，$i=0,1$，其中 $h(\boldsymbol x)$ 非负. 将\eqref{eq:8.3.4}的不等式乘以这一非负函数，就看到 $R$ 满足
$$\boldsymbol x\in R\quad\text{若 }f(\boldsymbol x\mid\theta_1)=g(T(\boldsymbol x)\mid\theta_1)h(\boldsymbol x)>kg(T(\boldsymbol x)\mid\theta_0)h(\boldsymbol x)=kf(\boldsymbol x\mid\theta_0),$$
以及
$$\boldsymbol x\in R^c\quad\text{若 }f(\boldsymbol x\mid\theta_1)=g(T(\boldsymbol x)\mid\theta_1)h(\boldsymbol x)<kg(T(\boldsymbol x)\mid\theta_0)h(\boldsymbol x)=kf(\boldsymbol x\mid\theta_0).$$
另外，由\eqref{eq:8.3.5}，
$$P_{\theta_0}(\boldsymbol X\in R)=P_{\theta_0}(T(\boldsymbol X)\in S)=\alpha.$$
所以，由 Neyman–Pearson 引理的充分性部分，基于 $T$ 的检验是水平 $\alpha$ 的一致最强检验.
\end{proof}
推导满足\eqref{eq:8.3.1}或\eqref{eq:8.3.4}、从而为水平 $\alpha$ 的一致最强检验时，通常把不等式改写为 $f(\boldsymbol x\mid\theta_1)/f(\boldsymbol x\mid\theta_0)>k$ 更方便（须留意除数为零的情形\mbox{）.} 以下例子采用这种方法.
\begin{example}[二项总体的一致最强检验]\label{ex:8.3.14}\index{SI0516@二项}
设 $X\sim\operatorname{binomial}(2,\theta)$. 要检验 $H_0:\theta=1/2$ 对 $H_1:\theta=3/4$. 计算概率质量函数之比，得到
$$\frac{f(0\mid\theta=3/4)}{f(0\mid\theta=1/2)}=\frac14,\qquad
\frac{f(1\mid\theta=3/4)}{f(1\mid\theta=1/2)}=\frac34,\qquad
\frac{f(2\mid\theta=3/4)}{f(2\mid\theta=1/2)}=\frac94.$$
选择 $3/4<k<9/4$，Neyman–Pearson 引理说明，在 $X=2$ 时拒绝 $H_0$ 的检验是水平 $\alpha=P(X=2\mid\theta=1/2)=1/4$ 的一致最强检验. 选择 $1/4<k<3/4$，引理说明，在 $X=1$ 或 $2$ 时拒绝 $H_0$ 的检验是水平 $\alpha=P(X=1\text{ 或 }2\mid\theta=1/2)=3/4$ 的一致最强检验. 取 $k<1/4$ 或 $k>9/4$，分别得到水平 $\alpha=1$ 或 $\alpha=0$ 的一致最强检验.

注意，若 $k=3/4$，\eqref{eq:8.3.1}要求在样本点 $x=2$ 时拒绝 $H_0$，在 $x=0$ 时接受 $H_0$，但未决定 $x=1$ 时采取何种行动. 若在 $x=1$ 时接受 $H_0$，就得到上面水平 $\alpha=1/4$ 的一致最强检验；若在 $x=1$ 时拒绝 $H_0$，就得到水平 $\alpha=3/4$ 的一致最强检验.
\end{example}
\autoref{ex:8.3.14}还说明，在离散分布情形，检验所能达到的 $\alpha$ 水平取决于具体的概率质量函数. （连续情形没有这种问题，任意 $\alpha$ 水平都可达到.）
\sourcepageanchor{348}
\begin{example}[正态总体的一致最强检验]\label{ex:8.3.15}\index{SI0517@正态}
设 $X_1,\ldots,X_n$ 是来自 $\mathrm n(\theta,\sigma^2)$ 总体的随机样本，$\sigma^2$ 已知. 样本均值 $\bar X$ 是 $\theta$ 的充分统计量. 考虑检验 $H_0:\theta=\theta_0$ 对 $H_1:\theta=\theta_1$，其中 $\theta_0>\theta_1$. \eqref{eq:8.3.4}中的不等式 $g(\bar x\mid\theta_1)>kg(\bar x\mid\theta_0)$ 等价于
$$\bar x<\frac{(2\sigma^2\log k)/n-\theta_0^2+\theta_1^2}{2(\theta_1-\theta_0)}.$$
这里利用了 $\theta_1-\theta_0<0$. 随着 $k$ 从 $0$ 增至 $\infty$，右端从 $-\infty$ 增至 $\infty$.\footnote{译注：按所写公式，分母 $2(\theta_1-\theta_0)<0$，故右端实际上随 $k$ 增大而从 $+\infty$ 减至 $-\infty$. 原文方向颠倒，但右端能遍历所有实数及后续检验结论不受影响.} 因此，由\autoref{cor:8.3.13}，拒绝域为 $\bar x<c$ 的检验是水平 $\alpha$ 的一致最强检验，其中 $\alpha=P_{\theta_0}(\bar X<c)$. 若指定某个 $\alpha$，则一致最强检验在 $\bar X<c=-\sigma z_\alpha/\sqrt n+\theta_0$ 时拒绝 $H_0$. 如此选择 $c$ 保证\eqref{eq:8.3.5}成立.
\end{example}
像 Neyman–Pearson 引理中的 $H_0,H_1$ 那样，仅为样本 $\boldsymbol X$ 指定一种可能分布的假设，称为简单假设（simple hypotheses\mbox{）.}\index{SI0480@简单}\semanticindexanchor{480:348} 多数实际问题中，所关心的假设为样本指定不止一种可能分布，这种假设称为复合假设（composite hypotheses\mbox{）.}\index{SI0476@复合}\semanticindexanchor{476:348} 由于\autoref{def:8.3.11}要求一致最强检验对每个单独的 $\theta\in\Theta_0^c$ 都最强，因此可以用 Neyman–Pearson 引理寻找复合假设问题中的一致最强检验.

特别地，断言一元参数较大（如 $H:\theta\ge\theta_0$）或较小（如 $H:\theta<\theta_0$）的假设，称为单侧假设（one-sided hypotheses\mbox{）.}\index{SI0478@单侧}\semanticindexanchor{478:348} 断言参数或者较大或者较小（如 $H:\theta\ne\theta_0$）的假设，称为双侧假设（two-sided hypotheses\mbox{）.}\index{SI0481@双侧}\semanticindexanchor{481:348} 一大类具有水平 $\alpha$ 一致最强检验的问题，涉及单侧假设以及具有单调似然比性质的密度或质量函数.
\begin{definition}\label{def:8.3.16}\index{SI0700@单调似然比}
对一元随机变量 $T$ 及实值参数 $\theta$，若密度或质量函数族 $\{g(t\mid\theta):\theta\in\Theta\}$ 满足：对每个 $\theta_2>\theta_1$，比值 $g(t\mid\theta_2)/g(t\mid\theta_1)$ 在 $\{t:g(t\mid\theta_1)>0\text{ 或 }g(t\mid\theta_2)>0\}$ 上是 $t$ 的单调函数（非增或非减），则称该族具有单调似然比（monotone likelihood ratio，MLR\mbox{）.} 注意，对 $c>0$，规定 $c/0=\infty$.
\end{definition}
许多常见分布族具有单调似然比，例如正态族（方差已知、均值未知）、泊松族和二项族. 事实上，若 $w(\theta)$ 是非减函数，则任意正则指数族 $g(t\mid\theta)=h(t)c(\theta)e^{w(\theta)t}$ 都具有单调似然比（见\autoref{exercise:8.25}\mbox{）.}
\begin{theorem}[Karlin–Rubin]\label{thm:8.3.17}\index{SI0594@Karlin–Rubin定理}
考虑检验 $H_0:\theta\le\theta_0$ 对 $H_1:\theta>\theta_0$. 设 $T$ 是 $\theta$ 的充分统计量，且 $T$ 的密度或质量函数族 $\{g(t\mid\theta):\theta\in\Theta\}$ 具有单调似然比. 则对任意 $t_0$，当且仅当 $T>t_0$ 时拒绝 $H_0$ 的检验，是水平 $\alpha=P_{\theta_0}(T>t_0)$ 的一致最强检验.\footnote{译注：本定理和下述证明实际需要似然比关于 $t$ 非减，而前面的定义同时允许非增. 若似然比非增，拒绝域方向应相应反转；例如 $X\sim\mathrm n(\theta,1)$ 而 $T=-X$ 时，似然比关于 $T$ 非增，$P_\theta(T>t_0)$ 随 $\theta$ 下降，不能直接套用本述. 正文保留原假设.}
\end{theorem}
\begin{proof}
设 $\beta(\theta)=P_\theta(T>t_0)$ 为检验的功效函数. 固定 $\theta'>\theta_0$，考虑检验 $H'_0:\theta=\theta_0$ 对 $H'_1:\theta=\theta'$. 由于 $T$ 的密度或质量函数族具有单调似然比，$\beta(\theta)$ 非减（\autoref{exercise:8.34}），因此
\begin{enumerate}[label=\roman*.]
\item\phantomsection\label{itm:proof:8.3.17:i} $\sup_{\theta\le\theta_0}\beta(\theta)=\beta(\theta_0)=\alpha$，故这是一个水平 $\alpha$ 的检验.
\item\phantomsection\label{itm:proof:8.3.17:ii} 定义
$$k'=\inf_{t\in\mathscr T}\frac{g(t\mid\theta')}{g(t\mid\theta_0)},$$
其中 $\mathscr T=\{t:t>t_0\text{ 且 }g(t\mid\theta')>0\text{ 或 }g(t\mid\theta_0)>0\}$，则\footnote{译注：这里的 $\mathscr T$ 按原文意思应读作 $t>t_0$ 且两个密度至少一个为正. 此外，“$\Longleftrightarrow$”及右端严格“$>$”一般不成立：离散分布可能在最小拒绝点达到该下确界，非严格单调的比值也可能有平台. 证明只需核验 Neyman–Pearson 的两个严格不等式所规定的方向，并在比值等于临界值处采用所需的拒绝规则；还应在同一自变量 $t=T$ 下读该式.}
$$T>t_0\ \Longleftrightarrow\ \frac{g(t\mid\theta')}{g(t\mid\theta_0)}>k'.$$
\end{enumerate}

结合\autoref{cor:8.3.13}，\hyperref[itm:proof:8.3.17:i]{(i)} 与\hyperref[itm:proof:8.3.17:ii]{(ii)} 蕴含 $\beta(\theta')\ge\beta^*(\theta')$，其中 $\beta^*(\theta)$ 是 $H'_0$ 的任一其他水平 $\alpha$ 检验的功效函数，即满足 $\beta(\theta_0)\le\alpha$ 的任一检验.\footnote{译注：此处原文写 $\beta(\theta_0)\le\alpha$，依上下文正在描述其他检验，宜为 $\beta^*(\theta_0)\le\alpha$；保留原记号.} 然而，$H_0$ 的任一水平 $\alpha$ 检验都满足 $\beta^*(\theta_0)\le\sup_{\theta\in\Theta_0}\beta^*(\theta)\le\alpha$. 因此对 $H_0$ 的任一水平 $\alpha$ 检验都有 $\beta(\theta')\ge\beta^*(\theta')$. 由于 $\theta'$ 任意，原检验是水平 $\alpha$ 的一致最强检验.
\end{proof}
\sourcepageanchor{349}
类似论证表明，在\autoref{thm:8.3.17}的条件下，当且仅当 $T<t_0$ 时拒绝 $H_0:\theta\ge\theta_0$、支持 $H_1:\theta<\theta_0$ 的检验，是水平 $\alpha=P_{\theta_0}(T<t_0)$ 的一致最强检验.
\begin{example}[续\autoref{ex:8.3.15}]\label{ex:8.3.18}\index{SI0517@正态}
考虑检验 $H'_0:\theta\ge\theta_0$ 对 $H'_1:\theta<\theta_0$，采用在
$$\bar X<-\frac{\sigma z_\alpha}{\sqrt n}+\theta_0$$
时拒绝 $H'_0$ 的检验. 因 $\bar X$ 充分，且其分布具有单调似然比（\autoref{exercise:8.25}），由\autoref{thm:8.3.17}可知，这是该问题中水平 $\alpha$ 的一致最强检验.

该检验的功效函数
$$\beta(\theta)=P_\theta\left(\bar X<-\frac{\sigma z_\alpha}{\sqrt n}+\theta_0\right)$$
是 $\theta$ 的减函数（$\theta$ 是 $\bar X$ 分布的位置参数），因此 $\alpha$ 由 $\sup_{\theta\ge\theta_0}\beta(\theta)=\beta(\theta_0)=\alpha$ 给出.
\end{example}
虽然多数实验者若知道一个水平 $\alpha$ 的一致最强检验，就会选择使用它，但遗憾的是许多问题中并不存在这样的检验. 原因是水平 $\alpha$ 的检验类太大，没有一个检验在功效上优于所有其他检验. 此时，继续寻找好检验的常见方法是考虑这一检验类的某个子集，尝试在子集中找到一致最强检验. 这应使人想起\autoref{sec:7}：为研究最优性，我们把范围限制在无偏点估计量中. 下面说明，把范围限制在无偏检验组成的子集，如何帮助找到最佳检验.

首先考虑一个例子，说明水平 $\alpha$ 的一致最强检验不存在的一种典型情况.
\begin{example}[一致最强检验的不存在性]\label{ex:8.3.19}
设 $X_1,\ldots,X_n$ 独立同分布于 $\mathrm n(\theta,\sigma^2)$，$\sigma^2$ 已知. 考虑检验 $H_0:\theta=\theta_0$ 对 $H_1:\theta\ne\theta_0$. 对指定的 $\alpha$，本问题中水平 $\alpha$ 的检验就是满足
\begin{equation}\label{eq:8.3.6}P_{\theta_0}(\text{拒绝 }H_0)\le\alpha\end{equation}
的任一检验.

考虑备择参数点 $\theta_1<\theta_0$. \autoref{ex:8.3.18}的分析表明，在满足\eqref{eq:8.3.6}的全部检验中，若 $\bar X<-\sigma z_\alpha/\sqrt n+\theta_0$ 则拒绝 $H_0$ 的检验，在 $\theta_1$ 处有最高功效. 称其为检验1. 而且，由 Neyman–Pearson 引理的\hyperref[itm:thm:8.3.12:b]{(b)（必要性）}，任何在 $\theta_1$ 处达到同样高功效的其他水平 $\alpha$ 检验，必须具有与检验1相同的拒绝域，至多相差一个满足 $\int_A f(\boldsymbol x\mid\theta_i)\,d\boldsymbol x=0$ 的集合 $A$. 所以，若本问题存在水平 $\alpha$ 的一致最强检验，它必须是检验1，因为没有其他检验在 $\theta_1$ 处具有同样高的功效.

现在考虑检验2：若 $\bar X>\sigma z_\alpha/\sqrt n+\theta_0$ 则拒绝 $H_0$. 检验2的水平也是 $\alpha$. 记 $\beta_i(\theta)$ 为检验 $i$ 的功效函数. 对任意 $\theta_2>\theta_0$，
$$\begin{aligned}
\beta_2(\theta_2)&=P_{\theta_2}\left(\bar X>\frac{\sigma z_\alpha}{\sqrt n}+\theta_0\right)\\
&=P_{\theta_2}\left(\frac{\bar X-\theta_2}{\sigma/\sqrt n}>z_\alpha+\frac{\theta_0-\theta_2}{\sigma/\sqrt n}\right)\\
&>P(Z>z_\alpha)&&\text{（$Z\sim\mathrm n(0,1)$，$\theta_0-\theta_2<0$）}\\
&=P(Z<-z_\alpha)\\
&>P_{\theta_2}\left(\frac{\bar X-\theta_2}{\sigma/\sqrt n}<-z_\alpha+\frac{\theta_0-\theta_2}{\sigma/\sqrt n}\right)&&\text{（仍因 $\theta_0-\theta_2<0$）}.
\end{aligned}$$
\sourcepageanchor{350}

\begin{figure}[htbp]\centering
\includegraphics[width=.83\textwidth]{assets/chapter8/fig8-3-3.pdf}
\caption{\autoref{ex:8.3.19}中三个检验的功效函数. $\beta_3(\theta)$ 是一个无偏的水平 $\alpha=.05$ 检验的功效函数.}\label{fig:8.3.3}
\end{figure}

上式最后一项等于
$$P_{\theta_2}\left(\bar X<-\frac{\sigma z_\alpha}{\sqrt n}+\theta_0\right)=\beta_1(\theta_2).$$
所以检验1不是水平 $\alpha$ 的一致最强检验，因为检验2在 $\theta_2$ 处的功效更高. 前面又已证明，如果存在水平 $\alpha$ 的一致最强检验，它必为检验1. 因此本问题不存在水平 $\alpha$ 的一致最强检验.\footnote{译注：这里的不存在性结论针对非平凡水平 $0<\alpha<1$，也是上述有限正态临界点论证的适用范围. 端点 $\alpha=0$ 或 $1$ 有相应的平凡一致最强检验.}
\end{example}
\autoref{ex:8.3.19}再次说明 Neyman–Pearson 引理的用途. 先前用其充分性部分构造水平 $\alpha$ 的一致最强检验；要说明这种检验不存在，则使用其必要性部分.
\begin{example}[无偏检验]\label{ex:8.3.20}
当全部检验类中没有水平 $\alpha$ 的一致最强检验时，可以尝试在无偏检验类中寻找. 检验3当且仅当
$$\bar X>\sigma z_{\alpha/2}/\sqrt n+\theta_0\quad\text{或}\quad\bar X<-\sigma z_{\alpha/2}/\sqrt n+\theta_0$$
时，拒绝 $H_0:\theta=\theta_0$、支持 $H_1:\theta\ne\theta_0$. 其功效函数 $\beta_3(\theta)$ 与\autoref{ex:8.3.19}中的 $\beta_1(\theta),\beta_2(\theta)$ 一起画在\autoref{fig:8.3.3}中. 检验3事实上是水平 $\alpha$ 的一致最强无偏检验，即在无偏检验类中一致最强.

注意，虽然检验1和检验2在某些参数点的功效略高于检验3，但在其他参数点，检验3的功效比检验1和检验2高得多. 例如，$\beta_3(\theta_2)$ 接近一，而 $\beta_1(\theta_2)$ 接近零. 若希望在 $\theta$ 偏大或偏小时都拒绝 $H_0$，\autoref{fig:8.3.3}表明，整体而言检验3优于检验1或检验2.
\end{example}



\subsubsection{并交检验与交并检验的大小}\label{sec:8.3.3}
由于构造方式简单，并交检验（UIT）与交并检验（IUT）的大小，常可用其他一些检验的大小给出上界. 当希望得到水平 $\alpha$ 的检验、却难以求出 UIT 或 IUT 的大小时，这类界很有用.
\sourcepageanchor{351}
本节讨论这些界，并给出界可达到的例子，即检验的大小恰好等于所给界.

先考虑并交检验. 回顾，此时检验的原假设形如 $H_0:\theta\in\Theta_0$，其中 $\Theta_0=\bigcap_{\gamma\in\Gamma}\Theta_\gamma$. 具体地，设 $\lambda_\gamma(\boldsymbol x)$ 是检验 $H_{0\gamma}:\theta\in\Theta_\gamma$ 对 $H_{1\gamma}:\theta\in\Theta_\gamma^c$ 的似然比统计量，$\lambda(\boldsymbol x)$ 是检验 $H_0:\theta\in\Theta_0$ 对 $H_1:\theta\in\Theta_0^c$ 的似然比统计量. 则总体似然比检验与基于 $\lambda_\gamma(\boldsymbol x)$ 的并交检验有如下关系.
\begin{theorem}\label{thm:8.3.21}\index{SI0524@水平}
考虑检验 $H_0:\theta\in\Theta_0$ 对 $H_1:\theta\in\Theta_0^c$，其中 $\Theta_0=\bigcap_{\gamma\in\Gamma}\Theta_\gamma$，$\lambda_\gamma(\boldsymbol x)$ 如上段定义. 令 $T(\boldsymbol x)=\inf_{\gamma\in\Gamma}\lambda_\gamma(\boldsymbol x)$，构造拒绝域为
$$\{\boldsymbol x:\lambda_\gamma(\boldsymbol x)<c\text{ 对某个 }\gamma\in\Gamma\}
=\{\boldsymbol x:T(\boldsymbol x)<c\}$$
的并交检验. 再考虑拒绝域为 $\{\boldsymbol x:\lambda(\boldsymbol x)<c\}$ 的通常似然比检验. 则
\begin{enumerate}
\item\phantomsection\label{itm:thm:8.3.21:a} 对每个 $\boldsymbol x$，$T(\boldsymbol x)\ge\lambda(\boldsymbol x)$；
\item\phantomsection\label{itm:thm:8.3.21:b} 若 $\beta_T(\theta)$ 和 $\beta_\lambda(\theta)$ 分别为基于 $T$ 与 $\lambda$ 的检验的功效函数，则对每个 $\theta\in\Theta$，$\beta_T(\theta)\le\beta_\lambda(\theta)$；
\item\phantomsection\label{itm:thm:8.3.21:c} 若似然比检验的水平为 $\alpha$，则并交检验的水平也为 $\alpha$.
\end{enumerate}
\end{theorem}
\begin{proof}
由于对任意 $\gamma$，$\Theta_0=\bigcap_{\gamma\in\Gamma}\Theta_\gamma\subset\Theta_\gamma$，由\autoref{def:8.2.1}可见，对任意 $\boldsymbol x$，
$$\lambda_\gamma(\boldsymbol x)\ge\lambda(\boldsymbol x),\qquad\text{对每个 }\gamma\in\Gamma,$$
因为单独的 $\lambda_\gamma$ 的最大化范围更大. 所以 $T(\boldsymbol x)=\inf_{\gamma\in\Gamma}\lambda_\gamma(\boldsymbol x)\ge\lambda(\boldsymbol x)$，证明了\hyperref[itm:thm:8.3.21:a]{(a)}. 由\hyperref[itm:thm:8.3.21:a]{(a)}，$\{\boldsymbol x:T(\boldsymbol x)<c\}\subset\{\boldsymbol x:\lambda(\boldsymbol x)<c\}$，故
$$\beta_T(\theta)=P_\theta(T(\boldsymbol X)<c)\le P_\theta(\lambda(\boldsymbol X)<c)=\beta_\lambda(\theta),$$
证明了\hyperref[itm:thm:8.3.21:b]{(b)}. 因\hyperref[itm:thm:8.3.21:b]{(b)}对每个 $\theta$ 都成立，$\sup_{\theta\in\Theta_0}\beta_T(\theta)\le\sup_{\theta\in\Theta_0}\beta_\lambda(\theta)\le\alpha$，所以\hyperref[itm:thm:8.3.21:c]{(c)}成立.
\end{proof}
\begin{example}[一种等价性]\label{ex:8.3.22}
在某些情形下，\autoref{thm:8.3.21}中的 $T(\boldsymbol x)=\lambda(\boldsymbol x)$，由单独似然比检验构成的并交检验，与总体似然比检验相同. \autoref{ex:8.2.8}就是这种情况：由两个单侧 $t$ 检验构成的并交检验，等价于双侧似然比检验.
\end{example}
既然\autoref{thm:8.3.21}中似然比检验的功效处处不低于并交检验，可能会问为什么还要使用并交检验. 一个原因是，对每个 $\theta\in\Theta_0$，并交检验的第一类错误概率更小. 此外，若 $H_0$ 被拒绝，我们可能希望查看各个 $H_{0\gamma}$ 的单独检验，以了解原因. 至今尚未讨论关于单独 $H_{0\gamma}$ 的推断；采用这种推断程序之前，必须考察其错误概率. 但显然，逐个查看 $H_{0\gamma}$，有可能比只查看总体似然比检验获得更多信息.

现在研究交并检验的大小. 交并检验大小的一个简单界，与构造它所用的各个单独检验的大小有关. 回顾，此时原假设可写成并集，即检验
$$H_0:\theta\in\Theta_0\quad\text{对}\quad H_1:\theta\in\Theta_0^c,\qquad
\Theta_0=\bigcup_{\gamma\in\Gamma}\Theta_\gamma.$$
交并检验的拒绝域形如 $R=\bigcap_{\gamma\in\Gamma}R_\gamma$，其中 $R_\gamma$ 是检验 $H_{0\gamma}:\theta\in\Theta_\gamma$ 的拒绝域.
\begin{theorem}\label{thm:8.3.23}\index{SI0496@水平}
设拒绝域为 $R_\gamma$ 的 $H_{0\gamma}$ 检验的大小为 $\alpha_\gamma$. 则拒绝域为 $R=\bigcap_{\gamma\in\Gamma}R_\gamma$ 的交并检验，是水平 $\alpha=\sup_{\gamma\in\Gamma}\alpha_\gamma$ 的检验.
\end{theorem}
\begin{proof}
取 $\theta\in\Theta_0$，则对某个 $\gamma$ 有 $\theta\in\Theta_\gamma$，并且
$$P_\theta(\boldsymbol X\in R)\le P_\theta(\boldsymbol X\in R_\gamma)\le\alpha_\gamma\le\alpha.$$
由于 $\theta\in\Theta_0$ 任意，交并检验的水平为 $\alpha$.
\end{proof}
\sourcepageanchor{352}
通常，单独拒绝域 $R_\gamma$ 被选成对所有 $\gamma$ 都有 $\alpha_\gamma=\alpha$. 此时\autoref{thm:8.3.23}说明，所得交并检验的水平为 $\alpha$.

给出交并检验大小上界的\autoref{thm:8.3.23}，比给出并交检验大小上界的\autoref{thm:8.3.21}更有用一些. 后者只适用于由似然比检验构造的并交检验，前者则适用于任意交并检验.

\autoref{thm:8.3.21}的界是似然比检验的大小，在复杂问题中可能难以计算. 但在\autoref{thm:8.3.23}中，求上界并不必使用似然比检验：任何大小 $\alpha_\gamma$ 已知的 $H_{0\gamma}$ 检验都可使用，交并检验大小的上界便用这些已知的 $\alpha_\gamma$，$\gamma\in\Gamma$ 表示.

\autoref{thm:8.3.23}的交并检验水平为 $\alpha$，但其大小可能远小于 $\alpha$，即检验可能很保守. 下列定理给出大小恰为 $\alpha$、因而检验不过于保守的条件.
\begin{theorem}\label{thm:8.3.24}
考虑检验 $H_0:\theta\in\bigcup_{j=1}^k\Theta_j$，其中 $k$ 为有限正整数. 对每个 $j=1,\ldots,k$，设 $R_j$ 为 $H_{0j}$ 的某水平 $\alpha$ 检验的拒绝域. 假设对某个 $i=1,\ldots,k$，存在参数点列 $\theta_l\in\Theta_i$，$l=1,2,\ldots$，满足
\begin{enumerate}[label=\roman*.]
\item\phantomsection\label{itm:thm:8.3.24:i} $\lim_{l\to\infty}P_{\theta_l}(\boldsymbol X\in R_i)=\alpha$；
\item\phantomsection\label{itm:thm:8.3.24:ii} 对每个 $j=1,\ldots,k$，$j\ne i$，$\lim_{l\to\infty}P_{\theta_l}(\boldsymbol X\in R_j)=1$.
\end{enumerate}
则拒绝域为 $R=\bigcap_{j=1}^kR_j$ 的交并检验，大小为 $\alpha$.
\end{theorem}
\begin{proof}
由\autoref{thm:8.3.23}，$R$ 是水平 $\alpha$ 的检验，即
\begin{equation}\label{eq:8.3.7}\sup_{\theta\in\Theta_0}P_\theta(\boldsymbol X\in R)\le\alpha.\end{equation}
但由于所有参数点 $\theta_l$ 都满足 $\theta_l\in\Theta_i\subset\Theta_0$，
$$\begin{aligned}
\sup_{\theta\in\Theta_0}P_\theta(\boldsymbol X\in R)&\ge\lim_{l\to\infty}P_{\theta_l}(\boldsymbol X\in R)\\
&=\lim_{l\to\infty}P_{\theta_l}\left(\boldsymbol X\in\bigcap_{j=1}^kR_j\right)\\
&\ge\lim_{l\to\infty}\left[\sum_{j=1}^kP_{\theta_l}(\boldsymbol X\in R_j)-(k-1)\right]&&\text{（Bonferroni 不等式）}\\
&=(k-1)+\alpha-(k-1)&&\text{（由\hyperref[itm:thm:8.3.24:i]{(i)}与\hyperref[itm:thm:8.3.24:ii]{(ii)}）}\\
&=\alpha.
\end{aligned}$$
结合\eqref{eq:8.3.7}，可知检验的大小恰为 $\alpha$.
\end{proof}
\begin{example}[交并检验]\label{ex:8.3.25}\index{SI0495@伯努利}\index{SI0497@正态}
在\autoref{ex:8.2.9}中，令 $n=m=58$，$t=1.672$，$b=57$. 则各个单独检验的大小均约为 $\alpha=.05$. 所以，由\autoref{thm:8.3.23}，交并检验是水平 $\alpha=.05$ 的检验，即产品实际上不合格却被判为合格的概率不超过 $.05$. 事实上，这个检验的大小为 $\alpha=.05$. 为说明这一点，考虑参数点列 $\theta_l=(\theta_{1l},\theta_2)$，其中 $l\to\infty$ 时 $\theta_{1l}\to\infty$，$\theta_2=.95$. 因 $\theta_2\le.95$，所有这些点都在 $\Theta_0$ 中. 而且 $\theta_{1l}\to\infty$ 时 $P_{\theta_l}(\boldsymbol X\in R_1)\to1$；因为 $\theta_2=.95$，对所有 $l$，$P_{\theta_l}(\boldsymbol X\in R_2)=.05$. 所以由\autoref{thm:8.3.24}，交并检验的大小为 $\alpha$.\footnote{译注：本例最初已说明大小为近似值，后面的“恰为”及“不超过”不能按精确的 $.05$ 理解. 对给定 $m=58,b=57$，第二项检验在 $58$ 次均成功时拒绝，其边界概率为 $.95^{58}\approx.05105>.05$；沿原文所选参数列，交并拒绝概率也趋于这一数值，故所给检验并非严格的水平 $.05$ 检验. 原参数和结论均保留，数值近似的影响在此说明.}
\end{example}
注意，\autoref{ex:8.3.25}只利用 $X_1,\ldots,X_n$ 与 $Y_1,\ldots,Y_m$ 的边缘分布求出检验大小. 这一点非常重要，直接关系到交并检验的实用性，因为联合分布往往难以获知，即使知道也往往难以处理. 例如，若 $X_i,Y_i$ 来自同一块织物的测量，二者可能有关联；要计算某个特定参数值下交并检验的精确功效，就必须建立并使用这种关联的模型.



\sourcepageanchor{353}
\subsubsection{\texorpdfstring{$p$}{p} 值}\label{sec:8.3.4}
假设检验完成后，必须以在统计上有意义的方式报告结论. 一种报告方式是给出所用检验的大小 $\alpha$，以及拒绝或接受 $H_0$ 的决定. 检验大小包含重要信息：若 $\alpha$ 小，拒绝 $H_0$ 的决定相当有说服力；若 $\alpha$ 大，这一决定就不太有说服力，因为检验错误作出该决定的概率较大. 另一种方式是报告称为 $p$ 值（$p$-value）的某种检验统计量的取值.
\begin{definition}\label{def:8.3.26}\index{SI0742@p值}
$p$ 值 $p(\boldsymbol X)$ 是一个检验统计量，对每个样本点 $\boldsymbol x$ 满足 $0\le p(\boldsymbol x)\le1$. $p(\boldsymbol X)$ 的小值为 $H_1$ 为真提供证据. 若对每个 $\theta\in\Theta_0$ 和每个 $0\le\alpha\le1$，都有
\begin{equation}\label{eq:8.3.8}P_\theta(p(\boldsymbol X)\le\alpha)\le\alpha,\end{equation}
则称该 $p$ 值有效（valid\mbox{）.}
\end{definition}
若 $p(\boldsymbol X)$ 有效，很容易基于它构造水平 $\alpha$ 的检验. 当且仅当 $p(\boldsymbol X)\le\alpha$ 时拒绝 $H_0$ 的检验，由\eqref{eq:8.3.8}可知水平为 $\alpha$. 用 $p$ 值报告检验结果的优点是，每位读者都能选择自己认为适当的 $\alpha$，将报告的 $p(\boldsymbol x)$ 与之比较，判断这些数据会导致接受还是拒绝 $H_0$. 此外，$p$ 值越小，拒绝 $H_0$ 的证据越强. 因而，$p$ 值在一个更连续的尺度上报告检验结果，而不只是“接受 $H_0$”或“拒绝 $H_0$”这样的二分决策.

\autoref{thm:8.3.27}给出定义有效 $p$ 值最常用的方法.
\begin{theorem}\label{thm:8.3.27}
设 $W(\boldsymbol X)$ 是一个检验统计量，其大值为 $H_1$ 为真提供证据. 对每个样本点 $\boldsymbol x$，定义
\begin{equation}\label{eq:8.3.9}p(\boldsymbol x)=\sup_{\theta\in\Theta_0}P_\theta(W(\boldsymbol X)\ge W(\boldsymbol x)).\end{equation}
则 $p(\boldsymbol X)$ 是有效 $p$ 值.
\end{theorem}
\begin{proof}
固定 $\theta\in\Theta_0$. 用 $F_\theta(w)$ 表示 $-W(\boldsymbol X)$ 的累积分布函数，定义
$$\begin{aligned}
p_\theta(\boldsymbol x)&=P_\theta(W(\boldsymbol X)\ge W(\boldsymbol x))\\
&=P_\theta(-W(\boldsymbol X)\le-W(\boldsymbol x))=F_\theta(-W(\boldsymbol x)).
\end{aligned}$$
则随机变量 $p_\theta(\boldsymbol X)$ 等于 $F_\theta(-W(\boldsymbol X))$. 因此，由概率积分变换或\autoref{exercise:2.10}，$p_\theta(\boldsymbol X)$ 的分布在随机意义下大于或等于 $\operatorname{uniform}(0,1)$ 分布，即对每个 $0\le\alpha\le1$，$P_\theta(p_\theta(\boldsymbol X)\le\alpha)\le\alpha$. 因为对每个 $\boldsymbol x$，$p(\boldsymbol x)=\sup_{\theta'\in\Theta_0}p_{\theta'}(\boldsymbol x)\ge p_\theta(\boldsymbol x)$，所以
$$P_\theta(p(\boldsymbol X)\le\alpha)\le P_\theta(p_\theta(\boldsymbol X)\le\alpha)\le\alpha.$$
这对每个 $\theta\in\Theta_0$ 和每个 $0\le\alpha\le1$ 成立，故 $p(\boldsymbol X)$ 有效.
\end{proof}
\eqref{eq:8.3.9}的上确界可能难以计算. 下面两个例子说明不太困难的常见情形：第一例无需取上确界；第二例则容易确定使上确界达到的 $\theta$ 值.
\begin{example}[双侧正态检验的 $p$ 值]\label{ex:8.3.28}
设 $X_1,\ldots,X_n$ 是来自 $\mathrm n(\mu,\sigma^2)$ 总体的随机样本. 考虑检验 $H_0:\mu=\mu_0$ 对 $H_1:\mu\ne\mu_0$. 由\autoref{exercise:8.38}，似然比检验在 $W(\boldsymbol X)=|\bar X-\mu_0|/(S/\sqrt n)$ 较大时拒绝 $H_0$. 若 $\mu=\mu_0$，不论 $\sigma$ 为何，$(\bar X-\mu_0)/(S/\sqrt n)$ 均服从自由度为 $n-1$ 的 Student $t$ 分布. 因而计算\eqref{eq:8.3.9}时，对所有 $\theta$ 值，也就是所有 $\sigma$ 值，概率都相同. 所以，这个双侧 $t$ 检验由\eqref{eq:8.3.9}给出的 $p$ 值为
$$p(\boldsymbol x)=2P\left(T_{n-1}\ge\frac{|\bar x-\mu_0|}{s/\sqrt n}\right),$$
其中 $T_{n-1}$ 服从自由度为 $n-1$ 的 Student $t$ 分布.
\end{example}
\sourcepageanchor{354}
\begin{example}[单侧正态检验的 $p$ 值]\label{ex:8.3.29}
再考虑\autoref{ex:8.3.28}的正态模型，但检验 $H_0:\mu\le\mu_0$ 对 $H_1:\mu>\mu_0$. 由\autoref{exercise:8.37}，似然比检验在 $W(\boldsymbol X)=(\bar X-\mu_0)/(S/\sqrt n)$ 较大时拒绝 $H_0$. 下列论证说明，对此统计量，\eqref{eq:8.3.9}的上确界总在参数 $(\mu_0,\sigma)$ 处达到，且所取 $\sigma$ 值无关紧要. 考虑任意 $\mu\le\mu_0$ 和任意 $\sigma$：
$$\begin{aligned}
P_{\mu,\sigma}(W(\boldsymbol X)\ge W(\boldsymbol x))
&=P_{\mu,\sigma}\left(\frac{\bar X-\mu_0}{S/\sqrt n}\ge W(\boldsymbol x)\right)\\
&=P_{\mu,\sigma}\left(\frac{\bar X-\mu}{S/\sqrt n}\ge W(\boldsymbol x)+\frac{\mu_0-\mu}{S/\sqrt n}\right)\\
&=P_{\mu,\sigma}\left(T_{n-1}\ge W(\boldsymbol x)+\frac{\mu_0-\mu}{S/\sqrt n}\right)\\
&\le P(T_{n-1}\ge W(\boldsymbol x)).
\end{aligned}$$
这里 $T_{n-1}$ 仍服从自由度为 $n-1$ 的 Student $t$ 分布. 最后一行的不等式成立，是因为 $\mu_0\ge\mu$，所以 $(\mu_0-\mu)/(S/\sqrt n)$ 是非负随机变量. 此处省去 $P$ 的下标，因为该概率不依赖 $(\mu,\sigma)$. 另外，
$$\begin{aligned}
P(T_{n-1}\ge W(\boldsymbol x))&=P_{\mu_0,\sigma}\left(\frac{\bar X-\mu_0}{S/\sqrt n}\ge W(\boldsymbol x)\right)\\
&=P_{\mu_0,\sigma}(W(\boldsymbol X)\ge W(\boldsymbol x)),
\end{aligned}$$
由于 $(\mu_0,\sigma)\in\Theta_0$，这一概率正是\eqref{eq:8.3.9}取上确界时所考虑的概率之一. 因此，单侧 $t$ 检验由\eqref{eq:8.3.9}给出的 $p$ 值为
$$p(\boldsymbol x)=P(T_{n-1}\ge W(\boldsymbol x))=P\left(T_{n-1}\ge\frac{\bar x-\mu_0}{s/\sqrt n}\right).$$
\end{example}
定义有效 $p$ 值的另一种方法，是对充分统计量取条件分布，从而替代\eqref{eq:8.3.9}. 设 $S(\boldsymbol X)$ 是模型 $\{f(\boldsymbol x\mid\theta):\theta\in\Theta_0\}$ 的充分统计量. （为避免检验功效低，重要的是 $S$ 仅对原假设模型充分，而不是对整个模型 $\{f(\boldsymbol x\mid\theta):\theta\in\Theta\}$ 充分.）若原假设为真，给定 $S=s$ 后 $\boldsymbol X$ 的条件分布不依赖 $\theta$. 再令 $W(\boldsymbol X)$ 为大值支持 $H_1$ 为真的检验统计量. 对每个样本点 $\boldsymbol x$ 定义
\begin{equation}\label{eq:8.3.10}\index{SI0743@条件}p(\boldsymbol x)=P(W(\boldsymbol X)\ge W(\boldsymbol x)\mid S=S(\boldsymbol x)).\end{equation}
像\autoref{thm:8.3.27}那样论证，但只考虑给定 $S=s$ 后 $\boldsymbol X$ 的这一个条件分布，可知对任意 $0\le\alpha\le1$，
$$P(p(\boldsymbol X)\le\alpha\mid S=s)\le\alpha.$$
于是，对任意 $\theta\in\Theta_0$，无条件地有
$$\begin{aligned}
P_\theta(p(\boldsymbol X)\le\alpha)&=\sum_sP(p(\boldsymbol X)\le\alpha\mid S=s)P_\theta(S=s)\\
&\le\sum_s\alpha P_\theta(S=s)\le\alpha.
\end{aligned}$$
所以，\eqref{eq:8.3.10}定义的 $p(\boldsymbol X)$ 有效. 对连续 $S$，求和可换成积分，但此法通常用于离散 $S$，如下例.
\begin{example}[Fisher 精确检验]\label{ex:8.3.30}
设 $S_1,S_2$ 是独立观测，$S_1\sim\operatorname{binomial}(n_1,p_1)$，$S_2\sim\operatorname{binomial}(n_2,p_2)$. 考虑检验 $H_0:p_1=p_2$ 对 $H_1:p_1>p_2$. 在 $H_0$ 下，用 $p$ 表示 $p_1=p_2$ 的共同值，$(S_1,S_2)$ 的联合质量函数为
$$\begin{aligned}
f(s_1,s_2\mid p)&=\binom{n_1}{s_1}p^{s_1}(1-p)^{n_1-s_1}\binom{n_2}{s_2}p^{s_2}(1-p)^{n_2-s_2}\\
&=\binom{n_1}{s_1}\binom{n_2}{s_2}p^{s_1+s_2}(1-p)^{n_1+n_2-(s_1+s_2)}.
\end{aligned}$$
\sourcepageanchor{355}
因此，$S=S_1+S_2$ 在 $H_0$ 下是充分统计量. 给定 $S=s$，合理的做法是用 $S_1$ 作检验统计量，当 $S_1$ 较大时拒绝 $H_0$、支持 $H_1$，因为较大的 $S_1$ 对应较小的 $S_2=s-S_1$. 给定 $S=s$ 后 $S_1$ 的条件分布为 $\operatorname{hypergeometric}(n_1+n_2,n_1,s)$（见\autoref{exercise:8.48}\mbox{）.} 所以，\eqref{eq:8.3.10}中的条件 $p$ 值为
$$p(s_1,s_2)=\sum_{j=s_1}^{\min\{n_1,s\}}f(j\mid s),$$
即超几何概率之和. 由这个 $p$ 值定义的检验称为 Fisher 精确检验（Fisher's Exact Test\mbox{）.}
\end{example}



\subsubsection{损失函数意义下的最优性}\label{sec:8.3.5}\index{SI0244@决策论}\index{SI0492@决策论}
可以像\autoref{sec:7.3.4}那样，通过决策论分析来比较假设检验，而不只通过功效函数比较. 为作这种分析，必须为假设检验问题指定行动空间和损失函数.

假设检验问题只允许两种行动：“接受 $H_0$”或“拒绝 $H_0$”，可分别记为 $a_0,a_1$. 行动空间是两点集合 $\mathscr A=\{a_0,a_1\}$.\index{SI0004@行动空间}\semanticindexanchor{4:355} 决策规则 $\delta(\boldsymbol x)$（一个假设检验）是样本空间 $\mathscr X$ 上只取 $a_0,a_1$ 两个值的函数. 集合 $\{\boldsymbol x:\delta(\boldsymbol x)=a_0\}$ 是接受域，$\{\boldsymbol x:\delta(\boldsymbol x)=a_1\}$ 是拒绝域，正如\autoref{def:8.1.3}.

假设检验问题的损失函数应反映：若 $\theta\in\Theta_0$ 却选择 $a_1$，或 $\theta\in\Theta_0^c$ 却选择 $a_0$，就犯了错误；另外两种可能情形则作出了正确决定. 由于只有两种可能行动，损失函数 $L(\theta,a)$ 只含两部分. $L(\theta,a_0)$ 是决定接受 $H_0$ 时，在不同 $\theta$ 值下产生的损失；$L(\theta,a_1)$ 则是决定拒绝 $H_0$ 时的损失.

检验问题中最简单的损失称为 $0$–$1$ 损失（$0$–$1$ loss），定义为
$$L(\theta,a_0)=\begin{cases}0,&\theta\in\Theta_0,\\1,&\theta\in\Theta_0^c,\end{cases}
\qquad L(\theta,a_1)=\begin{cases}1,&\theta\in\Theta_0,\\0,&\theta\in\Theta_0^c.\end{cases}$$
采用 $0$–$1$ 损失时，正确决定的损失为 $0$，错误决定的损失为 $1$. 这是两类错误后果相同的一种特别简单的情形. 稍微更现实的损失，是为两类错误赋予不同代价的广义 $0$–$1$ 损失（generalized $0$–$1$ loss）：
\begin{equation}\label{eq:8.3.11}\index{SI0655@广义零一}\index{SI0656@假设检验}\index{SI0663@零一}
L(\theta,a_0)=\begin{cases}0,&\theta\in\Theta_0,\\c_{\mathrm{II}},&\theta\in\Theta_0^c,\end{cases}
\qquad L(\theta,a_1)=\begin{cases}c_{\mathrm I},&\theta\in\Theta_0,\\0,&\theta\in\Theta_0^c.\end{cases}
\end{equation}
这里 $c_{\mathrm I}$ 是第一类错误、即错误拒绝 $H_0$ 的代价，$c_{\mathrm{II}}$ 是第二类错误、即错误接受 $H_0$ 的代价. （实际上，比较检验时真正重要的只是比值 $c_{\mathrm{II}}/c_{\mathrm I}$，不是两个单独的值. 若 $c_{\mathrm I}=c_{\mathrm{II}}$，实质上就是 $0$–$1$ 损失.）

决策论分析用风险函数（期望损失）评价假设检验程序. 检验的风险函数与功效函数密切相关，如下列分析所示.

设 $\beta(\theta)$ 是基于决策规则 $\delta$ 的检验的功效函数，即若用 $R=\{\boldsymbol x:\delta(\boldsymbol x)=a_1\}$ 表示拒绝域，则
$$\beta(\theta)=P_\theta(\boldsymbol X\in R)=P_\theta(\delta(\boldsymbol X)=a_1).$$
\sourcepageanchor{356}

\begin{figure}[htbp]\centering
\includegraphics[width=.8\textwidth]{assets/chapter8/fig8-3-4.pdf}
\caption{\autoref{ex:8.3.31}中检验的风险函数，其中 $z=\sqrt n(\theta-\theta_0)/\sigma$.}\label{fig:8.3.4}
\end{figure}

与\eqref{eq:8.3.11}相应、特别是与 $0$–$1$ 损失相应的风险函数非常简单. 对任意 $\theta\in\Theta$，$L(\theta,a)$ 只取两个值：若 $\theta\in\Theta_0$，为 $0$ 和 $c_{\mathrm I}$；若 $\theta\in\Theta_0^c$，为 $0$ 和 $c_{\mathrm{II}}$. 所以风险为
\begin{equation}\label{eq:8.3.12}\index{SI0505@与风险的关系}
\begin{aligned}
R(\theta,\delta)&=0P_\theta(\delta(\boldsymbol X)=a_0)+c_{\mathrm I}P_\theta(\delta(\boldsymbol X)=a_1)\\
&=c_{\mathrm I}\beta(\theta),&&\theta\in\Theta_0;\\
R(\theta,\delta)&=c_{\mathrm{II}}P_\theta(\delta(\boldsymbol X)=a_0)+0P_\theta(\delta(\boldsymbol X)=a_1)\\
&=c_{\mathrm{II}}(1-\beta(\theta)),&&\theta\in\Theta_0^c.
\end{aligned}
\end{equation}
决策论方法与更传统的功效方法之间的这种相似性，部分归因于损失函数的形式. 但如下所示，在所有假设检验问题中，功效函数都在风险函数中起重要作用.
\begin{example}[一致最强检验的风险]\label{ex:8.3.31}\index{SI0493@均值}
设 $X_1,\ldots,X_n$ 是来自 $\mathrm n(\mu,\sigma^2)$ 总体的随机样本，$\sigma^2$ 已知.\footnote{译注：本例原文开头用 $\mu$ 表示正态均值，随后检验与功效公式改用 $\theta$；应理解为同一均值参数，原记号保留.} 检验 $H_0:\theta\ge\theta_0$ 对 $H_1:\theta<\theta_0$ 的水平 $\alpha$ 一致最强检验，在 $(\bar X-\theta_0)/(\sigma/\sqrt n)<-z_\alpha$ 时拒绝 $H_0$（\autoref{ex:8.3.15}\mbox{）.} 其功效函数为
$$\beta(\theta)=P_\theta\left(Z<-z_\alpha-\frac{\theta-\theta_0}{\sigma/\sqrt n}\right),$$
其中 $Z\sim\mathrm n(0,1)$. 对 $\alpha=.10$，取 $c_{\mathrm I}=8$、$c_{\mathrm{II}}=3$ 时，\eqref{eq:8.3.12}的风险函数见\autoref{fig:8.3.4}. 注意风险函数在 $\theta=\theta_0$ 处不连续；这是因为在 $\theta_0$ 处表达式由 $\beta(\theta)$ 变为 $1-\beta(\theta)$，并且 $c_{\mathrm I}$ 与 $c_{\mathrm{II}}$ 不同.
\end{example}
\sourcepageanchor{357}
$0$–$1$ 损失只判断决定是对还是错. 某些错误决定可能比其他错误决定更严重，损失函数应反映这一点. 检验 $H_0:\theta\ge\theta_0$ 对 $H_1:\theta<\theta_0$ 时，若 $\theta$ 略大于 $\theta_0$ 而拒绝 $H_0$，虽犯第一类错误，却可能不很严重；若 $\theta$ 远大于 $\theta_0$，拒绝 $H_0$ 的不良后果可能严重得多. 反映这种情况的一个损失函数是
\begin{equation}\label{eq:8.3.13}
L(\theta,a_0)=\begin{cases}0,&\theta\ge\theta_0,\\b(\theta_0-\theta),&\theta<\theta_0,\end{cases}
\quad L(\theta,a_1)=\begin{cases}c(\theta-\theta_0)^2,&\theta\ge\theta_0,\\0,&\theta<\theta_0,\end{cases}
\end{equation}
其中 $b,c$ 为正常数. 例如，若实验者检验某药物是否降低胆固醇水平，可以这样设定 $H_0,H_1$，取 $\theta_0$ 为可接受胆固醇水平的标准. 因高胆固醇水平与心脏病相关，在 $\theta$ 很大时拒绝 $H_0$ 的后果相当严重. \eqref{eq:8.3.13}这样的损失函数反映了这一后果. \hyperref[bib:source320]{Vardeman (1987)} 主张使用类似类型的损失函数.

即使对\eqref{eq:8.3.13}这样的一般损失函数，风险函数与功效函数仍密切相关. 对任意固定 $\theta$，损失不是 $L(\theta,a_0)$，就是 $L(\theta,a_1)$，所以期望损失为
\begin{equation}\label{eq:8.3.14}
\begin{aligned}
R(\theta,\delta)&=L(\theta,a_0)P_\theta(\delta(\boldsymbol X)=a_0)+L(\theta,a_1)P_\theta(\delta(\boldsymbol X)=a_1)\\
&=L(\theta,a_0)(1-\beta(\theta))+L(\theta,a_1)\beta(\theta).
\end{aligned}
\end{equation}
评价假设检验时，功效函数总是重要的；但在决策论分析中，损失函数赋予的权重也同样重要.





\subsection{习题}\label{sec:8.4}
\setcounter{exercise}{0}
\begin{exercise}\label{exercise:8.1}
抛掷一枚硬币 $1000$ 次，出现 $560$ 次正面、$440$ 次反面. 假定这枚硬币公平是否合理？说明理由.
\end{exercise}
\begin{exercise}\label{exercise:8.2}
假设某城市一年内的汽车事故数服从泊松分布. 以往每年平均事故数为 $15$，今年为 $10$. 据此声称事故率下降，是否有充分依据？
\end{exercise}
\begin{exercise}\label{exercise:8.3}
本题推导\autoref{ex:8.2.9}中提到的似然比检验. 假设观测到 $m$ 个独立同分布的 $\operatorname{Bernoulli}(\theta)$ 随机变量 $Y_1,\ldots,Y_m$. 证明，检验 $H_0:\theta\le\theta_0$ 对 $H_1:\theta>\theta_0$ 的似然比检验，在 $\sum_{i=1}^mY_i>b$ 时拒绝 $H_0$.
\end{exercise}
\begin{exercise}\label{exercise:8.4}
证明\autoref{def:8.2.1}后正文的断言：若 $f(\boldsymbol x\mid\theta)$ 是离散随机变量的概率质量函数，则似然比检验统计量 $\lambda(\boldsymbol x)$ 的分子，是在原假设参数范围内计算的已观察样本的最大概率；其分母则是已观察样本在所有可能参数上的最大概率.
\end{exercise}
\begin{exercise}\label{exercise:8.5}\index{SI0641@帕累托}
随机样本 $X_1,\ldots,X_n$ 来自 Pareto 总体，其密度为
$$f(x\mid\theta,\nu)=\frac{\theta\nu^\theta}{x^{\theta+1}}I_{[\nu,\infty)}(x),\qquad\theta>0,\quad\nu>0.$$
\begin{enumerate}
\item 求 $\theta$ 与 $\nu$ 的最大似然估计量.
\item 证明，检验
$$H_0:\theta=1,\ \nu\text{ 未知}\quad\text{对}\quad H_1:\theta\ne1,\ \nu\text{ 未知}$$
的似然比检验具有形如 $\{\boldsymbol x:T(\boldsymbol x)\le c_1\text{ 或 }T(\boldsymbol x)\ge c_2\}$ 的临界域，其中 $0<c_1<c_2$，且
$$T=\log\left[\frac{\prod_{i=1}^nX_i}{(\min_iX_i)^n}\right].$$
\sourcepageanchor{358}
\item 证明，在 $H_0$ 下，$2T$ 服从卡方分布，并求自由度. （提示：给定 $\min_iX_i$，求 $n-1$ 个非平凡项 $X_i/(\min_iX_i)$ 的联合分布. 将这 $n-1$ 项合并，并注意到给定 $\min_iX_i$ 后 $T$ 的分布不依赖 $\min_iX_i$，因而也就是 $T$ 的无条件分布.）
\end{enumerate}
\end{exercise}
\begin{exercise}\label{exercise:8.6}\index{SI0634@指数}
设有两个独立随机样本：$X_1,\ldots,X_n$ 来自 $\operatorname{exponential}(\theta)$，$Y_1,\ldots,Y_m$ 来自 $\operatorname{exponential}(\mu)$.
\begin{enumerate}
\item\phantomsection\label{itm:exercise:8.6:a} 求 $H_0:\theta=\mu$ 对 $H_1:\theta\ne\mu$ 的似然比检验.
\item 证明，\hyperref[itm:exercise:8.6:a]{(a)} 中的检验可以基于统计量
$$T=\frac{\sum X_i}{\sum X_i+\sum Y_i}.$$
\item 求 $H_0$ 为真时 $T$ 的分布.
\end{enumerate}
\end{exercise}
\begin{exercise}\label{exercise:8.7}
我们已经看到似然比检验在处理干扰参数问题时的用处. 现在考察另外一些干扰参数问题.
\begin{enumerate}
\item 基于来自密度为 $f(x\mid\theta,\lambda)=\lambda^{-1}e^{-(x-\theta)/\lambda}I_{(\theta,\infty)}(x)$ 的总体的样本 $X_1,\ldots,X_n$，求
$$H_0:\theta\le0\quad\text{对}\quad H_1:\theta>0$$
的似然比检验，其中 $\theta,\lambda$ 都未知.
\item 前面已看到指数密度是伽马密度的特例. 沿另一方向推广，指数密度也是 $\operatorname{Weibull}(\gamma,\beta)$ 的特例：当 $\gamma=1$ 时，Weibull 密度退化为指数密度. Weibull 密度在系统可靠性建模中非常重要. 设 $X_1,\ldots,X_n$ 是来自 Weibull 总体的随机样本，$\gamma,\beta$ 均未知. 求 $H_0:\gamma=1$ 对 $H_1:\gamma\ne1$ 的似然比检验.
\end{enumerate}
\end{exercise}
\begin{exercise}\label{exercise:8.8}
正态族的一种特殊情况是均值与方差相联系的 $\mathrm n(\theta,a\theta)$ 族. 若希望不论 $\theta$ 的值为何，都检验这种关系，就再次遇到干扰参数问题.
\begin{enumerate}
\item\phantomsection\label{itm:exercise:8.8:a} 基于来自 $\mathrm n(\theta,a\theta)$ 族的样本 $X_1,\ldots,X_n$，求 $H_0:a=1$ 对 $H_1:a\ne1$ 的似然比检验，其中 $\theta$ 未知.
\item\phantomsection\label{itm:exercise:8.8:b} 对相关的 $\mathrm n(\theta,a\theta^2)$ 族也可提出类似问题. 即若 $X_1,\ldots,X_n$ 独立同分布于 $\mathrm n(\theta,a\theta^2)$，$\theta$ 未知，求 $H_0:a=1$ 对 $H_1:a\ne1$ 的似然比检验.\footnote{译注：原题未明示参数空间. 使用普通正态密度时，\hyperref[itm:exercise:8.8:a]{(a)} 至少须有 $a\theta>0$，\hyperref[itm:exercise:8.8:b]{(b)} 须有 $a>0$、$\theta\ne0$；是否还限定 $a,\theta$ 各自为正，会影响约束最大化的范围. 此处保留原题，不能默把零方差纳入同一密度公式.}
\end{enumerate}
\end{exercise}
\begin{exercise}\label{exercise:8.9}
\hyperref[bib:source300]{Stefanski (1996)} 用基于似然比检验的证明建立了算术—几何—调和平均不等式（见\autoref{ex:4.7.8}与\hyperref[sec:4.9.2]{杂录~4.9.2}\mbox{）.} 设 $Y_1,Y_2,\ldots,Y_n$ 相互独立，密度为 $\lambda_i e^{-\lambda_i y_i}$；希望检验 $H_0:\lambda_1=\lambda_2=\cdots=\lambda_n$ 对 $H_1:$ 各 $\lambda_i$ 不全相等.
\begin{enumerate}
\item 证明，似然比检验统计量为 $(\bar Y)^{-n}/(\prod_iY_i)^{-1}$，并由此推出算术—几何平均不等式.
\item 作变换 $X_i=1/Y_i$，证明基于 $X_1,X_2,\ldots,X_n$ 的似然比检验统计量为
$$\frac{[n/\sum_i(1/X_i)]^n}{\prod_iX_i},$$
并由此推出几何—调和平均不等式.
\end{enumerate}
\end{exercise}
\begin{exercise}\label{exercise:8.10}
设 $X_1,\ldots,X_n$ 独立同分布于 $\operatorname{Poisson}(\lambda)$，$\lambda$ 服从泊松分布的共轭族 $\operatorname{gamma}(\alpha,\beta)$. \autoref{exercise:7.24}已求出 $\lambda$ 的后验分布，包括后验均值和方差. 现考虑 $H_0:\lambda\le\lambda_0$ 对 $H_1:\lambda>\lambda_0$ 的贝叶斯检验.
\begin{enumerate}
\item 求 $H_0,H_1$ 的后验概率表达式.
\item 若 $\alpha=5/2$、$\beta=2$，先验分布为自由度 $5$ 的卡方分布. 说明如何利用卡方表进行贝叶斯检验.
\end{enumerate}
\end{exercise}
\begin{exercise}\label{exercise:8.11}
\autoref{exercise:7.23}使用 $\sigma^2$ 的共轭先验（参数为 $\alpha,\beta$ 的逆伽马密度），求出了给定大小为 $n$ 的样本的样本方差 $S^2$ 后，正态总体方差 $\sigma^2$ 的后验分布. 现要根据观测到的 $S^2$，对假设 $H_0:\sigma\le1$ 与 $H_1:\sigma>1$ 作出决定.
\begin{enumerate}
\item\phantomsection\label{itm:exercise:8.11:a} 求样本空间中满足 $P(\sigma\le1\mid s^2)>P(\sigma>1\mid s^2)$ 的区域，即贝叶斯检验判定 $\sigma\le1$ 的区域.
\item 将\hyperref[itm:exercise:8.11:a]{(a)} 的区域与似然比检验的接受域比较. 是否存在先验参数的选择，使两区域一致？
\end{enumerate}
\end{exercise}
\sourcepageanchor{359}
\begin{exercise}\label{exercise:8.12}
对来自均值为 $\mu$、已知方差为 $\sigma^2$ 的正态总体、大小分别为 $n=1,4,16,64,100$ 的样本，画出下列检验的功效函数. 取 $\alpha=.05$.
\begin{enumerate}
\item $H_0:\mu\le0$ 对 $H_1:\mu>0$.
\item $H_0:\mu=0$ 对 $H_1:\mu\ne0$.
\end{enumerate}
\end{exercise}
\begin{exercise}\label{exercise:8.13}
设 $X_1,X_2$ 独立同分布于 $\operatorname{uniform}(\theta,\theta+1)$. 对检验 $H_0:\theta=0$ 对 $H_1:\theta>0$，有两个候选检验：
$$\begin{aligned}\phi_1(X_1):&\quad X_1>.95\text{ 时拒绝 }H_0,\\
\phi_2(X_1,X_2):&\quad X_1+X_2>C\text{ 时拒绝 }H_0.\end{aligned}$$
\begin{enumerate}
\item 求 $C$，使 $\phi_2$ 与 $\phi_1$ 大小相同.
\item 计算各检验的功效函数，并画出标注清楚的图像.
\item 证明或否证：$\phi_2$ 比 $\phi_1$ 功效更高.
\item 说明如何得到与 $\phi_2$ 大小相同、但功效更高的检验.
\end{enumerate}
\end{exercise}
\begin{exercise}\label{exercise:8.14}
对 $\operatorname{Bernoulli}(p)$ 随机变量的随机样本 $X_1,\ldots,X_n$，希望检验
$$H_0:p=.49\quad\text{对}\quad H_1:p=.51.$$
利用中心极限定理，近似确定使两类错误概率都约为 $.01$ 所需的样本量. 采用在 $\sum_{i=1}^nX_i$ 较大时拒绝 $H_0$ 的检验函数.
\end{exercise}
\begin{exercise}\label{exercise:8.15}
证明，对来自 $\mathrm n(0,\sigma^2)$ 总体的随机样本 $X_1,\ldots,X_n$，检验 $H_0:\sigma=\sigma_0$ 对 $H_1:\sigma=\sigma_1$（$\sigma_0<\sigma_1$）的最强检验为
$$\phi\left(\sum X_i^2\right)=\begin{cases}1,&\sum X_i^2>c,\\0,&\sum X_i^2\le c.\end{cases}$$
给定第一类错误的大小 $\alpha$，说明如何明确确定 $c$.
\end{exercise}
\begin{exercise}\label{exercise:8.16}\index{SI0483@滥用}
对 $\alpha$ 水平的一种明显滥用，是在看过数据之后才选择它，并使选择强制导致拒绝（或接受）原假设. 为了解这种做法真正的第一类和第二类错误概率，计算以下两个平凡检验的大小与功效：
\begin{enumerate}
\item 无论取得什么数据，总拒绝 $H_0$（相当于选择 $\alpha$ 水平以强制拒绝 $H_0$\mbox{）.}
\item 无论取得什么数据，总接受 $H_0$（相当于选择 $\alpha$ 水平以强制接受 $H_0$\mbox{）.}
\end{enumerate}
\end{exercise}
\begin{exercise}\label{exercise:8.17}\index{SI0633@贝塔}
设 $X_1,\ldots,X_n$ 独立同分布于 $\operatorname{beta}(\mu,1)$，$Y_1,\ldots,Y_m$ 独立同分布于 $\operatorname{beta}(\theta,1)$，且 $X$ 样本与 $Y$ 样本独立.
\begin{enumerate}
\item\phantomsection\label{itm:exercise:8.17:a} 求 $H_0:\theta=\mu$ 对 $H_1:\theta\ne\mu$ 的一个似然比检验.
\item 证明，\hyperref[itm:exercise:8.17:a]{(a)} 中的检验可以基于统计量
$$T=\frac{\sum\log X_i}{\sum\log X_i+\sum\log Y_i}.$$
\item 求 $H_0$ 为真时 $T$ 的分布，再说明如何得到大小 $\alpha=.10$ 的检验.
\end{enumerate}
\end{exercise}
\begin{exercise}\label{exercise:8.18}
设 $X_1,\ldots,X_n$ 是来自 $\mathrm n(\theta,\sigma^2)$ 总体的随机样本，$\sigma^2$ 已知. 检验 $H_0:\theta=\theta_0$ 对 $H_1:\theta\ne\theta_0$ 的似然比检验，在 $|\bar X-\theta_0|/(\sigma/\sqrt n)>c$ 时拒绝 $H_0$.
\begin{enumerate}
\item 用标准正态概率表示该检验的功效函数.
\item 实验者希望第一类错误概率为 $.05$，并在 $\theta=\theta_0+\sigma$ 处使第二类错误概率至多为 $.25$. 求能达到要求的 $n,c$.
\end{enumerate}
\end{exercise}
\begin{exercise}\label{exercise:8.19}
随机变量 $X$ 的密度为 $f(x)=e^{-x}$，$x>0$. 观察随机变量 $Y=X^\theta$ 一次，需构造 $H_0:\theta=1$ 对 $H_1:\theta=2$ 的检验. 求水平 $\alpha=.10$ 的一致最强检验，并计算第二类错误概率.
\end{exercise}
\sourcepageanchor{360}
\begin{exercise}\label{exercise:8.20}
随机变量 $X$ 在 $H_0$ 与 $H_1$ 下的概率质量函数如下：
\begin{center}\begin{tabular}{c|rrrrrrr}\toprule
$x$&1&2&3&4&5&6&7\\\midrule
$f(x\mid H_0)$&.01&.01&.01&.01&.01&.01&.94\\
$f(x\mid H_1)$&.06&.05&.04&.03&.02&.01&.79\\\bottomrule
\end{tabular}\end{center}
利用 Neyman–Pearson 引理，求检验 $H_0$ 对 $H_1$ 的大小 $\alpha=.04$ 的最强检验，并计算其第二类错误概率.
\end{exercise}
\begin{exercise}\label{exercise:8.21}
\autoref{thm:8.3.12}（Neyman–Pearson 引理）的证明指出，对连续随机变量给出的证明很容易适配到离散情形. 补出细节，即证明离散随机变量的 Neyman–Pearson 引理. 假定给定的 $\alpha$ 水平可达到.
\end{exercise}
\begin{exercise}\label{exercise:8.22}
设 $X_1,\ldots,X_{10}$ 独立同分布于 $\operatorname{Bernoulli}(p)$.
\begin{enumerate}
\item\phantomsection\label{itm:exercise:8.22:a} 求检验 $H_0:p=1/2$ 对 $H_1:p=1/4$ 的大小 $\alpha=.0547$ 的最强检验，并求其功效.\footnote{译注：在非随机化检验下，原题 $.0547$ 应理解为舍入值；常用下尾拒绝域的大小为 $56/1024=.0546875$. 若坚持十进制 $.0547$ 为精确大小，则必须另考虑随机化；原数值保留.}
\item 对检验 $H_0:p\le1/2$ 对 $H_1:p>1/2$，求在 $\sum_{i=1}^{10}X_i\ge6$ 时拒绝 $H_0$ 的检验的大小，并画出功效函数.
\item 对哪些 $\alpha$ 水平，\hyperref[itm:exercise:8.22:a]{(a)} 中的假设存在一致最强检验？
\end{enumerate}
\end{exercise}
\begin{exercise}\label{exercise:8.23}
设 $X$ 是来自密度为 $\operatorname{beta}(\theta,1)$ 的总体的一次观测.
\begin{enumerate}
\item 对检验 $H_0:\theta\le1$ 对 $H_1:\theta>1$，求在 $X>1/2$ 时拒绝 $H_0$ 的检验的大小，并画出功效函数.
\item 求 $H_0:\theta=1$ 对 $H_1:\theta=2$ 的水平 $\alpha$ 最强检验.
\item 是否存在 $H_0:\theta\le1$ 对 $H_1:\theta>1$ 的一致最强检验？若存在，求出；否则证明不存在.
\end{enumerate}
\end{exercise}
\begin{exercise}\label{exercise:8.24}\index{SI0637@Neyman–Pearson引理}\index{SI0723@似然比检验}
求简单 $H_0$ 对简单 $H_1$ 的似然比检验. 它是否等价于由 Neyman–Pearson 引理得到的检验？（\hyperref[bib:source296]{Solomon (1975)} 对这一关系有较详细讨论.）
\end{exercise}
\begin{exercise}\label{exercise:8.25}\index{SI0701@二项}\index{SI0706@正态}\index{SI0707@泊松}
证明下列各族均具有单调似然比：
\begin{enumerate}
\item $\sigma^2$ 已知的 $\mathrm n(\theta,\sigma^2)$ 族；
\item $\operatorname{Poisson}(\theta)$ 族；
\item $n$ 已知的 $\operatorname{binomial}(n,\theta)$ 族.
\end{enumerate}
\end{exercise}
\begin{exercise}\label{exercise:8.26}\index{SI0708@随机序}\index{SI0993@单调似然比}\leavevmode
\begin{enumerate}
\item\phantomsection\label{itm:exercise:8.26:a} 证明，若密度族 $\{f(x\mid\theta):\theta\in\Theta\}$ 具有单调似然比，则相应的累积分布函数族关于 $\theta$ 随机递增（stochastically increasing\mbox{）.} （见\autoref{sec:8.5}的杂录.）\footnote{译注：同\autoref{thm:8.3.17}，此结论要求 $\theta_2>\theta_1$ 时 $f(x\mid\theta_2)/f(x\mid\theta_1)$ 关于 $x$ 非减；原定义允许的非增方向不能推出所写的随机递增结论.}
\item 证明\hyperref[itm:exercise:8.26:a]{(a)} 的逆命题不成立：给出一个关于 $\theta$ 随机递增、但相应密度族没有单调似然比的累积分布函数族.
\end{enumerate}
\end{exercise}
\begin{exercise}\label{exercise:8.27}\index{SI0703@指数族}
设 $g(t\mid\theta)=h(t)c(\theta)e^{w(\theta)t}$ 是随机变量 $T$ 的单参数指数族. 证明，若 $w(\theta)$ 是 $\theta$ 的增函数，则此族具有单调似然比. 给出三个这类分布族的例子.
\end{exercise}
\begin{exercise}\label{exercise:8.28}\index{SI0704@逻辑}
设 $f(x\mid\theta)$ 是 logistic 位置密度
$$f(x\mid\theta)=\frac{e^{x-\theta}}{(1+e^{x-\theta})^2},\qquad-\infty<x<\infty,\quad-\infty<\theta<\infty.$$
\begin{enumerate}
\item 证明此族具有单调似然比.
\item\phantomsection\label{itm:exercise:8.28:b} 基于一次观测 $X$，求 $H_0:\theta=0$ 对 $H_1:\theta=1$ 的大小 $\alpha$ 最强检验. 当 $\alpha=.2$ 时，求第二类错误的大小.
\item 证明\hyperref[itm:exercise:8.28:b]{(b)} 中的检验是 $H_0:\theta\le0$ 对 $H_1:\theta>0$ 的大小 $\alpha$ 一致最强检验. 对 logistic 位置族的一致最强检验，一般可以说些什么？
\end{enumerate}
\end{exercise}

\sourcepageanchor{361}
\begin{exercise}\label{exercise:8.29}\index{SI0702@柯西}
设 $X$ 是来自 $\operatorname{Cauchy}(\theta)$ 分布的一次观测.
\begin{enumerate}
\item 证明此族没有单调似然比.
\item\phantomsection\label{itm:exercise:8.29:b} 证明检验
$$\phi(x)=\begin{cases}1,&1<x<3,\\0,&\text{其他}\end{cases}$$
是检验 $H_0:\theta=0$ 对 $H_1:\theta=1$ 时、其自身大小下的最强检验. 计算第一类与第二类错误概率.
\item 证明或否证：\hyperref[itm:exercise:8.29:b]{(b)} 的检验是 $H_0:\theta\le0$ 对 $H_1:\theta>0$ 的一致最强检验. 对 Cauchy 位置族的一致最强检验，一般可以说些什么？
\end{enumerate}
\end{exercise}
\begin{exercise}\label{exercise:8.30}
设 $f(x\mid\theta)$ 是 Cauchy 尺度密度
$$f(x\mid\theta)=\frac\theta\pi\frac1{\theta^2+x^2},\qquad-\infty<x<\infty,\quad\theta>0.$$
\begin{enumerate}
\item 证明此族没有单调似然比.
\item 若 $X$ 是来自 $f(x\mid\theta)$ 的一次观测，证明 $|X|$ 对 $\theta$ 充分，且 $|X|$ 的分布确实具有单调似然比.
\end{enumerate}
\end{exercise}
\begin{exercise}\label{exercise:8.31}
设 $X_1,\ldots,X_n$ 独立同分布于 $\operatorname{Poisson}(\lambda)$.
\begin{enumerate}
\item 求 $H_0:\lambda\le\lambda_0$ 对 $H_1:\lambda>\lambda_0$ 的一个一致最强检验.
\item 考虑具体情形 $H_0:\lambda\le1$ 对 $H_1:\lambda>1$. 利用中心极限定理确定样本量 $n$，使一致最强检验满足 $P(\text{拒绝 }H_0\mid\lambda=1)=.05$ 和 $P(\text{拒绝 }H_0\mid\lambda=2)=.9$.
\end{enumerate}
\end{exercise}
\begin{exercise}\label{exercise:8.32}
设 $X_1,\ldots,X_n$ 独立同分布于 $\mathrm n(\theta,1)$，$\theta_0$ 为 $\theta$ 的指定值.
\begin{enumerate}
\item 求 $H_0:\theta\ge\theta_0$ 对 $H_1:\theta<\theta_0$ 的大小 $\alpha$ 一致最强检验.
\item 证明不存在 $H_0:\theta=\theta_0$ 对 $H_1:\theta\ne\theta_0$ 的大小 $\alpha$ 一致最强检验.
\end{enumerate}
\end{exercise}
\begin{exercise}\label{exercise:8.33}\index{SI0519@均匀}
设 $X_1,\ldots,X_n$ 是来自 $\operatorname{uniform}(\theta,\theta+1)$ 的随机样本. 检验 $H_0:\theta=0$ 对 $H_1:\theta>0$，采用规则
$$\text{若 }Y_n\ge1\text{ 或 }Y_1\ge k,\quad\text{则拒绝 }H_0,$$
其中 $k$ 为常数，$Y_1=\min\{X_1,\ldots,X_n\}$，$Y_n=\max\{X_1,\ldots,X_n\}$.
\begin{enumerate}
\item\phantomsection\label{itm:exercise:8.33:a} 确定 $k$，使检验的大小为 $\alpha$.
\item 求\hyperref[itm:exercise:8.33:a]{(a)} 中检验的功效函数表达式.
\item 证明该检验是大小 $\alpha$ 的一致最强检验.
\item 求 $n,k$，使水平 $.10$ 的一致最强检验在 $\theta>1$ 时功效至少为 $.8$.\footnote{译注：原题确实写 $\theta>1$. 在此范围，样本最大值以概率 $1$ 大于 $1$，故给定检验的功效已为 $1$；若原意是某个小于 $1$ 的正偏移量，则需要另一条件，不能据题面自行补写.}
\end{enumerate}
\end{exercise}
\begin{exercise}\label{exercise:8.34}\index{SI0652@随机序}\index{SI0992@位置尺度族}
在以下两种情形中，证明对任意数 $c$，若 $\theta_1\le\theta_2$，则
$$P_{\theta_1}(T>c)\le P_{\theta_2}(T>c).$$
\begin{enumerate}
\item $\theta$ 是随机变量 $T$ 分布中的位置参数.
\item\phantomsection\label{itm:exercise:8.34:b} $T$ 的密度族 $\{g(t\mid\theta):\theta\in\Theta\}$ 具有单调似然比.\footnote{译注：\hyperref[itm:exercise:8.34:b]{(b)} 需取似然比非减的方向，参见\autoref{exercise:8.26}与\autoref{thm:8.3.17}的译注；原定义包含非增方向，不能不加区分地推出本题不等式.}
\end{enumerate}
\end{exercise}
\begin{exercise}\label{exercise:8.35}\index{SI0382@非中心}\index{SI0705@非中心t}
\autoref{sec:5.3.2}推导的通常 $t$ 分布，也称为中心 $t$ 分布（central $t$ distribution\mbox{）.} 它可看成形如 $T=\mathrm n(0,1)/\sqrt{\chi_\nu^2/\nu}$ 的随机变量的密度，其中正态变量与卡方变量独立. 一种推广是非中心 $t$ 分布（noncentral $t$），形式为 $T'=\mathrm n(\mu,1)/\sqrt{\chi_\nu^2/\nu}$，正态与卡方变量独立，且允许 $\mu\ne0$. （我们已在\eqref{eq:4.4.3}见过非中心卡方密度.）正式地，若 $X\sim\mathrm n(\mu,1)$，$Y\sim\chi_\nu^2$ 且独立于 $X$，则 $T'=X/\sqrt{Y/\nu}$ 服从自由度为 $\nu$、非中心参数为 $\delta=\sqrt{\mu^2}$ 的非中心 $t$ 分布.\footnote{译注：对有符号的非中心 $t$ 随机变量，此参数通常应为 $\delta=\mu$，而非 $\sqrt{\mu^2}=|\mu|$. $\mu$ 与 $-\mu$ 的分布互为镜像，一般不相同；原式丢失符号，只有限制 $\mu\ge0$ 才一致. 下述密度级数也含奇次 $\delta t$ 项，体现了这种有符号依赖.}
\begin{enumerate}
\item 计算 $T'$ 的均值和方差.
\item 证明 $T'$ 的密度为
$$f_{T'}(t\mid\delta)=\frac{e^{-\delta^2/2}}{\Gamma(\frac12)\Gamma(\frac\nu2)\sqrt\nu}
\sum_{k=0}^\infty\frac{(2/\nu)^{k/2}(\delta t)^k}{k!}
\frac{\Gamma((\nu+k+1)/2)}{(1+t^2/\nu)^{(\nu+k+1)/2}}.$$
证明，当 $\delta=0$ 时，它退化为中心 $t$ 密度.
\item 证明 $T'$ 的密度关于其非中心参数具有单调似然比.
\end{enumerate}
\end{exercise}
\sourcepageanchor{362}
\begin{exercise}\label{exercise:8.36}
我们有来自 $\operatorname{beta}(1,\theta)$ 总体的一次观测.
\begin{enumerate}
\item 检验 $H_0:\theta_1\le\theta\le\theta_2$ 对 $H_1:\theta<\theta_1\text{ 或 }\theta>\theta_2$，其中 $\theta_1=1$、$\theta_2=2$. 某检验满足 $\mathrm E_{\theta_1}\phi=.5$ 与 $\mathrm E_{\theta_2}\phi=.3$. 求一个同样好的检验，并说明它为何同样好.
\item 对检验 $H_0:\theta=\theta_1$ 对 $H_1:\theta\ne\theta_1$，其中 $\theta_1=1$，求一个双侧检验（不同于 $\phi\equiv.1$），满足
$$\mathrm E_{\theta_1}\phi=.1,\qquad\left.\frac d{d\theta}\mathrm E_\theta(\phi)\right|_{\theta=\theta_1}=0.$$
\end{enumerate}
\end{exercise}
\begin{exercise}\label{exercise:8.37}\index{SI0517@正态}\index{SI0638@正态}\index{SI1027@无偏}
设 $X_1,\ldots,X_n$ 是来自 $\mathrm n(\theta,\sigma^2)$ 总体的随机样本. 考虑检验
$$H_0:\theta\le\theta_0\quad\text{对}\quad H_1:\theta>\theta_0.$$
\begin{enumerate}
\item\phantomsection\label{itm:exercise:8.37:a} 若 $\sigma^2$ 已知，证明在
$$\bar X>\theta_0+z_\alpha\sqrt{\sigma^2/n}$$
时拒绝 $H_0$ 的检验，大小为 $\alpha$. 证明它可由似然比检验导出.
\item 证明\hyperref[itm:exercise:8.37:a]{(a)} 中的检验是一致最强检验.
\item 若 $\sigma^2$ 未知，证明在
$$\bar X>\theta_0+t_{n-1,\alpha}\sqrt{S^2/n}$$
时拒绝 $H_0$ 的检验，大小为 $\alpha$. 证明它可由似然比检验导出.
\end{enumerate}
\end{exercise}
\begin{exercise}\label{exercise:8.38}
设 $X_1,\ldots,X_n$ 独立同分布于 $\mathrm n(\theta,\sigma^2)$，$\theta_0$ 为 $\theta$ 的指定值，$\sigma^2$ 未知. 希望检验
$$H_0:\theta=\theta_0\quad\text{对}\quad H_1:\theta\ne\theta_0.$$
\begin{enumerate}
\item\phantomsection\label{itm:exercise:8.38:a} 证明在
$$|\bar X-\theta_0|>t_{n-1,\alpha/2}\sqrt{S^2/n}$$
时拒绝 $H_0$ 的检验，大小为 $\alpha$.
\item 证明\hyperref[itm:exercise:8.38:a]{(a)} 中的检验可由似然比检验导出.
\end{enumerate}
\end{exercise}
\begin{exercise}\label{exercise:8.39}\index{SI1022@配对}
设 $(X_1,Y_1),\ldots,(X_n,Y_n)$ 是来自参数为 $\mu_X,\mu_Y,\sigma_X^2,\sigma_Y^2,\rho$ 的二元正态分布的随机样本. 希望检验
$$H_0:\mu_X=\mu_Y\quad\text{对}\quad H_1:\mu_X\ne\mu_Y.$$
\begin{enumerate}
\item 证明 $W_i=X_i-Y_i$ 独立同分布于 $\mathrm n(\mu_W,\sigma_W^2)$.
\item 证明可用统计量
$$T_W=\frac{\bar W}{\sqrt{S_W^2/n}}$$
检验上述假设，其中 $\bar W=n^{-1}\sum_{i=1}^nW_i$，$S_W^2=(n-1)^{-1}\sum_{i=1}^n(W_i-\bar W)^2$. 进一步证明，在 $H_0$ 下，$T_W$ 服从自由度为 $n-1$ 的 Student $t$ 分布. （这个检验称为配对样本 $t$ 检验（paired-sample $t$ test\mbox{）.}）
\end{enumerate}
\end{exercise}
\begin{exercise}\label{exercise:8.40}\index{SI1020@配对}
设 $(X_1,Y_1),\ldots,(X_n,Y_n)$ 是来自参数为 $\mu_X,\mu_Y,\sigma_X^2,\sigma_Y^2,\rho$ 的二元正态分布的随机样本.
\begin{enumerate}
\item\phantomsection\label{itm:exercise:8.40:a} 推导
$$H_0:\mu_X=\mu_Y\quad\text{对}\quad H_1:\mu_X\ne\mu_Y$$
的似然比检验，其中 $\sigma_X^2,\sigma_Y^2,\rho$ 未指定且未知.
\item 证明\hyperref[itm:exercise:8.40:a]{(a)} 的检验等价于\autoref{exercise:8.39}的配对 $t$ 检验.

（提示：可以直接最大化二元似然，但运算相当繁杂. 补全下面论证的空缺，可以得到更简洁的证明.）

作变换 $u=x-y$、$v=x+y$. 设 $f(x,y)$ 为二元正态密度，写成
$$f(x,y)=g(v\mid u)h(u),$$
\sourcepageanchor{363}
其中 $g(v\mid u)$ 是给定 $U$ 后 $V$ 的条件密度，$h(u)$ 为 $U$ 的边缘密度. 说明：(1) 似然可等价地因子分解；(2) 涉及 $g(v\mid u)$ 的部分，无论均值是否受约束，都有相同最大值. 因此这一部分可以忽略（它会约掉），似然比检验只基于 $h(u)$. 而 $h(u)$ 是均值为 $\mu_X-\mu_Y$ 的正态密度，所以所得似然比检验就是\autoref{exercise:8.38}推导的通常单样本 $t$ 检验.\footnote{译注：作为原坐标 $(x,y)$ 的密度，严格的变换公式有雅可比因子 $|\partial(u,v)/\partial(x,y)|=2$，故应为 $f_{X,Y}(x,y)=2g(v\mid u)h(u)$. 原文省略的常数与参数无关，会在似然比中相消，因而不影响所述检验.}
\end{enumerate}
\end{exercise}
\begin{exercise}\label{exercise:8.41}\index{SI0640@双样本}\index{SI1024@双样本}\index{SI1025@双侧}
设 $X_1,\ldots,X_n$ 是来自 $\mathrm n(\mu_X,\sigma_X^2)$ 的随机样本，$Y_1,\ldots,Y_m$ 是来自 $\mathrm n(\mu_Y,\sigma_Y^2)$ 的独立随机样本. 希望在 $\sigma_X^2=\sigma_Y^2=\sigma^2$ 的假设下检验
$$H_0:\mu_X=\mu_Y\quad\text{对}\quad H_1:\mu_X\ne\mu_Y.$$
\begin{enumerate}
\item 推导似然比检验，并证明它可基于统计量
$$T=\frac{\bar X-\bar Y}{\sqrt{S_p^2(1/n+1/m)}},$$
其中
$$S_p^2=\frac1{n+m-2}\left(\sum_{i=1}^n(X_i-\bar X)^2+\sum_{i=1}^m(Y_i-\bar Y)^2\right).$$
（$S_p^2$ 有时称为合并方差估计（pooled variance estimate），\autoref{sec:11.2}将广泛使用这种估计.）
\item 证明，在 $H_0$ 下，$T\sim t_{n+m-2}$. （这个检验称为双样本 $t$ 检验（two-sample $t$ test\mbox{）.}）
\item 从某座拜占庭教堂的核心和外围取得木材样本，测定年代得到以下数据：
\begin{center}\begin{tabular}{rr@{\qquad}rr}\toprule
\multicolumn{2}{c}{核心}&\multicolumn{2}{c}{外围}\\\cmidrule{1-2}\cmidrule{3-4}
1294&1251&1284&1274\\
1279&1248&1272&1264\\
1274&1240&1256&1256\\
1264&1232&1254&1250\\
1263&1220&1242&\\
1254&1218&&\\
1251&1210&&\\\bottomrule
\end{tabular}\end{center}
用双样本 $t$ 检验判断核心的平均年代是否与外围相同.
\end{enumerate}
\end{exercise}
\begin{exercise}\label{exercise:8.42}\index{SI0087@Behrens–Fisher问题}\index{SI0952@Satterthwaite近似}\index{SI1018@近似}
\autoref{exercise:8.41}的方差相等假设并非总能成立. 若不成立，统计量不再服从 $t$ 分布，计算合并方差估计是否明智也值得怀疑. （方差不等时对均值作推断，一般相当困难，称为 Behrens–Fisher 问题.）自然可尝试如下双样本 $t$ 检验的修正：不假定 $\sigma_X^2=\sigma_Y^2$，使用统计量
$$T'=\frac{\bar X-\bar Y}{\sqrt{S_X^2/n+S_Y^2/m}}$$
检验 $H_0:\mu_X=\mu_Y$ 对 $H_1:\mu_X\ne\mu_Y$，其中
$$S_X^2=\frac1{n-1}\sum_{i=1}^n(X_i-\bar X)^2,\qquad
S_Y^2=\frac1{m-1}\sum_{i=1}^m(Y_i-\bar Y)^2.$$
$T'$ 的精确分布不好处理，但可以用 Satterthwaite 近似（\autoref{ex:7.2.3}）来近似其分布.
\sourcepageanchor{364}
\begin{enumerate}
\item 证明近似地有
$$\frac{S_X^2/n+S_Y^2/m}{\sigma_X^2/n+\sigma_Y^2/m}\sim\frac{\chi_\nu^2}{\nu},$$
其中 $\nu$ 可由下式估计：
$$\hat\nu=\frac{(S_X^2/n+S_Y^2/m)^2}{S_X^4/[n^2(n-1)]+S_Y^4/[m^2(m-1)]}.$$
\item 说明 $T'$ 的分布可用自由度为 $\hat\nu$ 的 $t$ 分布近似.
\item 用本题的近似 $t$ 检验重新考察\autoref{exercise:8.41}的数据，即用 $T'$ 检验核心的平均年代是否与外围相同.
\item 是否有统计证据表明核心数据的方差可能不同于外围数据的方差？（回顾\hyperref[ex:5.3.5]{例~5.4.1}\footnote{译注：原文确为“例 5.4.1”，但本书该编号对应的是次序统计量的\autoref{def:5.4.1}，并无此编号的例. 这里所需的正态样本方差比及其 $F$ 分布见\autoref{ex:5.3.5}、\autoref{def:5.3.6}和\autoref{ex:5.3.7}. 正文保留原印编号，链接导向相关的方差比示例.}.）
\end{enumerate}
\end{exercise}
\begin{exercise}\label{exercise:8.43}
\hyperref[bib:source298]{Sprott and Farewell (1993)} 指出，在双样本 $t$ 检验中，只要方差比已知，就能导出有效的 $t$ 统计量. 设 $X_1,X_2,\ldots,X_{n_1}$ 是来自 $\mathrm N(\mu_1,\sigma^2)$ 的样本，$Y_1,Y_2,\ldots,Y_{n_2}$ 是来自 $\mathrm n(\mu_2,\rho^2\sigma^2)$ 的样本，$\rho^2$ 已知. 证明
$$\frac{(\bar X-\bar Y)-(\mu_1-\mu_2)}{
\sqrt{\frac1{n_1}+\frac{\rho^2}{n_2}}\,
\sqrt{\frac{n_1(n_1-1)s_X^2+n_2(n_2-1)s_Y^2/\mathrm{rho}^2}{n_1+n_2-2}}}$$
服从自由度为 $n_1+n_2-2$ 的 Student $t$ 分布，且 $n_2s_Y^2/(\rho^2n_1s_X^2)$ 服从自由度为 $n_1-1$ 与 $n_2-1$ 的 $F$ 分布.

Sprott 和 Farewell 还指出，$t$ 统计量在
$$\rho^2=\frac{n_1\sqrt{n_1-1}\,s_X^2}{n_2\sqrt{n_2-1}\,s_Y^2}$$
处取最大值，并建议对 $\rho^2$ 的合理取值（例如某个置信区间内的值）作出该统计量的图像.\footnote{译注：原式确实印有字母串“rho”，且本题未另外定义 $s_X^2,s_Y^2$. 若沿用通常的样本方差定义，$t$ 式合并方差中的额外因子 $n_1,n_2$ 应去掉，标准化方差比应为 $s_Y^2/(\rho^2s_X^2)$，其分子、分母自由度依次应为 $n_2-1,n_1-1$. 即使把 $s_X^2,s_Y^2$ 另释为样本均值的方差估计，原文所列 $F$ 自由度次序仍不符. 最后关于 $\rho^2$ 的极值公式也应另行推导核对；有符号 $t$ 的“最大”还取决于分子的符号. 这里保留原公式，避免静默采用不明方差约定.}
\end{exercise}
\begin{exercise}\label{exercise:8.44}
验证\autoref{ex:8.3.20}的检验3是无偏的水平 $\alpha$ 检验.
\end{exercise}
\begin{exercise}\label{exercise:8.45}
设 $X_1,\ldots,X_n$ 是来自 $\mathrm n(\theta,\sigma^2)$ 总体的随机样本. 考虑检验
$$H_0:\theta\le\theta_0\quad\text{对}\quad H_1:\theta>\theta_0.$$
对 $m=1,\ldots,n$，用 $\bar X_m$ 表示前 $m$ 个观测 $X_1,\ldots,X_m$ 的样本均值. 若 $\sigma^2$ 已知，证明对每个 $m=1,\ldots,n$，在
$$\bar X_m>\theta_0+z_\alpha\sqrt{\sigma^2/m}$$
时拒绝 $H_0$ 的检验是无偏的大小 $\alpha$ 检验. 当 $n=4$ 时，画出这些检验各自的功效函数.
\end{exercise}
\begin{exercise}\label{exercise:8.46}
设 $X_1,\ldots,X_n$ 是来自 $\mathrm n(\theta,\sigma^2)$ 总体的随机样本. 考虑检验
$$H_0:\theta_1\le\theta\le\theta_2\quad\text{对}\quad H_1:\theta<\theta_1\text{ 或 }\theta>\theta_2.$$
\begin{enumerate}
\item\phantomsection\label{itm:exercise:8.46:a} 证明检验
$$\text{若 }\bar X>\theta_2+t_{n-1,\alpha/2}\sqrt{S^2/n}\quad\text{或}\quad
\bar X<\theta_1-t_{n-1,\alpha/2}\sqrt{S^2/n},\quad\text{则拒绝 }H_0,$$
不是大小 $\alpha$ 的检验.
\item\phantomsection\label{itm:exercise:8.46:b} 证明，对恰当选择的常数 $k$，大小 $\alpha$ 的检验为
$$\text{若 }|\bar X-\bar\theta|>k\sqrt{S^2/n},\quad\text{则拒绝 }H_0,$$
其中 $\bar\theta=(\theta_1+\theta_2)/2$.
\item\phantomsection\label{itm:exercise:8.46:c} 证明\hyperref[itm:exercise:8.46:a]{(a)} 与\hyperref[itm:exercise:8.46:b]{(b)} 中的检验，就各自的大小而言都是无偏的.\footnote{译注：原题对干扰参数 $\sigma$ 的处理未说清楚，按通常在完整原假设参数空间 $(\theta,\sigma)$ 上取上确界的“大小”定义，题面存在问题. 若 $\theta_1<\theta_2$、$n\ge2$、$0<\alpha<1$，\hyperref[itm:exercise:8.46:a]{(a)} 的拒绝概率不超过 $\alpha$，而 $\sigma\to\infty$ 时趋于 $\alpha$，故大小恰为 $\alpha$；\hyperref[itm:exercise:8.46:b]{(b)} 对任意有限固定 $k>0$，取 $\theta=\theta_2$、$\sigma\downarrow0$，则拒绝概率趋于 $1$，故大小为 $1$. \hyperref[itm:exercise:8.46:c]{(c)} 的无偏性也不能在这一完整参数空间中直接成立：对 \hyperref[itm:exercise:8.46:a]{(a)}，在固定有限 $\sigma$ 下，原假设端点及其附近的某些备择点功效仍小于实际大小 $\alpha$；对 \hyperref[itm:exercise:8.46:b]{(b)} 的 $k>0$，有限参数值下功效小于实际大小 $1$. 若 $k\le0$，\hyperref[itm:exercise:8.46:b]{(b)} 是以概率 $1$ 拒绝的平凡检验，也不能给出所要求的大小 $\alpha$. 若作者意指固定 $\sigma$、仅对 $\theta$ 作比较，则需另外明确这一条件范围，并允许临界常数依赖该固定 $\sigma$；这里保留原题.} （假定非中心 $t$ 分布具有单调似然比.）
\end{enumerate}

\end{exercise}
\sourcepageanchor{365}
\begin{exercise}\label{exercise:8.47}\index{SI0490@生物等效性}
考虑\autoref{exercise:8.41}中的两个方差相等、相互独立的正态样本. 检验 $H_0:\mu_X-\mu_Y\le-\delta\text{ 或 }\mu_X-\mu_Y\ge\delta$ 对 $H_1:-\delta<\mu_X-\mu_Y<\delta$，其中 $\delta$ 为指定正常数. （这称为等效性检验（equivalence testing）问题.）
\begin{enumerate}
\item\phantomsection\label{itm:exercise:8.47:a} 证明，$H_0^-:\mu_X-\mu_Y\le-\delta$ 对 $H_1^-:\mu_X-\mu_Y>-\delta$ 的大小 $\alpha$ 似然比检验，在
$$T^-=\frac{\bar X-\bar Y-(-\delta)}{\sqrt{S_p^2(1/n+1/m)}}\ge t_{n+m-2,\alpha}$$
时拒绝 $H_0^-$.
\item\phantomsection\label{itm:exercise:8.47:b} 求 $H_0^+:\mu_X-\mu_Y\ge\delta$ 对 $H_1^+:\mu_X-\mu_Y<\delta$ 的大小 $\alpha$ 似然比检验.
\item\phantomsection\label{itm:exercise:8.47:c} 说明如何把\hyperref[itm:exercise:8.47:a]{(a)}、\hyperref[itm:exercise:8.47:b]{(b)} 的检验组合成 $H_0$ 对 $H_1$ 的水平 $\alpha$ 检验.
\item 证明\hyperref[itm:exercise:8.47:c]{(c)} 的检验大小为 $\alpha$. （提示：考虑 $\sigma\to0$.）
\end{enumerate}
这种程序有时称为两个单侧检验程序（two one-sided tests procedure），由 \hyperref[bib:source278]{Schuirmann (1987)} 为生物等效性检验问题推导（另见 Westlake 1972\mbox{）.} 亦可参见 \hyperref[bib:source27]{Berger and Hsu (1996)} 的综述，以及给出相应置信区间的\autoref{exercise:9.33}.
\end{exercise}
\begin{exercise}\label{exercise:8.48}
证明\autoref{ex:8.3.30}的断言：给定 $S$ 后 $S_1$ 的条件分布为超几何分布.
\end{exercise}
\begin{exercise}\label{exercise:8.49}
在以下各情形中，计算已观察数据的 $p$ 值.
\begin{enumerate}
\item 检验 $H_0:\theta\le1/2$ 对 $H_1:\theta>1/2$，在 $10$ 次伯努利试验中观察到 $7$ 次成功.
\item 检验 $H_0:\lambda\le1$ 对 $H_1:\lambda>1$，观测到 $X=3$，其中 $X\sim\operatorname{Poisson}(\lambda)$.
\item 检验 $H_0:\lambda\le1$ 对 $H_1:\lambda>1$，观测到 $X_1=3,X_2=5,X_3=1$，其中 $X_i\sim\operatorname{Poisson}(\lambda)$，相互独立.
\end{enumerate}
\end{exercise}
\begin{exercise}\label{exercise:8.50}
设 $X_1,\ldots,X_n$ 独立同分布于 $\mathrm n(\theta,\sigma^2)$，$\sigma^2$ 已知，$\theta$ 服从双指数分布，即 $\pi(\theta)=e^{-|\theta|/a}/(2a)$，$a$ 已知. 检验 $H_0:\theta\le0$ 对 $H_1:\theta>0$ 的贝叶斯检验，在 $H_1$ 的后验概率较大时支持 $H_1$.
\begin{enumerate}
\item 对给定常数 $K$，计算 $\theta>K$ 的后验概率，即 $P(\theta>K\mid x_1,\ldots,x_n,a)$.
\item\phantomsection\label{itm:exercise:8.50:b} 求 $\lim_{a\to\infty}P(\theta>K\mid x_1,\ldots,x_n,a)$ 的表达式.
\item 将\hyperref[itm:exercise:8.50:b]{(b)} 的答案与经典假设检验对应的 $p$ 值比较.
\end{enumerate}
\end{exercise}
\begin{exercise}\label{exercise:8.51}
这里给出 $p$ 值的另一种常见解释. 考虑检验 $H_0$ 对 $H_1$，设 $W(\boldsymbol X)$ 为检验统计量. 假设对每个 $0\le\alpha\le1$，都能选择临界值 $c_\alpha$，使 $\{\boldsymbol x:W(\boldsymbol x)\ge c_\alpha\}$ 为 $H_0$ 的某大小 $\alpha$ 检验的拒绝域. 利用这一检验族，证明\eqref{eq:8.3.9}定义的通常 $p$ 值 $p(\boldsymbol x)$，是在已观察数据 $\boldsymbol x$ 后能够拒绝 $H_0$ 的最小 $\alpha$ 水平.
\end{exercise}
\begin{exercise}\label{exercise:8.52}
考虑检验 $H_0:\theta\in\bigcup_{j=1}^k\Theta_j$. 对每个 $j=1,\ldots,k$，设 $p_j(\boldsymbol x)$ 是检验 $H_{0j}:\theta\in\Theta_j$ 的有效 $p$ 值. 令 $p(\boldsymbol x)=\max_{1\le j\le k}p_j(\boldsymbol x)$.
\begin{enumerate}
\item 证明 $p(\boldsymbol X)$ 是检验 $H_0$ 的有效 $p$ 值.
\item 证明 $p(\boldsymbol X)$ 定义的水平 $\alpha$ 检验，等同于由各 $p_j(\boldsymbol x)$ 所对应的单独检验构造的水平 $\alpha$ 交并检验.
\end{enumerate}
\end{exercise}
\begin{exercise}\label{exercise:8.53}\index{SI0489@p值}
\autoref{ex:8.2.7}给出了单侧贝叶斯假设检验的例子. 现在考虑类似情形，但使用双侧检验. 希望检验
$$H_0:\theta=0\quad\text{对}\quad H_1:\theta\ne0,$$
观察来自 $\mathrm n(\theta,\sigma^2)$ 总体的随机样本 $X_1,\ldots,X_n$，$\sigma^2$ 已知. 此时常用的一类先验是 $\theta=0$ 处的点质量与分散在 $H_1$ 上的密度的混合. 典型选择为 $P(\theta=0)=1/2$，当 $\theta\ne0$ 时，取先验分布为 $\frac12\mathrm n(0,\tau^2)$，$\tau^2$ 已知.\footnote{译注：原文“当 $\theta\ne0$ 时”之后仍写 $\tfrac12\mathrm n(0,\tau^2)$，若将其理解为条件分布，则总质量只有 $1/2$，不是概率分布. 应区分：无条件先验的连续部分密度权重为 $1/2$；给定非零分支后的条件分布则是完整的 $\mathrm n(0,\tau^2)$. 此处保留原写法，并按题意明确混合权重的含义.}
\begin{enumerate}
\item 证明上述先验正常，即 $P(-\infty<\theta<\infty)=1$.
\item 计算 $H_0$ 为真的后验概率 $P(\theta=0\mid x_1,\ldots,x_n)$.
\item 求对应于给定 $\bar x$ 值的 $p$ 值表达式.
\item 当 $\sigma^2=\tau^2=1$ 时，在一段 $\bar x$ 取值范围上比较 $P(\theta=0\mid x_1,\ldots,x_n)$ 与 $p$ 值. 特别地，
\sourcepageanchor{366}
\begin{enumerate}[label=\roman*.]
\item 对 $n=9$，把 $p$ 值和后验概率作为 $\bar x$ 的函数作图，证明在 $\bar x$ 适度偏大时，贝叶斯概率大于 $p$ 值.
\item 取 $\alpha=.05$，令 $\bar x=Z_{\alpha/2}/\sqrt n$，使所有 $n$ 下的 $p$ 值都固定为 $\alpha$. 证明当 $n\to\infty$ 时，在 $\bar x=Z_{\alpha/2}/\sqrt n$ 处的后验概率趋于 $1$. 这就是 Lindley 悖论（Lindley's Paradox\mbox{）.}
\end{enumerate}
注意，$P(\theta=0\mid x_1,\ldots,x_n)$ 的小值是反对 $H_0$ 的证据，因此这个量在意义上类似于 $p$ 值. \hyperref[bib:source204]{Lindley (1957)} 指出了这两个量可能差别很大，\hyperref[bib:source21]{Berger and Sellke (1987)} 也考察了这一事实（见\autoref{sec:8.5}的杂录\mbox{）.}
\end{enumerate}
\end{exercise}
\begin{exercise}\label{exercise:8.54}
\hyperref[bib:source66]{Casella and Berger (1987)} 指出，在单侧问题中，$p$ 值与贝叶斯后验概率的差异没有那么显著；杂录中也提到了这一点. 设 $X_1,\ldots,X_n$ 是来自 $\mathrm n(\theta,\sigma^2)$ 总体的随机样本，需检验
$$H_0:\theta\le0\quad\text{对}\quad H_1:\theta>0.$$
$\theta$ 的先验分布为 $\mathrm n(0,\tau^2)$，$\tau^2$ 已知. 它对两个假设是对称的，即 $P(\theta\le0)=P(\theta>0)=1/2$.
\begin{enumerate}
\item\phantomsection\label{itm:exercise:8.54:a} 计算 $H_0$ 为真的后验概率 $P(\theta\le0\mid x_1,\ldots,x_n)$.
\item\phantomsection\label{itm:exercise:8.54:b} 使用 $\bar X$ 较大时拒绝的检验，求对应于给定 $\bar x$ 值的 $p$ 值表达式.
\item 在 $\sigma^2=\tau^2=1$ 的特殊情形，对 $\bar x>0$ 比较 $P(\theta\le0\mid x_1,\ldots,x_n)$ 与 $p$ 值，证明贝叶斯概率总大于 $p$ 值.
\item 利用\hyperref[itm:exercise:8.54:a]{(a)}、\hyperref[itm:exercise:8.54:b]{(b)} 所得表达式，证明
$$\lim_{\tau^2\to\infty}P(\theta\le0\mid x_1,\ldots,x_n)=p\text{ 值},$$
而双侧问题不会出现这个等式.
\end{enumerate}
\end{exercise}
\begin{exercise}\label{exercise:8.55}
设 $X\sim\mathrm n(\theta,1)$，考虑检验 $H_0:\theta\ge\theta_0$ 对 $H_1:\theta<\theta_0$. 使用损失函数\eqref{eq:8.3.13}，研究分别取 $\alpha=.1,.3,.5$、在 $X<-z_\alpha+\theta_0$ 时拒绝 $H_0$ 的三个检验.
\begin{enumerate}
\item\phantomsection\label{itm:exercise:8.55:a} 对 $b=c=1$，画图比较它们的风险函数.
\item\phantomsection\label{itm:exercise:8.55:b} 对 $b=3,c=1$，画图比较它们的风险函数.
\item 画出三个检验的功效函数，并与\hyperref[itm:exercise:8.55:a]{(a)}、\hyperref[itm:exercise:8.55:b]{(b)} 的风险函数比较.
\end{enumerate}
\end{exercise}
\begin{exercise}\label{exercise:8.56}
设 $X\sim\operatorname{binomial}(5,p)$，在 $0$–$1$ 损失下考虑检验 $H_0:p\le1/3$ 对 $H_1:p>1/3$. 画图比较以下两个检验的风险函数：检验 I 在 $X=0$ 或 $1$ 时拒绝 $H_0$；检验 II 在 $X=4$ 或 $5$ 时拒绝 $H_0$.
\end{exercise}
\begin{exercise}\label{exercise:8.57}
设 $X\sim\mathrm n(\mu,1)$，在 $0$–$1$ 损失下考虑检验 $H_0:\mu\le0$ 对 $H_1:\mu>0$. 设 $\delta_c$ 为在 $X>c$ 时拒绝 $H_0$ 的检验. 对本问题的每个检验，类 $\{\delta_c:-\infty\le c\le\infty\}$ 中都存在一个 $\delta_c$，其风险函数关于 $\mu$ 一致更小. 设 $\delta$ 为在 $1<X<2$ 时拒绝 $H_0$ 的检验. 求一个优于 $\delta$ 的 $\delta_c$. （可以证明其优越性，也可以画出 $\delta$ 与 $\delta_c$ 的风险函数，并仔细说明为什么所提出的检验应当更好.）
\end{exercise}
\begin{exercise}\label{exercise:8.58}\index{SI0500@极小极大}
考虑\autoref{ex:8.3.31}的假设检验问题和损失函数，令 $\sigma=n=1$. 考虑在 $X<-z_\alpha+\theta_0$ 时拒绝 $H_0$ 的检验. 求使风险函数最大值最小的 $\alpha$，即得到极小极大检验（minimax test）的 $\alpha$.
\end{exercise}





\subsection{杂录}\label{sec:8.5}
\subsubsection{单调功效函数}\label{sec:8.5.1}\index{SI0503@单调}\index{SI0994@功效}
本章广泛使用了单调似然比性质，特别是用它研究检验功效函数的性质. 随机序（stochastic ordering）的概念也能用于推导功效函数的性质. （回顾，前几章已遇到随机序，例如\autoref{exercise:1.49}、\autoref{exercise:3.41}至\autoref{exercise:3.43}、\autoref{exercise:5.19}. 若对所有 $x$，$F(x)\le G(x)$，且至少某个 $x$ 取严格不等式，则称累积分布函数 $F$ 在随机意义下大于 $G$. 这意味着，若 $X\sim F,Y\sim G$，则对所有 $x$，$P(X>x)\ge P(Y>x)$，且至少某个 $x$ 取严格不等式. 换言之，$F$ 把更多概率赋予较大的值.）
\sourcepageanchor{367}
从假设检验的角度看，备择假设下的分布常在随机意义下大于原假设下的分布. 例如，若有来自 $\mathrm n(\theta,\sigma^2)$ 总体的随机样本，希望检验 $H_0:\theta\le\theta_0$ 对 $H_1:\theta>\theta_0$，则备择假设中的所有分布确实都在随机意义下大于原假设中的所有分布. \hyperref[bib:source142]{Gilat (1977)} 利用随机序而非单调似然比，在一般条件下证明功效函数的单调性.

\subsubsection{似然比作为证据}\label{sec:8.5.2}\index{SI0487@贝叶斯}
似然比 $L(\theta_1\mid\boldsymbol x)/L(\theta_0\mid\boldsymbol x)=f(\boldsymbol x\mid\theta_1)/f(\boldsymbol x\mid\theta_0)$，在检验 $H_0:\theta=\theta_0$ 对 $H_1:\theta=\theta_1$ 时起重要作用. 对产生较小 $\lambda$ 值的那些 $\boldsymbol x$，这一比值等于似然比检验统计量 $\lambda(\boldsymbol x)$.\footnote{译注：此处原文把比值与其倒数混淆. 在只含 $\theta_0,\theta_1$ 的简单对简单模型中，按\autoref{def:8.2.1}，$\lambda=f_0/\max(f_0,f_1)$；当 $\lambda<1$ 时，$f_1/f_0=1/\lambda$，并非 $\lambda$. 后面的后验优势比公式使用 $f_1/f_0$ 是正确的.} 此外，Neyman–Pearson 引理说明，$H_0$ 对 $H_1$ 的水平 $\alpha$ 一致最强检验可用该比值定义. 这一似然比还有重要的贝叶斯解释. 设 $\pi_0,\pi_1$ 分别为 $\theta_0,\theta_1$ 的先验概率，则支持 $\theta_1$ 的后验优势比（posterior odds）为
$$\frac{P(\theta=\theta_1\mid\boldsymbol x)}{P(\theta=\theta_0\mid\boldsymbol x)}
=\frac{f(\boldsymbol x\mid\theta_1)\pi_1/m(\boldsymbol x)}{f(\boldsymbol x\mid\theta_0)\pi_0/m(\boldsymbol x)}
=\frac{f(\boldsymbol x\mid\theta_1)}{f(\boldsymbol x\mid\theta_0)}\cdot\frac{\pi_1}{\pi_0}.$$
$\pi_1/\pi_0$ 是支持 $\theta_1$ 的先验优势比（prior odds\mbox{）.} 观察到数据 $\boldsymbol X=\boldsymbol x$ 后，似然比给出将先验优势比调整为后验优势比所需的倍数. 若似然比等于二，先验优势比就加倍. 似然比不依赖先验概率，所以可解释为数据中支持 $H_1$ 而非 $H_0$ 的证据. \hyperref[bib:source265]{Royall (1997)} 讨论了这种解释.

\subsubsection{\texorpdfstring{$p$}{p} 值与后验概率}\label{sec:8.5.3}\index{SI0744@与贝叶斯检验的关系}
在讨论贝叶斯检验的\autoref{sec:8.2.2}中，我们看到 $H_0$ 为真的后验概率可以度量数据反对（或支持）原假设的证据. \autoref{sec:8.3.4}又说明，$p$ 值可以度量数据反对 $H_0$ 的证据. 自然要问，这两种不同度量是否会一致，即能否协调起来？\hyperref[bib:source21]{Berger（James，而非 Roger）与 Sellke (1987)} 认为，在双侧检验问题中，这些度量无法协调，贝叶斯度量更优. \hyperref[bib:source66]{Casella 与 Berger（Roger，1987）}则认为，双侧贝叶斯问题是人为构造的，在更自然的单侧问题中，这些证据度量可以协调. 这种协调对 \hyperref[bib:source277]{Schervish (1996)} 意义不大；他认为，把 $p$ 值当作证据度量存在严重的逻辑缺陷.

\subsubsection{置信集 \texorpdfstring{$p$}{p} 值}\label{sec:8.5.4}
\hyperref[bib:source25]{Berger and Boos (1994)} 提出了计算 $p$ 值的另一种方法. 在\autoref{thm:8.3.27}的通常定义中，上确界取遍整个原假设空间 $\Theta_0$. Berger 与 Boos 建议在 $\Theta_0$ 的一个子集 $C$ 上取上确界. 集合 $C=C(\boldsymbol X)$ 由数据确定，并具有性质：若 $\theta\in\Theta_0$，则 $P_\theta(\theta\in C(\boldsymbol X))\ge1-\beta$. （关于 $C$ 这样的置信集，见\autoref{sec:9}.）置信集 $p$ 值（confidence set $p$-value）定义为
$$p_C(\boldsymbol x)=\sup_{\theta\in C(\boldsymbol x)}P_\theta(W(\boldsymbol X)\ge W(\boldsymbol x))+\beta.$$
Berger 与 Boos 证明 $p_C$ 是有效 $p$ 值.\footnote{译注：原式可能大于 $1$，而\autoref{def:8.3.26}要求 $p$ 值在 $[0,1]$ 内. 为与该定义严格一致，通常取 $\min\{1,p_C(\boldsymbol x)\}$；这一截断不改变 $0\le\alpha<1$ 时的拒绝事件. 原公式保留.}

$p_C$ 有两个潜在优点. 计算上的优点是，在较小集合 $C$ 上求上确界，可能比在较大的 $\Theta_0$ 上容易. 统计上的优点是，观察 $\boldsymbol X$ 后，我们对 $\theta$ 的值已有一些认识，$\theta\in C$ 有较大机会成立；考虑那些看起来不可能为真的 $\theta$ 值似乎没有必要. 置信集 $p$ 值只考察 $\Theta_0$ 中看起来可信的取值. \hyperref[bib:source25]{Berger and Boos (1994)} 与 \hyperref[bib:source291]{Silvapulle (1996)} 给出了许多例子. \hyperref[bib:source24]{Berger (1996)} 指出，在比较两个二项概率的问题中，置信集 $p$ 值能产生功效有所提高的检验.





\section{区间估计}\label{sec:9}
\chapterdoi{9}
\sourcepageanchor{368}
\begin{quote}
“恐怕，”福尔摩斯说，“如果这件事超出了人类的能力，那当然也超出了我的能力. 不过，在求助于这样的理论之前，我们必须穷尽一切自然的解释.”
\begin{flushright}夏洛克·福尔摩斯\\《魔鬼之足》\end{flushright}
\end{quote}

\subsection{引言}\label{sec:9.1}
在\autoref{sec:7}，我们讨论了参数$\theta$的点估计，即以单个数值作为$\theta$的值所作的推断. 本章讨论区间估计（interval estimation），以及更一般的集合估计（set estimation）. 集合估计问题中的推断是断言“$\theta\in C$”，其中$C\subset\Theta$，且$C=C(\boldsymbol x)$是由观测到的数据值$\boldsymbol X=\boldsymbol x$确定的集合. 如果$\theta$取实数值，我们通常希望集合估计$C$是一个区间. 区间估计量是本章的主要论题.

与前两章一样，本章分为两部分，第一部分研究如何寻找区间估计量，第二部分研究如何评价估计量的优劣. 我们从区间估计量的形式定义开始，这个定义与点估计量的定义一样宽泛.
\begin{definition}\label{def:9.1.1}\index{SI0179@置信区间（另见置信集；置信程序；可信集）}
实值参数$\theta$的区间估计，是样本的任意一对函数$L(x_1,\ldots,x_n)$和$U(x_1,\ldots,x_n)$，它们对所有$\boldsymbol x\in\mathcal X$均满足$L(\boldsymbol x)\le U(\boldsymbol x)$. 如果观测到$\boldsymbol X=\boldsymbol x$，便作出推断$L(\boldsymbol x)\le\theta\le U(\boldsymbol x)$. 随机区间$[L(\boldsymbol X),U(\boldsymbol X)]$称为区间估计量（interval estimator）.
\end{definition}
我们沿用前面规定的记号，以$[L(\boldsymbol X),U(\boldsymbol X)]$表示基于随机样本$\boldsymbol X=(X_1,\ldots,X_n)$的$\theta$的区间估计量，以$[L(\boldsymbol x),U(\boldsymbol x)]$表示该区间的实现值. 虽然大多数情形下我们研究的$L$和$U$均取有限值，但有时也关注单侧区间估计.\index{SI0165@单侧}\semanticindexanchor{165:368} 例如，若$L(\boldsymbol x)=-\infty$，就得到单侧区间$(-\infty,U(\boldsymbol x)]$，相应断言为“$\theta\le U(\boldsymbol x)$”，并不涉及下界. 类似地，令$U(\boldsymbol x)=\infty$便得到单侧区间$[L(\boldsymbol x),\infty)$.

虽然定义中使用闭区间$[L(\boldsymbol x),U(\boldsymbol x)]$，但有时使用开区间$(L(\boldsymbol x),U(\boldsymbol x))$，甚至像上一段那样使用半开半闭区间，会更加自然. 我们将采用最适合具体问题的形式，不过通常优先采用闭区间.
\begin{example}[区间估计量]\label{ex:9.1.2}\index{SI0163@正态}
对于来自$\mathrm n(\mu,1)$的样本$X_1,X_2,X_3,X_4$，$\mu$的一个区间估计量是$[\bar X-1,\bar X+1]$. 这意味着我们将断言$\mu$在这个区间内.
\end{example}
这时，自然要问使用区间估计量能带来什么好处. 以前我们以$\bar X$估计$\mu$，现在却采用精度较低的估计量$[\bar X-1,\bar X+1]$. 我们肯定应当有所收获！牺牲估计（或关于$\mu$的断言）的一部分精度，使我们对断言的正确性获得了某种置信或保证.
\sourcepageanchor{369}
\begin{example}[续\autoref{ex:9.1.2}]\label{ex:9.1.3}
以$\bar X$估计$\mu$时，估计恰好正确的概率，即$P(\bar X=\mu)$，等于零. 但采用区间估计量，正确的概率就是正的. 区间$[\bar X-1,\bar X+1]$覆盖$\mu$的概率可计算如下：
$$\begin{aligned}
P(\mu\in[\bar X-1,\bar X+1])
&=P(\bar X-1\le\mu\le\bar X+1)\\
&=P(-1\le\bar X-\mu\le1)\\
&=P\left(-2\le\frac{\bar X-\mu}{\sqrt{1/4}}\le2\right)\\
&=P(-2\le Z\le2)\qquad\left(\frac{\bar X-\mu}{\sqrt{1/4}}\text{服从标准正态分布}\right)\\
&=.9544.
\end{aligned}$$
因此，区间估计量覆盖未知参数的机会超过$95\%$. 从点估计改为区间估计，牺牲一些估计精度，使我们对断言正确性的信心增加了.
\end{example}
使用区间估计量而非点估计量，目的是使我们对捕获感兴趣的参数有某种保证. 以下定义对这种保证的确定程度加以量化.
\begin{definition}\label{def:9.1.4}\index{SI0224@覆盖概率}
对于参数$\theta$的区间估计量$[L(\boldsymbol X),U(\boldsymbol X)]$，其覆盖概率（coverage probability）是随机区间$[L(\boldsymbol X),U(\boldsymbol X)]$覆盖真实参数$\theta$的概率. 用符号表示为$P_\theta(\theta\in[L(\boldsymbol X),U(\boldsymbol X)])$或$P(\theta\in[L(\boldsymbol X),U(\boldsymbol X)]\mid\theta)$.
\end{definition}
\begin{definition}\label{def:9.1.5}\index{SI0127@置信系数}
对于参数$\theta$的区间估计量$[L(\boldsymbol X),U(\boldsymbol X)]$，其置信系数（confidence coefficient）是覆盖概率的下确界$\inf_\theta P_\theta(\theta\in[L(\boldsymbol X),U(\boldsymbol X)])$.
\end{definition}
关于这些定义，有几点需要注意. 首先，务必牢记随机的是区间，而不是参数. 因此，当我们写出$P_\theta(\theta\in[L(\boldsymbol X),U(\boldsymbol X)])$这样的概率陈述时，它针对的是$\boldsymbol X$，而不是$\theta$. 换言之，应把看似针对随机$\theta$的$P_\theta(\theta\in[L(\boldsymbol X),U(\boldsymbol X)])$，理解为与之在代数上等价、针对随机$\boldsymbol X$的陈述$P_\theta(L(\boldsymbol X)\le\theta,U(\boldsymbol X)\ge\theta)$.

区间估计量连同置信程度的度量（通常是置信系数），有时称为置信区间（confidence interval）.\index{SI0171@与集合估计量的关系}\semanticindexanchor{171:369}\index{SI0958@与置信区间的关系}\semanticindexanchor{958:369} 我们经常将这个术语与区间估计量互换使用. 虽然我们主要关注置信区间，但有时也研究更一般的集合. 当进行一般讨论、还不完全确定集合的具体形式时，我们将称之为置信集（confidence set）. 置信系数等于某个数值（比如$1-\alpha$）的置信集，简称为$1-\alpha$置信集.

另一个要点涉及覆盖概率与置信系数. 因为不知道$\theta$的真实值，我们只能保证覆盖概率达到其下确界，即置信系数. 在一些情形下，这没有关系，因为覆盖概率是$\theta$的常数函数. 但在另一些情形下，覆盖概率可能是随$\theta$变动颇大的函数.
\begin{example}[尺度均匀分布的区间估计量]\label{ex:9.1.6}\index{SI0177@均匀}
设$X_1,\ldots,X_n$为来自$\operatorname{uniform}(0,\theta)$总体的随机样本，并令$Y=\max\{X_1,\ldots,X_n\}$. 我们关心$\theta$的区间估计量，考虑两个候选估计量：$[aY,bY]$，$1\le a<b$，以及$[Y+c,Y+d]$，$0\le c<d$，其中
\sourcepageanchor{370}
$a,b,c,d$是给定的常数.（注意，$\theta$必然大于$y$.）对于第一个区间，有
$$\begin{aligned}
P_\theta(\theta\in[aY,bY])&=P_\theta(aY\le\theta\le bY)\\
&=P_\theta\left(\frac1b\le\frac Y\theta\le\frac1a\right)\\
&=P_\theta\left(\frac1b\le T\le\frac1a\right).\qquad(T=Y/\theta)
\end{aligned}$$
此前已经看到（\autoref{ex:7.3.13}），$f_Y(y)=ny^{n-1}/\theta^n$，$0\le y\le\theta$，所以$T$的pdf为$f_T(t)=nt^{n-1}$，$0\le t\le1$. 因此，
$$P_\theta\left(\frac1b\le T\le\frac1a\right)=\int_{1/b}^{1/a}nt^{n-1}\,dt=\left(\frac1a\right)^n-\left(\frac1b\right)^n.$$
第一个区间的覆盖概率与$\theta$的值无关，故$(1/a)^n-(1/b)^n$就是该区间的置信系数.

对于另一个区间，当$\theta\ge d$时，类似计算给出
$$\begin{aligned}
P_\theta(\theta\in[Y+c,Y+d])&=P_\theta(Y+c\le\theta\le Y+d)\\
&=P_\theta\left(1-\frac d\theta\le T\le1-\frac c\theta\right)\qquad(T=Y/\theta)\\
&=\int_{1-d/\theta}^{1-c/\theta}nt^{n-1}\,dt
=\left(1-\frac c\theta\right)^n-\left(1-\frac d\theta\right)^n.
\end{aligned}$$
此时覆盖概率依赖于$\theta$. 而且，直接计算可知，对任意常数$c,d$，
$$\lim_{\theta\to\infty}\left(1-\frac c\theta\right)^n-\left(1-\frac d\theta\right)^n=0,$$
这说明这个区间估计量的置信系数为零.
\end{example}



\subsection{寻找区间估计量的方法}\label{sec:9.2}
我们用五个小节介绍寻找估计量的方法. 这似乎意味着有五种不同的方法来寻找区间估计量. 实际并非如此：在操作上，随后四个小节介绍的所有方法都相同，均以反演检验统计量为策略. 最后一小节讨论贝叶斯区间，介绍一种不同的构造方法.\footnote{译注：原文写“五个小节”和“随后四个小节”，但本节实际只有四个小节，其中前三个讨论频率学派方法，第四个讨论贝叶斯区间. 此处保留原文措辞.}
\subsubsection{反演检验统计量}\label{sec:9.2.1}
假设检验与区间估计之间有非常紧密的对应关系. 事实上，一般地说，每个置信集都对应一个检验，反之亦然. 考虑下面的例子.
\begin{example}[反演正态检验]\label{ex:9.2.1}\index{SI0163@正态}
设$X_1,\ldots,X_n$独立同分布于$\mathrm n(\mu,\sigma^2)$，考虑检验$H_0:\mu=\mu_0$对$H_1:\mu\ne\mu_0$. 对于给定的水平$\alpha$，一个合理的检验（事实上是最强无偏检验）的拒绝域为$\{\boldsymbol x:|\bar x-\mu_0|>z_{\alpha/2}\sigma/\sqrt n\}$. 注意，对于满足$|\bar x-\mu_0|\le z_{\alpha/2}\sigma/\sqrt n$的样本点，我们接受$H_0$；等价地，
$$\bar x-z_{\alpha/2}\frac\sigma{\sqrt n}\le\mu_0\le\bar x+z_{\alpha/2}\frac\sigma{\sqrt n}.$$
\sourcepageanchor{371}
\par
\begin{figure}[htbp]\centering
\includegraphics[width=.64\textwidth]{assets/chapter9/fig-9.2.1.pdf}
\caption{置信区间与检验接受域之间的关系. 上方直线为$\bar x=\mu+z_{\alpha/2}\sigma/\sqrt n$，下方直线为$\bar x=\mu-z_{\alpha/2}\sigma/\sqrt n$.}\label{fig:9.2.1}
\end{figure}
\par
由于检验的大小为$\alpha$，这意味着$P(H_0\text{被拒绝}\mid\mu=\mu_0)=\alpha$；换一种说法，$P(H_0\text{被接受}\mid\mu=\mu_0)=1-\alpha$. 结合上述接受域的刻画，可写为
$$P\left(\bar X-z_{\alpha/2}\frac\sigma{\sqrt n}\le\mu_0\le\bar X+z_{\alpha/2}\frac\sigma{\sqrt n}\,\middle|\,\mu=\mu_0\right)=1-\alpha.$$
但这个概率陈述对\emph{每一个}$\mu_0$都成立. 因此，陈述
$$P_\mu\left(\bar X-z_{\alpha/2}\frac\sigma{\sqrt n}\le\mu\le\bar X+z_{\alpha/2}\frac\sigma{\sqrt n}\right)=1-\alpha$$
成立. 反演水平$\alpha$检验的接受域所得的区间$[\bar x-z_{\alpha/2}\sigma/\sqrt n,\bar x+z_{\alpha/2}\sigma/\sqrt n]$是一个$1-\alpha$置信区间.
\end{example}
我们已说明了置信集与检验之间的对应关系. 假设检验的接受域，即样本空间中接受$H_0:\mu=\mu_0$的集合，为
$$A(\mu_0)=\left\{(x_1,\ldots,x_n):\mu_0-z_{\alpha/2}\frac\sigma{\sqrt n}\le\bar x\le\mu_0+z_{\alpha/2}\frac\sigma{\sqrt n}\right\},$$
而置信区间，即参数空间中由$\mu$的合理取值构成的集合，为
$$C(x_1,\ldots,x_n)=\left\{\mu:\bar x-z_{\alpha/2}\frac\sigma{\sqrt n}\le\mu\le\bar x+z_{\alpha/2}\frac\sigma{\sqrt n}\right\}.$$
这两个集合由如下恒真式联系起来：
$$(x_1,\ldots,x_n)\in A(\mu_0)\iff\mu_0\in C(x_1,\ldots,x_n).$$
\autoref{fig:9.2.1}展示了正态双侧问题中检验与区间估计的对应关系. 从图中或许更容易看出，检验与区间所问的是同一问题，只是角度略有不同. 两种程序都在寻找样本统计量与总体参数之间的一致性. 假设检验固定参数，问哪些样本值（接受域）与这个固定值相容. 置信集固定样本值，问哪些参数值（置信区间）使这个样本值最为合理.

检验的接受域与置信集之间的对应关系一般都成立. 下面的定理给出这种对应关系的形式表述.
\sourcepageanchor{372}
\begin{theorem}\label{thm:9.2.2}\index{SI0187@与假设检验的关系}\index{SI0507@与置信集的关系}
对每个$\theta_0\in\Theta$，设$A(\theta_0)$为$H_0:\theta=\theta_0$的水平$\alpha$检验的接受域. 对每个$\boldsymbol x\in\mathcal X$，在参数空间中定义集合$C(\boldsymbol x)$为
\begin{equation}\label{eq:9.2.1}
C(\boldsymbol x)=\{\theta_0:\boldsymbol x\in A(\theta_0)\}.
\end{equation}
则随机集合$C(\boldsymbol X)$是一个$1-\alpha$置信集. 反之，设$C(\boldsymbol X)$是一个$1-\alpha$置信集. 对任意$\theta_0\in\Theta$，定义
$$A(\theta_0)=\{\boldsymbol x:\theta_0\in C(\boldsymbol x)\}.$$
则$A(\theta_0)$是$H_0:\theta=\theta_0$的水平$\alpha$检验的接受域.
\end{theorem}
\begin{proof}
对于第一部分，因为$A(\theta_0)$是水平$\alpha$检验的接受域，
$$P_{\theta_0}(\boldsymbol X\notin A(\theta_0))\le\alpha,\qquad\text{从而}\quad P_{\theta_0}(\boldsymbol X\in A(\theta_0))\ge1-\alpha.$$
由于$\theta_0$任意，可将$\theta_0$写为$\theta$. 上式结合\eqref{eq:9.2.1}表明，集合$C(\boldsymbol X)$的覆盖概率为
$$P_\theta(\theta\in C(\boldsymbol X))=P_\theta(\boldsymbol X\in A(\theta))\ge1-\alpha,$$
这说明$C(\boldsymbol X)$是一个$1-\alpha$置信集.

对于第二部分，以$A(\theta_0)$为接受域的$H_0:\theta=\theta_0$检验，其第一类错误概率为
$$P_{\theta_0}(\boldsymbol X\notin A(\theta_0))=P_{\theta_0}(\theta_0\notin C(\boldsymbol X))\le\alpha.$$
所以这是一个水平$\alpha$的检验.
\end{proof}
虽然人们常说反演一个检验来得到置信集，但\autoref{thm:9.2.2}清楚表明，实际反演的是一族检验，对每个$\theta_0\in\Theta$各有一个检验，从而得到一个置信集.

检验可以反演为置信集、反之亦然，这一事实在理论上很有趣，但\autoref{thm:9.2.2}真正有用的是第一部分. 构造水平$\alpha$接受域相对容易，构造置信集则较难. 所以，通过反演接受域获得置信集的方法相当有用. 我们已有的所有构造检验的技术，都可以立即应用于构造置信集.

在\autoref{thm:9.2.2}中，我们仅给出了原假设$H_0:\theta=\theta_0$. 对接受域的全部要求是
$$P_{\theta_0}(\boldsymbol X\in A(\theta_0))\ge1-\alpha.$$
实际上，通过检验反演构造置信集时，我们还会考虑某个备择假设，例如$H_1:\theta\ne\theta_0$或$H_1:\theta>\theta_0$. 备择假设决定$A(\theta_0)$的合理形式，而$A(\theta_0)$的形式又决定$C(\boldsymbol x)$的形状. 但请注意，我们谨慎地使用了“集合”而不是“区间”一词，因为不能保证检验反演所得的置信集是区间. 不过，在大多数情形下，单侧检验给出单侧区间，双侧检验给出双侧区间，形状奇特的接受域则给出形状奇特的置信集. 后面的例子将展示这一点.

被反演的检验的性质也会转移到置信集上（有时需适当修改）. 例如，无偏检验经反演会产生无偏置信集. 此外，更重要的是，既然我们知道寻找良好检验时可以将注意力限于充分统计量，那么寻找良好置信集时也可以将注意力限于充分统计量.

当直觉失效、我们不知道怎样的集合才算合理时，检验反演法尤其有用. 我们只须求助于构造合理检验的通用方法.
\sourcepageanchor{373}
\begin{example}[反演似然比检验]\label{ex:9.2.3}\index{SI0160@指数}
假设我们希望为$\operatorname{exponential}(\lambda)$总体的均值$\lambda$构造置信区间. 可通过反演$H_0:\lambda=\lambda_0$对$H_1:\lambda\ne\lambda_0$的水平$\alpha$检验得到这样的区间.

若取随机样本$X_1,\ldots,X_n$，似然比检验（LRT）统计量为
$$\frac{\lambda_0^{-n}e^{-\sum x_i/\lambda_0}}{\sup_\lambda\lambda^{-n}e^{-\sum x_i/\lambda}}
=\frac{\lambda_0^{-n}e^{-\sum x_i/\lambda_0}}{(\sum x_i/n)^{-n}e^{-n}}
=\left(\frac{\sum x_i}{n\lambda_0}\right)^n e^ne^{-\sum x_i/\lambda_0}.$$
固定$\lambda_0$时，接受域为
\begin{equation}\label{eq:9.2.2}
A(\lambda_0)=\left\{\boldsymbol x:\left(\frac{\sum x_i}{\lambda_0}\right)^ne^{-\sum x_i/\lambda_0}\ge k^*\right\},
\end{equation}
其中常数$k^*$选取为使$P_{\lambda_0}(\boldsymbol X\in A(\lambda_0))=1-\alpha$.（常数$e^n/n^n$已并入$k^*$.）这是样本空间中的集合，如\autoref{fig:9.2.2}所示. 反演这个接受域得到$1-\alpha$置信集
$$C(\boldsymbol x)=\left\{\lambda:\left(\frac{\sum x_i}\lambda\right)^ne^{-\sum x_i/\lambda}\ge k^*\right\}.$$
这是参数空间中的一个区间，如\autoref{fig:9.2.2}所示.
\par
\begin{figure}[htbp]\centering
\includegraphics[width=\textwidth]{assets/chapter9/fig-9.2.2.pdf}
\caption{\autoref{ex:9.2.3}的接受域和置信区间. 接受域为$A(\lambda_0)=\{\boldsymbol x:(\sum x_i/\lambda_0)^ne^{-\sum x_i/\lambda_0}\ge k^*\}$，置信区间为$C(\boldsymbol x)=\{\lambda:(\sum x_i/\lambda)^ne^{-\sum x_i/\lambda}\ge k^*\}$.}\label{fig:9.2.2}
\end{figure}
\par
定义$C(\boldsymbol x)$的表达式仅通过$\sum x_i$依赖$\boldsymbol x$. 因而置信区间可写为
\begin{equation}\label{eq:9.2.3}
C(\textstyle\sum x_i)=\{\lambda:L(\textstyle\sum x_i)\le\lambda\le U(\textstyle\sum x_i)\},
\end{equation}
其中函数$L,U$由以下约束确定：集合\eqref{eq:9.2.2}的概率为$1-\alpha$，并且
\begin{equation}\label{eq:9.2.4}
\left(\frac{\sum x_i}{L(\sum x_i)}\right)^ne^{-\sum x_i/L(\sum x_i)}
=\left(\frac{\sum x_i}{U(\sum x_i)}\right)^ne^{-\sum x_i/U(\sum x_i)}.
\end{equation}
\sourcepageanchor{374}
如果令
\begin{equation}\label{eq:9.2.5}
\frac{\sum x_i}{L(\sum x_i)}=a\quad\text{且}\quad\frac{\sum x_i}{U(\sum x_i)}=b,
\end{equation}
其中$a>b$为常数，则\eqref{eq:9.2.4}变成
\begin{equation}\label{eq:9.2.6}a^ne^{-a}=b^ne^{-b},\end{equation}
这个方程易于数值求解. 为具体说明，令$n=2$，并注意$\sum X_i\sim\operatorname{gamma}(2,\lambda)$以及$\sum X_i/\lambda\sim\operatorname{gamma}(2,1)$. 因此，由\eqref{eq:9.2.5}，置信区间变为$\{\lambda:\frac1a\sum x_i\le\lambda\le\frac1b\sum x_i\}$，其中$a,b$满足
$$P_\lambda\left(\frac1a\sum X_i\le\lambda\le\frac1b\sum X_i\right)=P\left(b\le\frac{\sum X_i}\lambda\le a\right)=1-\alpha,$$
而由\eqref{eq:9.2.6}，$a^2e^{-a}=b^2e^{-b}$. 于是
$$\begin{aligned}
P\left(b\le\frac{\sum X_i}\lambda\le a\right)&=\int_b^a te^{-t}\,dt\\
&=e^{-b}(b+1)-e^{-a}(a+1).\qquad\text{（分部积分）}
\end{aligned}$$
例如，为得到$90\%$置信区间，必须同时满足概率条件和约束. 保留三位小数，得到$a=5.480,b=.441$，置信系数为$.90006$. 因而
$$P_\lambda\left(\frac1{5.480}\sum X_i\le\lambda\le\frac1{.441}\sum X_i\right)=.90006.$$
\end{example}
反演$H_0:\theta=\theta_0$对$H_1:\theta\ne\theta_0$的LRT（\autoref{def:8.2.1}）得到区域，检验的形式为
$$\text{若 }\frac{L(\theta_0\mid\boldsymbol x)}{L(\hat\theta\mid\boldsymbol x)}\le k(\theta_0),\text{ 则接受 }H_0,$$
所得置信区域为
\begin{equation}\label{eq:9.2.7}\index{SI0156@似然}
\{\theta:L(\theta\mid\boldsymbol x)\ge k'(\boldsymbol x,\theta)\},
\end{equation}
其中某个函数$k'$给出$1-\alpha$的置信程度.\footnote{译注：上面“接受$H_0$”的似然比不等号原文为$\le$. 按\autoref{def:8.2.1}以及紧接的置信区域\eqref{eq:9.2.7}，此处应为$\ge$；这里保留原式.}

某些情形下（如正态分布和伽马分布），函数$k'$不依赖$\theta$. 此时似然区域有一种特别令人满意的解释：它由使似然最高的那些$\theta$值组成. 我们还将看到，在频率学派（\autoref{thm:9.3.2}）和贝叶斯学派（\autoref{cor:9.3.10}）中，最优性考虑也会导出这样的区间.

检验反演法完全一般，因为我们可以反演任意检验而得到置信集. 在\autoref{ex:9.2.3}中反演的是LRT，但也可以采用任何方法构造的检验. 此外，注意双侧检验的反演给出了双侧区间. 在接下来的例子中，我们反演单侧检验得到单侧区间.
\begin{example}[正态单侧置信界]\label{ex:9.2.4}\index{SI0167@正态}
设$X_1,\ldots,X_n$为来自$\mathrm n(\mu,\sigma^2)$总体的随机样本. 考虑为$\mu$构造$1-\alpha$置信上界，即希望得到形式为$C(\boldsymbol x)=(-\infty,U(\boldsymbol x)]$的置信区间. 为用\autoref{thm:9.2.2}获得这样的区间，我们反演$H_0:\mu=\mu_0$对$H_1:\mu<\mu_0$的单侧检验.（注意，这里利用$H_1$的规定确定置信区间的形式. $H_1$规定了“较大”的$\mu_0$值，
\sourcepageanchor{375}
所以置信集将包含“较小”的值，即小于某个界的值. 因而得到的是置信上界.）检验$H_0$对$H_1$的大小为$\alpha$的LRT在
$$\frac{\bar X-\mu_0}{S/\sqrt n}<-t_{n-1,\alpha}$$
时拒绝$H_0$（类似\autoref{ex:8.2.6}）. 因此该检验的接受域为
$$A(\mu_0)=\left\{\boldsymbol x:\bar x\ge\mu_0-t_{n-1,\alpha}\frac s{\sqrt n}\right\},$$
且$\boldsymbol x\in A(\mu_0)\iff\bar x+t_{n-1,\alpha}s/\sqrt n\ge\mu_0$. 根据\eqref{eq:9.2.1}，定义
$$C(\boldsymbol x)=\{\mu_0:\boldsymbol x\in A(\mu_0)\}=\left\{\mu_0:\bar x+t_{n-1,\alpha}\frac s{\sqrt n}\ge\mu_0\right\}.$$
由\autoref{thm:9.2.2}，随机集合$C(\boldsymbol X)=(-\infty,\bar X+t_{n-1,\alpha}S/\sqrt n]$是$\mu$的$1-\alpha$置信集. 我们看到，它确实具有置信上界所需的形式. 反演单侧检验得到了单侧置信区间.
\end{example}
\begin{example}[二项分布单侧置信界]\label{ex:9.2.5}\index{SI0166@二项}
作为较难的单侧置信区间例子，考虑对一系列伯努利试验的成功概率$p$给出$1-\alpha$置信下界. 即观测$X_1,\ldots,X_n$，其中$X_i\sim\operatorname{Bernoulli}(p)$，希望区间具有$(L(x_1,\ldots,x_n),1]$的形式，且$P_p(p\in(L(X_1,\ldots,X_n),1])\ge1-\alpha$.（正如将要看到的，所得到的区间在左端是开的.）

由于希望得到给出置信下界的单侧区间，我们考虑反演下列检验的接受域：
$$H_0:p=p_0\qquad\text{对}\qquad H_1:p>p_0.$$
为简化问题，我们知道可以基于$T=\sum_{i=1}^nX_i\sim\operatorname{binomial}(n,p)$进行检验，因为$T$对$p$充分.（参见杂录一节.）由于二项分布具有单调似然比（\autoref{exercise:8.25}），由Karlin--Rubin定理（\autoref{thm:8.3.17}），在$T>k(p_0)$时拒绝$H_0$的检验是其大小下的UMP检验. 对每个$p_0$，希望选择常数$k(p_0)$（可取整数），使检验的水平为$\alpha$. 除某些$p_0$值外，由于$T$的离散性，不能使检验大小恰好等于$\alpha$. 但我们选择$k(p_0)$，使检验的大小在不超过$\alpha$的条件下尽可能接近$\alpha$. 因而，$k(p_0)$定义为$0$与$n$之间同时满足以下不等式的整数：
\begin{equation}\label{eq:9.2.8}
\begin{gathered}
\sum_{y=0}^{k(p_0)}\binom ny p_0^y(1-p_0)^{n-y}\ge1-\alpha,\\
\text{且}\qquad\sum_{y=0}^{k(p_0)-1}\binom ny p_0^y(1-p_0)^{n-y}<1-\alpha.
\end{gathered}
\end{equation}
由二项分布的MLR性质，对任意$k=0,\ldots,n$，量
$$f(p_0\mid k)=\sum_{y=0}^k\binom ny p_0^y(1-p_0)^{n-y}$$
是$p_0$的递减函数（\autoref{exercise:8.26}）. 当然，$f(0\mid0)=1$，所以$k(0)=0$；在一段$p_0$值上，$f(p_0\mid0)$保持高于$1-\alpha$. 随后在某一点，$f(p_0\mid0)=1-\alpha$；对更大的$p_0$，有$f(p_0\mid0)<1-\alpha$. 因而在这一点，$k(p_0)$增至$1$.
\sourcepageanchor{376}
这一模式不断继续. 因此$k(p_0)$是整数值阶梯函数：它在一段$p_0$取值上保持不变，然后跳到下一个更大的整数. 由于$k(p_0)$是$p_0$的非递减函数，这就给出了置信下界.（置信上界参见\autoref{exercise:9.5}.）对$k(p_0)$求解\eqref{eq:9.2.8}中的不等式，便同时得到检验的接受域与置信集.

对每个$p_0$，接受域为$A(p_0)=\{t:t\le k(p_0)\}$，其中$k(p_0)$满足\eqref{eq:9.2.8}. 对每个$t$，置信集为$C(t)=\{p_0:t\le k(p_0)\}$. 不过，这个集合目前的形式在实际中用处不大. 虽然它在形式上正确，而且是$1-\alpha$置信集，但它由$p_0$隐式定义，而我们希望得到关于$p_0$的显式定义.

由于$k(p_0)$非递减，对于给定观测$T=t$，所有小于或等于某个值（记为$k^{-1}(t)$）的$p_0$均有$k(p_0)<t$. 在$k^{-1}(t)$处，$k(p_0)$向上跳至$t$，而对所有$p_0>k^{-1}(t)$都有$k(p_0)\ge t$.（注意，在$p_0=k^{-1}(t)$时，$f(p_0\mid t-1)=1-\alpha$. 所以$k(p_0)=t-1$仍满足\eqref{eq:9.2.8}. 只有$p_0>k^{-1}(t)$时才有$k(p_0)\ge t$.）于是置信集为
\begin{equation}\label{eq:9.2.9}
C(t)=\{p_0:t\le k(p_0)\}=\{p_0:p_0>k^{-1}(t)\},
\end{equation}
我们已构造出形式为$C(T)=(k^{-1}(T),1]$的$1-\alpha$置信下界.

数$k^{-1}(t)$可定义为
\begin{equation}\label{eq:9.2.10}
k^{-1}(t)=\sup\left\{p:\sum_{y=0}^{t-1}\binom ny p^y(1-p)^{n-y}\ge1-\alpha\right\}.
\end{equation}
应注意，$k^{-1}(t)$实际上并非$k(p_0)$的反函数，因为$k(p_0)$不是一一函数. 不过，\eqref{eq:9.2.8}和\eqref{eq:9.2.10}中的表达式给出了良好定义的$k$与$k^{-1}$.

\hyperref[bib:source79]{Clopper和Pearson（1934）}最早研究了二项分布置信界问题，对双侧区间得到了类似的答案（\autoref{exercise:9.21}），并由此开创了一条至今仍活跃的研究方向. 参见\autoref{sec:9.5.2}.
\end{example}
\subsubsection{枢轴量}\label{sec:9.2.2}
\autoref{ex:9.1.6}中的两个置信区间在许多方面不同. 一个重要区别是，区间$\{aY,bY\}$的覆盖概率不依赖参数$\theta$的值，而$\{Y+c,Y+d\}$的覆盖概率却依赖它. 原因在于$\{aY,bY\}$的覆盖概率能够通过量$Y/\theta$表示；这个随机变量的分布不依赖参数，称为枢轴量（pivotal quantity），简称枢轴（pivot）.

利用枢轴量构造置信集，由此产生所谓枢轴推断（pivotal inference），主要归功于G. A. Barnard（\hyperref[bib:source10]{1949}，\hyperref[bib:source11]{1980}），但也可追溯到\hyperref[bib:source124]{Fisher（1930）}；Fisher使用的术语是逆概率（inverse probability）. 与之密切相关的是D. A. S. Fraser的结构推断（structural inference）理论\index{SI0996@结构推断}\semanticindexanchor{996:376}（Fraser \hyperref[bib:source127]{1968}，\hyperref[bib:source128]{1979}）. \hyperref[bib:source22]{Berger和Wolpert（1984）}对这些方法的优缺点作了有趣的讨论.
\begin{definition}\label{def:9.2.6}\index{SI0756@枢轴量}
若随机变量$Q(\boldsymbol X,\theta)=Q(X_1,\ldots,X_n,\theta)$的分布不依赖任何参数，就称其为枢轴量（或枢轴）. 即若$\boldsymbol X\sim F(\boldsymbol x\mid\theta)$，则$Q(\boldsymbol X,\theta)$对所有$\theta$值都有相同分布.
\end{definition}
函数$Q(\boldsymbol x,\theta)$通常显式地同时包含参数和统计量，但对任意集合$\mathcal A$，$P_\theta(Q(\boldsymbol X,\theta)\in\mathcal A)$不能依赖$\theta$. 从枢轴量构造置信集的技术，依赖于能找到一个枢轴量及一个集合$\mathcal A$，使得$\{\theta:Q(\boldsymbol x,\theta)\in\mathcal A\}$是$\theta$的集合估计.
\begin{example}[位置尺度枢轴量]\label{ex:9.2.7}\index{SI0646@枢轴量}\index{SI0649@枢轴量}\index{SI0759@位置族}\index{SI0760@位置尺度族}\index{SI0761@尺度族}\index{SI0954@枢轴量}
在位置与尺度情形中有很多枢轴量. 这里展示几个，更多例子见\autoref{exercise:9.8}. 设$X_1,\ldots,X_n$为来自所列pdf的随机样本，$\bar X,S$分别为样本均值和标准差. 为证明上述量是枢轴量，只须证明它们的pdf不依赖参数（详见\autoref{exercise:9.9}）.
\sourcepageanchor{377}
\par
\begin{table}[htbp]\centering
\caption{位置尺度枢轴量}\label{tab:9.2.1}
\begin{tabular}{lll}\toprule
pdf的形式&pdf的类型&枢轴量\\\midrule
$f(x-\mu)$&位置&$\bar X-\mu$\\
$\frac1\sigma f(\frac x\sigma)$&尺度&$\bar X/\sigma$\\
$\frac1\sigma f(\frac{x-\mu}\sigma)$&位置尺度&$(\bar X-\mu)/S$\\\bottomrule
\end{tabular}
\end{table}
\par
特别地，注意若$X_1,\ldots,X_n$为来自$\mathrm n(\mu,\sigma^2)$总体的随机样本，则$t$统计量$(\bar X-\mu)/(S/\sqrt n)$是枢轴量，因为$t$分布不依赖参数$\mu$和$\sigma^2$.
\end{example}
在\autoref{sec:9.2.1}用检验反演法构造的区间中，有些最终基于枢轴量（\autoref{ex:9.2.3}和\autoref{ex:9.2.4}），有些则不是（\autoref{ex:9.2.5}）. 寻找枢轴量没有通用策略. 不过，我们可以略加巧思，不必完全依赖猜测. 例如，寻找位置参数或尺度参数的枢轴量相对容易. 一般地，位置问题中的\emph{差}是枢轴量，尺度问题中的\emph{比}（或积）是枢轴量.
\begin{example}[伽马枢轴量]\label{ex:9.2.8}\index{SI0757@伽马}
设$X_1,\ldots,X_n$独立同分布于$\operatorname{exponential}(\lambda)$. 则$T=\sum X_i$是$\lambda$的充分统计量，且$T\sim\operatorname{gamma}(n,\lambda)$. 在伽马pdf中，$t$与$\lambda$以$t/\lambda$的形式共同出现；事实上，$\operatorname{gamma}(n,\lambda)$的pdf $(\Gamma(n)\lambda^n)^{-1}t^{n-1}e^{-t/\lambda}$是尺度族. 因此，若$Q(T,\lambda)=2T/\lambda$，则
$$Q(T,\lambda)\sim\operatorname{gamma}(n,\lambda(2/\lambda))=\operatorname{gamma}(n,2),$$
不依赖$\lambda$. 量$Q(T,\lambda)=2T/\lambda$是服从$\operatorname{gamma}(n,2)$即$\chi^2_{2n}$分布的枢轴量.
\end{example}
有时可以根据pdf的形式判断是否存在枢轴量. 上例中的量$t/\lambda$出现在pdf中，而它确实是枢轴量. 正态pdf中出现量$(x-\mu)/\sigma$，它也为枢轴量. 一般地，假设统计量$T$的pdf $f(t\mid\theta)$可表示为
\begin{equation}\label{eq:9.2.11}\index{SI0758@一般形式}
f(t\mid\theta)=g(Q(t,\theta))\left|\frac\partial{\partial t}Q(t,\theta)\right|,
\end{equation}
其中$g$为某个函数，$Q$为某个单调函数（对每个$\theta$关于$t$单调）. 则可用\autoref{thm:2.1.5}证明（\autoref{exercise:9.10}），$Q(T,\theta)$是枢轴量.

得到枢轴量后，如何用它构造置信集？这一步实际上很简单. 若$Q(\boldsymbol X,\theta)$是枢轴量，则对于给定的$\alpha$，可以找到不依赖$\theta$的数$a,b$，使
$$P_\theta(a\le Q(\boldsymbol X,\theta)\le b)\ge1-\alpha.$$
于是，对每个$\theta_0\in\Theta$，
\begin{equation}\label{eq:9.2.12}
A(\theta_0)=\{\boldsymbol x:a\le Q(\boldsymbol x,\theta_0)\le b\}
\end{equation}
是$H_0:\theta=\theta_0$的水平$\alpha$检验的接受域. 我们将用检验反演法构造置信集，不过这里使用枢轴量来规定接受域的具体形式. 应用\autoref{thm:9.2.2}反演这些检验，得到
\begin{equation}\label{eq:9.2.13}
C(\boldsymbol x)=\{\theta_0:a\le Q(\boldsymbol x,\theta_0)\le b\},
\end{equation}
且$C(\boldsymbol X)$是$\theta$的$1-\alpha$置信集. 如果$\theta$是实值参数，且对每个$\boldsymbol x\in\mathcal X$，$Q(\boldsymbol x,\theta)$都是$\theta$的单调函数，那么$C(\boldsymbol x)$将是区间. 事实上，若$Q(\boldsymbol x,\theta)$关于$\theta$递增，则$C(\boldsymbol x)$具有$L(\boldsymbol x,a)\le\theta\le U(\boldsymbol x,b)$的形式. 若$Q(\boldsymbol x,\theta)$关于$\theta$递减（这是典型情形），则$C(\boldsymbol x)$具有$L(\boldsymbol x,b)\le\theta\le U(\boldsymbol x,a)$的形式.
\sourcepageanchor{378}
\begin{example}[续\autoref{ex:9.2.8}]\label{ex:9.2.9}\index{SI0152@枢轴}
在\autoref{ex:9.2.3}中，我们反演$H_0:\lambda=\lambda_0$对$H_1:\lambda\ne\lambda_0$的水平$\alpha$ LRT，得到$\operatorname{exponential}(\lambda)$密度的均值$\lambda$的置信区间. 现在还看到，对于样本$X_1,\ldots,X_n$，可定义$T=\sum X_i$和$Q(T,\lambda)=2T/\lambda\sim\chi^2_{2n}$.

若选择常数$a,b$使$P(a\le\chi^2_{2n}\le b)=1-\alpha$，则
$$P_\lambda\left(a\le\frac{2T}\lambda\le b\right)=P_\lambda(a\le Q(T,\lambda)\le b)=P(a\le\chi^2_{2n}\le b)=1-\alpha.$$
反演集合$A(\lambda)=\{t:a\le2t/\lambda\le b\}$得到$C(t)=\{\lambda:2t/b\le\lambda\le2t/a\}$，这是一个$1-\alpha$置信区间.（注意，如上所述，下端点依赖$b$，上端点依赖$a$，因为$Q(t,\lambda)=2t/\lambda$关于$\lambda$递减.）例如，若$n=10$，查卡方分布表可知，一个$95\%$置信区间为$\{\lambda:2T/34.17\le\lambda\le2T/9.59\}$.
\end{example}
对于位置问题，即使方差未知，枢轴区间的构造和计算也很容易. 事实上，我们已经用过这些思想，只是尚未赋予它们正式名称.
\begin{example}[正态枢轴区间]\label{ex:9.2.10}\index{SI0164@枢轴}
由\autoref{thm:5.3.1}可知，如果$X_1,\ldots,X_n$独立同分布于$\mathrm n(\mu,\sigma^2)$，则$(\bar X-\mu)/(\sigma/\sqrt n)$是枢轴量. 若$\sigma^2$已知，可用这个枢轴量计算$\mu$的置信区间. 对任意常数$a$，
$$P\left(-a\le\frac{\bar X-\mu}{\sigma/\sqrt n}\le a\right)=P(-a\le Z\le a)\qquad\text{（$Z$服从标准正态分布）},$$
而如今已熟悉的代数变换给出置信区间
$$\left\{\mu:\bar x-a\frac\sigma{\sqrt n}\le\mu\le\bar x+a\frac\sigma{\sqrt n}\right\}.$$
若$\sigma^2$未知，可用位置尺度枢轴量$(\bar X-\mu)/(S/\sqrt n)$. 由于$(\bar X-\mu)/(S/\sqrt n)$服从Student $t$分布，
$$P\left(-a\le\frac{\bar X-\mu}{S/\sqrt n}\le a\right)=P(-a\le T_{n-1}\le a).$$
所以对任意给定$\alpha$，取$a=t_{n-1,\alpha/2}$，便得到$1-\alpha$置信区间
\begin{equation}\label{eq:9.2.14}
\left\{\mu:\bar x-t_{n-1,\alpha/2}\frac s{\sqrt n}\le\mu\le\bar x+t_{n-1,\alpha/2}\frac s{\sqrt n}\right\},
\end{equation}
即基于Student $t$分布的$\mu$的经典$1-\alpha$置信区间.

继续考虑这一情形，假设还希望得到$\sigma$的区间估计. 由于$(n-1)S^2/\sigma^2\sim\chi^2_{n-1}$，$(n-1)S^2/\sigma^2$也是枢轴量. 因而，若选取$a,b$满足
$$P\left(a\le\frac{(n-1)S^2}{\sigma^2}\le b\right)=P(a\le\chi^2_{n-1}\le b)=1-\alpha,$$
则可反演这个集合，得到$1-\alpha$置信区间
$$\left\{\sigma^2:\frac{(n-1)s^2}b\le\sigma^2\le\frac{(n-1)s^2}a\right\},$$
或等价地，
$$\left\{\sigma:\sqrt{\frac{(n-1)s^2}b}\le\sigma\le\sqrt{\frac{(n-1)s^2}a}\right\}.$$
\sourcepageanchor{379}
一种能产生所需区间的选择是$a=\chi^2_{n-1,1-\alpha/2}$和$b=\chi^2_{n-1,\alpha/2}$. 这种选择平分了概率，在分布两侧尾部各放置$\alpha/2$. 不过，$\chi^2_{n-1}$是偏斜分布，对于偏斜分布，平分概率是否最优并不显然.（对于对称分布，平分概率是否最优也不显然，只是直觉让后一种情形显得更合理.）事实上，对于卡方分布，平分概率并非最优，\autoref{sec:9.3}将说明这一点.（另见\autoref{exercise:9.52}.）

关于这个问题，最后再说明一点. 我们现在分别构造了$\mu$和$\sigma$的置信区间. 对$\mu$和$\sigma$的\emph{同时}置信集感兴趣完全合情合理. Bonferroni不等式提供了实现这一目标的简单（且相当不错的）方法.\index{SI0560@Bonferroni}\semanticindexanchor{560:379}（参见\autoref{exercise:9.14}.）
\end{example}
\subsubsection{对累积分布函数作枢轴变换}\label{sec:9.2.3}
上一小节看到，枢轴量$Q$导出\eqref{eq:9.2.13}形式的置信集，即
$$C(\boldsymbol x)=\{\theta_0:a\le Q(\boldsymbol x,\theta_0)\le b\}.$$
如果对于每个$\theta$，函数$Q(\boldsymbol x,\theta)$都是$\theta$的单调函数，那么就能保证置信集$C(\boldsymbol x)$是一个区间. 目前所见的枢轴量主要通过位置与尺度变换构造，所得的$Q$函数单调，因而产生置信区间.

本小节研究另一个枢轴量，它完全一般，且在一些轻微假设下可保证得到区间.

如果感到困惑，或遇到陌生情形，我们建议在可能时通过反演LRT来构造置信集. 尽管无法保证这样的集合最优，但它绝不会很差. 不过，在某些情形下，这种策略无论在解析上还是计算上都太困难；接受域的反演有时是件相当费力的事. 如果本小节的方法适用，则实施起来相当直接，而且通常会产生合理的集合.

若不对所用接受域的具体类型施加额外条件，检验反演法可能带来某些麻烦. 为说明这一点，考虑下例，它展示了早期构造二项成功概率置信集的一种方法.
\begin{example}[最短长度二项置信集]\label{ex:9.2.11}\index{SI0145@构造}\index{SI0173@二项}
\hyperref[bib:source305]{Sterne（1954）}提出了下面构造二项置信集的方法，它产生最短长度的集合. 给定$\alpha$，对于每个$p$值，找出由概率最大的$x$值组成的大小为$\alpha$的接受域. 即对于每个$p$，将$x=0,\ldots,n$按概率从大到小排序，依次将值放入接受域$A(p)$，直到其概率为$1-\alpha$. 然后利用\eqref{eq:9.2.1}反演这些接受域，得到$1-\alpha$置信集；Sterne声称它具有长度最优性质.

为看清这个看似合理的构造中出人意料的问题，考虑一个小例子. 设$X\sim\operatorname{binomial}(3,p)$，使用置信系数$1-\alpha=.442$. \autoref{tab:9.2.2}给出了Sterne构造得到的接受域，以及反演这族检验得到的置信集.

令人惊讶的是，置信集并不是置信区间. 这种看似合理的构造导致了不合理的程序. 应当归咎于pmf，因为它并未按我们的预期变化.（参见\autoref{exercise:9.18}.）
\end{example}
我们基于实值统计量$T$来构造参数$\theta$的置信区间，$T$的cdf为$F_T(t\mid\theta)$.（实际中通常取$T$为$\theta$的充分统计量，但下面的理论并不需要这一点.）首先假设$T$是连续随机变量.
\sourcepageanchor{380}
\par
\begin{table}[htbp]\centering
\caption{Sterne构造的接受域与置信集，其中$X\sim\operatorname{binomial}(3,p)$，$1-\alpha=.442$}\label{tab:9.2.2}
\begin{tabular}{cc|cl}\toprule
$p$&接受域$=A(p)$&$x$&置信集$=C(x)$\\\midrule
$[.000,.238]$&$\{0\}$&&\\
$(.238,.305)$&$\{0,1\}$&$0$&$[.000,.305)\cup(.362,.366)$\\
$[.305,.362]$&$\{1\}$&&\\
$(.362,.366)$&$\{0,1\}$&$1$&$(.238,.634]$\\
$[.366,.634]$&$\{1,2\}$&&\\
$(.634,.638)$&$\{2,3\}$&$2$&$[.366,.762)$\\
$[.638,.695]$&$\{2\}$&&\\
$(.695,.762)$&$\{2,3\}$&$3$&$(.634,.638)\cup(.695,1.00]$\\
$[.762,1.00]$&$\{3\}$&&\\\bottomrule
\end{tabular}
\end{table}
\par
当$T$离散时情形类似，但需考虑一些额外的技术细节，因此我们将离散情形另列为一个定理.

首先，回顾概率积分变换\autoref{thm:2.1.10}：随机变量$F_T(T\mid\theta)$服从$\operatorname{uniform}(0,1)$分布，因而是枢轴量. 所以，若$\alpha_1+\alpha_2=\alpha$，则假设$H_0:\theta=\theta_0$的一个水平$\alpha$的接受域为（\autoref{exercise:9.11}）
$$\{t:\alpha_1\le F_T(t\mid\theta_0)\le1-\alpha_2\},$$
对应的置信集为
$$\{\theta:\alpha_1\le F_T(t\mid\theta)\le1-\alpha_2\}.$$
为了保证置信集是区间，需要$F_T(t\mid\theta)$关于$\theta$单调. 但我们在随机递增和随机递减的定义中已经见过这一点.\index{SI0989@递增}\semanticindexanchor{989:380}（参见\autoref{sec:8}的杂录一节与\autoref{exercise:8.26}，或\autoref{exercise:3.41}至\autoref{exercise:3.43}.）若对$T$的样本空间$\mathcal T$中的每个$t$，$F(t\mid\theta)$是$\theta$的递减（递增）函数，则称cdf族$F(t\mid\theta)$关于$\theta$随机递增（随机递减）. 以下仅需$F$单调这一事实，无论递增还是递减. 随机递增和随机递减这些统计意味更浓的概念，仅作为解释工具.
\begin{theorem}[对连续cdf作枢轴变换]\label{thm:9.2.12}\index{SI0848@置信区间}
设统计量$T$具有连续cdf $F_T(t\mid\theta)$. 固定$\alpha_1+\alpha_2=\alpha$，其中$0<\alpha<1$. 假设对于每个$t\in\mathcal T$，函数$\theta_L(t)$和$\theta_U(t)$可以按如下方式定义.
\begin{enumerate}[label=\roman*.]
\item\label{item:9.2.12.i}若对每个$t$，$F_T(t\mid\theta)$是$\theta$的递减函数，则以
$$F_T(t\mid\theta_U(t))=\alpha_1,\qquad F_T(t\mid\theta_L(t))=1-\alpha_2$$
定义$\theta_L(t)$和$\theta_U(t)$.
\item\label{item:9.2.12.ii}若对每个$t$，$F_T(t\mid\theta)$是$\theta$的递增函数，则以
$$F_T(t\mid\theta_U(t))=1-\alpha_2,\qquad F_T(t\mid\theta_L(t))=\alpha_1$$
定义$\theta_L(t)$和$\theta_U(t)$.
\end{enumerate}
则随机区间$[\theta_L(T),\theta_U(T)]$是$\theta$的$1-\alpha$置信区间.
\end{theorem}
\begin{proof}
只证明\hyperref[item:9.2.12.i]{(a)部分}. \hyperref[item:9.2.12.ii]{(b)部分}的证明类似，留作\autoref{exercise:9.19}.

假设已经构造了$1-\alpha$接受域
$$\{t:\alpha_1\le F_T(t\mid\theta_0)\le1-\alpha_2\}.$$
由于对每个$t$，$F_T(t\mid\theta)$是$\theta$的递减函数，且$1-\alpha_2>\alpha_1$，所以$\theta_L(t)<\theta_U(t)$，而且$\theta_L(t)$与$\theta_U(t)$的值唯一. 此外，
$$\begin{aligned}
F_T(t\mid\theta)<\alpha_1&\iff\theta>\theta_U(t),\\
F_T(t\mid\theta)>1-\alpha_2&\iff\theta<\theta_L(t),
\end{aligned}$$
故$\{\theta:\alpha_1\le F_T(t\mid\theta)\le1-\alpha_2\}=\{\theta:\theta_L(T)\le\theta\le\theta_U(T)\}$.
\end{proof}
\sourcepageanchor{381}
注意，没有额外信息时，通常取$\alpha_1=\alpha_2=\alpha/2$. 尽管这不总是最优（参见\autoref{thm:9.3.2}），但在多数情形下肯定是合理的策略. 如果希望得到单侧区间，则只须令$\alpha_1$或$\alpha_2$为$0$，很容易实现.

随机递增情形的方程
\begin{equation}\label{eq:9.2.15}
F_T(t\mid\theta_U(t))=\alpha_1,\qquad F_T(t\mid\theta_L(t))=1-\alpha_2
\end{equation}
也可用统计量$T$的pdf表示. 函数$\theta_U(t),\theta_L(t)$可定义为满足
$$\int_{-\infty}^t f_T(u\mid\theta_U(t))\,du=\alpha_1\quad\text{且}\quad\int_t^\infty f_T(u\mid\theta_L(t))\,du=\alpha_2.$$
随机递减情形也有类似的一组方程.
\begin{example}[位置指数分布区间]\label{ex:9.2.13}\index{SI0151@指数}
这个方法可用于获得位置指数pdf的置信区间.（\autoref{exercise:9.25}将这里的答案与似然法和枢轴法得到的答案比较. 另见\autoref{exercise:9.41}.）

若$X_1,\ldots,X_n$独立同分布，pdf为$f(x\mid\mu)=e^{-(x-\mu)}I_{[\mu,\infty)}(x)$，则$Y=\min\{X_1,\ldots,X_n\}$对$\mu$充分，其pdf为
$$f_Y(y\mid\mu)=ne^{-n(y-\mu)}I_{[\mu,\infty)}(y).$$
固定$\alpha$，定义$\mu_L(y),\mu_U(y)$使其满足
$$\int_{\mu_U(y)}^y ne^{-n(u-\mu_U(y))}\,du=\frac\alpha2,\qquad\int_y^\infty ne^{-n(u-\mu_L(y))}\,du=\frac\alpha2.$$
计算积分得到方程
$$1-e^{-n(y-\mu_U(y))}=\frac\alpha2,\qquad e^{-n(y-\mu_L(y))}=\frac\alpha2,$$
其解为
$$\mu_U(y)=y+\frac1n\log\left(1-\frac\alpha2\right),\qquad\mu_L(y)=y+\frac1n\log\left(\frac\alpha2\right).$$
因此，随机区间
$$C(Y)=\left\{\mu:Y+\frac1n\log\left(\frac\alpha2\right)\le\mu\le Y+\frac1n\log\left(1-\frac\alpha2\right)\right\}$$
是$\mu$的$1-\alpha$置信区间.
\end{example}
使用这个方法时有两点需要注意.\index{SI0145@构造}\semanticindexanchor{145:381} 第一，只须对实际观测到的统计量值求解方程\eqref{eq:9.2.15}. 若观测到$T=t_0$，则$\theta$的置信区间实现值为$[\theta_L(t_0),\theta_U(t_0)]$. 因而只须解以下两个方程，求得$\theta_L(t_0),\theta_U(t_0)$：
$$\int_{-\infty}^{t_0}f_T(u\mid\theta_U(t_0))\,du=\alpha_1\quad\text{且}\quad\int_{t_0}^\infty f_T(u\mid\theta_L(t_0))\,du=\alpha_2.$$
第二，即使这些方程不能解析求解，实际也只需进行数值求解，因为证明所得区间是$1-\alpha$置信区间时并不需要解析解.

现在考虑离散情形.
\begin{theorem}[对离散cdf作枢轴变换]\label{thm:9.2.14}\index{SI0849@离散}
设$T$为离散统计量，cdf为$F_T(t\mid\theta)=P(T\le t\mid\theta)$. 固定$\alpha_1+\alpha_2=\alpha$，其中$0<\alpha<1$. 假设对每个$t\in\mathcal T$，$\theta_L(t),\theta_U(t)$可定义如下.
\sourcepageanchor{382}
\begin{enumerate}[label=\roman*.]
\item\label{item:9.2.14.i}若对每个$t$，$F_T(t\mid\theta)$是$\theta$的递减函数，则以
$$P(T\le t\mid\theta_U(t))=\alpha_1,\qquad P(T\ge t\mid\theta_L(t))=\alpha_2$$
定义$\theta_L(t),\theta_U(t)$.
\item\label{item:9.2.14.ii}若对每个$t$，$F_T(t\mid\theta)$是$\theta$的递增函数，则以
$$P(T\ge t\mid\theta_U(t))=\alpha_1,\qquad P(T\le t\mid\theta_L(t))=\alpha_2$$
定义$\theta_L(t),\theta_U(t)$.
\end{enumerate}
则随机区间$[\theta_L(T),\theta_U(T)]$是$\theta$的$1-\alpha$置信区间.
\end{theorem}
\begin{proof}
仅概述\hyperref[item:9.2.14.i]{(a)部分}的证明. 细节以及\hyperref[item:9.2.14.ii]{(b)部分}的证明留作\autoref{exercise:9.20}.

首先回顾\autoref{exercise:2.10}，其中已证明$F_T(T\mid\theta)$随机大于均匀随机变量，即$P_\theta(F_T(T\mid\theta)\le x)\le x$. 此外，$\overline F_T(T\mid\theta)=P(T\ge t\mid\theta)$也具有这个性质，这意味着集合
$$\{\theta:F_T(T\mid\theta)\le\alpha_1\text{ 且 }\overline F_T(T\mid\theta)\le\alpha_2\}$$
是一个$1-\alpha$置信集.

对于每个$t$，$F_T(t\mid\theta)$是$\theta$的递减函数这一事实，意味着对于每个$t$，$\overline F(t\mid\theta)$是$\theta$的非递减函数. 因此，
$$\begin{aligned}
\theta>\theta_U(t)&\implies F_T(t\mid\theta)<\frac\alpha2,\\
\theta<\theta_L(t)&\implies\overline F_T(t\mid\theta)<\frac\alpha2,
\end{aligned}$$
故$\{\theta:F_T(T\mid\theta)\le\alpha_1\text{ 且 }\overline F_T(T\mid\theta)\le\alpha_2\}=\{\theta:\theta_L(T)\le\theta\le\theta_U(T)\}$.\footnote{译注：本证明原文的两个集合均使用“$\le\alpha_1$且$\le\alpha_2$”，与排除两侧尾部并构造中央置信集的论证方向不符；应考虑相反方向的不等式，并审慎处理端点. 两个蕴含式的右端原文均为$\alpha/2$，但定理未假设$\alpha_1=\alpha_2$，一般应分别为$\alpha_1,\alpha_2$. 原文还将$\overline F_T(T\mid\theta)$的定义右侧写为含$t$的式子. 此处均保留原式.}
\end{proof}
本小节最后给出一个例子，说明\autoref{thm:9.2.14}中的构造. 注意，也可通过反演LRT构造另一个区间（\autoref{exercise:9.23}）.
\begin{example}[泊松区间估计量]\label{ex:9.2.15}\index{SI0168@泊松}
设$X_1,\ldots,X_n$为来自参数为$\lambda$的泊松总体的随机样本，并定义$Y=\sum X_i$. $Y$对$\lambda$充分，且$Y\sim\operatorname{Poisson}(n\lambda)$. 应用上述方法，取$\alpha_1=\alpha_2=\alpha/2$；如果观测到$Y=y_0$，就需要从方程
\begin{equation}\label{eq:9.2.16}
\sum_{k=0}^{y_0}e^{-n\lambda}\frac{(n\lambda)^k}{k!}=\frac\alpha2\quad\text{且}\quad\sum_{k=y_0}^\infty e^{-n\lambda}\frac{(n\lambda)^k}{k!}=\frac\alpha2
\end{equation}
中解出$\lambda$. 回顾\autoref{ex:3.3.1}中联系泊松族与伽马族的恒等式. 将该恒等式用于\eqref{eq:9.2.16}中的和式，可以写为（牢记$y_0$是$Y$的观测值）
$$\frac\alpha2=\sum_{k=0}^{y_0}e^{-n\lambda}\frac{(n\lambda)^k}{k!}=P(Y\le y_0\mid\lambda)=P(\chi^2_{2(y_0+1)}>2n\lambda),$$
其中$\chi^2_{2(y_0+1)}$是自由度为$2(y_0+1)$的卡方随机变量. 因而上面方程的解为
$$\lambda=\frac1{2n}\chi^2_{2(y_0+1),\alpha/2}.$$
类似地，将恒等式应用于\eqref{eq:9.2.16}的另一个方程，得到
$$\frac\alpha2=\sum_{k=y_0}^\infty e^{-n\lambda}\frac{(n\lambda)^k}{k!}=P(Y\ge y_0\mid\lambda)=P(\chi^2_{2y_0}<2n\lambda).$$
作一些代数变换，得到$\lambda$的$1-\alpha$置信区间
\begin{equation}\label{eq:9.2.17}
\left\{\lambda:\frac1{2n}\chi^2_{2y_0,1-\alpha/2}\le\lambda\le\frac1{2n}\chi^2_{2(y_0+1),\alpha/2}\right\}.
\end{equation}
（在$y_0=0$时，定义$\chi^2_{0,1-\alpha/2}=0$.）
\sourcepageanchor{383}
这些区间最初由\hyperref[bib:source133]{Garwood（1936）}导出. \autoref{fig:9.2.5}给出了覆盖概率的图像. 注意，图形相当参差不齐. 跳跃发生在各置信区间的端点处，此时构成覆盖概率的和式中有项被加入或减去.（参见\autoref{exercise:9.24}.）

举一个数值例子：考虑$n=10$，观测到$y_0=\sum x_i=6$. $\lambda$的一个$90\%$置信区间为
$$\frac1{20}\chi^2_{12,.95}\le\lambda\le\frac1{20}\chi^2_{14,.05},$$
查卡方分布表得到
$$.262\le\lambda\le1.184.$$
涉及负二项分布和二项分布的类似推导见习题.
\end{example}
\subsubsection{贝叶斯区间}\label{sec:9.2.4}\index{SI0184@解释}
到目前为止，在描述置信区间与参数的关系时，我们一直谨慎地说区间\emph{覆盖}参数，而不说参数\emph{位于}区间内. 这是有意为之，目的是强调随机量是区间，不是参数. 因此，我们尽量使动作动词作用于区间，而不是参数.

在\autoref{ex:9.2.15}中看到，若$y_0=\sum_{i=1}^{10}x_i=6$，则$\lambda$的一个$90\%$置信区间为$.262\le\lambda\le1.184$. 人们很容易说（许多实验工作者确实这样说）：“$\lambda$位于区间$[.262,1.184]$内的概率为$90\%$.” 但在经典统计中，这样的说法无效，因为参数被假定为固定的. 形式上，区间$[.262,1.184]$是随机区间$[\frac1{2n}\chi^2_{2Y,.95},\frac1{2n}\chi^2_{2(Y+1),.05}]$的一个可能实现值；由于参数$\lambda$不变，$\lambda$位于实现区间$[.262,1.184]$内的概率只能为$0$或$1$. 说实现区间$[.262,1.184]$有$90\%$的覆盖机会时，仅意味着我们知道随机区间的样本点中有$90\%$覆盖真实参数.

相比之下，贝叶斯框架允许我们说$\lambda$以某个不同于$0$或$1$的概率位于$[.262,1.184]$内. 这是因为贝叶斯模型中的$\lambda$是具有概率分布的随机变量. 贝叶斯关于覆盖的全部断言，都是针对参数的后验分布作出的.

由于贝叶斯集合与经典集合所作的概率评价很不相同，为明确区分两者，贝叶斯集合估计称为可信集（credible set），而不是置信集.

因此，若$\pi(\theta\mid\boldsymbol x)$是在给定$\boldsymbol X=\boldsymbol x$下$\theta$的后验分布，则任意集合$A\subset\Theta$的可信概率为
\begin{equation}\label{eq:9.2.18}\index{SI0235@可信集（另见置信集；置信区间；置信程序）}
P(\theta\in A\mid\boldsymbol x)=\int_A\pi(\theta\mid\boldsymbol x)\,d\theta,
\end{equation}
且$A$是$\theta$的可信集. 若$\pi(\theta\mid\boldsymbol x)$是pmf，则将上式中的积分替换为求和.

注意，贝叶斯可信集的解释和构造均比经典置信集更为直接. 但要记住，任何便利都不是无代价的. 构造与解释的简便伴随着额外假设. 贝叶斯模型所需的输入比经典模型更多.
\begin{example}[泊松可信集]\label{ex:9.2.16}\index{SI0232@泊松}
现在为\autoref{ex:9.2.15}的问题构造可信集. 设$X_1,\ldots,X_n$独立同分布于$\operatorname{Poisson}(\lambda)$，并假定$\lambda$具有伽马先验pdf，$\lambda\sim\operatorname{gamma}(a,b)$. $\lambda$的后验pdf（\autoref{exercise:7.24}）为
\begin{equation}\label{eq:9.2.19}
\pi(\lambda\mid\textstyle\sum X=\sum x)=\operatorname{gamma}(a+\textstyle\sum x,[n+(1/b)]^{-1}).
\end{equation}
\sourcepageanchor{384}
可用许多不同方式为$\lambda$构造可信集，因为任何满足\eqref{eq:9.2.18}的集合$A$都可以. 一种简单方法是在上、下端点间平分$\alpha$. 由\eqref{eq:9.2.19}可得$2(nb+1)\lambda/b\sim\chi^2_{2(a+\sum x_i)}$（假设$a$是整数），因而一个$1-\alpha$可信区间为
\begin{equation}\label{eq:9.2.20}
\left\{\lambda:\frac b{2(nb+1)}\chi^2_{2(\sum x+a),1-\alpha/2}\le\lambda\le\frac b{2(nb+1)}\chi^2_{2(\sum x+a),\alpha/2}\right\}.
\end{equation}
若取$a=b=1$，则给定$\sum X=\sum x$时$\lambda$的后验分布可表示为$2(n+1)\lambda\sim\chi^2_{2(\sum x+1)}$. 与\autoref{ex:9.2.15}一样，假设$n=10$，$\sum x=6$. 因为$\chi^2_{14,.95}=6.571$，$\chi^2_{14,.05}=23.685$，所以$\lambda$的一个$90\%$可信集为$[.299,1.077]$.

$90\%$可信集的实现值与\autoref{ex:9.2.15}得到的$90\%$置信集$[.262,1.184]$不同. 为更清楚地看出差异，请看\autoref{fig:9.2.3}，它给出一系列$x$值对应的$90\%$可信区间和$90\%$置信区间. 注意，可信集中的区间稍短，上端点也更接近零. 这反映了先验的作用：先验将区间拉向零.
\end{example}
\par
\begin{figure}[htbp]\centering
\includegraphics[width=.68\textwidth]{assets/chapter9/fig-9.2.3.pdf}
\caption{\autoref{ex:9.2.16}中的$90\%$可信区间（虚线）和$90\%$置信区间（实线）.}\label{fig:9.2.3}
\end{figure}
\par
务必不要混淆可信概率（贝叶斯后验概率）与覆盖概率（经典概率）.\index{SI0225@频率学派与贝叶斯解释}\semanticindexanchor{225:384} 它们是很不同的量，有不同的含义与解释. 可信概率来自后验分布，后验分布又从先验分布获得其概率. 因而可信概率反映实验者的主观信念：这些信念用先验分布表达，再结合数据更新为后验分布. 贝叶斯关于$90\%$覆盖的断言，意味着实验者结合先验知识与数据后，对覆盖有$90\%$的把握.

另一方面，覆盖概率反映抽样程序的不确定性，其概率来自重复实验试验这一客观机制. 经典的$90\%$覆盖断言，意味着在一长串相同试验中，实现的置信集将有$90\%$覆盖真实参数.

统计学家有时争论，经典方法与贝叶斯方法中，哪一种是更好的统计方法. 我们不打算争论，甚至不想为某一方辩护. 事实上，我们认为不存在一种唯一最佳的统计方法；一些问题最好用经典统计解决，另一些问题最好用贝叶斯统计解决. 重要的是认识到，所得解答可能很不相同. 在经典评价下，贝叶斯解往往不合理，反之亦然.
\sourcepageanchor{385}
\begin{example}[泊松可信概率与覆盖概率]\label{ex:9.2.17}
\autoref{ex:9.2.16}中的$90\%$置信集和可信集保持了各自的概率保证，但用另一种准则评价时，表现如何？首先来看置信集\eqref{eq:9.2.17}的可信概率：
\begin{equation}\label{eq:9.2.21}
P\left\{\frac1{2n}\chi^2_{2\sum x,1-\alpha/2}\le\lambda\le\frac1{2n}\chi^2_{2(\sum x+1),\alpha/2}\right\},
\end{equation}
其中$\lambda$服从分布\eqref{eq:9.2.19}. \autoref{fig:9.2.4}给出集合\eqref{eq:9.2.20}恒等于$1-\alpha$的可信概率，同时给出置信集的可信概率\eqref{eq:9.2.21}.
\par
\begin{figure}[htbp]\centering
\includegraphics[width=.78\textwidth]{assets/chapter9/fig-9.2.4.pdf}
\caption{\autoref{ex:9.2.16}中的$90\%$可信区间（虚线）和$90\%$置信区间（实线）的可信概率.}\label{fig:9.2.4}
\end{figure}
\par
后一个概率似乎持续下降，我们希望知道它是否对所有$\sum x_i$值都保持大于$0$（对每个固定$n$而言）. 为此，考察$\sum x_i\to\infty$时的概率. 细节留作\autoref{exercise:9.30}，但结果是：当$\sum x_i\to\infty$时，除非$b=1/n$，否则概率\eqref{eq:9.2.21}$\to0$. 因此，这个置信区间无法保持非零的可信概率.\footnote{译注：原文中的例外条件“除非$b=1/n$”可能有误. 对固定的有限$b>0$，后验均值约为$\sum x_i/(n+1/b)$，而置信区间端点约为$\sum x_i/n$，二者的差为$\sum x_i$的量级，超过后验标准差的平方根量级；因此$b=1/n$似乎也不能构成例外. 原文保留，供对照\autoref{exercise:9.30}核实.}

把可信集\eqref{eq:9.2.20}作为置信集评价时，它的表现也好不了多少. \autoref{fig:9.2.5}提示，随着$\lambda\to\infty$，可信集的覆盖概率趋于零. 为评价覆盖概率，写成
$$\lambda=\frac\lambda{\chi^2_{2Y}}\chi^2_{2Y},$$
其中$\chi^2_{2Y}$是自由度为$2Y$的卡方随机变量，且$Y\sim\operatorname{Poisson}(n\lambda)$. 于是，当$\lambda\to\infty$时，$\lambda/\chi^2_{2Y}\to1/(2n)$，而\eqref{eq:9.2.20}的覆盖概率成为
\begin{equation}\label{eq:9.2.22}
P\left\{\frac{nb}{nb+1}\chi^2_{2(Y+a),1-\alpha/2}\le\chi^2_{2Y}\le\frac{nb}{nb+1}\chi^2_{2(Y+a),\alpha/2}\right\}.
\end{equation}
该概率在$\lambda\to\infty$时趋于零，见\autoref{exercise:9.31}.
\end{example}
\autoref{ex:9.2.17}展示的行为具有一定典型性. 下面给出一个能够显式完成计算的例子.
\begin{example}[正态可信集的覆盖]\label{ex:9.2.18}\index{SI0231@正态}
设$X_1,\ldots,X_n$独立同分布于$\mathrm n(\theta,\sigma^2)$，$\theta$的先验pdf为$\mathrm n(\mu,\tau^2)$，其中$\mu,\sigma,\tau$均已知. 在\autoref{ex:7.2.16}中看到
$$\pi(\theta\mid\bar x)\sim\mathrm n(\delta^{\mathrm B}(\bar x),\operatorname{Var}(\theta\mid\bar x)),$$
\sourcepageanchor{386}
\par
\begin{figure}[htbp]\centering
\includegraphics[width=.8\textwidth]{assets/chapter9/fig-9.2.5.pdf}
\caption{\autoref{ex:9.2.16}中的$90\%$可信区间（虚线）和$90\%$置信区间（实线）的覆盖概率.}\label{fig:9.2.5}
\end{figure}
\par
其中
$$\delta^{\mathrm B}(\bar x)=\frac{\sigma^2}{\sigma^2+n\tau^2}\mu+\frac{n\tau^2}{\sigma^2+n\tau^2}\bar x,\qquad\operatorname{Var}(\theta\mid\bar x)=\frac{\sigma^2\tau^2}{\sigma^2+n\tau^2}.$$
因此，在后验分布下，
$$\frac{\theta-\delta^{\mathrm B}(\bar x)}{\sqrt{\operatorname{Var}(\theta\mid\bar x)}}\sim\mathrm n(0,1),$$
而$\theta$的一个$1-\alpha$可信集为
\begin{equation}\label{eq:9.2.23}
\delta^{\mathrm B}(\bar x)-z_{\alpha/2}\sqrt{\operatorname{Var}(\theta\mid\bar x)}\le\theta\le\delta^{\mathrm B}(\bar x)+z_{\alpha/2}\sqrt{\operatorname{Var}(\theta\mid\bar x)}.
\end{equation}
现在计算贝叶斯区域\index{SI0234@正态}\semanticindexanchor{234:386}\eqref{eq:9.2.23}的覆盖概率. 经典模型中$\bar X$是随机变量，$\theta$固定，且$\bar X\sim\mathrm n(\theta,\sigma^2/n)$. 为使记号简便，定义$\gamma=\sigma^2/(n\tau^2)$. 利用$\delta^{\mathrm B}(\bar X)$和$\operatorname{Var}(\theta\mid\bar X)$的定义并略作代数变换，\eqref{eq:9.2.23}的覆盖概率为
$$\begin{aligned}
&P_\theta\left(|\theta-\delta^{\mathrm B}(\bar X)|\le z_{\alpha/2}\sqrt{\operatorname{Var}(\theta\mid\bar X)}\right)\\
&=P_\theta\left(\left|\theta-\left(\frac\gamma{1+\gamma}\mu+\frac1{1+\gamma}\bar X\right)\right|\le z_{\alpha/2}\sqrt{\frac{\sigma^2}{n(1+\gamma)}}\right)\\
&=P_\theta\left(-\sqrt{1+\gamma}\,z_{\alpha/2}+\frac{\gamma(\theta-\mu)}{\sigma/\sqrt n}\le Z\le\sqrt{1+\gamma}\,z_{\alpha/2}+\frac{\gamma(\theta-\mu)}{\sigma/\sqrt n}\right),
\end{aligned}$$
最后一个等号利用了$\sqrt n(\bar X-\theta)/\sigma=Z\sim\mathrm n(0,1)$这一事实.

尽管从$1-\alpha$可信集出发，我们得到的并非$1-\alpha$置信集，考虑如下参数配置即可看出. 固定$\theta\ne\mu$，令$\tau=\sigma/\sqrt n$，于是$\gamma=1$. 再令$\sigma/\sqrt n$非常小（$\to0$）. 容易看出，上述概率趋于零，因为若$\theta>\mu$，下界趋于无穷大；若$\theta<\mu$，上界趋于$-\infty$. 不过，若$\theta=\mu$，覆盖概率则有一个远离零的下界.
\sourcepageanchor{387}
另一方面，$\theta$的通常的$1-\alpha$置信集为$\{\theta:|\theta-\bar x|\le z_{\alpha/2}\sigma/\sqrt n\}$. 这个集合的可信概率（现在$\theta\sim\pi(\theta\mid\bar x)$）为
$$\begin{aligned}
&P_{\bar x}\left(|\theta-\bar x|\le z_{\alpha/2}\frac\sigma{\sqrt n}\right)\\
&=P_{\bar x}\left(\left|[\theta-\delta^{\mathrm B}(\bar x)]+[\delta^{\mathrm B}(\bar x)-\bar x]\right|\le z_{\alpha/2}\frac\sigma{\sqrt n}\right)\\
&=P_{\bar x}\left(-\sqrt{1+\gamma}\,z_{\alpha/2}+\frac{\gamma(\bar x-\mu)}{\sqrt{1+\gamma}\,\sigma/\sqrt n}\le Z\le\sqrt{1+\gamma}\,z_{\alpha/2}+\frac{\gamma(\bar x-\mu)}{\sqrt{1+\gamma}\,\sigma/\sqrt n}\right),
\end{aligned}$$
最后一个等号使用了$(\theta-\delta^{\mathrm B}(\bar x))/\sqrt{\operatorname{Var}(\theta\mid\bar x)}=Z\sim\mathrm n(0,1)$. 同样很容易证明，该概率没有远离零的下界，这说明置信集一般也不是可信集. 细节见\autoref{exercise:9.32}.
\end{example}



\subsection{评价区间估计量的方法}\label{sec:9.3}
我们已经见到许多导出置信集的方法，而且对于同一问题，实际上可以导出不同的置信集. 此时当然希望选择最好的一个. 因此，现在考察评价集合估计量的一些方法和准则.

集合估计中，有两个量相互竞争：大小和覆盖概率. 自然希望集合大小较小、覆盖概率较大，但这样的集合通常难以构造.（显然，增大集合就能获得较大的覆盖概率. 区间$(-\infty,\infty)$的覆盖概率为$1$！）在根据大小和覆盖概率优化集合之前，必须决定如何度量这些量.

除特殊情形外，置信集的覆盖概率是参数的函数，因此要考虑的不是一个值，而是无穷多个值. 不过，大多数时候，我们将通过置信系数，即覆盖概率的下确界，来度量覆盖概率的表现.\index{SI0127@置信系数}\semanticindexanchor{127:387} 这是概括覆盖概率信息的一种方式，但并非唯一方式.（例如，也可以计算平均覆盖概率.）

说到置信集的大小，若集合为区间，通常指其长度.\index{SI0188@大小}\semanticindexanchor{188:387} 若集合不是区间，或者处理的是多维集合，那么长度通常变为体积.（也有一些情形，采用长度之外的大小度量更自然，特别是在考虑等变性时. \hyperref[bib:source276]{Schervish 1995}，第6章，以及\hyperref[bib:source18]{Berger 1985}，第6章，讨论了这个主题.）
\subsubsection{大小与覆盖概率}\label{sec:9.3.1}
现在考虑一个看似简单的约束极小化问题：对于给定的覆盖概率，寻找长度最短的置信区间. 先看一个例子.
\begin{example}[长度的优化]\label{ex:9.3.1}
设$X_1,\ldots,X_n$独立同分布于$\mathrm n(\mu,\sigma^2)$，其中$\sigma$已知. 利用\autoref{sec:9.2.2}的方法和
$$Z=\frac{\bar X-\mu}{\sigma/\sqrt n}$$
是服从标准正态分布的枢轴量这一事实，任意满足
$$P(a\le Z\le b)=1-\alpha$$
的$a,b$
\sourcepageanchor{388}
均给出$1-\alpha$置信区间
$$\left\{\mu:\bar x-b\frac\sigma{\sqrt n}\le\mu\le\bar x-a\frac\sigma{\sqrt n}\right\}.$$
哪一对$a,b$最好？更正式地问，怎样选择$a,b$，可以在保持$1-\alpha$覆盖率的同时使置信区间长度最小？注意，置信区间的长度为$(b-a)\sigma/\sqrt n$，但由于每个区间长度都含因子$\sigma/\sqrt n$，可以忽略这一因子，而根据$b-a$的值比较长度. 因而，希望找到一对数$a,b$，使$P(a\le Z\le b)=1-\alpha$，并使$b-a$最小.

在\autoref{ex:9.2.1}中，我们取$a=-z_{\alpha/2},b=z_{\alpha/2}$，但没有提及最优性. 若取$1-\alpha=.90$，则下列任一对数都给出$90\%$区间. \autoref{tab:9.3.3}中的数值研究提示，选择$a=-1.65,b=1.65$给出最好的区间，事实也确实如此. 在这种情形下，平分概率$\alpha$是最优策略.
\par
\begin{table}[htbp]\centering
\setcounter{table}{2}
\caption{三个$90\%$正态置信区间}\label{tab:9.3.3}
\begin{tabular}{rrccr}\toprule
$a$&$b$&\multicolumn{2}{c}{概率}&$b-a$\\\midrule
$-1.34$&$2.33$&$P(Z<a)=.09$&$P(Z>b)=.01$&$3.67$\\
$-1.44$&$1.96$&$P(Z<a)=.075$&$P(Z>b)=.025$&$3.40$\\
$-1.65$&$1.65$&$P(Z<a)=.05$&$P(Z>b)=.05$&$3.30$\\\bottomrule
\end{tabular}
\end{table}
\par
\end{example}
平分$\alpha$的策略在上面的情形中最优，但并不总是最优. 上述情形中，平分$\alpha$之所以最优，是因为pdf在$-z_{\alpha/2}$和$z_{\alpha/2}$处高度相同. 下面证明一个说明此事实的定理；它具有一定的一般性，只需要假设pdf单峰. 回顾单峰（unimodal）的定义：若存在$x^*$，使$f(x)$在$x\le x^*$时非递减，在$x\ge x^*$时非递增，则称pdf $f(x)$单峰.（这是相当弱的要求.）
\begin{theorem}\label{thm:9.3.2}\index{SI0132@最优}\index{SI0135@单峰}
设$f(x)$是单峰pdf. 若区间$[a,b]$满足
\begin{enumerate}[label=\roman*.]
\item\label{item:9.3.2.i}$\int_a^b f(x)\,dx=1-\alpha$，
\item\label{item:9.3.2.ii}$f(a)=f(b)>0$，且
\item\label{item:9.3.2.iii}$a\le x^*\le b$，其中$x^*$为$f(x)$的一个众数，
\end{enumerate}
则$[a,b]$是所有满足\hyperref[item:9.3.2.i]{(a)}的区间中最短的.
\end{theorem}
\begin{proof}
设$[a',b']$是任意满足$b'-a'<b-a$的区间. 将证明这蕴含$\int_{a'}^{b'}f(x)\,dx<1-\alpha$. 只证明$a'\le a$的情形，$a<a'$时证明类似. 还需分两种情形：$b'\le a$和$b'>a$.

若$b'\le a$，则$a'\le b'\le a\le x^*$，且
$$\begin{aligned}
\int_{a'}^{b'}f(x)\,dx
&\le f(b')(b'-a')&&\text{（$x\le b'\le x^*\implies f(x)\le f(b')$）}\\
&\le f(a)(b'-a')&&\text{（$b'\le a\le x^*\implies f(b')\le f(a)$）}\\
&<f(a)(b-a)&&\text{（$b'-a'<b-a$且$f(a)>0$）}\\
&\le\int_a^b f(x)\,dx&&\text{（由\hyperref[item:9.3.2.ii]{(ii)}、\hyperref[item:9.3.2.iii]{(iii)}及单峰性，$a\le x\le b$时$f(x)\ge f(a)$）}\\
&=1-\alpha,&&\text{（\hyperref[item:9.3.2.i]{(i)}）}
\end{aligned}$$
第一种情形得证.
\sourcepageanchor{389}
若$b'>a$，则$a'\le a<b'<b$，因为如果$b'\ge b$，则$b'-a'\ge b-a$. 此时可以写为
$$\begin{aligned}
\int_{a'}^{b'}f(x)\,dx&=\int_a^b f(x)\,dx+\left[\int_{a'}^a f(x)\,dx-\int_{b'}^b f(x)\,dx\right]\\
&=(1-\alpha)+\left[\int_{a'}^a f(x)\,dx-\int_{b'}^b f(x)\,dx\right].
\end{aligned}$$
只要证明方括号内的表达式为负，定理即得证. 利用$f$的单峰性、次序$a'\le a<b'<b$及\hyperref[item:9.3.2.ii]{(b)}，得到
$$\int_{a'}^a f(x)\,dx\le f(a)(a-a')\quad\text{且}\quad\int_{b'}^b f(x)\,dx\ge f(b)(b-b').$$
因而
$$\begin{aligned}
\int_{a'}^a f(x)\,dx-\int_{b'}^b f(x)\,dx
&\le f(a)(a-a')-f(b)(b-b')\\
&=f(a)[(a-a')-(b-b')]\qquad(f(a)=f(b))\\
&=f(a)[(b'-a')-(b-a)],
\end{aligned}$$
当$(b'-a')<(b-a)$且$f(a)>0$时为负.
\end{proof}
若愿意对$f$施加更多假设，例如$f$连续，则可以简化\autoref{thm:9.3.2}的证明. 参见\autoref{exercise:9.38}.

回顾\autoref{ex:9.2.3}之后关于似然区域形式的讨论，\index{SI0134@与最高后验密度的关系}\semanticindexanchor{134:389}由\autoref{thm:9.3.2}现在知道这是一种最优构造. \autoref{cor:9.3.10}给出的类似论证将说明，这种构造如何产生最优贝叶斯区域. 此外，我们现在也可以看出，\autoref{ex:9.3.1}中最优的平分$\alpha$方法，对于任何对称单峰pdf都最优（\autoref{exercise:9.39}）. 即使最优性准则与最小长度有所不同，\autoref{thm:9.3.2}也可能适用.
\begin{example}[期望长度的优化]\label{ex:9.3.3}\index{SI0131@正态}\index{SI0150@正态}
对于基于枢轴量$(\bar X-\mu)/(S/\sqrt n)$的正态区间，我们知道，具有形式
$$\bar x-b\frac s{\sqrt n}\le\mu\le\bar x-a\frac s{\sqrt n}$$
的最短$1-\alpha$置信区间取$a=-t_{n-1,\alpha/2},b=t_{n-1,\alpha/2}$. 区间长度是$s$的函数，一般形式为
$$\operatorname{Length}(s)=(b-a)\frac s{\sqrt n}.$$
容易看出，如果考虑\emph{期望长度}准则，希望找到使
$$E_\sigma(\operatorname{Length}(S))=(b-a)\frac{E_\sigma S}{\sqrt n}=(b-a)c(n)\frac\sigma{\sqrt n}$$
最小的$1-\alpha$区间，那么\autoref{thm:9.3.2}仍适用，选择$a=-t_{n-1,\alpha/2},b=t_{n-1,\alpha/2}$再次给出最优区间.（量$c(n)$是只依赖$n$的常数，参见\autoref{exercise:7.50}.）
\end{example}
在某些情形下，特别是在位置问题以外的场合，应用\autoref{thm:9.3.2}必须谨慎. 尤其在尺度情形中，定理可能无法直接应用，但某个变形可能适用.
\sourcepageanchor{390}
\begin{example}[最短枢轴区间]\label{ex:9.3.4}\index{SI0155@伽马}
设$X\sim\operatorname{gamma}(k,\beta)$. 量$Y=X/\beta$是枢轴量，$Y\sim\operatorname{gamma}(k,1)$，因此可以寻找满足
\begin{equation}\label{eq:9.3.1}P(a\le Y\le b)=1-\alpha\end{equation}
的常数$a,b$来获得置信区间. 不过，盲目应用\autoref{thm:9.3.2}并不能得到最短置信区间. 即选择满足\eqref{eq:9.3.1}及$f_Y(a)=f_Y(b)$的$a,b$并非最优. 因为基于\eqref{eq:9.3.1}的$\beta$区间具有形式
$$\left\{\beta:\frac xb\le\beta\le\frac xa\right\},$$
所以区间长度为$(1/a-1/b)x$，即与$(1/a)-(1/b)$成正比，而不是与$b-a$成正比.

虽然\autoref{thm:9.3.2}不能直接应用，但修改论证可以解决这个问题. \autoref{thm:9.3.2}的\hyperref[item:9.3.2.i]{条件(i)}将$b$定义为$a$的函数，记为$b(a)$. 必须解以下约束极小化问题：
$$\begin{aligned}
\text{关于$a$极小化：}\quad&\frac1a-\frac1{b(a)}\\
\text{约束条件：}\quad&\int_a^{b(a)}f_Y(y)\,dy=1-\alpha.
\end{aligned}$$
对第一个式子关于$a$求导，并令导数为零，得到恒等式$db/da=b^2/a^2$. 将其代入第二个式子的导数（该导数必为零），得到$f(b)b^2=f(a)a^2$（\autoref{exercise:9.42}）. 正态分布方差的区间估计中也会出现这样的方程，参见\autoref{ex:9.2.10}和\autoref{exercise:9.52}. 注意，上述方程定义的不是总体意义下最短的区间，而是最短枢轴区间，即基于枢轴量$X/\beta$的最短区间. 涉及Neyman--Pearson引理的推广见\autoref{exercise:9.43}.
\end{example}
\subsubsection{与检验相关的最优性}\label{sec:9.3.2}
由于置信集与假设检验一一对应（\autoref{thm:9.2.2}），检验的最优性与置信集的最优性也存在某种对应关系. 通常，置信集与检验相关的最优性质并不直接涉及集合大小，而是涉及集合覆盖错误值的概率.

覆盖错误值的概率，或称错误覆盖概率（probability of false coverage），间接度量置信集的大小. 直觉上，较小的集合覆盖的值较少，因此覆盖错误值的可能性也较小. 此外，后面还将看到联系大小与错误覆盖概率的方程.

先考虑一般情形，$\boldsymbol X\sim f(\boldsymbol x\mid\theta)$，通过反演接受域$A(\theta)$，构造$\theta$的$1-\alpha$置信集$C(\boldsymbol x)$. $C(\boldsymbol x)$的覆盖概率，即正确覆盖概率，为$\theta$的函数$P_\theta(\theta\in C(\boldsymbol X))$. 错误覆盖概率是$\theta$和$\theta'$的函数，定义为
\begin{equation}\label{eq:9.3.2}\index{SI0858@错误覆盖概率}
\begin{aligned}
P_\theta(\theta'\in C(\boldsymbol X)),\ \theta\ne\theta',&\qquad\text{若 }C(\boldsymbol X)=[L(\boldsymbol X),U(\boldsymbol X)],\\
P_\theta(\theta'\in C(\boldsymbol X)),\ \theta'<\theta,&\qquad\text{若 }C(\boldsymbol X)=[L(\boldsymbol X),\infty),\\
P_\theta(\theta'\in C(\boldsymbol X)),\ \theta'>\theta,&\qquad\text{若 }C(\boldsymbol X)=(-\infty,U(\boldsymbol X)],
\end{aligned}
\end{equation}
即$\theta$是真实参数时覆盖$\theta'$的概率.

对单侧和双侧区间采用不同的错误覆盖概率定义是合理的. 例如，有置信下界时，我们断言$\theta$大于某个值，只有当覆盖了过小的$\theta$值时才会发生错误覆盖. 类似论证导出置信上界与双侧界所用的定义.
\sourcepageanchor{391}
若一个$1-\alpha$置信集在某类$1-\alpha$置信集中使错误覆盖概率最小，则称其为一致最精确置信集（uniformly most accurate，UMA）.\index{SI0189@一致最准确}\semanticindexanchor{189:391} 例如，可以考虑在$[L(\boldsymbol x),\infty)$形式的集合中寻找UMA置信集. UMA置信集可通过反演UMP检验的接受域构造，下面将证明这一点. 遗憾的是，虽然UMA置信集很理想，但它只在相当少见的情形下存在（UMP检验也如此）. 特别地，由于UMP检验通常是单侧的，UMA区间也通常为单侧. 不过，它们构成了优美的理论. 下一个定理说明，$H_0:\theta=\theta_0$对$H_1:\theta>\theta_0$的UMP检验给出UMA置信下界.
\begin{theorem}\label{thm:9.3.5}\index{SI0191@与一致最强检验的关系}\index{SI0518@与一致最准确集的关系}
设$\boldsymbol X\sim f(\boldsymbol x\mid\theta)$，其中$\theta$为实值参数. 对每个$\theta_0\in\Theta$，设$A^*(\theta_0)$是$H_0:\theta=\theta_0$对$H_1:\theta>\theta_0$的水平$\alpha$ UMP检验的接受域. 设$C^*(\boldsymbol x)$为反演这些UMP接受域所得的$1-\alpha$置信集. 则对于任何其他$1-\alpha$置信集$C$，
$$P_\theta(\theta'\in C^*(\boldsymbol X))\le P_\theta(\theta'\in C(\boldsymbol X)),\qquad\text{对所有 }\theta'<\theta.$$
\end{theorem}
\begin{proof}
设$\theta'$为小于$\theta$的任意值. 设$A(\theta')$是反演$C$得到的$H_0:\theta=\theta'$的水平$\alpha$检验的接受域. 因为$A^*(\theta')$是检验$H_0:\theta=\theta'$对$H_1:\theta>\theta'$的UMP接受域，且$\theta>\theta'$，所以
$$\begin{aligned}
P_\theta(\theta'\in C^*(\boldsymbol X))&=P_\theta(\boldsymbol X\in A^*(\theta'))&&\text{（反演置信集）}\\
&\le P_\theta(\boldsymbol X\in A(\theta'))&&\text{（$A^*$为UMP，故对任意$A$成立）}\\
&=P_\theta(\theta'\in C(\boldsymbol X)).&&\text{（反演$A$得到$C$）}
\end{aligned}$$
注意，上述不等号为“$\le$”，因为我们处理的是接受域的概率. 它等于$1-\text{功效}$，所以UMP检验使接受域的概率最小. 因而已证明，当$\theta'<\theta$时，反演UMP检验所得区间使错误覆盖概率最小.
\end{proof}
回顾\autoref{sec:9.2.1}中的讨论. 上述定理的UMA置信集通过反演下列假设的一族检验构造：
$$H_0:\theta=\theta_0\qquad\text{对}\qquad H_1:\theta>\theta_0,$$
其中置信集的形式由备择假设决定. 上述备择假设规定$\theta_0$小于某个特定值，因此产生置信下界；即如果集合是区间，它们就具有$[L(\boldsymbol X),\infty)$的形式.
\begin{example}[UMA置信界]\label{ex:9.3.6}\index{SI0190@正态}
设$X_1,\ldots,X_n$独立同分布于$\mathrm n(\mu,\sigma^2)$，其中$\sigma^2$已知. 区间
$$C(\bar x)=\left\{\mu:\mu\ge\bar x-z_\alpha\frac\sigma{\sqrt n}\right\}$$
是$1-\alpha$ UMA置信下界，因为它可通过反演$H_0:\mu=\mu_0$对$H_1:\mu>\mu_0$的UMP检验得到.

更常见的双侧区间
$$C(\bar x)=\left\{\mu:\bar x-z_{\alpha/2}\frac\sigma{\sqrt n}\le\mu\le\bar x+z_{\alpha/2}\frac\sigma{\sqrt n}\right\}$$
不是UMA，因为它来自$H_0:\mu=\mu_0$对$H_1:\mu\ne\mu_0$检验的双侧接受域反演，而对于这组假设，不存在UMP检验.
\end{example}
\sourcepageanchor{392}
在检验问题中考虑双侧检验时，我们发现无偏性既令人信服又有用. 置信区间问题中也有类似思想. 处理双侧置信区间时，将考虑范围限于无偏置信集是合理的. 记住，无偏检验的备择假设下的功效总大于原假设下的功效. 阅读下面定义时请牢记这一点.
\begin{definition}\label{def:9.3.7}\index{SI0192@无偏}
若一个$1-\alpha$置信集$C(\boldsymbol x)$对所有$\theta\ne\theta'$均满足$P_\theta(\theta'\in C(\boldsymbol X))\le1-\alpha$，则称其无偏（unbiased）.
\end{definition}
因此，对于无偏置信集，错误覆盖概率决不会大于最小正确覆盖概率. 无偏置信集可通过反演无偏检验得到.\index{SI0194@与无偏检验的关系}\semanticindexanchor{194:392}\index{SI0521@与无偏集的关系}\semanticindexanchor{521:392} 即若$A(\theta_0)$是检验$H_0:\theta=\theta_0$对$H_1:\theta\ne\theta_0$的无偏水平$\alpha$接受域，且$C(\boldsymbol x)$是反演这些接受域形成的$1-\alpha$置信集，则$C(\boldsymbol x)$是无偏$1-\alpha$置信集（\autoref{exercise:9.46}）.
\begin{example}[续\autoref{ex:9.3.6}]\label{ex:9.3.8}\index{SI0193@正态}
双侧正态区间
$$C(\bar x)=\left\{\mu:\bar x-z_{\alpha/2}\frac\sigma{\sqrt n}\le\mu\le\bar x+z_{\alpha/2}\frac\sigma{\sqrt n}\right\}$$
是无偏区间. 它可通过反演\autoref{ex:8.3.20}给出的$H_0:\mu=\mu_0$对$H_1:\mu\ne\mu_0$的无偏检验得到. 类似地，基于$t$分布的区间\eqref{eq:9.2.14}也是无偏区间，因为它也可以通过反演无偏检验得到（\autoref{exercise:8.38}）.
\end{example}
使错误覆盖概率最小的集合也称为Neyman最短（Neyman-shortest）集合.\index{SI0721@Neyman最短}\semanticindexanchor{721:392} 这个名称包含长度的意味，\hyperref[bib:source246]{Pratt（1961）}的下述定理为其提供了某种依据.
\begin{theorem}[Pratt]\label{thm:9.3.9}\index{SI0133@与错误覆盖的关系}\index{SI0822@Pratt定理}
设$X$是实值随机变量，$X\sim f(x\mid\theta)$，其中$\theta$是实值参数. 设$C(x)=[L(x),U(x)]$是$\theta$的置信区间. 若$L(x)$和$U(x)$均为$x$的递增函数，则对任意值$\theta^*$，
\begin{equation}\label{eq:9.3.3}
E_{\theta^*}(\operatorname{Length}[C(X)])=\int_{\theta\ne\theta^*}P_{\theta^*}(\theta\in C(X))\,d\theta.
\end{equation}
\end{theorem}
\autoref{thm:9.3.9}表明，$C(x)$的期望长度等于错误覆盖概率之和（积分），其中积分遍及参数的所有错误值.
\begin{proof}
根据期望值的定义，可以写为
$$\begin{aligned}
E_{\theta^*}(\operatorname{Length}[C(X)])
&=\int_{\mathcal X}\operatorname{Length}[C(x)]f(x\mid\theta^*)\,dx\\
&=\int_{\mathcal X}[U(x)-L(x)]f(x\mid\theta^*)\,dx&&\text{（长度的定义）}\\
&=\int_{\mathcal X}\left[\int_{L(x)}^{U(x)}d\theta\right]f(x\mid\theta^*)\,dx&&\text{（用$\theta$作积分变量）}\\
&=\int_\Theta\left[\int_{U^{-1}(\theta)}^{L^{-1}(\theta)}f(x\mid\theta^*)\,dx\right]d\theta&&\text{（交换积分次序，见下文）}\\
&=\int_\Theta[P_{\theta^*}(U^{-1}(\theta)\le X\le L^{-1}(\theta))]\,d\theta&&\text{（定义）}\\
&=\int_\Theta[P_{\theta^*}(\theta\in C(X))]\,d\theta&&\text{（反演接受域）}\\
&=\int_{\theta\ne\theta^*}[P_{\theta^*}(\theta\in C(X))]\,d\theta.&&\text{（一个点不改变积分值）}
\end{aligned}$$
\sourcepageanchor{393}
这串等式建立了所需恒等式，定理得证. 积分次序的交换形式上由Fubini定理保证（\hyperref[bib:source200]{Lehmann和Casella 1998}，第1.2节），但只要所有被积函数有限，就容易看出交换是合理的.\footnote{译注：一般而言，仅要求被积函数处处有限不足以保证交换积分次序；这里的被积函数非负，交换可由Tonelli定理得到（容许积分为无穷）. 原文表述保留.} 置信区间的反演是标准的，其中使用关系
$$\theta\in\{\theta:L(x)\le\theta\le U(x)\}\iff x\in\{x:U^{-1}(\theta)\le x\le L^{-1}(\theta)\},$$
它因$L,U$递增的假设而成立. 注意，修改定理后也可适用于端点递减的区间.
\end{proof}
\autoref{thm:9.3.9}说明，置信区间的长度与错误覆盖概率之间存在形式上的关系. 在双侧情形中，这意味着使错误覆盖概率最小也带来某种长度最优保证. 不过，在单侧情形中，这种类比并不完全成立. 此时，为使错误覆盖概率最小而构造的区间，只关注参数空间一部分中的参数，因而可能没有长度最优性. \hyperref[bib:source213]{Madansky（1962）}给出了一个可被改进的$1-\alpha$ UMA单侧区间例子：可以构造另一个更短的$1-\alpha$区间.（参见\autoref{exercise:9.45}.）此外，\hyperref[bib:source212]{Maatta和Casella（1987）}证明，反演UMP检验所得区间，按其他合理准则衡量时可能并非最优.
\subsubsection{贝叶斯最优性}\label{sec:9.3.3}
获得具有指定覆盖概率的最小置信集这一目标，也可借助贝叶斯准则实现. 若有后验分布$\pi(\theta\mid\boldsymbol x)$，即给定$\boldsymbol X=\boldsymbol x$时$\theta$的后验分布，我们希望找到集合$C(\boldsymbol x)$，使
\begin{enumerate}[label=\roman*.]
\item $\displaystyle\int_{C(\boldsymbol x)}\pi(\theta\mid\boldsymbol x)\,d\boldsymbol x=1-\alpha$；
\item $\operatorname{Size}(C(\boldsymbol x))\le\operatorname{Size}(C'(\boldsymbol x))$，
\end{enumerate}
其中$C'(\boldsymbol x)$为任意满足$\int_{C'(\boldsymbol x)}\pi(\theta\mid\boldsymbol x)\,d\boldsymbol x\ge1-\alpha$的集合.\footnote{译注：原文以上两处后验积分的微分均写作$d\boldsymbol x$；积分区域是参数集合，因此应为$d\theta$. 此处保留原式.}

如果以长度作为大小的度量，则可以应用\autoref{thm:9.3.2}，得到如下结果.
\begin{corollary}\label{cor:9.3.10}\index{SI0471@最高后验密度区域}
若后验密度$\pi(\theta\mid\boldsymbol x)$单峰，则对于给定$\alpha$，$\theta$的最短可信区间为
$$\{\theta:\pi(\theta\mid\boldsymbol x)\ge k\},\qquad\text{其中}\quad\int_{\{\theta:\pi(\theta\mid\boldsymbol x)\ge k\}}\pi(\theta\mid\boldsymbol x)\,d\theta=1-\alpha.$$
\end{corollary}
\autoref{cor:9.3.10}所述的可信集称为最高后验密度区域（highest posterior density region，HPD区域），因为它由后验密度最高的参数值组成. 注意HPD区域与似然区域在形式上的相似性.
\begin{example}[泊松HPD区域]\label{ex:9.3.11}
在\autoref{ex:9.2.16}中，我们导出了泊松参数的$1-\alpha$可信集. 现在构造HPD区域. 由\autoref{cor:9.3.10}，此区域为$\{\lambda:\pi(\lambda\mid\sum x)\ge k\}$，其中$k$选取为使
$$1-\alpha=\int_{\{\lambda:\pi(\lambda\mid\sum x)\ge k\}}\pi(\lambda\mid\textstyle\sum x)\,d\lambda.$$
回顾$\lambda$的后验pdf为$\operatorname{gamma}(a+\sum x,[n+(1/b)]^{-1})$，因此需要找到$\lambda_L,\lambda_U$使\footnote{译注：原文在两个端点密度$\pi(\lambda_L\mid\sum x)$和$\pi(\lambda_U\mid\sum x)$之间漏印关系符号；由HPD区域的等高端点条件，此处应为等号. 原式保留.}
$$\pi(\lambda_L\mid\textstyle\sum x)\pi(\lambda_U\mid\textstyle\sum x)\quad\text{且}\quad\int_{\lambda_L}^{\lambda_U}\pi(\lambda\mid\textstyle\sum x)\,d\lambda=1-\alpha.$$

\sourcepageanchor{394}
若取$a=b=1$（与\autoref{ex:9.2.16}相同），则给定$\sum X=\sum x$时$\lambda$的后验分布可表示为$2(n+1)\lambda\sim\chi^2_{2(\sum x+1)}$；若$n=10$且$\sum x=6$，则$\lambda$的$90\%$ HPD可信集为$[.253,1.005]$.

\autoref{fig:9.3.1}展示了$\lambda$的三个$1-\alpha$区间：\autoref{ex:9.2.16}中的$1-\alpha$等尾贝叶斯可信集、这里导出的HPD区域，以及\autoref{ex:9.2.15}中的经典$1-\alpha$置信集.
\par
\begin{figure}[htbp]\centering
\includegraphics[width=.72\textwidth]{assets/chapter9/fig-9.3.1.pdf}
\caption{\autoref{ex:9.2.16}中的三个区间估计量.}\label{fig:9.3.1}
\end{figure}
\par
\end{example}
HPD区域的形状由后验分布的形状决定. 一般地，HPD区域并不关于某个贝叶斯点估计量对称，而是像似然区域一样相当不对称. 对泊松分布，这显然成立，如上例所示. 尽管不总如此，但通常可预期：尺度参数问题产生不对称HPD区域，位置参数问题产生对称HPD区域.
\begin{example}[正态HPD区域]\label{ex:9.3.12}
\autoref{ex:9.2.18}导出的等尾可信集实际上是HPD区域. 因为$\theta$的后验分布是均值为$\delta^{\mathrm B}$的正态分布，所以对某个$k'$有（\autoref{exercise:9.40}）$\{\theta:\pi(\theta\mid\bar x)\ge k\}=\{\theta:\theta\in\delta^{\mathrm B}\pm k'\}$. 因而HPD区域关于均值$\delta^{\mathrm B}(\bar x)$对称.
\end{example}
\subsubsection{损失函数意义下的最优性}\label{sec:9.3.4}
前两小节研究区间估计量最优性时，先要求它们具有某个最小覆盖概率，然后寻找最短区间. 不过，
\sourcepageanchor{395}
也可以将这些要求合并到一个损失函数中，利用决策论寻找最优估计量. 在区间估计中，行动空间$\mathcal A$由参数空间$\Theta$的子集构成；更正式地说，我们或许应称之为“集合估计”，因为最优规则未必是区间. 然而，出于实际考虑，我们主要研究以区间为集合估计量的情形；幸运的是，许多最优程序最终确实是区间.

用$C$（取自置信区间）表示$\mathcal A$中的元素，行动$C$的含义是作出区间估计“$\theta\in C$\mbox{”.} 决策规则$\delta(\boldsymbol x)$仅仅规定：对每个$\boldsymbol x\in\mathcal X$，若观测到$\boldsymbol X=\boldsymbol x$，将用哪个集合$C\in\mathcal A$估计$\theta$. 因此，仍沿用记号$C(\boldsymbol x)$.

区间估计问题的损失函数通常包括两个量：集合估计是否正确包含真值$\theta$的度量，以及集合估计大小的度量. 大多数时候只考虑作为区间的集合$C$，因而一个自然的大小度量是$\operatorname{Len}(C)=C\text{的长度}$. 为表示正确性度量，通常用
$$I_C(\theta)=\begin{cases}1,&\theta\in C,\\0,&\theta\notin C.\end{cases}$$
即估计正确时$I_C(\theta)=1$，否则为零. 事实上，$I_C(\theta)$就是集合$C$的示性函数. 但应注意，$C$将是由数据$\boldsymbol X$的值决定的随机集合.

损失函数应反映良好估计具有较小$\operatorname{Len}(C)$和较大$I_C(\theta)$这一事实. 一种这样的损失函数是
\begin{equation}\label{eq:9.3.4}\index{SI0146@决策论}\index{SI0657@区间估计}
L(\theta,C)=b\operatorname{Length}(C)-I_C(\theta),
\end{equation}
其中$b$为正常数，反映我们希望赋予两个准则的相对权重；由于这两个量很不相同，这种权衡是必要的. 若更关注估计正确，应取较小的$b$；若更关注区间长度，则应取较大的$b$.

与\eqref{eq:9.3.4}对应的风险函数特别简单：
$$\begin{aligned}
R(\theta,C)&=bE_\theta[\operatorname{Length}(C(\boldsymbol X))]-E_\theta I_{C(\boldsymbol X)}(\theta)\\
&=bE_\theta[\operatorname{Length}(C(\boldsymbol X))]-P_\theta(I_{C(\boldsymbol X)}(\theta)=1)\\
&=bE_\theta[\operatorname{Length}(C(\boldsymbol X))]-P_\theta(\theta\in C(\boldsymbol X)).
\end{aligned}$$
风险有两部分：区间的期望长度和区间估计量的覆盖概率.\index{SI0128@与风险的关系}\semanticindexanchor{128:395}\index{SI0226@与风险的关系}\semanticindexanchor{226:395} 风险反映了与前面各小节相同的愿望：同时希望期望长度小、覆盖概率高. 但现在不再先要求最低覆盖概率、再极小化长度，而是在风险函数中规定这两个量的权衡. 如果能显著缩短长度，较小的覆盖概率也许可以接受.

改变损失\eqref{eq:9.3.4}中$b$的大小，就能改变区间估计量的大小与覆盖概率的相对重要性，这是此前无法做到的. 为展示当前框架的灵活性，考虑某些极限情形. 若$b=0$，则大小不重要，只有覆盖概率重要，因此覆盖概率为$1$的区间估计量$C=(-\infty,\infty)$是最佳决策规则. 类似地，若$b=\infty$，则覆盖概率不重要，因此单点集合最优. 于是，一整系列决策规则都可能成为候选. 下例中，对于某个指定的有限$b$值范围，选择良好规则归结为用风险函数确定置信系数；如果$b$在此范围外，则最优决策规则是点估计量.
\begin{example}[正态区间估计量]\label{ex:9.3.13}\index{SI0148@均值}
设$X\sim\mathrm n(\mu,\sigma^2)$，假设$\sigma^2$已知. $X$通常是样本均值，而$\sigma^2$具有$\tau^2/n$的形式，其中$\tau^2$为已知总体方差，$n$为样本量. 对每个$c\ge0$，定义$\mu$的区间估计量$C(x)=[x-c\sigma,x+c\sigma]$. 将利用\eqref{eq:9.3.4}中的损失比较这些估计量. 区间的
\sourcepageanchor{396}
长度$\operatorname{Len}(C(x))=2c\sigma$不依赖$x$. 因此，风险第一项为$b(2c\sigma)$. 第二项为
$$\begin{aligned}
P_\mu(\mu\in C(X))&=P_\mu(X-c\sigma\le\mu\le X+c\sigma)\\
&=P_\mu\left(-c\le\frac{X-\mu}\sigma\le c\right)\\
&=2P(Z\le c)-1,
\end{aligned}$$
其中$Z\sim\mathrm n(0,1)$. 因而，这一类区间估计量的风险函数为
\begin{equation}\label{eq:9.3.5}
R(\mu,C)=b(2c\sigma)-[2P(Z\le c)-1].
\end{equation}
风险函数不依赖$\mu$，所以是常数；此类中最佳区间估计量对应于使\eqref{eq:9.3.5}最小的$c$值.

若$b\sigma>1/\sqrt{2\pi}$，可证明$R(\mu,C)$在$c=0$时最小. 即损失中长度部分完全压倒了覆盖概率部分，最佳区间估计量是点估计量$C(x)=[x,x]$. 但若$b\sigma\le1/\sqrt{2\pi}$，则风险在$c=\sqrt{-2\log(b\sigma\sqrt{2\pi})}$处最小. 若将$c$写成某个$\alpha$对应的$z_{\alpha/2}$，则使风险最小的区间估计量恰是通常的$1-\alpha$置信区间.（细节见\autoref{exercise:9.53}.）
\end{example}
决策论在区间估计问题中的应用，不像在点估计或假设检验问题中那样广泛.\index{SI0147@困难}\semanticindexanchor{147:396} 一个原因是\eqref{eq:9.3.4}（或\autoref{ex:9.3.13}）中的$b$难以选择. 上例说明，一个看似合理的选择可能产生不符合直觉的结果，表明\eqref{eq:9.3.4}中的损失可能不适当. 有些人在其他问题中使用决策论分析，但仍倾向于仅使用置信系数固定为$1-\alpha$的区间估计量，然后用风险函数评价集合大小等其他品质.

另一个困难在于对$\mathcal A$中允许集合形状的限制. 理想地，应由损失和风险函数判断哪些形状最好. 但总可以向区间估计量加入孤立点，在大小方面不受损失惩罚，却提高覆盖概率. 上例中本可以使用估计量
$$C(x)=[x-c\sigma,x+c\sigma]\cup\{\mu\text{的所有整数值}\}.$$
这些集合的“长度”与此前相同，但现在对$\mu$的所有整数值，覆盖概率都是$1$. 必须采用更精细的大小度量才能避免此类反常现象.（\hyperref[bib:source183]{Joshi 1969}通过定义估计量的等价类处理了这个问题.）



\subsection{习题}\label{sec:9.4}
\begin{exercise}\label{exercise:9.1}
若$L(x),U(x)$满足$P_\theta(L(X)\le\theta)=1-\alpha_1$和$P_\theta(U(X)\ge\theta)=1-\alpha_2$，且对所有$x$均有$L(x)\le U(x)$，证明$P_\theta(L(X)\le\theta\le U(X))=1-\alpha_1-\alpha_2$.
\end{exercise}
\begin{exercise}\label{exercise:9.2}
设$X_1,\ldots,X_n$独立同分布于$\mathrm n(\theta,1)$. $\theta$的一个$95\%$置信区间为$\bar x\pm1.96/\sqrt n$. 令$p$表示额外的独立观测$X_{n+1}$落入这个区间的概率. $p$大于、小于还是等于$.95$？证明你的答案.
\end{exercise}
\begin{exercise}\label{exercise:9.3}
独立随机变量$X_1,\ldots,X_n$具有相同分布
$$P(X_i\le x)=\begin{cases}0,&x\le0,\\(x/\beta)^\alpha,&0<x<\beta,\\1,&x\ge\beta.\end{cases}$$
\begin{enumerate}
\item 在\autoref{exercise:7.10}中已求出$\alpha,\beta$的MLE. 若$\alpha$是已知常数$\alpha_0$，求$\beta$的置信系数为$.95$的置信上限.
\item 利用\autoref{exercise:7.10}的数据构造$\beta$的区间估计. 假设$\alpha$已知且等于其MLE.
\end{enumerate}
\end{exercise}
\sourcepageanchor{397}
\begin{exercise}\label{exercise:9.4}
设$X_1,\ldots,X_n$为来自$\mathrm n(0,\sigma_X^2)$的随机样本，$Y_1,\ldots,Y_m$为来自$\mathrm n(0,\sigma_Y^2)$且与各$X$独立的随机样本. 定义$\lambda=\sigma_Y^2/\sigma_X^2$.
\begin{enumerate}
\item\label{item:9.4.a}求$H_0:\lambda=\lambda_0$对$H_1:\lambda\ne\lambda_0$的水平$\alpha$ LRT.
\item 用$F$随机变量表示\hyperref[item:9.4.a]{(a)}中LRT的拒绝域.
\item 求$\lambda$的$1-\alpha$置信区间.
\end{enumerate}
\end{exercise}
\begin{exercise}\label{exercise:9.5}
\autoref{ex:9.2.5}对一系列伯努利试验的成功概率$p$给出了置信下界. 本题将导出置信上界. 即观测$X_1,\ldots,X_n$，其中$X_i\sim\operatorname{Bernoulli}(p)$，希望得到形式为$[0,U(x_1,\ldots,x_n))$的区间，使$P_p(p\in[0,U(X_1,\ldots,X_n)))\ge1-\alpha$.
\begin{enumerate}
\item 证明反演检验
$$H_0:p=p_0\qquad\text{对}\qquad H_1:p<p_0$$
的接受域，将给出具有所需置信水平和形式的置信区间.
\item 求类似\eqref{eq:9.2.8}的方程，用以构造置信区间.
\end{enumerate}
\end{exercise}
\begin{exercise}\label{exercise:9.6}\index{SI0159@二项}\leavevmode
\begin{enumerate}
\item 通过反演$H_0:p=p_0$对$H_1:p\ne p_0$的LRT，导出二项参数$p$的置信区间.
\item 证明该区间是来自$p^y(1-p)^{n-y}$的最高密度区域，且不等于\eqref{eq:10.4.4}中的区间.
\end{enumerate}
\end{exercise}
\begin{exercise}\label{exercise:9.7}\leavevmode
\begin{enumerate}
\item 基于来自$\mathrm n(\theta,a\theta)$族的样本$X_1,\ldots,X_n$（其中$\theta$未知），求反演$H_0:a=a_0$对$H_1:a\ne a_0$的LRT所得的$a$的$1-\alpha$置信集.
\item 对相关的$\mathrm n(\theta,a\theta^2)$族也可提出类似问题. 若$X_1,\ldots,X_n$独立同分布于$\mathrm n(\theta,a\theta^2)$，其中$\theta$未知，求基于反演$H_0:a=a_0$对$H_1:a\ne a_0$的LRT的$1-\alpha$置信集.
\end{enumerate}
\end{exercise}
\begin{exercise}\label{exercise:9.8}
给定来自形式为$\sigma^{-1}f((x-\theta)/\sigma)$的pdf的样本$X_1,\ldots,X_n$，列出至少五个不同的枢轴量.
\end{exercise}
\begin{exercise}\label{exercise:9.9}
证明\autoref{ex:9.2.7}所列三个量均为枢轴量.
\end{exercise}
\begin{exercise}\label{exercise:9.10}\leavevmode
\begin{enumerate}
\item 设$T$为实值统计量，并假设对于每个$\theta\in\Theta$，$Q(t,\theta)$都是$t$的单调函数. 证明：若$T$的pdf $f(t\mid\theta)$对某个函数$g$可表示为\eqref{eq:9.2.11}的形式，则$Q(T,\theta)$为枢轴量.
\item 证明取$g=1$及$Q(t,\theta)=F_\theta(t)$（$T$的cdf）即可满足\eqref{eq:9.2.11}.（这就是概率积分变换.）
\end{enumerate}
\end{exercise}
\begin{exercise}\label{exercise:9.11}
若$T$为cdf为$F_T(t\mid\theta)$的连续随机变量，且$\alpha_1+\alpha_2=\alpha$，证明假设$H_0:\theta=\theta_0$的一个水平$\alpha$接受域为$\{t:\alpha_1\le F_T(t\mid\theta_0)\le1-\alpha_2\}$，其对应的$1-\alpha$置信集为$\{\theta:\alpha_1\le F_T(t\mid\theta)\le1-\alpha_2\}$.
\end{exercise}
\begin{exercise}\label{exercise:9.12}
基于来自$\mathrm n(\theta,\theta)$总体、容量为$n$的随机样本寻找一个枢轴量，其中$\theta>0$. 利用该枢轴量建立$\theta$的$1-\alpha$置信区间.
\end{exercise}
\begin{exercise}\label{exercise:9.13}
设$X$为来自$\operatorname{beta}(\theta,1)$密度的单个观测.
\begin{enumerate}
\item\label{item:9.13.a}令$Y=-(\log X)^{-1}$. 计算集合$[y/2,y]$的置信系数.
\item 寻找一个枢轴量，并用它建立一个置信区间，使其置信系数与\hyperref[item:9.13.a]{(a)}中的区间相同.
\item 比较两个置信区间.
\end{enumerate}
\end{exercise}
\begin{exercise}\label{exercise:9.14}
设$X_1,\ldots,X_n$独立同分布于$\mathrm n(\mu,\sigma^2)$，两个参数均未知. 可以用多种方式通过Bonferroni不等式对$\mu,\sigma$进行同时推断.
\begin{enumerate}
\item\label{item:9.14.a}利用Bonferroni不等式，将两个置信集
$$\left\{\mu:\bar x-\frac{ks}{\sqrt n}\le\mu\le\bar x+\frac{ks}{\sqrt n}\right\}\quad\text{与}\quad\left\{\sigma^2:\frac{(n-1)s^2}b\le\sigma^2\le\frac{(n-1)s^2}a\right\}$$
合并为$(\mu,\sigma)$的一个置信集. 说明如何选择$a,b,k$，使该同时置信集为$1-\alpha$置信集.
\sourcepageanchor{398}
\item\label{item:9.14.b}利用Bonferroni不等式，将两个置信集
$$\left\{\mu:\bar x-\frac{k\sigma}{\sqrt n}\le\mu\le\bar x+\frac{k\sigma}{\sqrt n}\right\}\quad\text{与}\quad\left\{\sigma^2:\frac{(n-1)s^2}b\le\sigma^2\le\frac{(n-1)s^2}a\right\}$$
合并为$(\mu,\sigma)$的一个置信集. 说明如何选择$a,b,k$，使该同时置信集为$1-\alpha$置信集.
\item 比较\hyperref[item:9.14.a]{(a)}和\hyperref[item:9.14.b]{(b)}中的置信集.
\end{enumerate}
\end{exercise}
\begin{exercise}\label{exercise:9.15}
求解定义正态均值比的Fieller置信集的二次方程之根（见\autoref{sec:9.5.3}）. 求出随机变量满足何种条件时，
\begin{enumerate}
\item 抛物线开口向上（置信集为一个区间）.
\item 抛物线开口向下（置信集为一个区间的补集）.
\item 抛物线没有实根.
\end{enumerate}
在每种情形下解释置信集的含义. 例如，若某位实验者的数据导致抛物线没有实根，你会如何向他解释？
\end{exercise}
\begin{exercise}\label{exercise:9.16}
设$X_1,\ldots,X_n$独立同分布于$\mathrm n(\theta,\sigma^2)$，其中$\sigma^2$已知. 对下列每组假设，写出一个水平$\alpha$检验的接受域，以及反演该检验得到的$1-\alpha$置信区间.
\begin{enumerate}
\item $H_0:\theta=\theta_0$对$H_1:\theta\ne\theta_0$.
\item $H_0:\theta\ge\theta_0$对$H_1:\theta<\theta_0$.
\item $H_0:\theta\le\theta_0$对$H_1:\theta>\theta_0$.
\end{enumerate}
\end{exercise}
\begin{exercise}\label{exercise:9.17}\index{SI0145@构造}
给定独立同分布的$X_1,\ldots,X_n$，其pdf如下，求$\theta$的$1-\alpha$置信区间.
\begin{enumerate}
\item $f(x\mid\theta)=1$，$\theta-\tfrac12<x<\theta+\tfrac12$.
\item $f(x\mid\theta)=2x/\theta^2$，$0<x<\theta$，$\theta>0$.
\end{enumerate}
\end{exercise}
\begin{exercise}\label{exercise:9.18}\index{SI0173@二项}
本题进一步研究二项置信集的性质，特别是\hyperref[bib:source305]{Sterne（1954）}的构造. 与\autoref{ex:9.2.11}一样，仍考虑$\operatorname{binomial}(3,p)$分布.
\begin{enumerate}
\item 画出四个概率函数$P_p(X=x)$，$x=0,\ldots,3$，随$p$变化的图像. 确定$P_p(X=1)$和$P_p(X=2)$的最大值.
\item 证明对于小的$\varepsilon$，当$p=\frac13+\varepsilon$时有$P_p(X=0)>P_p(X=2)$.
\item 证明以下接受域可被反演为$1-\alpha=.442$置信\emph{区间}，从而说明Sterne集合中的困难应归咎于\emph{最大概率构造}.
\begin{center}\begin{tabular}{cc}\toprule
$p$&接受域\\\midrule
$[.000,.238]$&$\{0\}$\\
$(.238,.305)$&$\{0,1\}$\\
$[.305,.362]$&$\{1\}$\\
$(.362,.634)$&$\{1,2\}$\\
$[.634,.695]$&$\{2\}$\\
$(.695,.762)$&$\{2,3\}$\\
$[.762,1.00]$&$\{3\}$\\\bottomrule
\end{tabular}\end{center}
（这实质上是\hyperref[bib:source88]{Crow（1956）}对Sterne构造的修改，参见\autoref{sec:9.5.2}.）
\end{enumerate}
\end{exercise}
\begin{exercise}\label{exercise:9.19}
证明\autoref{thm:9.2.12}的\hyperref[item:9.2.12.ii]{(b)部分}.
\end{exercise}
\begin{exercise}\label{exercise:9.20}
需要补全\autoref{thm:9.2.14}证明中的一些细节，并证明定理的第二部分.
\begin{enumerate}
\item 证明：若$F_T(T\mid\theta)$随机大于或等于均匀随机变量，则$\overline F_T(T\mid\theta)$也如此. 即若对每个$0\le x\le1$都有$P_\theta(F_T(T\mid\theta)\le x)\le x$，则对每个$0\le x\le1$也有$P_\theta(\overline F_T(T\mid\theta)\le x)\le x$.
\item 证明：对于$\alpha_1+\alpha_2=\alpha$，集合$\{\theta:F_T(T\mid\theta)\le\alpha_1\text{ 且 }\overline F_T(T\mid\theta)\le\alpha_2\}$是$1-\alpha$置信集.\footnote{译注：本分题沿用了\autoref{thm:9.2.14}证明中的不等号疑误；要构造中央置信集，应排除下尾或上尾概率过小的参数值，原文两个“$\le$”的方向与此不符. 原式保留.}
\item 若cdf $F_T(t\mid\theta)$对每个$t$均为$\theta$的递减函数，证明由$\overline F_T(t\mid\theta)=P(T\ge t\mid\theta)$定义的函数$\overline F_T(t\mid\theta)$，对每个$t$均为$\theta$的非递减函数.
\item 证明\autoref{thm:9.2.14}的\hyperref[item:9.2.14.ii]{(b)部分}.
\end{enumerate}
\end{exercise}
\sourcepageanchor{399}
\begin{exercise}\label{exercise:9.21}\index{SI0143@与F分布的关系}
\autoref{ex:9.2.15}表明，泊松参数的置信区间可用卡方临界点表示. 用类似技术证明，若$X\sim\operatorname{binomial}(n,p)$，则$p$的一个$1-\alpha$置信区间为
$$\frac1{1+\frac{n-x+1}xF_{2(n-x+1),2x,\alpha/2}}\le p\le\frac{\frac{x+1}{n-x}F_{2(x+1),2(n-x),\alpha/2}}{1+\frac{x+1}{n-x}F_{2(x+1),2(n-x),\alpha/2}},$$
其中$F_{\nu_1,\nu_2,\alpha}$为自由度$\nu_1,\nu_2$的$F$分布的上侧$\alpha$临界点，并作如下端点调整：若$x=0$，下端点为$0$；若$x=n$，上端点为$1$. 这些就是\hyperref[bib:source79]{Clopper和Pearson（1934）}区间.

（提示：回顾\autoref{exercise:2.40}中的恒等式，它可解释为：若$X\sim\operatorname{binomial}(n,\theta)$，则$P_\theta(X\ge x)=P(Y\le\theta)$，其中$Y\sim\operatorname{beta}(x,n-x+1)$. 使用\autoref{sec:5}中$F$分布与贝塔分布的性质.）
\end{exercise}
\begin{exercise}\label{exercise:9.22}\index{SI0162@与F分布的关系}
若$X\sim\operatorname{negative\ binomial}(r,p)$，利用二项分布与负二项分布的关系，证明$p$的一个$1-\alpha$置信区间为
$$\frac1{1+\frac{x+1}rF_{2(x+1),2r,\alpha/2}}\le p\le\frac{\frac rxF_{2r,2x,\alpha/2}}{1+\frac rxF_{2r,2x,\alpha/2}},$$
在$x=0$时作适当修改.
\end{exercise}
\begin{exercise}\label{exercise:9.23}\leavevmode
\begin{enumerate}
\item 设$X_1,\ldots,X_n$为来自参数为$\lambda$的泊松总体的随机样本，定义$Y=\sum X_i$. \autoref{ex:9.2.15}用\autoref{sec:9.2.3}的方法求得了$\lambda$的置信区间. 通过反演LRT构造$\lambda$的另一个区间，并比较这些区间.
\item 以下数据是某块马铃薯田中九行作物每行的蚜虫数，可以假定它们服从泊松分布：
$$155,\quad104,\quad66,\quad50,\quad36,\quad40,\quad30,\quad35,\quad42.$$
用这些数据为每行的平均蚜虫数构造$90\%$ LRT置信区间. 再用\autoref{ex:9.2.15}的方法构造一个区间.
\end{enumerate}
\end{exercise}
\begin{exercise}\label{exercise:9.24}
对于$X\sim\operatorname{Poisson}(\lambda)$，证明\autoref{ex:9.2.15}中置信区间$[L(X),U(X)]$的覆盖概率为
$$P_\lambda(\lambda\in[L(X),U(X)])=\sum_{x=0}^\infty I_{[L(x),U(x)]}(\lambda)\frac{e^{-\lambda}\lambda^x}{x!},$$
并且可以定义函数$x_l(\lambda),x_u(\lambda)$使
$$P_\lambda(\lambda\in[L(X),U(X)])=\sum_{x=x_l(\lambda)}^{x_u(\lambda)}\frac{e^{-\lambda}\lambda^x}{x!}.$$
由此解释为什么\autoref{fig:9.2.5}所给泊松区间的覆盖概率图像，在各个置信区间的端点处发生跳跃.
\end{exercise}
\begin{exercise}\label{exercise:9.25}
若$X_1,\ldots,X_n$独立同分布，pdf为$f(x\mid\mu)=e^{-(x-\mu)}I_{[\mu,\infty)}(x)$，则$Y=\min\{X_1,\ldots,X_n\}$对$\mu$充分，其pdf为
$$f_Y(y\mid\mu)=ne^{-n(y-\mu)}I_{[\mu,\infty)}(y).$$
\autoref{ex:9.2.13}用\autoref{sec:9.2.3}的方法求得$\mu$的一个$1-\alpha$置信区间. 将它与似然法、枢轴法所得的$1-\alpha$区间比较.
\end{exercise}
\begin{exercise}\label{exercise:9.26}
设$X_1,\ldots,X_n$为来自$\operatorname{beta}(\theta,1)$密度的独立同分布观测，并假设$\theta$具有$\operatorname{gamma}(r,\lambda)$先验pdf. 求$\theta$的一个$1-\alpha$贝叶斯可信集.
\end{exercise}
\begin{exercise}\label{exercise:9.27}\leavevmode
\begin{enumerate}
\item 设$X_1,\ldots,X_n$为来自$\operatorname{exponential}(\lambda)$密度的独立同分布观测，其中$\lambda$具有共轭$\operatorname{IG}(a,b)$先验，即逆伽马分布，其pdf为
$$\pi(\lambda\mid a,b)=\frac1{\Gamma(a)b^a}\left(\frac1\lambda\right)^{a+1}e^{-1/(b\lambda)},\qquad0<\lambda<\infty.$$
说明如何寻找$\lambda$的$1-\alpha$贝叶斯HPD可信集.
\sourcepageanchor{400}
\item\label{item:9.27.b}基于样本方差$s^2$，并对正态分布方差$\sigma^2$采用共轭$\operatorname{IG}(a,b)$先验，求$\sigma^2$的$1-\alpha$贝叶斯HPD可信集.
\item 从\hyperref[item:9.27.b]{(b)}的区间出发，求$a\to0$且$b\to\infty$时所得的$\sigma^2$的极限$1-\alpha$贝叶斯HPD可信集.
\end{enumerate}
\end{exercise}
\begin{exercise}\label{exercise:9.28}
设$X_1,\ldots,X_n$独立同分布于$\mathrm n(\theta,\sigma^2)$，其中$\theta,\sigma^2$均未知，但只关心关于$\theta$的推断. 考虑先验pdf
$$\pi(\theta,\sigma^2\mid\mu,\tau^2,a,b)=\frac1{\sqrt{2\pi\tau^2\sigma^2}}e^{-(\theta-\mu)^2/(2\tau^2\sigma^2)}\frac1{\Gamma(a)b^a}\left(\frac1{\sigma^2}\right)^{a+1}e^{-1/(b\sigma^2)},$$
即$\mathrm n(\mu,\tau^2\sigma^2)$乘以$\operatorname{IG}(a,b)$.
\begin{enumerate}
\item 证明这个先验是该问题的共轭先验.
\item\label{item:9.28.b}求$\theta$的后验分布，并用它构造$\theta$的$1-\alpha$可信集.
\item $\theta$的经典$1-\alpha$置信集可以表示为
$$\left\{\theta:|\theta-\bar x|^2\le F_{1,n-1,\alpha/2}\frac{s^2}n\right\}.$$
是否存在$\tau^2,a,b$的某个（极限）序列，使\hyperref[item:9.28.b]{(b)}中的贝叶斯集合逼近这个集合？\footnote{译注：原式的$F$分布上侧概率写为$\alpha/2$；由$t_{n-1,\alpha/2}^2=F_{1,n-1,\alpha}$，通常的$1-\alpha$双侧$t$区间平方后应使用$F_{1,n-1,\alpha}$. 原式保留.}
\end{enumerate}
\end{exercise}
\begin{exercise}\label{exercise:9.29}
若$X_1,\ldots,X_n$是一系列$n$次$\operatorname{Bernoulli}(p)$试验，
\begin{enumerate}
\item 采用共轭$\operatorname{beta}(a,b)$先验，计算$p$的$1-\alpha$可信集.
\item 利用贝塔分布和$F$分布的关系，将可信集写成可与\autoref{exercise:9.21}的区间形式比较的形式. 比较这些区间.
\end{enumerate}
\end{exercise}
\begin{exercise}\label{exercise:9.30}
完成\autoref{ex:9.2.17}所需的可信概率计算.
\begin{enumerate}
\item\label{item:9.30.a}假设$a$是整数，证明$T=2(nb+1)\lambda/b\sim\chi^2_{2(a+\sum x)}$.
\item 证明当$\nu\to\infty$时，
$$\frac{\chi^2_\nu-\nu}{\sqrt{2\nu}}\to\mathrm n(0,1).$$
（使用矩母函数. 这个极限难以计算：先取对数，再作Taylor展开. 或参见\autoref{ex:12.6.8}.）
\item 将\hyperref[item:9.30.a]{(a)}中的随机变量$T$标准化，并用这个变量写出可信概率\eqref{eq:9.2.21}. 证明当$\sum x_i\to\infty$时，标准化后的下侧临界点$\to\infty$，从而可信概率趋于零.
\end{enumerate}
\end{exercise}
\begin{exercise}\label{exercise:9.31}
完成\autoref{ex:9.2.17}所需的覆盖概率计算.
\begin{enumerate}
\item 若$\chi^2_{2Y}$为卡方随机变量，且$Y\sim\operatorname{Poisson}(\lambda)$，证明$E(\chi^2_{2Y})=2\lambda$，$\operatorname{Var}(\chi^2_{2Y})=8\lambda$，$\chi^2_Y$的mgf为$\exp(-\lambda+\lambda/(1-2t))$，并且当$\lambda\to\infty$时，
$$\frac{\chi^2_{2Y}-2\lambda}{\sqrt{8\lambda}}\to\mathrm n(0,1).$$
（使用矩母函数.）\footnote{译注：本分题的mgf所对应随机变量原文写为$\chi_Y^2$，而题首和后续公式均为$\chi_{2Y}^2$；由给定mgf，其自由度下标应为$2Y$. 原文保留.}
\item 现在先将$\chi^2_{2Y}$标准化，再计算$\lambda\to\infty$时的\eqref{eq:9.2.22}. 证明标准化后的上限在$\lambda\to\infty$时趋于$-\infty$，从而覆盖概率趋于零.
\end{enumerate}
\end{exercise}
\begin{exercise}\label{exercise:9.32}\index{SI0234@正态}
本题计算\eqref{eq:9.2.23}中HPD区域的经典覆盖概率，即在概率模型$\bar X\sim\mathrm n(\theta,\sigma^2/n)$下贝叶斯HPD区域的覆盖概率.
\sourcepageanchor{401}
\begin{enumerate}
\item 利用\autoref{ex:9.3.12}中的定义，证明
$$\begin{aligned}
&P_\theta\left(|\theta-\delta^{\mathrm B}(\bar X)|\le z_{\alpha/2}\sqrt{\operatorname{Var}(\theta\mid\bar X)}\right)\\
&=P_\theta\left[-\sqrt{1+\gamma}\,z_{\alpha/2}+\frac{\gamma(\theta-\mu)}{\sigma/\sqrt n}\le Z\le\sqrt{1+\gamma}\,z_{\alpha/2}+\frac{\gamma(\theta-\mu)}{\sigma/\sqrt n}\right].
\end{aligned}$$
\item 证明上述集合虽然是$1-\alpha$可信集，却不是$1-\alpha$置信集.（固定$\theta\ne\mu$，令$\tau=\sigma/\sqrt n$，使$\gamma=1$. 证明当$\sigma^2/n\to0$时，上述概率趋于零.）
\item 不过，若$\theta=\mu$，证明覆盖概率有一个远离零的下界. 求此覆盖概率的最小值与最大值.
\item 现在来看另一面. $\theta$的通常$1-\alpha$置信集为$\{\theta:|\theta-\bar x|\le z_{\alpha/2}\sigma/\sqrt n\}$. 证明这个集合的可信概率为
$$\begin{aligned}
&P_{\bar x}(|\theta-\bar x|\le z_{\alpha/2}\sigma/\sqrt n)\\
&=P_{\bar x}\left[-\sqrt{1+\gamma}\,z_{\alpha/2}+\frac{\gamma(\bar x-\mu)}{\sqrt{1+\gamma}\,\sigma/\sqrt n}\le Z\le\sqrt{1+\gamma}\,z_{\alpha/2}+\frac{\gamma(\bar x-\mu)}{\sqrt{1+\gamma}\,\sigma/\sqrt n}\right],
\end{aligned}$$
且此概率没有远离零的下界. 因而这个$1-\alpha$置信集不是$1-\alpha$可信集.
\end{enumerate}
\end{exercise}
\begin{exercise}\label{exercise:9.33}\index{SI0144@生物等效性}
设$X\sim\mathrm n(\mu,1)$，考虑置信区间
$$C_a(x)=\{\mu:\min\{0,(x-a)\}\le\mu\le\max\{0,(x+a)\}\}.$$
\begin{enumerate}
\item 对$a=1.645$，证明除$\mu=0$时覆盖概率为$1$以外，$C_a(x)$的覆盖概率对所有$\mu$都恰为$.95$.
\item 现在考虑所谓无信息先验$\pi(\mu)=1$. 使用此先验并仍取$a=1.645$，证明$C_a(x)$的后验可信概率在$-1.645\le x\le1.645$时恰为$.90$，且当$|x|\to\infty$时上升至$.95$.
\end{enumerate}
这类区间出现在生物等效性（bioequivalence）问题中，其目标是确定两种处理（药物的不同制剂、治疗的不同给药系统）是否产生相同效应. 这个问题的表述导致原假设与备择假设的角色“互换”（\autoref{exercise:8.47}），从而产生一些有趣的统计问题. 生物等效性的综述见\hyperref[bib:source27]{Berger和Hsu（1996）}，置信集的推广见\hyperref[bib:source53]{Brown、Casella和Hwang（1995）}.
\end{exercise}
\begin{exercise}\label{exercise:9.34}
设$X_1,\ldots,X_n$是来自$\mathrm n(\mu,\sigma^2)$总体的随机样本.
\begin{enumerate}
\item 若$\sigma^2$已知，求保证$\mu$的$.95$置信区间长度不超过$\sigma/4$的最小$n$.
\item 若$\sigma^2$未知，求以$.90$的概率保证$\mu$的$.95$置信区间长度不超过$\sigma/4$的最小$n$.
\end{enumerate}
\end{exercise}
\begin{exercise}\label{exercise:9.35}
设$X_1,\ldots,X_n$为来自$\mathrm n(\mu,\sigma^2)$总体的随机样本. 比较在以下假设下计算的$\mu$的$1-\alpha$置信区间的期望长度：
\begin{enumerate}
\item $\sigma^2$已知.
\item $\sigma^2$未知.
\end{enumerate}
\end{exercise}
\begin{exercise}\label{exercise:9.36}
设$X_1,\ldots,X_n$相互独立，pdf为$f_{X_i}(x\mid\theta)=e^{i\theta-x}I_{[i\theta,\infty)}(x)$. 证明$T=\min_i(X_i/i)$是$\theta$的充分统计量. 基于$T$，求形式为$[T+a,T+b]$且长度最短的$\theta$的$1-\alpha$置信区间.
\end{exercise}
\begin{exercise}\label{exercise:9.37}
设$X_1,\ldots,X_n$独立同分布于$\operatorname{uniform}(0,\theta)$，$Y$为最大次序统计量. 证明$Y/\theta$是枢轴量，并证明区间
$$\left\{\theta:y\le\theta\le\frac y{\alpha^{1/n}}\right\}$$
是最短的$1-\alpha$枢轴区间.
\end{exercise}
\sourcepageanchor{402}
\begin{exercise}\label{exercise:9.38}
若在\autoref{thm:9.3.2}中假设$f$连续，则可以简化证明. 固定$c$，考虑积分$\int_a^{a+c}f(x)\,dx$.
\begin{enumerate}
\item 证明$\frac d{da}\int_a^{a+c}f(x)\,dx=f(a+c)-f(a)$.
\item 证明$f$的单峰性蕴含：当$a$满足$f(a+c)-f(a)=0$时，$\int_a^{a+c}f(x)\,dx$达到最大.
\item 假设给定$\alpha$，选择$c^*,a^*$满足$\int_{a^*}^{a^*+c^*}f(x)\,dx=1-\alpha$及$f(a^*+c^*)-f(a^*)=0$. 证明这是最短$1-\alpha$区间.
\end{enumerate}
\end{exercise}
\begin{exercise}\label{exercise:9.39}
证明\autoref{thm:9.3.2}的一个特例. 设$X\sim f(x)$，其中$f$为对称单峰pdf. 对固定的$1-\alpha$，在所有满足$\int_a^b f(x)\,dx=1-\alpha$的区间$[a,b]$中，最短区间通过选择满足$\int_{-\infty}^a f(x)\,dx=\alpha/2$及$\int_b^\infty f(x)\,dx=\alpha/2$的$a,b$得到.
\end{exercise}
\begin{exercise}\label{exercise:9.40}
在\autoref{exercise:9.39}的基础上，证明若$f$对称，则最优区间具有$m\pm k$的形式，其中$m$为$f$的众数，$k$为常数. 由此证明：
\begin{enumerate}
\item 若$k'$不依赖参数（见\eqref{eq:9.2.7}），则对称似然函数产生关于MLE对称的似然区域.
\item 对称后验密度产生关于后验均值对称的HPD区域.
\end{enumerate}
\end{exercise}
\begin{exercise}\label{exercise:9.41}\leavevmode
\begin{enumerate}
\item\label{item:9.41.a}证明下述与\autoref{thm:9.3.2}有关的结论. 设$X\sim f(x)$，其中$f$是$[0,\infty)$上严格递减的pdf. 对固定$1-\alpha$，在所有满足$\int_a^b f(x)\,dx=1-\alpha$的区间$[a,b]$中，最短区间由$a=0$且$\int_0^b f(x)\,dx=1-\alpha$确定.
\item 利用\hyperref[item:9.41.a]{(a)}的结果，求\autoref{ex:9.2.13}中最短的$1-\alpha$置信区间.
\end{enumerate}
\end{exercise}
\begin{exercise}\label{exercise:9.42}
回顾\autoref{ex:9.3.4}，为寻找伽马尺度参数的最短枢轴区间，我们必须解一个约束极小化问题.
\begin{enumerate}
\item 证明解由满足$\int_a^b f_Y(y)\,dy=1-\alpha$以及$f(b)b^2=f(a)a^2$的$a,b$给出.
\item 对来自$\operatorname{gamma}(k,\beta)$密度的一个观测，已知形状参数$k$，求形式为$\{\beta:x/b\le\beta\le x/a\}$的最短$1-\alpha$（枢轴）置信区间.
\end{enumerate}
\end{exercise}
\begin{exercise}\label{exercise:9.43}
\hyperref[bib:source184]{Juola（1993）}作了如下观察. 若有枢轴量$Q(\boldsymbol X,\theta)$，构造$1-\alpha$置信区间涉及寻找$a,b$使$P(a<Q<b)=1-\alpha$. 通常区间长度是$a,b$的某个函数，例如$b-a$或$1/b^2-1/a^2$. 若$Q$的密度为$f$，且长度可表示为$\int_a^b g(t)\,dt$，则最短枢轴区间是下述问题的解：
$$\min_{\{a,b\}}\int_a^b g(t)\,dt\qquad\text{满足}\quad\int_a^b f(t)\,dt=1-\alpha,$$
或更一般地，
$$\min_C\int_C g(t)\,dt\qquad\text{满足}\quad\int_C f(t)\,dt\ge1-\alpha.$$
\begin{enumerate}
\item\label{item:9.43.a}证明解为$C=\{t:g(t)<\lambda f(t)\}$，其中$\lambda$选为使$\int_C f(t)\,dt=1-\alpha$.（提示：可以改编Neyman--Pearson引理\autoref{thm:8.3.12}的证明.）
\item 应用\hyperref[item:9.43.a]{(a)}的结果，得到\autoref{exercise:9.37}和\autoref{exercise:9.42}中的最短区间.
\end{enumerate}
\end{exercise}
\begin{exercise}\label{exercise:9.44}\leavevmode
\begin{enumerate}
\item 设$X_1,\ldots,X_n$独立同分布于$\operatorname{Poisson}(\lambda)$. 反演$H_0:\lambda=\lambda_0$对$H_1:\lambda>\lambda_0$的水平$\alpha$ UMP检验，求一个UMA $1-\alpha$置信区间.
\item 设$f(x\mid\theta)$为$\operatorname{logistic}(\theta,1)$位置pdf. 基于一个观测$x$，求形式为$\{\theta:\theta\le U(x)\}$的UMA单侧$1-\alpha$置信区间.
\end{enumerate}
\end{exercise}
\begin{exercise}\label{exercise:9.45}\index{SI0175@长度优势}
设$X_1,\ldots,X_n$独立同分布于$\operatorname{exponential}(\lambda)$.
\begin{enumerate}
\item\label{item:9.45.a}求$H_0:\lambda=\lambda_0$对$H_1:\lambda<\lambda_0$的大小为$\alpha$的UMP假设检验.
\item 反演\hyperref[item:9.45.a]{(a)}中的检验，求一个UMA $1-\alpha$置信区间. 证明该区间可表示为
$$C^*(x_1,\ldots,x_n)=\left\{\lambda:0\le\lambda\le\frac{2\sum x_i}{\chi^2_{2n,\alpha}}\right\}.$$
\item 求$C^*(x_1,\ldots,x_n)$的期望长度.
\sourcepageanchor{403}
\item \hyperref[bib:source213]{Madansky（1962）}展示了一个期望长度比UMA区间更短的$1-\alpha$区间. 一般而言，Madansky区间难以计算，但在以下情形中计算相对简单. 令$1-\alpha=.3$，$n=120$. Madansky区间为
$$C^{\mathrm M}(x_1,\ldots,x_n)=\left\{\lambda:0\le\lambda\le-\frac{x_{(1)}}{\log(.99)}\right\},$$
它是$30\%$置信区间. 利用$\chi^2_{240,.7}=251.046$这一事实，证明$30\%$ UMA区间满足\footnote{译注：原文的卡方上侧概率下标$.7$与数值$251.046$不符；$251.046$约为$\chi^2_{240}$的上侧$.3$临界点. 本题关于$30\%$ UMA区间的说明以及前述临界点公式因此需一并核查，原式保留.}
$$\begin{aligned}
E[\operatorname{Length}(C^*(x_1,\ldots,x_n))]&=.956\lambda\\
&>E[\operatorname{Length}(C^{\mathrm M}(x_1,\ldots,x_n))]=.829\lambda.
\end{aligned}$$

\end{enumerate}
\end{exercise}
\begin{exercise}\label{exercise:9.46}
证明：若$A(\theta_0)$是$H_0:\theta=\theta_0$对$H_1:\theta\ne\theta_0$的无偏水平$\alpha$检验接受域，而$C(\boldsymbol x)$为反演这些接受域形成的$1-\alpha$置信集，则$C(\boldsymbol x)$是无偏$1-\alpha$置信集.
\end{exercise}
\begin{exercise}\label{exercise:9.47}\index{SI0176@正态}
设$X_1,\ldots,X_n$为来自$\mathrm n(\theta,\sigma^2)$总体的随机样本，其中$\sigma^2$已知. 证明通常的单侧$1-\alpha$置信上界$\{\theta:\theta\le\bar x+z_\alpha\sigma/\sqrt n\}$无偏，相应置信下界也无偏.
\end{exercise}
\begin{exercise}\label{exercise:9.48}
设$X_1,\ldots,X_n$为来自$\mathrm n(\theta,\sigma^2)$总体的随机样本，其中$\sigma^2$未知.
\begin{enumerate}
\item\label{item:9.48.a}证明区间$\theta\le\bar x+t_{n-1,\alpha}s/\sqrt n$可由反演LRT接受域导出.
\item\label{item:9.48.b}证明\eqref{eq:9.2.14}中相应的双侧区间也可由反演LRT接受域导出.
\item 证明\hyperref[item:9.48.a]{(a)}和\hyperref[item:9.48.b]{(b)}中的区间无偏.
\end{enumerate}
\end{exercise}
\begin{exercise}[Cox悖论]\label{exercise:9.49}\index{SI0227@Cox悖论}
我们要检验
$$H_0:\theta=\theta_0\qquad\text{对}\qquad H_1:\theta>\theta_0,$$
其中$\theta$为两个正态分布之一的均值，$\theta_0$是$\theta$的一个固定但任意的值. 观测到随机变量$X$，其分布为
$$X\sim\begin{cases}\mathrm n(\theta,100),&\text{概率为 }p,\\\mathrm n(\theta,1),&\text{概率为 }1-p.\end{cases}$$
\begin{enumerate}
\item 证明以下检验是水平$\alpha$检验：
$$\text{若 }X>\theta_0+z_\alpha\sigma,\text{ 则拒绝 }H_0,$$
其中$\sigma=1$或$10$取决于从哪个总体抽样. 反演这个检验的接受域，导出$1-\alpha$置信集.
\item 证明一个功效更高的水平$\alpha$检验（对$\alpha>p$）为
$$\begin{gathered}
\text{若 }X>\theta_0+z_{(\alpha-p)/(1-p)}\text{ 且 }\sigma=1,\text{ 则拒绝 }H_0,\\
\text{否则总是拒绝 }H_0.
\end{gathered}$$
反演该检验的接受域导出$1-\alpha$置信集，并证明它以正概率为空集.（Cox悖论指出，经典最优程序有时忽略条件分布的信息，给出一个虽最优、却在某种意义上不合理的程序；参见\hyperref[bib:source84]{Cox 1958}或\hyperref[bib:source83]{Cornfield 1969}.）\footnote{译注：原文“否则总是拒绝$H_0$”若按字面理解，会使检验对全部样本都拒绝. 结合所给水平计算，合理意图应是：当$\sigma=1$时使用所列阈值，而当$\sigma=10$时总是拒绝. 此处保留原文表述.}
\end{enumerate}
\end{exercise}
\begin{exercise}\label{exercise:9.50}
设$X\sim f(x\mid\theta)$，假设区间$\{\theta:a(X)\le\theta\le b(X)\}$为$\theta$的UMA置信集.
\begin{enumerate}
\item 求$1/\theta$的UMA置信集. 注意，若$a(x)<0<b(x)$，则此集合为$\{1/\theta:1/b(x)\le1/\theta\}\cup\{1/\theta:1/\theta\le1/a(x)\}$. 因而UMA置信集可能既非区间，也无界.
\item 证明：若$h$严格递增，则集合$\{h(\theta):h(a(X))\le h(\theta)\le h(b(X))\}$为$h(\theta)$的UMA置信集. 对$h$的条件能否放宽？
\end{enumerate}
\end{exercise}
\begin{exercise}\label{exercise:9.51}\index{SI0157@位置}
若$X_1,\ldots,X_n$独立同分布于位置pdf $f(x-\theta)$，证明置信集
$$C(x_1,\ldots,x_n)=\{\theta:\bar x-k_1\le\theta\le\bar x+k_2\},$$
其中$k_1,k_2$为常数，具有常数覆盖概率.（提示：$\bar X$的pdf具有$f_{\bar X}(\bar x-\theta)$的形式.）
\end{exercise}
\sourcepageanchor{404}
\begin{exercise}\label{exercise:9.52}\index{SI0178@方差}
设$X_1,\ldots,X_n$为来自$\mathrm n(\mu,\sigma^2)$总体的随机样本，其中$\mu,\sigma^2$均未知. 以下每种寻找$\sigma^2$置信区间的方法都给出形式为
$$\left\{\sigma^2:\frac{(n-1)s^2}b\le\sigma^2\le\frac{(n-1)s^2}a\right\}$$
的区间，但每种情形中$a,b$满足不同约束. 本题给出的区间由\hyperref[bib:source313]{Tate和Klett（1959）}导出，他们还列出了部分临界点.

定义$f_p(t)$为自由度为$p$的$\chi_p^2$随机变量的pdf. 为获得$1-\alpha$置信区间，$a,b$必须满足
$$\int_a^b f_{n-1}(t)\,dt=1-\alpha,$$
但还需要额外约束才能唯一确定$a,b$. 验证下列每个约束确可按所述方式导出.
\begin{enumerate}
\item 似然比区间：反演$H_0:\sigma=\sigma_0$对$H_1:\sigma\ne\sigma_0$的LRT所得的$1-\alpha$置信区间具有上述形式，且$a,b$还满足$f_{n+2}(a)=f_{n+2}(b)$.
\item 最小长度区间：对于上述形式的区间，通过极小化区间长度所得的$1-\alpha$置信区间约束$a,b$满足$f_{n+3}(a)=f_{n+3}(b)$.
\item 最短无偏区间：对于上述形式的区间，在所有无偏区间中极小化错误覆盖概率所得的$1-\alpha$置信区间约束$a,b$满足$f_{n+1}(a)=f_{n+1}(b)$. 这个区间也可通过极小化端点之比导出.
\item 等尾区间：对于上述形式的区间，要求区间上方和下方概率相等所得的$1-\alpha$置信区间约束$a,b$满足
$$\int_0^a f_{n-1}(t)\,dt=\frac\alpha2,\qquad\int_b^\infty f_{n-1}(t)\,dt=\frac\alpha2.$$
（这个区间虽然非常常见，但无论采用什么长度准则，显然都不是最优的.）
\item 对于$\alpha=.1$且$n=3$，求上述各情形中$a,b$的数值，并比较这些区间的长度.
\end{enumerate}
\end{exercise}
\begin{exercise}\label{exercise:9.53}
设$X\sim\mathrm n(\mu,\sigma^2)$，$\sigma^2$已知. 对每个$c\ge0$，定义$\mu$的区间估计量$C(x)=[x-c\sigma,x+c\sigma]$，并考虑\eqref{eq:9.3.4}中的损失.
\begin{enumerate}
\item 证明风险函数$R(\mu,C)$为
$$R(\mu,C)=b(2c\sigma)-P(-c\le Z\le c).$$
\item 利用微积分基本定理证明
$$\frac d{dc}R(\mu,C)=2b\sigma-\frac2{\sqrt{2\pi}}e^{-c^2/2},$$
从而导数在$c\ge0$上是$c$的递增函数.
\item 证明若$b\sigma>1/\sqrt{2\pi}$，则导数对所有$c\ge0$均为正，因而$R(\mu,C)$在$c=0$处最小. 即最优区间估计量是点估计量$C(x)=[x,x]$.
\item 证明若$b\sigma\le1/\sqrt{2\pi}$，则使风险最小的$c$为$c=\sqrt{-2\log(b\sigma\sqrt{2\pi})}$. 因而，若选取$b$使$c=z_{\alpha/2}$（对某个$\alpha$），则使风险最小的区间估计量就是通常的$1-\alpha$置信区间.
\end{enumerate}
\end{exercise}
\begin{exercise}\label{exercise:9.54}
设$X\sim\mathrm n(\mu,\sigma^2)$，但现在$\sigma^2$未知. 对每个$c\ge0$，定义$\mu$的区间估计量$C(x)=[x-cs,x+cs]$，其中$s^2$是与$X$独立的$\sigma^2$估计量，$\nu S^2/\sigma^2\sim\chi^2_\nu$（例如通常的样本方差）. 考虑对\eqref{eq:9.3.4}中损失的修改：
$$L((\mu,\sigma),C)=\frac b\sigma\operatorname{Length}(C)-I_C(\mu).$$
\sourcepageanchor{405}
\begin{enumerate}
\item 证明风险函数$R((\mu,\sigma),C)$为
$$R((\mu,\sigma),C)=b(2cM)-[2P(T\le c)-1],$$
其中$T\sim t_\nu$，$M=ES/\sigma$.
\item 若$b\le1/\sqrt{2\pi}$，证明使风险最小的$c$满足
$$b=\frac1{\sqrt{2\pi}}\left(\frac\nu{\nu+c^2}\right)^{(\nu+1)/2}.$$
\item 将这个问题与$\sigma^2$已知的情形协调起来. 证明当$\nu\to\infty$时，这里的解收敛到$\sigma^2$已知问题的解.（注意对损失函数所作的重新缩放.）
\end{enumerate}
\end{exercise}
\begin{exercise}\label{exercise:9.55}
集合估计的决策论方法可能很有用（\autoref{exercise:9.56}），但也可能给出令人不安的结果，表明需要审慎实施. 再次考虑$X\sim\mathrm n(\mu,\sigma^2)$、$\sigma^2$未知的情形，假设$\mu$的区间估计量为$C(x)=[x-cs,x+cs]$，其中$s^2$是与$X$独立的$\sigma^2$估计量，$\nu S^2/\sigma^2\sim\chi^2_\nu$. 当然，这就是通常的$t$区间，是经受住时间考验的伟大统计程序之一. 考虑损失
$$L((\mu,\sigma),C)=b\operatorname{Length}(C)-I_C(\mu),$$
它类似于\autoref{exercise:9.54}所用的损失，只是没有对长度缩放. 构造另一个程序$C'$：
$$C'=\begin{cases}[x-cs,x+cs],&s<K,\\\varnothing,&s\ge K,\end{cases}$$
其中$K$为正常数. 注意，$C'$做的恰恰是错误的事：当$s^2$较大、不确定性很大时，我们希望区间宽一些，而$C'$却为空！证明能找到一个$K$值，使
$$R((\mu,\sigma),C')\le R((\mu,\sigma),C),\qquad\text{对每个 }(\mu,\sigma),$$
并且对某些$(\mu,\sigma)$严格不等.
\end{exercise}
\begin{exercise}\label{exercise:9.56}
设$X\sim f(x\mid\theta)$，希望在\eqref{eq:9.3.4}的损失下，用区间估计量$C$估计$\theta$. 若$\theta$具有先验pdf $\pi(\theta)$，证明贝叶斯规则为
$$C^\pi=\{\theta:\pi(\theta\mid x)\ge b\}.$$
（提示：写成$\operatorname{Len}(C)=\int_C1\,d\theta$，并使用Neyman--Pearson引理.）
\end{exercise}
以下两题与\autoref{sec:9.5.4}有关.
\begin{exercise}\label{exercise:9.57}
设$X_1,X_2,\ldots,X_n$独立同分布于$\mathrm n(\mu,\sigma^2)$，其中$\sigma^2$已知. 已知$\mu$的一个$1-\alpha$置信区间为$\bar x\pm z_{\alpha/2}\sigma/\sqrt n$.
\begin{enumerate}
\item 证明$X_{n+1}$的一个$1-\alpha$预测区间为$\bar x\pm z_{\alpha/2}\sigma\sqrt{1+1/n}$.
\item 证明基础总体中$100p\%$的一个$1-\alpha$容忍区间为$\bar x\pm z_{\alpha/2}\sigma(1+1/\sqrt n)$.\footnote{译注：原式完全不含$p$，因此不能对任意指定的总体比例$p$成立. 一般容忍区间的宽度需要依赖$p$；此处疑有漏项或遗漏对$p$的限定，保留原式待核实.}
\item 若$\sigma^2$未知，求$X_{n+1}$的$1-\alpha$预测区间.
\end{enumerate}
（若$\sigma^2$未知，$1-\alpha$容忍区间的计算相当复杂.）
\end{exercise}
\begin{exercise}\label{exercise:9.58}
设$X_1,X_2,\ldots,X_n$为来自中位数为$m$的总体的独立同分布观测. 可按以下方式，基于次序统计量$X_{(1)}\le X_{(2)}\le\cdots\le X_{(n)}$构造无分布区间.
\begin{enumerate}
\item\label{item:9.58.a}证明单侧区间$(-\infty,x_{(n)}]$和$[x_{(1)},\infty)$均为$m$的置信区间，置信系数为$1-(1/2)^n$；区间$[x_{(1)},x_{(n)}]$的置信系数为$1-2(1/2)^n$.
\item 证明\hyperref[item:9.58.a]{(a)}中的单侧区间是系数为$n/(n+1)$的预测区间，而双侧区间是系数为$(n-1)/(n+1)$的预测区间.
\item \hyperref[item:9.58.a]{(a)}中的区间也可作为覆盖基础总体比例$p$的容忍区间. 证明，将它们视为容忍区间时，单侧区间的系数为$1-p^n$，双侧区间的系数为$1-p^n-n(1-p)p^{n-1}$. \hyperref[bib:source321]{Vardeman（1992）}将最后这一计算称为“一道漂亮的次序统计量习题”.\footnote{译注：所列精确系数通常以总体分布连续、样本无并列为前提；题干未明示这一条件. 对含原子的分布，相关覆盖率可能更高，不能直接视为逐分布成立的等式. 原题保留.}
\end{enumerate}

\end{exercise}



\sourcepageanchor{406}
\subsection{杂录}\label{sec:9.5}
\subsubsection{置信程序}\label{sec:9.5.1}\index{SI0180@置信程序（另见置信集；置信区间）}
通过定义一种称为置信程序（confidence procedure）的对象，可以在形式上联系置信集与检验（\hyperref[bib:source183]{Joshi 1969）}. 若$X\sim f(x\mid\theta)$，其中$x\in\mathcal X$且$\theta\in\Theta$，则置信程序是笛卡尔积空间$\mathcal X\times\Theta$中的集合. 对集合$\boldsymbol C\in\mathcal X\times\Theta$，定义为\footnote{译注：原文将作为集合的$\boldsymbol C$写成$\boldsymbol C\in\mathcal X\times\Theta$；按上下文应为$\boldsymbol C\subseteq\mathcal X\times\Theta$. 原记号保留.}
$$\{(x,\theta):(x,\theta)\in\boldsymbol C\}.$$

由置信程序，固定其中一个变量即可定义两个切片，或称截面. 固定$x$，定义$\theta$截面或置信集为
$$C(x)=\{\theta:(x,\theta)\in\boldsymbol C\}.$$
固定$\theta$，定义$x$截面或接受域为
$$A(\theta)=\{x:(x,\theta)\in\boldsymbol C\}.$$
虽然这种展开需要在积空间$\mathcal X\times\Theta$中工作（这也是我们这里不采用它的一个原因），但它确实提供了一种更直接的方式来观察检验与集合的关系. \autoref{fig:9.2.1}展示了正态情形中的这种对应关系.
\subsubsection{离散分布中的置信区间}\label{sec:9.5.2}
为离散分布的参数构造最优（或至少有所改进的）置信区间有很长的历史，\autoref{ex:9.2.11}已有所说明，其中考察了\hyperref[bib:source305]{Sterne（1954）}对\hyperref[bib:source79]{Clopper和Pearson（1934）}区间的修改. 当然，Sterne构造存在困难，但基本思想是正确的；\hyperref[bib:source88]{Crow（1956）}以及\hyperref[bib:source40]{Blyth和Still（1983）}修改了Sterne构造，后者产生了最短的一组精确区间. \hyperref[bib:source64]{Casella（1986）}给出了寻找一类最短二项置信区间的算法.

泊松区间研究（通常也包括其他离散分布）的历史类似. \hyperref[bib:source133]{Garwood（1936）}的构造恰好是将Clopper--Pearson论证用于二项分布，\hyperref[bib:source89]{Crow和Gardner（1959）}对这些区间作了改进. \hyperref[bib:source69]{Casella和Robert（1989）}找出了一类最短泊松区间.\footnote{译注：本段讨论泊松区间，但原文称Garwood的构造将论证“用于二项分布”；结合\autoref{ex:9.2.15}，这里可能应为“用于泊松分布”. 原文保留.}

\hyperref[bib:source39]{Blyth（1986）}给出了二项参数的非常精确的近似区间；\hyperref[bib:source196]{Leemis和Trivedi（1996）}比较了正态近似与泊松近似；\hyperref[bib:source3]{Agresti和Coull（1998）}认为，要求离散区间的覆盖率始终高于名义水平可能过于严格. \hyperref[bib:source37]{Blaker（2000）}为二项、泊松及其他离散分布构造了改进区间，并具有嵌套性质：若$\alpha<\alpha'$，则$1-\alpha$区间包含相应的$1-\alpha'$区间.
\subsubsection{Fieller定理}\label{sec:9.5.3}\index{SI0153@Fieller}\index{SI0169@均值之比}\index{SI0445@Fieller定理}
Fieller定理（\hyperref[bib:source121]{Fieller 1954）}用巧妙的论证，得到正态均值比的精确置信集.

给定来自参数为$(\mu_X,\allowbreak\mu_Y,\allowbreak\sigma_X^2,\allowbreak\sigma_Y^2,\allowbreak\rho)$的二元正态分布的随机样本$(X_1,Y_1),\allowbreak\ldots,\allowbreak(X_n,Y_n)$，可按如下方式构造$\theta=\mu_Y/\mu_X$的置信集. 对$i=1,\ldots,n$，定义$Z_{\theta i}=Y_i-\theta X_i$以及$\bar Z_\theta=\bar Y-\theta\bar X$. 可证明$\bar Z_\theta$为正态随机变量，均值为$0$，方差为
$$V_\theta=\frac1n(\sigma_Y^2-2\theta\rho\sigma_Y\sigma_X+\theta^2\sigma_X^2).$$
可用$\hat V_\theta$估计$V_\theta$，其中
$$\hat V_\theta=\frac1{n(n-1)}\sum_{i=1}^n(Z_{\theta i}-\bar Z_\theta)^2=\frac1{n-1}(S_Y^2-2\theta S_{YX}+\theta^2S_X^2),$$
\sourcepageanchor{407}
这里
$$S_Y^2=\frac1n\sum_{i=1}^n(Y_i-\bar Y)^2,\quad S_X^2=\frac1n\sum_{i=1}^n(X_i-\bar X)^2,\quad S_{YX}=\frac1n\sum_{i=1}^n(Y_i-\bar Y)(X_i-\bar X).$$
此外，还可证明$E\hat V_\theta=V_\theta$，$\hat V_\theta$与$\bar Z_\theta$独立，且$(n-1)\hat V_\theta/V_\theta\sim\chi^2_{n-1}$. 因此，$\bar Z_\theta/\sqrt{\hat V_\theta}\sim t_{n-1}$，集合
$$\left\{\theta:\frac{\bar z_\theta^2}{\hat v_\theta}\le t^2_{n-1,\alpha/2}\right\}$$
定义了均值比$\theta$的$1-\alpha$置信集. 此集合定义了关于$\theta$的抛物线，其根给出置信集的端点. 用原变量写出该集合，得到
$$\left\{\theta:\left(\bar x^2-\frac{t^2_{n-1,\alpha/2}}{n-1}S_x^2\right)\theta^2-2\left(\bar x\bar y-\frac{t^2_{n-1,\alpha/2}}{n-1}S_{yx}\right)\theta+\left(\bar y^2-\frac{t^2_{n-1,\alpha/2}}{n-1}S_y^2\right)\le0\right\}.$$
此集合有一个有趣特点：取决于抛物线的根，它可以是一个区间、一个区间的补集，或整条实数轴（参见\autoref{exercise:9.15}）. 此外，要保持$1-\alpha$的置信程度，该区间必须以正概率为无穷区间. 基于自助法的替代方法见\hyperref[bib:source175]{Hwang（1995）}；另一种置信方法见Tsao和Hwang（\hyperref[bib:source316]{1998}，\hyperref[bib:source317]{1999}）.
\subsubsection{其他区间呢？}\label{sec:9.5.4}
\hyperref[bib:source321]{Vardeman（1992）}提出本条杂录标题中的问题，主张主流统计学应花更多时间研究双侧置信区间以外的区间. 特别地，他列出：（a）单侧区间；（b）无分布区间；（c）预测区间；（d）容忍区间.

我们已经见过单侧区间；无分布区间（distribution-free interval）是指其概率保证对基础cdf只需很少假设（或无需假设）就成立的区间（参见\autoref{exercise:9.58}）. 另外两种区间的定义，连同通常的置信区间，提供了一个推断层级，后者逐一比前者更严格.

若$X_1,X_2,\ldots,X_n$独立同分布于cdf为$F(x\mid\theta)$的总体，且$C(\boldsymbol x)=[l(\boldsymbol x),u(\boldsymbol x)]$为区间，则对给定的$1-\alpha$，它是
\begin{enumerate}[label=\roman*.]
\item 置信区间，若$P_\theta[l(\boldsymbol X)\le\theta\le u(\boldsymbol X)]\ge1-\alpha$；
\item 预测区间（prediction interval），若$P_\theta[l(\boldsymbol X)\le X_{n+1}\le u(\boldsymbol X)]\ge1-\alpha$；
\item 容忍区间（tolerance interval），若对于指定的$p$，$P_\theta[F(u(\boldsymbol X)\mid\theta)-F(l(\boldsymbol X)\mid\theta)\ge p]\ge1-\alpha$.
\end{enumerate}
所以，置信区间覆盖一个均值，预测区间覆盖一个新的随机变量，而容忍区间覆盖总体的某个比例. 因而，它们分别给出不同推断，应根据具体问题选择恰当的一种.





\section{渐近评价}\label{sec:10}
\chapterdoi{10}
\sourcepageanchor{408}
\begin{quote}
我知道，亲爱的华生，你和我一样，喜爱一切离奇的、超出日常生活惯例与单调常规的事物.
\begin{flushright}夏洛克·福尔摩斯\\《红发会》\end{flushright}
\end{quote}

迄今为止，我们考虑的所有准则都是有限样本准则. 相比之下，也可以考虑渐近性质（asymptotic properties），即描述样本量趋于无穷时某一程序行为的性质. 本节考察其中一些性质，分别讨论点估计、假设检验与区间估计. 我们将特别强调最大似然程序的渐近性质.

渐近评价的力量在于，令样本量趋于无穷会使计算简化. 有限样本情形下无法完成的评价，变成了常规计算. 这种简化也让我们得以考察其他一些技术（如自助法和$M$估计），它们通常只能进行渐近评价.

让样本量无限增大（有时称为“渐近乌托邦”，asymptopia），不应被嘲笑为仅仅是一种异想天开的练习. 相反，渐近理论揭示一个程序最根本的性质，为我们提供极有力且普遍的评价工具.

\subsection{点估计}\label{sec:10.1}
\subsubsection{相合性}\label{sec:10.1.1}

相合性（consistency）似乎是一项十分基本的性质，它要求样本量趋于无穷时，估计量收敛到“正确”值. 这项性质如此基本，以至于不相合估计量的价值应受到质疑（或至少应受到认真研究）.

相合性（以及所有渐近性质）涉及的是估计量序列，而非单个估计量，尽管人们常说“相合估计量”. 若按照分布$f(x\mid\theta)$观测$X_1,X_2,\ldots$，只须对每个样本量$n$执行同一种估计程序，就能构造估计量序列$W_n=W_n(X_1,\ldots,X_n)$. 例如，$\bar X_1=X_1$，$\bar X_2=(X_1+X_2)/2$，$\bar X_3=(X_1+X_2+X_3)/3$，等等. 现在可以定义相合序列.

\begin{definition}\label{def:10.1.1}\index{SI0770@相合}
若对每个$\varepsilon>0$及每个$\theta\in\Theta$，估计量序列$W_n=W_n(X_1,\ldots,X_n)$满足
\begin{equation}\label{eq:10.1.1}
\lim_{n\to\infty}P_\theta(|W_n-\theta|<\varepsilon)=1,
\end{equation}
则称它为参数$\theta$的相合估计量序列.
\end{definition}

非正式地说，\eqref{eq:10.1.1}意味着，当样本量趋于无穷（样本信息越来越好）时，估计量会以很高的概率任意接近参数，这显然是理想的性质. 或者反过来说，相合估计量序列偏离真实参数的概率很小.
\sourcepageanchor{409}
与\eqref{eq:10.1.1}等价的表述如下：对每个$\varepsilon>0$及每个$\theta\in\Theta$，相合序列$W_n$满足
\begin{equation}\label{eq:10.1.2}
\lim_{n\to\infty}P_\theta(|W_n-\theta|\geq\varepsilon)=0.
\end{equation}
应将\autoref{def:10.1.1}与定义依概率收敛的\autoref{def:5.5.1}比较. \autoref{def:10.1.1}说，相合估计量序列依概率收敛到它所估计的参数$\theta$. \autoref{def:5.5.1}处理的是一个概率结构下的一列随机变量，而\autoref{def:10.1.1}处理的是由$\theta$标记的一整族概率结构. 对不同的$\theta$值，与序列$W_n$相关的概率结构不同. 定义要求，对每个$\theta$值，该概率结构均使序列依概率收敛到真实的$\theta$. 这就是概率定义与统计定义之间通常的区别：概率定义处理一个概率结构，统计定义处理一整族.

\begin{example}[$\bar X$的相合性]\label{ex:10.1.2}\index{SI0772@正态}
设$X_1,X_2,\ldots$独立同分布于$\mathrm n(\theta,1)$，考虑序列
$$\bar X_n=\frac1n\sum_{i=1}^nX_i.$$
回顾$\bar X_n\sim\mathrm n(\theta,1/n)$，所以
\begin{align*}
P_\theta(|\bar X_n-\theta|<\varepsilon)
&=\int_{\theta-\varepsilon}^{\theta+\varepsilon}\left(\frac n{2\pi}\right)^{1/2}e^{-(n/2)(\bar x_n-\theta)^2}\,d\bar x_n &&\text{（定义）}\\
&=\int_{-\varepsilon}^{\varepsilon}\left(\frac n{2\pi}\right)^{1/2}e^{-(n/2)y^2}\,dy &&\text{（代入$y=\bar x_n-\theta$）}\\
&=\int_{-\varepsilon\sqrt n}^{\varepsilon\sqrt n}\left(\frac1{2\pi}\right)^{1/2}e^{-t^2/2}\,dt &&\text{（代入$t=y\sqrt n$）}\\
&=P(-\varepsilon\sqrt n<Z<\varepsilon\sqrt n) &&(Z\sim\mathrm n(0,1))\\
&\longrightarrow1\quad\text{当$n\to\infty$}.
\end{align*}
因此，$\bar X_n$是$\theta$的相合估计量序列.
\end{example}

一般而言，验证相合性并不需要上面那样的详细计算.\index{SI0565@切比雪夫}\semanticindexanchor{565:409} 回顾切比雪夫不等式，对估计量$W_n$有
$$P_\theta(|W_n-\theta|\geq\varepsilon)\leq\frac{E_\theta[(W_n-\theta)^2]}{\varepsilon^2},$$
因此，若对每个$\theta\in\Theta$，
$$\lim_{n\to\infty}E_\theta[(W_n-\theta)^2]=0,$$
则该估计量序列相合. 此外，由\eqref{eq:7.3.1}，
\begin{equation}\label{eq:10.1.3}
E_\theta[(W_n-\theta)^2]=\operatorname{Var}_\theta W_n+[\operatorname{Bias}_\theta W_n]^2.
\end{equation}
综合这些事实，可以陈述如下定理.

\begin{theorem}\label{thm:10.1.3}
若$W_n$是参数$\theta$的估计量序列，且对每个$\theta\in\Theta$满足
\begin{enumerate}[label=\roman*.]
\item $\lim_{n\to\infty}\operatorname{Var}_\theta W_n=0$，
\item $\lim_{n\to\infty}\operatorname{Bias}_\theta W_n=0$，
\end{enumerate}
则$W_n$是$\theta$的相合估计量序列.
\end{theorem}

\sourcepageanchor{410}
\begin{example}[续\autoref{ex:10.1.2}]\label{ex:10.1.4}
因为
$$E_\theta\bar X_n=\theta\quad\text{且}\quad\operatorname{Var}_\theta\bar X_n=\frac1n,$$
所以满足\autoref{thm:10.1.3}的条件，序列$\bar X_n$相合. 此外，由\autoref{thm:5.2.6}，若从任意均值为$\theta$的总体中独立同分布抽样，只要总体方差有限，$\bar X_n$就相合地估计$\theta$.
\end{example}

本节开头说过，不相合估计量序列的价值应受到质疑. 这一说法的部分依据在于，相合序列如此之多，下面的定理说明了这一点. 证明留作\autoref{exercise:10.2}.

\begin{theorem}\label{thm:10.1.5}
设$W_n$是参数$\theta$的相合估计量序列，$a_1,a_2,\ldots$与$b_1,b_2,\ldots$是满足以下条件的常数序列：
\begin{enumerate}[label=\roman*.]
\item $\lim_{n\to\infty}a_n=1$，且
\item $\lim_{n\to\infty}b_n=0$.
\end{enumerate}
则序列$U_n=a_nW_n+b_n$是$\theta$的相合估计量序列.
\end{theorem}

本节最后概述一个关于最大似然估计量相合性的更一般的结果. 它表明，MLE是其参数的相合估计量；这是我们首次看到，寻找估计量的方法本身能够保证一种最优性性质.

要使MLE相合，基础密度（似然函数）必须满足某些“正则条件”（regularity conditions）. 这里不展开，详见\hyperref[sec:10.6.2]{杂录10.6.2}.

\begin{theorem}[MLE的相合性]\label{thm:10.1.6}\index{SI0674@相合性}\index{SI0771@最大似然估计量}
设$X_1,\ldots,X_n$独立同分布于$f(x\mid\theta)$，$L(\theta\mid\boldsymbol x)=\prod_{i=1}^nf(x_i\mid\theta)$为似然函数，$\hat\theta$为$\theta$的MLE，$\tau(\theta)$为$\theta$的连续函数. 在\hyperref[sec:10.6.2]{杂录10.6.2}对$f(x\mid\theta)$、从而对$L(\theta\mid\boldsymbol x)$施加的正则条件下，对每个$\varepsilon>0$及每个$\theta\in\Theta$，
$$\lim_{n\to\infty}P_\theta(|\tau(\hat\theta)-\tau(\theta)|\geq\varepsilon)=0.$$
即$\tau(\hat\theta)$是$\tau(\theta)$的相合估计量.
\end{theorem}
\begin{proof}
证明先说明，对每个$\theta\in\Theta$，$\frac1n\log L(\hat\theta\mid\boldsymbol x)$几乎处处收敛到$E_\theta(\log f(X\mid\theta))$. 在$f(x\mid\theta)$满足某些条件时，这意味着$\hat\theta$依概率收敛到$\theta$，从而$\tau(\hat\theta)$依概率收敛到$\tau(\theta)$. 详见\hyperref[bib:source311]{Stuart、Ord和Arnold（1999}，第18章）.
\end{proof}

\subsubsection{效率}\label{sec:10.1.2}\index{SI0394@效率（另见点估计）}\index{SI0779@效率}
相合性关心估计量的渐近准确性：它是否收敛到所估计的参数.\index{SI0071@渐近方差（另见方差）}\semanticindexanchor{71:410}\index{SI1063@渐近}\semanticindexanchor{1063:410} 本节考察相关的效率（efficiency）性质，它关心估计量的渐近方差.

计算渐近方差时，人们或许很容易采用如下做法. 给定基于样本量$n$的估计量$T_n$，先计算有限样本方差$\operatorname{Var}T_n$，再求$\lim_{n\to\infty}k_n\operatorname{Var}T_n$，其中$k_n$是某个规范化常数. （注意，在许多情形中，$n\to\infty$时$\operatorname{Var}T_n\to0$，所以需要因子$k_n$使它趋于一个极限.）

\begin{definition}\label{def:10.1.7}\index{SI1065@极限}
对估计量$T_n$，若$\lim_{n\to\infty}k_n\operatorname{Var}T_n=\tau^2<\infty$，其中$\{k_n\}$是常数序列，则称$\tau^2$为极限方差（limiting variance），或方差的极限（limit of the variances）.
\end{definition}

\begin{example}[极限方差]\label{ex:10.1.8}
对于$n$个独立同分布正态观测的均值$\bar X_n$，其中$EX=\mu$、$\operatorname{Var}X=\sigma^2$，若取$T_n=\bar X_n$，则$\lim\sqrt n\operatorname{Var}\bar X_n=\sigma^2$是$T_n$的极限方差.\footnote{译注：原文系数为$\sqrt n$；由$\operatorname{Var}\bar X_n=\sigma^2/n$可知，欲得到$\sigma^2$，此处应为$n$，而非$\sqrt n$. 原式保留.}

\sourcepageanchor{411}
但若改为用$1/\bar X_n$估计$1/\mu$，例如，就会出现麻烦. 此时取$T_n=1/\bar X_n$，会发现方差$\operatorname{Var}(T_n)=\infty$，因此方差的极限为无穷. 不过，回顾\autoref{ex:5.5.23}，其中说$1/\bar X_n$的“近似”均值和方差为
$$E\left(\frac1{\bar X_n}\right)\approx\frac1\mu,\qquad
\operatorname{Var}\left(\frac1{\bar X_n}\right)\approx\left(\frac1\mu\right)^4\operatorname{Var}\bar X_n,$$
于是，按第二种计算，方差为$\operatorname{Var}(T_n)\approx\sigma^2/(n\mu^4)<\infty$.
\end{example}

这个例子指出了把方差的极限作为大样本度量的问题. 当然，$1/\bar X$的精确有限样本方差为$\infty$. 然而，若$\mu\ne0$，$1/\bar X$变得非常大的区域，其概率趋于0. 所以\autoref{ex:10.1.8}中的第二种近似更符合实际（也更有用）. 我们采用的正是这第二种计算大样本方差的方法.

\begin{definition}\label{def:10.1.9}
对估计量$T_n$，假设$k_n(T_n-\mu)\to\mathrm n(0,\sigma^2)$. 参数$\sigma^2$称为$T_n$的渐近方差（asymptotic variance），或极限分布的方差（variance of the limit distribution）.
\end{definition}

计算样本均值及其他平均量的方差时，极限方差与渐近方差通常具有相同的值. 但在较复杂的情形中，极限方差有时会失效. 另一个有趣的事实是，渐近方差总是小于极限方差（\hyperref[bib:source200]{Lehmann和Casella，第6.1节}）. 下面给出说明.

\begin{example}[大样本混合方差]\label{ex:10.1.10}\index{SI1065@极限}
分层模型
$$\begin{aligned}
Y_n\mid W_n=w_n&\sim\mathrm n\bigl(0,w_n+(1-w_n)\sigma_n^2\bigr),\\
W_n&\sim\operatorname{Bernoulli}(p_n)
\end{aligned}$$
能够使渐近方差与极限方差出现很大差异. （有时也把它描述为混合模型，即以概率$p_n$观测到$Y_n\sim\mathrm n(0,1)$，以概率$1-p_n$观测到$Y_n\sim\mathrm n(0,\sigma_n^2)$.）

首先，由\autoref{thm:4.4.7}，
$$\operatorname{Var}(Y_n)=p_n+(1-p_n)\sigma_n^2.$$
由此可知，仅当$\lim_n(1-p_n)\sigma_n^2<\infty$时，$Y_n$的极限方差才有限.

另一方面，可利用
$$P(Y_n<a)=p_nP(Z<a)+(1-p_n)P(Z<a/\sigma_n)$$
直接计算$Y_n$的渐近分布. 现令$p_n\to1$且$\sigma_n\to\infty$，使$(1-p_n)\sigma_n^2\to\infty$. 则$P(Y_n<a)\to P(Z<a)$，即$Y_n\to\mathrm n(0,1)$，且
$$\begin{aligned}
\text{极限方差}&=\lim_n p_n+(1-p_n)\sigma_n^2=\infty,\\
\text{渐近方差}&=1.
\end{aligned}$$
更多细节见\autoref{exercise:10.6}.
\end{example}

仿照Cramér--Rao下界（\autoref{thm:7.3.9}）的思路，可以给出一个最优渐近方差.

\sourcepageanchor{412}
\begin{definition}\label{def:10.1.11}\index{SI0065@渐近效率}\index{SI0779@效率}
若估计量序列$W_n$满足$\sqrt n[W_n-\tau(\theta)]$依分布收敛到$\mathrm n[0,v(\theta)]$，且
$$v(\theta)=\frac{[\tau'(\theta)]^2}{E_\theta\left(\left(\frac\partial{\partial\theta}\log f(X\mid\theta)\right)^2\right)},$$
则称它对于参数$\tau(\theta)$渐近有效（asymptotically efficient）. 即$W_n$的渐近方差达到Cramér--Rao下界.
\end{definition}

回顾\autoref{thm:10.1.6}，它表明在一般条件下MLE相合. 在略强一些的正则条件下，关于渐近有效性也有同类定理成立，因此一般可以认为MLE既相合又渐近有效. 正则条件的细节仍见\hyperref[sec:10.6.2]{杂录10.6.2}.

\begin{theorem}[MLE的渐近有效性]\label{thm:10.1.12}\index{SI0668@渐近效率}\index{SI0796@渐近效率}
设$X_1,\ldots,X_n$独立同分布于$f(x\mid\theta)$，$\hat\theta$为$\theta$的MLE，$\tau(\theta)$为$\theta$的连续函数. 在\hyperref[sec:10.6.2]{杂录10.6.2}对$f(x\mid\theta)$、从而对$L(\theta\mid\boldsymbol x)$施加的正则条件下，
$$\sqrt n[\tau(\hat\theta)-\tau(\theta)]\longrightarrow\mathrm n[0,v(\theta)],$$
其中$v(\theta)$为Cramér--Rao下界. 即$\tau(\hat\theta)$是$\tau(\theta)$的相合且渐近有效的估计量.
\end{theorem}
\begin{proof}
这个定理的证明有趣之处在于使用泰勒级数，并利用MLE被定义为似然函数的零点这一事实.\footnote{译注：原文称“似然函数的零点”；按接下来的计算，应为对数似然函数的一阶导数的零点，即似然方程的根. 原文保留.} 我们概述$\hat\theta$渐近有效的证明；推广到$\tau(\hat\theta)$留作\autoref{exercise:10.7}.

回顾$l(\theta\mid\boldsymbol x)=\sum\log f(x_i\mid\theta)$为对数似然函数，将关于$\theta$的导数记为$l',l'',\ldots$. 现在在真值$\theta_0$处展开对数似然函数的一阶导数：
\begin{equation}\label{eq:10.1.4}
l'(\theta\mid\boldsymbol x)=l'(\theta_0\mid\boldsymbol x)+(\theta-\theta_0)l''(\theta_0\mid\boldsymbol x)+\cdots,
\end{equation}
我们将忽略高阶项（正则条件保证这种做法成立）.

现用MLE $\hat\theta$代替$\theta$，注意到\eqref{eq:10.1.4}左侧为零. 整理并乘以$\sqrt n$，得到
\begin{equation}\label{eq:10.1.5}
\sqrt n(\hat\theta-\theta_0)=\sqrt n\frac{-l'(\theta_0\mid\boldsymbol x)}{l''(\theta_0\mid\boldsymbol x)}
=\frac{\frac1{\sqrt n}l'(\theta_0\mid\boldsymbol x)}{-\frac1n l''(\theta_0\mid\boldsymbol x)}.
\end{equation}
若令$I(\theta_0)=E[l'(\theta_0\mid\boldsymbol X)]^2=1/v(\theta)$表示信息量，应用中心极限定理与大数定律可知（细节见\autoref{exercise:10.8}）\footnote{译注：这里$l$已定义为全样本对数似然，故其得分平方的期望为$n$倍的单个观测信息量. 要与下式及$1/v(\theta_0)$一致，定义中应使用单个观测的对数密度导数，或将$E\{[l'(\theta_0\mid\boldsymbol X)]^2\}$除以$n$. 原文还将$v$的自变量写成$\theta$，此处照录.}
\begin{equation}\label{eq:10.1.6}
\begin{aligned}
\frac1{\sqrt n}l'(\theta_0\mid\boldsymbol X)&\longrightarrow\mathrm n[0,I(\theta_0)]\quad\text{依分布},\\
\text{且}\qquad-\frac1n l''(\theta_0\mid\boldsymbol X)&\longrightarrow I(\theta_0)\quad\text{依概率}.
\end{aligned}
\end{equation}
因此，若令$W\sim\mathrm n[0,I(\theta_0)]$，则$\sqrt n(\hat\theta-\theta_0)$依分布收敛到$W/I(\theta_0)\sim\mathrm n[0,1/I(\theta_0)]$，定理得证.
\end{proof}

\begin{example}[渐近正态性与相合性]\label{ex:10.1.13}
上述定理表明，MLE通常有效且相合. 要指出的是，这种说法有些重复，因为效率仅在估计量渐近正态时才有定义，而正如下面所说明的，渐近正态性蕴含相合性. 假设
$$\sqrt n\frac{W_n-\mu}{\sigma}\longrightarrow Z\quad\text{依分布},$$
\sourcepageanchor{413}
其中$Z\sim\mathrm n(0,1)$. 应用Slutsky定理（\autoref{thm:5.5.17}），得到
$$W_n-\mu=\left(\frac\sigma{\sqrt n}\right)\left(\sqrt n\frac{W_n-\mu}{\sigma}\right)
\longrightarrow\lim_{n\to\infty}\left(\frac\sigma{\sqrt n}\right)Z=0,$$
因此$W_n-\mu$依分布收敛到0. 由\autoref{thm:5.5.13}，依分布收敛到一个点等价于依概率收敛，所以$W_n$是$\mu$的相合估计量.
\end{example}

\subsubsection{计算与比较}\label{sec:10.1.3}
前面各节导出的渐近公式，可以提供适合大样本使用的近似方差. 我们仍需关注正则条件（\hyperref[sec:10.6.2]{杂录10.6.2}），但这些条件相当一般，在通常情形下几乎总能满足. 不过，有一个条件值得特别强调；正如在\autoref{ex:7.3.13}中所见，违反它会造成复杂情况. 要使以下近似成立，pdf或pmf、从而似然函数的支撑集必须与参数无关.

假设MLE渐近有效，则\autoref{thm:10.1.6}中的渐近方差就是\autoref{thm:5.5.24}的Delta方法方差（去掉$1/n$项）.\footnote{译注：这里及本小节后文多次把渐近有效性归于\autoref{thm:10.1.6}；该定理只陈述相合性，渐近有效性在\autoref{thm:10.1.12}中陈述. 保留原书所引编号，后续同类引用不再逐一注释.} 因此，可以用Cramér--Rao下界近似MLE的真实方差. 假设$X_1,\ldots,X_n$独立同分布于$f(x\mid\theta)$，$\hat\theta$是$\theta$的MLE，而$I_n(\theta)=E_\theta\left(-\frac\partial{\partial\theta}\log L(\theta\mid\boldsymbol X)\right)^2$是样本信息量. 利用Delta方法与MLE的渐近有效性，$h(\hat\theta)$的方差可近似为
\begin{equation}\label{eq:10.1.7}\index{SI0582@信息量}
\begin{aligned}
\operatorname{Var}(h(\hat\theta)\mid\theta)
&\approx\frac{[h'(\theta)]^2}{I_n(\theta)}\\
&\approx\frac{[h'(\theta)]^2}{E_\theta\left(-\frac{\partial^2}{\partial\theta^2}\log L(\theta\mid\boldsymbol X)\right)}
&&\text{（使用\autoref{lem:7.3.11}的恒等式）}\\
&\approx\frac{[h'(\theta)]^2|_{\theta=\hat\theta}}{-\frac{\partial^2}{\partial\theta^2}\log L(\theta\mid\boldsymbol X)|_{\theta=\hat\theta}}
&&\text{（分母为观测信息量$\hat I_n(\hat\theta)$）}.
\end{aligned}
\end{equation}
此外，已有结果表明（\hyperref[bib:source110]{Efron和Hinkley，1978）}，使用观测信息量（observed information number）优于使用期望信息量（expected information number），后者正是Cramér--Rao下界中的信息量.

注意，方差估计分为两步，只是在\eqref{eq:10.1.7}中这一点不太明显. 为估计$\operatorname{Var}_\theta h(\hat\theta)$，先近似$\operatorname{Var}_\theta h(\hat\theta)$，再估计所得近似，通常用$\hat\theta$代入$\theta$. 所得估计可记为$\operatorname{Var}_{\hat\theta}h(\hat\theta)$或$\widehat{\operatorname{Var}}_\theta h(\hat\theta)$.

由\autoref{thm:10.1.6}，$-\frac1n\frac{\partial^2}{\partial\theta^2}\log L(\theta\mid\boldsymbol X)|_{\theta=\hat\theta}$是$I(\theta)$的相合估计量，故$\widehat{\operatorname{Var}}_\theta h(\hat\theta)$是$\operatorname{Var}_\theta h(\hat\theta)$的相合估计量.

\begin{example}[近似二项方差]\label{ex:10.1.14}\index{SI0670@近似方差}
在\autoref{ex:7.2.7}中看到，$\hat p=\sum X_i/n$是$p$的MLE，其中$X_1,\ldots,X_n$为来自Bernoulli$(p)$总体的随机样本. 直接计算也知道
$$\operatorname{Var}_p\hat p=\frac{p(1-p)}n,$$
而$\operatorname{Var}_p\hat p$的一个合理估计为
\begin{equation}\label{eq:10.1.8}
\widehat{\operatorname{Var}}_p\hat p=\frac{\hat p(1-\hat p)}n.
\end{equation}
\sourcepageanchor{414}
若应用\eqref{eq:10.1.7}的近似，取$h(p)=p$，则得到$\operatorname{Var}_p\hat p$的估计
$$\widehat{\operatorname{Var}}_p\hat p\approx\frac1{-\frac{\partial^2}{\partial p^2}\log L(p\mid\boldsymbol x)|_{p=\hat p}}.$$
回顾
$$\log L(p\mid\boldsymbol x)=n\hat p\log(p)+n(1-\hat p)\log(1-p),$$
因此\footnote{译注：下面二阶导数的第二项原文分子为$n(1-p)$；对上式求导应为$n(1-\hat p)$. 随后的$p=\hat p$代入式是相符的. 此处保留原式.}
$$\frac{\partial^2}{\partial p^2}\log L(p\mid\boldsymbol x)=-\frac{n\hat p}{p^2}-\frac{n(1-p)}{(1-p)^2}.$$
在$p=\hat p$处计算二阶导数，得到
$$\left.\frac{\partial^2}{\partial p^2}\log L(p\mid\boldsymbol x)\right|_{p=\hat p}
=-\frac{n\hat p}{\hat p^2}-\frac{n(1-\hat p)}{(1-\hat p)^2}=-\frac n{\hat p(1-\hat p)},$$
给出的方差近似与\eqref{eq:10.1.8}完全相同. 现在可以应用\autoref{thm:10.1.6}断言$\hat p$渐近有效，特别是
$$\sqrt n(\hat p-p)\longrightarrow\mathrm n[0,p(1-p)]$$
依分布成立. 若再用\autoref{thm:5.5.17}（Slutsky定理），则有
$$\sqrt n\frac{\hat p-p}{\sqrt{\hat p(1-\hat p)}}\longrightarrow\mathrm n[0,1].$$

估计$\hat p$的方差其实并不难，无须动用这些近似的全部工具. 不过，若转向稍复杂一些的函数，事情就可能棘手起来. 回顾\hyperref[ex:5.5.22]{习题5.5.22}\footnote{译注：原文称“习题5.5.22”，实际对应\autoref{ex:5.5.22}. 原称呼保留，链接指向对应例.}，其中用Delta方法近似$\hat p/(1-\hat p)$的方差，该量估计优势比$p/(1-p)$.\index{SI0728@优势}\semanticindexanchor{728:414} 现在可知，这个估计量实际上就是优势比的MLE，而且其方差可估计为
$$\begin{aligned}
\widehat{\operatorname{Var}}\left(\frac{\hat p}{1-\hat p}\right)
&\approx\frac{\left[\frac\partial{\partial p}\left(\frac p{1-p}\right)\right]^2\big|_{p=\hat p}}{-\frac{\partial^2}{\partial p^2}\log L(p\mid\boldsymbol x)|_{p=\hat p}}\\
&=\frac{\left[\frac{(1-p)+p}{(1-p)^2}\right]^2\big|_{p=\hat p}}{\frac n{p(1-p)}\big|_{p=\hat p}}
=\frac{\hat p}{n(1-\hat p)^3}.
\end{aligned}$$
此外，也知道该估计量渐近有效.
\end{example}

MLE方差近似在许多情形下表现很好，但并非万无一失. 特别是当函数$h(\hat\theta)$不单调时，必须谨慎. 这种情形下，导数$h'$会变号，可能导致低估方差的近似. 应认识到，这个近似基于Cramér--Rao下界，因此很可能本来就是低估；非单调函数会使问题更加严重.

\begin{example}[续\autoref{ex:10.1.14}]\label{ex:10.1.15}
现在假设要估计伯努利分布的方差$p(1-p)$. 该方差的MLE为$\hat p(1-\hat p)$，应用\eqref{eq:10.1.7}可估计这个估计量的方差：
$$\begin{aligned}
\widehat{\operatorname{Var}}\bigl(\hat p(1-\hat p)\bigr)
&\approx\frac{\left[\frac\partial{\partial p}(p(1-p))\right]^2\big|_{p=\hat p}}{-\frac{\partial^2}{\partial p^2}\log L(p\mid\boldsymbol x)|_{p=\hat p}}\\
&=\frac{(1-2p)^2|_{p=\hat p}}{\frac n{p(1-p)}\big|_{p=\hat p}}
=\frac{\hat p(1-\hat p)(1-2\hat p)^2}n,
\end{aligned}$$
当$\hat p=1/2$时它可能为零，显然低估了$\hat p(1-\hat p)$的方差. 函数$p(1-p)$不单调是造成这个问题的一个原因.

\sourcepageanchor{415}
利用\autoref{thm:10.1.6}，只要$p\ne1/2$，就能断言该估计量渐近有效. 若$p=1/2$，则需采用\autoref{thm:5.5.26}给出的二阶近似（见\autoref{exercise:10.10}）.
\end{example}

渐近有效性为我们希望达到的渐近方差提供一个基准\index{SI0780@有效}\semanticindexanchor{780:415}（不过见\hyperref[sec:10.6.1]{杂录10.6.1}）. 也可以通过渐近相对效率（asymptotic relative efficiency）的概念，利用渐近方差来比较估计量.

\begin{definition}\label{def:10.1.16}\index{SI0066@渐近相对效率}
若两个估计量$W_n$与$V_n$满足
$$\begin{aligned}
\sqrt n[W_n-\mu]&\longrightarrow\mathrm n[0,\sigma_W^2],\\
\sqrt n[V_n-\mu]&\longrightarrow\mathrm n[0,\sigma_V^2]
\end{aligned}$$
依分布成立，则$V_n$相对于$W_n$的渐近相对效率为
$$\operatorname{ARE}(V_n,W_n)=\frac{\sigma_W^2}{\sigma_V^2}.$$
\end{definition}

\begin{example}[泊松估计量的ARE]\label{ex:10.1.17}
假设$X_1,X_2,\ldots,X_n$独立同分布于Poisson$(\lambda)$，我们希望估计取零的概率. 例如，有时将给定时间段内进入银行的顾客人数建模为泊松随机变量，取零的概率就是一个时间段内无人进入银行的概率. 若$X\sim\operatorname{Poisson}(\lambda)$，则$P(X=0)=e^{-\lambda}$；一个自然（但有些朴素）的估计量来自定义$Y_i=I(X_i=0)$并采用
$$\hat\tau=\frac1n\sum_{i=1}^nY_i.$$
$Y_i$服从Bernoulli$(e^{-\lambda})$，所以
$$E(\hat\tau)=e^{-\lambda}\quad\text{且}\quad\operatorname{Var}(\hat\tau)=\frac{e^{-\lambda}(1-e^{-\lambda})}n.$$
另一种选择是$e^{-\lambda}$的MLE $e^{-\hat\lambda}$，其中$\hat\lambda=\sum_iX_i/n$是$\lambda$的MLE. 利用Delta方法近似，有
$$E(e^{-\hat\lambda})\approx e^{-\lambda}\quad\text{且}\quad\operatorname{Var}(e^{-\hat\lambda})\approx\frac{\lambda e^{-2\lambda}}n.$$
因为
$$\begin{aligned}
\sqrt n(\hat\tau-e^{-\lambda})&\longrightarrow\mathrm n[0,e^{-\lambda}(1-e^{-\lambda})],\\
\sqrt n(e^{-\hat\lambda}-e^{-\lambda})&\longrightarrow\mathrm n[0,\lambda e^{-2\lambda}]
\end{aligned}$$
依分布成立，所以$\hat\tau$相对于MLE $e^{-\hat\lambda}$的ARE为
$$\operatorname{ARE}(\hat\tau,e^{-\hat\lambda})=\frac{\lambda e^{-2\lambda}}{e^{-\lambda}(1-e^{-\lambda})}.$$
考察这个函数可知，它严格递减，在$\lambda=0$达到最大值1（这是$\hat\tau$所能达到的最佳情形），并随着$\lambda\to\infty$迅速下降、渐近趋于0，在$\lambda=4$时已低于10\%. （见\autoref{exercise:10.9}.）
\end{example}

\sourcepageanchor{416}
因为MLE通常渐近有效，另一个估计量无法指望在渐近方差上胜过它. 然而，其他估计量可能具有其他理想性质（计算方便、对基础假设稳健），使它们值得采用. 此时，MLE的效率就成为衡量使用替代估计量所付代价的重要标准.

我们再看最后一个例子，将计算方便与最优方差作比较. 下一节讨论稳健性问题.

\begin{example}[估计伽马均值]\label{ex:10.1.18}
尽管可能难以置信，估计伽马分布的均值并不是一件容易的事. 回顾伽马pdf $f(x\mid\alpha,\beta)$为
$$f(x\mid\alpha,\beta)=\frac1{\Gamma(\alpha)\beta^\alpha}x^{\alpha-1}e^{-x/\beta}.$$
其均值为$\alpha\beta$；要计算最大似然估计量，就必须处理伽马函数的导数（称为双伽马函数，digamma function），这从来不是令人愉快的事. 相比之下，矩方法给出容易计算的估计.

具体地，假设有来自上述伽马密度的随机样本$X_1,X_2,\ldots,X_n$，但重新参数化，使均值$\mu=\alpha\beta$显式出现. 此时
$$f(x\mid\mu,\beta)=\frac1{\Gamma(\mu/\beta)\beta^{\mu/\beta}}x^{\mu/\beta-1}e^{-x/\beta},$$
而$\mu$的矩估计量为$\bar X$，方差为$\beta\mu/n$.

为计算MLE，采用对数似然
$$l(\mu,\beta\mid\boldsymbol x)=\sum_{i=1}^n\log f(x_i\mid\mu,\beta).$$
为简化计算，假设$\beta$已知，故解$\frac d{d\mu}l(\mu,\beta\mid\boldsymbol x)=0$求MLE $\hat\mu$. 没有显式解，所以采用数值方法.

由\autoref{thm:10.1.6}，知道$\hat\mu$渐近有效. 关注的问题是，使用较容易计算的矩估计量会损失多少. 为作比较，计算渐近相对效率\footnote{译注：此处ARE公式与\autoref{def:10.1.16}及随后“大于1”的说法互有不一致. 按定义，$\operatorname{ARE}(\bar X,\hat\mu)$应为MLE渐近方差除以$\bar X$的渐近方差；若以单个观测信息量$I(\mu)$表示，应为$1/[I(\mu)\beta\mu]\leq1$. 原文公式使用全样本信息量作分子，图中又显示其倒数方向的效率损失. 保留原公式、文字与原图，待对照作者勘误核实.}
$$\operatorname{ARE}(\bar X,\hat\mu)=\frac{E\left(-\frac{d^2}{d\mu^2}l(\mu,\beta\mid\boldsymbol X)\right)}{\beta\mu},$$
并在\autoref{fig:10.1.1}中对若干$\beta$值画出它. 当然，知道ARE一定大于1，但从图中可以看到，对较大的$\beta$，进行更复杂的计算并使用MLE是值得的. （推广见\autoref{exercise:10.11}，计算细节见\autoref{ex:12.6.7}.）
\end{example}

\subsubsection{自助法标准误}\label{sec:10.1.4}\index{SI0094@自助法}
我们在\autoref{ex:1.2.20}中首次见到的自助法（bootstrap），提供了另一种计算标准误的方法. （它还能提供更多内容，见\hyperref[sec:10.6.3]{杂录10.6.3}.）

自助法基于一个简单而有力的思想（其数学细节可能相当复杂\footnote{易读的入门介绍见\hyperref[bib:source199]{Lehmann（1999}，第6.5节）.}）. 在统计学中，通过抽样了解总体的特征. 因为样本代表总体，样本的相应特征应能提供总体特征的信息. 自助法通过重抽样（resamples，即从原样本中再次抽样）帮助我们了解样本特征，并利用这些信息推断总体.
\sourcepageanchor{417}
\par\begin{figure}[htbp]
\centering\includegraphics[width=.9\textwidth]{assets/chapter10/fig-10.1.1.pdf}
\caption{伽马均值的矩估计量与MLE的渐近相对效率. 四条曲线对应尺度参数值$(1,3,5,10)$，较高的曲线对应较大的尺度参数.}\label{fig:10.1.1}
\end{figure}\par
自助法由Efron于20世纪70年代末提出，最初的思想见Efron（\hyperref[bib:source106]{1979a}\hyperref[bib:source107]{b}），随后有专著\hyperref[bib:source108]{Efron（1982）}. 较新的思考与发展另见\hyperref[bib:source109]{Efron（1998）}.

先看一个其实并不需要自助法的简单例子.

\begin{example}[方差的自助法]\label{ex:10.1.19}\index{SI0095@非参数}\index{SI0913@重抽样}\index{SI0950@有放回}
在\autoref{ex:1.2.20}中，我们从四个数
$$2,4,9,12$$
中有放回地抽取四个数，计算所有可能的平均值. 这是最简单形式的自助法，有时称为非参数自助法（nonparametric bootstrap）. \autoref{fig:1.2.2}以直方图显示这些值.

所构造的是样本均值可能取值的重抽样. 已经看到，有$\binom{4+4-1}{4}=35$个不同的可能值，但它们并不等概率（所以不能视为随机样本）. $4^4=256$个（不全相异的）重抽样等可能，可以视为随机样本. 对第$i$个重抽样，令$\bar x_i^*$为其均值. 则样本均值$\bar X$的方差可估计为
\begin{equation}\label{eq:10.1.9}
\operatorname{Var}^*(\bar X)=\frac1{n^n-1}\sum_{i=1}^{n^n}(\bar x_i^*-\bar{\bar x}^{\,*})^2,
\end{equation}
其中$\bar{\bar x}^{\,*}=\frac1{n^n}\sum_{i=1}^{n^n}\bar x_i^*$是各个重抽样的均值. （通常以记号$*$表示自助法所得的、即重抽样所得的值.）

在本例中，自助法均值和方差分别为$\bar{\bar x}^{\,*}=6.75$和$\operatorname{Var}^*(\bar X)=3.94$. 事实表明，就均值和方差而言，自助法估计与通常估计几乎相同（见\autoref{exercise:10.13}）.
\end{example}

\sourcepageanchor{418}
我们已经看到如何计算自助法标准误，不过是在一个实际上不需要它的问题中. 自助法真正的优势在于，与Delta方法一样，方差公式\eqref{eq:10.1.9}几乎适用于任何估计量. 因此，对任意估计量$\hat\theta(\boldsymbol x)=\hat\theta$，可写为
\begin{equation}\label{eq:10.1.10}
\operatorname{Var}^*(\hat\theta)=\frac1{n^n-1}\sum_{i=1}^{n^n}(\hat\theta_i^*-\bar{\hat\theta}^{\,*})^2,
\end{equation}
其中$\hat\theta_i^*$是从第$i$个重抽样计算的估计量，$\bar{\hat\theta}^{\,*}=\frac1{n^n}\sum_{i=1}^{n^n}\hat\theta_i^*$是重抽样值的均值.

\begin{example}[二项方差的自助法]\label{ex:10.1.20}
在\autoref{ex:10.1.15}中，用Delta方法估计$\hat p(1-\hat p)$的方差. 对样本量$n$，也可将这个方差估计为
$$\operatorname{Var}^*(\hat p(1-\hat p))=\frac1{n^n-1}\sum_{i=1}^{n^n}\left((\hat p(1-\hat p))_i^*-\overline{\hat p(1-\hat p)}^{\,*}\right)^2.$$
\end{example}

但出现了新问题. 在\autoref{ex:10.1.19}中，$n=4$，自助法求和有256项. 对更典型的样本量，这个数会大到无法计算. （当$n>15$时，枚举全部可能的重抽样几乎不可能，至少对作者们如此.）不过，现在记起我们是统计学家：对重抽样再抽样即可.

因此，对样本$\boldsymbol x=(x_1,x_2,\ldots,x_n)$及估计$\hat\theta(x_1,x_2,\ldots,x_n)=\hat\theta$，选取$B$个重抽样（或自助样本，bootstrap samples），并计算
\begin{equation}\label{eq:10.1.11}\index{SI0913@重抽样}
\operatorname{Var}_B^*(\hat\theta)=\frac1{B-1}\sum_{i=1}^B(\hat\theta_i^*-\bar{\hat\theta}^{\,*})^2.
\end{equation}

\begin{example}[\autoref{ex:10.1.20}的结尾]\label{ex:10.1.21}
对样本量$n=24$，取$B=1{,}000$，计算$\hat p(1-\hat p)$的Delta方法方差估计和自助法方差估计. 当$\hat p\ne1/2$时，采用\autoref{ex:10.1.15}的一阶Delta方法方差估计；当$\hat p=1/2$时，采用\autoref{thm:5.5.26}的二阶方差估计（见\autoref{exercise:10.16}）. \autoref{tab:10.1.1}显示，在所有情形下，自助法方差估计都更接近真实方差，而Delta方法方差都低估了它. （这不应令人意外，因为\eqref{eq:10.1.7}表明Delta方法方差估计基于一个下界.）\footnote{译注：表中“真实值”与所述$n=24$二项模型不符. 按该模型对$Y=0,\ldots,24$逐项求和，$\operatorname{Var}[(Y/24)(1-Y/24)]$在$p=1/4,1/2,2/3$时约为$.001911,.000208,.001109$，并非表中的$.00484,.00531,.00519$. 原表及其比较结论保留；此核算仅针对表述的这三个参数值.}

\par\begin{table}[htbp]
\centering
\caption{$\hat p(1-\hat p)$的自助法（上行）与Delta方法（下行）方差. 当$\hat p=1/2$时使用二阶Delta方法（见\autoref{thm:5.5.26}）. 真实方差在假设$\hat p=p$下数值计算.}\label{tab:10.1.1}
\begin{tabular}{lrrr}\toprule
&$\hat p=1/4$&$\hat p=1/2$&$\hat p=2/3$\\\midrule
自助法&.00508&.00555&.00561\\
Delta方法&.00195&.00022&.00102\\\midrule
真实值&.00484&.00531&.00519\\\bottomrule
\end{tabular}
\end{table}\par

Delta方法是“一阶”近似，因为它基于泰勒级数展开的第一项. 当该项为零（例如$\hat p=1/2$）时，必须使用二阶Delta方法. 相比之下，自助法常有“二阶”准确性，能使展开中的不止第一项正确（见\hyperref[sec:10.6.3]{杂录10.6.3}）. 所以这里，自助法会自动修正$\hat p=1/2$的情形. （注意$24^{24}\approx1.33\times10^{13}$，这是极大的数，所以无法生成全部自助样本.）\footnote{译注：原文数值为$1.33\times10^{13}$，但$24^{24}\approx1.33\times10^{33}$. 保留原值；两者均支持这里无法枚举全部样本的结论.}
\end{example}

迄今讨论的自助法称为非参数自助法，因为没有假设总体pdf的函数形式.\index{SI0096@参数}\semanticindexanchor{96:418} 相比之下，也可以有参数自助法（parametric bootstrap）.

\sourcepageanchor{419}
假设样本$X_1,X_2,\ldots,X_n$来自pdf为$f(x\mid\theta)$的分布，其中$\theta$可以是参数向量. 可以用MLE $\hat\theta$估计$\theta$，并抽取样本
$$X_1^*,X_2^*,\ldots,X_n^*\sim f(x\mid\hat\theta).$$
若抽取$B$个这样的样本，就能用\eqref{eq:10.1.11}估计$\hat\theta$的方差. 注意，这些样本并非对数据的重抽样，而是从$f(x\mid\hat\theta)$实际抽出的随机样本；这个分布有时称为代入分布（plug-in distribution）.

\begin{example}[参数自助法]\label{ex:10.1.22}
假设有样本
$$-1.81,0.63,2.22,2.41,2.95,4.16,4.24,4.53,5.09,$$
其$\bar x=2.71$，$s^2=4.82$. 若假设基础分布为正态分布，则参数自助法会抽取样本
$$X_1^*,X_2^*,\ldots,X_n^*\sim\mathrm n(2.71,4.82).$$
基于$B=1000$个样本，算得$\operatorname{Var}_B^*(S^2)=4.33$. 根据正态理论，$S^2$的方差为$2(\sigma^2)^2/8$，可以用MLE $2(4.82)^2/8=5.81$估计它. 数据值模拟自方差为4的正态分布，因此这里参数自助法给出更好的估计. （在\autoref{ex:5.6.6}中，我们已经用现在所知的参数自助法估计过$S^2$的分布.）\footnote{译注：这里把以$s^2=4.82$代入的方差估计称为MLE；一般正态模型的方差MLE为$(n-1)s^2/n$，故该称呼与通常MLE定义不符. 原书数值与称呼均保留.}
\end{example}

现在有了一种计算标准误的通用方法，如何知道它是好方法呢？在\autoref{ex:10.1.21}中，它似乎优于Delta方法，而后者已知具有一些良好性质. 特别是，基于最大似然估计的Delta方法通常产生相合估计量. 对自助法能作相同的断言吗？

虽然无法十分一般地回答这个问题，但可以说，在许多情形下，自助法的确提供合理且相合的估计量.

为更精确，将计算自助法估计量的两个不同部分分开.
\begin{enumerate}[label=\alph*.]
\item\label{item:10.1.bootstrap.a} 证明当$B\to\infty$时，\eqref{eq:10.1.11}收敛到\eqref{eq:10.1.10}，即
$$\operatorname{Var}_B^*(\hat\theta)\xrightarrow{B\to\infty}\operatorname{Var}^*(\hat\theta).$$
\item\label{item:10.1.bootstrap.b}\index{SI0773@相合的自助法} 证明使用全部自助样本的估计量\eqref{eq:10.1.10}相合，即
$$\operatorname{Var}^*(\hat\theta)\xrightarrow{n\to\infty}\operatorname{Var}(\hat\theta).$$
\end{enumerate}
\hyperref[item:10.1.bootstrap.a]{(a)部分}可由大数定律建立（\autoref{exercise:10.15}）. 还应注意，\hyperref[item:10.1.bootstrap.a]{(a)部分}全部发生在样本内部. （\hyperref[bib:source199]{Lehmann，1999}，第6.5节把$\operatorname{Var}_B^*(\hat\theta)$称为近似器，而不是估计量.）

证明\hyperref[item:10.1.bootstrap.b]{(b)部分}稍微微妙一些，相合性正是在这里建立的. 独立同分布抽样时通常能够得到相合性，但在更一般的情形中未必如此（\hyperref[bib:source199]{Lehmann，1999}，第6.5节给出一个例子）. 关于相合性的更多细节（必然需要较高深的层次），见\hyperref[bib:source288]{Shao和Tu（1995}，第3.2.2节）或\hyperref[bib:source287]{Shao（1999}，第5.5.3节）.



\subsection{稳健性}\label{sec:10.2}
迄今为止，我们总是假设基础模型正确来评价估计量的表现. 在这个假设下，已经导出在某种意义下最优的估计量. 然而，若基础模型不正确，估计量的最优性就无法保证.

\sourcepageanchor{420}
我们无法防范所有可能的情形；而且，如果模型经过仔细考虑才确定，也不应有这种必要.\index{SI0813@稳健}\semanticindexanchor{813:420} 但我们可能关心与所假设模型的小幅或中等程度偏离. 这会引导我们考虑稳健估计量（robust estimators）. 这类估计量放弃在假设模型下的最优性，换取假设模型不是真实模型时的合理表现. 因此存在取舍，最优性与稳健性哪个更重要，最好视具体情形决定.

“稳健性”一词可以有多种解释；\hyperref[bib:source171]{Huber（1981}，第1.2节）的说法或许是最好的概括：
\begin{quote}
“……任何统计程序都应具有以下理想特征：
\begin{enumerate}[label=(\arabic*)]
\item\label{item:10.2.huber.1} 在假设模型下，应有相当好的（最优或近乎最优的）效率.
\item\label{item:10.2.huber.2} 它应当稳健，即模型假设的小幅偏离只应轻微损害表现……
\item\label{item:10.2.huber.3} 模型较大一些的偏离不应造成灾难.”
\end{enumerate}
\end{quote}
先看一些简单例子以更好地理解这些要求，然后再看更一般的稳健估计量和稳健性度量.

\subsubsection{均值与中位数}\label{sec:10.2.1}
样本均值是稳健估计量吗？这可能取决于如何精确地形式化稳健性的度量.

\begin{example}[样本均值的稳健性]\label{ex:10.2.1}
设$X_1,X_2,\ldots,X_n$独立同分布于$\mathrm n(\mu,\sigma^2)$. 已知$\bar X$的方差$\operatorname{Var}(\bar X)=\sigma^2/n$就是Cramér--Rao下界. 因而$\bar X$满足\hyperref[item:10.2.huber.1]{(1)}，在假设模型下达到最佳方差.

为研究\hyperref[item:10.2.huber.2]{(2)}，即$\bar X$在模型小幅偏离时的表现，先要明确这种偏离的含义. 通常的解释是采用$\delta$污染模型（$\delta$-contamination model），即对很小的$\delta$，假设观测到
$$X_i\sim\begin{cases}\mathrm n(\mu,\sigma^2),&\text{概率为$1-\delta$},\\f(x),&\text{概率为$\delta$},\end{cases}$$
其中$f(x)$是另一个分布.

假设取$f(x)$为任意均值为$\mu$、方差为$\tau^2$的密度. 则\footnote{译注：原文先说污染分布的均值为$\mu$，随后却使用未定义的$\theta$；若意指污染均值为$\theta$，独立样本均值的混合方差最后一项还应除以$n$. 此处原文与原式均保留.}
$$\operatorname{Var}(\bar X)=(1-\delta)\frac{\sigma^2}n+\delta\frac{\tau^2}n+\delta(1-\delta)(\theta-\mu)^2.$$
这看起来对$\bar X$相当有利，因为若$\theta\approx\mu$且$\sigma\approx\tau$，$\bar X$就接近最优. 不过，可以进一步扰动模型，使情况相当糟糕. 考虑$f(x)$为柯西pdf的情形，立即可得$\operatorname{Var}(\bar X)=\infty$. （细节见\autoref{exercise:10.18}，另一个情形见\autoref{exercise:10.19}.）
\end{example}

转向要求\hyperref[item:10.2.huber.3]{(3)}，我们问：若有一个异常偏离的观测，会发生什么？设想一组固定的样本值，然后考虑增大其中最大观测值的影响. 例如，假设$X_{(n)}=x$，其中$x\to\infty$. 这种观测的影响可被认为是“灾难性的”. 虽然$\bar X$的分布性质并未改变，观测值却会变得“毫无意义”. 这说明了崩溃值（breakdown value）的概念，该思想归于\hyperref[bib:source156]{Hampel（1974）}.

\begin{definition}\label{def:10.2.2}\index{SI0100@崩溃值}\index{SI0769@崩溃值}
设$X_{(1)}<\cdots<X_{(n)}$为容量$n$的有序样本，$T_n$为基于该样本的统计量. 若对每个$\varepsilon>0$，
$$\lim_{X_{(\{(1-b)n\})}\to\infty}T_n<\infty\quad\text{且}\quad
\lim_{X_{(\{(1-(b+\varepsilon))n\})}\to\infty}T_n=\infty,$$
则称$T_n$的崩溃值为$b$，$0\leq b\leq1$.
\end{definition}
（百分位数记号见\autoref{def:5.4.2}.）

\sourcepageanchor{421}
容易看出，$\bar X$的崩溃值是0，也就是说，将样本中任意比例的观测推向无穷，$\bar X$也随之趋于无穷. 与此形成鲜明对比，样本中位数不因样本值的这种变化而改变. 对极端观测的不敏感有时被视为样本中位数的优点，其崩溃值为50\%. （关于崩溃值的更多内容见\autoref{exercise:10.20}.）\footnote{译注：\autoref{def:10.2.2}的第一个极限在$b=0$时要求最大观测趋于无穷而统计量仍有限，与这里样本均值的说法在有限样本边界上不相容. 需明确崩溃比例的上确界或渐近约定，并限制$\varepsilon$以保证次序下标有效；原定义和结论照录，待核实.}

既然中位数比均值更稳健，我们可以问，改用更稳健的估计量是否会有所损失（当然一定会！）. 例如，在\autoref{ex:10.2.1}的简单正态模型中，如果模型真实，均值是最佳无偏估计量. 因此，在正态模型下（及其附近），均值是更好的估计量. 但关键问题是，在正态模型下均值究竟好多少. 若能回答，就可以有根据地选择估计量，并判断最优性和稳健性哪个准则更重要. 为在一定一般性下回答这个问题，使用渐近相对效率准则.

计算中位数相对于均值的ARE，首先必须建立中位数的渐近正态性，计算渐近分布的方差.

\begin{example}[中位数的渐近正态性]\label{ex:10.2.3}\index{SI0804@极限分布}
为求中位数的极限分布，采用与\autoref{thm:5.4.3}和\autoref{thm:5.4.4}证明相似的论证，即基于二项分布的论证.

设$X_1,X_2,\ldots,X_n$来自pdf为$f$、cdf为$F$（假设可微）的总体，且$P(X_i\leq\mu)=1/2$，所以$\mu$是总体中位数. 令$M_n$为样本中位数，考虑对某个$a$计算
$$\lim_{n\to\infty}P\bigl(\sqrt n(M_n-\mu)\leq a\bigr).$$
若定义随机变量$Y_i$为
$$Y_i=\begin{cases}1,&X_i\leq\mu+a/\sqrt n,\\0,&\text{其他},\end{cases}$$
则$Y_i$为成功概率$p_n=F(\mu+a/\sqrt n)$的伯努利随机变量. 为避免复杂情况，假设$n$为奇数，于是事件$\{M_n\leq\mu+a/\sqrt n\}$等价于$\{\sum_iY_i\geq(n+1)/2\}$.

作一些代数运算得到
$$P\bigl(\sqrt n(M_n-\mu)\leq a\bigr)=P\left(\frac{\sum_iY_i-np_n}{\sqrt{np_n(1-p_n)}}\geq\frac{(n+1)/2-np_n}{\sqrt{np_n(1-p_n)}}\right).$$
现在$p_n\to p=F(\mu)=1/2$，因此预期应用中心极限定理可证明$(\sum_iY_i-np_n)/\sqrt{np_n(1-p_n)}$依分布收敛到标准正态随机变量$Z$. 直接计算极限也可得到
$$\frac{(n+1)/2-np_n}{\sqrt{np_n(1-p_n)}}\longrightarrow-2aF'(\mu)=-2af(\mu).$$
综合起来，得到
$$P\bigl(\sqrt n(M_n-\mu)\leq a\bigr)\longrightarrow P(Z\geq-2af(\mu)),$$
因此$\sqrt n(M_n-\mu)$渐近正态，均值为0，方差为$1/[2f(\mu)]^2$. （细节见\autoref{exercise:10.22}；严格且更一般的推导见\hyperref[bib:source287]{Shao，1999}，第5.3节.）
\end{example}

\sourcepageanchor{422}
\begin{example}[中位数相对于均值的ARE]\label{ex:10.2.4}\index{SI0068@中位数相对于均值}
均值与中位数的渐近方差都有简单表达式，因此很快就能计算ARE. 这里考察三种对称分布的ARE，发现正如预期的那样，分布尾部越重，ARE越大. 也就是说，在重尾分布下中位数表现更好. 更多比较见\autoref{exercise:10.23}.
\end{example}
\par\begin{table}[htbp]
\centering\setcounter{table}{1}
\caption{中位数相对于均值的渐近相对效率}\label{tab:10.2.2}
\begin{tabular}{ccc}\toprule 正态分布&逻辑斯蒂分布&双指数分布\\\midrule .64&.82&2\\\bottomrule\end{tabular}
\end{table}\par

\subsubsection{\texorpdfstring{$M$}{M}估计量}\label{sec:10.2.2}\index{SI0816@M估计量}
许多估计量都是通过极小化某种准则得到的. 例如，若$X_1,X_2,\ldots,X_n$独立同分布于$f(x\mid\theta)$，可采用的估计量包括：均值，它极小化$\sum(x_i-a)^2$；中位数，它极小化$\sum|x_i-a|$；或MLE，它极大化$\prod_{i=1}^nf(x_i\mid\theta)$（或极小化负似然）. 作为系统获得稳健估计量的方法，可以尝试写出某种准则函数，使其极小值对应具有理想稳健性质的估计量.

为定义稳健准则，\hyperref[bib:source170]{Huber（1964）}考虑了均值与中位数之间的折中. 均值准则为平方，这赋予它敏感性，但在“尾部”，平方给予大观测过多的权重. 相比之下，中位数的绝对值准则不会过分加重大观测或小观测的权重. 折中方案是极小化准则函数
\begin{equation}\label{eq:10.2.1}
\sum_{i=1}^n\rho(x_i-a),
\end{equation}
其中$\rho$为
\begin{equation}\label{eq:10.2.2}\index{SI1043@调节参数}
\rho(x)=\begin{cases}\frac12x^2,&|x|\leq k,\\k|x|-\frac12k^2,&|x|\geq k.\end{cases}
\end{equation}
函数$\rho(x)$在$|x|\leq k$时类似$x^2$，在$|x|>k$时类似$|x|$. 此外，因为$\frac12k^2=k|k|-\frac12k^2$，该函数连续（\autoref{exercise:10.28}）. 常数$k$也称调节参数（tuning parameter），它控制两者的混合；$k$越小，得到的估计量越“类似中位数”.

\begin{example}[Huber估计量]\label{ex:10.2.5}\index{SI0814@Huber估计量}
以\eqref{eq:10.2.1}和\eqref{eq:10.2.2}的极小化点定义的估计量称为Huber估计量. 为了解它如何工作以及$k$的选择为何重要，考虑以下数据集，其中有八个标准正态偏差值和三个“离群值”：
$$\boldsymbol x=-1.28,-0.96,-0.46,-0.44,-0.26,-0.21,-0.063,0.39,3,6,9.$$
这些数据的均值为1.33，中位数为$-.21$. 随$k$变化，Huber估计值如\autoref{tab:10.2.3}所示. 可见当$k$增大时，Huber估计在中位数与均值之间变化，因此可以把增大$k$理解为降低对离群值的稳健性.

\par\begin{table}[htbp]\centering
\caption{Huber估计量}\label{tab:10.2.3}
\begin{tabular}{c|rrrrrrrrr}\toprule
$k$&0&1&2&3&4&5&6&8&10\\\midrule
估计值&$-.21$&.03&$-.04$&.29&.41&.52&.87&.97&1.33\\\bottomrule
\end{tabular}
\end{table}\par
\end{example}

极小化\eqref{eq:10.2.2}的估计量是Huber所研究估计量的特例. 对一般函数$\rho$，把极小化$\sum_i\rho(x_i-\theta)$的估计量称为$M$估计量（$M$-estimator），这个名称提醒我们它们属于最大似然型估计量（maximum-likelihood-type estimators）. 注意，若选择$\rho$为负对数似然$-l(\theta\mid\boldsymbol x)$，则$M$估计量就是通常的MLE.
\sourcepageanchor{423}
不过，若在选择待极小化的函数时更灵活，就能导出具有不同性质的估计量.

由于极小化函数通常通过求导数的零点来完成（在可求导时），定义$\psi=\rho'$，可见$M$估计量是方程
\begin{equation}\label{eq:10.2.3}
\sum_{i=1}^n\psi(x_i-\theta)=0
\end{equation}
的解.

将估计量刻画为方程的根，对研究其性质特别有用，因为似然估计中使用的论证可以推广.\index{SI1029@泰勒级数}\semanticindexanchor{1029:423} 特别参见\autoref{sec:10.1.2}，尤其是\autoref{thm:10.1.6}的证明. 假设$\rho(x)$对称，其导数$\psi(x)$单调递增（这保证\eqref{eq:10.2.3}的根是唯一极小值点）. 然后，像\autoref{thm:10.1.6}的证明一样，对$\phi$作泰勒展开如下：\footnote{译注：本段原文写$\phi$，但此前定义的是$\psi$；所引\autoref{thm:10.1.6}应核对为\autoref{thm:10.1.12}. 此外，对$\psi(x_i-\theta)$关于$\theta$展开，一阶项前应为负号；原文以下展开及整理后的负号均保留. 最后的方差因分母平方而不受此符号问题影响.}
$$\sum_{i=1}^n\psi(x_i-\theta)=\sum_{i=1}^n\psi(x_i-\theta_0)+(\theta-\theta_0)\sum_{i=1}^n\psi'(x_i-\theta_0)+\cdots,$$
其中$\theta_0$是真值，并忽略高阶项. 令$\hat\theta_M$为\eqref{eq:10.2.3}的解，代替$\theta$得到
$$0=\sum_{i=1}^n\psi(x_i-\theta_0)+(\hat\theta_M-\theta_0)\sum_{i=1}^n\psi'(x_i-\theta_0)+\cdots,$$
左侧为零是因为$\hat\theta_M$为解. 现在再次类似\autoref{thm:10.1.6}的证明，整理各项，除以$\sqrt n$，忽略余项，得到
$$\sqrt n(\hat\theta_M-\theta_0)=\frac{-\frac1{\sqrt n}\sum_{i=1}^n\psi(x_i-\theta_0)}{\frac1n\sum_{i=1}^n\psi'(x_i-\theta_0)}.$$
现假设$\theta_0$满足$E_{\theta_0}\psi(X-\theta_0)=0$（通常也将此作为$\theta_0$的定义）. 于是
\begin{equation}\label{eq:10.2.4}
-\frac1{\sqrt n}\sum_{i=1}^n\psi(X_i-\theta_0)
=\sqrt n\left[-\frac1n\sum_{i=1}^n\psi(X_i-\theta_0)\right]
\longrightarrow\mathrm n\bigl(0,E_{\theta_0}[\psi(X-\theta_0)]^2\bigr)
\end{equation}
依分布成立，而大数定律给出
\begin{equation}\label{eq:10.2.5}
\frac1n\sum_{i=1}^n\psi'(x_i-\theta_0)\longrightarrow E_{\theta_0}\psi'(X-\theta_0)
\end{equation}
依概率成立. 综合起来有
\begin{equation}\label{eq:10.2.6}\index{SI0815@极限分布}
\sqrt n(\hat\theta_M-\theta_0)\longrightarrow\mathrm n\left(0,\frac{E_{\theta_0}\psi(X-\theta_0)^2}{[E_{\theta_0}\psi'(X-\theta_0)]^2}\right).
\end{equation}

\begin{example}[Huber估计量的极限分布]\label{ex:10.2.6}
若$X_1,X_2,\ldots,X_n$独立同分布于pdf $f(x-\theta)$，其中$f$关于零对称，则对\eqref{eq:10.2.2}给出的$\rho$，有
\begin{equation}\label{eq:10.2.7}
\psi(x)=\begin{cases}x,&|x|\leq k,\\k,&x>k,\\-k,&x<-k,\end{cases}
\end{equation}
\sourcepageanchor{424}
所以\footnote{译注：按\eqref{eq:10.2.7}直接积分，下一式第一项应为正号，左右尾的系数应分别为$-k$与$k$；原文的负号及$\pm2k$与该定义不符. 对称性仍使两种写法均为零，原式保留.}
\begin{equation}\label{eq:10.2.8}
\begin{aligned}
E_\theta\psi(X-\theta)
&=-\int_{\theta-k}^{\theta+k}(x-\theta)f(x-\theta)\,dx\\
&\quad+2k\int_{-\infty}^{\theta-k}f(x-\theta)\,dx-2k\int_{\theta+k}^\infty f(x-\theta)\,dx\\
&=-\int_{-k}^kyf(y)\,dy+2k\int_{-\infty}^{-k}f(y)\,dy-2k\int_k^\infty f(y)\,dy=0,
\end{aligned}
\end{equation}
其中代入$y=x-\theta$. 由$f$的对称性，各积分相加为零. 因此Huber估计量具有正确的均值. （推广见\autoref{exercise:10.25}.）

为计算方差，需要$\psi'$的期望. 虽然$\psi$并非处处可微，但超出不可微点$x=\pm k$之后，$\psi'$为零. 因此只需处理$|x|\leq k$内的期望，有
$$\begin{aligned}
E_\theta\psi'(X-\theta)&=\int_{\theta-k}^{\theta+k}f(x-\theta)\,dx=P_0(|X|\leq k),\\
E_\theta[\psi(X-\theta)]^2
&=\int_{\theta-k}^{\theta+k}(x-\theta)^2f(x-\theta)\,dx\\
&\quad+k^2\int_{\theta+k}^\infty f(x-\theta)\,dx+k^2\int_{-\infty}^{\theta-k}f(x-\theta)\,dx\\
&=\int_{-k}^kx^2f(x)\,dx+2k^2\int_k^\infty f(x)\,dx.
\end{aligned}$$
因此，Huber估计量渐近正态，均值为$\theta$，渐近方差为
$$\frac{\int_{-k}^kx^2f(x)\,dx+k^2P_0(|X|>k)}{[P_0(|X|\leq k)]^2}.$$
\end{example}

像\autoref{ex:10.2.4}一样，现在考察多种分布下Huber估计量的ARE.

\begin{example}[Huber估计量的ARE]\label{ex:10.2.7}\index{SI0067@Huber估计量}
Huber估计量在某种意义上是均值与中位数的折中，因此考察它相对于这两个估计量的相对效率. 对正态分布和逻辑斯蒂分布，Huber估计量与均值表现相近，且优于中位数. 对双指数分布，它优于均值，但不如中位数. 回顾，均值是正态分布的MLE，中位数是双指数分布的MLE（因此预期ARE小于1）. 对这些分布，Huber估计量具有接近MLE的表现，同时在其他情形下似乎也保持合理表现.

\par\begin{table}[htbp]\centering
\caption{Huber估计量的渐近相对效率，$k=1.5$}\label{tab:10.2.4}
\begin{tabular}{l|ccc}\toprule
&正态分布&逻辑斯蒂分布&双指数分布\\\midrule
相对于均值&.96&1.08&1.37\\
相对于中位数&1.51&1.31&.68\\\bottomrule
\end{tabular}
\end{table}\par
\end{example}

可见，$M$估计量在稳健性与效率之间折中. 下面更仔细地看看，为获得稳健性，可能在效率方面放弃多少.

更仔细地考察\eqref{eq:10.2.6}中的渐近方差. 分母包含$E_{\theta_0}\psi'(X-\theta)$，可写为\footnote{译注：原文把$\psi'(x-\theta)$等同于$\partial\psi(x-\theta)/\partial\theta$；若撇号是对函数自变量求导，二者相差负号. 原式保留；下一页的乘积求导恒等式则显式处理了该负号.}
$$E_\theta\psi'(X-\theta)=\int\psi'(x-\theta)f(x-\theta)\,dx
=\int\left[\frac\partial{\partial\theta}\psi(x-\theta)\right]f(x-\theta)\,dx.$$

\sourcepageanchor{425}
现用乘积求导法则得到
$$\begin{aligned}
\frac d{d\theta}\int\psi(x-\theta)f(x-\theta)\,dx
&=\int\left[\frac d{d\theta}\psi(x-\theta)\right]f(x-\theta)\,dx\\
&\quad+\int\psi(x-\theta)\left[\frac d{d\theta}f(x-\theta)\right]\,dx.
\end{aligned}$$
若左侧为零——合理的$\psi$函数应当如此，仿照$l'$的行为——则有
$$\begin{aligned}
-\int\left[\frac d{d\theta}\psi(x-\theta)\right]f(x-\theta)\,dx
&=\int\psi(x-\theta)\left[\frac d{d\theta}f(x-\theta)\right]\,dx\\
&=\int\psi(x-\theta)\left[\frac d{d\theta}\log f(x-\theta)\right]f(x-\theta)\,dx,
\end{aligned}$$
其中利用了$\frac d{dy}g(y)/g(y)=\frac d{dy}\log g(y)$. 最后一个表达式可写为$E_\theta[\psi(X-\theta)l'(\theta\mid X)]$，其中$l(\theta\mid X)$是对数似然，从而得到恒等式
$$-E_\theta\left[\frac d{d\theta}\psi(X-\theta)\right]=E_\theta[\psi(X-\theta)l'(\theta\mid X)].$$
（取$\psi=l'$时，就得到我们希望读者已经熟悉的等式$-E_\theta[l''(\theta\mid X)]=E_\theta l'(\theta\mid X)^2$，见\autoref{lem:7.3.11}.）

现在比较$M$估计量与MLE的渐近方差就很简单. 回顾MLE $\hat\theta$的渐近方差为$E_\theta l'(\theta\mid X)^2$，因此\footnote{译注：原文此句将MLE渐近方差写为信息量；正确的通常形式是其倒数$1/E_\theta[l'(\theta\mid X)]^2$. 接下来的ARE公式与倒数形式一致. 原句保留.}
\begin{equation}\label{eq:10.2.9}
\operatorname{ARE}(\hat\theta_M,\hat\theta)=\frac{[E_\theta(\psi(X-\theta_0)l'(\theta\mid X)]^2}{E_\theta\psi(X-\theta)^2E_\theta l'(\theta\mid X)^2}\leq1,
\end{equation}
其中应用Cauchy--Schwarz不等式. 因此，$M$估计量的效率总低于MLE，且只有当$\psi$与$l'$成比例时才能达到相同效率（见\autoref{exercise:10.29}）.

本节并未试图将所有稳健估计量分类，而是满足于若干例子. 有许多优秀著作详细讨论稳健性；感兴趣的读者可阅读\hyperref[bib:source299]{Staudte和Sheather（1990）}或\hyperref[bib:source165]{Hettmansberger和McKean（1998）}\footnote{译注：原书正文将第一作者姓氏拼作Hettmansberger，书后对应条目[165]作Hettmansperger；正文拼写保留，链接至该条目.}.



\subsection{假设检验}\label{sec:10.3}
与\autoref{sec:10.1}一样，本节介绍几种在复杂问题中导出某些检验的方法. 所考虑的问题，是先前定义的最优检验不存在（例如不存在UMP无偏检验），或尚不知道最优检验的问题. 此时，导出任何合理的检验都可能有用. 以下两小节讨论似然比检验的大样本性质，以及其他近似大样本检验.

\subsubsection{似然比检验的渐近分布}\label{sec:10.3.1}
对复杂模型，构造检验最有用的方法之一是似然比法，因为它明确地定义检验统计量
$$\lambda(\boldsymbol x)=\frac{\sup_{\Theta_0}L(\theta\mid\boldsymbol x)}{\sup_\Theta L(\theta\mid\boldsymbol x)},$$
并明确给出拒绝域的形式$\{\boldsymbol x:\lambda(\boldsymbol x)\leq c\}$. 观测到数据$\boldsymbol X=\boldsymbol x$后，似然函数$L(\theta\mid\boldsymbol x)$就是变量$\theta$的完全确定的函数.
\sourcepageanchor{426}
即使不能解析地求出$L(\theta\mid\boldsymbol x)$在$\Theta_0$和$\Theta$上的两个上确界，通常也能数值计算. 因此，即使没有定义$\lambda(\boldsymbol x)$的方便公式，也能对观测数据点求得检验统计量$\lambda(\boldsymbol x)$.

为定义水平$\alpha$的检验，必须选择常数$c$使
\begin{equation}\label{eq:10.3.1}
\sup_{\theta\in\Theta_0}P_\theta\bigl(\lambda(\boldsymbol X)\leq c\bigr)\leq\alpha.
\end{equation}
若不能导出$\lambda(\boldsymbol x)$的简单公式，似乎就无望求出$\lambda(\boldsymbol X)$的抽样分布，从而知道如何选$c$以保证\eqref{eq:10.3.1}. 但借助渐近理论，可以获得近似答案.

类似\autoref{thm:10.1.12}，有如下结果.

\begin{theorem}[LRT的渐近分布：简单$H_0$]\label{thm:10.3.1}
为检验$H_0:\theta=\theta_0$对$H_1:\theta\ne\theta_0$，假设$X_1,\ldots,X_n$独立同分布于$f(x\mid\theta)$，$\hat\theta$为$\theta$的MLE，且$f(x\mid\theta)$满足\hyperref[sec:10.6.2]{杂录10.6.2}中的正则条件. 则在$H_0$下，当$n\to\infty$时，
$$-2\log\lambda(\boldsymbol X)\longrightarrow\chi_1^2\quad\text{依分布},$$
其中$\chi_1^2$是自由度为1的卡方随机变量.
\end{theorem}
\begin{proof}
首先在$\hat\theta$处对$\log L(\theta\mid\boldsymbol x)=l(\theta\mid\boldsymbol x)$作泰勒展开，得到
$$l(\theta\mid\boldsymbol x)=l(\hat\theta\mid\boldsymbol x)+l'(\hat\theta\mid\boldsymbol x)(\theta-\hat\theta)+l''(\hat\theta\mid\boldsymbol x)\frac{(\theta-\hat\theta)^2}{2!}+\cdots.$$
现将$l(\theta_0\mid\boldsymbol x)$的展开式代入$-2\log\lambda\boldsymbol x=-2l(\theta_0\mid\boldsymbol x)+2l(\hat\theta\mid\boldsymbol x)$，得到\footnote{译注：原文下式把平方差除以$-l''$；由上面的泰勒展开应为$-l''(\hat\theta\mid\boldsymbol x)(\theta_0-\hat\theta)^2$，或除以观测信息量的倒数. 分子原文还写$\theta$而非$\theta_0$. 原式及随后“分母”一词保留.}
$$-2\log\lambda(\boldsymbol x)\approx\frac{(\theta-\hat\theta)^2}{-l''(\hat\theta\mid\boldsymbol x)},$$
其中使用$l'(\hat\theta\mid\boldsymbol x)=0$. 由于分母是观测信息量$\hat I_n(\hat\theta)$，且$\frac1n\hat I_n(\hat\theta)\to I(\theta_0)$，由\autoref{thm:10.1.12}和Slutsky定理（\autoref{thm:5.5.17}），得到$-2\log\lambda(\boldsymbol X)\to\chi_1^2$.
\end{proof}

\begin{example}[泊松LRT]\label{ex:10.3.2}
基于独立同分布于Poisson$(\lambda)$的观测$X_1,\ldots,X_n$，检验$H_0:\lambda=\lambda_0$对$H_1:\lambda\ne\lambda_0$，有
$$-2\log\lambda(\boldsymbol X)=-2\log\left(\frac{e^{-n\lambda_0}\lambda_0^{\sum x_i}}{e^{-n\hat\lambda}\hat\lambda^{\sum x_i}}\right)
=2n\bigl[(\lambda_0-\hat\lambda)-\hat\lambda\log(\lambda_0/\hat\lambda)\bigr],$$
其中$\hat\lambda=\sum x_i/n$是$\lambda$的MLE. 应用\autoref{thm:10.3.1}，当$-2\log\lambda(\boldsymbol x)>\chi_{1,\alpha}^2$时，在水平$\alpha$下拒绝$H_0$.

为了解渐近近似的准确程度，对检验统计量作一个小型模拟. 取$\lambda_0=5$、$n=25$，\autoref{fig:10.3.1}画出$-2\log\lambda(\boldsymbol x)$的10,000个值的直方图，以及$\chi_1^2$的pdf. 二者吻合似乎相当合理. 此外，\autoref{tab:10.3.5}比较模拟（精确）临界点与$\chi_1^2$（近似）临界点，显示它们非常接近.
\end{example}

\autoref{thm:10.3.1}可推广到原假设涉及参数向量的情形. 下面的推广不予证明，它至少对大样本保证\eqref{eq:10.3.1}成立. 完整讨论见\hyperref[bib:source311]{Stuart、Ord和Arnold（1999}，第22章）.

\sourcepageanchor{427}
\par\begin{figure}[htbp]\centering
\includegraphics[width=.73\textwidth]{assets/chapter10/fig-10.3.1.pdf}
\caption{$-2\log\lambda(\boldsymbol x)$的10,000个值的直方图，以及$\chi_1^2$的pdf；$\lambda_0=5$，$n=25$.}\label{fig:10.3.1}
\end{figure}\par
\par\begin{table}[htbp]\centering\setcounter{table}{4}
\caption{泊松LRT统计量的模拟（精确）百分位数与近似百分位数}\label{tab:10.3.5}
\begin{tabular}{lrrrr}\toprule
百分位&.80&.90&.95&.99\\\midrule
模拟&1.630&2.726&3.744&6.304\\
$\chi^2$&1.642&2.706&3.841&6.635\\\bottomrule
\end{tabular}
\end{table}\par

\begin{theorem}\label{thm:10.3.3}\index{SI0632@渐近}
设$X_1,\ldots,X_n$是来自pdf或pmf $f(x\mid\theta)$的随机样本. 在\hyperref[sec:10.6.2]{杂录10.6.2}的正则条件下，若$\theta\in\Theta_0$，则统计量$-2\log\lambda(\boldsymbol X)$的分布随着样本量$n\to\infty$收敛到卡方分布. 极限分布的自由度，是$\theta\in\Theta_0$所规定的自由参数个数与$\theta\in\Theta$所规定的自由参数个数之差.
\end{theorem}

当$\lambda(\boldsymbol X)$较小时拒绝$H_0:\theta\in\Theta_0$，等价于当$-2\log\lambda(\boldsymbol X)$较大时拒绝. 因此，
$$\text{拒绝$H_0$当且仅当}\quad-2\log\lambda(\boldsymbol X)\geq\chi_{\nu,\alpha}^2,$$
其中$\nu$为\autoref{thm:10.3.3}规定的自由度. 若$\theta\in\Theta_0$且样本量大，第一类错误概率近似为$\alpha$.\index{SI0509@渐近}\semanticindexanchor{509:427} 这样，样本量大时\eqref{eq:10.3.1}近似满足，便定义了渐近大小为$\alpha$的检验. 注意，定理实际上仅意味着
$$\lim_{n\to\infty}P_\theta(\text{拒绝$H_0$})=\alpha\quad\text{对每个$\theta\in\Theta_0$},$$
并不意味着$\sup_{\theta\in\Theta_0}P_\theta(\text{拒绝$H_0$})$收敛到$\alpha$. 渐近大小为$\alpha$的检验通常如此.

检验统计量自由度的计算通常很直接. 最常见的是，$\Theta$可表示为$q$维欧氏空间的子集，并包含$\mathbb R^q$中的一个开子集；$\Theta_0$可表示为$p$维欧氏空间的子集，并包含$\mathbb R^p$中的一个开子集，其中$p<q$. 则$q-p=\nu$是检验统计量的自由度.

\begin{example}[多项LRT]\label{ex:10.3.4}\index{SI0636@多项}
设$\theta=(p_1,p_2,p_3,p_4,p_5)$，其中$p_j$非负且和为1. 假设$X_1,\ldots,X_n$为独立同分布的离散随机变量，$P_\theta(X_i=j)=p_j$，$j=1,\ldots,5$. 因此$X_i$的pmf为$f(j\mid\theta)=p_j$，似然函数为
$$L(\theta\mid\boldsymbol x)=\prod_{i=1}^nf(x_i\mid\theta)=p_1^{y_1}p_2^{y_2}p_3^{y_3}p_4^{y_4}p_5^{y_5},$$
\sourcepageanchor{428}
其中$y_j$为$x_1,\ldots,x_n$中等于$j$的个数. 考虑检验
$$H_0:p_1=p_2=p_3\text{且}p_4=p_5\qquad\text{对}\qquad H_1:H_0\text{不成立}.$$
完整参数空间$\Theta$实际上是四维集合. 因为$p_5=1-p_1-p_2-p_3-p_4$，只有四个自由参数. 参数集合由
$$\sum_{j=1}^4p_j\leq1\quad\text{且}\quad p_j\geq0,\quad j=1,\ldots,4$$
定义，是包含$\mathbb R^4$中开子集的$\mathbb R^4$子集. 因此$q=4$. $H_0$规定的集合只有一个自由参数，因为一旦固定$p_1$，$0\leq p_1\leq1/3$，则$p_2=p_3$必须等于$p_1$，而$p_4=p_5$必须等于$(1-3p_1)/2$. 因而$p=1$，自由度$\nu=4-1=3$.

计算$\lambda(\boldsymbol x)$，必须确定$\theta$在$\Theta_0$和$\Theta$下的MLE. 令
$$\frac\partial{\partial p_j}\log L(\theta\mid\boldsymbol x)=0\quad\text{对每个$j=1,\ldots,4$},$$
并利用$p_5=1-p_1-p_2-p_3-p_4$及$y_5=n-y_1-y_2-y_3-y_4$，可验证$p_j$在$\Theta$下的MLE为$\hat p_j=y_j/n$. 在$H_0$下，似然函数化为
$$L(\theta\mid\boldsymbol x)=p_1^{y_1+y_2+y_3}\left(\frac{1-3p_1}2\right)^{y_4+y_5}.$$
同样，以令导数为零的通常方法可得，$H_0$下$p_1$的MLE为$\hat p_{10}=(y_1+y_2+y_3)/(3n)$. 此时$\hat p_{10}=\hat p_{20}=\hat p_{30}$，且$\hat p_{40}=\hat p_{50}=(1-3\hat p_{10})/2$. 将这些值和$\hat p_j$代入$L(\theta\mid\boldsymbol x)$，合并相同指数的项，得到
$$\begin{aligned}
\lambda(\boldsymbol x)&=\left(\frac{y_1+y_2+y_3}{3y_1}\right)^{y_1}\left(\frac{y_1+y_2+y_3}{3y_2}\right)^{y_2}\left(\frac{y_1+y_2+y_3}{3y_3}\right)^{y_3}\\
&\quad\times\left(\frac{y_4+y_5}{2y_4}\right)^{y_4}\left(\frac{y_4+y_5}{2y_5}\right)^{y_5}.
\end{aligned}$$
因此检验统计量为
\begin{equation}\label{eq:10.3.2}
-2\log\lambda(\boldsymbol x)=2\sum_{i=1}^5y_i\log\left(\frac{y_i}{m_i}\right),
\end{equation}
其中$m_1=m_2=m_3=(y_1+y_2+y_3)/3$，$m_4=m_5=(y_4+y_5)/2$. 渐近大小为$\alpha$的检验在$-2\log\lambda(\boldsymbol x)\geq\chi_{3,\alpha}^2$时拒绝$H_0$. 本例属于广泛使用似然比检验渐近理论的一大类检验问题.
\end{example}

\subsubsection{其他大样本检验}\label{sec:10.3.2}\index{SI0204@依分布}
构造大样本检验统计量的另一常用方法，基于具有渐近正态分布的估计量. 假设要检验关于实值参数$\theta$的假设，$W_n=W(X_1,\ldots,X_n)$是通过某种方法导出的、基于容量$n$样本的$\theta$的点估计量. 例如，$W_n$可以是$\theta$的MLE. 基于正态近似的近似检验，可按如下方式论证. 若$\sigma_n^2$表示$W_n$的方差，且能用某种形式的中心极限定理证明，当$n\to\infty$时，$(W_n-\theta)/\sigma_n$依分布收敛到标准正态随机变量，那么就可以将$(W_n-\theta)/\sigma_n$与$\mathrm n(0,1)$分布比较，由此得到近似检验的基础.

当然，上段论证有许多细节需要验证，但这一思想确实适用于许多情形. 例如，若$W_n$为MLE，可以用\autoref{thm:10.1.12}来确立上述论证. 注意，$W_n$的分布以及可能的$\sigma_n$值依赖于$\theta$. 因而，更正式地说，对每个固定的$\theta\in\Theta$，若采用$W_n$相应的分布和$\sigma_n$相应的值，则$(W_n-\theta)/\sigma_n$收敛到标准正态分布.
\sourcepageanchor{429}
若对每个$n$，$\sigma_n$是可计算的常数（可依赖$\theta$，但不依赖其他未知参数），便可能导出基于$(W_n-\theta)/\sigma_n$的检验.

在一些情形中，$\sigma_n$还依赖未知参数.\index{SI0979@Slutsky定理}\semanticindexanchor{979:429} 此时寻找$\sigma_n$的估计$S_n$，使$\sigma_n/S_n$依概率收敛到1. 再由Slutsky定理（像\autoref{ex:5.5.18}一样），推出$(W_n-\theta)/S_n$也依分布收敛到标准正态分布. 大样本检验可据此构造.

假设要检验双侧假设$H_0:\theta=\theta_0$对$H_1:\theta\ne\theta_0$. 近似检验可基于统计量$Z_n=(W_n-\theta_0)/S_n$，当且仅当$Z_n<-z_{\alpha/2}$或$Z_n>z_{\alpha/2}$时拒绝$H_0$. 若$H_0$成立，则$\theta=\theta_0$，$Z_n$依分布收敛到$Z\sim\mathrm n(0,1)$. 因此，第一类错误概率
$$P_{\theta_0}(Z_n<-z_{\alpha/2}\text{或}Z_n>z_{\alpha/2})\longrightarrow P(Z<-z_{\alpha/2}\text{或}Z>z_{\alpha/2})=\alpha,$$
所以它是渐近大小为$\alpha$的检验.

现考虑备择参数值$\theta\ne\theta_0$. 可写为
\begin{equation}\label{eq:10.3.3}
Z_n=\frac{W_n-\theta_0}{S_n}=\frac{W_n-\theta}{S_n}+\frac{\theta-\theta_0}{S_n}.
\end{equation}
无论$\theta$取何值，都有$(W_n-\theta)/S_n\to\mathrm n(0,1)$. 通常还满足$n\to\infty$时$\sigma_n\to0$. （回顾$\sigma_n=\operatorname{Var}W_n$，而估计量通常随着$n\to\infty$越来越精确.）\footnote{译注：原文括号内写$\sigma_n=\operatorname{Var}W_n$，与前文$\sigma_n^2$为方差的定义不符；应为$\sigma_n^2=\operatorname{Var}W_n$. 原文保留.} 因而$S_n$依概率收敛到0，$(\theta-\theta_0)/S_n$根据$\theta-\theta_0$为正或负，分别依概率收敛到$+\infty$或$-\infty$. 于是$Z_n$也依概率收敛到$+\infty$或$-\infty$，且
$$P_\theta(\text{拒绝$H_0$})=P_\theta(Z_n<-z_{\alpha/2}\text{或}Z_n>z_{\alpha/2})\longrightarrow1,\qquad n\to\infty.$$
这样就能构造渐近大小为$\alpha$、渐近功效为1的检验.

若要检验单侧假设$H_0:\theta\leq\theta_0$对$H_1:\theta>\theta_0$，也能构造类似检验. 仍使用统计量$Z_n=(W_n-\theta_0)/S_n$，当且仅当$Z_n>z_\alpha$时拒绝$H_0$. 与上面类似的推理表明，检验的功效函数在$\theta<\theta_0$、$\theta=\theta_0$、$\theta>\theta_0$时分别收敛到0、$\alpha$、1. 所以该检验也有合理的渐近功效性质.

一般地，Wald检验（Wald test）是基于如下形式统计量的检验：
$$Z_n=\frac{W_n-\theta_0}{S_n},$$
其中$\theta_0$是参数$\theta$的假设值，$W_n$是$\theta$的估计量，$S_n$是$W_n$的标准误，即$W_n$标准差的估计. 若$W_n$是$\theta$的MLE，则如\autoref{sec:10.1.3}所述，$1/\sqrt{I(W_n)}$是合理的标准误. 也常使用$1/\sqrt{\hat I(W_n)}$，其中
$$\hat I(W_n)=-\left.\frac{\partial^2}{\partial\theta^2}\log L(\theta\mid\boldsymbol X)\right|_{\theta=W_n}$$
是观测信息量（见\eqref{eq:10.1.7}）.

\begin{example}[大样本二项检验]\label{ex:10.3.5}\index{SI0486@伯努利}
设$X_1,\ldots,X_n$为Bernoulli$(p)$总体的随机样本. 考虑检验$H_0:p\leq p_0$对$H_1:p>p_0$，其中$0<p_0<1$为指定值. 容量$n$样本中$p$的MLE为$\hat p_n=\sum_{i=1}^nX_i/n$. 因为$\hat p_n$就是样本均值，中心极限定理适用，且说明对任意$0<p<1$，$(\hat p_n-p)/\sigma_n$收敛到标准正态随机变量. 这里$\sigma_n=\sqrt{p(1-p)/n}$，依赖未知参数$p$.
\sourcepageanchor{430}
$\sigma_n$的合理估计为$S_n=\sqrt{\hat p_n(1-\hat p_n)/n}$，而且可以证明（\autoref{exercise:5.32}）$\sigma_n/S_n$依概率收敛到1. 因而，对任意$0<p<1$，
$$\frac{\hat p_n-p}{\sqrt{\hat p_n(1-\hat p_n)/n}}\longrightarrow\mathrm n(0,1).$$
将$p$替换为$p_0$，即定义了Wald检验统计量$Z_n$；大样本Wald检验在$Z_n>z_\alpha$时拒绝$H_0$. 作为$\sigma_n$的另一种估计，容易验证$1/I(\hat p_n)=\hat p_n(1-\hat p_n)/n$. 所以用信息量导出$\hat p_n$的标准误，也得到相同统计量$Z_n$.

若要检验双侧假设$H_0:p=p_0$对$H_1:p\ne p_0$，其中$0<p_0<1$为指定值，上述策略仍适用. 但此时还有另一种近似检验. 由中心极限定理，对任意$0<p<1$，
$$\frac{\hat p_n-p}{\sqrt{p(1-p)/n}}\longrightarrow\mathrm n(0,1).$$
因此，若原假设成立，统计量
\begin{equation}\label{eq:10.3.4}
Z'_n=\frac{\hat p_n-p_0}{\sqrt{p_0(1-p_0)/n}}\sim\mathrm n(0,1)\qquad\text{（近似地）}.
\end{equation}
近似水平$\alpha$的检验在$|Z'_n|>z_{\alpha/2}$时拒绝$H_0$.

当两个检验均适用时，例如检验$H_0:p=p_0$时，并不清楚应优先选择哪个. 它们的（实际而非近似）功效函数彼此交叉，所以每个检验都在参数空间的某一部分更有功效. （\hyperref[bib:source139]{Ghosh，1979}，对这个问题提供了一些见解. \hyperref[bib:source254]{Robbins，1977}，以及\hyperref[bib:source105]{Eberhardt和Fligner，1977}，讨论了相关的双样本二项问题争论. \autoref{exercise:10.31}给出该问题的两个不同检验统计量.）

当然，比较功效函数时会受到一个事实的干扰：这些是近似检验，不一定保持水平$\alpha$. 连续性校正（\autoref{ex:3.3.2}）有助于解决这个问题. 许多情形下，使用连续性校正的近似程序是保守的，即保持其名义$\alpha$水平（见\autoref{ex:10.4.6}）.
\end{example}

\eqref{eq:10.3.4}是另一种有用的大样本检验——得分检验（score test）——的特例. 定义得分统计量为
$$S(\theta)=\frac\partial{\partial\theta}\log f(\boldsymbol X\mid\theta)=\frac\partial{\partial\theta}\log L(\theta\mid\boldsymbol X).$$
由\eqref{eq:7.3.8}，对所有$\theta$，$E_\theta S(\theta)=0$. 特别地，若检验$H_0:\theta=\theta_0$且$H_0$成立，则$S(\theta_0)$均值为零. 此外，由\eqref{eq:7.3.10}，
$$\operatorname{Var}_\theta S(\theta)=E_\theta\left(\left(\frac\partial{\partial\theta}\log L(\theta\mid\boldsymbol X)\right)^2\right)
=-E_\theta\left(\frac{\partial^2}{\partial\theta^2}\log L(\theta\mid\boldsymbol X)\right)=I(\theta);$$
信息量就是得分统计量的方差. 得分检验的检验统计量为
$$Z_S=S(\theta_0)/\sqrt{I(\theta_0)}.$$
若$H_0$成立，$Z_S$均值为0、方差为1. 由\autoref{thm:10.1.12}，在$H_0$成立时，$Z_S$收敛到标准正态随机变量. 因此，近似水平$\alpha$的得分检验在$|Z_S|>z_{\alpha/2}$时拒绝$H_0$. 若$H_0$为复合假设，则在$Z_S$中用假定$H_0$成立时$\theta$的估计$\hat\theta_0$替代$\theta_0$. 若$\hat\theta_0$是受限MLE，受限极大化可借助拉格朗日乘子完成. 因而得分检验有时称为拉格朗日乘子检验（Lagrange multiplier test）.

\sourcepageanchor{431}
\begin{example}[二项得分检验]\label{ex:10.3.6}
再次考虑\autoref{ex:10.3.5}的伯努利模型，检验$H_0:p=p_0$对$H_1:p\ne p_0$. 直接计算得到
$$S(p)=\frac{\hat p_n-p}{p(1-p)/n}\quad\text{且}\quad I(p)=\frac n{p(1-p)}.$$
所以得分统计量为
$$Z_S=\frac{S(p_0)}{\sqrt{I(p_0)}}=\frac{\hat p_n-p_0}{\sqrt{p_0(1-p_0)/n}},$$
与\eqref{eq:10.3.4}相同.
\end{example}

最后考虑的一类近似检验是稳健检验（robust tests，见\hyperref[sec:10.6.6]{杂录10.6.6}）. 在\autoref{sec:10.2}看到，若$X_1,X_2,\ldots,X_n$是来自位置族的独立同分布样本，$\hat\theta_M$为$M$估计量，则
\begin{equation}\label{eq:10.3.5}
\sqrt n(\hat\theta_M-\theta_0)\longrightarrow\mathrm n\bigl(0,\operatorname{Var}_{\theta_0}(\hat\theta_M)\bigr),
\end{equation}
其中$\operatorname{Var}_{\theta_0}(\hat\theta_M)=E_{\theta_0}\psi(X-\theta_0)^2/[E_{\theta_0}\psi'(X-\theta_0)]^2$为渐近方差. 因此，可以构造“广义”得分统计量
$$Z_{GS}=\sqrt n\frac{\hat\theta_M-\theta_0}{\sqrt{\operatorname{Var}_{\theta_0}(\hat\theta_M)}},$$
或广义Wald统计量
$$Z_{GW}=\sqrt n\frac{\hat\theta_M-\theta_0}{\sqrt{\widehat{\operatorname{Var}}_{\theta_0}(\hat\theta_M)}},$$
其中$\widehat{\operatorname{Var}}_{\theta_0}(\hat\theta_M)$可为任意相合估计量. 例如，可以使用标准误的自助法估计，或直接将估计量代入\eqref{eq:10.2.6}而使用
\begin{equation}\label{eq:10.3.6}
\widehat{\operatorname{Var}}_1(\hat\theta_M)=\frac{\frac1n\sum_{i=1}^n[\psi(x_i-\hat\theta_M)]^2}{\left[\frac1n\sum_{i=1}^n\psi'(x_i-\hat\theta_M)\right]^2}.
\end{equation}
方差估计的选择可能很重要；相关指导见\hyperref[bib:source41]{Boos（1992）}或\hyperref[bib:source62]{Carroll等（1995}，附录A.3）.

\begin{example}[基于Huber估计量的检验]\label{ex:10.3.7}
若$X_1,X_2,\ldots,X_n$独立同分布于pdf $f(x-\theta)$，其中$f$关于零对称，则采用\eqref{eq:10.2.2}的$\rho$函数和\eqref{eq:10.2.7}的$\psi$函数的Huber $M$估计量，其渐近方差为
\begin{equation}\label{eq:10.3.7}
\frac{\int_{-k}^kx^2f(x)\,dx+k^2P_0(|X|>k)}{[P_0(|X|\leq k)]^2}.
\end{equation}
因此，根据$M$估计量的渐近正态性，例如可以在$|Z_{GS}|>z_{\alpha/2}$时拒绝$H_0$，以水平$\alpha$检验$H_0:\theta=\theta_0$对$H_1:\theta\ne\theta_0$. 为更切合实际，考察使用估计标准误的近似检验. 采用统计量$Z_{GW}$，但以\eqref{eq:10.3.7}为方差估计的基础，即\footnote{译注：下面分母原文为“1减去区间内比例”的平方；由\eqref{eq:10.3.7}应是区间内比例本身的平方，或将指标改为区间外. 原式保留.}
\begin{equation}\label{eq:10.3.8}
\widehat{\operatorname{Var}}_2(\hat\theta_M)=
\frac{\frac1n\sum_{i=1}^n(x_i-\hat\theta_M)^2I(|x_i-\hat\theta_M|<k)+k^2\left(\frac1n\sum_{i=1}^nI(|x_i-\hat\theta_M|>k)\right)}{\left(1-\frac1n\sum_{i=1}^nI(|x_i-\hat\theta_M|<k)\right)^2}.
\end{equation}
\sourcepageanchor{432}
另外加入一个“朴素”检验$Z_N$，使用简单的方差估计
\begin{equation}\label{eq:10.3.9}
\widehat{\operatorname{Var}}_3(\hat\theta_M)=\frac1n\sum_{i=1}^n(x_i-\hat\theta_M)^2.
\end{equation}
这些检验表现如何？解析评价很困难，但下面的小型模拟显示，$z_{\alpha/2}$临界值通常太小（忽略了方差估计的变动），实际大小一般大于名义大小. 不过，在一系列分布中的表现相当一致，双指数分布的情形最好. （最后这一点并不十分意外，因为Huber估计量对具有指数尾的分布享有一种最优性性质，见\hyperref[bib:source171]{Huber，1981}，第4章.）
\end{example}

\par\begin{table}[htbp]\centering
\caption{指定参数值处，基于$Z_{GW}$与$Z_N$的名义$\alpha=.1$检验的功效；样本量$n=15$，10,000次模拟.}\label{tab:10.3.6}
\small\begin{tabular}{lrrrrrrrr}\toprule
&\multicolumn{8}{c}{基础pdf}\\\cmidrule(lr){2-9}
&\multicolumn{2}{c}{正态}&\multicolumn{2}{c}{$t_5$}&\multicolumn{2}{c}{逻辑斯蒂}&\multicolumn{2}{c}{双指数}\\
&$Z_{GW}$&$Z_N$&$Z_{GW}$&$Z_N$&$Z_{GW}$&$Z_N$&$Z_{GW}$&$Z_N$\\\midrule
$\theta_0$&0.16&0.16&0.14&0.13&0.15&0.15&0.11&0.09\\
$\theta_0+.25\sigma$&0.27&0.29&0.29&0.27&0.27&0.27&0.31&0.26\\
$\theta_0+.5\sigma$&0.58&0.60&0.65&0.63&0.59&0.60&0.70&0.64\\
$\theta_0+.75\sigma$&0.85&0.87&0.89&0.89&0.85&0.87&0.92&0.90\\
$\theta_0+1\sigma$&0.96&0.97&0.98&0.97&0.96&0.97&0.98&0.98\\
$\theta_0+2\sigma$&1.&1.&1.&1.&1.&1.&1.&1.\\\bottomrule
\end{tabular}
\end{table}\par



\subsection{区间估计}\label{sec:10.4}
像前两节一样，现在探讨置信集的若干近似和渐近版本. 目的仍是说明一些在更复杂情形中有用、能够得到某种答案的方法. 这里所得的答案几乎肯定不是最好的，但也肯定不是最坏的. 不过，在许多情形下，它们是我们所能做到的最好结果.

与前面一样，从基于MLE的近似开始.

\subsubsection{近似最大似然区间}\label{sec:10.4.1}
由\autoref{sec:10.1.2}的讨论，使用\autoref{thm:10.1.12}，有一般方法求得MLE的渐近分布. 因而也有一般方法构造置信区间.

若$X_1,\ldots,X_n$独立同分布于$f(x\mid\theta)$，$\hat\theta$是$\theta$的MLE，则由\hyperref[eq:10.1.7]{(7.4.1)}，函数$h(\hat\theta)$的方差可近似为\footnote{译注：原文引用公式(7.4.1)，但与下式直接对应的是\eqref{eq:10.1.7}. 保留原引用，待对照作者勘误核实.}
$$\widehat{\operatorname{Var}}(h(\hat\theta)\mid\theta)\approx\frac{[h'(\theta)]^2|_{\theta=\hat\theta}}{-\frac{\partial^2}{\partial\theta^2}\log L(\theta\mid\boldsymbol x)|_{\theta=\hat\theta}}.$$

\sourcepageanchor{433}
现在，对固定但任意的$\theta$，关心
$$\frac{h(\hat\theta)-h(\theta)}{\sqrt{\widehat{\operatorname{Var}}(h(\hat\theta)\mid\theta)}}$$
的渐近分布. 由\autoref{thm:10.1.12}及Slutsky定理（\autoref{thm:5.5.17}；见\autoref{exercise:10.33}），
$$\frac{h(\hat\theta)-h(\theta)}{\sqrt{\widehat{\operatorname{Var}}(h(\hat\theta)\mid\theta)}}\longrightarrow\mathrm n(0,1),$$
从而给出近似置信区间
$$h(\hat\theta)-z_{\alpha/2}\sqrt{\widehat{\operatorname{Var}}(h(\hat\theta)\mid\theta)}\leq h(\theta)\leq h(\hat\theta)+z_{\alpha/2}\sqrt{\widehat{\operatorname{Var}}(h(\hat\theta)\mid\theta)}.$$

\begin{example}[续\autoref{ex:10.1.14}]\label{ex:10.4.1}\index{SI0138@二项}
设$X_1,\ldots,X_n$为Bernoulli$(p)$总体的随机样本. 已经看到，可以用MLE $\hat p/(1-\hat p)$估计优势比$p/(1-p)$，且这个估计的近似方差为
$$\widehat{\operatorname{Var}}\left(\frac{\hat p}{1-\hat p}\right)\approx\frac{\hat p}{n(1-\hat p)^3}.$$
因此，可以构造近似置信区间
$$\frac{\hat p}{1-\hat p}-z_{\alpha/2}\sqrt{\widehat{\operatorname{Var}}\left(\frac{\hat p}{1-\hat p}\right)}
\leq\frac p{1-p}\leq\frac{\hat p}{1-\hat p}+z_{\alpha/2}\sqrt{\widehat{\operatorname{Var}}\left(\frac{\hat p}{1-\hat p}\right)}.$$
\end{example}

一种限制较多的似然近似形式基于得分统计量（\autoref{sec:10.3.2}）；但在适用时，它能给出更好的区间. 随机量
\begin{equation}\label{eq:10.4.1}
Q(\boldsymbol X\mid\theta)=\frac{\frac\partial{\partial\theta}\log L(\theta\mid\boldsymbol X)}{\sqrt{-E_\theta\left(\frac{\partial^2}{\partial\theta^2}\log L(\theta\mid\boldsymbol X)\right)}}
\end{equation}
在$n\to\infty$时渐近服从$\mathrm n(0,1)$分布. 因此集合
\begin{equation}\label{eq:10.4.2}\index{SI0139@最大似然}
\{\theta:|Q(\boldsymbol x\mid\theta)|\leq z_{\alpha/2}\}
\end{equation}
是近似$1-\alpha$置信集. 注意，应用\autoref{sec:7.3.2}的结果，有
$$E_\theta(Q(\boldsymbol X\mid\theta))=\frac{E_\theta\left(\frac\partial{\partial\theta}\log L(\theta\mid\boldsymbol X)\right)}{\sqrt{-E_\theta\left(\frac{\partial^2}{\partial\theta^2}\log L(\theta\mid\boldsymbol X)\right)}}=0,$$
且
\begin{equation}\label{eq:10.4.3}
\operatorname{Var}_\theta(Q(\boldsymbol X\mid\theta))=\frac{\operatorname{Var}_\theta\left(\frac\partial{\partial\theta}\log L(\theta\mid\boldsymbol X)\right)}{-E_\theta\left(\frac{\partial^2}{\partial\theta^2}\log L(\theta\mid\boldsymbol X)\right)}=1.
\end{equation}
所以这个近似精确匹配$\mathrm n(0,1)$随机变量的前两阶矩. \hyperref[bib:source330]{Wilks（1938）}证明，这些区间具有渐近最优性：渐近地，它们是某类区间中最短的.

当然，这些区间并非完全一般，未必总能适用于函数$h(\theta)$. 必须能够把\eqref{eq:10.4.2}表示为$h(\theta)$的函数.

\sourcepageanchor{434}
\begin{example}[二项得分区间]\label{ex:10.4.2}
再次使用二项例子，若$Y=\sum_{i=1}^nX_i$，每个$X_i$都是独立的Bernoulli$(p)$随机变量，则有
$$\begin{aligned}
Q(Y\mid p)&=\frac{\frac\partial{\partial p}\log L(p\mid Y)}{\sqrt{-E_p\left(\frac{\partial^2}{\partial p^2}\log L(p\mid Y)\right)}}\\
&=\frac{\frac yp-\frac{n-y}{1-p}}{\sqrt{\frac n{p(1-p)}}}
=\frac{\hat p-p}{\sqrt{p(1-p)/n}},
\end{aligned}$$
其中$\hat p=y/n$. 利用\eqref{eq:10.4.2}，得到近似$1-\alpha$置信区间
\begin{equation}\label{eq:10.4.4}
\left\{p:\left|\frac{\hat p-p}{\sqrt{p(1-p)/n}}\right|\leq z_{\alpha/2}\right\}.
\end{equation}
这是反演得分统计量（\autoref{ex:10.3.6}）所得的区间. 计算它需要解一个关于$p$的二次方程，细节见\autoref{ex:10.4.6}.
\end{example}

在\autoref{sec:10.3}中，根据$-2\log\lambda(\boldsymbol X)$渐近服从卡方分布这一事实，导出了另一种似然检验. 这提示：若$X_1,\ldots,X_n$独立同分布于$f(x\mid\theta)$，$\hat\theta$是$\theta$的MLE，则集合
\begin{equation}\label{eq:10.4.5}
\left\{\theta:-2\log\left(\frac{L(\theta\mid\boldsymbol x)}{L(\hat\theta\mid\boldsymbol x)}\right)\leq\chi_{1,\alpha}^2\right\}
\end{equation}
是近似$1-\alpha$置信区间. 确实如此，这又给出了一个近似似然区间.

当然，\eqref{eq:10.4.5}就是最初通过反演LRT统计量导出的最高似然区域\eqref{eq:9.2.7}. 不过，现在有了自动赋予近似置信水平的方法.

\begin{example}[二项LRT区间]\label{ex:10.4.3}
若$Y=\sum_{i=1}^nX_i$，每个$X_i$都是独立的Bernoulli$(p)$随机变量，则近似$1-\alpha$置信集为
$$\left\{p:-2\log\left(\frac{p^y(1-p)^{n-y}}{\hat p^y(1-\hat p)^{n-y}}\right)\leq\chi_{1,\alpha}^2\right\}.$$
\autoref{ex:10.4.7}将比较这个置信集与基于得分检验和Wald检验的区间.
\end{example}

\subsubsection{其他大样本区间}\label{sec:10.4.2}
多数近似置信区间基于寻找近似（或渐近）枢轴量，或者反演近似水平$\alpha$的检验统计量. 若有任意统计量$W$、$V$与参数$\theta$，使$n\to\infty$时
$$\frac{W-\theta}V\longrightarrow\mathrm n(0,1),$$
则可以构造$\theta$的近似置信区间
$$W-z_{\alpha/2}V\leq\theta\leq W+z_{\alpha/2}V,$$
本质上是Wald型区间. 直接应用中心极限定理与Slutsky定理，通常就会给出近似置信区间. （注意，上一小节的近似最大似然区间均体现这一策略.）

\sourcepageanchor{435}
\begin{example}[近似区间]\label{ex:10.4.4}\index{SI0137@近似}
若$X_1,\ldots,X_n$独立同分布，均值为$\mu$、方差为$\sigma^2$，则由中心极限定理，
$$\frac{\bar X-\mu}{\sigma/\sqrt n}\longrightarrow\mathrm n(0,1).$$
此外，由Slutsky定理，若$S^2$依概率收敛到$\sigma^2$，则
$$\frac{\bar X-\mu}{S/\sqrt n}\longrightarrow\mathrm n(0,1),$$
给出近似$1-\alpha$置信区间
\begin{equation}\label{eq:10.4.6}
\bar x-z_{\alpha/2}s/\sqrt n\leq\mu\leq\bar x+z_{\alpha/2}s/\sqrt n.
\end{equation}
为了解近似有多好，对多种pdf进行小型模拟，计算近似区间的精确覆盖概率. 注意，因为区间基于枢轴量，覆盖概率不依赖参数值；它是常数，因此就是置信系数. \autoref{tab:10.4.7}显示，即使样本量小至$n=15$，枢轴置信区间也有合理表现，但显然达不到名义置信系数. 这无疑是由于使用$z_{\alpha/2}$临界值过于乐观，没有计入$S$的变动. 样本量增大时，近似将改善.

\par\begin{table}[htbp]\centering\setcounter{table}{6}
\caption{枢轴区间\eqref{eq:10.4.6}的置信系数，$n=15$，基于10,000次模拟}\label{tab:10.4.7}
\begin{tabular}{l|rrrr}\toprule
&\multicolumn{4}{c}{基础pdf}\\\cmidrule(lr){2-5}
名义水平&正态&$t_5$&逻辑斯蒂&双指数\\\midrule
$1-\alpha=.90$&0.879&0.864&0.880&0.876\\
$1-\alpha=.95$&0.931&0.924&0.931&0.933\\\bottomrule
\end{tabular}
\end{table}\par
\end{example}

上例无需指定抽样分布形式，就能得到近似置信区间. 当指定分布形式时，应能做得更好.

\begin{example}[近似泊松区间]\label{ex:10.4.5}\index{SI0141@泊松}
若$X_1,\ldots,X_n$独立同分布于Poisson$(\lambda)$，则知道
$$\frac{\bar X-\lambda}{S/\sqrt n}\longrightarrow\mathrm n(0,1).$$
不过，即使不是从泊松总体抽样，这也成立. 利用泊松假设，知道$\operatorname{Var}(X)=\lambda=EX$，且$\bar X$是$\lambda$的良好估计量（\autoref{sec:7}）. 因此利用泊松假设，还能由
$$\frac{\bar X-\lambda}{\sqrt{\bar X/n}}\longrightarrow\mathrm n(0,1)$$
得到近似置信区间，它是反演Wald检验所得的区间. 还可以用另一种方式利用泊松假设. 因为$\operatorname{Var}(X)=\lambda$，所以
$$\frac{\bar X-\lambda}{\sqrt{\lambda/n}}\longrightarrow\mathrm n(0,1),$$
由此得到与得分检验对应的区间，也就是\eqref{eq:10.4.2}的似然区间；按\hyperref[bib:source330]{Wilks（1938）}的结果，它是最优的（见\autoref{exercise:10.40}）.
\end{example}

\sourcepageanchor{436}
一般而言，合理的经验法则是，在近似中尽可能少用估计值、多用参数. 理由很简单：参数固定，不给近似引入额外变动，而每个统计量都会带来更多变动.

\begin{example}[再论二项得分区间]\label{ex:10.4.6}
对Bernoulli$(p)$总体的随机样本$X_1,\ldots,X_n$，在\autoref{ex:10.3.5}中看到，当$n\to\infty$时，
$$\frac{\hat p-p}{\sqrt{\hat p(1-\hat p)/n}}\quad\text{与}\quad\frac{\hat p-p}{\sqrt{p(1-p)/n}}$$
均依分布收敛到标准正态随机变量，其中$\hat p=\sum x_i/n$. 在\autoref{ex:10.3.5}中看到，可以基于任一近似建立检验，前者为Wald检验，后者为得分检验. 也知道可以用任一近似构造$p$的置信区间. 不过，得分检验近似使用较少的统计量和较多的参数值，会给出\autoref{ex:10.4.2}中的区间\eqref{eq:10.4.4}，它是渐近最优的. 即
$$\left\{p:\left|\frac{\hat p-p}{\sqrt{p(1-p)/n}}\right|\leq z_{\alpha/2}\right\}$$
是更好的近似区间.

这个区间的样子并不一目了然，但可以显式解出其取值集合. 两边平方并整理，就要寻找满足
$$\left\{p:(\hat p-p)^2\leq z_{\alpha/2}^2\frac{p(1-p)}n\right\}$$
的$p$值. 这是关于$p$的二次不等式，进一步整理可写成熟悉的形式：
$$\left\{p:\left(1+\frac{z_{\alpha/2}^2}n\right)p^2-\left(2\hat p+\frac{z_{\alpha/2}^2}n\right)p+\hat p^2\leq0\right\}.$$
由于二次项系数为正，抛物线开口向上，因此当$p$位于两个根之间时不等式成立. 两个根为
\begin{equation}\label{eq:10.4.7}
\frac{2\hat p+z_{\alpha/2}^2/n\ \pm\sqrt{(2\hat p+z_{\alpha/2}^2/n)^2-4\hat p^2(1+z_{\alpha/2}^2/n)}}{2(1+z_{\alpha/2}^2/n)},
\end{equation}
它们定义$p$的置信区间端点. 虽然根的表达式有些繁琐，这实际上是$p$的一个很好区间. 不过，使用连续性校正（\autoref{ex:3.3.2}）还可进一步改进. 为此，需要分别解两个二次不等式（见\autoref{exercise:10.45}）：
$$\begin{aligned}
\left|\frac{\hat p+\frac1{2n}-p}{\sqrt{p(1-p)/n}}\right|&\leq z_{\alpha/2},&&\text{（较大根为区间上端点）},\\
\left|\frac{\hat p-\frac1{2n}-p}{\sqrt{p(1-p)/n}}\right|&\leq z_{\alpha/2}.&&\text{（较小根为区间下端点）}.
\end{aligned}$$
在端点处有显然的修正：若$\sum x_i=0$，则区间下端点取0；若$\sum x_i=n$，则区间上端点取1. 一些良好近似见\hyperref[bib:source39]{Blyth（1986）}.
\end{example}

\sourcepageanchor{437}
现在已经看到二项比例的三个区间：基于Wald统计量、得分统计量的区间，以及\autoref{ex:10.4.3}的LRT区间. 通常最不优先选择Wald区间，不过比较三者仍很有意思.

\begin{example}[二项区间的比较]\label{ex:10.4.7}
对$Y=\sum_{i=1}^nX_i$，其中$X_1,\ldots,X_n$独立同分布于Bernoulli$(p)$总体，Wald区间为
\begin{equation}\label{eq:10.4.8}
\hat p-z_{\alpha/2}\sqrt{\frac{\hat p(1-\hat p)}n}\leq p\leq\hat p+z_{\alpha/2}\sqrt{\frac{\hat p(1-\hat p)}n}.
\end{equation}
带连续性校正的得分区间见\autoref{ex:10.4.6}，近似LRT见\autoref{ex:10.4.3}. 为比较它们，来看一个例子.

\par\begin{figure}[htbp]\centering
\includegraphics[width=.9\textwidth]{assets/chapter10/fig-10.4.1.pdf}
\caption{LRT程序（实线）、得分程序（长虚线）和修正Wald程序（短虚线）给出的二项比例区间.}\label{fig:10.4.1}
\end{figure}\par

对$n=12$，\autoref{fig:10.4.1}显示三种程序的区间实现值. LRT程序给出的区间最短，得分区间最长. 在这幅图中，对Wald区间作了两项修正. 首先，当$y=0$时，未修正区间为$(0,0)$，故将上端点改为$1-(\alpha/2)^{1/n}$；在$y=n$时对下端点作相似修正. 另外，Wald区间端点有时会超出$[0,1]$，这些部分已被截断.

得分区间较长，反映在其覆盖概率更高. 实际上，三者中只有得分区间保持覆盖概率在$.9$以上，因此只有它的置信系数为$.9$. LRT与Wald区间似乎太短，覆盖概率远低于$.9$，难以接受. 当然，随着$n$增大，它们的表现会改善. 其他选择见\hyperref[sec:9.5.2]{杂录9.5.2}.

所以，经过连续性校正的得分区间虽然更长，但似乎是小$n$时应选择的区间（另一种选择见\autoref{exercise:10.44}）. LRT和Wald程序在小$n$时产生的区间实在太短，Wald区间还存在端点问题.
\end{example}

像\autoref{sec:10.3.2}一样，简要考察基于稳健估计量的区间.

\begin{example}[基于Huber估计量的区间]\label{ex:10.4.8}
通过与\autoref{ex:10.3.7}类似的推导，可以构造基于Huber $M$估计量的渐近置信区间. 若$X_1,X_2,\ldots,X_n$独立同分布于pdf $f(x-\theta)$，其中$f$关于零对称，则$\theta$的近似区间为
\sourcepageanchor{438}
\par\begin{figure}[htbp]\centering
\includegraphics[width=.9\textwidth]{assets/chapter10/fig-10.4.2.pdf}
\caption{二项比例的名义$.9$置信程序的覆盖概率：LRT程序（细实线，灰色阴影）、得分程序（虚线）和修正Wald程序（粗实线）.}\label{fig:10.4.2}
\end{figure}\par
$$\hat\theta_M\pm z_{\alpha/2}\sqrt{\frac{\operatorname{Var}(\hat\theta_M)}n},$$
其中$\operatorname{Var}(\hat\theta_M)$由\eqref{eq:10.3.7}给出. 现在用估计\eqref{eq:10.3.8}与\eqref{eq:10.3.9}替换$\operatorname{Var}(\hat\theta_M)$，得到Wald型区间. 为评价它们，制作与\autoref{tab:10.4.7}相似的表. 有趣的是，除双指数分布外，\autoref{tab:10.4.8}中的区间表现都不如\autoref{tab:10.4.7}中基于通常均值与方差的区间. 对此没有很好的解释，只能再次归咎于$z_{\alpha/2}$临界值过于乐观.
\end{example}

\par\begin{table}[htbp]\centering
\caption{基于Huber $M$估计量、名义$1-\alpha=.9$的区间置信系数；$n=15$，基于10,000次模拟.}\label{tab:10.4.8}
\begin{tabular}{l|rrrr}\toprule
&\multicolumn{4}{c}{基础pdf}\\\cmidrule(lr){2-5}
名义水平&正态&$t_5$&逻辑斯蒂&双指数\\\midrule
方差估计\eqref{eq:10.3.8}&.844&.856&.855&.889\\
方差估计\eqref{eq:10.3.9}&.837&.867&.855&.910\\\bottomrule
\end{tabular}
\end{table}\par

迄今提到的近似都基于令$n\to\infty$. 不过，还有其他适合使用近似区间的情形. 在\autoref{ex:9.2.17}中，需要参数趋于无穷时的近似. 另一种情形见\autoref{ex:2.3.13}：对某些参数配置，泊松分布可以近似二项分布. 这提示，如果认为这种参数配置很可能成立，就能以泊松分布为基础构造近似二项区间. 本着这种想法，介绍下面这个略不寻常的例子.

\sourcepageanchor{439}
\begin{example}[负二项区间]\label{ex:10.4.9}\index{SI0140@负二项}
设$X_1,\ldots,X_n$独立同分布于negative binomial$(r,p)$. 假设$r$已知，希望得到$p$的置信区间. 利用$Y=\sum X_i\sim\operatorname{negative\ binomial}(nr,p)$，可以用多种方法构造区间. 利用二项分布与$F$分布关系的一个变形，可构造精确置信区间（\autoref{exercise:9.22}），也可使用正态近似（\autoref{exercise:10.41}）. 还有另一种近似，不依赖$n$大，而依赖$p$小.

\autoref{exercise:2.38}证明，当$p\to0$时，
$$2pY\longrightarrow\chi_{2nr}^2\quad\text{依分布}.$$
所以，对小$p$，$2pY$是枢轴量！据此可构造适用于小$p$的枢轴$1-\alpha$置信区间：
$$\left\{p:\frac{\chi_{2nr,1-\alpha/2}^2}{2y}\leq p\leq\frac{\chi_{2nr,\alpha/2}^2}{2y}\right\}.$$
细节见\autoref{exercise:10.47}.
\end{example}



\subsection{习题}\label{sec:10.5}

\begin{exercise}\label{exercise:10.1}
从pdf为
$$f(x\mid\theta)=\frac12(1+\theta x),\qquad-1<x<1,\quad-1<\theta<1$$
的总体中抽取随机样本$X_1,\ldots,X_n$. 求$\theta$的一个相合估计量，并证明其相合性.
\end{exercise}

\begin{exercise}\label{exercise:10.2}
证明\autoref{thm:10.1.5}.
\end{exercise}

\begin{exercise}\label{exercise:10.3}
从$\mathrm n(\theta,\theta)$总体中抽取随机样本$X_1,\ldots,X_n$，其中$\theta>0$.
\begin{enumerate}[label=(\alph*)]
\item 证明$\theta$的MLE $\hat\theta$是二次方程$\theta^2+\theta-W=0$的根，其中$W=(1/n)\sum_{i=1}^nX_i^2$，并确定哪个根等于MLE.
\item 用\autoref{sec:10.1.3}的技术求$\hat\theta$的近似方差.
\end{enumerate}
\end{exercise}

\begin{exercise}\label{exercise:10.4}
\autoref{exercise:7.19}模型的一个变形是令随机变量$Y_1,\ldots,Y_n$满足
$$Y_i=\beta X_i+\varepsilon_i,\qquad i=1,\ldots,n,$$
其中$X_1,\ldots,X_n$为独立的$\mathrm n(\mu,\tau^2)$随机变量，$\varepsilon_1,\ldots,\varepsilon_n$独立同分布于$\mathrm n(0,\sigma^2)$，且$X$与$\varepsilon$相互独立. 精确方差计算相当困难，因此可以采用近似. 用$\mu,\tau^2,\sigma^2$表示以下各量的近似均值与方差：
\begin{enumerate}[label=(\alph*)]
\item $\sum X_iY_i/\sum X_i^2$.
\item $\sum Y_i/\sum X_i$.
\item $\sum(Y_i/X_i)/n$.
\end{enumerate}
\end{exercise}

\begin{exercise}\label{exercise:10.5}
对\autoref{ex:10.1.8}的情形，令$T_n=\sqrt n/\bar X_n$，证明：
\begin{enumerate}[label=(\alph*)]
\item $\operatorname{Var}(T_n)=\infty$.
\item 若$\mu\ne0$且从样本空间删除区间$(-\delta,\delta)$，则$\operatorname{Var}(T_n)<\infty$.
\item 若$\mu\ne0$，区间$(-\delta,\delta)$的概率含量趋于0.\footnote{译注：此处需限制$\delta$相对$|\mu|$的大小；若区间包含$\mu$，则$\bar X_n$落在其中的概率反而趋于1. 原题条件照录，待核实所指样本空间及$\delta$约定.}
\end{enumerate}
\end{exercise}

\begin{exercise}\label{exercise:10.6}
对\autoref{ex:10.1.10}的情形，证明：
\begin{enumerate}[label=(\alph*)]
\item $EY_n=0$，且$\operatorname{Var}(Y_n)=p_n+(1-p_n)\sigma_n^2$；
\item $P(Y_n<a)\to P(Z<a)$，因而$Y_n\to\mathrm n(0,1)$（回顾$p_n\to1$、$\sigma_n\to\infty$且$(1-p_n)\sigma_n^2\to\infty$）.
\end{enumerate}
\end{exercise}

\begin{exercise}\label{exercise:10.7}
\autoref{thm:10.1.12}的证明说明了MLE $\hat\theta$是$\theta$的渐近有效估计量. 证明：若$\tau(\theta)$是$\theta$的连续函数，则$\tau(\hat\theta)$是$\tau(\theta)$的相合且渐近有效的估计量.
\end{exercise}

\sourcepageanchor{440}
\begin{exercise}\label{exercise:10.8}
通过建立\eqref{eq:10.1.6}中的两个收敛结果，完成\autoref{thm:10.1.6}的证明.\footnote{译注：原题引\autoref{thm:10.1.6}，但\eqref{eq:10.1.6}位于\autoref{thm:10.1.12}的证明中. 下列导数应在$\theta=\theta_0$处求值，原题未显式写出该代入.}
\begin{enumerate}[label=(\alph*)]
\item 证明
$$\frac1{\sqrt n}l'(\theta_0\mid\boldsymbol X)=\sqrt n\left[\frac1n\sum_iW_i\right],$$
其中$W_i=\frac{\frac d{d\theta}f(X_i\mid\theta)}{f(X_i\mid\theta)}$均值为0、方差为$I(\theta_0)$. 再用中心极限定理建立向$\mathrm n[0,I(\theta_0)]$的收敛.
\item 证明
$$-\frac1n l''(\theta_0\mid\boldsymbol X)=\frac1n\sum_iW_i^2-\frac1n\sum_i\frac{\frac{d^2}{d\theta^2}f(X_i\mid\theta)}{f(X_i\mid\theta)},$$
且第一部分均值为$I(\theta_0)$，第二部分均值为0. 应用弱大数定律.
\end{enumerate}
\end{exercise}

\begin{exercise}\label{exercise:10.9}
假设$X_1,\ldots,X_n$独立同分布于Poisson$(\lambda)$. 求以下各量的最佳无偏估计量：
\begin{enumerate}[label=(\alph*)]
\item\label{item:10.9.a} $e^{-\lambda}$，即$X=0$的概率.
\item\label{item:10.9.b} $\lambda e^{-\lambda}$，即$X=1$的概率.
\item 对\hyperref[item:10.9.a]{(a)}和\hyperref[item:10.9.b]{(b)}的最佳无偏估计量，计算相对于MLE的渐近相对效率. 你更倾向于哪些估计量？为什么？
\item 测量暴露于某种化合物的微生物突变率，可以对该化合物可能的致癌性进行初步检验. 实验者将化合物放入15个培养皿，记录到以下突变菌落数：
$$10,7,8,13,8,9,5,7,6,8,3,6,6,3,5.$$
估计没有突变菌落出现的概率$e^{-\lambda}$，以及出现一个突变菌落的概率$\lambda e^{-\lambda}$. 分别计算最佳无偏估计与MLE.
\end{enumerate}
\end{exercise}

\begin{exercise}\label{exercise:10.10}
继续\autoref{ex:10.1.14}的计算，那里考察了$p(1-p)$估计量的性质.\footnote{译注：直接讨论$p(1-p)$估计量的是\autoref{ex:10.1.15}，原题所引\autoref{ex:10.1.14}保留.}
\begin{enumerate}[label=(\alph*)]
\item 证明若$p\ne1/2$，MLE $\hat p(1-\hat p)$渐近有效.
\item 若$p=1/2$，用\autoref{thm:5.5.26}求$\hat p(1-\hat p)$的极限分布.
\item 计算$\operatorname{Var}[\hat p(1-\hat p)]$的精确表达式. 近似失效的原因是否更清楚了？
\end{enumerate}
\end{exercise}

\begin{exercise}\label{exercise:10.11}
本题考察\autoref{ex:10.1.18}计算的若干细节与推广.
\begin{enumerate}[label=(\alph*)]
\item 复现\autoref{fig:10.1.1}，计算$\beta$已知时的ARE. （可沿用\autoref{ex:12.6.7}的计算，或自行编程.）
\item 验证，无论$\beta$已知还是未知，$\operatorname{ARE}(\bar X,\hat\mu)$的比较都相同.
\item 对$\mu$已知时$\beta$的估计，证明矩估计与最大似然估计相同. （使用$(\alpha,\beta)$参数化可能更容易.）
\item 对$\mu$未知时$\beta$的估计，矩估计与最大似然估计不同. 用渐近相对效率比较它们，并作出类似\autoref{fig:10.1.1}的图，不同曲线对应不同$\mu$值.
\end{enumerate}
\end{exercise}

\begin{exercise}\label{exercise:10.12}
验证\hyperref[sec:10.6.1]{杂录10.6.1}的超有效估计量$d_n$渐近正态，其方差在$\theta\ne0$时为$v(\theta)=1$，在$\theta=0$时为$v(\theta)=a^2$. （超有效估计量的更多内容见\hyperref[bib:source200]{Lehmann和Casella，1998}，第6.2节.）
\end{exercise}

\needspace{6\baselineskip}
\begin{exercise}\label{exercise:10.13}
参照\autoref{ex:10.1.19}：
\begin{enumerate}[label=(\alph*)]
\item 验证样本$2,4,9,12$的自助法均值和方差分别为6.75与3.94.
\item\label{item:10.13.b} 验证6.75是原样本的均值.
\item\label{item:10.13.c} 验证，当除以$n$而非$n-1$时，均值的自助法方差与均值方差的通常估计相同.
\item 说明如何利用$\binom{4+4-1}{4}=35$个不同的可能重抽样计算自助法均值和标准误.
\item 对一般样本$X_1,X_2,\ldots,X_n$证明\hyperref[item:10.13.b]{(b)}和\hyperref[item:10.13.c]{(c)}.
\end{enumerate}
\end{exercise}

\sourcepageanchor{441}
\begin{exercise}\label{exercise:10.14}
在下列各情形中考察参数与非参数自助法. 比较估计结果，讨论这些方法的优缺点.
\begin{enumerate}[label=(\alph*)]
\item 参照\autoref{ex:10.1.22}，用非参数自助法估计$S^2$的方差.
\item 在\autoref{ex:5.6.6}中，实质上已经对泊松样本的$S^2$分布进行了参数自助法. 使用非参数自助法给出该分布的另一幅直方图.
\item 在\autoref{ex:10.1.18}中考察了伽马均值的估计. 假设来自gamma$(\alpha,\beta)$分布的随机样本为
$$0.28,0.98,1.36,1.38,2.4,7.42.$$
使用最大似然法和自助法估计分布的均值与方差.
\end{enumerate}
\end{exercise}

\begin{exercise}\label{exercise:10.15}\leavevmode
\begin{enumerate}[label=(\alph*)]
\item 证明当$B\to\infty$时，\eqref{eq:10.1.11}中的$\operatorname{Var}_B^*(\hat\theta)$收敛到\eqref{eq:10.1.10}中的$\operatorname{Var}^*(\hat\theta)$.
\item 固定$B$，对$i=1,2,\ldots$，计算自助法方差$\operatorname{Var}_{B_i}^*(\hat\theta)$. 用大数定律证明当$m\to\infty$时，$(1/m)\sum_{i=1}^m\operatorname{Var}_{B_i}^*(\hat\theta)\to\operatorname{Var}^*(\hat\theta)$.
\end{enumerate}
\end{exercise}

\begin{exercise}\label{exercise:10.16}
对\autoref{ex:10.1.21}的情形，若观测到$\hat p=1/2$，可以使用\autoref{thm:5.5.26}的方差估计. 证明这个方差估计等于$2[\operatorname{Var}(\hat p)]^2$.\footnote{译注：\hyperref[itm:exercise:10.16:a]{(a)}的$.00007$对应一阶Delta近似，而非题干所列二阶形式；\hyperref[itm:exercise:10.16:b]{(b)}的$.00529$也与$n=24,p=11/24$下对二项概率逐项求和所得方差（约$.000271$）不符. 原题数值保留，须核查原书的数值计算或统计量表述.}
\begin{enumerate}[label=(\alph*)]
\item\phantomsection\label{itm:exercise:10.16:a} 若观测到$\hat p=11/24$，验证此方差估计为$.00007$.
\item\phantomsection\label{itm:exercise:10.16:b} 使用模拟，计算$n=24$且$p=11/24$时$\hat p(1-\hat p)$的“精确方差”. 验证它等于$.00529$.
\item 你认为Delta方法在此情形下为何如此差？二阶Delta方法会更好吗？自助法估计又如何？
\end{enumerate}
\end{exercise}

\begin{exercise}\label{exercise:10.17}
\hyperref[bib:source108]{Efron（1982）}分析法学院录取数据，考察LSAT（法学院入学考试，Law School Admission Test）与第一年GPA（平均绩点，Grade Point Average）之间的相关性. 对15所法学院中的每所，得到一对数据（平均LSAT，平均GPA）.\footnote{译注：原表第三对数据写$(5.58,2.81)$，第二行第一对写$(580.3.07)$，均照录；可能分别为$(558,2.81)$与$(580,3.07)$，需与数据来源核实. \hyperref[itm:exercise:10.17:c]{(c)}原文写GRE，但本题分析的是GPA；亦照录.}
\begin{center}\small
\begin{tabular}{ccccc}
$(576,3.39)$&$(635,3.30)$&$(5.58,2.81)$&$(578,3.03)$&$(666,3.44)$\\
$(580.3.07)$&$(555,3.00)$&$(661,3.43)$&$(651,3.36)$&$(605,3.13)$\\
$(653,3.12)$&$(575,2.74)$&$(545,2.76)$&$(572,2.88)$&$(594,2.96)$
\end{tabular}
\end{center}
\begin{enumerate}[label=(\alph*)]
\item 计算LSAT与GPA之间的相关系数.
\item 使用非参数自助法估计相关系数的标准差，取$B=1000$次重抽样，并画出其直方图.
\item\phantomsection\label{itm:exercise:10.17:c} 使用参数自助法估计相关系数的标准差. 假设(LSAT, GRE)服从二元正态分布，估计五个参数. 再从这个二元正态分布生成1,000个各含15对数据的样本.
\item\label{item:10.17.d}\index{SI0274@相关}\index{SI0775@极限分布} 若$(X,Y)$服从相关系数为$\rho$的二元正态分布，样本相关系数为$r$，则可用Delta方法证明
$$\sqrt n(r-\rho)\longrightarrow\mathrm n(0,(1-\rho^2)^2).$$
用此事实估计$r$的标准差. 它与自助法估计相比如何？画出$r$的近似pdf.
\item\label{item:10.17.e}\index{SI1039@Fisher的z变换} Fisher的$z$变换是相关系数的方差稳定变换（variance stabilizing transformation，见\autoref{exercise:11.4}）. 若$(X,Y)$服从相关系数为$\rho$的二元正态分布，样本相关系数为$r$，则
$$\frac12\left[\log\left(\frac{1+r}{1-r}\right)-\log\left(\frac{1+\rho}{1-\rho}\right)\right]$$
近似正态. 用此事实画出$r$的近似pdf.
\end{enumerate}
（建立\hyperref[item:10.17.d]{(d)}的正态性结果需要一些繁琐的矩阵计算，见\hyperref[bib:source200]{Lehmann和Casella，1998}，例6.5. \hyperref[item:10.17.e]{(e)}的$z$变换比\hyperref[item:10.17.d]{(d)}的Delta方法更快地收敛到正态分布. \hyperref[bib:source99]{Diaconis和Holmes，1994}，对此问题做了穷尽的自助法计算，枚举全部77,558,760个相关系数.）
\end{exercise}

\sourcepageanchor{442}
\begin{exercise}\label{exercise:10.18}
对\hyperref[ex:10.2.1]{习题10.2.1}的情形，即$X_1,X_2,\ldots,X_n$独立同分布，其中$X_i$以概率$1-\delta$服从$\mathrm n(\theta,\sigma^2)$，以概率$\delta$服从$f(x)$，而$f(x)$为任意均值$\mu$、方差$\tau^2$的密度，证明
$$\operatorname{Var}(\bar X)=(1-\delta)\frac{\sigma^2}n+\delta\frac{\tau^2}n+\frac{\delta(1-\delta)(\theta-\mu)^2}n.$$
并推出柯西pdf的污染总会导致无穷方差. （提示：将混合模型写成分层模型. 令$Y$以概率$1-\delta$为0、以概率$\delta$为1. 则$\operatorname{Var}(X_i)=E[\operatorname{Var}(X_i)\mid Y]+\operatorname{Var}(E[X_i\mid Y])$.）\footnote{译注：原题称“习题10.2.1”，对应正文实际为\autoref{ex:10.2.1}. 提示第一项条件位置原文如上，按全方差公式应为$E[\operatorname{Var}(X_i\mid Y)]$. 此处保留原文.}
\end{exercise}

\begin{exercise}\label{exercise:10.19}
基础假设还可能因抽样中的相关性而被违反，这会严重影响样本均值的性质. 假设在\hyperref[ex:10.2.1]{习题10.2.1}所讨论的情形中引入相关性，即观测$X_1,X_2,\ldots,X_n$，各$X_i\sim\mathrm n(\theta,\sigma^2)$，但不再独立.\footnote{译注：原题“习题10.2.1”应指\autoref{ex:10.2.1}. \hyperref[item:10.19.b]{(b)}方差式中括号内的加号与协方差求和不符，通常应为$n-(1-\rho^n)/(1-\rho)$；\hyperref[itm:exercise:10.19:c]{(c)}还需初始变量处于平稳分布，且与后续创新独立，才能得出所列相关式. 原题均保留.}
\begin{enumerate}[label=(\alph*)]
\item 对等相关情形，即$i\ne j$时$\operatorname{Corr}(X_i,X_j)=\rho$，证明
$$\operatorname{Var}(\bar X)=\frac{\sigma^2}n+\frac{n-1}n\rho\sigma^2,$$
因此$n\to\infty$时$\operatorname{Var}(\bar X)\not\to0$.
\item\label{item:10.19.b} 若沿时间（或距离）观测$X_i$，有时假设相关性随时间（或距离）减小，一个具体模型是$\operatorname{Corr}(X_i,X_j)=\rho^{|i-j|}$. 证明在这种情形下
$$\operatorname{Var}(\bar X)=\frac{\sigma^2}n+\frac{2\sigma^2}{n^2}\frac\rho{1-\rho}\left(n+\frac{1-\rho^n}{1-\rho}\right),$$
所以$n\to\infty$时$\operatorname{Var}(\bar X)\to0$. （相关性的另一影响见\hyperref[sec:5.8.2]{杂录5.8.2}.）
\item\phantomsection\label{itm:exercise:10.19:c}\index{SI0072@自回归模型} \hyperref[item:10.19.b]{(b)}的相关结构出现在自回归AR(1)模型（autoregressive AR(1) model）中，假设$X_{i+1}=\rho X_i+\delta_i$，其中$\delta_i$独立同分布于$\mathrm n(0,1)$. 若$|\rho|<1$，定义$\sigma^2=1/(1-\rho^2)$，证明$\operatorname{Corr}(X_1,X_i)=\rho^{i-1}$.
\end{enumerate}
\end{exercise}

\begin{exercise}\label{exercise:10.20}\index{SI0100@崩溃值}
参照关于崩溃值的\autoref{def:10.2.2}，证明：
\begin{enumerate}[label=(\alph*)]
\item 若$T_n=\bar X_n$为样本均值，证明$b=0$.
\item 若$T_n=M_n$为样本中位数，证明$b=.5$.

在敏感性方面介于均值与中位数之间的估计量是\index{SI0929@截尾}\semanticindexanchor{929:442}\index{SI1042@截尾均值}\semanticindexanchor{1042:442}$\alpha$截尾均值（$\alpha$-trimmed mean），$0<\alpha<1/2$，定义如下：删去最小的$\alpha n$个观测和最大的$\alpha n$个观测，取剩余观测的算术平均，得到$\bar X_n^\alpha$.
\item 若$T_n=\bar X_n^\alpha$为样本的$\alpha$截尾均值，$0<\alpha<1/2$，证明$0<b<1/2$.
\end{enumerate}
\end{exercise}

\begin{exercise}\label{exercise:10.21}\index{SI0805@中位数绝对偏差}
均值与中位数的崩溃表现也延续到对应的尺度估计量. 对样本$X_1,\allowbreak X_2,\allowbreak\ldots,\allowbreak X_n$：
\begin{enumerate}[label=(\alph*)]
\item 证明样本方差$S^2=\sum(X_i-\bar X)^2/(n-1)$的崩溃值为0.
\item 一种稳健替代量是中位数绝对偏差（median absolute deviation，MAD），即$|X_1-m|,|X_2-m|,\ldots,|X_n-m|$的中位数，其中$m$为样本中位数. 证明该估计量的崩溃值为50\%.
\end{enumerate}
\end{exercise}

\begin{exercise}\label{exercise:10.22}
本题考察\autoref{ex:10.2.3}的若干细节.
\begin{enumerate}[label=(\alph*)]
\item 验证若$n$为奇数，则
$$P\bigl(\sqrt n(M_n-\mu)\leq a\bigr)=P\left(\frac{\sum_iY_i-np_n}{\sqrt{np_n(1-p_n)}}\geq\frac{(n+1)/2-np_n}{\sqrt{np_n(1-p_n)}}\right).$$
\item 验证$p_n\to p=F(\mu)=1/2$，且
$$\frac{(n+1)/2-np_n}{\sqrt{np_n(1-p_n)}}\longrightarrow-2aF'(\mu)=-2af(\mu).$$
（提示：证明$[(n+1)/2-np_n]/\sqrt n$是导数的极限形式.）
\sourcepageanchor{443}
\item 解释如何从
$$P\bigl(\sqrt n(M_n-\mu)\leq a\bigr)\longrightarrow P(Z\geq-2af(\mu))$$
得到$\sqrt n(M_n-\mu)$渐近正态、均值为0、方差为$1/[2f(\mu)]^2$的结论.
\end{enumerate}
（注意，只有$p_n$不依赖$n$时，中心极限定理才可直接应用. 由于这里确实依赖$n$，要严格得出极限正态性还需更多工作. 完成这些工作后，结果与预期一致.）
\end{exercise}

\begin{exercise}\label{exercise:10.23}\index{SI0386@Tukey}\index{SI0916@Tukey模型}
本题进一步考察中位数相对于均值的ARE，即$\operatorname{ARE}(M_n,\bar X)$.
\begin{enumerate}[label=(\alph*)]
\item 验证\autoref{ex:10.2.4}给出的三个ARE.
\item 证明$\operatorname{ARE}(M_n,\bar X)$不受尺度变化影响，即基础pdf为$f(x)$还是$(1/\sigma)f(x/\sigma)$没有区别.
\item 当基础分布为自由度$\nu$的Student $t$分布时，对$\nu=3,5,10,25,50,\infty$计算$\operatorname{ARE}(M_n,\bar X)$. 关于ARE与分布尾部，你能得出什么结论？
\item 当基础pdf为Tukey模型
$$X\sim\begin{cases}\mathrm n(0,1),&\text{概率为$1-\delta$},\\\mathrm n(0,\sigma^2),&\text{概率为$\delta$},\end{cases}$$
时，计算$\operatorname{ARE}(M_n,\bar X)$. 对一系列$\delta$和$\sigma$计算ARE. 关于均值和中位数的相对表现，你能得出什么结论？
\end{enumerate}
\end{exercise}

\begin{exercise}\label{exercise:10.24}
假设$\theta_0$满足$E_{\theta_0}\psi(X-\theta_0)=0$，证明\eqref{eq:10.2.4}与\eqref{eq:10.2.5}蕴含\eqref{eq:10.2.6}.
\end{exercise}

\begin{exercise}\label{exercise:10.25}
若$f(x)$是关于0对称的pdf，$\rho$为对称函数，证明$\int\psi(x-\theta)f(x-\theta)\,dx=0$，其中$\psi=\rho'$. 再证明，这意味着若$X_1,X_2,\ldots,X_n$独立同分布于$f(x-\theta)$，$\hat\theta_M$为$\sum_i\rho(x_i-\theta)$的极小化点，则$\hat\theta_M$渐近正态，其均值等于$\theta$的真值.
\end{exercise}

\begin{exercise}\label{exercise:10.26}
考察\autoref{ex:10.2.6}计算的若干细节.
\begin{enumerate}[label=(\alph*)]
\item 验证$E_\theta\psi'(X-\theta)$与$E_\theta[\psi(X-\theta)]^2$的表达式，从而验证$\hat\theta_M$的方差公式.
\item 计算$\psi'$的期望时，注意到$\psi$并非处处可微，但可处理其中可微的部分. 另一种做法是注意$\psi$的期望可微，在\eqref{eq:10.2.5}中可以写为
$$\frac1n\sum_{i=1}^n\psi'(x_i-\theta_0)\longrightarrow\left.\frac d{d\theta_0}E_\theta\psi(X-\theta_0)\right|_{\theta_0=\theta}.$$
证明这与\eqref{eq:10.2.5}的极限相同.\footnote{译注：按函数自变量求导的撇号约定，右侧关于$\theta_0$的导数应带负号；这与\autoref{sec:10.2.2}的符号疑误相同. 原式保留.}
\end{enumerate}
\end{exercise}

\begin{exercise}\label{exercise:10.27}
对\autoref{ex:10.6.2}的情形：\footnote{译注：\hyperref[itm:exercise:10.27:b]{(b)}中的中位数分段式未单列中位数恰等于污染点$x$的中间情形；原式照录，待核实端点约定.}
\begin{enumerate}[label=(\alph*)]
\item 验证$IF(\bar X,x)=x-\mu$.
\item\phantomsection\label{itm:exercise:10.27:b} 对中位数，若$P(X\leq m)=1/2$，则$T(F)=m$，或$m=F^{-1}(1/2)$. 若$X\sim F_\delta$，证明
$$P(X\leq a)=\begin{cases}(1-\delta)F(a),&x>a,\\(1-\delta)F(a)+\delta,&\text{其他},\end{cases}$$
从而
$$T(F_\delta)=\begin{cases}F^{-1}\left(\dfrac1{2(1-\delta)}\right),&x>F^{-1}\left(\dfrac1{2(1-\delta)}\right),\\F^{-1}\left(\dfrac{1/2-\delta}{1-\delta}\right),&\text{其他}.
\end{cases}$$
\item 证明
$$\frac1\delta\left[F^{-1}\left(\frac1{2(1-\delta)}\right)-F^{-1}\left(\frac12\right)\right]\longrightarrow\frac1{2f(m)},$$
并完成计算$IF(M,x)$的论证.

（提示：写成$a_\delta=F^{-1}\left(\frac1{2(1-\delta)}\right)$，说明此极限为$a'_\delta|_{\delta=0}$. 后者可用隐函数求导及$(1-\delta)^{-1}=2F(a_\delta)$计算.）
\end{enumerate}
\end{exercise}

\sourcepageanchor{444}
\begin{exercise}\label{exercise:10.28}
证明：若$\rho$由\eqref{eq:10.2.2}定义，则$\rho$与$\rho'$均连续.
\end{exercise}

\begin{exercise}\label{exercise:10.29}\index{SI0794@渐近有效}
由\eqref{eq:10.2.9}，知道$M$估计量不可能比最大似然估计量更有效. 不过，也知道何时能同样有效.
\begin{enumerate}[label=(\alph*)]
\item 证明：若取$\psi(x-\theta)=cl'(\theta\mid x)$，其中$l$为对数似然，$c$为常数，则\eqref{eq:10.2.9}取等号.
\item 对下列各分布，验证对应的$\psi$函数给出渐近有效的$M$估计量.
\begin{enumerate}[label=(\roman*)]
\item 正态：$f(x)=e^{-x^2/2}/\sqrt{2\pi}$，$\psi(x)=x$.
\item 逻辑斯蒂：$f(x)=e^{-x}/(1+e^{-x})^2$，$\psi(x)=\tanh(x)$，其中$\tanh(x)$为双曲正切（hyperbolic tangent）.\footnote{译注：所列逻辑斯蒂密度的位置得分与$\tanh(x/2)$成比例，而非$\tanh(x)$. 此处原式保留.}
\item 柯西：$f(x)=[\pi(1+x^2)]^{-1}$，$\psi(x)=2x/(1+x^2)$.
\item 最少信息分布（least informative distribution）：
$$f(x)=\begin{cases}Ce^{-x^2/2},&|x|\leq c,\\Ce^{-c|x|+c^2/2},&|x|>c,\end{cases}$$
对应$\psi(x)=\max\{-c,\min(c,x)\}$，$C$与$c$为常数.
\end{enumerate}
\end{enumerate}
（更多细节见\hyperref[bib:source171]{Huber，1981}，第3.5节.）
\end{exercise}

\begin{exercise}\label{exercise:10.30}
对$M$估计量，$\psi$函数与崩溃值之间存在联系. 细节相当复杂（\hyperref[bib:source171]{Huber，1981}，第3.2节），但可以概括为：若$\psi$有界，则对应$M$估计量的崩溃值为
$$b^*=\frac\eta{1+\eta},\quad\text{其中}\quad\eta=\min\left\{-\frac{\psi(-\infty)}{\psi(\infty)},-\frac{\psi(\infty)}{\psi(-\infty)}\right\}.$$
\begin{enumerate}[label=(\alph*)]
\item 计算\autoref{exercise:10.29}中有效$M$估计量的崩溃值. 哪些既有效又稳健？
\item\label{item:10.30.b} 计算下列其他$M$估计量的崩溃值：
\begin{enumerate}[label=(\roman*)]
\item \eqref{eq:10.2.1}给出的Huber估计量.
\item Tukey双权函数（Tukey's biweight）：$|x|\leq c$时$\psi(x)=x(c^2-x^2)$，其他情形为0，其中$c$为常数.
\item Andrews正弦波（Andrew's sine wave）：$|x|\leq c\pi$时$\psi(x)=c\sin(x/c)$，其他情形为0.
\end{enumerate}
\item 当基础分布为(i)正态、(ii)双指数时，评价\hyperref[item:10.30.b]{(b)}中估计量相对于MLE的ARE.
\end{enumerate}
\end{exercise}

\begin{exercise}\label{exercise:10.31}\index{SI0106@独立性卡方检验}\index{SI0197@列联表}\index{SI0491@列联表}
从多个总体采集的二项数据常以列联表（contingency table）呈现. 对两个总体，表格可能如下：
\begin{center}
\begin{tabular}{lccr}
&\multicolumn{2}{c}{总体}&\\
&1&2&合计\\\cline{2-3}
成功&\multicolumn{1}{|c|}{$S_1$}&\multicolumn{1}{c|}{$S_2$}&$S=S_1+S_2$\\
失败&\multicolumn{1}{|c|}{$F_1$}&\multicolumn{1}{c|}{$F_2$}&$F=F_1+F_2$\\\cline{2-3}
合计&$n_1$&$n_2$&$n=n_1+n_2$
\end{tabular}
\end{center}
其中总体1为binomial$(n_1,p_1)$，有$S_1$次成功与$F_1$次失败；总体2为binomial$(n_2,p_2)$，有$S_2$次成功与$F_2$次失败. 通常感兴趣的假设为
$$H_0:p_1=p_2\qquad\text{对}\qquad H_1:p_1\ne p_2.$$
\begin{enumerate}[label=(\alph*)]
\item\label{item:10.31.a} 证明可以基于统计量
$$T=\frac{(\hat p_1-\hat p_2)^2}{\left(\frac1{n_1}+\frac1{n_2}\right)\bigl(\hat p(1-\hat p)\bigr)}$$
建立检验，其中$\hat p_1=S_1/n_1$、$\hat p_2=S_2/n_2$、$\hat p=(S_1+S_2)/(n_1+n_2)$. 再证明当$n_1,n_2\to\infty$时，$T$的分布趋于$\chi_1^2$. （这是称为独立性卡方检验，chi squared test of independence，的一类检验的特例.）
\sourcepageanchor{445}
\item 另一种衡量偏离$H_0$程度的方法是计算期望频数表（expected frequency table）. 以边际合计为条件，按照$H_0:p_1=p_2$填表，即
\begin{center}
\renewcommand{\arraystretch}{1.7}
\begin{tabular}{lccr}
&\multicolumn{2}{c}{期望频数}&\\
&1&2&合计\\\cline{2-3}
成功&\multicolumn{1}{|c|}{$\dfrac{n_1S}{n_1+n_2}$}&\multicolumn{1}{c|}{$\dfrac{n_2S}{n_1+n_2}$}&$S=S_1+S_2$\\\cline{2-3}
失败&\multicolumn{1}{|c|}{$\dfrac{n_1F}{n_1+n_2}$}&\multicolumn{1}{c|}{$\dfrac{n_2F}{n_1+n_2}$}&$F=F_1+F_2$\\\cline{2-3}
合计&$n_1$&$n_2$&$n=n_1+n_2$
\end{tabular}
\end{center}
利用期望频数表，逐个单元计算并求和，得到统计量$T^*$：
$$\begin{aligned}
T^*&=\sum\frac{(\text{观测值}-\text{期望值})^2}{\text{期望值}}\\
&=\frac{\left(S_1-\frac{n_1S}{n_1+n_2}\right)^2}{\frac{n_1S}{n_1+n_2}}+\cdots+\frac{\left(F_2-\frac{n_2F}{n_1+n_2}\right)^2}{\frac{n_2F}{n_1+n_2}}.
\end{aligned}$$
用代数方法证明$T^*=T$，从而$T^*$渐近服从卡方分布.
\item 另一个可用来检验$p_1$和$p_2$相等的统计量为
$$T^{**}=\frac{\hat p_1-\hat p_2}{\sqrt{\frac{\hat p_1(1-\hat p_1)}{n_1}+\frac{\hat p_2(1-\hat p_2)}{n_2}}}.$$
证明在$H_0$下，$T^{**}$渐近服从$\mathrm n(0,1)$，因此其平方渐近服从$\chi_1^2$. 再证明$(T^{**})^2\ne T^*$.
\item 在什么情况下，一个统计量优于另一个？
\item 19世纪后期，Joseph Lister进行了一项著名的医学实验. 当时手术相关死亡率很高，Lister猜想使用消毒剂石炭酸会有所帮助. 数年间，Lister完成75例使用或不使用石炭酸的截肢手术，数据如下：
\begin{center}
\begin{tabular}{llcc}
&&\multicolumn{2}{c}{是否使用石炭酸？}\\
&&是&否\\\cline{3-4}
患者是否存活？&是&\multicolumn{1}{|c|}{34}&\multicolumn{1}{c|}{19}\\\cline{3-4}
&否&\multicolumn{1}{|c|}{6}&\multicolumn{1}{c|}{16}\\\cline{3-4}
\end{tabular}
\end{center}
用这些数据检验使用石炭酸是否与患者死亡有关.
\end{enumerate}
\end{exercise}

\begin{exercise}\label{exercise:10.32}\index{SI0485@渐近}\index{SI0499@McNemar检验}\leavevmode
\begin{enumerate}[label=(\alph*)]
\item 设$(X_1,\ldots,X_n)\sim\operatorname{multinomial}(m,p_1,\ldots,p_n)$. 考虑检验$H_0:p_1=p_2$对$H_1:p_1\ne p_2$. 一个常用的检验称为McNemar检验（McNemar's Test），在
$$\frac{(X_1-X_2)^2}{X_1+X_2}>\chi_{1,\alpha}^2$$
时拒绝$H_0$. 证明这个检验统计量具有如下形式（如\autoref{exercise:10.31}）：
$$\sum_1^n\frac{(\text{观测值}-\text{期望值})^2}{\text{期望值}},$$
其中$X_i$是观测单元频数，期望单元频数是在$p_1=p_2$假设下$mp_i$的MLE.
\sourcepageanchor{446}
\item McNemar检验常用于以下类型的问题. 先问受试者赞成还是反对某个陈述，再让他们阅读与该陈述有关的信息，随后再次询问赞成还是反对. 各类别回答人数汇总于如下$2\times2$表：
\begin{center}
\begin{tabular}{llcc}
&&\multicolumn{2}{c}{之前}\\
&&赞成&反对\\\cline{3-4}
之后&赞成&\multicolumn{1}{|c|}{$X_3$}&\multicolumn{1}{c|}{$X_2$}\\\cline{3-4}
&反对&\multicolumn{1}{|c|}{$X_1$}&\multicolumn{1}{c|}{$X_4$}\\\cline{3-4}
\end{tabular}
\end{center}
假设$H_0:p_1=p_2$意味着，从赞成改为反对者的比例，与从反对改为赞成者的比例相同. 另一种可检验的假设是：最初赞成的人中改变意见的比例，与最初反对的人中改变意见的比例相同. 用条件概率表达这个假设，证明它不同于上述$H_0$. （此假设可用类似\autoref{exercise:10.31}的$\chi^2$检验检验.）
\end{enumerate}
\end{exercise}

\begin{exercise}\label{exercise:10.33}
补全\autoref{thm:10.3.1}中的空缺. 用\autoref{thm:10.1.12}及Slutsky定理（\autoref{thm:5.5.17}）证明$(\theta-\hat\theta)/\sqrt{-l''(\hat\theta\mid\boldsymbol x)}\to\mathrm n(0,1)$，从而$-2\log\lambda(\boldsymbol X)\to\chi_1^2$.\footnote{译注：原题除以观测信息量的平方根，与\autoref{thm:10.3.1}证明中的同一疑误相连；通常应乘以该平方根，或除以其倒数. 原式保留.}
\end{exercise}

\begin{exercise}\label{exercise:10.34}
检验$H_0:p=p_0$对$H_1:p\ne p_0$，假设观测到独立同分布于Bernoulli$(p)$的$X_1,\ldots,X_n$.
\begin{enumerate}[label=(\alph*)]
\item 导出$-2\log\lambda(\boldsymbol x)$的表达式，其中$\lambda(\boldsymbol x)$为LRT统计量.
\item 像\autoref{ex:10.3.2}一样，模拟$-2\log\lambda(\boldsymbol x)$的分布，与$\chi^2$近似比较.
\end{enumerate}
\end{exercise}

\begin{exercise}\label{exercise:10.35}
设$X_1,\ldots,X_n$为$\mathrm n(\mu,\sigma^2)$总体的随机样本.
\begin{enumerate}[label=(\alph*)]
\item 若$\mu$未知、$\sigma^2$已知，证明$Z=\sqrt n(\bar X-\mu_0)/\sigma$是检验$H_0:\mu=\mu_0$的Wald统计量.
\item 若$\sigma^2$未知、$\mu$已知，求检验$H_0:\sigma=\sigma_0$的Wald统计量.
\end{enumerate}
\end{exercise}

\begin{exercise}\label{exercise:10.36}
设$X_1,\ldots,X_n$为gamma$(\alpha,\beta)$总体的随机样本，假设$\alpha$已知、$\beta$未知. 考虑检验$H_0:\beta=\beta_0$.
\begin{enumerate}[label=(\alph*)]
\item $\beta$的MLE是什么？
\item\label{item:10.36.b} 导出检验$H_0$的Wald统计量，在统计量的分子与分母中都使用MLE.
\item 重做\hyperref[item:10.36.b]{(b)}，但在标准误中使用样本标准差.
\end{enumerate}
\end{exercise}

\begin{exercise}\label{exercise:10.37}
设$X_1,\ldots,X_n$为$\mathrm n(\mu,\sigma^2)$总体的随机样本.
\begin{enumerate}[label=(\alph*)]
\item 若$\mu$未知、$\sigma^2$已知，证明$Z=\sqrt n(\bar X-\mu_0)/\sigma$是检验$H_0:\mu=\mu_0$的得分统计量.
\item 若$\sigma^2$未知、$\mu$已知，求检验$H_0:\sigma=\sigma_0$的得分统计量.
\end{enumerate}
\end{exercise}

\begin{exercise}\label{exercise:10.38}
设$X_1,\ldots,X_n$为gamma$(\alpha,\beta)$总体的随机样本，假设$\alpha$已知、$\beta$未知. 考虑检验$H_0:\beta=\beta_0$，导出检验$H_0$的得分统计量.
\end{exercise}

\begin{exercise}\label{exercise:10.39}
扩展\autoref{ex:10.3.7}中的比较.
\begin{enumerate}[label=(\alph*)]
\item 另一种基于Huber $M$估计量的检验，可以使用\eqref{eq:10.3.6}的方差估计. 考察这种检验统计量的表现，评价它作为\eqref{eq:10.3.8}或\eqref{eq:10.3.9}替代量的优点（或不足）.
\item 另一种基于Huber $M$估计量的检验，可以使用自助法计算的方差. 考察这种检验统计量的表现.
\item $\hat\theta_M$的一个稳健竞争者是中位数. 考察基于中位数的位置参数检验的表现.
\end{enumerate}
\end{exercise}

\begin{exercise}\label{exercise:10.40}
在\autoref{ex:10.4.5}中看到，泊松假设结合中心极限定理，可以基于
$$\frac{\bar X-\lambda}{\sqrt{\lambda/n}}\longrightarrow\mathrm n(0,1)$$
构造近似区间.
\sourcepageanchor{447}
证明，根据\hyperref[bib:source330]{Wilks（1938）}，该近似最优. 即证明
$$\frac{\bar X-\lambda}{\sqrt{\lambda/n}}=\frac{\frac\partial{\partial\lambda}\log L(\lambda\mid\boldsymbol X)}{\sqrt{-E_\lambda\left(\frac{\partial^2}{\partial\lambda^2}\log L(\lambda\mid\boldsymbol X)\right)}}.$$
\end{exercise}

\begin{exercise}\label{exercise:10.41}
设$X_1,\ldots,X_n$独立同分布于negative binomial$(r,p)$. 希望为负二项分布参数构造近似置信区间.
\begin{enumerate}[label=(\alph*)]
\item 计算Wilks近似\eqref{eq:10.4.3}，说明如何用此表达式构造置信区间.
\item\label{item:10.41.b} 求负二项分布均值的一个近似$1-\alpha$置信区间. 说明如何在区间中加入连续性校正.
\item \autoref{exercise:9.23}的蚜虫数据也可以用负二项分布建模. 利用\hyperref[item:10.41.b]{(b)}的结果，对蚜虫数据构造近似90\%置信区间，与\autoref{exercise:9.23}基于泊松分布的区间比较.
\end{enumerate}
\end{exercise}

\begin{exercise}\label{exercise:10.42}
证明\eqref{eq:10.4.5}等价于最高似然区域\eqref{eq:9.2.7}，即对任意固定$\alpha$水平，它们产生相同置信集.
\end{exercise}

\begin{exercise}\label{exercise:10.43}
在\autoref{ex:10.4.7}中，对Wald区间作了两项修正.
\begin{enumerate}[label=(\alph*)]
\item\label{item:10.43.a} 当$y=0$时，将上端点改为$1-(\alpha/2)^{1/n}$；当$y=n$时，将下端点改为$(\alpha/2)^{1/n}$. 说明选取这些端点的理由. （提示：\autoref{sec:9.2.3}.）
\item 第二项修正是将所有区间截断在$[0,1]$内. 证明这个改变结合\hyperref[item:10.43.a]{(a)}的改变，改进了原Wald区间.
\end{enumerate}
\end{exercise}

\begin{exercise}\label{exercise:10.44}
\hyperref[bib:source3]{Agresti和Coull（1998）}“强烈推荐”二项参数的得分区间，但担心像\eqref{eq:10.4.7}这样的公式对初等统计课程而言有些令人生畏. 为以更简单公式产生合理的二项区间，他们建议对Wald区间作如下修正：增加两次成功与两次失败，再使用原Wald公式\eqref{eq:10.4.8}. 即用$\hat p=(y+2)/(n+4)$代替$\hat p=y/n$. 从长度与覆盖概率两方面，将此区间与二项得分区间比较. 你是否同意它是得分区间的合理替代？

（\hyperref[bib:source271]{Samuels和Lu，1992}，建议了另一种基于样本量的Wald区间修正.）
\end{exercise}

\begin{exercise}\label{exercise:10.45}
解出\autoref{ex:10.4.6}所给带连续性校正的近似二项置信区间端点. 证明该区间比不带连续性校正的对应区间更宽，且经过校正的区间具有一致更高的覆盖概率. （事实上，未校正区间的覆盖概率不能保持$1-\alpha$，会在一些参数值处低于此水平. 经过校正的区间则对所有参数值都保持大于$1-\alpha$的覆盖概率.）
\end{exercise}

\begin{exercise}\label{exercise:10.46}
扩展\autoref{ex:10.4.8}中的比较.
\begin{enumerate}[label=(\alph*)]
\item 制作类似\autoref{tab:10.4.8}的表，考察基于中位数的位置参数区间的稳健性. （基于均值的区间已在\autoref{tab:10.4.7}中完成.）
\item 另一种基于Huber $M$估计量的区间，可以使用自助法计算的方差. 考察这种区间的稳健性.
\end{enumerate}
\end{exercise}

\begin{exercise}\label{exercise:10.47}
设$X_1,\ldots,X_n$独立同分布于negative binomial$(r,p)$.
\begin{enumerate}[label=(\alph*)]
\item 补全\autoref{ex:10.4.9}的细节，即证明对小$p$，区间
$$\left\{p:\frac{\chi_{2nr,1-\alpha/2}^2}{2\sum x}\leq p\leq\frac{\chi_{2nr,\alpha/2}^2}{2\sum x}\right\}$$
是近似$1-\alpha$置信区间.
\item 说明如何选择端点以获得最短的$1-\alpha$区间.
\end{enumerate}
\end{exercise}

\begin{exercise}\label{exercise:10.48}
对Fieller置信集的情形（见\hyperref[sec:9.5.3]{杂录9.5.3}），即给定来自参数为$(\mu_X,\allowbreak\mu_Y,\allowbreak\sigma_X^2,\allowbreak\sigma_Y^2,\allowbreak\rho)$的二元正态分布的随机样本$(X_1,Y_1),\allowbreak\ldots,\allowbreak(X_n,Y_n)$，求$\theta=\mu_Y/\mu_X$的近似置信区间. 使用\autoref{ex:5.5.27}中的近似矩计算并应用中心极限定理.
\end{exercise}



\subsection{杂录}\label{sec:10.6}
\sourcepageanchor{448}
\subsubsection{超有效性}\label{sec:10.6.1}\index{SI1014@超效率}
虽然\autoref{thm:7.3.9}的Cramér--Rao下界确实是方差的下界，但\autoref{def:10.1.11}与\autoref{thm:10.1.6}中关于渐近方差的下界可以被突破. Hodges给出了一个胜过\autoref{def:10.1.11}下界的估计量例子（见\hyperref[bib:source194]{LeCam，1953）}.

若$X_1,\ldots,X_n$独立同分布于$\mathrm n(\theta,1)$，$\theta$的无偏估计量的Cramér--Rao下界为$v(\theta)=1/n$. 估计量
$$d_n=\begin{cases}\bar X,&|\bar X|\geq1/n^{1/4},\\a\bar X,&|\bar X|<1/n^{1/4},\end{cases}$$
满足
$$\sqrt n(d_n-\theta)\longrightarrow\mathrm n[0,v(\theta)]\quad\text{依分布},$$
其中$\theta\ne0$时$v(\theta)=1$，$\theta=0$时$v(\theta)=a^2$. 若$a<1$，则不等式\eqref{eq:7.2.5}在$\theta=0$处被违反.\footnote{译注：要保证$a^2<1$，需$|a|<1$，或事先限定$0\leq a<1$；原文仅写$a<1$. 此段所引\autoref{thm:10.1.6}应核对为\autoref{thm:10.1.12}，所引不等式(7.2.5)也需核对. 原文均保留.}

像$d_n$这样的估计量称为超有效估计量（superefficient estimator）. 虽然可以在一定一般性下构造它们，但它们更像理论上的奇特现象，而非实际关注的对象. 原因是，使方差低于下界的$\theta$值构成零概率集合，即Lebesgue测度为零的集合. 不过，超有效估计量的存在提醒我们，在建立估计量性质时，始终应仔细检查假设（在一般情况下也要谨慎！）.

\subsubsection{适当的正则条件}\label{sec:10.6.2}\index{SI0912@正则条件}
“在适当的正则条件下”这句话有些被滥用，因为只要假设足够多，大概就能证明想要的任何结论. 不过，“正则条件”通常十分技术性、相当乏味，而且大多数合理问题中通常都满足. 但它们不可避免，我们应该处理它们. 为完整起见，给出一组足以严格建立\autoref{thm:10.1.6}与\autoref{thm:10.1.12}的正则条件. 这些不是最一般的条件，但对许多应用已经足够一般（一个显著例外是MLE位于参数空间边界的情况）. 预先提醒，以下内容不适合畏难者；跳过也不会对理解造成太大损失.

这些条件主要涉及密度的可微性，以及交换微分与积分的能力（如\autoref{thm:7.3.9}的条件）. 更多细节与更一般的讨论，见\hyperref[bib:source311]{Stuart、Ord和Arnold（1999}，第18章）、\hyperref[bib:source120]{Ferguson（1996}，第4部分），或\hyperref[bib:source200]{Lehmann和Casella（1998}，第6.3节）.

以下四个假设足以证明\autoref{thm:10.1.6}，即MLE的相合性：
\begin{enumerate}[label=(A\arabic*)]
\item\label{item:10.6.A1} 观测$X_1,X_2,\ldots,X_n$，其中$X_i\sim f(x\mid\theta)$独立同分布.
\item\label{item:10.6.A2} 观测密度$f(x\mid\theta)$彼此不同，即若$\theta\ne\theta'$，则$f(x\mid\theta)\ne f(x\mid\theta')$.
\item\label{item:10.6.A3} 密度$f(x\mid\theta)$具有共同支撑集，且$f(x\mid\theta)$关于$\theta$可微.
\item\label{item:10.6.A4} 参数空间$\Omega$包含开集$\omega$，真实参数值$\theta_0$是其内点.
\end{enumerate}
以下两个假设与\hyperref[item:10.6.A1]{(A1)}至\hyperref[item:10.6.A4]{(A4)}一起，足以证明\autoref{thm:10.1.12}，即MLE的渐近正态性与有效性.
\begin{enumerate}[label=(A\arabic*),start=5]
\item\label{item:10.6.A5} 对每个$x\in\mathcal X$，密度$f(x\mid\theta)$关于$\theta$三次可微，三阶导数关于$\theta$连续，且$\int f(x\mid\theta)\,dx$可在积分号下对$\theta$求导三次.
\sourcepageanchor{449}
\item\label{item:10.6.A6} 对任意$\theta_0\in\Omega$，存在正数$c$与函数$M(x)$（二者均可依赖$\theta_0$），使
$$\left|\frac{\partial^3}{\partial\theta^3}\log f(x\mid\theta)\right|\leq M(x)\quad\text{对所有$x\in\mathcal X$、$\theta_0-c<\theta<\theta_0+c$},$$
且$E_{\theta_0}[M(X)]<\infty$.
\end{enumerate}

\subsubsection{再论自助法}\label{sec:10.6.3}
\paragraph{理论}
自助法背后的理论相当精深，以Edgeworth展开（Edgeworth expansions）为基础. 它们是在正态分布附近对分布函数进行的展开，其精神类似泰勒级数展开. 例如，若$X_1,X_2,\ldots,X_n$独立同分布，密度为$f$，均值和方差分别为$\mu$与$\sigma^2$，则$\sqrt n(\bar X-\mu)/\sigma$的cdf的一个Edgeworth展开为（\hyperref[bib:source154]{Hall，1992}，式2.17）
$$P\left(\frac{\sqrt n(\bar X-\mu)}\sigma\leq w\right)=\Phi(w)+\phi(w)\left[-\frac1{6\sqrt n}\kappa(w^2-1)+O\left(\frac1n\right)\right],$$
其中$\Phi$与$\phi$分别为标准正态分布函数与密度函数，$\kappa=E(X_1-\mu)^3$为偏度.\footnote{译注：对于此标准化后的展开，偏度通常应为$E[(X_1-\mu)^3]/\sigma^3$；原文定义缺少$\sigma^3$，除非另有$\sigma=1$的约定. 保留原式.} 展开的第一项是“通常”的正态近似；增加项数会提高展开的准确性.

自助法令人惊讶之处在于，在某些情形下，它会自动使展开的第二项正确（从而达到“二阶”准确性）. 并非所有情形都如此，但对枢轴量进行自助法就是其中一种. \hyperref[bib:source154]{Hall（1992）}全面讨论了自助法的Edgeworth理论，另见\hyperref[bib:source288]{Shao和Tu（1995）}.

\paragraph{实践}
我们仅用自助法计算过标准误，它还有许多用途，其中最流行的可能是构造置信区间. 此外，还有针对不同情形发展的许多自助法变体. 特别是，处理相依数据较为微妙. 关于自助法诸多用途及更多内容的入门介绍，见\hyperref[bib:source114]{Efron和Tibshirani（1993）}.

\paragraph{局限}
自助法或许是近年统计方法学中最重要的单项发展，但它也有局限，也有人持批评态度. 超出独立同分布抽样与枢轴量的情形，自助法就不那么自动化了，但仍可能极有用. 对这些问题的有趣讨论，见\hyperref[bib:source203]{LePage和Billard（1992）}或\hyperref[bib:source336]{Young（1994）}.

\subsubsection{影响函数}\label{sec:10.6.4}\index{SI0580@影响函数}
影响函数（influence function）是一种考虑分布性质的灾难性事件度量，也用于衡量异常观测的影响. 影响函数可以解释为导数，这还带来一些有趣的结果.

统计量的影响函数实际上用其总体对应量计算. 例如，样本均值的影响函数使用总体均值计算，因为它试图度量扰动总体的影响. 类似地，样本中位数的影响函数使用总体中位数计算. 为一致地处理这个思想，可以把估计量看成作用于cdf $F$或其样本对应量——经验cdf $F_n$（\autoref{def:1.5.1}）——的函数. 这种以其他函数为自变量的函数，称为泛函（functionals）.\index{SI0447@泛函}\semanticindexanchor{447:449}

\sourcepageanchor{450}
注意，对样本$X_1,X_2,\ldots,X_n$，知道样本等价于知道经验cdf $F_n$，因为$F_n$在每个$X_i$处跳跃$1/n$. 因此统计量$T=T(X_1,X_2,\ldots,X_n)$可等价写成$T(F_n)$，其总体对应量便可记为$T(F)$.

\begin{definition}\label{def:10.6.1}
对来自cdf为$F$的总体的样本$X_1,X_2,\ldots,X_n$，统计量$T=T(F_n)$在点$x$处的影响函数为
$$IF(T,x)=\lim_{\delta\to0}\frac1\delta[T(F_\delta)-T(F)],$$
其中，当
$$X\sim\begin{cases}F,&\text{概率为$1-\delta$},\\x,&\text{概率为$\delta$},\end{cases}$$
时，记$X\sim F_\delta$；即$F_\delta$是$F$与点$x$的混合.
\end{definition}

\begin{example}[均值与中位数的影响函数]\label{ex:10.6.2}\index{SI0798@影响函数}
假设总体具有连续cdf $F$与pdf $f$. 令$\mu$为总体均值，$\bar X$为样本均值，$T(\cdot)$为计算总体均值的泛函. 因此$T(F_n)=\bar X$，$T(F)=\mu$，且
$$T(F_\delta)=(1-\delta)\mu+\delta x,$$
所以$IF(\bar X,x)=x-\mu$；随着$x$增大，它对$\bar X$的影响越来越大.

对中位数$M$，有（\autoref{exercise:10.27}）
$$IF(M,x)=\begin{cases}\dfrac1{2f(m)},&x>m,\\-\dfrac1{2f(m)},&\text{其他}.\end{cases}$$
因此，与均值不同，中位数的影响函数有界.\index{SI0803@影响函数}\semanticindexanchor{803:450}
\end{example}

为什么有界影响函数很重要？为回答这个问题，考察$M$估计量的影响函数，均值与中位数都是其特例.

令$\hat\theta_M$为方程$\sum_i\psi(x_i-\theta)=0$的$M$估计量解，其中$X_1,X_2,\ldots,X_n$独立同分布，cdf为$F$. 在\autoref{sec:10.2.2}中看到，$\hat\theta_M$是满足$E_{\theta_0}\psi(X-\theta_0)=0$的$\theta_0$的相合估计量. $\hat\theta_M$的影响函数为\footnote{译注：若$\psi'$表示对函数自变量求导，则下面分母应为正的$E_{\theta_0}\psi'(X-\theta_0)$；取$\psi(u)=u$时，原式给出$-(x-\theta_0)$，与本节均值影响函数相反. 原文还写下标$E_0$，此处保留. 影响函数平方及后面的渐近方差不受整体符号影响.}
$$IF(\hat\theta_M,x)=\frac{\psi(x-\theta_0)}{-\int\psi'(t-\theta_0)f(t)\,dt}
=\frac{\psi(x-\theta_0)}{-E_0(\psi'(X-\theta_0))}.$$
现在回顾\eqref{eq:10.2.6}，可见影响函数平方的期望给出\index{SI0070@与影响函数的关系}\semanticindexanchor{70:450}\index{SI0581@与渐近方差的关系}\semanticindexanchor{581:450}$\hat\theta_M$的渐近方差，即
$$\sqrt n(\hat\theta_M-\theta_0)\longrightarrow\mathrm n\bigl(0,E_{\theta_0}[IF(\hat\theta_M,X)]^2\bigr)\quad\text{依分布}.$$
因此，影响函数与渐近方差直接相关.

\subsubsection{自助法区间}\label{sec:10.6.5}
在\autoref{sec:10.1.4}中看到，自助法是一种简单、通用的技术，可求任意统计量的标准误. 计算这些标准误时，实际上构造了统计量的一个分布，即自助分布（bootstrap distribution\mbox{）.} 自然会问：是否有一种简单、通用的方法，利用自助分布作置信陈述？自助法确实能够构造很好的置信区间，可惜它在计算标准误时具有的应用简便性，并未延续到置信区间中.

\sourcepageanchor{451}
使用自助分布百分位数，或对$t$统计量（枢轴量）进行自助法的方法，似乎可能具有普遍适用性. 然而，\hyperref[bib:source114]{Efron和Tibshirani（1993}，第13.4节）指出，“这两种区间一般都表现不好”. \hyperref[bib:source154]{Hall（1992}，第3章）更倾向于$t$统计量方法，并指出对枢轴量进行自助法通常是更优的技术.

百分位数区间和百分位数$t$区间，仅仅是自助法置信区间庞大发展的一角，其中许多方法表现优秀. 但无法用一个简单配方概括这些程序：不同问题需要不同技术.

\subsubsection{稳健区间}\label{sec:10.6.6}
虽然在\autoref{sec:10.2}中较详细讨论了点估计量的稳健性，但除\autoref{ex:10.3.7}与\autoref{ex:10.4.8}外，对稳健检验与置信区间并未给出很多细节. 这并非对主题重要性的评价，更多是篇幅所限.

考察点估计量的稳健性质时，主要关心基础假设发生小幅或大幅偏离后的表现. 同样的关切延续到检验与区间，期望稳健点估计量能够导出稳健检验与区间. 特别地，希望稳健检验在基础模型的一系列偏离下保持功效，希望稳健区间保持覆盖概率. 以下事实支持这种联系（见\hyperref[bib:source299]{Staudte和Sheather，1990}，第5.3.3节）：检验的功效函数可以与其所基于的点估计量的影响函数联系起来. 当然，这立即意味着相关区间估计的覆盖性质也可与影响函数联系起来.

\hyperref[bib:source41]{Boos（1992）}通过估计方程与得分检验，对稳健检验作了很好的介绍. \hyperref[bib:source299]{Staudte和Sheather（1990）}以及\hyperref[bib:source165]{Hettmansberger和McKean（1998）}的著作也是优秀来源，\hyperref[bib:source171]{Huber（1981）}这本现已成为经典的著作亦然.





\section{方差分析与回归}\label{sec:11}
\chapterdoi{11}
\sourcepageanchor{452}
\begin{quote}
“我已经浪费够多时间了，”雷斯垂德站起身说，“我相信的是踏实工作，而不是坐在炉火旁编织精巧的理论.”
\begin{flushright}
雷斯垂德探长\\《贵族单身汉探案》
\end{flushright}
\end{quote}

\subsection{引言}\label{sec:11.1}
到目前为止，我们一直用依赖待估参数的pdf或pmf为随机变量建模. 在许多情形中，包括下面的一些情形，随机变量的模型不仅包含未知参数，还包含已知的（有时可以控制的）协变量（covariates）. 这就是方差分析（analysis of variance，ANOVA）和回归分析（regression analysis）的方法. 它们以线性关系这一基础假设为依据，构成实际应用中统计方法的一大核心部分.

ANOVA是应用最广泛的统计技术之一. ANOVA的一个基本思想——分解变异——是实验统计学的基本思想. ANOVA的名称有些名不副实，因为它关心的并不是分析方差，而是分析\emph{均值的变异}.

我们将研究一种常见的ANOVA，即单因素ANOVA（oneway ANOVA）.\index{SI0043@单因素}\semanticindexanchor{43:452} 关于ANOVA设计各方面的全面论述，可参阅\hyperref[bib:source82]{Cochran和Cox（1957）}的经典教材，或较现代但仍带有一定经典色彩的\hyperref[bib:source94]{Dean和Voss（1999）}以及\hyperref[bib:source192]{Kuehl（2000）}的论述. \hyperref[bib:source232]{Neter等（1993）}的教材介绍了实验统计学的总体策略.

回归技术，特别是线性回归，很可能当选最受欢迎的统计工具. 回归有各种形式：线性、非线性、简单、多元、参数、非参数，等等. 本章考察最简单的情形：只有一个预测变量的线性回归. （这通常称为简单线性回归（simple linear regression），以区别于处理多个预测变量的多元线性回归（multiple linear regression）.）

回归的一个主要目的是探索一个变量对其他变量的依赖.\index{SI0878@回归}\semanticindexanchor{878:452}\index{SI0901@总体}\semanticindexanchor{901:452} 在简单线性回归中，通过关系$EY=\alpha+\beta x$，将随机变量$Y$的均值建模为另一个可观测变量$x$的函数. 一般而言，把$EY$表示为$x$的函数的那个函数，称为总体回归函数（population regression function）.

关于回归模型，较好的综合参考书是\hyperref[bib:source75]{Christensen（1996）}以及\hyperref[bib:source102]{Draper和Smith（1998\mbox{）.}} \hyperref[bib:source311]{Stuart、Ord和Arnold（1999}，第27章）给出了更偏重理论的论述.



\subsection{单因素方差分析}\label{sec:11.2}
在最简单的形式下，ANOVA是一种估计多个总体均值的方法，通常假设这些总体服从正态分布. 不过，ANOVA的核心在于统计设计：怎样用最少的观测，从尽可能多的总体中获得最多的信息？但ANOVA的设计问题不是我们主要关心的内容；我们将关注ANOVA中的推断，也就是估计与检验.

\sourcepageanchor{453}
经典ANOVA以检验为主要目标，特别是检验所谓“ANOVA原假设”. 但近年来，尤其随着计算能力的提高，实验人员已经认识到，检验一个假设（而且正如我们将看到的，这是个有些荒唐的假设）并不能形成良好的实验推断. 因此，虽然我们将推导ANOVA原假设的检验，它却远不是方差分析中最重要的部分. 更重要的是估计，包括点估计和区间估计. 特别是，基于对比（稍后定义）的推断极为重要.

在单因素方差分析（也称单向分类）中，假设观测到的数据$Y_{ij}$服从模型
\begin{equation}\label{eq:11.2.1}
Y_{ij}=\theta_i+\varepsilon_{ij},\qquad i=1,\ldots,k,\quad j=1,\ldots,n_i,
\end{equation}
其中$\theta_i$为未知参数，$\varepsilon_{ij}$为误差随机变量.

\begin{example}[单因素ANOVA]\label{ex:11.2.1}
单因素ANOVA的数据$y_{ij}$可示意如下：
\begin{center}
\begin{tabular}{ccccc}
\toprule
\multicolumn{5}{c}{处理}\\
1&2&3&$\cdots$&$k$\\\midrule
$y_{11}$&$y_{21}$&$y_{31}$&$\cdots$&$y_{k1}$\\
$y_{12}$&$y_{22}$&$y_{32}$&$\cdots$&$y_{k2}$\\
$\vdots$&$\vdots$&$\vdots$&$\cdots$&$y_{k3}$\\
&&$y_{3n_3}$&&$\vdots$\\
$y_{1n_1}$&&&&\\
&$y_{2n_2}$&&&$y_{kn_k}$\\\bottomrule
\end{tabular}
\end{center}
注意，我们并未假设各处理组的观测数相同.

例如，考虑下面一项实验，它旨在评价三种毒素和一个对照处理对某种鳟鱼肝脏的相对影响. 数据为每条被处死的鱼的肝脏损伤程度（用标准单位表示）.
\begin{center}
\begin{tabular}{rrrr}
\toprule
毒素1&毒素2&毒素3&对照\\\midrule
28&33&18&11\\
23&36&21&14\\
14&34&20&11\\
27&29&22&16\\
&31&24&\\
&34&&\\\bottomrule
\end{tabular}
\end{center}
\end{example}

不失一般性，可假设$E\varepsilon_{ij}=0$；否则可以重新调整$\varepsilon_{ij}$，把剩余的均值吸收到$\theta_i$中. 因而
$$
EY_{ij}=\theta_i,\qquad j=1,\ldots,n_i,
$$
所以$\theta_i$是$Y_{ij}$的均值. $\theta_i$通常称为处理均值（treatment means），因为下标常对应不同的处理，或某一处理的不同水平，例如某种药物的剂量水平.\index{SI0039@因子的水平}\semanticindexanchor{39:453}\index{SI0055@处理均值}\semanticindexanchor{55:453}

\eqref{eq:11.2.1}还有另一种模型形式，有时称为过参数化模型（overparameterized model），可写成
\begin{equation}\label{eq:11.2.2}\index{SI0046@过参数化模型}\index{SI0741@方差分析}
Y_{ij}=\mu+\tau_i+\varepsilon_{ij},\qquad i=1,\ldots,k,\quad j=1,\ldots,n_i,
\end{equation}
其中仍有$E\varepsilon_{ij}=0$. 由此模型得到
$$EY_{ij}=\mu+\tau_i.$$

\sourcepageanchor{454}
在这种表述中，将$\mu$看作总均值（grand mean），即各处理共同的均值水平. 参数$\tau_i$表示处理$i$特有的效应，即该处理引起的偏离均值水平的程度. 但由于存在可识别性问题，我们不能分别估计$\tau_i$与$\mu$.

\begin{definition}\label{def:11.2.2}\index{SI0038@可识别性}\index{SI0526@可识别性}\index{SI0527@方差分析}
对于分布族$\{f(x\mid\theta):\theta\in\Theta\}$，如果不同的$\theta$值对应不同的pdf或pmf，就称参数$\theta$是可识别的（identifiable）. 也就是说，若$\theta\ne\theta'$，则$f(x\mid\theta)$与$f(x\mid\theta')$作为$x$的函数并不相同.
\end{definition}

可识别性是模型的性质，不是估计量或估计程序的性质. 但模型不可识别时，推断会遇到困难. 例如，若$f(x\mid\theta)=f(x\mid\theta')$，则来自这两个分布的观测看起来完全相同，我们无法知道参数真值究竟是$\theta$还是$\theta'$. 特别地，$\theta$与$\theta'$都会使似然函数取相同的值.

应当认识到，可识别性问题通常可以通过重新定义模型来解决. 之前没有遇到可识别性问题的一个原因，是我们的模型不仅直观合理，而且可识别（例如，用均值和方差为正态总体建模）. 但这里的模型\eqref{eq:11.2.2}虽然直观合理，却不可识别. 在\autoref{sec:12}将看到二元正态分布的一种参数化，它很好地描述了某种情形，却不可识别.

\eqref{eq:11.2.2}的参数化有$k+1$个参数$(\mu,\tau_1,\ldots,\tau_k)$，却只有$k$个均值$EY_{ij}$，$i=1,\ldots,k$. 若不进一步限制参数，就有不止一组$(\mu,\tau_1,\ldots,\tau_k)$值产生同一个分布. 通常在此模型中加上限制$\sum_{i=1}^k\tau_i=0$，从而实际上把参数个数减少为$k$，使模型可识别. 这一限制还使$\tau_i$可以解释为相对于总体均值水平的偏差. （见\autoref{exercise:11.5}.）

对于单因素ANOVA，我们更愿意使用解释较直接的模型\eqref{eq:11.2.1}，即单元均值模型（cell means model）.\index{SI0029@单元均值模型}\semanticindexanchor{29:454} 但在更复杂的ANOVA中，模型\eqref{eq:11.2.2}有时在解释上更有优势.

\subsubsection{模型与分布假设}\label{sec:11.2.1}\index{SI0027@假设}
在模型\eqref{eq:11.2.1}下，进行任何估计之前所需的最低限度假设，是对所有$i,j$都有$E\varepsilon_{ij}=0$且$\operatorname{Var}\varepsilon_{ij}<\infty$. 在这些假设下，可以对$\theta_i$作一些估计（如\autoref{exercise:7.41}）. 但要作置信区间估计或检验，就需要分布假设. 下面是经典ANOVA假设.

\paragraph{单因素ANOVA假设}
按照模型观测随机变量$Y_{ij}$：
$$
Y_{ij}=\theta_i+\varepsilon_{ij},\qquad i=1,\ldots,k,\quad j=1,\ldots,n_i,
$$
其中
\begin{enumerate}[label=(\roman*)]
\item\label{item:11.2.assumption.i} 对所有$i,j$，$E\varepsilon_{ij}=0$，$\operatorname{Var}\varepsilon_{ij}=\sigma_i^2<\infty$. 对所有$i,i',j,j'$，除非$i=i'$且$j=j'$，否则$\operatorname{Cov}(\varepsilon_{ij},\varepsilon_{i'j'})=0$.
\item\label{item:11.2.assumption.ii} $\varepsilon_{ij}$相互独立，且服从正态分布（正态误差）.
\item\label{item:11.2.assumption.iii}\index{SI0034@方差相等}\index{SI0472@同方差性} 对所有$i$，$\sigma_i^2=\sigma^2$（方差相等，也称同方差性，homoscedasticity）.
\end{enumerate}
没有假设\hyperref[item:11.2.assumption.ii]{(ii)}，只能作点估计，或许还可寻找某一类中方差最小的估计量，却不能作区间估计或检验. 若假设某种非正态分布，区间与检验的推导会很困难（但仍可能做到）. 当然，当样本量适当且总体不太偏斜时，可以依靠中心极限定理.

\sourcepageanchor{455}
方差相等假设也相当重要. 有意思的是，它的重要性与正态性假设相联系. 一般而言，如果怀疑数据严重违反ANOVA假设，首先通常尝试对数据作非线性变换，\index{SI0054@数据变换}\semanticindexanchor{54:455}以便更接近满足ANOVA假设；这通常比为未经变换的数据寻找另一种模型容易. 若干常见变换可见\hyperref[bib:source294]{Snedecor和Cochran（1989）}，另见\autoref{exercise:11.1}和\autoref{exercise:11.2}. （关于变换的其他研究涉及Box--Cox幂变换族，见\autoref{exercise:11.3}.）

\hyperref[bib:source43]{Box（1954）}的经典论文表明，ANOVA对正态性假设的稳健性取决于各方差接近相等的程度（相等更好）. 方差不等时的均值估计问题称为Behrens--Fisher问题，\index{SI0087@Behrens–Fisher问题}\semanticindexanchor{87:455}有着丰富的统计学历史，可追溯到Fisher（\hyperref[bib:source125]{1935}，\hyperref[bib:source126]{1939}）. \hyperref[bib:source311]{Stuart、Ord和Arnold（1999）}对Behrens--Fisher问题作了完整论述.

本章余下部分将采用多数实验情形中的做法，假设这三个经典假设成立. 若数据需要变换或应用中心极限定理，则假设已经采取了这些措施.

\subsubsection{经典的方差分析假设}\label{sec:11.2.2}
经典ANOVA检验针对的原假设是
$$H_0:\quad\theta_1=\theta_2=\cdots=\theta_k,$$
在许多情况下，这个假设显得愚蠢、无趣，而且并不成立. 实验人员通常不会相信不同处理具有\emph{完全相同}的均值. 更合理地说，实验是为了查明哪些处理更好（例如均值更大）；ANOVA真正关心的是估计，而不是检验. （在某些特殊情形中，人们确实关心ANOVA原假设本身.）多数情形类似下面的例子.

\begin{example}[ANOVA假设]\label{ex:11.2.3}
ANOVA起初是分析农业实验的方法. 例如，在研究不同肥料对菠菜植株含锌量（$y_{ij}$）的影响时，考察五种处理. 每种处理由肥料物质（镁、钾和锌）的混合物组成，数据的排列类似\autoref{ex:11.2.1}. 五种处理如下，单位为磅/英亩：
\begin{center}
\begin{tabular}{rrrr}
\toprule 处理&镁&钾&锌\\\midrule
1&0&0&0\\2&0&200&0\\3&50&200&0\\4&200&200&0\\5&0&200&15\\\bottomrule
\end{tabular}
\end{center}
经典ANOVA原假设其实没有什么吸引力，因为实验人员确信不同肥料混合物的效果有所不同. 真正关心的是将这些效果量化.
\end{example}

我们会花一些时间讨论ANOVA原假设，但主要把它作为达到另一目的的手段. 回顾\autoref{sec:9}建立的检验与区间估计之间的联系. 利用这一联系，可以先推导适当的检验，再将其反演，得到置信域（这里这条路线更容易）.

ANOVA原假设的备择假设就是各均值不全相等，即检验
\begin{equation}\label{eq:11.2.3}\index{SI0042@零假设}
H_0:\quad\theta_1=\theta_2=\cdots=\theta_k
\quad\text{对}\quad H_1:\quad\theta_i\ne\theta_j\text{，对某些 }i,j.
\end{equation}

\sourcepageanchor{456}
等价地，可把$H_1$表述为$H_1$：非$H_0$. 要认识到，拒绝$H_0$时，只能断定$\theta_i$之间存在某种差异，却不能推断差异究竟在哪里. （注意，接受$H_1$并不意味着所有$\theta_i$都不同，而只意味着至少有两个不同.）

ANOVA假设的一个问题，也是许多多元假设的共同问题，是假设不容易解释. 相比仅仅断言某些$\theta_i$不同，更有用的是给出$\theta_i$的统计描述. 通过把ANOVA假设分解为较小、较容易描述的部分，就能得到这种描述.

我们已经遇到过把复杂假设分解成较小、较容易理解部分的方法——\autoref{sec:8}的并交法与交并法. 对ANOVA而言，并交法最合适，因为ANOVA原假设是一些较容易理解的单变量假设的交，而这些假设可用对比来表示. 此外，在我们将考察的情形中，基于并交法得到的检验与LRT相同（\autoref{exercise:11.13}），因而具有似然检验的所有性质.

\begin{definition}\label{def:11.2.4}\index{SI0032@对比}\index{SI0199@对比}\index{SI0644@线性组合}
设$\boldsymbol t=(t_1,\ldots,t_k)$是一组变量，可以是参数，也可以是统计量，且$\boldsymbol a=(a_1,\ldots,a_k)$为已知常数. 函数
\begin{equation}\label{eq:11.2.4}
\sum_{i=1}^k a_it_i
\end{equation}
称为$t_i$的线性组合（linear combination）. 若还有$\sum a_i=0$，就称为对比（contrast）.
\end{definition}

对比之所以重要，是因为它能用来比较处理均值. 例如，对均值$\theta_1,\ldots,\theta_k$和常数$\boldsymbol a=(1,-1,0,\ldots,0)$，
$$\sum_{i=1}^k a_i\theta_i=\theta_1-\theta_2$$
就是比较$\theta_1$与$\theta_2$的对比. （更多关于对比的内容见\autoref{exercise:11.10}.）

并交方法的力量在于增进理解. ANOVA原假设是各个原假设的交，而这些单独的原假设很容易直观理解.

\begin{theorem}\label{thm:11.2.5}\index{SI0056@并交检验}\index{SI0523@方差分析}
设$\boldsymbol\theta=(\theta_1,\ldots,\theta_k)$为任意参数，则
$$
\theta_1=\theta_2=\cdots=\theta_k
\quad\Longleftrightarrow\quad
\sum_{i=1}^k a_i\theta_i=0\quad\text{对所有 }\boldsymbol a\in\mathcal A,
$$
其中$\mathcal A$为满足$\mathcal A=\{\boldsymbol a=(a_1,\ldots,a_k):\sum a_i=0\}$的常数组成的集合，也就是说，所有对比都必须满足$\sum a_i\theta_i=0$.
\end{theorem}
\begin{proof}
若$\theta_1=\cdots=\theta_k=\theta$，则
$$
\sum_{i=1}^k a_i\theta_i=\sum_{i=1}^k a_i\theta
=\theta\sum_{i=1}^k a_i=0,
\qquad\text{（因为 }\boldsymbol a\text{ 满足 }\sum a_i=0\text{）},
$$
这就证明了一个方向（$\Rightarrow$）. 为证明另一个方向，考虑下列$\boldsymbol a_i\in\mathcal A$：
$$
\boldsymbol a_1=(1,-1,0,\ldots,0),\quad
\boldsymbol a_2=(0,1,-1,0,\ldots,0),\quad\ldots,\quad
\boldsymbol a_{k-1}=(0,\ldots,0,1,-1).
$$
（集合$(\boldsymbol a_1,\ldots,\boldsymbol a_{k-1})$张成$\mathcal A$中的元素，即任意$\boldsymbol a\in\mathcal A$都能写成$(\boldsymbol a_1,\ldots,\boldsymbol a_{k-1})$的线性组合.）用这些$\boldsymbol a_i$构成对比，得到
$$
\boldsymbol a_1\Rightarrow\theta_1=\theta_2,\quad
\boldsymbol a_2\Rightarrow\theta_2=\theta_3,\quad\ldots,\quad
\boldsymbol a_{k-1}\Rightarrow\theta_{k-1}=\theta_k,
$$
合在一起就推出$\theta_1=\cdots=\theta_k$，定理得证.
\end{proof}

\sourcepageanchor{457}
由\autoref{thm:11.2.5}立即得到，ANOVA原假设可以表示为关于对比的假设. 即原假设成立当且仅当以下假设成立：
$$
H_0:\quad\sum_{i=1}^k a_i\theta_i=0
\quad\text{对所有满足 }\sum_{i=1}^k a_i=0\text{ 的 }(a_1,\ldots,a_k).
$$
此外，如果$H_0$不成立，现在就知道一定至少有一个非零对比. 也就是说，ANOVA备择假设$H_1$：$\theta_i$不全相等，等价于
$$
H_1:\quad\sum_{i=1}^k a_i\theta_i\ne0
\quad\text{对某些满足 }\sum_{i=1}^k a_i=0\text{ 的 }(a_1,\ldots,a_k).
$$
因此，使用对比的好处是，得到的假设稍容易理解，或许也稍容易解释. 但真正的好处是，使用对比后，我们可以按单变量的方式思考和处理问题.

\subsubsection{关于均值线性组合的推断}\label{sec:11.2.3}\index{SI0040@线性组合}
线性组合，特别是对比，在方差分析中起着极其重要的作用. 通过理解和分析对比，可以对$\theta_i$作有意义的推断. 上一节说明，ANOVA原假设实际上是关于对比的陈述. 事实上，ANOVA中大多数有意义的推断都能表示为一个或一组对比. 先从简单的情形开始，考虑单个线性组合的推断.

在单因素ANOVA假设下，有
$$Y_{ij}\sim\mathrm n(\theta_i,\sigma^2),\qquad i=1,\ldots,k,\quad j=1,\ldots,n_i.$$
因此
$$\bar Y_{i\cdot}=\frac1{n_i}\sum_{j=1}^{n_i}Y_{ij}\sim\mathrm n(\theta_i,\sigma^2/n_i),\qquad i=1,\ldots,k.$$
关于记号的说明：通常约定，若一个下标用点$\cdot$代替，就表示已对这个下标求和.\index{SI0393@点记号}\semanticindexanchor{393:457}\index{SI0462@总均值}\semanticindexanchor{462:457} 因而$Y_{i\cdot}=\sum_{j=1}^{n_i}Y_{ij}$，$Y_{\cdot j}=\sum_{i=1}^kY_{ij}$. 加上一条横线表示取均值，如上面的$\bar Y_{i\cdot}$. 若对两个下标都求和并计算总体均值（称为总均值），为了使记号简单些，将破例写作$\bar{\bar Y}=(1/N)\sum_{i=1}^k\sum_{j=1}^{n_i}Y_{ij}$，其中$N=\sum_{i=1}^kn_i$.

对任意常数$\boldsymbol a=(a_1,\ldots,a_k)$，$\sum_{i=1}^ka_i\bar Y_{i\cdot}$也服从正态分布（\autoref{exercise:11.8}），且
$$
E\left(\sum_{i=1}^ka_i\bar Y_{i\cdot}\right)=\sum_{i=1}^ka_i\theta_i,
\qquad
\operatorname{Var}\left(\sum_{i=1}^ka_i\bar Y_{i\cdot}\right)=\sigma^2\sum_{i=1}^k\frac{a_i^2}{n_i},
$$
并且
$$
\frac{\sum_{i=1}^ka_i\bar Y_{i\cdot}-\sum_{i=1}^ka_i\theta_i}
{\sqrt{\sigma^2\sum_{i=1}^ka_i^2/n_i}}\sim\mathrm n(0,1).
$$
虽然这个结果很好，但通常希望在不知道$\sigma$的情况下对$\theta_i$作推断，所以需要用估计值代替$\sigma$. 若将各总体的样本方差记为$S_i^2$，即
$$S_i^2=\frac1{n_i-1}\sum_{j=1}^{n_i}(Y_{ij}-\bar Y_{i\cdot})^2,\qquad i=1,\ldots,k,$$
\sourcepageanchor{458}
则$S_i^2$估计$\sigma^2$，且$(n_i-1)S_i^2/\sigma^2\sim\chi^2_{n_i-1}$. 此外，在ANOVA假设下，各个$S_i^2$都估计同一个$\sigma^2$，所以可以把它们合并来改进估计量. 因而使用$\sigma^2$的合并估计量（pooled estimator）$S_p^2$：
\begin{equation}\label{eq:11.2.5}\index{SI0057@方差估计量}
S_p^2=\frac1{N-k}\sum_{i=1}^k(n_i-1)S_i^2
=\frac1{N-k}\sum_{i=1}^k\sum_{j=1}^{n_i}(Y_{ij}-\bar Y_{i\cdot})^2.
\end{equation}
注意$N-k=\sum(n_i-1)$. 由于各$S_i^2$独立，\autoref{lem:5.3.2}表明$(N-k)S_p^2/\sigma^2\sim\chi^2_{N-k}$. $S_p^2$也与每个$\bar Y_{i\cdot}$独立（\autoref{exercise:11.6}），故
\begin{equation}\label{eq:11.2.6}
\frac{\sum_{i=1}^ka_i\bar Y_{i\cdot}-\sum_{i=1}^ka_i\theta_i}{\sqrt{S_p^2\sum_{i=1}^ka_i^2/n_i}}\sim t_{N-k},
\end{equation}
即自由度为$N-k$的Student $t$分布.

要在水平$\alpha$检验
$$H_0:\quad\sum_{i=1}^ka_i\theta_i=0\quad\text{对}\quad H_1:\quad\sum_{i=1}^ka_i\theta_i\ne0,$$
则在
\begin{equation}\label{eq:11.2.7}\index{SI0051@t检验}\index{SI1017@方差分析}
\left|\frac{\sum_{i=1}^ka_i\bar Y_{i\cdot}}{\sqrt{S_p^2\sum_{i=1}^ka_i^2/n_i}}\right|>t_{N-k,\alpha/2}
\end{equation}
时拒绝$H_0$. （\autoref{exercise:11.9}给出了涉及线性组合的一些其他检验.）此外，\eqref{eq:11.2.6}定义了一个枢轴量，反演后得到$\sum a_i\theta_i$的区间估计量. 以下不等式以概率$1-\alpha$成立：
\begin{equation}\label{eq:11.2.8}\index{SI0030@置信区间}\index{SI0136@方差分析}
\begin{aligned}
\sum_{i=1}^ka_i\bar Y_{i\cdot}-t_{N-k,\alpha/2}\sqrt{S_p^2\sum_{i=1}^k\frac{a_i^2}{n_i}}
&\le\sum_{i=1}^ka_i\theta_i\\
&\le\sum_{i=1}^ka_i\bar Y_{i\cdot}+t_{N-k,\alpha/2}\sqrt{S_p^2\sum_{i=1}^k\frac{a_i^2}{n_i}}.
\end{aligned}
\end{equation}

\begin{example}[ANOVA对比]\label{ex:11.2.6}\index{SI0032@对比}
$\boldsymbol a$的特定取值产生特定的检验或置信区间. 例如，要比较处理1和2，取$\boldsymbol a=(1,-1,0,\ldots,0)$. 利用\eqref{eq:11.2.6}，检验$H_0:\theta_1=\theta_2$对$H_1:\theta_1\ne\theta_2$时，在
$$\left|\frac{\bar Y_{1\cdot}-\bar Y_{2\cdot}}{\sqrt{S_p^2(1/n_1+1/n_2)}}\right|>t_{N-k,\alpha/2}$$
时拒绝$H_0$.

注意，这个检验与两样本$t$检验（\autoref{exercise:8.41}）的区别在于，这里估计$\sigma^2$时，不仅利用处理1和2，也利用处理3到$k$的信息.

另一种情形是把处理1与处理2和3的平均效果比较（例如，处理1为对照，2和3为实验处理，考察某种总体效应）. 此时取$\boldsymbol a=(1,-\tfrac12,-\tfrac12,0,\ldots,0)$，在
$$\left|\frac{\bar Y_{1\cdot}-\tfrac12\bar Y_{2\cdot}-\tfrac12\bar Y_{3\cdot}}
{\sqrt{S_p^2(1/n_1+1/(4n_2)+1/(4n_3))}}\right|>t_{N-k,\alpha/2}$$
时拒绝$H_0:\theta_1=\tfrac12(\theta_2+\theta_3)$.

\sourcepageanchor{459}
利用\eqref{eq:11.2.6}或\eqref{eq:11.2.8}，便有办法检验或估计ANOVA中的任意线性组合. 明智地选择线性组合，可以使我们了解处理均值的许多信息. 例如，考察对比$\theta_1-\theta_2$、$\theta_2-\theta_3$和$\theta_1-\theta_3$，可以了解$\theta_i$的大小顺序. （当然，进行多个检验或构造多个区间时，必须注意总体$\alpha$水平，但可以使用Bonferroni不等式. 见\autoref{ex:11.2.9}.）

由多个对比的组合得出正式结论时，也必须谨慎.\index{SI0032@对比}\semanticindexanchor{32:459} 考虑假设
$$H_0:\quad\theta_1=\tfrac12(\theta_2+\theta_3)\quad\text{对}\quad H_1:\quad\theta_1<\tfrac12(\theta_2+\theta_3)$$
和
$$H_0:\quad\theta_2=\theta_3\quad\text{对}\quad H_1:\quad\theta_2<\theta_3.$$
若两个原假设都被拒绝，可以得出$\theta_3$大于$\theta_1$与$\theta_2$的结论，但仅由这两个检验，不能对$\theta_2$与$\theta_1$的大小顺序作正式结论. （见\autoref{exercise:11.10}.）
\end{example}

现在利用这些关于线性组合的单变量结果，以及\autoref{thm:11.2.5}给出的ANOVA原假设与对比的关系，推导ANOVA原假设的检验.

\subsubsection{方差分析的\texorpdfstring{$F$}{F}检验}\label{sec:11.2.4}\index{SI0056@并交检验}
上一节介绍了如何处理ANOVA中的单个线性组合，特别是对比. 在\autoref{sec:11.2}还看到，ANOVA原假设等价于一个关于对比的假设. 本节利用这一等价性，结合\autoref{sec:8}的并交方法，推导ANOVA假设的检验.

由\autoref{thm:11.2.5}，ANOVA假设检验可以写成
$$
H_0:\ \sum_{i=1}^ka_i\theta_i=0\text{，对所有 }\boldsymbol a\in\mathcal A
\quad\text{对}\quad
H_1:\ \sum_{i=1}^ka_i\theta_i\ne0\text{，对某些 }\boldsymbol a\in\mathcal A,
$$
其中$\mathcal A=\{\boldsymbol a=(a_1,\ldots,a_k):\sum_{i=1}^ka_i=0\}$. 为更清楚地将其视作并交检验，对每个$\boldsymbol a$定义集合
$$\Theta_{\boldsymbol a}=\left\{\boldsymbol\theta=(\theta_1,\ldots,\theta_k):\sum_{i=1}^ka_i\theta_i=0\right\}.$$
则有
$$\boldsymbol\theta\in\{\boldsymbol\theta:\theta_1=\cdots=\theta_k\}
\Longleftrightarrow\boldsymbol\theta\in\Theta_{\boldsymbol a}\text{，对所有 }\boldsymbol a\in\mathcal A
\Longleftrightarrow\boldsymbol\theta\in\bigcap_{\boldsymbol a\in\mathcal A}\Theta_{\boldsymbol a},$$
这说明ANOVA原假设可以写成一个交集.

回顾\autoref{sec:8.2.3}的并交方法，只要对任意一个$\boldsymbol a$，能够拒绝
$$H_{0\boldsymbol a}:\quad\boldsymbol\theta\in\Theta_{\boldsymbol a}
\quad\text{对}\quad H_{1\boldsymbol a}:\quad\boldsymbol\theta\notin\Theta_{\boldsymbol a},$$
就拒绝$H_0:\boldsymbol\theta\in\bigcap_{\boldsymbol a\in\mathcal A}\Theta_{\boldsymbol a}$，因而也拒绝ANOVA原假设. 用\eqref{eq:11.2.6}的$t$统计量检验$H_{0\boldsymbol a}$：
\begin{equation}\label{eq:11.2.9}
T_{\boldsymbol a}=\left|\frac{\sum_{i=1}^ka_i\bar Y_{i\cdot}-\sum_{i=1}^ka_i\theta_i}
{\sqrt{S_p^2\sum_{i=1}^ka_i^2/n_i}}\right|.
\end{equation}

\sourcepageanchor{460}
于是，当$T_{\boldsymbol a}>k$时拒绝$H_{0\boldsymbol a}$，其中$k$为某个常数. 按照并交方法，如果能够对任意一个$\boldsymbol a$拒绝，就能够对使$T_{\boldsymbol a}$最大的$\boldsymbol a$拒绝. 因此，ANOVA原假设的并交检验是在$\sup_{\boldsymbol a}T_{\boldsymbol a}>k$时拒绝$H_0$，其中$k$选择为满足$P_{H_0}(\sup_{\boldsymbol a}T_{\boldsymbol a}>k)=\alpha$.

计算$\sup_{\boldsymbol a}T_{\boldsymbol a}$并非直接明了，但稍加仔细也不困难. 这是一个约束最大化问题，类似前面遇到的问题（例如\autoref{exercise:7.41}计算了约束最小值）. 我们将采用类似此前的办法，使用Cauchy--Schwarz不等式. （也可用Lagrange乘子等方法，但那样就必须用二阶条件验证所得确为最大值.）

最大化的大部分技术论证放在下面的引理中，再应用引理求$T_{\boldsymbol a}$的上确界. 这个引理只是关于二次函数约束最大值的陈述. 不想深入技术细节的读者可以跳过证明.

\begin{lemma}\label{lem:11.2.7}\index{SI0681@二次函数的最大值}
设$(v_1,\ldots,v_k)$为常数，$(c_1,\ldots,c_k)$为正常数. 对$\mathcal A=\{\boldsymbol a=(a_1,\ldots,a_k):\sum a_i=0\}$，有
\begin{equation}\label{eq:11.2.10}
\max_{\boldsymbol a\in\mathcal A}\left\{\frac{(\sum_{i=1}^ka_iv_i)^2}{\sum_{i=1}^ka_i^2/c_i}\right\}
=\sum_{i=1}^kc_i(v_i-\bar v_c)^2,
\end{equation}
其中$\bar v_c=\sum c_iv_i/\sum c_i$. 最大值在任何形如$a_i=Kc_i(v_i-\bar v_c)$的$\boldsymbol a$处取得，其中$K$为非零常数.\footnote{译注：分母要求$\boldsymbol a\ne\boldsymbol0$，原集合未排除零向量. 所列取等向量还需$v_i$不全相同；若全相同，则对每个非零对比向量商均为0，而所列向量本身为零. 原文保留.}
\end{lemma}
\begin{proof}
定义$\mathcal B=\{\boldsymbol b=(b_1,\ldots,b_k):\sum b_i=0\text{ 且 }\sum b_i^2/c_i=1\}$. 对任意$\boldsymbol a\in\mathcal A$，定义$\boldsymbol b=(b_1,\ldots,b_k)$为
$$b_i=\frac{a_i}{\sqrt{\sum_{i=1}^ka_i^2/c_i}},$$
注意$\boldsymbol b\in\mathcal B$. 对任意$\boldsymbol a\in\mathcal A$，
$$\frac{(\sum_{i=1}^ka_iv_i)^2}{\sum_{i=1}^ka_i^2/c_i}=\left(\sum_{i=1}^kb_iv_i\right)^2.$$
我们先求$\boldsymbol b\in\mathcal B$时$(\sum b_iv_i)^2$的上界，再证明引理给出的最大化向量$\boldsymbol a$能达到该上界.

由于处理的是乘积之和，自然可以尝试Cauchy--Schwarz不等式（\autoref{sec:4.7}），但必须仔细纳入涉及$c_i$的约束. 可如下进行. 定义$C=\sum c_i$，写成
$$\frac1{C^2}\left(\sum_{i=1}^kb_iv_i\right)^2
=\left\{\sum_{i=1}^k\left(\frac{b_i}{c_i}\right)(v_i)\left(\frac{c_i}{C}\right)\right\}^2.$$
这是由比值$c_i/C$定义的概率测度下某个协方差的平方. 形式上，若定义随机变量$B,V$满足
$$P\left(B=\frac{b_i}{c_i},V=v_i\right)=\frac{c_i}{C},\qquad i=1,\ldots,k,$$

\sourcepageanchor{461}
则$EB=\sum(b_i/c_i)(c_i/C)=\sum b_i/C=0$. 因而
\begin{align*}
&\left\{\sum_{i=1}^k\left(\frac{b_i}{c_i}\right)(v_i)\left(\frac{c_i}{C}\right)\right\}^2\\
&\quad=(EBV)^2\\
&\quad=(\operatorname{Cov}(B,V))^2\qquad(EB=0)\\
&\quad\le(\operatorname{Var}B)(\operatorname{Var}V)\qquad\text{（Cauchy--Schwarz不等式）}\\
&=\left(\sum_{i=1}^k\left(\frac{b_i}{c_i}\right)^2\frac{c_i}{C}\right)
\left(\sum_{i=1}^k(v_i-\bar v_c)^2\frac{c_i}{C}\right),
\end{align*}
其中$\bar v_c=\sum c_iv_i/\sum c_i$. 利用$\sum b_i^2/c_i=1$并消去公共因子，得到
\begin{equation}\label{eq:11.2.11}
\left(\sum_{i=1}^kb_iv_i\right)^2\le\sum_{i=1}^kc_i(v_i-\bar v_c)^2,
\qquad\text{对任意 }\boldsymbol b\in\mathcal B.
\end{equation}
最后，若$a_i=Kc_i(v_i-\bar v_c)$，其中$K$为任意非零常数，则$\boldsymbol a\in\mathcal A$，且\footnote{译注：原式第二个等号把$K/|K|$省去了；当$K<0$时须补一个负号. 后续最终比较的是平方，取等结论不受影响.}
$$b_i=\frac{Kc_i(v_i-\bar v_c)}{\sqrt{\sum_{i=1}^k(Kc_i(v_i-\bar v_c))^2/c_i}}
=\frac{c_i(v_i-\bar v_c)}{\sqrt{\sum_{i=1}^kc_i(v_i-\bar v_c)^2}}.$$
由于$\sum c_i(v_i-\bar v_c)=0$，
\begin{align*}
\sum_{i=1}^kb_iv_i
&=\frac{\sum_{i=1}^kc_i(v_i-\bar v_c)v_i}{\sqrt{\sum_{i=1}^kc_i(v_i-\bar v_c)^2}}\\
&=\frac{\sum_{i=1}^kc_i(v_i-\bar v_c)^2}{\sqrt{\sum_{i=1}^kc_i(v_i-\bar v_c)^2}}
=\sqrt{\sum_{i=1}^kc_i(v_i-\bar v_c)^2},
\end{align*}
于是\eqref{eq:11.2.11}取等. 因此上界可以达到，函数在这样的$\boldsymbol a$处最大.
\end{proof}

回到\eqref{eq:11.2.9}的$T_{\boldsymbol a}$，显然最大化$T_{\boldsymbol a}$等价于最大化$T_{\boldsymbol a}^2$. 有
$$T_{\boldsymbol a}^2=
\frac{(\sum_{i=1}^ka_i\bar Y_{i\cdot}-\sum_{i=1}^ka_i\theta_i)^2}{S_p^2\sum_{i=1}^ka_i^2/n_i}
=\frac{(\sum_{i=1}^ka_i\bar U_i)^2}{S_p^2\sum_{i=1}^ka_i^2/n_i},
\qquad(\bar U_i=\bar Y_{i\cdot}-\theta_i).$$
注意$S_p^2$不影响最大化，对上式应用\autoref{lem:11.2.7}，得到以下定理.

\begin{theorem}\label{thm:11.2.8}
对\eqref{eq:11.2.9}定义的$T_{\boldsymbol a}$，有
\begin{equation}\label{eq:11.2.12}
\sup_{\boldsymbol a:\sum a_i=0}T_{\boldsymbol a}^2
=\frac{\sum_{i=1}^kn_i\bigl((\bar Y_{i\cdot}-\bar{\bar Y})-(\theta_i-\bar\theta)\bigr)^2}{S_p^2},
\end{equation}
\sourcepageanchor{462}
其中$\bar{\bar Y}=\sum n_i\bar Y_{i\cdot}/\sum n_i$，$\bar\theta=\sum n_i\theta_i/\sum n_i$. 此外，在ANOVA假设下，
\begin{equation}\label{eq:11.2.13}
\sup_{\boldsymbol a:\sum a_i=0}T_{\boldsymbol a}^2\sim(k-1)F_{k-1,N-k},
\end{equation}
即$\sup_{\boldsymbol a:\sum a_i=0}T_{\boldsymbol a}^2/(k-1)$服从自由度为$k-1$和$N-k$的$F$分布. （回顾$N=\sum n_i$.）
\end{theorem}
\begin{proof}
为证明\eqref{eq:11.2.12}，使用\hyperref[lem:11.2.7]{引理11.2.1}，将$v_i$对应为$\bar U_i$，$c_i$对应为$n_i$，结果立即得到.\footnote{译注：原文写引理11.2.1，此处实际应为\autoref{lem:11.2.7}；11.2.1是例. 保留原显示编号.}

为证明\eqref{eq:11.2.13}，必须说明\eqref{eq:11.2.12}的分子、分母是独立的卡方随机变量各除以其自由度. 由ANOVA假设得到两点：分子、分母独立，且$S_p^2\sim\sigma^2\chi^2_{N-k}/(N-k)$. 还需一些工作才能证明
$$\frac1{\sigma^2}\sum_{i=1}^kn_i\bigl((\bar Y_{i\cdot}-\bar{\bar Y})-(\theta_i-\bar\theta)\bigr)^2\sim\chi^2_{k-1}.$$
不过这是可以做到的，留作习题. （见\autoref{exercise:11.7}.）
\end{proof}

若$H_0:\theta_1=\theta_2=\cdots=\theta_k$成立，则对所有$i=1,\ldots,k$都有$\theta_i=\bar\theta$，\eqref{eq:11.2.12}中的$\theta_i-\bar\theta$项消失. 因此，对ANOVA假设
$$H_0:\ \theta_1=\theta_2=\cdots=\theta_k\quad\text{对}\quad H_1:\ \theta_i\ne\theta_j\text{，对某些 }i,j,$$
其水平为$\alpha$的检验在
\begin{equation}\label{eq:11.2.14}\index{SI0036@F检验}
\frac{\sum_{i=1}^kn_i(\bar Y_{i\cdot}-\bar{\bar Y})^2}{S_p^2}>(k-1)F_{k-1,N-k,\alpha}
\end{equation}
时拒绝$H_0$. 这一拒绝域通常写作
$$\text{若 }F=\frac{\sum_{i=1}^kn_i(\bar Y_{i\cdot}-\bar{\bar Y})^2/(k-1)}{S_p^2}>F_{k-1,N-k,\alpha}\text{，则拒绝 }H_0,$$
检验统计量$F$称为ANOVA $F$统计量（ANOVA $F$ statistic）.

\subsubsection{对比的同时估计}\label{sec:11.2.5}\index{SI0032@对比}
已经介绍了如何估计和检验ANOVA中的单个对比；$t$统计量见\eqref{eq:11.2.6}，区间见\eqref{eq:11.2.8}. 然而，ANOVA中经常希望作不止一个推断，而多个水平为$\alpha$的检验所形成的同时推断，未必仍处于水平$\alpha$. 在ANOVA的背景下，前面已提到过这个问题.

\begin{example}[两两差值]\label{ex:11.2.9}\index{SI0028@Bonferroni不等式}\index{SI0561@方差分析}
很多时候关心的是均值的两两差值. 因此，若ANOVA的均值为$\theta_1,\ldots,\theta_k$，就可能关心$\theta_1-\theta_2$、$\theta_2-\theta_3$、$\theta_3-\theta_4$等的区间估计. 利用Bonferroni不等式，可建立同时推断的陈述. 定义
$$C_{ij}=\left\{\theta_i-\theta_j:\theta_i-\theta_j\in\bar Y_{i\cdot}-\bar Y_{j\cdot}
\pm t_{N-k,\alpha/2}\sqrt{S_p^2\left(\frac1{n_i}+\frac1{n_j}\right)}\right\}.$$
则对每个$C_{ij}$，$P(C_{ij})=1-\alpha$，但例如$P(C_{12}\text{ 且 }C_{23})<1-\alpha$. 然而，后一类推断正是ANOVA中希望作出的推断.

\sourcepageanchor{463}
回顾\eqref{eq:1.2.10}给出的Bonferroni不等式：对任意集合$A_1,\ldots,A_n$，
$$P\left(\bigcap_{i=1}^nA_i\right)\ge\sum_{i=1}^nP(A_i)-(n-1).$$
这里希望为$P(\bigcap_{i,j}C_{ij})$给出界，即所有两两比较区间都覆盖各自差值的概率.

若希望对$m$个置信集的覆盖作同时$1-\alpha$的陈述，那么由Bonferroni不等式，可以将每个置信集构造成水平为$\gamma$，其中$\gamma$满足
$$1-\alpha=\sum_{i=1}^m\gamma-(m-1),$$
等价地，
$$\gamma=1-\frac\alpha m.$$
还可以稍作推广，不必要求每个单独推断的水平相同. 可将每个置信集构造成水平为$\gamma_i$，其中$\gamma_i$满足
$$1-\alpha=\sum_{i=1}^m\gamma_i-(m-1).$$
在有$k$个处理的ANOVA中，若各$t$区间的置信度为$1-2\alpha/[k(k-1)]$，便可对全部$k(k-1)/2$个两两差值作置信度为$1-\alpha$的同时推断.
\end{example}

\hyperref[bib:source275]{Scheff\'e（1959）}给出了另一种相当优美的同时推断方法. Scheff\'e方法有时称为$S$方法，可以对所有对比构造同时置信区间（或检验）. （\autoref{exercise:11.14}表明，Scheff\'e方法也可用来构造任意线性组合的同时区间，而不只限于对比.）该方法允许设定一个置信系数，对所有对比区间同时有效，而不只是对某个指定的组有效. 若要考察大量对比，更宜使用Scheff\'e方法；若对比数较少，Bonferroni界几乎肯定更小. （其他多重比较方法的讨论见\autoref{sec:11.5}.）

Scheff\'e方法对所有对比具有同时$1-\alpha$覆盖率，这一结论容易由ANOVA检验的并交性质证明.

\begin{theorem}\label{thm:11.2.10}\index{SI0047@Scheffé方法}\index{SI0955@Scheffé方法}
在ANOVA假设下，若$M=\sqrt{(k-1)F_{k-1,N-k,\alpha}}$，则以下不等式对所有$\boldsymbol a\in\mathcal A=\{\boldsymbol a=(a_1,\ldots,a_k):\sum a_i=0\}$同时成立的概率为$1-\alpha$：
$$\sum_{i=1}^ka_i\bar Y_{i\cdot}-M\sqrt{S_p^2\sum_{i=1}^k\frac{a_i^2}{n_i}}
\le\sum_{i=1}^ka_i\theta_i
\le\sum_{i=1}^ka_i\bar Y_{i\cdot}+M\sqrt{S_p^2\sum_{i=1}^k\frac{a_i^2}{n_i}}.$$
\end{theorem}
\begin{proof}
同时概率陈述要求$M$满足
$$P\left(\left|\sum_{i=1}^ka_i\bar Y_{i\cdot}-\sum_{i=1}^ka_i\theta_i\right|
\le M\sqrt{S_p^2\sum_{i=1}^k\frac{a_i^2}{n_i}}\text{，对所有 }\boldsymbol a\in\mathcal A\right)=1-\alpha,$$
等价地，
$$P(T_{\boldsymbol a}^2\le M^2\text{，对所有 }\boldsymbol a\in\mathcal A)=1-\alpha,$$
\sourcepageanchor{464}
其中$T_{\boldsymbol a}$由\eqref{eq:11.2.9}定义. 但由于
$$P(T_{\boldsymbol a}^2\le M^2\text{，对所有 }\boldsymbol a\in\mathcal A)
=P\left(\sup_{\boldsymbol a:\sum a_i=0}T_{\boldsymbol a}^2\le M^2\right),$$
\autoref{thm:11.2.8}表明，选取$M^2=(k-1)F_{k-1,N-k,\alpha}$就满足所需的概率要求.
\end{proof}

Scheff\'e方法的一项真正优势，是允许合法的“数据窥探”（data snooping）.\index{SI0033@数据窥探}\semanticindexanchor{33:464} 在经典统计学中，检验由数据启发的假设是禁忌，因为这会使结果产生偏倚，进而使推断失效. （通常不会仅仅因为发现$\bar Y_{1\cdot}$与$\bar Y_{2\cdot}$不同，就检验$H_0:\theta_1=\theta_2$. 见\autoref{exercise:11.18}.）但使用Scheff\'e方法时，这一策略是合法的. 区间或检验对所有对比都有效，是否由数据启发并无区别，因为Scheff\'e方法已经将它们全部考虑在内.

当然，Scheff\'e方法提供如此大的推断能力，必须付出代价，代价表现为区间长度. 为保证同时置信水平，区间可能相当长. 例如，可以证明（\autoref{exercise:11.15}），比较$t$和$F$分布时，对任意$\nu,\alpha,k$，临界点满足
$$t_{\nu,\alpha/2}\le\sqrt{(k-1)F_{k-1,\nu,\alpha}},$$
所以Scheff\'e区间总比单个对比的区间宽，有时宽得多（这也支持一种观点：实验中，没有什么能代替细致的规划和准备）. 区间长度的现象也延伸到检验. 由上述不等式还可知，Scheff\'e检验的功效低于$t$检验.

\subsubsection{平方和的分解}\label{sec:11.2.6}
ANOVA为思考不同处理如何影响测量变量提供了一种有用的方式，即把变异分配给不同来源. 分配变异的基本思想可概括为下面的恒等式.

\begin{theorem}\label{thm:11.2.11}\index{SI0050@平方和}
对任意数$y_{ij}$，$i=1,\ldots,k$，$j=1,\ldots,n_i$，有
\begin{equation}\label{eq:11.2.15}
\sum_{i=1}^k\sum_{j=1}^{n_i}(y_{ij}-\bar{\bar y})^2
=\sum_{i=1}^kn_i(\bar y_{i\cdot}-\bar{\bar y})^2
+\sum_{i=1}^k\sum_{j=1}^{n_i}(y_{ij}-\bar y_{i\cdot})^2,
\end{equation}
其中$\bar y_{i\cdot}=\frac1{n_i}\sum_jy_{ij}$，$\bar{\bar y}=\sum_i n_i\bar y_{i\cdot}/\sum_i n_i$.
\end{theorem}
\begin{proof}
证明非常简单，仅依赖于处理均值时交叉项常常消失这一事实. 写成
$$\sum_{i=1}^k\sum_{j=1}^{n_i}(y_{ij}-\bar{\bar y})^2
=\sum_{i=1}^k\sum_{j=1}^{n_i}\bigl((y_{ij}-\bar y_{i\cdot})+(\bar y_{i\cdot}-\bar{\bar y})\bigr)^2,$$
展开右端并重新合并各项即可. （见\autoref{exercise:11.21}.）
\end{proof}

\eqref{eq:11.2.15}中的各个和称为平方和（sums of squares），用来度量数据中可归因于不同来源的变异. （有时也称校正平方和，corrected sums of squares，其中“校正”指减去了均值.）特别地，单因素ANOVA模型
$$Y_{ij}=\theta_i+\varepsilon_{ij}$$
\sourcepageanchor{465}
中的各项与\eqref{eq:11.2.15}中的各项一一对应. \eqref{eq:11.2.15}说明如何将变异分配给处理（处理\emph{之间}的变异）和随机误差（处理\emph{内部}的变异）. \eqref{eq:11.2.15}的左端不区分处理类别，度量全部变异；右端两项则分别度量仅由处理引起的变异和仅由随机误差引起的变异. 这些变异来源满足上述恒等式，表明由平方和度量的数据变异，与ANOVA模型一样具有可加性.

平方和较容易处理的一个原因是，在正态性下，校正平方和是卡方随机变量；而已经看到，独立卡方变量相加可得到新的卡方变量.

在ANOVA假设下，特别是若$Y_{ij}\sim\mathrm n(\theta_i,\sigma^2)$，容易证明
\begin{equation}\label{eq:11.2.16}
\frac1{\sigma^2}\sum_{i=1}^k\sum_{j=1}^{n_i}(Y_{ij}-\bar Y_{i\cdot})^2\sim\chi^2_{N-k},
\end{equation}
因为对每个$i=1,\ldots,k$，$\frac1{\sigma^2}\sum_{j=1}^{n_i}(Y_{ij}-\bar Y_{i\cdot})^2\sim\chi^2_{n_i-1}$，且它们相互独立，而对独立卡方随机变量，$\sum_{i=1}^k\chi^2_{n_i-1}\sim\chi^2_{N-k}$. 此外，若对所有$i,j$都有$\theta_i=\theta_j$，则
\begin{equation}\label{eq:11.2.17}
\frac1{\sigma^2}\sum_{i=1}^kn_i(\bar Y_{i\cdot}-\bar{\bar Y})^2\sim\chi^2_{k-1}
\quad\text{且}\quad
\frac1{\sigma^2}\sum_{i=1}^k\sum_{j=1}^{n_i}(Y_{ij}-\bar{\bar Y})^2\sim\chi^2_{N-1}.
\end{equation}
所以在$H_0:\theta_1=\cdots=\theta_k$下，\eqref{eq:11.2.15}的平方和分解就是卡方随机变量的分解. 缩放后，左端服从$\chi^2_{N-1}$分布，右端是两个独立随机变量的和，分别服从$\chi^2_{k-1}$和$\chi^2_{N-k}$. 注意，卡方分解只有在\eqref{eq:11.2.15}右端各项独立时才成立；这里的独立性来自ANOVA假设中的正态性. 卡方分解也在稍一般的背景下成立，其刻画有时称为Cochran定理（Cochran's Theorem）.\index{SI0107@Cochran定理}\semanticindexanchor{107:465} （见\hyperref[bib:source282]{Searle（1971）}，以及\autoref{sec:11.5}.）

一般而言，可以把一个平方和分解成互不相关的对比的平方和，各有一个自由度. 若平方和有$\nu$个自由度且服从$\chi^2_\nu$，就可分解成$\nu$个独立项，每项服从$\chi^2_1$.

量$(\sum a_i\bar Y_{i\cdot})^2/(\sum a_i^2/n_i)$称为处理对比$\sum a_i\bar Y_{i\cdot}$的对比平方和（contrast sum of squares）. 在单因素ANOVA中，总能找到常数组$\boldsymbol a^{(l)}=(a_1^{(l)},\ldots,a_k^{(l)})$，$l=1,\ldots,k-1$，使其满足\footnote{译注：原书下一式各分子排作$\sum_i a_i^{(l)}\bar Y_{i\cdot}^{\,2}$，与前文对比平方和定义不一致；按定义应是$(\sum_i a_i^{(l)}\bar Y_{i\cdot})^2$. 此处照录原式. 紧接着关于\eqref{eq:11.2.18}左、右端卡方分布的说明，实际指这条平方和分解式，并须在相应对比均值为零时取中心卡方分布.}
$$\sum_{i=1}^kn_i(\bar Y_{i\cdot}-\bar{\bar Y})^2
=\frac{\sum_{i=1}^ka_i^{(1)}\bar Y_{i\cdot}^{\,2}}{\sum_{i=1}^k(a_i^{(1)})^2/n_i}
+\frac{\sum_{i=1}^ka_i^{(2)}\bar Y_{i\cdot}^{\,2}}{\sum_{i=1}^k(a_i^{(2)})^2/n_i}
+\cdots+\frac{\sum_{i=1}^ka_i^{(k-1)}\bar Y_{i\cdot}^{\,2}}{\sum_{i=1}^k(a_i^{(k-1)})^2/n_i}$$
以及
\begin{equation}\label{eq:11.2.18}\index{SI0031@对比平方和}
\sum_{i=1}^k\frac{a_i^{(l)}a_i^{(l')}}{n_i}=0\qquad\text{对所有 }l\ne l'.
\end{equation}
所以各个对比平方和互不相关，从而在正态性下相互独立（\autoref{lem:5.3.3}）. 适当归一化后，\eqref{eq:11.2.18}左端服从$\chi^2_{k-1}$，右端是$k-1$个$\chi^2_1$. （这种对比称为正交对比，orthogonal contrasts.\index{SI0044@正交对比}\semanticindexanchor{44:465} 见\autoref{exercise:11.10}和\autoref{exercise:11.11}.）

通常将ANOVA $F$检验的结果汇总成一种标准形式，称为ANOVA表. 表中还给出一些有用的中间统计量，列标题应当不言自明.

\begin{example}[\autoref{ex:11.2.1}续]\label{ex:11.2.12}\index{SI0053@表}
鱼类毒素数据的ANOVA表如下. $F$统计量为26.09，高度显著，表明有强有力的证据说明各毒素产生不同的效应.
\end{example}
\sourcepageanchor{466}
\begin{center}
单向分类的ANOVA表

\small
\begin{tabular}{p{.15\textwidth}cp{.26\textwidth}p{.21\textwidth}c}
\toprule 变异来源&自由度&平方和&均方&$F$统计量\\\midrule
处理组之间&$k-1$&$\mathrm{SSB}=\sum n_i(\bar y_i-\bar{\bar y})^2$&$\mathrm{MSB}=\mathrm{SSB}/(k-1)$&$F=\dfrac{\mathrm{MSB}}{\mathrm{MSW}}$\\[1ex]
处理组内部&$N-k$&$\mathrm{SSW}=\sum\sum(y_{ij}-\bar y_{i\cdot})^2$&$\mathrm{MSW}=$\newline$\mathrm{SSW}/(N-k)$&\\[1ex]
总计&$N-1$&$\mathrm{SST}=\sum\sum(y_{ij}-\bar{\bar y})^2$&&\\\bottomrule
\end{tabular}

\medskip
\begin{tabular}{lrrrr}
\toprule 变异来源&自由度&平方和&均方&$F$统计量\\\midrule
处理&3&995.90&331.97&26.09\\
组内&15&190.83&12.72&\\
总计&18&1,186.73&&\\\bottomrule
\end{tabular}
\end{center}

由\eqref{eq:11.2.15}可知，平方和这一列可以“相加”，即$\mathrm{SSB}+\mathrm{SSW}=\mathrm{SST}$. 类似地，自由度列也可以相加. 但均方列不能，因为它们是平均值，而不是总和.

ANOVA表没有引入新的统计量，只是为计算和呈现提供一种有序的形式. $F$统计量与之前推导的完全相同；此外，MSW就是通常的$\sigma^2$合并无偏估计量，即\eqref{eq:11.2.5}的$S_p^2$（\autoref{exercise:11.22}）.



\subsection{简单线性回归}\label{sec:11.3}\index{SI0878@回归}
在方差分析中，考察了一个因子（变量）怎样影响响应变量的均值. 现在转向简单线性回归，试图更好地理解一个变量对另一个变量的函数依赖. 特别地，在简单线性回归中，有如下形式的关系：
\begin{equation}\label{eq:11.3.1}
Y_i=\alpha+\beta x_i+\varepsilon_i,
\end{equation}
其中$Y_i$为随机变量，$x_i$为另一个可观测变量. $\alpha$与$\beta$分别为回归的截距（intercept）和斜率（slope），假设为固定的未知参数，而$\varepsilon_i$必然是随机变量. 通常还假设$E\varepsilon_i=0$（否则只需将多出的均值调整到$\alpha$中），因此由\eqref{eq:11.3.1}有
\begin{equation}\label{eq:11.3.2}\index{SI0901@总体}
EY_i=\alpha+\beta x_i.
\end{equation}
一般而言，把$EY$表示为$x$的函数的函数称为总体回归函数. \eqref{eq:11.3.2}定义了简单线性回归的总体回归函数.

回归的一个主要目的是利用\eqref{eq:11.3.2}这样的关系，根据已知的$x_i$预测$Y_i$. 日常用语中常将其解释为$Y_i$依赖于$x_i$，并常把$Y_i$称为因变量（dependent variable），$x_i$称为自变量（independent variable）.\index{SI0889@因变量}\semanticindexanchor{889:466}\index{SI0892@自变量}\semanticindexanchor{892:466}
\sourcepageanchor{467}
但这一术语容易引起混淆，因为这里“independent”的含义不同于此前的用法. （$x_i$未必是随机变量，所以不能按通常意义称其为统计上“独立”.）我们不使用这种容易混淆的术语，而采用描述性更强的另一组术语，将$Y_i$称为响应变量（response variable），$x_i$称为预测变量（predictor variable）.\index{SI0903@预测变量}\semanticindexanchor{903:467}\index{SI0906@响应变量}\semanticindexanchor{906:467}

实际上，为明确关于$Y_i$与$x_i$关系的推断以已知$x_i$为前提，可以把\eqref{eq:11.3.2}写为
\begin{equation}\label{eq:11.3.3}\index{SI0884@条件期望}
E(Y_i\mid x_i)=\alpha+\beta x_i.
\end{equation}
我们倾向于使用\eqref{eq:11.3.3}，以强调推断的条件性质.

回顾\autoref{sec:4}，在讨论条件期望时遇到过“回归”一词（见\autoref{exercise:4.13}）. 那里把$Y$对$X$的回归定义为$E(Y\mid x)$，即给定$X=x$时$Y$的条件期望. 更一般地，统计学中“回归”表示变量之间的一种关系. 说回归是线性的，可以指给定$X=x$时$Y$的条件期望是$x$的线性函数. 注意，在\eqref{eq:11.3.3}中，$x_i$是固定已知的，还是可观测随机变量$X_i$的一次实现，并无区别；两种情况下\eqref{eq:11.3.3}的解释相同. 但在\autoref{sec:11.3.4}中将利用$X_i,Y_i$的联合分布进行推断，那时情况就不同了.

术语线性回归是指模型设定\emph{对参数线性}.\index{SI0895@线性}\semanticindexanchor{895:467}\index{SI0896@关于x线性}\semanticindexanchor{896:467} 因而，$E(Y_i\mid x_i)=\alpha+\beta x_i^2$和$E(\log Y_i\mid x_i)=\alpha+\beta(1/x_i)$都规定了线性回归. 前者规定$Y_i$与$x_i^2$之间的线性关系，后者规定$\log Y_i$与$1/x_i$之间的线性关系. 相反，$E(Y_i\mid x_i)=\alpha+\beta^2x_i$并不规定线性回归.

“回归”一词有一段有趣的历史，可追溯至19世纪Francis Galton爵士的工作. （更多细节见\hyperref[bib:source129]{Freedman等（1991）}，深入的历史论述见\hyperref[bib:source308]{Stigler（1986）}.）Galton研究了父亲身高与儿子身高的关系. 不出所料，他发现高个父亲倾向于有高个儿子，矮个父亲倾向于有矮个儿子. 但他也发现，特别高的父亲倾向于有较矮的儿子，特别矮的父亲倾向于有较高的儿子. （想一想——这是有道理的.）Galton把这种现象称为向均值回归（regression toward the mean），\index{SI0910@向均值回归}\semanticindexanchor{910:467}采用“regression”的通常含义“回去”；今天“回归”一词的用法便源于此.

\begin{example}[预测葡萄产量]\label{ex:11.3.1}
回归的一项较现代应用是预测葡萄收成. 七月，葡萄藤结出果串，果串数量可用于预测采收时的最终产量. 典型数据如下，列出若干年的果串数和产量（吨/英亩）.
\begin{center}
\begin{tabular}{rrr}
\toprule 年份&产量（$Y$）&果串数（$x$）\\\midrule
1971&5.6&116.37\\1973&3.2&82.77\\1974&4.5&110.68\\1975&4.2&97.50\\
1976&5.2&115.88\\1977&2.7&80.19\\1978&4.8&125.24\\1979&4.9&116.15\\
1980&4.7&117.36\\1981&4.1&93.31\\1982&4.4&107.46\\1983&5.4&122.30\\\bottomrule
\end{tabular}
\end{center}
1972年的数据缺失，因为那年的作物被飓风摧毁. 这些数据的图形会显示出很强的线性关系.
\end{example}

\sourcepageanchor{468}
写下\eqref{eq:11.3.3}这样的方程时，隐含地假设$Y$对$X$的回归是线性的，即给定$X=x$时$Y$的条件期望是$x$的线性函数. 这一假设未必有依据，因为可能没有支持线性关系的基础理论. 但线性关系处理起来非常方便，所以可能愿意假设$Y$对$X$的回归能用线性函数充分近似. 因而，实际上并不指望\eqref{eq:11.3.3}严格成立，而是希望
\begin{equation}\label{eq:11.3.4}
E(Y_i\mid x_i)\approx\alpha+\beta x_i
\end{equation}
是一种合理的近似. 若从一个（相当强的）假设出发，即$(X_i,Y_i)$服从二元正态分布，就立即得到$Y$对$X$的回归是线性的. 此时条件期望$E(Y\mid x)$对参数线性（见\autoref{def:4.5.10}及其后的讨论）.

最后还需作一个区分. 回归分析，即考察预测变量与响应变量之间的关系，包含两个步骤. 第一步完全以数据为对象，只试图概括观测到的数据.\index{SI0888@数据拟合}\semanticindexanchor{888:468} （这一步总会进行，因为几乎总要计算样本均值、方差或其他汇总统计量；但现在这部分分析往往变得更复杂.）必须记住，这一步“数据拟合”不属于统计推断. 因为只关心手中的数据，不必对参数作任何假设.

回归分析的第二步是统计步骤，试图对总体中的关系，也就是总体回归函数，作推断.\index{SI0893@推断}\semanticindexanchor{893:468} 为此必须对总体作假设. 特别是，若希望对总体线性关系的斜率和截距作推断，就必须假设存在与这些量对应的参数.

简单线性回归问题中，观测数据由$n$对观测$(x_1,y_1),\allowbreak\ldots,\allowbreak(x_n,y_n)$组成. 本节将考虑这些数据的若干不同模型. 各模型对$x$、$y$或二者是否为随机变量$X$、$Y$的观测值有不同假设.

每种模型都关注$x$与$y$之间的线性关系. $n$个数据点不会恰好落在一条直线上，但我们希望通过对观测点拟合直线来概括样本信息. 将看到，许多不同的方法都导向同一条直线.

基于数据$(x_1,y_1),\allowbreak\ldots,\allowbreak(x_n,y_n)$，定义以下各量. 样本均值为
\begin{equation}\label{eq:11.3.5}
\bar x=\frac1n\sum_{i=1}^nx_i\quad\text{且}\quad\bar y=\frac1n\sum_{i=1}^ny_i.
\end{equation}
平方和为
\begin{equation}\label{eq:11.3.6}\index{SI0907@平方和}
S_{xx}=\sum_{i=1}^n(x_i-\bar x)^2\quad\text{且}\quad S_{yy}=\sum_{i=1}^n(y_i-\bar y)^2,
\end{equation}
交叉乘积和为
\begin{equation}\label{eq:11.3.7}
S_{xy}=\sum_{i=1}^n(x_i-\bar x)(y_i-\bar y).
\end{equation}
\eqref{eq:11.3.4}中$\alpha,\beta$最常用的估计分别记为$a,b$，如下给出；稍后将在不同模型下论证它们：
\begin{equation}\label{eq:11.3.8}
b=\frac{S_{xy}}{S_{xx}}\quad\text{且}\quad a=\bar y-b\bar x.
\end{equation}

\sourcepageanchor{469}
\begin{figure}[htbp]
\centering\includegraphics[width=.77\textwidth]{assets/chapter11/fig-11.3.1.pdf}
\caption{\autoref{tab:11.3.1}的数据：RSS所度量的竖直距离}\label{fig:11.3.1}
\end{figure}

\subsubsection{最小二乘：一种数学解法}\label{sec:11.3.1}\index{SI0606@最小二乘估计}
第一次推导$\alpha,\beta$的估计时，不对观测$(x_i,y_i)$作任何统计假设. 只把$(x_1,y_1),\allowbreak\ldots,\allowbreak(x_n,y_n)$看成绘制于散点图中的$n$对数，如\autoref{fig:11.3.1}. （图中24个数据点列于\autoref{tab:11.3.1}.）想象穿过这团点画一条直线，使它“尽可能接近”所有点.

对任意直线$y=c+dx$，定义残差平方和（residual sum of squares，RSS）为\index{SI1013@残差}\semanticindexanchor{1013:469}
$$\mathrm{RSS}=\sum_{i=1}^n\bigl(y_i-(c+dx_i)\bigr)^2.$$
RSS度量每个数据点到直线$c+dx$的\emph{竖直}距离，再将这些距离的平方相加. （\autoref{fig:11.3.1}画出了其中两个距离.）

\sourcepageanchor{470}
\begin{table}[htbp]
\centering\caption{\autoref{fig:11.3.1}中的数据}\label{tab:11.3.1}
\begin{tabular}{*{8}{>{\centering\arraybackslash}p{\dimexpr(\textwidth-16\tabcolsep)/8\relax}}}
\toprule $x$&$y$&$x$&$y$&$x$&$y$&$x$&$y$\\\midrule
3.74&3.22&0.20&2.81&1.22&1.23&1.76&4.12\\
3.66&4.87&2.50&3.71&1.00&3.13&0.51&3.16\\
0.78&0.12&3.50&3.11&1.29&4.05&2.17&4.40\\
2.40&2.31&1.35&0.90&0.95&2.28&1.99&1.18\\
2.18&4.25&2.36&4.39&1.05&3.60&1.53&2.54\\
1.93&2.24&3.13&4.36&2.92&5.39&2.60&4.89\\\midrule
\multicolumn{8}{c}{$\bar x=1.95\quad\bar y=3.18\quad S_{xx}=22.82\quad S_{yy}=43.62\quad S_{xy}=15.48$}\\\bottomrule
\end{tabular}
\end{table}
$\alpha,\beta$的最小二乘估计（least squares estimates）定义为使直线$a+bx$的RSS最小的$a,b$值. 即最小二乘估计$a,b$满足
$$\min_{c,d}\sum_{i=1}^n\bigl(y_i-(c+dx_i)\bigr)^2
=\sum_{i=1}^n\bigl(y_i-(a+bx_i)\bigr)^2.$$

这个关于$c,d$两个变量的函数可如下最小化. 对任意固定的$d$，写成
$$\sum_{i=1}^n\bigl(y_i-(c+dx_i)\bigr)^2=\sum_{i=1}^n\bigl((y_i-dx_i)-c\bigr)^2,$$
便可求出给出最小值的$c$. 由\autoref{thm:5.2.4}，使其最小的$c$值为
\begin{equation}\label{eq:11.3.9}\index{SI0614@截距的}
c=\frac1n\sum_{i=1}^n(y_i-dx_i)=\bar y-d\bar x.
\end{equation}
因此，给定$d$时，RSS的最小值为
\begin{align*}
\sum_{i=1}^n\bigl((y_i-dx_i)-(\bar y-d\bar x)\bigr)^2
&=\sum_{i=1}^n\bigl((y_i-\bar y)-d(x_i-\bar x)\bigr)^2\\
&=S_{yy}-2dS_{xy}+d^2S_{xx}.
\end{align*}
将这个关于$d$的二次函数求导并令导数为零，即得使RSS整体最小的$d$：
\begin{equation}\label{eq:11.3.10}\index{SI0615@斜率的}
d=\frac{S_{xy}}{S_{xx}}.
\end{equation}
由于$d^2$的系数为正，该值确实给出最小值. 因此，由\eqref{eq:11.3.9}和\eqref{eq:11.3.10}，\eqref{eq:11.3.8}的$a,b$就是使残差平方和最小的$c,d$值.

RSS只是度量直线$c+dx$到数据点距离的许多合理方法之一. 例如，也可使用水平距离来代替竖直距离. 这等价于将$y$变量画在横轴、$x$变量画在纵轴，然后像上面那样使用竖直距离. 利用上述结果，交换$x,y$的角色，得到最小二乘直线$\hat x=a'+b'y$，其中
$$b'=\frac{S_{xy}}{S_{yy}}\quad\text{且}\quad a'=\bar x-b'\bar y.$$
把直线重新表达为$y$是$x$的函数，得$\hat y=-(a'/b')+(1/b')x$.

通常，由水平距离得到的直线不同于由竖直距离得到的直线. 利用\autoref{tab:11.3.1}的数值，$y$对$x$的回归（竖直距离）
\sourcepageanchor{471}
为$\hat y=1.86+.68x$. $x$对$y$的回归（水平距离）为$\hat y=-2.31+2.82x$. \autoref{fig:12.2.1}画出这两条直线，以及\autoref{sec:12.2}讨论的第三条直线. 若两条直线相同，则斜率相同，$b/(1/b')$应等于1. 但实际上$b/(1/b')\le1$，只在特殊情况下取等. 注意
$$\frac{b}{1/b'}=bb'=\frac{(S_{xy})^2}{S_{xx}S_{yy}}.$$
在\eqref{eq:4.7.9}的H\"older不等式中取$p=q=2$、$a_i=x_i-\bar x$、$b_i=y_i-\bar y$，得$(S_{xy})^2\le S_{xx}S_{yy}$，所以该比值小于1.\index{SI0568@Hölder}\semanticindexanchor{568:471}

若$x$是预测变量，$y$是响应变量，且考虑由$x$预测$y$，则RSS所度量的竖直距离是合理的，它度量$y_i$与其预测值$\hat y_i=c+dx_i$之间的距离. 但若不对$x,y$作这种区分，另一个合理准则——水平距离——却给出不同的直线，就令人不安.

最小二乘法应当仅被视为对一组数据“拟合直线”的方法，\index{SI0888@数据拟合}\semanticindexanchor{888:471}而不应被视为统计推断方法. 在本节没有对数据使用任何统计模型，因此没有构造置信区间或检验假设的依据. 在本节背景下谈论$a,b$，称它们为最小二乘\emph{解}可能比称最小二乘\emph{估计}更合适，因为它们是最小化RSS这一数学问题的解，而不是从统计模型推导出的估计. 不过，正如将看到的，这些最小二乘解在某些统计模型中具有最优性.

\subsubsection{最优线性无偏估计量：一种统计解法}\label{sec:11.3.2}
本节说明，在相当一般的统计模型下，\eqref{eq:11.3.8}的估计$a,b$在线性无偏估计类中最优. 模型如下：假设$x_1,\ldots,x_n$是已知的固定值. （可把它们看作实验人员在实验室实验中选定并设定的值.）$y_1,\ldots,y_n$是互不相关随机变量$Y_1,\ldots,Y_n$的观测值. 假设$x$与$y$之间的线性关系为
\begin{equation}\label{eq:11.3.11}
EY_i=\alpha+\beta x_i,\qquad i=1,\ldots,n,
\end{equation}
并假设
\begin{equation}\label{eq:11.3.12}
\operatorname{Var}Y_i=\sigma^2.
\end{equation}
$\sigma^2$没有下标，因为假设所有$Y_i$都有相同的（未知）方差. 关于$Y_i$前两阶矩的这些假设，就是本小节推导所需的全部假设. 例如，不必指定$Y_1,\ldots,Y_n$的概率分布.

\eqref{eq:11.3.11}和\eqref{eq:11.3.12}的模型也可以这样表述：假设
\begin{equation}\label{eq:11.3.13}
Y_i=\alpha+\beta x_i+\varepsilon_i,\qquad i=1,\ldots,n,
\end{equation}
其中$\varepsilon_1,\ldots,\varepsilon_n$是互不相关的随机变量，满足
\begin{equation}\label{eq:11.3.14}
E\varepsilon_i=0\quad\text{且}\quad\operatorname{Var}\varepsilon_i=\sigma^2.
\end{equation}
$\varepsilon_1,\ldots,\varepsilon_n$称为随机误差（random errors）. 由于$Y_i$只依赖$\varepsilon_i$，且各$\varepsilon_i$互不相关，所以$Y_i$互不相关. 由\eqref{eq:11.3.13}和\eqref{eq:11.3.14}也容易验证\eqref{eq:11.3.11}、\eqref{eq:11.3.12}中$EY_i$和$\operatorname{Var}Y_i$的表达式.

为推导参数$\alpha,\beta$的估计量，将注意力限制在线性估计量类. 若一个估计量具有形式
\begin{equation}\label{eq:11.3.15}\index{SI0792@线性}
\sum_{i=1}^nd_iY_i,
\end{equation}
\sourcepageanchor{472}
其中$d_1,\ldots,d_n$是已知的固定常数，则称为线性估计量（linear estimator）. （\autoref{exercise:7.39}讨论了总体均值的线性估计量.）在线性估计量类中，再限制为无偏估计量，这就限制了$d_1,\ldots,d_n$允许取的值.

无论参数$\alpha,\beta$的真值为何，斜率$\beta$的无偏估计量必须满足
$$E\sum_{i=1}^nd_iY_i=\beta.$$
这意味着
\begin{align*}
\beta&=E\sum_{i=1}^nd_iY_i=\sum_{i=1}^nd_iEY_i=\sum_{i=1}^nd_i(\alpha+\beta x_i)\\
&=\alpha\left(\sum_{i=1}^nd_i\right)+\beta\left(\sum_{i=1}^nd_ix_i\right).
\end{align*}
这一等式对所有$\alpha,\beta$成立，当且仅当
\begin{equation}\label{eq:11.3.16}
\sum_{i=1}^nd_i=0\quad\text{且}\quad\sum_{i=1}^nd_ix_i=1.
\end{equation}
所以，要使该估计量成为$\beta$的无偏估计量，$d_1,\ldots,d_n$必须满足\eqref{eq:11.3.16}.

在\autoref{sec:7}中，如果一个无偏估计量在所有无偏估计量中方差最小，就称其为“最优”. 类似地，若一个线性无偏估计量在线性无偏估计量中方差最小，就称为最优线性无偏估计量（best linear unbiased estimator，BLUE）. 现在说明，定义估计量$b=S_{xY}/S_{xx}$的选择$d_i=(x_i-\bar x)/S_{xx}$是最优的，因为它产生$\beta$的方差最小的线性无偏估计量. （$d_i$必须是已知的固定常数，而$x_i$是已知固定常数，所以这一选择合法.）

关于记号的说明：记号$S_{xY}$强调它是随机变量$Y_1,\ldots,Y_n$的函数，因而是随机变量. $S_{xY}$也依赖非随机的量$x_1,\ldots,x_n$.

由于$Y_1,\ldots,Y_n$互不相关且方差同为$\sigma^2$，任意线性估计量的方差为
$$\operatorname{Var}\sum_{i=1}^nd_iY_i
=\sum_{i=1}^nd_i^2\operatorname{Var}Y_i
=\sum_{i=1}^nd_i^2\sigma^2=\sigma^2\sum_{i=1}^nd_i^2.$$
因此，$\beta$的BLUE由满足\eqref{eq:11.3.16}并使$\sum_{i=1}^nd_i^2$最小的常数$d_1,\ldots,d_n$定义. （$\sigma^2$不影响在线性估计量上的最小化，因为每个线性估计量的方差都以它为乘数.）

现在可用\autoref{lem:11.2.7}求使其最小的$d_1,\ldots,d_n$. 将引理用于这个最小化问题时，作如下对应，其中左端是\hyperref[lem:11.2.7]{引理11.2.1}中的记号，右端是当前记号.\footnote{译注：原文此处再次写引理11.2.1，实际使用的是\autoref{lem:11.2.7}. 保留原显示编号.} 令
$$k=n,\quad v_i=x_i,\quad c_i=1,\quad a_i=d_i,$$
则$\bar v_c=\bar x$. 若$d_i$具有形式
\begin{equation}\label{eq:11.3.17}
d_i=Kc_i(v_i-\bar v_c)=K(x_i-\bar x),\qquad i=1,\ldots,n,
\end{equation}
则由\autoref{lem:11.2.7}，$d_1,\ldots,d_n$在所有满足$\sum d_i=0$的向量中使
\begin{equation}\label{eq:11.3.18}\index{SI0681@二次函数的最大值}
\frac{(\sum_{i=1}^nd_ix_i)^2}{\sum_{i=1}^nd_i^2}
\end{equation}
达到最大.

\sourcepageanchor{473}
此外，由于
$$\{(d_1,\ldots,d_n):\sum d_i=0,\ \sum d_ix_i=1\}
\subset\{(d_1,\ldots,d_n):\sum d_i=0\},$$
如果形如\eqref{eq:11.3.17}的$d_i$还满足\eqref{eq:11.3.16}，它们当然也在所有满足\eqref{eq:11.3.16}的$d_1,\ldots,d_n$中使\eqref{eq:11.3.18}最大. （取最大值的集合变小，最大值不可能增大.）现在利用\eqref{eq:11.3.17}，有
$$\sum_{i=1}^nd_ix_i=\sum_{i=1}^nK(x_i-\bar x)x_i=KS_{xx}.$$
若$K=1/S_{xx}$，就满足\eqref{eq:11.3.16}的第二个约束. 因此，定义
\begin{equation}\label{eq:11.3.19}\index{SI0882@斜率的最优线性无偏估计量}
d_i=\frac{x_i-\bar x}{S_{xx}},\qquad i=1,\ldots,n,
\end{equation}
就满足\eqref{eq:11.3.16}的两个约束，且这一组$d_i$达到最大值. 最后注意，对所有满足\eqref{eq:11.3.16}的$d_1,\ldots,d_n$，
$$\frac{(\sum_{i=1}^nd_ix_i)^2}{\sum_{i=1}^nd_i^2}=\frac1{\sum_{i=1}^nd_i^2}.$$
因此，对满足\eqref{eq:11.3.16}的$d_1,\ldots,d_n$，最大化\eqref{eq:11.3.18}等价于最小化$\sum d_i^2$. 所以可断定，\eqref{eq:11.3.19}定义的$d_i$在所有满足\eqref{eq:11.3.16}的$d_i$中使$\sum d_i^2$最小，而它们定义的线性无偏估计量
$$b=\sum_{i=1}^n\frac{x_i-\bar x}{S_{xx}}y_i=\frac{S_{xy}}{S_{xx}}$$
就是$\beta$的BLUE.

\autoref{fig:11.3.2}给出了构造$\beta$的BLUE的几何描述，其中取$n=3$. 图中是以$d_1,d_2,d_3$为坐标的三维空间. 两个平面分别代表满足\eqref{eq:11.3.16}两个线性约束的向量$(d_1,d_2,d_3)$，两平面的交线代表同时满足两个等式的向量. 对直线上任意点，$\sum_{i=1}^nd_i^2$是该点到原点$\boldsymbol0$距离的平方. 定义BLUE的向量$(d_1,d_2,d_3)$就是直线上距原点最近的点. 图中的球面是与直线相交的最小球面，交点就是定义$\beta$的BLUE的点$(d_1,d_2,d_3)$. 已经证明，该点满足$d_i=(x_i-\bar x)/S_{xx}$.

$b$的方差为
\begin{equation}\label{eq:11.3.20}
\operatorname{Var}b=\sigma^2\sum_{i=1}^nd_i^2=\frac{\sigma^2}{S_{xx}}
=\frac{\sigma^2}{\sum_{i=1}^n(x_i-\bar x)^2}.
\end{equation}
由于$x_1,\ldots,x_n$是实验人员选定的值，\index{SI0890@设计}\semanticindexanchor{890:473}可以选取它们，使$S_{xx}$较大而估计量方差较小. 也就是说，实验人员可通过设计实验使估计更精确. 假设所有$x_1,\ldots,x_n$都必须在区间$[e,f]$中选取. 若$n$为偶数，使$S_{xx}$尽可能大的选择是令一半$x_i$等于$e$，另一半等于$f$（\autoref{exercise:11.26}）. 若实验人员确信\eqref{eq:11.3.11}和\eqref{eq:11.3.12}的模型正确，这就是最优设计，因为它给出斜率$\beta$最精确的估计. 然而实践中很少使用这种设计，因为实验人员几乎从不确信模型. 这种两点设计只提供$x=e$和$x=f$两个值处的$E(Y\mid x)$信息. 如果将$Y$的均值表示为$x$函数的总体回归函数$E(Y\mid x)$是非线性的，从这种“最优”两点设计所得数据中永远无法发现这一点.

\sourcepageanchor{474}
\begin{figure}[htbp]
\centering\includegraphics[width=.65\textwidth]{assets/chapter11/fig-11.3.2.pdf}
\caption{BLUE的几何描述}\label{fig:11.3.2}
\end{figure}
已经证明$b$是$\beta$的BLUE. 类似分析可说明$a$是截距$\alpha$的BLUE. 定义$\alpha$的线性估计量的常数$d_1,\ldots,d_n$必须满足
\begin{equation}\label{eq:11.3.21}\index{SI0881@截距的最优线性无偏估计量}
\sum_{i=1}^nd_i=1\quad\text{且}\quad\sum_{i=1}^nd_ix_i=0.
\end{equation}
推导细节留作\autoref{exercise:11.27}. 最小二乘估计量为BLUE这一事实也在其他线性模型中成立.\index{SI0451@Gauss–Markov定理}\semanticindexanchor{451:474} 这个一般结果称为Gauss--Markov定理（Gauss--Markov Theorem）（见\hyperref[bib:source75]{Christensen，1996}；\hyperref[bib:source200]{Lehmann和Casella，1998}，第3.4节；更一般的论述见\hyperref[bib:source160]{Harville，1981）}.

\subsubsection{模型与分布假设}\label{sec:11.3.3}\index{SI0895@线性}
本节再为配对数据$(x_1,y_1),\allowbreak\ldots,\allowbreak(x_n,y_n)$引入两个模型，称为简单线性回归模型.

为了得到\hyperref[sec:11.3.1]{第12.2.1节}中的最小二乘估计，我们没有使用统计模型，而只是解了一个数学上的最小化问题. 因为没有概率模型可用，所以不能推导这种方法所得估计量的任何统计性质. 实际上也没有可为之构造假设检验或置信区间的参数.\footnote{译注：原文写第12.2.1节，但这里回顾的是\autoref{sec:11.3.1}的最小二乘推导. 保留原显示节号.}

在\autoref{sec:11.3.2}，对数据作了一些统计假设. 具体而言，对数据前两阶矩——均值、方差和协方差——作了假设. 它们都是与数据概率模型有关的统计假设，据此推导了估计量的统计性质. 对参数$\alpha,\beta$的估计量$a,b$证明的无偏性和最小方差性质，都属于统计性质.

为得到这些性质，无须指定数据的完整概率模型，只需关于前两阶矩的假设. 在这些最低限度的假设下，得到了一种一般最优性，但只限于受限估计量类——线性无偏估计量. 在该模型下无法推导精确检验与置信区间，因为模型对数据概率分布规定得还不够. 下面给出两个完整规定数据概率结构的统计模型.

\paragraph{11.3.3.1\quad 条件正态模型}\label{sec:11.3.3.1}\index{SI0885@条件正态}
条件正态模型（conditional normal model）是最常见、分析起来最直接的简单线性回归模型. 观测数据为$n$对$(x_1,y_1),\allowbreak\ldots,\allowbreak(x_n,y_n)$. 预测变量的值$x_1,\ldots,x_n$视为已知的固定常数. 如\autoref{sec:11.3.2}，可认为它们是实验人员选定并设定的值.
\sourcepageanchor{475}
响应变量的值$y_1,\ldots,y_n$是随机变量$Y_1,\ldots,Y_n$的观测值. 假设$Y_1,\ldots,Y_n$相互独立，并且$Y_i$服从正态分布，具体为
\begin{equation}\label{eq:11.3.22}
Y_i\sim\mathrm n(\alpha+\beta x_i,\sigma^2),\qquad i=1,\ldots,n.
\end{equation}
所以总体回归函数是$x$的线性函数，即$E(Y\mid x)=\alpha+\beta x$，所有$Y_i$的方差同为$\sigma^2$. 条件正态模型可表示成类似\eqref{eq:11.3.13}、\eqref{eq:11.3.14}的形式：
\begin{equation}\label{eq:11.3.23}
Y_i=\alpha+\beta x_i+\varepsilon_i,\qquad i=1,\ldots,n,
\end{equation}
其中$\varepsilon_1,\ldots,\varepsilon_n$独立同分布于$\mathrm n(0,\sigma^2)$.

条件正态模型是\autoref{sec:11.3.2}所考虑模型的特例. 总体回归函数$E(Y\mid x)=\alpha+\beta x$和方差$\operatorname{Var}Y=\sigma^2$都与该模型相同. $Y_1,\ldots,Y_n$（等价地，$\varepsilon_1,\ldots,\varepsilon_n$）的不相关性被加强为独立性. 当然，现在不仅指定$Y_1,\ldots,Y_n$分布的前两阶矩，还指定了概率分布的精确形式.

由于独立性，$Y_1,\ldots,Y_n$的联合pdf为边缘pdf的乘积，即
\begin{equation}\label{eq:11.3.24}
\begin{aligned}
f(\boldsymbol y\mid\alpha,\beta,\sigma^2)
&=f(y_1,\ldots,y_n\mid\alpha,\beta,\sigma^2)\\
&=\prod_{i=1}^nf(y_i\mid\alpha,\beta,\sigma^2)\\
&=\prod_{i=1}^n\frac1{\sqrt{2\pi}\sigma}\exp\left[-\frac{(y_i-(\alpha+\beta x_i))^2}{2\sigma^2}\right]\\
&=\frac1{(2\pi)^{n/2}\sigma^n}\exp\left[-\frac{\sum_{i=1}^n(y_i-\alpha-\beta x_i)^2}{2\sigma^2}\right].
\end{aligned}
\end{equation}
\autoref{sec:11.3.4}和\autoref{sec:11.3.5}将使用这一联合概率分布建立统计程序. 例如，用\eqref{eq:11.3.24}求$\alpha,\beta,\sigma^2$的MLE.

\paragraph{11.3.3.2\quad 二元正态模型}\label{sec:11.3.3.2}\index{SI0880@二元正态}
此前讨论的所有模型中，预测变量的值$x_1,\ldots,x_n$都是固定的已知常数. 但有时这些值实际上是随机变量$X_1,\ldots,X_n$的观测值. 在\autoref{sec:11.3}的Galton例子中，$x_1,\ldots,x_n$是观测到的父亲身高，实验人员显然没有在收集数据前选定这些身高. 因此必须考虑预测变量和响应变量都随机的模型. 一种相当简单的模型是二元正态模型（bivariate normal model）. \autoref{sec:12.2}将讨论更复杂的模型.

在二元正态模型中，数据$(x_1,y_1),\allowbreak\ldots,\allowbreak(x_n,y_n)$是二元随机向量$(X_1,Y_1),\allowbreak\ldots,\allowbreak(X_n,Y_n)$的观测值. 各随机向量独立，且假设$(X_i,Y_i)$的联合分布为二元正态. 具体假设为
$$ (X_i,Y_i)\sim\text{二元正态}(\mu_X,\mu_Y,\sigma_X^2,\sigma_Y^2,\rho). $$
二元正态分布的联合pdf及各种性质，见\autoref{def:4.5.10}及其后的讨论. 全部数据$(X_1,Y_1),\allowbreak\ldots,\allowbreak(X_n,Y_n)$的联合pdf是这些二元pdf的乘积.

在简单线性回归分析中，仍把$x$看成预测变量，$y$看成响应变量，即最关心的是观测$x$值以后预测$y$值.
\sourcepageanchor{476}
这自然引导我们基于给定$X=x$时$Y$的条件分布进行推断. 对二元正态模型，给定$X=x$时$Y$的条件分布为正态. 总体回归函数现在确实是一个条件期望，正如记号所提示的那样，且为
\begin{equation}\label{eq:11.3.25}
E(Y\mid x)=\mu_Y+\rho\frac{\sigma_Y}{\sigma_X}(x-\mu_X)
=\left[\mu_Y-\rho\frac{\sigma_Y}{\sigma_X}\mu_X\right]
+\left[\rho\frac{\sigma_Y}{\sigma_X}\right]x.
\end{equation}
二元正态模型\emph{蕴含}总体回归是$x$的线性函数，无须像此前模型那样另外假设. 这里$E(Y\mid x)=\alpha+\beta x$，其中$\beta=\rho\sigma_Y/\sigma_X$，$\alpha=\mu_Y-\rho(\sigma_Y/\sigma_X)\mu_X$. 与条件正态模型一样，响应变量$Y$的条件方差也不依赖$x$：
\begin{equation}\label{eq:11.3.26}
\operatorname{Var}(Y\mid x)=\sigma_Y^2(1-\rho^2).
\end{equation}

对二元正态模型，线性回归分析几乎总是利用给定$X_1=x_1,\ldots,X_n=x_n$时$(Y_1,\ldots,Y_n)$的条件分布，而非$(X_1,Y_1),\allowbreak\ldots,\allowbreak(X_n,Y_n)$的无条件分布. 这样就回到上面描述的条件正态模型的情形. 只要以$x_1,\ldots,x_n$为条件，它们原本是随机变量观测值这一事实就无关紧要. 一般而言，简单线性回归中除用二元正态性确定条件分布外，并不使用二元正态性. （事实上，在大多数情况下，$X$的边缘分布完全无关紧要；线性回归中重要的是条件分布.）两种模型基于点估计量、区间或检验的推断相同. 另一种观点见\hyperref[bib:source52]{Brown（1990b）}.

\subsubsection{正态误差下的估计与检验}\label{sec:11.3.4}\index{SI0898@最大似然估计量}
本小节及下一小节在条件正态模型，即\eqref{eq:11.3.22}或\eqref{eq:11.3.23}定义的回归模型下建立推断程序.

首先求三个参数$\alpha,\beta,\sigma^2$的最大似然估计. 由\eqref{eq:11.3.24}的联合pdf，对数似然函数为
$$\log L(\alpha,\beta,\sigma^2\mid\boldsymbol x,\boldsymbol y)
=-\frac n2\log(2\pi)-\frac n2\log\sigma^2
-\frac{\sum_{i=1}^n(y_i-\alpha-\beta x_i)^2}{2\sigma^2}.$$
对任意固定的$\sigma^2$，使$\log L$作为$\alpha,\beta$的函数最大的值$\hat\alpha,\hat\beta$，就是使
$$\sum_{i=1}^n(y_i-\alpha-\beta x_i)^2$$
最小的值. 但这个函数就是\autoref{sec:11.3.1}中的RSS！那里已求得使其最小的值为
$$\hat\beta=b=\frac{S_{xy}}{S_{xx}}\quad\text{且}\quad
\hat\alpha=a=\bar y-b\bar x=\bar y-\hat\beta\bar x.$$
所以$\alpha,\beta$的最小二乘估计也是它们的MLE. 对任意固定的$\sigma^2$，$\hat\alpha,\hat\beta$都是使其最大的值. 现在将它们代入对数似然，为求$\sigma^2$的MLE，只需最大化
$$-\frac n2\log(2\pi)-\frac n2\log\sigma^2
-\frac{\sum_{i=1}^n(y_i-\hat\alpha-\hat\beta x_i)^2}{2\sigma^2}.$$
这与通常正态抽样中求$\sigma^2$的MLE类似（\autoref{ex:7.2.11}），细节留作\autoref{exercise:11.28}. 条件正态模型下$\sigma^2$的MLE为
$$\hat\sigma^2=\frac1n\sum_{i=1}^n(y_i-\hat\alpha-\hat\beta x_i)^2,$$
\sourcepageanchor{477}
即在最小二乘直线处计算的RSS除以样本量. 此后提到RSS，均指在最小二乘直线处计算的RSS.

\autoref{sec:11.3.2}说明$\hat\alpha,\hat\beta$是$\alpha,\beta$的线性无偏估计量. 但$\hat\sigma^2$并不是$\sigma^2$的无偏估计量. 计算$E\hat\sigma^2$以及后续许多计算时，下面的引理会有用.

\begin{lemma}\label{lem:11.3.2}
设$Y_1,\ldots,Y_n$是互不相关的随机变量，且对所有$i=1,\ldots,n$有$\operatorname{Var}Y_i=\sigma^2$. 设$c_1,\ldots,c_n$与$d_1,\ldots,d_n$是两组常数，则
$$\operatorname{Cov}\left(\sum_{i=1}^nc_iY_i,\sum_{i=1}^nd_iY_i\right)
=\left(\sum_{i=1}^nc_id_i\right)\sigma^2.$$
\end{lemma}
\begin{proof}
此前已遇到过此类结果，它类似于\autoref{lem:5.3.3}和\autoref{exercise:11.11}. 不过，这里既不需要$Y_1,\ldots,Y_n$正态，也不需要它们独立.
\end{proof}

接下来求$\sigma^2$中的偏倚.\index{SI0899@有偏}\semanticindexanchor{899:477}\footnote{译注：原文此句写“$\sigma^2$中的偏倚”，实际计算的是估计量$\hat\sigma^2$的偏倚.} 由\eqref{eq:11.3.23}有
$$\varepsilon_i=Y_i-\alpha-\beta x_i.$$
定义回归残差（residuals）为
\begin{equation}\label{eq:11.3.27}\index{SI0905@残差}
\hat\varepsilon_i=Y_i-\hat\alpha-\hat\beta x_i,
\end{equation}
因而
$$\hat\sigma^2=\frac1n\sum_{i=1}^n\hat\varepsilon_i^2=\frac1n\mathrm{RSS}.$$
经计算（\autoref{exercise:11.29}）可得
$$E\hat\varepsilon_i=0,$$
再经过较长的计算（同见\autoref{exercise:11.29}），得
\begin{equation}\label{eq:11.3.28}
\operatorname{Var}\hat\varepsilon_i=E\hat\varepsilon_i^2
=\left(\frac{n-2}n+\frac1{S_{xx}}\left(\frac1n\sum_{j=1}^nx_j^2+x_i^2-2(x_i-\bar x)^2-2x_i\bar x\right)\right)\sigma^2.
\end{equation}
因此
\begin{align*}
E\hat\sigma^2
&=\frac1n\sum_{i=1}^nE\hat\varepsilon_i^2\\
&=\frac1n\sum_{i=1}^n\left[\frac{n-2}n+\frac1{S_{xx}}\left(\frac1n\sum_{j=1}^nx_j^2+x_i^2-2(x_i-\bar x)^2-2x_i\bar x\right)\right]\sigma^2\\
&=\left[\frac{n-2}n+\frac1{nS_{xx}}\left\{\sum_{j=1}^nx_j^2+\sum_{i=1}^nx_i^2-2S_{xx}-\frac2n\left(\sum_{i=1}^nx_i\right)^2\right\}\right]\sigma^2\\
&\hspace{7em}\text{（因为 }\sum x_i\bar x=\tfrac1n(\sum x_i)^2\text{）}\\
&=\left(\frac{n-2}n+0\right)\sigma^2
\qquad\text{（因为 }\sum x_i^2-\tfrac1n(\sum x_i)^2=S_{xx}\text{）}\\
&=\frac{n-2}n\sigma^2.
\end{align*}

\sourcepageanchor{478}
MLE $\hat\sigma^2$是$\sigma^2$的有偏估计量. 更常用的$\sigma^2$估计量是无偏的，定义为
\begin{equation}\label{eq:11.3.29}
S^2=\frac n{n-2}\hat\sigma^2
=\frac1{n-2}\sum_{i=1}^n(y_i-\hat\alpha-\hat\beta x_i)^2
=\frac1{n-2}\sum_{i=1}^n\hat\varepsilon_i^2.
\end{equation}
要基于这些估计量建立估计和检验程序，需要知道其抽样分布. 下面的定理将其汇总.

\begin{theorem}\label{thm:11.3.3}\index{SI0900@抽样分布}
在条件正态回归模型\eqref{eq:11.3.22}下，估计量$\hat\alpha,\hat\beta,S^2$的抽样分布为
$$\hat\alpha\sim\mathrm n\left(\alpha,\frac{\sigma^2}{nS_{xx}}\sum_{i=1}^nx_i^2\right),
\qquad\hat\beta\sim\mathrm n\left(\beta,\frac{\sigma^2}{S_{xx}}\right),$$
并有
$$\operatorname{Cov}(\hat\alpha,\hat\beta)\frac{-\sigma^2\bar x}{S_{xx}}.$$
\footnote{译注：原书协方差式中漏排等号；应读作$\operatorname{Cov}(\hat\alpha,\hat\beta)=-\sigma^2\bar x/S_{xx}$. 此处保留原式.}
此外，$(\hat\alpha,\hat\beta)$与$S^2$独立，且
$$\frac{(n-2)S^2}{\sigma^2}\sim\chi^2_{n-2}.$$
\end{theorem}
\begin{proof}
首先说明$\hat\alpha,\hat\beta$具有所述正态分布. 二者都是独立正态随机变量$Y_1,\ldots,Y_n$的线性函数，所以由\autoref{cor:4.6.10}，二者均为正态分布. 具体而言，\autoref{sec:11.3.2}证明了$\hat\beta=\sum_{i=1}^nd_iY_i$，其中$d_i$由\eqref{eq:11.3.19}给出，并证明了
$$E\hat\beta=\beta\quad\text{且}\quad\operatorname{Var}\hat\beta=\frac{\sigma^2}{S_{xx}}.$$
估计量$\hat\alpha=\bar Y-\hat\beta\bar x$可表示为$\hat\alpha=\sum_{i=1}^nc_iY_i$，其中
$$c_i=\frac1n-\frac{(x_i-\bar x)\bar x}{S_{xx}},$$
因此容易验证
$$E\hat\alpha=\sum_{i=1}^nc_iEY_i
=\sum_{i=1}^n\left(\frac1n-\frac{(x_i-\bar x)\bar x}{S_{xx}}\right)(\alpha+\beta x_i)=\alpha,$$
以及
$$\operatorname{Var}\hat\alpha=\sigma^2\sum_{i=1}^nc_i^2
=\sigma^2\left[\frac1{nS_{xx}}\sum_{i=1}^nx_i^2\right],$$
从而$\hat\alpha,\hat\beta$具有指定的分布. 利用\autoref{lem:11.3.2}也容易求出$\operatorname{Cov}(\hat\alpha,\hat\beta)$. 细节留作\autoref{exercise:11.30}.

接下来说明$\hat\alpha,\hat\beta$与$S^2$独立，这将由\autoref{lem:11.3.2}和\autoref{lem:5.3.3}得到. 根据\eqref{eq:11.3.27}对$\hat\varepsilon_i$的定义，可写为\footnote{译注：原式求和指标是$j$，但最后的随机变量写成$Y_i$；由残差定义应为$Y_j$. 原式照录.}
\begin{equation}\label{eq:11.3.30}
\hat\varepsilon_i=\sum_{j=1}^n[\delta_{ij}-(c_j+d_jx_i)]Y_i,
\end{equation}
其中
$$\delta_{ij}=\begin{cases}1,&i=j,\\0,&i\ne j,\end{cases}
\qquad c_j=\frac1n-\frac{(x_j-\bar x)\bar x}{S_{xx}},\qquad d_j=\frac{x_j-\bar x}{S_{xx}}.$$

\sourcepageanchor{479}
由于$\hat\alpha=\sum c_iY_i$、$\hat\beta=\sum d_iY_i$，应用\autoref{lem:11.3.2}再作一些代数运算，可得
$$\operatorname{Cov}(\hat\varepsilon_i,\hat\alpha)=\operatorname{Cov}(\hat\varepsilon_i,\hat\beta)=0,\qquad i=1,\ldots,n.$$
细节留作\autoref{exercise:11.31}. 因此，由\autoref{lem:5.3.3}，在正态抽样下，$S^2=\sum\hat\varepsilon_i^2/(n-2)$与$\hat\alpha,\hat\beta$独立.

为证明$(n-2)S^2/\sigma^2\sim\chi^2_{n-2}$，将$(n-2)S^2$写成$n-2$个相互独立、各服从$\chi^2_1$分布的随机变量之和.\footnote{译注：此句遗漏除以$\sigma^2$的标准化. 下述系数除已列条件外，还应取为与$(x_1,\ldots,x_n)$正交且各列单位长度，使对应线性组合均值为0、方差为$\sigma^2$；原文明确略去细节，所列内容照录.} 即找常数$a_{ij}$，$i=1,\ldots,n$，$j=1,\ldots,n-2$，满足
\begin{equation}\label{eq:11.3.31}
\sum_{i=1}^n\hat\varepsilon_i^2
=\sum_{j=1}^{n-2}\left(\sum_{i=1}^na_{ij}Y_i\right)^2,
\end{equation}
其中
$$\sum_{i=1}^na_{ij}=0,\quad j=1,\ldots,n-2,
\qquad\text{且}\qquad\sum_{i=1}^na_{ij}a_{ij'}=0,\quad j\ne j'.$$
由于$x_i$具有一般性，细节较为繁复，故略去.
\end{proof}

线性回归的RSS包含了高阶多项式拟合相较线性拟合是否有价值的信息. 由于本模型假设总体回归线性，高阶拟合中的变异只是随机变异. \hyperref[bib:source257]{Robson（1959）}给出了寻找这类高阶多项式拟合系数的一般递推公式，经调整可明确求出\eqref{eq:11.3.31}中的$a_{ij}$. 另外，也可以用Cochran定理（\hyperref[sec:11.5.1]{杂录11.5.1}）证明$\sum\hat\varepsilon_i^2/\sigma^2\sim\chi^2_{n-2}$.

关于两个参数$\alpha,\beta$的推断通常基于以下两个Student $t$分布，它们立即由\autoref{thm:11.3.3}中的正态分布、卡方分布及独立性得到：
\begin{equation}\label{eq:11.3.32}
\frac{\hat\alpha-\alpha}{S\sqrt{(\sum_{i=1}^nx_i^2)/(nS_{xx})}}\sim t_{n-2},
\end{equation}
以及
\begin{equation}\label{eq:11.3.33}\index{SI0908@t统计量}
\frac{\hat\beta-\beta}{S/\sqrt{S_{xx}}}\sim t_{n-2}.
\end{equation}
这两个$t$统计量的联合分布称为二元Student $t$分布\index{SI0275@二元t}\semanticindexanchor{275:479}（bivariate Student's $t$ distribution），其推导方式类似单变量情形. 这里利用了$\hat\alpha,\hat\beta$的联合分布为二元正态，以及两个单变量$t$统计量使用同一个方差估计$S$. 若要对$\alpha,\beta$作同时推断，就要使用这个联合分布. 不过，这里只处理一次关于一个参数的推断.

通常人们对$\beta$比对$\alpha$更感兴趣. 参数$\alpha$是$x=0$时$Y$的期望，即$E(Y\mid x=0)$. 依具体问题而定，这可能是有意义的量，也可能不是. 特别地，$x=0$未必是预测变量的合理取值. 但$\beta$是$E(Y\mid x)$作为$x$函数的变化率，即$x$改变一个单位时$E(Y\mid x)$的改变量. 所以这个参数涉及整个$x$值范围，包含$Y$与$x$之间一切线性关系的信息. （见\autoref{exercise:11.33}.）此外，$\beta=0$这一值尤其值得关注.

若$\beta=0$，则$E(Y\mid x)=\alpha+\beta x=\alpha$，且$Y\sim\mathrm n(\alpha,\sigma^2)$，不依赖$x$. 对一个经过周密思考、准备进行回归分析的实验，我们并不预期会如此，但若事实确实如此，也希望知道.

检验$\beta=0$与ANOVA中检验所有处理相等十分相似. ANOVA的原假设表示，处理与响应在任何方面都无关系；线性回归的原假设$\beta=0$则表示，处理（$x$）与响应之间没有线性关系.

\sourcepageanchor{480}
要检验
\begin{equation}\label{eq:11.3.34}
H_0:\quad\beta=0\quad\text{对}\quad H_1:\quad\beta\ne0,
\end{equation}
利用\eqref{eq:11.3.33}，当
$$\left|\frac{\hat\beta-0}{S/\sqrt{S_{xx}}}\right|>t_{n-2,\alpha/2}$$
时在水平$\alpha$拒绝$H_0$；等价地，当
\begin{equation}\label{eq:11.3.35}\index{SI0909@t检验}
\frac{\hat\beta^2}{S^2/S_{xx}}>F_{1,n-2,\alpha}
\end{equation}
时拒绝. 回顾$\hat\beta$的公式以及$\mathrm{RSS}=\sum\hat\varepsilon_i^2$，有
$$\frac{\hat\beta^2}{S^2/S_{xx}}
=\frac{S_{xy}^2/S_{xx}}{\mathrm{RSS}/(n-2)}
=\frac{\text{回归平方和}}{\text{残差平方和}/\text{自由度}}.$$
最后这个公式汇总在回归ANOVA表中，类似\autoref{sec:11.2}中的ANOVA表. 简单线性回归对应的表见\autoref{tab:11.3.2}，产生\eqref{eq:11.3.35}给出的检验. 注意，表中只涉及关于$\beta$的假设. 参数$\alpha$与估计$\hat\alpha$在这里起的作用，与\autoref{sec:11.2}的总均值相同：只是确定数据的总体水平，在平方和中已被“校正”掉.

\begin{table}[htbp]
\centering\caption{简单线性回归的ANOVA表}\label{tab:11.3.2}\index{SI0879@方差分析表}
\small
\begin{tabular}{lcp{.23\textwidth}p{.22\textwidth}c}
\toprule 变异来源&自由度&平方和&均方&$F$统计量\\\midrule
回归（斜率）&1&回归SS=$S_{xy}^2/S_{xx}$&MS(回归)=回归SS&$F=\dfrac{\mathrm{MS}(\text{回归})}{\mathrm{MS}(\text{残差})}$\\[1ex]
残差&$n-2$&$\mathrm{RSS}=\sum\hat\varepsilon_i^2$&MS(残差)=\newline$\mathrm{RSS}/(n-2)$&\\[1ex]
总计&$n-1$&$\mathrm{SST}=\sum(y_i-\bar y)^2$&&\\\bottomrule
\end{tabular}
\end{table}

\begin{example}[\autoref{ex:11.3.1}续]\label{ex:11.3.4}
葡萄产量数据的回归ANOVA如下.
\begin{center}
葡萄数据的ANOVA表

\begin{tabular}{lrrrr}
\toprule 变异来源&自由度&平方和&均方&$F$统计量\\\midrule
回归&1&6.66&6.66&50.23\\
残差&10&1.33&.133&\\
总计&11&7.99&&\\\bottomrule
\end{tabular}
\end{center}
这表明回归直线的斜率高度显著.
\end{example}

最后再指出与方差分析的一个相似之处. 从\autoref{tab:11.3.2}未必明显看出，但ANOVA的平方和分解在回归中也有对应形式：
\begin{equation}\label{eq:11.3.36}\index{SI0907@平方和}
\begin{aligned}
\text{总平方和}&=\text{回归平方和}+\text{残差平方和},\\
\sum_{i=1}^n(y_i-\bar y)^2&=\sum_{i=1}^n(\hat y_i-\bar y)^2+\sum_{i=1}^n(y_i-\hat y_i)^2,
\end{aligned}
\end{equation}
\sourcepageanchor{481}
其中$\hat y_i=\hat\alpha+\hat\beta x_i$. 注意这些平方和与ANOVA中平方和的相似性. 总平方和当然相同. RSS度量拟合直线与观测值的偏离，回归平方和类似ANOVA的处理平方和，度量预测值（“处理均值”）与总均值的偏离. 也与ANOVA一样，平方和恒等式成立是因为交叉项消失（\autoref{exercise:11.34}）. \eqref{eq:11.3.36}中的总平方和与残差平方和显然与\autoref{tab:11.3.2}相同，回归平方和看起来却不同. 不过，二者相等（\autoref{exercise:11.34}），即
$$\sum_{i=1}^n(\hat y_i-\bar y)^2=\frac{S_{xy}^2}{S_{xx}}.$$
表达式$S_{xy}^2/S_{xx}$更便于计算，并建立了与$t$检验的联系；而$\sum_{i=1}^n(\hat y_i-\bar y)^2$更容易解释.

用来量化拟合直线对数据描述程度的统计量是决定系数（coefficient of determination）.\index{SI0108@决定系数}\semanticindexanchor{108:481}\index{SI0883@决定系数}\semanticindexanchor{883:481}\index{SI0904@r${}^2$}\semanticindexanchor{904:481} 它定义为回归平方和与总平方和之比，通常记为$r^2$，可以写成几种形式：
$$r^2=\frac{\text{回归平方和}}{\text{总平方和}}
=\frac{\sum_{i=1}^n(\hat y_i-\bar y)^2}{\sum_{i=1}^n(y_i-\bar y)^2}
=\frac{S_{xy}^2}{S_{xx}S_{yy}}.$$
决定系数度量$y_1,\ldots,y_n$的总变异（由$S_{yy}$度量）中，由拟合直线解释的部分（由回归平方和度量）所占的比例. 由\eqref{eq:11.3.36}，$0\le r^2\le1$. 若$y_1,\ldots,y_n$恰好全部落在拟合直线上，则对所有$i$有$y_i=\hat y_i$，且$r^2=1$. 若$y_1,\ldots,y_n$并不接近拟合直线，残差平方和就大，$r^2$接近0. 决定系数也可以（或许更直接地）推导为$n$对$(y_1,x_1),\allowbreak\ldots,\allowbreak(y_n,x_n)$，或$n$对$(y_1,\hat y_1),\allowbreak\ldots,\allowbreak(y_n,\hat y_n)$的样本相关系数的平方.

\eqref{eq:11.3.33}可用于构造$\beta$的$100(1-\alpha)\%$置信区间：
\begin{equation}\label{eq:11.3.37}\index{SI0170@回归}\index{SI0887@置信区间}
\hat\beta-t_{n-2,\alpha/2}\frac S{\sqrt{S_{xx}}}<\beta
<\hat\beta+t_{n-2,\alpha/2}\frac S{\sqrt{S_{xx}}}.
\end{equation}
此外，$H_0:\beta=\beta_0$对$H_1:\beta\ne\beta_0$的水平$\alpha$检验，在
\begin{equation}\label{eq:11.3.38}\index{SI0909@t检验}\index{SI1023@回归}
\left|\frac{\hat\beta-\beta_0}{S/\sqrt{S_{xx}}}\right|>t_{n-2,\alpha/2}
\end{equation}
时拒绝$H_0$. 如前所述，通常检验$H_0:\beta=0$对$H_1:\beta\ne0$，判断预测变量与响应变量之间是否存在某种线性关系. 但上述检验更一般，因为可以指定任意$\beta_0$值. 固守“程式”的回归ANOVA只能检验$H_0:\beta=0$.

\subsubsection{在给定\texorpdfstring{$x=x_0$}{x=x₀}处的估计与预测}\label{sec:11.3.5}
预测变量的一个指定值，比如$x=x_0$，对应一个$Y$值总体. 事实上，根据条件正态模型，从该总体取得的随机观测满足$Y\sim\mathrm n(\alpha+\beta x_0,\sigma^2)$. 在观测回归数据$(x_1,y_1),\allowbreak\ldots,\allowbreak(x_n,y_n)$并估计参数$\alpha,\beta,\sigma^2$后，实验人员可能将设定$x=x_0$并取得新观测，记为$Y_0$. 可能希望估计这一观测所属总体的均值，甚至预测这一观测的值. 现在讨论这些推断.

\sourcepageanchor{482}
假设$(x_1,Y_1),\allowbreak\ldots,\allowbreak(x_n,Y_n)$满足条件正态回归模型，并基于这$n$个观测得到估计$\hat\alpha,\hat\beta,S^2$. 设$x_0$为预测变量的指定值. 先考虑估计与$x_0$对应的$Y$总体的均值，即$E(Y\mid x_0)=\alpha+\beta x_0$. 点估计量的显然选择是$\hat\alpha+\hat\beta x_0$. 它是无偏估计量，因为$E(\hat\alpha+\hat\beta x_0)=E\hat\alpha+(E\hat\beta)x_0=\alpha+\beta x_0$. 利用\autoref{thm:11.3.3}给出的矩，还可计算
\begin{align*}
\operatorname{Var}(\hat\alpha+\hat\beta x_0)
&=\operatorname{Var}\hat\alpha+(\operatorname{Var}\hat\beta)x_0^2+2x_0\operatorname{Cov}(\hat\alpha,\hat\beta)\\
&=\frac{\sigma^2}{nS_{xx}}\sum_{i=1}^nx_i^2+\frac{\sigma^2x_0^2}{S_{xx}}-\frac{2\sigma^2x_0\bar x}{S_{xx}}\\
&=\frac{\sigma^2}{S_{xx}}\left(\frac1n\sum_{i=1}^nx_i^2-\bar x^2+\bar x^2-2x_0\bar x+x_0^2\right)
\quad(\pm\bar x)\\
&=\frac{\sigma^2}{S_{xx}}\left(\frac1n\left[\sum_{i=1}^nx_i^2-\frac1n\left(\sum_{i=1}^nx_i\right)^2\right]+(x_0-\bar x)^2\right)
\quad\text{（重新组合各项）}\\
&=\sigma^2\left(\frac1n+\frac{(x_0-\bar x)^2}{S_{xx}}\right)
\quad\text{（因为 }\sum x_i^2-\tfrac1n(\sum x_i)^2=S_{xx}\text{）}.
\end{align*}
最后，由于$\hat\alpha,\hat\beta$都是$Y_1,\ldots,Y_n$的线性函数，$\hat\alpha+\hat\beta x_0$也是，因而服从正态分布，具体为
\begin{equation}\label{eq:11.3.39}
\hat\alpha+\hat\beta x_0\sim\mathrm n\left(\alpha+\beta x_0,\sigma^2\left(\frac1n+\frac{(x_0-\bar x)^2}{S_{xx}}\right)\right).
\end{equation}
由\autoref{thm:11.3.3}，$(\hat\alpha,\hat\beta)$与$S^2$独立，所以$S^2$也与$\hat\alpha+\hat\beta x_0$独立（\autoref{thm:4.6.12}），且
$$\frac{\hat\alpha+\hat\beta x_0-(\alpha+\beta x_0)}{S\sqrt{1/n+(x_0-\bar x)^2/S_{xx}}}\sim t_{n-2}.$$
反演这个枢轴量，得到$\alpha+\beta x_0$的$100(1-\alpha)\%$置信区间：
\begin{equation}\label{eq:11.3.40}
\begin{aligned}
\hat\alpha+\hat\beta x_0-t_{n-2,\alpha/2}S\sqrt{\frac1n+\frac{(x_0-\bar x)^2}{S_{xx}}}
&\le\alpha+\beta x_0\\
&\le\hat\alpha+\hat\beta x_0+t_{n-2,\alpha/2}S\sqrt{\frac1n+\frac{(x_0-\bar x)^2}{S_{xx}}}.
\end{aligned}
\end{equation}
$\alpha+\beta x_0$的置信区间长度通过$(x_0-\bar x)^2/S_{xx}$依赖$x_1,\ldots,x_n$的值. 显然，$x_0$接近$\bar x$时区间较短，在$x_0=\bar x$时最短. 因此，实验设计时应选择$x_1,\ldots,x_n$，使待估均值所在的$x_0$等于或接近$\bar x$. 在观测数据的中心附近估计更精确，是合乎情理的.

一种尚未讨论的推断是对还未观测到的随机变量$Y$进行预测（prediction），这在回归情境中很有意义. 例如，假设$x$度量大学申请者的高中表现. 大学招生办公室可能希望用$x$预测$Y$，即该学生大学一年后的平均绩点. 显然，学生尚未被录取，$Y$还没有观测到！学校拥有以往学生的数据$(x_1,y_1),\allowbreak\ldots,\allowbreak(x_n,y_n)$，记录高中表现和大学一年后的平均绩点，可用这些数据预测新生的平均绩点.

\sourcepageanchor{483}
\begin{definition}\label{def:11.3.5}\index{SI0823@预测区间}
基于已观测数据$\boldsymbol X$，未观测随机变量$Y$的$100(1-\alpha)\%$预测区间（prediction interval）是随机区间$[L(\boldsymbol X),U(\boldsymbol X)]$，对参数$\theta$的所有取值都满足
$$P_\theta\bigl(L(\boldsymbol X)\le Y\le U(\boldsymbol X)\bigr)\ge1-\alpha.$$
\end{definition}

注意预测区间与置信区间定义的相似性. 区别在于预测区间针对随机变量，而不是参数. 直观上，随机变量比参数（常数）更具有变动性，所以预期预测区间比同水平的置信区间更宽. 在线性回归这一特例中，确实如此.

假设将在$x=x_0$处取得的新观测$Y_0$服从$\mathrm n(\alpha+\beta x_0,\sigma^2)$，并与此前的数据$(x_1,Y_1),\allowbreak\ldots,\allowbreak(x_n,Y_n)$独立. 估计量$\hat\alpha,\hat\beta,S^2$由此前数据计算，因而$Y_0$与它们独立. 利用\eqref{eq:11.3.39}，得到$Y_0-(\hat\alpha+\hat\beta x_0)$服从正态分布，均值为
$$E\bigl(Y_0-(\hat\alpha+\hat\beta x_0)\bigr)=\alpha+\beta x_0-(\alpha+\beta x_0)=0,$$
方差为
\begin{align*}
\operatorname{Var}\bigl(Y_0-(\hat\alpha+\hat\beta x_0)\bigr)
&=\operatorname{Var}Y_0+\operatorname{Var}(\hat\alpha+\hat\beta x_0)\\
&=\sigma^2+\sigma^2\left(\frac1n+\frac{(x_0-\bar x)^2}{S_{xx}}\right).
\end{align*}
利用$S^2$与$Y_0-(\hat\alpha+\hat\beta x_0)$独立，有
$$T=\frac{Y_0-(\hat\alpha+\hat\beta x_0)}{S\sqrt{1+1/n+(x_0-\bar x)^2/S_{xx}}}\sim t_{n-2}.$$
按通常方式整理，得到$100(1-\alpha)\%$预测区间：
\begin{equation}\label{eq:11.3.41}\index{SI0902@预测区间}
\begin{aligned}
\hat\alpha+\hat\beta x_0-t_{n-2,\alpha/2}S\sqrt{1+\frac1n+\frac{(x_0-\bar x)^2}{S_{xx}}}
&<Y_0\\
&<\hat\alpha+\hat\beta x_0+t_{n-2,\alpha/2}S\sqrt{1+\frac1n+\frac{(x_0-\bar x)^2}{S_{xx}}}.
\end{aligned}
\end{equation}
由于区间端点只依赖已观测数据，\eqref{eq:11.3.41}定义了新观测$Y_0$的预测区间.

\subsubsection{同时估计与置信带}\label{sec:11.3.6}\index{SI0170@回归}
上一节考察了单个$x_0$值处的预测. 但有些情况下可能希望在许多$x_0$处预测. 例如，前述平均绩点预测问题中，招生办公室可能希望预测许多申请者的平均绩点，自然就需要在多个$x_0$处预测.

遇到的问题是现在已很熟悉的同时推断问题：如何控制同时推断的总体置信水平？上一节得到，与$x_0$对应的$Y$总体均值$E(Y\mid x_0)=\alpha+\beta x_0$的$1-\alpha$置信区间为
$$\begin{aligned}
\hat\alpha+\hat\beta x_0-t_{n-2,\alpha/2}S\sqrt{\frac1n+\frac{(x_0-\bar x)^2}{S_{xx}}}
&<\alpha+\beta x_0\\
&<\hat\alpha+\hat\beta x_0+t_{n-2,\alpha/2}S\sqrt{\frac1n+\frac{(x_0-\bar x)^2}{S_{xx}}}.
\end{aligned}$$

\sourcepageanchor{484}
现在假设希望在多个$x_0$值处对$Y$总体均值作推断，例如希望构造$E(Y\mid x_{0i})$的区间，$i=1,\ldots,m$. 已经知道，若像上面那样构造$m$个各为$1-\alpha$水平的区间，总体推断就不会处于$1-\alpha$水平.

一种简单且相当好的办法，是像\autoref{ex:11.2.9}那样使用Bonferroni不等式. 利用它可断言，下列不等式对$i=1,\ldots,m$同时成立的概率至少为$1-\alpha$：
\begin{equation}\label{eq:11.3.42}\index{SI0562@回归}
\begin{aligned}
\hat\alpha+\hat\beta x_{0i}-t_{n-2,\alpha/(2m)}S\sqrt{\frac1n+\frac{(x_{0i}-\bar x)^2}{S_{xx}}}
&<\alpha+\beta x_{0i}\\
&<\hat\alpha+\hat\beta x_{0i}+t_{n-2,\alpha/(2m)}S\sqrt{\frac1n+\frac{(x_{0i}-\bar x)^2}{S_{xx}}}.
\end{aligned}
\end{equation}
（见\autoref{exercise:11.39}.）

回归中的同时推断还可再向前推进一步. 关于总体回归直线的假设意味着$E(Y\mid x)=\alpha+\beta x$对所有$x$成立，所以应能在所有$x$处作推断. 因而，希望作出类似\eqref{eq:11.3.42}的陈述，但要求对所有$x$成立. 正如所预料的，Scheff\'e像在ANOVA中一样，给出了这个问题的解. 下面的定理总结简单线性回归情形的结果.

\begin{theorem}\label{thm:11.3.6}\index{SI0681@二次函数的最大值}
在条件正态回归模型\eqref{eq:11.3.22}下，下列不等式对所有$x$同时成立的概率至少为$1-\alpha$：
\begin{equation}\label{eq:11.3.43}
\begin{aligned}
\hat\alpha+\hat\beta x-M_\alpha S\sqrt{\frac1n+\frac{(x-\bar x)^2}{S_{xx}}}
&<\alpha+\beta x\\
&<\hat\alpha+\hat\beta x+M_\alpha S\sqrt{\frac1n+\frac{(x-\bar x)^2}{S_{xx}}},
\end{aligned}
\end{equation}
其中$M_\alpha=\sqrt{2F_{2,n-2,\alpha}}$.
\end{theorem}
\begin{proof}
重新整理各项可知，只要能找到常数$M_\alpha$满足
$$P\left(\frac{\bigl((\hat\alpha+\hat\beta x)-(\alpha+\beta x)\bigr)^2}
{S^2[1/n+(x-\bar x)^2/S_{xx}]}\le M_\alpha^2\text{，对所有 }x\right)=1-\alpha,$$
或等价地，
$$P\left(\max_x\frac{\bigl((\hat\alpha+\hat\beta x)-(\alpha+\beta x)\bigr)^2}
{S^2[1/n+(x-\bar x)^2/S_{xx}]}\le M_\alpha^2\right)=1-\alpha,$$
定理结论就成立. \autoref{exercise:11.32}给出的参数化使$\alpha,\beta$的估计量独立，能使上述最大化更容易. 写成
\begin{align*}
\hat\alpha+\hat\beta x&=\bar Y+\hat\beta(x-\bar x),\\
\alpha+\beta x&=\mu_{\bar Y}+\beta(x-\bar x),
\qquad(\mu_{\bar Y}=E\bar Y=\alpha+\beta\bar x),
\end{align*}
并为简便定义$t=x-\bar x$. 则
$$\frac{\bigl((\hat\alpha+\hat\beta x)-(\alpha+\beta x)\bigr)^2}{S^2[1/n+(x-\bar x)^2/S_{xx}]}
=\frac{\bigl((\bar Y-\mu_{\bar Y})+(\hat\beta-\beta)t\bigr)^2}{S^2[1/n+t^2/S_{xx}]},$$
\sourcepageanchor{485}
希望求$M_\alpha$使
$$P\left(\max_t\frac{\bigl((\bar Y-\mu_{\bar Y})+(\hat\beta-\beta)t\bigr)^2}{S^2[1/n+t^2/S_{xx}]}\le M_\alpha^2\right)=1-\alpha.$$
注意$S^2$只作为常数，不影响最大化. 应用直接由微积分得到的\autoref{exercise:11.40}的结果，得
\begin{equation}\label{eq:11.3.44}
\begin{aligned}
\max_t\frac{\bigl((\bar Y-\mu_{\bar Y})+(\hat\beta-\beta)t\bigr)^2}{S^2[1/n+t^2/S_{xx}]}
&=\frac{n(\bar Y-\mu_{\bar Y})^2+S_{xx}(\hat\beta-\beta)^2}{S^2}\\
&=\frac{\dfrac{(\bar Y-\mu_{\bar Y})^2}{\sigma^2/n}+\dfrac{(\hat\beta-\beta)^2}{\sigma^2/S_{xx}}}{S^2/\sigma^2}.
\quad\text{（乘以 }\sigma^2/\sigma^2\text{）}
\end{aligned}
\end{equation}
由\autoref{thm:11.3.3}和\autoref{exercise:11.32}，最后这个表达式是独立卡方随机变量之商，分母还除以其自由度. 分子是两个独立的$\chi^2_1$随机变量之和，所以分子服从$\chi^2_2$，商的分布为
$$\frac{\dfrac{(\bar Y-\mu_{\bar Y})^2}{\sigma^2/n}+\dfrac{(\hat\beta-\beta)^2}{\sigma^2/S_{xx}}}{S^2/\sigma^2}\sim2F_{2,n-2},$$
并且当$M_\alpha=\sqrt{2F_{2,n-2}}$时，\footnote{译注：原书这里的$F$下标漏了临界概率$\alpha$；应与定理陈述一致，为$\sqrt{2F_{2,n-2,\alpha}}$. 原式保留.}
$$P\left(\max_t\frac{\bigl((\bar Y-\mu_{\bar Y})+(\hat\beta-\beta)t\bigr)^2}{S^2[1/n+t^2/S_{xx}]}\le M_\alpha^2\right)=1-\alpha,$$
定理得证.
\end{proof}

由于\eqref{eq:11.3.43}对所有$x$都成立，它实际上给出了整条总体回归直线的置信带（confidence band）.\index{SI0886@置信带}\semanticindexanchor{886:485} 也就是说，置信区间覆盖一个单值参数，而置信带以一个带状区域覆盖整条直线. \autoref{fig:11.3.3}给出了Scheff\'e置信带的例子，同时画出两个Bonferroni区间和一个单独的$t$区间. 注意，虽然\autoref{fig:11.3.3}并非如此，但Bonferroni区间可能比Scheff\'e置信带更宽，尽管Bonferroni推断必然只涉及较少的区间. 只要
$$t_{n-2,\alpha/(2m)}>2F_{2,n-2,\alpha},$$
就会出现这种情形，其中$m$如\eqref{eq:11.3.42}定义.\footnote{译注：比较两种区间的实际临界宽度时，右端应为$\sqrt{2F_{2,n-2,\alpha}}$；原书此处漏了平方根. 保留原式.} 对足够大的$m$，这个不等式总会成立，所以即使只关心有限多个$x$值，也总会达到一个从Bonferroni改用Scheff\'e更合算的点. 这种看起来似乎“不付代价便有所得”的现象，是因为Bonferroni不等式是一个通用界，而Scheff\'e置信带是当前问题的精确解. （Bonferroni区间的实际覆盖概率高于$1-\alpha$.）Scheff\'e置信带有许多变体，有的形状不同，有的只保证某个特定$x$区间的覆盖. 这些其他置信带的讨论见\autoref{sec:11.5}.

理论上，对\autoref{thm:11.3.6}的证明作适当修改，也可得到同时预测区间. （事实上，\autoref{exercise:11.40}中函数的最大化几乎立即给出结果.）但问题是，所得统计量的分布并不特别理想.

\sourcepageanchor{486}
\begin{figure}[htbp]
\centering\includegraphics[width=.8\textwidth]{assets/chapter11/fig-11.3.3.pdf}
\caption{\autoref{tab:11.3.1}数据的Scheff\'e置信带、$t$区间（$x=3.5$处）及Bonferroni区间（$x=1$和$x=3$处）}\label{fig:11.3.3}
\end{figure}
最后指出，用Scheff\'e置信带等程序对观测$x$范围之外的$x$值作推断存在一个问题. 这些程序基于这样的假设：我们\emph{知道}总体回归函数对所有$x$都是线性的. 虽然假设回归函数在观测$x$的范围内线性可能合理，但向该范围之外的$x$作外推（extrapolation）通常不明智.\index{SI0891@外推}\semanticindexanchor{891:486} （观测范围之外没有数据，无法检查回归是否变成非线性.）这一提醒同样适用于\autoref{sec:11.3.5}的方法.



\subsection{习题}\label{sec:11.4}
\begin{exercise}\label{exercise:11.1}\index{SI0054@数据变换}
ANOVA中的方差稳定变换（variance-stabilizing transformation）以下述近似方式稳定方差. 设$Y$的均值为$\theta$，方差为$v(\theta)$.
\sourcepageanchor{487}
\begin{enumerate}[label=(\alph*)]
\item\label{item:11.1.a} 用类似\autoref{sec:10.1.3}的论证，说明单项Taylor级数近似给出$g(y)$的方差为
$$\operatorname{Var}(g(Y))=\left[\frac d{d\theta}g(\theta)\right]^2v(\theta).$$
\item\label{item:11.1.b} 证明$g^*(Y)$的近似方差与$\theta$无关，其中$g^*(y)=\int[1/\sqrt{v(y)}]\,dy$.
\end{enumerate}
\end{exercise}

\begin{exercise}\label{exercise:11.2}
验证下列变换在\autoref{exercise:11.1}的意义下近似稳定方差.
\begin{enumerate}[label=(\alph*)]
\item\label{item:11.2.a} $Y\sim\operatorname{Poisson}$，$g^*(y)=\sqrt y$.
\item\label{item:11.2.b} $Y\sim\operatorname{binomial}(n,p)$，$g^*(y)=\sin^{-1}(\sqrt{y/n})$.
\item\label{item:11.2.c} $Y$的方差为$v(\theta)=K\theta^2$，其中$K$为常数，$g^*(y)=\log y$.
\end{enumerate}
（方差稳定变换的存在条件至少可追溯至\hyperref[bib:source90]{Curtiss，1943}；\hyperref[bib:source7]{Bar-Lev和Enis，1988}、1990，给出了改进.）
\end{exercise}

\begin{exercise}\label{exercise:11.3}
Box--Cox幂变换族（\hyperref[bib:source44]{Box和Cox，1964）}定义为
$$g_\lambda^*(y)=\begin{cases}(y^\lambda-1)/\lambda,&\lambda\ne0,\\\log y,&\lambda=0,\end{cases}$$
其中$\lambda$为自由参数.
\begin{enumerate}[label=(\alph*)]
\item\label{item:11.3.a} 证明，对每个$y$，$g_\lambda^*(y)$关于$\lambda$连续. 特别地，证明
$$\lim_{\lambda\to0}(y^\lambda-1)/\lambda=\log y.$$
\item\label{item:11.3.b} 求函数$v(\theta)$，即由$g_\lambda^*(y)$稳定的$Y$的近似方差. （注意$v(\theta)$很可能还依赖$\lambda$.）
\end{enumerate}
变换后数据的分析，特别是Box--Cox幂变换，在统计文献中曾引起一些争论. 见\hyperref[bib:source31]{Bickel和Doksum（1981）}、\hyperref[bib:source45]{Box和Cox（1982）}，以及\hyperref[bib:source168]{Hinkley和Runger（1984）}.
\end{exercise}

\begin{exercise}\label{exercise:11.4}
一种极著名且有用的方差稳定变换是Fisher的$z$变换，在\autoref{exercise:10.17}已遇到. 这里再考察一些细节. 假设$(X,Y)$服从二元正态分布，相关系数为$\rho$，样本相关系数为$r$.
\begin{enumerate}[label=(\alph*)]
\item\label{item:11.4.a} 从\hyperref[item:10.17.d]{习题10.17(d)}出发，用Delta方法证明
$$\frac12\left[\log\left(\frac{1+r}{1-r}\right)-\log\left(\frac{1+\rho}{1-\rho}\right)\right]$$
近似正态，均值为0，方差为$1/n$.
\item\label{item:11.4.b} Fisher实际上使用了稍精确的展开（\hyperref[bib:source309]{Stuart和Ord，1987}，第16.33节），证明\hyperref[item:11.4.a]{(a)}中的量近似正态，且
$$\text{均值}=\frac\rho{2(n-1)},\qquad\text{方差}=\frac1{n-1}+\frac{4-\rho^2}{2(n-1)^2}.$$
证明，对较小的$\rho$和中等的$n$，可将此均值、方差近似为0和$1/(n-3)$，这就是Fisher $z$变换最常用的形式.
\end{enumerate}
\end{exercise}

\begin{exercise}\label{exercise:11.5}
假设按照\eqref{eq:11.2.2}的过参数化单因素ANOVA模型观测随机变量$Y_{ij}$. 给出两组不同的参数，使$Y_{ij}$具有完全相同的分布，从而说明：若不对参数施加某种限制，该模型不可识别.
\end{exercise}

\begin{exercise}\label{exercise:11.6}
在单因素ANOVA假设下：
\begin{enumerate}[label=(\alph*)]
\item\label{item:11.6.a}\index{SI0049@充分统计量}\index{SI1001@方差分析} 证明统计量组$(\bar Y_{1\cdot},\bar Y_{2\cdot},\ldots,\bar Y_{k\cdot},S_p^2)$对$(\theta_1,\theta_2,\ldots,\theta_k,\sigma^2)$充分.
\item\label{item:11.6.b} 证明$S_p^2=\frac1{N-k}\sum_{i=1}^k(n_i-1)S_i^2$与每个$\bar Y_{i\cdot}$独立，$i=1,\ldots,k$. （见\autoref{lem:5.3.3}.）
\item\label{item:11.6.c} 若$\sigma^2$已知，说明ANOVA数据如何等价于\hyperref[sec:11.5.6]{杂录11.5.6}中的典范形式.
\end{enumerate}
\end{exercise}

\sourcepageanchor{488}
\begin{exercise}\label{exercise:11.7}
证明
$$\frac1{\sigma^2}\sum_{i=1}^kn_i\bigl((\bar Y_{i\cdot}-\bar{\bar Y})-(\theta_i-\bar\theta)\bigr)^2\sim\chi^2_{k-1},$$
从而完成\autoref{thm:11.2.8}的证明.

（提示：定义$\bar U_i=\bar Y_{i\cdot}-\theta_i$，$i=1,\ldots,k$. 证明$\bar U_i$独立服从$\mathrm n(0,\sigma^2/n_i)$. 再调整\autoref{lem:5.3.2}的归纳论证，证明$\sum n_i(\bar U_i-\bar{\bar U})^2/\sigma^2\sim\chi^2_{k-1}$，其中$\bar{\bar U}=\sum n_i\bar U_i/\sum n_i$.）
\end{exercise}

\begin{exercise}\label{exercise:11.8}
证明：在单因素ANOVA假设下，对任意常数组$\boldsymbol a=(a_1,\ldots,a_k)$，量$\sum a_i\bar Y_{i\cdot}$服从均值为$\sum a_i\theta_i$、方差为$\sigma^2\sum a_i^2/n_i$的正态分布. （见\autoref{cor:4.6.10}.）
\end{exercise}

\begin{exercise}\label{exercise:11.9}
利用类似推导\eqref{eq:11.2.7}的$t$检验的论证，说明如何对下列假设构造$t$检验：
\begin{enumerate}[label=(\alph*)]
\item\label{item:11.9.a} $H_0:\sum a_i\theta_i=\delta$对$H_1:\sum a_i\theta_i\ne\delta$.
\item\label{item:11.9.b} $H_0:\sum a_i\theta_i\le\delta$对$H_1:\sum a_i\theta_i>\delta$，其中$\delta$为指定常数.
\end{enumerate}
\end{exercise}

\begin{exercise}\label{exercise:11.10}
假设有一个五处理的单因素ANOVA. 处理均值记为$\theta_1,\ldots,\theta_5$，其中$\theta_1$为对照，$\theta_2,\ldots,\theta_5$为备选新处理，并假设每个处理取相同数量的观测. 考虑由下列向量定义的四个对比$\sum a_i\theta_i$：
$$\begin{aligned}
\boldsymbol a_1&=(1,-\tfrac14,-\tfrac14,-\tfrac14,-\tfrac14),\\
\boldsymbol a_2&=(0,1,-\tfrac13,-\tfrac13,-\tfrac13),\\
\boldsymbol a_3&=(0,0,1,-\tfrac12,-\tfrac12),\\
\boldsymbol a_4&=(0,0,0,1,-1).
\end{aligned}$$
\begin{enumerate}[label=(\alph*)]
\item\label{item:11.10.a} 说明使用这些对比进行四个$t$检验，其结果可以对$\theta_1,\ldots,\theta_5$的排序得出结论. 可能得到哪些结论？
\item\label{item:11.10.b} 证明由\hyperref[item:11.10.a]{(a)}中的四个$\boldsymbol a_i$形成的任意两个对比$\sum a_i\bar Y_{i\cdot}$互不相关. （回顾它们称为正交对比.）
\item\label{item:11.10.c} 对\autoref{ex:11.2.3}的肥料实验，预先安排了以下对比：
$$\begin{aligned}
\boldsymbol a_1&=(-1,1,0,0,0),\\
\boldsymbol a_2&=(0,-1,\tfrac12,\tfrac12,0),\\
\boldsymbol a_3&=(0,0,1,-1,0),\\
\boldsymbol a_4&=(0,-1,0,0,1,).
\end{aligned}$$
证明这些对比并不正交. 在肥料实验的背景下解释这些对比，说明它们是一组合理的对比.
\end{enumerate}
\end{exercise}

\begin{exercise}\label{exercise:11.11}
对任意常数组$\boldsymbol a=(a_1,\ldots,a_k)$和$\boldsymbol b=(b_1,\ldots,b_k)$，证明在单因素ANOVA假设下
$$\operatorname{Cov}\left(\sum a_i\bar Y_{i\cdot},\sum b_i\bar Y_{i\cdot}\right)
=\sigma^2\sum\frac{a_ib_i}{n_i}.$$
因此，单因素ANOVA中，若$\sum a_ib_i/n_i=0$，则对比互不相关（正交）.
\end{exercise}

\begin{exercise}\label{exercise:11.12}\index{SI0037@与t检验的关系}\index{SI0052@与F检验的关系}
假设单因素ANOVA中每个处理的观测数相同，即$n_i=n$，$i=1,\ldots,k$. 此时$F$检验可看成平均的$t$检验.
\begin{enumerate}[label=(\alph*)]
\item\label{item:11.12.a} 证明，$H_0:\theta_i=\theta_{i'}$对$H_1:\theta_i\ne\theta_{i'}$的$t$检验可以基于统计量
$$t_{ii'}^2=\frac{(\bar Y_{i\cdot}-\bar Y_{i'\cdot})^2}{S_p^2(2/n)}.$$
\item\label{item:11.12.b} 证明
$$\frac1{k(k-1)}\sum_{i,i'}t_{ii'}^2=F,$$
其中$F$为通常的ANOVA $F$统计量. （提示：见\hyperref[item:5.8.a]{习题5.8(a)}.）（由George McCabe提供，他从John Tukey处获知此结果.）
\end{enumerate}
\end{exercise}

\sourcepageanchor{489}
\begin{exercise}\label{exercise:11.13}\index{SI0041@似然比检验}\index{SI0631@方差分析}
在单因素ANOVA假设下，证明$H_0:\theta_1=\theta_2=\cdots=\theta_k$的似然比检验就是\eqref{eq:11.2.14}的$F$检验.
\end{exercise}

\begin{exercise}\label{exercise:11.14}\index{SI0047@Scheffé方法}\index{SI0681@二次函数的最大值}\index{SI0956@方差分析}
Scheff\'e同时区间方法实际上适用于所有线性组合，不仅适用于对比. 证明：在单因素ANOVA假设下，若$M=\sqrt{kF_{k,N-k,\alpha}}$（注意分子自由度的改变），则下列不等式对所有$\boldsymbol a=(a_1,\ldots,a_k)$同时成立的概率为$1-\alpha$：
$$\begin{aligned}
\sum_{i=1}^ka_i\bar Y_{i\cdot}-M\sqrt{S_p^2\sum_{i=1}^k\frac{a_i^2}{n_i}}
&\le\sum_{i=1}^ka_i\theta_i\\
&\le\sum_{i=1}^ka_i\bar Y_{i\cdot}+M\sqrt{S_p^2\sum_{i=1}^k\frac{a_i^2}{n_i}}.
\end{aligned}$$
（最容易的做法可能是，仿照\autoref{lem:11.2.7}先证明：若$v_1,\ldots,v_k$为常数，$c_1,\ldots,c_k$为正常数，则
$$\max_{\boldsymbol a}\left\{\frac{(\sum_{i=1}^ka_iv_i)^2}{\sum_{i=1}^ka_i^2/c_i}\right\}=\sum_{i=1}^kc_iv_i^2.$$
然后调整\autoref{thm:11.2.10}的证明，建立所需结果.）
\end{exercise}

\begin{exercise}\label{exercise:11.15}\leavevmode
\begin{enumerate}[label=(\alph*)]
\item\label{item:11.15.a} 证明，对$t$分布和$F$分布，以及任意$\nu,\alpha,k$，
$$t_{\nu,\alpha/2}\le\sqrt{(k-1)F_{k-1,\nu,\alpha}}.$$
（回顾$t$与$F$的关系. 固定$\nu$时，$kF_{k,\nu}$的分布随$k$随机递增，该不等式是其推论，但实际上是一个较弱的陈述. 见\autoref{exercise:5.19}.）
\item\label{item:11.15.b} 解释上述不等式如何说明，同时Scheff\'e区间总比单个对比区间宽.
\item\label{item:11.15.c} 证明，由上述不等式还可得到Scheff\'e检验的功效低于$t$检验.
\end{enumerate}
\end{exercise}

\begin{exercise}\label{exercise:11.16}
\autoref{thm:11.2.5}表明ANOVA原假设等价于所有对比为零. 也可将ANOVA原假设写成另一组假设的交.
\begin{enumerate}[label=(\alph*)]
\item\label{item:11.16.a} 证明假设
$$H_0:\theta_1=\theta_2=\cdots=\theta_k\quad\text{对}\quad H_1:\theta_i\ne\theta_j\text{，对某些 }i,j$$
与假设
$$H_0:\theta_i-\theta_j=0\text{，对所有 }i,j\quad\text{对}\quad H_1:\theta_i-\theta_j\ne0\text{，对某些 }i,j$$
等价.
\item\label{item:11.16.b} 把ANOVA检验的$H_0,H_1$表示为集合
$$\Theta_{ij}=\{\boldsymbol\theta=(\theta_1,\ldots,\theta_k):\theta_i-\theta_j=0\}$$
的并与交. 描述如何利用这些表达式，构造ANOVA原假设的另一种不同的并交检验. （见\hyperref[sec:11.5.2]{杂录11.5.2}.）
\end{enumerate}
\end{exercise}

\begin{exercise}\label{exercise:11.17}\index{SI0717@受保护的最小显著差}
一种称为受保护LSD（Protected Least Significant Difference，受保护的最小显著差）的多重比较程序如下进行：若ANOVA $F$检验在水平$\alpha$拒绝$H_0$，则对每对均值$\theta_i,\theta_{i'}$，在
$$\frac{|\bar Y_{i\cdot}-\bar Y_{i'\cdot}|}{\sqrt{S_p^2(1/n_i+1/n_{i'})}}>t_{\alpha/2,N-k}$$
时判定均值不同.\footnote{译注：原题$t$分位数的下标次序与本书通常记号相反，按通常记号应为$t_{N-k,\alpha/2}$. 原式保留.} 注意，每个$t$检验都在与ANOVA $F$检验相同的$\alpha$水平进行. 这里使用全实验$\alpha$水平（experimentwise $\alpha$），其中
$$\text{全实验 }\alpha=P\left(\text{至少作出一个错误的差异断言}\mid\text{所有均值相等}\right).$$
\sourcepageanchor{490}
\begin{enumerate}[label=(\alph*)]
\item\label{item:11.17.a} 证明，无论实验中有多少均值，受保护LSD所作的同时推断都处于水平$\alpha$.
\item\label{item:11.17.b}\index{SI0715@逐次比较错误率}\index{SI0716@整体试验错误率} 普通（或不受保护的）LSD不论ANOVA $F$检验结果如何，都直接进行水平$\alpha$的各个$t$检验. 证明，普通LSD的全实验错误率可能大于$\alpha$. （不受保护LSD确实维持了每次比较的错误率$\alpha$.）
\item\label{item:11.17.c} 对\autoref{ex:11.2.1}的鱼类毒素数据进行LSD程序. 结论是什么？
\end{enumerate}
\end{exercise}

\needspace{6\baselineskip}
\begin{exercise}\label{exercise:11.18}\index{SI0033@数据窥探}
为说明“数据窥探”，即检验由数据启发的假设，通常不是良好做法：
\begin{enumerate}[label=(\alph*)]
\item\label{item:11.18.a} 证明，对任意随机变量$Y$和常数$a,b$，若$a>b$且$P(Y>b)<1$，则$P(Y>a\mid Y>b)>P(Y>a)$.\footnote{译注：原题还需$P(Y>b)>0$以定义条件概率，且严格不等号需$P(Y>a)>0$；否则可能取等. 原条件与结论保留.}
\item\label{item:11.18.b} 令$Y$为检验统计量，$a$为临界点，把\hyperref[item:11.18.a]{(a)}的不等式用于由数据启发的假设检验的大小.
\end{enumerate}
\end{exercise}

\begin{exercise}\label{exercise:11.19}
设$X_i\sim\operatorname{gamma}(\lambda_i,1)$相互独立，$i=1,\ldots,n$. 定义$Y_i=X_{i+1}/(\sum_{j=1}^iX_j)$，$i=1,\ldots,n-1$，以及$Y_n=\sum_{i=1}^nX_i$.
\begin{enumerate}[label=(\alph*)]
\item\label{item:11.19.a} 求$Y_i$，$i=1,\ldots,n$，的联合分布和边缘分布.
\item\label{item:11.19.b} 将结果与ANOVA中常用的分布联系起来.
\end{enumerate}
\end{exercise}

\begin{exercise}\label{exercise:11.20}
证明：若单因素ANOVA原假设成立，则
\begin{enumerate}[label=(\alph*)]
\item\label{item:11.20.a} $\sum n_i(\bar Y_{i\cdot}-\bar{\bar Y})^2/(k-1)$是$\sigma^2$的无偏估计量.
\item\label{item:11.20.b} 说明如何用\autoref{ex:5.3.5}的方法推导ANOVA $F$检验.
\end{enumerate}
\end{exercise}

\begin{exercise}\label{exercise:11.21}\leavevmode
\begin{enumerate}[label=(\alph*)]
\item\label{item:11.21.a} 计算下列数据的完整ANOVA表，说明ANOVA中的平方和分解. 为确定饲料质量，对雄性断奶大鼠喂食不同蛋白质水平的饲料. 将15只大鼠分别随机分配至三种饲料之一，记录其体重增加量，单位为克.
\begin{center}
\begin{tabular}{rrr}
\toprule\multicolumn{3}{c}{饲料蛋白质水平}\\
低&中&高\\\midrule
3.89&8.54&20.39\\3.87&9.32&24.22\\3.26&8.76&30.91\\2.70&9.30&22.78\\3.82&10.45&26.33\\\bottomrule
\end{tabular}
\end{center}
\item\label{item:11.21.b} 完成\autoref{thm:11.2.11}的证明，以解析方式验证ANOVA平方和分解.
\item\label{item:11.21.c} 用\hyperref[item:11.21.a]{(a)}的数据说明\hyperref[item:11.12.b]{习题11.12(b)}中$t$统计量与$F$统计量的关系.
\end{enumerate}
\end{exercise}

\begin{exercise}\label{exercise:11.22}\index{SI0035@均方的期望}
计算单因素ANOVA表中MSB与MSW的期望. （这些期望正式称为期望均方，expected mean squares，可用于帮助确定复杂ANOVA中的$F$检验. 有算法可计算期望均方，算法细节例如可见\hyperref[bib:source190]{Kirk，1982}.）
\end{exercise}

\begin{exercise}\label{exercise:11.23}
对\hyperref[sec:11.5.3]{杂录11.5.3}的模型：
\begin{enumerate}[label=(\alph*)]
\item\label{item:11.23.a} 证明$Y_{ij}$的均值、方差为$EY_{ij}=\mu+\tau_i$和$\operatorname{Var}Y_{ij}=\sigma_B^2+\sigma^2$.
\item\label{item:11.23.b} 若$\sum a_i=0$，证明$\sum a_i\bar Y_{i\cdot}$的无条件方差为
$$\operatorname{Var}\left(\sum a_i\bar Y_{i\cdot}\right)=\frac1r(\sigma^2+\sigma_B^2)(1-\rho)\sum a_i^2,$$
其中$\rho$为组内相关系数.
\end{enumerate}
\end{exercise}

\begin{exercise}\label{exercise:11.24}
\hyperref[sec:11.5.6]{杂录11.5.6}的Stein估计量形式，可以通过\hyperref[bib:source111]{Efron和Morris（1972）}给出的经验贝叶斯论证得到一定解释，该论证在数据分析中相当有用. \hyperref[bib:source301]{Stein（1956）}可能已经知道这样的论证，但未提及. 设$X_i\sim\mathrm n(\theta_i,1)$，$i=1,\ldots,p$，且$\theta_i$独立同分布于$\mathrm n(0,\tau^2)$.
\begin{enumerate}[label=(\alph*)]
\item\label{item:11.24.a} 证明$X_i$的边缘分布独立同为$\mathrm n(0,\tau^2+1)$，从而$\sum X_i^2/(\tau^2+1)\sim\chi^2_p$.
\sourcepageanchor{491}
\item\label{item:11.24.b} 利用边缘分布，证明若$p\ge3$，则
$$E\left(1-\frac{p-2}{\sum_{j=1}^pX_j^2}\right)=\frac{\tau^2}{\tau^2+1}.$$
因此，\hyperref[sec:11.5.6]{杂录11.5.6}的Stein估计量是贝叶斯估计量$\delta_\pi^i(\boldsymbol X)=[\tau^2/(\tau^2+1)]X_i$的经验贝叶斯版本.
\item\label{item:11.24.c} 证明：若$Y\sim\chi^2_p$且$p<3$，则$E(1/Y)=\infty$，从而说明$p<3$时该论证失效.
\end{enumerate}
\end{exercise}

\begin{exercise}\label{exercise:11.25}\index{SI0613@正规方程}\index{SI0726@正规方程}
\autoref{sec:11.3.1}通过两阶段最小化求出了$\alpha,\beta$的最小二乘估计. 这个最小化也可用偏导数进行.
\begin{enumerate}[label=(\alph*)]
\item\label{item:11.25.a} 计算$\partial\mathrm{RSS}/\partial c$和$\partial\mathrm{RSS}/\partial d$，并令其为零. 证明所得两个方程可写成
$$nc+\left(\sum_{i=1}^nx_i\right)d=\sum_{i=1}^ny_i,
\qquad
\left(\sum_{i=1}^nx_i\right)c+\left(\sum_{i=1}^nx_i^2\right)d=\sum_{i=1}^nx_iy_i.$$
（这些方程称为该最小化问题的正规方程，normal equations.）
\item\label{item:11.25.b} 证明$c=a$和$d=b$是正规方程的解.
\item\label{item:11.25.c} 检查二阶偏导数条件，验证$c=a,d=b$这一点确实使RSS最小.
\end{enumerate}
\end{exercise}

\begin{exercise}\label{exercise:11.26}
假设$n$为偶数. 预测变量的值$x_1,\ldots,x_n$都必须在区间$[e,f]$中选取. 证明，使$S_{xx}$最大的选择，是令一半$x_i$等于$e$，另一半等于$f$. （这是\autoref{sec:11.3.2}中提到的使$\operatorname{Var}b$最小的选择.）
\end{exercise}

\begin{exercise}\label{exercise:11.27}
观测$(x_i,Y_i)$，$i=1,\ldots,n$，服从模型$Y_i=\alpha+\beta x_i+\varepsilon_i$，其中$E\varepsilon_i=0$、$\operatorname{Var}\varepsilon_i=\sigma^2$，且$i\ne j$时$\operatorname{Cov}(\varepsilon_i,\varepsilon_j)=0$. 求$\alpha$的最优线性无偏估计量.
\end{exercise}

\begin{exercise}\label{exercise:11.28}
证明，在简单线性回归的条件正态模型中，$\sigma^2$的MLE为
$$\hat\sigma^2=\frac1n\sum_{i=1}^n(y_i-\hat\alpha-\hat\beta x_i)^2.$$
\end{exercise}

\begin{exercise}\label{exercise:11.29}
考虑\autoref{sec:11.3.4}中由$\hat\varepsilon_i=Y_i-\hat\alpha-\hat\beta x_i$定义的残差$\hat\varepsilon_1,\ldots,\hat\varepsilon_n$.
\begin{enumerate}[label=(\alph*)]
\item\label{item:11.29.a} 证明$E\hat\varepsilon_i=0$.
\item\label{item:11.29.b} 验证
$$\begin{aligned}
\operatorname{Var}\hat\varepsilon_i
={}&\operatorname{Var}Y_i+\operatorname{Var}\hat\alpha+x_i^2\operatorname{Var}\hat\beta
-2\operatorname{Cov}(Y_i,\hat\alpha)\\
&-2x_i\operatorname{Cov}(Y_i,\hat\beta)+2x_i\operatorname{Cov}(\hat\alpha,\hat\beta).
\end{aligned}$$
\item\label{item:11.29.c} 用\autoref{lem:11.3.2}证明
$$\operatorname{Cov}(Y_i,\hat\alpha)=\sigma^2\left(\frac1n+\frac{(x_i-\bar x)\bar x}{S_{xx}}\right)
\quad\text{且}\quad\operatorname{Cov}(Y_i,\hat\beta)=\sigma^2\frac{x_i-\bar x}{S_{xx}},$$
并用它们验证\eqref{eq:11.3.28}.\footnote{译注：第一个协方差式中的加号应为减号，由$\hat\alpha=\sum_i[1/n-(x_i-\bar x)\bar x/S_{xx}]Y_i$可直接核对. 原题照录.}
\end{enumerate}
\end{exercise}

\begin{exercise}\label{exercise:11.30}
补全\autoref{thm:11.3.3}证明中略去的关于$\hat\alpha$分布的细节.
\begin{enumerate}[label=(\alph*)]
\item\label{item:11.30.a} 证明估计量$\hat\alpha=\bar y-\hat\beta\bar x$可表示为$\hat\alpha=\sum_{i=1}^nc_iY_i$，其中
$$c_i=\frac1n-\frac{(x_i-\bar x)\bar x}{S_{xx}}.$$
\item\label{item:11.30.b} 验证
$$E\hat\alpha=\alpha\quad\text{且}\quad\operatorname{Var}\hat\alpha=\sigma^2\left[\frac1{nS_{xx}}\sum_{i=1}^nx_i^2\right].$$
\item\label{item:11.30.c} 验证
$$\operatorname{Cov}(\hat\alpha,\hat\beta)=-\frac{\sigma^2\bar x}{S_{xx}}.$$
\end{enumerate}
\end{exercise}

\begin{exercise}\label{exercise:11.31}
验证\autoref{thm:11.3.3}中的断言：$\hat\varepsilon_i$与$\hat\alpha,\hat\beta$互不相关. （证明$\hat\varepsilon_i=\sum e_jY_j$，其中$e_j$由\eqref{eq:11.3.30}给出. 再利用$\hat\alpha=\sum c_jY_j$和$\hat\beta=\sum d_jY_j$，验证$\sum e_jc_j=\sum e_jd_j=0$，并应用\autoref{lem:11.3.2}.）
\end{exercise}

\sourcepageanchor{492}
\begin{exercise}\label{exercise:11.32}
按照模型
$$Y_i=\alpha+\beta x_i+\varepsilon_i$$
取得观测$(x_i,Y_i)$，$i=1,\ldots,n$，其中$x_1,\ldots,x_n$为固定常数，$\varepsilon_1,\ldots,\varepsilon_n$独立同分布于$\mathrm n(0,\sigma^2)$. 将模型重新参数化为
$$Y_i=\alpha'+\beta'(x_i-\bar x)+\varepsilon_i.$$
设$\hat\alpha,\hat\beta$分别为$\alpha,\beta$的MLE，$\hat\alpha',\hat\beta'$分别为$\alpha',\beta'$的MLE.
\begin{enumerate}[label=(\alph*)]
\item\label{item:11.32.a} 证明$\hat\beta'=\hat\beta$.
\item\label{item:11.32.b} 证明$\hat\alpha'\ne\hat\alpha$. 事实上，证明$\hat\alpha'=\bar Y$. 求$\hat\alpha'$的分布.\footnote{译注：原题不等号应理解为两个估计量一般不同；若$\bar x=0$，二者相同，个别样本也可能取相同值. 原文保留.}
\item\label{item:11.32.c} 证明$\hat\alpha',\hat\beta'$互不相关，从而在正态性下独立.
\end{enumerate}
\end{exercise}

\begin{exercise}\label{exercise:11.33}
从参数为$(\mu_X,\mu_Y,\sigma_X^2,\sigma_Y^2,\rho)$的二元正态总体取得观测$(X_i,Y_i)$，$i=1,\ldots,n$，拟合模型$Y_i=\alpha+\beta x_i+\varepsilon_i$.
\begin{enumerate}[label=(\alph*)]
\item\label{item:11.33.a} 说明假设$H_0:\beta=0$成立当且仅当假设$H_0:\rho=0$成立. （见\eqref{eq:11.3.25}.）
\item\label{item:11.33.b} 用代数方法证明
$$\frac{\hat\beta}{S/\sqrt{S_{xx}}}=\sqrt{n-2}\frac r{\sqrt{1-r^2}},$$
其中$r$为样本相关系数，也是$\rho$的MLE.
\item\label{item:11.33.c} 说明如何只给定$r^2$和$n$，使用自由度为$n-2$的Student $t$分布检验$H_0:\rho=0$（见\eqref{eq:11.3.33}）. （Fisher利用方差稳定变换推导了$\rho$的近似置信区间，见\autoref{exercise:11.4}.）
\end{enumerate}
\end{exercise}

\begin{exercise}\label{exercise:11.34}\leavevmode
\begin{enumerate}[label=(\alph*)]
\item\label{item:11.34.a} 计算下列数据的回归ANOVA表，说明简单线性回归的平方和分解. 父母常希望预测孩子成年后的身高. 以下数据来自一项可能受到Galton分析启发的研究，是其中的一部分.
\begin{center}
\begin{tabular}{lrrrrrrrr}
\toprule 两岁身高（英寸）（$x$）&39&30&32&34&35&36&36&30\\\midrule
成年身高（英寸）（$y$）&71&63&63&67&68&68&70&64\\\bottomrule
\end{tabular}
\end{center}
\item\label{item:11.34.b} 验证\eqref{eq:11.3.36}，从解析上建立简单线性回归的平方和分解.
\item\label{item:11.34.c} 证明回归平方和的两个表达式确实相等，即
$$\sum_{i=1}^n(\hat y_i-\bar y)^2=\frac{S_{xy}^2}{S_{xx}}.$$
\item\label{item:11.34.d} 证明决定系数
$$r^2=\frac{\sum_{i=1}^n(\hat y_i-\bar y)^2}{\sum_{i=1}^n(y_i-\bar y)^2}$$
既可以推导为$n$对$(y_1,x_1),\allowbreak\ldots,\allowbreak(y_n,x_n)$的样本相关系数平方，也可以推导为$n$对$(y_1,\hat y_1),\allowbreak\ldots,\allowbreak(y_n,\hat y_n)$的样本相关系数平方.
\end{enumerate}
\end{exercise}

\begin{exercise}\label{exercise:11.35}
观测$Y_1,\ldots,Y_n$由关系$Y_i=\theta x_i^2+\varepsilon_i$描述，其中$x_1,\ldots,x_n$为固定常数，$\varepsilon_1,\ldots,\varepsilon_n$独立同分布于$\mathrm n(0,\sigma^2)$.
\begin{enumerate}[label=(\alph*)]
\item\label{item:11.35.a} 求$\theta$的最小二乘估计量.
\item\label{item:11.35.b} 求$\theta$的MLE.
\item\label{item:11.35.c} 求$\theta$的最优无偏估计量.
\end{enumerate}
\end{exercise}

\begin{exercise}\label{exercise:11.36}
按照模型$Y_i=\alpha+\beta x_i+\varepsilon_i$观测$Y_1,\ldots,Y_n$，其中$x_1,\ldots,x_n$为固定常数，$\varepsilon_1,\ldots,\varepsilon_n$独立同分布于$\mathrm n(0,\sigma^2)$. 设$\hat\alpha,\hat\beta$为$\alpha,\beta$的MLE.
\sourcepageanchor{493}
\begin{enumerate}[label=(\alph*)]
\item\label{item:11.36.a} 假设$x_1,\ldots,x_n$是独立同分布随机变量$X_1,\ldots,X_n$的观测值，其分布为$\mathrm n(\mu_X,\sigma_X^2)$. 证明，对$X,Y$的联合分布取期望时，仍有$E\hat\alpha=\alpha$和$E\hat\beta=\beta$.
\item\label{item:11.36.b} \hyperref[item:11.36.a]{(a)}的现象不能推广到协方差. 利用$X,Y$的联合分布，计算$\hat\alpha,\hat\beta$的无条件协方差.
\end{enumerate}
\end{exercise}

\begin{exercise}\label{exercise:11.37}
观测随机变量$Y_1,\ldots,Y_n$，它们相互独立，各服从方差为$\sigma^2$的正态分布. 此外，$EY_i=\beta x_i$，其中$\beta$为未知参数，$x_1,\ldots,x_n$为不全为零的固定常数.
\begin{enumerate}[label=(\alph*)]
\item\label{item:11.37.a} 求$\beta$的MLE，并计算其均值和方差.
\item\label{item:11.37.b} 计算$\beta$的无偏估计量方差的Cram\'er--Rao下界.
\item\label{item:11.37.c} 求$\beta$的最优无偏估计量.
\item\label{item:11.37.d} 若可将$x_1,\ldots,x_n$放在给定非退化闭区间$[A,B]$内任意位置，应如何放置？说明理由.
\item\label{item:11.37.e}\index{SI0682@最大概率估计量}\index{SI0897@最大概率估计量} 对给定正数$r$，$\beta$关于$r$的最大概率估计量（maximum probability estimator）是使积分
$$\int_{D-r}^{D+r}f(y_1,\ldots,y_n\mid\beta)\,d\beta$$
最大的$D$值，其中$f(y_1,\ldots,y_n\mid\beta)$是$Y_1,\ldots,Y_n$的联合pdf. 求此估计量.
\end{enumerate}
\end{exercise}

\begin{exercise}\label{exercise:11.38}
一位生态学家取得数据$(x_i,Y_i)$，$i=1,\ldots,n$，其中$x_i$为区域面积，$Y_i$为该区域的苔藓植株数. 用模型$Y_i\sim\operatorname{Poisson}(\theta x_i)$描述数据，各$Y_i$独立.
\begin{enumerate}[label=(\alph*)]
\item\label{item:11.38.a} 证明$\theta$的最小二乘估计量为$\sum x_iY_i/\sum x_i^2$. 证明它的方差为$\theta\sum x_i^3/(\sum x_i^2)^2$，并计算其偏倚.
\item\label{item:11.38.b} 证明$\theta$的MLE为$\sum Y_i/\sum x_i$，方差为$\theta/\sum x_i$. 计算其偏倚.
\item\label{item:11.38.c} 求$\theta$的最优无偏估计量，并证明其方差达到Cram\'er--Rao下界.
\end{enumerate}
\end{exercise}

\begin{exercise}\label{exercise:11.39}
验证\eqref{eq:11.3.42}的同时置信区间具有所声称的覆盖概率.
\end{exercise}

\begin{exercise}\label{exercise:11.40}\leavevmode
\begin{enumerate}[label=(\alph*)]
\item\label{item:11.40.a}\index{SI0681@二次函数的最大值} 证明：若$a,b,c,d$为常数，$c>0,d>0$，则\footnote{译注：若$a=0,b\ne0$，右端仅为$t\to\pm\infty$时的上确界，不在有限$t$处达到；将max改为sup可覆盖此情形. 原题照录.}
$$\max_t\frac{(a+bt)^2}{c+dt^2}=\frac{a^2}c+\frac{b^2}d.$$
\item\label{item:11.40.b} 用\hyperref[item:11.40.a]{(a)}验证\eqref{eq:11.3.44}，从而补全\autoref{thm:11.3.6}中的空缺.
\item\label{item:11.40.c} 用\hyperref[item:11.40.a]{(a)}，根据\eqref{eq:11.3.41}的预测区间求一个Scheff\'e型同时带. 即仿照\autoref{thm:11.3.6}重新表示预测区间，证明
$$\max_t\frac{\bigl((\bar Y-\mu_{\bar Y})+(\hat\beta-\beta)t\bigr)^2}{S^2[1+1/n+t^2/S_{xx}]}
=\frac{\dfrac n{n+1}(\bar Y-\mu_{\bar Y})^2+S_{xx}(\hat\beta-\beta)^2}{S^2}.$$
\item\label{item:11.40.d} 最大值的分布不容易写出，但可以近似. 像\autoref{ex:7.2.3}那样使用矩匹配，近似这个统计量.
\end{enumerate}
\end{exercise}

\begin{exercise}\label{exercise:11.41}
注意\autoref{ex:12.4.2}的讨论中，Potoroo数据有一个观测含缺失值. 假设第24只动物观测到$\mathrm O_2=16.3$.
\begin{enumerate}[label=(\alph*)]
\item\label{item:11.41.a} 写出观测数据的对数似然函数，以及完整数据对数似然函数的期望.
\item\label{item:11.41.b} 描述求MLE的EM算法的E步和M步.
\item\label{item:11.41.c} 用全部24个观测求MLE.
\item\label{item:11.41.d} 实际上，第24只动物未观测到$\mathrm O_2$，而是观测到$\mathrm{CO}_2=4.2$，$\mathrm O_2$缺失. 对此情形建立EM算法并求MLE. （这个问题困难得多，因为现在必须对$x$取期望，这意味着必须用二元正态分布表述回归问题.）
\end{enumerate}
\end{exercise}



\subsection{杂录}\label{sec:11.5}
\sourcepageanchor{494}
\subsubsection{Cochran定理}\label{sec:11.5.1}\index{SI0107@Cochran定理}
正态随机变量的平方和经适当缩放与中心化后，服从卡方分布. 这类结果最早由\hyperref[bib:source80]{Cochran（1934）}给出. Cochran定理给出了缩放所需的充要条件，使独立同分布正态随机变量平方后求和所得的量服从卡方分布. 条件并不难，但最好用矩阵性质来陈述，这里不作讨论. Cochran定理的直接推论是，单因素ANOVA中的卡方随机变量按\autoref{sec:11.2.6}所述方式分解. 另一个推论是，在随机完全区组ANOVA（\hyperref[sec:11.5.3]{杂录11.5.3}）中，各均方都具有卡方分布.

Cochran定理已有推广，现在已知正态变量（未必独立同分布）平方所得量服从卡方分布的充要条件. 细节见\hyperref[bib:source309]{Stuart和Ord（1987}，第15章）.

\subsubsection{多重比较}\label{sec:11.5.2}\index{SI0714@多重比较}
本章介绍了两种同时推断的方法：Scheff\'e方法和使用Bonferroni不等式. 还有大量其他同时推断程序，多数涉及两两比较的推断，即均值之差. 这些程序可用于估计单因素ANOVA中的处理均值.

Tukey提出的一种方法（见\hyperref[bib:source221]{Miller，1981）}有时称为$Q$方法，\index{SI0718@Tukey的Q方法}\semanticindexanchor{718:494}它使用Scheff\'e式的最大化论证，但只对两两差值最大化，而不是对所有对比. $Q$分布是以下统计量的分布：
$$Q=\max_{i,j}\left|\frac{(\bar Y_{i\cdot}-\bar Y_{j\cdot})-(\theta_i-\theta_j)}{\sqrt{S_p^2(1/n+1/n)}}\right|,$$
其中对所有$i$有$n_i=n$. （\hyperref[bib:source164]{Hayter，1984}，证明：若$n_i\ne n_j$，并将上面的$n$替换为调和均值$n_h$，其中$1/n_h=\tfrac12((1/n_i)+(1/n_j))$，所得程序是保守的.）如果只关心两两差值，$Q$方法比Scheff\'e的$S$方法更有功效（区间更短），因此是一种改进. 这很容易看出，因为按定义，$Q$的最大化产生的最大值小于$S$方法.

其他处理两两差值的多重比较程序也比$S$方法更有功效. 其中包括LSD（最小显著差，Least Significant Difference）程序、受保护LSD、Duncan程序和Student--Neumann--Keuls程序. 后两者是多重极差程序（multiple range procedures）：各次比较所用的临界点随比较而改变.

充分理解多重比较程序的一个困难，是第一类错误的定义并非不可变动. 有些程序改变了多重比较中第一类错误的定义，所以“$\alpha$水平”的确切含义未必清楚. 所考虑的一些错误率称为全实验错误率（experimentwise error rate）、每次比较错误率（comparisonwise error rate）和族错误率（familywise error rate）. \hyperref[bib:source221]{Miller（1981）}和\hyperref[bib:source169]{Hsu（1996）}是这方面的良好参考书. \hyperref[bib:source58]{Carmer和Walker（1982）}对此有一篇幽默而富有启发性的论述.

\subsubsection{随机完全区组设计}\label{sec:11.5.3}
\autoref{sec:11.2}讨论数据的单向分类，即实验中只有一种分类（处理）. 一般而言，ANOVA允许多种分类，其中最常用的ANOVA之一是随机完全区组（Randomized Complete Block，RCB）ANOVA.

\sourcepageanchor{495}
区组（block）或区组因子（blocking factor）是在实验中专为消除变异而设置的分类. 与处理不同，通常并不关心区组间差异. 区组化实践起源于农业，实验人员利用相似的生长条件控制实验方差. 为对此建模，把实验中的实际区组看成来自一个大区组总体的随机样本，因而它们成为随机因子（random factor）.

\paragraph{RCB ANOVA假设}
按照下列模型观测随机变量$Y_{ij}$：
$$Y_{ij}\mid\boldsymbol b=\mu+\tau_i+b_j+\varepsilon_{ij},\qquad i=1,\ldots,k,\quad j=1,\ldots,r,$$
其中：
\begin{enumerate}[label=(\roman*)]
\item\label{item:11.5.3.i} 随机变量$\varepsilon_{ij}$独立同分布于$\mathrm n(0,\sigma^2)$，$i=1,\ldots,k$，$j=1,\ldots,r$（方差相等的正态误差）.
\item\label{item:11.5.3.ii} 随机变量$B_1,\ldots,B_r$的实现值（未观测）为区组$b_1,\ldots,b_r$；它们独立同分布于$\mathrm n(0,\sigma_B^2)$，并与所有$i,j$的$\varepsilon_{ij}$独立.
\end{enumerate}
$Y_{ij}$的均值、方差为
$$EY_{ij}=\mu+\tau_i,\qquad\operatorname{Var}Y_{ij}=\sigma_B^2+\sigma^2.$$
此外，虽然$Y_{ij}$在给定条件下互不相关，但不取条件时，区组内部存在相关性. 区组$j$中$Y_{ij}$与$Y_{i'j}$（$i\ne i'$）的相关系数为
$$\frac{\operatorname{Cov}(Y_{ij},Y_{i'j})}{\sqrt{(\operatorname{Var}Y_{ij})(\operatorname{Var}Y_{i'j})}}
=\frac{\sigma_B^2}{\sigma_B^2+\sigma^2},$$
这个量称为组内相关系数（intraclass correlation）. 因此，模型不仅意味着区组内有相关，而且意味着正相关. 这是加性模型以及$\varepsilon$与$B$独立的假设所带来的结果（\autoref{exercise:11.23}\mbox{）.} 即使$Y_{ij}$并不独立，这种组内相关结构仍导出均方之比具有$F$分布的方差分析（\hyperref[sec:11.5.1]{杂录11.5.1}）.

\subsubsection{其他类型的方差分析}\label{sec:11.5.4}\index{SI0045@其他模型}
所考虑的两种ANOVA——单因素ANOVA和RCB ANOVA——是最简单的类型. 例如，完全区组设计的一种推广是不完全区组设计（incomplete block design）. 有时实际条件限制，无法将所有处理放入每个区组，就需要不完全区组设计. 决定怎样在这种设计中安排处理，既困难又关键. 当然，设计越复杂，分析也越复杂.

统计设计研究如何用最少的观测获得最多的信息，这一主题在许多情形下导向更复杂、更有效率的ANOVA. 基于部分因子设计（fractional factorials）、拉丁方（Latin squares）、平衡不完全区组（balanced incomplete blocks）等设计的ANOVA，可以有效收集关于某一现象的大量信息. 这一主题较好的综合参考书是\hyperref[bib:source82]{Cochran和Cox（1957）}、\hyperref[bib:source94]{Dean和Voss（1999）}以及\hyperref[bib:source192]{Kuehl（2000）}.

\subsubsection{置信带的形状}\label{sec:11.5.5}
置信带有许多形状，不仅限于Scheff\'e置信带规定的双曲线形状. 例如，\hyperref[bib:source131]{Gafarian（1964）}说明了如何在有限区间上构造直线置信带. Gafarian型置信带允许作如下形式的陈述：
$$P\bigl(\hat\alpha+\hat\beta x-d_\alpha\le\alpha+\beta x\le\hat\alpha+\hat\beta x+d_\alpha
\text{，对所有 }x\in[a,b]\bigr)=1-\alpha.$$

\sourcepageanchor{496}
Gafarian给出了$d_\alpha$的数值表. 有限宽度的置信带必然只能适用于有限的$x$范围. 任意$1-\alpha$水平的置信带，在$|x|\to\infty$时都必须具有无限长度.

\hyperref[bib:source71]{Casella和Strawderman（1980）}等人说明了如何在有限区间上构造Scheff\'e型置信带，在维持与无限区间Scheff\'e置信带相同置信度的同时缩小宽度. \hyperref[bib:source227]{Naiman（1983）}比较了有限区间上直线置信带与Scheff\'e置信带的表现. 按他的平均宽度准则，Scheff\'e置信带更优. 不过，有些情形中，实验人员可能觉得直线置信带的解释更自然.

除了直线和双曲线，还有其他可能的形状. \hyperref[bib:source241]{Piegorsch（1985）}研究并刻画了在概率陈述无法再改进意义下的容许形状，得到了容许置信带必须满足的“增长条件”. Naiman（\hyperref[bib:source227]{1983}，\hyperref[bib:source228]{1984}，\hyperref[bib:source229]{1987}）及Naiman和Wynn（\hyperref[bib:source230]{1992}，\hyperref[bib:source231]{1997}）将这一理论发展到很高水平，建立有用的不等式和几何恒等式，进一步改进推断.

\subsubsection{Stein悖论}\label{sec:11.5.6}\index{SI0048@Stein估计}\index{SI0983@方差分析}
方差分析的一部分涉及一组正态均值的同时估计. 自\hyperref[bib:source301]{Stein（1956）}开始，对这一特定问题的发展深刻影响了点估计的理论和应用.

方差分析的一个典范形式是观测$\boldsymbol X=(X_1,\ldots,X_p)$，其中各$X_i\sim\mathrm n(\theta_i,1)$相互独立，$i=1,\ldots,p$，目标是估计$\boldsymbol\theta=(\theta_1,\ldots,\theta_p)$. 通常用$X_i$估计$\theta_i$，但\hyperref[bib:source301]{Stein（1956）}建立了一个令人惊讶的结果：若$p\ge3$，则下面的$\theta_i$估计量
$$\delta_S^i(\boldsymbol X)=\left(1-\frac{p-2}{\sum_{i=1}^pX_i^2}\right)X_i$$
在下述意义下是$\theta_i$的更好估计量：
$$\sum_{i=1}^pE_{\boldsymbol\theta}(X_i-\theta_i)^2
\ge\sum_{i=1}^pE_{\boldsymbol\theta}\bigl(\delta_S^i(\boldsymbol X)-\theta_i\bigr)^2.$$
即Stein估计量的均方误差之和总较小，通常严格小于$\boldsymbol X$的均方误差之和.

注意，比较估计量时使用的是各分量均方误差之和，而每个$\delta_S^i$都可以是整个向量$(X_1,\ldots,X_p)$的函数，因此估计每个均值时都可以使用全部数据. 由于$X_i$相互独立，可能以为把$\delta_S^i$限制为仅是$X_i$的函数就已足够. 但对均方误差求和，就把各分量联系在一起.

在单因素ANOVA中，观测到
$$\bar Y_{i\cdot}\sim\mathrm n\left(\theta_i,\frac{\sigma^2}{n_i}\right),\qquad i=1,\ldots,k,\quad\text{相互独立},$$
其中$\bar Y_{i\cdot}$为单元均值. Stein估计量具有形式
$$\delta_S^i(\bar Y_{1\cdot},\ldots,\bar Y_{k\cdot})
=\left(1-\frac{(k-2)\sigma^2}{\sum n_j\bar Y_{j\cdot}^{\,2}}\right)^+\bar Y_{i\cdot},
\qquad i=1,\ldots,k.$$
还可以选择有意义的收缩目标，进一步改进这一Stein型估计量（上述估计量向零收缩\mbox{）.} \hyperref[bib:source205]{Lindley（1962）}提出的一种估计量向观测的总均值收缩，其形式为
$$\delta_L^i(\bar Y_{1\cdot},\ldots,\bar Y_{k\cdot})
=\bar{\bar Y}+\left(1-\frac{(k-3)\sigma^2}{\sum n_j(\bar Y_{j\cdot}-\bar{\bar Y})^2}\right)^+
(\bar Y_{i\cdot}-\bar{\bar Y}),\qquad i=1,\ldots,k.$$

\sourcepageanchor{497}
选择其他收缩目标可能更合适. \hyperref[bib:source68]{Casella和Hwang（1987）}对此作了讨论，包括改进Scheff\'e $S$方法等置信陈述的方法. \hyperref[bib:source225]{Morris（1983）}也讨论了这些估计量的应用.

使用Stein型估计量已有许多理论发展，不仅涉及点估计，还涉及置信集估计；已证明，将中心移到Stein估计量处，能提高覆盖概率并缩小置信集. Stein估计量与经验贝叶斯估计量（\hyperref[sec:7.5.6]{杂录7.5.6}）也有密切联系，这最早由Efron和Morris在20世纪70年代的一系列论文（\hyperref[bib:source111]{1972}，\hyperref[bib:source112]{1973}，\hyperref[bib:source113]{1975}）中揭示：利用共同的先验分布把$\boldsymbol\theta$的各分量联系起来. \hyperref[bib:source200]{Lehmann和Casella（1998}，第5章）介绍了这一理论以及Stein估计量的一些应用.





\section{回归模型}\label{sec:12}
\chapterdoi{12}
\sourcepageanchor{498}
\begin{quote}
他的结果会使不懂内情的人如此惊讶，以至于在了解他取得这些结果的过程之前，他们很可能把他当作巫师.
\begin{flushright}华生医生，谈夏洛克·福尔摩斯\\《血字的研究》\end{flushright}
\end{quote}

\subsection{引言}\label{sec:12.1}
\autoref{sec:11}讨论了可称为“经典线性模型”的内容. 方差分析和简单线性回归都基于具有正态误差的线性模型. 本章考察这一模型的一些推广，它们在实际问题中已证明很有用.

在\autoref{sec:12.2}中，将线性模型推广到预测变量含误差的模型，称为变量含误差的回归（regression with errors in variables，EIV）. 此时预测变量$X$像响应变量$Y$一样成为随机变量. 这一模型中的估计会遇到许多意想不到的困难，与简单线性回归模型可能很不相同.

\autoref{sec:12.3}进一步推广线性模型，考察逻辑回归（logistic regression）. 此时响应变量是离散的Bernoulli变量. Bernoulli均值是有界函数，在有界函数上建立线性模型可能遇到问题，尤其在边界附近. 因此，用logit变换将均值变为无界参数，再把变换后的参数建模为预测变量的线性函数. 当对响应均值的某个函数建立线性模型时，就得到广义线性模型（generalized linear model）.

最后，在\hyperref[sec:12.4]{节11.3}中考察线性回归中的稳健性.\footnote{译注：原文这里写“Section 11.3”，按本章内容应为\autoref{sec:12.4}. 原编号保留.} 本章其他节改变模型，这里则改变拟合准则. 其发展与\autoref{sec:10.2.2}关于稳健点估计的讨论平行. 即以基于$\rho$函数的准则代替最小二乘准则，使估计对观测值较不敏感，同时保留一定效率.



\subsection{变量含误差的回归}\label{sec:12.2}\index{SI0407@变量含误差的回归}
变量含误差的回归（EIV），又称测量误差模型（measurement error model），与\autoref{sec:11.3}的简单线性回归有根本差异，最好把它看成完全不同的主题. 这里将其作为通常回归模型的推广，主要出于传统原因. 然而，这一模型引出的问题很不相同.

本节模型在以下意义上推广简单线性回归：仍处理形如
\begin{equation}\label{eq:12.2.1}Y_i=\alpha+\beta x_i+\varepsilon_i,\end{equation}
的模型，但不再假定$x$已知. 相反，可以测量一个均值为$x_i$的随机变量. （为符合既定符号约定，将把所测随机变量记为$X_i$，其均值记为$\xi_i$，而非$x_i$.）

\sourcepageanchor{499}
这里旨在展示处理EIV模型的不同方法，给出一些标准解法及其有时令人意外的困难. 更全面的介绍可参阅\hyperref[bib:source147]{Gleser（1991）}的综述、\hyperref[bib:source130]{Fuller（1987）}及\hyperref[bib:source62]{Carroll、Ruppert和Stefanski（1995）}的专著，以及\hyperref[bib:source55]{Brown和Fuller（1991）}主编的文集. \hyperref[bib:source189]{Kendall和Stuart（1979}，第29章）也较详细地讨论了这个主题.

一般EIV模型假设观测到成对数据$(x_i,y_i)$，它们来自随机变量$(X_i,Y_i)$，其均值满足线性关系
\begin{equation}\label{eq:12.2.2}EY_i=\alpha+\beta(EX_i).\end{equation}
若定义
$$EY_i=\eta_i\quad\text{和}\quad EX_i=\xi_i,$$
则\eqref{eq:12.2.2}成为
\begin{equation}\label{eq:12.2.3}\eta_i=\alpha+\beta\xi_i,\end{equation}
即随机变量均值之间的线性关系.

变量$\xi_i,\eta_i$有时称为潜变量（latent variables），指不能直接测量的量.\index{SI0602@潜在变量}\semanticindexanchor{602:499} 潜变量不仅可能无法直接测量，甚至可能根本无法测量. 例如，人的智商IQ无法测量. 我们可以测量智商测试的分数，但永远无法测量IQ这个变量. 然而，人们常常假定IQ与其他变量存在关系.

\eqref{eq:12.2.2}所规定的模型实际上没有区分$X$与$Y$. 但若关心回归，就应有理由选择$Y$为响应、$X$为预测变量. 牢记这一关于$Y$对$X$回归的规定，把变量含误差模型或测量误差模型定义如下：按照
\begin{equation}\label{eq:12.2.4}
\begin{aligned}Y_i&=\alpha+\beta\xi_i+\varepsilon_i,&\varepsilon_i&\sim\mathrm n(0,\sigma_\varepsilon^2),\\
X_i&=\xi_i+\delta_i,&\delta_i&\sim\mathrm n(0,\sigma_\delta^2)
\end{aligned}
\end{equation}
观测$(X_i,Y_i)$，$i=1,\ldots,n$.

注意，正态性假设虽然常见，却并非必要. 也可采用其他分布. 实际上，这一模型遇到的一些问题正由正态性假设造成. （例如见\hyperref[bib:source295]{Solari，1969}.）

\begin{example}[估计大气压]\label{ex:12.2.1}
当$x$变量与$y$变量一起被观测，而不是受控制时，EIV回归模型会很自然地出现. 例如，19世纪苏格兰物理学家J. D. Forbes试图用水的沸点测量值估计海拔. 为此，他同时测量沸点和大气压，后者可用于得到海拔. 由于19世纪的气压计十分易损，若能从温度估计气压，更确切地说，估计$\log(\text{气压})$，就会很有用. 在九个地点观测的数据如下，对这种情形采用EIV模型是合理的.
\begin{center}\begin{tabular}{rr}\toprule
沸点（$^\circ\mathrm F$）&$\log(\text{气压})$（$\log(\mathrm{Hg})$）\\\midrule
194.5&1.3179\\197.9&1.3502\\199.4&1.3646\\200.9&1.3782\\201.4&1.3806\\203.6&1.4004\\209.5&1.4547\\210.7&1.4630\\212.2&1.4780\\\bottomrule
\end{tabular}\end{center}
\end{example}

\sourcepageanchor{500}
已经见过模型\eqref{eq:12.2.4}的若干特殊情形. 若$\delta_i=0$，模型就成为简单线性回归，因为此时没有测量误差，可以直接观测$\xi_i$. 若$\alpha=0$，则有
$$\begin{aligned}Y_i&\sim\mathrm n(\eta_i,\sigma_\varepsilon^2),&i&=1,\ldots,n,\\
X_i&\sim\mathrm n(\xi_i,\sigma_\delta^2),&i&=1,\ldots,n,
\end{aligned}$$
其中可能$\sigma_\delta^2\ne\sigma_\varepsilon^2$，这是Behrens--Fisher问题的一个版本.\index{SI0087@Behrens–Fisher问题}\semanticindexanchor{87:500}

\subsubsection{函数关系与结构关系}\label{sec:12.2.1}
EIV模型可规定两类不同的关系：函数线性关系（functional linear relationship）和结构线性关系（structural linear relationship）. 不同的关系规定可能产生不同估计量及不同性质. 正如\hyperref[bib:source222]{Moran（1971）}所言：“这些术语不太令人满意，但我们仍沿用它们，因为这种区别至关重要……” 杂录中给出了这些术语的一些解释. 这里暂且只介绍两个模型.

\paragraph{线性函数关系模型}
这就是\eqref{eq:12.2.4}给出的模型：随机变量$X_i,Y_i$满足$EX_i=\xi_i$、$EY_i=\eta_i$，并假设函数关系
$$\eta_i=\alpha+\beta\xi_i.$$
按照
\begin{equation}\label{eq:12.2.5}\index{SI0410@函数关系}
\begin{aligned}Y_i&=\alpha+\beta\xi_i+\varepsilon_i,&\varepsilon_i&\sim\mathrm n(0,\sigma_\varepsilon^2),\\
X_i&=\xi_i+\delta_i,&\delta_i&\sim\mathrm n(0,\sigma_\delta^2),
\end{aligned}
\end{equation}
观测$(X_i,Y_i)$，$i=1,\ldots,n$，其中$\xi_i$为固定的未知参数，$\varepsilon_i$与$\delta_i$相互独立. 主要关心的参数为$\alpha,\beta$；利用给定$\xi_1,\ldots,\xi_n$时$((X_1,Y_1),\allowbreak\ldots,\allowbreak(X_n,Y_n))$的联合分布，对这些参数作推断.

\paragraph{线性结构关系模型}
可把这一模型看成通过以下层次结构对函数关系模型的推广. 与函数关系模型一样，随机变量$X_i,Y_i$满足$EX_i=\xi_i$、$EY_i=\eta_i$，并假设函数关系$\eta_i=\alpha+\beta\xi_i$. 但现在假定参数$\xi_1,\ldots,\xi_n$本身是来自同一总体的随机样本. 因此，给定$\xi_1,\ldots,\xi_n$时，按照
\begin{equation}\label{eq:12.2.6}\index{SI0417@结构关系}
\begin{aligned}Y_i&=\alpha+\beta\xi_i+\varepsilon_i,&\varepsilon_i&\sim\mathrm n(0,\sigma_\varepsilon^2),\\
X_i&=\xi_i+\delta_i,&\delta_i&\sim\mathrm n(0,\sigma_\delta^2),
\end{aligned}
\end{equation}
观测$(X_i,Y_i)$，$i=1,\ldots,n$，并且
$$\xi_i\sim\text{iid }\mathrm n(\xi,\sigma_\xi^2).$$
如前，$\varepsilon_i$与$\delta_i$相互独立，并且也独立于$\xi_i$. 与函数关系模型一样，主要关心的参数为$\alpha,\beta$. 但这里对它们的推断使用$((X_1,Y_1),\allowbreak\ldots,\allowbreak(X_n,Y_n))$关于$\xi_1,\ldots,\xi_n$的无条件联合分布. （即按\eqref{eq:12.2.6}中的分布积分消去$\xi_1,\ldots,\xi_n$.）

两个模型十分相似，一个模型中估计量的统计性质，例如相合性，常能转移到另一个模型. 更确切地说，在函数模型中相合的估计量，在结构模型中也相合（\hyperref[bib:source236]{Nussbaum，1976}；或\hyperref[bib:source144]{Gleser，1983）}.
\sourcepageanchor{501}
这很合理，因为函数模型是结构模型的“条件版本”. 在函数模型中相合的估计量，必须对$\xi_i$的所有取值都相合，所以在对$\xi_i$平均的结构模型中必然也相合. 逆命题不成立. 不过，从结构关系模型到函数关系模型有一个有用推论：若参数在结构模型中不可识别，则在函数模型中也不可识别.\index{SI0411@可识别性}\semanticindexanchor{411:501}\index{SI0528@变量含误差}\semanticindexanchor{528:501} （见\autoref{def:11.2.2}.）

\begin{figure}[htbp]\centering
\includegraphics[width=.60\textwidth]{assets/chapter12/fig-12.2.1.pdf}
\caption{正交最小二乘所最小化的距离}\label{fig:12.2.1}
\end{figure}

后面将看到，两个模型有类似问题，并且在某些情形下有类似的似然解. 在结构模型中建立统计理论可能更容易，但在许多情形下函数模型往往显得更合理. 因而，这些内在相似性很有用.

如前所述，两个模型的主要差异之一在于对描述回归关系的参数$\alpha,\beta$的推断. 这种差异极为重要，怎样强调也不过分. 函数关系模型中，给定$\xi_1,\ldots,\xi_n$，利用$X,Y$的条件联合分布进行推断. 结构关系模型中，则利用积分消去$\xi_1,\ldots,\xi_n$后$X,Y$的边际分布，进行不以$\xi_1,\ldots,\xi_n$为条件的推断.

\subsubsection{一种最小二乘解法}\label{sec:12.2.2}\index{SI0609@变量含误差}\index{SI0739@正交最小二乘}
像\autoref{sec:11.3.1}那样，暂时抛开统计，尝试寻找通过观测点$(x_i,y_i)$，$i=1,\ldots,n$的“最佳”直线. 先前假定$x$测量无误差，因此考虑最小化竖直距离是合理的. 这种距离度量隐含假定$x$值正确，得到普通最小二乘（ordinary least squares）. 但这里$x$也有误差，没有理由只考虑竖直距离. 事实上，从统计角度看，普通最小二乘在EIV模型中存在一些问题（见杂录）.

考虑$x$测量也有误差的一种方法，是进行正交最小二乘（orthogonal least squares），即寻找使正交于直线的距离而非竖直距离最小的直线（见\autoref{fig:12.2.1}）. 这种距离度量不像普通最小二乘那样偏向$x$，而是平等对待两个变量. 它也称为总体最小二乘法（total least squares）.\index{SI1035@全最小二乘}\semanticindexanchor{1035:501} 参看\autoref{fig:12.2.1}，对于给定数据点$(x',y')$，在直线$y=a+bx$上按正交距离衡量最近的点由下式给出（\autoref{exercise:12.1}）：
\sourcepageanchor{502}
\begin{equation}\label{eq:12.2.7}
\hat x'=\frac{by'+x'-ab}{1+b^2},\qquad
\hat y'=a+\frac b{1+b^2}(by'+x'-ab).
\end{equation}
现假设有数据$(x_i,y_i)$，$i=1,\ldots,n$. 观测点$(x_i,y_i)$到直线$y=a+bx$上最近点的距离平方为$(x_i-\hat x_i)^2+(y_i-\hat y_i)^2$，其中$\hat x_i,\hat y_i$由\eqref{eq:12.2.7}定义. 总体最小二乘问题是在所有$a,b$上最小化
$$\sum_{i=1}^n\bigl[(x_i-\hat x_i)^2+(y_i-\hat y_i)^2\bigr].$$
容易证明
\begin{equation}\label{eq:12.2.8}
\begin{aligned}
&\sum_{i=1}^n\bigl[(x_i-\hat x_i)^2+(y_i-\hat y_i)^2\bigr]\\
&=\sum_{i=1}^n\left\{\frac{b^2}{(1+b^2)^2}[y_i-(a+bx_i)]^2
+\frac1{(1+b^2)^2}[y_i-(a+bx_i)]^2\right\}\\
&=\frac1{1+b^2}\sum_{i=1}^n(y_i-(a+bx_i))^2.
\end{aligned}
\end{equation}
固定$b$时，和式前的因子为常数. 因而，与\eqref{eq:11.3.9}一样，使和式最小的$a$为$a=\bar y-b\bar x$.\index{SI0610@截距}\semanticindexanchor{610:502} 代回\eqref{eq:12.2.8}，总体最小二乘解就是使
\begin{equation}\label{eq:12.2.9}
\frac1{1+b^2}\sum_{i=1}^n\bigl((y_i-\bar y)-b(x_i-\bar x)\bigr)^2
\end{equation}
在所有$b$上最小的解.

如\eqref{eq:11.3.6}和\eqref{eq:11.3.6}那样，定义平方和与乘积和为\footnote{译注：原文将同一式号11.3.6重复两次；平方和见\eqref{eq:11.3.6}，乘积和见\eqref{eq:11.3.7}. 原重复保留.}
\begin{equation}\label{eq:12.2.10}
\begin{aligned}S_{xx}&=\sum_{i=1}^n(x_i-\bar x)^2,\qquad S_{yy}=\sum_{i=1}^n(y_i-\bar y)^2,\\
S_{xy}&=\sum_{i=1}^n(x_i-\bar x)(y_i-\bar y).
\end{aligned}
\end{equation}
展开平方并求和，\eqref{eq:12.2.9}成为
$$\frac1{1+b^2}\bigl[S_{yy}-2bS_{xy}+b^2S_{xx}\bigr].$$
通常的微积分方法给出最小值（\hyperref[exercise:12.2]{习题12.20}），\footnote{译注：原文引用习题12.20，但该题讨论逻辑回归；这里对应\autoref{exercise:12.2}. 原编号保留.} 得到正交最小二乘直线$y=a+bx$，其中
\begin{equation}\label{eq:12.2.11}\index{SI0612@斜率}
a=\bar y-b\bar x,\qquad b=\frac{-(S_{xx}-S_{yy})+\sqrt{(S_{xx}-S_{yy})^2+4S_{xy}^2}}{2S_{xy}}.
\end{equation}
可以预料，这条直线与最小二乘直线不同. 事实上，后面将看到，它总位于$y$对$x$的普通回归与$x$对$y$的普通回归之间. \autoref{fig:12.2.2}说明了这一点：用\autoref{tab:11.3.1}的数据算得正交最小二乘直线$\hat y=-.49+1.88x$.

简单线性回归中，已看到在正态性假设下，$\alpha,\beta$的普通最小二乘解与MLE相同. 这里，正交最小二乘解只在对参数作出某些假定的特殊情形下才是MLE.

似然估计将遇到的困难，再次表明数学解与统计解的区别. 直线拟合问题的数学最小二乘解并不难获得，似然解却不会如此.

\sourcepageanchor{503}
\begin{figure}[htbp]\centering
\includegraphics[width=.85\textwidth]{assets/chapter12/fig-12.2.2.pdf}
\caption{\autoref{tab:11.3.1}数据的三条回归直线}\label{fig:12.2.2}
\end{figure}
\subsubsection{最大似然估计}\label{sec:12.2.3}
先考虑函数线性关系模型的最大似然解，结构关系模型的情形类似，并且在某些方面更容易. 利用正态性假设，函数关系模型可表示为
$$Y_i\sim\mathrm n(\alpha+\beta\xi_i,\sigma_\varepsilon^2)
\quad\text{和}\quad X_i\sim\mathrm n(\xi_i,\sigma_\delta^2),\quad i=1,\ldots,n,$$
各$X_i,Y_i$相互独立. 给定观测$(\boldsymbol x,\boldsymbol y)=((x_1,y_1),\allowbreak\ldots,\allowbreak(x_n,y_n))$，似然函数为
\begin{equation}\label{eq:12.2.12}
\begin{aligned}
&L\bigl(\alpha,\beta,\xi_1,\ldots,\xi_n,\sigma_\delta^2,\sigma_\varepsilon^2\mid(\boldsymbol x,\boldsymbol y)\bigr)\\
&=\frac1{(2\pi)^n}\frac1{(\sigma_\delta^2\sigma_\varepsilon^2)^{n/2}}
\exp\left(-\sum_{i=1}^n\frac{(x_i-\xi_i)^2}{2\sigma_\delta^2}\right)
\exp\left(-\sum_{i=1}^n\frac{(y_i-(\alpha+\beta\xi_i))^2}{2\sigma_\varepsilon^2}\right).
\end{aligned}
\end{equation}
\sourcepageanchor{504}
这个似然函数的问题在于没有有限的最大值. 取参数配置$\xi_i=x_i$，再令$\sigma_\delta^2\to0$，函数值便趋于无穷，说明最大似然解不存在. 事实上，\hyperref[bib:source295]{Solari（1969）}证明，若将$L$的一阶导数方程置零并求解，结果是鞍点而非最大值. 注意，只要参数完全由我们控制，就总能迫使似然函数趋于无穷. 特别地，总能使某个方差趋于零，同时让指数项保持有界.

我们作一个常见假设，它既合理，又可缓解许多问题：$\sigma_\delta^2=\lambda\sigma_\varepsilon^2$，其中$\lambda>0$固定且已知. （关于方差的其他假设，见\hyperref[bib:source189]{Kendall和Stuart，1979}，第29章.）这一假设属于限制最少的假设之一，只说已知方差之比，而不知各方差的值. 而且，由此得到的模型具有较好的性质.

在这一假设下，似然函数可写为
\begin{equation}\label{eq:12.2.13}
\begin{aligned}
&L\bigl(\alpha,\beta,\xi_1,\ldots,\xi_n,\sigma_\delta^2\mid(\boldsymbol x,\boldsymbol y)\bigr)\\
&=\frac1{(2\pi)^n}\frac{\lambda^{n/2}}{(\sigma_\delta^2)^n}
\exp\left[-\sum_{i=1}^n\frac{(x_i-\xi_i)^2+\lambda(y_i-(\alpha+\beta\xi_i))^2}{2\sigma_\delta^2}\right],
\end{aligned}
\end{equation}
现在可以将它最大化. 将分阶段最大化，确保每一步都得到最大值后再进行下一步. 考察\eqref{eq:12.2.13}，可以确定合理的最大化顺序.

首先，对每个$\alpha,\beta,\sigma_\delta^2$值，为使$L$关于$\xi_1,\ldots,\xi_n$最大，最小化$\sum_{i=1}^n\bigl((x_i-\xi_i)^2+\lambda(y_i-(\alpha+\beta\xi_i))^2\bigr)$. （细节见\autoref{exercise:12.3}.）对每个$i$，这是$\xi_i$的二次式，最小值在
$$\xi_i^*=\frac{x_i+\lambda\beta(y_i-\alpha)}{1+\lambda\beta^2}$$
处取得. 代回得到
$$\sum_{i=1}^n\bigl((x_i-\xi_i^*)^2+\lambda(y_i-(\alpha+\beta\xi_i^*))^2\bigr)
=\frac\lambda{1+\lambda\beta^2}\sum_{i=1}^n(y_i-(\alpha+\beta x_i))^2.$$
似然函数现在成为
\begin{equation}\label{eq:12.2.14}
\begin{aligned}
&\max_{\xi_1,\ldots,\xi_n}L(\alpha,\beta,\xi_1,\ldots,\xi_n,\sigma_\delta^2\mid(\boldsymbol x,\boldsymbol y))\\
&=\frac1{(2\pi)^n}\frac{\lambda^{n/2}}{(\sigma_\delta^2)^n}
\exp\left\{-\frac1{2\sigma_\delta^2}\left[\frac\lambda{1+\lambda\beta^2}\sum_{i=1}^n(y_i-(\alpha+\beta x_i))^2\right]\right\}.
\end{aligned}
\end{equation}
现在可以关于$\alpha,\beta$最大化，但稍加计算就会发现，已经在正交最小二乘解中做过这件事了！确实，EIV模型的正交最小二乘与最大似然之间存在某种对应，下面就利用它. 定义
\begin{equation}\label{eq:12.2.15}
\alpha^*=\sqrt\lambda\alpha,\quad\beta^*=\sqrt\lambda\beta,\quad
y_i^*=\sqrt\lambda y_i,\quad i=1,\ldots,n.
\end{equation}
\eqref{eq:12.2.14}中的指数部分成为
$$\frac\lambda{1+\lambda\beta^2}\sum_{i=1}^n(y_i-(\alpha+\beta x_i))^2
=\frac1{1+\beta^{*2}}\sum_{i=1}^n(y_i^*-(\alpha^*+\beta^*x_i))^2,$$
\sourcepageanchor{505}
这与正交最小二乘问题中的表达式相同. 由\eqref{eq:12.2.11}知道使其最小的$\alpha^*,\beta^*$，再利用\eqref{eq:12.2.15}，得到截距和斜率的MLE：
\begin{equation}\label{eq:12.2.16}\index{SI0413@函数型}
\hat\alpha=\bar y-\hat\beta\bar x\quad\text{和}\quad
\hat\beta=\frac{-(S_{xx}-\lambda S_{yy})+\sqrt{(S_{xx}-\lambda S_{yy})^2+4\lambda S_{xy}^2}}{2\lambda S_{xy}}.
\end{equation}
从公式可见，当$\lambda=1$时，MLE与正交最小二乘解一致. 这很合理：正交最小二乘把$x,y$当作具有相同大小的误差，即方差比为1. 将这一论证推进一步，可把此解与假定$x$固定时的普通最小二乘或最大似然联系起来. 若$x$固定，其方差为零，因而$\lambda=0$. 一般$\lambda$下的最大似然解在这种情形下确实退化为普通最小二乘解. \autoref{exercise:12.4}探讨这一关系及其他关系.

将\eqref{eq:12.2.16}与\eqref{eq:12.2.14}结合，似然几乎已完全最大化. 有
\begin{equation}\label{eq:12.2.17}
\begin{aligned}
&\max_{\alpha,\beta,\xi_1,\ldots,\xi_n}L(\alpha,\beta,\xi_1,\ldots,\xi_n,\sigma_\delta^2\mid(\boldsymbol x,\boldsymbol y))\\
&=\frac1{(2\pi)^n}\frac{\lambda^{n/2}}{(\sigma_\delta^2)^n}
\exp\left[-\frac1{2\sigma_\delta^2}\frac\lambda{1+\lambda\hat\beta^2}\sum_{i=1}^n\bigl(y_i-(\hat\alpha+\hat\beta x_i)\bigr)^2\right].
\end{aligned}
\end{equation}
现在关于$\sigma_\delta^2$最大化$L$，与通常正态抽样中求$\sigma^2$的MLE非常相似（\autoref{ex:7.2.11}），主要差异在于$\sigma_\delta^2$的指数为$n$而非$n/2$. 细节留给\autoref{exercise:12.5}. 得到$\sigma_\delta^2$的MLE为
\begin{equation}\label{eq:12.2.18}
\hat\sigma_\delta^2=\frac1{2n}\frac\lambda{1+\lambda\hat\beta^2}\sum_{i=1}^n\bigl(y_i-(\hat\alpha+\hat\beta x_i)\bigr)^2.
\end{equation}
由MLE的性质，$\sigma_\varepsilon^2$的MLE为$\hat\sigma_\varepsilon^2=\hat\sigma_\delta^2/\lambda$，并且$\hat\xi_i=\hat\alpha+\hat\beta x_i$.\footnote{译注：原文此处将$\hat\xi_i$写成拟合响应. 根据前一页对$\xi_i$的最大化，实际应为$[x_i+\lambda\hat\beta(y_i-\hat\alpha)]/(1+\lambda\hat\beta^2)$. 原式保留.} 尽管通常不关心$\hat\xi_i$，但需要预测时它们有时会有用. 此外，$\hat\xi_i$也可用于检验拟合是否充分（见\hyperref[bib:source130]{Fuller，1987）}.

有趣的是，虽然$\hat\alpha,\hat\beta$是相合估计量，$\sigma_\delta^2$却不是.\index{SI0409@相合性}\semanticindexanchor{409:505}\index{SI0675@变量含误差}\semanticindexanchor{675:505}\footnote{译注：此处讨论估计量的相合性，原文$\sigma_\delta^2$应指$\hat\sigma_\delta^2$，下式也如此表明. 原符号保留.} 更确切地说，当$n\to\infty$时，
$$\hat\alpha\to\alpha\quad\text{依概率},\qquad\hat\beta\to\beta\quad\text{依概率},$$
但
$$\hat\sigma_\delta^2\to\tfrac12\sigma_\delta^2\quad\text{依概率}.$$
\hyperref[bib:source143]{Gleser（1981）}给出了EIV函数关系模型中相合性的一般结果.

现在转向线性结构关系模型.\index{SI0415@结构型}\semanticindexanchor{415:505} 回顾这里假设按照
$$\begin{aligned}Y_i&\sim\mathrm n(\alpha+\beta\xi_i,\sigma_\varepsilon^2),\\
X_i&\sim\mathrm n(\xi_i,\sigma_\delta^2),\\
\xi_i&\sim\mathrm n(\xi,\sigma_\xi^2),
\end{aligned}$$
观测$(X_i,Y_i)$，$i=1,\ldots,n$，其中$\xi_i$相互独立，给定$\xi_i$时，各$X_i,Y_i$相互独立. 如前，对$\alpha,\beta$的推断将使用$X_i,Y_i$的边际分布，即积分消去$\xi_i$所得的分布.
\sourcepageanchor{506}
积分消去$\xi_i$，得到$(X_i,Y_i)$的边际分布（\autoref{exercise:12.6}）：
\begin{equation}\label{eq:12.2.19}
(X_i,Y_i)\sim\text{二元正态}\bigl(\xi,\alpha+\beta\xi,\sigma_\delta^2+\sigma_\xi^2,
\sigma_\varepsilon^2+\beta^2\sigma_\xi^2,\beta\sigma_\xi^2\bigr).
\end{equation}
注意这一相关结构与RCB ANOVA\index{SI0416@与随机完全区组的关系}\semanticindexanchor{416:506}\index{SI0875@与变量含误差的关系}\semanticindexanchor{875:506}（\hyperref[sec:11.5.3]{杂录11.5.3}）的相似之处. 在那里，给定区组时观测不相关，但无条件时存在组内相关. 这里，以$\xi_i$为条件的函数关系模型的观测不相关，但在不以$\xi_i$为条件作推断的结构关系模型中，观测相关. $\xi_i$起着类似区组的作用，这里的相关类似于组内相关. （事实上，若$\beta=1$且$\sigma_\delta^2=\sigma_\varepsilon^2$，它就等于组内相关.）

在这种情形下作似然估计，给定观测$(\boldsymbol x,\boldsymbol y)=((x_1,y_1),\allowbreak\ldots,\allowbreak(x_n,y_n))$，似然函数是二元正态的似然，如\autoref{exercise:7.18}中所见. 在那里，二元正态的似然估计可通过令样本量等于总体量求得. 因此，为求$\alpha,\beta,\xi,\sigma_\varepsilon^2,\sigma_\delta^2,\sigma_\xi^2$的MLE，解方程
\begin{equation}\label{eq:12.2.20}
\begin{aligned}
\bar y&=\hat\alpha+\hat\beta\hat\xi,&\bar x&=\hat\xi,\\
\frac1nS_{yy}&=\hat\sigma_\varepsilon^2+\hat\beta^2\hat\sigma_\xi^2,\\
\frac1nS_{xx}&=\hat\sigma_\delta^2+\hat\sigma_\xi^2,\\
\frac1nS_{xy}&=\hat\beta\hat\sigma_\xi^2.
\end{aligned}
\end{equation}
五个方程却有六个未知数，所以方程组不定. 即方程组没有唯一解，使似然最大的参数向量$(\alpha,\beta,\xi,\sigma_\varepsilon^2,\sigma_\delta^2,\sigma_\xi^2)$也没有唯一取值.

继续之前，须认识到这里$X_i,Y_i$的方差与函数关系模型中的方差不同. 在那里，以$\xi_1,\ldots,\xi_n$为条件；这里则对$\xi_i$取边际分布. 例如，函数关系模型中写$\operatorname{Var}X_i=\sigma_\delta^2$，隐含这是给定$\xi_1,\ldots,\xi_n$的条件方差；结构模型中写$\operatorname{Var}X_i=\sigma_\delta^2+\sigma_\xi^2$，隐含这是不以$\xi_1,\ldots,\xi_n$为条件的方差. 不应因此产生混淆.

\eqref{eq:12.2.20}的解对$\hat\beta$施加一个限制，函数关系情形中已经遇到这个限制（\autoref{exercise:12.4}）. 由上述方差与协方差方程，容易推得
$$\begin{aligned}
\hat\sigma_\delta^2\ge0&\quad\text{仅当}\quad S_{xx}\ge\frac1{\hat\beta}S_{xy},\\
\hat\sigma_\varepsilon^2\ge0&\quad\text{仅当}\quad S_{yy}\ge\hat\beta S_{xy},
\end{aligned}$$
二者共同推出
$$\frac{|S_{xy}|}{S_{xx}}\le|\hat\beta|\le\frac{S_{yy}}{|S_{xy}|}.$$
（\autoref{exercise:12.9}建立$\hat\beta$的界.）

现在处理结构关系情形中的可识别性问题.\index{SI0411@可识别性}\semanticindexanchor{411:506}\index{SI0528@变量含误差}\semanticindexanchor{528:506} 这个问题并不意外，因为\eqref{eq:12.2.19}中的参数数目多于确定分布所需的数目. 要使结构线性关系模型可识别，必须作一个假设，把参数数目减少到五个. 幸运的是，对函数关系所作的方差假设也能解决这里的可识别性问题.
\sourcepageanchor{507}
因此，假设$\sigma_\delta^2=\lambda\sigma_\varepsilon^2$，其中$\lambda$已知. 这使未知参数数目降到五个，使模型可识别（见\autoref{exercise:12.8}）. 更强的限制，如假设$\sigma_\delta^2$已知，可能导致方差的MLE取零值. \hyperref[bib:source189]{Kendall和Stuart（1979}，第29章）对此有完整讨论.

一旦假设$\sigma_\delta^2=\lambda\sigma_\varepsilon^2$，本模型中$\hat\alpha,\hat\beta$的最大似然估计便与函数关系模型相同，由\eqref{eq:12.2.16}给出. 但方差估计不同，为
\begin{equation}\label{eq:12.2.21}
\begin{aligned}
\hat\sigma_\delta^2&=\frac1n\left(S_{xx}-\frac{S_{xy}}{\hat\beta}\right),\\
\hat\sigma_\varepsilon^2&=\frac{\hat\sigma_\delta^2}\lambda=\frac1n(S_{yy}-\hat\beta S_{xy}),\\
\hat\sigma_\xi^2&=\frac1n\frac{S_{xy}}{\hat\beta}.
\end{aligned}
\end{equation}
（\autoref{exercise:12.10}证明这些结果，并探讨这里的方差估计与函数模型中的方差估计之间的关系.）注意，与函数关系模型不同，在线性结构关系模型中，这些估计量全都相合\index{SI0409@相合性}\semanticindexanchor{409:507}\index{SI0675@变量含误差}\semanticindexanchor{675:507}（当$\sigma_\delta^2=\lambda\sigma_\varepsilon^2$时）.

\subsubsection{置信集}\label{sec:12.2.4}\index{SI0182@变量含误差}\index{SI0408@置信集}
可以预料，EIV模型中构造置信集是困难的. 完整讨论需要本书尚未建立的工具. 这里仅集中于斜率$\beta$的置信集.

首先，可用\autoref{sec:10.4.1}的近似似然方法构造近似置信区间. 实践中这可能是最常见的做法，也并非全无道理. 但是，这些近似区间不能保持名义$1-\alpha$置信水平. 事实上，\hyperref[bib:source150]{Gleser和Hwang（1987）}的结果给出一个令人不安的结论：任何长度始终有限的斜率区间估计量，其置信系数都等于零！

为明确起见，本节余下部分假定采用EIV模型的结构关系情形. 所给置信集结果对结构与函数两种情形都有效，尤其公式保持相同. 继续假定$\sigma_\delta^2=\lambda\sigma_\varepsilon^2$，其中$\lambda$已知.

\hyperref[bib:source150]{Gleser和Hwang（1987）}指出，参数
$$\tau^2=\frac{\sigma_\xi^2}{\sigma_\delta^2}$$
决定数据中为确定斜率$\beta$而可能提供的信息量. 他们证明，当$\tau^2\to0$时，任何有限长度的$\beta$置信区间，其覆盖概率也必趋于0. 要理解这为何合理，注意$\tau^2=0$意味着$\xi_i$没有变动，从而不可能拟合唯一的直线.

可利用以下事实构造$\beta$的近似置信区间：估计量
$$\hat\sigma_\beta^2=\frac{(1+\lambda\hat\beta^2)^2(S_{xx}S_{yy}-S_{xy}^2)}{(S_{xx}-\lambda S_{yy})^2+4\lambda S_{xy}^2}$$
是$\hat\beta$的真实方差$\sigma_\beta^2$的相合估计量.\footnote{译注：结合下式除以$\sqrt n$的标准误，这里的$\sigma_\beta^2$应理解为$\sqrt n(\hat\beta-\beta)$的渐近方差，而非$\hat\beta$本身的有限样本方差. 原表述保留.} 因而，结合中心极限定理与Slutsky定理（\autoref{sec:5.5}），可证明区间
$$\hat\beta-\frac{z_{\alpha/2}\hat\sigma_\beta}{\sqrt n}\le\beta\le
\hat\beta+\frac{z_{\alpha/2}\hat\sigma_\beta}{\sqrt n}$$
\sourcepageanchor{508}
是$\beta$的近似$1-\alpha$置信区间. 但因为长度有限，它无法对所有参数值保持$1-\alpha$覆盖率.

\hyperref[bib:source145]{Gleser（1987）}考虑了这一区间的一种修正，并将其覆盖概率的下确界报告为$\tau^2$的函数. Gleser的修正$C_G(\hat\beta)$为
\begin{equation}\label{eq:12.2.22}
\hat\beta-\frac{t_{n-2,\alpha/2}\hat\sigma_\beta}{\sqrt{n-2}}\le\beta\le
\hat\beta+\frac{t_{n-2,\alpha/2}\hat\sigma_\beta}{\sqrt{n-2}}.
\end{equation}
再利用中心极限定理与Slutsky定理，可证明这是$\beta$的近似$1-\alpha$置信区间. 因为这个区间的长度也有限，同样无法对所有参数值保持$1-\alpha$覆盖率. Gleser作了若干有限样本数值计算，给出覆盖概率下确界作为$\tau^2$函数的界. 对合理大小的$n$（$\ge10$），若$\tau^2\ge.25$，名义90\%区间的覆盖概率至少为80\%. 随$\tau^2$或$n$增加，表现会改善.

\eqref{eq:12.2.22}的$C_G(\hat\beta)$长度有限，却不保证覆盖概率；现在考察一个精确置信集，按其必须具有的性质，它的长度无穷.\footnote{译注：结合下文可由两个有限区间组成的情形，这里应理解为置信集不能始终具有有限长度，而非每个样本得到的集合都长为无穷. 原表述保留.} 这个集合称为Creasy--Williams置信集，由\hyperref[bib:source87]{Creasy（1956）}及\hyperref[bib:source331]{Williams（1959）}提出，依据以下事实（见\autoref{exercise:12.11}）：若$\sigma_\delta^2=\lambda\sigma_\varepsilon^2$，则
$$\operatorname{Cov}(\beta\lambda Y_i+X_i,Y_i-\beta X_i)=0.$$
定义$r_\lambda(\beta)$为$\beta\lambda Y_i+X_i$与$Y_i-\beta X_i$的样本相关系数，即\footnote{译注：原式第二行分母最后一项为$-\beta^2S_{xx}$；展开第一行第二个平方和可知应为$+\beta^2S_{xx}$. 原式保留.}
\begin{equation}\label{eq:12.2.23}
\begin{aligned}
r_\lambda(\beta)
&=\frac{\sum_{i=1}^n((\beta\lambda y_i+x_i)-(\beta\lambda\bar y+\bar x))((y_i-\beta x_i)-(\bar y-\beta\bar x))}
{\sqrt{\sum_{i=1}^n((\beta\lambda y_i+x_i)-(\beta\lambda\bar y+\bar x))^2\sum_{i=1}^n((y_i-\beta x_i)-(\bar y-\beta\bar x))^2}}\\
&=\frac{\beta\lambda S_{yy}+(1-\beta^2\lambda)S_{xy}-\beta S_{xx}}
{\sqrt{(\beta^2\lambda^2S_{yy}+2\beta\lambda S_{xy}+S_{xx})(S_{yy}-2\beta S_{xy}-\beta^2S_{xx})}}.
\end{aligned}
\end{equation}
因为$\beta\lambda Y_i+X_i$与$Y_i-\beta X_i$服从相关系数为零的二元正态分布，由\autoref{exercise:11.33}有
$$\frac{\sqrt{n-2}\,r_\lambda(\beta)}{\sqrt{1-r_\lambda^2(\beta)}}\sim t_{n-2},$$
对任何$\beta$值都成立. 因而找到一个枢轴量，并得出集合
\begin{equation}\label{eq:12.2.24}\index{SI0186@变量含误差}
\left\{\beta:\frac{(n-2)r_\lambda^2(\beta)}{1-r_\lambda^2(\beta)}\le F_{1,n-2,\alpha}\right\}
\end{equation}
是$\beta$的$1-\alpha$置信集（见\autoref{exercise:12.11}）.

尽管这是$1-\alpha$置信集，它有类似于Fieller区间的缺点. 描述集合\eqref{eq:12.2.24}的函数有两个极小值，函数值均为零. 置信集可以由两个有限的不相交区间、一个有限和两个无限的不相交区间组成，也可能是整个实轴. 例如，\autoref{tab:11.3.1}的数据在$\lambda=1$时，$F$统计量函数的图形见\autoref{fig:12.2.3}. 置信集包含使函数值小于等于$F_{1,22,\alpha}$的所有$\beta$. 当$\alpha=.05$、$F_{1,22,.05}=4.30$时，置信集为$[-1.13,-.14]\cup[.89,7.38]$. 当$\alpha=.01$、$F_{1,22,.01}=7.95$时，置信集为$(-\infty,-18.18]\cup[-1.68,.06]\cup[.60,\infty)$.

此外，对每个$\beta$，$-r_\lambda(\beta)=r_\lambda(-1/(\lambda\beta))$（见\autoref{exercise:12.12}），所以若$\beta$在置信集中，$-1/(\lambda\beta)$也在其中. 这一置信集不能区分$\beta$与$-1/(\lambda\beta)$，始终同时包含正值和负值. 我们永远不能由这个置信集确定斜率的符号！

\eqref{eq:12.2.24}的置信集并非\hyperref[bib:source87]{Creasy（1956）}原来讨论的那个，而是其修正. 她实际关心估计的是$\beta$与$x$轴所成的角$\phi$，即$\beta=\tan(\phi)$，相应置信集的问题较少.
\sourcepageanchor{509}
在EIV模型中，估计$\phi$也许更自然（例如见\hyperref[bib:source4]{Anderson，1976）}，但我们似乎更倾向于估计$\alpha,\beta$.

\begin{figure}[htbp]\centering
\includegraphics[width=.78\textwidth]{assets/chapter12/fig-12.2.3.pdf}
\caption{定义Creasy--Williams置信集的$F$统计量，$\lambda=1$}\label{fig:12.2.3}
\end{figure}
普通线性回归中能进行的其他标准统计分析，大多数在EIV模型中也有对应方法. 例如，可以检验关于$\beta$的假设，或估计$EY_i$的值. 更多内容见\hyperref[bib:source130]{Fuller（1987）}或\hyperref[bib:source189]{Kendall和Stuart（1979}，第29章）.



\subsection{逻辑回归}\label{sec:12.3}
\autoref{sec:11.3.3}的条件正态模型是广义线性模型（GLM）的一个例子. GLM描述响应变量$Y$的均值与自变量$x$之间的关系，但这种关系可能比\eqref{eq:11.3.2}的$EY_i=\alpha+\beta x_i$更复杂. 许多不同模型都可表示为GLM. 本节专门讨论一种GLM：逻辑回归模型.

\subsubsection{模型}\label{sec:12.3.1}
GLM由三部分组成：随机成分、系统成分和连接函数.
\begin{enumerate}[label=(\arabic*)]
\item 响应变量$Y_1,\ldots,Y_n$是随机成分（random component）. 假定它们相互独立，各自的分布都属于指定的指数族. 各$Y_i$不同分布，但都属于同一个分布族，如二项、Poisson、正态等.
\item 系统成分（systematic component）就是模型. 它是预测变量$x_i$的函数，对参数为线性，并与$Y_i$的均值有关. 例如，系统成分可以是$\alpha+\beta x_i$或$\alpha+\beta/x_i$. 这里只考虑$\alpha+\beta x_i$.
\item 最后，连接函数（link function）$g(\mu)$通过规定$g(\mu_i)=\alpha+\beta x_i$联系这两部分，其中$\mu_i=EY_i$.
\end{enumerate}

\sourcepageanchor{510}
\autoref{sec:11.3.3}的条件正态回归模型就是GLM. 在该模型中，各响应$Y_i$都服从正态分布. 正态族当然是指数族，构成随机成分. 回归函数的形式为$\alpha+\beta x_i$，这就是系统成分. 最后，假设关系$\mu_i=EY_i=\alpha+\beta x_i$，意味着连接函数为$g(\mu)=\mu$. 这个简单连接函数称为恒等连接（identity link）.

另一个很有用的GLM是逻辑回归模型. 此时响应$Y_1,\ldots,Y_n$相互独立，$Y_i\sim\operatorname{Bernoulli}(\pi_i)$. （Bernoulli族是指数族.）回顾$EY_i=\pi_i=P(Y_i=1)$. 在这一模型中，假定$\pi_i$与$x_i$满足
\begin{equation}\label{eq:12.3.1}
\log\left(\frac{\pi_i}{1-\pi_i}\right)=\alpha+\beta x_i.
\end{equation}
左边是$Y_i$成功的优势（odds）的对数. 模型假定这一对数优势（log-odds，或logit）是预测变量$x$的线性函数. Bernoulli的pmf可写成指数族形式
$$\pi^y(1-\pi)^{1-y}=(1-\pi)\exp\left\{y\log\left(\frac{\pi_i}{1-\pi_i}\right)\right\}.$$
$\log(\pi/(1-\pi))$是这个指数族的自然参数，\eqref{eq:12.3.1}采用连接函数$g(\pi)=\log(\pi/(1-\pi))$. 以这种方式使用自然参数时，称为典范连接（canonical link）.

\eqref{eq:12.3.1}可改写为
$$\pi_i=\frac{e^{\alpha+\beta x_i}}{1+e^{\alpha+\beta x_i}},$$
或更一般地，
\begin{equation}\label{eq:12.3.2}
\pi(x)=\frac{e^{\alpha+\beta x}}{1+e^{\alpha+\beta x}}.
\end{equation}
可见$0<\pi(x)<1$，由于$\pi(x)$是概率，这似乎合理. 但若某些$x$可能有$\pi(x)=0$或1，这个模型就不适用. 更仔细地考察$\pi(x)$，其导数可写为
\begin{equation}\label{eq:12.3.3}
\frac{d\pi(x)}{dx}=\beta\pi(x)(1-\pi(x)).
\end{equation}
因为$\pi(x)(1-\pi(x))$始终为正，$\pi(x)$的导数随$\beta$为正、零、负而分别为正、零、负. 若$\beta>0$，$\pi(x)$关于$x$严格递增；若$\beta<0$，则严格递减. 若$\beta=0$，对所有$x$有$\pi(x)=e^\alpha/(1+e^\alpha)$. 与简单线性回归一样，$\beta=0$表示$\pi$与$x$之间没有关系. 此外，在逻辑回归模型中，$\pi(-\alpha/\beta)=1/2$. 逻辑回归函数具有如下对称性：对任意$c$，$\pi((-\alpha/\beta)+c)=1-\pi((-\alpha/\beta)-c)$.

参数$\alpha,\beta$的含义与简单线性回归相似. 在\eqref{eq:12.3.1}中令$x=0$，得$\alpha$是$x=0$时成功的对数优势. 分别在$x,x+1$处计算\eqref{eq:12.3.1}，对任意$x$有
$$\log\left(\frac{\pi(x+1)}{1-\pi(x+1)}\right)-\log\left(\frac{\pi(x)}{1-\pi(x)}\right)
=\alpha+\beta(x+1)-\alpha-\beta(x)=\beta.$$
所以$\beta$是$x$增加一个单位所对应的成功对数优势变化量. 简单线性回归中，$\beta$是$x$增加一个单位所对应的$Y$均值变化量. 对等式两边取指数得到
\begin{equation}\label{eq:12.3.4}
e^\beta=\frac{\pi(x+1)/(1-\pi(x+1))}{\pi(x)/(1-\pi(x))}.
\end{equation}
\sourcepageanchor{511}
右边是比较$x+1$与$x$处成功优势的优势比（odds ratio）. （回顾\autoref{ex:5.5.19}和\autoref{ex:5.5.22}讨论过优势的估计.）逻辑回归模型中，这个比值作为$x$的函数为常数. 最后，
\begin{equation}\label{eq:12.3.5}
\frac{\pi(x+1)}{1-\pi(x+1)}=e^\beta\frac{\pi(x)}{1-\pi(x)};
\end{equation}
即$e^\beta$是$x$增加一个单位所对应的成功优势的乘法变化因子.

\eqref{eq:12.3.2}提示了将Bernoulli成功概率$\pi(x)$建模为预测变量$x$的函数的其他方法. 回顾$F(w)=e^w/(1+e^w)$是logistic$(0,1)$分布的cdf. 在\eqref{eq:12.3.2}中假定$\pi(x)=F(\alpha+\beta x)$. 使用其他连续cdf可定义$\pi(x)$的其他模型. 若$F(w)$为标准正态cdf，称为probit回归（probit regression，见\autoref{exercise:12.17}）. 若采用Gumbel的cdf，连接函数称为log-log连接.

\subsubsection{估计}\label{sec:12.3.2}
在线性回归中，采用$Y_i=\alpha+\beta x_i+\varepsilon_i$这样的模型时，可以选择最小二乘方法计算$\alpha,\beta$的估计. 这里已不再如此. 在模型\eqref{eq:12.3.1}中，$Y_i\sim\operatorname{Bernoulli}(\pi_i)$，$Y_i$与$\alpha+\beta x_i$不再直接相连，这也正是需要连接函数的原因. 因而，最小二乘已不再是一种选择.

最常用的估计方法是最大似然. 一般模型为$Y_i\sim\operatorname{Bernoulli}(\pi_i)$，其中$\pi(x)=F(\alpha+\beta x)$. 若记$F_i=F(\alpha+\beta x_i)$，似然函数为
$$L(\alpha,\beta\mid\boldsymbol y)=\prod_{i=1}^n\pi(x_i)^{y_i}(1-\pi(x_i))^{1-y_i}
=\prod_{i=1}^nF_i^{y_i}(1-F_i)^{1-y_i},$$
对数似然为
$$\log L(\alpha,\beta\mid\boldsymbol y)=\sum_{i=1}^n\left\{\log(1-F_i)+y_i\log\left(\frac{F_i}{1-F_i}\right)\right\}.$$
对$\alpha,\beta$求导便得到似然方程. 记$dF(w)/dw=f(w)$，即$F(w)$对应的pdf，再记$f_i=f(\alpha+\beta x_i)$. 则
$$\frac{\partial\log(1-F_i)}{\partial\alpha}=-\frac{f_i}{1-F_i}=-\frac{F_if_i}{F_i(1-F_i)},$$
并且
\begin{equation}\label{eq:12.3.6}
\frac\partial{\partial\alpha}\log\left(\frac{F_i}{1-F_i}\right)=\frac{f_i}{F_i(1-F_i)}.
\end{equation}
因此，
\begin{equation}\label{eq:12.3.7}
\frac\partial{\partial\alpha}\log L(\alpha,\beta\mid\boldsymbol y)
=\sum_{i=1}^n(y_i-F_i)\frac{f_i}{F_i(1-F_i)}.
\end{equation}
类似计算得到
\begin{equation}\label{eq:12.3.8}
\frac\partial{\partial\beta}\log L(\alpha,\beta\mid\boldsymbol y)
=\sum_{i=1}^n(y_i-F_i)\frac{f_i}{F_i(1-F_i)}x_i.
\end{equation}
逻辑回归中$F(w)=e^w/(1+e^w)$，$f_i/[F_i(1-F_i)]=1$，所以\eqref{eq:12.3.7}与\eqref{eq:12.3.8}有所简化.

\sourcepageanchor{512}
令\eqref{eq:12.3.7}、\eqref{eq:12.3.8}等于零，并对$\alpha,\beta$求解，得到MLE. 这些是$\alpha,\beta$的非线性方程，必须用数值方法求解，后面将讨论. 对逻辑和probit回归，对数似然严格凹. 因而，若似然方程有解，解就唯一，并且是MLE. 但对某些极端数据，似然方程无解. 似然最大值在参数趋于$\pm\infty$的某个极限中取得，例子见\autoref{exercise:12.16}. 原因是逻辑模型假定$0<\pi(x)<1$，但对某些数据集，逻辑似然的最大值出现在$\pi(x)=0$或1的极限中. 若逻辑模型成立，获得这类数据的概率趋于零.

\needspace{12\baselineskip}
\begin{example}[挑战者号数据]\label{ex:12.3.1}
航天飞机O形密封圈失效的数据集如今已因悲剧而广为人知，这种失效与温度有联系. \autoref{tab:12.3.1}给出起飞时的温度，以及O形圈是否失效.

使用$F(\alpha+\beta x_i)=e^{\alpha+\beta x_i}/(1+e^{\alpha+\beta x_i})$求解\eqref{eq:12.3.6}和\eqref{eq:12.3.7}，得到MLE $\hat\alpha=15.81$、$\hat\beta=-.243$.\footnote{译注：这里实际应解\eqref{eq:12.3.7}、\eqref{eq:12.3.8}；原引用照录.} \autoref{fig:12.3.1}给出拟合曲线和数据.

挑战者号航天飞机在起飞时爆炸，机上七名宇航员遇难. 爆炸由O形圈失效导致，据信与发射时异常寒冷的天气（31$^\circ\mathrm F$）有关. 31度时O形圈失效概率的MLE为.9997. （完整叙述见\hyperref[bib:source91]{\hyperref[bib:source91]{Dalal等，1989}}.）
\end{example}
\begin{table}[htbp]\centering
\caption{飞行时温度（华氏度）及O形圈失效情况（1为失效，0为成功）}\label{tab:12.3.1}
\setlength{\tabcolsep}{3.5pt}
\begin{tabular}{lrrrrrrrrrrrr}\toprule
飞行编号&14&9&23&10&1&5&13&15&4&3&8&17\\
失效&1&1&1&1&0&0&0&0&0&0&0&0\\
温度&53&57&58&63&66&67&67&68&69&70&70&70\\\midrule
飞行编号&2&11&6&7&16&21&19&22&12&20&18&\\
失效&1&1&0&0&0&1&0&0&0&0&0&\\
温度&70&70&72&73&75&75&76&76&78&79&81&\\\bottomrule
\end{tabular}
\end{table}
\begin{figure}[htbp]\centering
\includegraphics[width=.92\textwidth]{assets/chapter12/fig-12.3.1.pdf}
\caption{\autoref{tab:12.3.1}的数据及拟合的逻辑曲线}\label{fig:12.3.1}
\end{figure}
\sourcepageanchor{513}
到目前为止，假定在每个$x_i$值只观测一次Bernoulli试验的结果. 虽然经常如此，但很多情形会在每个$x$值处观测多个Bernoulli结果. 现在重新考察这一更一般情形的似然解.

设数据集中预测变量$x$有$J$个不同取值$x_1,\ldots,x_J$. 用$n_j$表示$x_j$处的Bernoulli观测数，$Y_j^*$表示这些$n_j$次观测中的成功数. 因而$Y_j^*\sim\operatorname{binomial}(n_j,\pi(x_j))$. 似然为\footnote{译注：原式第一项使用$x_i$而乘积下标为$j$，按定义应为$x_j$；还省略了不依赖参数的二项系数，故应理解为与似然成比例. 原式保留.}
$$L(\alpha,\beta\mid\boldsymbol y^*)=\prod_{j=1}^J\pi(x_i)^{y_j^*}(1-\pi(x_i))^{n_j-y_j^*}
=\prod_{j=1}^JF_j^{y_j^*}(1-F_j)^{n_j-y_j^*},$$
似然方程为
$$\begin{aligned}
0&=\sum_{j=1}^J(y_j^*-n_jF_j)\frac{f_j}{F_j(1-F_j)},\\
0&=\sum_{j=1}^J(y_j^*-n_jF_j)\frac{f_j}{F_j(1-F_j)}x_j.
\end{aligned}$$
既然用最大似然估计了逻辑回归参数，接下来就可利用MLE的渐近理论得到近似方差. 但必须以更一般的方式处理. \autoref{sec:10.1.3}曾利用信息数近似MLE的方差. 这里采用同一策略，但由于有两个参数，需要由以下$2\times2$矩阵给出的信息矩阵（information matrix）：
\begin{equation}\label{eq:12.3.9}
I(\theta_1,\theta_2)=\begin{pmatrix}
-\dfrac{\partial^2}{\partial\theta_1^2}\log L(\theta_1,\theta_2\mid\boldsymbol y)&
-\dfrac{\partial^2}{\partial\theta_1\partial\theta_2}\log L(\theta_1,\theta_2\mid\boldsymbol y)\\
-\dfrac{\partial^2}{\partial\theta_1\partial\theta_2}\log L(\theta_1,\theta_2\mid\boldsymbol y)&
-\dfrac{\partial^2}{\partial\theta_2^2}\log L(\theta_1,\theta_2\mid\boldsymbol y)
\end{pmatrix}.
\end{equation}
逻辑回归的信息矩阵为
\begin{equation}\label{eq:12.3.10}
I(\alpha,\beta)=\begin{pmatrix}
\sum_{j=1}^Jn_jF_j(1-F_j)&\sum_{j=1}^Jx_jn_jF_j(1-F_j)\\
\sum_{j=1}^Jx_jn_jF_j(1-F_j)&\sum_{j=1}^Jx_j^2n_jF_j(1-F_j)
\end{pmatrix},
\end{equation}
通常利用这一矩阵近似MLE $\hat\alpha,\hat\beta$的方差. 注意$I(\alpha,\beta)$的元素不依赖$Y_1^*,\ldots,Y_J^*$. 因而，在这种情形下，观测信息与信息相同.

\autoref{sec:10.1.3}使用近似\eqref{eq:10.1.7}，即$\operatorname{Var}(h(\hat\theta)\mid\theta)\approx[h'(\hat\theta)]^2/I(\hat\theta)$，其中$I(\cdot)$为信息数. 对信息矩阵不能完全照搬，而要先求逆，再用逆矩阵元素近似方差. 回顾$2\times2$矩阵的逆为
$$\begin{pmatrix}a&b\\c&d\end{pmatrix}^{-1}
=\frac1{ad-bc}\begin{pmatrix}d&-b\\-c&a\end{pmatrix}.$$
为得到近似方差，用MLE估计矩阵\eqref{eq:12.3.10}中的参数；方差估计$[\operatorname{se}(\hat\alpha)]^2$、$[\operatorname{se}(\hat\beta)]^2$就是$I(\hat\alpha,\hat\beta)$的逆矩阵的对角元素.

\enlargethispage{1pt}
\begin{example}[挑战者号数据续]\label{ex:12.3.2}
挑战者号数据所得估计量的估计信息矩阵为
$$\begin{aligned}
I(\hat\alpha,\hat\beta)&=\begin{pmatrix}
\sum_{j=1}^J\hat F_j(1-\hat F_j)&\sum_{j=1}^Jx_j\hat F_j(1-\hat F_j)\\
\sum_{j=1}^Jx_j\hat F_j(1-\hat F_j)&\sum_{j=1}^Jx_j^2\hat F_j(1-\hat F_j)
\end{pmatrix}\\
&=\begin{pmatrix}3.15&214.75\\214.75&14728.5\end{pmatrix},
\end{aligned}$$
\sourcepageanchor{514}
其中$\hat F_j=e^{\hat\alpha+\hat\beta x_j}/(1+e^{\hat\alpha+\hat\beta x_j})$，其逆为
$$I(\hat\alpha,\hat\beta)^{-1}=\begin{pmatrix}59.19&-.86\\-.86&.013\end{pmatrix}.$$
例如，似然渐近理论说明，大样本时$\hat\beta\pm z_{\alpha/2}\operatorname{se}(\hat\beta)$是$\beta$的近似$100(1-\alpha)\%$置信区间. 因而，挑战者号数据给出的95\%置信区间为
$$\beta\in-.243\pm1.96\times\sqrt{.013}\quad\Rightarrow\quad-.466\le\beta\le-.02,$$
支持$\beta<0$的结论.
\end{example}

这一模型中最常见的检验也许是假设$H_0:\beta=0$，因为和简单线性回归一样，这一假设表示预测变量与响应变量没有关系. 若$H_0$成立且样本量很大，Wald检验统计量$Z=\hat\beta/\operatorname{se}(\hat\beta)$近似服从标准正态分布. 因而，若$|Z|\ge z_{\alpha/2}$，可拒绝$H_0$. 也可用对数似然比检验统计量检验$H_0$：\footnote{译注：括号中第二项按对数似然比定义应为$\log L(\hat\alpha_0,0\mid\boldsymbol y^*)$，原文漏log，照录.}
$$-2\log\lambda(\boldsymbol y^*)=2[\log L(\hat\alpha,\hat\beta\mid\boldsymbol y^*)-L(\hat\alpha_0,0\mid\boldsymbol y^*)],$$
其中$\hat\alpha_0$是$\beta=0$时$\alpha$的MLE. 利用通常的二项分布论证（\autoref{exercise:12.20}），可证明$\hat\alpha_0=\sum_{i=1}^ny_i/n=\sum_{j=1}^Jy_j^*/\sum_{j=1}^Jn_j$.\footnote{译注：右边估计的是共同成功概率$p$，而截距满足$\alpha_0=\log(p/(1-p))$，所以应对所给比例取logit. 原式保留.} 因此，在$H_0$下，$-2\log\lambda$近似服从$\chi_1^2$分布，若$-2\log\lambda\ge\chi_{1,\alpha}^2$，可在水平$\alpha$拒绝$H_0$.

这里只介绍了最简单的逻辑回归与广义线性模型. 更多内容可见\hyperref[bib:source1]{Agresti（1990）}等标准教材.



\subsection{稳健回归}\label{sec:12.4}
如\autoref{sec:10.2}，现在考察基础模型不正确时所用方法的表现，并研究最小二乘估计的一些稳健替代方法，先作一个类似于均值与中位数的比较.

回顾观测$x_1,x_2,\ldots,x_n$时，可将均值和中位数定义为以下量的最小化点：
$$\text{均值：}\min_m\left\{\sum_{i=1}^n(x_i-m)^2\right\},\qquad
\text{中位数：}\min_m\left\{\sum_{i=1}^n|x_i-m|\right\}.$$
对于简单线性回归，观测$(y_1,x_1),\allowbreak(y_2,x_2),\allowbreak\ldots,\allowbreak(y_n,x_n)$，最小二乘回归估计满足
$$\text{最小二乘：}\min_{a,b}\left\{\sum_{i=1}^n[y_i-(a+bx_i)]^2\right\},$$
类似地，定义最小绝对偏差（least absolute deviation，LAD）回归估计为\index{SI0894@最小绝对偏差}\semanticindexanchor{894:514}
$$\text{最小绝对偏差：}\min_{a,b}\left\{\sum_{i=1}^n|y_i-(a+bx_i)|\right\}.$$
（LAD估计可能不唯一，见\autoref{exercise:12.25}.）

可见，最小二乘估计量是样本均值在回归中的对应物. 因此，应当考察它们在\autoref{sec:10.2}第\hyperref[item:10.2.huber.1]{(1)}至\hyperref[item:10.2.huber.3]{(3)}项所列意义下的稳健性表现.

\sourcepageanchor{515}
\begin{example}[最小二乘估计的稳健性]\label{ex:12.4.1}
若观测$(y_1,x_1),\allowbreak(y_2,x_2),\allowbreak\ldots,\allowbreak(y_n,x_n)$，其中
$$Y_i=\alpha+\beta x_i+\varepsilon_i,$$
各$\varepsilon_i$互不相关，$E\varepsilon_i=0$、$\operatorname{Var}\varepsilon_i=\sigma^2$，则已知最小二乘估计量$b$的方差为$\sigma^2/\sum(x_i-\bar x)^2$，是$\beta$的BLUE，满足\autoref{sec:10.2}的第\hyperref[item:10.2.huber.1]{(1)}项.

为考察小偏离下$b$的表现，假定
$$\operatorname{Var}(\varepsilon_i)=\begin{cases}\sigma^2,&\text{概率为 }1-\delta,\\\tau^2,&\text{概率为 }\delta.\end{cases}$$
写$b=\sum d_iY_i$，其中$d_i=(x_i-\bar x)/\sum(x_i-\bar x)^2$，现在有
$$\operatorname{Var}(b)=\sum_{i=1}^nd_i^2\operatorname{Var}(\varepsilon_i)
=\frac{(1-\delta)\sigma^2+\delta\tau^2}{\sum_{i=1}^n(x_i-\bar x)^2}.$$
这说明，和样本均值一样，小扰动下$b$表现相当好. （当然，例如用Cauchy的pdf污染，仍可使情况失控.）最小二乘截距$a$的表现类似（\autoref{exercise:12.22}）. 偏倚污染的影响另见\autoref{exercise:12.23}.
\end{example}

接下来考察一个“灾难性”观测的影响，并将最小二乘与对应于中位数的替代方法——最小绝对偏差回归——比较.

\begin{example}[灾难性观测]\label{ex:12.4.2}
\hyperref[bib:source217]{McPherson（1990）}描述了一项实验，测量24只长鼻袋鼠（potoroo，一种有袋动物）育儿袋中的二氧化碳（$\mathrm{CO}_2$）和氧气（$\mathrm O_2$）浓度. 关心$\mathrm{CO}_2$对$\mathrm O_2$的回归，实验者预期斜率为$-1$. 23只动物的数据见\autoref{tab:12.4.2}，另一只有缺失值. 对原始数据，最小二乘与LAD直线很接近：
$$\begin{aligned}\text{最小二乘}&:\quad y=18.67-0.89x,\\
\text{最小绝对偏差}&:\quad y=18.59-0.89x.
\end{aligned}$$
然而，异常观测可能扰乱最小二乘. 输入数据时，第15只动物的$\mathrm O_2$值18被错录为10，我们确实这样做了. 对这一新的错误数据集，有
$$\begin{aligned}\text{最小二乘}&:\quad y=6.41-0.23x,\\
\text{最小绝对偏差}&:\quad y=15.95-0.75x,
\end{aligned}$$
说明异常观测对LAD的影响小得多. 回归直线见\autoref{fig:12.4.1}.

这些计算说明，与最小二乘相比，LAD具有抵抗性. 由于均值与中位数的类比，可以推测这种表现反映在崩溃值中：最小二乘为0\%，LAD为50\%.\footnote{译注：回归LAD的崩溃值不能一般地照搬中位数的50\%；若允许污染预测变量$x$，单个高杠杆点也可使LAD斜率失控. 上述数值须有额外污染范围和设计条件，原文未列明，照录.}
\end{example}

不过，均值与中位数的类比仍在继续. 虽然LAD估计量能稳健地对待灾难性观测，但相对于最小二乘估计量，它损失了很多效率（另见\autoref{exercise:12.25}）.

\begin{example}[LAD估计量的渐近正态性]\label{ex:12.4.3}
调整推导\eqref{eq:10.2.6}的论证，求LAD估计量的渐近分布. 为简化，仅考虑模型
$$Y_i=\beta x_i+\varepsilon_i,$$
即令$\alpha=0$，从而避免处理二元极限分布.

\sourcepageanchor{516}
\begin{table}[htbp]\centering
\setcounter{table}{1}
\caption{23只长鼻袋鼠育儿袋中的$\mathrm{CO}_2$与$\mathrm O_2$值（\hyperref[bib:source217]{McPherson，1990）}}\label{tab:12.4.2}
\begin{tabular}{lrrrrrrrr}\toprule
动物&1&2&3&4&5&6&7&8\\\midrule
$\%\,\mathrm O_2$&20&19.6&19.6&19.4&18.4&19&19&18.3\\
$\%\,\mathrm{CO}_2$&1&1.2&1.1&1.4&2.3&1.7&1.7&2.4\\\midrule
动物&9&10&11&12&13&14&15&16\\\midrule
$\%\,\mathrm O_2$&18.2&18.6&19.2&18.2&18.7&18.5&18&17.4\\
$\%\,\mathrm{CO}_2$&2.1&2.1&1.2&2.3&1.9&2.4&2.6&2.9\\\midrule
动物&17&18&19&20&21&22&23&\\\midrule
$\%\,\mathrm O_2$&16.5&17.2&17.3&17.8&17.3&18.4&16.9&\\
$\%\,\mathrm{CO}_2$&4.0&3.3&3.0&3.4&2.9&1.9&3.9&\\\bottomrule
\end{tabular}
\end{table}
\begin{figure}[htbp]\centering
\includegraphics[width=.75\textwidth]{assets/chapter12/fig-12.4.1.pdf}
\caption{\autoref{tab:12.4.2}数据的最小二乘、LAD及M估计拟合，分别使用原始数据与将$(18,2.6)$误录为$(10,2.6)$的数据. LAD与M估计直线十分相似，而最小二乘直线会对数据变化作出反应.}\label{fig:12.4.1}
\end{figure}
用M估计量的术语说，LAD估计量由最小化下式得到：
\begin{equation}\label{eq:12.4.1}
\begin{aligned}
\sum_{i=1}^n\rho(y_i-\beta x_i)&=\sum_{i=1}^n|y_i-\beta x_i|\\
&=\sum_{i=1}^n\bigl[(y_i-\beta x_i)I(y_i>\beta x_i)-(y_i-\beta x_i)I(y_i<\beta x_i)\bigr].
\end{aligned}
\end{equation}
然后计算$\psi=\rho'$，求解$\sum_i\psi(y_i-\beta x_i)=0$得到$\beta$，其中\footnote{译注：原文$\psi=\rho'$与随后含$x_i$的式子并非同一变量下的导数；后式实际是目标函数对$\beta$的负导数所对应的加权符号函数，并在残差为零处不可微. 原表述保留.}
$$\psi(y_i-\beta x_i)=x_iI(y_i>\beta x_i)-x_iI(y_i<\beta x_i).$$

\sourcepageanchor{517}
若$\hat\beta_L$是解，在$\beta$附近将$\psi$展开为Taylor级数：
$$\begin{aligned}
\sum_{i=1}^n\psi(y_i-\hat\beta_Lx_i)
&=\sum_{i=1}^n\psi(y_i-\beta x_i)\\
&\quad+(\hat\beta_L-\beta)\left.\frac d{d\hat\beta_L}\sum_{i=1}^n\psi(y_i-\hat\beta_Lx_i)\right|_{\hat\beta_L=\beta}+\cdots.
\end{aligned}$$
虽然等式左边不等于零，假定它在$n\to\infty$时趋于零（见\autoref{exercise:12.27}）. 重新整理得
\begin{equation}\label{eq:12.4.2}
\sqrt n(\hat\beta_L-\beta)=
\frac{-\dfrac1{\sqrt n}\sum_{i=1}^n\psi(y_i-\beta x_i)}
{\left.\dfrac1n\dfrac d{d\hat\beta_L}\sum_{i=1}^n\psi(y_i-\hat\beta_Lx_i)\right|_{\hat\beta_L=\beta}}.
\end{equation}
先看分子. 因为$E_\beta\psi(Y_i-\hat\beta_Lx_i)=0$且$\operatorname{Var}\psi(Y_i-\hat\beta_Lx_i)=x_i^2$，有\footnote{译注：这些均值、方差应在真实参数$\beta$处计算，并需误差中位数为零等条件；原文此处写估计量$\hat\beta_L$，下一式又写$\hat\beta$，与\eqref{eq:12.4.2}分子的$\beta$不一致. 原符号均保留.}
\begin{equation}\label{eq:12.4.3}
-\frac1{\sqrt n}\sum_{i=1}^n\psi(Y_i-\hat\beta x_i)
=\sqrt n\left[-\frac1n\sum_{i=1}^n\psi(Y_i-\hat\beta x_i)\right]
\to\mathrm n(0,\sigma_x^2),
\end{equation}
其中$\sigma_x^2=\lim_{n\to\infty}\frac1n\sum_{i=1}^nx_i^2$. 转向分母，因为$\psi$有不可微点，必须小心. 因此，先应用大数定律，再求导，并采用近似
\begin{equation}\label{eq:12.4.4}
\begin{aligned}
\frac1n\frac d{d\beta_0}\sum_{i=1}^n\psi(y_i-\beta_0x_i)
&\approx\frac1n\sum_{i=1}^n\frac d{d\beta_0}E_\beta[\psi(Y_i-\beta_0x_i)]\\
&=\frac1n\sum_{i=1}^n\frac d{d\beta_0}\bigl[x_iP_\beta(Y_i>\beta_0x_i)-x_iP_\beta(Y_i<\beta_0x_i)\bigr]\\
&=\frac1n\sum_{i=1}^n\bigl[x_i^2f(\beta_0x_i-\beta x_i)+x_i^2f(\beta_0x_i-\beta x_i)\bigr].
\end{aligned}
\end{equation}
在$\beta_0=\beta$处计算导数，得\footnote{译注：按上文$\psi$的符号约定，对两个概率求导后均应带负号，故\eqref{eq:12.4.4}最后一行及此近似右端应为负；最终方差使用平方，不受这一符号影响. 原式保留.}
$$\left.\frac1n\frac d{d\beta_0}\sum_{i=1}^n\psi(y_i-\beta_0x_i)\right|_{\beta_0=\beta}
\approx2f(0)\frac1n\sum_{i=1}^nx_i^2,$$
结合\eqref{eq:12.4.2}、\eqref{eq:12.4.3}，有
\begin{equation}\label{eq:12.4.5}
\sqrt n(\hat\beta_L-\beta)\to\mathrm n\left(0,\frac1{4f(0)^2\sigma_x^2}\right).
\end{equation}
最后，$\alpha=0$时最小二乘估计量为$\hat\beta=\sum_{i=1}^nx_iy_i/\sum_{i=1}^nx_i^2$，满足
$$\sqrt n(\hat\beta-\beta)\to\mathrm n\left(0,\frac1{\sigma_x^2}\right),$$
所以$\hat\beta_L$相对于$\hat\beta$的渐近相对效率为
$$\operatorname{ARE}(\hat\beta_L,\hat\beta)=\frac{1/\sigma_x^2}{1/(4f(0)^2\sigma_x^2)}=4f(0)^2,$$
其值与\autoref{tab:10.2.2}中中位数相对于均值的值相同. 因此，对正态误差，LAD估计量相对于最小二乘的ARE只有64\%，说明LAD牺牲了相当多的效率.\footnote{译注：这里的最小二乘极限方差及ARE采用误差方差为1的标准化；一般误差方差$\sigma^2$下，应分别为$\sigma^2/\sigma_x^2$与$4\sigma^2f(0)^2$. 原式未明示此标准化，照录.}
\end{example}
\sourcepageanchor{518}
于是又遇到\autoref{sec:10.2.1}的情形：若误差确为正态，以LAD代替最小二乘似乎损失了过多效率. 折中方法仍然是M估计量. 可以最小化类似\eqref{eq:10.2.2}的函数$\sum_i\rho_i(\alpha,\beta)$来构造，其中
\begin{equation}\label{eq:12.4.6}
\rho_i(\alpha,\beta)=\begin{cases}
\tfrac12(y_i-\alpha-\beta x_i)^2,&|y_i-\alpha-\beta x_i|\le k,\\
k|y_i-\alpha-\beta x_i|-\tfrac12k^2,&|y_i-\alpha-\beta x_i|\ge k,
\end{cases}
\end{equation}
$k$为调节参数.

\begin{example}[回归M估计量]\label{ex:12.4.4}
使用\eqref{eq:12.4.6}，取$k=1.5\sigma$，对\autoref{tab:12.4.2}的数据拟合$\alpha,\beta$的M估计量，结果为
$$\begin{aligned}\text{原始数据的M估计}&:\quad y=18.5-0.89x,\\
\text{误录数据的M估计}&:\quad y=14.67-0.68x,
\end{aligned}$$
其中以最小二乘拟合残差的标准差.23估计$\sigma$.

可见，M估计比最小二乘直线更有抵抗性，有离群值时的表现更像LAD拟合.
\end{example}

如\autoref{sec:10.2}，预期M估计量的ARE优于LAD. 确实如此，但计算变得十分复杂，比LAD还复杂，因此不在这里给出细节. \hyperref[bib:source171]{Huber（1981}，第7章）详细讨论了M估计量的渐近理论，另见\hyperref[bib:source243]{Portnoy（1987）}. 这里满足于通过一个小型模拟研究评价M估计量，重现类似\autoref{tab:10.2.4}的表.

\begin{example}[回归ARE的模拟]\label{ex:12.4.5}
对于模型$Y_i=\alpha+\beta x_i+\varepsilon_i$，$i=1,2,\ldots,5$，取$x_i$为$(-2,-1,0,1,2)$，$\alpha=0$、$\beta=1$. 从正态、双指数和Laplace分布生成$\varepsilon_i$，计算最小二乘、LAD和M估计量的方差，列于\autoref{tab:12.4.3}.\footnote{译注：双指数分布与Laplace分布通常是同一分布；原文列举与表中“正态、Logistic、双指数”不一致，可能将Logistic误写为Laplace. 原文照录.}
\begin{table}[htbp]\centering
\caption{回归M估计量的渐近相对效率，$k=1.5$（基于10,000次模拟）}\label{tab:12.4.3}
\begin{tabular}{lrrr}\toprule
&正态&Logistic&双指数\\\midrule
相对于最小二乘&0.98&1.03&1.07\\
相对于LAD&1.39&1.27&1.14\\\bottomrule
\end{tabular}
\end{table}
三种分布下，M估计量的方差都与最小二乘相近，并且都优于LAD. M估计量相对于LAD的优势，比\autoref{tab:10.2.4}中Huber估计量相对于中位数的优势更显著.
\end{example}



\subsection{习题}\label{sec:12.5}
\begin{exercise}\label{exercise:12.1}
验证\eqref{eq:12.2.7}的表达式. （提示：使用勾股定理.）
\end{exercise}

\begin{exercise}\label{exercise:12.2}
证明函数
$$f(b)=\frac1{1+b^2}\bigl[S_{yy}-2bS_{xy}+b^2S_{xx}\bigr]$$
的极值点为
$$b=\frac{-(S_{xx}-S_{yy})\pm\sqrt{(S_{xx}-S_{yy})^2+4S_{xy}^2}}{2S_{xy}}.$$
证明“$+$”解给出$f(b)$的最小值.
\end{exercise}

\sourcepageanchor{519}
\begin{center}\begin{tabular}{lrr}\toprule
物种&体重（kg）&脑重（g）\\&$(x)$&$(y)$\\\midrule
北极狐&3.385&44.50\\夜猴&.480&15.50\\山河狸&1.350&8.10\\豚鼠&1.040&5.50\\
毛丝鼠&.425&6.40\\地松鼠&.101&4.00\\树蹄兔&2.000&12.30\\大棕蝠&.023&.30\\\bottomrule
\end{tabular}\end{center}

\begin{exercise}\label{exercise:12.3}
最大化似然\eqref{eq:12.2.13}时，对每个$\alpha,\beta,\sigma_\delta^2$值，先关于$\xi_1,\ldots,\xi_n$最小化函数
$$f(\xi_1,\ldots,\xi_n)=\sum_{i=1}^n\bigl((x_i-\xi_i)^2+\lambda(y_i-(\alpha+\beta\xi_i))^2\bigr).$$
\begin{enumerate}[label=(\alph*)]
\item\label{item:12.3.a} 证明这一函数在
$$\xi_i^*=\frac{x_i+\lambda\beta(y_i-\alpha)}{1+\lambda\beta^2}$$
处取得最小值.
\item\label{item:12.3.b}\index{SI0687@度量} 证明函数
$$D_\lambda((x,y),(\xi,\alpha+\beta\xi))=(x-\xi)^2+\lambda(y-(\alpha+\beta\xi))^2$$
定义了点$(x,y)$与$(\xi,\alpha+\beta\xi)$之间的度量（metric）. 度量是一种距离度量，即测量两点$A,B$间距离的函数$D$. 度量满足以下四个性质：\footnote{译注：所给$D_\lambda$为加权距离的平方，一般不满足三角不等式；例如同一直线上间距为1的三点有$4>1+1$. 对$D_\lambda$取平方根才得到度量. 第\hyperref[item:12.3.b.iii]{(iii)}项原称“reflexive”，实际为对称性. 原题保留.}
\begin{enumerate}[label=\roman*.]
\item $D(A,A)=0$.
\item 若$A\ne B$，则$D(A,B)>0$.
\item\phantomsection\label{item:12.3.b.iii} $D(A,B)=D(B,A)$（自反性）.
\item $D(A,B)\le D(A,C)+D(C,B)$（三角不等式）.
\end{enumerate}
\end{enumerate}
\end{exercise}

\begin{exercise}\label{exercise:12.4}\index{SI0414@与最小二乘的关系}\index{SI0611@与最大似然估计量的关系}
考虑EIV模型中斜率的MLE
$$\hat\beta(\lambda)=\frac{-(S_{xx}-\lambda S_{yy})+\sqrt{(S_{xx}-\lambda S_{yy})^2+4\lambda S_{xy}^2}}{2\lambda S_{xy}},$$
其中假定$\lambda=\sigma_\delta^2/\sigma_\varepsilon^2$已知.
\begin{enumerate}[label=(\alph*)]
\item\label{item:12.4.a} 证明$\lim_{\lambda\to0}\hat\beta(\lambda)=S_{xy}/S_{xx}$，即$y$对$x$普通回归的斜率.
\item\label{item:12.4.b} 证明$\lim_{\lambda\to\infty}\hat\beta(\lambda)=S_{yy}/S_{xy}$，即$x$对$y$普通回归斜率的倒数.
\item\label{item:12.4.c} 证明$\hat\beta(\lambda)$关于$\lambda$单调；当$S_{xy}>0$时递增，当$S_{xy}<0$时递减.
\item\label{item:12.4.d} 证明正交最小二乘直线（$\lambda=1$）总在$y$对$x$及$x$对$y$的普通回归直线之间.
\item\label{item:12.4.e} 上表数据来自一项研究，旨在考察若干动物物种的脑重与体重的关系. 假设EIV模型，计算斜率的MLE. 再计算$y$对$x$及$x$对$y$回归的最小二乘斜率，说明这些量如何给出MLE的界.
\end{enumerate}
\end{exercise}

\begin{exercise}\label{exercise:12.5}
在EIV函数关系模型中，假定$\lambda=\sigma_\delta^2/\sigma_\varepsilon^2$已知，证明$\sigma_\delta^2$的MLE由\eqref{eq:12.2.18}给出.
\end{exercise}

\sourcepageanchor{520}
\begin{exercise}\label{exercise:12.6}
证明在线性结构关系模型\eqref{eq:12.2.6}中，积分消去$\xi_i$后，$(X_i,Y_i)$的边际分布由\eqref{eq:12.2.19}给出.
\end{exercise}

\needspace{6\baselineskip}
\begin{exercise}\label{exercise:12.7}
考虑线性结构关系模型，假定$\xi_i$具有非正常分布$\xi_i\sim\operatorname{uniform}(-\infty,\infty)$.
\begin{enumerate}[label=(\alph*)]
\item\label{item:12.7.a} 证明对每个$i$，
$$\begin{aligned}
&\int_{-\infty}^\infty\frac1{2\pi}\frac1{\sigma_\delta\sigma_\varepsilon}
\exp\left[-\frac{(x_i-\xi_i)^2}{2\sigma_\delta^2}\right]
\exp\left[-\frac{(y_i-(\alpha+\beta\xi_i))^2}{2\sigma_\varepsilon^2}\right]d\xi_i\\
&=\frac1{\sqrt{2\pi}}\frac1{\sqrt{\beta^2\sigma_\delta^2+\sigma_\varepsilon^2}}
\exp\left[-\frac12\frac{(y_i-(\alpha+\beta x_i))^2}{\beta^2\sigma_\delta^2+\sigma_\varepsilon^2}\right].
\end{aligned}$$
（在指数中配方便容易积分.）
\item\label{item:12.7.b} \hyperref[item:12.7.a]{(a)}的积分结果看起来像pdf；若把它当作给定$X$时$Y$的pdf，似乎就有$X,Y$之间的线性关系. 因此，有时称结构关系的这一“极限情形”导致简单线性回归和普通最小二乘. 解释为何对上述函数的这一解释是错误的.
\end{enumerate}
\end{exercise}

\begin{exercise}\label{exercise:12.8}
用以下方式验证结构关系模型的不可识别性问题.
\begin{enumerate}[label=(\alph*)]
\item\label{item:12.8.a} 给出两组不同参数，使$(X_i,Y_i)$具有相同的边际分布.
\item\label{item:12.8.b} 证明至少存在两个不同的参数向量，对\eqref{eq:12.2.20}的方程给出相同的解.
\end{enumerate}
\end{exercise}

\begin{exercise}\label{exercise:12.9}
结构关系模型中，\eqref{eq:12.2.20}的解对$\hat\beta$施加一个限制，与函数关系情形的限制相同（\autoref{exercise:12.4}）.
\begin{enumerate}[label=(\alph*)]
\item\label{item:12.9.a} 证明\eqref{eq:12.2.20}中，$\sigma_\delta^2$的MLE非负仅当$S_{xx}\ge(1/\hat\beta)S_{xy}$；$\sigma_\varepsilon^2$的MLE非负仅当$S_{yy}\ge\hat\beta S_{xy}$.
\item\label{item:12.9.b} 证明\hyperref[item:12.9.a]{(a)}的限制与\eqref{eq:12.2.20}其余方程共同推出
$$\frac{|S_{xy}|}{S_{xx}}\le|\hat\beta|\le\frac{S_{yy}}{|S_{xy}|}.$$
\end{enumerate}
\end{exercise}

\begin{exercise}\label{exercise:12.10}\leavevmode
\begin{enumerate}[label=(\alph*)]
\item\label{item:12.10.a} 假设$\sigma_\delta^2=\lambda\sigma_\varepsilon^2$，解\eqref{eq:12.2.20}，推导结构关系模型中$(\alpha,\beta,\sigma_\varepsilon^2,\sigma_\delta^2,\sigma_\xi^2)$的MLE.
\item\label{item:12.10.b} 假设结构关系模型成立且$\sigma_\delta^2=\lambda\sigma_\varepsilon^2$，计算\autoref{exercise:12.4}数据的$(\alpha,\beta,\sigma_\varepsilon^2,\sigma_\delta^2,\sigma_\xi^2)$的MLE.
\item\label{item:12.10.c} 验证函数关系模型与结构关系模型中方差估计的关系. 特别地，证明
$$\widehat{\operatorname{Var}}X_i(\text{结构})=2\widehat{\operatorname{Var}}X_i(\text{函数}).$$
即验证\footnote{译注：后一等式比较的是两模型中的测量误差方差$\sigma_\delta^2$估计；结构模型的无条件$\operatorname{Var}X_i=\sigma_\delta^2+\sigma_\xi^2$，故前一式的方差名称易造成混淆. 原题保留.}
$$\left(S_{xx}-\frac{S_{xy}}{\hat\beta}\right)=\frac\lambda{1+\lambda\hat\beta^2}\sum_{i=1}^n\bigl(y_i-(\hat\alpha+\hat\beta x_i)\bigr)^2.$$
\item\label{item:12.10.d} 验证\eqref{eq:12.2.21}的MLE方差估计中隐含的等式：
$$S_{xx}-\frac{S_{xy}}{\hat\beta}=\lambda(S_{yy}-\hat\beta S_{xy}).$$
\end{enumerate}
\end{exercise}

\begin{exercise}\label{exercise:12.11}\leavevmode
\begin{enumerate}[label=(\alph*)]
\item\label{item:12.11.a}\index{SI0222@线性函数}\index{SI0534@协方差} 证明对随机变量$X,Y$和常数$a,b,c,d$，
$$\operatorname{Cov}(aY+bX,cY+dX)=ac\operatorname{Var}Y+(bc+ad)\operatorname{Cov}(X,Y)+bd\operatorname{Var}X.$$
\item\label{item:12.11.b} 用\hyperref[item:12.11.a]{(a)}验证，在$\sigma_\delta^2=\lambda\sigma_\varepsilon^2$的结构关系模型中，
$$\operatorname{Cov}(\beta\lambda Y_i+X_i,Y_i-\beta X_i)=0,$$
这就是Creasy--Williams置信集所依据的恒等式.
\sourcepageanchor{521}
\item\label{item:12.11.c} 用\hyperref[item:12.11.b]{(b)}证明，对任意$\beta$，
$$\frac{\sqrt{n-2}\,r_\lambda(\beta)}{\sqrt{1-r_\lambda^2(\beta)}}\sim t_{n-2},$$
其中$r_\lambda(\beta)$由\eqref{eq:12.2.23}给出. 再证明\eqref{eq:12.2.24}定义的置信集具有恒定覆盖概率$1-\alpha$.
\end{enumerate}
\end{exercise}

\begin{exercise}\label{exercise:12.12}
验证关于$\hat\beta$（假设$\sigma_\delta^2=\lambda\sigma_\varepsilon^2$时$\beta$的MLE）、\eqref{eq:12.2.23}的$r_\lambda(\beta)$以及\eqref{eq:12.2.24}的Creasy--Williams置信集$C_\lambda(\hat\beta)$的下列事实.
\begin{enumerate}[label=(\alph*)]
\item\label{item:12.12.a} $\hat\beta$与$-1/(\lambda\hat\beta)$是定义似然函数\eqref{eq:12.2.14}一阶导数零点的二次方程的两个根.
\item\label{item:12.12.b} 对每个$\beta$，$r_\lambda(\beta)=-r_\lambda(-1/(\lambda\beta))$.
\item\label{item:12.12.c} 若$\beta\in C_\lambda(\hat\beta)$，则$-1/(\lambda\beta)\in C_\lambda(\hat\beta)$.
\end{enumerate}
\end{exercise}

\begin{exercise}\label{exercise:12.13}
\eqref{eq:12.2.24}的Creasy--Williams置信集与\eqref{eq:12.2.22}的区间$C_G(\hat\beta)$之间有一个有趣联系.
\begin{enumerate}[label=(\alph*)]
\item\label{item:12.13.a} 证明
$$C_G(\hat\beta)=\left\{\beta:\frac{(\beta-\hat\beta)^2}{\hat\sigma_\beta^2/(n-2)}\le F_{1,n-2,\alpha}\right\},$$
其中$\hat\beta$是$\beta$的MLE，$\hat\sigma_\beta^2$是先前定义的$\sigma_\beta^2$的相合估计量.
\item\label{item:12.13.b} 证明Creasy--Williams集合可写为
$$\left\{\beta:\frac{(\beta-\hat\beta)^2}{\hat\sigma_\beta^2/(n-2)}
\left[\frac{(1+\lambda\beta\hat\beta)^2}{(1+\lambda\beta^2)^2}\right]\le F_{1,n-2,\alpha}\right\}.$$
因此，将方括号中的项替换为其概率极限1，即得到$C_G(\hat\beta)$. （推导这一表示时，$\hat\beta$与$-1/(\lambda\hat\beta)$是$r_\lambda(\beta)$分子的根这一事实很有帮助. 特别地，容易建立
$$\frac{r_\lambda^2(\beta)}{1-r_\lambda^2(\beta)}
=\frac{\lambda^2S_{xy}^2(\beta-\hat\beta)^2(\beta+(1/\lambda\hat\beta))^2}
{(1+\lambda\beta^2)^2(S_{xx}S_{yy}-S_{xy}^2)}.\quad\text{）}$$
\end{enumerate}
\end{exercise}

\begin{exercise}\label{exercise:12.14}
画出\eqref{eq:12.3.2}的逻辑回归函数$\pi(x)$在以下三种情形的图像：$\alpha=\beta=1$、$\alpha=\beta=2$、$\alpha=\beta=3$.
\end{exercise}

\begin{exercise}\label{exercise:12.15}
对\eqref{eq:12.3.2}的逻辑回归函数，验证以下关系.
\begin{enumerate}[label=(\alph*)]
\item\label{item:12.15.a} $\pi(-\alpha/\beta)=1/2$.
\item\label{item:12.15.b} 对任意$c$，$\pi((-\alpha/\beta)+c)=1-\pi((-\alpha/\beta)-c)$.
\item\label{item:12.15.c} 验证关于$d\pi(x)/dx$的\eqref{eq:12.3.3}.
\item\label{item:12.15.d} 验证关于优势比的\eqref{eq:12.3.4}.
\item\label{item:12.15.e} 验证关于优势乘法变化的\eqref{eq:12.3.5}.
\item\label{item:12.15.f} 验证Bernoulli GLM似然方程的\eqref{eq:12.3.6}、\eqref{eq:12.3.8}.
\item\label{item:12.15.g} 验证逻辑回归中，\eqref{eq:12.3.7}、\eqref{eq:12.3.8}内$f_i/(F_i(1-F_i)=1$.\footnote{译注：原题左边漏一个右括号，应为$f_i/[F_i(1-F_i)]=1$. 原式保留.}
\end{enumerate}
\end{exercise}

\begin{exercise}\label{exercise:12.16}
考虑以下逻辑回归数据，只观测两个取值$x=0,1$. 在$x=0$处，十次试验有十次成功；在$x=1$处，十次试验有五次成功. 验证以下结论，证明这些数据的逻辑回归MLE $\hat\alpha,\hat\beta$不存在.
\begin{enumerate}[label=(\alph*)]
\item\label{item:12.16.a} 不受\eqref{eq:12.3.2}限制时，$\pi(0),\pi(1)$的MLE为$\hat\pi(0)=1$、$\hat\pi(1)=.5$.
\item\label{item:12.16.b} \hyperref[item:12.16.a]{(a)}的估计所给似然函数整体最大值，不能在逻辑回归参数$\alpha,\beta$的任何有限值处取得，但可在$\beta\to-\infty$、$\alpha=-\beta$的极限中达到.
\end{enumerate}
\end{exercise}
\sourcepageanchor{522}
\begin{exercise}\label{exercise:12.17}
在probit回归中，连接函数为标准正态cdf $\Phi(x)=P(Z\le x)$，其中$Z\sim\mathrm n(0,1)$. 因而，在这一模型中观测$(Y_1,x_1),\allowbreak(Y_2,x_2),\allowbreak\ldots,\allowbreak(Y_n,x_n)$，其中$Y_i\sim\operatorname{Bernoulli}(\pi_i)$，且$\pi_i=\Phi(\alpha+\beta x_i)$.\footnote{译注：按\autoref{sec:12.3.1}的定义，连接函数应为$\Phi^{-1}$；$\Phi$是逆连接函数. 原文名称保留.}
\begin{enumerate}[label=(\alph*)]
\item\label{item:12.17.a} 写出似然函数，并说明如何求解$\alpha,\beta$的MLE.
\item\label{item:12.17.b} 对\autoref{tab:12.3.1}的数据拟合probit模型，评论它与逻辑回归拟合的任何差异.
\end{enumerate}
\end{exercise}

\begin{exercise}\label{exercise:12.18}
\hyperref[bib:source49]{Brown和Rothery（1993}，第4章）讨论了将线性逻辑模型推广为二次模型：
$$\log\left(\frac{\pi_i}{1-\pi_i}\right)=\alpha+\beta x_i+\gamma x_i^2.$$
\begin{enumerate}[label=(\alph*)]
\item\label{item:12.18.a} 写出似然函数，并说明如何求解$\alpha,\beta,\gamma$的MLE.
\item\label{item:12.18.b} 用对数似然比检验，说明如何检验$H_0:\gamma=0$，即模型实际为线性逻辑模型.
\item\label{item:12.18.c} 对\autoref{tab:12.5.4}中不同年龄雀鹰的存活数据拟合二次逻辑模型.
\item\label{item:12.18.d} 判断雀鹰数据用线性还是二次模型更好，即检验$H_0:\gamma=0$.
\end{enumerate}
\begin{table}[htbp]\centering\setcounter{table}{3}
\caption{按年龄分组的雀鹰存活情况}\label{tab:12.5.4}
\begin{tabular}{lrrrrrrrrr}\toprule
年龄&1&2&3&4&5&6&7&8&9\\
鸟数&77&149&182&118&78&46&27&10&4\\
存活数&35&89&130&79&52&28&14&3&1\\\bottomrule
\end{tabular}
\end{table}
\end{exercise}

\begin{exercise}\label{exercise:12.19}
对于逻辑回归模型：
\begin{enumerate}[label=(\alph*)]
\item\label{item:12.19.a} 证明$\bigl(\sum_{j=1}^JY_j^*,\sum_{j=1}^JY_j^*x_j\bigr)$是$(\alpha,\beta)$的充分统计量.
\item\label{item:12.19.b} 验证\eqref{eq:12.3.10}的逻辑回归信息矩阵公式.
\end{enumerate}
\end{exercise}

\begin{exercise}\label{exercise:12.20}
考虑逻辑回归模型，并假定$\beta=0$.
\begin{enumerate}[label=(\alph*)]
\item\label{item:12.20.a} 若$0<\sum_{i=1}^ny_i<n$，证明$\pi(x)$的MLE为$\hat\pi=\sum_{i=1}^ny_i/n$，此时它不依赖$x$.
\item\label{item:12.20.b} 若$0<\sum_{i=1}^ny_i<n$，证明$\alpha$的MLE为$\hat\alpha_0=\log\bigl((\sum_{i=1}^ny_i)/(n-\sum_{i=1}^ny_i)\bigr)$.
\item\label{item:12.20.c} 证明若$\sum_{i=1}^ny_i=0$或$n$，$\hat\alpha_0$不存在，但检验$H_0:\beta=0$的似然比统计量仍有良好定义.
\end{enumerate}
\end{exercise}

\begin{exercise}\label{exercise:12.21}
设$Y\sim\operatorname{binomial}(n,\pi)$，用$\hat\pi=Y/n$表示$\pi$的MLE. 令$W=\log(\hat\pi/(1-\hat\pi))$为样本logit，即$\log(\pi/(1-\pi))$的MLE. 用Delta方法证明$1/(n\hat\pi(1-\hat\pi))$是$\operatorname{Var}W$的合理估计.
\end{exercise}

\begin{exercise}\label{exercise:12.22}
\autoref{ex:12.4.1}考察了小扰动对斜率最小二乘估计的影响. 作类似计算，评价截距最小二乘估计对小扰动的稳健性.
\end{exercise}

\begin{exercise}\label{exercise:12.23}
\autoref{ex:12.4.1}与\hyperref[ex:10.2.1]{习题10.2.1}不同，引入$\varepsilon_i$的污染分布时没有引入偏倚.\footnote{译注：原文称“Exercise 10.2.1”，本书习题没有这一编号；所指应为\autoref{ex:10.2.1}. 原称呼保留.} 证明即使引入偏倚也不会有影响. 即假设
$$\bigl(E\varepsilon_i,\operatorname{Var}\varepsilon_i\bigr)=
\begin{cases}(0,\sigma^2),&\text{概率为 }1-\delta,\\(\mu,\tau^2),&\text{概率为 }\delta,
\end{cases}$$
则
\begin{enumerate}[label=(\alph*)]
\item\label{item:12.23.a} 最小二乘估计量$b$仍为$\beta$的无偏估计量.
\item\label{item:12.23.b} 最小二乘估计量$a$的期望为$\alpha+\delta\mu$，所以也可将模型假设为$Y_i=\alpha+\delta\mu+\beta x_i+\varepsilon_i$.
\end{enumerate}
\end{exercise}

\begin{exercise}\label{exercise:12.24}
对模型$Y_i=\beta x_i+\varepsilon_i$，证明LAD估计量为$t_{(k^*+1)}$，其中$t_i=y_i/x_i$，$t_{(1)}\le t_{(2)}\le\cdots\le t_{(n)}$，若$x_{(i)}$为与$t_{(i)}$配对的$x$值，则$k^*$满足
$$\sum_{i=1}^{k^*}|x_{(i)}|\le\sum_{i=k^*+1}^n|x_{(i)}|\quad\text{和}\quad
\sum_{i=1}^{k^*+1}|x_{(i)}|>\sum_{i=k^*+2}^n|x_{(i)}|.$$
\end{exercise}

\sourcepageanchor{523}
\begin{exercise}\label{exercise:12.25}\index{SI0894@最小绝对偏差}
LAD回归直线的一个问题是，它不总有唯一定义.
\begin{enumerate}[label=(\alph*)]
\item\label{item:12.25.a} 对含三个观测$(x_1,y_1)$、$(x_1,y_2)$、$(x_3,y_3)$的数据集，注意前两个$x$相同，证明任何通过$(x_3,y_3)$并位于$(x_1,y_1)$与$(x_1,y_2)$之间的直线，都是最小绝对偏差直线.
\item\label{item:12.25.b} 对三名个体测量心率$x$，单位为每分钟心跳次数，以及耗氧量$y$，单位为ml/kg. $(x,y)$数据为$(127,14.4)$、$(127,11.9)$、$(136,17.9)$. 计算最小二乘直线的斜率、截距及最小绝对偏差直线的范围.
\end{enumerate}
对于最小绝对偏差直线的价值，似乎存在一些分歧. 它肯定比最小二乘稳健，但可能很难计算，不过\hyperref[bib:source244]{Portnoy和Koenker（1997）}给出了高效算法. \hyperref[bib:source115]{Ellis（1998）}似乎质疑其稳健性，而\hyperref[bib:source245]{Portnoy和Misera（1998）}\footnote{译注：原正文写 Misera，书末对应文献写 Mizera；本段 Ellis 文献的书末条目没有列出年份. 保留原正文和书目各自的写法.}在一篇讨论中又质疑Ellis.
\end{exercise}

\autoref{exercise:12.26}至\autoref{exercise:12.28}考察\autoref{ex:12.4.3}的一些细节.
\begin{exercise}\label{exercise:12.26}\leavevmode
\begin{enumerate}[label=(\alph*)]
\item\label{item:12.26.a} \autoref{ex:12.4.3}始终假定$\frac1n\sum_{i=1}^nx_i^2\to\sigma_x^2<\infty$. 证明以下情形满足此条件：(i) $x_i=1$，即普通中位数的情形；(ii) $|x_i|\le1$，即$x_i$有界的情形.\footnote{译注：有界性只保证这些平均值有界，不保证其收敛；(ii)还需平方均值存在极限等条件. 可由长短交替的0、1区块构造均值不收敛的有界序列. 原题照录.}
\item\label{item:12.26.b} 证明在\hyperref[item:12.26.a]{(a)}关于$x_i$的条件下，$\frac1n\sum_{i=1}^n\psi(y_i-\hat\beta_Lx_i)\to0$依概率成立.
\end{enumerate}
\end{exercise}

\begin{exercise}\label{exercise:12.27}\leavevmode
\begin{enumerate}[label=(\alph*)]
\item\label{item:12.27.a} 验证$-\frac1{\sqrt n}\sum_{i=1}^n\psi(Y_i-\hat\beta x_i)\to\mathrm n(0,\sigma_x^2)$.
\item\label{item:12.27.b} 验证
$$\left.\frac1n\sum_{i=1}^n\frac d{d\beta_0}E_\beta[\psi(Y_i-\beta_0x_i)]\right|_{\beta_0=\beta}
=2f(0)\frac1n\sum_{i=1}^nx_i^2,$$
并结合\hyperref[item:12.27.a]{(a)}得出$\sqrt n(\hat\beta_L-\beta)\to\mathrm n\bigl(0,1/(4f(0)^2\sigma_x^2)\bigr)$.\footnote{译注：本题重复\autoref{ex:12.4.3}中的符号疑点：\hyperref[item:12.27.a]{(a)}的$\hat\beta$应核为真实参数$\beta$；按所给$\psi$定义，\hyperref[item:12.27.b]{(b)}导数右端应带负号. 原题保留，依据见该例译注.}
\end{enumerate}
\end{exercise}

\begin{exercise}\label{exercise:12.28}
证明最小二乘估计量为$\hat\beta=\sum_{i=1}^nx_iy_i/\sum_{i=1}^nx_i^2$，并且$\sqrt n(\hat\beta-\beta)\to\mathrm n(0,1/\sigma_x^2)$.
\end{exercise}

\begin{exercise}\label{exercise:12.29}
用\autoref{ex:12.4.3}那样的Taylor级数论证，推导独立同分布抽样中中位数的渐近分布.
\end{exercise}

\begin{exercise}\label{exercise:12.30}
对\autoref{tab:12.4.2}的数据，用参数自助法评价LAD和M估计拟合的标准误. 具体地：
\begin{enumerate}[label=(\alph*)]
\item\label{item:12.30.a} 拟合直线$y=\alpha+\beta x$，得到估计$\tilde\alpha,\tilde\beta$.
\item\label{item:12.30.b} 计算残差均方误差$\hat\sigma^2=\frac1{n-2}\sum_{i=1}^n[y_i-(\tilde\alpha+\tilde\beta x_i)]^2$.
\item\label{item:12.30.c} 从$\mathrm n(0,\hat\sigma^2)$生成新残差，重新估计$\alpha,\beta$.
\item\label{item:12.30.d} 将\hyperref[item:12.30.c]{(c)}重复$B$次，计算$\tilde\alpha,\tilde\beta$的标准差.
\item\label{item:12.30.e} 分别用双指数和Laplace分布作为误差，重复\hyperref[item:12.30.a]{(a)}至\hyperref[item:12.30.d]{(d)}，并与正态情形比较.\footnote{译注：双指数分布与Laplace分布通常同义；此处可能沿袭\autoref{ex:12.4.5}中的名称疑误，是否原拟指Logistic分布尚待核实. 原题照录.}
\end{enumerate}
\end{exercise}

\begin{exercise}\label{exercise:12.31}
对\autoref{tab:12.4.2}的数据，也可以用非参数自助法评价LAD和M估计拟合的标准误. 此时：
\begin{enumerate}[label=(\alph*)]
\item\label{item:12.31.a} 拟合直线$y=\alpha+\beta x$，得到估计$\tilde\alpha,\tilde\beta$.
\item\label{item:12.31.b} 对拟合残差重抽样生成新残差，并重新估计$\alpha,\beta$.
\item\label{item:12.31.c} 将\hyperref[item:12.31.c]{(c)}重复$B$次，计算$\tilde\alpha,\tilde\beta$的标准差.\footnote{译注：原题(c)自我引用；按自助法步骤，应重复\hyperref[item:12.31.b]{(b)}. 原引用保留.}
\end{enumerate}
\end{exercise}



\subsection{杂录}\label{sec:12.6}
\subsubsection{函数与结构的含义}\label{sec:12.6.1}\index{SI0410@函数关系}\index{SI0417@结构关系}
“函数”和“结构”这两个名称本身就是EIV模型中混淆的主要来源. \hyperref[bib:source189]{Kendall和Stuart（1979}，第29章）详细讨论了这些概念，区分数学变量（非随机变量）间的关系与随机变量间的关系. 理解二者联系的一种方式，是把模型写成层次结构：对函数模型的参数赋予分布，就得到结构关系模型：\footnote{译注：原图将含随机误差的等式左边写为条件期望；由\eqref{eq:12.2.5}、\eqref{eq:12.2.6}，这里应直接为$Y_i,X_i$，或在保留期望时去掉误差项. 原表达式保留.}
$$\underbrace{\begin{aligned}
\text{函数关系模型}\quad\left\{\begin{aligned}
E(Y_i\mid\xi_i)&=\alpha+\beta\xi_i+\varepsilon_i,&\varepsilon_i&\sim\mathrm n(0,\sigma_\varepsilon^2),\\
E(X_i\mid\xi_i)&=\xi_i+\delta_i,&\delta_i&\sim\mathrm n(0,\sigma_\delta^2),
\end{aligned}\right.\\
\xi_i\sim\mathrm n(\xi,\sigma_\xi^2)
\end{aligned}}_{\text{结构关系模型}}$$

\sourcepageanchor{524}
可通过以下区别理解这两个词，尽管这种区别并非普遍接受. 例如，微积分中常见方程$y=f(x)$，它描述函数关系，即假定变量间存在的关系. 因此，若把函数关系理解为两变量间假定的关系，那么在函数关系模型和结构关系模型中，方程$\eta_i=\alpha+\beta\xi_i$都是函数关系，即假设的关系，其中$\eta_i=E(Y_i\mid\xi_i)$.

另一方面，结构关系是从问题的假定结构中产生的关系. 因而，在结构关系模型中，$\eta=EY_i=\alpha+\beta\xi=\alpha+\beta EX_i$可由模型结构推出，所以是结构关系.

为使这些想法更清楚，考虑假定$x$没有误差的简单线性回归. 方程$E(Y_i\mid x_i)=\alpha+\beta x_i$是函数关系，即假定$E(Y_i\mid x_i)$与$x_i$间存在的关系. 然而，也可以假定$(X_i,Y_i)$服从二元正态分布，并给定$x_i$作简单线性回归. 此时，关系$E(Y_i\mid x_i)=\alpha+\beta x_i$由所假设模型的结构推出，因而是结构关系.

按这些含义，术语的区分成了偏好问题. 任一模型中，都可以由函数关系推出结构关系，反之亦然. 重要区别在于：推断之前，是否积分消去滋扰参数$\xi_i$.

\subsubsection{EIV模型中普通最小二乘的相合性}\label{sec:12.6.2}\index{SI0409@相合性}\index{SI0608@变量含误差}
一般而言，用普通最小二乘估计量估计EIV回归斜率并不是好主意，因为它不相合. 假设线性结构关系\eqref{eq:12.2.6}成立，则
$$\begin{aligned}
\hat\beta&=\frac{\sum_{i=1}^n(X_i-\bar X)(Y_i-\bar Y)}{\sum_{i=1}^n(X_i-\bar X)^2}\\
&=\frac{\frac1n\sum_{i=1}^n(X_i-\bar X)(Y_i-\bar Y)}{\frac1n\sum_{i=1}^n(X_i-\bar X)^2}\\
&\to\frac{\operatorname{Cov}(X,Y)}{\operatorname{Var}X}\qquad(n\to\infty\text{，利用弱大数定律})\\
&=\beta\frac{\sigma_\xi^2}{\sigma_\delta^2+\sigma_\xi^2}\qquad\text{（由\eqref{eq:12.2.19}）},
\end{aligned}$$
说明$\hat\beta$不可能相合. 函数关系情形也可得到同类结果.

\hyperref[bib:source81]{Cochran（1968）}讨论了EIV模型中$\hat\beta$的表现. \hyperref[bib:source59]{Carroll、Gallo和Gleser（1985）}以及\hyperref[bib:source148]{Gleser、Carroll和Gallo（1987）}研究了普通最小二乘估计量的函数何时相合.

\subsubsection{EIV模型中的工具变量}\label{sec:12.6.3}\index{SI0583@工具变量}
工具变量（instrumental variables）的概念至少可追溯到\hyperref[bib:source324]{Wald（1940）}，他借此构造了斜率的相合估计量. 为理解工具变量，将EIV模型写为
$$\begin{aligned}Y_i&=\alpha+\beta\xi_i+\varepsilon_i,\\X_i&=\xi_i+\delta_i,\end{aligned}$$
作代数变换，得
$$Y_i=\alpha+\beta X_i+[\varepsilon_i-\beta\delta_i].$$
\sourcepageanchor{525}
工具变量$Z_i$是能良好预测$X_i$但与$\nu_i=\varepsilon_i-\beta\delta_i$不相关的随机变量. 若能找到这样的变量，就可用来改进预测；特别地，可以构造$\beta$的相合估计量.

\hyperref[bib:source324]{Wald（1940）}证明，在相当一般的条件下，估计量
$$\hat\beta_{\mathrm W}=\frac{\bar Y_{(1)}-\bar Y_{(2)}}{\bar X_{(1)}-\bar X_{(2)}}$$
在可识别模型中是$\beta$的相合估计量，下标表示数据的两个分组. 只取两个值以定义分组的变量$Z_i$就是工具变量. 关于Wald估计量的讨论，见\hyperref[bib:source222]{Moran（1971）}.

工具变量虽然很有帮助，使用时也可能有问题. 例如，\hyperref[bib:source117]{Feldstein（1974）}给出使用工具变量反而有害的情形. \hyperref[bib:source222]{Moran（1971）}讨论了验证$\hat\beta_{\mathrm W}$这样简单的估计量相合所需条件的困难. \hyperref[bib:source130]{Fuller（1987）}深入讨论了工具变量. \hyperref[bib:source28]{Berkson（1950）}提出的模型利用了类似于工具变量的相关结构.

\subsubsection{逻辑回归的似然方程}\label{sec:12.6.4}
逻辑回归模型的似然方程关于参数为非线性，必须数值求解. 最常用的方法为Newton--Raphson方法. 它从MLE的初始猜测$(\hat\alpha^{(1)},\hat\beta^{(1)})$开始，用对数似然在该点的二阶Taylor展开作二次函数近似. MLE的下一次猜测$(\hat\alpha^{(2)},\hat\beta^{(2)})$就是这个二次函数的最大化点. 再以$(\hat\alpha^{(2)},\hat\beta^{(2)})$为中心作另一个二次近似，其最大化点给出下一次猜测. Taylor近似涉及对数似然的一阶、二阶导数，在当前猜测$(\hat\alpha^{(t)},\hat\beta^{(t)})$处计算. 这些二阶导数正是\eqref{eq:12.3.10}的信息矩阵中出现的导数. 因此，这种求解似然方程方法的一个副产品，是$\hat\alpha,\hat\beta$的方差与协方差估计. 对逻辑回归模型，猜测$(\hat\alpha^{(t)},\hat\beta^{(t)})$通常很快收敛到MLE $(\hat\alpha,\hat\beta)$，往往只需几次迭代就得到满意近似.

Newton--Raphson方法也称为迭代重加权最小二乘（iteratively reweighted least squares）. 每一步对$(\hat\alpha,\hat\beta)$的下一次猜测，都可表示为某个最小二乘问题的解. 但这个问题对平方和中的不同项给予不同权重. 此时权重为$n_jF_j^{(t)}(1-F_j^{(t)})$，其中$F_j^{(t)}=F(\hat\alpha^{(t)}+\hat\beta^{(t)}x_j)$，$F$为logistic的cdf. 这是第$j$个样本logit方差近似值的倒数（\autoref{exercise:12.21}）. 每一步都使用当前MLE猜测重新计算权重，因而称为“迭代重加权”. 所以，Newton--Raphson方法近似等于以样本logit为数据，进行加权最小二乘估计参数所得的结果.

\subsubsection{稳健回归的进一步讨论}\label{sec:12.6.5}
多年来，最小二乘的稳健替代方法一直是研究对象，已有大量文献讨论各种问题. \autoref{sec:12.4}只简要介绍了稳健回归，但其一些优点和困难应已显现. 有许多优秀专著详细讨论稳健回归，包括\hyperref[bib:source165]{Hettmansberger和McKean（1998）}、\hyperref[bib:source299]{Staudte和Sheather（1990）}及\hyperref[bib:source171]{Huber（1981）}.\footnote{译注：原正文作者姓氏写为Hettmansberger，而本书书目相应条目写为Hettmansperger. 原正文拼法保留.} 以下其他主题也受到大量关注.

\paragraph{截尾与变换}
\sourcepageanchor{526}
Carroll、Ruppert及其合作者的工作涵盖了稳健回归的许多方面. \hyperref[bib:source269]{Ruppert和Carroll（1979）}细致研究了截尾的渐近理论，截尾均值见\autoref{exercise:10.20}；\hyperref[bib:source60]{Carroll和Ruppert（1985）}考察误差不同分布时最小二乘的替代方法. 后续工作研究了回归中变换的优点（Carroll和Ruppert，\hyperref[bib:source60]{1985}，\hyperref[bib:source61]{1988}）.

\paragraph{其他稳健替代方法}
这里只考察了LAD估计量和一个M估计量，当然还有许多其他稳健估计量. 一个常用替代方法是\hyperref[bib:source262]{Rousseeuw（1984）}的最小平方中位数（Least Median of Squares，LMS）估计量，另见\hyperref[bib:source264]{Rousseeuw和Leroy（1987）}. 此外还有基于秩的回归估计量，即R估计量（R-estimators；见\hyperref[bib:source101]{Draper，1988}的综述）. 较近期还有关于数据深度（data depth）的研究（\hyperref[bib:source209]{Liu，1990}；\hyperref[bib:source210]{Liu和Singh，1992）}，在回归中应用于寻找最深直线（deepest line；\hyperref[bib:source263]{Rousseeuw和Hubert，1999）}.

\paragraph{计算}
从实际角度看，计算稳健估计可能颇具挑战，因为经常面对困难的最小化问题. \hyperref[bib:source244]{Portnoy和Koenker（1997）}的综述讨论LAD估计的计算. Hawkins（\hyperref[bib:source161]{1993}，\hyperref[bib:source162]{1994}，\hyperref[bib:source163]{1995}）给出了计算LMS及相关估计的一些算法.




\appendix

% 原PDF528–533／印刷497–502；全部原页已目视核对。
\clearpage\section*{计算机代数}\phantomsection\label{app:computer-algebra}
\addcontentsline{toc}{section}{计算机代数}
\begingroup
\renewcommand{\thesection}{\arabic{section}}
\setcounter{section}{12}\setcounter{subsection}{6}\setcounter{theorem}{0}
\sourcepageanchor{528}
计算机代数系统（computer algebra system）能够对表达式进行符号运算.面对繁琐、机械的计算时（例如对一个比值求二阶导数），它们尤其有用.本附录通过各种问题说明这类系统的使用.理解和使用统计学绝不以使用计算机代数系统为必要条件；不过，它不仅能减轻一些枯燥的计算工作，还能引导我们得到新的认识和更实用的答案.

需要认识到，计算机代数系统有很多种，它们可以帮助完成的计算也多种多样（求和、积分、模拟等等）.本附录的目的既不是教授这些系统的用法，也不是展示它们的一切可能用途，而只是说明其中一些可能性.

我们使用Mathematica软件包演示计算.其他软件包，例如Maple，也能用于这些计算.

\needspace{5\baselineskip}\medskip\noindent\textit{\hyperref[sec:1]{第1章}}
\begin{example}[无序抽样]\label{ex:12.6.1}
我们给出Mathematica代码，枚举\autoref{ex:1.2.20}中从$\{2,4,9,12\}$有放回抽样得到的无序结果.枚举结果并计算多项式权重后，将结果与权重排序.注意，要画出\autoref{fig:1.2.2}中的直方图，还需要做一些工作.例如，有两个不同结果的平均值都是8，因此要得到\autoref{fig:1.2.2}那样的图形，需要将结果$\{8,\frac{3}{128}\}$、$\{8,\frac{3}{64}\}$合并成$\{8,\frac{9}{128}\}$.

集合中的数超过七个时，这种枚举会变得非常耗时，因为这会产生$\binom{13}{7}=27,132$个无序结果.\footnote{译注：原文等式照录，但$\binom{13}{7}=1716$，与所列$27,132$不一致.}
\begin{enumerate}[label=(\arabic*)]
\item “DiscreteMath”包包含用于计算排列与组合的函数.
\begin{sourcecode}
In[1]:= Needs["DiscreteMath`Combinatorica`"]
\end{sourcecode}
\item 令$x=$数值的集合.不同样本的个数为$\operatorname{NumberofCompositions}[n,m]=\binom{n+m-1}{n}$.
\begin{sourcecode}
In[2]:= x ={2,4,9,12};
n=Length[x];
ncomp=NumberOfCompositions[n,n]
\end{sourcecode}
\noindent$\textbf{Out[4]}=35$.
\item 枚举样本（w），计算每个样本的平均值（avg），并计算各个值的权重（wt）.权重就是对应于该配置的多项式系数.
\begin{sourcecode}
In[5]:= w = Compositions[n, n];
wt = n!/(Apply[Times, Factorial /@ w, 1]*n^n);
avg = w . x/n;
Sort[Transpose[{avg, wt}]]
Out[8] = {2, 1/256}, {5/2, 1/64}, {3, 3/128}, {7/2, 1/64}, {15/4, 1/64},
{4, 1/256}, {17/4, 3/64}, {9/2, 1/64}, {19/4, 3/64}, {5, 3/64},
{21/4, 1/64}, {11/2, 3/128}, {11/2, 3/64}, {6, 1/64}, {6, 3/64},
{25/4, 3/64}, {13/2, 3/128}, {27/4, 3/32}, {7, 3/128}, {29/4, 1/64},
{29/4, 3/64}, {15/2, 3/64}, {31/4, 1/64}, {8, 3/128}, {8, 3/64},
{17/2, 3/64}, {35/4, 3/64}, {9, 1/256}, {37/4, 3/64}, {19/2, 1/64},
{39/4, 1/64}, {10, 1/64}, {21/2, 3/128}, {45/4, 1/64}, {12, 1/256}
\end{sourcecode}
\end{enumerate}
\end{example}
\sourcepageanchor{529}
\needspace{5\baselineskip}\medskip\noindent\textit{\hyperref[sec:2]{第2章}}
\begin{example}[一元变换]\label{ex:12.6.2}
\hyperref[exercise:2.1]{习题2.1(a)}是一个标准的一元变量变换.对计算机代数程序来说，这类计算通常很容易.
\begin{enumerate}[label=(\arabic*)]
\item 输入$f(x)$，并解出变换后的变量.\footnote{译注：代码中的反斜线以及最后一个复根的符号、$y$的幂次均按原页保留.最后一项一般不满足$x^3=y$；本例后续仅使用第一个实根$x=y^{1/3}$.本译稿没有执行或修补这些历史代码.}
\begin{sourcecode}
In[1]:=f[x_] := 42*(x^5)*(1 - x)
sol = Solve[y == x\^3, x]
\end{sourcecode}
\begin{equation*}
\textbf{Out[2]}=\{\{x\mathbin{->}y^{1/3}\},\{x\mathbin{->}-(-1)^{1/3}y^{1/3}\},\{x\mathbin{->}-(-1)^{2/3}y^{2/3}\}\}.
\end{equation*}

\item 计算雅可比因子.
\begin{sourcecode}
In[3]:= D[x/.sol[[1]],y]
\end{sourcecode}
$$\textbf{Out[3]}=\frac{1}{3y^{2/3}}.$$
\item 计算变换后变量的密度.
\begin{sourcecode}
In[4]:= f[x/.sol[[1]]]* D[x/.sol[[1]],y]
\end{sourcecode}
$$\textbf{Out[4]}=14(1-y^{1/3})y.$$
\end{enumerate}
\end{example}

\needspace{5\baselineskip}\medskip\noindent\textit{\hyperref[sec:4]{第4章}}
\begin{example}[二元变换]\label{ex:12.6.3}
下面演示一些二元变换.使用类似的代码，也可以完成多元变换.第一个计算演示\autoref{ex:4.3.4}，求正态变量之和的分布.然后计算\autoref{ex:4.3.3}：先对两个贝塔密度的乘积进行二元变换，再求边缘分布.
\begin{enumerate}[label=(\arabic*)]
\item 正态变量之和.Out[4]是联合分布，Out[5]是边缘分布.\footnote{译注：原代码对$v$从0积到$\infty$，只覆盖其支持的一半，因此所列Out[5]是完整边缘密度的一半.若求$u=X+Y$的密度，需对$v$从$-\infty$积到$\infty$.}
\begin{sourcecode}
In[1]:=f[x_,y_]:=(1/(2*Pi))*E^(-x^2/2)*E^(-y^2/2)
So:=Solve[{u==x+y,v==x-y},{x,y}]
g:=f[x/.So,y/.So]*Abs[Det[Outer[D,First[{x,y}/.So],{u,v}]]]
Simplify[g]
\end{sourcecode}
$$\textbf{Out[4]}=\left\{\frac{e^{-u^2/4-v^2/4}}{4\pi}\right\}]$$
\begin{sourcecode}
In[5]:= Integrate[g,{v,0,Infinity}]
\end{sourcecode}
$$\textbf{Out[5]}=\left\{\frac{e^{-u^2/4}}{4\sqrt{\pi}}\right\}]$$
\item 贝塔变量的乘积.（“ContinuousDistributions”包包含许多标准分布的pdf和cdf.）Out[10]是贝塔变量乘积的联合密度，Out[11]是$u$的密度.If语句读作If(检验条件,真时的值,假时的值)，因此检验条件为真时，取中间的值.在大多数情形下，检验条件为真，边缘密度就是所给的贝塔密度.\footnote{译注：由变换$u=xy,v=x$及$0<x,y<1$，联合支持应为$0<u<v<1$；原代码的积分写成$v\in(0,1)$，没有显式编码这一支持限制.原输出中还存在括号不配对，以下均照录，不将其当作已核验的可执行代码.}
\sourcepageanchor{530}
\begin{sourcecode}
In[6]:= Needs["Statistics`ContinuousDistributions`"]
Clear[f, g, u, v, x, y, a, b, c]
f[x_, y_] := PDF[BetaDistribution[a, b],x]
*PDF[BetaDistribution[a + b, c], y]
So := Solve[{u == x*y, v == x}, {x, y}]
g[u_, v_] = f[x /. So, y /. So]
*Abs[Det[Outer[D, First[{x, y} /. So], {u, v}]]]
Integrate[g[u, v] ,{v,0,1}]
\end{sourcecode}
$$\textbf{Out[10]}=\left\{\frac{(1-u/v)^{-1+c}(1-v)^{-1+b}(u/v)^{-1+a+b}v^{-1+a}}{\operatorname{Abs}[v]\operatorname{Beta}[a,b]\operatorname{Beta}[a+b,c]}\right\}$$
\begin{equation*}\begin{aligned}
\textbf{Out[11]}={}&\{\operatorname{If}[\operatorname{Re}[b]>0\mathbin{\&\&}\operatorname{Re}[b+c]<1\mathbin{\&\&}\operatorname{Im}[u]==0\mathbin{\&\&}u>0\mathbin{\&\&}u<1),\\
&\frac{1}{\operatorname{Gamma}[1-c]}\Bigl((1-u)^{-1+b+c}(-1/u)^{b+c}(-u)^cu^{-1+a+b}\\
&\hspace{3em}\operatorname{Gamma}[b]\operatorname{Gamma}[1-b-c]\Bigr),\\
&\left.\int_0^1\frac{(1-u/v)^{-1+c}(1-v)^{-1+b}(u/v)^{-1+a+b}v^{-1+a}}{\operatorname{Abs}[v]\operatorname{Beta}[a,b]\operatorname{Beta}[a+b,c]}\,dv\right\}.
\end{aligned}\end{equation*}
\end{enumerate}
\end{example}

\begin{example}[正态概率]\label{ex:12.6.4}
\hyperref[exercise:4.14]{习题4.14(a)}所要求的计算很容易处理.首先直接计算；Mathematica容易以数值方式求出积分，却没有找到闭式表达式.注意，答案用erf函数表示，该函数定义为
$$\operatorname{erf}(z)=\frac{2}{\sqrt{\pi}}\int_0^z e^{-t^2}\,dt.$$
如果用一个卡方随机变量重新表达此概率，就能找到闭式表达式.

为求积分，先设置被积函数与积分限.
\begin{sourcecode}
In[1]:= Needs["Statistics`ContinuousDistributions`"]
Clear[f, g, x, y]
f[x_, y_] = PDF[NormalDistribution[0, 1],x]
*PDF[NormalDistribution[0, 1],y]
g[x_] = Sqrt[1 - x^2]
\end{sourcecode}
$$\textbf{Out[3]}=\frac{e^{-x^2/2-y^2/2}}{2\pi}.$$
现在计算二重积分，得到Erf函数.命令N[\%]对前一行作数值求值.
\sourcepageanchor{531}
\begin{sourcecode}
In[5]:= Integrate[f[x, y], {x, -1, 1}, {y, -g[x], g[x]}]
N[%]
\end{sourcecode}
$$\textbf{Out[5]}=\frac{\int_{-1}^1e^{-x^2/2}\operatorname{Erf}\!\left(\frac{\sqrt{1-x^2}}{\sqrt2}\right)\,dx}{\sqrt{2\pi}},\qquad\textbf{Out[6]}=0.393469.$$
当然，我们知道$X^2+Y^2$是自由度为2的卡方随机变量.利用这一事实，就得到闭式答案.
\begin{sourcecode}
In[7]: = Clear[f, t]
f[t_] = PDF[ChiSquareDistribution[2], t];
Integrate[f[t], {t, 0, 1}]
N[%]
\end{sourcecode}
$$\textbf{Out[10]}=\frac12\left(2-\frac{2}{\sqrt e}\right),\qquad\textbf{Out[11]}=0.393469.$$
\end{example}

\needspace{5\baselineskip}\medskip\noindent\textit{\hyperref[sec:5]{第5章}}
\begin{example}[和的密度]\label{ex:12.6.5}
用于说明\autoref{thm:5.2.9}的\autoref{ex:5.2.10}，其中的计算相当复杂.下面在正态、柯西和Student的$t$这三种情形中演示这种计算.

有两点需要注意.
\begin{enumerate}[label=(\arabic*)]
\item 为正确解释答案，可能需要一些复分析知识.正态情形下，报告的答案带有关于变量$z$实部取值的条件（$z$可能取复数值）.在柯西例子中，重要的是知道$I^2=-1$，因此有
$$\frac{2}{\pi(-2I+z)(2I+z)}=\frac{2}{\pi(4+z^2)}.$$
\item Student的$t$变量相加时，似乎当自由度之和为偶数时存在闭式表达式，否则积分就必须用数值方法求解.（这是我们试用计算机代数系统时发现的一条经验观察.）
\end{enumerate}
还应指出，较新版本的计算机代数程序也许会避免这里出现的复数.不过，其他计算中仍会出现复数，因此最好做好处理它们的准备.
\begin{enumerate}[label=(\arabic*)]
\item 两个正态变量之和的密度.\footnote{译注：原页中$f$的代码括号不配对，分母写为$2\pi$；所列Out[3]也没有给出通常的$N(0,2)$密度$e^{-z^2/4}/(2\sqrt\pi)$.这里保留原代码和输出，不据其报告实际运行结果.}
\begin{sourcecode}
In[1]:= Clear[f, x, y, z]
f[x_] = Exp[(-x^2/2]/(2 Pi);
Integrate[f[y]* f[z - y], {y,-Infinity,Infinity}]
\end{sourcecode}
$$\textbf{Out[3]}=\operatorname{If}\left[\operatorname{Re}[z]<0,\frac{E^{-z^2/4}}{2\pi},\int_{-\infty}^{\infty}\frac{E^{-y^2/2-\frac12(-y+z)^2}}{2\pi}\,dy\right].$$
\item 两个柯西变量之和的密度.
\begin{sourcecode}
In[4]:= Clear[f, x, y, z]
f[x_] = 1/(Pi*(1+x^2));
Integrate[f[y]* f[z - y], {y,-Infinity,Infinity}]
\end{sourcecode}
$$\textbf{Out[6]}=\frac{2}{\pi(-2I+z)(2I+z)}.$$
\sourcepageanchor{532}
\item 两个自由度为5的$t$变量之和的密度.\footnote{译注：下面原代码的PDF函数少了一个右方括号；规范的函数调用应为\texttt{PDF[StudentTDistribution[5], x]}.原排印保留.}
\begin{sourcecode}
In[7]:= Needs["Statistics`ContinuousDistributions`"]
Clear[f, x, y, z]
f[x_] = PDF[StudentTDistribution[5, x]
Integrate[f[y]* f[z - y], {y,-Infinity,Infinity}]
\end{sourcecode}
$$\textbf{Out[10]}=\frac{400\sqrt5(8400+120z^2+z^4)}{3\pi(20+z^2)^5}.$$
\end{enumerate}
\end{example}

\begin{example}[均匀变量之和的四阶矩]\label{ex:12.6.6}
\autoref{exercise:5.51}要求计算12个均匀随机变量之和的四阶矩.由于密度具有分段形式，推导它有些费力（不过也可以用计算机代数完成）.然而，使用矩母函数可以简化事情.
\begin{enumerate}[label=(\arabic*)]
\item 先计算均匀随机变量$X_1$的矩母函数，再计算$\sum_{i=1}^{12}X_i$的矩母函数，其中各$X_i$相互独立.
\begin{sourcecode}
In[1]:= M[t_] = Integrate[Exp[t*x], {x, 0, 1}]
In[2]:= Msum[t_] = M[t]^12
\end{sourcecode}
$$\textbf{Out[1]}=\frac{-1+e^t}{t},\qquad\textbf{Out[2]}=\frac{(-1+e^t)^{12}}{t^{12}}.$$
\item 计算$\sum_{i=1}^{12}X_i-6$的矩母函数的四阶导数.它太长，不适合打印出来.
\begin{sourcecode}
In[3]:= g[t_]=D[Exp[-6*t]*Msum[t],{t,4}];
\end{sourcecode}
\item 值g[0]就是四阶矩.但是直接代入0会导致除以零，因此必须用极限进行计算.
\begin{sourcecode}
In[4]:= g[0]
\end{sourcecode}
\noindent Power:infy: Infinite expression $\frac{1}{0^{16}}$ encountered.
\begin{sourcecode}
In[5]:= Limit[g[t],t->0]
\end{sourcecode}
$$\textbf{Out[5]}=\frac{29}{10}.$$
\end{enumerate}
\end{example}

\needspace{5\baselineskip}\medskip\noindent\textit{\hyperref[sec:7]{第7章}}\footnote{译注：此处原小标题为“Chapter 7”，但正文引用的是第10章的例与图，以下保留原小标题.}
\begin{example}[伽马均值的渐近相对效率]\label{ex:12.6.7}
\autoref{ex:10.1.18}中得到\autoref{fig:10.1.1}的计算是在Mathematica中完成的.下面的代码会生成其中一条渐近相对效率（ARE）曲线.
\begin{enumerate}[label=(\arabic*)]
\item 用符号方法求对数似然的二阶导数.
\begin{sourcecode}
In[1]:= Needs["Statistics`ContinuousDistributions`"]
Clear[m, b, x]
f[x_, m_, b_] = PDF[GammaDistribution[m/b, b], x];
loglike2[m_, b_, x_] = D[Log[f[x, m, b]], {m, 2}];
\end{sourcecode}
\item 对二阶导数按该密度积分，求出渐近方差.
\begin{sourcecode}
In[5]:= var[m_, b_] := 1/(-Integrate[loglike2[m, b, x]*f[x, m, b],
 {x, 0, Infinity}])
\end{sourcecode}
\sourcepageanchor{533}
\item 以下代码设置绘图.
\begin{sourcecode}
In[6]:=
mu = {1, 2, 3, 4, 6, 8, 10};
beta = 5;
mlevar = Table[var[mu[[i]], beta], {i, 1, 7}];
momvar = Table[mu[[i]]*beta, {i, 1, 7}];
ARE = momvar/mlevar
Out[10] = {5.25348, 2.91014, 2.18173, 1.83958, 1.52085, 1.37349, 1.28987}
In[11]:=ListPlot[Transpose[{mu, ARE}], PlotJoined -> True,
 PlotRange -> {{0, mu[[7]]}, {0, 8}}, AxesLabel -> {"Gamma mean",
 "ARE"}]
\end{sourcecode}
\end{enumerate}
\end{example}

\needspace{5\baselineskip}\medskip\noindent\textit{\hyperref[sec:9]{第9章}}
\begin{example}[卡方矩母函数的极限]\label{ex:12.6.8}
计算\hyperref[exercise:9.30]{习题9.30(b)}中的极限分布需要小心，但在Mathematica中相当直接.
\begin{enumerate}[label=(\arabic*)]
\item 先计算自由度为$n$的卡方随机变量的矩母函数.（当然，这一步其实不必要.）\footnote{译注：此处Out[3]中外部的$2^{n/2}$因子和中间项的正指数均与卡方矩母函数不符，原页照录.满足条件时应为$(1-2t)^{-n/2}$；下面作者定义M[t\_]时使用了这个正确形式.}
\begin{sourcecode}
In[1]:= Needs["Statistics`ContinuousDistributions`"]
f[x_] = PDF[ChiSquareDistribution[n], x];
Integrate[Exp[t*x]*f[x], {x, 0, Infinity}]
\end{sourcecode}
\begin{equation*}\begin{aligned}
\textbf{Out[3]}={}&\frac{2^{n/2}}{\operatorname{Gamma}[n/2]}\operatorname{If}\Biggl[\operatorname{Re}[n]>0\mathbin{\&\&}\operatorname{Re}[t]<\frac12,\\
&\left(\frac12-t\right)^{n/2}\operatorname{Gamma}[n/2],\int_0^\infty e^{-x/2+tx}x^{-1+n/2}\,dx\Biggr].
\end{aligned}\end{equation*}

\item 因为检验条件得到满足，$\chi_n^2$的矩母函数就是中间项.现在求$\frac{\chi_n^2-n}{\sqrt{2n}}$的矩母函数的极限.
\begin{sourcecode}
In[4]:M[t_] = (1 - 2*t)^(-n/2);
Limit[Exp[-n*t/Sqrt[2*n]]*M[t/Sqrt[2*n]], n -> Infinity]
\end{sourcecode}
$$\textbf{Out[5]}=e^{t^2/2}.$$
\end{enumerate}
\end{example}


\endgroup



% 原PDF534–539／印刷503–508；全部原页已目视核对。
\clearpage\section*{常见分布表}\phantomsection\label{app:distributions}
\addcontentsline{toc}{section}{常见分布表}
\begingroup
\newcommand{\distheading}[1]{\par\medskip\needspace{10\baselineskip}\noindent\rule{\linewidth}{.3pt}\par\smallskip\noindent\textbf{#1}\par\smallskip}
\sourcepageanchor{534}
\subsection*{离散分布}
\distheading{伯努利分布（Bernoulli）$(p)$}
概率质量函数：$P(X=x\mid p)=p^x(1-p)^{1-x}$；$x=0,1$；$0\le p\le1$.
$$\text{均值与方差：}\quad EX=p,\qquad\operatorname{Var}X=p(1-p).$$
$$\text{矩母函数：}\quad M_X(t)=(1-p)+pe^t.$$

\distheading{二项分布（Binomial）$(n,p)$}
概率质量函数：$P(X=x\mid n,p)=\binom nx p^x(1-p)^{n-x}$；$x=0,1,2,\ldots,n$；$0\le p\le1$.
$$\text{均值与方差：}\quad EX=np,\qquad\operatorname{Var}X=np(1-p).$$
$$\text{矩母函数：}\quad M_X(t)=[pe^t+(1-p)]^n.$$
注：与二项式定理（\autoref{thm:3.2.2}）有关.多项分布（\autoref{def:4.6.2}）是二项分布的多元版本.

\distheading{离散均匀分布（Discrete Uniform）}
概率质量函数：$P(X=x\mid N)=\frac1N$；$x=1,2,\ldots,N$；$N=1,2,\ldots$.
$$\text{均值与方差：}\quad EX=\frac{N+1}{2},\qquad\operatorname{Var}X=\frac{(N+1)(N-1)}{12}.$$
$$\text{矩母函数：}\quad M_X(t)=\frac1N\sum_{i=1}^N e^{it}.$$

\distheading{几何分布（Geometric）$(p)$}
概率质量函数：$P(X=x\mid p)=p(1-p)^{x-1}$；$x=1,2,\ldots$；$0\le p\le1$.\footnote{译注：原表允许$p=0$，但在上述正整数支持上，此时总质量为0；作为通常的几何概率分布须有$0<p\le1$.下面的无记忆等式还需$s\ge t$为相应整数且条件事件概率为正.}
$$\text{均值与方差：}\quad EX=\frac1p,\qquad\operatorname{Var}X=\frac{1-p}{p^2}.$$
$$\text{矩母函数：}\quad M_X(t)=\frac{pe^t}{1-(1-p)e^t},\qquad t<-\log(1-p).$$
注：$Y=X-1$服从负二项$(1,p)$分布.该分布具有无记忆性：$P(X>s\mid X>t)=P(X>s-t)$.

\sourcepageanchor{535}
\distheading{超几何分布（Hypergeometric）}
\begin{equation*}\begin{aligned}
\text{概率质量函数：}\quad&P(X=x\mid N,M,K)=\frac{\binom Mx\binom{N-M}{K-x}}{\binom NK};\quad x=0,1,2,\ldots,K;\\
&M-(N-K)\le x\le M;\quad N,M,K\ge0.
\end{aligned}\end{equation*}
$$\text{均值与方差：}\quad EX=\frac{KM}{N},\qquad\operatorname{Var}X=\frac{KM}{N}\frac{(N-M)(N-K)}{N(N-1)}.$$
注：若$K\ll M$且$K\ll N$，则范围$x=0,1,2,\ldots,K$是适当的.\footnote{译注：原表只写$N,M,K\ge0$，通常还要求它们为整数、$M,K\le N$；所列方差公式在$N=1$等退化情形须另作解释.}

\distheading{负二项分布（Negative binomial）$(r,p)$}
概率质量函数：$P(X=x\mid r,p)=\binom{r+x-1}{x}p^r(1-p)^x$；$x=0,1,\ldots$；$0\le p\le1$.\footnote{译注：原表对$p$的端点约定有与几何分布相同的问题；通常取$r>0$、$0<p\le1$，按失败次数解释时$r$为正整数.}
$$\text{均值与方差：}\quad EX=\frac{r(1-p)}p,\qquad\operatorname{Var}X=\frac{r(1-p)}{p^2}.$$
$$\text{矩母函数：}\quad M_X(t)=\left(\frac{p}{1-(1-p)e^t}\right)^r,\qquad t<-\log(1-p).$$
注：概率质量函数的另一种形式为$P(Y=y\mid r,p)=\binom{y-1}{r-1}p^r(1-p)^{y-r}$，$y=r,r+1,\ldots$.该随机变量满足$Y=X+r$.负二项分布可以由泊松分布的伽马混合导出（见\autoref{exercise:4.32}\mbox{）.}

\distheading{泊松分布（Poisson）$(\lambda)$}
概率质量函数：$P(X=x\mid\lambda)=e^{-\lambda}\lambda^x/x!$；$x=0,1,\ldots$；$0\le\lambda<\infty$.
$$\text{均值与方差：}\quad EX=\lambda,\qquad\operatorname{Var}X=\lambda.$$
$$\text{矩母函数：}\quad M_X(t)=e^{\lambda(e^t-1)}.$$

\subsection*{连续分布}
\distheading{贝塔分布（Beta）$(\alpha,\beta)$}
概率密度函数：$f(x\mid\alpha,\beta)=\frac1{B(\alpha,\beta)}x^{\alpha-1}(1-x)^{\beta-1}$，$0\le x\le1$，$\alpha>0$，$\beta>0$.
$$\text{均值与方差：}\quad EX=\frac\alpha{\alpha+\beta},\qquad\operatorname{Var}X=\frac{\alpha\beta}{(\alpha+\beta)^2(\alpha+\beta+1)}.$$
$$\text{矩母函数：}\quad M_X(t)=1+\sum_{k=1}^{\infty}\left(\prod_{r=0}^{k-1}\frac{\alpha+r}{\alpha+\beta+r}\right)\frac{t^k}{k!}.$$
注：贝塔密度中的常数可以用伽马函数定义，$B(\alpha,\beta)=\frac{\Gamma(\alpha)\Gamma(\beta)}{\Gamma(\alpha+\beta)}$.式\eqref{eq:3.3.18}给出了矩的一般表达式.

\sourcepageanchor{536}
\distheading{柯西分布（Cauchy）$(\theta,\sigma)$}
$$f(x\mid\theta,\sigma)=\frac1{\pi\sigma}\frac1{1+((x-\theta)/\sigma)^2},\quad -\infty<x<\infty,\quad -\infty<\theta<\infty,\quad\sigma>0.$$
均值与方差：不存在.\qquad 矩母函数：不存在.

注：是自由度等于1的Student的$t$分布的特例.此外，若$X$与$Y$相互独立且均服从$n(0,1)$，则$X/Y$服从柯西分布.\footnote{译注：自由度为 $1$ 的通常 $t$ 分布和所写比值 $X/Y$ 对应的是标准柯西分布，即 $\theta=0,\sigma=1$. 表中一般的 $(\theta,\sigma)$ 分布须作位置尺度变换 $\theta+\sigma X/Y$.}

\distheading{卡方分布（Chi squared）}
$$f(x\mid p)=\frac1{\Gamma(p/2)2^{p/2}}x^{(p/2)-1}e^{-x/2};\quad0\le x<\infty;\quad p=1,2,\ldots.$$
$$\text{均值与方差：}\quad EX=p,\qquad\operatorname{Var}X=2p.$$
$$\text{矩母函数：}\quad M_X(t)=\left(\frac1{1-2t}\right)^{p/2},\qquad t<\frac12.$$
注：是伽马分布的特例.

\distheading{双指数分布（Double exponential）$(\mu,\sigma)$}
$$f(x\mid\mu,\sigma)=\frac1{2\sigma}e^{-|x-\mu|/\sigma},\quad-\infty<x<\infty,\quad-\infty<\mu<\infty,\quad\sigma>0.$$
$$\text{均值与方差：}\quad EX=\mu,\qquad\operatorname{Var}X=2\sigma^2.$$
$$\text{矩母函数：}\quad M_X(t)=\frac{e^{\mu t}}{1-(\sigma t)^2},\qquad |t|<\frac1\sigma.$$
注：也称为拉普拉斯分布（Laplace distribution）.

\distheading{指数分布（Exponential）$(\beta)$}
概率密度函数：$f(x\mid\beta)=\frac1\beta e^{-x/\beta}$，$0\le x<\infty$，$\beta>0$.
$$\text{均值与方差：}\quad EX=\beta,\qquad\operatorname{Var}X=\beta^2.$$
$$\text{矩母函数：}\quad M_X(t)=\frac1{1-\beta t},\qquad t<\frac1\beta.$$
注：是伽马分布的特例，具有无记忆性.有许多特殊变换：$Y=X^{1/\gamma}$服从威布尔分布，$Y=\sqrt{2X/\beta}$服从瑞利分布（Rayleigh），$Y=\alpha-\gamma\log(X/\beta)$服从耿贝尔分布\mbox{（Gumbel）.}

\sourcepageanchor{537}
\distheading{$F$分布}
\begin{equation*}\begin{aligned}
f(x\mid\nu_1,\nu_2)&=\frac{\Gamma((\nu_1+\nu_2)/2)}{\Gamma(\nu_1/2)\Gamma(\nu_2/2)}\left(\frac{\nu_1}{\nu_2}\right)^{\nu_1/2}\frac{x^{(\nu_1-2)/2}}{(1+(\nu_1/\nu_2)x)^{(\nu_1+\nu_2)/2}};\\
&\hspace{4em}0\le x<\infty;\qquad\nu_1,\nu_2=1,\ldots.
\end{aligned}\end{equation*}
均值与方差：
$$EX=\frac{\nu_2}{\nu_2-2},\quad\nu_2>2;\qquad\operatorname{Var}X=2\left(\frac{\nu_2}{\nu_2-2}\right)^2\frac{\nu_1+\nu_2-2}{\nu_1(\nu_2-4)},\quad\nu_2>4.$$
矩（矩母函数不存在）：
$$EX^n=\frac{\Gamma((\nu_1+2n)/2)\Gamma((\nu_2-2n)/2)}{\Gamma(\nu_1/2)\Gamma(\nu_2/2)}\left(\frac{\nu_2}{\nu_1}\right)^n,\qquad n<\frac{\nu_2}{2}.$$
注：与卡方分布有关（$F_{\nu_1,\nu_2}=(\chi_{\nu_1}^2/\nu_1)/(\chi_{\nu_2}^2/\nu_2)$，其中两个卡方变量相互独立），也与$t$分布有关（$F_{1,\nu}=t_\nu^2$）.

\distheading{伽马分布（Gamma）$(\alpha,\beta)$}
$$f(x\mid\alpha,\beta)=\frac1{\Gamma(\alpha)\beta^\alpha}x^{\alpha-1}e^{-x/\beta},\qquad0\le x<\infty,\quad\alpha,\beta>0.$$
$$\text{均值与方差：}\quad EX=\alpha\beta,\qquad\operatorname{Var}X=\alpha\beta^2.$$
$$\text{矩母函数：}\quad M_X(t)=\left(\frac1{1-\beta t}\right)^\alpha,\qquad t<\frac1\beta.$$
注：一些特例为指数分布（$\alpha=1$）与卡方分布（$\alpha=p/2,\beta=2$）.若$\alpha=3/2$，则$Y=\sqrt{X/\beta}$服从麦克斯韦分布（Maxwell）.$Y=1/X$服从逆伽马分布（inverted gamma distribution）.还可以与泊松分布联系起来（\autoref{ex:3.3.1}）.

\distheading{逻辑分布（Logistic）$(\mu,\beta)$}
$$f(x\mid\mu,\beta)=\frac1\beta\frac{e^{-(x-\mu)/\beta}}{[1+e^{-(x-\mu)/\beta}]^2},\quad-\infty<x<\infty,\quad-\infty<\mu<\infty,\quad\beta>0.$$
$$\text{均值与方差：}\quad EX=\mu,\qquad\operatorname{Var}X=\frac{\pi^2\beta^2}{3}.$$
$$\text{矩母函数：}\quad M_X(t)=e^{\mu t}\Gamma(1-\beta t)\Gamma(1+\beta t),\qquad |t|<\frac1\beta.$$
注：累积分布函数为$F(x\mid\mu,\beta)=\frac1{1+e^{-(x-\mu)/\beta}}$.

\distheading{对数正态分布（Lognormal）$(\mu,\sigma^2)$}
$$f(x\mid\mu,\sigma^2)=\frac1{\sqrt{2\pi}\sigma}\frac{e^{-(\log x-\mu)^2/(2\sigma^2)}}x,\quad0\le x<\infty,\quad-\infty<\mu<\infty,\quad\sigma>0.$$
$$\text{均值与方差：}\quad EX=e^{\mu+\sigma^2/2},\qquad\operatorname{Var}X=e^{2(\mu+\sigma^2)}-e^{2\mu+\sigma^2}.$$
$$\text{矩（矩母函数不存在）：}\quad EX^n=e^{n\mu+n^2\sigma^2/2}.$$
注：\autoref{ex:2.3.10}给出了具有相同各阶矩的另一个分布.

\sourcepageanchor{538}
\distheading{正态分布（Normal）$(\mu,\sigma^2)$}
$$f(x\mid\mu,\sigma^2)=\frac1{\sqrt{2\pi}\sigma}e^{-(x-\mu)^2/(2\sigma^2)},\quad-\infty<x<\infty,\quad-\infty<\mu<\infty,\quad\sigma>0.$$
$$\text{均值与方差：}\quad EX=\mu,\qquad\operatorname{Var}X=\sigma^2.$$
$$\text{矩母函数：}\quad M_X(t)=e^{\mu t+\sigma^2t^2/2}.$$
注：有时称为高斯分布（Gaussian distribution）.

\distheading{帕累托分布（Pareto）$(\alpha,\beta)$}
$$f(x\mid\alpha,\beta)=\frac{\beta\alpha^\beta}{x^{\beta+1}},\qquad\alpha<x<\infty,\quad\alpha>0,\quad\beta>0.$$
$$\text{均值与方差：}\quad EX=\frac{\beta\alpha}{\beta-1},\quad\beta>1;\qquad\operatorname{Var}X=\frac{\beta\alpha^2}{(\beta-1)^2(\beta-2)},\quad\beta>2.$$
矩母函数：不存在.

\distheading{$t$分布}
$$f(x\mid\nu)=\frac{\Gamma((\nu+1)/2)}{\Gamma(\nu/2)}\frac1{\sqrt{\nu\pi}}\frac1{(1+x^2/\nu)^{(\nu+1)/2}},\quad-\infty<x<\infty,\quad\nu=1,\ldots.$$
$$\text{均值与方差：}\quad EX=0,\quad\nu>1;\qquad\operatorname{Var}X=\frac\nu{\nu-2},\quad\nu>2.$$
矩（矩母函数不存在）：
\begin{equation*}\begin{aligned}
EX^n&=\frac{\Gamma((n+1)/2)\Gamma((\nu-n)/2)}{\sqrt\pi\Gamma(\nu/2)}\nu^{n/2},&&\text{若$n<\nu$且为偶数};\\
EX^n&=0,&&\text{若$n<\nu$且为奇数}.
\end{aligned}\end{equation*}
注：与$F$分布有关（$F_{1,\nu}=t_\nu^2$）.

\distheading{均匀分布（Uniform）$(a,b)$}
概率密度函数：$f(x\mid a,b)=\frac1{b-a}$，$a\le x\le b$.
$$\text{均值与方差：}\quad EX=\frac{b+a}{2},\qquad\operatorname{Var}X=\frac{(b-a)^2}{12}.$$
$$\text{矩母函数：}\quad M_X(t)=\frac{e^{bt}-e^{at}}{(b-a)t}.$$
注：若$a=0$且$b=1$，这是贝塔分布（$\alpha=\beta=1$）的特例.\footnote{译注：此处密度需$a<b$；矩母函数在$t=0$处按连续延拓取1.原表未显写这两项约定.}

\sourcepageanchor{539}
\distheading{威布尔分布（Weibull）$(\gamma,\beta)$}
$$f(x\mid\gamma,\beta)=\frac\gamma\beta x^{\gamma-1}e^{-x^\gamma/\beta},\qquad0\le x<\infty,\quad\gamma>0,\quad\beta>0.$$
均值与方差：
$$EX=\beta^{1/\gamma}\Gamma\!\left(1+\frac1\gamma\right),\qquad\operatorname{Var}X=\beta^{2/\gamma}\left[\Gamma\!\left(1+\frac2\gamma\right)-\Gamma^2\!\left(1+\frac1\gamma\right)\right].$$
矩：\footnote{译注：原页的矩公式漏印了伽马函数符号；对非负整数$n$应为$EX^n=\beta^{n/\gamma}\Gamma(1+n/\gamma)$，这也与上方的一阶矩公式一致.}
$$EX^n=\beta^{n/\gamma}\left(1+\frac n\gamma\right).$$

注：矩母函数仅在$\gamma\ge1$时存在，其形式并不十分有用.一个特例是指数分布\mbox{（$\gamma=1$）.}

\clearpage
\begin{center}
\includegraphics[width=.94\textwidth,height=.80\textheight,keepaspectratio]{assets/backmatter/distribution-relations.pdf}
\end{center}
\noindent\textbf{常见分布之间的关系.}实线表示变换与特例，虚线表示极限.据\hyperref[bib:source195]{Leemis（1986）}改编.

\smallskip\noindent\small 图中名称中英对照：Geometric—几何；Discrete uniform—离散均匀；Negative binomial—负二项；Beta-binomial—贝塔二项；Hypergeometric—超几何；Poisson—泊松；Binomial—二项；Bernoulli—伯努利；Lognormal—对数正态；Normal—正态；Beta—贝塔；Gamma—伽马；Chi-squared—卡方；Uniform—均匀；Cauchy—柯西；Exponential—指数；Weibull—威布尔；Double exponential—双指数.
\endgroup



\clearpage\renewcommand{\refname}{参考文献}
\phantomsection\label{app:references}\addcontentsline{toc}{section}{参考文献}
\begin{thebibliography}{338}
\small\setlength{\itemsep}{3pt plus 1pt minus 1pt}
\sourcepageanchor{540}
\bibitem[1]{source1}\label{bib:source1} Agresti, A. (1990). Categorical Data Analysis. New York: Wiley.
\bibitem[2]{source2}\label{bib:source2} Agresti, A. and Caffo, B. (2000). Simple and Effective Confidence Intervals for Proportions and Differences of Proportions Result from Adding Two Successes and Two Failures. Amer. Statist. 54 280–288.
\bibitem[3]{source3}\label{bib:source3} Agresti, A. and Coull, B. A. (1998). Approximate Is Better Than “Exact” for Interval Estimation of Binomial Proportions. Amer. Statist. 52 119–126.
\bibitem[4]{source4}\label{bib:source4} Anderson, T. W. (1976). Estimation of Linear Functional Relationships: Approximate Distributions and Connection with Simultaneous Equations in Econometrics (with discussion). J. Roy. Statist. Soc. Ser. B 38 1–36.
\bibitem[5]{source5}\label{bib:source5} Anderson, T. W. (1984). An Introduction to Multivariate Statistical Analysis, 2nd edition. New York: Wiley.
\bibitem[6]{source6}\label{bib:source6} Balanda, K. P and MacGillivray, H. L. (1988). Kurtosis: A Critical Review. Amer. Statist. 42 111–119.
\bibitem[7]{source7}\label{bib:source7} Bar-Lev, S. K. and Enis, P. (1988). On the Classical Choice of Variance Stabilizing Transformations and an Application for a Poisson Variate. Biometrika 75 803–804.
\bibitem[8]{source8}\label{bib:source8} Bar-Lev, S. K. and Enis, P. (1990). On the Construction of Classes of Variance Stabilizing Transformations. Statist. Prob. Let. 10 95–100.
\bibitem[9]{source9}\label{bib:source9} Barlow, R. and Proschan, F. (1975). Statistical Theory of Life Testing. New York: Holt, Rinehart and Winston.
\bibitem[10]{source10}\label{bib:source10} Barnard, G. A. (1949). Statistical Inference (with discussion). J. Roy. Statist. Soc. Ser. B 11, 115–139.
\bibitem[11]{source11}\label{bib:source11} Barnard, G. A. (1980). Pivotal Inference and the Bayesian Controversy (with discussion). Bayesian Statistics (J. M. Bernardo, M. H. DeGroot, D. V. Lindley, and A. F. M. Smith, eds.). Valencia: University Press.
\bibitem[12]{source12}\label{bib:source12} Barr, D. R. and Zehna, P. W. (1983). Probability: Modeling Uncertainty. Reading, MA: Addison-Wesley.
\bibitem[13]{source13}\label{bib:source13} Basu, D. (1959). The Family of Ancillary Statistics. Sankhya, Ser. A 21 247–256.
\bibitem[14]{source14}\label{bib:source14} Bechhofer, R. E. (1954). A Single-Sample Multiple Decision Procedure for Ranking Means of Normal Populations with Known Variances. Ann. Math. Statist. 25 16–39.
\bibitem[15]{source15}\label{bib:source15} Behboodian, J. (1990). Examples of Uncorrelated Dependent Random Variables Using a Bivariate Mixture. Amer. Statist. 44 218.
\bibitem[16]{source16}\label{bib:source16} Berg, C. (1988). The Cube of a Normal Distribution Is Indeterminate. Ann. Prob. 16 910–913.
\bibitem[17]{source17}\label{bib:source17} Berger, J. O. (1984). The Robust Bayesian Viewpoint (with discussion). Robustness of Bayesian Analysis (J. Kadane, ed.), 63–144. Amsterdam: North-Holland.
\bibitem[18]{source18}\label{bib:source18} Berger, J. O. (1985). Statistical Decision Theory and Bayesian Analysis, 2nd edition. New York: Springer-Verlag.
\bibitem[19]{source19}\label{bib:source19} Berger, J. O. (1990). Robust Bayesian Analysis: Sensitivity to the Prior. J. Statist. Plan. Inf. 25 303–328.
\sourcepageanchor{541}
\bibitem[20]{source20}\label{bib:source20} Berger, J. O. (1994). An Overview of Robust Bayesian Analysis (with discussion). Test 3 5–124.
\bibitem[21]{source21}\label{bib:source21} Berger, J. O. and Sellke, T. (1987). Testing a Point Null Hypothesis: The Irreconcilability of P Values and Evidence (with discussion). J. Amer. Statist. Assoc. 82 112–122.
\bibitem[22]{source22}\label{bib:source22} Berger, J. O. and Wolpert, R. W. (1984). The Likelihood Principle. Institute of Mathematical Statistics Lecture Notes—Monograph Series. Hayward, CA: IMS.
\bibitem[23]{source23}\label{bib:source23} Berger, R. L. (1982). Multiparameter Hypothesis Testing and Acceptance Sampling. Technometrics 24 295–300.
\bibitem[24]{source24}\label{bib:source24} Berger, R. L. (1996). More Powerful Tests from Confidence Interval P Values. Amer. Statist. 50 314–318.
\bibitem[25]{source25}\label{bib:source25} Berger, R. L. and Boos, D. D. (1994). P Values Maximized over a Confidence Set for the Nuisance Parameter. J. Amer. Statist. Assoc. 89 1012–1016.
\bibitem[26]{source26}\label{bib:source26} Berger, R. L. and Casella, G. (1992). Deriving Generalized Means as Least Squares and Maximum Likelihood Estimates. Amer. Statist. 46 279–282.
\bibitem[27]{source27}\label{bib:source27} Berger, R. L. and Hsu, J. C. (1996). Bioequivalence Trials, Intersection-Union Tests and Equivalence Confidence Sets (with discussion). Statist. Sci. 11 283–319.
\bibitem[28]{source28}\label{bib:source28} Berkson, J. (1950). Are There Two Regressions? J. Amer. Statist. Assoc. 45 164–180.
\bibitem[29]{source29}\label{bib:source29} Bernardo, J. M and Smith, A. F. M. (1994). Bayesian Theory. New York: Wiley.
\bibitem[30]{source30}\label{bib:source30} Betteley, I. G. (1977). The Addition Law for Expectations. Amer. Statist. 31 33–35.
\bibitem[31]{source31}\label{bib:source31} Bickel, P. J. and Doksum, K. A. (1981). An Analysis of Transformations Revisited. J. Amer. Statist. Assoc. 76 296–311.
\bibitem[32]{source32}\label{bib:source32} Bickel, P. J. and Mallows, C. L. (1988). A Note on Unbiased Bayes Estimates. Amer. Statist. 42 132–134.
\bibitem[33]{source33}\label{bib:source33} Billingsley P. (1995). Probability and Measure, 3rd edition. New York: Wiley.
\bibitem[34]{source34}\label{bib:source34} Binder, D. A. (1993). Letter to the Editor about the “Exchange Paradox”. Amer. Statist. 47 160.
\bibitem[35]{source35}\label{bib:source35} Birnbaum, A. (1962). On the Foundations of Statistical Inference (with discussion). J. Amer. Statist. Assoc. 57 269–306.
\bibitem[36]{source36}\label{bib:source36} Blachman, N. M. (1996) Letter to the Editor about the “Exchange Paradox”. Amer. Statist. 50 98–99.
\bibitem[37]{source37}\label{bib:source37} Blaker, H. (2000). Confidence Curves and Improved Exact Confidence Intervals for Discrete Distributions. Can. J. Statist. 28 783–798.
\bibitem[38]{source38}\label{bib:source38} Bloch, D. A. and Moses, L. E. (1988). Nonoptimally Weighted Least Squares. Amer. Statist. 42 50–53.
\bibitem[39]{source39}\label{bib:source39} Blyth, C. R. (1986). Approximate Binomial Confidence Limits. J. Amer. Statist. Assoc. 81 843–855; correction 84 636.
\bibitem[40]{source40}\label{bib:source40} Blyth, C. R. and Still, H. A. (1983). Binomial Confidence Intervals. J. Amer. Statist. Assoc. 78 108–116.
\bibitem[41]{source41}\label{bib:source41} Boos, D. D. (1992). On Generalized Score Tests (Com: 93V47 P311-312). Amer. Statist. 46 327–333.
\bibitem[42]{source42}\label{bib:source42} Boos, D. D. and Hughes-Oliver, J. M. (1998). Applications of Basu’s Theorem. Amer. Statist. 52 218–221.
\bibitem[43]{source43}\label{bib:source43} Box, G. E. P. (1954). Some Theorems on Quadratic Forms Applied in the Study of Analysis of Variances Problems I: Effect of Inequality of Variance in the One-Way Classification. Ann. Math. Statist. 25 290–302.
\sourcepageanchor{542}
\bibitem[44]{source44}\label{bib:source44} Box, G. E. P. and Cox, D. R. (1964). An Analysis of Transformations (with discussion). J. Roy. Statist. Soc. Ser. B 26 211–252.
\bibitem[45]{source45}\label{bib:source45} Box, G. E. P. and Cox, D. R. (1982). An Analysis of Transformations Revisited, Rebutted. J. Amer. Statist. Assoc. 77 209–210.
\bibitem[46]{source46}\label{bib:source46} Box, G.E.P. and Muller, M. (1958) A Note on the Generation of Random Normal Variates. Ann. Math. Statist. 29 610–611.
\bibitem[47]{source47}\label{bib:source47} Boyles, R. A. (1983). On the convergence of the EM algorithm. J. Roy. Statist. Soc. Ser. B 45 47–50.
\bibitem[48]{source48}\label{bib:source48} Brewster, J. F. and Zidek, J. V. (1974). Improving on Equivariant Estimators. Ann. Statist. 2 21–38.
\bibitem[49]{source49}\label{bib:source49} Brown, D. and Rothery, P. (1993). Models in Biology: Mathematics, Statistics and Computing. New York: Wiley
\bibitem[50]{source50}\label{bib:source50} Brown L. D. (1986). Fundamentals of Statistical Exponential Families with Applications to Statistical Decision Theory. Institute of Mathematical Statistics Lecture Notes—Monograph Series. Hayward, CA: IMS.
\bibitem[51]{source51}\label{bib:source51} Brown L. D. (1990a). Comment on the Paper by Maatta and Casella. Statist. Sci. 5 103–106.
\bibitem[52]{source52}\label{bib:source52} Brown, L. D. (1990b). An Ancillarity Paradox Which Occurs in Multiple Linear Regression (with discussion). Ann. Statist. 18 471–492.
\bibitem[53]{source53}\label{bib:source53} Brown, L. D., Casella, G. and Hwang, J. T. G. (1995). Optimal Confidence Sets, Bioequivalence, and the Limaçon of Pascal. J. Amer. Statist. Assoc. 90 880–889.
\bibitem[54]{source54}\label{bib:source54} Brown, L. D. and Purves, R. (1973). Measurable Selections of Extrema. Ann. Statist. 1 902–912.
\bibitem[55]{source55}\label{bib:source55} Brown, P. J. and Fuller, W. A. (1991). Statistical Analysis of Measurement Error Models and Applications. Providence, RI: American Mathematical Society.
\bibitem[56]{source56}\label{bib:source56} Buehler, R. J. (1982). Some Ancillary Statistics and Their Properties (with discussion). J. Amer. Statist. Assoc. 77 581–594.
\bibitem[57]{source57}\label{bib:source57} Carlin, B.P. and Louis, T.A. (1996). Bayes and Empirical Bayes Methods for Data Analysis. London: Chapman and Hall.
\bibitem[58]{source58}\label{bib:source58} Carmer, S. G. and Walker, W. M. (1982). Baby Bear’s Dilemma: A Statistical Tale. Agronomy Journal 74 122–124.
\bibitem[59]{source59}\label{bib:source59} Carroll, R. J., Gallo, P., and Gleser, L. J. (1985). Comparison of Least Squares and Errorsin-Variables Regression, with Special Reference to Randomized Analysis of Covariance. J. Amer. Statist. Assoc. 80 929–932.
\bibitem[60]{source60}\label{bib:source60} Carroll, R. J. and Ruppert, D. (1985). Transformations: A Robust Analysis. Technometrics 27 1–12.
\bibitem[61]{source61}\label{bib:source61} Carroll, R. J. and Ruppert, D. (1988). Transformation and Weighting in Regression. London: Chapman and Hall.
\bibitem[62]{source62}\label{bib:source62} Carroll, R. J., Ruppert, D., and Stefanski, L. A. (1995). Measurement Error in Nonlinear Models. London: Chapman and Hall.
\bibitem[63]{source63}\label{bib:source63} Casella, G. (1985). An Introduction to Empirical Bayes Data Analysis. Amer. Statist. 39 83–87.
\bibitem[64]{source64}\label{bib:source64} Casella, G. (1986). Refining Binomial Confidence Intervals. Can. J. Statist. 14 113–129.
\bibitem[65]{source65}\label{bib:source65} Casella, G. (1992). Illustrating Empirical Bayes Methods. Chemolab 16 107–125.
\bibitem[66]{source66}\label{bib:source66} Casella, G. and Berger, R. L. (1987). Reconciling Bayesian and Frequentist Evidence in the One-Sided Testing Problem (with discussion). J. Amer. Statist. Assoc. 82 106–111.
\sourcepageanchor{543}
\bibitem[67]{source67}\label{bib:source67} Casella, G. and George, E. I. (1992). Explaining the Gibbs Sampler. Amer. Statist. 46 167–174.
\bibitem[68]{source68}\label{bib:source68} Casella, G. and Hwang, J. T. (1987). Employing Vague Prior Information in the Construction of Confidence Sets. J. Mult. Analysis. 21 79–104.
\bibitem[69]{source69}\label{bib:source69} Casella, G. and Robert, C. (1989). Refining Poisson Confidence Intervals. Can. J. Statist. 17 45–57.
\bibitem[70]{source70}\label{bib:source70} Casella, G. and Robert, C. P. (1996). Rao-Blackwellisation of Sampling Schemes. Biometrika 83 81–94.
\bibitem[71]{source71}\label{bib:source71} Casella, G. and Strawderman, W. E. (1980). Confidence Bands for Linear Regression with Restricted Predictor Variables. J. Amer. Statist. Assoc. 75 862–868.
\bibitem[72]{source72}\label{bib:source72} Chapman, D. G. and Robbins, H. (1951). Minimum Variance Estimation Without Regularity Assumptions. Ann. Math. Statist. 22 581–586.
\bibitem[73]{source73}\label{bib:source73} Chib, S. and Greenberg, E. (1995). Understanding the Metropolis–Hastings Algorithm. Amer. Statist. 49 327–335.
\bibitem[74]{source74}\label{bib:source74} Chikkara, R. S. and Folks, J. L. (1989). The Inverse Gaussian Distribution: Theory, Methodology, and Applications. New York: Marcel Dekker.
\bibitem[75]{source75}\label{bib:source75} Christensen, R. (1996). Plane Answers to Complex Questions. The Theory of Linear Models, 2nd edition. New York: Springer-Verlag.
\bibitem[76]{source76}\label{bib:source76} Christensen, R. and Utts, J. (1992). Bayesian Resolution of the “Exchange Paradox.” Amer. Statist. 46 274–278.
\bibitem[77]{source77}\label{bib:source77} Chun, Y. H. (1999). On the Information Economics Approach to the Generalized Game Show Problem. Amer. Statist. 53 43–51.
\bibitem[78]{source78}\label{bib:source78} Chung, K. L. (1974). A Course in Probability Theory. New York: Academic Press.
\bibitem[79]{source79}\label{bib:source79} Clopper, C. J. and Pearson, E. S. (1934). The Use of Confidence or Fiducial Limits Illustrated in the Case of the Binomial. Biometrika 26 404–413. (Also in The Selected Papers of E. S. Pearson, New York: Cambridge University Press, 1966.)
\bibitem[80]{source80}\label{bib:source80} Cochran, W. G. (1934). The Distribution of Quadratic Forms in a Normal System with Applications to the Analysis of Variance. Proceedings of the Cambridge Philosophical Society 30 178–191.
\bibitem[81]{source81}\label{bib:source81} Cochran, W. G. (1968). Errors of Measurement in Statistics. Technometrics 10 637–666.
\bibitem[82]{source82}\label{bib:source82} Cochran, W. G. and Cox, G. M. (1957). Experimental Designs, 2nd edition. New York: Wiley.
\bibitem[83]{source83}\label{bib:source83} Cornfield, J. (1969). The Bayesian Outlook and Its Application (with discussion). Biometrika 25 617–657.
\bibitem[84]{source84}\label{bib:source84} Cox, D. R. (1958). Some Problems Connected with Statistical Inference. Ann. Math. Statist. 29 357–372.
\bibitem[85]{source85}\label{bib:source85} Cox, D. R. (1971). The Choice Between Ancillary Statistics. J. Roy. Statist. Soc. Ser. B 3, 251–255.
\bibitem[86]{source86}\label{bib:source86} Cox, D. R. and Hinkley, D. V. (1974). Theoretical Statistics. London: Chapman and Hall.
\bibitem[87]{source87}\label{bib:source87} Creasy, M. A. (1956). Confidence Limits for the Gradient in the Linear Functional Relationship. J. Roy. Statist. Soc. Ser. B 18 65–69.
\bibitem[88]{source88}\label{bib:source88} Crow, E. L. (1956). Confidence Intervals for a Proportion. Biometrika 43 423–425.
\bibitem[89]{source89}\label{bib:source89} Crow, E. L. and Gardner, R. S. (1959). Confidence Intervals for the Expectation of a Poisson Variable. Biometrika 46 441–453.
\bibitem[90]{source90}\label{bib:source90} Curtiss, J. H. (1943). On Transformations Used in the Analysis of Variance. Ann. Math. Statist. 14 107–122.
\sourcepageanchor{544}
\bibitem[91]{source91}\label{bib:source91} Dalal, S. R., Fowlkes, E. B., and Hoadley, B. (1989). Risk Analysis of the Space Shuttle: Pre-Challenger Prediction of Failure. J. Amer. Statist. Assoc. 84 945–957.
\bibitem[92]{source92}\label{bib:source92} David, H. A. (1985). Bias of $S^2$ Under Dependence. Amer. Statist. 39 201.
\bibitem[93]{source93}\label{bib:source93} Davidson, R. R. and Solomon, D. L. (1974). Moment-Type Estimation in the Exponential Family. Communications in Statistics 3 1101–1108.
\bibitem[94]{source94}\label{bib:source94} Dean, A. and Voss, D. (1999) Design and Analysis of Experiments New York: Springer- Verlag.
\bibitem[95]{source95}\label{bib:source95} deFinetti, B. (1972). Probability, Induction, and Statistics. New York: Wiley.
\bibitem[96]{source96}\label{bib:source96} DeGroot, M. H. (1986). Probability and Statistics, 2nd edition. New York: Addison-Wesley.
\bibitem[97]{source97}\label{bib:source97} Dempster, A. P., Laird, N. M., and Rubin, D. B. (1977). Maximum Likelihood from Incomplete Data via the EM Algorithm. J. Roy. Statist. Soc. Ser. B 39 1–22.
\bibitem[98]{source98}\label{bib:source98} Devroye, L. (1985). Non-Uniform Random Variate Generation. New York: Springer-Verlag.
\bibitem[99]{source99}\label{bib:source99} Diaconis, P. and Holmes, S. (1994). Gray Codes for Randomization Procedures. Statistics and Computing 4 287–302.
\bibitem[100]{source100}\label{bib:source100} Diaconis, P. and Mosteller, F. (1989). Methods for Studying Coincidences. J. Amer. Statist. Assoc. 84 853–861.
\bibitem[101]{source101}\label{bib:source101} Draper, D. (1988). Rank-Based Robust Analysis of Linear Models. I. Exposition and Review (with discussion). Statist. Sci. 3 239–271.
\bibitem[102]{source102}\label{bib:source102} Draper, N. R. and Smith, H. (1998). Applied Regression Analysis, 3rd edition. New York: Wiley.
\bibitem[103]{source103}\label{bib:source103} Durbin, J. (1970). On Birnbaum’s Theorem and the Relation Between Sufficiency, Conditionality, and Likelihood. J. Amer. Statist. Assoc. 65 395–398.
\bibitem[104]{source104}\label{bib:source104} Dynkin, E. B. (1951). Necessary and Sufficient Statistics for a Family of Probability Distributions. English translation in Selected Translations in Mathematical Statistics and Probability 1 (1961), 23–41.
\bibitem[105]{source105}\label{bib:source105} Eberhardt, K. R. and Fligner, M. A. (1977). A Comparison of Two Tests for Equality of Two Proportions. Amer. Statist. 31 151–155.
\bibitem[106]{source106}\label{bib:source106} Efron, B. F. (1979a). Bootstrap Methods: Another Look at the Jackknife. Ann. Statist. 7 1–26.
\bibitem[107]{source107}\label{bib:source107} Efron, B. F. (1979b). Computers and the Theory of Statistics: Thinking the Unthinkable. SIAM Review 21 460–480.
\bibitem[108]{source108}\label{bib:source108} Efron, B. F. (1982). The Jackknife, the Bootstrap, and Other Resampling Plans. Philadelphia: Society for Industrial and Applied Mathematics.
\bibitem[109]{source109}\label{bib:source109} Efron, B. F. (1998). R. A. Fisher in the 21st Century (with discussion). Statist. Sci. 13 95–122.
\bibitem[110]{source110}\label{bib:source110} Efron, B. F. and Hinkley, D. V. (1978). Assessing the Accuracy of the Maximum Likelihood Estimator: Observed Versus Expected Fisher Information. Biometrika 65 457–487.
\bibitem[111]{source111}\label{bib:source111} Efron, B. F. and Morris, C. N. (1972). Limiting the Risk of Bayes and Empirical Bayes Estimators Part II: The Empirical Bayes Case. J. Amer. Statist. Assoc. 67 130–139.
\bibitem[112]{source112}\label{bib:source112} Efron, B. F. and Morris, C. N. (1973). Stein’s Estimation Rule and Its Competitors—An Empirical Bayes Approach. J. Amer. Statist. Assoc. 68 117–130.
\bibitem[113]{source113}\label{bib:source113} Efron, B. F. and Morris, C. N. (1975). Data Analysis Using Stein’s Estimator and Its Generalizations. J. Amer. Statist. Assoc. 70 311–319.
\bibitem[114]{source114}\label{bib:source114} Efron, B. and Tibshirani, R. J. (1993). An Introduction to the Bootstrap. London: Chapman and Hall..
\bibitem[115]{source115}\label{bib:source115} Ellis, S. Instability of Least Squares, Least Absolute Deviation and Least Median of Squares Linear Regression (with discussion). Statist. Sci. 13 337–350.
\sourcepageanchor{545}
\bibitem[116]{source116}\label{bib:source116} Feldman, D. and Fox, M. (1968). Estimation of the Parameter n in the Binomial Distribution. J. Amer. Statist. Assoc. 63 150–158.
\bibitem[117]{source117}\label{bib:source117} Feldstein, M. (1974). Errors in Variables: A Consistent Estimator with Smaller MSE in Finite Samples. J. Amer. Statist. Assoc. 69 990–996.
\bibitem[118]{source118}\label{bib:source118} Feller, W. (1968). An Introduction to Probability Theory and Its Applications, Volume I. New York: Wiley.
\bibitem[119]{source119}\label{bib:source119} Feller, W. (1971). An Introduction to Probability Theory and Its Applications, Volume II. New York: Wiley.
\bibitem[120]{source120}\label{bib:source120} Ferguson, T. S. (1996). A Course in Large Sample Theory. London: Chapman and Hall.
\bibitem[121]{source121}\label{bib:source121} Fieller, E. C. (1954). Some Problems in Interval Estimation. J. Roy. Statist. Soc. Ser. B 16 175–185.
\bibitem[122]{source122}\label{bib:source122} Finch, S. J., Mendell, N. R., and Thode, H. C. (1989). Probabilistic Measures of Adequacy of a Numerical Search for a Global Maximum. J. Amer. Statist. Assoc. 84 1020–1023.
\bibitem[123]{source123}\label{bib:source123} Fisher, R. A. (1925). Theory of Statistical Estimation. Proceedings of the Cambridge Philosophical Society 22 700–725.
\bibitem[124]{source124}\label{bib:source124} Fisher, R. A. (1930). Inverse Probability. Proceedings of the Cambridge Philosophical Society 26 528–535.
\bibitem[125]{source125}\label{bib:source125} Fisher, R. A. (1935). The Fiducial Argument in Statistical Inference. Annals of Eugenics 6 391–398. (Also in R. A. Fisher, Contributions to Mathematical Statistics, New York: Wiley, 1950.)
\bibitem[126]{source126}\label{bib:source126} Fisher, R. A. (1939). The Comparison of Samples with Possibly Unequal Variances. Annals of Eugenics 9 174–180. (Also in R. A. Fisher, Contributions to Mathematical Statistics, New York: Wiley, 1950.)
\bibitem[127]{source127}\label{bib:source127} Fraser, D. A. S. (1968). The Structure of Inference. New York: Wiley.
\bibitem[128]{source128}\label{bib:source128} Fraser, D. A. S. (1979). Inference and Linear Models. New York: McGraw-Hill.
\bibitem[129]{source129}\label{bib:source129} Freedman, D., Pisani, R., Purves, R., and Adhikari, A. (1991). Statistics, 2nd edition. New York: Norton.
\bibitem[130]{source130}\label{bib:source130} Fuller, W. A. (1987). Measurement Error Models. New York: Wiley.
\bibitem[131]{source131}\label{bib:source131} Gafarian, A. V. (1964). Confidence Bands in Straight Line Regression. J. Amer. Statist. Assoc. 59 182–213.
\bibitem[132]{source132}\label{bib:source132} Gardner, M. (1961). The Second Scientific American Book of Mathematical Puzzles and Diversions. New York: Simon \& Schuster.
\bibitem[133]{source133}\label{bib:source133} Garwood, F. (1936). Fiducial Limits for the Poisson Distribution. Biometrika 28 437–442.
\bibitem[134]{source134}\label{bib:source134} Gelfand, A. E. and Smith, A. F. M. (1990). Sampling-Based Approaches to Calculating Marginal Densities. J. Amer. Statist. Assoc. 85 398–409.
\bibitem[135]{source135}\label{bib:source135} Gelman, A. and Meng, X-L. (1991). A Note on Bivariate Distributions That Are Conditionally Normal. Amer. Statist. 45 125–126.
\bibitem[136]{source136}\label{bib:source136} Gelman, A. and Rubin, D.B. (1992). Inference from Iterative Simulation Using Multiple Sequences (with discussion). Statist. Sci. 7 457–511.
\bibitem[137]{source137}\label{bib:source137} Geman, S. and Geman, D. (1984) Stochastic Relaxation, Gibbs Distributions and the Bayesian Restoration of Images. IEEE Trans. Pattern Anal. Mach. Intell. 6 721–741.
\bibitem[138]{source138}\label{bib:source138} Geyer, C.J. and Thompson, E.A. (1992) Constrained Monte Carlo Maximum Likelihood for Dependent Data (with discussion). J. Roy. Statist. Soc. Ser. B 54 657–699.
\bibitem[139]{source139}\label{bib:source139} Ghosh, B. K. (1979). A Comparison of Some Approximate Confidence Intervals for the Binomial Parameter. J. Amer. Statist. Assoc. 74 894–900.
\sourcepageanchor{546}
\bibitem[140]{source140}\label{bib:source140} Ghosh, M. and Meeden, G. (1977). On the Non-Attainability of Chebychev Bounds. Amer. Statist. 31 35–36.
\bibitem[141]{source141}\label{bib:source141} Gianola, D. and Fernando, R. L. (1986). Bayesian Methods in Animal Breeding Theory. Journal of Animal Science 63 217–244.
\bibitem[142]{source142}\label{bib:source142} Gilat, D. (1977). Monotonicity of a Power Function: An Elementary Probabilistic Proof. Amer. Statist. 31 91–93.
\bibitem[143]{source143}\label{bib:source143} Gleser, L. J. (1981). Estimation in a Multivariate Errors-in-Variables Regression Model: Large-Sample Results. Ann. Statist. 9 24–44.
\bibitem[144]{source144}\label{bib:source144} Gleser, L. J. (1983). Functional, Structural, and Ultrastructural Errors-in-Variables Models. Proceedings of the Business and Economic Statistics Section 57–66. Alexandria, VA: American Statistical Association.
\bibitem[145]{source145}\label{bib:source145} Gleser, L. J. (1987). Confidence Intervals for the Slope in a Linear Errors-in-Variables Regression Model. Advances in Multivariate Statistical Analysis (K. Gupta, ed.) 85–109. Dordrecht: D. Reidel.
\bibitem[146]{source146}\label{bib:source146} Gleser, L. J. (1989). The Gamma Distribution as a Mixture of Exponential Distributions. Amer. Statist. 43 115–117.
\bibitem[147]{source147}\label{bib:source147} Gleser, L. J. (1991). Measurement error models (with discussion). Chem. Int. Lab. Sys. 10 45–67.
\bibitem[148]{source148}\label{bib:source148} Gleser, L. J., Carroll, R. J., and Gallo, P. (1987). The Limiting Distribution of Least Squares in an Errors-in-Variables Regression Model. Ann. Statist. 15 220–233.
\bibitem[149]{source149}\label{bib:source149} Gleser, L. J. and Healy, J. D. (1976). Estimating the Mean of a Normal Distribution with Known Coefficient of Variation. J. Amer. Statist. Assoc. 71 977–981.
\bibitem[150]{source150}\label{bib:source150} Gleser, L. J. and Hwang, J. T. (1987). The Nonexistence of 100(1 − \ensuremath{\alpha})\% Confidence Sets of Finite Expected Diameter in Errors-in-Variables and Related Models. Ann. Statist. 15 1351–1362.
\bibitem[151]{source151}\label{bib:source151} Gnedenko, B. V. (1978). The Theory of Probability. Moscow: MIR Publishers.
\bibitem[152]{source152}\label{bib:source152} Groeneveld, R. A. (1991). An Influence Function Approach to Describing the Skewness of a Distribution. Amer. Statist. 45 97–102.
\bibitem[153]{source153}\label{bib:source153} Guenther, W. C. (1978). Some Easily Found Minimum Variance Unbiased Estimators. Amer. Statist. 32 29–33.
\bibitem[154]{source154}\label{bib:source154} Hall, P. (1992). The Bootstrap and Edgeworth Expansion. New York: Springer-Verlag.
\bibitem[155]{source155}\label{bib:source155} Halmos, P. R. and Savage, L. J. (1949). Applications of the Radon–Nikodym Theorem to the Theory of Sufficient Statistics. Ann. Math. Statist. 20 225–241.
\bibitem[156]{source156}\label{bib:source156} Hampel, F. R. (1974). The Influence Curve and Its Role in Robust Estimation. J. Amer. Statist. Assoc. 69 383–393.
\bibitem[157]{source157}\label{bib:source157} Hanley, J. A. (1992). Jumping to Coincidences: Defying Odds in the Realm of the Preposterous. Amer. Statist. 46 197–202.
\bibitem[158]{source158}\label{bib:source158} Hardy, G. H., Littlewood, J. E., and Polya, G. (1952). Inequalities, 2nd edition. London: Cambridge University Press.
\bibitem[159]{source159}\label{bib:source159} Hartley, H. O. (1958). Maximum Likelihood Estimation from Incomplete Data. Biometrics 14 174–194.
\bibitem[160]{source160}\label{bib:source160} Harville, D. A. (1981). Unbiased and Minimum-Variance Unbiased Estimation of Estimable Functions for Fixed Linear Models with Arbitrary Covariance Structure. Ann. Statist. 9 633–637.
\bibitem[161]{source161}\label{bib:source161} Hawkins, D. M. (1993). The Feasible Set Algorithm for Least Median of Squares Regression. Computational Statistics and Data Analysis 16 81–101.
\sourcepageanchor{547}
\bibitem[162]{source162}\label{bib:source162} Hawkins, D. M. (1994). The Feasible Solution Algorithm for Least Trimmed Squares Regression. Computational Statistics and Data Analysis 17 186–196.
\bibitem[163]{source163}\label{bib:source163} Hawkins, D. M. (1995). Convergence of the Feasible Solution Algorithm for Least Median of Squares Regressionn. Computational Statistics and Data Analysis 19 519–538.
\bibitem[164]{source164}\label{bib:source164} Hayter, A. J. (1984). A Proof of the Conjecture that the Tukey-Kramer Multiple Comparison Procedure Is Conservative. Ann. Statist. 12 61–75.
\bibitem[165]{source165}\label{bib:source165} Hettmansperger, T. P. and McKean, J. W. (1998). Robust Nonparametric Statistical Methods. London: Kendall’s Library of Statistics, 5, Edward Arnold; New York: Wiley.
\bibitem[166]{source166}\label{bib:source166} Hinkley, D. V. (1980). Likelihood. Can. J. Statist. 8 151–163.
\bibitem[167]{source167}\label{bib:source167} Hinkley, D. V., Reid, N, and Snell, L. (1991). Statistical Theory and Modelling. In honor of Sir David Cox. London: Chapman and Hall.
\bibitem[168]{source168}\label{bib:source168} Hinkley, D. V. and Runger, G. (1984). The Analysis of Transformed Data (with discussion). J. Amer. Statist. Assoc. 79 302–320.
\bibitem[169]{source169}\label{bib:source169} Hsu, J. C. (1996). Multiple Comparisons: Theory and Methods. London: Chapman and Hall..
\bibitem[170]{source170}\label{bib:source170} Huber, P. J. (1964). Robust Estimation of a Location Parameter. Ann. Math. Statist. 35 73–101.
\bibitem[171]{source171}\label{bib:source171} Huber, P. J. (1981). Robust Statistics. New York: Wiley.
\bibitem[172]{source172}\label{bib:source172} Hudson, H. M. (1978). A Natural Identity for Exponential Families with Applications in Multiparameter Estimation. Ann. Statist. 6 473–484.
\bibitem[173]{source173}\label{bib:source173} Huzurbazar, V. S. (1949). On a Property of Distributions Admitting Sufficient Statistics. Biometrika 36 71–74.
\bibitem[174]{source174}\label{bib:source174} Hwang, J. T. (1982). Improving on Standard Estimators in Discrete Exponential Families with Applications to Poisson and Negative Binomial Cases. Ann. Statist. 10 857–867.
\bibitem[175]{source175}\label{bib:source175} Hwang, J. T. (1995). Fieller’s Problems and Resampling Techniques. Statistica Sinica 5 161–171.
\bibitem[176]{source176}\label{bib:source176} James, W. and Stein, C. (1961). Estimation with Quadratic Loss. Proc. Fourth Berkeley Symp. Math. Statist. Prob. 1. University of California Press.
\bibitem[177]{source177}\label{bib:source177} Johnson, N. L. and Kotz, S. (1969-1972). Distributions in Statistics (4 vols.). New York: Wiley.
\bibitem[178]{source178}\label{bib:source178} Johnson, N. L., Kotz. S., and Balakrishnan, N. (1994). Continuous Univariate Distributions, Volume 1, 2nd edition. New York: Wiley.
\bibitem[179]{source179}\label{bib:source179} Johnson, N. L., Kotz. S., and Balakrishnan, N. (1995). Continuous Univariate Distributions, Volume 2, 2nd edition. New York: Wiley.
\bibitem[180]{source180}\label{bib:source180} Johnson, N. L., Kotz. S., and Kemp, A. W. (1992). Univariate Discrete Distributions, 2nd edition. New York: Wiley.
\bibitem[181]{source181}\label{bib:source181} Jones, M. C. (1999). Distributional Relationships Arising from Simple Trigonometric Formulas. Amer. Statist. 53 99–102
\bibitem[182]{source182}\label{bib:source182} Joshi, S. M. and Nabar, S. P. (1989). Linear Estimators for the Parameter in the Problem of the Nile. Amer. Statist. 43 40–41.
\bibitem[183]{source183}\label{bib:source183} Joshi, V. M. (1969). Admissibility of the Usual Confidence Sets for the Mean of a Univariate or Bivariate Normal Population. Ann. Math. Statist. 40 1042–1067.
\bibitem[184]{source184}\label{bib:source184} Juola, R. C. (1993). More on Shortest Confidence Intervals. Amer. Statist. 47 117–119.
\bibitem[185]{source185}\label{bib:source185} Kalbfleisch, J. D. (1975). Sufficiency and Conditionality. Biometrika 62 251–268.
\bibitem[186]{source186}\label{bib:source186} Kalbfleisch, J. D. and Prentice, R. L. (1980). The Statistical Analysis of Failure Time Data. New York: Wiley.
\sourcepageanchor{548}
\bibitem[187]{source187}\label{bib:source187} Karlin, S. and Ost, F. (1988). Maximal Length of Common Words among Random Letter Sequences. Ann. Prob. 16 535–563.
\bibitem[188]{source188}\label{bib:source188} Kelker, D. (1970). Distribution Theory of Spherical Distributions and a Location–Scale Parameter Generalization. Sankhyā, Ser. A 32 419–430.
\bibitem[189]{source189}\label{bib:source189} Kendall, M. and Stuart, A. (1979). The Advanced Theory of Statistics, Volume II: Inference and Relationship, 4th edition. New York: Macmillan.
\bibitem[190]{source190}\label{bib:source190} Kirk, R. E. (1982). Experimental Design: Procedures for the Behavorial Sciences, 2nd edition. Pacific Grove, CA: Brooks/Cole.
\bibitem[191]{source191}\label{bib:source191} Koopmans, L. H. (1993). A Note on Using the Moment Generating Function to Teach the Laws of Large Numbers. Amer. Statist. 47 199–202.
\bibitem[192]{source192}\label{bib:source192} Kuehl, R. O. (2000). Design of Experiments : Statistical Principles of Research Design and Analysis, 2nd edition. Pacific Grove, CA : Duxbury.
\bibitem[193]{source193}\label{bib:source193} Lange, N., Billard, L., Conquest, L., Ryan, L., Brillinger, D., Greenhouse, J. (eds). (1994), Case Studies in Biometry. New York: Wiley-Interscience.
\bibitem[194]{source194}\label{bib:source194} Le Cam, L. (1953). On Some Asymptotic Properties of Maximum Likelihood Estimates and Related Bayes’ Estimates. Univ. of Calif. Publ. in Statist. 1 277–330.
\bibitem[195]{source195}\label{bib:source195} Leemis, L. M. (1986). Relationships Among Common Univariate Distributions. Amer. Statist. 40 143–146.
\bibitem[196]{source196}\label{bib:source196} Leemis, L. M. and Trivedi, K. S. (1996). A Comparison of Approximate Interval Estimators for the Bernoulli Parameter. Amer. Statist. 50 63–68.
\bibitem[197]{source197}\label{bib:source197} Lehmann, E. L. (1981). An Interpretation of Completeness and Basu’s Theorem. J. Amer. Statist. Assoc. 76 335–340.
\bibitem[198]{source198}\label{bib:source198} Lehmann, E. L. (1986). Testing Statistical Hypotheses, 2nd edition. New York: Wiley.
\bibitem[199]{source199}\label{bib:source199} Lehmann, E. L. (1999). Introduction to Large-Sample Theory. New York: Springer-Verlag.
\bibitem[200]{source200}\label{bib:source200} Lehmann, E. L. and Casella, G. (1998). Theory of Point Estimation, 2nd edition. New York: Springer-Verlag.
\bibitem[201]{source201}\label{bib:source201} Lehmann, E. L. and Scheffé, H. (1950, 1955, 1956). Completeness, Similar Regions, and Unbiased Estimation. Sankhyā, Ser. A 10 305–340; 15 219–236; correction 17 250.
\bibitem[202]{source202}\label{bib:source202} Lehmann, E. L. and Scholz, F. W. (1992). Ancillarity. Current Issues in Statistical Inference: Essays in Honor of D. Basu (M. Ghosh and P. K. Pathak, eds.). Hayward, CA: IMS Monograph Series, 32–51.
\bibitem[203]{source203}\label{bib:source203} LePage, R. and Billard, L. (eds.) (1992). Exploring the Limits of Bootstrap. New York: Wiley.
\bibitem[204]{source204}\label{bib:source204} Lindley, D. V. (1957). A Statistical Paradox. Biometrika 44 187–192.
\bibitem[205]{source205}\label{bib:source205} Lindley, D. V. (1962). Discussion of the Article by Stein. J. Roy. Statist. Soc. Ser. B 24 265–296.
\bibitem[206]{source206}\label{bib:source206} Lindley, D. V. and Phillips, L. D. (1976). Inference for a Bernoulli Process (a Bayesian View). Amer. Statist. 30 112–119.
\bibitem[207]{source207}\label{bib:source207} Lindley, D. V. and Smith, A. F. M. (1972). Bayes Estimates for the Linear Model. J. Roy. Statist. Soc. Ser. B 34 1–41.
\bibitem[208]{source208}\label{bib:source208} Little, R. J. A. and Rubin, D. B. (1987). Statistical Analysis with Missing Data. New York: Wiley.
\bibitem[209]{source209}\label{bib:source209} Liu, R. Y. (1990). On a Notion of Data Depth Based on Random Simplices. Ann. Statist. 18 405–414.
\bibitem[210]{source210}\label{bib:source210} Liu, R. Y. and Singh, K. (1992). Ordering Directional Data: Concepts of Data Depth on Circles and Spheres. Ann. Statist. 20 1468–1484.
\sourcepageanchor{549}
\bibitem[211]{source211}\label{bib:source211} Luceño, A. (1997). Further Evidence Supporting the Numerical Usefulness of Characteristic Functions. Amer. Statist. 51 233–234.
\bibitem[212]{source212}\label{bib:source212} Maatta, J. M. and Casella, G. (1987). Conditional Properties of Interval Estimators of the Normal Variance. Ann. Statist. 15 1372–1388.
\bibitem[213]{source213}\label{bib:source213} Madansky, A. (1962). More on Length of Confidence Intervals. J. Amer. Statist. Assoc. 57 586–589.
\bibitem[214]{source214}\label{bib:source214} Marshall, A. W. and Olkin, I. (1979). Inequalities: Theory of Majorization and Its Applications. New York: Academic Press.
\bibitem[215]{source215}\label{bib:source215} McCullagh, P. (1994). Does the Moment Generating Function Characterize a Distribution? Amer. Statist. 48 208.
\bibitem[216]{source216}\label{bib:source216} McLachlan, G., and Krishnan, T. (1997). The EM Algotithm and Extensions. New York: Wiley.
\bibitem[217]{source217}\label{bib:source217} McPherson, G. (1990). Statistics in Scientific Inverstigation. New York: Springer-Verlag.
\bibitem[218]{source218}\label{bib:source218} Mengersen, K. L. and Tweedie, R. L. (1996). Rates of Convergence of the Hastings and Metropolis Algorithms. Ann. Statist. 24 101–121.
\bibitem[219]{source219}\label{bib:source219} Metropolis, N., Rosenbluth, A. W., Rosenbluth, M. N., Teller, A. H., Teller, E. (1953). Equations of State Calculations by Fast Computing Machines. J. Chem. Phys. 21 1087–1092.
\bibitem[220]{source220}\label{bib:source220} Miller, R. G. (1974). The Jackknife—A Review. Biometrika 61 1–15.
\bibitem[221]{source221}\label{bib:source221} Miller, R. G. (1981). Simultaneous Statistical Inference, 2nd edition. New York: Springer- Verlag.
\bibitem[222]{source222}\label{bib:source222} Moran, P. A. P. (1971). Estimating Structural and Functional Relationships. J. Mult. Analysis. 1 232–255.
\bibitem[223]{source223}\label{bib:source223} Morgan, J. P., Chaganty, N. R., Dahiya, R. C. and Doviak, M. J. (1991). Let’s Make a Deal: The Player’s Dilemma (with discussion). Amer. Statist. 45 284–289.
\bibitem[224]{source224}\label{bib:source224} Morris, C. N. (1982). Natural Exponential Families with Quadratic Variance Functions. Ann. Statist. 10 65–80.
\bibitem[225]{source225}\label{bib:source225} Morris, C. N. (1983). Parametric Empirical Bayes Inference: Theory and Applications (with discussion). J. Amer. Statist. Assoc. 78 47–65.
\bibitem[226]{source226}\label{bib:source226} Morrison, D. G. (1978). A Probability Model for Forced Binary Choices. Amer. Statist. 32 23–25.
\bibitem[227]{source227}\label{bib:source227} Naiman, D. Q. (1983). Comparing Scheffé-Type to Constant Width Confidence Bounds in Regression. J. Amer. Statist. Assoc. 78 906–912.
\bibitem[228]{source228}\label{bib:source228} Naiman, D. Q. (1984). Optimal Simultaneous Confidence Bounds. Ann. Statist. 12 702–715.
\bibitem[229]{source229}\label{bib:source229} Naiman, D. Q. (1987). Simultaneous Confidence Bounds in Multiple Regression Using Predictor Variable Constraints. J. Amer. Statist. Assoc. 82 214–219.
\bibitem[230]{source230}\label{bib:source230} Naiman, D. Q. and Wynn, H. P. (1992). Inclusion-Exclusion-Bonferroni Identities and Inequalities for Discrete Tube-like Problems via Euler Characteristics. Ann. Statist. 20 43–76.
\bibitem[231]{source231}\label{bib:source231} Naiman, D. Q. and Wynn, H. P. (1997). Abstract Tubes, Improved Inclusion-Exclusion Identities and Inequalities and Importance Sampling. Ann. Statist. 25 1954–1983.
\bibitem[232]{source232}\label{bib:source232} Neter, J., Wasserman, W., and Whitmore, G. A. (1993). Applied Statistics. Boston : Allyn and Bacon.
\bibitem[233]{source233}\label{bib:source233} Neyman, J. (1935). Su un teorema concernente le cosiddette statistiche sufficienti. Inst. Ital. Atti. Giorn. 6 320–334.
\bibitem[234]{source234}\label{bib:source234} Noorbaloochi, S. and Meeden, G. (1983). Unbiasedness as the Dual of Being Bayes. J. Amer. Statist. Assoc. 78 619–623.
\sourcepageanchor{550}
\bibitem[235]{source235}\label{bib:source235} Norton, R. M. (1984). The Double Exponential Distribution: Using Calculus to Find an MLE. Amer. Statist. 38 135–136.
\bibitem[236]{source236}\label{bib:source236} Nussbaum, M. (1976). Maximum Likelihood and Least Squares Estimation of Linear Functional Relationships. Mathematische Operationsforschung und Statistik, Ser. Statistik 7 23–49.
\bibitem[237]{source237}\label{bib:source237} Olkin, I., Petkau, A. J., and Zidek, J. V. (1981). A Comparison of n Estimators for the Binomial Distribution. J. Amer. Statist. Assoc. 76 637–642.
\bibitem[238]{source238}\label{bib:source238} Pal, N. and Berry, J. (1992). On Invariance and Maximum Likelihood Estimation. Amer. Statist. 46 209–212.
\bibitem[239]{source239}\label{bib:source239} Park, C. G., Park, T., and Shin, D. W. (1996). A Simple Method for Generating Correlated Binary Variates. Amer. Statist. 50 306–310.
\bibitem[240]{source240}\label{bib:source240} Pena, E. A. and Rohatgi, V. (1994). Some Comments about Sufficiency and Unbiased Estimation. Amer. Statist. 48 242–243.
\bibitem[241]{source241}\label{bib:source241} Piegorsch, W. W. (1985). Admissible and Optimal Confidence Bands in Simple Linear Regression. Ann. Statist. 13 801–810.
\bibitem[242]{source242}\label{bib:source242} Pitman, E. J. G. (1939). The Estimation of the Location and Scale Parameters of a Continuous Population of Any Given Form. Biometrika 30 200–215.
\bibitem[243]{source243}\label{bib:source243} Portnoy, S. (1987). A Central Limit Theorem Applicable to Robust Regression Estimators. J. Mult. Analysis. 22 24–50.
\bibitem[244]{source244}\label{bib:source244} Portnoy, S. and Koenker, R. (1997). The Gaussian Hare and the Laplacian Tortoise: Computability of Squared-Error versus Absolute-Error Estimators. Statist. Sci. 12 279–300.
\bibitem[245]{source245}\label{bib:source245} Portnoy, S. and Mizera, I. (1998). Discussion of the Paper by Ellis. Statist. Sci. 13 344–347.
\bibitem[246]{source246}\label{bib:source246} Pratt, J. W. (1961). Length of Confidence Intervals. J. Amer. Statist. Assoc. 56 549–567.
\bibitem[247]{source247}\label{bib:source247} Proschan, M. A. and Presnell, B. (1998). Expect the Unexpected from Conditional Expectation. Amer. Statist. 52 248–252.
\bibitem[248]{source248}\label{bib:source248} Pukelsheim, F. (1994). The Three Sigma Rule. Amer. Statist. 48 88–91.
\bibitem[249]{source249}\label{bib:source249} Quenouille, M. H. (1956). Notes on Bias in Estimation. Biometrika 43 353–360.
\bibitem[250]{source250}\label{bib:source250} Reid, N. (1995). The Role of Conditioning in Inference (with discussion). Statist. Sci. 10 138–166.
\bibitem[251]{source251}\label{bib:source251} Resnick, S. I. (1999). A Probability Path. Basel: Birkhauser.
\bibitem[252]{source252}\label{bib:source252} Ridgeway, T. (1993). Letter to the Editor about the “Exchange Paradox.” Amer. Statist. 47 311.
\bibitem[253]{source253}\label{bib:source253} Ripley, B. D. (1987) Stochastic Simulation. New York: Wiley.
\bibitem[254]{source254}\label{bib:source254} Robbins, H. (1977). A Fundamental Question of Practical Statistics (letter to the editor). Amer. Statist. 31 97.
\bibitem[255]{source255}\label{bib:source255} Robert, C. P. (1994) The Bayesian Choice. New York: Springer-Verlag.
\bibitem[256]{source256}\label{bib:source256} Robert, C. P. and Casella, G. (1999). Monte Carlo Statistical Methods New York: Springer- Verlag.
\bibitem[257]{source257}\label{bib:source257} Robson, D. S. (1959). A Simple Method for Constructing Orthogonal Polynomials when the Independent Variable Is Unequally Spaced. Biometrics 15 187–191.
\bibitem[258]{source258}\label{bib:source258} Romano, J. P. and Siegel, A. F. (1986). Counterexamples in Probability and Statistics. Pacific Grove, CA: Wadsworth and Brooks/Cole.
\bibitem[259]{source259}\label{bib:source259} Ross, S. M. (1988) A First Course in Probability Theory, 3rd edition. New York: MacMillan.
\bibitem[260]{source260}\label{bib:source260} Ross, S. M. (1994) Letter to the Editor about the “Exchange Paradox.” Amer. Statist. 48 267.
\sourcepageanchor{551}
\bibitem[261]{source261}\label{bib:source261} Ross, S. M. (1996). Bayesians Should Not Resample a Prior Sample to Learn about the Posterior. Amer. Statist. 50 116.
\bibitem[262]{source262}\label{bib:source262} Rousseeuw, P. J. (1984). Least Median of Squares Regression. J. Amer. Statist. Assoc. 79 871–880.
\bibitem[263]{source263}\label{bib:source263} Rousseeuw, P. J. and Hubert, M. (1999). Regression Depth (with discussion). J. Amer. Statist. Assoc. 94 388–433
\bibitem[264]{source264}\label{bib:source264} Rouseeuw, P. J. and Leroy, A. M. (1987). Robust Regression and Outlier Detection. New York: Wiley.
\bibitem[265]{source265}\label{bib:source265} Royall, R. M. (1997). Statistical Evidence: A Likelihood Paradigm. London: Chapman and Hall..
\bibitem[266]{source266}\label{bib:source266} Rubin, D. B. (1988). Using the SIR Algorithm to Simulate Posterior Distributions. Bayesian Statistics 3 (J. M. Bernardo, M. H. DeGroot, D. V. Lindley, and A. F. M. Smith, eds.) Cambridge, MA: Oxford University Press 395–402.
\bibitem[267]{source267}\label{bib:source267} Rudin, W. (1976). Principles of Real Analysis. New York: McGraw-Hill.
\bibitem[268]{source268}\label{bib:source268} Ruppert, D. (1987). What Is Kurtosis? Amer. Statist. 41 1–5.
\bibitem[269]{source269}\label{bib:source269} Ruppert, D. and Carroll, R. J. (1979). Trimmed Least Squares Estimation in the Linear Model. J. Amer. Statist. Assoc. 75 828–838.
\bibitem[270]{source270}\label{bib:source270} Russell, K. G. (1991). Estimating the value of e by Simulation. Amer. Statist. 45 66–68
\bibitem[271]{source271}\label{bib:source271} Samuels, M. L. and Lu, T.-F. C. (1992). Sample Size Requirements for the Back-of-the- Envelope Binomial Confidence Interval. Amer. Statist. 46 228–231.
\bibitem[272]{source272}\label{bib:source272} Satterthwaite, F. E. (1946). An Approximate Distribution of Estimates of Variance Components. Biometrics Bulletin (now called Biometrics) 2 110–114.
\bibitem[273]{source273}\label{bib:source273} Saw, J. G., Yang, M. C. K., and Mo, T. C. (1984). Chebychev’s Inequality with Estimated Mean and Variance. Amer. Statist. 38 130–132.
\bibitem[274]{source274}\label{bib:source274} Schafer, J. L. (1997). Analysis of Incomplete Multivariate Data. London: Chapman and Hall.
\bibitem[275]{source275}\label{bib:source275} Scheffé, H. (1959). The Analysis of Variance. New York: Wiley.
\bibitem[276]{source276}\label{bib:source276} Schervish, M. J. (1995). Theory of Statistics. New York: Springer-Verlag.
\bibitem[277]{source277}\label{bib:source277} Shervisch, M. J. (1996). P Values: What They Are and What They Are Not. Amer. Statist. 50 203–206.
\bibitem[278]{source278}\label{bib:source278} Schuirmann, D. J. (1987). A Comparison of the Two One-Sided Test Procedures and the Power Approach for Assessing the Equivalence of Average Bioavailability. J. Pharmacokinetics and Biopharmaceutics 15 657–680.
\bibitem[279]{source279}\label{bib:source279} Schwager, S. J. (1984). Bonferroni Sometimes Loses. Amer. Statist. 38 192–197.
\bibitem[280]{source280}\label{bib:source280} Schwager, S. J. (1985). Reply to Worsley. Amer. Statist. 39 236.
\bibitem[281]{source281}\label{bib:source281} Schwarz, C. J. and Samanta, M. (1991). An Inductive Proof of the Sampling Distributions for the MLEs of the Parameters in an Inverse Gaussian Distribution. Amer. Statist. 45 223–225.
\bibitem[282]{source282}\label{bib:source282} Searle, S. R. (1971). Linear Models. New York: Wiley.
\bibitem[283]{source283}\label{bib:source283} Searle, S. R. (1982). Matrix Algebra Useful for Statistics. New York: Wiley.
\bibitem[284]{source284}\label{bib:source284} Searls, D. T. and Intarapanich, P. (1990). A Note on an Estimator for the Variance that Utilizes the Kurtosis. Amer. Statist. 44 295–296.
\bibitem[285]{source285}\label{bib:source285} Selvin, S. (1975). A Problem in Probability (letter to the editor). Amer. Statist. 29 67.
\bibitem[286]{source286}\label{bib:source286} Seshadri, V. (1993). The Inverse Gaussian Distribution. A Case Study in Exponential Families. New York: Clarendon Press.
\bibitem[287]{source287}\label{bib:source287} Shao, J. (1999). Mathematical Statistics. New York: Springer-Verlag.
\sourcepageanchor{552}
\bibitem[288]{source288}\label{bib:source288} Shao, J. and Tu, D. (1995). The Jackknife and the Bootstrap. New York: Springer-Verlag.
\bibitem[289]{source289}\label{bib:source289} Shier, D. P. (1988). The Monotonicity of Power Means Using Entropy. Amer. Statist. 42 203–204.
\bibitem[290]{source290}\label{bib:source290} Shuster, J. J. (1991). The Statistician in a Reverse Cocaine Sting. Amer. Statist. 45 123–124.
\bibitem[291]{source291}\label{bib:source291} Silvapulle, M. J. (1996). A Test in the Presence of Nuisance Parameters. J. Amer. Statist. Assoc. 91 1690–1693.
\bibitem[292]{source292}\label{bib:source292} Smith, A. F. M. and Gelfand, A. E. (1992). Bayesian Statistics Without Tears: A Sampling- Resampling Perspective. Amer. Statist. 46 84–88.
\bibitem[293]{source293}\label{bib:source293} Smith, A. F. M. and Roberts, G. O. (1993). Bayesian Computation via the Gibbs Sampler and Related Markov Chain Monte Carlo Methods (with discussion). J. Roy. Statist. Soc. Ser. B 55 3–24.
\bibitem[294]{source294}\label{bib:source294} Snedecor, G. W. and Cochran, W. G. (1989). Statistical Methods, 8th edition. Ames, IA: Iowa State University Press.
\bibitem[295]{source295}\label{bib:source295} Solari, M. E. (1969). The “Maximum Likelihood Solution” of the Problem of Estimating a Linear Functional Relationship. J. Roy. Statist. Soc. Ser. B 31 372–375.
\bibitem[296]{source296}\label{bib:source296} Solomon, D. L. (1975). A Note on the Non-Equivalence of the Neyman–Pearson and Generalized Likelihood Ratio Tests for Testing a Simple Null Versus a Simple Alternative Hypothesis. Amer. Statist. 29 101–102.
\bibitem[297]{source297}\label{bib:source297} Solomon, D. L. (1983). The Spatial Distribution of Cabbage Butterfly Eggs. Life Science Models, Volume 4 (H. Marcus-Roberts and M. Thompson, eds.) New York: Springer-Verlag 350–366.
\bibitem[298]{source298}\label{bib:source298} Sprott, D. A. and Farewell, V. T. (1993). The Difference Between Two Normal Means. Amer. Statist. 47 126–128.
\bibitem[299]{source299}\label{bib:source299} Staudte, R. G. and Sheather, S. J. (1990). Robust Estimation and Testing. New York: Wiley.
\bibitem[300]{source300}\label{bib:source300} Stefanski, L. A. (1996). A Note on the Arithmetic-Geometric-Harmonic Mean Inequalities. Amer. Statist. 50 246–247.
\bibitem[301]{source301}\label{bib:source301} Stein, C. (1956). Inadmissibility of the Usual Estimator for the Mean of a Multivariate Normal Distribution. Proc. Third Berkeley Symp. Math. Statist. Prob. 1. University of California Press 197–206.
\bibitem[302]{source302}\label{bib:source302} Stein, C. (1964). Inadmissibility of the Usual Estimator for the Variance of a Normal Distribution with Unknown Mean. Ann. Inst. Statist. Math. 16 155–160.
\bibitem[303]{source303}\label{bib:source303} Stein, C. (1973). Estimation of the Mean of a Multivariate Distribution. Proceedings of the Prague Symposium on Asymptotic Statistics. Prague: Charles Univ. 345–381.
\bibitem[304]{source304}\label{bib:source304} Stein, C. (1981). Estimation of the Mean of a Multivariate Normal Distribution. Ann. Statist. 9 1135–1151.
\bibitem[305]{source305}\label{bib:source305} Sterne, T. E. (1954). Some Remarks on Confidence or Fiducial Limits. Biometrika 41 275–278.
\bibitem[306]{source306}\label{bib:source306} Stigler, S. M. (1983). Who Discovered Bayes’ Theorem? Amer. Statist. 37 290–296.
\bibitem[307]{source307}\label{bib:source307} Stigler, S. M. (1984). Kruskal’s Proof of the Joint Distribution of $\bar X$ and $S^2$ . Amer. Statist. 38 134–135.
\bibitem[308]{source308}\label{bib:source308} Stigler, S. M. (1986). The History of Statistics: The Measurement of Uncertainty Before 1900. Cambridge, MA: Harvard University Press.
\bibitem[309]{source309}\label{bib:source309} Stuart, A. and Ord, J. K. (1987). Kendall’s Advanced Theory of Statistics, Volume I: Distribution Theory, 5th edition. New York: Oxford University Press.
\bibitem[310]{source310}\label{bib:source310} Stuart, A. and Ord, J. K. (1991). Kendall’s Advanced Theory of Statistics, Volume II, 5th edition. New York: Oxford University Press.
\sourcepageanchor{553}
\bibitem[311]{source311}\label{bib:source311} Stuart, A., Ord, J. K., and Arnold, S. (1999). Advanced Theory of Statistics, Volume 2A: Classical Inference and the Linear Model, 6th edition. London: Oxford University Press.
\bibitem[312]{source312}\label{bib:source312} Tanner, M. A. (1996) Tools for Statistical Inference: Observed Data and Data Augmentation Methods, 3rd edition. New York: Springer-Verlag.
\bibitem[313]{source313}\label{bib:source313} Tate, R. F. and Klett, G. W. (1959). Optimal Confidence Intervals for the Variance of a Normal Distribution. J. Amer. Statist. Assoc. 54 674–682.
\bibitem[314]{source314}\label{bib:source314} Tierney, J. (1991). Behind Monte Hall’s Door: Puzzle, Debate and Answer? The New York Times, July 21, 1991.
\bibitem[315]{source315}\label{bib:source315} Tierney, L. (1994). Markov Chains for Exploring Posterior Distributions (with discussion). Ann. Statist. 22 1701–1786.
\bibitem[316]{source316}\label{bib:source316} Tsao, C. A. and Hwang, J. T. (1998). Improved Confidence Estimators for Fieller’s Confidence Sets. Can. J. Statist. 26 299–310.
\bibitem[317]{source317}\label{bib:source317} Tsao, C. A. and Hwang, J. T. (1999). Generalized Bayes Confidence Estimators for Fieller’s Confidence Sets. Statistica Sinica 9 795–810.
\bibitem[318]{source318}\label{bib:source318} Tukey, J. W. (1977). Exploratory Data Analysis. Reading, MA: Addison-Wesley.
\bibitem[319]{source319}\label{bib:source319} Tweedie, M. C. K. (1957). Statistical Properties of Inverse Gaussian Distributions I. Ann. Math. Statist. 28 362–377.
\bibitem[320]{source320}\label{bib:source320} Vardeman, S. B. (1987). Discussion of the Articles by Casella and Berger and Berger and Sellke. J. Amer. Statist. Assoc. 82 130–131.
\bibitem[321]{source321}\label{bib:source321} Vardeman, S. B. (1992). What About Other Intervals. Amer. Statist. 46 193–197.
\bibitem[322]{source322}\label{bib:source322} vos Savant, M. (1990). Ask Marilyn. Parade Magazine, September 9, 15.
\bibitem[323]{source323}\label{bib:source323} vos Savant, M. (1991). Letter to the Editor. Amer. Statist. 45 347.
\bibitem[324]{source324}\label{bib:source324} Wald, A. (1940). The Fitting of Straight Lines when Both Variables Are Subject to Error. Ann. Math. Statist. 11 284–300.
\bibitem[325]{source325}\label{bib:source325} Waller, L. A. (1995). Does the Characteristic Function Numerically Distinguish Distributions? Amer. Statist. 49 150–151.
\bibitem[326]{source326}\label{bib:source326} Waller, L. A., Turnbull, B. W., and Hardin, J. M. (1995). Obtaining Distribution Functions by Numerical Inversion of Characteristic Functions with Applications. Amer. Statist. 49 346–350.
\bibitem[327]{source327}\label{bib:source327} Wasserman, L. (1992). Recent Methodological Advances in Robust Bayesian Inference (with discussion). Bayesian Statistics 4. J. M. Bernardo, J. O. Berger, A. P. Dawid and A. F. M. Smith, eds. Oxford: Clarendon Press, 483–502.
\bibitem[328]{source328}\label{bib:source328} Westlake, W. J. (1981). Bioequivalence Testing – A Need to Rethink. Biometrics 37 591–593.
\bibitem[329]{source329}\label{bib:source329} Widder, D. V. (1946). The Laplace Transform. Princeton, NJ: Princeton University Press.
\bibitem[330]{source330}\label{bib:source330} Wilks, S. S. (1938). Shortest Average Confidence Intervals from Large Samples. Ann. Math. Statist. 9 166–175.
\bibitem[331]{source331}\label{bib:source331} Williams, E. J. (1959). Regression Analysis. New York: Wiley.
\bibitem[332]{source332}\label{bib:source332} Worsley, K. J. (1982). An Improved Bonferroni Inequality and Applications. Biometrika 69 297–302.
\bibitem[333]{source333}\label{bib:source333} Worsley, K. J. (1985). Bonferroni (Improved) Wins Again. Amer. Statist. 39 235.
\bibitem[334]{source334}\label{bib:source334} Wright, T. (1992). Lagrange’s Identity Reveals Correlation Coefficient and Straight-Line Connection. Amer. Statist. 46 106–107.
\bibitem[335]{source335}\label{bib:source335} Wu, C. F. J. (1983). On the Convergence of the EM Algorithm. Ann. Statist. 11 95–103.
\sourcepageanchor{554}
\bibitem[336]{source336}\label{bib:source336} Young, G. A. (1994). Bootstrap: More Than a Stab in the Dark? (with discussion). Statist. Sci. 9 382–415.
\bibitem[337]{source337}\label{bib:source337} Zehna, P. W. (1966). Invariance of Maximum Likelihood Estimators. Ann. Math. Statist. 37 744.
\bibitem[338]{source338}\label{bib:source338} Zellner, A. (1986). Bayesian Estimation and Prediction Using Asymmetric Loss Functions. J. Amer. Statist. Assoc. 81 446–451.
\end{thebibliography}



\clearpage\section*{索引}\phantomsection\label{app:index}
\addcontentsline{toc}{section}{索引}
按原书英文词条顺序保留主项与子项，附中文译名.下列页码由本译稿中相关概念、定理、公式或习题的位置生成；同一译稿页的重复定位合并.原词条及原书页码的完整映射保存在核验资料中.\footnote{译注：原索引的“$t$ 检验／一致最强／无偏”（原印刷331页）与“置信区间／一致最准确／正态”（原印刷372页）分类比所引内容的实际陈述更宽. 此处分别定位到\autoref{exercise:8.37}、\autoref{exercise:9.47}的相关讨论，保留原分类，不据索引额外断言其结论. “每次比较错误率”原指印刷458页，明确定义实际在下一页的\hyperref[item:11.17.b]{习题11.17(b)}；“回归／残差”的明确公式同样在原索引所指页的下一页. 页码已按实际语义位置生成.}
\begingroup\small\raggedright\interlinepenalty=10000\setlength{\parskip}{0pt}\setlength{\columnsep}{1.5em}
\begin{multicols}{2}
\sourcepageanchor{556}
\par\addvspace{2pt}
\par\noindent\hangindent=1em\hspace*{0em}\phantomsection\label{idx:1}接受／拒绝算法（Accept-Reject algorithm），\hyperref[ex:5.6.7]{\pageref*{ex:5.6.7}}, \hyperref[thm:5.6.8]{\pageref*{thm:5.6.8}}\par
\par\addvspace{2pt}
\par\noindent\hangindent=1em\hspace*{0em}\phantomsection\label{idx:2}验收抽样（Acceptance sampling），\hyperref[idxocc:2:333]{\pageref*{idxocc:2:333}}, \hyperref[ex:8.2.9]{\pageref*{ex:8.2.9}}\par
\par\noindent\hangindent=2em\hspace*{1em}\phantomsection\label{idx:3}超几何（hypergeometric），\hyperref[ex:3.2.1]{\pageref*{ex:3.2.1}}\par
\par\addvspace{2pt}
\par\noindent\hangindent=1em\hspace*{0em}\phantomsection\label{idx:4}行动空间（Action space），\hyperref[idxocc:4:355]{\pageref*{idxocc:4:355}}\par
\par\addvspace{2pt}
\needspace{3\baselineskip}
\par\noindent\hangindent=1em\hspace*{0em}\phantomsection\label{idx:5}算法（Algorithm）\par
\par\noindent\hangindent=2em\hspace*{1em}\phantomsection\label{idx:6}接受／拒绝（accept-reject），\hyperref[ex:5.6.7]{\pageref*{ex:5.6.7}}, \hyperref[thm:5.6.8]{\pageref*{thm:5.6.8}}\par
\par\noindent\hangindent=2em\hspace*{1em}\phantomsection\label{idx:7}EM，\hyperref[sec:7.2.4]{\pageref*{sec:7.2.4}}, \hyperref[sec:7.5.4]{\pageref*{sec:7.5.4}}\par
\par\noindent\hangindent=2.8em\hspace*{1.8em}\phantomsection\label{idx:8}完整数据（complete-data），\hyperref[idxocc:8:295]{\pageref*{idxocc:8:295}}\par
\par\noindent\hangindent=2.8em\hspace*{1.8em}\phantomsection\label{idx:9}完整数据似然（complete-data likelihood），\hyperref[eq:7.2.11]{\pageref*{eq:7.2.11}}\par
\par\noindent\hangindent=2.8em\hspace*{1.8em}\phantomsection\label{idx:10}收敛（convergence），\hyperref[thm:7.2.20]{\pageref*{thm:7.2.20}}, \hyperref[thm:7.5.2]{\pageref*{thm:7.5.2}}\par
\par\noindent\hangindent=2.8em\hspace*{1.8em}\phantomsection\label{idx:11}不完整数据（incomplete-data），\hyperref[idxocc:11:295]{\pageref*{idxocc:11:295}}\par
\par\noindent\hangindent=2.8em\hspace*{1.8em}\phantomsection\label{idx:12}不完整数据似然（incomplete-data likelihood），\hyperref[eq:7.2.13]{\pageref*{eq:7.2.13}}\par
\par\noindent\hangindent=2.8em\hspace*{1.8em}\phantomsection\label{idx:13}缺失数据（missing data），\hyperref[sec:7.2.4]{\pageref*{sec:7.2.4}}\par
\par\noindent\hangindent=2em\hspace*{1em}\phantomsection\label{idx:14}重要性抽样（importance sampling），\hyperref[exercise:5.64]{\pageref*{exercise:5.64}}\par
\par\noindent\hangindent=2em\hspace*{1em}\phantomsection\label{idx:15}Metropolis，\hyperref[idxocc:15:238]{\pageref*{idxocc:15:238}}\par
\par\addvspace{2pt}
\needspace{3\baselineskip}
\par\noindent\hangindent=1em\hspace*{0em}\phantomsection\label{idx:16}算法（Algorithms）\par
\par\noindent\hangindent=2em\hspace*{1em}\phantomsection\label{idx:17}马尔可夫链蒙特卡洛（Markov chain Monte Carlo），\hyperref[sec:5.8.5]{\pageref*{sec:5.8.5}}\par
\par\addvspace{2pt}
\par\noindent\hangindent=1em\hspace*{0em}\phantomsection\label{idx:18}辅助统计量（Ancillary statistic），\hyperref[def:6.2.16]{\pageref*{def:6.2.16}}, \hyperref[sec:6.6.2]{\pageref*{sec:6.6.2}}\par
\par\noindent\hangindent=2em\hspace*{1em}\phantomsection\label{idx:19}指数（exponential），\hyperref[ex:6.2.26]{\pageref*{ex:6.2.26}}\par
\par\noindent\hangindent=2em\hspace*{1em}\phantomsection\label{idx:20}位置族（location family），\hyperref[ex:6.2.18]{\pageref*{ex:6.2.18}}, \hyperref[exercise:6.11]{\pageref*{exercise:6.11}}\par
\par\noindent\hangindent=2em\hspace*{1em}\phantomsection\label{idx:21}位置尺度族（location–scale family），\hyperref[exercise:6.40]{\pageref*{exercise:6.40}}\par
\par\noindent\hangindent=2em\hspace*{1em}\phantomsection\label{idx:22}样本量（sample size），\hyperref[exercise:6.12]{\pageref*{exercise:6.12}}\par
\par\noindent\hangindent=2em\hspace*{1em}\phantomsection\label{idx:23}尺度族（scale family），\hyperref[ex:6.2.19]{\pageref*{ex:6.2.19}}\par
\par\noindent\hangindent=2em\hspace*{1em}\phantomsection\label{idx:24}均匀（uniform），\hyperref[ex:6.2.17]{\pageref*{ex:6.2.17}}\par
\par\noindent\hangindent=2em\hspace*{1em}\phantomsection\label{idx:25}用途（usefulness），\hyperref[ex:6.2.20]{\pageref*{ex:6.2.20}}\par
\par\addvspace{2pt}
\needspace{3\baselineskip}
\par\noindent\hangindent=1em\hspace*{0em}\phantomsection\label{idx:26}方差分析（ANOVA）\par
\par\noindent\hangindent=2em\hspace*{1em}\phantomsection\label{idx:27}假设（assumptions），\hyperref[sec:11.2.1]{\pageref*{sec:11.2.1}}\par
\par\noindent\hangindent=2em\hspace*{1em}\phantomsection\label{idx:28}Bonferroni不等式（Bonferroni Inequality），\hyperref[ex:11.2.9]{\pageref*{ex:11.2.9}}\par
\par\noindent\hangindent=2em\hspace*{1em}\phantomsection\label{idx:29}单元均值模型（cell means model），\hyperref[idxocc:29:454]{\pageref*{idxocc:29:454}}\par
\par\noindent\hangindent=2em\hspace*{1em}\phantomsection\label{idx:30}置信区间（confidence interval），\hyperref[eq:11.2.8]{\pageref*{eq:11.2.8}}\par
\par\noindent\hangindent=2em\hspace*{1em}\phantomsection\label{idx:31}对比平方和（contrast sum of squares），\hyperref[eq:11.2.18]{\pageref*{eq:11.2.18}}\par
\par\noindent\hangindent=2em\hspace*{1em}\phantomsection\label{idx:32}对比（contrasts），\hyperref[def:11.2.4]{\pageref*{def:11.2.4}}, \hyperref[ex:11.2.6]{\pageref*{ex:11.2.6}}, \hyperref[idxocc:32:459]{\pageref*{idxocc:32:459}}, \hyperref[sec:11.2.5]{\pageref*{sec:11.2.5}}\par
\par\noindent\hangindent=2em\hspace*{1em}\phantomsection\label{idx:33}数据窥探（data snooping），\hyperref[idxocc:33:464]{\pageref*{idxocc:33:464}}, \hyperref[exercise:11.18]{\pageref*{exercise:11.18}}\par
\par\noindent\hangindent=2em\hspace*{1em}\phantomsection\label{idx:34}方差相等（equality of variance），\hyperref[item:11.2.assumption.iii]{\pageref*{item:11.2.assumption.iii}}\par
\par\noindent\hangindent=2em\hspace*{1em}\phantomsection\label{idx:35}均方的期望（expected mean squares），\hyperref[exercise:11.22]{\pageref*{exercise:11.22}}\par
\par\noindent\hangindent=2em\hspace*{1em}\phantomsection\label{idx:36}F检验（F test），\hyperref[eq:11.2.14]{\pageref*{eq:11.2.14}}\par
\par\noindent\hangindent=2.8em\hspace*{1.8em}\phantomsection\label{idx:37}与t检验的关系（relation to t test），\hyperref[exercise:11.12]{\pageref*{exercise:11.12}}\par
\par\noindent\hangindent=2em\hspace*{1em}\phantomsection\label{idx:38}可识别性（identifiability），\hyperref[def:11.2.2]{\pageref*{def:11.2.2}}\par
\par\noindent\hangindent=2em\hspace*{1em}\phantomsection\label{idx:39}因子的水平（levels of a factor），\hyperref[idxocc:39:453]{\pageref*{idxocc:39:453}}\par
\par\noindent\hangindent=2em\hspace*{1em}\phantomsection\label{idx:40}线性组合（linear combinations），\hyperref[sec:11.2.3]{\pageref*{sec:11.2.3}}\par
\par\noindent\hangindent=2em\hspace*{1em}\phantomsection\label{idx:41}似然比检验（LRT），\hyperref[exercise:11.13]{\pageref*{exercise:11.13}}\par
\par\noindent\hangindent=2em\hspace*{1em}\phantomsection\label{idx:42}零假设（null hypothesis），\hyperref[eq:11.2.3]{\pageref*{eq:11.2.3}}\par
\par\noindent\hangindent=2em\hspace*{1em}\phantomsection\label{idx:43}单因素（oneway），\hyperref[idxocc:43:452]{\pageref*{idxocc:43:452}}\par
\par\noindent\hangindent=2em\hspace*{1em}\phantomsection\label{idx:44}正交对比（orthogonal contrast），\hyperref[idxocc:44:465]{\pageref*{idxocc:44:465}}\par
\par\noindent\hangindent=2em\hspace*{1em}\phantomsection\label{idx:45}其他模型（other models），\hyperref[sec:11.5.4]{\pageref*{sec:11.5.4}}\par
\par\noindent\hangindent=2em\hspace*{1em}\phantomsection\label{idx:46}过参数化模型（overparameterized model），\hyperref[eq:11.2.2]{\pageref*{eq:11.2.2}}\par
\par\noindent\hangindent=2em\hspace*{1em}\phantomsection\label{idx:47}Scheffé方法（Scheffé’s method），\hyperref[thm:11.2.10]{\pageref*{thm:11.2.10}}, \hyperref[exercise:11.14]{\pageref*{exercise:11.14}}\par
\par\noindent\hangindent=2em\hspace*{1em}\phantomsection\label{idx:48}Stein估计（Stein estimation），\hyperref[sec:11.5.6]{\pageref*{sec:11.5.6}}\par
\par\noindent\hangindent=2em\hspace*{1em}\phantomsection\label{idx:49}充分统计量（sufficient statistic），\hyperref[item:11.6.a]{\pageref*{item:11.6.a}}\par
\par\noindent\hangindent=2em\hspace*{1em}\phantomsection\label{idx:50}平方和（sum of squares），\hyperref[thm:11.2.11]{\pageref*{thm:11.2.11}}\par
\par\noindent\hangindent=2em\hspace*{1em}\phantomsection\label{idx:51}t检验（t test），\hyperref[eq:11.2.7]{\pageref*{eq:11.2.7}}\par
\par\noindent\hangindent=2.8em\hspace*{1.8em}\phantomsection\label{idx:52}与F检验的关系（relation to F test），\hyperref[exercise:11.12]{\pageref*{exercise:11.12}}\par
\par\noindent\hangindent=2em\hspace*{1em}\phantomsection\label{idx:53}表（table），\hyperref[ex:11.2.12]{\pageref*{ex:11.2.12}}\par
\par\noindent\hangindent=2em\hspace*{1em}\phantomsection\label{idx:54}数据变换（transformation of data），\hyperref[idxocc:54:455]{\pageref*{idxocc:54:455}}, \hyperref[exercise:11.1]{\pageref*{exercise:11.1}}\par
\par\noindent\hangindent=2em\hspace*{1em}\phantomsection\label{idx:55}处理均值（treatment means），\hyperref[idxocc:55:453]{\pageref*{idxocc:55:453}}\par
\par\noindent\hangindent=2em\hspace*{1em}\phantomsection\label{idx:56}并交检验（union–intersection test），\hyperref[thm:11.2.5]{\pageref*{thm:11.2.5}}, \hyperref[sec:11.2.4]{\pageref*{sec:11.2.4}}\par
\par\noindent\hangindent=2em\hspace*{1em}\phantomsection\label{idx:57}方差估计量（variance estimator），\hyperref[eq:11.2.5]{\pageref*{eq:11.2.5}}\par
\par\addvspace{2pt}
\needspace{3\baselineskip}
\par\noindent\hangindent=1em\hspace*{0em}\phantomsection\label{idx:58}近似（Approximation）\par
\par\noindent\hangindent=2em\hspace*{1em}\phantomsection\label{idx:59}以正态近似二项（binomial by normal），\hyperref[ex:3.3.2]{\pageref*{ex:3.3.2}}\par
\par\noindent\hangindent=2em\hspace*{1em}\phantomsection\label{idx:60}以泊松近似二项（binomial by Poisson），\hyperref[ex:2.3.13]{\pageref*{ex:2.3.13}}, \hyperref[eq:3.2.8]{\pageref*{eq:3.2.8}}\par
\par\noindent\hangindent=2em\hspace*{1em}\phantomsection\label{idx:61}以二项近似超几何（hypergeometric by binomial），\hyperref[idxocc:61:137]{\pageref*{idxocc:61:137}}\par
\par\noindent\hangindent=2em\hspace*{1em}\phantomsection\label{idx:62}以泊松近似超几何（hypergeometric by Poisson），\hyperref[itm:exercise:3.11:b]{\pageref*{itm:exercise:3.11:b}}\par
\par\noindent\hangindent=2em\hspace*{1em}\phantomsection\label{idx:63}以正态近似负二项（negative binomial by normal），\hyperref[ex:5.5.16]{\pageref*{ex:5.5.16}}\par
\par\addvspace{2pt}
\par\noindent\hangindent=1em\hspace*{0em}\phantomsection\label{idx:64}算术平均数（Arithmetic mean），\hyperref[ex:4.7.8]{\pageref*{ex:4.7.8}}\par
\par\addvspace{2pt}
\par\noindent\hangindent=1em\hspace*{0em}\phantomsection\label{idx:65}渐近效率（Asymptotic efficiency），\hyperref[def:10.1.11]{\pageref*{def:10.1.11}}\par
\par\addvspace{2pt}
\par\noindent\hangindent=1em\hspace*{0em}\phantomsection\label{idx:66}渐近相对效率（Asymptotic Relative Efficiency），\hyperref[def:10.1.16]{\pageref*{def:10.1.16}}\par
\par\noindent\hangindent=2em\hspace*{1em}\phantomsection\label{idx:67}Huber估计量（Huber estimator），\hyperref[ex:10.2.7]{\pageref*{ex:10.2.7}}\par
\par\noindent\hangindent=2em\hspace*{1em}\phantomsection\label{idx:68}中位数相对于均值（median to mean），\hyperref[ex:10.2.4]{\pageref*{ex:10.2.4}}\par
\par\addvspace{2pt}
\needspace{3\baselineskip}
\par\noindent\hangindent=1em\hspace*{0em}\phantomsection\label{idx:69}渐近方差（Asymptotic variance）\par
\par\noindent\hangindent=2em\hspace*{1em}\phantomsection\label{idx:70}与影响函数的关系（relation to influence function），\hyperref[idxocc:70:450]{\pageref*{idxocc:70:450}}\par
\par\addvspace{2pt}
\par\noindent\hangindent=1em\hspace*{0em}\phantomsection\label{idx:71}渐近方差（另见\hyperref[idx:1062]{方差}）（Asymptotic variance (see also Variance)），\hyperref[idxocc:71:410]{\pageref*{idxocc:71:410}}\par
\par\addvspace{2pt}
\par\noindent\hangindent=1em\hspace*{0em}\phantomsection\label{idx:72}自回归模型（Autoregressive model），\hyperref[itm:exercise:10.19:c]{\pageref*{itm:exercise:10.19:c}}\par
\par\addvspace{2pt}
\par\noindent\hangindent=1em\hspace*{0em}\phantomsection\label{idx:73}连续性公理（Axiom of Continuity），\hyperref[exercise:1.12]{\pageref*{exercise:1.12}}\par
\par\addvspace{2pt}
\par\noindent\hangindent=1em\hspace*{0em}\phantomsection\label{idx:74}可数可加性公理（Axiom of Countable Additivity），\hyperref[idxocc:74:38]{\pageref*{idxocc:74:38}}, \hyperref[exercise:1.12]{\pageref*{exercise:1.12}}\par
\par\addvspace{2pt}
\par\noindent\hangindent=1em\hspace*{0em}\phantomsection\label{idx:75}有限可加性公理（Axiom of Finite Additivity），\hyperref[idxocc:75:39]{\pageref*{idxocc:75:39}}, \hyperref[exercise:1.12]{\pageref*{exercise:1.12}}\par
\par\addvspace{2pt}
\par\noindent\hangindent=1em\hspace*{0em}\phantomsection\label{idx:76}Basu定理（Basu’s Theorem），\hyperref[thm:6.2.24]{\pageref*{thm:6.2.24}}\par
\par\noindent\hangindent=2em\hspace*{1em}\phantomsection\label{idx:77}逆命题（converse），\hyperref[sec:6.6.1]{\pageref*{sec:6.6.1}}\par
\par\addvspace{2pt}
\needspace{3\baselineskip}
\par\noindent\hangindent=1em\hspace*{0em}\phantomsection\label{idx:78}贝叶斯估计量（Bayes estimator）\par
\par\noindent\hangindent=2em\hspace*{1em}\phantomsection\label{idx:79}伯努利（Bernoulli），\hyperref[ex:7.2.14]{\pageref*{ex:7.2.14}}, \hyperref[ex:7.3.5]{\pageref*{ex:7.3.5}}\par
\par\noindent\hangindent=2em\hspace*{1em}\phantomsection\label{idx:80}正态（normal），\hyperref[ex:7.2.16]{\pageref*{ex:7.2.16}}, \hyperref[exercise:7.22]{\pageref*{exercise:7.22}}\par
\par\noindent\hangindent=2em\hspace*{1em}\phantomsection\label{idx:81}正态方差（normal variance），\hyperref[exercise:7.23]{\pageref*{exercise:7.23}}\par
\par\noindent\hangindent=2em\hspace*{1em}\phantomsection\label{idx:82}与无偏估计量的关系（relation to unbiased estimator），\hyperref[sec:7.5.2]{\pageref*{sec:7.5.2}}\par
\par\addvspace{2pt}
\par\noindent\hangindent=1em\hspace*{0em}\phantomsection\label{idx:83}贝叶斯风险（Bayes risk），\hyperref[idxocc:83:316]{\pageref*{idxocc:83:316}}\par
\par\addvspace{2pt}
\needspace{3\baselineskip}
\par\noindent\hangindent=1em\hspace*{0em}\phantomsection\label{idx:84}贝叶斯规则（Bayes rule）\par
\par\noindent\hangindent=2em\hspace*{1em}\phantomsection\label{idx:85}先验分布（Prior distribution），\hyperref[idxocc:85:316]{\pageref*{idxocc:85:316}}\par
\par\addvspace{2pt}
\par\noindent\hangindent=1em\hspace*{0em}\phantomsection\label{idx:86}贝叶斯公式（Bayes’ Rule），\hyperref[thm:1.3.5]{\pageref*{thm:1.3.5}}\par
\par\addvspace{2pt}
\par\noindent\hangindent=1em\hspace*{0em}\phantomsection\label{idx:87}Behrens–Fisher问题（Behrens–Fisher problem），\hyperref[exercise:8.42]{\pageref*{exercise:8.42}}, \hyperref[idxocc:87:455]{\pageref*{idxocc:87:455}}, \hyperref[idxocc:87:500]{\pageref*{idxocc:87:500}}\par
\sourcepageanchor{557}
\par\addvspace{2pt}
\par\noindent\hangindent=1em\hspace*{0em}\phantomsection\label{idx:88}伯努利试验（见\hyperref[idx:252]{分布：伯努利}）（Bernoulli trials (see Distribution Bernoulli)），\hyperref[idxocc:88:104]{\pageref*{idxocc:88:104}}\par
\par\addvspace{2pt}
\par\noindent\hangindent=1em\hspace*{0em}\phantomsection\label{idx:89}贝塔函数（另见\hyperref[idx:254]{分布：贝塔}）（Beta function (see also Distribution beta)），\hyperref[eq:3.3.17]{\pageref*{eq:3.3.17}}\par
\par\addvspace{2pt}
\par\noindent\hangindent=1em\hspace*{0em}\phantomsection\label{idx:90}偏差（Bias），\hyperref[def:7.3.2]{\pageref*{def:7.3.2}}\par
\par\addvspace{2pt}
\par\noindent\hangindent=1em\hspace*{0em}\phantomsection\label{idx:91}二项式系数（Binomial coefficients），\hyperref[def:1.2.17]{\pageref*{def:1.2.17}}\par
\par\addvspace{2pt}
\par\noindent\hangindent=1em\hspace*{0em}\phantomsection\label{idx:92}二项式定理（Binomial Theorem），\hyperref[thm:3.2.2]{\pageref*{thm:3.2.2}}\par
\par\addvspace{2pt}
\par\noindent\hangindent=1em\hspace*{0em}\phantomsection\label{idx:93}Birnbaum定理（Birnbaum’s Theorem），\hyperref[thm:6.3.6]{\pageref*{thm:6.3.6}}, \hyperref[exercise:6.33]{\pageref*{exercise:6.33}}\par
\par\addvspace{2pt}
\par\noindent\hangindent=1em\hspace*{0em}\phantomsection\label{idx:94}自助法（Bootstrap），\hyperref[sec:10.1.4]{\pageref*{sec:10.1.4}}\par
\par\noindent\hangindent=2em\hspace*{1em}\phantomsection\label{idx:95}非参数（nonparametric），\hyperref[ex:10.1.19]{\pageref*{ex:10.1.19}}\par
\par\noindent\hangindent=2em\hspace*{1em}\phantomsection\label{idx:96}参数（parametric），\hyperref[ex:5.6.6]{\pageref*{ex:5.6.6}}, \hyperref[idxocc:96:418]{\pageref*{idxocc:96:418}}\par
\par\addvspace{2pt}
\par\noindent\hangindent=1em\hspace*{0em}\phantomsection\label{idx:97}博雷尔域（见\hyperref[idx:972]{\ensuremath{\sigma}代数}）（Borel field (see Sigma algebra)），\hyperref[def:1.2.1]{\pageref*{def:1.2.1}}\par
\par\addvspace{2pt}
\par\noindent\hangindent=1em\hspace*{0em}\phantomsection\label{idx:98}博雷尔悖论（Borel Paradox），\hyperref[sec:4.9.3]{\pageref*{sec:4.9.3}}\par
\par\addvspace{2pt}
\par\noindent\hangindent=1em\hspace*{0em}\phantomsection\label{idx:99}Box–Muller变换（Box-Muller transformation），\hyperref[ex:5.6.4]{\pageref*{ex:5.6.4}}\par
\par\addvspace{2pt}
\par\noindent\hangindent=1em\hspace*{0em}\phantomsection\label{idx:100}崩溃值（Breakdown value），\hyperref[def:10.2.2]{\pageref*{def:10.2.2}}, \hyperref[exercise:10.20]{\pageref*{exercise:10.20}}\par
\par\addvspace{2pt}
\par\noindent\hangindent=1em\hspace*{0em}\phantomsection\label{idx:101}Carleman条件（Carleman’s Condition），\hyperref[idxocc:101:99]{\pageref*{idxocc:101:99}}\par
\par\addvspace{2pt}
\par\noindent\hangindent=1em\hspace*{0em}\phantomsection\label{idx:102}中心极限定理（Central Limit Theorem (CLT)），\hyperref[thm:5.5.14]{\pageref*{thm:5.5.14}}, \hyperref[thm:5.5.15]{\pageref*{thm:5.5.15}}, \hyperref[sec:5.8.1]{\pageref*{sec:5.8.1}}\par
\par\addvspace{2pt}
\par\noindent\hangindent=1em\hspace*{0em}\phantomsection\label{idx:103}特征函数（Characteristic function），\hyperref[idxocc:103:100]{\pageref*{idxocc:103:100}}\par
\par\noindent\hangindent=2em\hspace*{1em}\phantomsection\label{idx:104}收敛（convergence），\hyperref[thm:2.6.1]{\pageref*{thm:2.6.1}}\par
\par\noindent\hangindent=2em\hspace*{1em}\phantomsection\label{idx:105}在中心极限定理中的用途（use in CLT），\hyperref[thm:5.5.15]{\pageref*{thm:5.5.15}}\par
\par\addvspace{2pt}
\par\noindent\hangindent=1em\hspace*{0em}\phantomsection\label{idx:106}独立性卡方检验（Chi squared test of independence），\hyperref[exercise:10.31]{\pageref*{exercise:10.31}}\par
\par\addvspace{2pt}
\par\noindent\hangindent=1em\hspace*{0em}\phantomsection\label{idx:107}Cochran定理（Cochran’s Theorem），\hyperref[idxocc:107:465]{\pageref*{idxocc:107:465}}, \hyperref[sec:11.5.1]{\pageref*{sec:11.5.1}}\par
\par\addvspace{2pt}
\par\noindent\hangindent=1em\hspace*{0em}\phantomsection\label{idx:108}决定系数（Coefficient of determination），\hyperref[idxocc:108:481]{\pageref*{idxocc:108:481}}\par
\par\addvspace{2pt}
\par\noindent\hangindent=1em\hspace*{0em}\phantomsection\label{idx:109}巧合（Coincidences），\hyperref[exercise:3.9]{\pageref*{exercise:3.9}}\par
\par\addvspace{2pt}
\par\noindent\hangindent=1em\hspace*{0em}\phantomsection\label{idx:110}完备性（Completeness），\hyperref[def:6.2.21]{\pageref*{def:6.2.21}}\par
\par\noindent\hangindent=2em\hspace*{1em}\phantomsection\label{idx:111}辅助性（ancillarity），\hyperref[thm:6.2.24]{\pageref*{thm:6.2.24}}, \hyperref[thm:6.6.2]{\pageref*{thm:6.6.2}}\par
\par\noindent\hangindent=2em\hspace*{1em}\phantomsection\label{idx:112}Basu定理（Basu’s Theorem），\hyperref[thm:6.2.24]{\pageref*{thm:6.2.24}}\par
\par\noindent\hangindent=2em\hspace*{1em}\phantomsection\label{idx:113}最优无偏估计量（best unbiased estimator），\hyperref[thm:7.3.23]{\pageref*{thm:7.3.23}}\par
\par\noindent\hangindent=2em\hspace*{1em}\phantomsection\label{idx:114}二项（binomial），\hyperref[ex:6.2.22]{\pageref*{ex:6.2.22}}\par
\par\noindent\hangindent=2em\hspace*{1em}\phantomsection\label{idx:115}指数（exponential），\hyperref[ex:6.2.26]{\pageref*{ex:6.2.26}}, \hyperref[exercise:6.30]{\pageref*{exercise:6.30}}\par
\par\noindent\hangindent=2em\hspace*{1em}\phantomsection\label{idx:116}指数族（exponential family），\hyperref[thm:6.2.25]{\pageref*{thm:6.2.25}}\par
\par\noindent\hangindent=2em\hspace*{1em}\phantomsection\label{idx:117}独立性（independence），\hyperref[thm:6.2.24]{\pageref*{thm:6.2.24}}\par
\par\noindent\hangindent=2em\hspace*{1em}\phantomsection\label{idx:118}泊松（Poisson），\hyperref[exercise:6.18]{\pageref*{exercise:6.18}}\par
\par\noindent\hangindent=2em\hspace*{1em}\phantomsection\label{idx:119}充分性（sufficiency），\hyperref[thm:6.2.24]{\pageref*{thm:6.2.24}}\par
\par\noindent\hangindent=2em\hspace*{1em}\phantomsection\label{idx:120}均匀（uniform），\hyperref[ex:6.2.23]{\pageref*{ex:6.2.23}}\par
\par\addvspace{2pt}
\par\noindent\hangindent=1em\hspace*{0em}\phantomsection\label{idx:121}凹函数（Concave function），\hyperref[def:4.7.6]{\pageref*{def:4.7.6}}\par
\par\addvspace{2pt}
\par\noindent\hangindent=1em\hspace*{0em}\phantomsection\label{idx:122}条件概率（Conditional probability），\hyperref[def:1.3.2]{\pageref*{def:1.3.2}}\par
\par\noindent\hangindent=2em\hspace*{1em}\phantomsection\label{idx:123}乘法法则（multiplication rule），\hyperref[eq:1.3.3]{\pageref*{eq:1.3.3}}\par
\par\noindent\hangindent=2em\hspace*{1em}\phantomsection\label{idx:124}性质（properties），\hyperref[exercise:1.38]{\pageref*{exercise:1.38}}\par
\par\addvspace{2pt}
\par\noindent\hangindent=1em\hspace*{0em}\phantomsection\label{idx:125}条件性原则（Conditionality Principle），\hyperref[idxocc:125:269]{\pageref*{idxocc:125:269}}\par
\par\noindent\hangindent=2em\hspace*{1em}\phantomsection\label{idx:126}解释（interpretations），\hyperref[ex:6.3.7]{\pageref*{ex:6.3.7}}\par
\par\addvspace{2pt}
\par\noindent\hangindent=1em\hspace*{0em}\phantomsection\label{idx:127}置信系数（Confidence coefficient），\hyperref[def:9.1.5]{\pageref*{def:9.1.5}}, \hyperref[idxocc:127:387]{\pageref*{idxocc:127:387}}\par
\par\noindent\hangindent=2em\hspace*{1em}\phantomsection\label{idx:128}与风险的关系（relation to risk），\hyperref[idxocc:128:395]{\pageref*{idxocc:128:395}}\par
\par\addvspace{2pt}
\needspace{3\baselineskip}
\par\noindent\hangindent=1em\hspace*{0em}\phantomsection\label{idx:129}置信区间（Confidence interval）\par
\needspace{3\baselineskip}
\par\noindent\hangindent=2em\hspace*{1em}\phantomsection\label{idx:130}长度（length）\par
\par\noindent\hangindent=2.8em\hspace*{1.8em}\phantomsection\label{idx:131}正态（normal），\hyperref[ex:9.3.3]{\pageref*{ex:9.3.3}}\par
\par\noindent\hangindent=2.8em\hspace*{1.8em}\phantomsection\label{idx:132}最优（optimal），\hyperref[thm:9.3.2]{\pageref*{thm:9.3.2}}\par
\par\noindent\hangindent=2.8em\hspace*{1.8em}\phantomsection\label{idx:133}与错误覆盖的关系（relation to false coverage），\hyperref[thm:9.3.9]{\pageref*{thm:9.3.9}}\par
\par\noindent\hangindent=2.8em\hspace*{1.8em}\phantomsection\label{idx:134}与最高后验密度的关系（relation to HPD），\hyperref[idxocc:134:389]{\pageref*{idxocc:134:389}}\par
\par\noindent\hangindent=2.8em\hspace*{1.8em}\phantomsection\label{idx:135}单峰（unimodal），\hyperref[thm:9.3.2]{\pageref*{thm:9.3.2}}\par
\par\noindent\hangindent=2em\hspace*{1em}\phantomsection\label{idx:136}方差分析（ANOVA），\hyperref[eq:11.2.8]{\pageref*{eq:11.2.8}}\par
\par\noindent\hangindent=2em\hspace*{1em}\phantomsection\label{idx:137}近似（approximate），\hyperref[ex:10.4.4]{\pageref*{ex:10.4.4}}\par
\par\noindent\hangindent=2.8em\hspace*{1.8em}\phantomsection\label{idx:138}二项（binomial），\hyperref[ex:10.4.1]{\pageref*{ex:10.4.1}}\par
\par\noindent\hangindent=2.8em\hspace*{1.8em}\phantomsection\label{idx:139}最大似然（maximum likelihood），\hyperref[eq:10.4.2]{\pageref*{eq:10.4.2}}\par
\par\noindent\hangindent=2.8em\hspace*{1.8em}\phantomsection\label{idx:140}负二项（negative binomial），\hyperref[ex:10.4.9]{\pageref*{ex:10.4.9}}\par
\par\noindent\hangindent=2.8em\hspace*{1.8em}\phantomsection\label{idx:141}泊松（Poisson），\hyperref[ex:10.4.5]{\pageref*{ex:10.4.5}}\par
\needspace{3\baselineskip}
\par\noindent\hangindent=2em\hspace*{1em}\phantomsection\label{idx:142}二项（binomial）\par
\par\noindent\hangindent=2.8em\hspace*{1.8em}\phantomsection\label{idx:143}与F分布的关系（relation to F distribution），\hyperref[exercise:9.21]{\pageref*{exercise:9.21}}\par
\par\noindent\hangindent=2em\hspace*{1em}\phantomsection\label{idx:144}生物等效性（bioequivalence），\hyperref[exercise:9.33]{\pageref*{exercise:9.33}}\par
\par\noindent\hangindent=2em\hspace*{1em}\phantomsection\label{idx:145}构造（construction），\hyperref[ex:9.2.11]{\pageref*{ex:9.2.11}}, \hyperref[idxocc:145:381]{\pageref*{idxocc:145:381}}, \hyperref[exercise:9.17]{\pageref*{exercise:9.17}}\par
\par\noindent\hangindent=2em\hspace*{1em}\phantomsection\label{idx:146}决策论（decision theoretic），\hyperref[eq:9.3.4]{\pageref*{eq:9.3.4}}\par
\par\noindent\hangindent=2.8em\hspace*{1.8em}\phantomsection\label{idx:147}困难（difficulty），\hyperref[idxocc:147:396]{\pageref*{idxocc:147:396}}\par
\par\noindent\hangindent=2.8em\hspace*{1.8em}\phantomsection\label{idx:148}均值（mean），\hyperref[ex:9.3.13]{\pageref*{ex:9.3.13}}\par
\needspace{3\baselineskip}
\par\noindent\hangindent=2em\hspace*{1em}\phantomsection\label{idx:149}期望长度（expected length）\par
\par\noindent\hangindent=2.8em\hspace*{1.8em}\phantomsection\label{idx:150}正态（normal），\hyperref[ex:9.3.3]{\pageref*{ex:9.3.3}}\par
\par\noindent\hangindent=2em\hspace*{1em}\phantomsection\label{idx:151}指数（exponential），\hyperref[ex:9.2.13]{\pageref*{ex:9.2.13}}\par
\par\noindent\hangindent=2.8em\hspace*{1.8em}\phantomsection\label{idx:152}枢轴（pivotal），\hyperref[ex:9.2.9]{\pageref*{ex:9.2.9}}\par
\par\noindent\hangindent=2em\hspace*{1em}\phantomsection\label{idx:153}Fieller，\hyperref[sec:9.5.3]{\pageref*{sec:9.5.3}}\par
\needspace{3\baselineskip}
\par\noindent\hangindent=2em\hspace*{1em}\phantomsection\label{idx:154}长度（length）\par
\par\noindent\hangindent=2.8em\hspace*{1.8em}\phantomsection\label{idx:155}伽马（gamma），\hyperref[ex:9.3.4]{\pageref*{ex:9.3.4}}\par
\par\noindent\hangindent=2em\hspace*{1em}\phantomsection\label{idx:156}似然（likelihood），\hyperref[eq:9.2.7]{\pageref*{eq:9.2.7}}\par
\par\noindent\hangindent=2em\hspace*{1em}\phantomsection\label{idx:157}位置（location），\hyperref[exercise:9.51]{\pageref*{exercise:9.51}}\par
\needspace{3\baselineskip}
\par\noindent\hangindent=2em\hspace*{1em}\phantomsection\label{idx:158}似然比检验（LRT）\par
\par\noindent\hangindent=2.8em\hspace*{1.8em}\phantomsection\label{idx:159}二项（binomial），\hyperref[exercise:9.6]{\pageref*{exercise:9.6}}\par
\par\noindent\hangindent=2.8em\hspace*{1.8em}\phantomsection\label{idx:160}指数（exponential），\hyperref[ex:9.2.3]{\pageref*{ex:9.2.3}}\par
\needspace{3\baselineskip}
\par\noindent\hangindent=2em\hspace*{1em}\phantomsection\label{idx:161}负二项（negative binomial）\par
\par\noindent\hangindent=2.8em\hspace*{1.8em}\phantomsection\label{idx:162}与F分布的关系（relation to F distribution），\hyperref[exercise:9.22]{\pageref*{exercise:9.22}}\par
\par\noindent\hangindent=2em\hspace*{1em}\phantomsection\label{idx:163}正态（normal），\hyperref[ex:9.1.2]{\pageref*{ex:9.1.2}}, \hyperref[ex:9.2.1]{\pageref*{ex:9.2.1}}\par
\par\noindent\hangindent=2.8em\hspace*{1.8em}\phantomsection\label{idx:164}枢轴（pivotal），\hyperref[ex:9.2.10]{\pageref*{ex:9.2.10}}\par
\par\noindent\hangindent=2em\hspace*{1em}\phantomsection\label{idx:165}单侧（one-sided），\hyperref[idxocc:165:368]{\pageref*{idxocc:165:368}}\par
\par\noindent\hangindent=2.8em\hspace*{1.8em}\phantomsection\label{idx:166}二项（binomial），\hyperref[ex:9.2.5]{\pageref*{ex:9.2.5}}\par
\par\noindent\hangindent=2.8em\hspace*{1.8em}\phantomsection\label{idx:167}正态（normal），\hyperref[ex:9.2.4]{\pageref*{ex:9.2.4}}\par
\par\noindent\hangindent=2em\hspace*{1em}\phantomsection\label{idx:168}泊松（Poisson），\hyperref[ex:9.2.15]{\pageref*{ex:9.2.15}}\par
\par\noindent\hangindent=2em\hspace*{1em}\phantomsection\label{idx:169}均值之比（ratio of means），\hyperref[sec:9.5.3]{\pageref*{sec:9.5.3}}\par
\par\noindent\hangindent=2em\hspace*{1em}\phantomsection\label{idx:170}回归（regression），\hyperref[eq:11.3.37]{\pageref*{eq:11.3.37}}, \hyperref[sec:11.3.6]{\pageref*{sec:11.3.6}}\par
\par\noindent\hangindent=2em\hspace*{1em}\phantomsection\label{idx:171}与集合估计量的关系（relation to set estimator），\hyperref[idxocc:171:369]{\pageref*{idxocc:171:369}}\par
\needspace{3\baselineskip}
\par\noindent\hangindent=2em\hspace*{1em}\phantomsection\label{idx:172}最短（shortest）\par
\par\noindent\hangindent=2.8em\hspace*{1.8em}\phantomsection\label{idx:173}二项（binomial），\hyperref[ex:9.2.11]{\pageref*{ex:9.2.11}}, \hyperref[exercise:9.18]{\pageref*{exercise:9.18}}\par
\needspace{3\baselineskip}
\par\noindent\hangindent=2em\hspace*{1em}\phantomsection\label{idx:174}一致最准确（UMA）\par
\par\noindent\hangindent=2.8em\hspace*{1.8em}\phantomsection\label{idx:175}长度优势（length dominance），\hyperref[exercise:9.45]{\pageref*{exercise:9.45}}\par
\par\noindent\hangindent=2.8em\hspace*{1.8em}\phantomsection\label{idx:176}正态（normal），\hyperref[exercise:9.47]{\pageref*{exercise:9.47}}\par
\par\noindent\hangindent=2em\hspace*{1em}\phantomsection\label{idx:177}均匀（uniform），\hyperref[ex:9.1.6]{\pageref*{ex:9.1.6}}\par
\par\noindent\hangindent=2em\hspace*{1em}\phantomsection\label{idx:178}方差（variance），\hyperref[exercise:9.52]{\pageref*{exercise:9.52}}\par
\par\addvspace{2pt}
\par\noindent\hangindent=1em\hspace*{0em}\phantomsection\label{idx:179}置信区间（另见\hyperref[idx:181]{置信集}；\hyperref[idx:180]{置信程序}；\hyperref[idx:230]{可信集}）（Confidence interval (see also Confidence set; Confidence procedure; Credible Set)），\hyperref[def:9.1.1]{\pageref*{def:9.1.1}}\par
\par\addvspace{2pt}
\par\noindent\hangindent=1em\hspace*{0em}\phantomsection\label{idx:180}置信程序（另见\hyperref[idx:181]{置信集}；\hyperref[idx:129]{置信区间}）（Confidence procedure (see also Confidence set; Confidence interval)），\hyperref[sec:9.5.1]{\pageref*{sec:9.5.1}}\par
\par\addvspace{2pt}
\needspace{3\baselineskip}
\par\noindent\hangindent=1em\hspace*{0em}\phantomsection\label{idx:181}置信集（Confidence set）\par
\par\noindent\hangindent=2em\hspace*{1em}\phantomsection\label{idx:182}变量含误差（EIV），\hyperref[sec:12.2.4]{\pageref*{sec:12.2.4}}\par
\needspace{3\baselineskip}
\par\noindent\hangindent=2em\hspace*{1em}\phantomsection\label{idx:183}频率学派与贝叶斯（frequentist–Bayesian）\par
\par\noindent\hangindent=2.8em\hspace*{1.8em}\phantomsection\label{idx:184}解释（interpretations），\hyperref[sec:9.2.4]{\pageref*{sec:9.2.4}}\par
\needspace{3\baselineskip}
\par\noindent\hangindent=2em\hspace*{1em}\phantomsection\label{idx:185}枢轴（pivotal）\par
\par\noindent\hangindent=2.8em\hspace*{1.8em}\phantomsection\label{idx:186}变量含误差（EIV），\hyperref[eq:12.2.24]{\pageref*{eq:12.2.24}}\par
\par\noindent\hangindent=2em\hspace*{1em}\phantomsection\label{idx:187}与假设检验的关系（relation to hypothesis test），\hyperref[thm:9.2.2]{\pageref*{thm:9.2.2}}\par
\sourcepageanchor{558}
\par\noindent\hangindent=2em\hspace*{1em}\phantomsection\label{idx:188}大小（size），\hyperref[idxocc:188:387]{\pageref*{idxocc:188:387}}\par
\par\noindent\hangindent=2em\hspace*{1em}\phantomsection\label{idx:189}一致最准确（UMA），\hyperref[idxocc:189:391]{\pageref*{idxocc:189:391}}\par
\par\noindent\hangindent=2.8em\hspace*{1.8em}\phantomsection\label{idx:190}正态（normal），\hyperref[ex:9.3.6]{\pageref*{ex:9.3.6}}\par
\par\noindent\hangindent=2.8em\hspace*{1.8em}\phantomsection\label{idx:191}与一致最强检验的关系（relation to UMP test），\hyperref[thm:9.3.5]{\pageref*{thm:9.3.5}}\par
\par\noindent\hangindent=2em\hspace*{1em}\phantomsection\label{idx:192}无偏（unbiased），\hyperref[def:9.3.7]{\pageref*{def:9.3.7}}\par
\par\noindent\hangindent=2.8em\hspace*{1.8em}\phantomsection\label{idx:193}正态（normal），\hyperref[ex:9.3.8]{\pageref*{ex:9.3.8}}\par
\par\noindent\hangindent=2.8em\hspace*{1.8em}\phantomsection\label{idx:194}与无偏检验的关系（relation to unbiased test），\hyperref[idxocc:194:392]{\pageref*{idxocc:194:392}}\par
\par\addvspace{2pt}
\par\noindent\hangindent=1em\hspace*{0em}\phantomsection\label{idx:195}共轭族（Conjugate family），\hyperref[def:7.2.15]{\pageref*{def:7.2.15}}\par
\par\addvspace{2pt}
\par\noindent\hangindent=1em\hspace*{0em}\phantomsection\label{idx:196}相合性（见\hyperref[idx:770]{点估计量的相合性}）（Consistency (see Point estimator consistent)），\hyperref[ex:5.5.3]{\pageref*{ex:5.5.3}}\par
\par\addvspace{2pt}
\par\noindent\hangindent=1em\hspace*{0em}\phantomsection\label{idx:197}列联表（Contingency table），\hyperref[exercise:10.31]{\pageref*{exercise:10.31}}\par
\par\addvspace{2pt}
\par\noindent\hangindent=1em\hspace*{0em}\phantomsection\label{idx:198}连续性校正（Continuity correction），\hyperref[ex:3.3.2]{\pageref*{ex:3.3.2}}\par
\par\addvspace{2pt}
\par\noindent\hangindent=1em\hspace*{0em}\phantomsection\label{idx:199}对比（Contrasts），\hyperref[def:11.2.4]{\pageref*{def:11.2.4}}\par
\par\addvspace{2pt}
\needspace{3\baselineskip}
\par\noindent\hangindent=1em\hspace*{0em}\phantomsection\label{idx:200}收敛（Convergence）\par
\par\noindent\hangindent=2em\hspace*{1em}\phantomsection\label{idx:201}几乎处处（almost sure），\hyperref[def:5.5.6]{\pageref*{def:5.5.6}}\par
\par\noindent\hangindent=2em\hspace*{1em}\phantomsection\label{idx:202}几乎处处（almost surely），\hyperref[sec:5.8.4]{\pageref*{sec:5.8.4}}\par
\par\noindent\hangindent=2em\hspace*{1em}\phantomsection\label{idx:203}遍历定理（ergodic theorem），\hyperref[idxocc:203:250]{\pageref*{idxocc:203:250}}\par
\par\noindent\hangindent=2em\hspace*{1em}\phantomsection\label{idx:204}依分布（in distribution），\hyperref[def:5.5.10]{\pageref*{def:5.5.10}}, \hyperref[thm:5.5.17]{\pageref*{thm:5.5.17}}, \hyperref[sec:10.3.2]{\pageref*{sec:10.3.2}}\par
\par\noindent\hangindent=2.8em\hspace*{1.8em}\phantomsection\label{idx:205}与依概率收敛的关系（relation to convergence in probability），\hyperref[thm:5.5.12]{\pageref*{thm:5.5.12}}\par
\par\noindent\hangindent=2em\hspace*{1em}\phantomsection\label{idx:206}依概率（in probability），\hyperref[def:5.5.1]{\pageref*{def:5.5.1}}\par
\par\noindent\hangindent=2.8em\hspace*{1.8em}\phantomsection\label{idx:207}与依分布收敛的关系（relation to convergence in distribution），\hyperref[thm:5.5.12]{\pageref*{thm:5.5.12}}\par
\par\noindent\hangindent=2.8em\hspace*{1.8em}\phantomsection\label{idx:208}收敛到常数（to a constant），\hyperref[thm:5.5.13]{\pageref*{thm:5.5.13}}\par
\par\noindent\hangindent=2em\hspace*{1em}\phantomsection\label{idx:209}强大数定律（strong law），\hyperref[thm:5.5.9]{\pageref*{thm:5.5.9}}\par
\par\noindent\hangindent=2em\hspace*{1em}\phantomsection\label{idx:210}弱大数定律（weak law），\hyperref[thm:5.5.2]{\pageref*{thm:5.5.2}}\par
\par\noindent\hangindent=2em\hspace*{1em}\phantomsection\label{idx:211}以概率1（with probability one），\hyperref[def:5.5.6]{\pageref*{def:5.5.6}}\par
\par\addvspace{2pt}
\par\noindent\hangindent=1em\hspace*{0em}\phantomsection\label{idx:212}凸函数（Convex function），\hyperref[def:4.7.6]{\pageref*{def:4.7.6}}\par
\par\addvspace{2pt}
\par\noindent\hangindent=1em\hspace*{0em}\phantomsection\label{idx:213}卷积（Convolution），\hyperref[thm:5.2.9]{\pageref*{thm:5.2.9}}\par
\par\addvspace{2pt}
\needspace{3\baselineskip}
\par\noindent\hangindent=1em\hspace*{0em}\phantomsection\label{idx:214}相关（Correlation）\par
\par\noindent\hangindent=2em\hspace*{1em}\phantomsection\label{idx:215}二元（binary），\hyperref[exercise:5.57]{\pageref*{exercise:5.57}}\par
\par\noindent\hangindent=2em\hspace*{1em}\phantomsection\label{idx:216}系数（coefficient），\hyperref[def:4.5.2]{\pageref*{def:4.5.2}}\par
\par\addvspace{2pt}
\par\noindent\hangindent=1em\hspace*{0em}\phantomsection\label{idx:217}相关（另见\hyperref[idx:220]{协方差}）（Correlation (see also Covariance)），\hyperref[def:4.5.2]{\pageref*{def:4.5.2}}, \hyperref[thm:4.5.7]{\pageref*{thm:4.5.7}}\par
\par\addvspace{2pt}
\par\noindent\hangindent=1em\hspace*{0em}\phantomsection\label{idx:218}计数公式（Counting formulas），\hyperref[tab:1.2.1]{\pageref*{tab:1.2.1}}\par
\par\addvspace{2pt}
\par\noindent\hangindent=1em\hspace*{0em}\phantomsection\label{idx:219}计数方法（Counting methods），\hyperref[sec:1.2.3]{\pageref*{sec:1.2.3}}\par
\par\addvspace{2pt}
\needspace{3\baselineskip}
\par\noindent\hangindent=1em\hspace*{0em}\phantomsection\label{idx:220}协方差（Covariance）\par
\par\noindent\hangindent=2em\hspace*{1em}\phantomsection\label{idx:221}独立性（independence），\hyperref[thm:4.5.5]{\pageref*{thm:4.5.5}}\par
\par\noindent\hangindent=2em\hspace*{1em}\phantomsection\label{idx:222}线性函数（linear functions），\hyperref[item:12.11.a]{\pageref*{item:12.11.a}}\par
\par\addvspace{2pt}
\par\noindent\hangindent=1em\hspace*{0em}\phantomsection\label{idx:223}协方差（另见\hyperref[idx:214]{相关}）（Covariance (see also Correlation)），\hyperref[def:4.5.1]{\pageref*{def:4.5.1}}\par
\par\addvspace{2pt}
\par\noindent\hangindent=1em\hspace*{0em}\phantomsection\label{idx:224}覆盖概率（Coverage probability），\hyperref[def:9.1.4]{\pageref*{def:9.1.4}}\par
\par\noindent\hangindent=2em\hspace*{1em}\phantomsection\label{idx:225}频率学派与贝叶斯解释（frequentist–Bayesian interpretation），\hyperref[idxocc:225:384]{\pageref*{idxocc:225:384}}\par
\par\noindent\hangindent=2em\hspace*{1em}\phantomsection\label{idx:226}与风险的关系（relation to risk），\hyperref[idxocc:226:395]{\pageref*{idxocc:226:395}}\par
\par\addvspace{2pt}
\par\noindent\hangindent=1em\hspace*{0em}\phantomsection\label{idx:227}Cox悖论（Cox’s paradox），\hyperref[exercise:9.49]{\pageref*{exercise:9.49}}\par
\par\addvspace{2pt}
\par\noindent\hangindent=1em\hspace*{0em}\phantomsection\label{idx:228}Cramér–Rao下界（Cramér–Rao Lower Bound），\hyperref[thm:7.3.9]{\pageref*{thm:7.3.9}}\par
\par\noindent\hangindent=2em\hspace*{1em}\phantomsection\label{idx:229}必要条件（necessary condition），\hyperref[cor:7.3.15]{\pageref*{cor:7.3.15}}\par
\par\addvspace{2pt}
\needspace{3\baselineskip}
\par\noindent\hangindent=1em\hspace*{0em}\phantomsection\label{idx:230}可信集（Credible set）\par
\par\noindent\hangindent=2em\hspace*{1em}\phantomsection\label{idx:231}正态（normal），\hyperref[ex:9.2.18]{\pageref*{ex:9.2.18}}\par
\par\noindent\hangindent=2em\hspace*{1em}\phantomsection\label{idx:232}泊松（Poisson），\hyperref[ex:9.2.16]{\pageref*{ex:9.2.16}}\par
\needspace{3\baselineskip}
\par\noindent\hangindent=2em\hspace*{1em}\phantomsection\label{idx:233}与置信集的关系（relation to confidence set）\par
\par\noindent\hangindent=2.8em\hspace*{1.8em}\phantomsection\label{idx:234}正态（normal），\hyperref[idxocc:234:386]{\pageref*{idxocc:234:386}}, \hyperref[exercise:9.32]{\pageref*{exercise:9.32}}\par
\par\addvspace{2pt}
\par\noindent\hangindent=1em\hspace*{0em}\phantomsection\label{idx:235}可信集（另见\hyperref[idx:181]{置信集}；\hyperref[idx:129]{置信区间}；\hyperref[idx:180]{置信程序}）（Credible set (see also Confidence set, Confidence interval, Confidence procedure），\hyperref[eq:9.2.18]{\pageref*{eq:9.2.18}}\par
\par\addvspace{2pt}
\par\noindent\hangindent=1em\hspace*{0em}\phantomsection\label{idx:236}直积集（Cross product set），\hyperref[idxocc:236:157]{\pageref*{idxocc:236:157}}\par
\par\addvspace{2pt}
\par\noindent\hangindent=1em\hspace*{0em}\phantomsection\label{idx:237}累积分布函数（Cumulative distribution function (cdf)），\hyperref[def:1.5.1]{\pageref*{def:1.5.1}}\par
\par\noindent\hangindent=2em\hspace*{1em}\phantomsection\label{idx:238}独立性（independence），\hyperref[exercise:4.9]{\pageref*{exercise:4.9}}\par
\par\noindent\hangindent=2em\hspace*{1em}\phantomsection\label{idx:239}联合（joint），\hyperref[idxocc:239:152]{\pageref*{idxocc:239:152}}\par
\par\noindent\hangindent=2em\hspace*{1em}\phantomsection\label{idx:240}混合（mixed），\hyperref[ex:1.5.6]{\pageref*{ex:1.5.6}}\par
\par\noindent\hangindent=2em\hspace*{1em}\phantomsection\label{idx:241}单调变换（monotone transformation），\hyperref[thm:2.1.3]{\pageref*{thm:2.1.3}}\par
\par\noindent\hangindent=2em\hspace*{1em}\phantomsection\label{idx:242}性质（properties），\hyperref[thm:1.5.3]{\pageref*{thm:1.5.3}}\par
\par\addvspace{2pt}
\par\noindent\hangindent=1em\hspace*{0em}\phantomsection\label{idx:243}临界点（Cutoff point），\hyperref[idxocc:243:344]{\pageref*{idxocc:243:344}}\par
\par\addvspace{2pt}
\par\noindent\hangindent=1em\hspace*{0em}\phantomsection\label{idx:244}决策论（Decision theory），\hyperref[sec:7.3.4]{\pageref*{sec:7.3.4}}, \hyperref[sec:8.3.5]{\pageref*{sec:8.3.5}}\par
\par\addvspace{2pt}
\par\noindent\hangindent=1em\hspace*{0em}\phantomsection\label{idx:245}自由度（Degrees of freedom），\hyperref[eq:3.3.10]{\pageref*{eq:3.3.10}}\par
\par\addvspace{2pt}
\par\noindent\hangindent=1em\hspace*{0em}\phantomsection\label{idx:246}Delta方法（Delta method），\hyperref[sec:5.5.4]{\pageref*{sec:5.5.4}}\par
\par\addvspace{2pt}
\par\noindent\hangindent=1em\hspace*{0em}\phantomsection\label{idx:247}德摩根律（见\hyperref[idx:959]{集合运算}）（DeMorgan’s Laws (see Set operation)），\hyperref[idxocc:247:36]{\pageref*{idxocc:247:36}}, \hyperref[exercise:1.9]{\pageref*{exercise:1.9}}\par
\par\addvspace{2pt}
\par\noindent\hangindent=1em\hspace*{0em}\phantomsection\label{idx:248}行列式（Determinant），\hyperref[idxocc:248:160]{\pageref*{idxocc:248:160}}, \hyperref[idxocc:248:182]{\pageref*{idxocc:248:182}}\par
\par\addvspace{2pt}
\par\noindent\hangindent=1em\hspace*{0em}\phantomsection\label{idx:249}不交集（Disjoint sets），\hyperref[def:1.1.5]{\pageref*{def:1.1.5}}\par
\par\noindent\hangindent=2em\hspace*{1em}\phantomsection\label{idx:250}独立性（independence），\hyperref[exercise:1.39]{\pageref*{exercise:1.39}}\par
\par\addvspace{2pt}
\needspace{3\baselineskip}
\par\noindent\hangindent=1em\hspace*{0em}\phantomsection\label{idx:251}分布（Distribution）\par
\par\noindent\hangindent=2em\hspace*{1em}\phantomsection\label{idx:252}伯努利（Bernoulli），\hyperref[eq:3.2.3]{\pageref*{eq:3.2.3}}\par
\par\noindent\hangindent=2.8em\hspace*{1.8em}\phantomsection\label{idx:253}极限（limit），\hyperref[exercise:5.44]{\pageref*{exercise:5.44}}\par
\par\noindent\hangindent=2em\hspace*{1em}\phantomsection\label{idx:254}贝塔（beta），\hyperref[eq:3.3.16]{\pageref*{eq:3.3.16}}\par
\par\noindent\hangindent=2.8em\hspace*{1.8em}\phantomsection\label{idx:255}生成（generating），\hyperref[eq:5.6.5]{\pageref*{eq:5.6.5}}, \hyperref[ex:5.6.7]{\pageref*{ex:5.6.7}}, \hyperref[ex:5.6.9]{\pageref*{ex:5.6.9}}, \hyperref[exercise:5.61]{\pageref*{exercise:5.61}}\par
\par\noindent\hangindent=2.8em\hspace*{1.8em}\phantomsection\label{idx:256}乘积（product），\hyperref[ex:4.3.3]{\pageref*{ex:4.3.3}}\par
\par\noindent\hangindent=2.8em\hspace*{1.8em}\phantomsection\label{idx:257}与二项的关系（relation to binomial），\hyperref[exercise:2.40]{\pageref*{exercise:2.40}}\par
\par\noindent\hangindent=2.8em\hspace*{1.8em}\phantomsection\label{idx:258}与F的关系（relation to F），\hyperref[thm:5.3.8]{\pageref*{thm:5.3.8}}\par
\par\noindent\hangindent=2.8em\hspace*{1.8em}\phantomsection\label{idx:259}与伽马的关系（relation to gamma），\hyperref[exercise:4.24]{\pageref*{exercise:4.24}}\par
\par\noindent\hangindent=2em\hspace*{1em}\phantomsection\label{idx:260}贝塔二项（beta–binomial），\hyperref[exercise:4.34]{\pageref*{exercise:4.34}}\par
\par\noindent\hangindent=2em\hspace*{1em}\phantomsection\label{idx:261}二项（binomial），\hyperref[ex:2.1.1]{\pageref*{ex:2.1.1}}, \hyperref[idxocc:261:104]{\pageref*{idxocc:261:104}}\par
\par\noindent\hangindent=2.8em\hspace*{1.8em}\phantomsection\label{idx:262}指数族（exponential family），\hyperref[ex:3.4.1]{\pageref*{ex:3.4.1}}\par
\par\noindent\hangindent=2.8em\hspace*{1.8em}\phantomsection\label{idx:263}生成（generating），\hyperref[ex:5.6.5]{\pageref*{ex:5.6.5}}\par
\par\noindent\hangindent=2.8em\hspace*{1.8em}\phantomsection\label{idx:264}均值（mean），\hyperref[ex:2.2.3]{\pageref*{ex:2.2.3}}\par
\par\noindent\hangindent=2.8em\hspace*{1.8em}\phantomsection\label{idx:265}矩母函数（mgf），\hyperref[ex:2.3.9]{\pageref*{ex:2.3.9}}\par
\par\noindent\hangindent=2.8em\hspace*{1.8em}\phantomsection\label{idx:266}矩（moments），\hyperref[ex:3.4.3]{\pageref*{ex:3.4.3}}\par
\par\noindent\hangindent=2.8em\hspace*{1.8em}\phantomsection\label{idx:267}递推关系（recursion relation），\hyperref[eq:3.2.7]{\pageref*{eq:3.2.7}}\par
\par\noindent\hangindent=2.8em\hspace*{1.8em}\phantomsection\label{idx:268}与贝塔的关系（relation to beta），\hyperref[exercise:2.40]{\pageref*{exercise:2.40}}\par
\par\noindent\hangindent=2.8em\hspace*{1.8em}\phantomsection\label{idx:269}与多项的关系（relation to multinomial），\hyperref[idxocc:269:179]{\pageref*{idxocc:269:179}}\par
\par\noindent\hangindent=2.8em\hspace*{1.8em}\phantomsection\label{idx:270}与负二项的关系（relation to negative binomial），\hyperref[exercise:3.12]{\pageref*{exercise:3.12}}\par
\par\noindent\hangindent=2.8em\hspace*{1.8em}\phantomsection\label{idx:271}与泊松的关系（relation to Poisson），\hyperref[exercise:4.15]{\pageref*{exercise:4.15}}\par
\par\noindent\hangindent=2.8em\hspace*{1.8em}\phantomsection\label{idx:272}方差（variance），\hyperref[ex:2.3.5]{\pageref*{ex:2.3.5}}\par
\par\noindent\hangindent=2em\hspace*{1em}\phantomsection\label{idx:273}二元正态（bivariate normal），\hyperref[def:4.5.10]{\pageref*{def:4.5.10}}, \hyperref[exercise:4.45]{\pageref*{exercise:4.45}}\par
\par\noindent\hangindent=2.8em\hspace*{1.8em}\phantomsection\label{idx:274}相关（correlation），\hyperref[item:10.17.d]{\pageref*{item:10.17.d}}\par
\par\noindent\hangindent=2em\hspace*{1em}\phantomsection\label{idx:275}二元t（bivariate t），\hyperref[idxocc:275:479]{\pageref*{idxocc:275:479}}\par
\par\noindent\hangindent=2em\hspace*{1em}\phantomsection\label{idx:276}柯西（Cauchy），\hyperref[eq:3.3.19]{\pageref*{eq:3.3.19}}, \hyperref[exercise:3.39]{\pageref*{exercise:3.39}}\par
\par\noindent\hangindent=2.8em\hspace*{1.8em}\phantomsection\label{idx:277}生成（generation），\hyperref[exercise:2.12]{\pageref*{exercise:2.12}}\par
\par\noindent\hangindent=2.8em\hspace*{1.8em}\phantomsection\label{idx:278}均值（mean），\hyperref[ex:2.2.4]{\pageref*{ex:2.2.4}}\par
\par\noindent\hangindent=2.8em\hspace*{1.8em}\phantomsection\label{idx:279}与正态的关系（relation to normal），\hyperref[ex:4.3.6]{\pageref*{ex:4.3.6}}, \hyperref[exercise:4.28]{\pageref*{exercise:4.28}}\par
\par\noindent\hangindent=2em\hspace*{1em}\phantomsection\label{idx:280}卡方（chi squared），\hyperref[eq:3.3.10]{\pageref*{eq:3.3.10}}\par
\par\noindent\hangindent=2.8em\hspace*{1.8em}\phantomsection\label{idx:281}生成（generating），\hyperref[eq:5.6.5]{\pageref*{eq:5.6.5}}\par
\par\noindent\hangindent=2.8em\hspace*{1.8em}\phantomsection\label{idx:282}性质（properties），\hyperref[lem:5.3.2]{\pageref*{lem:5.3.2}}\par
\par\noindent\hangindent=2.8em\hspace*{1.8em}\phantomsection\label{idx:283}与负二项的关系（relation to negative binomial），\hyperref[exercise:2.38]{\pageref*{exercise:2.38}}\par
\par\noindent\hangindent=2.8em\hspace*{1.8em}\phantomsection\label{idx:284}与正态的关系（relation to normal），\hyperref[ex:2.1.9]{\pageref*{ex:2.1.9}}\par
\sourcepageanchor{559}
\needspace{3\baselineskip}
\par\noindent\hangindent=2em\hspace*{1em}\phantomsection\label{idx:285}离散（discrete）\par
\par\noindent\hangindent=2.8em\hspace*{1.8em}\phantomsection\label{idx:286}生成（generating），\hyperref[ex:5.6.5]{\pageref*{ex:5.6.5}}\par
\par\noindent\hangindent=2em\hspace*{1em}\phantomsection\label{idx:287}离散均匀（discrete uniform），\hyperref[eq:3.2.1]{\pageref*{eq:3.2.1}}\par
\par\noindent\hangindent=2em\hspace*{1em}\phantomsection\label{idx:288}双指数（double exponential），\hyperref[eq:3.3.23]{\pageref*{eq:3.3.23}}\par
\par\noindent\hangindent=2em\hspace*{1em}\phantomsection\label{idx:289}指数（exponential），\hyperref[eq:3.3.11]{\pageref*{eq:3.3.11}}\par
\par\noindent\hangindent=2.8em\hspace*{1.8em}\phantomsection\label{idx:290}生成（generating），\hyperref[ex:5.6.3]{\pageref*{ex:5.6.3}}\par
\par\noindent\hangindent=2.8em\hspace*{1.8em}\phantomsection\label{idx:291}风险率函数（hazard function），\hyperref[idxocc:291:139]{\pageref*{idxocc:291:139}}\par
\par\noindent\hangindent=2.8em\hspace*{1.8em}\phantomsection\label{idx:292}均值（mean），\hyperref[ex:2.2.2]{\pageref*{ex:2.2.2}}\par
\par\noindent\hangindent=2.8em\hspace*{1.8em}\phantomsection\label{idx:293}无记忆性（memoryless），\hyperref[idxocc:293:114]{\pageref*{idxocc:293:114}}\par
\par\noindent\hangindent=2.8em\hspace*{1.8em}\phantomsection\label{idx:294}递推关系（recursion relation），\hyperref[eq:2.4.7]{\pageref*{eq:2.4.7}}\par
\par\noindent\hangindent=2.8em\hspace*{1.8em}\phantomsection\label{idx:295}与伽马的关系（relation to gamma），\hyperref[exercise:4.38]{\pageref*{exercise:4.38}}\par
\par\noindent\hangindent=2.8em\hspace*{1.8em}\phantomsection\label{idx:296}与均匀的关系（relation to uniform），\hyperref[ex:2.1.4]{\pageref*{ex:2.1.4}}\par
\par\noindent\hangindent=2.8em\hspace*{1.8em}\phantomsection\label{idx:297}方差（variance），\hyperref[ex:2.3.3]{\pageref*{ex:2.3.3}}\par
\needspace{3\baselineskip}
\par\noindent\hangindent=2em\hspace*{1em}\phantomsection\label{idx:298}指数族（exponential family）\par
\par\noindent\hangindent=2.8em\hspace*{1.8em}\phantomsection\label{idx:299}曲线（curved），\hyperref[idxocc:299:208]{\pageref*{idxocc:299:208}}\par
\par\noindent\hangindent=2em\hspace*{1em}\phantomsection\label{idx:300}极值（extreme value），\hyperref[idxocc:300:139]{\pageref*{idxocc:300:139}}\par
\par\noindent\hangindent=2em\hspace*{1em}\phantomsection\label{idx:301}F，\hyperref[def:5.3.6]{\pageref*{def:5.3.6}}\par
\par\noindent\hangindent=2.8em\hspace*{1.8em}\phantomsection\label{idx:302}倒数（reciprocal），\hyperref[thm:5.3.8]{\pageref*{thm:5.3.8}}\par
\par\noindent\hangindent=2.8em\hspace*{1.8em}\phantomsection\label{idx:303}与贝塔的关系（relation to beta），\hyperref[thm:5.3.8]{\pageref*{thm:5.3.8}}\par
\par\noindent\hangindent=2.8em\hspace*{1.8em}\phantomsection\label{idx:304}与t的关系（relation to t），\hyperref[thm:5.3.8]{\pageref*{thm:5.3.8}}\par
\par\noindent\hangindent=2em\hspace*{1em}\phantomsection\label{idx:305}折叠正态（folded normal），\hyperref[exercise:3.20]{\pageref*{exercise:3.20}}\par
\par\noindent\hangindent=2em\hspace*{1em}\phantomsection\label{idx:306}伽马（gamma），\hyperref[eq:3.3.6]{\pageref*{eq:3.3.6}}\par
\par\noindent\hangindent=2.8em\hspace*{1.8em}\phantomsection\label{idx:307}生成（generating），\hyperref[eq:5.6.5]{\pageref*{eq:5.6.5}}, \hyperref[exercise:5.61]{\pageref*{exercise:5.61}}\par
\par\noindent\hangindent=2.8em\hspace*{1.8em}\phantomsection\label{idx:308}矩母函数（mgf），\hyperref[ex:2.3.8]{\pageref*{ex:2.3.8}}\par
\par\noindent\hangindent=2.8em\hspace*{1.8em}\phantomsection\label{idx:309}矩（moments），\hyperref[exercise:3.17]{\pageref*{exercise:3.17}}\par
\par\noindent\hangindent=2.8em\hspace*{1.8em}\phantomsection\label{idx:310}递推关系（recursion relation），\hyperref[thm:3.6.4]{\pageref*{thm:3.6.4}}\par
\par\noindent\hangindent=2.8em\hspace*{1.8em}\phantomsection\label{idx:311}与贝塔的关系（relation to beta），\hyperref[exercise:4.24]{\pageref*{exercise:4.24}}\par
\par\noindent\hangindent=2.8em\hspace*{1.8em}\phantomsection\label{idx:312}与指数的关系（relation to exponential），\hyperref[exercise:4.38]{\pageref*{exercise:4.38}}\par
\par\noindent\hangindent=2.8em\hspace*{1.8em}\phantomsection\label{idx:313}与逆伽马的关系（relation to inverted gamma），\hyperref[ex:2.1.6]{\pageref*{ex:2.1.6}}\par
\par\noindent\hangindent=2.8em\hspace*{1.8em}\phantomsection\label{idx:314}与负二项的关系（relation to negative binomial），\hyperref[exercise:3.18]{\pageref*{exercise:3.18}}\par
\par\noindent\hangindent=2.8em\hspace*{1.8em}\phantomsection\label{idx:315}与泊松的关系（relation to Poisson），\hyperref[ex:3.3.1]{\pageref*{ex:3.3.1}}, \hyperref[exercise:3.19]{\pageref*{exercise:3.19}}\par
\par\noindent\hangindent=2.8em\hspace*{1.8em}\phantomsection\label{idx:316}和（sum），\hyperref[ex:4.6.8]{\pageref*{ex:4.6.8}}\par
\par\noindent\hangindent=2em\hspace*{1em}\phantomsection\label{idx:317}高斯（Gaussian），\hyperref[idxocc:317:114]{\pageref*{idxocc:317:114}}\par
\par\noindent\hangindent=2em\hspace*{1em}\phantomsection\label{idx:318}几何（geometric），\hyperref[idxocc:318:111]{\pageref*{idxocc:318:111}}\par
\par\noindent\hangindent=2.8em\hspace*{1.8em}\phantomsection\label{idx:319}累积分布函数（cdf），\hyperref[ex:1.5.4]{\pageref*{ex:1.5.4}}\par
\par\noindent\hangindent=2.8em\hspace*{1.8em}\phantomsection\label{idx:320}无记忆性（memoryless），\hyperref[eq:3.2.11]{\pageref*{eq:3.2.11}}\par
\par\noindent\hangindent=2.8em\hspace*{1.8em}\phantomsection\label{idx:321}概率质量函数（pmf），\hyperref[ex:1.6.2]{\pageref*{ex:1.6.2}}\par
\par\noindent\hangindent=2em\hspace*{1em}\phantomsection\label{idx:322}耿贝尔（Gumbel），\hyperref[idxocc:322:139]{\pageref*{idxocc:322:139}}\par
\par\noindent\hangindent=2em\hspace*{1em}\phantomsection\label{idx:323}超几何（hypergeometric），\hyperref[eq:3.2.2]{\pageref*{eq:3.2.2}}\par
\par\noindent\hangindent=2em\hspace*{1em}\phantomsection\label{idx:324}逆高斯（inverse Gaussian），\hyperref[exercise:6.27]{\pageref*{exercise:6.27}}, \hyperref[exercise:7.15]{\pageref*{exercise:7.15}}\par
\needspace{3\baselineskip}
\par\noindent\hangindent=2em\hspace*{1em}\phantomsection\label{idx:325}逆伽马（inverted gamma）\par
\par\noindent\hangindent=2.8em\hspace*{1.8em}\phantomsection\label{idx:326}与伽马的关系（relation to gamma），\hyperref[ex:2.1.6]{\pageref*{ex:2.1.6}}\par
\par\noindent\hangindent=2em\hspace*{1em}\phantomsection\label{idx:327}对数级数（logarithmic series），\hyperref[exercise:3.14]{\pageref*{exercise:3.14}}\par
\needspace{3\baselineskip}
\par\noindent\hangindent=2em\hspace*{1em}\phantomsection\label{idx:328}逻辑（logistic）\par
\par\noindent\hangindent=2.8em\hspace*{1.8em}\phantomsection\label{idx:329}累积分布函数（cdf），\hyperref[ex:1.5.5]{\pageref*{ex:1.5.5}}\par
\par\noindent\hangindent=2.8em\hspace*{1.8em}\phantomsection\label{idx:330}风险率函数（hazard function），\hyperref[idxocc:330:139]{\pageref*{idxocc:330:139}}\par
\par\noindent\hangindent=2.8em\hspace*{1.8em}\phantomsection\label{idx:331}概率密度函数（pdf），\hyperref[ex:1.6.4]{\pageref*{ex:1.6.4}}\par
\par\noindent\hangindent=2em\hspace*{1em}\phantomsection\label{idx:332}对数正态（lognormal），\hyperref[eq:3.3.21]{\pageref*{eq:3.3.21}}\par
\par\noindent\hangindent=2.8em\hspace*{1.8em}\phantomsection\label{idx:333}矩母函数（mgf），\hyperref[exercise:2.36]{\pageref*{exercise:2.36}}\par
\par\noindent\hangindent=2.8em\hspace*{1.8em}\phantomsection\label{idx:334}矩（moments），\hyperref[ex:2.3.10]{\pageref*{ex:2.3.10}}\par
\par\noindent\hangindent=2em\hspace*{1em}\phantomsection\label{idx:335}麦克斯韦（Maxwell），\hyperref[idxocc:335:139]{\pageref*{idxocc:335:139}}\par
\par\noindent\hangindent=2em\hspace*{1em}\phantomsection\label{idx:336}混合（mixture），\hyperref[def:4.4.4]{\pageref*{def:4.4.4}}\par
\par\noindent\hangindent=2.8em\hspace*{1.8em}\phantomsection\label{idx:337}伯努利–贝塔（Bernoulli–beta），\hyperref[ex:4.4.6]{\pageref*{ex:4.4.6}}, \hyperref[ex:4.4.8]{\pageref*{ex:4.4.8}}\par
\par\noindent\hangindent=2.8em\hspace*{1.8em}\phantomsection\label{idx:338}二项–贝塔（binomial–beta），\hyperref[exercise:4.34]{\pageref*{exercise:4.34}}\par
\par\noindent\hangindent=2.8em\hspace*{1.8em}\phantomsection\label{idx:339}二项–负二项（binomial–negative binomial），\hyperref[idxocc:339:167]{\pageref*{idxocc:339:167}}\par
\par\noindent\hangindent=2.8em\hspace*{1.8em}\phantomsection\label{idx:340}二项–泊松（binomial–Poisson），\hyperref[ex:4.4.1]{\pageref*{ex:4.4.1}}\par
\par\noindent\hangindent=2.8em\hspace*{1.8em}\phantomsection\label{idx:341}二项–泊松–指数（binomial–Poisson–exponential），\hyperref[ex:4.4.5]{\pageref*{ex:4.4.5}}\par
\par\noindent\hangindent=2.8em\hspace*{1.8em}\phantomsection\label{idx:342}卡方–泊松（chi squared–Poisson），\hyperref[eq:4.4.3]{\pageref*{eq:4.4.3}}\par
\par\noindent\hangindent=2.8em\hspace*{1.8em}\phantomsection\label{idx:343}指数–伽马（exponential-gamma），\hyperref[exercise:4.38]{\pageref*{exercise:4.38}}\par
\par\noindent\hangindent=2.8em\hspace*{1.8em}\phantomsection\label{idx:344}对数级数–泊松（logarithmic series–Poisson），\hyperref[exercise:4.33]{\pageref*{exercise:4.33}}\par
\par\noindent\hangindent=2.8em\hspace*{1.8em}\phantomsection\label{idx:345}正态–伽马（normal–gamma），\hyperref[exercise:5.20]{\pageref*{exercise:5.20}}\par
\par\noindent\hangindent=2.8em\hspace*{1.8em}\phantomsection\label{idx:346}泊松–伽马（Poisson–gamma），\hyperref[idxocc:346:167]{\pageref*{idxocc:346:167}}, \hyperref[exercise:4.32]{\pageref*{exercise:4.32}}\par
\par\noindent\hangindent=2.8em\hspace*{1.8em}\phantomsection\label{idx:347}t–F，\hyperref[exercise:5.20]{\pageref*{exercise:5.20}}\par
\par\noindent\hangindent=2em\hspace*{1em}\phantomsection\label{idx:348}多项（multinomial），\hyperref[def:4.6.2]{\pageref*{def:4.6.2}}\par
\par\noindent\hangindent=2.8em\hspace*{1.8em}\phantomsection\label{idx:349}协方差（covariance），\hyperref[idxocc:349:180]{\pageref*{idxocc:349:180}}, \hyperref[exercise:4.39]{\pageref*{exercise:4.39}}\par
\par\noindent\hangindent=2.8em\hspace*{1.8em}\phantomsection\label{idx:350}与二项的关系（relation to binomial），\hyperref[idxocc:350:179]{\pageref*{idxocc:350:179}}\par
\par\noindent\hangindent=2em\hspace*{1em}\phantomsection\label{idx:351}负二项（negative binomial），\hyperref[eq:3.2.10]{\pageref*{eq:3.2.10}}, \hyperref[exercise:4.32]{\pageref*{exercise:4.32}}\par
\par\noindent\hangindent=2.8em\hspace*{1.8em}\phantomsection\label{idx:352}二次方差（quadratic variance），\hyperref[idxocc:352:110]{\pageref*{idxocc:352:110}}\par
\par\noindent\hangindent=2.8em\hspace*{1.8em}\phantomsection\label{idx:353}与二项的关系（relation to binomial），\hyperref[exercise:3.12]{\pageref*{exercise:3.12}}\par
\par\noindent\hangindent=2.8em\hspace*{1.8em}\phantomsection\label{idx:354}与卡方的关系（relation to chi squared），\hyperref[exercise:2.38]{\pageref*{exercise:2.38}}\par
\par\noindent\hangindent=2.8em\hspace*{1.8em}\phantomsection\label{idx:355}与伽马的关系（relation to gamma），\hyperref[exercise:3.18]{\pageref*{exercise:3.18}}\par
\par\noindent\hangindent=2.8em\hspace*{1.8em}\phantomsection\label{idx:356}与泊松的关系（relation to Poisson），\hyperref[idxocc:356:110]{\pageref*{idxocc:356:110}}\par
\par\noindent\hangindent=2em\hspace*{1em}\phantomsection\label{idx:357}非中心卡方（noncentral chi squared），\hyperref[eq:4.4.3]{\pageref*{eq:4.4.3}}\par
\par\noindent\hangindent=2em\hspace*{1em}\phantomsection\label{idx:358}正态（normal），\hyperref[idxocc:358:114]{\pageref*{idxocc:358:114}}\par
\par\noindent\hangindent=2.8em\hspace*{1.8em}\phantomsection\label{idx:359}二元（bivariate），\hyperref[exercise:4.47]{\pageref*{exercise:4.47}}\par
\par\noindent\hangindent=2.8em\hspace*{1.8em}\phantomsection\label{idx:360}指数族（exponential family），\hyperref[ex:3.4.4]{\pageref*{ex:3.4.4}}\par
\par\noindent\hangindent=2.8em\hspace*{1.8em}\phantomsection\label{idx:361}生成（generating），\hyperref[eq:5.6.5]{\pageref*{eq:5.6.5}}\par
\par\noindent\hangindent=2.8em\hspace*{1.8em}\phantomsection\label{idx:362}线性组合（linear combination），\hyperref[cor:4.6.10]{\pageref*{cor:4.6.10}}, \hyperref[lem:5.3.3]{\pageref*{lem:5.3.3}}\par
\par\noindent\hangindent=2.8em\hspace*{1.8em}\phantomsection\label{idx:363}比值（ratio），\hyperref[ex:4.3.6]{\pageref*{ex:4.3.6}}\par
\par\noindent\hangindent=2.8em\hspace*{1.8em}\phantomsection\label{idx:364}递推关系（recursion relation），\hyperref[idxocc:364:90]{\pageref*{idxocc:364:90}}\par
\par\noindent\hangindent=2.8em\hspace*{1.8em}\phantomsection\label{idx:365}与柯西的关系（relation to Cauchy），\hyperref[ex:4.3.6]{\pageref*{ex:4.3.6}}, \hyperref[exercise:4.28]{\pageref*{exercise:4.28}}\par
\par\noindent\hangindent=2.8em\hspace*{1.8em}\phantomsection\label{idx:366}与卡方的关系（relation to chi squared），\hyperref[ex:2.1.9]{\pageref*{ex:2.1.9}}\par
\par\noindent\hangindent=2.8em\hspace*{1.8em}\phantomsection\label{idx:367}与t的关系（relation to t），\hyperref[exercise:5.18]{\pageref*{exercise:5.18}}\par
\par\noindent\hangindent=2.8em\hspace*{1.8em}\phantomsection\label{idx:368}标准（standard），\hyperref[idxocc:368:115]{\pageref*{idxocc:368:115}}\par
\par\noindent\hangindent=2.8em\hspace*{1.8em}\phantomsection\label{idx:369}和（sum），\hyperref[ex:4.2.13]{\pageref*{ex:4.2.13}}, \hyperref[ex:4.3.4]{\pageref*{ex:4.3.4}}\par
\par\noindent\hangindent=2em\hspace*{1em}\phantomsection\label{idx:370}帕累托（Pareto），\hyperref[exercise:3.23]{\pageref*{exercise:3.23}}\par
\par\noindent\hangindent=2em\hspace*{1em}\phantomsection\label{idx:371}泊松（Poisson），\hyperref[eq:3.2.5]{\pageref*{eq:3.2.5}}\par
\par\noindent\hangindent=2.8em\hspace*{1.8em}\phantomsection\label{idx:372}矩母函数（mgf），\hyperref[lem:2.3.14]{\pageref*{lem:2.3.14}}\par
\par\noindent\hangindent=2.8em\hspace*{1.8em}\phantomsection\label{idx:373}正态近似（normal approximation），\hyperref[ex:3.4.9]{\pageref*{ex:3.4.9}}\par
\par\noindent\hangindent=2.8em\hspace*{1.8em}\phantomsection\label{idx:374}基本假设（postulates），\hyperref[thm:3.8.1]{\pageref*{thm:3.8.1}}\par
\par\noindent\hangindent=2.8em\hspace*{1.8em}\phantomsection\label{idx:375}递推关系（recursion relation），\hyperref[eq:3.2.6]{\pageref*{eq:3.2.6}}, \hyperref[eq:3.6.2]{\pageref*{eq:3.6.2}}\par
\par\noindent\hangindent=2.8em\hspace*{1.8em}\phantomsection\label{idx:376}与二项的关系（relation to binomial），\hyperref[exercise:4.15]{\pageref*{exercise:4.15}}\par
\par\noindent\hangindent=2.8em\hspace*{1.8em}\phantomsection\label{idx:377}与伽马的关系（relation to gamma），\hyperref[ex:3.3.1]{\pageref*{ex:3.3.1}}, \hyperref[exercise:3.19]{\pageref*{exercise:3.19}}\par
\par\noindent\hangindent=2.8em\hspace*{1.8em}\phantomsection\label{idx:378}与负二项的关系（relation to negative binomial），\hyperref[idxocc:378:110]{\pageref*{idxocc:378:110}}, \hyperref[exercise:3.15]{\pageref*{exercise:3.15}}\par
\par\noindent\hangindent=2.8em\hspace*{1.8em}\phantomsection\label{idx:379}和（sum），\hyperref[thm:4.3.2]{\pageref*{thm:4.3.2}}\par
\par\noindent\hangindent=2em\hspace*{1em}\phantomsection\label{idx:380}瑞利（Rayleigh），\hyperref[idxocc:380:139]{\pageref*{idxocc:380:139}}\par
\par\noindent\hangindent=2em\hspace*{1em}\phantomsection\label{idx:381}t，\hyperref[def:5.3.4]{\pageref*{def:5.3.4}}\par
\par\noindent\hangindent=2.8em\hspace*{1.8em}\phantomsection\label{idx:382}非中心（noncentral），\hyperref[exercise:8.35]{\pageref*{exercise:8.35}}\par
\par\noindent\hangindent=2.8em\hspace*{1.8em}\phantomsection\label{idx:383}与F的关系（relation to F），\hyperref[thm:5.3.8]{\pageref*{thm:5.3.8}}\par
\par\noindent\hangindent=2.8em\hspace*{1.8em}\phantomsection\label{idx:384}与正态的关系（relation to normal），\hyperref[exercise:5.18]{\pageref*{exercise:5.18}}\par
\par\noindent\hangindent=2em\hspace*{1em}\phantomsection\label{idx:385}截断（truncated），\hyperref[exercise:3.13]{\pageref*{exercise:3.13}}\par
\par\noindent\hangindent=2em\hspace*{1em}\phantomsection\label{idx:386}Tukey，\hyperref[exercise:10.23]{\pageref*{exercise:10.23}}\par
\par\noindent\hangindent=2em\hspace*{1em}\phantomsection\label{idx:387}均匀（uniform），\hyperref[eq:3.3.1]{\pageref*{eq:3.3.1}}\par
\par\noindent\hangindent=2.8em\hspace*{1.8em}\phantomsection\label{idx:388}最大值（maximum），\hyperref[ex:5.5.11]{\pageref*{ex:5.5.11}}\par
\sourcepageanchor{560}
\par\noindent\hangindent=2.8em\hspace*{1.8em}\phantomsection\label{idx:389}次序统计量（order statistic），\hyperref[ex:5.4.5]{\pageref*{ex:5.4.5}}\par
\par\noindent\hangindent=2.8em\hspace*{1.8em}\phantomsection\label{idx:390}与指数的关系（relation to exponential），\hyperref[ex:2.1.4]{\pageref*{ex:2.1.4}}\par
\par\noindent\hangindent=2em\hspace*{1em}\phantomsection\label{idx:391}威布尔（Weibull），\hyperref[eq:3.3.12]{\pageref*{eq:3.3.12}}, \hyperref[idxocc:391:139]{\pageref*{idxocc:391:139}}\par
\par\noindent\hangindent=2.8em\hspace*{1.8em}\phantomsection\label{idx:392}风险率函数（hazard function），\hyperref[idxocc:392:139]{\pageref*{idxocc:392:139}}\par
\par\addvspace{2pt}
\par\noindent\hangindent=1em\hspace*{0em}\phantomsection\label{idx:393}点记号（Dot notation），\hyperref[idxocc:393:457]{\pageref*{idxocc:393:457}}\par
\par\addvspace{2pt}
\par\noindent\hangindent=1em\hspace*{0em}\phantomsection\label{idx:394}效率（另见\hyperref[idx:763]{点估计}）（Efficiency (see also Point estimation)），\hyperref[sec:10.1.2]{\pageref*{sec:10.1.2}}\par
\par\addvspace{2pt}
\par\noindent\hangindent=1em\hspace*{0em}\phantomsection\label{idx:395}EM算法（见\hyperref[idx:7]{算法：EM}）（EM Algorithm (see Algorithm, EM)），\hyperref[sec:7.2.4]{\pageref*{sec:7.2.4}}, \hyperref[sec:7.5.4]{\pageref*{sec:7.5.4}}\par
\par\addvspace{2pt}
\par\noindent\hangindent=1em\hspace*{0em}\phantomsection\label{idx:396}经验贝叶斯估计量（Empirical Bayes estimator），\hyperref[idxocc:396:331]{\pageref*{idxocc:396:331}}\par
\par\addvspace{2pt}
\par\noindent\hangindent=1em\hspace*{0em}\phantomsection\label{idx:397}空集（Empty set），\hyperref[idxocc:397:33]{\pageref*{idxocc:397:33}}\par
\par\addvspace{2pt}
\par\noindent\hangindent=1em\hspace*{0em}\phantomsection\label{idx:398}熵（Entropy），\hyperref[idxocc:398:196]{\pageref*{idxocc:398:196}}\par
\par\addvspace{2pt}
\needspace{3\baselineskip}
\par\noindent\hangindent=1em\hspace*{0em}\phantomsection\label{idx:399}等变性（Equivariance）\par
\par\noindent\hangindent=2em\hspace*{1em}\phantomsection\label{idx:400}二项（binomial），\hyperref[ex:6.4.1]{\pageref*{ex:6.4.1}}\par
\par\noindent\hangindent=2em\hspace*{1em}\phantomsection\label{idx:401}测量（measurement），\hyperref[idxocc:401:272]{\pageref*{idxocc:401:272}}, \hyperref[exercise:6.39]{\pageref*{exercise:6.39}}\par
\par\noindent\hangindent=2.8em\hspace*{1.8em}\phantomsection\label{idx:402}位置（location），\hyperref[exercise:6.41]{\pageref*{exercise:6.41}}\par
\par\addvspace{2pt}
\par\noindent\hangindent=1em\hspace*{0em}\phantomsection\label{idx:403}等变性（另见\hyperref[idx:587]{不变性}）（Equivariance (see also Invariance)），\hyperref[sec:6.4]{\pageref*{sec:6.4}}\par
\par\addvspace{2pt}
\par\noindent\hangindent=1em\hspace*{0em}\phantomsection\label{idx:404}等变性原则（Equivariance Principle），\hyperref[idxocc:404:272]{\pageref*{idxocc:404:272}}\par
\par\noindent\hangindent=2em\hspace*{1em}\phantomsection\label{idx:405}解释（interpretations），\hyperref[ex:6.4.6]{\pageref*{ex:6.4.6}}\par
\par\addvspace{2pt}
\par\noindent\hangindent=1em\hspace*{0em}\phantomsection\label{idx:406}遍历定理（Ergodic theorem），\hyperref[idxocc:406:250]{\pageref*{idxocc:406:250}}\par
\par\addvspace{2pt}
\par\noindent\hangindent=1em\hspace*{0em}\phantomsection\label{idx:407}变量含误差的回归（Errors in variables regression (EIV)），\hyperref[sec:12.2]{\pageref*{sec:12.2}}\par
\par\noindent\hangindent=2em\hspace*{1em}\phantomsection\label{idx:408}置信集（confidence set），\hyperref[sec:12.2.4]{\pageref*{sec:12.2.4}}\par
\par\noindent\hangindent=2em\hspace*{1em}\phantomsection\label{idx:409}相合性（consistency），\hyperref[idxocc:409:505]{\pageref*{idxocc:409:505}}, \hyperref[idxocc:409:507]{\pageref*{idxocc:409:507}}, \hyperref[sec:12.6.2]{\pageref*{sec:12.6.2}}\par
\par\noindent\hangindent=2em\hspace*{1em}\phantomsection\label{idx:410}函数关系（functional relationship），\hyperref[eq:12.2.5]{\pageref*{eq:12.2.5}}, \hyperref[sec:12.6.1]{\pageref*{sec:12.6.1}}\par
\par\noindent\hangindent=2em\hspace*{1em}\phantomsection\label{idx:411}可识别性（identifiability），\hyperref[idxocc:411:501]{\pageref*{idxocc:411:501}}, \hyperref[idxocc:411:506]{\pageref*{idxocc:411:506}}\par
\needspace{3\baselineskip}
\par\noindent\hangindent=2em\hspace*{1em}\phantomsection\label{idx:412}最大似然估计量（MLE）\par
\par\noindent\hangindent=2.8em\hspace*{1.8em}\phantomsection\label{idx:413}函数型（functional），\hyperref[eq:12.2.16]{\pageref*{eq:12.2.16}}\par
\par\noindent\hangindent=2.8em\hspace*{1.8em}\phantomsection\label{idx:414}与最小二乘的关系（relation to least squares），\hyperref[exercise:12.4]{\pageref*{exercise:12.4}}\par
\par\noindent\hangindent=2.8em\hspace*{1.8em}\phantomsection\label{idx:415}结构型（structural），\hyperref[idxocc:415:505]{\pageref*{idxocc:415:505}}\par
\par\noindent\hangindent=2em\hspace*{1em}\phantomsection\label{idx:416}与随机完全区组的关系（relation to RCB），\hyperref[idxocc:416:506]{\pageref*{idxocc:416:506}}\par
\par\noindent\hangindent=2em\hspace*{1em}\phantomsection\label{idx:417}结构关系（structural relationship），\hyperref[eq:12.2.6]{\pageref*{eq:12.2.6}}, \hyperref[sec:12.6.1]{\pageref*{sec:12.6.1}}\par
\par\addvspace{2pt}
\par\noindent\hangindent=1em\hspace*{0em}\phantomsection\label{idx:418}事件（Event），\hyperref[def:1.1.2]{\pageref*{def:1.1.2}}\par
\par\addvspace{2pt}
\par\noindent\hangindent=1em\hspace*{0em}\phantomsection\label{idx:419}证据（Evidence），\hyperref[idxocc:419:269]{\pageref*{idxocc:419:269}}\par
\par\addvspace{2pt}
\par\noindent\hangindent=1em\hspace*{0em}\phantomsection\label{idx:420}交换悖论（Exchange paradox），\hyperref[sec:4.9.1]{\pageref*{sec:4.9.1}}\par
\par\addvspace{2pt}
\par\noindent\hangindent=1em\hspace*{0em}\phantomsection\label{idx:421}可交换（Exchangeable），\hyperref[exercise:5.4]{\pageref*{exercise:5.4}}\par
\par\addvspace{2pt}
\par\noindent\hangindent=1em\hspace*{0em}\phantomsection\label{idx:422}期望值（Expected value），\hyperref[def:2.2.1]{\pageref*{def:2.2.1}}\par
\par\noindent\hangindent=2em\hspace*{1em}\phantomsection\label{idx:423}二元（bivariate），\hyperref[idxocc:423:147]{\pageref*{idxocc:423:147}}, \hyperref[idxocc:423:149]{\pageref*{idxocc:423:149}}\par
\par\noindent\hangindent=2.8em\hspace*{1.8em}\phantomsection\label{idx:424}性质（properties），\hyperref[exercise:4.2]{\pageref*{exercise:4.2}}\par
\par\noindent\hangindent=2em\hspace*{1em}\phantomsection\label{idx:425}条件（conditional），\hyperref[idxocc:425:154]{\pageref*{idxocc:425:154}}\par
\par\noindent\hangindent=2.8em\hspace*{1.8em}\phantomsection\label{idx:426}最小化（minimizing），\hyperref[exercise:4.13]{\pageref*{exercise:4.13}}\par
\par\noindent\hangindent=2em\hspace*{1em}\phantomsection\label{idx:427}独立性（independence），\hyperref[thm:4.2.10]{\pageref*{thm:4.2.10}}, \hyperref[thm:4.6.6]{\pageref*{thm:4.6.6}}\par
\par\noindent\hangindent=2em\hspace*{1em}\phantomsection\label{idx:428}迭代（iterated），\hyperref[thm:4.4.3]{\pageref*{thm:4.4.3}}\par
\par\noindent\hangindent=2em\hspace*{1em}\phantomsection\label{idx:429}联合（joint），\hyperref[eq:4.6.3]{\pageref*{eq:4.6.3}}\par
\par\noindent\hangindent=2em\hspace*{1em}\phantomsection\label{idx:430}最小化性质（minimizing property），\hyperref[ex:2.2.6]{\pageref*{ex:2.2.6}}\par
\par\noindent\hangindent=2em\hspace*{1em}\phantomsection\label{idx:431}极大值与极小值的（of maxima and minima），\hyperref[exercise:2.15]{\pageref*{exercise:2.15}}\par
\par\noindent\hangindent=2em\hspace*{1em}\phantomsection\label{idx:432}性质（properties），\hyperref[thm:2.2.5]{\pageref*{thm:2.2.5}}\par
\par\noindent\hangindent=2em\hspace*{1em}\phantomsection\label{idx:433}和（sum），\hyperref[lem:5.2.5]{\pageref*{lem:5.2.5}}\par
\par\noindent\hangindent=2em\hspace*{1em}\phantomsection\label{idx:434}两种计算方法（two calculation methods），\hyperref[eq:2.2.5]{\pageref*{eq:2.2.5}}\par
\par\addvspace{2pt}
\par\noindent\hangindent=1em\hspace*{0em}\phantomsection\label{idx:435}指数族（Exponential family），\hyperref[eq:3.4.1]{\pageref*{eq:3.4.1}}, \hyperref[exercise:3.28]{\pageref*{exercise:3.28}}\par
\par\noindent\hangindent=2em\hspace*{1em}\phantomsection\label{idx:436}曲线（curved），\hyperref[def:3.4.7]{\pageref*{def:3.4.7}}, \hyperref[idxocc:436:265]{\pageref*{idxocc:436:265}}\par
\par\noindent\hangindent=2.8em\hspace*{1.8em}\phantomsection\label{idx:437}正态近似（normal approximations），\hyperref[exercise:3.34]{\pageref*{exercise:3.34}}\par
\par\noindent\hangindent=2em\hspace*{1em}\phantomsection\label{idx:438}对数正态（lognormal），\hyperref[sec:3.8.3]{\pageref*{sec:3.8.3}}\par
\par\noindent\hangindent=2em\hspace*{1em}\phantomsection\label{idx:439}矩（moments），\hyperref[thm:3.4.2]{\pageref*{thm:3.4.2}}\par
\par\noindent\hangindent=2em\hspace*{1em}\phantomsection\label{idx:440}自然参数（natural parameter），\hyperref[idxocc:440:125]{\pageref*{idxocc:440:125}}\par
\par\noindent\hangindent=2em\hspace*{1em}\phantomsection\label{idx:441}和（sums），\hyperref[thm:5.2.11]{\pageref*{thm:5.2.11}}\par
\par\addvspace{2pt}
\par\noindent\hangindent=1em\hspace*{0em}\phantomsection\label{idx:442}阶乘（Factorial），\hyperref[def:1.2.16]{\pageref*{def:1.2.16}}\par
\par\addvspace{2pt}
\par\noindent\hangindent=1em\hspace*{0em}\phantomsection\label{idx:443}因子分解定理（Factorization Theorem），\hyperref[thm:6.2.6]{\pageref*{thm:6.2.6}}\par
\par\addvspace{2pt}
\par\noindent\hangindent=1em\hspace*{0em}\phantomsection\label{idx:444}信仰推断（Fiducial inference），\hyperref[idxocc:444:268]{\pageref*{idxocc:444:268}}\par
\par\addvspace{2pt}
\par\noindent\hangindent=1em\hspace*{0em}\phantomsection\label{idx:445}Fieller定理（Fieller’s Theorem），\hyperref[sec:9.5.3]{\pageref*{sec:9.5.3}}\par
\par\addvspace{2pt}
\par\noindent\hangindent=1em\hspace*{0em}\phantomsection\label{idx:446}强制二元选择（Forced binary choice），\hyperref[exercise:4.37]{\pageref*{exercise:4.37}}\par
\par\addvspace{2pt}
\par\noindent\hangindent=1em\hspace*{0em}\phantomsection\label{idx:447}泛函（Functional），\hyperref[idxocc:447:449]{\pageref*{idxocc:447:449}}\par
\par\addvspace{2pt}
\par\noindent\hangindent=1em\hspace*{0em}\phantomsection\label{idx:448}微积分基本定理（Fundamental Theorem of Calculus），\hyperref[eq:1.6.2]{\pageref*{eq:1.6.2}}, \hyperref[eq:4.1.3]{\pageref*{eq:4.1.3}}\par
\par\addvspace{2pt}
\par\noindent\hangindent=1em\hspace*{0em}\phantomsection\label{idx:449}计数基本定理（Fundamental Theorem of Counting），\hyperref[thm:1.2.14]{\pageref*{thm:1.2.14}}\par
\par\addvspace{2pt}
\par\noindent\hangindent=1em\hspace*{0em}\phantomsection\label{idx:450}伽马函数（另见\hyperref[idx:306]{分布：伽马}）（Gamma function (See also Distribution gamma)），\hyperref[eq:3.3.2]{\pageref*{eq:3.3.2}}\par
\par\addvspace{2pt}
\par\noindent\hangindent=1em\hspace*{0em}\phantomsection\label{idx:451}Gauss–Markov定理（Gauss–Markov Theorem），\hyperref[idxocc:451:474]{\pageref*{idxocc:451:474}}\par
\par\addvspace{2pt}
\par\noindent\hangindent=1em\hspace*{0em}\phantomsection\label{idx:452}广义均值（generalized means），\hyperref[exercise:7.16]{\pageref*{exercise:7.16}}\par
\par\addvspace{2pt}
\needspace{3\baselineskip}
\par\noindent\hangindent=1em\hspace*{0em}\phantomsection\label{idx:453}母函数（Generating function）\par
\par\noindent\hangindent=2em\hspace*{1em}\phantomsection\label{idx:454}累积量（cumulant），\hyperref[idxocc:454:99]{\pageref*{idxocc:454:99}}\par
\par\noindent\hangindent=2em\hspace*{1em}\phantomsection\label{idx:455}阶乘矩（factorial moment），\hyperref[idxocc:455:99]{\pageref*{idxocc:455:99}}\par
\par\noindent\hangindent=2em\hspace*{1em}\phantomsection\label{idx:456}概率（probability），\hyperref[idxocc:456:99]{\pageref*{idxocc:456:99}}\par
\par\addvspace{2pt}
\par\noindent\hangindent=1em\hspace*{0em}\phantomsection\label{idx:457}几何平均数（Geometric mean），\hyperref[ex:4.7.8]{\pageref*{ex:4.7.8}}\par
\par\addvspace{2pt}
\needspace{3\baselineskip}
\par\noindent\hangindent=1em\hspace*{0em}\phantomsection\label{idx:458}等比级数（Geometric series）\par
\par\noindent\hangindent=2em\hspace*{1em}\phantomsection\label{idx:459}求导（differentiation），\hyperref[ex:2.4.7]{\pageref*{ex:2.4.7}}\par
\par\noindent\hangindent=2em\hspace*{1em}\phantomsection\label{idx:460}部分和（partial sum），\hyperref[eq:1.5.4]{\pageref*{eq:1.5.4}}\par
\par\addvspace{2pt}
\par\noindent\hangindent=1em\hspace*{0em}\phantomsection\label{idx:461}Gosset，W. S.（Gosset, W. S.），\hyperref[idxocc:461:212]{\pageref*{idxocc:461:212}}\par
\par\addvspace{2pt}
\par\noindent\hangindent=1em\hspace*{0em}\phantomsection\label{idx:462}总均值（Grand mean），\hyperref[idxocc:462:457]{\pageref*{idxocc:462:457}}\par
\par\addvspace{2pt}
\needspace{3\baselineskip}
\par\noindent\hangindent=1em\hspace*{0em}\phantomsection\label{idx:463}群（Group）\par
\par\noindent\hangindent=2em\hspace*{1em}\phantomsection\label{idx:464}位置（location），\hyperref[ex:6.4.6]{\pageref*{ex:6.4.6}}, \hyperref[ex:7.3.6]{\pageref*{ex:7.3.6}}\par
\par\noindent\hangindent=2em\hspace*{1em}\phantomsection\label{idx:465}位置尺度（location–scale），\hyperref[exercise:6.42]{\pageref*{exercise:6.42}}\par
\par\noindent\hangindent=2em\hspace*{1em}\phantomsection\label{idx:466}变换的（of transformations），\hyperref[def:6.4.2]{\pageref*{def:6.4.2}}\par
\par\noindent\hangindent=2em\hspace*{1em}\phantomsection\label{idx:467}尺度（scale），\hyperref[exercise:6.43]{\pageref*{exercise:6.43}}\par
\par\addvspace{2pt}
\par\noindent\hangindent=1em\hspace*{0em}\phantomsection\label{idx:468}调和平均数（Harmonic mean），\hyperref[ex:4.7.8]{\pageref*{ex:4.7.8}}\par
\par\addvspace{2pt}
\par\noindent\hangindent=1em\hspace*{0em}\phantomsection\label{idx:469}风险率函数（Hazard function），\hyperref[idxocc:469:114]{\pageref*{idxocc:469:114}}, \hyperref[exercise:3.25]{\pageref*{exercise:3.25}}\par
\par\addvspace{2pt}
\par\noindent\hangindent=1em\hspace*{0em}\phantomsection\label{idx:470}分层贝叶斯估计量（Hierarchical Bayes estimator），\hyperref[idxocc:470:332]{\pageref*{idxocc:470:332}}\par
\par\addvspace{2pt}
\par\noindent\hangindent=1em\hspace*{0em}\phantomsection\label{idx:471}最高后验密度区域（Highest posterior density (HPD) region），\hyperref[cor:9.3.10]{\pageref*{cor:9.3.10}}\par
\par\addvspace{2pt}
\par\noindent\hangindent=1em\hspace*{0em}\phantomsection\label{idx:472}同方差性（Homoscedasticity），\hyperref[item:11.2.assumption.iii]{\pageref*{item:11.2.assumption.iii}}\par
\par\addvspace{2pt}
\par\noindent\hangindent=1em\hspace*{0em}\phantomsection\label{idx:473}Hwang引理（Hwang’s Lemma），\hyperref[thm:3.6.8]{\pageref*{thm:3.6.8}}\par
\par\addvspace{2pt}
\par\noindent\hangindent=1em\hspace*{0em}\phantomsection\label{idx:474}假设（Hypothesis），\hyperref[def:8.1.1]{\pageref*{def:8.1.1}}\par
\par\noindent\hangindent=2em\hspace*{1em}\phantomsection\label{idx:475}备择（alternative），\hyperref[def:8.1.2]{\pageref*{def:8.1.2}}\par
\par\noindent\hangindent=2em\hspace*{1em}\phantomsection\label{idx:476}复合（composite），\hyperref[idxocc:476:348]{\pageref*{idxocc:476:348}}\par
\par\noindent\hangindent=2em\hspace*{1em}\phantomsection\label{idx:477}零（null），\hyperref[def:8.1.2]{\pageref*{def:8.1.2}}\par
\par\noindent\hangindent=2em\hspace*{1em}\phantomsection\label{idx:478}单侧（one-sided），\hyperref[idxocc:478:348]{\pageref*{idxocc:478:348}}\par
\par\noindent\hangindent=2em\hspace*{1em}\phantomsection\label{idx:479}研究（research），\hyperref[idxocc:479:344]{\pageref*{idxocc:479:344}}\par
\par\noindent\hangindent=2em\hspace*{1em}\phantomsection\label{idx:480}简单（simple），\hyperref[idxocc:480:345]{\pageref*{idxocc:480:345}}, \hyperref[idxocc:480:348]{\pageref*{idxocc:480:348}}\par
\par\noindent\hangindent=2em\hspace*{1em}\phantomsection\label{idx:481}双侧（two-sided），\hyperref[idxocc:481:348]{\pageref*{idxocc:481:348}}\par
\par\addvspace{2pt}
\par\noindent\hangindent=1em\hspace*{0em}\phantomsection\label{idx:482}假设检验（Hypothesis test），\hyperref[def:8.1.3]{\pageref*{def:8.1.3}}\par
\par\noindent\hangindent=2em\hspace*{1em}\phantomsection\label{idx:483}滥用（abuse），\hyperref[exercise:8.16]{\pageref*{exercise:8.16}}\par
\par\noindent\hangindent=2em\hspace*{1em}\phantomsection\label{idx:484}接受域（acceptance region），\hyperref[idxocc:484:333]{\pageref*{idxocc:484:333}}\par
\par\noindent\hangindent=2em\hspace*{1em}\phantomsection\label{idx:485}渐近（asymptotic），\hyperref[exercise:10.32]{\pageref*{exercise:10.32}}\par
\par\noindent\hangindent=2.8em\hspace*{1.8em}\phantomsection\label{idx:486}伯努利（Bernoulli），\hyperref[ex:10.3.5]{\pageref*{ex:10.3.5}}\par
\par\noindent\hangindent=2em\hspace*{1em}\phantomsection\label{idx:487}贝叶斯（Bayes），\hyperref[sec:8.5.2]{\pageref*{sec:8.5.2}}\par
\sourcepageanchor{561}
\par\noindent\hangindent=2.8em\hspace*{1.8em}\phantomsection\label{idx:488}正态（normal），\hyperref[ex:8.2.7]{\pageref*{ex:8.2.7}}\par
\par\noindent\hangindent=2.8em\hspace*{1.8em}\phantomsection\label{idx:489}p值（p-value），\hyperref[exercise:8.53]{\pageref*{exercise:8.53}}\par
\par\noindent\hangindent=2em\hspace*{1em}\phantomsection\label{idx:490}生物等效性（bioequivalence），\hyperref[exercise:8.47]{\pageref*{exercise:8.47}}\par
\par\noindent\hangindent=2em\hspace*{1em}\phantomsection\label{idx:491}列联表（contingency table），\hyperref[exercise:10.31]{\pageref*{exercise:10.31}}\par
\par\noindent\hangindent=2em\hspace*{1em}\phantomsection\label{idx:492}决策论（decision theoretic），\hyperref[sec:8.3.5]{\pageref*{sec:8.3.5}}\par
\par\noindent\hangindent=2.8em\hspace*{1.8em}\phantomsection\label{idx:493}均值（mean），\hyperref[ex:8.3.31]{\pageref*{ex:8.3.31}}\par
\par\noindent\hangindent=2em\hspace*{1em}\phantomsection\label{idx:494}交并（intersection–union），\hyperref[eq:8.2.6]{\pageref*{eq:8.2.6}}\par
\par\noindent\hangindent=2.8em\hspace*{1.8em}\phantomsection\label{idx:495}伯努利（Bernoulli），\hyperref[ex:8.2.9]{\pageref*{ex:8.2.9}}, \hyperref[ex:8.3.25]{\pageref*{ex:8.3.25}}\par
\par\noindent\hangindent=2.8em\hspace*{1.8em}\phantomsection\label{idx:496}水平（level），\hyperref[thm:8.3.23]{\pageref*{thm:8.3.23}}\par
\par\noindent\hangindent=2.8em\hspace*{1.8em}\phantomsection\label{idx:497}正态（normal），\hyperref[ex:8.2.9]{\pageref*{ex:8.2.9}}, \hyperref[ex:8.3.25]{\pageref*{ex:8.3.25}}\par
\par\noindent\hangindent=2em\hspace*{1em}\phantomsection\label{idx:498}检验水平（level of test），\hyperref[def:8.3.6]{\pageref*{def:8.3.6}}\par
\par\noindent\hangindent=2em\hspace*{1em}\phantomsection\label{idx:499}McNemar检验（McNemar’s test），\hyperref[exercise:10.32]{\pageref*{exercise:10.32}}\par
\par\noindent\hangindent=2em\hspace*{1em}\phantomsection\label{idx:500}极小极大（minimax），\hyperref[exercise:8.58]{\pageref*{exercise:8.58}}\par
\par\noindent\hangindent=2em\hspace*{1em}\phantomsection\label{idx:501}功效（power），\hyperref[def:8.3.1]{\pageref*{def:8.3.1}}\par
\par\noindent\hangindent=2.8em\hspace*{1.8em}\phantomsection\label{idx:502}二项（binomial），\hyperref[ex:8.3.2]{\pageref*{ex:8.3.2}}\par
\par\noindent\hangindent=2.8em\hspace*{1.8em}\phantomsection\label{idx:503}单调（monotone），\hyperref[sec:8.5.1]{\pageref*{sec:8.5.1}}\par
\par\noindent\hangindent=2.8em\hspace*{1.8em}\phantomsection\label{idx:504}正态（normal），\hyperref[ex:8.3.3]{\pageref*{ex:8.3.3}}\par
\par\noindent\hangindent=2.8em\hspace*{1.8em}\phantomsection\label{idx:505}与风险的关系（relation to risk），\hyperref[eq:8.3.12]{\pageref*{eq:8.3.12}}\par
\par\noindent\hangindent=2em\hspace*{1em}\phantomsection\label{idx:506}拒绝域（rejection region），\hyperref[idxocc:506:333]{\pageref*{idxocc:506:333}}\par
\par\noindent\hangindent=2em\hspace*{1em}\phantomsection\label{idx:507}与置信集的关系（relation to confidence set），\hyperref[thm:9.2.2]{\pageref*{thm:9.2.2}}\par
\needspace{3\baselineskip}
\par\noindent\hangindent=2em\hspace*{1em}\phantomsection\label{idx:508}样本量（sample size）\par
\par\noindent\hangindent=2.8em\hspace*{1.8em}\phantomsection\label{idx:509}渐近（asymptotic），\hyperref[idxocc:509:427]{\pageref*{idxocc:509:427}}\par
\par\noindent\hangindent=2.8em\hspace*{1.8em}\phantomsection\label{idx:510}正态（normal），\hyperref[ex:8.3.4]{\pageref*{ex:8.3.4}}\par
\par\noindent\hangindent=2em\hspace*{1em}\phantomsection\label{idx:511}检验大小（size of test），\hyperref[def:8.3.5]{\pageref*{def:8.3.5}}\par
\par\noindent\hangindent=2em\hspace*{1em}\phantomsection\label{idx:512}充分性（sufficiency），\hyperref[cor:8.3.13]{\pageref*{cor:8.3.13}}\par
\par\noindent\hangindent=2em\hspace*{1em}\phantomsection\label{idx:513}第一类错误（type I error），\hyperref[tab:8.3.1]{\pageref*{tab:8.3.1}}, \hyperref[ex:8.3.4]{\pageref*{ex:8.3.4}}\par
\par\noindent\hangindent=2em\hspace*{1em}\phantomsection\label{idx:514}第二类错误（type II error），\hyperref[tab:8.3.1]{\pageref*{tab:8.3.1}}, \hyperref[ex:8.3.4]{\pageref*{ex:8.3.4}}\par
\par\noindent\hangindent=2em\hspace*{1em}\phantomsection\label{idx:515}一致最强（UMP），\hyperref[def:8.3.11]{\pageref*{def:8.3.11}}\par
\par\noindent\hangindent=2.8em\hspace*{1.8em}\phantomsection\label{idx:516}二项（binomial），\hyperref[ex:8.3.14]{\pageref*{ex:8.3.14}}\par
\par\noindent\hangindent=2.8em\hspace*{1.8em}\phantomsection\label{idx:517}正态（normal），\hyperref[ex:8.3.15]{\pageref*{ex:8.3.15}}, \hyperref[ex:8.3.18]{\pageref*{ex:8.3.18}}, \hyperref[exercise:8.37]{\pageref*{exercise:8.37}}\par
\par\noindent\hangindent=2.8em\hspace*{1.8em}\phantomsection\label{idx:518}与一致最准确集的关系（relation to UMA set），\hyperref[thm:9.3.5]{\pageref*{thm:9.3.5}}\par
\par\noindent\hangindent=2.8em\hspace*{1.8em}\phantomsection\label{idx:519}均匀（uniform），\hyperref[exercise:8.33]{\pageref*{exercise:8.33}}\par
\par\noindent\hangindent=2em\hspace*{1em}\phantomsection\label{idx:520}无偏（unbiased），\hyperref[def:8.3.9]{\pageref*{def:8.3.9}}\par
\par\noindent\hangindent=2.8em\hspace*{1.8em}\phantomsection\label{idx:521}与无偏集的关系（relation to unbiased set），\hyperref[idxocc:521:392]{\pageref*{idxocc:521:392}}\par
\par\noindent\hangindent=2em\hspace*{1em}\phantomsection\label{idx:522}并交（union–intersection），\hyperref[eq:8.2.4]{\pageref*{eq:8.2.4}}\par
\par\noindent\hangindent=2.8em\hspace*{1.8em}\phantomsection\label{idx:523}方差分析（ANOVA），\hyperref[thm:11.2.5]{\pageref*{thm:11.2.5}}\par
\par\noindent\hangindent=2.8em\hspace*{1.8em}\phantomsection\label{idx:524}水平（level），\hyperref[thm:8.3.21]{\pageref*{thm:8.3.21}}\par
\par\noindent\hangindent=2.8em\hspace*{1.8em}\phantomsection\label{idx:525}正态（normal），\hyperref[ex:8.2.8]{\pageref*{ex:8.2.8}}, \hyperref[ex:8.3.8]{\pageref*{ex:8.3.8}}\par
\par\addvspace{2pt}
\par\noindent\hangindent=1em\hspace*{0em}\phantomsection\label{idx:526}可识别性（Identifiability），\hyperref[def:11.2.2]{\pageref*{def:11.2.2}}\par
\par\noindent\hangindent=2em\hspace*{1em}\phantomsection\label{idx:527}方差分析（ANOVA），\hyperref[def:11.2.2]{\pageref*{def:11.2.2}}\par
\par\noindent\hangindent=2em\hspace*{1em}\phantomsection\label{idx:528}变量含误差（EIV），\hyperref[idxocc:528:501]{\pageref*{idxocc:528:501}}, \hyperref[idxocc:528:506]{\pageref*{idxocc:528:506}}\par
\par\addvspace{2pt}
\needspace{3\baselineskip}
\par\noindent\hangindent=1em\hspace*{0em}\phantomsection\label{idx:529}恒等式（Identity）\par
\par\noindent\hangindent=2em\hspace*{1em}\phantomsection\label{idx:530}贝塔（beta），\hyperref[idxocc:530:142]{\pageref*{idxocc:530:142}}\par
\par\noindent\hangindent=2em\hspace*{1em}\phantomsection\label{idx:531}二项式系数（binomial coefficients），\hyperref[exercise:1.27]{\pageref*{exercise:1.27}}, \hyperref[idxocc:531:103]{\pageref*{idxocc:531:103}}\par
\par\noindent\hangindent=2em\hspace*{1em}\phantomsection\label{idx:532}卡方（chi squared），\hyperref[thm:3.6.7]{\pageref*{thm:3.6.7}}\par
\par\noindent\hangindent=2em\hspace*{1em}\phantomsection\label{idx:533}条件协方差（conditional covariance），\hyperref[exercise:4.59]{\pageref*{exercise:4.59}}\par
\par\noindent\hangindent=2em\hspace*{1em}\phantomsection\label{idx:534}协方差（covariance），\hyperref[item:12.11.a]{\pageref*{item:12.11.a}}\par
\needspace{3\baselineskip}
\par\noindent\hangindent=2em\hspace*{1em}\phantomsection\label{idx:535}期望（expectation）\par
\par\noindent\hangindent=2.8em\hspace*{1.8em}\phantomsection\label{idx:536}连续（continuous），\hyperref[exercise:2.14]{\pageref*{exercise:2.14}}\par
\par\noindent\hangindent=2.8em\hspace*{1.8em}\phantomsection\label{idx:537}离散（discrete），\hyperref[exercise:2.14]{\pageref*{exercise:2.14}}\par
\par\noindent\hangindent=2em\hspace*{1em}\phantomsection\label{idx:538}指数族（exponential family），\hyperref[thm:3.4.2]{\pageref*{thm:3.4.2}}\par
\par\noindent\hangindent=2em\hspace*{1em}\phantomsection\label{idx:539}伽马（gamma），\hyperref[idxocc:539:142]{\pageref*{idxocc:539:142}}\par
\par\noindent\hangindent=2em\hspace*{1em}\phantomsection\label{idx:540}容斥（inclusion-exclusion），\hyperref[exercise:1.42]{\pageref*{exercise:1.42}}, \hyperref[idxocc:540:68]{\pageref*{idxocc:540:68}}\par
\par\noindent\hangindent=2em\hspace*{1em}\phantomsection\label{idx:541}分部积分（integration–by–parts），\hyperref[lem:3.6.5]{\pageref*{lem:3.6.5}}\par
\par\noindent\hangindent=2em\hspace*{1em}\phantomsection\label{idx:542}拉格朗日（Lagrange），\hyperref[exercise:5.9]{\pageref*{exercise:5.9}}\par
\par\noindent\hangindent=2em\hspace*{1em}\phantomsection\label{idx:543}负二项（negative binomial），\hyperref[eq:3.6.6]{\pageref*{eq:3.6.6}}\par
\par\noindent\hangindent=2em\hspace*{1em}\phantomsection\label{idx:544}泊松（Poisson），\hyperref[eq:3.6.5]{\pageref*{eq:3.6.5}}\par
\par\addvspace{2pt}
\par\noindent\hangindent=1em\hspace*{0em}\phantomsection\label{idx:545}独立同分布（Iid），\hyperref[def:5.1.1]{\pageref*{def:5.1.1}}\par
\par\addvspace{2pt}
\needspace{3\baselineskip}
\par\noindent\hangindent=1em\hspace*{0em}\phantomsection\label{idx:546}独立性（Independence）\par
\par\noindent\hangindent=2em\hspace*{1em}\phantomsection\label{idx:547}完备性（completeness），\hyperref[thm:6.2.24]{\pageref*{thm:6.2.24}}\par
\par\noindent\hangindent=2em\hspace*{1em}\phantomsection\label{idx:548}协方差（covariance），\hyperref[thm:4.5.5]{\pageref*{thm:4.5.5}}\par
\par\noindent\hangindent=2em\hspace*{1em}\phantomsection\label{idx:549}不交（disjoint），\hyperref[exercise:1.39]{\pageref*{exercise:1.39}}\par
\par\noindent\hangindent=2em\hspace*{1em}\phantomsection\label{idx:550}事件（events），\hyperref[def:1.3.7]{\pageref*{def:1.3.7}}\par
\par\noindent\hangindent=2.8em\hspace*{1.8em}\phantomsection\label{idx:551}相互（mutual），\hyperref[def:1.3.12]{\pageref*{def:1.3.12}}\par
\par\noindent\hangindent=2.8em\hspace*{1.8em}\phantomsection\label{idx:552}性质（properties），\hyperref[thm:1.3.9]{\pageref*{thm:1.3.9}}\par
\par\noindent\hangindent=2em\hspace*{1em}\phantomsection\label{idx:553}期望值（expected value），\hyperref[thm:4.2.10]{\pageref*{thm:4.2.10}}, \hyperref[thm:4.6.6]{\pageref*{thm:4.6.6}}\par
\par\noindent\hangindent=2em\hspace*{1em}\phantomsection\label{idx:554}随机变量（random variables），\hyperref[def:4.2.5]{\pageref*{def:4.2.5}}, \hyperref[thm:4.3.5]{\pageref*{thm:4.3.5}}, \hyperref[def:4.6.5]{\pageref*{def:4.6.5}}, \hyperref[def:5.1.1]{\pageref*{def:5.1.1}}\par
\par\noindent\hangindent=2.8em\hspace*{1.8em}\phantomsection\label{idx:555}和（sums），\hyperref[thm:4.2.12]{\pageref*{thm:4.2.12}}\par
\par\addvspace{2pt}
\par\noindent\hangindent=1em\hspace*{0em}\phantomsection\label{idx:556}示性函数（Indicator function），\hyperref[def:3.4.5]{\pageref*{def:3.4.5}}\par
\par\addvspace{2pt}
\needspace{3\baselineskip}
\par\noindent\hangindent=1em\hspace*{0em}\phantomsection\label{idx:557}不等式（Inequality）\par
\par\noindent\hangindent=2em\hspace*{1em}\phantomsection\label{idx:558}算术–几何–调和平均数（arithmetic-geometric-harmonic mean），\hyperref[sec:4.9.2]{\pageref*{sec:4.9.2}}\par
\par\noindent\hangindent=2em\hspace*{1em}\phantomsection\label{idx:559}Bonferroni与Boole的相似性（Bonferroni - Boole similarity），\hyperref[idxocc:559:41]{\pageref*{idxocc:559:41}}\par
\par\noindent\hangindent=2em\hspace*{1em}\phantomsection\label{idx:560}Bonferroni（Bonferroni’s），\hyperref[eq:1.2.9]{\pageref*{eq:1.2.9}}, \hyperref[eq:1.2.10]{\pageref*{eq:1.2.10}}, \hyperref[sec:1.8.1]{\pageref*{sec:1.8.1}}, \hyperref[idxocc:560:379]{\pageref*{idxocc:560:379}}\par
\par\noindent\hangindent=2.8em\hspace*{1.8em}\phantomsection\label{idx:561}方差分析（ANOVA），\hyperref[ex:11.2.9]{\pageref*{ex:11.2.9}}\par
\par\noindent\hangindent=2.8em\hspace*{1.8em}\phantomsection\label{idx:562}回归（regression），\hyperref[eq:11.3.42]{\pageref*{eq:11.3.42}}\par
\par\noindent\hangindent=2em\hspace*{1em}\phantomsection\label{idx:563}Boole，\hyperref[thm:1.2.11]{\pageref*{thm:1.2.11}}, \hyperref[sec:1.8.1]{\pageref*{sec:1.8.1}}\par
\par\noindent\hangindent=2em\hspace*{1em}\phantomsection\label{idx:564}Cauchy–Schwarz，\hyperref[thm:4.5.7]{\pageref*{thm:4.5.7}}, \hyperref[thm:4.7.3]{\pageref*{thm:4.7.3}}\par
\par\noindent\hangindent=2em\hspace*{1em}\phantomsection\label{idx:565}切比雪夫（Chebychev’s），\hyperref[thm:3.6.1]{\pageref*{thm:3.6.1}}, \hyperref[thm:3.8.2]{\pageref*{thm:3.8.2}}, \hyperref[thm:5.5.2]{\pageref*{thm:5.5.2}}, \hyperref[sec:5.8.3]{\pageref*{sec:5.8.3}}, \hyperref[idxocc:565:409]{\pageref*{idxocc:565:409}}\par
\par\noindent\hangindent=2em\hspace*{1em}\phantomsection\label{idx:566}协方差（covariance），\hyperref[thm:4.7.9]{\pageref*{thm:4.7.9}}\par
\par\noindent\hangindent=2em\hspace*{1em}\phantomsection\label{idx:567}高斯（Gauss），\hyperref[thm:3.8.4]{\pageref*{thm:3.8.4}}\par
\par\noindent\hangindent=2em\hspace*{1em}\phantomsection\label{idx:568}Hölder（Hölder’s），\hyperref[thm:4.7.2]{\pageref*{thm:4.7.2}}, \hyperref[idxocc:568:471]{\pageref*{idxocc:568:471}}\par
\par\noindent\hangindent=2em\hspace*{1em}\phantomsection\label{idx:569}Jensen（Jensen’s），\hyperref[thm:4.7.7]{\pageref*{thm:4.7.7}}\par
\par\noindent\hangindent=2em\hspace*{1em}\phantomsection\label{idx:570}Liapounov（Liapounov’s），\hyperref[eq:4.7.6]{\pageref*{eq:4.7.6}}\par
\par\noindent\hangindent=2em\hspace*{1em}\phantomsection\label{idx:571}对数（logarithm），\hyperref[exercise:1.28]{\pageref*{exercise:1.28}}\par
\par\noindent\hangindent=2em\hspace*{1em}\phantomsection\label{idx:572}马尔可夫（Markov’s），\hyperref[lem:3.8.3]{\pageref*{lem:3.8.3}}\par
\par\noindent\hangindent=2em\hspace*{1em}\phantomsection\label{idx:573}均值（means），\hyperref[ex:4.7.8]{\pageref*{ex:4.7.8}}\par
\par\noindent\hangindent=2em\hspace*{1em}\phantomsection\label{idx:574}矩母函数（mgf），\hyperref[exercise:3.45]{\pageref*{exercise:3.45}}, \hyperref[exercise:5.38]{\pageref*{exercise:5.38}}\par
\par\noindent\hangindent=2em\hspace*{1em}\phantomsection\label{idx:575}Minkowski（Minkowski’s），\hyperref[thm:4.7.5]{\pageref*{thm:4.7.5}}\par
\par\noindent\hangindent=2em\hspace*{1em}\phantomsection\label{idx:576}正态（normal），\hyperref[ex:3.6.3]{\pageref*{ex:3.6.3}}, \hyperref[exercise:3.47]{\pageref*{exercise:3.47}}\par
\par\noindent\hangindent=2em\hspace*{1em}\phantomsection\label{idx:577}幂平均数（power mean），\hyperref[sec:4.9.2]{\pageref*{sec:4.9.2}}\par
\par\noindent\hangindent=2em\hspace*{1em}\phantomsection\label{idx:578}三角（triangle），\hyperref[eq:4.7.8]{\pageref*{eq:4.7.8}}, \hyperref[exercise:4.64]{\pageref*{exercise:4.64}}\par
\par\noindent\hangindent=2em\hspace*{1em}\phantomsection\label{idx:579}Vysochanskiĭ–Petunin（Vysochanskiĭ-Petunin），\hyperref[thm:3.8.5]{\pageref*{thm:3.8.5}}\par
\par\addvspace{2pt}
\par\noindent\hangindent=1em\hspace*{0em}\phantomsection\label{idx:580}影响函数（Influence function），\hyperref[sec:10.6.4]{\pageref*{sec:10.6.4}}\par
\par\noindent\hangindent=2em\hspace*{1em}\phantomsection\label{idx:581}与渐近方差的关系（relation to asymptotic variance），\hyperref[idxocc:581:450]{\pageref*{idxocc:581:450}}\par
\par\addvspace{2pt}
\par\noindent\hangindent=1em\hspace*{0em}\phantomsection\label{idx:582}信息量（Information number），\hyperref[idxocc:582:305]{\pageref*{idxocc:582:305}}, \hyperref[eq:10.1.7]{\pageref*{eq:10.1.7}}\par
\par\addvspace{2pt}
\par\noindent\hangindent=1em\hspace*{0em}\phantomsection\label{idx:583}工具变量（Instrumental variable），\hyperref[sec:12.6.3]{\pageref*{sec:12.6.3}}\par
\par\addvspace{2pt}
\needspace{3\baselineskip}
\par\noindent\hangindent=1em\hspace*{0em}\phantomsection\label{idx:584}换序（Interchange）\par
\par\noindent\hangindent=2em\hspace*{1em}\phantomsection\label{idx:585}求导与求和（differentiation and summation），\hyperref[thm:2.4.8]{\pageref*{thm:2.4.8}}\par
\par\noindent\hangindent=2em\hspace*{1em}\phantomsection\label{idx:586}积分与求和（integration and summation），\hyperref[thm:2.4.10]{\pageref*{thm:2.4.10}}\par
\par\addvspace{2pt}
\needspace{3\baselineskip}
\par\noindent\hangindent=1em\hspace*{0em}\phantomsection\label{idx:587}不变性（Invariance）\par
\par\noindent\hangindent=2em\hspace*{1em}\phantomsection\label{idx:588}分布（distribution），\hyperref[def:6.4.4]{\pageref*{def:6.4.4}}\par
\par\noindent\hangindent=2em\hspace*{1em}\phantomsection\label{idx:589}形式化（formal），\hyperref[idxocc:589:272]{\pageref*{idxocc:589:272}}\par
\par\noindent\hangindent=2.8em\hspace*{1.8em}\phantomsection\label{idx:590}位置（location），\hyperref[exercise:6.41]{\pageref*{exercise:6.41}}\par
\sourcepageanchor{562}
\par\noindent\hangindent=2em\hspace*{1em}\phantomsection\label{idx:591}正态（normal），\hyperref[ex:6.4.6]{\pageref*{ex:6.4.6}}\par
\par\addvspace{2pt}
\par\noindent\hangindent=1em\hspace*{0em}\phantomsection\label{idx:592}逆二项抽样（Inverse binomial sampling），\hyperref[ex:3.2.6]{\pageref*{ex:3.2.6}}\par
\par\addvspace{2pt}
\par\noindent\hangindent=1em\hspace*{0em}\phantomsection\label{idx:593}雅可比行列式（Jacobian），\hyperref[idxocc:593:160]{\pageref*{idxocc:593:160}}\par
\par\addvspace{2pt}
\par\noindent\hangindent=1em\hspace*{0em}\phantomsection\label{idx:594}Karlin–Rubin定理（Karlin–Rubin Theorem），\hyperref[thm:8.3.17]{\pageref*{thm:8.3.17}}\par
\par\addvspace{2pt}
\par\noindent\hangindent=1em\hspace*{0em}\phantomsection\label{idx:595}核（Kernel），\hyperref[idxocc:595:83]{\pageref*{idxocc:595:83}}\par
\par\addvspace{2pt}
\par\noindent\hangindent=1em\hspace*{0em}\phantomsection\label{idx:596}柯尔莫哥洛夫公理（Kolmogorov’s Axioms），\hyperref[def:1.2.4]{\pageref*{def:1.2.4}}\par
\par\noindent\hangindent=2em\hspace*{1em}\phantomsection\label{idx:597}条件概率（conditional probability），\hyperref[idxocc:597:49]{\pageref*{idxocc:597:49}}\par
\par\noindent\hangindent=2em\hspace*{1em}\phantomsection\label{idx:598}诱导概率（induced probability），\hyperref[eq:1.4.1]{\pageref*{eq:1.4.1}}, \hyperref[idxocc:598:55]{\pageref*{idxocc:598:55}}, \hyperref[idxocc:598:70]{\pageref*{idxocc:598:70}}\par
\par\addvspace{2pt}
\par\noindent\hangindent=1em\hspace*{0em}\phantomsection\label{idx:599}峰度（Kurtosis），\hyperref[exercise:2.28]{\pageref*{exercise:2.28}}, \hyperref[exercise:7.45]{\pageref*{exercise:7.45}}\par
\par\addvspace{2pt}
\par\noindent\hangindent=1em\hspace*{0em}\phantomsection\label{idx:600}Lévy定理（Lévy Theorem），\hyperref[sec:5.8.1]{\pageref*{sec:5.8.1}}\par
\par\addvspace{2pt}
\par\noindent\hangindent=1em\hspace*{0em}\phantomsection\label{idx:601}拉普拉斯变换（Laplace transform），\hyperref[eq:2.3.8]{\pageref*{eq:2.3.8}}, \hyperref[idxocc:601:265]{\pageref*{idxocc:601:265}}\par
\par\addvspace{2pt}
\par\noindent\hangindent=1em\hspace*{0em}\phantomsection\label{idx:602}潜在变量（Latent variables），\hyperref[idxocc:602:499]{\pageref*{idxocc:602:499}}\par
\par\addvspace{2pt}
\needspace{3\baselineskip}
\par\noindent\hangindent=1em\hspace*{0em}\phantomsection\label{idx:603}大数定律（Law of Large Numbers）\par
\par\noindent\hangindent=2em\hspace*{1em}\phantomsection\label{idx:604}强（Strong），\hyperref[thm:5.5.9]{\pageref*{thm:5.5.9}}\par
\par\noindent\hangindent=2em\hspace*{1em}\phantomsection\label{idx:605}弱（Weak），\hyperref[thm:5.5.2]{\pageref*{thm:5.5.2}}\par
\par\addvspace{2pt}
\par\noindent\hangindent=1em\hspace*{0em}\phantomsection\label{idx:606}最小二乘估计（Least squares estimate），\hyperref[sec:11.3.1]{\pageref*{sec:11.3.1}}\par
\needspace{3\baselineskip}
\par\noindent\hangindent=2em\hspace*{1em}\phantomsection\label{idx:607}相合性（consistency）\par
\par\noindent\hangindent=2.8em\hspace*{1.8em}\phantomsection\label{idx:608}变量含误差（EIV），\hyperref[sec:12.6.2]{\pageref*{sec:12.6.2}}\par
\par\noindent\hangindent=2em\hspace*{1em}\phantomsection\label{idx:609}变量含误差（EIV），\hyperref[sec:12.2.2]{\pageref*{sec:12.2.2}}\par
\par\noindent\hangindent=2.8em\hspace*{1.8em}\phantomsection\label{idx:610}截距（intercept），\hyperref[idxocc:610:502]{\pageref*{idxocc:610:502}}\par
\par\noindent\hangindent=2.8em\hspace*{1.8em}\phantomsection\label{idx:611}与最大似然估计量的关系（relation to MLE），\hyperref[exercise:12.4]{\pageref*{exercise:12.4}}\par
\par\noindent\hangindent=2.8em\hspace*{1.8em}\phantomsection\label{idx:612}斜率（slope），\hyperref[eq:12.2.11]{\pageref*{eq:12.2.11}}\par
\par\noindent\hangindent=2em\hspace*{1em}\phantomsection\label{idx:613}正规方程（normal equations），\hyperref[exercise:11.25]{\pageref*{exercise:11.25}}\par
\par\noindent\hangindent=2em\hspace*{1em}\phantomsection\label{idx:614}截距的（of intercept），\hyperref[eq:11.3.9]{\pageref*{eq:11.3.9}}\par
\par\noindent\hangindent=2em\hspace*{1em}\phantomsection\label{idx:615}斜率的（of slope），\hyperref[eq:11.3.10]{\pageref*{eq:11.3.10}}\par
\par\addvspace{2pt}
\par\noindent\hangindent=1em\hspace*{0em}\phantomsection\label{idx:616}勒贝格控制收敛定理（Lebesgue’s Dominated Convergence Theorem），\hyperref[idxocc:616:88]{\pageref*{idxocc:616:88}}\par
\par\addvspace{2pt}
\par\noindent\hangindent=1em\hspace*{0em}\phantomsection\label{idx:617}Lehmann–Scheffé定理（Lehmann–Scheffé Theorem），\hyperref[thm:7.5.1]{\pageref*{thm:7.5.1}}\par
\par\addvspace{2pt}
\par\noindent\hangindent=1em\hspace*{0em}\phantomsection\label{idx:618}Leibnitz法则（Leibnitz’s Rule），\hyperref[thm:2.4.1]{\pageref*{thm:2.4.1}}\par
\par\addvspace{2pt}
\par\noindent\hangindent=1em\hspace*{0em}\phantomsection\label{idx:619}似然函数（Likelihood function），\hyperref[def:6.3.1]{\pageref*{def:6.3.1}}\par
\par\noindent\hangindent=2em\hspace*{1em}\phantomsection\label{idx:620}诱导（induced），\hyperref[eq:7.2.5]{\pageref*{eq:7.2.5}}\par
\par\noindent\hangindent=2em\hspace*{1em}\phantomsection\label{idx:621}对数似然（log likelihood），\hyperref[idxocc:621:288]{\pageref*{idxocc:621:288}}, \hyperref[exercise:7.3]{\pageref*{exercise:7.3}}\par
\par\noindent\hangindent=2em\hspace*{1em}\phantomsection\label{idx:622}负二项（negative binomial），\hyperref[ex:6.3.2]{\pageref*{ex:6.3.2}}\par
\par\noindent\hangindent=2em\hspace*{1em}\phantomsection\label{idx:623}充分性（sufficiency），\hyperref[sec:6.6.3]{\pageref*{sec:6.6.3}}\par
\par\addvspace{2pt}
\par\noindent\hangindent=1em\hspace*{0em}\phantomsection\label{idx:624}似然原则（Likelihood Principle），\hyperref[idxocc:624:268]{\pageref*{idxocc:624:268}}\par
\par\noindent\hangindent=2em\hspace*{1em}\phantomsection\label{idx:625}二项与负二项（binomial and negative binomial），\hyperref[ex:6.3.7]{\pageref*{ex:6.3.7}}\par
\par\noindent\hangindent=2em\hspace*{1em}\phantomsection\label{idx:626}Birnbaum定理（Birnbaum’s Theorem），\hyperref[thm:6.3.6]{\pageref*{thm:6.3.6}}, \hyperref[exercise:6.33]{\pageref*{exercise:6.33}}\par
\par\noindent\hangindent=2em\hspace*{1em}\phantomsection\label{idx:627}形式化（formal），\hyperref[sec:6.3.2]{\pageref*{sec:6.3.2}}, \hyperref[idxocc:627:270]{\pageref*{idxocc:627:270}}\par
\par\noindent\hangindent=2em\hspace*{1em}\phantomsection\label{idx:628}解释（interpretations），\hyperref[ex:6.3.7]{\pageref*{ex:6.3.7}}\par
\par\noindent\hangindent=2em\hspace*{1em}\phantomsection\label{idx:629}样本量（sample size），\hyperref[exercise:6.34]{\pageref*{exercise:6.34}}\par
\par\addvspace{2pt}
\par\noindent\hangindent=1em\hspace*{0em}\phantomsection\label{idx:630}似然比检验（Likelihood ratio test (LRT)），\hyperref[def:8.2.1]{\pageref*{def:8.2.1}}\par
\par\noindent\hangindent=2em\hspace*{1em}\phantomsection\label{idx:631}方差分析（ANOVA），\hyperref[exercise:11.13]{\pageref*{exercise:11.13}}\par
\par\noindent\hangindent=2em\hspace*{1em}\phantomsection\label{idx:632}渐近（asymptotic），\hyperref[thm:10.3.3]{\pageref*{thm:10.3.3}}\par
\par\noindent\hangindent=2em\hspace*{1em}\phantomsection\label{idx:633}贝塔（beta），\hyperref[exercise:8.17]{\pageref*{exercise:8.17}}\par
\par\noindent\hangindent=2em\hspace*{1em}\phantomsection\label{idx:634}指数（exponential），\hyperref[ex:8.2.3]{\pageref*{ex:8.2.3}}, \hyperref[ex:8.2.5]{\pageref*{ex:8.2.5}}, \hyperref[ex:8.3.7]{\pageref*{ex:8.3.7}}, \hyperref[exercise:8.6]{\pageref*{exercise:8.6}}\par
\par\noindent\hangindent=2em\hspace*{1em}\phantomsection\label{idx:635}最大似然估计量（MLE），\hyperref[idxocc:635:335]{\pageref*{idxocc:635:335}}\par
\par\noindent\hangindent=2em\hspace*{1em}\phantomsection\label{idx:636}多项（multinomial），\hyperref[ex:10.3.4]{\pageref*{ex:10.3.4}}\par
\par\noindent\hangindent=2em\hspace*{1em}\phantomsection\label{idx:637}Neyman–Pearson引理（Neyman–Pearson Lemma），\hyperref[exercise:8.24]{\pageref*{exercise:8.24}}\par
\par\noindent\hangindent=2em\hspace*{1em}\phantomsection\label{idx:638}正态（normal），\hyperref[ex:8.2.2]{\pageref*{ex:8.2.2}}, \hyperref[ex:8.2.6]{\pageref*{ex:8.2.6}}, \hyperref[ex:8.3.7]{\pageref*{ex:8.3.7}}, \hyperref[exercise:8.37]{\pageref*{exercise:8.37}}\par
\par\noindent\hangindent=2.8em\hspace*{1.8em}\phantomsection\label{idx:639}滋扰参数（nuisance parameter），\hyperref[ex:8.2.6]{\pageref*{ex:8.2.6}}\par
\par\noindent\hangindent=2.8em\hspace*{1.8em}\phantomsection\label{idx:640}双样本（two-sample），\hyperref[exercise:8.41]{\pageref*{exercise:8.41}}\par
\par\noindent\hangindent=2em\hspace*{1em}\phantomsection\label{idx:641}帕累托（Pareto），\hyperref[exercise:8.5]{\pageref*{exercise:8.5}}\par
\par\noindent\hangindent=2em\hspace*{1em}\phantomsection\label{idx:642}充分统计量（sufficient statistic），\hyperref[thm:8.2.4]{\pageref*{thm:8.2.4}}\par
\par\addvspace{2pt}
\par\noindent\hangindent=1em\hspace*{0em}\phantomsection\label{idx:643}Lindeberg–Feller条件（Lindeberg–Feller Condition），\hyperref[sec:5.8.1]{\pageref*{sec:5.8.1}}\par
\par\addvspace{2pt}
\par\noindent\hangindent=1em\hspace*{0em}\phantomsection\label{idx:644}线性组合（Linear combination），\hyperref[def:11.2.4]{\pageref*{def:11.2.4}}\par
\par\addvspace{2pt}
\needspace{3\baselineskip}
\par\noindent\hangindent=1em\hspace*{0em}\phantomsection\label{idx:645}位置族（Location family）\par
\par\noindent\hangindent=2em\hspace*{1em}\phantomsection\label{idx:646}枢轴量（pivotal quantity），\hyperref[ex:9.2.7]{\pageref*{ex:9.2.7}}\par
\par\addvspace{2pt}
\par\noindent\hangindent=1em\hspace*{0em}\phantomsection\label{idx:647}位置尺度族（Location–scale family），\hyperref[sec:3.5]{\pageref*{sec:3.5}}\par
\par\noindent\hangindent=2em\hspace*{1em}\phantomsection\label{idx:648}位置参数（location parameter），\hyperref[def:3.5.2]{\pageref*{def:3.5.2}}, \hyperref[def:3.5.5]{\pageref*{def:3.5.5}}\par
\par\noindent\hangindent=2em\hspace*{1em}\phantomsection\label{idx:649}枢轴量（pivotal quantity），\hyperref[ex:9.2.7]{\pageref*{ex:9.2.7}}\par
\par\noindent\hangindent=2em\hspace*{1em}\phantomsection\label{idx:650}尺度参数（scale parameter），\hyperref[def:3.5.4]{\pageref*{def:3.5.4}}\par
\par\noindent\hangindent=2em\hspace*{1em}\phantomsection\label{idx:651}标准概率密度函数（standard pdf），\hyperref[idxocc:651:126]{\pageref*{idxocc:651:126}}, \hyperref[def:3.5.5]{\pageref*{def:3.5.5}}\par
\par\noindent\hangindent=2em\hspace*{1em}\phantomsection\label{idx:652}随机序（stochastic order），\hyperref[exercise:8.34]{\pageref*{exercise:8.34}}\par
\par\addvspace{2pt}
\par\noindent\hangindent=1em\hspace*{0em}\phantomsection\label{idx:653}损失函数（Loss function），\hyperref[idxocc:653:313]{\pageref*{idxocc:653:313}}\par
\par\noindent\hangindent=2em\hspace*{1em}\phantomsection\label{idx:654}绝对误差（absolute error），\hyperref[ex:7.3.28]{\pageref*{ex:7.3.28}}\par
\par\noindent\hangindent=2em\hspace*{1em}\phantomsection\label{idx:655}广义零一（generalized zero–one），\hyperref[eq:8.3.11]{\pageref*{eq:8.3.11}}\par
\par\noindent\hangindent=2em\hspace*{1em}\phantomsection\label{idx:656}假设检验（hypothesis testing），\hyperref[eq:8.3.11]{\pageref*{eq:8.3.11}}\par
\par\noindent\hangindent=2em\hspace*{1em}\phantomsection\label{idx:657}区间估计（interval estimation），\hyperref[eq:9.3.4]{\pageref*{eq:9.3.4}}\par
\par\noindent\hangindent=2em\hspace*{1em}\phantomsection\label{idx:658}似然（likelihood），\hyperref[ex:7.3.27]{\pageref*{ex:7.3.27}}\par
\par\noindent\hangindent=2em\hspace*{1em}\phantomsection\label{idx:659}点估计（point estimation），\hyperref[ex:7.3.28]{\pageref*{ex:7.3.28}}\par
\par\noindent\hangindent=2em\hspace*{1em}\phantomsection\label{idx:660}后验期望损失（posterior expected loss），\hyperref[eq:7.3.19]{\pageref*{eq:7.3.19}}\par
\par\noindent\hangindent=2em\hspace*{1em}\phantomsection\label{idx:661}平方误差（squared error），\hyperref[ex:7.3.28]{\pageref*{ex:7.3.28}}\par
\par\noindent\hangindent=2em\hspace*{1em}\phantomsection\label{idx:662}Stein损失（Stein’s loss），\hyperref[ex:7.3.27]{\pageref*{ex:7.3.27}}\par
\par\noindent\hangindent=2em\hspace*{1em}\phantomsection\label{idx:663}零一（zero–one），\hyperref[eq:8.3.11]{\pageref*{eq:8.3.11}}\par
\par\addvspace{2pt}
\par\noindent\hangindent=1em\hspace*{0em}\phantomsection\label{idx:664}映射（Mapping），\hyperref[idxocc:664:69]{\pageref*{idxocc:664:69}}\par
\par\noindent\hangindent=2em\hspace*{1em}\phantomsection\label{idx:665}逆（inverse），\hyperref[idxocc:665:69]{\pageref*{idxocc:665:69}}\par
\par\addvspace{2pt}
\par\noindent\hangindent=1em\hspace*{0em}\phantomsection\label{idx:666}马尔可夫链蒙特卡洛（Markov chain Monte Carlo），\hyperref[sec:5.8.5]{\pageref*{sec:5.8.5}}\par
\par\addvspace{2pt}
\par\noindent\hangindent=1em\hspace*{0em}\phantomsection\label{idx:667}最大似然估计量（Maximum likelihood estimator (MLE)），\hyperref[def:7.2.4]{\pageref*{def:7.2.4}}\par
\par\noindent\hangindent=2em\hspace*{1em}\phantomsection\label{idx:668}渐近效率（asymptotic efficiency），\hyperref[thm:10.1.12]{\pageref*{thm:10.1.12}}\par
\par\noindent\hangindent=2em\hspace*{1em}\phantomsection\label{idx:669}伯努利（Bernoulli），\hyperref[ex:7.2.7]{\pageref*{ex:7.2.7}}\par
\par\noindent\hangindent=2.8em\hspace*{1.8em}\phantomsection\label{idx:670}近似方差（approximate variance），\hyperref[ex:10.1.14]{\pageref*{ex:10.1.14}}\par
\par\noindent\hangindent=2em\hspace*{1em}\phantomsection\label{idx:671}二项（binomial），\hyperref[ex:7.2.9]{\pageref*{ex:7.2.9}}, \hyperref[ex:7.2.13]{\pageref*{ex:7.2.13}}\par
\par\noindent\hangindent=2em\hspace*{1em}\phantomsection\label{idx:672}二元正态（bivariate normal），\hyperref[exercise:7.18]{\pageref*{exercise:7.18}}\par
\par\noindent\hangindent=2em\hspace*{1em}\phantomsection\label{idx:673}微积分（calculus），\hyperref[ex:7.2.12]{\pageref*{ex:7.2.12}}\par
\par\noindent\hangindent=2em\hspace*{1em}\phantomsection\label{idx:674}相合性（consistency），\hyperref[thm:10.1.6]{\pageref*{thm:10.1.6}}\par
\par\noindent\hangindent=2.8em\hspace*{1.8em}\phantomsection\label{idx:675}变量含误差（EIV），\hyperref[idxocc:675:505]{\pageref*{idxocc:675:505}}, \hyperref[idxocc:675:507]{\pageref*{idxocc:675:507}}\par
\par\noindent\hangindent=2em\hspace*{1em}\phantomsection\label{idx:676}不稳定性（instability），\hyperref[ex:7.2.13]{\pageref*{ex:7.2.13}}\par
\par\noindent\hangindent=2em\hspace*{1em}\phantomsection\label{idx:677}不变性（invariance），\hyperref[thm:7.2.10]{\pageref*{thm:7.2.10}}\par
\par\noindent\hangindent=2em\hspace*{1em}\phantomsection\label{idx:678}似然比检验（LRT），\hyperref[idxocc:678:335]{\pageref*{idxocc:678:335}}\par
\par\noindent\hangindent=2em\hspace*{1em}\phantomsection\label{idx:679}正态（normal），\hyperref[ex:7.2.5]{\pageref*{ex:7.2.5}}, \hyperref[ex:7.2.11]{\pageref*{ex:7.2.11}}\par
\par\noindent\hangindent=2em\hspace*{1em}\phantomsection\label{idx:680}与矩估计法的关系（relation to method of moments），\hyperref[sec:7.5.1]{\pageref*{sec:7.5.1}}\par
\par\addvspace{2pt}
\par\noindent\hangindent=1em\hspace*{0em}\phantomsection\label{idx:681}二次函数的最大值（Maximum of quadratic functions），\hyperref[lem:11.2.7]{\pageref*{lem:11.2.7}}, \hyperref[eq:11.3.18]{\pageref*{eq:11.3.18}}, \hyperref[thm:11.3.6]{\pageref*{thm:11.3.6}}, \hyperref[exercise:11.14]{\pageref*{exercise:11.14}}, \hyperref[item:11.40.a]{\pageref*{item:11.40.a}}\par
\par\addvspace{2pt}
\par\noindent\hangindent=1em\hspace*{0em}\phantomsection\label{idx:682}最大概率估计量（Maximum probability estimator），\hyperref[item:11.37.e]{\pageref*{item:11.37.e}}\par
\par\addvspace{2pt}
\par\noindent\hangindent=1em\hspace*{0em}\phantomsection\label{idx:683}均方误差（Mean squared error），\hyperref[def:7.3.1]{\pageref*{def:7.3.1}}, \hyperref[ex:7.3.28]{\pageref*{ex:7.3.28}}\par
\par\addvspace{2pt}
\par\noindent\hangindent=1em\hspace*{0em}\phantomsection\label{idx:684}中位数（Median），\hyperref[exercise:2.17]{\pageref*{exercise:2.17}}\par
\par\noindent\hangindent=2em\hspace*{1em}\phantomsection\label{idx:685}最小化性质（minimizing property），\hyperref[exercise:2.18]{\pageref*{exercise:2.18}}\par
\par\addvspace{2pt}
\par\noindent\hangindent=1em\hspace*{0em}\phantomsection\label{idx:686}德梅雷骰子问题（Meré dice problem），\hyperref[ex:1.3.8]{\pageref*{ex:1.3.8}}, \hyperref[ex:3.2.3]{\pageref*{ex:3.2.3}}\par
\par\addvspace{2pt}
\par\noindent\hangindent=1em\hspace*{0em}\phantomsection\label{idx:687}度量（Metric），\hyperref[item:12.3.b]{\pageref*{item:12.3.b}}\par
\par\addvspace{2pt}
\par\noindent\hangindent=1em\hspace*{0em}\phantomsection\label{idx:688}Metropolis算法（Metropolis Algorithm），\hyperref[idxocc:688:238]{\pageref*{idxocc:688:238}}\par
\sourcepageanchor{563}
\par\addvspace{2pt}
\par\noindent\hangindent=1em\hspace*{0em}\phantomsection\label{idx:689}矩母函数（Moment generating function (mgf)），\hyperref[def:2.3.6]{\pageref*{def:2.3.6}}, \hyperref[sec:2.6.3]{\pageref*{sec:2.6.3}}\par
\par\noindent\hangindent=2em\hspace*{1em}\phantomsection\label{idx:690}收敛（convergence），\hyperref[thm:2.3.12]{\pageref*{thm:2.3.12}}\par
\par\noindent\hangindent=2em\hspace*{1em}\phantomsection\label{idx:691}同分布（identically distributed），\hyperref[thm:2.3.11]{\pageref*{thm:2.3.11}}\par
\par\noindent\hangindent=2em\hspace*{1em}\phantomsection\label{idx:692}独立随机变量（independent random variables），\hyperref[thm:4.2.12]{\pageref*{thm:4.2.12}}\par
\par\noindent\hangindent=2em\hspace*{1em}\phantomsection\label{idx:693}性质（properties），\hyperref[thm:2.3.15]{\pageref*{thm:2.3.15}}\par
\par\noindent\hangindent=2em\hspace*{1em}\phantomsection\label{idx:694}样本均值（sample mean），\hyperref[thm:5.2.7]{\pageref*{thm:5.2.7}}\par
\par\noindent\hangindent=2em\hspace*{1em}\phantomsection\label{idx:695}和（sum），\hyperref[thm:4.6.7]{\pageref*{thm:4.6.7}}\par
\par\noindent\hangindent=2em\hspace*{1em}\phantomsection\label{idx:696}在中心极限定理中的用途（use in CLT），\hyperref[idxocc:696:224]{\pageref*{idxocc:696:224}}\par
\par\addvspace{2pt}
\par\noindent\hangindent=1em\hspace*{0em}\phantomsection\label{idx:697}矩（Moments），\hyperref[def:2.3.1]{\pageref*{def:2.3.1}}\par
\par\noindent\hangindent=2em\hspace*{1em}\phantomsection\label{idx:698}非唯一性（nonuniqueness），\hyperref[ex:2.3.10]{\pageref*{ex:2.3.10}}, \hyperref[exercise:2.35]{\pageref*{exercise:2.35}}\par
\par\noindent\hangindent=2em\hspace*{1em}\phantomsection\label{idx:699}唯一性（uniqueness），\hyperref[sec:2.6.1]{\pageref*{sec:2.6.1}}\par
\par\addvspace{2pt}
\par\noindent\hangindent=1em\hspace*{0em}\phantomsection\label{idx:700}单调似然比（Monotone likelihood ratio (MLR)），\hyperref[def:8.3.16]{\pageref*{def:8.3.16}}\par
\par\noindent\hangindent=2em\hspace*{1em}\phantomsection\label{idx:701}二项（binomial），\hyperref[exercise:8.25]{\pageref*{exercise:8.25}}\par
\par\noindent\hangindent=2em\hspace*{1em}\phantomsection\label{idx:702}柯西（Cauchy），\hyperref[exercise:8.29]{\pageref*{exercise:8.29}}\par
\par\noindent\hangindent=2em\hspace*{1em}\phantomsection\label{idx:703}指数族（exponential family），\hyperref[exercise:8.27]{\pageref*{exercise:8.27}}\par
\par\noindent\hangindent=2em\hspace*{1em}\phantomsection\label{idx:704}逻辑（logistic），\hyperref[exercise:8.28]{\pageref*{exercise:8.28}}\par
\par\noindent\hangindent=2em\hspace*{1em}\phantomsection\label{idx:705}非中心t（noncentral t），\hyperref[exercise:8.35]{\pageref*{exercise:8.35}}\par
\par\noindent\hangindent=2em\hspace*{1em}\phantomsection\label{idx:706}正态（normal），\hyperref[exercise:8.25]{\pageref*{exercise:8.25}}\par
\par\noindent\hangindent=2em\hspace*{1em}\phantomsection\label{idx:707}泊松（Poisson），\hyperref[exercise:8.25]{\pageref*{exercise:8.25}}\par
\par\noindent\hangindent=2em\hspace*{1em}\phantomsection\label{idx:708}随机序（stochastic order），\hyperref[exercise:8.26]{\pageref*{exercise:8.26}}\par
\par\addvspace{2pt}
\needspace{3\baselineskip}
\par\noindent\hangindent=1em\hspace*{0em}\phantomsection\label{idx:709}蒙特卡洛（Monte Carlo）\par
\par\noindent\hangindent=2em\hspace*{1em}\phantomsection\label{idx:710}巧算（swindle），\hyperref[exercise:6.31]{\pageref*{exercise:6.31}}\par
\par\addvspace{2pt}
\par\noindent\hangindent=1em\hspace*{0em}\phantomsection\label{idx:711}蒙提霍尔问题（Monte Hall problem），\hyperref[exercise:1.37]{\pageref*{exercise:1.37}}\par
\par\addvspace{2pt}
\par\noindent\hangindent=1em\hspace*{0em}\phantomsection\label{idx:712}多项式系数（Multinomial coefficient），\hyperref[idxocc:712:179]{\pageref*{idxocc:712:179}}\par
\par\addvspace{2pt}
\par\noindent\hangindent=1em\hspace*{0em}\phantomsection\label{idx:713}多项式定理（Multinomial Theorem），\hyperref[thm:4.6.4]{\pageref*{thm:4.6.4}}\par
\par\addvspace{2pt}
\par\noindent\hangindent=1em\hspace*{0em}\phantomsection\label{idx:714}多重比较（Multiple comparisons），\hyperref[sec:11.5.2]{\pageref*{sec:11.5.2}}\par
\par\noindent\hangindent=2em\hspace*{1em}\phantomsection\label{idx:715}逐次比较错误率（comparisonwise error rate），\hyperref[item:11.17.b]{\pageref*{item:11.17.b}}\par
\par\noindent\hangindent=2em\hspace*{1em}\phantomsection\label{idx:716}整体试验错误率（experimentwise error rate），\hyperref[item:11.17.b]{\pageref*{item:11.17.b}}\par
\par\noindent\hangindent=2em\hspace*{1em}\phantomsection\label{idx:717}受保护的最小显著差（protected LSD），\hyperref[exercise:11.17]{\pageref*{exercise:11.17}}\par
\par\noindent\hangindent=2em\hspace*{1em}\phantomsection\label{idx:718}Tukey的Q方法（Tukey’s Q method），\hyperref[idxocc:718:494]{\pageref*{idxocc:718:494}}\par
\par\addvspace{2pt}
\par\noindent\hangindent=1em\hspace*{0em}\phantomsection\label{idx:719}必要统计量（Necessary statistic），\hyperref[def:6.6.3]{\pageref*{def:6.6.3}}\par
\par\addvspace{2pt}
\par\noindent\hangindent=1em\hspace*{0em}\phantomsection\label{idx:720}邻域（Neighborhood），\hyperref[def:2.3.6]{\pageref*{def:2.3.6}}\par
\par\addvspace{2pt}
\par\noindent\hangindent=1em\hspace*{0em}\phantomsection\label{idx:721}Neyman最短（Neyman shortest），\hyperref[idxocc:721:392]{\pageref*{idxocc:721:392}}\par
\par\addvspace{2pt}
\par\noindent\hangindent=1em\hspace*{0em}\phantomsection\label{idx:722}Neyman–Pearson引理（Neyman–Pearson Lemma），\hyperref[thm:8.3.12]{\pageref*{thm:8.3.12}}\par
\par\noindent\hangindent=2em\hspace*{1em}\phantomsection\label{idx:723}似然比检验（LRT），\hyperref[exercise:8.24]{\pageref*{exercise:8.24}}\par
\par\addvspace{2pt}
\needspace{3\baselineskip}
\par\noindent\hangindent=1em\hspace*{0em}\phantomsection\label{idx:724}尼罗河（Nile）\par
\par\noindent\hangindent=2em\hspace*{1em}\phantomsection\label{idx:725}问题（problem），\hyperref[exercise:6.37]{\pageref*{exercise:6.37}}\par
\par\addvspace{2pt}
\par\noindent\hangindent=1em\hspace*{0em}\phantomsection\label{idx:726}正规方程（Normal equations），\hyperref[exercise:11.25]{\pageref*{exercise:11.25}}\par
\par\addvspace{2pt}
\par\noindent\hangindent=1em\hspace*{0em}\phantomsection\label{idx:727}滋扰参数（Nuisance parameter），\hyperref[idxocc:727:337]{\pageref*{idxocc:727:337}}\par
\par\addvspace{2pt}
\par\noindent\hangindent=1em\hspace*{0em}\phantomsection\label{idx:728}优势（Odds），\hyperref[idxocc:728:414]{\pageref*{idxocc:728:414}}\par
\par\addvspace{2pt}
\par\noindent\hangindent=1em\hspace*{0em}\phantomsection\label{idx:729}一一变换（One-to-one transformation），\hyperref[idxocc:729:71]{\pageref*{idxocc:729:71}}, \hyperref[idxocc:729:160]{\pageref*{idxocc:729:160}}\par
\par\addvspace{2pt}
\par\noindent\hangindent=1em\hspace*{0em}\phantomsection\label{idx:730}满射变换（Onto transformation），\hyperref[idxocc:730:71]{\pageref*{idxocc:730:71}}, \hyperref[idxocc:730:160]{\pageref*{idxocc:730:160}}\par
\par\addvspace{2pt}
\par\noindent\hangindent=1em\hspace*{0em}\phantomsection\label{idx:731}次序统计量（Order statistics），\hyperref[def:5.4.1]{\pageref*{def:5.4.1}}\par
\par\noindent\hangindent=2em\hspace*{1em}\phantomsection\label{idx:732}累积分布函数（cdf），\hyperref[thm:5.4.3]{\pageref*{thm:5.4.3}}\par
\par\noindent\hangindent=2em\hspace*{1em}\phantomsection\label{idx:733}联合累积分布函数（joint cdf），\hyperref[exercise:5.26]{\pageref*{exercise:5.26}}\par
\par\noindent\hangindent=2em\hspace*{1em}\phantomsection\label{idx:734}联合概率密度函数（joint pdf），\hyperref[thm:5.4.6]{\pageref*{thm:5.4.6}}\par
\par\noindent\hangindent=2em\hspace*{1em}\phantomsection\label{idx:735}概率密度函数（pdf），\hyperref[thm:5.4.4]{\pageref*{thm:5.4.4}}\par
\par\noindent\hangindent=2em\hspace*{1em}\phantomsection\label{idx:736}概率质量函数（pmf），\hyperref[thm:5.4.3]{\pageref*{thm:5.4.3}}\par
\par\noindent\hangindent=2em\hspace*{1em}\phantomsection\label{idx:737}间距（spacings），\hyperref[idxocc:737:183]{\pageref*{idxocc:737:183}}\par
\par\noindent\hangindent=2em\hspace*{1em}\phantomsection\label{idx:738}充分（sufficient），\hyperref[exercise:6.8]{\pageref*{exercise:6.8}}\par
\par\addvspace{2pt}
\par\noindent\hangindent=1em\hspace*{0em}\phantomsection\label{idx:739}正交最小二乘（Orthogonal least squares），\hyperref[sec:12.2.2]{\pageref*{sec:12.2.2}}\par
\par\addvspace{2pt}
\needspace{3\baselineskip}
\par\noindent\hangindent=1em\hspace*{0em}\phantomsection\label{idx:740}过参数化模型（Overparameterized model）\par
\par\noindent\hangindent=2em\hspace*{1em}\phantomsection\label{idx:741}方差分析（ANOVA），\hyperref[eq:11.2.2]{\pageref*{eq:11.2.2}}\par
\par\addvspace{2pt}
\par\noindent\hangindent=1em\hspace*{0em}\phantomsection\label{idx:742}p值（P-value），\hyperref[def:8.3.26]{\pageref*{def:8.3.26}}\par
\par\noindent\hangindent=2em\hspace*{1em}\phantomsection\label{idx:743}条件（conditional），\hyperref[eq:8.3.10]{\pageref*{eq:8.3.10}}\par
\par\noindent\hangindent=2em\hspace*{1em}\phantomsection\label{idx:744}与贝叶斯检验的关系（relation to Bayes test），\hyperref[sec:8.5.3]{\pageref*{sec:8.5.3}}\par
\par\addvspace{2pt}
\par\noindent\hangindent=1em\hspace*{0em}\phantomsection\label{idx:745}两两不交集（Pairwise disjoint sets），\hyperref[def:1.1.5]{\pageref*{def:1.1.5}}\par
\par\addvspace{2pt}
\needspace{3\baselineskip}
\par\noindent\hangindent=1em\hspace*{0em}\phantomsection\label{idx:746}悖论（Paradox）\par
\par\noindent\hangindent=2em\hspace*{1em}\phantomsection\label{idx:747}博雷尔（Borel），\hyperref[sec:4.9.3]{\pageref*{sec:4.9.3}}\par
\par\noindent\hangindent=2em\hspace*{1em}\phantomsection\label{idx:748}交换（exchange），\hyperref[sec:4.9.1]{\pageref*{sec:4.9.1}}\par
\par\addvspace{2pt}
\par\noindent\hangindent=1em\hspace*{0em}\phantomsection\label{idx:749}并联系统（Parallel system），\hyperref[exercise:4.55]{\pageref*{exercise:4.55}}\par
\par\addvspace{2pt}
\par\noindent\hangindent=1em\hspace*{0em}\phantomsection\label{idx:750}参数（Parameter），\hyperref[ex:2.1.1]{\pageref*{ex:2.1.1}}\par
\par\addvspace{2pt}
\par\noindent\hangindent=1em\hspace*{0em}\phantomsection\label{idx:751}参数空间（Parameter space），\hyperref[idxocc:751:313]{\pageref*{idxocc:751:313}}\par
\par\addvspace{2pt}
\par\noindent\hangindent=1em\hspace*{0em}\phantomsection\label{idx:752}参数自助法（Parametric bootstrap），\hyperref[ex:5.6.6]{\pageref*{ex:5.6.6}}\par
\par\addvspace{2pt}
\par\noindent\hangindent=1em\hspace*{0em}\phantomsection\label{idx:753}划分（Partition），\hyperref[def:1.1.6]{\pageref*{def:1.1.6}}, \hyperref[thm:1.3.5]{\pageref*{thm:1.3.5}}, \hyperref[idxocc:753:252]{\pageref*{idxocc:753:252}}, \hyperref[idxocc:753:259]{\pageref*{idxocc:753:259}}\par
\par\addvspace{2pt}
\par\noindent\hangindent=1em\hspace*{0em}\phantomsection\label{idx:754}百分位数（Percentile），\hyperref[idxocc:754:216]{\pageref*{idxocc:754:216}}\par
\par\addvspace{2pt}
\par\noindent\hangindent=1em\hspace*{0em}\phantomsection\label{idx:755}Pitman估计量（见\hyperref[idx:787]{点估计量：不变}）（Pitman Estimator (see Point estimator, invariant)），\hyperref[exercise:7.35]{\pageref*{exercise:7.35}}\par
\par\addvspace{2pt}
\par\noindent\hangindent=1em\hspace*{0em}\phantomsection\label{idx:756}枢轴量（Pivotal quantity），\hyperref[def:9.2.6]{\pageref*{def:9.2.6}}\par
\par\noindent\hangindent=2em\hspace*{1em}\phantomsection\label{idx:757}伽马（gamma），\hyperref[ex:9.2.8]{\pageref*{ex:9.2.8}}\par
\par\noindent\hangindent=2em\hspace*{1em}\phantomsection\label{idx:758}一般形式（general form），\hyperref[eq:9.2.11]{\pageref*{eq:9.2.11}}\par
\par\noindent\hangindent=2em\hspace*{1em}\phantomsection\label{idx:759}位置族（location family），\hyperref[ex:9.2.7]{\pageref*{ex:9.2.7}}\par
\par\noindent\hangindent=2em\hspace*{1em}\phantomsection\label{idx:760}位置尺度族（location–scale family），\hyperref[ex:9.2.7]{\pageref*{ex:9.2.7}}\par
\par\noindent\hangindent=2em\hspace*{1em}\phantomsection\label{idx:761}尺度族（scale family），\hyperref[ex:9.2.7]{\pageref*{ex:9.2.7}}\par
\par\addvspace{2pt}
\par\noindent\hangindent=1em\hspace*{0em}\phantomsection\label{idx:762}点估计值（Point estimate），\hyperref[idxocc:762:283]{\pageref*{idxocc:762:283}}\par
\par\addvspace{2pt}
\par\noindent\hangindent=1em\hspace*{0em}\phantomsection\label{idx:763}点估计量（Point estimator），\hyperref[def:7.1.1]{\pageref*{def:7.1.1}}\par
\par\noindent\hangindent=2em\hspace*{1em}\phantomsection\label{idx:764}近似（approximation），\hyperref[eq:5.5.9]{\pageref*{eq:5.5.9}}\par
\needspace{3\baselineskip}
\par\noindent\hangindent=2em\hspace*{1em}\phantomsection\label{idx:765}贝叶斯（Bayes）\par
\par\noindent\hangindent=2.8em\hspace*{1.8em}\phantomsection\label{idx:766}经验（empirical），\hyperref[idxocc:766:331]{\pageref*{idxocc:766:331}}\par
\par\noindent\hangindent=2.8em\hspace*{1.8em}\phantomsection\label{idx:767}分层（hierarchical），\hyperref[idxocc:767:332]{\pageref*{idxocc:767:332}}\par
\par\noindent\hangindent=2.8em\hspace*{1.8em}\phantomsection\label{idx:768}稳健（robust），\hyperref[idxocc:768:331]{\pageref*{idxocc:768:331}}\par
\par\noindent\hangindent=2em\hspace*{1em}\phantomsection\label{idx:769}崩溃值（breakdown value），\hyperref[def:10.2.2]{\pageref*{def:10.2.2}}\par
\par\noindent\hangindent=2em\hspace*{1em}\phantomsection\label{idx:770}相合（consistent），\hyperref[def:10.1.1]{\pageref*{def:10.1.1}}\par
\par\noindent\hangindent=2.8em\hspace*{1.8em}\phantomsection\label{idx:771}最大似然估计量（maximum likelihood estimator），\hyperref[thm:10.1.6]{\pageref*{thm:10.1.6}}\par
\par\noindent\hangindent=2.8em\hspace*{1.8em}\phantomsection\label{idx:772}正态（normal），\hyperref[ex:10.1.2]{\pageref*{ex:10.1.2}}\par
\par\noindent\hangindent=2em\hspace*{1em}\phantomsection\label{idx:773}相合的自助法（consistent bootstrap），\hyperref[item:10.1.bootstrap.b]{\pageref*{item:10.1.bootstrap.b}}\par
\needspace{3\baselineskip}
\par\noindent\hangindent=2em\hspace*{1em}\phantomsection\label{idx:774}相关（correlation）\par
\par\noindent\hangindent=2.8em\hspace*{1.8em}\phantomsection\label{idx:775}极限分布（limit distribution），\hyperref[item:10.17.d]{\pageref*{item:10.17.d}}\par
\par\noindent\hangindent=2em\hspace*{1em}\phantomsection\label{idx:776}决策论（decision theoretic），\hyperref[ex:7.3.28]{\pageref*{ex:7.3.28}}\par
\par\noindent\hangindent=2.8em\hspace*{1.8em}\phantomsection\label{idx:777}二项（binomial），\hyperref[ex:7.3.30]{\pageref*{ex:7.3.30}}\par
\par\noindent\hangindent=2.8em\hspace*{1.8em}\phantomsection\label{idx:778}方差（variance），\hyperref[ex:7.3.26]{\pageref*{ex:7.3.26}}\par
\par\noindent\hangindent=2em\hspace*{1em}\phantomsection\label{idx:779}效率（efficiency），\hyperref[sec:10.1.2]{\pageref*{sec:10.1.2}}, \hyperref[def:10.1.11]{\pageref*{def:10.1.11}}\par
\par\noindent\hangindent=2em\hspace*{1em}\phantomsection\label{idx:780}有效（efficient），\hyperref[idxocc:780:415]{\pageref*{idxocc:780:415}}\par
\needspace{3\baselineskip}
\par\noindent\hangindent=2em\hspace*{1em}\phantomsection\label{idx:781}经验贝叶斯（empirical Bayes）\par
\par\noindent\hangindent=2.8em\hspace*{1.8em}\phantomsection\label{idx:782}正态（normal），\hyperref[exercise:7.25]{\pageref*{exercise:7.25}}\par
\needspace{3\baselineskip}
\par\noindent\hangindent=2em\hspace*{1em}\phantomsection\label{idx:783}等变（equivariant）\par
\par\noindent\hangindent=2.8em\hspace*{1.8em}\phantomsection\label{idx:784}位置（location），\hyperref[exercise:6.41]{\pageref*{exercise:6.41}}\par
\par\noindent\hangindent=2.8em\hspace*{1.8em}\phantomsection\label{idx:785}位置尺度（location–scale），\hyperref[exercise:6.42]{\pageref*{exercise:6.42}}\par
\par\noindent\hangindent=2em\hspace*{1em}\phantomsection\label{idx:786}广义均值（generalized means），\hyperref[exercise:7.16]{\pageref*{exercise:7.16}}\par
\needspace{3\baselineskip}
\par\noindent\hangindent=2em\hspace*{1em}\phantomsection\label{idx:787}不变（invariant）\par
\par\noindent\hangindent=2.8em\hspace*{1.8em}\phantomsection\label{idx:788}形式化（formal），\hyperref[idxocc:788:301]{\pageref*{idxocc:788:301}}\par
\sourcepageanchor{564}
\par\noindent\hangindent=2.8em\hspace*{1.8em}\phantomsection\label{idx:789}均方误差（mean squared error），\hyperref[ex:7.3.6]{\pageref*{ex:7.3.6}}\par
\par\noindent\hangindent=2.8em\hspace*{1.8em}\phantomsection\label{idx:790}Pitman，\hyperref[exercise:7.35]{\pageref*{exercise:7.35}}\par
\par\noindent\hangindent=2em\hspace*{1em}\phantomsection\label{idx:791}刀切法（jackknife），\hyperref[exercise:7.66]{\pageref*{exercise:7.66}}\par
\par\noindent\hangindent=2em\hspace*{1em}\phantomsection\label{idx:792}线性（linear），\hyperref[exercise:7.41]{\pageref*{exercise:7.41}}, \hyperref[eq:11.3.15]{\pageref*{eq:11.3.15}}\par
\needspace{3\baselineskip}
\par\noindent\hangindent=2em\hspace*{1em}\phantomsection\label{idx:793}M估计量（M-estimator）\par
\par\noindent\hangindent=2.8em\hspace*{1.8em}\phantomsection\label{idx:794}渐近有效（asymptotically efficient），\hyperref[exercise:10.29]{\pageref*{exercise:10.29}}\par
\needspace{3\baselineskip}
\par\noindent\hangindent=2em\hspace*{1em}\phantomsection\label{idx:795}最大似然估计量（maximum likelihood estimator）\par
\par\noindent\hangindent=2.8em\hspace*{1.8em}\phantomsection\label{idx:796}渐近效率（asymptotic efficiency），\hyperref[thm:10.1.12]{\pageref*{thm:10.1.12}}\par
\needspace{3\baselineskip}
\par\noindent\hangindent=2em\hspace*{1em}\phantomsection\label{idx:797}均值（mean）\par
\par\noindent\hangindent=2.8em\hspace*{1.8em}\phantomsection\label{idx:798}影响函数（influence function），\hyperref[ex:10.6.2]{\pageref*{ex:10.6.2}}\par
\needspace{3\baselineskip}
\par\noindent\hangindent=2em\hspace*{1em}\phantomsection\label{idx:799}均方误差（mean squared error）\par
\par\noindent\hangindent=2.8em\hspace*{1.8em}\phantomsection\label{idx:800}伯努利（Bernoulli），\hyperref[ex:7.3.5]{\pageref*{ex:7.3.5}}\par
\par\noindent\hangindent=2.8em\hspace*{1.8em}\phantomsection\label{idx:801}方差（variance），\hyperref[eq:7.3.1]{\pageref*{eq:7.3.1}}\par
\needspace{3\baselineskip}
\par\noindent\hangindent=2em\hspace*{1em}\phantomsection\label{idx:802}中位数（median）\par
\par\noindent\hangindent=2.8em\hspace*{1.8em}\phantomsection\label{idx:803}影响函数（influence function），\hyperref[idxocc:803:450]{\pageref*{idxocc:803:450}}\par
\par\noindent\hangindent=2.8em\hspace*{1.8em}\phantomsection\label{idx:804}极限分布（limit distribution），\hyperref[ex:10.2.3]{\pageref*{ex:10.2.3}}\par
\par\noindent\hangindent=2em\hspace*{1em}\phantomsection\label{idx:805}中位数绝对偏差（median absolute deviation），\hyperref[exercise:10.21]{\pageref*{exercise:10.21}}\par
\par\noindent\hangindent=2em\hspace*{1em}\phantomsection\label{idx:806}矩估计法（method of moments），\hyperref[sec:7.2.1]{\pageref*{sec:7.2.1}}\par
\par\noindent\hangindent=2.8em\hspace*{1.8em}\phantomsection\label{idx:807}二项（binomial），\hyperref[ex:7.2.2]{\pageref*{ex:7.2.2}}\par
\par\noindent\hangindent=2.8em\hspace*{1.8em}\phantomsection\label{idx:808}卡方（chi squared），\hyperref[ex:7.2.3]{\pageref*{ex:7.2.3}}\par
\par\noindent\hangindent=2.8em\hspace*{1.8em}\phantomsection\label{idx:809}正态（normal），\hyperref[ex:7.2.1]{\pageref*{ex:7.2.1}}\par
\par\noindent\hangindent=2.8em\hspace*{1.8em}\phantomsection\label{idx:810}与最大似然估计量的关系（relation to MLE），\hyperref[sec:7.5.1]{\pageref*{sec:7.5.1}}\par
\par\noindent\hangindent=2em\hspace*{1em}\phantomsection\label{idx:811}幂平均数（power means），\hyperref[exercise:7.16]{\pageref*{exercise:7.16}}\par
\par\noindent\hangindent=2em\hspace*{1em}\phantomsection\label{idx:812}比值（ratio），\hyperref[ex:5.5.27]{\pageref*{ex:5.5.27}}\par
\par\noindent\hangindent=2em\hspace*{1em}\phantomsection\label{idx:813}稳健（robust），\hyperref[idxocc:813:420]{\pageref*{idxocc:813:420}}\par
\par\noindent\hangindent=2.8em\hspace*{1.8em}\phantomsection\label{idx:814}Huber估计量（Huber estimator），\hyperref[ex:10.2.5]{\pageref*{ex:10.2.5}}\par
\par\noindent\hangindent=2.8em\hspace*{1.8em}\phantomsection\label{idx:815}极限分布（limit distribution），\hyperref[eq:10.2.6]{\pageref*{eq:10.2.6}}\par
\par\noindent\hangindent=2.8em\hspace*{1.8em}\phantomsection\label{idx:816}M估计量（M-estimator），\hyperref[sec:10.2.2]{\pageref*{sec:10.2.2}}\par
\par\noindent\hangindent=2em\hspace*{1em}\phantomsection\label{idx:817}充分性（sufficiency），\hyperref[thm:7.3.17]{\pageref*{thm:7.3.17}}\par
\par\addvspace{2pt}
\par\noindent\hangindent=1em\hspace*{0em}\phantomsection\label{idx:818}泊松基本假设（Poisson postulates），\hyperref[thm:3.8.1]{\pageref*{thm:3.8.1}}\par
\par\addvspace{2pt}
\par\noindent\hangindent=1em\hspace*{0em}\phantomsection\label{idx:819}后验分布（Posterior distribution），\hyperref[idxocc:819:293]{\pageref*{idxocc:819:293}}\par
\par\addvspace{2pt}
\par\noindent\hangindent=1em\hspace*{0em}\phantomsection\label{idx:820}后验期望损失（另见\hyperref[idx:653]{损失函数}）（Posterior expected loss (see also Loss function)），\hyperref[eq:7.3.19]{\pageref*{eq:7.3.19}}\par
\par\addvspace{2pt}
\par\noindent\hangindent=1em\hspace*{0em}\phantomsection\label{idx:821}幂平均数（Power means），\hyperref[sec:4.9.2]{\pageref*{sec:4.9.2}}, \hyperref[exercise:7.16]{\pageref*{exercise:7.16}}\par
\par\addvspace{2pt}
\par\noindent\hangindent=1em\hspace*{0em}\phantomsection\label{idx:822}Pratt定理（Pratt’s Theorem），\hyperref[thm:9.3.9]{\pageref*{thm:9.3.9}}\par
\par\addvspace{2pt}
\par\noindent\hangindent=1em\hspace*{0em}\phantomsection\label{idx:823}预测区间（Prediction interval），\hyperref[def:11.3.5]{\pageref*{def:11.3.5}}\par
\par\addvspace{2pt}
\par\noindent\hangindent=1em\hspace*{0em}\phantomsection\label{idx:824}先验分布（Prior distribution），\hyperref[idxocc:824:293]{\pageref*{idxocc:824:293}}\par
\needspace{3\baselineskip}
\par\noindent\hangindent=2em\hspace*{1em}\phantomsection\label{idx:825}共轭（conjugate）\par
\par\noindent\hangindent=2.8em\hspace*{1.8em}\phantomsection\label{idx:826}伽马（gamma），\hyperref[exercise:7.23]{\pageref*{exercise:7.23}}\par
\par\noindent\hangindent=2.8em\hspace*{1.8em}\phantomsection\label{idx:827}泊松（Poisson），\hyperref[exercise:7.24]{\pageref*{exercise:7.24}}\par
\par\addvspace{2pt}
\par\noindent\hangindent=1em\hspace*{0em}\phantomsection\label{idx:828}概率密度函数（Probability density function (pdf)），\hyperref[def:1.6.3]{\pageref*{def:1.6.3}}\par
\par\noindent\hangindent=2em\hspace*{1em}\phantomsection\label{idx:829}二元变换（bivariate transformation），\hyperref[eq:4.3.2]{\pageref*{eq:4.3.2}}\par
\par\noindent\hangindent=2em\hspace*{1em}\phantomsection\label{idx:830}条件（conditional），\hyperref[def:4.2.3]{\pageref*{def:4.2.3}}, \hyperref[eq:4.6.6]{\pageref*{eq:4.6.6}}\par
\par\noindent\hangindent=2em\hspace*{1em}\phantomsection\label{idx:831}构造（construction），\hyperref[idxocc:831:61]{\pageref*{idxocc:831:61}}\par
\par\noindent\hangindent=2em\hspace*{1em}\phantomsection\label{idx:832}偶（even），\hyperref[exercise:2.25]{\pageref*{exercise:2.25}}\par
\par\noindent\hangindent=2em\hspace*{1em}\phantomsection\label{idx:833}联合（joint），\hyperref[def:4.1.10]{\pageref*{def:4.1.10}}, \hyperref[eq:4.6.2]{\pageref*{eq:4.6.2}}\par
\par\noindent\hangindent=2em\hspace*{1em}\phantomsection\label{idx:834}边缘（marginal），\hyperref[eq:4.1.2]{\pageref*{eq:4.1.2}}, \hyperref[idxocc:834:176]{\pageref*{idxocc:834:176}}\par
\par\noindent\hangindent=2em\hspace*{1em}\phantomsection\label{idx:835}单调变换（monotone transformation），\hyperref[thm:2.1.5]{\pageref*{thm:2.1.5}}\par
\par\noindent\hangindent=2em\hspace*{1em}\phantomsection\label{idx:836}多元变换（multivariate transformation），\hyperref[eq:4.6.7]{\pageref*{eq:4.6.7}}\par
\par\noindent\hangindent=2em\hspace*{1em}\phantomsection\label{idx:837}分段单调变换（piecewise monotone transformation），\hyperref[thm:2.1.8]{\pageref*{thm:2.1.8}}\par
\par\noindent\hangindent=2em\hspace*{1em}\phantomsection\label{idx:838}性质（properties），\hyperref[thm:1.6.5]{\pageref*{thm:1.6.5}}\par
\par\noindent\hangindent=2em\hspace*{1em}\phantomsection\label{idx:839}样本（sample），\hyperref[eq:5.1.1]{\pageref*{eq:5.1.1}}\par
\par\noindent\hangindent=2em\hspace*{1em}\phantomsection\label{idx:840}对称（symmetric），\hyperref[exercise:2.26]{\pageref*{exercise:2.26}}\par
\par\noindent\hangindent=2em\hspace*{1em}\phantomsection\label{idx:841}截断（truncated），\hyperref[exercise:1.52]{\pageref*{exercise:1.52}}\par
\par\noindent\hangindent=2em\hspace*{1em}\phantomsection\label{idx:842}单峰（unimodal），\hyperref[exercise:2.27]{\pageref*{exercise:2.27}}\par
\par\addvspace{2pt}
\par\noindent\hangindent=1em\hspace*{0em}\phantomsection\label{idx:843}概率函数（Probability function），\hyperref[def:1.2.4]{\pageref*{def:1.2.4}}\par
\par\noindent\hangindent=2em\hspace*{1em}\phantomsection\label{idx:844}有限集（finite set），\hyperref[thm:1.2.6]{\pageref*{thm:1.2.6}}\par
\par\noindent\hangindent=2em\hspace*{1em}\phantomsection\label{idx:845}诱导（induced），\hyperref[eq:1.4.1]{\pageref*{eq:1.4.1}}\par
\par\noindent\hangindent=2em\hspace*{1em}\phantomsection\label{idx:846}性质（properties），\hyperref[thm:1.2.8]{\pageref*{thm:1.2.8}}\par
\par\addvspace{2pt}
\par\noindent\hangindent=1em\hspace*{0em}\phantomsection\label{idx:847}概率积分变换（Probability integral transformation），\hyperref[thm:2.1.10]{\pageref*{thm:2.1.10}}, \hyperref[ex:5.6.3]{\pageref*{ex:5.6.3}}\par
\par\noindent\hangindent=2em\hspace*{1em}\phantomsection\label{idx:848}置信区间（confidence interval），\hyperref[thm:9.2.12]{\pageref*{thm:9.2.12}}\par
\par\noindent\hangindent=2em\hspace*{1em}\phantomsection\label{idx:849}离散（discrete），\hyperref[exercise:2.10]{\pageref*{exercise:2.10}}, \hyperref[thm:9.2.14]{\pageref*{thm:9.2.14}}\par
\par\addvspace{2pt}
\par\noindent\hangindent=1em\hspace*{0em}\phantomsection\label{idx:850}概率质量函数（Probability mass function (pmf)），\hyperref[def:1.6.1]{\pageref*{def:1.6.1}}\par
\par\noindent\hangindent=2em\hspace*{1em}\phantomsection\label{idx:851}条件（conditional），\hyperref[def:4.2.1]{\pageref*{def:4.2.1}}, \hyperref[eq:4.6.6]{\pageref*{eq:4.6.6}}\par
\par\noindent\hangindent=2em\hspace*{1em}\phantomsection\label{idx:852}构造（construction），\hyperref[idxocc:852:61]{\pageref*{idxocc:852:61}}\par
\par\noindent\hangindent=2em\hspace*{1em}\phantomsection\label{idx:853}联合（joint），\hyperref[eq:4.6.1]{\pageref*{eq:4.6.1}}\par
\par\noindent\hangindent=2.8em\hspace*{1.8em}\phantomsection\label{idx:854}二元（bivariate），\hyperref[def:4.1.3]{\pageref*{def:4.1.3}}\par
\par\noindent\hangindent=2em\hspace*{1em}\phantomsection\label{idx:855}边缘（marginal），\hyperref[thm:4.1.6]{\pageref*{thm:4.1.6}}, \hyperref[idxocc:855:176]{\pageref*{idxocc:855:176}}\par
\par\noindent\hangindent=2em\hspace*{1em}\phantomsection\label{idx:856}性质（properties），\hyperref[thm:1.6.5]{\pageref*{thm:1.6.5}}\par
\par\noindent\hangindent=2em\hspace*{1em}\phantomsection\label{idx:857}样本（sample），\hyperref[eq:5.1.1]{\pageref*{eq:5.1.1}}\par
\par\addvspace{2pt}
\par\noindent\hangindent=1em\hspace*{0em}\phantomsection\label{idx:858}错误覆盖概率（Probability of false coverage），\hyperref[eq:9.3.2]{\pageref*{eq:9.3.2}}\par
\par\addvspace{2pt}
\par\noindent\hangindent=1em\hspace*{0em}\phantomsection\label{idx:859}尼罗河问题（Problem of the Nile），\hyperref[exercise:7.54]{\pageref*{exercise:7.54}}\par
\par\addvspace{2pt}
\par\noindent\hangindent=1em\hspace*{0em}\phantomsection\label{idx:860}四分位数（Quartile），\hyperref[idxocc:860:216]{\pageref*{idxocc:860:216}}\par
\par\addvspace{2pt}
\par\noindent\hangindent=1em\hspace*{0em}\phantomsection\label{idx:861}随机数生成（Random number generation），\hyperref[idxocc:861:76]{\pageref*{idxocc:861:76}}, \hyperref[exercise:5.50]{\pageref*{exercise:5.50}}\par
\par\addvspace{2pt}
\needspace{3\baselineskip}
\par\noindent\hangindent=1em\hspace*{0em}\phantomsection\label{idx:862}随机样本（Random sample）\par
\par\noindent\hangindent=2em\hspace*{1em}\phantomsection\label{idx:863}生成（generating），\hyperref[sec:5.6]{\pageref*{sec:5.6}}\par
\par\addvspace{2pt}
\par\noindent\hangindent=1em\hspace*{0em}\phantomsection\label{idx:864}随机变量（Random variable），\hyperref[def:1.4.1]{\pageref*{def:1.4.1}}\par
\par\noindent\hangindent=2em\hspace*{1em}\phantomsection\label{idx:865}绝对连续（absolutely continuous），\hyperref[idxocc:865:62]{\pageref*{idxocc:865:62}}\par
\par\noindent\hangindent=2em\hspace*{1em}\phantomsection\label{idx:866}删失（censored），\hyperref[exercise:4.26]{\pageref*{exercise:4.26}}\par
\par\noindent\hangindent=2em\hspace*{1em}\phantomsection\label{idx:867}连续（continuous），\hyperref[def:1.5.7]{\pageref*{def:1.5.7}}\par
\par\noindent\hangindent=2em\hspace*{1em}\phantomsection\label{idx:868}离散（discrete），\hyperref[def:1.5.7]{\pageref*{def:1.5.7}}\par
\par\noindent\hangindent=2em\hspace*{1em}\phantomsection\label{idx:869}同分布（identically distributed），\hyperref[def:1.5.8]{\pageref*{def:1.5.8}}\par
\par\noindent\hangindent=2.8em\hspace*{1.8em}\phantomsection\label{idx:870}矩（moments），\hyperref[thm:2.3.11]{\pageref*{thm:2.3.11}}\par
\par\addvspace{2pt}
\par\noindent\hangindent=1em\hspace*{0em}\phantomsection\label{idx:871}随机向量（Random vector），\hyperref[def:4.1.1]{\pageref*{def:4.1.1}}\par
\par\noindent\hangindent=2em\hspace*{1em}\phantomsection\label{idx:872}连续（continuous），\hyperref[def:4.1.10]{\pageref*{def:4.1.10}}\par
\par\noindent\hangindent=2em\hspace*{1em}\phantomsection\label{idx:873}离散（discrete），\hyperref[idxocc:873:146]{\pageref*{idxocc:873:146}}\par
\par\addvspace{2pt}
\needspace{3\baselineskip}
\par\noindent\hangindent=1em\hspace*{0em}\phantomsection\label{idx:874}随机完全区组设计（Randomized complete block design (RCB)）\par
\par\noindent\hangindent=2em\hspace*{1em}\phantomsection\label{idx:875}与变量含误差的关系（relation to EIV），\hyperref[idxocc:875:506]{\pageref*{idxocc:875:506}}\par
\par\addvspace{2pt}
\par\noindent\hangindent=1em\hspace*{0em}\phantomsection\label{idx:876}随机停止的和（Randomly stopped sum），\hyperref[exercise:4.33]{\pageref*{exercise:4.33}}\par
\par\addvspace{2pt}
\par\noindent\hangindent=1em\hspace*{0em}\phantomsection\label{idx:877}Rao–Blackwell定理（Rao–Blackwell Theorem），\hyperref[thm:7.3.17]{\pageref*{thm:7.3.17}}\par
\par\addvspace{2pt}
\par\noindent\hangindent=1em\hspace*{0em}\phantomsection\label{idx:878}回归（Regression），\hyperref[idxocc:878:452]{\pageref*{idxocc:878:452}}, \hyperref[sec:11.3]{\pageref*{sec:11.3}}\par
\par\noindent\hangindent=2em\hspace*{1em}\phantomsection\label{idx:879}方差分析表（ANOVA table），\hyperref[tab:11.3.2]{\pageref*{tab:11.3.2}}\par
\par\noindent\hangindent=2em\hspace*{1em}\phantomsection\label{idx:880}二元正态（bivariate normal），\hyperref[sec:11.3.3.2]{\pageref*{sec:11.3.3.2}}\par
\par\noindent\hangindent=2em\hspace*{1em}\phantomsection\label{idx:881}截距的最优线性无偏估计量（BLUE of intercept），\hyperref[eq:11.3.21]{\pageref*{eq:11.3.21}}\par
\par\noindent\hangindent=2em\hspace*{1em}\phantomsection\label{idx:882}斜率的最优线性无偏估计量（BLUE of slope），\hyperref[eq:11.3.19]{\pageref*{eq:11.3.19}}\par
\par\noindent\hangindent=2em\hspace*{1em}\phantomsection\label{idx:883}决定系数（coefficient of determination），\hyperref[idxocc:883:481]{\pageref*{idxocc:883:481}}\par
\par\noindent\hangindent=2em\hspace*{1em}\phantomsection\label{idx:884}条件期望（conditional expectation），\hyperref[exercise:4.13]{\pageref*{exercise:4.13}}, \hyperref[eq:11.3.3]{\pageref*{eq:11.3.3}}\par
\par\noindent\hangindent=2em\hspace*{1em}\phantomsection\label{idx:885}条件正态（conditional normal），\hyperref[sec:11.3.3.1]{\pageref*{sec:11.3.3.1}}\par
\par\noindent\hangindent=2em\hspace*{1em}\phantomsection\label{idx:886}置信带（confidence band），\hyperref[idxocc:886:485]{\pageref*{idxocc:886:485}}\par
\par\noindent\hangindent=2em\hspace*{1em}\phantomsection\label{idx:887}置信区间（confidence interval），\hyperref[eq:11.3.37]{\pageref*{eq:11.3.37}}\par
\par\noindent\hangindent=2em\hspace*{1em}\phantomsection\label{idx:888}数据拟合（data fitting），\hyperref[idxocc:888:468]{\pageref*{idxocc:888:468}}, \hyperref[idxocc:888:471]{\pageref*{idxocc:888:471}}\par
\par\noindent\hangindent=2em\hspace*{1em}\phantomsection\label{idx:889}因变量（dependent variable），\hyperref[idxocc:889:466]{\pageref*{idxocc:889:466}}\par
\sourcepageanchor{565}
\par\noindent\hangindent=2em\hspace*{1em}\phantomsection\label{idx:890}设计（design），\hyperref[idxocc:890:473]{\pageref*{idxocc:890:473}}\par
\par\noindent\hangindent=2em\hspace*{1em}\phantomsection\label{idx:891}外推（extrapolation），\hyperref[idxocc:891:486]{\pageref*{idxocc:891:486}}\par
\par\noindent\hangindent=2em\hspace*{1em}\phantomsection\label{idx:892}自变量（independent variable），\hyperref[idxocc:892:466]{\pageref*{idxocc:892:466}}\par
\par\noindent\hangindent=2em\hspace*{1em}\phantomsection\label{idx:893}推断（inference），\hyperref[idxocc:893:468]{\pageref*{idxocc:893:468}}\par
\par\noindent\hangindent=2em\hspace*{1em}\phantomsection\label{idx:894}最小绝对偏差（least absolute deviation），\hyperref[idxocc:894:514]{\pageref*{idxocc:894:514}}, \hyperref[exercise:12.25]{\pageref*{exercise:12.25}}\par
\par\noindent\hangindent=2em\hspace*{1em}\phantomsection\label{idx:895}线性（linear），\hyperref[idxocc:895:467]{\pageref*{idxocc:895:467}}, \hyperref[sec:11.3.3]{\pageref*{sec:11.3.3}}\par
\par\noindent\hangindent=2em\hspace*{1em}\phantomsection\label{idx:896}关于x线性（linear in x），\hyperref[idxocc:896:467]{\pageref*{idxocc:896:467}}\par
\par\noindent\hangindent=2em\hspace*{1em}\phantomsection\label{idx:897}最大概率估计量（maximum probability estimator），\hyperref[item:11.37.e]{\pageref*{item:11.37.e}}\par
\par\noindent\hangindent=2em\hspace*{1em}\phantomsection\label{idx:898}最大似然估计量（MLE），\hyperref[sec:11.3.4]{\pageref*{sec:11.3.4}}\par
\par\noindent\hangindent=2.8em\hspace*{1.8em}\phantomsection\label{idx:899}有偏（biased），\hyperref[idxocc:899:477]{\pageref*{idxocc:899:477}}\par
\par\noindent\hangindent=2.8em\hspace*{1.8em}\phantomsection\label{idx:900}抽样分布（sampling distribution），\hyperref[thm:11.3.3]{\pageref*{thm:11.3.3}}\par
\par\noindent\hangindent=2em\hspace*{1em}\phantomsection\label{idx:901}总体（population），\hyperref[idxocc:901:452]{\pageref*{idxocc:901:452}}, \hyperref[eq:11.3.2]{\pageref*{eq:11.3.2}}\par
\par\noindent\hangindent=2em\hspace*{1em}\phantomsection\label{idx:902}预测区间（prediction interval），\hyperref[eq:11.3.41]{\pageref*{eq:11.3.41}}\par
\par\noindent\hangindent=2em\hspace*{1em}\phantomsection\label{idx:903}预测变量（predictor variable），\hyperref[idxocc:903:467]{\pageref*{idxocc:903:467}}\par
\par\noindent\hangindent=2em\hspace*{1em}\phantomsection\label{idx:904}$r^2$，\hyperref[idxocc:904:481]{\pageref*{idxocc:904:481}}\par
\par\noindent\hangindent=2em\hspace*{1em}\phantomsection\label{idx:905}残差（residuals），\hyperref[eq:11.3.27]{\pageref*{eq:11.3.27}}\par
\par\noindent\hangindent=2em\hspace*{1em}\phantomsection\label{idx:906}响应变量（response variable），\hyperref[idxocc:906:467]{\pageref*{idxocc:906:467}}\par
\par\noindent\hangindent=2em\hspace*{1em}\phantomsection\label{idx:907}平方和（sum of squares），\hyperref[eq:11.3.6]{\pageref*{eq:11.3.6}}, \hyperref[eq:11.3.36]{\pageref*{eq:11.3.36}}\par
\par\noindent\hangindent=2em\hspace*{1em}\phantomsection\label{idx:908}t统计量（t statistic），\hyperref[eq:11.3.33]{\pageref*{eq:11.3.33}}\par
\par\noindent\hangindent=2em\hspace*{1em}\phantomsection\label{idx:909}t检验（t test），\hyperref[eq:11.3.35]{\pageref*{eq:11.3.35}}, \hyperref[eq:11.3.38]{\pageref*{eq:11.3.38}}\par
\par\noindent\hangindent=2em\hspace*{1em}\phantomsection\label{idx:910}向均值回归（toward the mean），\hyperref[idxocc:910:467]{\pageref*{idxocc:910:467}}\par
\par\noindent\hangindent=2em\hspace*{1em}\phantomsection\label{idx:911}零截距（zero intercept），\hyperref[exercise:7.19]{\pageref*{exercise:7.19}}\par
\par\addvspace{2pt}
\par\noindent\hangindent=1em\hspace*{0em}\phantomsection\label{idx:912}正则条件（Regularity conditions），\hyperref[sec:10.6.2]{\pageref*{sec:10.6.2}}\par
\par\addvspace{2pt}
\par\noindent\hangindent=1em\hspace*{0em}\phantomsection\label{idx:913}重抽样（Resampling），\hyperref[ex:10.1.19]{\pageref*{ex:10.1.19}}, \hyperref[eq:10.1.11]{\pageref*{eq:10.1.11}}\par
\par\addvspace{2pt}
\par\noindent\hangindent=1em\hspace*{0em}\phantomsection\label{idx:914}风险函数（Risk function），\hyperref[eq:7.3.16]{\pageref*{eq:7.3.16}}\par
\par\addvspace{2pt}
\needspace{3\baselineskip}
\par\noindent\hangindent=1em\hspace*{0em}\phantomsection\label{idx:915}稳健（Robust）\par
\par\noindent\hangindent=2em\hspace*{1em}\phantomsection\label{idx:916}Tukey模型（Tukey model），\hyperref[exercise:10.23]{\pageref*{exercise:10.23}}\par
\par\addvspace{2pt}
\par\noindent\hangindent=1em\hspace*{0em}\phantomsection\label{idx:917}样本（Sample），\hyperref[def:5.1.1]{\pageref*{def:5.1.1}}\par
\par\noindent\hangindent=2em\hspace*{1em}\phantomsection\label{idx:918}指数族（exponential family），\hyperref[thm:5.2.11]{\pageref*{thm:5.2.11}}\par
\par\noindent\hangindent=2em\hspace*{1em}\phantomsection\label{idx:919}有限总体（finite population），\hyperref[idxocc:919:201]{\pageref*{idxocc:919:201}}\par
\par\noindent\hangindent=2em\hspace*{1em}\phantomsection\label{idx:920}无限总体（infinite population），\hyperref[idxocc:920:201]{\pageref*{idxocc:920:201}}\par
\par\noindent\hangindent=2em\hspace*{1em}\phantomsection\label{idx:921}均值（mean），\hyperref[def:5.2.2]{\pageref*{def:5.2.2}}\par
\par\noindent\hangindent=2.8em\hspace*{1.8em}\phantomsection\label{idx:922}柯西（Cauchy），\hyperref[ex:5.2.10]{\pageref*{ex:5.2.10}}\par
\par\noindent\hangindent=2.8em\hspace*{1.8em}\phantomsection\label{idx:923}期望值（expected value），\hyperref[thm:5.2.6]{\pageref*{thm:5.2.6}}\par
\par\noindent\hangindent=2.8em\hspace*{1.8em}\phantomsection\label{idx:924}位置尺度（location–scale），\hyperref[idxocc:924:208]{\pageref*{idxocc:924:208}}\par
\par\noindent\hangindent=2.8em\hspace*{1.8em}\phantomsection\label{idx:925}矩母函数（mgf），\hyperref[thm:5.2.7]{\pageref*{thm:5.2.7}}\par
\par\noindent\hangindent=2.8em\hspace*{1.8em}\phantomsection\label{idx:926}正态（normal），\hyperref[ex:5.2.8]{\pageref*{ex:5.2.8}}, \hyperref[thm:5.3.1]{\pageref*{thm:5.3.1}}\par
\par\noindent\hangindent=2.8em\hspace*{1.8em}\phantomsection\label{idx:927}概率密度函数（pdf），\hyperref[exercise:5.5]{\pageref*{exercise:5.5}}\par
\par\noindent\hangindent=2.8em\hspace*{1.8em}\phantomsection\label{idx:928}递推关系（recursion relation），\hyperref[exercise:5.15]{\pageref*{exercise:5.15}}\par
\par\noindent\hangindent=2.8em\hspace*{1.8em}\phantomsection\label{idx:929}截尾（trimmed），\hyperref[idxocc:929:442]{\pageref*{idxocc:929:442}}\par
\par\noindent\hangindent=2em\hspace*{1em}\phantomsection\label{idx:930}中位数（median），\hyperref[eq:5.4.1]{\pageref*{eq:5.4.1}}\par
\par\noindent\hangindent=2em\hspace*{1em}\phantomsection\label{idx:931}中程数（midrange），\hyperref[ex:5.4.7]{\pageref*{ex:5.4.7}}\par
\par\noindent\hangindent=2em\hspace*{1em}\phantomsection\label{idx:932}概率密度函数与概率质量函数（pdf, pmf），\hyperref[eq:5.1.1]{\pageref*{eq:5.1.1}}\par
\par\noindent\hangindent=2em\hspace*{1em}\phantomsection\label{idx:933}极差（range），\hyperref[idxocc:933:215]{\pageref*{idxocc:933:215}}\par
\par\noindent\hangindent=2em\hspace*{1em}\phantomsection\label{idx:934}简单随机（simple random），\hyperref[idxocc:934:202]{\pageref*{idxocc:934:202}}\par
\par\noindent\hangindent=2em\hspace*{1em}\phantomsection\label{idx:935}标准差（standard deviation），\hyperref[def:5.2.3]{\pageref*{def:5.2.3}}\par
\par\noindent\hangindent=2em\hspace*{1em}\phantomsection\label{idx:936}方差（variance），\hyperref[def:5.2.3]{\pageref*{def:5.2.3}}\par
\par\noindent\hangindent=2.8em\hspace*{1.8em}\phantomsection\label{idx:937}偏差（bias），\hyperref[sec:5.8.2]{\pageref*{sec:5.8.2}}\par
\par\noindent\hangindent=2.8em\hspace*{1.8em}\phantomsection\label{idx:938}期望值（expected value），\hyperref[thm:5.2.6]{\pageref*{thm:5.2.6}}\par
\par\noindent\hangindent=2.8em\hspace*{1.8em}\phantomsection\label{idx:939}恒等式（identity），\hyperref[exercise:5.8]{\pageref*{exercise:5.8}}\par
\par\noindent\hangindent=2.8em\hspace*{1.8em}\phantomsection\label{idx:940}正态（normal），\hyperref[thm:5.3.1]{\pageref*{thm:5.3.1}}\par
\par\noindent\hangindent=2.8em\hspace*{1.8em}\phantomsection\label{idx:941}递推关系（recursion relation），\hyperref[exercise:5.15]{\pageref*{exercise:5.15}}\par
\par\noindent\hangindent=2em\hspace*{1em}\phantomsection\label{idx:942}有放回（with replacement），\hyperref[idxocc:942:202]{\pageref*{idxocc:942:202}}\par
\par\noindent\hangindent=2em\hspace*{1em}\phantomsection\label{idx:943}无放回（without replacement），\hyperref[idxocc:943:202]{\pageref*{idxocc:943:202}}\par
\par\addvspace{2pt}
\needspace{3\baselineskip}
\par\noindent\hangindent=1em\hspace*{0em}\phantomsection\label{idx:944}样本极差（Sample range）\par
\par\noindent\hangindent=2em\hspace*{1em}\phantomsection\label{idx:945}均匀（uniform），\hyperref[ex:5.4.7]{\pageref*{ex:5.4.7}}\par
\par\addvspace{2pt}
\needspace{3\baselineskip}
\par\noindent\hangindent=1em\hspace*{0em}\phantomsection\label{idx:946}样本量（Sample size）\par
\par\noindent\hangindent=2em\hspace*{1em}\phantomsection\label{idx:947}假设检验：正态（hypothesis test – normal），\hyperref[ex:8.3.4]{\pageref*{ex:8.3.4}}\par
\par\addvspace{2pt}
\par\noindent\hangindent=1em\hspace*{0em}\phantomsection\label{idx:948}样本空间（Sample space），\hyperref[def:1.1.1]{\pageref*{def:1.1.1}}\par
\par\addvspace{2pt}
\needspace{3\baselineskip}
\par\noindent\hangindent=1em\hspace*{0em}\phantomsection\label{idx:949}抽样（Sampling）\par
\par\noindent\hangindent=2em\hspace*{1em}\phantomsection\label{idx:950}有放回（with replacement），\hyperref[ex:10.1.19]{\pageref*{ex:10.1.19}}\par
\par\addvspace{2pt}
\par\noindent\hangindent=1em\hspace*{0em}\phantomsection\label{idx:951}抽样分布（Sampling distribution），\hyperref[def:5.2.1]{\pageref*{def:5.2.1}}\par
\par\addvspace{2pt}
\par\noindent\hangindent=1em\hspace*{0em}\phantomsection\label{idx:952}Satterthwaite近似（Satterthwaite approximation），\hyperref[ex:7.2.3]{\pageref*{ex:7.2.3}}, \hyperref[exercise:8.42]{\pageref*{exercise:8.42}}\par
\par\addvspace{2pt}
\needspace{3\baselineskip}
\par\noindent\hangindent=1em\hspace*{0em}\phantomsection\label{idx:953}尺度族（Scale family）\par
\par\noindent\hangindent=2em\hspace*{1em}\phantomsection\label{idx:954}枢轴量（pivotal quantity），\hyperref[ex:9.2.7]{\pageref*{ex:9.2.7}}\par
\par\addvspace{2pt}
\par\noindent\hangindent=1em\hspace*{0em}\phantomsection\label{idx:955}Scheffé方法（Scheffé’s method），\hyperref[thm:11.2.10]{\pageref*{thm:11.2.10}}\par
\par\noindent\hangindent=2em\hspace*{1em}\phantomsection\label{idx:956}方差分析（ANOVA），\hyperref[exercise:11.14]{\pageref*{exercise:11.14}}\par
\par\addvspace{2pt}
\needspace{3\baselineskip}
\par\noindent\hangindent=1em\hspace*{0em}\phantomsection\label{idx:957}集合估计量（Set estimator）\par
\par\noindent\hangindent=2em\hspace*{1em}\phantomsection\label{idx:958}与置信区间的关系（relation to confidence interval），\hyperref[idxocc:958:369]{\pageref*{idxocc:958:369}}\par
\par\addvspace{2pt}
\needspace{3\baselineskip}
\par\noindent\hangindent=1em\hspace*{0em}\phantomsection\label{idx:959}集合运算（Set operations）\par
\par\noindent\hangindent=2em\hspace*{1em}\phantomsection\label{idx:960}结合律（associativity），\hyperref[thm:1.1.4]{\pageref*{thm:1.1.4}}\par
\par\noindent\hangindent=2em\hspace*{1em}\phantomsection\label{idx:961}交换律（commutativity），\hyperref[thm:1.1.4]{\pageref*{thm:1.1.4}}\par
\par\noindent\hangindent=2em\hspace*{1em}\phantomsection\label{idx:962}求补（complementation），\hyperref[idxocc:962:33]{\pageref*{idxocc:962:33}}\par
\par\noindent\hangindent=2em\hspace*{1em}\phantomsection\label{idx:963}德摩根律（DeMorgan’s Laws），\hyperref[thm:1.1.4]{\pageref*{thm:1.1.4}}\par
\par\noindent\hangindent=2em\hspace*{1em}\phantomsection\label{idx:964}分配律（Distributive Law），\hyperref[thm:1.1.4]{\pageref*{thm:1.1.4}}\par
\par\noindent\hangindent=2em\hspace*{1em}\phantomsection\label{idx:965}恒等式（identities），\hyperref[exercise:1.2]{\pageref*{exercise:1.2}}\par
\par\noindent\hangindent=2em\hspace*{1em}\phantomsection\label{idx:966}交（intersection），\hyperref[idxocc:966:33]{\pageref*{idxocc:966:33}}\par
\par\noindent\hangindent=2.8em\hspace*{1.8em}\phantomsection\label{idx:967}可数（countable），\hyperref[idxocc:967:34]{\pageref*{idxocc:967:34}}\par
\par\noindent\hangindent=2.8em\hspace*{1.8em}\phantomsection\label{idx:968}不可数（uncountable），\hyperref[idxocc:968:35]{\pageref*{idxocc:968:35}}\par
\par\noindent\hangindent=2em\hspace*{1em}\phantomsection\label{idx:969}并（union），\hyperref[idxocc:969:33]{\pageref*{idxocc:969:33}}\par
\par\noindent\hangindent=2.8em\hspace*{1.8em}\phantomsection\label{idx:970}可数（countable），\hyperref[idxocc:970:34]{\pageref*{idxocc:970:34}}\par
\par\noindent\hangindent=2.8em\hspace*{1.8em}\phantomsection\label{idx:971}不可数（uncountable），\hyperref[idxocc:971:35]{\pageref*{idxocc:971:35}}\par
\par\addvspace{2pt}
\par\noindent\hangindent=1em\hspace*{0em}\phantomsection\label{idx:972}\ensuremath{\sigma}代数（Sigma algebra），\hyperref[def:1.2.1]{\pageref*{def:1.2.1}}\par
\par\noindent\hangindent=2em\hspace*{1em}\phantomsection\label{idx:973}可数（countable），\hyperref[ex:1.2.2]{\pageref*{ex:1.2.2}}\par
\par\noindent\hangindent=2em\hspace*{1em}\phantomsection\label{idx:974}生成的（generated），\hyperref[idxocc:974:59]{\pageref*{idxocc:974:59}}\par
\par\noindent\hangindent=2em\hspace*{1em}\phantomsection\label{idx:975}性质（properties），\hyperref[exercise:1.11]{\pageref*{exercise:1.11}}\par
\par\noindent\hangindent=2em\hspace*{1em}\phantomsection\label{idx:976}平凡（trivial），\hyperref[idxocc:976:36]{\pageref*{idxocc:976:36}}\par
\par\noindent\hangindent=2em\hspace*{1em}\phantomsection\label{idx:977}不可数（uncountable），\hyperref[ex:1.2.3]{\pageref*{ex:1.2.3}}\par
\par\addvspace{2pt}
\par\noindent\hangindent=1em\hspace*{0em}\phantomsection\label{idx:978}偏度（Skewness），\hyperref[exercise:2.28]{\pageref*{exercise:2.28}}\par
\par\addvspace{2pt}
\par\noindent\hangindent=1em\hspace*{0em}\phantomsection\label{idx:979}Slutsky定理（Slutsky’s Theorem），\hyperref[thm:5.5.17]{\pageref*{thm:5.5.17}}, \hyperref[idxocc:979:429]{\pageref*{idxocc:979:429}}\par
\par\addvspace{2pt}
\par\noindent\hangindent=1em\hspace*{0em}\phantomsection\label{idx:980}标准差（Standard deviation），\hyperref[def:2.3.2]{\pageref*{def:2.3.2}}\par
\par\addvspace{2pt}
\par\noindent\hangindent=1em\hspace*{0em}\phantomsection\label{idx:981}统计量（Statistic），\hyperref[def:5.2.1]{\pageref*{def:5.2.1}}\par
\par\addvspace{2pt}
\needspace{3\baselineskip}
\par\noindent\hangindent=1em\hspace*{0em}\phantomsection\label{idx:982}Stein估计（Stein estimation）\par
\par\noindent\hangindent=2em\hspace*{1em}\phantomsection\label{idx:983}方差分析（ANOVA），\hyperref[sec:11.5.6]{\pageref*{sec:11.5.6}}\par
\par\addvspace{2pt}
\par\noindent\hangindent=1em\hspace*{0em}\phantomsection\label{idx:984}Stein引理（Stein’s Lemma），\hyperref[lem:3.6.5]{\pageref*{lem:3.6.5}}\par
\par\addvspace{2pt}
\par\noindent\hangindent=1em\hspace*{0em}\phantomsection\label{idx:985}Stirling公式（Stirling’s formula），\hyperref[exercise:1.28]{\pageref*{exercise:1.28}}, \hyperref[exercise:5.35]{\pageref*{exercise:5.35}}\par
\par\addvspace{2pt}
\needspace{3\baselineskip}
\par\noindent\hangindent=1em\hspace*{0em}\phantomsection\label{idx:986}随机序（Stochastic order）\par
\par\noindent\hangindent=2em\hspace*{1em}\phantomsection\label{idx:987}递减（decreasing），\hyperref[exercise:3.43]{\pageref*{exercise:3.43}}\par
\par\noindent\hangindent=2em\hspace*{1em}\phantomsection\label{idx:988}较大（greater），\hyperref[exercise:1.49]{\pageref*{exercise:1.49}}, \hyperref[exercise:2.10]{\pageref*{exercise:2.10}}\par
\par\noindent\hangindent=2em\hspace*{1em}\phantomsection\label{idx:989}递增（increasing），\hyperref[exercise:3.41]{\pageref*{exercise:3.41}}, \hyperref[idxocc:989:380]{\pageref*{idxocc:989:380}}\par
\par\noindent\hangindent=2.8em\hspace*{1.8em}\phantomsection\label{idx:990}卡方（chi squared），\hyperref[exercise:5.19]{\pageref*{exercise:5.19}}\par
\par\noindent\hangindent=2.8em\hspace*{1.8em}\phantomsection\label{idx:991}F，\hyperref[exercise:5.19]{\pageref*{exercise:5.19}}\par
\par\noindent\hangindent=2em\hspace*{1em}\phantomsection\label{idx:992}位置尺度族（location–scale family），\hyperref[exercise:8.34]{\pageref*{exercise:8.34}}\par
\sourcepageanchor{566}
\par\noindent\hangindent=2em\hspace*{1em}\phantomsection\label{idx:993}单调似然比（MLR），\hyperref[exercise:8.26]{\pageref*{exercise:8.26}}\par
\par\noindent\hangindent=2em\hspace*{1em}\phantomsection\label{idx:994}功效（power），\hyperref[sec:8.5.1]{\pageref*{sec:8.5.1}}\par
\par\addvspace{2pt}
\par\noindent\hangindent=1em\hspace*{0em}\phantomsection\label{idx:995}强大数定律（Strong Law of Large Numbers），\hyperref[sec:5.8.4]{\pageref*{sec:5.8.4}}\par
\par\addvspace{2pt}
\par\noindent\hangindent=1em\hspace*{0em}\phantomsection\label{idx:996}结构推断（Structural inference），\hyperref[idxocc:996:376]{\pageref*{idxocc:996:376}}\par
\par\addvspace{2pt}
\par\noindent\hangindent=1em\hspace*{0em}\phantomsection\label{idx:997}充分性原则（Sufficiency Principle），\hyperref[idxocc:997:252]{\pageref*{idxocc:997:252}}\par
\par\noindent\hangindent=2em\hspace*{1em}\phantomsection\label{idx:998}形式化（formal），\hyperref[idxocc:998:269]{\pageref*{idxocc:998:269}}\par
\par\noindent\hangindent=2em\hspace*{1em}\phantomsection\label{idx:999}解释（interpretations），\hyperref[ex:6.3.7]{\pageref*{ex:6.3.7}}\par
\par\addvspace{2pt}
\par\noindent\hangindent=1em\hspace*{0em}\phantomsection\label{idx:1000}充分统计量（Sufficient statistic），\hyperref[def:6.2.1]{\pageref*{def:6.2.1}}\par
\par\noindent\hangindent=2em\hspace*{1em}\phantomsection\label{idx:1001}方差分析（ANOVA），\hyperref[item:11.6.a]{\pageref*{item:11.6.a}}\par
\par\noindent\hangindent=2em\hspace*{1em}\phantomsection\label{idx:1002}伯努利（Bernoulli），\hyperref[ex:6.2.3]{\pageref*{ex:6.2.3}}\par
\par\noindent\hangindent=2em\hspace*{1em}\phantomsection\label{idx:1003}离散均匀（discrete uniform），\hyperref[ex:6.2.8]{\pageref*{ex:6.2.8}}\par
\par\noindent\hangindent=2em\hspace*{1em}\phantomsection\label{idx:1004}指数族（exponential family），\hyperref[thm:6.2.10]{\pageref*{thm:6.2.10}}\par
\par\noindent\hangindent=2em\hspace*{1em}\phantomsection\label{idx:1005}因子分解定理（Factorization Theorem），\hyperref[thm:6.2.6]{\pageref*{thm:6.2.6}}\par
\par\noindent\hangindent=2em\hspace*{1em}\phantomsection\label{idx:1006}最小（minimal），\hyperref[def:6.2.11]{\pageref*{def:6.2.11}}, \hyperref[exercise:6.26]{\pageref*{exercise:6.26}}\par
\par\noindent\hangindent=2.8em\hspace*{1.8em}\phantomsection\label{idx:1007}刻画（characterization），\hyperref[thm:6.2.13]{\pageref*{thm:6.2.13}}\par
\par\noindent\hangindent=2.8em\hspace*{1.8em}\phantomsection\label{idx:1008}必要（necessary），\hyperref[thm:6.6.4]{\pageref*{thm:6.6.4}}\par
\par\noindent\hangindent=2.8em\hspace*{1.8em}\phantomsection\label{idx:1009}正态（normal），\hyperref[ex:6.2.14]{\pageref*{ex:6.2.14}}\par
\par\noindent\hangindent=2.8em\hspace*{1.8em}\phantomsection\label{idx:1010}均匀（uniform），\hyperref[ex:6.2.15]{\pageref*{ex:6.2.15}}, \hyperref[exercise:6.10]{\pageref*{exercise:6.10}}\par
\par\noindent\hangindent=2em\hspace*{1em}\phantomsection\label{idx:1011}正态（normal），\hyperref[ex:6.2.4]{\pageref*{ex:6.2.4}}, \hyperref[ex:6.2.7]{\pageref*{ex:6.2.7}}, \hyperref[ex:6.2.9]{\pageref*{ex:6.2.9}}, \hyperref[ex:6.2.12]{\pageref*{ex:6.2.12}}\par
\par\addvspace{2pt}
\needspace{3\baselineskip}
\par\noindent\hangindent=1em\hspace*{0em}\phantomsection\label{idx:1012}平方和（Sum of squares）\par
\par\noindent\hangindent=2em\hspace*{1em}\phantomsection\label{idx:1013}残差（residual），\hyperref[idxocc:1013:469]{\pageref*{idxocc:1013:469}}\par
\par\addvspace{2pt}
\par\noindent\hangindent=1em\hspace*{0em}\phantomsection\label{idx:1014}超效率（Superefficiency），\hyperref[sec:10.6.1]{\pageref*{sec:10.6.1}}\par
\par\addvspace{2pt}
\par\noindent\hangindent=1em\hspace*{0em}\phantomsection\label{idx:1015}支持集（Support），\hyperref[idxocc:1015:71]{\pageref*{idxocc:1015:71}}\par
\par\addvspace{2pt}
\needspace{3\baselineskip}
\par\noindent\hangindent=1em\hspace*{0em}\phantomsection\label{idx:1016}t检验（t test）\par
\par\noindent\hangindent=2em\hspace*{1em}\phantomsection\label{idx:1017}方差分析（ANOVA），\hyperref[eq:11.2.7]{\pageref*{eq:11.2.7}}\par
\par\noindent\hangindent=2em\hspace*{1em}\phantomsection\label{idx:1018}近似（approximate），\hyperref[exercise:8.42]{\pageref*{exercise:8.42}}\par
\par\noindent\hangindent=2em\hspace*{1em}\phantomsection\label{idx:1019}似然比检验（LRT），\hyperref[ex:8.2.6]{\pageref*{ex:8.2.6}}\par
\par\noindent\hangindent=2.8em\hspace*{1.8em}\phantomsection\label{idx:1020}配对（paired），\hyperref[exercise:8.40]{\pageref*{exercise:8.40}}\par
\par\noindent\hangindent=2.8em\hspace*{1.8em}\phantomsection\label{idx:1021}双侧（two-sided），\hyperref[ex:8.2.8]{\pageref*{ex:8.2.8}}\par
\par\noindent\hangindent=2em\hspace*{1em}\phantomsection\label{idx:1022}配对（paired），\hyperref[exercise:8.39]{\pageref*{exercise:8.39}}\par
\par\noindent\hangindent=2em\hspace*{1em}\phantomsection\label{idx:1023}回归（regression），\hyperref[eq:11.3.38]{\pageref*{eq:11.3.38}}\par
\par\noindent\hangindent=2em\hspace*{1em}\phantomsection\label{idx:1024}双样本（two-sample），\hyperref[exercise:8.41]{\pageref*{exercise:8.41}}\par
\par\noindent\hangindent=2em\hspace*{1em}\phantomsection\label{idx:1025}双侧（two-sided），\hyperref[exercise:8.41]{\pageref*{exercise:8.41}}\par
\needspace{3\baselineskip}
\par\noindent\hangindent=2em\hspace*{1em}\phantomsection\label{idx:1026}一致最强（UMP）\par
\par\noindent\hangindent=2.8em\hspace*{1.8em}\phantomsection\label{idx:1027}无偏（unbiased），\hyperref[exercise:8.37]{\pageref*{exercise:8.37}}\par
\par\noindent\hangindent=2em\hspace*{1em}\phantomsection\label{idx:1028}并交（union–intersection），\hyperref[ex:8.3.8]{\pageref*{ex:8.3.8}}\par
\par\addvspace{2pt}
\par\noindent\hangindent=1em\hspace*{0em}\phantomsection\label{idx:1029}泰勒级数（Taylor series），\hyperref[def:5.5.20]{\pageref*{def:5.5.20}}, \hyperref[idxocc:1029:423]{\pageref*{idxocc:1029:423}}\par
\par\addvspace{2pt}
\par\noindent\hangindent=1em\hspace*{0em}\phantomsection\label{idx:1030}泰勒定理（Taylor’s Theorem），\hyperref[thm:5.5.21]{\pageref*{thm:5.5.21}}\par
\par\noindent\hangindent=2em\hspace*{1em}\phantomsection\label{idx:1031}在中心极限定理中的用途（use in CLT），\hyperref[idxocc:1031:224]{\pageref*{idxocc:1031:224}}\par
\par\addvspace{2pt}
\par\noindent\hangindent=1em\hspace*{0em}\phantomsection\label{idx:1032}检验函数（Test function），\hyperref[idxocc:1032:346]{\pageref*{idxocc:1032:346}}\par
\par\addvspace{2pt}
\par\noindent\hangindent=1em\hspace*{0em}\phantomsection\label{idx:1033}检验统计量（Test statistic），\hyperref[idxocc:1033:334]{\pageref*{idxocc:1033:334}}\par
\par\addvspace{2pt}
\par\noindent\hangindent=1em\hspace*{0em}\phantomsection\label{idx:1034}阈值参数（Threshold parameter），\hyperref[ex:3.5.3]{\pageref*{ex:3.5.3}}\par
\par\addvspace{2pt}
\par\noindent\hangindent=1em\hspace*{0em}\phantomsection\label{idx:1035}全最小二乘（Total least squares），\hyperref[idxocc:1035:501]{\pageref*{idxocc:1035:501}}\par
\par\addvspace{2pt}
\needspace{3\baselineskip}
\par\noindent\hangindent=1em\hspace*{0em}\phantomsection\label{idx:1036}变换（Transformation）\par
\par\noindent\hangindent=2em\hspace*{1em}\phantomsection\label{idx:1037}Box–Cox（Box-Cox），\hyperref[exercise:7.16]{\pageref*{exercise:7.16}}\par
\par\noindent\hangindent=2em\hspace*{1em}\phantomsection\label{idx:1038}Box–Muller（Box-Muller），\hyperref[ex:5.6.4]{\pageref*{ex:5.6.4}}\par
\par\noindent\hangindent=2em\hspace*{1em}\phantomsection\label{idx:1039}Fisher的z变换（Fisher’s z），\hyperref[item:10.17.e]{\pageref*{item:10.17.e}}\par
\par\noindent\hangindent=2em\hspace*{1em}\phantomsection\label{idx:1040}单调（monotone），\hyperref[idxocc:1040:71]{\pageref*{idxocc:1040:71}}\par
\par\noindent\hangindent=2em\hspace*{1em}\phantomsection\label{idx:1041}概率积分（probability integral），\hyperref[ex:5.6.3]{\pageref*{ex:5.6.3}}\par
\par\addvspace{2pt}
\par\noindent\hangindent=1em\hspace*{0em}\phantomsection\label{idx:1042}截尾均值（Trimmed mean），\hyperref[idxocc:1042:442]{\pageref*{idxocc:1042:442}}\par
\par\addvspace{2pt}
\par\noindent\hangindent=1em\hspace*{0em}\phantomsection\label{idx:1043}调节参数（Tuning parameter），\hyperref[eq:10.2.2]{\pageref*{eq:10.2.2}}\par
\par\addvspace{2pt}
\par\noindent\hangindent=1em\hspace*{0em}\phantomsection\label{idx:1044}无偏估计量（Unbiased estimator），\hyperref[idxocc:1044:206]{\pageref*{idxocc:1044:206}}, \hyperref[def:7.3.2]{\pageref*{def:7.3.2}}\par
\par\noindent\hangindent=2em\hspace*{1em}\phantomsection\label{idx:1045}最优（best），\hyperref[def:7.3.7]{\pageref*{def:7.3.7}}\par
\par\noindent\hangindent=2.8em\hspace*{1.8em}\phantomsection\label{idx:1046}二项（binomial），\hyperref[ex:7.3.24]{\pageref*{ex:7.3.24}}\par
\par\noindent\hangindent=2.8em\hspace*{1.8em}\phantomsection\label{idx:1047}刻画（characterization），\hyperref[thm:7.3.20]{\pageref*{thm:7.3.20}}\par
\par\noindent\hangindent=2.8em\hspace*{1.8em}\phantomsection\label{idx:1048}完备性（completeness），\hyperref[thm:7.3.23]{\pageref*{thm:7.3.23}}\par
\par\noindent\hangindent=2.8em\hspace*{1.8em}\phantomsection\label{idx:1049}泊松（Poisson），\hyperref[ex:7.3.12]{\pageref*{ex:7.3.12}}\par
\par\noindent\hangindent=2.8em\hspace*{1.8em}\phantomsection\label{idx:1050}充分性（sufficiency），\hyperref[thm:7.3.17]{\pageref*{thm:7.3.17}}\par
\par\noindent\hangindent=2.8em\hspace*{1.8em}\phantomsection\label{idx:1051}均匀（uniform），\hyperref[ex:7.3.22]{\pageref*{ex:7.3.22}}\par
\par\noindent\hangindent=2.8em\hspace*{1.8em}\phantomsection\label{idx:1052}唯一（unique），\hyperref[thm:7.3.19]{\pageref*{thm:7.3.19}}\par
\par\noindent\hangindent=2em\hspace*{1em}\phantomsection\label{idx:1053}线性（linear），\hyperref[exercise:7.41]{\pageref*{exercise:7.41}}\par
\par\noindent\hangindent=2em\hspace*{1em}\phantomsection\label{idx:1054}位置（location），\hyperref[exercise:6.41]{\pageref*{exercise:6.41}}\par
\par\noindent\hangindent=2em\hspace*{1em}\phantomsection\label{idx:1055}正态（normal），\hyperref[ex:7.3.3]{\pageref*{ex:7.3.3}}\par
\par\noindent\hangindent=2.8em\hspace*{1.8em}\phantomsection\label{idx:1056}标准差（standard deviation），\hyperref[exercise:7.50]{\pageref*{exercise:7.50}}\par
\par\noindent\hangindent=2em\hspace*{1em}\phantomsection\label{idx:1057}泊松（Poisson），\hyperref[ex:7.3.8]{\pageref*{ex:7.3.8}}\par
\par\noindent\hangindent=2em\hspace*{1em}\phantomsection\label{idx:1058}与贝叶斯估计量的关系（relation to Bayes estimator），\hyperref[sec:7.5.2]{\pageref*{sec:7.5.2}}\par
\par\noindent\hangindent=2em\hspace*{1em}\phantomsection\label{idx:1059}均匀（uniform），\hyperref[ex:7.3.13]{\pageref*{ex:7.3.13}}, \hyperref[ex:7.3.21]{\pageref*{ex:7.3.21}}\par
\par\noindent\hangindent=2em\hspace*{1em}\phantomsection\label{idx:1060}方差界（variance bound），\hyperref[ex:7.3.14]{\pageref*{ex:7.3.14}}\par
\par\noindent\hangindent=2em\hspace*{1em}\phantomsection\label{idx:1061}零（zero），\hyperref[ex:7.3.21]{\pageref*{ex:7.3.21}}\par
\par\addvspace{2pt}
\par\noindent\hangindent=1em\hspace*{0em}\phantomsection\label{idx:1062}方差（Variance），\hyperref[def:2.3.2]{\pageref*{def:2.3.2}}\par
\par\noindent\hangindent=2em\hspace*{1em}\phantomsection\label{idx:1063}渐近（asymptotic），\hyperref[idxocc:1063:410]{\pageref*{idxocc:1063:410}}\par
\par\noindent\hangindent=2em\hspace*{1em}\phantomsection\label{idx:1064}条件（conditional），\hyperref[idxocc:1064:155]{\pageref*{idxocc:1064:155}}, \hyperref[thm:4.4.7]{\pageref*{thm:4.4.7}}\par
\par\noindent\hangindent=2em\hspace*{1em}\phantomsection\label{idx:1065}极限（limiting），\hyperref[def:10.1.7]{\pageref*{def:10.1.7}}\par
\par\noindent\hangindent=2em\hspace*{1em}\phantomsection\label{idx:1066}性质（properties），\hyperref[thm:2.3.4]{\pageref*{thm:2.3.4}}\par
\par\noindent\hangindent=2em\hspace*{1em}\phantomsection\label{idx:1067}均值的二次函数（quadratic function of mean），\hyperref[idxocc:1067:110]{\pageref*{idxocc:1067:110}}\par
\par\noindent\hangindent=2em\hspace*{1em}\phantomsection\label{idx:1068}和（sum），\hyperref[thm:4.5.6]{\pageref*{thm:4.5.6}}, \hyperref[exercise:4.44]{\pageref*{exercise:4.44}}, \hyperref[lem:5.2.5]{\pageref*{lem:5.2.5}}\par
\par\noindent\hangindent=2em\hspace*{1em}\phantomsection\label{idx:1069}泰勒近似（Taylor approximation），\hyperref[eq:5.5.9]{\pageref*{eq:5.5.9}}\par
\par\addvspace{2pt}
\par\noindent\hangindent=1em\hspace*{0em}\phantomsection\label{idx:1070}韦恩图（Venn diagram），\hyperref[idxocc:1070:34]{\pageref*{idxocc:1070:34}}\par
\par\addvspace{2pt}
\par\noindent\hangindent=1em\hspace*{0em}\phantomsection\label{idx:1071}看守问题（Warden problem），\hyperref[ex:1.3.4]{\pageref*{ex:1.3.4}}, \hyperref[exercise:1.37]{\pageref*{exercise:1.37}}\par
\end{multicols}
\endgroup


\end{document}
